% Auto-generated from data/segments.json + layout.json (+ transforms) — do not edit by hand.
% volume: 02Vin02
% mode: reading
% segments: 3830
% transform_rules: 0
% toc_marks: 518

% Volume layout defaults (layout.json ``layout``).
\csromanlayoutapply{1}{1.5}{21.6pt}{5pt}{30pt}{65pt}{2.5em}%

% Reading physical-page layout (layout.json ``page_layout_reading_mode``).
\csromanreadingpagelayoutvolume{1}{1.5}{21.6pt}{5pt}{30pt}{65pt}{2.5em}%
\csromanreadingpagelayoutenable

\pitaka{\textbf{วินยปิฏก}}
\csromanmatikaanchor{mtk.0}
\csromantocmarknopage{chapter}{ปาจิตฺติยปาลิ}{mtk.0}
\bookmark[dest={mtk.0},level=0]{ปาจิตฺติยปาลิ}
\gambhira{\textbf{ปาจิตฺติยปาฬิ}}
\namakkaram{นโม ตสฺส ภควโต อรหโต สมฺมาสมฺพุทฺธสฺส.}
\csromanmatikaanchor{mtk.1}
\csromantocmarknopage{chapter}{5. ปาจิตฺติยกณฺฑ}{mtk.1}
\bookmark[dest={mtk.1},level=0]{5. ปาจิตฺติยกณฺฑ}
\setcsromanheadcha{5. ปาจิตฺติยกณฺฑ}
\setcsromanheadhclear{}
\chapterhead{5. \textbf{ปาจิตฺติยกณฺฑ}}
\csromanmatikaanchor{mtk.2}
\csromanmatikaanchor{mtk.3}
\csromantocmarknopage{section}{1. มุสาวาทวคฺค}{mtk.2}
\bookmark[dest={mtk.2},level=1]{1. มุสาวาทวคฺค}
\csromantocmarkref{subsection}{1. มุสาวาทสิกฺขาปท}{mtk.3}
\bookmark[dest={mtk.3},level=2]{1. มุสาวาทสิกฺขาปท}
\setcsromanheadh{1. มุสาวาทวคฺค}
\csromanheader{2}{1. \textbf{มุสาวาทวคฺค 1. มุสาวาทสิกฺขาปท}}
\prose{\csromanfolio{1}อิเม โข ปนายสฺมนฺโต เทฺวนวุติ ปาจิตฺติยา ธมฺมา อุทฺเทสํ อาคจฺฉนฺติ.}

\proseitem{1}{เตน สมเยน พุทฺโธ ภควา สาวตฺถิยํ วิหรติ เชตวเน อนาถปิณฺฑิกสฺส อาราเม.\csromanspacer{} เตน โข ปน สมเยน หตฺถโก สกฺยปุตฺโต วาทกฺขิตฺโต โหติ.\csromanspacer{} โส ติตฺถิเยหิ สทฺธิํ สลฺลปนฺโต อวชานิตฺวา ปฏิชานาติ, ปฏิชานิตฺวา อวชานาติ, อญฺเญนญฺญํ ปฏิจรติ, สมฺปชานมุสา ภาสติ, สงฺเกตํ กตฺวา วิสํวาเทติ.\csromanspacer{} ติตฺถิยา อุชฺฌายนฺติ ขิยฺยนฺติ วิปาเจนฺติ “กถํ หิ นาม หตฺถโก สกฺยปุตฺโต อเมฺหหิ สทฺธิํ สลฺลปนฺโต อวชานิตฺวา ปฏิชานิสฺสติ, ปฏิชานิตฺวา อวชานิสฺสติ, อญฺเญนญฺญํ ปฏิจริสฺสติ, สมฺปชานมุสา ภาสิสฺสติ, สงฺเกตํ กตฺวา วิสํวาเทสฺสตี”ติ.}

\prose{อสฺโสสุํ โข ภิกฺขู เตสํ ติตฺถิยานํ อุชฺฌายนฺตานํ ขิยฺยนฺตานํ วิปาเจนฺตานํ.\csromanspacer{} อถ โข เต ภิกฺขู เยน หตฺถโก สกฺยปุตฺโต เตนุปสงฺกมิํสุ, อุปสงฺกมิตฺวา หตฺถกํ สกฺยปุตฺตํ เอตทโวจุํ “สจฺจํ กิร ตฺวํ อาวุโส หตฺถก ติตฺถิเยหิ สทฺธิํ สลฺลปนฺโต อวชานิตฺวา ปฏิชานาสิ, ปฏิชานิตฺวา อวชานาสิ, อญฺเญนญฺญํ ปฏิจรสิ, สมฺปชานมุสา ภาสสิ, สงฺเกตํ กตฺวา วิสํวาเทสี”ติ.\csromanspacer{} เอเต โข อาวุโส \csromanfolio{2}ติตฺถิยา นาม เยน เกนจิ เชตพฺพา, เนว เตสํ ชโย ทาตพฺโพติ.\csromanspacer{} เย เต ภิกฺขู อปฺปิจฺฉา ฯเปฯ เต อุชฺฌายนฺติ ขิยฺยนฺติ วิปาเจนฺติ “กถํ หิ นาม หตฺถโก สกฺยปุตฺโต ติตฺถิเยหิ สทฺธิํ สลฺลปนฺโต อวชานิตฺวา ปฏิชานิสฺสติ, ปฏิชานิตฺวา อวชานิสฺสติ, อญฺเญนญฺญํ ปฏิจริสฺสติ, สมฺปชานมุสา ภาสิสฺสติ, สงฺเกตํ กตฺวา วิสํวาเทสฺสตี”ติ.}

\prose{อถ โข เต ภิกฺขู หตฺถกํ สกฺยปุตฺตํ อเนกปริยาเยน วิครหิตฺวา ภควโต เอตมตฺถํ อาโรเจสุํ.\csromanspacer{} อถ โข ภควา เอตสฺมิํ นิทาเน เอตสฺมิํ ปกรเณ ภิกฺขุสํฆํ สนฺนิปาตาเปตฺวา หตฺถกํ สกฺยปุตฺตํ ปฏิปุจฺฉิ “สจฺจํ กิร ตฺวํ หตฺถก ติตฺถิเยหิ สทฺธิํ สลฺลปนฺโต อวชานิตฺวา ปฏิชานาสิ, ปฏิชานิตฺวา อวชานาสิ, อญฺเญนญฺญํ ปฏิจรสิ, สมฺปชานมุสา ภาสสิ, สงฺเกตํ กตฺวา วิสํวาเทสี”ติ.\csromanspacer{} สจฺจํ ภควาติ.\csromanspacer{} วิครหิ พุทฺโธ ภควา ฯเปฯ.\csromanspacer{} กถํ หิ นาม ตฺวํ โมฆปุริส ติตฺถิเยหิ สทฺธิํ สลฺลปนฺโต อวชานิตฺวา ปฏิชานิสฺสสิ, ปฏิชานิตฺวา อวชานิสฺสสิ, อญฺเญนญฺญํ ปฏิจริสฺสสิ, สมฺปชานมุสา ภาสิสฺสสิ, สงฺเกตํ กตฺวา วิสํวาเทสฺสสิ.\csromanspacer{} เนตํ โมฆปุริส อปฺปสนฺนานํ วา ปสาทาย ฯเปฯ.\csromanspacer{} เอวญฺจ ปน ภิกฺขเว อิมํ สิกฺขาปทํ อุทฺทิเสยฺยาถ\csromandash{}}

\proseitem{2}{\textbf{“สมฺปชานมุสาวาเท ปาจิตฺติยนฺ”}ติ.}

\proseitem{3}{สมฺปชานมุสาวาโท นาม วิสํวาทน\-ปุเรกฺขารสฺส วาจา คิรา พฺยปฺปโถ วจีเภโท วาจสิกา วิญฺญตฺติ, อฏฺฐ อนริยโวหาราอทิฏฺฐํ ทิฏฺฐํ เมติ, อสฺสุตํ สุตํ เมติ, อมุตํ มุตํ เมติ, อวิญฺญาตํ วิญฺญาตํ เมติ, ทิฏฺฐํ อทิฏฺฐํ เมติ, สุตํ อสฺสุตํ เมติ, มุตํ อมุตํ เมติ, วิญฺญาตํ อวิญฺญาตํ เมติ.}

\prose{\textbf{อทิฏฺฐํ} นาม น จกฺขุนา ทิฏฺฐํ.\csromanspacer{} \textbf{อสฺสุตํ} นาม น โสเตน สุตํ.\csromanspacer{} \textbf{อมุตํ} นาม น ฆาเนน ฆายิตํ, น ชิวฺหาย สายิตํ, น กาเยน ผุฏฺฐํ.\csromanspacer{} \textbf{อวิญฺญาตํ} นาม น มนสา วิญฺญาตํ.\csromanspacer{} \textbf{ทิฏฺฐํ} นาม จกฺขุนา ทิฏฺฐํ.\csromanspacer{} \textbf{สุตํ} นาม โสเตน สุตํ.\csromanspacer{} \textbf{มุตํ} นาม ฆาเนน ฆายิตํ, ชิวฺหาย สายิตํ, กาเยน ผุฏฺฐํ.\csromanspacer{} \textbf{วิญฺญาตํ} นาม มนสา วิญฺญาตํ.}

\proseitem{4}{ตีหากาเรหิ อทิฏฺฐํ ทิฏฺฐํ เมติ สมฺปชานมุสา ภณนฺตสฺส อาปตฺติ ปาจิตฺติยสฺส, ปุพฺเพวสฺส โหติ “มุสา ภณิสฺสนฺ”ติ, ภณนฺตสฺส โหติ “มุสา ภณามี”ติ, ภณิตสฺส โหติ “มุสา มยา ภณิตนฺ”ติ.}

\prose{\csromanfolio{3}จตูหากาเรหิ อทิฏฺฐํ ทิฏฺฐํ เมติ สมฺปชานมุสา ภณนฺตสฺส อาปตฺติ ปาจิตฺติยสฺส, ปุพฺเพวสฺส โหติ “มุสา ภณิสฺสนฺ”ติ, ภณนฺตสฺส โหติ “มุสา ภณามี”ติ, ภณิตสฺส โหติ “มุสา มยา ภณิตนฺ”ติ, วินิธาย ทิฏฺฐิํ.}

\prose{ปญฺจหากาเรหิ อทิฏฺฐํ ทิฏฺฐํ เมติ สมฺปชานมุสา ภณนฺตสฺส อาปตฺติ ปาจิตฺติยสฺส, ปุพฺเพวสฺส โหติ “มุสา ภณิสฺสนฺ”ติ, ภณนฺตสฺส โหติ “มุสา ภณามี”ติ, ภณิตสฺส โหติ “มุสา มยา ภณิตนฺ”ติ, วินิธาย ทิฏฺฐิํ, วินิธาย ขนฺติํ.}

\prose{ฉหากาเรหิ อทิฏฺฐํ ทิฏฺฐํ เมติ สมฺปชานมุสา ภณนฺตสฺส อาปตฺติ ปาจิตฺติยสฺส, ปุพฺเพวสฺส โหติ “มุสา ภณิสฺสนฺ”ติ, ภณนฺตสฺส โหติ “มุสา ภณามี”ติ, ภณิตสฺส โหติ “มุสา มยา ภณิตนฺ”ติ, วินิธาย ทิฏฺฐิํ, วินิธาย ขนฺติํ, วินิธาย รุจิํ.}

\prose{สตฺตหากาเรหิ อทิฏฺฐํ ทิฏฺฐํ เมติ สมฺปชานมุสา ภณนฺตสฺส อาปตฺติ ปาจิตฺติยสฺส, ปุพฺเพวสฺส โหติ “มุสา ภณิสฺสนฺ”ติ, ภณนฺตสฺส โหติ “มุสา ภณามี”ติ, ภณิตสฺส โหติ “มุสา มยา ภณิตนฺ”ติ, วินิธาย ทิฏฺฐิํ, วินิธาย ขนฺติํ, วินิธาย รุจิํ, วินิธาย ภาวํ.}

\proseitem{5}{ตีหากาเรหิ อสฺสุตํ สุตํ เมติ ฯเปฯ อมุตํ มุตํ เมติ ฯเปฯ อวิญฺญาตํ วิญฺญาตํ เมติ สมฺปชานมุสา ภณนฺตสฺส อาปตฺติ ปาจิตฺติยสฺส, ปุพฺเพวสฺส โหติ “มุสา ภณิสฺสนฺ”ติ, ภณนฺตสฺส โหติ “มุสา ภณามี”ติ, ภณิตสฺส โหติ “มุสา มยา ภณิตนฺ”ติ.}

\prose{จตูหากาเรหิ ฯเปฯ.\csromanspacer{} ปญฺจหากาเรหิ ฯเปฯ.\csromanspacer{} ฉหากาเรหิ ฯเปฯ.\csromanspacer{} สตฺตหากาเรหิ อวิญฺญาตํ วิญฺญาตํ เมติ สมฺปชานมุสา ภณนฺตสฺส อาปตฺติ ปาจิตฺติยสฺส, ปุพฺเพวสฺส โหติ “มุสา ภณิสฺสนฺ”ติ, ภณนฺตสฺส โหติ “มุสา ภณามี”ติ, ภณิตสฺส โหติ “มุสา มยา ภณิตนฺ”ติ, วินิธาย ทิฏฺฐิํ, วินิธาย ขนฺติํ, วินิธาย รุจิํ, วินิธาย ภาวํ.}

\proseitem{6}{ตีหากาเรหิ อทิฏฺฐํ ทิฏฺฐญฺจ เม สุตญฺจาติ สมฺปชานมุสา ภณนฺตสฺส อาปตฺติ ปาจิตฺติยสฺส ฯเปฯ.\csromanspacer{} ตีหากาเรหิ อทิฏฺฐํ ทิฏฺฐญฺจ เม มุตญฺจาติ สมฺปชานมุสา ภณนฺตสฺส อาปตฺติ ปาจิตฺติยสฺส ฯเปฯ.\csromanspacer{} ตีหากาเรหิ อทิฏฺฐํ ทิฏฺฐญฺจ เม วิญฺญาตญฺจาติ สมฺปชานมุสา ภณนฺตสฺส อาปตฺติ \csromanfolio{4}ปาจิตฺติยสฺส ฯเปฯ.\csromanspacer{} ตีหากาเรหิ อทิฏฺฐํ ทิฏฺฐญฺจ เม สุตญฺจ มุตญฺจาติ ฯเปฯ.\csromanspacer{} ตีหากาเรหิ อทิฏฺฐํ ทิฏฺฐญฺจ เม สุตญฺจ วิญฺญาตญฺจาติ ฯเปฯ.\csromanspacer{} ตีหากาเรหิ อทิฏฺฐํ ทิฏฺฐญฺจ เม สุตญฺจ มุตญฺจ วิญฺญาตญฺจาติ สมฺปชานมุสา ภณนฺตสฺส อาปตฺติ ปาจิตฺติยสฺส ฯเปฯ.}

\prose{ตีหากาเรหิ อสฺสุตํ สุตญฺจ เม มุตญฺจาติ ฯเปฯ.\csromanspacer{} ตีหากาเรหิ อสฺสุตํ สุตญฺจ เม วิญฺญาตญฺจาติ ฯเปฯ.\csromanspacer{} ตีหากาเรหิ อสฺสุตํ สุตญฺจ เม ทิฏฺฐญฺจาติ สมฺปชานมุสา ภณนฺตสฺส อาปตฺติ ปาจิตฺติยสฺส ฯเปฯ.\csromanspacer{} ตีหากาเรหิ อสฺสุตํ สุตญฺจ เม มุตญฺจ วิญฺญาตญฺจาติ ฯเปฯ.\csromanspacer{} ตีหากาเรหิ อสฺสุตํ สุตญฺจ เม มุตญฺจ ทิฏฺฐญฺจาติ ฯเปฯ.\csromanspacer{} ตีหากาเรหิ อสฺสุตํ สุตญฺจ เม มุตญฺจ วิญฺญาตญฺจ ทิฏฺฐญฺจาติ สมฺปชานมุสา ภณนฺตสฺส อาปตฺติ ปาจิตฺติยสฺส ฯเปฯ.}

\prose{ตีหากาเรหิ อมุตํ มุตญฺจ เม วิญฺญาตญฺจาติ ฯเปฯ.\csromanspacer{} ตีหากาเรหิ อมุตํ มุตญฺจ เม ทิฏฺฐญฺจาติ ฯเปฯ.\csromanspacer{} ตีหากาเรหิ อมุตํ มุตญฺจ เม สุตญฺจาติ สมฺปชานมุสา ภณนฺตสฺส อาปตฺติ ปาจิตฺติยสฺส ฯเปฯ.\csromanspacer{} ตีหากาเรหิ อมุตํ มุตญฺจ เม วิญฺญาตญฺจ ทิฏฺฐญฺจาติ ฯเปฯ.\csromanspacer{} ตีหากาเรหิ อมุตํ มุตญฺจ เม วิญฺญาตญฺจ สุตญฺจาติ ฯเปฯ.\csromanspacer{} ตีหากาเรหิ อมุตํ มุตญฺจ เม วิญฺญาตญฺจ ทิฏฺฐญฺจ สุตญฺจาติ สมฺปชานมุสา ภณนฺตสฺส อาปตฺติ ปาจิตฺติยสฺส ฯเปฯ.}

\prose{ตีหากาเรหิ อวิญฺญาตํ วิญฺญาตญฺจ เม ทิฏฺฐญฺจาติ ฯเปฯ.\csromanspacer{} ตีหากาเรหิ อวิญฺญาตํ วิญฺญาตญฺจ เม สุตญฺจาติ ฯเปฯ.\csromanspacer{} ตีหากาเรหิ อวิญฺญาตํ วิญฺญาตญฺจ เม มุตญฺจาติ สมฺปชานมุสา ภณนฺตสฺส อาปตฺติ ปาจิตฺติยสฺส ฯเปฯ.\csromanspacer{} ตีหากาเรหิ อวิญฺญาตํ วิญฺญาตญฺจ เม ทิฏฺฐญฺจ สุตญฺจาติ ฯเปฯ.\csromanspacer{} ตีหากาเรหิ อวิญฺญาตํ วิญฺญาตญฺจ เม ทิฏฺฐญฺจ มุตญฺจาติ ฯเปฯ.\csromanspacer{} ตีหากาเรหิ อวิญฺญาตํ วิญฺญาตญฺจ เม ทิฏฺฐญฺจ สุตญฺจ มุตญฺจาติ สมฺปชานมุสา ภณนฺตสฺส อาปตฺติ ปาจิตฺติยสฺส ฯเปฯ.}

\proseitem{7}{ตีหากาเรหิ ทิฏฺฐํ อทิฏฺฐํ เมติ ฯเปฯ สุตํ อสฺสุตํ เมติ ฯเปฯ มุตํ อมุตํ เมติ ฯเปฯ วิญฺญาตํ อวิญฺญาตํ เมติ สมฺปชานมุสา ภณนฺตสฺส อาปตฺติ ปาจิตฺติยสฺส ฯเปฯ.}

\proseitem{8}{ตีหากาเรหิ ทิฏฺฐํ สุตํ เมติ ฯเปฯ.\csromanspacer{} ตีหากาเรหิ ทิฏฺฐํ มุตํ เมติ ฯเปฯ.\csromanspacer{} ตีหากาเรหิ ทิฏฺฐํ วิญฺญาตํ เมติ สมฺปชานมุสา ภณนฺตสฺส \csromanfolio{5}อาปตฺติ ปาจิตฺติยสฺส ฯเปฯ.\csromanspacer{} ตีหากาเรหิ ทิฏฺฐํ สุตญฺจ เม มุตญฺจาติ ฯเปฯ.\csromanspacer{} ตีหากาเรหิ ทิฏฺฐํ สุตญฺจ เม วิญฺญาตญฺจาติ ฯเปฯ.\csromanspacer{} ตีหากาเรหิ ทิฏฺฐํ สุตญฺจ เม มุตญฺจ วิญฺญาตญฺจาติ สมฺปชานมุสา ภณนฺตสฺส อาปตฺติ ปาจิตฺติยสฺส ฯเปฯ.}

\prose{ตีหากาเรหิ สุตํ มุตํ เมติ ฯเปฯ.\csromanspacer{} ตีหากาเรหิ สุตํ วิญฺญาตํ เมติ ฯเปฯ.\csromanspacer{} ตีหากาเรหิ สุตํ ทิฏฺฐํ เมติ สมฺปชานมุสา ภณนฺตสฺส อาปตฺติ ปาจิตฺติยสฺส ฯเปฯ.\csromanspacer{} ตีหากาเรหิ สุตํ มุตญฺจ เม วิญฺญาตญฺจาติ ฯเปฯ.\csromanspacer{} ตีหากาเรหิ สุตํ มุตญฺจ เม ทิฏฺฐญฺจาติ ฯเปฯ.\csromanspacer{} ตีหากาเรหิ สุตํ มุตญฺจ เม วิญฺญาตญฺจ ทิฏฺฐญฺจาติ สมฺปชานมุสา ภณนฺตสฺส อาปตฺติ ปาจิตฺติยสฺส ฯเปฯ.}

\prose{ตีหากาเรหิ มุตํ วิญฺญาตํ เมติ ฯเปฯ.\csromanspacer{} ตีหากาเรหิ มุตํ ทิฏฺฐํ เมติ ฯเปฯ.\csromanspacer{} ตีหากาเรหิ มุตํ สุตํ เมติ สมฺปชานมุสา ภณนฺตสฺส อาปตฺติ ปาจิตฺติยสฺส ฯเปฯ.\csromanspacer{} ตีหากาเรหิ มุตํ วิญฺญาตญฺจ เม ทิฏฺฐญฺจาติ ฯเปฯ.\csromanspacer{} ตีหากาเรหิ มุตํ วิญฺญาตญฺจ เม สุตญฺจาติ ฯเปฯ.\csromanspacer{} ตีหากาเรหิ มุตํ วิญฺญาตญฺจ เม ทิฏฺฐญฺจ สุตญฺจาติ สมฺปชานมุสา ภณนฺตสฺส อาปตฺติ ปาจิตฺติยสฺส ฯเปฯ.}

\prose{ตีหากาเรหิ วิญฺญาตํ ทิฏฺฐํ เมติ ฯเปฯ.\csromanspacer{} ตีหากาเรหิ วิญฺญาตํ สุตํ เมติ ฯเปฯ.\csromanspacer{} ตีหากาเรหิ วิญฺญาตํ มุตํ เมติ สมฺปชานมุสา ภณนฺตสฺส อาปตฺติ ปาจิตฺติยสฺส ฯเปฯ.\csromanspacer{} ตีหากาเรหิ วิญฺญาตํ ทิฏฺฐญฺจ เม สุตญฺจาติ ฯเปฯ.\csromanspacer{} ตีหากาเรหิ วิญฺญาตํ ทิฏฺฐญฺจ เม มุตญฺจาติ ฯเปฯ.\csromanspacer{} ตีหากาเรหิ วิญฺญาตํ ทิฏฺฐญฺจ เม สุตญฺจ มุตญฺจาติ สมฺปชานมุสา ภณนฺตสฺส อาปตฺติ ปาจิตฺติยสฺส ฯเปฯ.}

\proseitem{9}{ตีหากาเรหิ ทิฏฺเฐ เวมติโก ทิฏฺฐํ โนกปฺเปติ ทิฏฺฐํ นสฺสรติ ทิฏฺฐํ ปมุฏฺโฐ โหติ ฯเปฯ สุเต เวมติโก สุตํ โนกปฺเปติ สุตํ นสฺสรติ สุตํ ปมุฏฺโฐ โหติ ฯเปฯ มุเต เวมติโก มุตํ โนกปฺเปติ มุตํ นสฺสรติ มุตํ ปมุฏฺโฐ โหติ ฯเปฯ วิญฺญาเต เวมติโก วิญฺญาตํ โนกปฺเปติ วิญฺญาตํ นสฺสรติ วิญฺญาตํ ปมุฏฺโฐ โหติ, วิญฺญาตญฺจ เม ทิฏฺฐญฺจาติ ฯเปฯ วิญฺญาตํ ปมุฏฺโฐ โหติ, วิญฺญาตญฺจ เม สุตญฺจาติ ฯเปฯ วิญฺญาตํ ปมุฏฺโฐ โหติ, วิญฺญาตญฺจ เม มุตญฺจาติ ฯเปฯ วิญฺญาตํ ปมุฏฺโฐ โหติ, วิญฺญาตญฺจ เม ทิฏฺฐญฺจ สุตญฺจาติ ฯเปฯ วิญฺญาตํ \csromanfolio{6}ปมุฏฺโฐ โหติ, วิญฺญาตญฺจ เม ทิฏฺฐญฺจ มุตญฺจาติ ฯเปฯ วิญฺญาตํ ปมุฏฺโฐ โหติ, วิญฺญาตญฺจ เม ทิฏฺฐญฺจ สุตญฺจ มุตญฺจาติ สมฺปชานมุสา ภณนฺตสฺส อาปตฺติ ปาจิตฺติยสฺส.}

\proseitem{10}{จตูหากาเรหิ ฯเปฯ.\csromanspacer{} ปญฺจหากาเรหิ ฯเปฯ.\csromanspacer{} ฉหากาเรหิ ฯเปฯ.\csromanspacer{} สตฺตหากาเรหิ ฯเปฯ วิญฺญาตํ ปมุฏฺโฐ โหติ, วิญฺญาตญฺจ เม ทิฏฺฐญฺจ สุตญฺจ มุตญฺจาติ สมฺปชานมุสา ภณนฺตสฺส อาปตฺติ ปาจิตฺติยสฺส, ปุพฺเพวสฺส โหติ “มุสา ภณิสฺสนฺ”ติ, ภณนฺตสฺส โหติ “มุสา ภณามี”ติ, ภณิตสฺส โหติ “มุสา มยา ภณิตนฺ”ติ, วินิธาย ทิฏฺฐิํ, วินิธาย ขนฺติํ, วินิธาย รุจิํ, วินิธาย ภาวํ.}

\proseitemwithcloserruled{11}{อนาปตฺติ ทวา ภณติ, รวา ภณติ\footnote{ทวาย ภณติ, รวาย ภณติ (สฺยา)}. \textbf{“ทวา ภณติ} นาม สหสา ภณติ.\csromanspacer{} \textbf{รวา ภณติ} นาม ‘อญฺญํ ภณิสฺสามี’ติ อญฺญํ ภณติ”.\csromanspacer{} อุมฺมตฺตกสฺส อาทิกมฺมิกสฺสาติ.}{มุสาวาทสิกฺขาปทํ นิฏฺฐิตํ ปฐมํ.}
\csromanmatikaanchor{mtk.4}
\csromantocmarkref{subsection}{2. โอมสวาทสิกฺขาปท}{mtk.4}
\bookmark[dest={mtk.4},level=2]{2. โอมสวาทสิกฺขาปท}
\csromanheader{2}{1. \textbf{มุสาวาทวคฺค 2. โอมสวาทสิกฺขาปท}}
\proseitem{12}{เตน สมเยน พุทฺโธ ภควา สาวตฺถิยํ วิหรติ เชตวเน อนาถปิณฺฑิกสฺส อาราเม.\csromanspacer{} เตน โข ปน สมเยน ฉพฺพคฺคิยา ภิกฺขู เปสเลหิ ภิกฺขูหิ สทฺธิํ ภณฺฑนฺตา\footnote{ภณฺเฑนฺตา (อิติปิ)} เปสเล ภิกฺขู โอมสนฺติ, ชาติยาปิ นาเมนปิ โคตฺเตนปิ กมฺเมนปิ สิปฺเปนปิ อาพาเธนปิ ลิงฺเคนปิ กิเลเสนปิ อาปตฺติยาปิ หีเนนปิ อกฺโกเสน ขุํเสนฺติ วมฺเภนฺติ.\csromanspacer{} เย เต ภิกฺขู อปฺปิจฺฉา ฯเปฯ เต อุชฺฌายนฺติ ขิยฺยนฺติ วิปาเจนฺติ “กถํ หิ นาม ฉพฺพคฺคิยา ภิกฺขู เปสเลหิ ภิกฺขูหิ สทฺธิํ ภณฺฑนฺตา เปสเล ภิกฺขู โอมสิสฺสนฺติ, ชาติยาปิ นาเมนปิ โคตฺเตนปิ กมฺเมนปิ สิปฺเปนปิ อาพาเธนปิ ลิงฺเคนปิ กิเลเสนปิ อาปตฺติยาปิ หีเนนปิ อกฺโกเสน ขุํเสสฺสนฺติ วมฺเภสฺสนฺตี”ติ.}

\prose{อถ โข เต ภิกฺขู ฉพฺพคฺคิเย ภิกฺขู อเนกปริยาเยน วิครหิตฺวา ภควโต เอตมตฺถํ อาโรเจสุํ ฯเปฯ “สจฺจํ กิร ตุเมฺห ภิกฺขเว \csromanfolio{7}เปสเลหิ ภิกฺขูหิ สทฺธิํ ภณฺฑนฺตา เปสเล ภิกฺขู โอมสถ, ชาติยาปิ ฯเปฯ หีเนนปิ อกฺโกเสน ขุํเสถ วมฺเภถา”ติ.\csromanspacer{} สจฺจํ ภควาติ.\csromanspacer{} วิครหิ พุทฺโธ ภควา ฯเปฯ.\csromanspacer{} กถํ หิ นาม ตุเมฺห โมฆปุริสา เปสเลหิ ภิกฺขูหิ สทฺธิํ ภณฺฑนฺตา เปสเล ภิกฺขู โอมสิสฺสถ, ชาติยาปิ ฯเปฯ หีเนนปิ อกฺโกเสน ขุํเสสฺสถ วมฺเภสฺสถ.\csromanspacer{} เนตํ โมฆปุริสา อปฺปสนฺนานํ วา ปสาทาย ฯเปฯ วิครหิตฺวา ฯเปฯ ธมฺมิํ กถํ กตฺวา \textbf{ภิกฺขู อามนฺเตสิ\csromandash{}}}

\proseitem{13}{ภูตปุพฺพํ ภิกฺขเว ตกฺกสิลายํ\footnote{ตกฺกสีลายํ (ก) 2. พลิวทฺโท (สี), พลิพทฺโท (สฺยา) 3. อญฺฉ กูฏ (สี, สฺยา)} อญฺญตรสฺส พฺราหฺมณสฺส นนฺทิวิสาโล นาม พลีพทฺโท2 อโหสิ.\csromanspacer{} อถ โข ภิกฺขเว นนฺทิวิสาโล พลีพทฺโท ตํ พฺรหฺมณํ เอตทโวจ “คจฺฉ ตฺวํ พฺราหฺมณ, เสฏฺฐินา สทฺธิํ สหสฺเสน อพฺภูตํ กโรหิ มยฺหํ พลีพทฺโท สกฏสตํ อติพทฺธํ ปวฏฺเฏสฺสตี”ติ.\csromanspacer{} อถ โข ภิกฺขเว โส พฺราหฺมโณ เสฏฺฐินา สทฺธิํ สหสฺเสน อพฺภุตํ อกาสิ มยฺหํ พลีพทฺโท สกฏสตํ อติพทฺธํ ปวฏฺเฏสฺสตีติ.\csromanspacer{} อถ โข ภิกฺขเว โส พฺราหฺมโณ สกฏสตํ อติพนฺธิตฺวา นนฺทิวิสาลํ พลีพทฺทํ ยุญฺชิตฺวา เอตทโวจ “คจฺฉ กูฏ3, วหสฺสุ กูฏา”ติ.\csromanspacer{} อถ โข ภิกฺขเว นนฺทิวิสาโล พลีพทฺโท ตตฺเถว อฏฺฐาสิ.\csromanspacer{} อถ โข ภิกฺขเว โส พฺราหฺมโณ สหสฺเสน ปราชิโต ปชฺฌายิ.\csromanspacer{} อถ โข ภิกฺขเว นนฺทิวิสาโล พลีพทฺโท ตํ พฺราหฺมณํ เอตทโวจ “กิสฺส ตฺวํ พฺราหฺมณ ปชฺฌายสี”ติ.\csromanspacer{} ตถา หิ ปนาหํ โภ ตยา สหสฺเสน ปราชิโตติ.\csromanspacer{} กิสฺส ปน มํ ตฺวํ พฺราหฺมณ อกูฏํ กูฏวาเทน ปาเปสิ, คจฺฉ ตฺวํ พฺราหฺมณ, เสฏฺฐินา สทฺธิํ ทฺวีหิ สหสฺเสหิ อพฺภุตํ กโรหิ “มยฺหํ พลีพทฺโท สกฏสตํ อติพทฺธํ ปวฏฺเฏสฺสตี”ติ. “มา จ มํ อกูฏํ กูฏวาเทน ปาเปสี”ติ.\csromanspacer{} อถ โข ภิกฺขเว โส พฺราหฺมโณ เสฏฺฐินา สทฺธิํ ทฺวีหิ สหสฺเสหิ อพฺภุตํ อกาสิ “มยฺหํ พลีพทฺโท สกฏสตํ อติพทฺธํ ปวฏฺเฏสฺสตี”ติ.\csromanspacer{} อถ โข ภิกฺขเว โส พฺราหฺมโณ สกฏสตํ อติพนฺธิตฺวา นนฺทิวิสาลํ พลีพทฺทํ ยุญฺชิตฺวา เอตทโวจ “คจฺฉ ภทฺร วหสฺสุ ภทฺรา”ติ.\csromanspacer{} อถ โข ภิกฺขเว นนฺทิวิสาโล พลีพทฺโท สกฏสตํ อติพทฺธํ ปวฏฺเฏสิ.}

\setlength{\csromangathapagewidth}{0pt}
\csromangathameasuregroup{*มนาปเมว ภาเสยฺย, \\ มนาปํ ภาสมานสฺส, \\ ธนญฺจ นํ อลาเภสิ,}{\csromangathabat{*มนาปเมว ภาเสยฺย,}{นา’มนาปํ กุทาจนํ.} \\ \csromangathabat{มนาปํ ภาสมานสฺส,}{ครุํ ภารํ อุทพฺพหิ.} \\ \csromangathabat{ธนญฺจ นํ อลาเภสิ,}{เตน จ’ตฺตมโน อหูติ.}}
\csromangathasetleft{*มนาปเมว ภาเสยฺย, \\ มนาปํ ภาสมานสฺส, \\ ธนญฺจ นํ อลาเภสิ,}
\csromangathagroup{\csromangathastanza{\csromanfolio{8}\csromangathabat{\csromansymbolfootnote{*}{ขุ 5. 7 ปิฏฺเฐ นนฺทิวิสาลชาตเกปิ, ตตฺถ ปน มนุญฺญสทฺโท ทิสฺสติ.}มนาปเมว ภาเสยฺย,}{นา’มนาปํ กุทาจนํ.} \\ \csromangathabat{มนาปํ ภาสมานสฺส,}{ครุํ ภารํ อุทพฺพหิ.} \\ \csromangathabat{ธนญฺจ นํ อลาเภสิ,}{เตน จ’ตฺตมโน อหูติ.}}}

\prose{ตทาปิ เม ภิกฺขเว อมนาปา ขุํสนา วมฺภนา, กิมงฺคํ ปน เอตรหิ มนาปา ภวิสฺสติ ขุํสนา วมฺภนา.\csromanspacer{} เนตํ ภิกฺขเว อปฺปสนฺนานํ วา ปสาทาย ฯเปฯ.\csromanspacer{} เอวญฺจ ปน ภิกฺขเว อิมํ สิกฺขาปทํ อุทฺทิเสยฺยาถ\csromandash{}}

\proseitem{14}{\textbf{“โอมสวาเท ปาจิตฺติยนฺ”}ติ.}

\proseitem{15}{\textbf{โอมสวาโท} นาม ทสหิ อากาเรหิ โอมสติ ชาติยาปิ นาเมนปิ โคตฺเตนปิ กมฺเมนปิ สิปฺเปนปิ อาพาเธนปิ ลิงฺเคนปิ กิเลเสนปิ อาปตฺติยาปิ อกฺโกเสนปิ.}

\prose{\textbf{ชาติ} นาม เทฺว ชาติโย หีนา จ ชาติ อุกฺกฏฺฐา จ ชาติ.\csromanspacer{} \textbf{หีนา} นาม ชาติ จณฺฑาลชาติ เวนชาติ เนสาทชาติ รถการชาติ ปุกฺกุสชาติ, เอสา หีนา นาม ชาติ.\csromanspacer{} \textbf{อุกฺกฏฺฐา} นาม ชาติ ขตฺติยชาติ พฺราหฺมณชาติ, เอสา อุกฺกฏฺฐา นาม ชาติ.}

\prose{\textbf{นามํ} นาม เทฺว นามานิ หีนญฺจ นามํ อุกฺกฏฺฐญฺจ นามํ.\csromanspacer{} \textbf{หีนํ} นาม นามํ อวกณฺณกํ ชวกณฺณกํ ธนิฏฺฐกํ สวิฏฺฐกํ กุลวฑฺฒกํ, เตสุ เตสุ วา ปน ชนปเทสุ โอญฺญาตํ อวญฺญาตํ หีฬิตํ ปริภูตํ อจิตฺตีกตํ, เอตํ หีนํ นาม นามํ.\csromanspacer{} \textbf{อุกฺกฏฺฐํ} นาม นามํ พุทฺธปฺปฏิสํยุตฺตํ ธมฺมปฺปฏิสํยุตฺตํ สํฆปฺปฏิสํยุตฺตํ, เตสุ เตสุ วา ปน ชนปเทสุ อโนญฺญาตํ อนวญฺญาตํ อหีฬิตํ อปริภูตํ จิตฺตีกตํ, เอตํ อุกฺกฏฺฐํ นาม นามํ.}

\prose{\textbf{โคตฺตํ} นาม เทฺว โคตฺตานิ หีนญฺจ โคตฺตํ อุกฺกฏฺฐญฺจ โคตฺตํ.\csromanspacer{} \textbf{หีนํ} นาม โคตฺตํ โกสิยโคตฺตํ ภารทฺวาชโคตฺตํ, เตสุ เตสุ วา ปน ชนปเทสุ โอญฺญาตํ อวญฺญาตํ หีฬิตํ ปริภูตํ อจิตฺตีกตํ, เอตํ หีนํ นาม โคตฺตํ.\csromanspacer{} \textbf{อุกฺกฏฺฐํ} นาม โคตฺตํ โคตมโคตฺตํ โมคฺคลฺลานโคตฺตํ กจฺจานโคตฺตํ วาสิฏฺฐโคตฺตํ, เตสุ เตสุ วา ปน ชนปเทสุ อโนญฺญาตํ อนวญฺญาตํ อหีฬิตํ อปริภูตํ จิตฺตีกตํ, เอตํ อุกฺกฏฺฐํ นาม โคตฺตํ.}

\prose{\textbf{กมฺมํ} นาม เทฺว กมฺมานิ หีนญฺจ กมฺมํ อุกฺกฏฺฐญฺจ กมฺมํ.\csromanspacer{} \textbf{หีนํ} นาม กมฺมํ โกฏฺฐกกมฺมํ ปุปฺผฉฑฺฑกกมฺมํ, เตสุ เตสุ วา ปน ชนปเทสุ โอญฺญาตํ อวญฺญาตํ หีฬิตํ ปริภูตํ อจิตฺตีกตํ, เอตํ หีนํ นาม กมฺมํ.\csromanspacer{} \textbf{อุกฺกฏฺฐํ} นาม กมฺมํ \csromanfolio{9}กสิ วณิชฺชา โครกฺขา, เตสุ เตสุ วา ปน ชนปเทสุ อโนญฺญาตํ อนวญฺญาตํ อหีฬิตํ อปริภูตํ จิตฺตีกตํ, เอตํ อุกฺกฏฺฐํ นาม กมฺมํ.}

\prose{\textbf{สิปฺปํ} นาม เทฺว สิปฺปานิ หีนญฺจ สิปฺปํ อุกฺกฏฺฐญฺจ สิปฺปํ.\csromanspacer{} \textbf{หีนํ} นาม สิปฺปํ นฬการสิปฺปํ กุมฺภการสิปฺปํ เปสการสิปฺปํ จมฺมการสิปฺปํ นหาปิตสิปฺปํ, เตสุ เตสุ วา ปน ชนปเทสุ โอญฺญาตํ อวญฺญาตํ หีฬิตํ ปริภูตํ อจิตฺตีกตํ, เอตํ หีนํ นาม สิปฺปํ.\csromanspacer{} \textbf{อุกฺกฏฺฐํ} นาม สิปฺปํ มุทฺทา คณนา เลขา, เตสุ เตสุ วา ปน ชนปเทสุ อโนญฺญาตํ อนวญฺญาตํ อหีฬิตํ อปริภูตํ จิตฺตีกตํ, เอตํ อุกฺกฏฺฐํ นาม สิปฺปํ.}

\prose{สพฺเพปิ \textbf{อาพาธา} หีนา, อปิ จ มธุเมโห อาพาโธ อุกฺกฏฺโฐ.}

\prose{\textbf{ลิงฺคํ} นาม เทฺว ลิงฺคานิ หีนญฺจ ลิงฺคํ อุกฺกฏฺฐญฺจ ลิงฺคํ.\csromanspacer{} \textbf{หีนํ} นาม ลิงฺคํ อติทีฆํ อติรสฺสํ อติกณฺหํ อจฺโจทาตํ, เอตํ หีนํ นาม ลิงฺคํ.\csromanspacer{} \textbf{อุกฺกฏฺฐํ} นาม ลิงฺคํ นาติทีฆํ นาติรสฺสํ นาติกณฺหํ นาจฺโจทาตํ, เอตํ อุกฺกฏฺฐํ นาม ลิงฺคํ.}

\prose{สพฺเพปิ \textbf{กิเลสา} หีนา.}

\prose{สพฺพาปิ \textbf{อาปตฺติโย} หีนา, อปิ จ โสตาปตฺติสมาปตฺติ อุกฺกฏฺฐา.}

\prose{\textbf{อกฺโกโส} นาม เทฺว อกฺโกสา หีโน จ อกฺโกโส อุกฺกฏฺโฐ จ อกฺโกโส.\csromanspacer{} \textbf{หีโน} นาม อกฺโกโส โอฏฺโฐสิ เมณฺโฑสิ โคโณสิ คทฺรโภสิ ติรจฺฉานคโตสิ เนรยิโกสิ, นตฺถิ ตุยฺหํ สุคติ, ทุคฺคติ เยว ตุยฺหํ ปาฏิกงฺขาติ, ยกาเรน วา ภกาเรน วา กาฏโกฏจิกาย วา, เอโส หีโน นาม อกฺโกโส.\csromanspacer{} \textbf{อุกฺกฏฺโฐ} นาม อกฺโกโส ปณฺฑิโตสิ พฺยตฺโตสิ เมธาวีสิ พหุสฺสุโตสิ ธมฺมกถิโกสิ, นตฺถิ ตุยฺหํ ทุคฺคติ, สุคติ เยว ตุยฺหํ ปาฏิกงฺขาติ, เอโส อุกฺกฏฺโฐ นาม อกฺโกโส.}

\proseitem{16}{อุปสมฺปนฺโน อุปสมฺปนฺนํ ขุํเสตุกาโม วมฺเภตุกาโม มงฺกุกตฺตุกาโม หีเนน หีนํ วเทติ, จณฺฑาลํ เวนํ เนสาทํ รถการํ ปุกฺกุสํ “จณฺฑาโลสิ เวโนสิ เนสาโทสิ รถกาโรสิ ปุกฺกุโสสี”ติ ภณติ\footnote{วเทตีติ อุทฺเทโส. ภณตีติ วิตฺถาโร (วชิรพุทฺธิ)}, อาปตฺติ วาจาย วาจาย ปาจิตฺติยสฺส.}

\prose{\csromanfolio{10}อุปสมฺปนฺโน อุปสมฺปนฺนํ ขุํเสตุกาโม วมฺเภตุกาโม มงฺกุกตฺตุกาโม หีเนน อุกฺกฏฺฐํ วเทติ, ขตฺติยํ พฺราหฺมณํ “จณฺฑาโลสิ เวโนสิ เนสาโทสิ รถกาโรสิ ปุกฺกุโสสี”ติ ภณติ, อาปตฺติ วาจาย}

\prose{อุปสมฺปนฺโน อุปสมฺปนฺนํ ขุํเสตุกาโม วมฺเภตุกาโม มงฺกุกตฺตุกาโม อุกฺกฏฺเฐน หีนํ วเทติ, จณฺฑาลํ เวนํ เนสาทํ รถการํ ปุกฺกุสํ “ขตฺติโยสิ พฺราหฺมโณสี”ติ ภณติ, อาปตฺติ วาจาย}

\prose{อุปสมฺปนฺโน อุปสมฺปนฺนํ ขุํเสตุกาโม วมฺเภตุกาโม มงฺกุกตฺตุกาโม อุกฺกฏฺเฐน อุกฺกฏฺฐํ วเทติ, ขตฺติยํ พฺราหฺมณํ “ขตฺติโยสิ พฺราหฺมโณสี”ติ ภณติ, อาปตฺติ วาจาย วาจาย ปาจิตฺติยสฺส.}

\proseitem{17}{อุปสมฺปนฺโน อุปสมฺปนฺนํ ขุํเสตุกาโม วมฺเภตุกาโม มงฺกุกตฺตุกาโม หีเนน หีนํ วเทติ, อวกณฺณกํ ชวกณฺณกํ ธนิฏฺฐกํ สวิฏฺฐกํ กุลวฑฺฒกํ “อวกณฺณโกสิ ชวกณฺณโกสิ ธนิฏฺฐโกสิ สวิฏฺฐโกสิ กุลวฑฺฒโกสี”ติ ภณติ, อาปตฺติ วาจาย วาจาย}

\prose{อุปสมฺปนฺโน อุปสมฺปนฺนํ ขุํเสตุกาโม วมฺเภตุกาโม มงฺกุกตฺตุกาโม หีเนน อุกฺกฏฺฐํ วเทติ, พุทฺธรกฺขิตํ ธมฺมรกฺขิตํ สํฆรกฺขิตํ “อวกณฺณโกสิ ชวกณฺณโกสิ ธนิฏฺฐโกสิ สวิฏฺฐโกสิ กุลวฑฺฒโกสี”ติ ภณติ, อาปตฺติ วาจาย วาจาย}

\prose{อุปสมฺปนฺโน อุปสมฺปนฺนํ ขุํเสตุกาโม วมฺเภตุกาโม มงฺกุกตฺตุกาโม อุกฺกฏฺเฐน หีนํ วเทติ, อวกณฺณกํ ชวกณฺณกํ ธนิฏฺฐกํ สวิฏฺฐกํ กุลวฑฺฒกํ “พุทฺธรกฺขิโตสิ ธมฺมรกฺขิโตสิ สํฆรกฺขิโตสี”ติ ภณติ, อาปตฺติ วาจาย วาจาย}

\prose{อุปสมฺปนฺโน อุปสมฺปนฺนํ ขุํเสตุกาโม วมฺเภตุกาโม มงฺกุกตฺตุกาโม อุกฺกฏฺเฐน อุกฺกฏฺฐํ วเทติ, พุทฺธรกฺขิตํ ธมฺมรกฺขิตํ สํฆรกฺขิตํ “พุทฺธรกฺขิโตสิ ธมฺมรกฺขิโตสิ สํฆรกฺขิโตสี”ติ ภณติ, อาปตฺติ วาจาย วาจาย ปาจิตฺติยสฺส.}

\proseitem{18}{\csromanfolio{11}อุปสมฺปนฺโน อุปสมฺปนฺนํ ขุํเสตุกาโม วมฺเภตุกาโม มงฺกุกตฺตุกาโม หีเนน หีนํ วเทติ, โกสิยํ ภารทฺวาชํ “โกสิโยสิ ภารทฺวาโชสี”ติ ภณติ, อาปตฺติ วาจาย วาจาย ปาจิตฺติยสฺส.}

\prose{อุปสมฺปนฺโน อุปสมฺปนฺนํ ขุํเสตุกาโม วมฺเภตุกาโม มงฺกุกตฺตุกาโม หีเนน อุกฺกฏฺฐํ วเทติ, โคตมํ โมคฺคลฺลานํ กจฺจานํ วาสิฏฺฐํ “โกสิโยสิ ภารทฺวาโชสี”ติ ภณติ, อาปตฺติ วาจาย}

\prose{อุปสมฺปนฺโน อุปสมฺปนฺนํ ขุํเสตุกาโม วมฺเภตุกาโม มงฺกุกตฺตุกาโม อุกฺกฏฺเฐน หีนํ วเทติ, โกสิยํ ภารทฺวาชํ “โคตโมสิ โมคฺคลฺลาโนสิ กจฺจาโนสิ วาสิฏฺโฐสี”ติ ภณติ, อาปตฺติ วาจาย}

\prose{อุปสมฺปนฺโน อุปสมฺปนฺนํ ขุํเสตุกาโม วมฺเภตุกาโม มงฺกุกตฺตุกาโม อุกฺกฏฺเฐน อุกฺกฏฺฐํ วเทติ, โคตมํ โมคฺคลฺลานํ กจฺจานํ วาสิฏฺฐํ “โคตโมสิ โมคฺคลฺลาโนสิ กจฺจาโนสิ วาสิฏฺโฐสี”ติ ภณติ, อาปตฺติ วาจาย วาจาย ปาจิตฺติยสฺส.}

\proseitem{19}{อุปสมฺปนฺโน อุปสมฺปนฺนํ ขุํเสตุกาโม วมฺเภตุกาโม มงฺกุกตฺตุกาโม หีเนน หีนํ วเทติ, โกฏฺฐกํ ปุปฺผฉฑฺฑกํ “โกฏฺฐโกสิ ปุปฺผฉฑฺฑโกสี”ติ ภณติ, อาปตฺติ วาจาย วาจาย ปาจิตฺติยสฺส.}

\prose{อุปสมฺปนฺโน อุปสมฺปนฺนํ ขุํเสตุกาโม วมฺเภตุกาโม มงฺกุกตฺตุกาโม หีเนน อุกฺกฏฺฐํ วเทติ, กสฺสกํ วาณิชํ โครกฺขํ “โกฏฺฐโกสิ ปุปฺผฉฑฺฑโกสี”ติ ภณติ, อาปตฺติ วาจาย วาจาย ปาจิตฺติยสฺส.}

\prose{อุปสมฺปนฺโน อุปสมฺปนฺนํ ขุํเสตุกาโม วมฺเภตุกาโม มงฺกุกตฺตุกาโม อุกฺกฏฺเฐน หีนํ วเทติ, โกฏฺฐกํ ปุปฺผฉฑฺฑกํ “กสฺสโกสิ วาณิโชสิ โครกฺโขสี”ติ ภณติ, อาปตฺติ วาจาย วาจาย ปาจิตฺติยสฺส.}

\prose{อุปสมฺปนฺโน อุปสมฺปนฺนํ ขุํเสตุกาโม วมฺเภตุกาโม มงฺกุกตฺตุกาโม อุกฺกฏฺเฐน อุกฺกฏฺฐํ วเทติ, กสฺสกํ วาณิชํ โครกฺขํ “กสฺสโกสิ วาณิโชสิ โครกฺโขสี”ติ ภณติ, อาปตฺติ วาจาย วาจาย}

\proseitem{20}{\csromanfolio{12}อุปสมฺปนฺโน อุปสมฺปนฺนํ ขุํเสตุกาโม วมฺเภตุกาโม มงฺกุกตฺตุกาโม หีเนน หีนํ วเทติ, นฬการํ กุมฺภการํ เปสการํ จมฺมการํ นหาปิตํ “นฬกาโรสิ กุมฺภกาโรสิ เปสกาโรสิ จมฺมกาโรสิ นหาปิโตสี”ติ ภณติ, อาปตฺติ วาจาย วาจาย ปาจิตฺติยสฺส.}

\prose{อุปสมฺปนฺโน อุปสมฺปนฺนํ ขุํเสตุกาโม วมฺเภตุกาโม มงฺกุกตฺตุกาโม หีเนน อุกฺกฏฺฐํ วเทติ, มุทฺทิกํ คณกํ เลขกํ “นฬกาโรสิ กุมฺภกาโรสิ เปสกาโรสิ จมฺมกาโรสิ นหาปิโตสี”ติ ภณติ, อาปตฺติ วาจาย วาจาย ปาจิตฺติยสฺส.}

\prose{อุปสมฺปนฺโน อุปสมฺปนฺนํ ขุํเสตุกาโม วมฺเภตุกาโม มงฺกุกตฺตุกาโม อุกฺกฏฺเฐน หีนํ วเทติ, นฬการํ กุมฺภการํ เปสการํ จมฺมการํ นหาปิตํ “มุทฺทิโกสิ คณโกสิ เลขโกสี”ติ ภณติ, อาปตฺติ วาจาย วาจาย ปาจิตฺติยสฺส.}

\prose{อุปสมฺปนฺโน อุปสมฺปนฺนํ ขุํเสตุกาโม วมฺเภตุกาโม มงฺกุกตฺตุกาโม อุกฺกฏฺเฐน อุกฺกฏฺฐํ วเทติ, มุทฺทิกํ คณกํ เลขกํ “มุทฺทิโกสิ คณโกสิ เลขโกสี”ติ ภณติ, อาปตฺติ วาจาย วาจาย}

\proseitem{21}{อุปสมฺปนฺโน อุปสมฺปนฺนํ ขุํเสตุกาโม วมฺเภตุกาโม มงฺกุกตฺตุกาโม หีเนน หีนํ วเทติ, กุฏฺฐิกํ คณฺฑิกํ กิลาสิกํ โสสิกํ อปมาริกํ “กุฏฺฐิโกสิ คณฺฑิโกสิ กิลาสิโกสิ โสสิโกสิ อปมาริโกสี”ติ ภณติ, อาปตฺติ วาจาย วาจาย ปาจิตฺติยสฺส.}

\prose{อุปสมฺปนฺโน อุปสมฺปนฺนํ ขุํเสตุกาโม วมฺเภตุกาโม มงฺกุกตฺตุกาโม หีเนน อุกฺกฏฺฐํ วเทติ, มธุเมหิกํ “กุฏฺฐิโกสิ คณฺฑิโกสิ กิลาสิโกสิ โสสิโกสิ อปมาริโกสี”ติ ภณติ, อาปตฺติ วาจาย วาจาย}

\prose{อุปสมฺปนฺโน อุปสมฺปนฺนํ ขุํเสตุกาโม วมฺเภตุกาโม มงฺกุกตฺตุกาโม อุกฺกฏฺเฐน หีนํ วเทติ, กุฏฺฐิกํ คณฺฑิกํ กิลาสิกํ โสสิกํ อปมาริกํ “มธุเมหิโกสี”ติ ภณติ, อาปตฺติ วาจาย วาจาย}

\prose{\csromanfolio{13}อุปสมฺปนฺโน อุปสมฺปนฺนํ ขุํเสตุกาโม วมฺเภตุกาโม มงฺกุกตฺตุกาโม อุกฺกฏฺเฐน อุกฺกฏฺฐํ วเทติ, มธุเมหิกํ “มธุเมหิโกสี”ติ ภณติ, อาปตฺติ วาจาย วาจาย ปาจิตฺติยสฺส.}

\proseitem{22}{อุปสมฺปนฺโน อุปสมฺปนฺนํ ขุํเสตุกาโม วมฺเภตุกาโม มงฺกุกตฺตุกาโม หีเนน หีนํ วเทติ, อติทีฆํ อติรสฺสํ อติกณฺหํ อจฺโจทาตํ “อติทีโฆสิ อติรสฺโสสิ อติกโณฺหสิ อจฺโจทาโตสี”ติ ภณติ, อาปตฺติ วาจาย วาจาย ปาจิตฺติยสฺส.}

\prose{อุปสมฺปนฺโน อุปสมฺปนฺนํ ขุํเสตุกาโม วมฺเภตุกาโม มงฺกุกตฺตุกาโม หีเนน อุกฺกฏฺฐํ วเทติ, นาติทีฆํ นาติรสฺสํ นาติกณฺหํ นาจฺโจทาตํ “อติทีโฆสิ อติรสฺโสสิ อติกโณฺหสิ อจฺโจทาโตสี”ติ ภณติ, อาปตฺติ วาจาย วาจาย ปาจิตฺติยสฺส.}

\prose{อุปสมฺปนฺโน อุปสมฺปนฺนํ ขุํเสตุกาโม วมฺเภตุกาโม มงฺกุกตฺตุกาโม อุกฺกฏฺเฐน หีนํ วเทติ, อติทีฆํ อติรสฺสํ อติกณฺหํ อจฺโจทาตํ “นาติทีโฆสิ นาติรสฺโสสิ นาติกโณฺหสิ นาจฺโจทาโตสี”ติ ภณติ, อาปตฺติ วาจาย วาจาย ปาจิตฺติยสฺส.}

\prose{อุปสมฺปนฺโน อุปสมฺปนฺนํ ขุํเสตุกาโม วมฺเภตุกาโม มงฺกุกตฺตุกาโม อุกฺกฏฺเฐน อุกฺกฏฺฐํ วเทติ, นาติทีฆํ นาติรสฺสํ นาติกณฺหํ นาจฺโจทาตํ “นาติทีโฆสิ นาติรสฺโสสิ นาติกโณฺหสิ นาจฺโจทาโตสี”ติ ภณติ, อาปตฺติ วาจาย วาจาย ปาจิตฺติยสฺส.}

\proseitem{23}{อุปสมฺปนฺโน อุปสมฺปนฺนํ ขุํเสตุกาโม วมฺเภตุกาโม มงฺกุกตฺตุกาโม หีเนน หีนํ วเทติ, ราคปริยุฏฺฐิตํ โทสปริยุฏฺฐิตํ โมหปริยุฏฺฐิตํ “ราคปริยุฏฺฐิโตสิ โทสปริยุฏฺฐิโตสิ โมหปริยุฏฺฐิโตสี”ติ ภณติ, อาปตฺติ วาจาย วาจาย ปาจิตฺติยสฺส.}

\prose{อุปสมฺปนฺโน อุปสมฺปนฺนํ ขุํเสตุกาโม วมฺเภตุกาโม มงฺกุกตฺตุกาโม หีเนน อุกฺกฏฺฐํ วเทติ, วีตราคํ วีตโทสํ วีตโมหํ “ราคปริยุฏฺฐิโตสิ โทสปริยุฏฺฐิโตสิ โมหปติยุฏฺฐิโตสี”ติ ภณติ, อาปตฺติ วาจาย วาจาย ปาจิตฺติยสฺส.}

\prose{อุปสมฺปนฺโน อุปสมฺปนฺนํ ขุํเสภุกาโม วมฺเภตุกาโม มงฺกุกตฺตุกาโม อุกฺกฏฺเฐน หีนํ วเทติ, ราคปริยุฏฺฐิตํ โทสปริยุฏฺฐิตํ โมหปริยุฏฺฐิตํ \csromanfolio{14}“วีตราโคสิ วีตโทโสสิ วีตโมโหสี”ติ ภณติ, อาปตฺติ วาจาย}

\prose{อุปสมฺปนฺโน อุปสมฺปนฺนํ ขุํเสตุกาโม วมฺเภตุกาโม มงฺกุกตฺตุกาโม อุกฺกฏฺเฐน อุกฺกฏฺฐํ วเทติ, วีตราคํ วีตโทสํ วีตโมหํ “วีตราโคสิ วีตโทโสสิ วีตโมโหสี”ติ ภณติ, อาปตฺติ วาจาย}

\proseitem{24}{อุปสมฺปนฺโน อุปสมฺปนฺนํ ขุํเสตุกาโม วมฺเภตุกาโม มงฺกุกตฺตุกาโม หีเนน หีนํ วเทหิ, ปาราชิกํ อชฺฌาปนฺนํ สํฆาทิเสสํ อชฺฌาปนฺนํ ถุลฺลจฺจยํ อชฺฌาปนฺนํ ปาจิตฺติยํ อชฺฌาปนฺนํ ปาฏิเทสนียํ อชฺฌาปนฺนํ ทุกฺกฏํ อชฺฌาปนฺนํ ทุพฺภาสิตํ อชฺฌาปนฺนํ “ปาราชิกํ อชฺฌาปนฺโนสิ สํฆาทิเสสํ อชฺฌาปนฺโนสิ ถุลฺลจฺจยํ อชฺฌาปนฺโนสิ ปาจิตฺติยํ อชฺฌาปนฺโนสิ ปาฏิเทสนียํ อชฺฌาปนฺโนสิ ทุกฺกฏํ อชฺฌาปนฺโนสิ ทุพฺภาสิตํ อชฺฌาปนฺโนสี”ติ ภณติ, อาปตฺติ วาจาย วาจาย ปาจิตฺติยสฺส.}

\prose{อุปสมฺปนฺโน อุปสมฺปนฺนํ ขุํเสตุกาโม วมฺเภตุกาโม มงฺกุกตฺตุกาโม หีเนน อุกฺกฏฺฐํ วเทติ, โสตาปนฺนํ “ปาราชิกํ อชฺฌาปนฺโนสิ ฯเปฯ ทุพฺภาสิตํ อชฺฌาปนฺโนสี”ติ ภณติ, อาปตฺติ วาจาย วาจาย}

\prose{อุปสมฺปนฺโน อุปสมฺปนฺนํ ขุํเสตุกาโม วมฺเภตุกาโม มงฺกุกตฺตุกาโม อุกฺกฏฺเฐน หีนํ วเทติ, ปาราชิกํ อชฺฌาปนฺนํ ฯเปฯ ทุพฺภาสิตํ อชฺฌาปนฺนํ “โสตาปนฺโนสี”ติ ภณติ, อาปตฺติ วาจาย วาจาย}

\prose{อุปสมฺปนฺโน อุปสมฺปนฺนํ ขุํเสตุกาโม วมฺเภตุกาโม มงฺกุกตฺตุกาโม อุกฺกฏฺเฐน อุกฺกฏฺฐํ วเทติ, โสตาปนฺนํ “โสตาปนฺโนสี”ติ ภณติ, อาปตฺติ วาจาย วาจาย ปาจิตฺติยสฺส.}

\proseitem{25}{อุปสมฺปนฺโน อุปสมฺปนฺนํ ขุํเสตุกาโม วมฺเภตุกาโม มงฺกุกตฺตุกาโม หีเนน หีนํ วเทติ, โอฏฺฐํ เมณฺฑํ โคณํ คทฺรพํ ติรจฺฉานคตํ เนรยิกํ “โอฏฺโฐสิ เมณฺโฑสิ โคโณสิ คทฺรโภสิ ติรจฺฉานคโตสิ เนรยิโกสิ, นตฺถิ ตุยฺหํ สุคติ, ทุคฺคติ เยว ตุยฺหํ ปาฏิกงฺขา”ติ ภณติ, อาปตฺติ วาจาย วาจาย ปาจิตฺติยสฺส.}

\prose{\csromanfolio{15}อุปสมฺปนฺโน อุปสมฺปนฺนํ ขุํเสตุกาโม วมฺเภตุกาโม มงฺกุกตฺตุกาโม หีเนน อุกฺกฏฺฐํ วเทติ, ปณฺฑิตํ พฺยตฺตํ เมธาวิํ พหุสฺสุตํ ธมฺมกถิกํ “โอฏฺโฐสิ เมณฺโฑสิ โคโณสิ คทฺรโภสิ ติรจฺฉานคโตสิ เนรยิโกสิ, นตฺถิ ตุยฺหํ สุคติ, ทุคฺคติ เยว ตุยฺหํ ปาฏิกงฺขา”ติ ภณติ, อาปตฺติ วาจาย วาจาย ปาจิตฺติยสฺส.}

\prose{อุปสมฺปนฺโน อุปสมฺปนฺนํ ขุํเสตุกาโม วมฺเภตุกาโม มงฺกุกตฺตุกาโม อุกฺกฏฺเฐน หีนํ วเทติ, โอฏฺฐํ เมณฺฑํ โคณํ คทฺรภํ ติรจฺฉานคตํ เนรยิกํ “ปณฺฑิโตสิ พฺยตฺโตสิ เมธาวีสิ พหุสฺสุโตสิ ธมฺมกถิโกสิ, นตฺถิ ตุยฺหํ ทุคฺคติ, สุคติ เยว ตุยฺหํ ปาฏิกนฺขา”ติ ภณติ, อาปตฺติ วาจาย วาจาย ปาจิตฺติยสฺส.}

\prose{อุปสมฺปนฺโน อุปสมฺปนฺนํ ขุํเสตุกาโม วมฺเภตุกาโม มงฺกุกตฺตุกาโม อุกฺกฏฺฐน อุกฺกฏฺฐํ วเทติ, ปณฺฑิตํ พฺยตฺตํ เมธาวิํ พหุสฺสุตํ ธมฺมกถิกํ “ปณฺฑิโสสิ พฺยตฺโตสิ เมธาวีสิ พหุสฺสุโตสิ ธมฺมกถิโกสิ, นตฺถิ ตุยฺหํ ทุคฺคติ, สุคติ เยว ตุยฺหํ ปาฏิกงฺขา”ติ ภณติ, อาปตฺติ วาจาย วาจาย ปาจิตฺติยสฺส.}

\proseitem{26}{อุปสมฺปนฺโน อุปสมฺปนฺนํ ขุํเสตุกาโม วมฺเภตุกาโม มงฺกุกตฺตุกาโม เอวํ วเทติ, “สนฺติ อิเธกจฺเจ จณฺฑาลา เวนา เนสาทา รถการา ปุกฺกุสา”ติ ภณติ, อาปตฺติ วาจาย วาจาย ทุกฺกฏสฺส.}

\prose{อุปสมฺปนฺโน อุปสมฺปนฺนํ ขุํเสตุกาโม วมฺเภตุกาโม มงฺกุกตฺตุกาโม เอวํ วเทติ, “สนฺติ อิเธกจฺเจ ขตฺติยา พฺราหฺมณา”ติ ภณติ, อาปตฺติ วาจาย วาจาย ทุกฺกฏสฺส.}

\proseitem{27}{อุปสมฺปนฺโน อุปสมฺปนฺนํ ขุํเสตุกาโม วมฺเภตุกาโม มงฺกุกตฺตุกาโม เอวํ วเทติ, “สนฺติ อิเธกจฺเจ อวกณฺณกา ชวกณฺณกา ธนิฏฺฐกา สวิฏฺฐกา กุลวฑฺฒกา”ติ ภณติ ฯเปฯ “สนฺติ อิเธกจฺเจ พุทฺธรกฺขิตา ธมฺมรกฺขิตา สํฆรกฺขิตา”ติ ภณติ ฯเปฯ “สนฺติ อิเธกจฺเจ โกสิยา ภารทฺวาชา”ติ ภณติ ฯเปฯ “สนฺติ อิเธกจฺเจ โคตมา โมคฺคลฺลานา กจฺจานา วาสิฏฺฐา”ติ ภณติ ฯเปฯ “สนฺติ อิเธกจฺเจ โกฏฺฐกา ปุปฺผฉฑฺฑกา”ติ ภณติ ฯเปฯ “สนฺติ อิเธกจฺเจ กสฺสกา วาณิชา โครกฺขา”ติ ภณติ ฯเปฯ “สนฺติ อิเธกจฺเจ นฬการา \csromanfolio{16}กุมฺภการา เปสการา จมฺมการา นหาปิตา”ติ ภณติ ฯเปฯ “สนฺติ อิเธกจฺเจ มุทฺทิกา คณกา เลขกา”ติ ภณติ ฯเปฯ “สนฺติ อิเธกจฺเจ กุฏฺฐิกา คณฺฑิกา กิลาสิกา โสสิกา อปมาริกา”ติ ภณติ ฯเปฯ “สนฺติ อิเธกจฺเจ มธุเมหิกา”ติ ภณติ ฯเปฯ “สนฺติ อิเธกจฺเจ อติทีฆา อติรสฺสา อติกณฺหา อจฺโจทาตา”ติ ภณติ ฯเปฯ “สนฺติ อิเธกจฺเจ นาติทีฆา นาติรสฺสา นาติกณฺหา นาจฺโจทาตา”ติ ภณติ ฯเปฯ “สนฺติ อิเธกจฺเจ ราคปริยุฏฺฐิตา โทสปริยุฏฺฐิตา โมหปริยุฏฺฐิตา”ติ ภณติ ฯเปฯ “สนฺติ อิเธกจฺเจ วีตราคา วีตโทสา วีตโมหา”ติ ภณติ ฯเปฯ “สนฺติ อิเธกจฺเจ ปาราชิกํ อชฺฌาปนฺนา ฯเปฯ ทุพฺภาสิตํ อชฺฌาปนฺนา”ติ ภณติ ฯเปฯ “สนฺติ อิเธกจฺเจ โสตาปนฺนา”ติ ภณติ ฯเปฯ “สนฺติ อิเธกจฺเจ โอฏฺฐา เมณฺฑา โคณา คทฺรภา ติรจฺฉานคตา เนรยิกา, นตฺถิ เตสํ สุคติ, ทุคฺคติ เยว เตสํ ปาฏิกงฺขา”ติ ภณติ, อาปตฺติ วาจาย วาจาย ทุกฺกฏสฺส.}

\proseitem{28}{อุปสมฺปนฺโน อุปสมฺปนฺนํ ขุํเสตุกาโม วมฺเภตุกาโม มงฺกุกตฺตุกาโม เอวํ วเทติ, “สนฺติ อิเธกจฺเจ ปณฺฑิตา พฺยตฺตา เมธาวี พหุสฺสุตา ธมฺมกถิกา, นตฺถิ เตสํ ทุคฺคติ, สุคติ เยว เตสํ ปาฏิกงฺขา”ติ ภณติ, อาปตฺติ วาจาย วาจาย ทุกฺกฏสฺส.}

\proseitem{29}{อุปสมฺปนฺโน อุปสมฺปนฺนํ ขุํเสตุกาโม วมฺเภตุกาโม มงฺกุกตฺตุกาโม เอวํ วเทติ, “เย นูน จณฺฑาลา เวนา เนสาทา รถการา ปุกฺกุสา”ติ ภณติ, อาปตฺติ วาจาย วาจาย ทุกฺกฏสฺส ฯเปฯ.}

\prose{อุปสมฺปนฺโน อุปสมฺปนฺนํ ขุํเสตุกาโม วมฺเภตุกาโม มงฺกุกตฺตุกาโม เอวํ วเทติ, “เย นูน ปณฺฑิตา พฺยตฺตา เมธาวี พหุสฺสุตา ธมฺมกถิกา”ติ ภณติ, อาปตฺติ วาจาย วาจาย ทุกฺกฏสฺส.}

\proseitem{30}{อุปสมฺปนฺโน อุปสมฺปนฺนํ ขุํเสตุกาโม วมฺเภตุกาโม มงฺกุกตฺตุกาโม เอวํ วเทติ, “น มยํ จณฺฑาลา เวนา เนสาทา รถการา ปุกฺกุสา”ติ ภณติ ฯเปฯ “น มยํ ปณฺฑิตา พฺยตฺตา เมธาวี พหุสฺสุตา ธมฺมกถิกา, นตฺถมฺหากํ ทุคฺคติ, สุคติ เยว อมฺหากํ ปาฏิกงฺขา”ติ ภณติ, อาปตฺติ วาจาย วาจาย ทุกฺกฏสฺส.}

\proseitem{31}{\csromanfolio{17}อุปสมฺปนฺโน อนุปสมฺปนฺนํ ขุํเสตุกาโม วมฺเภตุกาโม มงฺกุกตฺตุกาโม หีเนน หีนํ วเทติ ฯเปฯ หีเนน อุกฺกฏฺฐํ วเทติ ฯเปฯ อุกฺกฏฺเฐน หีนํ วเทติ ฯเปฯ อุกฺกฏฺเฐน อุกฺกฏฺฐํ วเทติ, ปณฺฑิตํ พฺยตฺตํ เมธาวิํ พหุสฺสุตํ ธมฺมกถิกํ “ปณฺฑิโตสิ พฺยตฺโตสิ เมธาวีสิ พหุสฺสุโตสิ ธมฺมกถิโกสิ, นตฺถิ ตุยฺหํ ทุคฺคติ, สุคติ เยว ตุยฺหํ ปาฏิกงฺขา”ติ ภณติ, อาปตฺติ วาจาย วาจาย ทุกฺกฏสฺส.}

\prose{อุปสมฺปนฺโน อนุปสมฺปนฺนํ ขุํเสตุกาโม วมฺเภตุกาโม มงฺกุกตฺตุกาโม เอวํ วเทติ, “สนฺติ อิเธกจฺเจ จณฺฑาลา เวนา เนสาทา รถการา ปุกฺกุสา”ติ ภณติ ฯเปฯ “สนฺติ อิเธกจฺเจ ปณฺฑิตา พฺยตฺตา เมธาวี พหุสฺสุตา ธมฺมกถิกา, นตฺถิ เตสํ ทุคฺคติ, สุคติ เยว เตสํ ปาฏิกงฺขา”ติ ภณติ, อาปตฺติ วาจาย วาจาย ทุกฺกฏสฺส.}

\prose{อุปสมฺปนฺโน อนุปสมฺปนฺนํ ขุํเสตุกาโม วมฺเภตุกาโม มงฺกุกตฺตุกาโม เอวํ วเทติ, “เย นูน จณฺฑาลา เวนา เนสาทา รถการา ปุกฺกุสา”ติ ภณติ ฯเปฯ “เย นูน ปณฺฑิตา พฺยตฺตา เมธาวี พหุสฺสุตา ธมฺมกถิกา”ติ ภณติ, อาปตฺติ วาจาย วาจาย ทุกฺกฏสฺส.}

\prose{อุปสมฺปนฺโน อนุปสมฺปนฺนํ ขุํเสตุกาโม วมฺเภตุกาโม มงฺกุกตฺตุกาโม เอวํ วเทติ, “น มยํ จณฺฑาลา เวนา เนสาทา รถการา ปุกฺกุสา”ติ ภณติ ฯเปฯ “น มยํ ปณฺฑิตา พฺยตฺตา เมธาวี พหุสฺสุตา ธมฺมกถิกา, นตฺถมฺหากํ ทุคฺคติ, สุคติ เยว อมฺหากํ ปาฏิกงฺขา”ติ ภณติ, อาปตฺติ วาจาย วาจาย ทุกฺกฏสฺส.}

\proseitem{32}{อุปสมฺปนฺโน อุปสมฺปนฺนํ น ขุํเสตุกาโม น วมฺเภตุกาโม น มงฺกุกตฺตุกาโม ทวกมฺยตา หีเนน หีนํ วเทติ, จณฺฑาลํ เวนํ เนสาทํ รถการํ ปุกฺกุสํ “จณฺฑาโลสิ เวโนสิ เนสาโทสิ รถกาโรสิ ปุกฺกุโสสี”ติ ภณติ, อาปตฺติ วาจาย วาจาย ทุพฺภาสิตสฺส.}

\prose{อุปสมฺปนฺโน อุปสมฺปนฺนํ น ขุํเสตุกาโม น วมฺเภตุกาโม น มงฺกุกตฺตุกาโม ทวกมฺยตา หีเนน อุกฺกฏฺฐํ วเทติ, ขตฺติยํ พฺราหฺมณํ “จณฺฑาโลสิ เวโนสิ เนสาโทสิ รถกาโรสิ ปุกฺกุโสสี”ติ ภณติ, อาปตฺติ วาจาย วาจาย ทุพฺภาสิตสฺส.}

\prose{\csromanfolio{18}อุปสมฺปนฺโน อุปสมฺปนฺนํ น ขุํเสตุกาโม น วมฺเภตุกาโม น มงฺกุกตฺตุกาโม ทวกมฺยตา อุกฺกฏฺเฐน หีนํ วเทติ, จณฺฑาลํ เวนํ เนสาทํ รถการํ ปุกฺกุสํ “ขตฺติโยสิ พฺราหฺมโณสี”ติ ภณติ อาปตฺติ วาจาย วาจาย ทุพฺภาสิตสฺส.}

\prose{อุปสมฺปนฺโน อุปสมฺปนฺนํ น ขุํเสตุกาโม น วมฺเภตุกาโม น มงฺกุกตฺตุกาโม ทวกมฺยตา อุกฺกฏฺเฐน อุกฺกฏฺฐํ วเทติ, ขตฺติยํ พฺราหฺมณํ “ขตฺติโยสิ พฺราหฺมโณสี”ติ ภณติ, อาปตฺติ วาจาย วาจาย ทุพฺภาสิตสฺส.}

\prose{อุปสมฺปนฺโน อุปสมฺปนฺนํ น ขุํเสตุกาโม น วมฺเภตุกาโม น มงฺกุกตฺตุกาโม ทวกมฺยตา หีเนน หีนํ วเทติ ฯเปฯ หีเนน อุกฺกฏฺฐํ วเทติ ฯเปฯ อุกฺกฏฺเฐน หีนํ วเทติ ฯเปฯ อุกฺกฏฺเฐน อุกฺกฏฺฐํ วเทติ, ปณฺฑิตํ พฺยตฺตํ เมธาวิํ พหุสฺสุตํ ธมฺมกถิกํ “ปณฺฑิโตสิ พฺยตฺโตสิ เมธาวีสิ พหุสฺสุโตสิ ธมฺมกถิโกสิ, นตฺถิ ตุยฺหํ ทุคฺคติ, สุคติ เยว ตุยฺหํ ปาฏิกงฺขา”ติ ภณติ, อาปตฺติ วาจาย วาจาย ทุพฺภาสิตสฺส.}

\proseitem{33}{อุปสมฺปนฺโน อุปสมฺปนฺนํ น ขุํเสตุกาโม น วมฺเภตุกาโม น มงฺกุกตฺตุกาโม ทวกมฺยตา เอวํ วเทติ, “สนฺติ อิเธกจฺเจ จณฺฑาลา เวนา เนสาทา รถการา ปุกฺกุสา”ติ ภณติ ฯเปฯ “สนฺติ อิเธกจฺเจ ปณฺฑิตา พฺยตฺตา เมธาวี พหุสฺสุตา ธมฺมกถิกา, นตฺถิ เตสํ ทุคฺคติ, สุคติ เยว เตสํ ปาฏิกงฺขา”ติ ภณติ, อาปตฺติ วาจาย วาจาย ทุพฺภาสิตสฺส.}

\prose{อุปสมฺปนฺโน อุปสมฺปนฺนํ น ขุํเสตุกาโม น วมฺเภตุกาโม น มงฺกุกตฺตุกาโม ทวกมฺยตา เอวํ วเทติ, “เย นูน จณฺฑาลา เวนา เนสาทา รถการา ปุกฺกุสา”ติ ภณติ ฯเปฯ “เย นูน ปณฺฑิตา พฺยตฺตา เมธาวี พหุสฺสุตา ธมฺมกถิกา”ติ ภณติ, อาปตฺติ วาจาย วาจาย ทุพฺภาสิตสฺส.}

\prose{อุปสมฺปนฺโน อุปสมฺปนฺนํ น ขุํเสตุกาโม น วมฺเภตุกาโม น มงฺกุกตฺตุกาโม ทวกมฺยตา เอวํ วเทติ, “น มยํ จณฺฑาลา เวนา เนสาทา รถการา ปุกฺกุสา”ติ ภณติ ฯเปฯ “น มยํ ปณฺฑิตา พฺยตฺตา เมธาวี พหุสฺสุตา ธมฺมกถิกา, นตฺถมฺหากํ ทุคฺคติ, สุคติ เยว อมฺหากํ ปาฏิกงฺขา”ติ ภณติ, อาปตฺติ วาจาย วาจาย ทุพฺภาสิตสฺส.}

\proseitem{34}{\csromanfolio{19}อุปสมฺปนฺโน อนุปสมฺปนฺนํ น ขุํเสตุกาโม น วมฺเภตุกาโม น มงฺกุกตฺตุกาโม ทวกมฺยตา หีเนน หีนํ วเทติ ฯเปฯ หีเนน อุกฺกฏฺฐํ วเทติ ฯเปฯ อุกฺกฏฺเฐน หีนํ วเทติ ฯเปฯ อุกฺกฏฺเฐน อุกฺกฏฺฐํ วเทติ, ปณฺฑิตํ พฺยตฺตํ เมธาวิํ พหุสฺสุตํ ธมฺมกถิกํ “ปณฺฑิโตสิ พฺยตฺโตสิ เมธาวีสิ พหุสฺสุโตสิ ธมฺมกถิโกสิ, นตฺถิ ตุยฺหํ ทุคฺคติ, สุคติ เยว ตุยฺหํ ปาฏิกงฺขา”ติ ภณติ, อาปตฺติ วาจาย วาจาย ทุพฺภาสิตสฺส.}

\prose{อุปสมฺปนฺโน อนุปสมฺปนฺนํ น ขุํเสตุกาโม น วมฺเภตุกาโม น มงฺกุกตฺตุกาโม ทวกมฺยตา เอวํ วเทติ, “สนฺติ อิเธกจฺเจ จณฺฑาลา เวนา เนสาทา รถการา ปุกฺกุสา”ติ ภณติ ฯเปฯ “สนฺติ อิเธกจฺเจ ปณฺฑิตา พฺยตฺตา เมธาวี พหุสฺสุตา ธมฺมกถิกา, นตฺถิ เตสํ ทุคฺคติ, สุคติ เยว เตสํ ปาฏิกงฺขา”ติ ภณติ, อาปตฺติ วาจาย วาจาย ทุพฺภาสิตสฺส.}

\prose{อุปสมฺปนฺโน อนุปสมฺปนฺนํ น ขุํเสตุกาโม น วมฺเภตุกาโม น มงฺกุกตฺตุกาโม ทวกมฺยตา เอวํ วเทติ, “เย นูน จณฺฑาลา เวนา เนสาทา รถการา ปุกฺกุสา”ติ ภณติ ฯเปฯ “เย นูน ปณฺฑิตา พฺยตฺตา เมธาวี พหุสฺสุตา ธมฺมกถิกา”ติ ภณติ, อาปตฺติ วาจาย วาจาย ทุพฺภาสิตสฺส.}

\prose{อุปสมฺปนฺโน อนุปสมฺปนฺนํ น ขุํเสตุกาโม น วมฺเภตุกาโม น มงฺกุกตฺตุกาโม ทวกมฺยตา เอวํ วเทติ, “น มยํ จณฺฑาลา เวนา เนสาทา รถการา ปุกฺกุสา”ติ ภณติ ฯเปฯ “น มยํ ปณฺฑิตา พฺยตฺตา เมธาวี พหุสฺสุตา ธมฺมกถิกา, นตฺถมฺหากํ ทุคฺคติ, สุคติ เยว อมฺหากํ ปาฏิกงฺขา”ติ ภณติ, อาปตฺติ วาจาย วาจาย ทุพฺภาสิตสฺส.}

\proseitemwithcloserruled{35}{อนาปตฺติ อตฺถปุเรกฺขารสฺส ธมฺมปุเรกฺขารสฺส อนุสาสนิ\-ปุเรกฺขารสฺส อุมฺมตฺตกสฺส ขิตฺตจิตฺตสฺส เวทนาฏฺฏสฺส อาทิกมฺมิกสฺสาติ.}{โอมสวาทสิกฺขาปทํ นิฏฺฐิตํ ทุติยํ.}
\csromanmatikaanchor{mtk.5}
\csromantocmarkref{subsection}{3. เปสุญฺญสิกฺขาปท}{mtk.5}
\bookmark[dest={mtk.5},level=2]{3. เปสุญฺญสิกฺขาปท}
\csromanheader{2}{1. \textbf{มุสาวาทวคฺค 3. เปสุญฺญสิกฺขาปท}}
\proseitem{36}{เตน สมเยน พุทฺโธ ภควา สาวตฺถิยํ วิหรติ เชตวเน อนาถปิณฺฑิกสฺส อาราเม.\csromanspacer{} เตน โข ปน สมเยน ฉพฺพคฺคิยา ภิกฺขู \csromanfolio{20}ภิกฺขูนํ ภณฺฑนชาตานํ กลหชาตานํ วิวาทาปนฺนานํ เปสุญฺญํ อุปสํหรนฺติ, อิมสฺส สุตฺวา อมุสฺส อกฺขายนฺติ อิมสฺส เภทาย, อมุสฺส สุตฺวา อิมสฺส อกฺขายนฺติ อมุสฺส เภทาย.\csromanspacer{} เตน อนุปฺปนฺนานิ เจว ภณฺฑนานิ อุปฺปชฺชนฺติ, อุปฺปนฺนานิ จ ภณฺฑนานิ ภิยฺโยภาวาย เวปุลฺลาย สํวตฺตนฺติ.\csromanspacer{} เย เต ภิกฺขู อปฺปิจฺฉา ฯเปฯ เต อุชฺฌายนฺติ ขิยฺยนฺติ วิปาเจนฺติ “กถํ หิ นาม ฉพฺพคฺคิยา ภิกฺขู ภิกฺขูนํ ภณฺฑนชาตานํ กลหชาตานํ วิวาทาปนฺนานํ เปสุญฺญํ อุปสํหริสฺสนฺติ, อิมสฺส สุตฺวา อมุสฺส อกฺขายิสฺสนฺติ อิมสฺส เภทาย, อมุสฺส สุตฺวา อิมสฺส อกฺขายิสฺสนฺติ อมุสฺส เภทาย.\csromanspacer{} เตน อนุปฺปนฺนานิ เจว ภณฺฑนานิ อุปฺปชฺชนฺติ, อุปฺปนฺนานิ จ ภณฺฑนานิ ภิยฺโยภาวาย เวปุลฺลาย สํวตฺตนฺตี”ติ.\csromanspacer{} อถ โข เต ภิกฺขู ฉพฺพคฺคิเย ภิกฺขู อเนกปริยาเยน วิครหิตฺวา ภควโต เอตมตฺถํ อาโรเจสุํ ฯเปฯ “สจฺจํ กิร ตุเมฺห ภิกฺขเว ภิกฺขูนํ ภณฺฑนชาตานํ กลหชาตานํ วิวาทาปนฺนานํ เปสุญฺญํ อุปสํหรถ, อิมสฺส สุตฺวา อมุสฺส อกฺขายถ อิมสฺส เภทาย, อมุสฺส สุตฺวา อิมสฺส อกฺขายถ อมุสฺส เภทาย.\csromanspacer{} เตน อนุปฺปนฺนานิ เจว ภณฺฑนานิ อุปฺปชฺชนฺติ, อุปนฺนานิ จ ภณฺฑนานิ ภิยฺโยภาวาย เวปุลฺลาย สํวตฺตนฺตี”ติ.\csromanspacer{} สจฺจํ ภควาติ.\csromanspacer{} วิครหิ พุทฺโธ ภควา ฯเปฯ.\csromanspacer{} กถํ หิ นาม ตุเมฺห โมฆปุริสา ภิกฺขูนํ ภณฺฑนชาตานํ กลหชาตานํ วิวาทาปนฺนานํ เปสุญฺญํ อุปสํหริสฺสถ, อิมสฺส สุตฺวา อมุสฺส อกฺขายิสฺสถ อิมสฺส เภทาย, อมุสฺส สุตฺวา อิมสฺส อกฺขายิสฺสถ อมุสฺส เภทาย, เตน อนุปฺปนฺนานิ เจว ภณฺฑนานิ อุปฺปชฺชนฺติ, อุปฺปนฺนานิ จ ภณฺฑนานิ ภิยฺโยภาวาย เวปุลฺลาย สํวตฺตนฺติ.\csromanspacer{} เนตํ โมฆปุริสา อปฺปสนฺนานํ วา ปสาทาย.\csromanspacer{} ปสนฺนานํ วา ภิยฺโยภาวาย ฯเปฯ.\csromanspacer{} เอวญฺจ ปน ภิกฺขเว อิมํ สิกฺขาปทํ อุทฺทิเสยฺยาถ\csromandash{}}

\proseitem{37}{\textbf{“ภิกฺขุเปสุญฺเญ ปาจิตฺติยนฺ”}ติ.}

\proseitem{38}{\textbf{เปสุญฺญํ} นาม ทฺวีหากาเรหิ เปสุญฺญํ โหติ ปิยกมฺยสฺส วา เภทาธิปฺปายสฺส วา, ทสหากาเรหิ เปสุญฺญํ อุปสํหรติ ชาติโตปิ นามโตปิ โคตฺตโตปิ กมฺมโตปิ สิปฺปโตปิ อาพาธโตปิ ลิงฺคโตปิ กิเลสโตปิ อาปตฺติโตปิ อกฺโกสโตปิ.}

\prose{\textbf{ชาติ} นาม เทฺว ชาติโย หีนา จ ชาติ อุกฺกฏฺฐา จ ชาติ.\csromanspacer{} \textbf{หีนา} นาม ชาติ จณฺฑาลชาติ เวนชาติ เนสาทชาติ รถการชาติ \csromanfolio{21}ปุกฺกุสชาติ, เอสา หีนา นาม ชาติ.\csromanspacer{} \textbf{อุกฺกฏฺฐา} นาม ชาติ ขตฺติยชาติ พฺราหฺมณชาติ, เอสา อุกฺกฏฺฐา นาม ชาติ ฯเปฯ.}

\prose{\textbf{อกฺโกโส} นาม เทฺว อกฺโกสา หีโน จ อกฺโกโส อุกฺกฏฺโฐ จ อกฺโกโส.\csromanspacer{} \textbf{หีโน} นาม อกฺโกโส โอฏฺโฐสิ เมณฺโฑสิ โคโณสิ คทฺรโภสิ ติรจฺฉานคโตสิ เนรยิโกสิ, นตฺถิ ตุยฺหํ สุคติ, ทุคฺคติ เยว ตุยฺหํ ปาฏิกงฺขาติ, ยกาเรน วา ภกาเรน วา กาฏโกฏจิกาย วา, เอโส หีโน นาม อกฺโกโส.\csromanspacer{} \textbf{อุกฺกฏฺโฐ} นาม อกฺโกโส ปณฺฑิโตสิ พฺยตฺโตสิ เมธาวีสิ พหุสฺสุโตสิ ธมฺมกถิโกสิ, นตฺถิ ตุยฺหํ ทุคฺคติ, สุคติ เยว ตุยฺหํ ปาฏิกงฺขาติ, เอโส อุกฺกฏฺโฐ นาม อกฺโกโส.}

\proseitem{39}{อุปสมฺปนฺโน อุปสมฺปนฺนสฺส สุตฺวา อุปสมฺปนฺนสฺส เปสุญฺญํ อุปสํหรติ “อิตฺถนฺนาโม ตํ ‘จณฺฑาโล เวโน เนสาโท รถกาโร ปุกฺกุโส’ติ ภณตี”ติ, อาปตฺติ วาจาย วาจาย ปาจิตฺติยสฺส.}

\prose{อุปสมฺปนฺโน อุปสมฺปนฺนสฺส สุตฺวา อุปสมฺปนฺนสฺส เปสุญฺญํ อุปสํหรติ “อิตฺถนฺนาโม ตํ ‘ขตฺติโย พฺราหฺมโณ’ติ ภณตี”ติ, อาปตฺติ วาจาย วาจาย ปาจิตฺติยสฺส.}

\prose{อุปสมฺปนฺโน อุปสมฺปนฺนสฺส สุตฺวา อุปสมฺปนฺนสฺส เปสุญฺญํ อุปสํหรติ “อิตฺถนฺนาโม ตํ ‘อวกณฺณโก ชวกณฺณโก ธนิฏฺฐโก สวิฏฺฐโก กุลวฑฺฒโก’ติ ภณตี”ติ, อาปตฺติ วาจาย วาจาย ปาจิตฺติยสฺส.}

\prose{อุปสมฺปนฺโน อุปสมฺปนฺนสฺส สุตฺวา อุปสมฺปนฺนสฺส เปสุญฺญํ อุปสํหรติ “อิตฺถนฺนาโม ตํ ‘พุทฺธรกฺขิโต ธมฺมรกฺขิโต สํฆรกฺขิโต’ติ ภณตี”ติ, อาปตฺติ วาจาย วาจาย ปาจิตฺติยสฺส.}

\prose{อุปสมฺปนฺโน อุปสมฺปนฺนสฺส สุตฺวา อุปสมฺปนฺนสฺส เปสุญฺญํ อุปสํหรติ “อิตฺถนฺนาโม ตํ ‘โกสิโย ภารทฺวาโช’ติ ภณตี”ติ, อาปตฺติ วาจาย วาจาย ปาจิตฺติยสฺส.}

\prose{อุปสมฺปนฺโน อุปสมฺปนฺนสฺส สุตฺวา อุปสมฺปนฺนสฺส เปสุญฺญํ อุปสํหรติ “อิตฺถนฺนาโม ตํ ‘โคตโม โมคฺคลฺลาโน กจฺจาโน วาสิฏฺโฐ’ติ ภณตี”ติ, อาปตฺติ วาจาย วาจาย ปาจิตฺติยสฺส.}

\prose{\csromanfolio{22}อุปสมฺปนฺโน อุปสมฺปนฺนสฺส สุตฺวา อุปสมฺปนฺนสฺส เปสุญฺญํ อุปสํหรติ “อิตฺถนฺนาโม ตํ ‘โกฏฺฐโก ปุปฺผฉฑฺฑโก’ติ ภณตี”ติ, อาปตฺติ วาจาย วาจาย ปาจิตฺติยสฺส.}

\prose{อุปสมฺปนฺโน อุปสมฺปนฺนสฺส สุตฺวา อุปสมฺปนฺนสฺส เปสุญฺญํ อุปสํหรติ “อิตฺถนฺนาโม ตํ ‘กสฺสโก วาณิโช โครกฺโข’ติ ภณตี”ติ, อาปตฺติ วาจาย วาจาย ปาจิตฺติยสฺส.}

\prose{อุปสมฺปนฺโน อุปสมฺปนฺนสฺส สุตฺวา อุปสมฺปนฺนสฺส เปสุญฺญํ อุปสํหรติ “อิตฺถนฺนาโม ตํ ‘นฬกาโร กุมฺภกาโร เปสกาโร จมฺมกาโร นหาปิโต’ติ ภณตี”ติ, อาปตฺติ วาจาย วาจาย ปาจิตฺติยสฺส.}

\prose{อุปสมฺปนฺโน อุปสมฺปนฺนสฺส สุตฺวา อุปสมฺปนฺนสฺส เปสุญฺญํ อุปสํหรติ “อิตฺถนฺนาโม ตํ ‘มุทฺทิโก คณโก เลขโก’ติ ภณตี”ติ, อาปตฺติ วาจาย วาจาย ปาจิตฺติยสฺส.}

\proseitem{40}{อุปสมฺปนฺโน อุปสมฺปนฺนสฺส สุตฺวา อุปสมฺปนฺนสฺส เปสุญฺญํ อุปสํหรติ “อิตฺถนฺนาโม ตํ ‘กุฏฺฐิโก คณฺฑิโก กิลาสิโก โสสิโก อปมาริโก’ติ ภณตี”ติ, อาปตฺติ วาจาย วาจาย ปาจิตฺติยสฺส.}

\prose{อุปสมฺปนฺโน อุปสมฺปนฺนสฺส สุตฺวา อุปสมฺปนฺนสฺส เปสุญฺญํ อุปสํหรติ “อิตฺถนฺนาโม ตํ ‘มธุเมหิโก’ติ ภณตี”ติ, อาปตฺติ วาจาย}

\prose{อุปสมฺปนฺโน อุปสมฺปนฺนสฺส สุตฺวา อุปสมฺปนฺนสฺส เปสุญฺญํ อุปสํหรติ “อิตฺถนฺนาโม ตํ ‘อติทีโฆ อติรสฺโส อติกโณฺห อจฺโจทาโต’ติ ภณตี”ติ, อาปตฺติ วาจาย วาจาย ปาจิตฺติยสฺส.}

\prose{อุปสมฺปนฺโน อุปสมฺปนฺนสฺส สุตฺวา อุปสมฺปนฺนสฺส เปสุญฺญํ อุปสํหรติ “อิตฺถนฺนาโม ตํ ‘นาติทีโฆ นาติรสฺโส นาติกโณฺห นาจฺโจทาโต’ติ ภณตี”ติ, อาปตฺติ วาจาย วาจาย ปาจิตฺติยสฺส.}

\prose{อุปสมฺปนฺโน อุปสมฺปนฺนสฺส สุตฺวา อุปสมฺปนฺนสฺส เปสุญฺญํ อุปสํหรติ “อิตฺถนฺนาโม ตํ ‘ราคปริยุฏฺฐิโต โทสปริยุฏฺฐิโต โมหปริยุฏฺฐิโต’ติ ภณตี”ติ, อาปตฺติ วาจาย วาจาย ปาจิตฺติยสฺส.}

\prose{\csromanfolio{23}อุปสมฺปนฺโน อุปสมฺปนฺนสฺส สุตฺวา อุปสมฺปนฺนสฺส เปสุญฺญํ อุปสํหรติ “อิตฺถนฺนาโม ตํ ‘วีตราโค วีตโทโส วีตโมโห’ติ ภณตี”ติ, อาปตฺติ วาจาย วาจาย ปาจิตฺติยสฺส.}

\prose{อุปสมฺปนฺโน อุปสมฺปนฺนสฺส สุตฺวา อุปสมฺปนฺนสฺส เปสุญฺญํ อุปสํหรติ “อิตฺถนฺนาโม ตํ ‘ปาราชิกํ อชฺฌาปนฺโน, สํฆาทิเสสํ อชฺฌาปนฺโน, ถุลฺลจฺจยํ อชฺฌาปนฺโน, ปาจิตฺติยํ อชฺฌาปนฺโน, ปาฏิเทสนียํ อชฺฌาปนฺโน, ทุกฺกฏํ อชฺฌาปนฺโน, ทุพฺภาสิตํ อชฺฌาปนฺโน’ติ ภณตี”ติ, อาปตฺติ วาจาย วาจาย ปาจิตฺติยสฺส.}

\prose{อุปสมฺปนฺโน อุปสมฺปนฺนสฺส สุตฺวา อุปสมฺปนฺนสฺส เปสุญฺญํ อุปสํหรติ “อิตฺถนฺนาโม ตํ ‘โสตาปนฺโน’ติ ภณตี”ติ, อาปตฺติ วาจาย วาจาย}

\prose{อุปสมฺปนฺโน อุปสมฺปนฺนสฺส สุตฺวา อุปสมฺปนฺนสฺส เปสุญฺญํ อุปสํหรติ “อิตฺถนฺนาโม ตํ ‘โอฏฺโฐ เมณฺโฑ โคโณ คทฺรโภ ติรจฺฉานคโต เนรยิโก, นตฺถิ ตสฺส สุคติ, ทุคฺคติ เยว ตสฺส ปาฏิกงฺขา’ติ ภณตี”ติ, อาปตฺติ วาจาย วาจาย ปาจิตฺติยสฺส.}

\prose{อุปสมฺปนฺโน อุปสมฺปนฺนสฺส สุตฺวา อุปสมฺปนฺนสฺส เปสุญฺญํ อุปสํหรติ “อิตฺถนฺนาโม ตํ ‘ปณฺฑิโต พฺยตฺโต เมธาวี พหุสฺสุโต ธมฺมกถิโก, นตฺถิ ตสฺส ทุคฺคติ, สุคติ เยว ตสฺส ปาฏิกงฺขา’ติ ภณตี”ติ, อาปตฺติ วาจาย วาจาย ปาจิตฺติยสฺส.}

\proseitem{41}{อุปสมฺปนฺโน อุปสมฺปนฺนสฺส สุตฺวา อุปสมฺปนฺนสฺส เปสุญฺญํ อุปสํหรติ “อิตฺถนฺนาโม ‘สนฺติ อิเธกจฺเจ จณฺฑาลา เวนา เนสาทา รถการา ปุกฺกุสา’ติ ภณติ, น โส อญฺญํ ภณติ, ตญฺเญว ภณตี”ติ, อาปตฺติ วาจาย วาจาย ทุกฺกฏสฺส.}

\prose{อุปสมฺปนฺโน อุปสมฺปนฺนสฺส สุตฺวา อุปสมฺปนฺนสฺส เปสุญฺญํ อุปสํหรติ “อิตฺถนฺนาโม ‘สนฺติ อิเธกจฺเจ ขตฺติยา พฺราหฺมณา’ติ ภณติ, น โส อญฺญํ ภณติ, ตญฺเญว ภณตี”ติ, อาปตฺติ วาจาย วาจาย ทุกฺกฏสฺส ฯเปฯ.}

\prose{อุปสมฺปนฺโน อุปสมฺปนฺนสฺส สุตฺวา อุปสมฺปนฺนสฺส เปสุญฺญํ อุปสํหรติ “อิตฺถนฺนาโม ‘สนฺติ อิเธกจฺเจ ปณฺฑิตา พฺยตฺตา เมธาวี พหุสฺสุตา ธมฺมกถิกา, นตฺถิ เตสํ ทุคฺคติ, สุคติ เยว เตสํ ปาฏิกงฺขา’ติ ภณติ, \csromanfolio{24}\textbf{น โส อญฺญํ ภณติ, ตญฺเญว ภณตี”ติ, อาปตฺติ วาจาย วาจาย ทุกฺกฏสฺส.}}

\prose{อุปสมฺปนฺโน อุปสมฺปนฺนสฺส สุตฺวา อุปสมฺปนฺนสฺส เปสุญฺญํ อุปสํหรติ “อิตฺถนฺนาโม ‘เย นูน จณฺฑาลา เวนา เนสาทา รถการา ปุกฺกุสา’ติ ภณติ, น โส อญฺญํ ภณติ, ตญฺเญว ภณตี”ติ, อาปตฺติ วาจาย วาจาย ทุกฺกฏสฺส.}

\prose{อุปสมฺปนฺโน อุปสมฺปนฺนสฺส สุตฺวา อุปสมฺปนฺนสฺส เปสุญฺญํ อุปสํหรติ “อิตฺถนฺนาโม ‘เย นูน ปณฺฑิตา พฺยตฺตา เมธาวี พหุสฺสุตา ธมฺมกถิกา’ติ ภณติ, น โส อญฺญํ ภณติ, ตญฺเญว ภณตี”ติ, อาปตฺติ วาจาย วาจาย ทุกฺกฏสฺส.}

\prose{อุปสมฺปนฺโน อุปสมฺปนฺนสฺส สุตฺวา อุปสมฺปนฺนสฺส เปสุญฺญํ อุปสํหรติ “อิตฺถนฺนาโม ‘น มยํ จณฺฑาลา เวนา เนสาทา รถการา ปุกฺกุสา’ติ ภณติ, น โส อญฺญํ ภณติ, ตญฺเญว ภณตี”ติ, อาปตฺติ วาจาย วาจาย ทุกฺกฏสฺส.}

\prose{อุปสมฺปนฺโน อุปสมฺปนฺนสฺส สุตฺวา อุปสมฺปนฺนสฺส เปสุญฺญํ อุปสํหรติ “อิตฺถนฺนาโม ‘น มยํ ปณฺฑิตา พฺยตฺตา เมธาวี พหุสฺสุตา ธมฺมกถิกา, นตฺถมฺหากํ ทุคฺคติ, สุคติ เยว อมฺหากํ ปาฏิกงฺขา’ติ ภณติ, น โส อญฺญํ ภณติ, ตญฺเญว ภณตี”ติ, อาปตฺติ วาจาย วาจาย ทุกฺกฏสฺส.}

\proseitem{42}{อุปสมฺปนฺโน อุปสมฺปนฺนสฺส สุตฺวา อุปสมฺปนฺนสฺส เปสุญฺญํ อุปสํหรติ, อาปตฺติ วาจาย วาจาย ปาจิตฺติยสฺส.}

\prose{อุปสมฺปนฺโน อุปสมฺปนฺนสฺส สุตฺวา อนุปสมฺปนฺนสฺส เปสุญฺญํ อุปสํหรติ, อาปตฺติ ทุกฺกฏสฺส.}

\prose{อุปสมฺปนฺโน อนุปสมฺปนฺนสฺส สุตฺวา อุปสมฺปนฺนสฺส เปสุญฺญํ อุปสํหรติ, อาปตฺติ ทุกฺกฏสฺส.}

\prose{อุปสมฺปนฺโน อนุปสมฺปนฺนสฺส สุตฺวา อนุปสมฺปนฺนสฺส เปสุญฺญํ อุปสํหรติ, อาปตฺติ ทุกฺกฏสฺส.}

\proseitemwithcloser{43}{อนาปตฺติ นปิยกมฺยสฺส นเภทาธิปฺปายสฺส อุมฺมตฺตกสฺส อาทิกมฺมิกสฺสาติ.}{เปสุญฺญสิกฺขาปทํ นิฏฺฐิตํ ตติยํ.}
\csromanmatikaanchor{mtk.6}
\csromantocmarkref{subsection}{4. ปทโสธมฺมสิกฺขาปท}{mtk.6}
\bookmark[dest={mtk.6},level=2]{4. ปทโสธมฺมสิกฺขาปท}
\csromanheader{2}{1. \textbf{มุสาวาทวคฺค 4. ปทโสธมฺมสิกฺขาปท}}
\proseitem{44}{\csromanfolio{25}เตน สมเยน พุทฺโธ ภควา สาวตฺถิยํ วิหรติ เชตวเน อนาถปิณฺฑิกสฺส อาราเม.\csromanspacer{} เตน โข ปน สมเยน ฉพฺพคฺคิยา ภิกฺขู อุปาสเก ปทโส ธมฺมํ วาเจนฺติ.\csromanspacer{} อุปาสกา ภิกฺขูสุ อคารวา อปฺปติสฺสา อสภาควุตฺติกา วิหรนฺติ.\csromanspacer{} เย เต ภิกฺขู อปฺปิจฺฉา ฯเปฯ เต อุชฺฌายนฺติ ขิยฺยนฺติ วิปาเจนฺติ “กถํ หิ นาม ฉพฺพคฺคิยา ภิกฺขู อุปาสเก ปทโส ธมฺมํ วาเจสฺสนฺติ, อุปาสกา ภิกฺขูสุ อคารวา อปฺปติสฺสา อสภาควุตฺติกา วิหรนฺตี”ติ.\csromanspacer{} อถ โข เต ภิกฺขู ฉพฺพคฺคิเย ภิกฺขู อเนกปริยาเยน วิครหิตฺวา ภควโต เอตมตฺถํ อาโรเจสุํ ฯเปฯ “สจฺจํ กิร ตุเมฺห ภิกฺขเว อุปาสเก ปทโส ธมฺมํ วาเจถ, อุปาสกา ภิกฺขูสุ อคารวา อปฺปติสฺสา อสภาควุตฺติกา วิหรนฺตี”ติ.\csromanspacer{} สจฺจํ ภควาติ.\csromanspacer{} วิครหิ พุทฺโธ ภควา ฯเปฯ.\csromanspacer{} กถํ หิ นาม ตุเมฺห โมฆปุริสา อุปาสเก ปทโส ธมฺมํ วาเจสฺสถ, อุปาสกา ภิกฺขูสุ อคารวา อปฺปติสฺสา อสภาควุตฺติกา วิหรนฺติ.\csromanspacer{} เนตํ โมฆปุริสา อปฺปสนฺนานํ วา ปสาทาย ปสนฺนานํ วา ภิยฺโยภาวาย ฯเปฯ.\csromanspacer{} เอวญฺจ ปน ภิกฺขเว อิมํ สิกฺขาปทํ อุทฺทิเสยฺยาถ\csromandash{}}

\proseitem{45}{\textbf{“โย ปน ภิกฺขุ อนุปสมฺปนฺนํ ปทโส ธมฺมํ วาเจยฺย, ปาจิตฺติยนฺ”}ติ.}

\proseitem{46}{\textbf{โย ปนา}ติ โย ยาทิโส ฯเปฯ.\csromanspacer{} \textbf{ภิกฺขู}ติ ฯเปฯ อยํ อิมสฺมิํ อตฺเถ อธิปฺเปโต ภิกฺขูติ.}

\prose{\textbf{อนุปสมฺปนฺโน} นาม ภิกฺขุญฺจ ภิกฺขุนิญฺจ ฐเปตฺวา อวเสโส อนุปสมฺปนฺโน นาม.}

\prose{\textbf{ปทโส} นาม ปทํ อนุปทํ อนฺวกฺขรํ อนุพฺยญฺชนํ.}

\prose{\textbf{ปทํ} นาม เอกโต ปฏฺฐเปตฺวา เอกโต โอสาเปนฺติ.\csromanspacer{} \textbf{อนุปทํ} นาม ปาเฏกฺกํ ปฏฺฐเปตฺวา เอกโต โอสาเปนฺติ.\csromanspacer{} \textbf{อนฺวกฺขรํ} นาม “รูปํ อนิจฺจนฺ”ติ วุจฺจมาโน “รุนฺ”ติ โอปาเตติ.\csromanspacer{} \textbf{อนุพฺยญฺชนํ} นาม “รูปํ อนิจฺจนฺ”ติ วุจฺจมาโน “เวทนา อนิจฺจา”ติ สทฺทํ นิจฺฉาเรติ.}

\prose{\csromanfolio{26}ยญฺจ ปทํ ยญฺจ อนุปทํ ยญฺจ อนฺวกฺขรํ ยญฺจ อนุพฺยญฺชนํ, สพฺพเมตํ ปทโส\footnote{ปทโส ธมฺโม (อิติปิ)} นาม.}

\prose{\textbf{ธมฺโม} นาม พุทฺธภาสิโต สาวกภาสิโต อิสิภาสิโต เทวตาภาสิโต อตฺถูปสญฺหิโต ธมฺมูปสญฺหิโต.}

\prose{\textbf{วาเจยฺยา}ติ ปเทน วาเจติ, ปเท ปเท อาปตฺติ ปาจิตฺติยสฺส.\csromanspacer{} อกฺขราย วาเจติ.\csromanspacer{} อกฺขรกฺขราย อาปตฺติ ปาจิตฺติยสฺส.}

\proseitem{47}{อนุปสมฺปนฺเน อนุปสมฺปนฺนสญฺญี ปทโส ธมฺมํ วาเจติ, อาปตฺติ ปาจิตฺติยสฺส.\csromanspacer{} อนุปสมฺปนฺเน เวมติโก ปทโส ธมฺมํ วาเจติ, อาปตฺติ ปาจิตฺติยสฺส.\csromanspacer{} อนุปสมฺปนฺเน อุปสมฺปนฺนสญฺญี ปทโส ธมฺมํ วาเจติ, อาปตฺติ ปาจิตฺติยสฺส.}

\prose{อุปสมฺปนฺเน อนุปสมฺปนฺนสญฺญี, อาปตฺติ ทุกฺกฏสฺส.\csromanspacer{} อุปสมฺปนฺเน เวมติโก อาปตฺติ ทุกฺกฏสฺส.\csromanspacer{} อุปสมฺปนฺเน อุปสมฺปนฺนสญฺญี, อนาปตฺติ.}

\proseitemwithcloserruled{48}{อนาปตฺติ เอกโต อุทฺทิสาเปนฺโต, เอกโต สชฺฌายํ กโรนฺโต, เยภุยฺเยน ปคุณํ คนฺถํ ภณนฺตํ โอปาเตติ, โอสาเรนฺตํ โอปาเตติ, อุมฺมตฺตกสฺส อาทิกมฺมิกสฺสาติ.}{ปทโสธมฺมสิกฺขาปทํ นิฏฺฐิตํ จตุตฺถํ.}
\csromanmatikaanchor{mtk.7}
\csromantocmarkref{subsection}{5. สหเสยฺยสิกฺขาปท}{mtk.7}
\bookmark[dest={mtk.7},level=2]{5. สหเสยฺยสิกฺขาปท}
\csromanheader{2}{1. \textbf{มุสาวาทวคฺค 5. สหเสยฺยสิกฺขาปท}}
\proseitem{49}{เตน สมเยน พุทฺโธ ภควา อาฬวิยํ วิหรติ อคฺคาฬเว เจติเย.\csromanspacer{} เตน โข ปน สมเยน อุปาสกา อารามํ อาคจฺฉนฺติ ธมฺมสฺสวนาย.\csromanspacer{} ธมฺเม ภาสิเต เถรา ภิกฺขู ยถาวิหารํ คจฺฉนฺติ.\csromanspacer{} นวกา ภิกฺขู ตตฺเถว อุปฏฺฐานสาลายํ อุปาสเกหิ สทฺธิํ มุฏฺฐสฺสตี อสมฺปชานา นคฺคา วิกูชมานา กากจฺฉมานา เสยฺยํ กปฺเปนฺติ.\csromanspacer{} อุปาสกา อุชฺฌายนฺติ ขิยฺยนฺติ วิปาเจนฺติ “กถํ หิ นาม ภทนฺตา มุฏฺฐสฺสตี อสมฺปชานา นคฺคา วิกูชมานา กากจฺฉมานา เสยฺยํ กปฺเปสฺสนฺตี”ติ.\csromanspacer{} อสฺโสสุํ โข ภิกฺขู เตสํ อุปาสกานํ อุชฺฌายนฺตานํ ขิยฺยนฺตานํ วิปาเจนฺตานํ.\csromanspacer{} เย เต ภิกฺขู อปฺปิจฺฉา ฯเปฯ เต อุชฺฌายนฺติ ขิยฺยนฺติ วิปาเจนฺติ “กถํ หิ นาม ภิกฺขู \csromanfolio{27}อนุปสมฺปนฺเนน สหเสยฺยํ กปฺเปสฺสนฺตี”ติ.\csromanspacer{} อถ โข เต ภิกฺขู เต นวเก ภิกฺขู อเนกปริยาเยน วิครหิตฺวา ภควโต เอตมตฺถํ อาโรเจสุํ ฯเปฯ “สจฺจํ กิร ภิกฺขเว ภิกฺขู อนุปสมฺปนฺเนน สหเสยฺยํ กปฺเปนฺตี”ติ.\csromanspacer{} สจฺจํ ภควาติ.\csromanspacer{} วิครหิ พุทฺโธ ภควา ฯเปฯ.\csromanspacer{} กถํ หิ นาม เต ภิกฺขเว โมฆปุริสา อนุปสมฺปนฺเนน สหเสยฺยํ กปฺเปสฺสนฺติ.\csromanspacer{} เนตํ ภิกฺขเว อปฺปสนฺนานํ วา ปสาทาย ฯเปฯ.\csromanspacer{} เอวญฺจ ปน ภิกฺขเว อิมํ สิกฺขาปทํ อุทฺทิเสยฺยาถ\csromandash{}}

\prose{\textbf{“โย ปน ภิกฺขุ อนุปสมฺปนฺเนน สหเสยฺยํ กปฺเปยฺย, ปาจิตฺติยนฺ”}ติ.}

\prose{เอวญฺจิทํ ภควตา ภิกฺขูนํ สิกฺขาปทํ ปญฺญตฺตํ โหติ.}

\proseitem{50}{อถ โข ภควา อาฬวิยํ ยถาภิรนฺตํ วิหริตฺวา เยน โกสมฺพี เตน จาริกํ ปกฺกามิ.\csromanspacer{} อนุปุพฺเพน จาริกํ จรมาโน เยน โกสมฺพี ตทวสริ.\csromanspacer{} ตตฺร สุทํ ภควา โกสมฺพิยํ วิหรติ พทริการาเม.\csromanspacer{} ภิกฺขู อายสฺมนฺตํ ราหุลํ เอตทโวจุํ “ภควตา อาวุโส ราหุล สิกฺขาปทํ ปญฺญตฺตํ ‘น อนุปสมฺปนฺเนน สหเสยฺยา กปฺเปตพฺพา’ติ.\csromanspacer{} เสยฺยํ อาวุโส ราหุล ชานาหี”ติ.\csromanspacer{} อถ โข อายสฺมา ราหุโล เสยฺยํ อลภมาโน วจฺจกุฏิยา เสยฺยํ กปฺเปสิ.\csromanspacer{} อถ โข ภควา รตฺติยา ปจฺจูสสมยํ ปจฺจุฏฺฐาย เยน วจฺจกุฏิ เตนุปสงฺกมิ, อุปสงฺกมิตฺวา อุกฺกาสิ.\csromanspacer{} ราหุลํ อายสฺมาปิ ราหุโล อุกฺกาสิ.\csromanspacer{} โก เอตฺถาติ.\csromanspacer{} อหํ ภควา ราหุโลติ.\csromanspacer{} กิสฺส ตฺวํ ราหุล อิธ นิสินฺโนสีติ.\csromanspacer{} อถ โข อายสฺมา ราหุโล ภควโต เอตมตฺถํ อาโรเจสิ.\csromanspacer{} อถ โข ภควา เอตสฺมิํ นิทาเน เอตสฺมิํ ปกรเณ ธมฺมิํ กถํ กตฺวา ภิกฺขู อามนฺเตสิ อนุชานามิ ภิกฺขเว อนุปสมฺปนฺเนน ทิรตฺตติรตฺตํ สหเสยฺยํ กปฺเปตุํ.\csromanspacer{} เอวญฺจ ปน ภิกฺขเว อิมํ สิกฺขาปทํ อุทฺทิเสยฺยาถ\csromandash{}}

\proseitem{51}{“\textbf{โย ปน ภิกฺขุ อนุปฺ}อสมฺปนฺเนน อุตฺตริทิรตฺตติรตฺตํ สหเสยฺยํ กปฺเปยฺย, \textbf{ปาจิตฺติยนฺ”ติ.}}

\proseitem{52}{\textbf{โย ปนาติ} โย ยาทิโส ฯเปฯ.\csromanspacer{} \textbf{ภิกฺขู}ติ ฯเปฯ อยํ อิมสฺมิํ อตฺเถ อธิปฺเปโต ภิกฺขูติ.}

\prose{\textbf{อนุปสมฺปนฺโน} นาม ภิกฺขุํ ฐเปตฺวา อวเสโส อนุปสมฺปนฺโน นาม.}

\prose{\csromanfolio{28}\textbf{อุตฺตริทิรตฺตติรตฺตนฺ}ติ อติเรกทิรตฺตติรตฺตํ.}

\prose{\textbf{สหา}ติ เอกโต.}

\prose{\textbf{เสยฺยา} นาม สพฺพจฺฉนฺนา สพฺพปริจฺฉนฺนา, เยภุยฺเยนจฺฉนฺนา เยภุยฺเยน ปริจฺฉนฺนา.}

\prose{\textbf{เสยฺยํ กปฺเปยฺยา}ติ จตุตฺเถ ทิวเส อตฺถงฺคเต สูริเย อนุปสมฺปนฺเน นิปนฺเน ภิกฺขุ นิปชฺชติ, อาปตฺติ ปาจิตฺติยสฺส.\csromanspacer{} ภิกฺขุ นิปนฺเน อนุปสมฺปนฺโน นิปชฺชติ, อาปตฺติ ปาจิตฺติยสฺส.\csromanspacer{} อุโภ วา นิปชฺชนฺติ, อาปตฺติ ปาจิตฺติยสฺส.\csromanspacer{} อุฏฺฐหิตฺวา ปุนปฺปุนํ นิปชฺชนฺติ, อาปตฺติ ปาจิตฺติยสฺส.}

\proseitem{53}{อนุปสมฺปนฺเน อนุปสมฺปนฺนสญฺญี อุตฺตริทิรตฺตติรตฺตํ สหเสยฺยํ กปฺเปติ, อาปตฺติ ปาจิตฺติยสฺส.\csromanspacer{} อนุปสมฺปนฺเน เวมติโก อุตฺตริ ทิรตฺตติรตฺตํ สหเสยฺยํ กปฺเปติ, อาปตฺติ ปาจิตฺติยสฺส.\csromanspacer{} อนุปสมฺปนฺเน อุปสมฺปนฺนสญฺญี อุตฺตริทิรตฺตติรตฺตํ สหเสยฺยํ กปฺเปติ, อาปตฺติ}

\prose{อุปฑฺฒจฺฉนฺเน อุปฑฺฒปริจฺฉนฺเน อาปตฺติ ทุกฺกฏสฺส.\csromanspacer{} อุปสมฺปนฺเน อนุปสมฺปนฺนสญฺญี, อาปตฺติ ทุกฺกฏสฺส.\csromanspacer{} อุปสมฺปนฺเน เวมติโก, อาปตฺติ ทุกฺกฏสฺส.\csromanspacer{} อุปสมฺปนฺเน อุปสมฺปนฺนสญฺญี, อนาปตฺติ.}

\proseitemwithcloserruled{54}{อนาปตฺติ เทฺวติสฺโส รตฺติโย วสติ, อูนกเทฺวติสฺโส รตฺติโย วสติ, เทฺว รตฺติโย วสิตฺวา ตติยาย รตฺติยา ปุรารุณา นิกฺขมิตฺวา ปุน วสติ, สพฺพจฺฉนฺเน สพฺพอปริจฺฉนฺเน, สพฺพปริจฺฉนฺเน สพฺพอจฺฉนฺเน\footnote{วสติ, สพฺพอจฺฉนฺเน สพฺพอปริจฺฉนฺเน, (สี)}, เยภุยฺเยน อจฺฉนฺเน เยภุยฺเยน อปริจฺฉนฺเน, อนุปสมฺปนฺเน นิปนฺเน ภิกฺขุ นิสีทติ, ภิกฺขุ นิปนฺเน อนุปสมฺปนฺโน นิสีทติ, อุโภ วา นิสีทนฺติ, อุมฺมตฺตกสฺส อาทิกมฺมิกสฺสาติ.}{สหเสยฺยสิกฺขาปทํ นิฏฺฐิตํ ปญฺจมํ.}
\csromanmatikaanchor{mtk.8}
\csromantocmarkref{subsection}{6. ทุติยสหเสยฺยสิกฺขาปท}{mtk.8}
\bookmark[dest={mtk.8},level=2]{6. ทุติยสหเสยฺยสิกฺขาปท}
\csromanheader{2}{1. มุสาวาทวคฺค 6. ทุติย\-สหเสยฺย\-สิกฺขาปท}
\proseitem{55}{เตน สมเยน พุทฺโธ ภควา สาวตฺถิยํ วิหรติ เชตวเน อนาถปิณฺฑิกสฺส อาราเม.\csromanspacer{} เตน โข ปน สมเยน อายสฺมา อนุรุทฺโธ \csromanfolio{29}โกสเลสุ ชนปเท\footnote{เอตฺถ อมฺพฏฺฐสุตฺตาทิฏีกา โอโลเกตพฺพา.} สาวตฺถิํ คจฺฉนฺโต สายํ อญฺญตรํ คามํ อุปคจฺฉิ.\csromanspacer{} เตน โข ปน สมเยน ตสฺมิํ คาเม อญฺญตริสฺสา อิตฺถิยา อาวสถาคารํ ปญฺญตฺตํ โหติ.\csromanspacer{} อถ โข อายสฺมา อนุรุทฺโธ เยน สา อิตฺถี เตนุปสงฺกมิ, อุปสงฺกมิตฺวา ตํ อิตฺถิํ เอตทโวจ “สเจ เต ภคินิ อครุ, วเสยฺยาม เอกรตฺตํ อาวสถาคาเร”ติ.\csromanspacer{} วเสยฺยาถ ภนฺเตติ.\csromanspacer{} อญฺเญปิ อทฺธิกา เยน สา อิตฺถี เตนุปสงฺกมิํสุ, อุปสงฺกมิตฺวา ตํ อิตฺถิํ เอตทโวจุํ “สเจ เต อยฺเย อครุ, วเสยฺยาม เอกรตฺตํ อาวสถาคาเร”ติ.\csromanspacer{} เอโส โข อยฺโย สมโณ ปฐมํ อุปคโต, สเจ โส อนุชานาติ วเสยฺยาถาติ.\csromanspacer{} อถ โข เต อทฺธิกา เยนายสฺมา อนุรุทฺโธ เตนุปสงฺกมิํสุ, อุปสงฺกมิตฺวา อายสฺมนฺตํ อนุรุทฺธํ เอตทโวจุํ “สเจ เต ภนฺเต อครุ, วเสยฺยาม เอกรตฺตํ อาวสถาคาเร”ติ.\csromanspacer{} วเสยฺยาถ อาวุโสติ.\csromanspacer{} อถ โข สา อิตฺถี อายสฺมนฺเต อนุรุทฺเธ สห ทสฺสเนน ปฏิพทฺธจิตฺตา อโหสิ.\csromanspacer{} อถ โข สา อิตฺถี เยนายสฺมา อนุรุทฺโธ เตนุปสงฺกมิ, อุปสงฺกมิตฺวา อายสฺมนฺตํ อนุรุทฺธํ เอตทโวจ “อยฺโย ภนฺเต อิเมหิ มนุสฺเสหิ อากิณฺโณ น ผาสุ วิหริสฺสติ, สาธาหํ ภนฺเต อยฺยสฺส มญฺจกํ อพฺภนฺตรํ ปญฺญเปยฺยนฺ”ติ.\csromanspacer{} อธิวาเสสิ โข อายสฺมา อนุรุทฺโธ ตุณฺหีภาเวน.\csromanspacer{} อถ โข สา อิตฺถี อายสฺมโต อนุรุทฺธสฺส มญฺจกํ อพฺภนฺตรํ ปญฺญเปตฺวา อลงฺกตปฺปฏิยตฺตา คนฺธคนฺธินี เยนายสฺมา อนุรุทฺโธ เตนุปสงฺกมิ, อุปสงฺกมิตฺวา อายสฺมนฺตํ อนุรุทฺธํ เอตทโวจ “อยฺโย ภนฺเต อภิรูโป ทสฺสนีโย ปาสาทิโก, อหญฺจมฺหิ อภิรูปา ทสฺสนียา ปาสาทิกา, สาธาหํ ภนฺเต อยฺยสฺส ปชาปติ ภเวยฺยนฺ”ติ.\csromanspacer{} เอวํ วุตฺเต อายสฺมา อนุรุทฺโธ ตุณฺหี อโหสิ.\csromanspacer{} ทุติยมฺปิ โข ฯเปฯ.\csromanspacer{} ตติยมฺปิ โข สา อิตฺถี อายสฺมนฺตํ อนุรุทฺธํ เอตทโวจ “อยฺโย ภนฺเต อภิรูโป ทสฺสนีโย ปาสาทิโก, อหญฺจมฺหิ อภิรูปา ทสฺสนียา ปาสาทิกา, สาธุ ภนฺเต อยฺโย มญฺเจว ปฏิจฺฉตุ\footnote{สมฺปฏิจฺฉตุ (สฺยา)} สพฺพญฺจ สาปเตยฺยนฺ”ติ.\csromanspacer{} ตติยมฺปิ โข อายสฺมา อนุรุทฺโธ ตุณฺหี อโหสิ.\csromanspacer{} อถ โข สา อิตฺถี สาฏกํ นิกฺขิปิตฺวา อายสฺมโต อนุรุทฺธสฺส ปุรโต จงฺกมติปิ ติฏฺฐติปิ นิสีทติปิ เสยฺยมฺปิ กปฺเปติ.\csromanspacer{} อถ โข อายสฺมา อนุรุทฺโธ อินฺทฺริยานิ โอกฺขิปิตฺวา ตํ อิตฺถิํ เนว โอโลเกสิ, นปิ อาลปิ.\csromanspacer{} อถ \csromanfolio{30}โข สา อิตฺถี “อจฺฉริยํ วต โภ, อพฺภุตํ วต โภ, พหู เม มนุสฺสา สเตนปิ สหสฺเสนปิ ปหิณนฺติ, อยํ ปน สมโณ มยา สามํ ยาจิยมาโน น อิจฺฉติ มญฺเจว ปฏิจฺฉิตุํ สพฺพญฺจ สาปเตยฺยนฺ”ติ สาฏกํ นิวาเสตฺวา อายสฺมโต อนุรุทฺธสฺส ปาเทสุ สิรสา นิปติตฺวา อายสฺมนฺตํ \textbf{อนุรุทฺธํ เอตทโวจ “อจฺจโย มํ ภนฺเต อจฺจคมา ยถาพาลํ} ยถามูฬฺหํ ยถาอกุสลํ ยาหํ เอวมกาสิํ, ตสฺสา เม ภนฺเต อยฺโย อจฺจยํ อจฺจยโต ปฏิคฺคณฺหาตุ อายติํ สํวรายา”ติ.\csromanspacer{} ตคฺฆ ตฺวํ ภคินิ อจฺจโย อจฺจคมา ยถาพาลํ ยถามูฬฺหํ ยถาอกุสลํ ยา ตฺวํ เอวมกาสิ, ยโต จ โข ตฺวํ ภคินิ อจฺจยํ อจฺจยโต ทิสฺวา ยถาธมฺมํ ปฏิกโรสิ, ตํ เต มยํ ปฏิคฺคณฺหาม.\csromanspacer{} วุทฺธิ เหสา ภคินิ อริยสฺส วินเย, โย อจฺจยํ อจฺจยโต ทิสฺวา ยถาธมฺมํ ปฏิกโรติ, อายติํ จ สํวรํ อาปชฺชตีติ.}

\prose{อถ โข สา อิตฺถี ตสฺสา รตฺติยา อจฺจเยน อายสฺมนฺตํ อนุรุทฺธํ ปณีเตน ขาทนีเยน โภชนีเยน สหตฺถา สนฺตปฺเปตฺวา สมฺปวาเรตฺวา อายสฺมนฺตํ อนุรุทฺธํ ภุตฺตาวิํ โอนีตปตฺตปาณิํ อภิวาเทตฺวา เอกมนฺตํ นิสีทิ.\csromanspacer{} เอกมนฺตํ นิสินฺนํ โข ตํ อิตฺถิํ อายสฺมา อนุรุทฺโธ ธมฺมิยา กถาย สนฺทสฺเสสิ สมาทเปสิ สมุตฺเตเชสิ สมฺปหํเสสิ.\csromanspacer{} อถ โข สา อิตฺถี อายสฺมตา อนุรุทฺเธน ธมฺมิยา กถาย สนฺทสฺสิตา สมาทปิตา สมุตฺเตชิตา สมฺปหํสิตา อายสฺมนฺตํ อนุรุทฺธํ เอตทโวจ “อภิกฺกนฺตํ ภนฺเต, อภิกฺกนฺตํ ภนฺเต, เสยฺยถาปิ ภนฺเต นิกฺกุชฺชิตํ วา อุกฺกุชฺเชยฺย, ปฏิจฺฉนฺนํ วา วิวเรยฺย, มูฬฺหสฺส วา มคฺคํ อาจิกฺเขยฺย, อนฺธกาเร วา เตลปชฺโชตํ ธาเรยฺย ‘จกฺขุมนฺโต รูปานิ ทกฺขนฺตี’ติ, เอวเมวํ อยฺเยน อนุรุทฺเธน อเนกปริยาเยน ธมฺโม ปกาสิโต.\csromanspacer{} เอสาหํ ภนฺเต ตํ ภควนฺตํ สรณํ คจฺฉามิ ธมฺมญฺจ ภิกฺขุสํฆญฺจ, อุปาสิกํ มํ อยฺโย ธาเรตุ อชฺชตคฺเค ปาณุเปตํ สรณํ คตนฺ”ติ.}

\prose{อถ โข อายสฺมา อนุรุทฺโธ สาวตฺถิํ คนฺตฺวา ภิกฺขูนํ เอตมตฺถํ อาโรเจสิ.\csromanspacer{} เย เต ภิกฺขู อปฺปิจฺฉา ฯเปฯ เต อุชฺฌายนฺติ ขิยฺยนฺติ วิปาเจนฺติ “กถํ หิ นาม อายสฺมา อนุรุทฺโธ มาตุคาเมน สหเสยฺยํ กปฺเปสฺสตี”ติ.\csromanspacer{} อถ โข เต ภิกฺขู อายสฺมนฺตํ อนุรุทฺธํ อเนกปริยาเยน วิครหิตฺวา ภควโต เอตมตฺถํ อาโรเจสุํ ฯเปฯ “สจฺจํ กิร ตฺวํ อนุรุทฺธ มาตุคาเมน สหเสยฺยํ กปฺเปสี”ติ.\csromanspacer{} สจฺจํ \csromanfolio{31}ภควาติ.\csromanspacer{} วิครหิ พุทฺโธ ภควา ฯเปฯ.\csromanspacer{} กถํ หิ นาม ตฺวํ อนุรุทฺธ มาตุคาเมน สหเสยฺยํ กปฺเปสฺสสิ.\csromanspacer{} เนตํ อนุรุทฺธ อปฺปสนฺนานํ วา ปสาทาย ฯเปฯ.\csromanspacer{} เอวญฺจ ปน ภิกฺขเว อิมํ สิกฺขาปทํ อุทฺทิเสยฺยาถ\csromandash{}}

\proseitem{56}{\textbf{“โย ปน ภิกฺขุ มาตุคาเมน สหเสยฺยํ กปฺเปยฺย, ปาจิตฺติยนฺ”}ติ.}

\proseitem{57}{\textbf{โย ปนา}ติ โย ยาทิโส ฯเปฯ.\csromanspacer{} \textbf{ภิกฺขู}ติ ฯเปฯ อยํ อิมสฺมิํ อตฺเถ อธิปฺเปโต ภิกฺขูติ.}

\prose{\textbf{มาตุโคโม} นาม มนุสฺสิตฺถี, น ยกฺขี\footnote{ยกฺขินี (ก)}, น เปตี, น ติรจฺฉานคตา, อนฺตมโส ตทหุชาตาปิ ทาริกา, ปเคว มหตฺตรี.}

\prose{\textbf{สหา}ติ เอกโต.}

\prose{\textbf{เสยฺยา} นาม สพฺพจฺฉนฺนา สพฺพปริจฺฉนฺนา, เยภุยฺเยนจฺฉนฺนา เยภุยฺเยน ปริจฺฉนฺนา.}

\prose{\textbf{เสยฺยํ กปฺเปยฺยา}ติ อตฺถงฺคเต สูริเย มาตุคาเม นิปนฺเน ภิกฺขุ นิปชฺชติ, อาปตฺติ ปาจิตฺติยสฺส.\csromanspacer{} ภิกฺขุ นิปนฺเน มาตุคาโม นิปชฺชติ, อาปตฺติ ปาจิตฺติยสฺส.\csromanspacer{} อุโภ วา นิปชฺชนฺติ, อาปตฺติ ปาจิตฺติยสฺส.\csromanspacer{} อุฏฺฐหิตฺวา ปุนปฺปุนํ นิปชฺชนฺติ, อาปตฺติ ปาจิตฺติยสฺส.}

\proseitem{58}{มาตุคาเม มาตุคามสญฺญี สหเสยฺยํ กปฺเปติ, อาปตฺติ ปาจิตฺติยสฺส.\csromanspacer{} มาตุคาเม เวมติโก สหเสยฺยํ กปฺเปติ, อาปตฺติ ปาจิตฺติยสฺส.\csromanspacer{} มาตุคาเม อมาตุคามสญฺญี สหเสยฺยํ กปฺเปติ, อาปตฺติ ปาจิตฺติยสฺส.}

\prose{อุปฑฺฒจฺฉนฺเน อุปฑฺฒปริจฺฉนฺเน อาปตฺติ ทุกฺกฏสฺส.\csromanspacer{} ยกฺขิยา วา เปติยา วา ปณฺฑเกน วา ติรจฺฉานคติตฺถิยา วา สหเสยฺยํ กปฺเปติ, อาปตฺติ ทุกฺกฏสฺส.\csromanspacer{} อมาตุคาเม มาตุคามสญฺญี, อาปตฺติ ทุกฺกฏสฺส.\csromanspacer{} อมาตุคาเม เวมติโก, อาปตฺติ ทุกฺกฏสฺส.\csromanspacer{} อมาตุคาเม อมาตุคามสญฺญี, อนาปตฺติ.}

\proseitemwithcloserruled{59}{อนาปตฺติ สพฺพจฺฉนฺเน สพฺพอปริจฺฉนฺเน, สพฺพปริจฺฉนฺเน สพฺพอจฺฉนฺเน\footnote{อนาปตฺติ สพฺพอจฺฉนฺเน สพฺพอปริจฺฉนฺเน, (สี)}, เยภุยฺเยน อจฺฉนฺเน เยภุยฺเยน อปริจฺฉนฺเน, มาตุคาเม นิปนฺเน ภิกฺขุ \csromanfolio{32}นิสีทติ, ภิกฺขุ นิปนฺเน มาตุคาโม นิสีทติ, อุโภ วา นิสีทนฺติ, อุมฺมตฺตกสฺส อาทิกมฺมิกสฺสาติ.}{ทุติย\-สหเสยฺย\-สิกฺขาปทํ นิฏฺฐิตํ ฉฏฺฐํ.}
\csromanmatikaanchor{mtk.9}
\csromantocmarkref{subsection}{7. ธมฺมเทสนาสิกฺขาปท}{mtk.9}
\bookmark[dest={mtk.9},level=2]{7. ธมฺมเทสนาสิกฺขาปท}
\csromanheader{2}{1. \textbf{มุสาวาทวคฺค 7. ธมฺมเทสนาสิกฺขาปท}}
\proseitem{60}{เตน สมเยน พุทฺโธ ภควา สาวตฺถิยํ วิหรติ เชตวเน อนาถปิณฺฑิกสฺส อาราเม.\csromanspacer{} เตน โข ปน สมเยน อายสฺมา อุทายี สาวตฺถิยํ กุลูปโก โหติ, พหุกานิ กุลานิ อุปสงฺกมติ.\csromanspacer{} อถ โข อายสฺมา อุทายี ปุพฺพณฺหสมยํ นิวาเสตฺวา ปตฺตจีวรมาทาย เยน อญฺญตรํ กุลํ เตนุปสงฺกมิ, เตน โข ปน สมเยน ฆรณี นิเวสนทฺวาเร นิสินฺนา โหติ, ฆรสุณฺหา อาวสถทฺวาเร นิสินฺนา โหติ.\csromanspacer{} อถ โข อายสฺมา อุทายี เยน ฆรณี เตนุปสงฺกมิ, อุปสงฺกมิตฺวา ฆรณิยา อุปกณฺณเก ธมฺมํ เทเสสิ.\csromanspacer{} อถ โข ฆรสุณฺหาย เอตทโหสิ “กิํ นุ โข โส สมโณ สสฺสุยา ชาโร อุทาหุ โอภาสตี”ติ.}

\prose{อถ โข อายสฺมา อุทายี ฆรณิยา อุปกณฺณเก ธมฺมํ เทเสตฺวา เยน ฆรสุณฺหา เตนุปสงฺกมิ, อุปสงฺกมิตฺวา ฆรสุณฺหาย อุปกณฺณเก ธมฺมํ เทเสสิ.\csromanspacer{} อถ โข ฆรณิยา เอตทโหสิ “กิํ นุ โข โส สมโณ ฆรสุณฺหาย ชาโร อุทาหุ โอภาสตี”ติ.\csromanspacer{} อถ โข อายสฺมา อุทายี ฆรสุณฺหาย อุปกณฺณเก ธมฺมํ เทเสตฺวา ปกฺกามิ.\csromanspacer{} อถ โข ฆรณี ฆรสุณฺหํ เอตทโวจ “เห เช กิํ เต เอโส สมโณ อโวจา”ติ.\csromanspacer{} ธมฺมํ เม อยฺเย เทเสสิ, อยฺยาย ปน กิํ อโวจาติ.\csromanspacer{} มยฺหมฺปิ ธมฺมํ เทเสสีติ.\csromanspacer{} ตา อุชฺฌายนฺติ ขิยฺยนฺติ วิปาเจนฺติ “กถํ หิ นาม อยฺโย อุทายี อุปกณฺณเก ธมฺมํ เทเสสฺสติ, นนุ นาม วิสฺสฏฺเฐน วิวเฏน ธมฺโม เทเสตพฺโพ”ติ.}

\prose{อสฺโสสุํ โข ภิกฺขู ตาสํ อิตฺถีนํ อุชฺฌายนฺตีนํ ขิยฺยนฺตีนํ วิปาเจนฺตีนํ.\csromanspacer{} เย เต ภิกฺขู อปฺปิจฺฉา ฯเปฯ เต อุชฺฌายนฺติ ขิยฺยนฺติ วิปาเจนฺติ “กถํ หิ นาม อายสฺมา อุทายี มาตุคามสฺส ธมฺมํ เทเสสฺสตี”ติ.\csromanspacer{} อถ โข เต \csromanfolio{33}ภิกฺขู อายสฺมนฺตํ อุทายิํ อเนกปริยาเยน วิครหิตฺวา ภควโต เอตมตฺถํ อาโรเจสุํ ฯเปฯ “สจฺจํ กิร ตฺวํ อุทายิ มาตุคามสฺส ธมฺมํ เทเสสี”ติ.\csromanspacer{} สจฺจํ ภควาติ.\csromanspacer{} วิครหิ พุทฺโธ ภควา ฯเปฯ.\csromanspacer{} กถํ หิ นาม ตฺวํ โมฆปุริส มาตุคามสฺส ธมฺมํ เทเสสฺสสิ.\csromanspacer{} เนตํ โมฆปุริส อปฺปสนฺนานํ วา ปสาทาย ฯเปฯ.\csromanspacer{} เอวญฺจ ปน ภิกฺขเว อิมํ สิกฺขาปทํ อุทฺทิเสยฺยาถ\csromandash{}}

\prose{\textbf{“โย ปน ภิกฺขุ มาตุคามสฺส ธมฺมํ เทเสยฺย, ปาจิตฺติยนฺ”}ติ.}

\prose{เอวญฺจิทํ ภควตา ภิกฺขูนํ สิกฺขาปทํ ปญฺญตฺตํ โหติ.}

\proseitem{61}{เตน โข ปน สมเยน อุปาสิกา ภิกฺขู ปสฺสิตฺวา เอตทโวจุํ “อิงฺฆายฺยา ธมฺมํ เทเสถา”ติ.\csromanspacer{} น ภคินี กปฺปติ มาตุคามสฺส ธมฺมํ เทเสตุนฺติ.\csromanspacer{} อิงฺฆายฺยา ฉปฺปญฺจวาจาหิ ธมฺมํ เทเสถ, สกฺกา เอตฺตเกนปิ ธมฺโม อญฺญาตุนฺติ. “น ภคินี กปฺปติ มาตุคามสฺส ธมฺมํ เทเสตุนฺ”ติ กุกฺกุจฺจายนฺตา น เทเสสุํ.\csromanspacer{} อุปาสิกา อุชฺฌายนฺติ ขิยฺยนฺติ วิปาเจนฺติ “กถํ หิ นาม อยฺยา อเมฺหหิ ยาจียมานา ธมฺมํ น เทเสสฺสนฺตี”ติ.\csromanspacer{} อสฺโสสุํ โข ภิกฺขู ตาสํ อุปาสิกานํ อุชฺฌายนฺตีนํ ขิยฺยนฺตีนํ วิปาเจนฺตีนํ.\csromanspacer{} อถ โข เต ภิกฺขู ภควโต เอตมตฺถํ อาโรเจสุํ.\csromanspacer{} อถ โข ภควา เอตสฺมิํ นิทาเน เอตสฺมิํ ปกรเณ ธมฺมิํ กถํ กตฺวา ภิกฺขู อามนฺเตสิ อนุชานามิ ภิกฺขเว มาตุคามสฺส ฉปฺปญฺจวาจาหิ ธมฺมํ เทเสตุํ.\csromanspacer{} เอวญฺจ ปน ภิกฺขเว อิมํ สิกฺขาปทํ อุทฺทิเสยฺยาถ\csromandash{}}

\prose{“\textbf{โย ปน ภิกฺขุ มาตุคามฺ}อสฺส อุตฺตริฉปฺปญฺจวาจาหิ ธมฺมํ เทเสยฺย, ปาจิตฺติยนฺ”ติ.}

\prose{เอวญฺจิทํ ภควตา ภิกฺขูนํ สิกฺขาปทํ ปญฺญตฺตํ โหติ.}

\proseitem{62}{เตน โข ปน สมเยน ฉพฺพคฺคิยา ภิกฺขู “ภควตา อนุญฺญาตํ มาตุคามสฺส ฉปฺปญฺจวาจาหิ ธมฺมํ เทเสตุนฺ”ติ เต อวิญฺญุํ ปุริสวิคฺคหํ อุปนิสีทาเปตฺวา มาตุคามสฺส อุตฺตริฉปฺปญฺจวาจาหิ ธมฺมํ เทเสนฺติ.\csromanspacer{} เย เต ภิกฺขู อปฺปิจฺฉา ฯเปฯ เต อุชฺฌายนฺติ ขิยฺยนฺติ วิปาเจนฺติ “กถํ หิ นาม ฉพฺพคฺคิยา ภิกฺขู อวิญฺญุํ ปุริสวิคฺคหํ อุปนิสีทาเปตฺวา มาตุคามสฺส อุตฺตริฉปฺปญฺจวาจาหิ ธมฺมํ เทเสสฺสนฺตี”ติ.}

\prose{\csromanfolio{34}อถ โข เต ภิกฺขู ฉพฺพคฺคิเย ภิกฺขู อเนกปริยาเยน วิครหิตฺวา ภควโต เอตมตฺถํ อาโรเจสุํ ฯเปฯ “สจฺจํ กิร ตุเมฺห ภิกฺขเว อวิญฺญุํ ปุริสวิคฺคหํ อุปนิสีทาเปตฺวา มาตุคามสฺส อุตฺตริฉปฺปญฺจวาจาหิ ธมฺมํ เทเสถา”ติ.\csromanspacer{} สจฺจํ ภควาติ.\csromanspacer{} วิครหิ พุทฺโธ ภควา ฯเปฯ.\csromanspacer{} กถํ หิ นาม ตุเมฺห โมฆปุริสา อวิญฺญุํ ปุริสวิคฺคหํ อุปนิสีทาเปตฺวา มาตุคามสฺส อุตฺตริฉปฺปญฺจวาจาหิ ธมฺมํ เทเสสฺสถ.\csromanspacer{} เนตํ โมฆปุริสา อปฺปสนฺนานํ วา ปสาทาย ฯเปฯ.\csromanspacer{} เอวญฺจ ปน ภิกฺขเว อิมํ สิกฺขาปทํ อุทฺทิเสยฺยาถ\csromandash{}}

\proseitem{63}{\textbf{“โย ปน ภิกฺขุ มาตุคามสฺส อุตฺตริฉปฺปญฺจวาจาหิ ธมฺมํ เทเสยฺย อญฺญตฺร วิญฺญุนา ปุริสวิคฺคเหน, ปาจิตฺติยนฺ”}ติ.}

\proseitem{64}{\textbf{โย ปนา}ติ โย ยาทิโส ฯเปฯ.\csromanspacer{} \textbf{ภิกฺขู}ติ ฯเปฯ อยํ อิมสฺมิํ อตฺเถ อธิปฺเปโต ภิกฺขูติ.}

\prose{\textbf{มาตุคาโม} นาม มนุสฺสิตฺถี, น ยกฺขี, น เปตี, น ติรจฺฉานคตา, วิญฺญู ปฏิพลา สุภาสิตทุพฺภาสิตํ ทุฏฺฐุลฺลาทุฏฺฐุลฺลํ อาชานิตุํ.}

\prose{อุตฺตริ\-ฉปฺปญฺจ\-วาจาหีติ อติเรก\-ฉปฺปญฺจ\-วาจาหิ.}

\prose{\textbf{ธมฺโม} นาม พุทฺธภาสิโต สาวกภาสิโต อิสิภาสิโต เทวตาภาสิโต อตฺถูปสญฺหิโต ธมฺมูปสญฺหิโต.}

\prose{\textbf{เทเสยฺยา}ติ ปเทน เทเสติ, ปเท ปเท อาปตฺติ ปาจิตฺติยสฺส.\csromanspacer{} อกฺขราย เทเสติ, อกฺขรกฺขราย อาปตฺติ ปาจิตฺติยสฺส.}

\prose{\textbf{อญฺญตฺร วิญฺญุนา ปุริสวิคฺคเหนา}ติ ฐเปตฺวา วิญฺญุํ ปุริสวิคฺคหํ, วิญฺญู นาม ปุริสวิคฺคโห ปฏิพโล โหติ สุภาสิตทุพฺภาสิตํ ทุฏฺฐุลฺลาทุฏฺฐุลฺลํ อาชานิตุํ.}

\proseitem{65}{มาตุคาเม มาตุคามสญฺญี อุตฺตริฉปฺปญฺจวาจาหิ ธมฺมํ เทเสติ อญฺญตฺร วิญฺญุนา ปุริสวิคฺคเหน, อาปตฺติ ปาจิตฺติยสฺส.\csromanspacer{} มาตุคาเม เวมติโก อุตฺตริฉปฺปญฺจวาจาหิ ธมฺมํ เทเสติ อญฺญตฺร วิญฺญุนา ปุริสวิคฺคเหน, อาปตฺติ ปาจิตฺติยสฺส.\csromanspacer{} มาตุคาเม อมาตุคามสญฺญี อุตฺตริฉปฺปญฺจวาจาหิ ธมฺมํ เทเสติ อญฺญตฺร วิญฺญุนา ปุริสวิคฺคเหน, อาปตฺติ ปาจิตฺติยสฺส.}

\prose{\csromanfolio{35}ยกฺขิยา วา เปติยา วา ปณฺฑกสฺส วา ติรจฺฉานคตมนุสฺสวิคฺคหิตฺถิยา วา อุตฺตริฉปฺปญฺจวาจาหิ ธมฺมํ เทเสติ อญฺญตฺร วิญฺญุนา ปุริสวิคฺคเหน, อาปตฺติ ทุกฺกฏสฺส.\csromanspacer{} อมาตุคาเม มาตุคามสญฺญี, อาปตฺติ ทุกฺกฏสฺส.\csromanspacer{} อมาตุคาเม เวมติโก อาปตฺติ ทุกฺกฏสฺส.\csromanspacer{} อมาตุคาเม อมาตุคามสญฺญี, อนาปตฺติ.}

\proseitemwithcloserruled{66}{อนาปตฺติ วิญฺญุนา ปุริสวิคฺคเหน, ฉปฺปญฺจวาจาหิ ธมฺมํ เทเสติ, อูนกฉปฺปญฺจวาจาหิ ธมฺมํ เทเสติ, อุฏฺฐหิตฺวา ปุน นิสีทิตฺวา เทเสติ, มาตุคาโม อุฏฺฐหิตฺวา ปุน นิสีทติ ตสฺมิํ เทเสติ, อญฺญสฺส มาตุคามสฺส เทเสติ, ปญฺหํ ปุจฺฉติ, ปญฺหํ ปุฏฺโฐ กเถติ, อญฺญสฺสตฺถาย ภณนฺตํ มาตุคาโม สุณาติ, อุมฺมตฺตกสฺส อาทิกมฺมิกสฺสาติ.}{ธมฺมเทสนาสิกฺขาปทํ นิฏฺฐิตํ สตฺตมํ.}
\csromanmatikaanchor{mtk.10}
\csromantocmarkref{subsection}{8. ภูตาโรจนสิกฺขาปท}{mtk.10}
\bookmark[dest={mtk.10},level=2]{8. ภูตาโรจนสิกฺขาปท}
\csromanheader{2}{1. \textbf{มุสาวาทวคฺค 8. ภูตาโรจนสิกฺขาปท}}
\proseitem{67}{\csromansymbolfootnote{*}{วิ 1. 112 ปิฏฺเฐปิ.}เตน สมเยน พุทฺโธ ภควา เวสาลิยํ วิหรติ มหาวเน กูฏาคารสาลายํ.\csromanspacer{} เตน โข ปน สมเยน สมฺพหุลา สนฺทิฏฺฐา สมฺภตฺตา ภิกฺขู วคฺคุมุทาย นทิยา ตีเร วสฺสํ อุปคจฺฉิํสุ.\csromanspacer{} เตน โข ปน สมเยน วชฺชี ทุพฺภิกฺขา โหติ ทฺวิหิติกา เสตฏฺฐิกา สลากาวุตฺตา, น สุกรา อุญฺเฉน ปคฺคเหน ยาเปตุํ.\csromanspacer{} อถ โข เตสํ ภิกฺขูนํ เอตทโหสิ “เอตรหิ โข วชฺชี ทุพฺภิกฺขา ทฺวิหิติกา เสตฏฺฐิกา สลากาวุตฺตา, น สุกรา อุญฺเฉน ปคฺคเหน ยาเปตุํ, เกน นุ โข มยํ อุปาเยน สมคฺคา สมฺโมทมานา อวิวทมานา ผาสุกํ วสฺสํ วเสยฺยาม, น จ ปิณฺฑเกน กิลเมยฺยามา”ติ.\csromanspacer{} เอกจฺเจ เอวมาหํสุ “หนฺท มยํ อาวุโส คิหีนํ กมฺมนฺตํ อธิฏฺเฐม, เอวํ เต อมฺหากํ ทาตุํ มญฺญิสฺสนฺติ, เอวํ มยํ สมคฺคา สมฺโมทมานา อวิวทมานา ผาสุกํ วสฺสํ วสิสฺสาม, น จ ปิณฺฑเกน กิลมิสฺสามา”ติ.\csromanspacer{} เอกจฺเจ เอวมาหํสุ “อลํ อาวุโส กิํ คิหีนํ กมฺมนฺตํ อธิฏฺฐิเตน, หนฺท มยํ อาวุโส คิหีนํ ทูเตยฺยํ หราม, เอวํ เต อมฺหากํ ทาตุํ มญฺญิสฺสนฺติ, เอวํ มยํ สมคฺคา สมฺโมทมานา อวิวทมานา ผาสุกํ วสฺสํ วสิสฺสาม, น จ ปิณฺฑเกน กิลมิสฺสามา”ติ.\csromanspacer{} เอกจฺเจ เอวมาหํสุ “อลํ \csromanfolio{36}อาวุโส กิํ คิหีนํ กมฺมนฺตํ อธิฏฺฐิเตน, กิํ คิหีนํ ทูเตยฺยํ หเฏน, หนฺท มยํ อาวุโส คิหีนํ อญฺญมญฺญสฺส อุตฺตริมนุสฺสธมฺมสฺส วณฺณํ ภาสิสฺสาม ‘อสุโก ภิกฺขุ ปฐมสฺส ฌานสฺส ลาภี, อสุโก ภิกฺขุ ทุติยสฺส ฌานสฺส ลาภี, อสุโก ภิกฺขุ ตติยสฺส ฌานสฺส ลาภี, อสุโก ภิกฺขุ จตุตฺถสฺส ฌานสฺส ลาภี, อสุโก ภิกฺขุ โสตาปนฺโน, อสุโก ภิกฺขุ สกทาคามี, อสุโก ภิกฺขุ อนาคามี, อสุโก ภิกฺขุ อรหา, อสุโก ภิกฺขุ เตวิชฺโช, อสุโก ภิกฺขุ ฉฬภิญฺโญ’ติ, เอวํ เต อมฺหากํ ทาตุํ มญฺญิสฺสนฺติ, เอวํ มยํ สมคฺคา สมฺโมทมานา อวิวทมานา ผาสุกํ วสฺสํ วสิสฺสาม, น จ ปิณฺฑเกน กิลมิสฺสามา”ติ.\csromanspacer{} เอโส เยว โข อาวุโส เสยฺโย, โย อมฺหากํ คิหีนํ อญฺญมญฺญสฺส อุตฺตริมนุสฺสธมฺมสฺส วณฺโณ ภาสิโตติ.}

\prose{อถ โข เต ภิกฺขู คิหีนํ อญฺญมญฺญสฺส อุตฺตริมนุสฺสธมฺมสฺส วณฺณํ ภาสิํสุ “อสุโก ภิกฺขุ ปฐมสฺส ฌานสฺส ลาภี ฯเปฯ อสุโก ภิกฺขุ ฉฬภิญฺโญ”ติ.\csromanspacer{} อถ โข เต มนุสฺสา “ลาภา วต โน, สุลทฺธํ วต โน, เยสํ โน เอวรูปา ภิกฺขู วสฺสํ อุปคตา, น วต โน อิโต ปุพฺเพ เอวรูปา ภิกฺขู วสฺสํ อุปคตา, ยถยิเม ภิกฺขู สีลวนฺโต กลฺยาณธมฺมา”ติ, เต น ตาทิสานิ โภชนานิ อตฺตนา ภุญฺชนฺติ มาตาปิตูนํ เทนฺติ ปุตฺตทารสฺส เทนฺติ ทาสกมฺมกรโปริสสฺส เทนฺติ มิตฺตามจฺจานํ เทนฺติ ญาติสาโลหิตานํ เทนฺติ, ยาทิสานิ ภิกฺขูนํ เทนฺติ.\csromanspacer{} น ตาทิสานิ ขาทนียานิ สายนียานิ ปานานิ อตฺตนา ขาทนฺติ สายนฺติ ปิวนฺติ\footnote{อตฺตนา ปิวนฺติ (สฺยา, ก)} มาตาปิตูนํ เทนฺติ ปุตฺตทารสฺส เทนฺติ ทาสกมฺมกรโปริสสฺส เทนฺติ มิตฺตามจฺจานํ เทนฺติ ญาติสาโลหิตานํ เทนฺติ, ยาทิสานิ ภิกฺขูนํ เทนฺติ.\csromanspacer{} อถ โข เต ภิกฺขู วณฺณวา อเหสุํ ปีณินฺทฺริยา ปสนฺนมุขวณฺณา วิปฺปสนฺนฉวิวณฺณา.}

\proseitem{68}{อาจิณฺณํ โข ปเนตํ วสฺสํวุฏฺฐานํ ภิกฺขูนํ ภควนฺตํ ทสฺสนาย อุปสงฺกมิตุํ.\csromanspacer{} อถ โข เต ภิกฺขู วสฺสํวุฏฺฐา เตมาสจฺจเยน เสนาสนํ สํสาเมตฺวา ปตฺตจีวรมาทาย เยน เวสาลี เตน ปกฺกมิํสุ, อนุปุพฺเพน เยน เวสาลี มหาวนํ กูฏาคารสาลา เยน ภควา เตนุปสงฺกมิํสุ, อุปสงฺกมิตฺวา ภควนฺตํ อภิวาเทตฺวา เอกมนฺตํ นิสีทิํสุ.\csromanspacer{} เตน โข ปน สมเยน ทิสาสุ วสฺสํวุฏฺฐา ภิกฺขู กิสา โหนฺติ ลูขา \csromanfolio{37}ทุพฺพณฺณา อุปฺปณฺฑุปฺปณฺฑุกชาตา ธมนิสนฺถตคตฺตา.\csromanspacer{} วคฺคุมุทาตีริยา ปน ภิกฺขู วณฺณวา โหนฺติ ปีณินฺทฺริยา ปสนฺนมุขวณฺณา วิปฺปสนฺนฉวิวณฺณา.\csromanspacer{} อาจิณฺณํ โข ปเนตํ พุทฺธานํ ภควนฺตานํ อาคนฺตุเกหิ ภิกฺขูหิ สทฺธิํ ปฏิสมฺโมทิตุํ.\csromanspacer{} อถ โข ภควา วคฺคุมุทาตีริเย ภิกฺขู เอตทโวจ “กจฺจิ ภิกฺขเว ขมนียํ, กจฺจิ ยาปนียํ, กจฺจิ สมคฺคา สมฺโมทมานา อวิวทมานา ผาสุกํ วสฺสํ วสิตฺถ, น จ ปิณฺฑเกน กิลมิตฺถา”ติ.\csromanspacer{} ขมนียํ ภควา, ยาปนียํ ภควา, สมคฺคา จ มยํ ภนฺเต สมฺโมทมานา อวิวทมานา ผาสุกํ วสฺสํ วสิมฺหา, น จ ปิณฺฑเกน กิลมิมฺหาติ.\csromanspacer{} ชานนฺตาปิ ตถาคตา ปุจฺฉนฺติ, ชานนฺตาปิ น ปุจฺฉนฺติ.\csromanspacer{} กาลํ วิทิตฺวา ปุจฺฉนฺติ, กาลํ วิทิตฺวา น ปุจฺฉนฺติ.\csromanspacer{} อตฺถสญฺหิตํ ตถาคตา ปุจฺฉนฺติ, โน อนตฺถสญฺหิตํ, อนตฺถสญฺหิเต เสตุฆาโต ตถาคตานํ.\csromanspacer{} ทฺวีหากาเรหิ พุทฺธา ภควนฺโต ภิกฺขู ปฏิปุจฺฉนฺติ “ธมฺมํ วา เทเสสฺสาม, สาวกานํ วา \textbf{สิกฺขาปทํ ปญฺญเปสฺสามา”ติ.}}

\prose{อถ โข ภควา วคฺคุมุทาตีริเย ภิกฺขู เอตทโวจ “ยถา กถํ ปน ตุเมฺห ภิกฺขเว สมคฺคา สมฺโมทมานา อวิวทมานา ผาสุกํ วสฺสํ วสิตฺถ, น จ ปิณฺฑเกน กิลมิตฺถา”ติ.\csromanspacer{} อถ โข เต ภิกฺขู ภควโต เอตมตฺถํ อาโรเจสุํ.\csromanspacer{} กจฺจิ ปน โว ภิกฺขเว ภูตนฺติ.\csromanspacer{} ภูตํ ภควาติ.\csromanspacer{} วิครหิ พุทฺโธ ภควา ฯเปฯ.\csromanspacer{} กถํ หิ นาม ตุเมฺห ภิกฺขเว อุทรสฺส การณา คิหีนํ อญฺญมญฺญสฺส อุตฺตริมนุสฺสธมฺมสฺส วณฺณํ ภาสิสฺสถ.\csromanspacer{} เนตํ ภิกฺขเว อปฺปสนฺนานํ วา ปสาทาย ฯเปฯ.\csromanspacer{} เอวญฺจ ปน ภิกฺขเว อิมํ สิกฺขาปทํ อุทฺทิเสยฺยาถ\csromandash{}}

\proseitem{69}{\textbf{“โย ปน ภิกฺขุ อนุปสมฺปนฺนสฺส อุตฺตริมนุสฺสธมฺมํ อาโรเจยฺย, ภูตสฺมิํ ปาจิตฺติยนฺ”}ติ.}

\proseitem{70}{\textbf{โย ปนา} ติ โย ยาทิโส ฯเปฯ.\csromanspacer{} \textbf{ภิกฺขู}ติ ฯเปฯ อยํ อิมสฺมิํ อตฺเถ อธิปฺเปโต ภิกฺขูติ.}

\prose{\textbf{อนุปสมฺปนฺโน} นาม ภิกฺขุญฺจ ภิกฺขุนิญฺจ ฐเปตฺวา อวเสโส อนุปสมฺปนฺโน นาม.}

\prose{\csromansymbolmark{*}\textbf{อุตฺตริมนุสฺสธมฺโม} นาม ฌานํ วิโมกฺโข สมาธิ สมาปตฺติ ญาณทสฺสนํ มคฺคภาวนา ผลสจฺฉิกิริยา กิเลสปฺปหานํ วินีวรณตา จิตฺตสฺส สุญฺญาคาเร อภิรติ.}

\prose{\csromanfolio{38}\csromansymbolfootnote{*}{วิ 1. 118 ปิฏฺเฐปิ.}\textbf{ฌานนฺ}ติ ปฐมํ ฌานํ ทุติยํ ฌานํ ตติยํ ฌานํ จตุตฺถํ ฌานํ.}

\prose{\csromansymbolmark{*}\textbf{วิโมกฺโข}ติ สุญฺญโต วิโมกฺโข อนิมิตฺโต วิโมกฺโข อปฺปณิหิโต วิโมกฺโข.}

\prose{\csromansymbolmark{*}\textbf{สมาธี}ติ สุญฺญโต สมาธิ อนิมิตฺโต สมาธิ อปฺปณิหิโต สมาธิ.}

\prose{\csromansymbolmark{+}\textbf{สมาปตฺตี}ติ สุญฺญตา สมาปตฺติ อนิมิตฺตา สมาปตฺติ อปฺปณิหิตา สมาปตฺติ.}

\prose{\csromansymbolmark{+}\textbf{ญาณทสฺสนนฺ}ติ ติสฺโส วิชฺชา.}

\prose{\csromansymbolmark{+}\textbf{มคฺคภาวนา}ติ จตฺตาโร สติปฏฺฐานา จตฺตาโร สมฺมปฺปธานา จตฺตาโร อิทฺธิปาทา ปญฺจินฺทฺริยานิ ปญฺจ พลานิ สตฺต โพชฺฌงฺคา อริโย อฏฺฐงฺคิโก มคฺโค.}

\prose{\csromansymbolmark{+}\textbf{ผลสจฺฉิกิริยา}ติ โสตาปตฺติผลสฺส สจฺฉิกิริยา สกทาคามิผลสฺส สจฺฉิกิริยา อนาคามิผลสฺส สจฺฉิกิริยา อรหตฺตสฺส\footnote{อรหตฺตผลสฺส (สฺยา)} สจฺฉิกิริยา.}

\prose{\csromansymbolmark{+}\textbf{กิเลสปฺปหานนฺ}ติ ราคสฺส ปหานํ โทสสฺส ปหานํ โมหสฺส ปหานํ.}

\prose{\csromansymbolmark{+}\textbf{วินีวรณตา จิตฺตสฺสา}ติ ราคา จิตฺตํ วินีวรณตา โทสา จิตฺตํ วินีวรณตา โมหา จิตฺตํ วินีวรณตา.}

\prose{\csromansymbolmark{+}\textbf{สุญฺญาคาเร อภิรตี}ติ ปฐเมน ฌาเนน สุญฺญาคาเร อภิรติ, ทุติเยน ฌาเนน สุญฺญาคาเร อภิรติ, ตติเยน ฌาเนน สุญฺญาคาเร อภิรติ, จตุตฺเถน ฌาเนน สุญฺญาคาเร อภิรติ.}

\proseitem{71}{\textbf{อาโรเจยฺยา}ติ อนุปสมฺปนฺนสฺส ปฐมํ ฌานํ สมาปชฺชินฺติ ภณนฺตสฺส อาปตฺติ ปาจิตฺติยสฺส.}

\prose{\textbf{อาโรเจยฺยา}ติ อนุปสมฺปนฺนสฺส ปฐมํ ฌานํ สมาปชฺชามีติ ภณนฺตสฺส อาปตฺติ ปาจิตฺติยสฺส.}

\prose{\textbf{อาโรเจยฺยา}ติ อนุปสมฺปนฺนสฺส ปฐมํ ฌานํ สมาปนฺโนติ ภณนฺตสฺส อาปตฺติ ปาจิตฺติยสฺส.}

\prose{\csromanfolio{39}\textbf{อาโรเจยฺยา}ติ อนุปสมฺปนฺนสฺส ปฐมสฺส ฌานสฺส ลาภิมฺหีติ ภณนฺตสฺส อาปตฺติ ปาจิตฺติยสฺส.}

\prose{\textbf{อาโรเจยฺยา}ติ อนุปสมฺปนฺนสฺส ปฐมสฺส ฌานสฺส วสิมฺหีติ ภณนฺตสฺส อาปตฺติ ปาจิตฺติยสฺส.}

\prose{\textbf{อาโรเจยฺยา}ติ อนุปสมฺปนฺนสฺส ปฐมํ ฌานํ สจฺฉิกตํ มยาติ ภณนฺตสฺส อาปตฺติ ปาจิตฺติยสฺส.}

\prose{\textbf{อาโรเจยฺยา}ติ อนุปสมฺปนฺนสฺส ทุติยํ ฌานํ ฯเปฯ ตติยํ ฌานํ.\csromanspacer{} จตุตฺถํ ฌานํ สมาปชฺชิํ.\csromanspacer{} สมาปชฺชามิ.\csromanspacer{} สมาปนฺโน.\csromanspacer{} จตุตฺถสฺส ฌานสฺส ลาภิมฺหิ.\csromanspacer{} วสิมฺหิ.\csromanspacer{} จตุตฺถํ ฌานํ สจฺฉิกตํ มยาติ ภณนฺตสฺส อาปตฺติ ปาจิตฺติยสฺส.}

\prose{\textbf{อาโรเจยฺยา}ติ อนุปสมฺปนฺนสฺส สุญฺญตํ วิโมกฺขํ ฯเปฯ อนิมิตฺตํ วิโมกฺขํ.\csromanspacer{} อปฺปณิหิตํ วิโมกฺขํ.\csromanspacer{} สุญฺญตํ สมาธิํ.\csromanspacer{} อนิมิตฺตํ สมาธิํ.\csromanspacer{} อปฺปณิหิตํ สมาธิํ สมาปชฺชิํ.\csromanspacer{} สมาปชฺชามิ.\csromanspacer{} สมาปนฺโน.\csromanspacer{} อปฺปณิหิตสฺส สมาธิสฺส ลาภิมฺหิ.\csromanspacer{} วสิมฺหิ.\csromanspacer{} อปฺปณิหิโต สมาธิ สจฺฉิกโต มยาติ ภณนฺตสฺส อาปตฺติ ปาจิตฺติยสฺส.}

\prose{\textbf{อาโรเจยฺยา}ติ อนุปสมฺปนฺนสฺส สุญฺญตํ สมาปตฺติํ ฯเปฯ อนิมิตฺตํ สมาปตฺติํ.\csromanspacer{} อปฺปณิหิตํ สมาปตฺติํ สมาปชฺชิํ.\csromanspacer{} สมาปชฺชามิ.\csromanspacer{} สมาปนฺโน.\csromanspacer{} อปฺปณิหิตาย สมาปตฺติยา ลาภิมฺหิ.\csromanspacer{} วสิมฺหิ.\csromanspacer{} อปฺปณิหิตา สมาปตฺติ สจฺฉิกตา มยาติ ภณนฺตสฺส อาปตฺติ ปาจิตฺติยสฺส.}

\prose{\textbf{อาโรเจยฺยา}ติ อนุปสมฺปนฺนสฺส ติสฺโส วิชฺชา สมาปชฺชิํ.\csromanspacer{} สมาปชฺชามิ.\csromanspacer{} สมาปนฺโน.\csromanspacer{} ติสฺสนฺนํ วิชฺชานํ ลาภิมฺหิ.\csromanspacer{} วสิมฺหิ.\csromanspacer{} ติสฺโส วิชฺชา สจฺฉิกตา มยาติ ภณนฺตสฺส อาปตฺติ ปาจิตฺติยสฺส.}

\prose{\textbf{อาโรเจยฺยา}ติ อนุปสมฺปนฺนสฺส จตฺตาโร สติปฏฺฐาเน ฯเปฯ จตฺตาโร สมฺมปฺปธาเน.\csromanspacer{} จตฺตาโร อิทฺธิปาเท สมาปชฺชิํ.\csromanspacer{} สมาปชฺชามิ.\csromanspacer{} สมาปนฺโน.\csromanspacer{} จตุนฺนํ อิทฺธิปาทานํ ลาภิมฺหิ.\csromanspacer{} วสิมฺหิ.\csromanspacer{} จตฺตาโร อิทฺธิปาเท สจฺฉิกตา มยาติ ภณนฺตสฺส อาปตฺติ ปาจิตฺติยสฺส.}

\prose{\textbf{อาโรเจยฺยา}ติ อนุปสมฺปนฺนสฺส ปญฺจินฺทฺริยานิ ฯเปฯ ปญฺจ พลานิ สมาปชฺชิํ.\csromanspacer{} สมาปชฺชามิ.\csromanspacer{} สมาปนฺโน.\csromanspacer{} ปญฺจนฺนํ พลานํ ลาภิมฺหิ.\csromanspacer{} วสิมฺหิ.\csromanspacer{} ปญฺจ พลานิ สจฺฉิกตานิ มยาติ ภณนฺตสฺส อาปตฺติ ปาจิตฺติยสฺส.}

\prose{\csromanfolio{40}\textbf{อาโรเจยฺยา}ติ อนุปสมฺปนฺนสฺส สตฺต โพชฺฌงฺเค สมาปชฺชิํ.\csromanspacer{} สมาปชฺชามิ.\csromanspacer{} สมาปนฺโน.\csromanspacer{} สตฺตนฺนํ โพชฺฌงฺคานํ ลาภิมฺหิ.\csromanspacer{} วสิมฺหิ.\csromanspacer{} สตฺต โพชฺฌงฺคา สจฺฉิกตา มยาติ ภณนฺตสฺส อาปตฺติ ปาจิตฺติยสฺส.}

\prose{\textbf{อาโรเจยฺยา}ติ อนุปสมฺปนฺนสฺส อริยํ อฏฺฐงฺคิกํ มคฺคํ สมาปชฺชิํ.\csromanspacer{} สมาปชฺชามิ.\csromanspacer{} สมาปนฺโน.\csromanspacer{} อริยสฺส อฏฺฐงฺคิกสฺส มคฺคสฺส ลาภิมฺหิ.\csromanspacer{} วสิมฺหิ.\csromanspacer{} อริโย อฏฺฐงฺคิโก มคฺโค สจฺฉิกโต มยาติ ภณนฺตสฺส อาปตฺติ ปาจิตฺติยสฺส.}

\prose{\textbf{อาโรเจยฺยา}ติ อนุปสมฺปนฺนสฺส โสตาปตฺติผลํ ฯเปฯ สกทาคามิผลํ.\csromanspacer{} อนาคามิผลํ.\csromanspacer{} อรหตฺตํ\footnote{อรหตฺตผลํ (สฺยา)} สมาปชฺชิํ.\csromanspacer{} สมาปชฺชามิ.\csromanspacer{} สมาปนฺโน.\csromanspacer{} อรหตฺตสฺส ลาภิมฺหิ.\csromanspacer{} วสิมฺหิ.\csromanspacer{} อรหตฺตํ1 สจฺฉิกตํ มยาติ ภณนฺตสฺส อาปตฺติ ปาจิตฺติยสฺส.}

\prose{\textbf{อาโรเจยฺยา}ติ อนุปสมฺปนฺนสฺส ราโค เม จตฺโต ฯเปฯ โทโส เม จตฺโต.\csromanspacer{} โมโห เม จตฺโต วนฺโต มุตฺโต ปหีโน ปฏินิสฺสฏฺโฐ อุกฺเขฏิโต สมุกฺเขฏิโตติ ภณนฺตสฺส อาปตฺติ ปาจิตฺติยสฺส.}

\prose{\textbf{อาโรเจยฺยา}ติ อนุปสมฺปนฺนสฺส ราคา เม จิตฺตํ วินีวรณํ ฯเปฯ โทสา เม จิตฺตํ วินีวรณํ.\csromanspacer{} โมหา เม จิตฺตํ วินีวรณนฺติ ภณนฺตสฺส อาปตฺติ}

\prose{\textbf{อาโรเจยฺยา}ติ อนุปสมฺปนฺนสฺส สุญฺญาคาเร ปฐมํ ฌานํ ฯเปฯ ทุติยํ ฌานํ.\csromanspacer{} ตติยํ ฌานํ.\csromanspacer{} จตุตฺถํ ฌานํ สมาปชฺชิํ.\csromanspacer{} สมาปชฺชามิ.\csromanspacer{} สมาปนฺโน.\csromanspacer{} สุญฺญาคาเร จตุตฺถสฺส ฌานสฺส ลาภิมฺหิ.\csromanspacer{} วสิมฺหิ.\csromanspacer{} สุญฺญาคาเร จตุตฺถํ ฌานํ สจฺฉิกตํ มยาติ ภณนฺตสฺส อาปตฺติ ปาจิตฺติยสฺส.}

\proseitem{72}{\textbf{อาโรเจยฺยา}ติ อนุปสมฺปนฺนสฺส ปฐมญฺจ ฌานํ ทุติยญฺจ ฌานํ สมาปชฺชิํ.\csromanspacer{} สมาปชฺชามิ.\csromanspacer{} สมาปนฺโน.\csromanspacer{} ปฐมสฺส จ ฌานสฺส ทุติยสฺส จ ฌานสฺส ลาภิมฺหิ.\csromanspacer{} วสิมฺหิ.\csromanspacer{} ปฐมญฺจ ฌานํ ทุติยญฺจ ฌานํ สจฺฉิกตํ มยาติ ภณนฺตสฺส อาปตฺติ ปาจิตฺติยสฺส.}

\prose{\textbf{อาโรเจยฺยา}ติ อนุปสมฺปนฺนสฺส ปฐมญฺจ ฌานํ ตติยญฺจ ฌานํ ฯเปฯ ปฐมญฺจ ฌานํ จตุตฺถญฺจ ฌานํ สมาปชฺชิํ.\csromanspacer{} สมาปชฺชามิ.\csromanspacer{} สมาปนฺโน.\csromanspacer{} ปฐมสฺส จ ฌานสฺส จตุตฺถสฺส จ ฌานสฺส ลาภิมฺหิ.\csromanspacer{} วสิมฺหิ.\csromanspacer{} ปฐมญฺจ ฌานํ จตุตฺถญฺจ ฌานํ สจฺฉิกตํ มยาติ ภณนฺตสฺส อาปตฺติ ปาจิตฺติยสฺส.}

\prose{\csromanfolio{41}\textbf{อาโรเจยฺยา}ติ อนุปสมฺปนฺนสฺส ปฐมญฺจ ฌานํ สุญฺญตญฺจ วิโมกฺขํ ฯเปฯ อนิมิตฺตญฺจ วิโมกฺขํ.\csromanspacer{} อปฺปณิหิตญฺจ วิโมกฺขํ.\csromanspacer{} สุญฺญตญฺจ สมาธิํ.\csromanspacer{} อนิมิตฺตญฺจ สมาธิํ.\csromanspacer{} อปฺปณิหิตญฺจ สมาธิํ สมาปชฺชิํ.\csromanspacer{} สมาปชฺชามิ.\csromanspacer{} สมาปนฺโน.\csromanspacer{} ปฐมสฺส จ ฌานสฺส อปฺปณิหิตสฺส จ สมาธิสฺส ลาภิมฺหิ.\csromanspacer{} วสิมฺหิ.\csromanspacer{} ปฐมญฺจ ฌานํ อปฺปณิหิโต จ สมาธิ สจฺฉิกโต มยาติ ภณนฺตสฺส อาปตฺติ ปาจิตฺติยสฺส.}

\prose{\textbf{อาโรเจยฺยา}ติ อนุปสมฺปนฺนสฺส ปฐมญฺจ ฌานํ สุญฺญตญฺจ สมาปตฺติํ ฯเปฯ อนิมิตฺตญฺจ สมาปตฺติํ.\csromanspacer{} อปฺปณิหิตญฺจ สมาปตฺติํ สมาปชฺชิํ.\csromanspacer{} สมาปชฺชามิ.\csromanspacer{} สมาปนฺโน.\csromanspacer{} ปฐมสฺส จ ฌานสฺส อปฺปณิหิตาย จ สมาปตฺติยา ลาภิมฺหิ.\csromanspacer{} วสิมฺหิ.\csromanspacer{} ปฐมญฺจ ฌานํ อปฺปณิหิตา จ สมาปตฺติ สจฺฉิกตา มยาติ ภณนฺตสฺส อาปตฺติ ปาจิตฺติยสฺส.}

\prose{\textbf{อาโรเจยฺยา}ติ อนุปสมฺปนฺนสฺส ปฐมญฺจ ฌานํ ติสฺโส จ วิชฺชา สมาปชฺชิํ.\csromanspacer{} สมาปชฺชามิ.\csromanspacer{} สมาปนฺโน.\csromanspacer{} ปฐมสฺส จ ฌานสฺส ติสฺสนฺนญฺจ วิชฺชานํ ลาภิมฺหิ.\csromanspacer{} วสิมฺหิ.\csromanspacer{} ปฐมญฺจ ฌานํ ติสฺโส จ วิชฺชา สจฺฉิกตา มยาติ ภณนฺตสฺส อาปตฺติ ปาจิตฺติยสฺส.}

\prose{\textbf{อาโรเจยฺยา}ติ อนุปสมฺปนฺนสฺส ปฐมญฺจ ฌานํ จตฺตาโร จ สติปฏฺฐาเน ฯเปฯ จตฺตาโร จ สมฺมปฺปธาเน.\csromanspacer{} จตฺตาโร จ อิทฺธิปาเท สมาปชฺชิํ.\csromanspacer{} สมาปชฺชามิ.\csromanspacer{} สมาปนฺโน.\csromanspacer{} ปฐมสฺส จ ฌานสฺส จตุนฺนญฺจ อิทฺธิปาทานํ ลาภิมฺหิ.\csromanspacer{} วสิมฺหิ.\csromanspacer{} ปฐมญฺจ ฌานํ จตฺตาโร จ อิทฺธิปาทา สจฺฉิกตา มยาติ ภณนฺตสฺส อาปตฺติ ปาจิตฺติยสฺส.}

\prose{\textbf{อาโรเจยฺยา}ติ อนุปสมฺปนฺนสฺส ปฐมญฺจ ฌานํ ปญฺจ จ อินฺทฺริยานิ ฯเปฯ ปญฺจ จ พลานิ สมาปชฺชิํ.\csromanspacer{} สมาปชฺชามิ.\csromanspacer{} สมาปฺปนฺโน.\csromanspacer{} ปฐมสฺส จ ฌานสฺส ปญฺจนฺนญฺจ พลานํ ลาภิมฺหิ.\csromanspacer{} วสิมฺหิ.\csromanspacer{} ปฐมญฺจ ฌานํ ปญฺจ จ พลานิ สจฺฉิกตานิ มยาติ ภณนฺตสฺส อาปตฺติ ปาจิตฺติยสฺส.}

\prose{\textbf{อาโรเจยฺยา}ติ อนุปสมฺปนฺนสฺส ปฐมญฺจ ฌานํ สตฺต จ โพชฺฌงฺเค สมาปชฺชิํ.\csromanspacer{} สมาปชฺชามิ.\csromanspacer{} สมาปนฺโน.\csromanspacer{} ปฐมสฺส จ ฌานสฺส สตฺตนฺนญฺจ โพชฺฌงฺคานํ ลาภิมฺหิ.\csromanspacer{} วสิมฺหิ.\csromanspacer{} ปฐมญฺจ ฌานํ สตฺต จ โพชฺฌงฺคา สจฺฉิกตา มยาติ ภณนฺตสฺส อาปตฺติ ปาจิตฺติยสฺส.}

\prose{\textbf{อาโรเจยฺยา}ติ อนุปสมฺปนฺนสฺส ปฐมญฺจ ฌานํ อริยญฺจ อฏฺฐงฺคิกํ มคฺคํ สมาปชฺชิํ.\csromanspacer{} สมาปชฺชามิ.\csromanspacer{} สมาปนฺโน.\csromanspacer{} ปฐมสฺส จ ฌานสฺส อริยสฺส จ \csromanfolio{42}อฏฺฐงฺคิกสฺส มคฺคสฺส ลาภิมฺหิ.\csromanspacer{} วสิมฺหิ.\csromanspacer{} ปฐมญฺจ ฌานํ อริโย จ อฏฺฐงฺคิโก มคฺโค สจฺฉิกโต มยาติ ภณนฺตสฺส อาปตฺติ ปาจิตฺติยสฺส.}

\prose{\textbf{อาโรเจยฺยา}ติ อนุปสมฺปนฺนสฺส ปฐมญฺจ ฌานํ โสตาปตฺติผลญฺจ ฯเปฯ สกทาคามิผลญฺจ.\csromanspacer{} อนาคามิผลญฺจ.\csromanspacer{} อรหตฺตญฺจ\footnote{อรหตฺตผลญฺจ (สฺยา)} สมาปชฺชิํ.\csromanspacer{} สมาปชฺชามิ.\csromanspacer{} สมาปนฺโน.\csromanspacer{} ปฐมสฺส จ ฌานสฺส อรหตฺตสฺส\footnote{อรหตฺตผลสฺส (สฺยา)} จ ลาภิมฺหิ.\csromanspacer{} วสิมฺหิ.\csromanspacer{} ปฐมญฺจ ฌานํ อรหตฺตญฺจ1 สจฺฉิกตํ มยาติ ภณนฺตสฺส อาปตฺติ ปาจิตฺติยสฺส.}

\prose{\textbf{อาโรเจยฺยา}ติ อนุปสมฺปนฺนสฺส ปฐมญฺจ ฌานํ สมาปชฺชิํ.\csromanspacer{} สมาปชฺชามิ.\csromanspacer{} สมาปนฺโน ฯเปฯ ราโค จ เม จตฺโต.\csromanspacer{} โทโส จ เม จตฺโต.\csromanspacer{} โมโห จ เม จตฺโต วนฺโต มุตฺโต ปหีโน ปฏินิสฺสฏฺโฐ อุกฺเขฏิโต สมุกฺเขฏิโตติ ภณนฺตสฺส อาปตฺติ ปาจิตฺติยสฺส.}

\prose{\textbf{อาโรเจยฺยา}ติ อนุปสมฺปนฺนสฺส ปฐมญฺจ ฌานํ สมาปชฺชิํ.\csromanspacer{} สมาปชฺชามิ.\csromanspacer{} สมาปนฺโน ฯเปฯ ราคา จ เม จิตฺตํ วินีวรณํ.\csromanspacer{} โทสา จ เม จิตฺตํ วินีวรณํ.\csromanspacer{} โมหา จ เม จิตฺตํ วินีวรณนฺติ ภณนฺตสฺส อาปตฺติ}

\proseitem{73}{\textbf{อาโรเจยฺยา}ติ อนุปสมฺปนฺนสฺส ทุติยญฺจ ฌานํ ตติยญฺจ ฌานํ ฯเปฯ ทุติยญฺจ ฌานํ จตุตฺถญฺจ ฌานํ สมาปชฺชิํ.\csromanspacer{} สมาปชฺชามิ.\csromanspacer{} สมาปนฺโน.\csromanspacer{} ทุติยสฺส จ ฌานสฺส จตุตฺถสฺส จ ฌานสฺส ลาภิมฺหิ.\csromanspacer{} วสิมฺหิ.\csromanspacer{} ทุติยญฺจ ฌานํ จตุตฺถญฺจ ฌานํ สจฺฉิกตํ มยาติ ภณนฺตสฺส อาปตฺติ ปาจิตฺติยสฺส.}

\prose{\textbf{อาโรเจยฺยา}ติ อนุปสมฺปนฺนสฺส ทุติยญฺจ ฌานํ สุญฺญตญฺจ วิโมกฺขํ ฯเปฯ โมหา จ เม จิตฺตํ วินีวรณนฺติ ภณนฺตสฺส อาปตฺติ}

\prose{\textbf{อาโรเจยฺยา}ติ อนุปสมฺปนฺนสฺส ทุติยญฺจ ฌานํ ปฐมญฺจ ฌานํ สมาปชฺชิํ.\csromanspacer{} สมาปชฺชามิ.\csromanspacer{} สมาปนฺโน.\csromanspacer{} ทุติยสฺส จ ฌานสฺส ปฐมสฺส จ ฌานสฺส ลาภิมฺหิ.\csromanspacer{} วสิมฺหิ.\csromanspacer{} ทุติยญฺจ ฌานํ ปฐมญฺจ ฌานํ สจฺฉิกตํ มยาติ ภณนฺตสฺส อาปตฺติ ปาจิตฺติยสฺส ฯเปฯ.\csromanspacer{} มูลํ สํขิตฺตํ}

\prose{\textbf{อาโรเจยฺยา}ติ อนุปสมฺปนฺนสฺส โมหา จ เม จิตฺตํ วินีวรณํ ปฐมญฺจ ฌานํ สมาปชฺชิํ.\csromanspacer{} สมาปชฺชามิ.\csromanspacer{} สมาปนฺโน.\csromanspacer{} โมหา จ เม จิตฺตํ วินีวรณํ \csromanfolio{43}ปฐมสฺส จ ฌานสฺส ลาภิมฺหิ.\csromanspacer{} วสิมฺหิ.\csromanspacer{} โมหา จ เม จิตฺตํ วินีวรณํ ปฐมญฺจ ฌานํ สจฺฉิกตํ มยาติ ภณนฺตสฺส อาปตฺติ ปาจิตฺติยสฺส ฯเปฯ.}

\prose{\textbf{อาโรเจยฺยา}ติ อนุปสมฺปนฺนสฺส โมหา จ เม จิตฺตํ วินีวรณํ โทสา จ เม จิตฺตํ วินีวรณนฺติ ภณนฺตสฺส อาปตฺติ ปาจิตฺติยสฺส ฯเปฯ.}

\prose{\textbf{อาโรเจยฺยา}ติ อนุปสมฺปนฺนสฺส ปฐมญฺจ ฌานํ ทุติยญฺจ ฌานํ ตติยญฺจ ฌานํ จตุตฺถญฺจ ฌานํ สุญฺญตญฺจ วิโมกฺขํ อนิมิตฺตญฺจ วิโมกฺขํ อปฺปณิหิตญฺจ วิโมกฺขํ สุญฺญตญฺจ สมาธิํ อนิมิตฺตญฺจ สมาธิํ อปฺปณิหิตญฺจ สมาธิํ สุญฺญตญฺจ สมาปตฺติํ อนิมิตฺตญฺจ สมาปตฺติํ อปฺปณิหิตญฺจ สมาปตฺติํ ติสฺโส จ วิชฺชา จตฺตาโร จ สติปฏฺฐาเน จตฺตาโร จ สมฺมปฺปธาเน จตฺตาโร จ อิทฺธิปาเท ปญฺจ จ อินฺทฺริยานิ ปญฺจ จ พลานิ สตฺต จ โพชฺฌงฺเค อริยญฺจ อฏฺฐงฺคิกํ มคฺคํ โสตาปตฺติผลญฺจ สกทาคามิผลญฺจ อนาคามิผลญฺจ อรหตฺตญฺจ\footnote{อรหตฺตผลญฺจ (สฺยา)} สมาปชฺชิํ ฯเปฯ ราโค จ เม จตฺโต โทโส จ เม จตฺโต โมโห จ เม จตฺโต วนฺโต มุตฺโต ปหีโน ปฏินิสฺสฏฺโฐ อุกฺเขฏิโต สมุกฺเขฏิโต ราคา จ เม จิตฺตํ วินีวรณํ โทสา จ เม จิตฺตํ วินีวรณํ โมหา จ เม จิตฺตํ วินีวรณนฺติ ภณนฺตสฺส อาปตฺติ}

\proseitem{74}{\textbf{อาโรเจยฺยา}ติ อนุปสมฺปนฺนสฺส ปฐมํ ฌานํ สมาปชฺชินฺติ วตฺตุกาโม ทุติยํ ฌานํ สมาปชฺชินฺติ ภณนฺตสฺส ปฏิวิชานนฺตสฺส อาปตฺติ ปาจิตฺติยสฺส, น ปฏิวิชานนฺตสฺส อาปตฺติ ทุกฺกฏสฺส.}

\prose{\textbf{อาโรเจยฺยา}ติ อนุปสมฺปนฺนสฺส ปฐมํ ฌานํ สมาปชฺชินฺติ วตฺตุกาโม ตติยํ ฌานํ ฯเปฯ จตุตฺถํ ฌานํ.\csromanspacer{} สุญฺญตํ วิโมกฺขํ.\csromanspacer{} อนิมิตฺตํ วิโมกฺขํ.\csromanspacer{} อปฺปณิหิตํ วิโมกฺขํ.\csromanspacer{} สุญฺญตํ สมาธิํ.\csromanspacer{} อนิมิตฺตํ สมาธิํ.\csromanspacer{} อปฺปณิหิตํ สมาธิํ.\csromanspacer{} สุญฺญตํ สมาปตฺติํ.\csromanspacer{} อนิมิตฺตํ สมาปตฺติํ.\csromanspacer{} อปฺปณิหิตํ สมาปตฺติํ.\csromanspacer{} ติสฺโส วิชฺชา.\csromanspacer{} จตฺตาโร สติปฏฺฐาเน.\csromanspacer{} จตฺตาโร สมฺมปฺปธาเน.\csromanspacer{} จตฺตาโร อิทฺธิปาเท.\csromanspacer{} ปญฺจนฺทฺริยานิ.\csromanspacer{} ปญฺจ พลานิ.\csromanspacer{} สตฺต โพชฺฌงฺเค.\csromanspacer{} อริยํ อฏฺฐงฺคิกํ มคฺคํ.\csromanspacer{} โสตาปตฺติผลํ.\csromanspacer{} สกทาคามิผลํ.\csromanspacer{} อนาคามิผลํ.\csromanspacer{} อรหตฺตํ\footnote{อรหตฺตผลํ (สฺยา)} สมาปชฺชิํ ฯเปฯ ราโค เม จตฺโต.\csromanspacer{} โทโส เม จตฺโต.\csromanspacer{} โมโห เม จตฺโต วนฺโต มุตฺโต ปหีโน ปฏินิสฺสฏฺโฐ อุกฺเขฏิโต สมุกฺเขฏิโต.\csromanspacer{} ราคา เม จิตฺตํ วินีวรณํ.\csromanspacer{} โทสา เม จิตฺตํ วินีวรณํ.\csromanspacer{} โมหา เม จิตฺตํ วินีวรณนฺติ ภณนฺตสฺส ปฏิวิชานนฺตสฺส อาปตฺติ ปาจิตฺติยสฺส, น ปฏิชานนฺตสฺส อาปตฺติ ทุกฺกฏสฺส.}

\prose{\csromanfolio{44}\textbf{อาโรเจยฺยา}ติ อนุปสมฺปนฺนสฺส ทุติยํ ฌานํ สมาปชฺชินฺติ วตฺตุกาโม ฯเปฯ โมหา เม จิตฺตํ วินีวรณนฺติ ภณนฺตสฺส ปฏิวิชานนฺตสฺส อาปตฺติ ปาจิตฺติยสฺส, น ปฏิวิชานนฺตสฺส อาปตฺติ ทุกฺกฏสฺส ฯเปฯ.}

\prose{\textbf{อาโรเจยฺยา}ติ อนุปสมฺปนฺนสฺส ทุติยํ ฌานํ สมาปชฺชินฺติ วตฺตุกาโม ปฐมํ ฌานํ สมาปชฺชินฺติ ภณนฺตสฺส ปฏิวิชานนฺตสฺส อาปตฺติ ปาจิตฺติยสฺส, น ปฏิวิชานนฺตสฺส อาปตฺติ ทุกฺกฏสฺส ฯเปฯ.\csromanspacer{} มูลํ สํขิตฺตํ}

\prose{\textbf{อาโรเจยฺยา}ติ อนุปสมฺปนฺนสฺส โมหา เม จิตฺตํ วินีวรณนฺติ วตฺตุกาโม ปฐมํ ฌานํ สมาปชฺชินฺติ ภณนฺตสฺส ปฏิวิชานนฺตสฺส อาปตฺติ ปาจิตฺติยสฺส, น ปฏิวิชานนฺตสฺส อาปตฺติ ทุกฺกฏสฺส ฯเปฯ.}

\prose{\textbf{อาโรเจยฺยา}ติ อนุปสมฺปนฺนสฺส โมหา เม จิตฺตํ วินีวรณนฺติ วตฺตุกาโม โทสา เม จิตฺตํ วินีวรณนฺติ ภณนฺตสฺส ปฏิวิชานนฺตสฺส อาปตฺติ ปาจิตฺติยสฺส, น ปฏิวิชานนฺตสฺส อาปตฺติ ทุกฺกฏสฺส ฯเปฯ.}

\prose{\textbf{อาโรเจยฺยา}ติ อนุปสมฺปนฺนสฺส ปฐมญฺจ ฌานํ ทุติยญฺจ ฌานํ ตติยญฺจ ฌานํ จตุตฺถญฺจ ฌานํ ฯเปฯ โทสา จ เม จิตฺตํ วินีวรณนฺติ วตฺตุกาโม โมหา เม จิตฺตํ วินีวรณนฺติ ภณนฺตสฺส ปฏิวิชานนฺตสฺส อาปตฺติ ปาจิตฺติยสฺส, น ปฏิวิชานนฺตสฺส อาปตฺติ ทุกฺกฏสฺส.}

\prose{\textbf{อาโรเจยฺยา}ติ อนุปสมฺปนฺนสฺส ทุติยญฺจ ฌานํ ตติยญฺจ ฌานํ จตุตฺถญฺจ ฌานํ ฯเปฯ โมหา จ เม จิตฺตํ วินีวรณนฺติ วตฺตุกาโม ปฐมํ ฌานํ สมาปชฺชินฺติ ภณนฺตสฺส ปฏิวิชานนฺตสฺส อาปตฺติ ปาจิตฺติยสฺส, น ปฏิวิชานนฺตสฺส อาปตฺติ ทุกฺกฏสฺส ฯเปฯ.}

\proseitem{75}{\textbf{อาโรเจยฺยา}ติ อนุปสมฺปนฺนสฺส โย เต วิหาเร วสิ, โส ภิกฺขุ ปฐมํ ฌานํ สมาปชฺชิ.\csromanspacer{} สมาปชฺชติ.\csromanspacer{} สมาปนฺโน.\csromanspacer{} โส ภิกฺขุ ปฐมสฺส ฌานสฺส ลาภี.\csromanspacer{} วสี.\csromanspacer{} เตน ภิกฺขุนา ปฐมํ ฌานํ สจฺฉิกตนฺติ ภณนฺตสฺส อาปตฺติ ทุกฺกฏสฺส.}

\prose{\textbf{อาโรเจยฺยา}ติ อนุปสมฺปนฺนสฺส โย เต วิหาเร วสิ, โส ภิกฺขุ ทุติยํ ฌานํ ฯเปฯ ตติยํ ฌานํ.\csromanspacer{} จตุตฺถํ ฌานํ สมาปชฺชิ.\csromanspacer{} สมาปชฺชติ.\csromanspacer{} สมาปนฺโน.\csromanspacer{} โส ภิกฺขุ จตุตฺถสฺส ฌานสฺส ลาภี.\csromanspacer{} วสี.\csromanspacer{} เตน ภิกฺขุนา จตุตฺถํ ฌานํ สจฺฉิกตนฺติ ภณนฺตสฺส อาปตฺติ ทุกฺกฏสฺส.}

\prose{\csromanfolio{45}\textbf{อาโรเจยฺยา}ติ อนุปสมฺปนฺนสฺส โย เต วิหาเร วสิ, โส ภิกฺขุ สุญฺญตํ วิโมกฺขํ ฯเปฯ อนิมิตฺตํ วิโมกฺขํ.\csromanspacer{} อปฺปณิหิตํ วิโมกฺขํ.\csromanspacer{} สุญฺญตํ สมาธิํ.\csromanspacer{} อนิมิตฺตํ สมาธิํ.\csromanspacer{} อปฺปณิหิตํ สมาธิํ สมาปชฺชิ.\csromanspacer{} สมาปชฺชติ.\csromanspacer{} สมาปนฺโน.\csromanspacer{} โส ภิกฺขุ อปฺปณิหิตสฺส สมาธิสฺส ลาภี.\csromanspacer{} วสี.\csromanspacer{} เตน ภิกฺขุนา อปฺปณิหิโต สมาธิ สจฺฉิกโตติ ภณนฺตสฺส อาปตฺติ ทุกฺกฏสฺส.}

\prose{\textbf{อาโรเจยฺยา}ติ อนุปสมฺปนฺนสฺส โย เต วิหาเร วสิ, โส ภิกฺขุ สุญฺญตํ สมาปตฺติํ ฯเปฯ อนิมิตฺตํ สมาปตฺติํ.\csromanspacer{} อปฺปณิหิตํ สมาปตฺติํ สมาปชฺชิ.\csromanspacer{} สมาปชฺชติ.\csromanspacer{} สมาปนฺโน.\csromanspacer{} อปฺปณิหิตาย สมาปตฺติยา ลาภี.\csromanspacer{} วสี.\csromanspacer{} เตน ภิกฺขุนา อปฺปณิหิตา สมาปตฺติ สจฺฉิกตาติ ภณนฺตสฺส อาปตฺติ ทุกฺกฏสฺส.}

\prose{\textbf{อาโรเจยฺยา}ติ อนุปสมฺปนฺนสฺส โย เต วิหเร วสิ.\csromanspacer{} โส ภิกฺขุ ติสฺโส วิชฺชา ฯเปฯ จตฺตาโร สติปฏฺฐาเน.\csromanspacer{} จตฺตาโร สมฺมปฺปธาเน.\csromanspacer{} จตฺตาโร อิทฺธิปาเท.\csromanspacer{} ปญฺจินฺทฺริยานิ.\csromanspacer{} ปญฺจ พลานิ.\csromanspacer{} สตฺต โพชฺฌงฺเค.\csromanspacer{} อริยํ อฏฺฐงฺคิกํ มคฺคํ.\csromanspacer{} โสตาปตฺติผลํ.\csromanspacer{} สกทาคามิผลํ.\csromanspacer{} อนาคามิผลํ.\csromanspacer{} อรหตฺตํ\footnote{อรหตฺตผลํ (สฺยา)} สมาปชฺชิ.\csromanspacer{} สมาปชฺชติ.\csromanspacer{} สมาปนฺโน.\csromanspacer{} ตสฺส ภิกฺขุโน ราโค จตฺโต.\csromanspacer{} โทโส จตฺโต.\csromanspacer{} โมโห จตฺโต วนฺโต มุตฺโต ปหีโน ปฏินิสฺสฏฺโฐ อุกฺเขฏิโต สมุกฺเขฏิโต.\csromanspacer{} ตสฺส ภิกฺขุโน ราคา จิตฺตํ วินีวรณํ.\csromanspacer{} โทสา จิตฺตํ วิเนวรณํ.\csromanspacer{} โมหา จิตฺตํ วินีวรณนฺติ ภณนฺตสฺส อาปตฺติ ทุกฺกฏสฺส.}

\prose{\textbf{อาโรเจยฺยา}ติ อนุปสมฺปนฺนสฺส โย เต วิหาเร วสิ, โส ภิกฺขุ สุญฺญาคาเร ปฐมํ ฌานํ ฯเปฯ ทุติยํ ฌานํ.\csromanspacer{} ตติยํ ฌานํ.\csromanspacer{} จตุตฺถํ ฌานํ สมาปชฺชิ.\csromanspacer{} สมาปชฺชติ.\csromanspacer{} สมาปนฺโน.\csromanspacer{} โส ภิกฺขุ สุญฺญาคาเร จตุตฺถสฺส ฌานสฺส ลาภี.\csromanspacer{} วสี.\csromanspacer{} เตน ภิกฺขุนา สุญฺญาคาเร จตุตฺถํ ฌานํ สจฺฉิกตนฺติ ภณนฺตสฺส อาปตฺติ ทุกฺกฏสฺส.}

\prose{อาโรเจยฺยาติ อนุปสมฺปนฺนสฺส โย เต จีวรํ ปริภุญฺชิ.\csromanspacer{} โย เต ปิณฺฑปาตํ ปริภุญฺชิ.\csromanspacer{} โย เต เสนาสนํ ปริภุญฺชิ.\csromanspacer{} โย เต คิลานปฺปจฺจย\-เภสชฺช\-ปริกฺขารํ ปริภุญฺชิ, โส ภิกฺขุ สุญฺญาคาเร จตุตฺถํ ฌานํ สมาปชฺชิ.\csromanspacer{} สมาปชฺชติ.\csromanspacer{} สมาปนฺโน.\csromanspacer{} โส ภิกฺขุ สุญฺญาคาเร จตุตฺถสฺส ฌานสฺส ลาภี.\csromanspacer{} วสี.\csromanspacer{} เตน ภิกฺขุนา สุญฺญาคาเร จตุตฺถํ ฌานํ สจฺฉิกตนฺติ ภณนฺตสฺส อาปตฺติ ทุกฺกฏสฺส.}

\proseitem{76}{\csromanfolio{46}\textbf{อาโรเจยฺยา}ติ อนุปสมฺปนฺนสฺส เยน เต วิหาโร ปริภุตฺโต ฯเปฯ เยน เต จีวรํ ปริภุตฺตํ.\csromanspacer{} เยน เต ปิณฺฑปาโต ปริภุตฺโต.\csromanspacer{} เยน เต เสนาสนํ ปริภุตฺตํ.\csromanspacer{} เยน เต คิลานปฺปจฺจยเภสชฺชปริกฺขาโร ปริภุตฺโต, โส ภิกฺขุ สุญฺญาคาเร จตุตฺถํ ฌานํ สมาปชฺชิ.\csromanspacer{} สมาปชฺชติ.\csromanspacer{} สมาปนฺโน.\csromanspacer{} โส ภิกฺขุ สุญฺญาคาเร จตุตฺถสฺส ฌานสฺส ลาภี.\csromanspacer{} วสี.\csromanspacer{} เตน ภิกฺขุนา สุญฺญาคาเร จตุตฺถํ ฌานํ สจฺฉิกตนฺติ ภณนฺตสฺส อาปตฺติ ทุกฺกฏสฺส.}

\prose{อาโรเจยฺยาติ อนุปสมฺปนฺนสฺส ยํ ตฺวํ อาคมฺม วิหารํ อทาสิ ฯเปฯ จีวรํ อทาสิ.\csromanspacer{} ปิณฺฑปาตํ อทาสิ.\csromanspacer{} เสนาสนํ อทาสิ.\csromanspacer{} คิลานปฺปจฺจย\-เภสชฺช\-ปริกฺขารํ อทาสิ, โส ภิกฺขุ สุญฺญาคาเร จตุตฺถํ ฌานํ สมาปชฺชิ.\csromanspacer{} สมาปชฺชติ.\csromanspacer{} สมาปนฺโน.\csromanspacer{} โส ภิกฺขุ สุญฺญาคาเร จตุตฺถสฺส ฌานสฺส ลาภี.\csromanspacer{} วสี.\csromanspacer{} เตน ภิกฺขุนา สุญฺญาคาเร จตุตฺถํ ฌานํ สจฺฉิกตนฺติ ภณนฺตสฺส อาปตฺติ ทุกฺกฏสฺส.}

\proseitemwithcloserruled{77}{อนาปตฺติ อุปสมฺปนฺนสฺส ภูตํ อาโรเจติ, อาทิกมฺมิกสฺสาติ.}{ภูตาโรจนสิกฺขาปทํ นิฏฺฐิตํ อฏฺฐมํ.}
\csromanmatikaanchor{mtk.11}
\csromantocmarkref{subsection}{9. ทุฏฺฐุลฺลาโรจนสิกฺขาปท}{mtk.11}
\bookmark[dest={mtk.11},level=2]{9. ทุฏฺฐุลฺลาโรจนสิกฺขาปท}
\csromanheader{2}{1. มุสาวาทวคฺค 9. ทุฏฺฐุลฺลา\-โรจน\-สิกฺขาปท}
\proseitem{78}{เตน สมเยน พุทฺโธ ภควา สาวตฺถิยํ วิหรติ เชตวเน อนาถปิณฺฑิกสฺส อาราเม.\csromanspacer{} เตน โข ปน สมเยน อายสฺมา อุปนนฺโท สกฺยปุตฺโต ฉพฺพคฺคิเยหิ ภิกฺขูหิ สทฺธิํ ภณฺฑนกโต โหติ, โส สญฺเจตนิกํ สุกฺกวิสฺสฏฺฐิํ อาปตฺติํ อาปชฺชิตฺวา สํฆํ ตสฺสา อาปตฺติยา ปริวาสํ ยาจิ, ตสฺส สํโฆ ตสฺสา อาปตฺติยา ปริวาสํ อทาสิ.\csromanspacer{} เตน โข ปน สมเยน สาวตฺถิยํ อญฺญตรสฺส ปูคสฺส สํฆภตฺตํ โหติ, โส ปริวสนฺโต ภตฺตคฺเค อาสนปริยนฺเต นิสีทิ.\csromanspacer{} ฉพฺพคฺคิยา ภิกฺขู เต อุปาสเก เอตทโวจุํ “เอโส อาวุโส อายสฺมา อุปนนฺโท สกฺยปุตฺโต ตุมฺหากํ สมฺภาวิโต กุลูปโก เยนว หตฺเถน สทฺธาเทยฺยํ ภุญฺชติ, เตเนว หตฺเถน อุปกฺกมิตฺวา อสุจิํ โมเจสิ, โส สญฺเจตนิกํ สุกฺกวิสฺสฏฺฐิํ อาปตฺติํ อาปชฺชิตฺวา สํฆํ ตสฺสา อาปตฺติยา ปริวาสํ ยาจิ, ตสฺส สํโฆ ตสฺสา อาปตฺติยา ปริวาสํ \csromanfolio{47}อทาสิ, โส ปริวสนฺโต อาสนปริยนฺเต นิสินฺโน”ติ.\csromanspacer{} เย เต ภิกฺขู อปฺปิจฺฉา ฯเปฯ เต อุชฺฌายนฺติ ขิยฺยนฺติ วิปาเจนฺติ “กถํ หิ นาม ฉพฺพคฺคิยา ภิกฺขู \textbf{ภิกฺขุสฺส ทุฏฺฐุลฺลํ อาปตฺติํ อนุปสมฺปนฺนสฺส อาโรเจสฺสนฺตี”ติ ฯเปฯ} “\textbf{สจฺจํ กิร ตุเมฺห ภิกฺขเว ภิกฺขุสฺส ทุฏฺฐุลฺลํ อาปตฺติํ} \textbf{อนุปสมฺปนฺนสฺส อาโรเจถา”ติ.}\csromanspacer{}\textbf{ สจฺจํ ภควาติ.}\csromanspacer{}\textbf{ วิครหิ พุทฺโธ} ภควา ฯเปฯ.\csromanspacer{} กถํ หิ นาม ตุเมฺห โมฆปุริสา ภิกฺขุสฺส ทุฏฺฐุลฺลํ อาปตฺติํ อนุปสมฺปนฺนสฺส อาโรเจสฺสถ.\csromanspacer{} เนตํ โมฆปุริสา อปฺปสนฺนานํ วา ปสาทาย ฯเปฯ.\csromanspacer{} เอวญฺจ ปน ภิกฺขเว อิมํ สิกฺขาปทํ อุทฺทิเสยฺยาถ\csromandash{}}

\proseitem{79}{\textbf{“โย ปน ภิกฺขุ ภิกฺขุสฺส ทุฏฺฐุลฺลํ อาปตฺติํ อนุปสมฺปนฺนสฺส อาโรเจยฺย อญฺญตฺร ภิกฺขุสมฺมุติยา}\footnote{ภิกฺขุสมฺมติยา (สฺยา)}\textbf{, ปาจิตฺติยนฺ”}ติ.}

\proseitem{80}{\textbf{โย ปนา}ติ โย ยาทิโส ฯเปฯ.\csromanspacer{} \textbf{ภิกฺขู}ติ ฯเปฯ อยํ อิมสฺมิํ อตฺเถ อธิปฺเปโต ภิกฺขูติ.}

\prose{\textbf{ภิกฺขุสฺสา}ติ อญฺญสฺส ภิกฺขุสฺส.}

\prose{\textbf{ทุฏฺฐุลฺลา} นาม อาปตฺติ จตฺตาริ จ ปาราชิกานิ เตรส จ สํฆาทิเสสา.}

\prose{\textbf{อนุปสมฺปนฺโน} นาม ภิกฺขุญฺจ ภิกฺขุนิญฺจ ฐเปตฺวา อวเสโส อนุปสมฺปนฺโน นาม.}

\prose{\textbf{อาโรเจยฺยา}ติ อาโรเจยฺย อิตฺถิยา วา ปุริสสฺส วา คหฏฺฐสฺส วา ปพฺพชิตสฺส วา.}

\prose{\textbf{อญฺญตฺร ภิกฺขุสมฺมุติยา}ติ ฐเปตฺวา ภิกฺขุสมฺมุติํ.}

\prose{อตฺถิ ภิกฺขุสมฺมุติ อาปตฺติปริยนฺตา น กุลปริยนฺตา, อตฺถิ ภิกฺขุสมฺมุติ กุลปริยนฺตา น อาปตฺติปริยนฺตา, อตฺถิ ภิกฺขุสมฺมุติ อาปตฺติปริยนฺตา จ กุลปริยนฺตา จ, อตฺถิ ภิกฺขุสมฺมุติ เนว อาปตฺติปริยนฺตา น กุลปริยนฺตา.}

\prose{\textbf{อาปตฺติปริยนฺตา} นาม อาปตฺติโย ปริคฺคหิตาโย โหนฺติ “เอตฺตกาหิ อาปตฺตีติ อาโรเจตพฺโพ”ติ.}

\prose{\csromanfolio{48}\textbf{กุลปริยนฺตา} นาม กุลานิ ปริคฺคหิตานิ โหนฺติ “เอตฺตเกสุ กุเลสุ อาโรเจตพฺโพ”ติ.\csromanspacer{} \textbf{อาปตฺติปริยนฺตา จ กุลปริยนฺตา จ} นาม อาปตฺติโย จ ปริคฺคหิตาโย โหนฺติ, กุลานิ จ ปริคฺคหิตานิ โหนฺติ “เอตฺตกาหิ อาปตฺตีหิ เอตฺตเกสุ กุเลสุ อาโรเจตพฺโพ”ติ.\csromanspacer{} \textbf{เนว อาปตฺติปริยนฺตา น กุลปริยนฺตา} นาม อาปตฺติโย จ อปริคฺคหิตาโย โหนฺติ, กุลานิ จ อปริคฺคหิตานิ โหนฺติ “เอตฺตกาหิ อาปตฺตีหิ เอตฺตเกสุ กุเลสุ อาโรเจตพฺโพ”ติ.}

\proseitem{81}{อาปตฺติปริยนฺเต ยา อาปตฺติโย ปริคฺคหิตาโย โหนฺติ, ตา อาปตฺติโย ฐเปตฺวา อญฺญาหิ อาปตฺตีหิ อาโรเจติ, อาปตฺติ ปาจิตฺติยสฺส.}

\prose{กุลปริยนฺเต ยานิ กุลานิ ปริคฺคหิตานิ โหนฺติ, ตานิ กุลานิ ฐเปตฺวา อญฺเญสุ กุเลสุ อาโรเจติ, อาปตฺติ ปาจิตฺติยสฺส.}

\prose{อาปตฺติปริยนฺเต จ กุลปริยนฺเต จ ยา อาปตฺติโย ปริคฺคหิตาโย โหนฺติ, ตา อาปตฺติโย ฐเปตฺวา.\csromanspacer{} ยานิ กุลานิ ปริคฺคหิตานิ โหนฺติ, ตานิ กุลานิ ฐเปตฺวา อญฺญาหิ อาปตฺตีหิ อญฺเญสุ กุเลสุ อาโรเจติ, อาปตฺติ ปาจิตฺติยสฺส.}

\prose{เนว อาปตฺติปริยนฺเต น กุลปริยนฺเต อนาปตฺติ.}

\proseitem{82}{ทุฏฺฐุลฺลาย อาปตฺติยา ทุฏฺฐุลฺลาปตฺติสญฺญี อนุปสมฺปนฺนสฺส อาโรเจติ อญฺญตฺร ภิกฺขุสมฺมุติยา, อาปตฺติ ปาจิตฺติยสฺส.}

\prose{ทุฏฺฐุลฺลาย อาปตฺติยา เวมติโก อนุปสมฺปนฺนสฺส อาโรเจติ อญฺญตฺร ภิกฺขุสมฺมุติยา, อาปตฺติ ปาจิตฺติยสฺส.}

\prose{ทุฏฺฐุลฺลาย อาปตฺติยา อทุฏฺฐุลฺลาปตฺติสญฺญี อนุปสมฺปนฺนสฺส อาโรเจติ อญฺญตฺร ภิกฺขุสมฺมุติยา, อาปตฺติ ปาจิตฺติยสฺส.}

\prose{อทุฏฺฐุลฺลํ อาปตฺติํ อาโรเจติ, อาปตฺติ ทุกฺกฏสฺส.}

\prose{อนุปสมฺปนฺนสฺส ทุฏฺฐุลฺลํ วา อทุฏฺฐุลฺลํ วา อชฺฌาจารํ อาโรเจติ, อาปตฺติ ทุกฺกฏสฺส.}

\prose{\csromanfolio{49}อทุฏฺฐุลฺลาย อาปตฺติยา ทุฏฺฐุลฺลาปตฺติสญฺญี, อาปตฺติ ทุกฺกฏสฺส.}

\prose{อทุฏฺฐุลฺลาย อาปตฺติยา เวมติโก, อาปตฺติ ทุกฺกฏสฺส.}

\prose{อทุฏฺฐุลฺลาย อาปตฺติยา อทุฏฺฐุลฺลาปตฺติสญฺญี, อาปตฺติ ทุกฺกฏสฺส.}

\proseitemwithcloserruled{83}{อนาปตฺติ วตฺถุํ อาโรเจติ โน อาปตฺติํ, อาปตฺติํ อาโรเจติ โน วตฺถุํ, ภิกฺขุสมฺมุติยา อุมฺมตฺตกสฺส อาทิกมฺมิกสฺสาติ.}{ทุฏฺฐุลฺลา\-โรจน\-สิกฺขาปทํ นิฏฺฐิตํ นวมํ.}
\csromanmatikaanchor{mtk.12}
\csromantocmarkref{subsection}{10. ปถวีขณนสิกฺขาปท}{mtk.12}
\bookmark[dest={mtk.12},level=2]{10. ปถวีขณนสิกฺขาปท}
\csromanheader{2}{1. มุสาวาทวคฺค 10. \textbf{ปถวี}ขณนสิกฺขาปท}
\proseitem{84}{เตน สมเยน พุทฺโธ ภควา อาฬวิยํ วิหรติ อคฺคาฬเว เจติเย.\csromanspacer{} เตน โข ปน สมเยน อาฬวกา\footnote{อาฬวิกา (สฺยา)} ภิกฺขู นวกมฺมํ กโรนฺตา ปถวิํ ขณนฺติปิ ขณาเปนฺติปิ.\csromanspacer{} มนุสฺสา อุชฺฌายนฺติ ขิยฺยนฺติ วิปาเจนฺติ “กถํ หิ นาม สมณา สกฺยปุตฺติยา ปถวิํ ขณิสฺสนฺติปิ ขณาเปสฺสนฺติปิ.\csromanspacer{} เอกินฺทฺริยํ สมณา สกฺยปุตฺติยา ชีวํ วิเหฏฺเฐนฺตี”ติ.\csromanspacer{} อสฺโสสุํ โข ภิกฺขู เตสํ มนุสฺสานํ อุชฺฌายนฺตานํ ขิยฺยนฺตานํ วิปาเจนฺตานํ.\csromanspacer{} เย เต ภิกฺขู อปฺปิจฺฉา ฯเปฯ เต อุชฺฌายนฺติ ขิยฺยนฺติ วิปาเจนฺติ “กถํ หิ นาม อาฬวกา ภิกฺขู ปถวิํ ขณิสฺสนฺติปิ ขณาเปสฺสนฺติปี”ติ ฯเปฯ “สจฺจํ กิร ตุเมฺห ภิกฺขเว ปถวิํ ขณถปิ ขณาเปถปี”ติ.\csromanspacer{} สจฺจํ ภควาติ.\csromanspacer{} วิครหิ พุทฺโธ ภควา ฯเปฯ.\csromanspacer{} กถํ หิ นาม ตุเมฺห โมฆปุริสา ปถวิํ ขณิสฺสถปิ ขณาเปสฺสถปิ.\csromanspacer{} ชีวสญฺญิโน หิ โมฆปุริสา มนุสฺสา ปถวิยา.\csromanspacer{} เนตํ โมฆปุริสา อปฺปสนฺนานํ วา ปสาทาย ฯเปฯ.\csromanspacer{} เอวญฺจ ปน ภิกฺขเว อิมํ สิกฺขปทํ อุทฺทิเสยฺยาถ\csromandash{}}

\proseitem{85}{\textbf{“โย ปน ภิกฺขุ ปถวิํ ขเณยฺย วา ขณาเปยฺย วา ปาจิตฺติยนฺ”}ติ.}

\proseitem{86}{\textbf{โย ปนา}ติ โย ยาทิโส ฯเปฯ.\csromanspacer{} \textbf{ภิกฺขู}ติ ฯเปฯ อยํ อิมสฺมิํ อตฺเถ อธิปฺเปโต ภิกฺขูติ.}

\prose{\textbf{ปถวี} นาม เทฺว ปถวิโย ชาตา จ ปถวี อชาตา จ ปถวี.}

\prose{\csromanfolio{50}\textbf{ชาตา} นาม ปถวี สุทฺธปํสุ สุทฺธมตฺติกา อปฺปปาสาณา อปฺปสกฺขรา อปฺปกฐลา อปฺปมรุมฺพา อปฺปวาลิกา เยภุยฺเยน ปํสุกา เยภุยฺเยน มตฺติกา, อทฑฺฒาปิ วุจฺจติ ชาตา ปถวี.\csromanspacer{} โยปิ ปํสุปุญฺโช วา มตฺติกาปุญฺโช วา อติเรกจาตุมาสํ โอวฏฺโฐ, อยมฺปิ วุจฺจติ ชาตา ปถวี.}

\prose{\textbf{อชาตา} นาม ปถวี สุทฺธปาสาณา สุทฺธสกฺขรา สุทฺธกฐลา สุทฺธมรุมฺพา สุทฺธวาลิกา อปฺปปํสุกา อปฺปมตฺติกา เยภุยฺเยน ปาสาณา เยภุยฺเยน สกฺขรา เยภุยฺเยน กฐลา เยภุยฺเยน มรุมฺพา เยภุยฺเยน วาลิกา, ทฑฺฒาปิ วุจฺจติ อชาตา ปถวี.\csromanspacer{} โยปิ ปํสุปุญฺโช วา มตฺติกาปุญฺโช วา โอมกจาตุมาสํ โอวฏฺโฐ, อยมฺปิ วุจฺจติ อชาตา ปถวี.}

\prose{\textbf{ขเณยฺยา}ติ สยํ ขณติ, อาปตฺติ ปาจิตฺติยสฺส.}

\prose{\textbf{ขณาเปยฺยา}ติ อญฺญํ อาณาเปติ, อาปตฺติ ปาจิตฺติยสฺส.\csromanspacer{} สกิํ อาณตฺโต พหุกํปิ ขณติ, อาปตฺติ ปาจิตฺติยสฺส.}

\proseitem{87}{ปถวิยา ปถวิสญฺญี ขณติ วา ขณาเปติ วา, ภินฺทติ วา เภทาเปติ วา, ทหติ วา ทหาเปติ วา, อาปตฺติ ปาจิตฺติยสฺส.}

\prose{ปถวิยา เวมติโก ขณติ วา ขณาเปติ วา, ภินฺทติ วา เภทาเปติ วา, ทหติ วา ทหาเปติ วา, อาปตฺติ ทุกฺกฏสฺส.}

\prose{ปถวิยา อปถวิสญฺญี ขณติ วา ขณาเปติ วา, ภินฺทติ วา เภทาเปติ วา, ทหติ วา ทหาเปติ วา, อนาปตฺติ.}

\prose{อปถวิยา ปถวิสญฺญี, อาปตฺติ ทุกฺกฏสฺส.\csromanspacer{} อปถวิยา เวมติโก อาปตฺติ ทุกฺกฏสฺส.\csromanspacer{} อปถวิยา อปถวิสญฺญี, อนาปตฺติ.}

\proseitemwithcloser{88}{อนาปตฺติ “อิมํ ชาน, อิมํ เทหิ อิมํ อาหร, อิมินา อตฺโถ อิมํ กปฺปิยํ กโรหี”ติ ภณติ, อสญฺจิจฺจ อสฺสติยา อชานนฺตสฺส อุมฺมตฺตกสฺส อาทิกมฺมิกสฺสาติ.}{ปถวีขณนสิกฺขาปทํ นิฏฺฐิตํ ทสมํ.}
\vspace{\tipitakaparskip}
\csromancenter{มุสาวาทวคฺโค ปฐโม.}
\csromancenter{ตสฺสุทฺทานํ}
\vspace{0.5\baselineskip}
\setlength{\csromangathapagewidth}{0pt}
\csromangathameasuregroup{มุสา โอมสเปสุญฺญํ,}{\csromangathabat{มุสา โอมสเปสุญฺญํ,}{ปทเสยฺยาย เว ทุเว.} \\ อญฺญตฺร วิญฺญุนา ภูตา, ทุฏฺฐุลฺลาปตฺติ ขณนาติ.}
\csromangathasetleft{มุสา โอมสเปสุญฺญํ,}
\csromangathagroup{\csromangathastanza{\csromanfolio{51}\csromangathabat{มุสา โอมสเปสุญฺญํ,}{ปทเสยฺยาย เว ทุเว.} \\ อญฺญตฺร วิญฺญุนา ภูตา, ทุฏฺฐุลฺลาปตฺติ ขณนาติ.}}

\csromanmatikaanchor{mtk.13}
\csromanmatikaanchor{mtk.14}
\csromantocmarknopage{section}{2. ภูตคามวคฺค}{mtk.13}
\bookmark[dest={mtk.13},level=1]{2. ภูตคามวคฺค}
\csromantocmarkref{subsection}{1. ภูตคามสิกฺขาปท}{mtk.14}
\bookmark[dest={mtk.14},level=2]{1. ภูตคามสิกฺขาปท}
\setcsromanheadh{2. ภูตคามวคฺค}
\csromanheader{2}{2. \textbf{ภูตคามวคฺค 1. ภูตคามสิกฺขาปท}}
\proseitem{89}{เตน สมเยน พุทฺโธ ภควา อาฬวิยํ วิหรติ อคฺคาฬเว เจติเย, เตน โข ปน สมเยน อาฬวกา ภิกฺขู นวกมฺมํ กโรนฺตา รุกฺขํ ฉินฺทนฺติปิ เฉทาเปนฺติปิ, อญฺญตโรปิ อาฬวโก ภิกฺขุ รุกฺขํ ฉินฺทติ.\csromanspacer{} ตสฺมิํ รุกฺเข อธิวตฺถา เทวตา ตํ ภิกฺขุํ เอตทโวจ “มา ภนฺเต อตฺตโน ภวนํ กตฺตุกาโม มยฺหํ ภวนํ ฉินฺที”ติ.\csromanspacer{} โส ภิกฺขุ อนาทิยนฺโต ฉินฺทิ เยว, ตสฺสา จ เทวตาย ทารกสฺส พาหุํ อาโกเฏสิ.\csromanspacer{} อถ โข ตสฺสา เทวตาย เอตทโหสิ “ยนฺนูนาหํ อิมํ ภิกฺขุํ อิเธว ชีวิตา โวโรเปยฺยนฺ”ติ, อถ โข ตสฺสา เทวตาย เอตทโหสิ “น โข เมตํ ปติรูปํ ยาหํ อิมํ ภิกฺขุํ อิเธว ชีวิตา โวโรเปยฺยํ, ยนฺนูนาหํ ภควโต เอตมตฺถํ อาโรเจยฺยนฺ”ติ.\csromanspacer{} อถ โข สา เทวตา เยน ภควา เตนุปสงฺกมิ, อุปสงฺกมิตฺวา ภควโต เอตมตฺถํ อาโรเจสิ.\csromanspacer{} สาธุ สาธุ เทวเต, สาธุ โข ตฺวํ เทวเต ตํ ภิกฺขุํ ชีวิตา น โวโรเปสิ, สจชฺช ตฺวํ เทวเต ตํ ภิกฺขุํ ชีวิตา โวโรเปยฺยาสิ, พหุญฺจ ตฺวํ เทวเต อปุญฺญํ ปสเวยฺยาสิ.\csromanspacer{} คจฺฉ ตฺวํ เทวเต อมุกสฺมิํ โอกาเส รุกฺโข วิวิตฺโต, ตสฺมิํ อุปคจฺฉาติ.\csromanspacer{} มนุสฺสา อุชฺฌายนฺติ ขิยฺยนฺติ วิปาเจนฺติ “กถํ หิ นาม สมณา สกฺยปุตฺติยา รุกฺขํ ฉินฺทิสฺสนฺติปิ เฉทาเปสฺสนฺติปิ.\csromanspacer{} เอกินฺทฺริยํ สมณา สกฺยปุตฺติยา ชีวํ วิเหเฐนฺตี”ติ.}

\prose{อสฺโสสุํ โข ภิกฺขู เตสํ มนุสฺสานํ อุชฺฌายนฺตานํ ขิยฺยนฺตานํ วิปาเจนฺตานํ.\csromanspacer{} เย เต ภิกฺขู อปฺปิจฺฉา ฯเปฯ เต อุชฺฌายนฺติ ขิยฺยนฺติ วิปาเจนฺติ “กถํ หิ นาม อาฬาวกา ภิกฺขู รูกฺขํ ฉินฺทิสฺสนฺติปิ เฉทาเปสฺสนฺติปี”ติ ฯเปฯ “สจฺจํ กิร ตุเมฺห ภิกฺขเว รุกฺขํ ฉินฺทถาปิ เฉทาเปสาปี”ติ.\csromanspacer{} สจฺจํ ภควาติ.\csromanspacer{} วิครหิ พุทฺโธ ภควา ฯเปฯ.\csromanspacer{} กถํ หิ นาม ตุเมฺห โมฆปุริสา รุกฺขํ ฉินฺทิสฺสถาปิ เฉทาเปสฺสถาปิ.\csromanspacer{} ชีวสญฺญิโน หิ โมฆปุริสา มนุสฺสา \csromanfolio{52}\textbf{รุกฺขสฺมิํ.}\csromanspacer{}\textbf{ เนตํ โมฆปุริสา อปฺปสนฺนานํ วา ปสาทาย ฯเปฯ.}\csromanspacer{}\textbf{ เอวญฺจ} \textbf{ปน ภิกฺขเว อิมํ สิกฺขาปทํ อุทฺทิเสยฺยาถ\csromandash{}}}

\proseitem{90}{\textbf{“ภูตคามปาตพฺยตาย ปาจิตฺติยนฺ”}ติ.}

\proseitem{91}{\textbf{ภูตคาโม} นาม ปญฺจ พีชชาตานิ, มูลพีชํ ขนฺธพีชํ ผฬุพีชํ อคฺคพีชํ พีชพีชเมว\footnote{พีชพีชญฺเจว (อิติปิ)} ปญฺจมํ.}

\prose{\textbf{มูลพีชํ} นาม หลิทฺทิ สิงฺคิเวรํ วจา วจตฺตํ อติวิสา กฏุกโรหิณี อุสีรํ ภทฺทมุตฺตกํ, ยานิ วา ปนญฺญานิปิ อตฺถิ มูเล ชายนฺติ มูเล สญฺชายนฺติ, เอตํ มูลพีชํ นาม.}

\prose{\textbf{ขนฺธพีชํ} นาม อสฺสตฺโถ นิโคฺรโธ ปิลกฺโข อุทุมฺพโร กจฺฉโก กปิตฺถโน, ยานิ วา ปนญฺญานิปิ อตฺถิ ขนฺเธ ชายนฺติ ขนฺเธ สญฺชายนฺติ, เอตํ ขนฺธพีชํ นาม.}

\prose{\textbf{ผฬุพีชํ} นาม อุจฺฉุ เวฬุ นโฬ, ยานิ วา ปนญฺญานิปิ อตฺถิ ปพฺเพ ชายนฺติ ปพฺเพ สญฺชายนฺติ, เอตํ ผฬุพีชํ นาม.}

\prose{\textbf{อคฺคพีชํ} นาม อชฺชุกํ ผณิชฺชกํ หิริเวรํ, ยานิ วา ปนญฺญานิปิ อตฺถิ อคฺเค ชายนฺติ อคฺเค สญฺชายนฺติ, เอตํ อคฺคพีชํ นาม.}

\prose{\textbf{พีชพีชํ} นาม ปุพฺพณฺณํ อปรณฺณํ, ยานิ วา ปนญฺญานิปิ อตฺถิ พีเช ชายนฺติ พีเช สญฺชายนฺติ, เอตํ พีชพีชํ นาม.}

\proseitem{92}{พีเช พีชสญฺญี ฉินฺทติ วา เฉทาเปติ วา, ภินฺทติ วา เภทาเปติ วา, ปจติ วา ปจาเปติ วา, อาปตฺติ ปาจิตฺติยสฺส.\csromanspacer{} พีเช เวมติโก ฉินฺทติ วา เฉทาเปติ วา, ภินฺทติ วา เภทาเปติ วา, ปจติ วา ปจาเปติ วา, อาปตฺติ ทุกฺกฏสฺส.\csromanspacer{} พีเช อพีชสญฺญี ฉินฺทติ วา เฉทาเปติ วา, ภินฺทติ วา เภทาเปติ วา, ปจติ วา ปจาเปติ วา, อนาปตฺติ.\csromanspacer{} อพีเช พีชสญฺญี อาปตฺติ ทุกฺกฏสฺส.\csromanspacer{} อพีเช เวมติโก, อาปตฺติ ทุกฺกฏสฺส.\csromanspacer{} อพีเช อพีชสญฺญี, อนาปตฺติ.}

\proseitemwithcloser{93}{อนาปตฺติ “อิมํ ชาน, อิมํ เทหิ, อิมํ อาหร, อิมินา อตฺโถ, อิมํ กปฺปิยํ กโรหี”ติ ภณติ, อสญฺจิจฺจ อสฺสติยา อชานนฺตสฺส อุมฺมตฺตกสฺส อาทิกมฺมิกสฺสาติ.}{ภูตคามสิกฺขาปทํ นิฏฺฐิตํ ปฐมํ.}
\csromanmatikaanchor{mtk.15}
\csromantocmarkref{subsection}{2. อญฺญวาทกสิกฺขาปท}{mtk.15}
\bookmark[dest={mtk.15},level=2]{2. อญฺญวาทกสิกฺขาปท}
\csromanheader{2}{2. \textbf{ภูตคามวคฺค 2. อญฺญวาทกสิกฺขาปท}}
\proseitem{94}{\csromanfolio{53}เตน สมเยน พุทฺโธ ภควา โกสมฺพิยํ วิหรติ โฆสิตาราเม.\csromanspacer{} เตน โข ปน สมเยน อายสฺมา ฉนฺโน อนาจารํ อาจริตฺวา สํฆมชฺเฌ อาปตฺติยา อนุยุญฺชียมาโน อญฺเญนญฺญํ ปฏิจรติ “โก อาปนฺโน กิํ อาปนฺโน กิสฺมิํ อาปนฺโน กถํ อาปนฺโน กํ ภณถ กิํ ภณถา”ติ.\csromanspacer{} เย เต ภิกฺขู อปิจฺฉา ฯเปฯ เต อุชฺฌายนฺติ ขิยฺยนฺติ วิปาเจนฺติ “กถํ หิ นาม อายสฺมา ฉนฺโน สํฆมชฺเฌ อาปตฺติยา อนุยุญฺชียมาโน อญฺเญนญฺญํ ปฏิจริสฺสติ โก อาปนฺโน กิํ อาปนฺโน กิสฺมิํ อาปนฺโน กถํ อาปนฺโน กํ ภณถ กิํ ภณถา”ติ ฯเปฯ.\csromanspacer{} สจฺจํ กิร ตฺวํ ฉนฺน สํฆมชฺเฌ อาปตฺติยา อนุยุญฺชียมาโน อญฺเญนญฺญํ ปฏิจรสิ “โก อาปนฺโน กิํ อาปนฺโน กิสฺมิํ อาปนฺโน กถํ อาปนฺโน กํ ภณถ กิํ ภณถา”ติ.\csromanspacer{} สจฺจํ ภควาติ.\csromanspacer{} วิครหิ พุทฺโธ ภควา ฯเปฯ.\csromanspacer{} กถํ หิ นาม ตฺวํ โมฆปุริส สํฆมชฺเฌ อาปตฺติยา อนุยุญฺชียมาโน อญฺเญนญฺญํ ปฏิจริสฺสสิ “โก อาปนฺโน กิํ อาปนฺโน กิสฺมิํ อาปนฺโน กถํ อาปนฺโน กํ ภณถ กิํ ภณถา”ติ.\csromanspacer{} เนตํ โมฆปุริส อปฺปสนฺนานํ วา ปสาทาย ฯเปฯ วิครหิตฺวา ฯเปฯ ธมฺมิํ กถํ กตฺวา ภิกฺขู อามนฺเตสิ เตน หิ ภิกฺขเว สํโฆ ฉนฺนสฺส ภิกฺขุโน อญฺญวาทกํ โรเปตุ.\csromanspacer{} เอวญฺจ ปน ภิกฺขเว โรเปตพฺพํ, พฺยตฺเตน ภิกฺขุนา ปฏิพเลน สํโฆ ญาเปตพฺโพ\csromandash{}}

\proseitem{95}{“สุณาตุ เม ภนฺเต สํโฆ, อยํ ฉนฺโน ภิกฺขุ สํฆมชฺเฌ อาปตฺติยา อนุยุญฺชียมาโน อญฺเญนญฺญํ ปฏิจรติ, ยทิ สํฆสฺส ปตฺตกลฺลํ, สํโฆ ฉนฺนสฺส ภิกฺขุโน อญฺญวาทกํ โรเปยฺย, เอสา ญตฺติ.}

\prose{สุณาตุ เม ภนฺเต สํโฆ, อยํ ฉนฺโน ภิกฺขุ สํฆมชฺเฌ อาปตฺติยา อนุยุญฺชียมาโน อญฺเญนญฺญํ ปฏิจรติ, สํโฆ ฉนฺนสฺส ภิกฺขุโน อญฺญวาทกํ โรเปติ, ยสฺสายสฺมโต ขมติ ฉนฺนสฺส ภิกฺขุโน อญฺญวาทกสฺส โรปนา, โส ตุณฺหสฺส, ยสฺส นกฺขมติ, โส ภาเสยฺย.}

\prose{โรปิตํ สํเฆน ฉนฺนสฺส ภิกฺขุโน อญฺญวาทกํ, ขมติ สํฆสฺส, ตสฺมา ตุณฺหี, เอวเมตํ ธารยามี”ติ.}

\prose{\csromanfolio{54}อถ โข ภควา อายสฺมนฺตํ ฉนฺนํ อเนกปริยาเยน วิครหิตฺวา ทุพฺภรตาย ฯเปฯ.\csromanspacer{} เอวญฺจ ปน ภิกฺขเว อิมํ สิกฺขาปทํ อุทฺทิเสยฺยาถ\csromandash{}}

\prose{\textbf{“อญฺญวาทเก ปาจิตฺติยนฺ”}ติ.}

\prose{เอวญฺจิทํ ภควตา ภิกฺขูนํ สิกฺขาปทํ ปญฺญตฺตํ โหติ.}

\proseitem{96}{เตน โข ปน สมเยน อายสฺมา ฉนฺโน สํฆมชฺเฌ อาปตฺติยา อนุยุญฺชียมาโน “อญฺเญนญฺญํ ปฏิจรนฺโต อาปตฺติํ อาปชฺชิสฺสามี”ติ ตุณฺหีภูโต สํฆํ วิเหเสติ.\csromanspacer{} เย เต ภิกฺขู อปฺปิจฺฉา ฯเปฯ เต อุชฺฌายนฺติ ขิยฺยนฺติ วิปาเจนฺติ “กถํ หิ นาม อายสฺมา ฉนฺโน สํฆมชฺเฌ อาปตฺติยา อนุยุญฺชียมาโน ตุณฺหีภูโต สํฆํ วิเหเสสฺสตี”ติ ฯเปฯ “สจฺจํ กิร ตฺวํ ฉนฺน สํฆมชฺเฌ อาปตฺติยา อนุยุญฺชียมาโน ตุณฺหีภูโต สํฆํ วิเหเสสี”ติ.\csromanspacer{} สจฺจํ ภควาติ.\csromanspacer{} วิครหิ พุทฺโธ ภควา ฯเปฯ.\csromanspacer{} กถํ หิ นาม ตฺวํ โมฆปุริส สํฆมชฺเฌ อาปตฺติยา อนุยุญฺชียมาโน ตุณฺหีภูโต สํฆํ วิเหเสสฺสสิ.\csromanspacer{} เนตํ โมฆปุริส อปฺปสนฺนานํ วา ปสาทาย ฯเปฯ วิครหิตฺวา ฯเปฯ ธมฺมิํ กถํ กตฺวา ภิกฺขู อามนฺเตสิ เตน หิ ภิกฺขเว สํโฆ ฉนฺนสฺส ภิกฺขุโน วิเหสกํ โรเปตุ.\csromanspacer{} เอวญฺจ ปน ภิกฺขเว โรเปตพฺพํ พฺยตฺเตน ภิกฺขุนา ปฏิพเลน สํโฆ ญาเปตพฺโพ\csromandash{}}

\proseitem{97}{“สุณาตุ เม ภนฺเต สํโฆ, อยํ ฉนฺโน ภิกฺขุ สํฆมชฺเฌ อาปตฺติยา อนุยุญฺชียมาโน ตุณฺหีภูโต สํฆํ วิเหเสติ, ยทิ สํฆสฺส ปตฺตกลฺลํ, สํโฆ ฉนฺนสฺส ภิกฺขุโน วิเหสกํ โรเปยฺย, เอสา ญตฺติ.}

\prose{สุณาตุ เม ภนฺเต สํโฆ, อยํ ฉนฺโน ภิกฺขุ สํฆมชฺเฌ อาปตฺติยา อนุยุญฺชียมาโน ตุณฺหีภูโต สํฆํ วิเหเสติ, สํโฆ ฉนฺนสฺส ภิกฺขุโน วิเหสกํ โรเปติ, ยสฺสายสฺมาโต ขมติ ฉนฺนสฺส ภิกฺขุโน วิเหสกสฺส โรปนา, โส ตุณฺหสฺส, ยสฺส นกฺขมติ, โส ภาเสยฺย.}

\prose{โรปิตํ สํเฆน ฉนฺนสฺส ภิกฺขุโน วิเหสกํ, ขมติ สํฆสฺส ตสฺมา ตุณฺหี, เอวเมตํ ธารยามี”ติ.}

\prose{อถ โข ภควา อายสฺมนฺตํ ฉนฺนํ อเนกปริยาเยน วิครหิตฺวา ทุพฺภรตาย ฯเปฯ.\csromanspacer{} เอวญฺจ ปน ภิกฺขเว อิมํ สิกฺขาปทํ อุทฺทิเสยฺยาถ\csromandash{}}

\proseitem{98}{\csromanfolio{55}\textbf{“อญฺญวาทเก วิเหสเก ปาจิตฺติยนฺ”}ติ.}

\proseitem{99}{\textbf{อญฺญวาทโก} นาม สํฆมชฺเฌ วตฺถุสฺมิํ วา อาปตฺติยา วา อนุยุญฺชียมาโน ตํ น กเถตุกาโม ตํ น อุคฺฆาเฏตุกาโม อญฺเญนญฺญํ ปฏิจรติ “โก อาปนฺโน กิํ อาปนฺโน กิสฺมิํ อาปนฺโน กถํ อาปนฺโน กํ ภณถ กิํ ภณถา”ติ, เอโส อญฺญวาทโก นาม.}

\prose{\textbf{วิเหสโก} นาม สํฆมชฺเฌ วตฺถุสฺมิํ วา อาปตฺติยา วา อนุยุญฺชียมาโน ตํ น กเถตุกาโม ตํ น อุคฺฆาเฏตุกาโม ตุณฺหีภูโต สํฆํ วิเหเสติ, เอโส วิเหสโก นาม.}

\proseitem{100}{อโรปิเต อญฺญวาทเก สํฆมชฺเฌ วตฺถุสฺมิํ วา อาปตฺติยา วา อนุยุญฺชียมาโน ตํ น กเถตุกาโม ตํ น อุคฺฆาเฏตุกาโม อญฺเญนญฺญํ ปฏิจรติ “โก อาปนฺโน กิํ อาปนฺโน กิสฺมิํ อาปนฺโน กถํ อาปนฺโน กํ ภณถ กิํ ภณถา”ติ, อาปตฺติ ทุกฺกฏสฺส.\csromanspacer{} อโรปิเต วิเหสเก สํฆมชฺเฌ วตฺถุสฺมิํ วา อาปตฺติยา วา อนุยุญฺชียมาโน ตํ น กเถตุกาโม ตํ น อุคฺฆาเฏตุกาโม ตุณฺหีภูโต สํฆํ วิเหเสติ, อาปตฺติ ทุกฺกฏสฺส.\csromanspacer{} โรปิเต อญฺญวาทเก สํฆมชฺเฌ วตฺถุสฺมิํ วา อาปตฺติยา วา อนุยุญฺชียมาโน ตํ น กเถตุกาโม ตํ น อุคฺฆาเฏตุกาโม อญฺเญนญฺญํ ปฏิจรติ “โก อาปนฺโน กิํ อาปนฺโน กิสฺมิํ อาปนฺโน กถํ อาปนฺโน กํ ภณถ กิํ ภณถา”ติ, อาปตฺติ ปาจิตฺติยสฺส.\csromanspacer{} โรปิเต วิเหสเก สํฆมชฺเฌ วตฺถุสฺมิํ วา อาปตฺติยา วา อนุยุญฺชียมาโน ตํ น กเถตุกาโม ตํ น อุคฺฆาเฏตุกาโม ตุณฺหีภูโต สํฆํ วิเหเสติ, อาปตฺติ ปาจิตฺติยสฺส.}

\proseitem{101}{ธมฺมกมฺเม ธมฺมกมฺมสญฺญี อญฺญวาทเก วิเหสเก อาปตฺติ ปาจิตฺติยสฺส.\csromanspacer{} ธมฺมกมฺเม เวมติโก อญฺญวาทเก วิเหสเก อาปตฺติ ปาจิตฺติยสฺส.\csromanspacer{} ธมฺมกมฺเม อธมฺมกมฺมสญฺญี อญฺญวาทเก วิเหสเก อาปตฺติ ปาจิตฺติยสฺส.\csromanspacer{} อธมฺมกมฺเม ธมฺมกมฺมสญฺญี, อาปตฺติ ทุกฺกฏสฺส.\csromanspacer{} อธมฺมกมฺเม เวมติโก, อาปตฺติ ทุกฺกฏสฺส.\csromanspacer{} อธมฺมกมฺเม อธมฺมกมฺมสญฺญี, อาปตฺติ ทุกฺกฏสฺส.}

\proseitemwithcloserruled{102}{\csromanfolio{56}อนาปตฺติ อชานนฺโต ปุจฺฉติ, คิลาโน วา น กเถติ, “สํฆสฺส ภณฺฑนํ วา กลโห วา วิคฺคโห วา วิวาโท วา ภวิสฺสตี”ติ น กเถติ, “สํฆเภโท วา สํฆราชิ วา ภวิสฺสตี”ติ น กเถติ, “อธมฺเมน วา วคฺเคน วา นกมฺมารหสฺส วา กมฺมํ กริสฺสตี”ติ น กเถติ, อุมฺมตฺตกสฺส อาทิกมฺมิกสฺสาติ.}{อญฺญวาทกสิกฺขาปทํ นิฏฺฐิตํ ทุติยํ.}
\csromanmatikaanchor{mtk.16}
\csromantocmarkref{subsection}{3. อุชฺฌาปนกสิกฺขาปท}{mtk.16}
\bookmark[dest={mtk.16},level=2]{3. อุชฺฌาปนกสิกฺขาปท}
\csromanheader{2}{2. \textbf{ภูตคามวคฺค 3. อุชฺฌาปนกสิกฺขาปท}}
\proseitem{103}{เตน สมเยน พุทฺโธ ภควา ราชคเห วิหรติ เวฬุวเน กลนฺทกนิวาเป.\csromanspacer{} เตน โข ปน สมเยน อายสฺมา ทพฺโพ มลฺลปุตฺโต สํฆสฺส เสนาสนญฺจ ปญฺญเปติ, ภตฺตานิ จ อุทฺทิสติ.\csromanspacer{} เตน โข ปน สมเยน เมตฺติยภูมชกา\footnote{เมตฺติยภุมฺมชกา (สี, สฺยา)} ภิกฺขู นวกา เจว โหนฺติ อปฺปปุญฺญา จ.\csromanspacer{} ยานิ สํฆสฺส ลามกานิ เสนาสนานิ, ตานิ เตสํ ปาปุณนฺติ ลามกานิ จ ภตฺตานิ.\csromanspacer{} เต อายสฺมนฺตํ ทพฺพํ มลฺลปุตฺตํ ภิกฺขู อุชฺฌาเปนฺติ “ฉนฺทาย ทพฺโพ มลฺลปุตฺโต เสนาสนํ ปญฺญเปติ, ฉนฺทาย จ ภตฺตานิ อุทฺทิสตี”ติ.\csromanspacer{} เย เต ภิกฺขู อปฺปิจฺฉา ฯเปฯ เต อุชฺฌายนฺติ ขิยฺยนฺติ วิปาเจนฺติ “กถํ หิ นาม เมตฺติยภูมชกา ภิกฺขู อายสฺมนฺตํ ทพฺพํ มลฺลปุตฺตํ ภิกฺขู อุชฺฌาเปสฺสนฺตี”ติ ฯเปฯ “สจฺจํ กิร ตุเมฺห ภิกฺขเว ทพฺพํ มลฺลปุตฺตํ ภิกฺขู อุชฺฌาเปถา”ติ.\csromanspacer{} สจฺจํ ภควาติ.\csromanspacer{} วิครหิ พุทฺโธ ภควา ฯเปฯ.\csromanspacer{} กถํ หิ นาม ตุเมฺห โมฆปุริสา ทพฺพํ มลฺลปุตฺตํ ภิกฺขู อุชฺฌาเปสฺสถ.\csromanspacer{} เนตํ โมฆปุริสา อปฺปสนฺนานํ วา ปสาทาย ฯเปฯ.\csromanspacer{} เอวญฺจ ปน ภิกฺขเว อิมํ สิกฺขาปทํ อุทฺทิเสยฺยาถ\csromandash{}}

\prose{\textbf{“อุชฺฌาปนเก ปาจิตฺติยนฺ”}ติ.}

\prose{เอวญฺจิทํ ภควตา ภิกฺขูนํ สิกฺขาปทํ ปญฺญตฺตํ โหติ.}

\proseitem{104}{เตน โข ปน สมเยน เมตฺติยภูมชกา ภิกฺขู “ภควตา อุชฺฌาปนกํ ปฏิกฺขิตฺตนฺ”ติ “เอตฺตาวตา ภิกฺขู โสสฺสนฺตี”ติ\footnote{วิเหสิสฺสนฺตีติ (อิติปิ)} ภิกฺขูนํ สามนฺตา อายสฺมนฺตํ ทพฺพํ มลฺลปุตฺตํ ขิยฺยนฺติ “ฉนฺทาย ทพฺโพ มลฺลปุตฺโต \csromanfolio{57}\textbf{เสนาสนํ ปญฺญเปติ ฉนฺทาย จ ภตฺตานิ อุทฺทิสตี”ติ.}\csromanspacer{}\textbf{ เย เต ภิกฺขู} \textbf{อปฺปิจฺฉา ฯเปฯ เต อุชฺฌายนฺติ ขิยฺยนฺติ วิปาเจนฺติ “กถํ หิ นาม} เมตฺติยภูมชกา ภิกฺขู อายสฺมนฺตํ ทพฺพํ มลฺลปุตฺตํ ขิยฺยิสฺสนฺตี”ติ ฯเปฯ “สจฺจํ กิร ตุเมฺห ภิกฺขเว ทพฺพํ มลฺลปุตฺตํ ขิยฺยถา”ติ.\csromanspacer{} สจฺจํ ภควาติ.\csromanspacer{} วิครหิ พุทฺโธ ภควา ฯเปฯ.\csromanspacer{} กถํ หิ นาม ตุเมฺห โมฆปุริสา ทพฺพํ มลฺลปุตฺตํ ขิยฺยิสฺสถ.\csromanspacer{} เนตํ โมฆปุริสา อปฺปสนฺนานํ วา ปสาทาย ฯเปฯ.\csromanspacer{} เอวญฺจ ปน ภิกฺขเว อิมํ สิกฺขาปทํ อุทฺทิเสยฺยาถ\csromandash{}}

\proseitem{105}{\textbf{“อุชฺฌาปนเก ขิยฺยนเก ปาจิตฺติยนฺ”}ติ.}

\proseitem{106}{\textbf{อุชฺฌาปนกํ} นาม อุปสมฺปนฺนํ สํเฆน สมฺมตํ เสนาสนปญฺญาปกํ วา ภตฺตุทฺเทสกํ วา ยาคุภาชกํ วา ผลภาชกํ วา ขชฺชภาชกํ วา อปฺปมตฺตกวิสฺสชฺชกํ วา อวณฺณํ กตฺตุกาโม อยสํ กตฺตุกาโม มงฺกุกตฺตุกาโม อุปสมฺปนฺนํ อุชฺฌาเปติ วา ขิยฺยติ วา, อาปตฺติ ปาจิตฺติยสฺส.\csromanspacer{} ธมฺมกมฺเม ธมฺมกมฺมสญฺญี อุชฺฌาปนเก ขิยฺยนเก อาปตฺติ ปาจิตฺติยสฺส.\csromanspacer{} ธมฺมกมฺเม เวมติโก อุชฺฌาปนเก ขิยฺยนเก อาปตฺติ ปาจิตฺติยสฺส.\csromanspacer{} ธมฺมกมฺเม อธมฺมกมฺมสญฺญี อุชฺฌาปนเก ขิยฺยนเก อาปตฺติ ปาจิตฺติยสฺส.}

\prose{อนุปสมฺปนฺนํ อุชฺฌาเปติ วา ขิยฺยติ วา, อาปตฺติ ทุกฺกฏสฺส.\csromanspacer{} อุปสมฺปนฺนํ สํเฆน อสมฺมตํ เสนาสนปญฺญาปกํ วา ภตฺตุทฺเทสกํ วา ยาคุภาชกํ วา ผลภาชกํ วา ขชฺชภาชกํ วา อปฺปมตฺตกวิสฺสชฺชกํ วา อวณฺณํ กตฺตุกาโม อยสํ กตฺตุกาโม มงฺกุกตฺตุกาโม อุปสมฺปนฺนํ วา อนุปสมฺปนฺนํ วา อุชฺฌาเปติ วา ขิยฺยติ วา, อาปตฺติ ทุกฺกฏสฺส.\csromanspacer{} อนุปสมฺปนฺนํ สํเฆน สมฺมตํ วา อสมฺมตํ วา เสนาสนปญฺญาปกํ วา ภตฺตุทฺเทสกํ วา ยาคุภาชกํ วา ผลภาชกํ วา ขชฺชภาชกํ วา อปฺปมตฺตกวิสฺสชฺชกํ วา อวณฺณํ กตฺตุกาโม อยสํ กตฺตุกาโม มงฺกุกตฺตุกาโม อุปสมฺปนฺนํ วา อนุปสมฺปนฺนํ วา อุชฺฌาเปติ วา ขิยฺยติ วา, อาปตฺติ ทุกฺกฏสฺส.\csromanspacer{} อธมฺมกมฺเม ธมฺมกมฺมสญฺญี, อาปตฺติ ทุกฺกฏสฺส.\csromanspacer{} อธมฺมกมฺเม เวมติโก, อาปตฺติ ทุกฺกฏสฺส.\csromanspacer{} อธมฺมกมฺเม อธมฺมกมฺมสญฺญี, อาปตฺติ ทุกฺกฏสฺส.}

\proseitemwithcloser{107}{อนาปตฺติ ปกติยา ฉนฺทา โทสา โมหา ภยา กโรนฺตํ อุชฺฌาเปติ วา ขิยฺยติ วา, อุมฺมตฺตกสฺส อาทิกมฺมิกสฺสาติ.}{อุชฺฌาปนกสิกฺขาปทํ นิฏฺฐิตํ ตติยํ.}
\csromanmatikaanchor{mtk.17}
\csromantocmarkref{subsection}{4. ปฐมเสนาสนสิกฺขาปท}{mtk.17}
\bookmark[dest={mtk.17},level=2]{4. ปฐมเสนาสนสิกฺขาปท}
\csromanheader{2}{2. ภูตคามวคฺค 4. ปฐม\-เสนาสน\-สิกฺขาปท}
\proseitem{108}{\csromanfolio{58}เตน สมเยน พุทฺโธ ภควา สาวตฺถิยํ วิหรติ เชตวเน อนาถปิณฺฑิกสฺส อาราเม.\csromanspacer{} เตน โข ปน สมเยน ภิกฺขู เหมนฺติเก กาเล อชฺโฌกาเส เสนาสนํ ปญฺญเปตฺวา กายํ โอตาเปนฺตา กาเล อาโรจิเต ตํ ปกฺกมนฺตา เนว อุทฺธริํสุ น อุทฺธราเปสุํ อนาปุจฺฉา ปกฺกมิํสุ, เสนาสนํ โอวฏฺฐํ โหติ.\csromanspacer{} เย เต ภิกฺขู อปฺปิจฺฉา ฯเปฯ เต อุชฺฌายนฺติ ขิยฺยนฺติ วิปาเจนฺติ “กถํ หิ นาม ภิกฺขู อชฺโฌกาเส เสนาสนํ ปญฺญเปตฺวา ตํ ปกฺกมนฺตา เนว อุทฺธริสฺสนฺติ น อุทฺธราเปสฺสนฺติ อนาปุจฺฉา ปกฺกมิสฺสนฺติ, เสนาสนํ โอวฏฺฐนฺ”ติ.\csromanspacer{} อถ โข เต ภิกฺขู เต อเนปริยาเยน วิครหิตฺวา ภควโต เอตมตฺถํ อาโรเจสุํ ฯเปฯ.\csromanspacer{} สจฺจํ กิร ภิกฺขเว ภิกฺขู อชฺโฌกาเส ฯเปฯ.\csromanspacer{} เอวญฺจ ปน ภิกฺขเว อิมํ สิกฺขาปทํ อุทฺทิเสยฺยาถ\csromandash{}}

\proseitem{109}{\textbf{“โยปน ภิกฺขุ สํฆิกํ มญฺจํ วา ปีฐํ วา ภิสิํ วา โกจฺฉํ วา อชฺโฌกาเส สนฺถริตฺวา วา สนฺถราเปตฺวา วา ตํ ปกฺกมนฺโต เนว อุทฺธเรยฺย น อุทฺธราเปยฺย, อนาปุจฺฉํ วา คจฺเฉยฺย, ปาจิตฺติยนฺ”}ติ.}

\prose{เอวญฺจิทํ ภควตา ภิกฺขูนํ สิกฺขาปทํ ปญฺญตฺตํ โหติ.}

\proseitem{110}{เตน โข ปน สมเยน ภิกฺขู อชฺโฌกาเส วสิตฺวา กาลสฺเสว เสนาสนํ อภิหรนฺติ.\csromanspacer{} อทฺทสา โข ภควา เต ภิกฺขู กาลสฺเสว เสนาสนํ อภิหรนฺเต, ทิสฺวาน เอตสฺมิํ นิทาเน เอตสฺมิํ ปกรเณ ธมฺมิํ กถํ กตฺวา ภิกฺขู อามนฺเตสิ “อนุชานามิ ภิกฺขเว อฏฺฐ มาเส อวสฺสิกสงฺเกเต\footnote{วสิตุํ อวสฺสิกสํเกเต (อิติปิ)} มณฺฑเป วา รุกฺขมูเล วา ยตฺถ กากา วา กุลลา วา น อูหทนฺติ, ตตฺถ เสนาสนํ นิกฺขิปิตุนฺ”ติ.}

\proseitem{111}{\textbf{โย ปนา}ติ โย ยาทิโส ฯเปฯ.\csromanspacer{} \textbf{ภิกฺขู}ติ ฯเปฯ อยํ อิมสฺมิํ อตฺเถ อธิปฺเปโต ภิกฺขูติ.}

\prose{\textbf{สํฆิกํ} นาม สํฆสฺส ทินฺนํ โหติ ปริจฺจตฺตํ.}

\prose{\textbf{มญฺโจ} นาม จตฺตาโร มญฺจา มสารโก พุนฺทิกาพทฺโธ กุฬีรปาทโก อาหจฺจปาทโก.}

\prose{\csromanfolio{59}\textbf{ปีฐํ} นาม จตฺตาริ ปีฐานิ มสารกํ พุนฺทิกาพทฺธํ กุฬีรปาทกํ อาหจฺจปาทกํ.}

\prose{\textbf{ภิสิ} นาม ปญฺจ ภิสิโย อุณฺณภิสิ โจฬภิภิ วากภิสิ ติณภิสิ ปณฺณภิสิ.}

\prose{\textbf{โกจฺฉํ} นาม วากมยํ วา อุสีรมยํ วา มุญฺชมยํ วา ปพฺพชมยํ\footnote{พพฺพชมยํ (สี)} วา อนฺโต สํเวเฐตฺวา พทฺธํ โหติ.}

\prose{\textbf{สนฺถริตฺวา}ติ สยํ สนฺถริตฺวา.}

\prose{\textbf{สนฺถราเปตฺวา}ติ อญฺญํ สนฺถราเปตฺวา.\csromanspacer{} อนุปสมฺปนฺนํ สนฺถราเปติ, ตสฺส ปลิโพโธ.\csromanspacer{} อุปสมฺปนฺนํ สนฺถราเปติ, สนฺถารกสฺส\footnote{สนฺถตสฺส (อิติปิ)} ปลิโพโธ.}

\prose{\textbf{ตํ ปกฺกมนฺโต เนว อุทฺธเรยฺยา}ติ น สยํ อุทฺธเรยฺย.}

\prose{\textbf{น อุทฺธราเปยฺยา}ติ น อญฺญํ อุทฺธราเปยฺย.}

\prose{\textbf{อนาปุจฺฉํ วา คจฺเฉยฺยา}ติ ภิกฺขุํ วา สามเณรํ วา อารามิกํ วา อนาปุจฺฉา มชฺฌิมสฺส ปุริสสฺส เลฑฺฑุปาตํ อติกฺกมนฺตสฺส อาปตฺติ}

\proseitem{112}{สํฆิเก สํฆิกสญฺญี อชฺโฌกาเส สนฺถริตฺวา วา สนฺถราเปตฺวา วา ตํ ปกฺกมนฺโต เนว อุทฺธเรยฺย น อุทฺธราเปยฺย อนาปุจฺฉํ วา คจฺเฉยฺย, อาปตฺติ ปาจิตฺติยสฺส.\csromanspacer{} สํฆิเก เวมติโก ฯเปฯ สํฆิเก ปุคฺคเลกสญฺญี อชฺโฌกาเส สนฺถริตฺวา วา สนฺถราเปตฺวา วา ตํ ปกฺกมนฺโต เนว อุทฺธเรยฺย น อุทฺธราเปยฺย อนาปุจฺฉํ วา คจฺเฉยฺย, อาปตฺติ}

\prose{จิมิลิกํ วา อุตฺตรตฺถรณํ วา ภูมตฺถรณํ วา ตฏฺฏิกํ วา จมฺมขณฺฑํ วา ปาทปุญฺฉนิํ วา ผลกปีฐํ วา อชฺโฌกาเส สนฺถริตฺวา วา สนฺถราเปตฺวา วา ตํ ปกฺกมนฺโต เนว อุทฺธเรยฺย น อุทฺธราเปยฺย อนาปุจฺฉํ วา คจฺเฉยฺย, อาปตฺติ ทุกฺกฏสฺส.\csromanspacer{} ปุคฺคลิเก สํฆิกสญฺญี, อาปตฺติ ทุกฺกฏสฺส.\csromanspacer{} ปุคฺคลิเก เวมติโก, อาปตฺติ ทุกฺกฏสฺส.\csromanspacer{} ปุคฺคลิเก ปุคฺคลิกสญฺญี อญฺญสฺส ปุคฺคลิเก อาปตฺติ ทุกฺกฏสฺส.\csromanspacer{} อตฺตโน ปุคฺคลิเก อนาปตฺติ.}

\proseitemwithcloserruled{113}{\csromanfolio{60}อนาปตฺติ อุทฺธริตฺวา คจฺฉติ, อุทฺธราเปตฺวา คจฺฉติ, อาปุจฺฉํ คจฺฉติ, โอตาเปนฺโต คจฺฉติ, เกนจิ ปลิพุทฺธํ โหติ, อาปทาสุ อุมฺมตฺตกสฺส อาทิกมฺมิกสฺสาติ.}{ปฐม\-เสนาสน\-สิกฺขาปทํ นิฏฺฐิตํ จตุตฺถํ.}
\csromanmatikaanchor{mtk.18}
\csromantocmarkref{subsection}{5. ทุติยเสนาสนสิกฺขาปท}{mtk.18}
\bookmark[dest={mtk.18},level=2]{5. ทุติยเสนาสนสิกฺขาปท}
\csromanheader{2}{2. ภูตคามวคฺค 5. ทุติย\-เสนาสน\-สิกฺขาปท}
\proseitem{114}{เตน สมเยน พุทฺโธ ภควา สาวตฺถิยํ วิหรติ เชตวเน อนาถปิณฺฑิกสฺส อาราเม.\csromanspacer{} เตน โข ปน สมเยน สตฺตรสวคฺคิยา ภิกฺขู สหายกา โหนฺติ.\csromanspacer{} เต วสนฺตาปิ เอกโตว วสนฺติ, ปกฺกมนฺตาปิ เอกโตว ปกฺกมนฺติ.\csromanspacer{} เต อญฺญตรสฺมิํ สํฆิเก วิหาเร เสยฺยํ สนฺถริตฺวา ตํ ปกฺกมนฺตา เนว อุทฺธริํสุ น อุทฺธราเปสุํ อนาปุจฺฉา ปกฺกมิํสุ, เสนาสนํ อุปจิกาหิ ขายิตํ โหติ.\csromanspacer{} เย เต ภิกฺขู อปฺปิจฺฉา ฯเปฯ เต อุชฺฌายนฺติ ขิยฺยนฺติ วิปาเจนฺติ “กถํ หิ นาม สตฺตรสวคฺคิยา ภิกฺขู สํฆิเก วิหาเร เสยฺยํ สนฺถริตฺวา ตํ ปกฺกมนฺตา เนว อุทฺธริสฺสนฺติ น อุทฺธราเปสฺสนฺติ อนาปุจฺฉา ปกฺกมิสฺสนฺติ, เสนาสนํ อุปจิกาหิ ขายิตนฺ”ติ.\csromanspacer{} อถ โข เต ภิกฺขู สตฺตรสวคฺคิเย ภิกฺขู อเนกปริยาเยน วิครหิตฺวา ภควโต เอตมตฺถํ อาโรเจสุํ ฯเปฯ.\csromanspacer{} สจฺจํ กิร ภิกฺขเว สตฺตรสวคฺคิยา ภิกฺขู สํฆิเก วิหาเร เสยฺยํ สนฺถริตฺวา ตํ ปกฺกมนฺตา เนว อุทฺธริํสุ น อุทฺธราเปสุํ อนาปุจฺฉา ปกฺกมิํสุ, เสนาสนํ อุปจิกาหิ ขายิตนฺติ.\csromanspacer{} สจฺจํ ภควาติ.\csromanspacer{} วิครหิ พุทฺโธ ภควา ฯเปฯ.\csromanspacer{} กถํ หิ นาม เต ภิกฺขเว โมฆปุริสา สํฆิเก วิหาเร เสยฺยํ สนฺถริตฺวา ตํ ปกฺกมนฺตา เนว อุทฺธริสฺสนฺติ น อุทฺธราเปสฺสนฺติ อนาปุจฺฉา ปกฺกมิสฺสนฺติ.\csromanspacer{} เสนาสนํ อุปจิกาหิ ขายิตํ.\csromanspacer{} เนตํ ภิกฺขเว อปฺปสนฺนานํ วา ปสาทาย ฯเปฯ.\csromanspacer{} เอวญฺจ ปน ภิกฺขเว อิมํ สิกฺขาปทํ อุทฺทิเสยฺยาถ\csromandash{}}

\proseitem{115}{\textbf{“โย ปน ภิกฺขุ สํฆิเก วิหาเร เสยฺยํ สนฺถริตฺวา วา สนฺถราเปตฺวา วา ตํ ปกฺกมนฺโต เนว อุทฺธเรยฺย น อุทฺธราเปยฺย, อนาปุจฺฉํ วา คจฺเฉยฺย, ปาจิตฺติยนฺ”}ติ.}

\proseitem{116}{\csromanfolio{61}\textbf{โย ปนา}ติ โย ยาทิโส ฯเปฯ.\csromanspacer{} \textbf{ภิกฺกู}ติ ฯเปฯ อยํ อิมสฺมิํ อตฺเถ อธิปฺเปโต ภิกฺขูติ.}

\prose{\textbf{สํฆิโก} นาม วิหาโร สํฆสฺส ทินฺโน โหติ ปริจฺจตฺโต.}

\prose{\textbf{เสยฺยํ} นาม ภิสิ จิมิลิกา อุตฺตรตฺถรณํ ภูมตฺถรณํ ตฏฺฏิกา จมฺมขณฺโฑ นิสีทนํ ปจฺจตฺถรณํ ติณสนฺถาโร ปณฺณสนฺถาโร.}

\prose{\textbf{สนฺถริตฺวา}ติ สยํ สนฺถริตฺวา.}

\prose{\textbf{สนฺถราเปตฺวา}ติ อญฺญํ สนฺถราเปตฺวา.}

\prose{\textbf{ตํ ปกฺกมนฺโต เนว อุทฺธเรยฺยา}ติ น สยํ อุทฺธเรยฺย.}

\prose{\textbf{น อุทฺธราเปยฺยา}ติ น อญฺญํ อุทฺธราเปยฺย.}

\prose{\textbf{อนาปุจฺฉํ วา คจฺเฉยฺยา}ติ ภิกฺขุํ วา สามเณรํ วา อารามิกํ วา อนาปุจฺฉา ปริกฺขิตฺตสฺส อารามสฺส ปริกฺเขปํ อติกฺกมนฺตสฺส อาปตฺติ ปาจิตฺติยสฺส.\csromanspacer{} อปริกฺขิตฺตสฺส อารามสฺส อุปจารํ อติกฺกมนฺตสฺส อาปตฺติ ปาจิตฺติยสฺส.\csromanspacer{} สํฆิเก สํฆิกสญฺญี เสยฺยํ สนฺถริตฺวา วา สนฺถราเปตฺวา วา ตํ ปกฺกมนฺโต เนว อุทฺธเรยฺย น อุทฺธราเปยฺย อนาปุจฺฉํ วา คจฺเฉยฺย, อาปตฺติ ปาจิตฺติยสฺส.\csromanspacer{} สํฆิเก เวมติโก เสยฺยํ สนฺถริตฺวา วา สนฺถราเปตฺวา วา ตํ ปกฺกมนฺโต เนว อุทฺธเรยฺย น อุทฺธราเปยฺย อนาปุจฺฉํ วา คจฺเฉยฺย, อาปตฺติ ปาจิตฺติยสฺส.\csromanspacer{} สํฆิเก ปุคฺคลิกสญฺญี เสยฺยํ สนฺถริตฺวา วา สนฺถราเปตฺวา วา ตํ ปกฺกมนฺโต เนว อุทฺธเรยฺย น อุทฺธราเปยฺย อนาปุจฺฉํ วา คจฺเฉยฺย, อาปตฺติ ปาจิตฺติยสฺส.}

\proseitem{117}{วิหารสฺส อุปจาเร วา อุปฏฺฐานสาลายํ วา มณฺฑเป วา รุกฺขมูเล วา เสยฺยํ สนฺถริตฺวา วา สนฺถราเปตฺวา วา ตํ ปกฺกมนฺโต เนว อุทฺธเรยฺย น อุทฺธราเปยฺย อนาปุจฺฉํ วา คจฺเฉยฺย, อาปตฺติ ทุกฺกฏสฺส.\csromanspacer{} มญฺจํ วา ปีฐํ วา วิหาเร วา วิหารสฺสูปจาเร วา อุปฏฺฐานสาลายํ วา มณฺฑเป วา รุกฺขมูเล วา สนฺถริตฺวา วา สนฺถราเปตฺวา วา ตํ ปกฺกมนฺโต เนว อุทฺธเรยฺย น อุทฺธราเปยฺย อนาปุจฺฉํ วา คจฺเฉยฺย, อาปตฺติ ทุกฺกฏสฺส.}

\prose{ปุคฺคลิเก สํฆิกสญฺญี, อาปตฺติ ทุกฺกฏสฺส.\csromanspacer{} ปุคฺคลิเก เวมติโก, อาปตฺติ ทุกฺกฏสฺส.\csromanspacer{} ปุคฺคลิเก ปุคฺคลิกสญฺญี อญฺญสฺส ปุคฺคลิเก อาปตฺติ ทุกฺกฏสฺส.\csromanspacer{} อตฺตโน ปุคฺคลิเก อนาปตฺติ.}

\proseitemwithcloserruled{118}{\csromanfolio{62}อนาปตฺติ อุทฺธริตฺวา คจฺฉติ, อุทฺธราเปตฺวา คจฺฉติ, อาปุจฺฉํ คจฺฉติ, เกนจิ ปลิพุทฺธํ โหติ, สาเปกฺโข คนฺตฺวา ตตฺถ ฐิโต อาปุจฺฉติ, เกนจิ ปลิพุทฺโธ โหติ, อาปทาสุ อุมฺมตฺตกสฺส อาทิกมฺมิกสฺสาติ.}{ทุติย\-เสนาสน\-สิกฺขาปทํ นิฏฺฐิตํ ปญฺจมํ.}
\csromanmatikaanchor{mtk.19}
\csromantocmarkref{subsection}{6. อนุปขชฺชสิกฺขาปท}{mtk.19}
\bookmark[dest={mtk.19},level=2]{6. อนุปขชฺชสิกฺขาปท}
\csromanheader{2}{2. \textbf{ภูตคามวคฺค 6. อนุปขชฺชสิกฺขาปท}}
\proseitem{119}{เตน สมเยน พุทฺโธ ภควา สาวตฺถิยํ วิหรติ เชตวเน อนาถปิณฺฑิกสฺส อาราเม.\csromanspacer{} เตน โข ปน สมเยน ฉพฺพคฺคิยา ภิกฺขู วรเสยฺยาโย ปลิพุนฺเธนฺติ, เถรา ภิกฺขู วุฏฺฐาเปนฺติ\footnote{เอตฺถ เต อิติ กมฺมปทํ อูนํ วิย ทิสฺสติ.}.\csromanspacer{} อถ โข ฉพฺพคฺคิยานํ ภิกฺขูนํ เอตทโหสิ “เกน นุ โข มยํ อุปาเยน อิเธว วสฺสํ วเสยฺยามา”ติ.\csromanspacer{} อถ โข ฉพฺพคฺคิยา ภิกฺขู เถเร ภิกฺขู อนุปขชฺช เสยฺยํ กปฺเปนฺติ “ยสฺส สมฺพาโธ ภวิสฺสติ, โส ปกฺกมิสฺสตี”ติ.\csromanspacer{} เย เต ภิกฺขู อปฺปิจฺฉา ฯเปฯ เต อุชฺฌายนฺติ ขิยฺยนฺติ วิปาเจนฺติ “กถํ หิ นาม ฉพฺพคฺคิยา ภิกฺขู เถเร ภิกฺขู อนุปขชฺช เสยฺยํ กปฺเปสฺสนฺตี”ติ.\csromanspacer{} อถ โข เต ภิกฺขู ฉพฺพคฺคิเย ภิกฺขู อเนกปริยาเยน วิครหิตฺวา ภควโต เอตมตฺถํ อาโรเจสุํ ฯเปฯ.\csromanspacer{} สจฺจํ กิร ตุเมฺห ภิกฺขเว เถเร ภิกฺขู อนุปขชฺช เสยฺยํ กปฺเปถาติ.\csromanspacer{} สจฺจํ ภควาติ.\csromanspacer{} วิครหิ พุทฺโธ ภควา ฯเปฯ.\csromanspacer{} กถํ หิ นาม ตุเมฺห โมฆปุริสา เถเร ภิกฺขู อนุปขชฺช เสยฺยํ กปฺเปสฺสถ.\csromanspacer{} เนตํ โมฆปุริสา อปฺปสนฺนานํ วา ปสาทาย ฯเปฯ.\csromanspacer{} เอวญฺจ ปน ภิกฺขเว อิมํ สิกฺขาปทํ อุทฺทิเสยฺยาถ\csromandash{}}

\proseitem{120}{\textbf{“โย ปน ภิกฺขุ สํฆิเก วิหาเร ชานํ ปุพฺพุปคตํ ภิกฺขุํ อนุปขชฺช เสยฺยํ กปฺเปยฺย ‘ยสฺส สมฺพาโธ ภวิสฺสติ, โส ปกฺกมิสฺสตี’ติ.}\csromanspacer{} เอตเทว ปจฺจยํ กริตฺวา อนญฺญํ, ปาจิตฺติยนฺ”ติ.}

\proseitem{121}{\textbf{โย ปนา}ติ โย ยาทิโส ฯเปฯ.\csromanspacer{} \textbf{ภิกฺขู}ติ ฯเปฯ อยํ อิมสฺมิํ อตฺเถ อธิปฺเปโต ภิกฺขูติ.}

\prose{\textbf{สํฆิโก} นาม วิหาโร สํฆสฺส ทินฺโน โหติ ปริจฺจตฺโต.}

\prose{\csromanfolio{63}\textbf{ชานาติ} นาม วุฑฺโฒติ ชานาติ, คิลาโนติ ชานาติ, สํเฆน ทินฺโนติ ชานาติ.}

\prose{\textbf{อนุปขชฺชา}ติ อนุปวิสิตฺวา.}

\prose{\textbf{เสยฺยํ กปฺเปยฺยา}ติ มญฺจสฺส วา ปีฐสฺส วา ปวิสนฺตสฺส วา นิกฺขมนฺตสฺส วา อุปจาเร เสยฺยํ สนฺถรติ วา สนฺถราเปติ วา, อาปตฺติ ทุกฺกฏสฺส.\csromanspacer{} อภินิสีทติ วา อภินิปชฺชติ วา, อาปตฺติ ปาจิตฺติยสฺส.}

\prose{\textbf{เอตเทว ปจฺจยํ กริตฺวา อนญฺญนฺ}ติ น อญฺโญ โกจิ ปจฺจโย โหติ อนุปขชฺช เสยฺยํ กปฺเปตุํ.}

\proseitem{122}{สํฆิเก สํฆิกสญฺญี อนุปขชฺช เสยฺยํ กปฺเปติ, อาปตฺติ ปาจิตฺติยสฺส.\csromanspacer{} สํฆิเก เวมติโก อนุปขชฺช เสยฺยํ กปฺเปติ, อาปตฺติ ปาจิตฺติยสฺส.\csromanspacer{} สํฆิเก ปุคฺคลิกสญฺญี อนุปขชฺช เสยฺยํ กปฺเปติ, อาปตฺติ}

\prose{มญฺจสฺส วา ปีฐสฺส วา ปวิสนฺตสฺส วา นิกฺขมนฺตสฺส วา อุปจารํ ฐเปตฺวา เสยฺยํ สนฺถรติ วา สนฺถราเปติ วา, อาปตฺติ ทุกฺกฏสฺส.\csromanspacer{} อภินิสีทติ วา อภินิปชฺชติ วา, อาปตฺติ ทุกฺกฏสฺส.\csromanspacer{} วิหารสฺส อุปจาเร วา อุปฏฺฐานสาลายํ วา มณฺฑเป วา รุกฺขมูเล วา อชฺโฌกาเส วา เสยฺยํ สนฺถรติ วา สนฺถราเปติ วา, อาปตฺติ ทุกฺกฏสฺส.\csromanspacer{} อภินิสีทติ วา อภินิปฺปชฺชติ วา, อาปตฺติ ทุกฺกฏสฺส.\csromanspacer{} ปุคฺคลิเก สํฆิกสญฺญี, อาปตฺติ ทุกฺกฏสฺส.\csromanspacer{} ปุคฺคลิเก เวมติโก, อาปตฺติ ทุกฺกฏสฺส.\csromanspacer{} ปุคฺคลิเก ปุคฺคลิกสญฺญี อญฺญสฺส ปุคฺคลิเก อาปตฺติ ทุกฺกฏสฺส.\csromanspacer{} อตฺตโน ปุคฺคลิเก อนาปตฺติ.}

\proseitemwithcloserruled{123}{อนาปตฺติ คิลาโน ปวิสติ, สีเตน วา อุเณฺหน วา ปีฬิโต ปวิสติ, อาปทาสุ อุมฺมตฺตกสฺส อาทิกมฺมิกสฺสาติ.}{อนุปขชฺชสิกฺขาปทํ นิฏฺฐิตํ ฉฏฺฐํ.}
\csromanmatikaanchor{mtk.20}
\csromantocmarkref{subsection}{7. นิกฺกฑฺฒนสิกฺขาปท}{mtk.20}
\bookmark[dest={mtk.20},level=2]{7. นิกฺกฑฺฒนสิกฺขาปท}
\csromanheader{2}{2. \textbf{ภูตคามวคฺค 7. นิกฺกฑฺฒนสิกฺขาปท}}
\proseitem{124}{เตน สมเยน พุทฺโธ ภควา สาวตฺถิยํ วิหรติ เชตวเน อนาถปิณฺฑิกสฺส อาราเม.\csromanspacer{} เตน โข ปน สมเยน สตฺตรสวคฺคิยา ภิกฺขู อญฺญตรํ ปจฺจนฺติมํ มหาวิหารํ ปฏิสงฺขโรนฺติ “อิธ มยํ วสฺสํ \csromanfolio{64}วสิสฺสามา”ติ.\csromanspacer{} อทฺทสํสุ โข ฉพฺพคฺคิยา ภิกฺขู สตฺตรสวคฺคิเย ภิกฺขู วิหารํ ปฏิสงฺขโรนฺเต, ทิสฺวาน เอวมาหํสุ “อิเม อาวุโส สตฺตรสวคฺคิยา ภิกฺขู วิหารํ ปฏิสงฺขโรนฺติ.\csromanspacer{} หนฺท เน วุฏฺฐาเปสฺสามา”ติ.\csromanspacer{} เอกจฺเจ เอวมาหํสุ “อาคเมถาวุโส ยาว ปฏิสงฺขโรนฺติ, ปฏิสงฺขเต วุฏฺฐาเปสฺสามา”ติ.}

\prose{อถ โข ฉพฺพคฺคิยา ภิกฺขู สตฺตรสวคฺคิเย ภิกฺขู เอตทโวจุํ “อุฏฺเฐถาวุโส อมฺหากํ วิหาโร ปาปุณาตี”ติ.\csromanspacer{} นนุ อาวุโส ปฏิกจฺเจว\footnote{ปฏิคจฺเจว (สี)} อาจิกฺขิตพฺพํ, มยญฺจญฺญํ ปฏิสงฺขเรยฺยามาติ.\csromanspacer{} นนุ อาวุโส สํฆิโก วิหาโรติ.\csromanspacer{} อามาวุโส สํฆิโก วิหาโรติ.\csromanspacer{} อุฏฺเฐถาวุโส อมฺหากํ วิหาโร ปาปุณาตีติ.\csromanspacer{} มหลฺลโก อาวุโส วิหาโร, ตุเมฺหปิ วสถ, มยมฺปิ วสิสฺสามาติ. “อุฏฺเฐถาวุโส อมฺหากํ วิหาโร ปาปุณาตี”ติ กุปิตา อนตฺตมนา คีวายํ คเหตฺวา นิกฺกฑฺฒนฺติ.\csromanspacer{} เต นิกฺกฑฺฒียมานา โรทนฺติ.\csromanspacer{} ภิกฺขู เอวมาหํสุ “กิสฺส ตุเมฺห อาวุโส โรทถา”ติ.\csromanspacer{} อิเม อาวุโส ฉพฺพคฺคิยา ภิกฺขู กุปิตา อนตฺตมนา อเมฺห สํฆิกา วิหารา นิกฺกฑฺฒนฺตีติ.\csromanspacer{} เย เต ภิกฺขู อปฺปิจฺฉา ฯเปฯ เต อุชฺฌายนฺติ ขิยฺยนฺติ วิปาเจนฺติ “กถํ หิ นาม ฉพฺพคฺคิยา ภิกฺขู กุปิตา อนตฺตมนา ภิกฺขู สํฆิกา วิหารา นิกฺกฑฺฒิสฺสนฺตี”ติ.\csromanspacer{} อถ โข เต ภิกฺขู ฉพฺพคฺคิเย ภิกฺขู อเนกปริยาเยน วิครหิตฺวา ภควโต เอตมตฺถํ อาโรเจสุํ ฯเปฯ.\csromanspacer{} สจฺจํ กิร ตุเมฺห ภิกฺขเว กุปิตา อนตฺตมนา ภิกฺขู สํฆิกา วิหารา นิกฺกฑฺฒถาติ.\csromanspacer{} สจฺจํ ภควาติ.\csromanspacer{} วิครหิ พุทฺโธ ภควา ฯเปฯ.\csromanspacer{} กถํ หิ นาม ตุเมฺห โมฆปุริสา กุปิตา อนตฺตมนา ภิกฺขู สํฆิกา วิหารา นิกฺกฑฺฒิสฺสถ.\csromanspacer{} เนตํ โมฆปุริสา อปฺปสนฺนานํ วา ปสาทาย ฯเปฯ.\csromanspacer{} เอวญฺจ ปน ภิกฺขเว อิมํ สิกฺขาปทํ อุทฺทิเสยฺยาถ\csromandash{}}

\proseitem{125}{\textbf{“โย ปน ภิกฺขุ ภิกฺขุํ กุปิโต อนตฺตมโน สํฆิกา วิหารา นิกฺกฑฺเฒยฺย วา นิกฺกฑฺฒาเปยฺย วา, ปาจิตฺติยนฺ”}ติ.}

\proseitem{126}{\textbf{โย ปนา}ติ โย ยาทิโส ฯเปฯ.\csromanspacer{} \textbf{ภิกฺขู}ติ ฯเปฯ อยํ อิมสฺมิํ อตฺเถ อธิปฺเปโต ภิกฺขูติ.}

\prose{\textbf{ภิกฺขุนฺ}ติ อญฺญํ ภิกฺขุํ.}

\prose{\csromanfolio{65}\textbf{กุปิโต อนตฺตมโน}ติ อนภิรทฺโธ อาหตจิตฺโต ขิลชาโต.}

\prose{\textbf{สํฆิโก} นาม วิหาโร สํฆสฺส ทินฺโน โหติ ปริจฺจตฺโต.}

\prose{\textbf{นิกฺกฑฺเฒยฺยา}ติ คพฺเภ คเหตฺวา ปมุขํ นิกฺกฑฺฒติ, อาปตฺติ ปาจิตฺติยสฺส.\csromanspacer{} ปมุเข คเหตฺวา พหิ นิกฺกฑฺฒติ, อาปตฺติ ปาจิตฺติยสฺส.\csromanspacer{} เอเกน ปโยเคน พหุเกปิ ทฺวาเร อติกฺกาเมติ, อาปตฺติ ปาจิตฺติยสฺส.}

\prose{\textbf{นิกฺกฑฺฒาเปยฺยา}ติ อญฺญํ อาณาเปติ, อาปตฺติ ปาจิตฺติยสฺส\footnote{ทุกฺกฏสฺส (กตฺถจิ)}.\csromanspacer{} สกิํ อาณตฺโต พหุเกปิ ทฺวาเร อติกฺกาเมติ, อาปตฺติ ปาจิตฺติยสฺส.}

\proseitem{127}{สํฆิเก สํฆิกสญฺญี กุปิโต อนตฺตมโน นิกฺกฑฺฒติ วา นิกฺกฑฺฒาเปติ วา, อาปตฺติ ปาจิตฺติยสฺส.\csromanspacer{} สํฆิเก เวมติโก กุปิโต อนตฺตมโน นิกฺกฑฺฒติ วา นิกฺกฑฺฒาเปติ วา, อาปตฺติ ปาจิตฺติยสฺส.\csromanspacer{} สํฆิเก ปุคฺคลิกสญฺญี กุปิโต อนตฺตมโน นิกฺกฑฺฒติ วา นิกฺกฑฺฒาเปติ วา, อาปตฺติ}

\prose{ตสฺส ปริกฺขารํ นิกฺกฑฺฒติ วา นิกฺกฑฺฒาเปติ วา, อาปตฺติ ทุกฺกฏสฺส.\csromanspacer{} วิหารสฺส อุปจารา วา อุปฏฺฐานสาลาย วา มณฺฑปา วา รุกฺขมูลา วา อชฺโฌกาสา วา นิกฺกฑฺฒติ วา นิกฺกฑฺฒาเปติ วา, อาปตฺติ ทุกฺกฏสฺส.\csromanspacer{} ตสฺส ปริกฺขารํ นิกฺกฑฺฒติ วา นิกฺกฑฺฑาเปติ วา, อาปตฺติ ทุกฺกฏสฺส.\csromanspacer{} อนุปสมฺปนฺนํ วิหารา วา วิหารสฺส อุปจารา วา อุปฏฺฐานสาลาย วา มณฺฑปา วา รุกฺขมูลา วา อชฺโฌกาสา วา นิกฺกฑฺฒติ วา นิกฺกฑฺฒาเปติ วา, อาปตฺติ ทุกฺกฏสฺส.\csromanspacer{} ตสฺส ปริกฺขารํ นิกฺกฑฺฒติ วา นิกฺกฑฺฒาเปติ วา, อาปตฺติ ทุกฺกฏสฺส.}

\prose{ปุคฺคลิเก สํฆิกสญฺญี, อาปตฺติ ทุกฺกฏสฺส.\csromanspacer{} ปุคฺคลิเก เวมติโก, อาปตฺติ ทุกฺกฏสฺส.\csromanspacer{} ปุคฺคลิเก ปุคฺคลิกสญฺญี อญฺญสฺส ปุคฺคลิเก อาปตฺติ ทุกฺกฏสฺส.\csromanspacer{} อตฺตโน ปุคฺคลิเก อนาปตฺติ.}

\proseitemwithcloserruled{128}{อนาปตฺติ อลชฺชิํ นิกฺกฑฺฒติ วา นิกฺกฑฺฒาเปติ วา, ตสฺส ปริกฺขารํ นิกฺกฑฺฒติ วา นิกฺกฑฺฒาเปติ วา, อุมฺมตฺตกํ นิกฺกฑฺฒติ วา นิกฺกฑฺฒาเปติ วา, ตสฺส ปริกฺขารํ นิกฺกฑฺฒติ วา นิกฺกฑฺฒาเปติ วา, ภณฺฑนการกํ กลหการกํ วิวาทการกํ ภสฺสการกํ สํเฆ อธิกรณการกํ นิกฺกฑฺฒติ วา \csromanfolio{66}นิกฺกฑฺฒาเปติ วา, ตสฺส ปริกฺขารํ นิกฺกฑฺฒติ วา นิกฺกฑฺฒาเปติ วา, อนฺเตวาสิกํ วา สทฺธิวิหาริกํ วา น สมฺมา วตฺตนฺตํ นิกฺกฑฺฒติ วา นิกฺกฑฺฒาเปติ วา, ตสฺส ปริกฺขารํ นิกฺกฑฺฒติ วา นิกฺกฑฺฒาเปติ วา, อุมฺมตฺตกสฺส อาทิกมฺมิกสฺสาติ.}{นิกฺกฑฺฒนสิกฺขาปทํ นิฏฺฐิตํ สตฺตมํ.}
\csromanmatikaanchor{mtk.21}
\csromantocmarkref{subsection}{8. เวหาสกุฏิสิกฺขาปท}{mtk.21}
\bookmark[dest={mtk.21},level=2]{8. เวหาสกุฏิสิกฺขาปท}
\csromanheader{2}{2. \textbf{ภูตคามวคฺค 8. เวหาสกุฏิสิกฺขาปท}}
\proseitem{129}{เตน สมเยน พุทฺโธ ภควา สาวตฺถิยํ วิหรติ เชตวเน อนาถปิณฺฑิกสฺส อาราเม.\csromanspacer{} เตน โข ปน สมเยน เทฺว ภิกฺขู สํฆิเก วิหาเร อุปริเวหาสกุฏิยา เอโก เหฏฺฐา วิหรติ\footnote{วิหรนฺติ, เอโก เหฏฺฐา (?)} เอโก อุปริ, อุปริโม ภิกฺขุ อาหจฺจปาทกํ มญฺจํ สหสา อภินิสีทิ, มญฺจปาโท นิปฺปติตฺวา\footnote{ปติตฺวา (สฺยา)} เหฏฺฐิมสฺส ภิกฺขุโน มตฺถเก อวตฺถาสิ, โส ภิกฺขุ วิสฺสรมกาสิ.\csromanspacer{} ภิกฺขู อุปธาวิตฺวา ตํ ภิกฺขุํ เอตทโวจุํ “กิสฺส ตฺวํ อาวุโส วิสฺสรมกาสี”ติ.\csromanspacer{} อถ โข โส ภิกฺขุ ภิกฺขูนํ เอตมตฺถํ อาโรเจสิ.\csromanspacer{} เย เต ภิกฺขู อปฺปิจฺฉา ฯเปฯ เต อุชฺฌายนฺติ ขิยฺยนฺติ วิปาเจนฺติ “กถํ หิ นาม ภิกฺขุ สํฆิเก วิหาเร อุปริเวหาสกุฏิยา อาหจฺจปาทกํ มญฺจํ สหสา อภินิสีทิสฺสตี”ติ.\csromanspacer{} อถ โข เต ภิกฺขู ตํ ภิกฺขุํ อเนกปริยาเยน วิครหิตฺวา ภควโต เอตมตฺถํ อาโรเจสุํ ฯเปฯ.\csromanspacer{} สจฺจํ กิร ตฺวํ ภิกฺขุ สํฆิเก วิหาเร อุปริเวหาสกุฏิยา อาหจฺจปาทกํ มญฺจํ สหสา อภินิสีทสีติ.\csromanspacer{} สจฺจํ ภควาติ.\csromanspacer{} วิครหิ พุทฺโธ ภควา ฯเปฯ.\csromanspacer{} กถํ หิ นาม ตฺวํ โมฆปุริส สํฆิเก วิหาเร อุปริเวหาสกุฏิยา อาหจฺจปาทกํ มญฺจํ สหสา อภินิสีทิสฺสสิ.\csromanspacer{} เนตํ โมฆปุริส อปฺปสนฺนานํ วา ปสาทาย ฯเปฯ.\csromanspacer{} เอวญฺจ ปน ภิกฺขเว อิมํ สิกฺขาปทํ อุทฺทิเสยฺยาถ\csromandash{}}

\proseitem{130}{\textbf{“โย ปน ภิกฺขุ สํฆิเก วิหาเร อุปริเวหาสกุฏิยา อาหจฺจปาทกํ มญฺจํ วา ปีฐํ วา อภินิสีเทยฺย วา อภินิปชฺเชยฺย วา, ปาจิตฺติยนฺ”}ติ.}

\proseitem{131}{\csromanfolio{67}\textbf{โย ปนา}ติ โย ยาทิโส ฯเปฯ.\csromanspacer{} \textbf{ภิกฺขู}ติ ฯเปฯ อยํ อิมสฺมิํ อตฺเถ อธิปฺเปโต ภิกฺขูติ.}

\prose{\textbf{สํฆิโก} นาม วิหาโร สํฆสฺส ทินฺโน โหติ ปริจฺจตฺโต.}

\prose{\textbf{เวหาสกุฏิ} นาม มชฺฌิมสฺส ปุริสสฺส อสีสฆฏฺฏา.}

\prose{\textbf{อาหจฺจปาทโก} นาม \textbf{มญฺโจ} องฺเค วิชฺฌิตฺวา ฐิโต โหติ.\csromanspacer{} \textbf{อาหจฺจปาทกํ} นาม \textbf{ปีฐํ} องฺเค วิชฺฌิตฺวา ฐิตํ โหติ.}

\prose{\textbf{อภินิสีเทยฺยา}ติ ตสฺมิํ อภินิสีทติ, อาปตฺติ ปาจิตฺติยสฺส.}

\prose{\textbf{อภินิปชฺเชยฺยา}ติ ตสฺมิํ อภินิปชฺชติ, อาปตฺติ ปาจิตฺติยสฺส.}

\proseitem{132}{สํฆิเก สํฆิกสญฺญี อุปริเวหาสกุฏิยา อาหจฺจปาทกํ มญฺจํ วา ปีฐํ วา อภินิสีทติ วา อภินิปชฺชติ วา, อาปตฺติ ปาจิตฺติยสฺส.\csromanspacer{} สํฆิเก เวมติโก ฯเปฯ.\csromanspacer{} สํฆิเก ปุคฺคลิกสญฺญี อุปริเวหาสกุฏิยา อาหจฺจปาทกํ มญฺจํ วา ปีฐํ วา อภินิสีทติ วา อภินิปชฺชติ วา, อาปตฺติ}

\prose{ปุคฺคลิเก สํฆิกสญฺญี, อาปตฺติ ทุกฺกฏสฺส.\csromanspacer{} ปุคฺคลิเก เวมติโก, อาปตฺติ ทุกฺกฏสฺส.\csromanspacer{} ปุคฺคลิเก ปุคฺคลิกสญฺญี อญฺญสฺส ปุคฺคลิเก อาปตฺติ ทุกฺกฏสฺส.\csromanspacer{} อตฺตโน ปุคฺคลิเก อนาปตฺติ.}

\proseitemwithcloserruled{133}{อนาปตฺติ อเวหาสกุฏิยา, สีสฆฏฺฏาย, เหฏฺฐา อปริโภคํ โหติ, ปทรสญฺจิตํ โหติ, ปฏาณิ ทินฺนา โหติ, ตสฺมิํ ฐิโต คณฺหติ วา ลคฺเคติ วา, อุมฺมตฺตกสฺส อาทิกมฺมิกสฺสาติ.}{เวหาสกุฏิสิกฺขาปทํ นิฏฺฐิตํ อฏฺฐมํ.}
\csromanmatikaanchor{mtk.22}
\csromantocmarkref{subsection}{9. มหลฺลกวิหารสิกฺขาปท}{mtk.22}
\bookmark[dest={mtk.22},level=2]{9. มหลฺลกวิหารสิกฺขาปท}
\csromanheader{2}{2. ภูตคามวคฺค 9. มหลฺลก\-วิหาร\-สิกฺขาปท}
\proseitem{134}{เตน สมเยน พุทฺโธ ภควา โกสมฺพิยํ วิหรติ โฆสิตาราเม.\csromanspacer{} เตน โข ปน สมเยน อายสฺมโต ฉนฺนสฺส อุปฏฺฐาโก มหามตฺโต อายสฺมโต ฉนฺนสฺส วิหารํ การาเปติ.\csromanspacer{} อถ โข อายสฺมา ฉนฺโน กตปริโยสิตํ วิหารํ ปุนปฺปุนํ ฉาทาเปติ \csromanfolio{68}\textbf{ปุนปฺปุนํ เลปาเปติ, อติภาริโต วิหาโร ปริปติ.}\csromanspacer{}\textbf{ อถ โข อายสฺมา ฉนฺโน} \textbf{ติณญฺจ กฏฺฐญฺจ สํกฑฺฒนฺโต อญฺญตรสฺส พฺราหฺมณสฺส ยวเขตฺตํ} ทูเสสิ.\csromanspacer{} อถ โข โส พฺราหฺมโณ อุชฺฌายติ ขิยฺยติ วิปาเจติ “กถํ หิ นาม ภทนฺตา อมฺหากํ ยวเขตฺตํ ทูเสสฺสนฺตี”ติ.\csromanspacer{} อสฺโสสุํ โข ภิกฺขู ตสฺส พฺราหฺมณสฺส อุชฺฌายนฺตสฺส ขิยฺยนฺตสฺส วิปาเจนฺตสฺส.\csromanspacer{} เย เต ภิกฺขู อปฺปิจฺฉา ฯเปฯ เต อุชฺฌายนฺติ ขิยฺยนฺติ วิปาเจนฺติ “กถํ หิ นาม อายสฺมา \textbf{ฉนฺโน กตปริโยสิตํ วิหารํ ปุนปฺปุนํ ฉาทาเปสฺสติ ปุนปฺปุนํ} \textbf{เลปาเปสฺสติ.}\csromanspacer{}\textbf{ อติภาริโต วิหาโร ปริปตี”ติ.}\csromanspacer{}\textbf{ อถ โข เต ภิกฺขู อายสฺมนฺตํ} ฉนฺนํ อเนกปริยาเยน วิครหิตฺวา ภควโต เอตมตฺถํ อาโรเจสุํ ฯเปฯ.\csromanspacer{} สจฺจํ กิร ตฺวํ ฉนฺน กตปริโยสิตํ วิหารํ ปุนปฺปุนํ ฉาทาเปสิ ปุนปฺปุนํ เลปาเปสิ.\csromanspacer{} อติภาริโต วิหาโร ปริปตีติ.\csromanspacer{} สจฺจํ ภควาติ.\csromanspacer{} วิครหิ พุทฺโธ ภควา ฯเปฯ.\csromanspacer{} กถํ หิ นาม ตฺวํ โมฆปุริส กตปริโยสิตํ วิหารํ ปุนปฺปุนํ ฉาทาเปสฺสสิ ปุนปฺปุนํ เลปาเปสฺสสิ.\csromanspacer{} อติภาริโต วิหาโร ปริปติ.\csromanspacer{} เนตํ โมฆปุริส อปฺปสนฺนานํ วา ปสาทาย ฯเปฯ.\csromanspacer{} เอวญฺจ ปน ภิกฺขเว อิมํ สิกฺขาปทํ อุทฺทิเสยฺยาถ\csromandash{}}

\proseitem{135}{“มฺอหลฺลกํ ปน ภิกฺขุนา วิหารํ การยมาเนน ยาว ทฺวารโกสา อคฺคฬฏฺฐปนาย อาโลก\-สนฺธิ\-ปริกมฺมาย ทฺวตฺติจฺฉทนสฺส ปริยายํ อปฺปหริเต ฐิเตน อธิฏฺฐาตพฺพํ, ตโต เจ อุตฺตริ อปฺปหริเตปิ ฐิโต อธิฏฺฐเหยฺย, ปาจิตฺติยนฺ”ติ.}

\proseitem{136}{\textbf{มหลฺลโก} นาม วิหาโร สสฺสามิโก วุจฺจติ.}

\prose{\textbf{วิหาโร} นาม อุลฺลิตฺโต วา โหติ อวลิตฺโต วา อุลฺลิตฺตาวลิตฺโต วา.}

\prose{\textbf{การยมาเนนา}ติ กโรนฺโต วา การาเปนฺโต วา.}

\prose{\textbf{ยาว ทฺวารโกสา}ติ ปิฏฺฐสํฆาฏสฺส\footnote{ปีฐสํฆาฏสฺส (อิติปิ), ปิฏฺฐิสงฺฆาฏสฺส (สฺยา)} สมนฺตา หตฺถปาสา.}

\prose{\textbf{อคฺคฬฏฺฐปนายา}ติ ทฺวารฏฺฐปนาย.}

\prose{อาโลก\-สนฺธิ\-ปริกมฺมายาติ วาตปานปริกมฺมาย เสตวณฺณํ กาฬวณฺณํ เครุกปริกมฺมํ มาลากมฺมํ ลตากมฺมํ มกรทนฺตกํ ปญฺจปฏิกํ.}

\prose{\csromanfolio{69}\textbf{ทฺวตฺติจฺฉทนสฺส ปริยายํ อปฺปหริเต ฐิเตน อธิฏฺฐาตพฺพนฺ}ติ หริตํ นาม ปุพฺพณฺณํ อปรณฺณํ.\csromanspacer{} สเจ หริเต ฐิโต อธิฏฺฐาติ, อาปตฺติ ทุกฺกฏสฺส.\csromanspacer{} มคฺเคน ฉาเทนฺตสฺส เทฺว มคฺเค อธิฏฺฐหิตฺวา ตติยํ มคฺคํ อาณาเปตฺวา ปกฺกมิตพฺพํ.\csromanspacer{} ปริยาเยน ฉาเทนฺตสฺส เทฺว ปริยาเย อธิฏฺฐหิตฺวา ตติยํ ปริยายํ อาณาเปตฺวา ปกฺกมิตพฺพํ.}

\proseitem{137}{\textbf{ตโต เจ อุตฺตริ อปฺปหริเตปิ ฐิโต อธิฏฺฐเหยฺยา}ติ อิฏฺฐกาย ฉาเทนฺตสฺส อิฏฺฐกิฏฺฐกาย อาปตฺติ ปาจิตฺติยสฺส.\csromanspacer{} สิลาย ฉาเทนฺตสฺส สิลาย สิลาย อาปตฺติ ปาจิตฺติยสฺส.\csromanspacer{} สุธาย ฉาเทนฺตสฺส ปิณฺเฑ ปิณฺเฑ อาปตฺติ ปาจิตฺติยสฺส.\csromanspacer{} ติเณน ฉาเทนฺตสฺส กรเฬ กรเฬ อาปตฺติ ปาจิตฺติยสฺส.\csromanspacer{} ปณฺเณน ฉาเทนฺตสฺส ปณฺเณ ปณฺเณ อาปตฺติ ปาจิตฺติยสฺส.}

\prose{อติเรกทฺวตฺติปริยาเย อติเรกสญฺญี อธิฏฺฐาติ, อาปตฺติ ปาจิตฺติยสฺส.\csromanspacer{} อติเรกทฺวตฺติปริยาเย เวมติโก อธิฏฺฐาติ, อาปตฺติ ปาจิตฺติยสฺส.\csromanspacer{} อติเรกทฺวตฺติปริยาเย อูนกสญฺญี อธิฏฺฐาติ, อาปตฺติ ปาจิตฺติยสฺส.}

\prose{อูนกทฺวตฺติปริยาเย อติเรกสญฺญี, อาปตฺติ ทุกฺกฏสฺส.\csromanspacer{} อูนกทฺวตฺติปริยาเย เวมติโก, อาปตฺติ ทุกฺกฏสฺส.\csromanspacer{} อูนกทฺวตฺติปริยาเย อูนกสญฺญี, อนาปตฺติ.}

\proseitemwithcloserruled{138}{อนาปตฺติ ทฺวตฺติปริยาเย อูนกทฺวตฺติปริยาเย เลเณ คุหาย ติณกุฏิกาย อญฺญสฺสตฺถาย อตฺตโน ธเนน, วาสาคารํ ฐเปตฺวา สพฺพตฺถ, อนาปตฺติ อุมฺมตฺตกสฺส อาทิกมฺมิกสฺสาติ.}{มหลฺลก\-วิหาร\-สิกฺขาปทํ นิฏฺฐิตํ นวมํ.}
\csromanmatikaanchor{mtk.23}
\csromantocmarkref{subsection}{10. สปฺปาณกสิกฺขาปท}{mtk.23}
\bookmark[dest={mtk.23},level=2]{10. สปฺปาณกสิกฺขาปท}
\csromanheader{2}{2. \textbf{ภูตคามวคฺค} 10. สปฺปาณกสิกฺขาปท}
\proseitem{139}{เตน สมเยน พุทฺโธ ภควา อาฬวิยํ วิหรติ อคฺคาฬเว เจติเย.\csromanspacer{} เตน โข ปน สมเยน อาฬวกา ภิกฺขู นวกมฺมํ กโรนฺตา ชานํ สปฺปาณกํ อุทกํ ติณมฺปิ มตฺติกมฺปิ สิญฺจนฺติปิ สิญฺจาเปนฺติปิ.\csromanspacer{} เย เต ภิกฺขู อปฺปิจฺฉา ฯเปฯ เต อุชฺฌายนฺติ ขิยฺยนฺติ วิปาเจนฺติ “กถํ หิ นาม \csromanfolio{70}อาฬวกา ภิกฺขู ชานํ สปฺปาณกํ อุทกํ ติณมฺปิ มตฺติกมฺปิ สิญฺจิสฺสนฺติปิ สิญฺจาเปสฺสนฺติปี”ติ.\csromanspacer{} อถ โข เต ภิกฺขู อาฬวเก ภิกฺขู อเนกปริยาเยน วิครหิตฺวา ภควโต เอตมตฺถํ อาโรเจสุํ ฯเปฯ “สจฺจํ กิร ตุเมฺห ภิกฺขเว ชานํ สปฺปาณกํ อุทกํ ติณมฺปิ มตฺติกมฺปิ สิญฺจถปิ สิญฺจาเปถปี”ติ.\csromanspacer{} สจฺจํ ภควาติ.\csromanspacer{} วิครหิ พุทฺโธ ภควา \textbf{ฯเปฯ.}\csromanspacer{}\textbf{ กถํ หิ นาม ตุเมฺห โมฆปุริสา ชานํ สปฺปาณกํ อุทกํ} \textbf{ติณมฺปิ มตฺติกมฺปิ สิญฺจิสฺสถปิ สิญฺจาเปสฺสถปิ.}\csromanspacer{}\textbf{ เนตํ โมฆปุริสา} อปฺปสนฺนานํ วา ปสาทาย ฯเปฯ.\csromanspacer{} เอวญฺจ ปน ภิกฺขเว อิมํ สิกฺขาปทํ อุทฺทิเสยฺยาถ\csromandash{}}

\proseitem{140}{\textbf{“โย ปน ภิกฺขุ ชานํ สปฺปาณกํ อุทกํ ติณํ วา มตฺติกํ วา สิญฺเจยฺย วา สิญฺจาเปยฺย วา, ปาจิตฺติยนฺ”}ติ.}

\proseitem{141}{โย ปนาติ โย ยาทิโส ฯเปฯ.\csromanspacer{} \textbf{ภิกฺขู}ติ ฯเปฯ อยํ อิมสฺมิํ อตฺเถ อธิปฺเปโต ภิกฺขูติ.}

\prose{\textbf{ชานา}ติ นาม\footnote{ชานํ นาม (สฺยา)} สามํ วา ชานาติ.\csromanspacer{} อญฺเญ วา ตสฺส อาโรเจนฺติ.}

\prose{\textbf{สิญฺเจยฺยา}ติ สยํ สิญฺจติ, อาปตฺติ ปาจิตฺติยสฺส.}

\prose{\textbf{สิญฺจาเปยฺยา}ติ อญฺญํ อาณาเปติ, อาปตฺติ ปาจิตฺติยสฺส\footnote{อาปตฺติ ทุกฺกฏสฺส (กตฺถจิ)}.\csromanspacer{} สกิํ อาณตฺโต พหุกมฺปิ สิญฺจติ, อาปตฺติ ปาจิตฺติยสฺส.}

\proseitem{142}{สปฺปาณเก สปฺปาณกสญฺญี ติณํ วา มตฺติกํ สิญฺจติ วา สิญฺจาเปติ วา, อาปตฺติ ปาจิตฺติยสฺส.\csromanspacer{} สปฺปาณเก เวมติโก ติณํ วา มตฺติกํ วา สิญฺจติ วา สิญฺจาเปติ วา, อาปตฺติ ทุกฺกฏสฺส.\csromanspacer{} สปฺปาณเก อปฺปาณกสญฺญี ติณํ วา มตฺติกํ วา สิญฺจติ วา สิญฺจาเปติ วา, อนาปตฺติ.\csromanspacer{} อปฺปาณเก สปฺปาณกสญฺญี, อาปตฺติ ทุกฺกฏสฺส.\csromanspacer{} อปฺปาณเก เวมติโก, อาปตฺติ ทุกฺกฏสฺส.\csromanspacer{} อปฺปาณเก อปฺปาณกสญฺญี, อนาปตฺติ.}

\proseitemwithcloser{143}{อนาปตฺติ อสญฺจิจฺจ อสฺสติยา อชานนฺตสฺส อุมฺมตฺตกสฺส อาทิกมฺมิกสฺสาติ.}{สปฺปาณกสิกฺขาปทํ นิฏฺฐิตํ ทสมํ.}
\vspace{\tipitakaparskip}
\csromancenter{ภูตคามวคฺโค ทุติโย.}
\csromancenter{ตสฺสุทฺทานํ}
\vspace{0.5\baselineskip}
\setlength{\csromangathapagewidth}{0pt}
\csromangathameasuregroup{ภูตํ อญฺญาย อุชฺฌายํ,}{\csromangathabat{ภูตํ อญฺญาย อุชฺฌายํ,}{ปกฺกมนฺเตน เต ทุเว.} \\ ปุพฺเพ นิกฺกฑฺฒนาหจฺจ, ทฺวารํ สปฺปาณเกน จาติ.}
\csromangathasetleft{ภูตํ อญฺญาย อุชฺฌายํ,}
\csromangathagroup{\csromangathastanza{\csromanfolio{71}\csromangathabat{ภูตํ อญฺญาย อุชฺฌายํ,}{ปกฺกมนฺเตน เต ทุเว.} \\ ปุพฺเพ นิกฺกฑฺฒนาหจฺจ, ทฺวารํ สปฺปาณเกน จาติ.}}

\csromanmatikaanchor{mtk.24}
\csromanmatikaanchor{mtk.25}
\csromantocmarknopage{section}{3. โอวาทวคฺค}{mtk.24}
\bookmark[dest={mtk.24},level=1]{3. โอวาทวคฺค}
\csromantocmarkref{subsection}{1. โอวาทสิกฺขาปท}{mtk.25}
\bookmark[dest={mtk.25},level=2]{1. โอวาทสิกฺขาปท}
\setcsromanheadh{3. โอวาทวคฺค}
\csromanheader{2}{3. \textbf{โอวาทวคฺค 1. โอวาทสิกฺขาปท}}
\proseitem{144}{เตน สมเยน พุทฺโธ ภควา สาวตฺถิยํ วิหรติ เชตวเน อนาถปิณฺฑิกสฺส อาราเม.\csromanspacer{} เตน โข ปน สมเยน เถรา ภิกฺขู ภิกฺขุนิโย โอวทนฺตา ลาภิโน โหนฺติ จีวร\-ปิณฺฑปาต\-เสนาสน\-คิลานปฺปจฺจย\-เภสชฺช\-ปริกฺขารานํ.\csromanspacer{} อถ โข ฉพฺพคฺคิยานํ ภิกฺขูนํ เอตทโหสิ “เอตรหิ โข อาวุโส เถรา ภิกฺขู ภิกฺขุนิโย โอวทนฺตา ลาภิโน โหนฺติ จีวร\-ปิณฺฑปาต\-เสนาสน\-คิลานปฺปจฺจย\-เภสชฺช\-ปริกฺขารานํ, หนฺทาวุโส มยมฺปิ ภิกฺขุนิโย โอวทามา”ติ.\csromanspacer{} อถ โข ฉพฺพคฺคิยา ภิกฺขู ภิกฺขุนิโย อุปสงฺกมิตฺวา เอตทโวจุํ “อเมฺหปิ ภคินิโย อุปสงฺกมถ, มยมฺปิ โอวทิสฺสามา”ติ.}

\prose{อถ โข ตา ภิกฺขุนิโย เยน ฉพฺพคฺคิยา ภิกฺขู เตนุปสงฺกมิํสุ, อุปสงฺกมิตฺวา ฉพฺพคฺคิเย ภิกฺขู อภิวาเทตฺวา เอกมนฺตํ นิสีทิํสุ.\csromanspacer{} อถ โข ฉพฺพคฺคิย ภิกฺขู ภิกฺขุนีนํ ปริตฺตญฺเญว ธมฺมิํ กถํ กตฺวา ทิวสํ ติรจฺฉานกถาย วีตินาเมตฺวา อุยฺโยเชสุํ “คจฺฉถ ภคินิโย”ติ.\csromanspacer{} อถ โข ตา ภิกฺขุนิโย เยน ภควา เตนุปสงฺกมิํสุ, อุปสงฺกมิตฺวา ภควนฺตํ อภิวาเทตฺวา เอกมนฺตํ อฏฺฐํสุ.\csromanspacer{} เอกมนฺตํ ฐิตา โข ตา ภิกฺขุนิโย ภควา เอตทโวจ “กจฺจิ ภิกฺขุนิโย โอวาโท อิทฺโธ อโหสี”ติ.\csromanspacer{} กุโต ภนฺเต โอวาโท อิทฺโธ ภวิสฺสติ, อยฺยา ฉพฺพคฺคิยา ปริตฺตญฺเญว ธมฺมิํ กถํ กตฺวา ทิวสํ ติรจฺฉานกถาย วีตินาเมตฺวา อุยฺโยเชสุนฺติ.\csromanspacer{} อถ โข ภควา ตา ภิกฺขุนิโย ธมฺมิยา กถาย สนฺทสฺเสสิ สมาทเปสิ สมุตฺเตเชสิ สมฺปหํเสสิ.\csromanspacer{} อถ โข ตา ภิกฺขุนิโย ภควตา ธมฺมิยา กถาย สนฺทสฺสิตา สมาทปิตา สมุตฺเตชิตา สมฺปหํสิตา ภควนฺตํ อภิวาเทตฺวา ปทกฺขิณํ กตฺวา ปกฺกมิํสุ.}

\prose{อถ โข ภควา เอตสฺมิํ นิทาเน เอตสฺมิํ ปกรเณ ภิกฺขุสํฆํ สนฺนิปาตาเปตฺวา ฉพฺพคฺคิเย ภิกฺขู ปฏิปุจฺฉิ “สจฺจํ กิร ตุเมฺห ภิกฺขเว \csromanfolio{72}ภิกฺขูนีนํ ปริตฺตญฺเญว ธมฺมิํ กถํ กตฺวา ทิวสํ ติรจฺฉานกถาย วีตินาเมตฺวา อุยฺโยเชถา”ติ.\csromanspacer{} สจฺจํ ภควาติ.\csromanspacer{} วิครหิ พุทฺโธ ภควา ฯเปฯ.\csromanspacer{} กถํ หิ นาม ตุเมฺห โมฆปุริสา ภิกฺขุนีนํ ปริตฺตญฺเญว ธมฺมิํ กถํ กตฺวา ทิวสํ ติรจฺฉานกถาย วีตินาเมตฺวา อุยฺโยเชสฺสถ.\csromanspacer{} เนตํ โมฆปุริสา อปฺปสนฺนานํ วา ปสาทาย ฯเปฯ วิครหิตฺวา ฯเปฯ ธมฺมิํ กถํ กตฺวา ภิกฺขู อามนฺเตสิ “อนุชานามิ ภิกฺขเว ภิกฺขุโนวาทกํ สมฺมนฺนิตุํ.\csromanspacer{} เอวญฺจ ปน ภิกฺขเว สมฺมนฺนิตพฺโพ, ปฐมํ ภิกฺขุ ยาจิตพฺโพ, ยาจิตฺวา พฺยตฺเตน ภิกฺขุนา ปฏิพเลน สํโฆ ญาเปตพฺโพ\csromandash{}}

\proseitem{145}{“สุณาตุ เม ภนฺเต สํโฆ, ยทิ สํฆสฺส ปตฺตกลฺลํ, สํโฆ อิตฺถนฺนามํ ภิกฺขุํ ภิกฺขุโนวาทกํ สมฺมนฺเนยฺย, เอสา ญตฺติ.}

\prose{สุณาตุ เม ภนฺเต สํโฆ, สํโฆ อิตฺถนฺนามํ ภิกฺขุํ ภิกฺขุโนวาทกํ สมฺมนฺนติ, ยสฺสายสฺมโต ขมติ อิตฺถนฺนามสฺส ภิกฺขุโน ภิกฺขุโนวาทกสฺส สมฺมุติ, โส ตุณฺหสฺส, ยสฺส นกฺขมติ, โส ภาเสยฺย.}

\prose{ทุติยมฺปิ เอตมตฺถํ วทามิ ฯเปฯ.\csromanspacer{} ตติยมฺปิ เอตมตฺถํ วทามิ.\csromanspacer{} สุณาตุ เม ภนฺเต สํโฆ, สํโฆ อิตฺถนฺนามํ ภิกฺขุํ ภิกฺขุโนวาทกํ สมฺมนฺนติ, ยสฺสายสฺมโต ขมติ อิตฺถนฺนามสฺส ภิกฺขุโน ภิกฺขุโนวาทกสฺส สมฺมุติ, โส ตุณฺหสฺส, ยสฺส นกฺขมติ, โส ภาเสยฺย.}

\prose{สมฺมโต สํเฆน อิตฺถนฺนาโม ภิกฺขุ ภิกฺขุโนวาทโก, ขมติ สํฆสฺส, ตสฺมา ตุณฺหี, เอวเมตํ ธารยามี”ติ.}

\prose{อถ โข ภควา ฉพฺพคฺคิเย ภิกฺขู อเนกปริยาเยน วิครหิตฺวา ทุพฺภรตาย ฯเปฯ.\csromanspacer{} เอวญฺจ ปน ภิกฺขเว อิมํ สิกฺขาปทํ อุทฺทิเสยฺยาถ\csromandash{}}

\proseitem{146}{\textbf{“โย ปน ภิกฺขุ อสมฺมโต ภิกฺขุนิโย โอวเทยฺย, ปาจิตฺติยนฺ”}ติ.}

\prose{เอวญฺจิทํ ภควตา ภิกฺขูนํ สิกฺขาปทํ ปญฺญตฺตํ โหติ.}

\proseitem{147}{เตน โข ปน สมเยน เถรา ภิกฺขู สมฺมตา ภิกฺขุนิโย โอวทนฺตา ตเถว ลาภิโน โหนฺติ จีวรปิณฺฑปาตเสนาสน\-คิลานปฺปจฺจย- เภสชฺชปริกฺขารานํ.\csromanspacer{} อถ โข ฉพฺพคฺคิยานํ ภิกฺขูนํ เอตทโหสิ “เอตรหิ โข อาวุโส เถรา ภิกฺขู สมฺมตา ภิกฺขุนิโย โอวทนฺตา \csromanfolio{73}ตเถว ลาภิโน โหนฺติ จีวรปิณฺฑปาตเสนาสน\-คิลานปฺปจฺจย\-เภสชฺช- ปริกฺขารานํ, หนฺทาวุโส มยมฺปิ นิสฺสีมํ คนฺตฺวา อญฺญมญฺญํ ภิกฺขุโนวาทกํ สมฺมนฺนิตฺวา ภิกฺขุนิโย โอวทามา”ติ.\csromanspacer{} อถ โข ฉพฺพคฺคิยา ภิกฺขู นิสฺสีมํ คนฺตฺวา อญฺญมญฺญํ ภิกฺขุโนวาทกํ สมฺมนฺนิตฺวา ภิกฺขุนิโย อุปสงฺกมิตฺวา เอตทโวจุํ “มยมฺปิ ภคินิโย สมฺมตา อเมฺหปิ อุปสงฺกมถ มยมฺปิ โอวทิสฺสามา”ติ.}

\prose{อถ โข ตา ภิกฺขุนิโย เยน ฉพฺพคฺคิยา ภิกฺขู เตนุปสงฺกมิํสุ, อุปสงฺกมิตฺวา ฉพฺพคฺคิเย ภิกฺขู อภิวาเทตฺวา เอกมนฺตํ นิสีทิํสุ.\csromanspacer{} อถ โข ฉพฺพคฺคิยา ภิกฺขู ภิกฺขุนีนํ ปริตฺตญฺเญว ธมฺมิํ กถํ กตฺวา ทิวสํ ติรจฺฉานกถาย วีตินาเมตฺวา อุยฺโยเชสุํ “คจฺฉถ ภคินิโย”ติ.\csromanspacer{} อถ โข ตา ภิกฺขุนิโย เยน ภควา เตนุปสงฺกมิํสุ, อุปสงฺกมิตฺวา ภควนฺตํ อภิวาเทตฺวา เอกมนฺตํ อฏฺฐํสุ.\csromanspacer{} เอกมนฺตํ ฐิตา โข ตา ภิกฺขุนิโย ภควา เอตทโวจ “กจฺจิ ภิกฺขุนิโย โอวาโท อิทฺโธ อโหสี”ติ.\csromanspacer{} กุโต ภนฺเต โอวาโท อิทฺโธ ภวิสฺสติ, อยฺยา ฉพฺพคฺคิยา ปริตฺตญฺเญว ธมฺมิํ กถํ กตฺวา ทิวสํ ติรจฺฉานกถาย วีตินาเมตฺวา อุยฺโยเชสุนฺติ.}

\prose{อถ โข ภควา ตา ภิกฺขุนิโย ธมฺมิยา กถาย สนฺทสฺเสสิ ฯเปฯ.\csromanspacer{} อถ โข ตา ภิกฺขุนิโย ภควตา ธมฺมิยา กถาย สนฺทสฺสิตา สมาทปิตา สมุตฺเตชิตา สมฺปหํสิตา ภควนฺตํ อภิวาเทตฺวา ปทกฺขิณํ กตฺวา ปกฺกมิํสุ.\csromanspacer{} อถ โข ภควา เอตสฺมิํ นิทาเน เอตสฺมิํ ปกรเณ ภิกฺขุสํฆํ สนฺนิปาตาเปตฺวา ฉพฺพคฺคิเย ภิกฺขู ปฏิปุจฺฉิ “สจฺจํ กิร ตุเมฺห ภิกฺขเว ภิกฺขุนีนํ ปริตฺตญฺเญว ธมฺมิํ กถํ กตฺวา ทิวสํ ติรจฺฉานกถาย วีตินาเมตฺวา อุยฺโยเชถา”ติ.\csromanspacer{} สจฺจํ ภควาติ.\csromanspacer{} วิครหิ พุทฺโธ ภควา ฯเปฯ.\csromanspacer{} กถํ หิ นาม ตุเมฺห โมฆปุริสา ภิกฺขุนีนํ ปริตฺตญฺเญว ธมฺมิํ กถํ กตฺวา ทิวสํ ติรจฺฉานกถาย วีตินาเมตฺวา อุยฺโยเชสฺสถ.\csromanspacer{} เนตํ โมฆปุริสา อปฺปสนฺนานํ วา ปสาทาย ฯเปฯ วิครหิตฺวา ฯเปฯ ธมฺมิํ กถํ กตฺวา ภิกฺขู อามนฺเตสิ “อนุชานามิ ภิกฺขเว อฏฺฐหงฺเคหิ สมนฺนาคตํ ภิกฺขุํ ภิกฺขุโนวาทกํ สมฺมนฺนิตุํ,\csromandash{}สีลวา โหติ ปาติโมกฺขสํวรสํวุโต วิหรติ อาจารโคจรสมฺปนฺโน, อณุมตฺเตสุ วชฺเชสุ ภยทสฺสาวี สมาทาย สิกฺขติ สิกฺขาปเทสุ.\csromanspacer{} พหุสฺสุโต โหติ สุตธโร สุตสนฺนิจโย, เย เต ธมฺมา อาทิกลฺยาณา มชฺเฌ \csromanfolio{74}กลฺยาณา ปริโยสานกลฺยาณา สาตฺถํ สพฺยญฺชนํ เกวลปริปุณฺณํ ปริสุทฺธํ พฺรหฺมจริยํ อภิวทนฺติ, ตถารูปาสฺส ธมฺมา พหุสฺสุตา โหนฺติ ธาตา วจสา ปริจิตา มนสา อนุเปกฺขิตา ทิฏฺฐิยา สุปฺปฏิวิทฺธา.\csromanspacer{} อุภยานิ โข ปนสฺส ปาติโมกฺขานิ วิตฺถาเรน สฺวาคตานิ โหนฺติ สุวิภตฺตานิ สุปฺปวตฺตีนิ สุวินิจฺฉิตานิ สุตฺตโส อนุพฺยญฺชนโส.\csromanspacer{} กลฺยาณวาโจ โหติ กลฺยาณวากฺกรโณ.\csromanspacer{} เยภุยฺเยน ภิกฺขุนีนํ ปิโย โหติ มนาโป.\csromanspacer{} ปฏิพโล โหติ ภิกฺขุนิโย โอวทิตุํ.\csromanspacer{} น โข ปเนตํ\footnote{ปน ตํ (?)} ภควนฺตํ อุทฺทิสฺส ปพฺพชิตาย กาสายวตฺถวสนาย ครุธมฺมํ อชฺฌาปนฺนปุพฺโพ โหติ.\csromanspacer{} วีสติวสฺโส วา โหติ อติเรกวีสติวสฺโส วา.\csromanspacer{} อนุชานามิ ภิกฺขเว อิเมหิ อฏฺฐหงฺเคหิ สมนฺนาคตํ ภิกฺขุํ ภิกฺขุโนวาทกํ สมฺมนฺนิตุนฺ”ติ.}

\proseitem{148}{\textbf{โย ปนา}ติ โย ยาทิโส ฯเปฯ.\csromanspacer{} \textbf{ภิกฺขู}ติ ฯเปฯ อยํ อิมสฺมิํ อตฺเถ อธิปฺเปโต ภิกฺขูติ.}

\prose{\textbf{อสมฺมโต} นาม ญตฺติจตุตฺเถน กมฺเมน อสมฺมโต.}

\prose{\textbf{ภิกฺขุนิโย} นาม อุภโตสํเฆ อุปสมฺปนฺนา.}

\prose{\textbf{โอวเทยฺยา}ติ อฏฺฐหิ ครุธมฺเมหิ โอวทติ, อาปตฺติ ปาจิตฺติยสฺส.\csromanspacer{} อญฺเญน ธมฺเมน โอวทติ, อาปตฺติ ทุกฺกฏสฺส.\csromanspacer{} เอกโต อุปสมฺปนฺนํ โอวทติ, อาปตฺติ ทุกฺกฏสฺส.}

\proseitem{149}{เตน สมฺมเตน ภิกฺขุนา ปริเวณํ สมฺมชฺชิตฺวา ปานียํ ปริโภชนียํ อุปฏฺฐาเปตฺวา อาสนํ ปญฺญเปตฺวา ทุติยํ คเหตฺวา นิสีทิตพฺพํ.\csromanspacer{} ภิกฺขุนีหิ ตตฺถ คนฺตฺวา ตํ ภิกฺขุํ อภิวาเทตฺวา เอกมนฺตํ นิสีทิตพฺพํ.\csromanspacer{} เตน ภิกฺขุนา ปุจฺฉิตพฺพา “สมคฺคาตฺถ ภคินิโย”ติ.\csromanspacer{} สเจ “สมคฺคามฺหายฺยา”ติ ภณนฺติ, วตฺตนฺติ ภคินิโย อฏฺฐ ครุธมฺมาติ.\csromanspacer{} สเจ “วตฺตนฺตายฺยา”ติ ภณนฺติ, “เอโส ภคินิโย โอวาโท”ติ นิยฺยาเทตพฺโพ\footnote{นิยฺยาเตตพฺโพ (อิติปิ)}.\csromanspacer{} สเจ “น วตฺตนฺตายฺยา”ติ ภณนฺติ, โอสาเรตพฺพา “วสฺสสตูปสมฺปนฺนาย ภิกฺขุนิยา ตทหุปสมฺปนฺนสฺส ภิกฺขุโน อภิวาทนํ ปจฺจุฏฺฐานํ อญฺชลิกมฺมํ สามีจิกมฺมํ กาตพฺพํ, อยมฺปิ ธมฺโม สกฺกตฺวา ครุกตฺวา มาเนตฺวา \csromanfolio{75}ปูเชตฺวา ยาวชีวํ อนติกฺกมนีโย.\csromanspacer{} น ภิกฺขุนิยา อภิกฺขุเก อาวาเส วสฺสํ วสิตพฺพํ, อยมฺปิ ธมฺโม สกฺกตฺวา ครุกตฺวา มาเนตฺวา ปูเชตฺวา ยาวชีวํ อนติกฺกมนีโย.\csromanspacer{} อนฺวทฺธมาสํ ภิกฺขุนิยา ภิกฺขุสํฆโต เทฺว ธมฺมา ปจฺจาสีสิตพฺพา\footnote{ปจฺจาสิํ สิตพฺพา (อิติปิ)} อุโปสถปุจฺฉกญฺจ โอวาทุปสงฺกมนญฺจ, อยมฺปิ ธมฺโม ฯเปฯ.\csromanspacer{} วสฺสํวุฏฺฐาย ภิกฺขุนิยา อุภโตสํเฆ ตีหิ ฐาเนหิ ปวาเรตพฺพํ ทิฏฺเฐน วา สุเตน วา ปริสงฺกาย วา, อยมฺปิ ธมฺโม ฯเปฯ.\csromanspacer{} ครุธมฺมํ อชฺฌาปนฺนาย ภิกฺขุนิยา อุภโตสํเฆ ปกฺขมานตฺตํ จริตพฺพํ, อยมฺปิ ธมฺโม ฯเปฯ.\csromanspacer{} เทฺว วสฺสานิ ฉสุ ธมฺเมสุ สิกฺขิตสิกฺขาย สิกฺขมานาย อุภโตสํเฆ อุปสมฺปทา ปริเยสิตพฺพา, อยมฺปิ ธมฺโม ฯเปฯ.\csromanspacer{} น ภิกฺขุนิยา เกนจิ ปริยาเยน ภิกฺขุ อกฺโกสิตพฺโพ ปริภาสิตพฺโพ, อยมฺปิ ธมฺโม ฯเปฯ.\csromanspacer{} อชฺชตคฺเค โอวโฏ ภิกฺขุนีนํ ภิกฺขูสุ วจนปโถ, อโนวโฏ ภิกฺขูนํ ภิกฺขุนีสุ วจนปโถ, อยมฺปิ ธมฺโม สกฺกตฺวา ครุกตฺวา มาเนตฺวา ปูเชตฺวา ยาวชีวํ อนติกฺกมนีโย”ติ.}

\prose{สเจ “สมคฺคามฺหายฺยา”ติ ภณนฺตํ อญฺญํ ธมฺมํ ภณติ, อาปตฺติ ทุกฺกฏสฺส.\csromanspacer{} สเจ “วคฺคามฺหายฺยา”ติ ภณนฺตํ อฏฺฐ ครุธมฺเม ภณติ, อาปตฺติ ทุกฺกฏสฺส.\csromanspacer{} โอวาทํ อนิยฺยาเทตฺวา อญฺญํ ธมฺมํ ภณติ, อาปตฺติ ทุกฺกฏสฺส.}

\proseitem{150}{อธมฺมกมฺเม อธมฺมกมฺมสญฺญี วคฺคํ ภิกฺขุนิสํฆํ วคฺคสญฺญี โอวทติ, อาปตฺติ ปาจิตฺติยสฺส.\csromanspacer{} อธมฺมกมฺเม อธมฺมกมฺมสญฺญี วคฺคํ ภิกฺขุนิสํฆํ เวมติโก โอวทติ, อาปตฺติ ปาจิตฺติยสฺส.\csromanspacer{} อธมฺมกมฺเม อธมฺมกมฺมสญฺญี วคฺคํ ภิกฺขุนิสํฆํ สมคฺคสญฺญี โอวทติ, อาปตฺติ ปาจิตฺติยสฺส.}

\prose{อธมฺมกมฺเม เวมติโก วคฺคํ ภิกฺขุนิสํฆํ วคฺคสญฺญี โอวทติ, อาปตฺติ ปาจิตฺติยสฺส.\csromanspacer{} อธมฺมกมฺเม เวมติโก วคฺคํ ภิกฺขุนิสํฆํ เวมติโก โอวทติ, อาปตฺติ ปาจิตฺติยสฺส.\csromanspacer{} อธมฺมกมฺเม เวมติโก วคฺคํ ภิกฺขุนิสํฆํ สมคฺคสญฺญี โอวทติ, อาปตฺติ}

\prose{อธมฺมกมฺเม ธมฺมกมฺมสญฺญี วคฺคํ ภิกฺขุนิสํฆํ วคฺคสญฺญี โอวทติ, อาปตฺติ ปาจิตฺติยสฺส.\csromanspacer{} อธมฺมกมฺเม ธมฺมกมฺมสญฺญี วคฺคํ ภิกฺขุนิสํฆํ เวมติโก โอวทติ, อาปตฺติ ปาจิตฺติยสฺส.\csromanspacer{} อธมฺมกมฺเม ธมฺมกมฺมสญฺญี วคฺคํ ภิกฺขุนิสํฆํ สมคฺคสญฺญี โอวทติ, อาปตฺติ ปาจิตฺติยสฺส.}

\prose{\csromanfolio{76}อธมฺมกมฺเม อธมฺมกมฺมสญฺญี สมคฺคํ ภิกฺขุนิสํฆํ วคฺคสญฺญี โอวทติ, อาปตฺติ ปาจิตฺติยสฺส.\csromanspacer{} อธมฺมกมฺเม อธมฺมกมฺมสญฺญี สมคฺคํ ภิกฺขุนิสํฆํ เวมติโก โอวทติ, อาปตฺติ ปาจิตฺติยสฺส.\csromanspacer{} อธมฺมกมฺเม อธมฺมกมฺมสญฺญี สมคฺคํ ภิกฺขุนิสํฆํ สมคฺคสญฺญี โอวทติ, อาปตฺติ ปาจิตฺติยสฺส.}

\prose{อธมฺมกมฺเม เวมติโก สมคฺคํ ภิกฺขุนิสํฆํ วคฺคสญฺญี โอวทติ ฯเปฯ เวมติโก โอวทติ ฯเปฯ สมคฺคสญฺญี โอวทติ, อาปตฺติ ปาจิตฺติยสฺส.}

\prose{อธมฺมกมฺเม ธมฺมกมฺมสญฺญี สมคฺคํ ภิกฺขุนิสํฆํ วคฺคสญฺญี โอวทติ ฯเปฯ เวมติโก โอวทติ ฯเปฯ สมคฺคสญฺญี โอวทติ, อาปตฺติ}

\proseitem{151}{ธมฺมกมฺเม อธมฺมกมฺมสญฺญี วคฺคํ ภิกฺขุนิสํฆํ วคฺคสญฺญี โอวทติ, อาปตฺติ ทุกฺกฏสฺส.\csromanspacer{} ธมฺมกมฺเม อธมฺมกมฺมสญฺญี วคฺคํ ภิกฺขุนิสํฆํ เวมติโก โอวทติ ฯเปฯ สมคฺคสญฺญี โอวทติ, อาปตฺติ ทุกฺกฏสฺส.}

\prose{ธมฺมกมฺเม เวมติโก วคฺคํ ภิกฺขุนิสํฆํ วคฺคสญฺญี โอวทติ ฯเปฯ เวมติโก โอวทติ ฯเปฯ สมคฺคสญฺญี โอวทติ, อาปตฺติ ทุกฺกฏสฺส.}

\prose{ธมฺมกมฺเม ธมฺมกมฺมสญฺญี วคฺคํ ภิกฺขุนิสํฆํ วคฺคสญฺญี โอวทติ ฯเปฯ เวมติโก โอวทติ ฯเปฯ สมคฺคสญฺญี โอวทติ, อาปตฺติ ทุกฺกฏสฺส.}

\prose{ธมฺมกมฺเม อธมฺมกมฺมสญฺญี สมคฺคํ ภิกฺขุนิสํฆํ วคฺคสญฺญี โอวทติ ฯเปฯ เวมติโก โอวทติ ฯเปฯ สมคฺคสญฺญี โอวทติ, อาปตฺติ ทุกฺกฏสฺส.}

\prose{ธมฺมกมฺเม เวมติโก สมคฺคํ ภิกฺขุนิสํฆํ วคฺคสญฺญี โอวทติ ฯเปฯ เวมติโก โอวทติ ฯเปฯ สมคฺคสญฺญี โอวทติ, อาปตฺติ ทุกฺกฏสฺส.}

\prose{ธมฺมกมฺเม ธมฺมกมฺมสญฺญี สมคฺคํ ภิกฺขุนิสํฆํ วคฺคสญฺญี โอวทติ, อาปตฺติ ทุกฺกฏสฺส.\csromanspacer{} ธมฺมกมฺเม ธมฺมกมฺมสญฺญี สมคฺคํ ภิกฺขุนิสํฆํ เวมติโก โอวทติ, อาปตฺติ ทุกฺกฏสฺส.\csromanspacer{} ธมฺมกมฺเม ธมฺมกมฺมสญฺญี สมคฺคํ ภิกฺขุนิสํฆํ สมคฺคสญฺญี โอวทติ, อนาปตฺติ.}

\proseitemwithcloser{152}{อนาปตฺติ อุทฺเทสํ เทนฺโต, ปริปุจฺฉํ เทนฺโต, “โอสาเรหิ อยฺยา”ติ วุจฺจมาโน โอสาเรติ, ปญฺหํ ปุจฺฉติ, ปญฺหํ ปุฏฺโฐ กเถติ, อญฺญสฺสตฺถาย ภณนฺตํ ภิกฺขุนิโย สุณนฺติ, สิกฺขมานาย สามเณริยา อุมฺมตฺตกสฺส อาทิกมฺมิกสฺสาติ.}{โอวาทสิกฺขาปทํ นิฏฺฐิตํ ปฐมํ.}
\csromanmatikaanchor{mtk.26}
\csromantocmarkref{subsection}{2. อตฺถงฺคตสิกฺขาปท}{mtk.26}
\bookmark[dest={mtk.26},level=2]{2. อตฺถงฺคตสิกฺขาปท}
\csromanheader{2}{3. \textbf{โอวาทวคฺค 2. อตฺถงฺคตสิกฺขาปท}}
\proseitem{153}{\csromanfolio{77}เตน สมเยน พุทฺโธ ภควา สาวตฺถิยํ วิหรติ เชตวเน อนาถปิณฺฑิกสฺส อาราเม.\csromanspacer{} เตน โข ปน สมเยน เถรา ภิกฺขู ภิกฺขุนิโย โอวทนฺติ ปริยาเยน.\csromanspacer{} เตน โข ปน สมเยน อายสฺมโต จูฬปนฺถกสฺส ปริยาโย โหติ ภิกฺขุนิโย โอวทิตุํ.\csromanspacer{} ภิกฺขุนิโย เอวมาหํสุ “น ทานิ อชฺช โอวาโท อิทฺโธ ภวิสฺสติ, ตญฺเญว ทานิ อุทานํ อยฺโย จูฬปนฺถโก ปุนปฺปุนํ ภณิสฺสตี”ติ.\csromanspacer{} อถ โข ตา ภิกฺขุนิโย เยนายสฺมา จูฬปนฺถโก เตนุปสงฺกมิํสุ, อุปสงฺกมิตฺวา อายสฺมนฺตํ จูฬปนฺถกํ อภิวาเทตฺวา เอกมนฺตํ นิสีทิํสุ.\csromanspacer{} เอกมนฺตํ นิสินฺนา โข ตา ภิกฺขุนิโย อายสฺมา จูฬปนฺถโก เอตทโวจ “สมคฺคาตฺถ ภคินิโย”ติ.\csromanspacer{} สมคฺคามฺหายฺยาติ.\csromanspacer{} วตฺตนฺติ ภคินิโย อฏฺฐ ครุธมฺมาติ.\csromanspacer{} วตฺตนฺตายฺยาติ. “เอโส ภคินิโย โอวาโท”ติ นิยฺยาเทตฺวา อิมํ อุทานํ ปุนปฺปุนํ อภาสิ\csromandash{}}

\setlength{\csromangathapagewidth}{0pt}
\csromangathameasuregroup{}{*“อธิเจตโส อปฺปมชฺชโต, \\ มุนิโน โมนปเถสุ สิกฺขโต. \\ โสกา น ภวนฺติ ตาทิโน, \\ อุปสนฺถสฺส สทา สตีมโต”ติ.}
\setlength{\csromangathaleftcol}{0pt}
\csromangathagroup{\csromangathastanza{\csromansymbolfootnote{*}{ขุ 1. 128 ปิฏฺเฐ อุทาเนปิ.}“อธิเจตโส อปฺปมชฺชโต, \\ มุนิโน โมนปเถสุ สิกฺขโต. \\ โสกา น ภวนฺติ ตาทิโน, \\ อุปสนฺถสฺส สทา สตีมโต”ติ.}}

\prose{ภิกฺขุนิโย เอวมาหํสุ “นนุ อโวจุมฺหา น ทานิ อชฺช โอวาโท อิทฺโธ ภวิสฺสติ, ตญฺเญว ทานิ อุทานํ อยฺโย จูฬปนฺถโก ปุนปฺปุนํ ภณิสฺสตี”ติ.\csromanspacer{} อสฺโสสิ โข อายสฺมา จูฬปนฺถโก ตาสํ ภิกฺขูนีนํ อิมํ กถาสลฺลาปํ.\csromanspacer{} อถ โข อายสฺมา จูฬปนฺถโก เวหาสํ อพฺภุคฺคนฺตฺวา อากาเส อนฺตลิกฺเข จงฺกมติปิ ติฏฺฐติปิ นิสีทติปิ เสยฺยมฺปิ กปฺเปติ ธูมายติปิ ปชฺชลติปิ อนฺตรธายติปิ.\csromanspacer{} ตญฺเจว\footnote{ตญฺเญว (อิติปิ)} อุทานํ ภณติ อญฺญญฺจ พหุํ พุทฺธวจนํ.\csromanspacer{} ภิกฺขุนิโย เอวมาหํสุ “อจฺฉริยํ วต โภ, อพฺภุตํ วต โภ, น วต โน อิโต ปุพฺเพ โอวาโท เอวํ อิทฺโธ ภูตปุพฺโพ ยถา อยฺยสฺส จูฬปนฺถกสฺสา”ติ.\csromanspacer{} อถ โข อายสฺมา จูฬปนฺถโก ตา ภิกฺขุนิโย ยาว สมนฺธการา โอวทิตฺวา อุยฺโยเชสิ “คจฺฉถ ภคินิโย”ติ.}

\prose{\csromanfolio{78}อถ โข ตา ภิกฺขุนิโย นครทฺวาเร ถกิเต พหินคเร วสิตฺวา กาลสฺเสว นครํ ปวิสนฺติ.\csromanspacer{} มนุสฺสา อุชฺฌายนฺติ ขิยฺยนฺติ วิปาเจนฺติ “อพฺรหฺมจารินิโย อิมา ภิกฺขุนิโย อาราเม ภิกฺขูหิ สทฺธิํ วสิตฺวา อิทานิ นครํ ปวิสนฺตี”ติ.\csromanspacer{} อสฺโสสุํ โข ภิกฺขู เตสํ มนุสฺสานํ อุชฺฌายนฺตานํ ขิยฺยนฺตานํ วิปาเจนฺตานํ.\csromanspacer{} เย เต ภิกฺขู อปฺปิจฺฉา ฯเปฯ เต อุชฺฌายนฺติ ขิยฺยนฺติ วิปาเจนฺติ “กถํ หิ นาม อายสฺมา จูฬปนฺถโก อตฺถงฺคเต สูริเย ภิกฺขุนิโย โอวทิสฺสตี”ติ ฯเปฯ “สจฺจํ กิร ตฺวํ จูฬปนฺถก อตฺถงฺคเต สูริเย ภิกฺขุนิโย โอวทสี”ติ.\csromanspacer{} สจฺจํ ภควาติ.\csromanspacer{} วิครหิ พุทฺโธ ภควา ฯเปฯ.\csromanspacer{} กถํ หิ นาม ตฺวํ จูฬปนฺถก อตฺถงฺคเต สูริเย ภิกฺขุนิโย โอวทิสฺสสิ.\csromanspacer{} เนตํ จูฬปนฺถก อปฺปสนฺนานํ วา ปสาทาย ฯเปฯ.\csromanspacer{} เอวญฺจ ปน ภิกฺขเว อิมํ สิกฺขาปทํ อุทฺทิเสยฺยาถ\csromandash{}}

\proseitem{154}{\textbf{“สมฺมโตปิ เจ ภิกฺขุ อตฺถงฺคเต สูริเย ภิกฺขุนิโย โอวเทยฺย, ปาจิตฺติยนฺ”}ติ.}

\proseitem{155}{\textbf{สมฺมโต} นาม ญตฺติจตุตฺเถน กมฺเมน สมฺมโต.}

\prose{\textbf{อตฺถงฺคเต สูริเย}ติ โอคฺคเต สูริเย.}

\prose{\textbf{ภิกฺขุนี} นาม อุภโตสํเฆ อุปสมฺปนฺนา.}

\prose{\textbf{โอวเทยฺยา}ติ อฏฺฐหิ วา ครุธมฺเมหิ อญฺเญน วา ธมฺเมน โอวทติ, อาปตฺติ ปาจิตฺติยสฺส.}

\proseitem{156}{อตฺถงฺคเต อตฺถงฺคตสญฺญี โอวทติ, อาปตฺติ ปาจิตฺติยสฺส.\csromanspacer{} อตฺถงฺคเต เวมติโก โอวทติ, อาปตฺติ ปาจิตฺติยสฺส.\csromanspacer{} อตฺถงฺคเต อนตฺถงฺคตสญฺญี โอวทติ, อาปตฺติ ปาจิตฺติยสฺส.}

\prose{เอกโตอุปสมฺปนฺนาย โอวทติ, อาปตฺติ ทุกฺกฏสฺส.\csromanspacer{} อนตฺถงฺคเต อตฺถงฺคตสญฺญี, อาปตฺติ ทุกฺกฏสฺส.\csromanspacer{} อนตฺตงฺคเต เวมติโก, อาปตฺติ ทุกฺกฏสฺส.\csromanspacer{} อนตฺถงฺคเต อนตฺถงฺคตสญฺญี, อนาปตฺติ.}

\proseitemwithcloser{157}{อนาปตฺติ อุทฺเทสํ เทนฺโต, ปริปุจฺฉํ เทนฺโต, “โอสาเรหิ อยฺยา”ติ วุจฺจมาโน โอสาเรติ, ปญฺหํ ปุจฺฉติ, ปญฺหํ ปุฏฺโฐ กเถติ, อญฺญสฺสตฺถาย ภณนฺตํ ภิกฺขุนิโย สุณนฺติ, สิกฺขมานาย สามเณริยา อุมฺมตฺตกสฺส อาทิกมฺมิกสฺสาติ.}{อตฺเถงฺคตสิกฺขาปทํ นิฏฺฐิตํ ทุติยํ.}
\csromanmatikaanchor{mtk.27}
\csromantocmarkref{subsection}{3. ภิกฺขุนุปสฺสยสิกฺขาปท}{mtk.27}
\bookmark[dest={mtk.27},level=2]{3. ภิกฺขุนุปสฺสยสิกฺขาปท}
\csromanheader{2}{3. \textbf{โอวาทวคฺค 3. ภิกฺขุนุปสฺสยสิกฺขาปท}}
\proseitem{158}{\csromanfolio{79}เตน สมเยน พุทฺโธ ภควา สกฺเกสุ วิหรติ กปิลวตฺถุสฺมิํ นิโคฺรธาราเม.\csromanspacer{} เตน โข ปน สมเยน ฉพฺพคฺคิยา ภิกฺขู ภิกฺขุนุปสฺสยํ อุปสงฺกมิตฺวา ฉพฺพคฺคิยา ภิกฺขุนิโย โอวทนฺติ.\csromanspacer{} ภิกฺขุนิโย ฉพฺพคฺคิยา ภิกฺขุนิโย เอตทโวจุํ “เอถายฺเย โอวาทํ คมิสฺสามา”ติ.\csromanspacer{} ยมฺปิ\footnote{ยํ หิ (ก)} มยํ อยฺเย คจฺเฉยฺยาม โอวาทสฺส การณา, อยฺยา ฉพฺพคฺคิยา อิเธว อาคนฺตฺวา อเมฺห โอวทนฺตีติ.\csromanspacer{} ภิกฺขุนิโย อุชฺฌายนฺติ ขิยฺยนฺติ วิปาเจนฺติ “กถํ หิ นาม ฉพฺพคฺคิยา ภิกฺขู ภิกฺขุนุปสฺสยํ อุปสงฺกมิตฺวา ภิกฺขุนิโย โอวทิสฺสนฺตี”ติ\footnote{ฉพฺพคฺคิยา ภิกฺขุนิโย โอวาทํ น คจฺฉิสฺสนฺตีติ (สี)}.\csromanspacer{} อถ โข ตา ภิกฺขุนิโย ภิกฺขูนํ เอตมตฺถํ อาโรเจสุํ.\csromanspacer{} เย เต ภิกฺขู อปฺปิจฺฉา ฯเปฯ เต อุชฺฌายนฺติ ขิยฺยนฺติ วิปาเจนฺติ “กถํ หิ นาม ฉพฺพคฺคิยา ภิกฺขู ภิกฺขุนุปสฺสยํ อุปสงฺกมิตฺวา ภิกฺขุนิโย โอวทิสฺสนฺตี”ติ ฯเปฯ.\csromanspacer{} สจฺจํ กิร ตุเมฺห ภิกฺขเว ภิกฺขุนุปสฺสยํ อุปสงฺกมิตฺวา ภิกฺขุนิโย โอวทถาติ.\csromanspacer{} สจฺจํ ภควาติ.\csromanspacer{} วิครหิ พุทฺโธ ภควา ฯเปฯ.\csromanspacer{} กถํ หิ นาม ตุเมฺห โมฆปุริสา ภิกฺขุนุปสฺสยํ อุปสงฺกมิตฺวา ภิกฺขุนิโย โอวทิสฺสถ.\csromanspacer{} เนตํ โมฆปุริสา อปฺปสนฺนานํ วา ปสาทาย ฯเปฯ.\csromanspacer{} เอวญฺจ ปน ภิกฺขเว อิมํ สิกฺขาปทํ อุทฺทิเสยฺยาถ\csromandash{}}

\prose{\textbf{“โย ปน ภิกฺขุ ภิกฺขุนุปสฺสยํ อุปสงฺกมิตฺวา ภิกฺขุนิโย โอวเทยฺย, ปาจิตฺติยนฺ”}ติ.}

\prose{เอวญฺจิทํ ภควตา ภิกฺขูนํ สิกฺขาปทํ ปญฺญตฺตํ โหติ.}

\proseitem{159}{เตน โข ปน สมเยน มหาปชาปติ โคตมี คิลานา โหติ.\csromanspacer{} เถรา ภิกฺขู เยน มหาปชาปติ โคตมี เตนุปสงฺกมิํสุ, อุปสงฺกมิตฺวา มหาปชาปติํ โคตมิํ เอตทโวจุํ “กจฺจิ โคตมิ ขมนียํ, กจฺจิ ยาปนียนฺ”ติ.\csromanspacer{} น เม อยฺยา ขมนียํ, น ยาปนียํ, อิงฺฆยฺยา ธมฺมํ เทเสถาติ. “น ภคินิ กปฺปติ ภิกฺขุนุปสฺสยํ อุปสงฺกมิตฺวา ภิกฺขุนิโย ธมฺมํ เทเสตุนฺ”ติ กุกฺกุจฺจายนฺตา น เทเสสุํ.\csromanspacer{} อถ โข ภควา ปุพฺพณฺหสมยํ นิวาเสตฺวา ปตฺตจีวรมาทาย เยน มหาปชาปติ โคตมี เตนุปสงฺกมิ, อุปสงฺกมิตฺวา ปญฺญตฺเต อาสเน นิสีทิ.\csromanspacer{} นิสชฺช โข ภควา มหาปชาปติํ โคตมิํ เอตทโวจ “กจฺจิ โคตมิ ขมนียํ, กจฺจิ \csromanfolio{80}ยาปนียนฺ”ติ.\csromanspacer{} ปุพฺเพ เม ภนฺเต เถรา ภิกฺขู อาคนฺตฺวา ธมฺมํ เทเสนฺติ, เตน เม ผาสุ โหติ, อิทานิ ปน “ภควตา ปฏิกฺขิตฺตนฺ”ติ กุกฺกุจฺจายนฺตา น เทเสนฺติ, เตน เม น ผาสุ โหตีติ.\csromanspacer{} อถ โข ภควา มหาปชาปติํ โคตมิํ ธมฺมิยา กถาย สนฺทสฺเสตฺวา สมาทเปตฺวา สมุตฺเตเชตฺวา สมฺปหํเสตฺวา อุฏฺฐายาสนา ปกฺกามิ.\csromanspacer{} อถ โข ภควา เอตสฺมิํ นิทาเน เอตสฺมิํ ปกรเณ ธมฺมิํ กถํ กตฺวา ภิกฺขู อามนฺเตสิ อนุชานามิ ภิกฺขเว ภิกฺขุนุปสฺสยํ อุปสงฺกมิตฺวา คิลานํ ภิกฺขุนิํ โอวทิตุํ.\csromanspacer{} เอวญฺจ ปน ภิกฺขเว อิมํ สิกฺขาปทํ อุทฺทิเสยฺยาถ\csromandash{}}

\proseitem{160}{\textbf{“โย ปน ภิกฺขุ ภิกฺขุนุปสฺสยํ อุปสงฺกมิตฺวา ภิกฺขุนิโย โอวเทยฺย, อญฺญตฺร สมยา ปาจิตฺติยํ.}\csromanspacer{} ตตฺถายํ สมโย, คิลานา โหติ ภิกฺขุนี, อยํ ตตฺถ สมโย”ติ.}

\proseitem{161}{\textbf{โย ปนา}ติ โย ยาทิโส ฯเปฯ.\csromanspacer{} \textbf{ภิกฺขู}ติ ฯเปฯ อยํ อิมสฺมิํ อตฺเถ อธิปฺเปโต ภิกฺขูติ.}

\prose{\textbf{ภิกฺขุนุปสฺสโย} นาม ยตฺถ ภิกฺขุนิโย เอกรตฺตํปิ วสนฺติ.}

\prose{\textbf{อุปสงฺกมิตฺวา}ติ ตตฺถ คนฺตฺวา.}

\prose{\textbf{ภิกฺขุนี} นาม อุภโตสํเฆ อุปสมฺปนฺนา.}

\prose{\textbf{โอวเทยฺยา}ติ อฏฺฐหิ ครุธมฺเมหิ โอวทติ, อาปตฺติ ปาจิตฺติยสฺส.}

\prose{\textbf{อญฺญตฺร สมยา}ติ ฐเปตฺวา สมยํ.}

\prose{\textbf{คิลานา} นาม ภิกฺขูนี น สกฺโกติ โอวาทาย วา สํวาสาย วา คนฺตุํ.}

\proseitem{162}{อุปสมฺปนฺนาย อุปสมฺปนฺนสญฺญี ภิกฺขุนุปสฺสยํ อุปสงฺกมิตฺวา อญฺญตฺร สมยา โอวทติ, อาปตฺติ ปาจิตฺติยสฺส.\csromanspacer{} อุปสมฺปนฺนาย เวมติโก ภิกฺขุนุปสฺสยํ อุปสงฺกมิตฺวา อญฺญตฺร สมยา โอวทติ, อาปตฺติ ปาจิตฺติยสฺส.\csromanspacer{} อุปสมฺปนฺนาย อนุปสมฺปนฺนสญฺญี ภิกฺขุนุปสฺสยํ อุปสงฺกมิตฺวา อญฺญตฺร สมยา โอวทติ, อาปตฺติ ปาจิตฺติยสฺส.}

\prose{อญฺเญน ธมฺเมน โอวทติ, อาปตฺติ ทุกฺกฏสฺส.\csromanspacer{} เอกโตอุปสมฺปนฺนาย โอวทติ, อาปตฺติ ทุกฺกฏสฺส.\csromanspacer{} อนุปสมฺปนฺนาย อุปสมฺปนฺนสญฺญี, อาปตฺติ \csromanfolio{81}ทุกฺกฏสฺส.\csromanspacer{} อนุปสมฺปนฺนาย เวมติโก, อาปตฺติ ทุกฺกฏสฺส.\csromanspacer{} อนุปสมฺปนฺนาย อนุปสมฺปนฺนสญฺญี, อนาปตฺติ.}

\proseitemwithcloserruled{163}{อนาปตฺติ สมเย, อุทฺเทสํ เทนฺโต, ปริปุจฺฉํ เทนฺโต, “โอสาเรหิ อยฺยา”ติ วุจฺจมาโน โอสาเรติ, ปญฺหํ ปุจฺฉติ, ปญฺหํ ปุฏฺโฐ กเถติ, อญฺญสฺสตฺถาย ภณนฺตํ ภิกฺขุนิโย สุณนฺติ, สิกฺขมานาย สามเณริยา อุมฺมตฺตกสฺส อาทิกมฺมิกสฺสาติ.}{ภิกฺขุนุปสฺสยสิกฺขาปทํ นิฏฺฐิตํ ตติยํ.}
\csromanmatikaanchor{mtk.28}
\csromantocmarkref{subsection}{4. อามิสสิกฺขาปท}{mtk.28}
\bookmark[dest={mtk.28},level=2]{4. อามิสสิกฺขาปท}
\csromanheader{2}{3. \textbf{โอวาทวคฺค 4. อามิสสิกฺขาปท}}
\proseitem{164}{เตน สมเยน พุทฺโธ ภควา สาวตฺถิยํ วิหรติ เชตวเน อนาถปิณฺฑิกสฺส อาราเม.\csromanspacer{} เตน โข ปน สมเยน เถรา ภิกฺขู ภิกฺขุนิโย โอวทนฺตา ลาภิโน โหนฺติ จีวร\-ปิณฺฑปาต\-เสนาสน\-คิลานปฺปจฺจย\-เภสชฺช\-ปริกฺขารานํ.\csromanspacer{} ฉพฺพคฺคิยา ภิกฺขู เอวํ วทนฺติ “น พหุกตา เถรา ภิกฺขู ภิกฺขุนิโย โอวทิตุํ, อามิสเหตุ เถรา ภิกฺขู ภิกฺขุนิโย โอวทนฺตี”ติ.\csromanspacer{} เย เต ภิกฺขู อปฺปิจฺฉา ฯเปฯ เต อุชฺฌายนฺติ ขิยฺยนฺติ วิปาเจนฺติ “กถํ หิ นาม ฉพฺพคฺคิยา ภิกฺขู เอวํ วกฺขนฺติ น พหุกตา เถรา ภิกฺขู ภิกฺขุนิโย โอวทิตุํ, อามิสเหตุ เถรา ภิกฺขู ภิกฺขุนิโย โอวทนฺตี”ติ ฯเปฯ.\csromanspacer{} สจฺจํ กิร ตุเมฺห ภิกฺขเว เอวํ วเทถ “น พหุกตา เถรา ภิกฺขู ภิกฺขุนิโย โอวทิตุํ, อามิสเหตุ เถรา ภิกฺขู ภิกฺขุนิโย โอวทนฺตี”ติ.\csromanspacer{} สจฺจํ ภควาติ.\csromanspacer{} วิครหิ พุทฺโธ ภควา ฯเปฯ.\csromanspacer{} กถํ หิ นาม ตุเมฺห โมฆปุริสา เอวํ วกฺขถ “น พหุกตา เถรา ภิกฺขู ภิกฺขุนิโย โอวทิตุํ, อามิสเหตุ เถรา ภิกฺขู ภิกฺขุนิโย โอวทนฺตี”ติ.\csromanspacer{} เนตํ โมฆปุริสา อปฺปสนฺนานํ วา ปสาทาย ฯเปฯ.\csromanspacer{} เอวญฺจ ปน ภิกฺขเว อิมํ สิกฺขาปทํ อุทฺทิเสยฺยาถ\csromandash{}}

\proseitem{165}{\textbf{“โย ปน ภิกฺขุ เอวํ วเทยฺย อามิสเหตุ เถรา ภิกฺขู ภิกฺขุนิโย โอวทนฺตีติ, ปาจิตฺติยนฺ”}ติ.}

\proseitem{166}{\csromanfolio{82}\textbf{โย ปนา}ติ โย ยาทิโส ฯเปฯ.\csromanspacer{} \textbf{ภิกฺขู}ติ ฯเปฯ อยํ อิมสฺมิํ อตฺเถ อธิปฺเปโต ภิกฺขูติ.}

\prose{อามิสเหตูติ จีวรเหตุ ปิณฺฑปาตเหตุ เสนาสนเหตุ คิลานปฺปจฺจย\-เภสชฺช\-ปริกฺขาร\-เหตุ สกฺการเหตุ ครุการเหตุ มานนเหตุ วนฺทนเหตุ ปูชนเหตุ.}

\prose{เอวํ วเทยฺยาติ อุปสมฺปนฺนํ สํเฆน สมฺมตํ ภิกฺขุโนวาทกํ อวณฺณํ กตฺตุกาโม อยสํ กตฺตุกาโม มงฺกุกตฺตุกาโม เอวํ วเทติ “จีวรเหตุ ปิณฺฑปาตเหตุ เสนาสนเหตุ คิลานปฺปจฺจย\-เภสชฺช\-ปริกฺขาร\-เหตุ สกฺการเหตุ ครุการเหตุ มานนเหตุ วนฺทนเหตุ ปูชนเหตุ โอวทตี”ติ ภณติ, อาปตฺติ ปาจิตฺติยสฺส.}

\proseitem{167}{ธมฺมกมฺเม ธมฺมกมฺมสญฺญี เอวํ วเทติ, อาปตฺติยสฺส.\csromanspacer{} ธมฺมกมฺเม เวมติโก เอวํ วเทติ, อาปตฺติ ปาจิตฺติยสฺส.\csromanspacer{} ธมฺมกมฺเม อธมฺมกมฺมสญฺญี เอวํ วเทติ, อาปตฺติ ปาจิตฺติยสฺส.}

\prose{อุปสมฺปนฺนํ สํเฆน อสมฺมตํ ภิกฺขุโนวาทกํ อวณฺณํ กตฺตุกาโม อยสํ กตฺตุกาโม มงฺกุกตฺตุกาโม เอวํ วเทติ “จีวรเหตุ ฯเปฯ ปูชนเหตุ โอวทตี”ติ ภณติ, อาปตฺติ ทุกฺกฏสฺส.\csromanspacer{} อนุปสมฺปนฺนํ สํเฆน สมฺมตํ วา อสมฺมตํ วา ภิกฺขุโนวาทกํ อวณฺณํ กตฺตุกาโม อยสํ กตฺตุกาโม มงฺกุกตฺตุกาโม เอวํ วเทติ “จีวรเหตุ ปิณฺฑปาตเหตุ เสนาสนเหตุ คิลานปฺปจฺจย\-เภสชฺช\-ปริกฺขาร\-เหตุ สกฺการเหตุ ครุการเหตุ มานนเหตุ วนฺทนเหตุ ปูชนเหตุ โอวทตี”ติ ภณติ, อาปตฺติ ทุกฺกฏสฺส.\csromanspacer{} อธมฺมกมฺเม ธมฺมกมฺมสญฺญี, อาปตฺติ ทุกฺกฏสฺส.\csromanspacer{} อธมฺมกมฺเม เวมติโก, อาปตฺติ ทุกฺกฏสฺส.\csromanspacer{} อธมฺมกมฺเม อธมฺมกมฺมสญฺญี, อาปตฺติ ทุกฺกฏสฺส.}

\proseitemwithcloser{168}{อนาปตฺติ ปกติยา จีวรเหตุ ปิณฺฑปาตเหตุ เสนาสนเหตุ คิลานปฺปจฺจย\-เภสชฺช\-ปริกฺขาร\-เหตุ สกฺการเหตุ ครุการเหตุ มานนเหตุ วนฺทนเหตุ ปูชนเหตุ โอวทนฺตํ ภณติ, อุมฺมตฺตกสฺส อาทิกมฺมิกสฺสาติ.}{อามิสสิกฺขาปทํ นิฏฺฐิตํ จตุตฺถํ.}
\csromanmatikaanchor{mtk.29}
\csromantocmarkref{subsection}{5. จีวรทานสิกฺขาปท}{mtk.29}
\bookmark[dest={mtk.29},level=2]{5. จีวรทานสิกฺขาปท}
\csromanheader{2}{3. \textbf{โอวาทวคฺค 5. จีวรทานสิกฺขาปท}}
\proseitem{169}{\csromanfolio{83}เตน สมเยน พุทฺโธ ภควา สาวตฺถิยํ วิหรติ เชตวเน อนาถปิณฺฑิกสฺส อาราเม.\csromanspacer{} เตน โข ปน สมเยน อญฺญตโร ภิกฺขุ สาวตฺถิยํ อญฺญตริสฺสา วิสิขาย ปิณฺฑาย จรติ, อญฺญตราปิ ภิกฺขุนี ตสฺสา วิสิขาย ปิณฺฑาย จรติ.\csromanspacer{} อถ โข โส ภิกฺขุ ตํ ภิกฺขุนิํ เอตทโวจ “คจฺฉ ภคินิ อมุกสฺมิํ โอกาเส ภิกฺขา ทิยฺยตี”ติ, สาปิ โข เอวมาห “คจฺฉายฺย อมุกสฺมิํ โอกาเส ภิกฺขา ทิยฺยตี”ติ.\csromanspacer{} เต อภิณฺหทสฺสเนน สนฺทิฏฺฐา อเหสุํ.\csromanspacer{} เตน โข ปน สมเยน สํฆสฺส จีวรํ ภาชียติ.\csromanspacer{} อถ โข สา ภิกฺขุนี โอวาทํ คนฺตฺวา เยน โส ภิกฺขุ เตนุปสงฺกมิ, อุปสงฺกมิตฺวา ตํ ภิกฺขุํ อภิวาเทตฺวา เอกมนฺตํ อฏฺฐาสิ, เอกมนฺตํ ฐิตํ โข ตํ ภิกฺขุนิํ โส ภิกฺขุ เอตทโวจ “อยํ เม ภคินิ จีวรปฏิวีโส\footnote{ปฏิวิํ โส (สี), ปฏิวิโส (อิติปิ)} สาทิยิสฺสสี”ติ.\csromanspacer{} อามายฺย ทุพฺพลจีวรามฺหีติ.}

\prose{อถ โข โส ภิกฺขุ ตสฺสา ภิกฺขุนิยา จีวรํ อทาสิ.\csromanspacer{} โส ปิ โข ภิกฺขุ ทุพฺพลจีวโร โหติ.\csromanspacer{} ภิกฺขู ตํ ภิกฺขุํ เอตทโวจุํ “กโรหิ ทานิ เต อาวุโส จีวรนฺ”ติ.\csromanspacer{} อถ โข โส ภิกฺขุ ภิกฺขูนํ เอตมตฺถํ อาโรเจสิ.\csromanspacer{} เย เต ภิกฺขู อปฺปิจฺฉา ฯเปฯ เต อุชฺฌายนฺติ ขิยฺยนฺติ วิปาเจนฺติ “กถํ หิ นาม ภิกฺขุ ภิกฺขุนิยา จีวรํ ทสฺสตี”ติ ฯเปฯ.\csromanspacer{} สจฺจํ กิร ตฺวํ ภิกฺขุ ภิกฺขุนิยา จีวรํ อทาสีติ.\csromanspacer{} สจฺจํ ภควาติ.\csromanspacer{} ญาติกา เต ภิกฺขุ อญฺญาติกาติ.\csromanspacer{} อญฺญาติกา ภควาติ.\csromanspacer{} อญฺญาตโก โมฆปุริส อญฺญาติกาย น ชานาติ ปติรูปํ วา อปฺปติรูปํ วา สนฺตํ วา อสนฺตํ วา, กถํ หิ นาม ตฺวํ โมฆปุริส อญฺญาติกาย ภิกฺขุนิยา จิวรํ ทสฺสสิ.\csromanspacer{} เนตํ โมฆปุริส อปฺปสนฺนานํ วา ปสาทาย ฯเปฯ.\csromanspacer{} เอวญฺจ ปน ภิกฺขเว อิมํ สิกฺขาปทํ อุทฺทิเสยฺยาถ\csromandash{}}

\prose{\textbf{“โย ปน ภิกฺขุ อญฺญาติกาย ภิกฺขุนิยา จีวรํ ทเทยฺย, ปาจิตฺติยนฺ”}ติ.}

\prose{เอวญฺจิทํ ภควตา ภิกฺขูนํ สิกฺขาปทํ ปญฺญตฺตํ โหติ.}

\proseitem{170}{เตน โข ปน สมเยน ภิกฺขู กุกฺกุจฺจายนฺตา ภิกฺขุนีนํ ปาริวตฺตกํ\footnote{ปาริวฏฺฏกํ (อิติปิ)} จีวรํ น เทนฺติ.\csromanspacer{} ภิกฺขุนิโย อุชฺฌายฺยนฺติ ขิยฺยนฺติ วิปาเจนฺติ “กถํ หิ นาม อยฺยา อมฺหากํ ปริวตฺตกํ จิวรํ น ทสฺสนฺตี”ติ.\csromanspacer{} อสฺโสสุํ โข \csromanfolio{84}\textbf{ภิกฺขู ตาสํ ภิกฺขุนีนํ อุชฺฌายนฺตีนํ ขิยฺยนฺตีนํ วิปาเจนฺตีนํ.}\csromanspacer{} \textbf{อถ โข เต ภิกฺขู ภควโต เอตมตฺถํ อาโรเจสุํ.}\csromanspacer{}\textbf{ อถ โข ภควา} เอตสฺมิํ นิทาเน เอตสฺมิํ ปกรเณ ธมฺมิํ กถํ กตฺวา ภิกฺขู อามนฺเตสิ อนุชานามิ ภิกฺขเว ปญฺจนฺนํ ปาริวตฺตกํ ทาตุํ ภิกฺขุสฺส ภิกฺขุนิยา สิกฺขมานาย สามเณรสฺส สามเณริยา, อนุชานามิ ภิกฺขเว อิเมสํ ปญฺจนฺนํ ปาริวตฺตกํ ทาตุํ.\csromanspacer{} เอวญฺจ ปน ภิกฺขเว อิมํ สิกฺขาปทํ อุทฺทิเสยฺยาถ\csromandash{}}

\proseitem{171}{\textbf{“โย ปน ภิกฺขุ อญฺญาติกาย ภิกฺขุนิยา จิวรํ ทเทยฺย อญฺญตฺร ปาริวตฺตกา, ปาจิตฺติยนฺ”}ติ.}

\proseitem{172}{\textbf{โย ปนา}ติ โย ยาทิโส ฯเปฯ.\csromanspacer{} \textbf{ภิกฺขู}ติ ฯเปฯ อยํ อิมสฺมิํ อตฺเถ อธิปฺเปโต ภิกฺขูติ.}

\prose{\textbf{อญฺญาติกา} นาม มาติโต วา ปิติโต วา ยาว สตฺตมา ปิตามหยุคา อสมฺพทฺธา.}

\prose{\textbf{ภิกฺขุนี} นาม อุภโตสํเฆ อุปสมฺปนฺนา.}

\prose{\textbf{จีวรํ} นาม ฉนฺนํ จีวรานํ อญฺญตรํ จีวรํ วิกปฺปนุปคํ ปจฺฉิมํ\footnote{วิกปฺปนุปคปจฺฉิมํ (สี)}.}

\prose{\textbf{อญฺญตฺร ปาริวตฺตกา}ติ ฐเปตฺวา ปาริวตฺตกํ เทติ, อาปตฺติ ปาจิตฺติยสฺส.}

\proseitem{173}{อญฺญาติกาย อญฺญาติกสญฺญี จีวรํ เทติ อญฺญตฺร ปาริวตฺตกา, อาปตฺติ ปาจิตฺติยสฺส.\csromanspacer{} อญฺญาติกาย เวมติโก จิวรํ เทติ อญฺญตฺร ปาริวตฺตกา, อาปตฺติ ปาจิตฺติยสฺส.\csromanspacer{} อญฺญาติกาย ญาติกสญฺญี จิวรํ เทติ อญฺญตฺร ปาริวตฺตกา, อาปตฺติ}

\prose{เอกโต อุปสมฺปนฺนาย จีวรํ เทติ อญฺญตฺร ปาริวตฺตกา, อาปตฺติ ทุกฺกฏสฺส.\csromanspacer{} ญาติกาย อญฺญาติกสญฺญี, อาปตฺติ ทุกฺกฏสฺส.\csromanspacer{} ญาติกาย เวมติโก, อาปตฺติ ทุกฺกฏสฺส.\csromanspacer{} ญาติกาย ญาติกสญฺญี, อนาปตฺติ.}

\proseitemwithcloser{174}{อนาปตฺติ ญาติกาย, ปาริวตฺตกํ ปริตฺเตน วา วิปุลํ, วิปุเลน วา ปริตฺตํ, ภิกฺขุนี วิสฺสาสํ คณฺหาติ, ตาวกาลิกํ คณฺหาติ, จีวรํ ฐเปตฺวา อญฺญํ ปริกฺขารํ เทติ, สิกฺขมานาย สามเณริยา อุมฺมตฺตกสฺส อาทิกมฺมิกสฺสาติ.}{จีวรทานสิกฺขาปทํ นิฏฺฐิตํ ปญฺจมํ.}
\csromanmatikaanchor{mtk.30}
\csromantocmarkref{subsection}{6. จีวรสิพฺพนสิกฺขาปท}{mtk.30}
\bookmark[dest={mtk.30},level=2]{6. จีวรสิพฺพนสิกฺขาปท}
\csromanheader{2}{3. \textbf{โอวาทวคฺค 6. จีวรสิพฺพนสิกฺขาปท}}
\proseitem{175}{\csromanfolio{85}เตน สมเยน พุทฺโธ ภควา สาวตฺถิยํ วิหรติ เชตวเน อนาถปิณฺฑิกสฺส อาราเม.\csromanspacer{} เตน โข ปน สมเยน อายสฺมา อุทายี ปฏฺโฏ\footnote{ปฏฺโฐ (สี, สฺยา)} โหติ จีวรกมฺมํ กาตุํ.\csromanspacer{} อญฺญตรา ภิกฺขุนี เยนายสฺมา อุทายี เตนุปสงฺกมิ, อุปสงฺกมิตฺวา อายสฺมนฺตํ อุทายิํ เอตทโวจ “สาธุ เม ภนฺเต อยฺโย จีวรํ สิพฺพตู”ติ.\csromanspacer{} อถ โข อายสฺมา อุทายี ตสฺสา ภิกฺขุนิยา จีวรํ สิพฺพิตฺวา สุรตฺตํ สุปริกมฺมกตํ กตฺวา มชฺเฌ ปฏิภานจิตฺตํ วุฏฺฐาเปตฺวา สํหริตฺวา นิกฺขิปิ.\csromanspacer{} อถ โข สา ภิกฺขุนี เยนายสฺมา อุทายี เตนุปสงฺกมิ, อุปสงฺกมิตฺวา อายสฺมนฺตํ อุทายิํ เอตทโวจ “กหํ ตํ ภนฺเต จีวรนฺ”ติ.\csromanspacer{} หนฺท ภคินิ, อิมํ จีวรํ ยถาสํหฏํ หริตฺวา นิกฺขิปิตฺวา ยทา ภิกฺขุนิสํโฆ โอวาทํ อาคจฺฉติ, ตทา อิมํ จีวรํ ปารุปิตฺวา ภิกฺขุนิสํฆสฺส ปิฏฺฐิโต ปิฏฺฐิโต อาคจฺฉา”ติ.\csromanspacer{} อถ โข สา ภิกฺขุนี ตํ จีวรํ ยถาสํหฏํ หริตฺวา นิกฺขิปิตฺวา ยทา ภิกฺขุนิสํโฆ โอวาทํ อาคจฺฉติ, ตทา ตํ จิวรํ ปารุปิตฺวา ภิกฺขุนิสํฆสฺส ปิฏฺฐิโต ปิฏฺฐิโต อาคจฺฉติ.\csromanspacer{} มนุสฺสา อุชฺฌายนฺติ ขิยฺยนฺติ วิปาเจนฺติ “ยาว ฉินฺนิกา อิมา ภิกฺขุนิโย ธุตฺติกา อหิริกาโย, ยตฺร หิ นาม จีวเร ปฏิภานจิตฺตํ วุฏฺฐาเปสฺสนฺตี”ติ.}

\prose{ภิกฺขุนิโย เอวมาหํสุ “กสฺสิทํ กมฺมนฺ”ติ.\csromanspacer{} อยฺยสฺส อุทายิสฺสาติ.\csromanspacer{} เยปิ เต ฉินฺนกา ธุตฺตกา อหิริกา, เตสํปิ เอวรูปํ น โสเภยฺย, กิํ ปน อยฺยสฺส อุทายิสฺสาติ.\csromanspacer{} อถ โข ตา ภิกฺขุนิโย ภิกฺขูนํ เอตมตฺถํ อาโรเจสุํ.\csromanspacer{} เย เต ภิกฺขู อปฺปิจฺฉา ฯเปฯ เต อุชฺฌายนฺติ ขิยฺยนฺติ วิปาเจนฺติ “กถํ หิ นาม อายสฺมา อุทายี ภิกฺขุนิยา จีวรํ สิพฺพิสฺสตี”ติ ฯเปฯ.\csromanspacer{} สจฺจํ กิร ตฺวํ อุทายิ ภิกฺขุนิยา จีวรํ สิพฺพสีติ.\csromanspacer{} สจฺจํ ภควาติ.\csromanspacer{} ญาติกา เต อุทายิ อญฺญาติกาติ.\csromanspacer{} อญฺญาติกา ภควาติ.\csromanspacer{} อญฺญาตโก โมฆปุริส อญฺญาติกาย น ชานาติ ปติรูปํ วา อปฺปติรูปํ วา ปาสาทิกํ วา อปาสาทิกํ วา, กถํ หิ นาม ตฺวํ โมฆปุริส อญฺญาติกาย ภิกฺขุนิยา จีวรํ สิพฺพิสฺสสิ.\csromanspacer{} เนตํ โมฆปุริส อปฺปสนฺนานํ วา ปสาทาย ฯเปฯ.\csromanspacer{} เอวญฺจ ปน ภิกฺขเว อิมํ สิกฺขาปทํ อุทฺทิเสยฺยาถ\csromandash{}}

\proseitem{176}{\csromanfolio{86}\textbf{“โย ปน ภิกฺขุ อญฺญาติกาย ภิกฺขุนิยา จีวรํ สิพฺเพยฺย วา สิพฺพาเปยฺยวา, ปาจิตฺติยนฺ”}ติ.}

\proseitem{177}{\textbf{โย ปนา}ติ โย ยาทิโส ฯเปฯ.\csromanspacer{} \textbf{ภิกฺขู}ติ ฯเปฯ อยํ อิมสฺมิํ อตฺเถ อธิปฺเปโต ภิกฺขูติ.}

\prose{\textbf{อญฺญาติกา} นาม มาติโต วา ปิติโต วา ยาว สตฺตมา ปิตามหยุคา อสมฺพทฺธา.}

\prose{\textbf{ภิกฺขุนี} นาม อุภโตสํเฆ อุปสมฺปนฺนา.}

\prose{\textbf{จีวรํ} นาม ฉนฺนํ จีวรานํ อญฺญตรํ จีวรํ.}

\prose{\textbf{สิพฺเพยฺยา}ติ สยํ สิพฺพติ, อาราปเถ อาราปเถ อาปตฺติ ปาจิตฺติยสฺส.}

\prose{\textbf{สิพฺพาเปยฺยา}ติ อญฺญํ อาณาเปติ, อาปตฺติ ปาจิตฺติยสฺส.\csromanspacer{} สกิํ อาณตฺโต พหุกมฺปิ สิพฺพติ, อาปตฺติ ปาจิตฺติยสฺส.}

\proseitem{178}{อญฺญาติกาย อญฺญาติกสญฺญี จีวรํ สิพฺพติ วา สิพฺพาเปติ วา, อาปตฺติ ปาจิตฺติยสฺส.\csromanspacer{} อญฺญาติกาย เวมติโก จีวรํ สิพฺพติ วา สิพฺพาเปติ วา, อาปตฺติ ปาจิตฺติยสฺส.\csromanspacer{} อญฺญาติกาย ญาติกสญฺญี จิวรํ สิพฺพติ วา สิพฺพาเปติ วา, อาปตฺติ}

\prose{เอกโต อุปสมฺปนฺนาย จีวรํ สิพฺพติ วา สิพฺพาเปติ วา, อาปตฺติ ทุกฺกฏสฺส.\csromanspacer{} ญาติกาย อญฺญาติกสญฺญี, อาปตฺติ ทุกฺกฏสฺส.\csromanspacer{} ญาติกาย เวมติโก, อาปตฺติ ทุกฺกฏสฺส.\csromanspacer{} ญาติกาย ญาติกสญฺญี, อนาปตฺติ.}

\proseitemwithcloserruled{179}{อนาปตฺติ ญาติกาย, จีวรํ ฐเปตฺวา อญฺญํ ปริกฺขารํ สิพฺพติ วา สิพฺพาเปติ วา, สิกฺขมานาย สามเณริยา อุมฺมตฺตกสฺส อาทิกมฺมิกสฺสาติ.}{จีวรสิพฺพนสิกฺขาปทํ นิฏฺฐิตํ ฉฏฺฐํ.}
\csromanmatikaanchor{mtk.31}
\csromantocmarkref{subsection}{7. สํวิธานสิกฺขาปท}{mtk.31}
\bookmark[dest={mtk.31},level=2]{7. สํวิธานสิกฺขาปท}
\csromanheader{2}{3. \textbf{โอวาทวคฺค 7. สํวิธานสิกฺขาปท}}
\proseitem{180}{เตน สมเยน พุทฺโธ ภควา สาวตฺถิยํ วิหรติ เชตวเน อนาถปิณฺฑิกสฺส อาราเม.\csromanspacer{} เตน โข ปน สมเยน ฉพฺพคฺคิยา ภิกฺขู \csromanfolio{87}\textbf{ภิกฺขุนีหิ สทฺธิํ สํวิธาย เอกทฺธานมคฺคํ ปฏิปชฺชนฺติ.}\csromanspacer{}\textbf{ มนุสฺสา} \textbf{อุชฺฌายนฺติ ขิยฺยนฺติ วิปาเจนฺติ “ยเถว มยํ สปชาปติกา อาหิณฺฑาม,} \textbf{เอวเมวิเม สมณา สกฺยปุตฺติยา ภิกฺขุนีหิ สทฺธิํ สํวิธาย} อาหิณฺฑนฺตี”ติ.\csromanspacer{} อสฺโสสุํ โข ภิกฺขู เตสํ มนุสฺสานํ อุชฺฌายนฺตานํ ขิยฺยนฺตานํ วิปาเจนฺตานํ.\csromanspacer{} เย เต ภิกฺขู อปฺปิจฺฉา ฯเปฯ เต อุชฺฌายนฺติ ขิยฺยนฺติ วิปาเจนฺติ “กถํ หิ นาม ฉพฺพคฺคิยา ภิกฺขู ภิกฺขุนีหิ สทฺธิํ สํวิธาย เอกทฺธานมคฺคํ ปฏิปชฺชิสฺสนฺตี”ติ ฯเปฯ.\csromanspacer{} สจฺจํ กิร ตุเมฺห ภิกฺขเว ภิกฺขุนีหิ สทฺธิํ สํวิธาย เอกทฺธานมคฺคํ ปฏิปชฺชถาติ.\csromanspacer{} สจฺจํ ภควาติ.\csromanspacer{} วิครหิ พุทฺโธ ภควา ฯเปฯ.\csromanspacer{} กถํ หิ นาม ตุเมฺห โมฆปุริสา ภิกฺขุนีหิ สทฺธิํ สํวิธาย เอกทฺธานมคฺคํ ปฏิปชฺชิสฺสถ.\csromanspacer{} เนตํ โมฆปุริสา อปฺปสนฺนานํ วา ปสาทาย ฯเปฯ.\csromanspacer{} เอวญฺจ ปน ภิกฺขเว อิมํ สิกฺขาปทํ อุทฺทิเสยฺยาถ\csromandash{}}

\prose{\textbf{“โย ปน ภิกฺขุ ภิกฺขุนิยา สทฺธิํ สํวิธาย เอกทฺธานมคฺคํ ปฏิปชฺเชยฺย อนฺตมโส คามนฺตรมฺปิ, ปาจิตฺติยนฺ”}ติ.}

\prose{เอวญฺจิทํ ภควตา ภิกฺขูนํ สิกฺขาปทํ ปญฺญตฺตํ โหติ.}

\proseitem{181}{เตน โข ปน สมเยน สมฺพหุลา ภิกฺขู จ ภิกฺขุนิโย จ สาเกตา สาวตฺถิํ อทฺธานมคฺคปฺปฏิปนฺนา โหนฺติ.\csromanspacer{} อถ โข ตา ภิกฺขุนิโย เต ภิกฺขู เอตทโวจุํ “มยมฺปิ อยฺเยหิ สทฺธิํ คมิสฺสามา”ติ.\csromanspacer{} น ภคินี กปฺปติ ภิกฺขุนิยา สทฺธิํ สํวิธาย เอกทฺธานมคฺคํ ปฏิปชฺชิตุํ.\csromanspacer{} ตุเมฺห วา ปฐมํ คจฺฉถ, มยํ วา คมิสฺสามาติ.\csromanspacer{} อยฺยา ภนฺเต อคฺคปุริสา อยฺยาว ปฐมํ คจฺฉนฺตูติ.\csromanspacer{} อถ โข ตาสํ ภิกฺขุนีนํ ปจฺฉา คจฺฉนฺตีนํ อนฺตรามคฺเค โจรา อจฺฉินฺทิํสุ จ ทูเสสุํ จ.\csromanspacer{} อถ โข ตา ภิกฺขุนิโย สาวตฺถิํ คนฺตฺวา ภิกฺขุนีนํ เอตมตฺถํ อาโรเจสุํ.\csromanspacer{} ภิกฺขุนิโย ภิกฺขูนํ เอตมตฺถํ อาโรเจสุํ.\csromanspacer{} ภิกฺขู ภควโต เอตมตฺถํ อาโรเจสุํ.\csromanspacer{} อถ โข ภควา เอตสฺมิํ นิทาเน เอตสฺมิํ ปกรเณ ธมฺมิํ กถํ กตฺวา ภิกฺขู อามนฺเตสิ อนุชานามิ ภิกฺขเว สตฺถคมนีเย มคฺเค สาสงฺกสมฺมเต สปฺปฏิภเย ภิกฺขุนิยา สทฺธิํ สํวิธาย เอกทฺธานมคฺคํ ปฏิปชฺชิตุํ.\csromanspacer{} เอวญฺจ ปน ภิกฺขเว อิมํ สิกฺขาปทํ อุทฺทิเสยฺยาถ\csromandash{}}

\proseitem{182}{“\textbf{โย ปน ภิกฺขุ ภิกฺขุนิยา สทฺธิํ สํวิธาย} \textbf{เอกทฺธานมคฺคํ ปฏิปชฺเชยฺย อนฺตมโส คามนฺตรฺ}อมฺปิ อญฺญตฺร สมยา, ปาจิตฺติยํ.\csromanspacer{} ตตฺถายํ \csromanfolio{88}\textbf{สมโย, สตฺถคมนีโย โหติ มคฺโค สาสงฺกสมฺมโต สปฺปฏิภโย, อยํ ตตฺถ สมโย”}ติ.}

\proseitem{183}{\textbf{โย ปนา}ติ โย ยาทิโส ฯเปฯ.\csromanspacer{} \textbf{ภิกฺขู}ติ ฯเปฯ อยํ อิมสฺมิํ อตฺเถ อธิปฺเปโต ภิกฺขูติ.}

\prose{\textbf{ภิกฺขุนี} นาม อุภโตสํเฆ อุปสมฺปนฺนา.}

\prose{\textbf{สทฺธินฺ}ติ เอกโต.}

\prose{\textbf{สํวิธายา}ติ “คจฺฉาม ภคินิ คจฺฉามายฺย, คจฺฉามายฺย คจฺฉาม ภคินิ, อชฺช วา หิยฺโย วา ปเร วา คจฺฉามา”ติ สํวิทหติ, อาปตฺติ ทุกฺกฏสฺส.}

\prose{\textbf{อนฺตมโส คามนฺตรมฺปี}ติ กุกฺกุฏสมฺปาเต คาเม คามนฺตเร คามนฺตเร อาปตฺติ ปาจิตฺติยสฺส.\csromanspacer{} อคามเก อรญฺเญ อทฺธโยชเน อทฺธโยชเน อาปตฺติ}

\prose{\textbf{อญฺญตฺร สมยา}ติ ฐเปตฺวา สมยํ.}

\prose{\textbf{สตฺถคมนีโย} นาม มคฺโค น สกฺกา โหติ วินา สตฺเถน คนฺตุํ.}

\prose{\textbf{สาสงฺกํ} นาม ตสฺมิํ มคฺเค\footnote{ยสฺมิํ มคฺเค (?)} โจรานํ นิวิฏฺโฐกาโส ทิสฺสติ, ภุตฺโตกาโส ทิสฺสติ, ฐิโตกาโส ทิสฺสติ, นิสินฺโนกาโส ทิสฺสติ, นิปนฺโนกาโส ทิสฺสติ.}

\prose{\textbf{สปฺปฏิภยํ} นาม ตสฺมิํ มคฺเค โจเรหิ มนุสฺสา หตา ทิสฺสนฺติ, วิลุตฺตา ทิสฺสนฺติ, อาโกฏิตา ทิสฺสนฺติ.\csromanspacer{} สปฺปฏิภยํ คนฺตฺวา อปฺปฏิภยํ ทสฺเสตฺวา อุยฺโยเชตพฺพา “คจฺฉถ ภคินิโย”ติ.}

\proseitem{184}{สํวิทหิเต สํวิทหิตสญฺญี เอกทฺธานมคฺคํ ปฏิปชฺชติ อนฺตมโส คามนฺตรมฺปิ อญฺญตฺร สมยา, อาปตฺติ ปาจิตฺติยสฺส.\csromanspacer{} สํวิทหิเต เวมติโก เอกทฺธานมคฺคํ ปฏิปชฺชติ อนฺตมโส คามนฺตรมฺปิ อญฺญตฺร สมยา, อาปตฺติ ปาจิตฺติยสฺส.\csromanspacer{} สํวิทหิเต อสํวิทหิตสญฺญี เอกทฺธานมคฺคํ ปฏิปชฺชติ อนฺตมโส คามนฺตรมฺปิ อญฺญตฺร สมยา, อาปตฺติ ปาจิตฺติยสฺส.}

\prose{ภิกฺขุ สํวิทหติ ภิกฺขุนี น สํวิทหติ, อาปตฺติ ทุกฺกฏสฺส.\csromanspacer{} อสํวิทหิเต สํวิทหิตสญฺญี, อาปตฺติ ทุกฺกฏสฺส.\csromanspacer{} อสํวิทหิเต เวมติโก อาปตฺติ ทุกฺกฏสฺส.\csromanspacer{} อสํวิทหิเต อสํวิทหิตสญฺญี, อนาปตฺติ.}

\proseitemwithcloserruled{185}{\csromanfolio{89}อนาปตฺติ สมเย, อสํวิทหิตฺวา คจฺฉติ, ภิกฺขุนี สํวิทหติ ภิกฺขุ น สํวิทหติ, วิสงฺเกเตน คจฺฉนฺติ, อาปทาสุ อุมฺมตฺตกสฺส อาทิกมฺมิกสฺสาติ.}{สํวิธานสิกฺขาปทํ นิฏฺฐิตํ สตฺตมํ.}
\csromanmatikaanchor{mtk.32}
\csromantocmarkref{subsection}{8. นาวาภิรุหนสิกฺขาปท}{mtk.32}
\bookmark[dest={mtk.32},level=2]{8. นาวาภิรุหนสิกฺขาปท}
\csromanheader{2}{3. \textbf{โอวาทวคฺค 8. นาวาภิรุหนสิกฺขาปท}}
\proseitem{186}{เตน สมเยน พุทฺโธ ภควา สาวตฺถิยํ วิหรติ เชตวเน อนาถปิณฺฑิกสฺส อาราเม.\csromanspacer{} เตน โข ปน สมเยน ฉพฺพคฺคิยา ภิกฺขู ภิกฺขุนีหิ สทฺธิํ สํวิธาย เอกํ นาวํ\footnote{เอกนาวํ (สฺยา)} อภิรุหนฺติ.\csromanspacer{} มนุสฺสา อุชฺฌายนฺติ ขิยฺยนฺติ วิปาเจนฺติ “ยเถว มยํ สปชาปติกา นาวาย\footnote{เอกนาวาย (สฺยา)} กีฬาม, เอวเมวิเม สมณา สกฺยปุตฺติยา ภิกฺขุนีหิ สทฺธิํ สํวิธาย นาวาย กีฬนฺตี”ติ.\csromanspacer{} อสฺโสสุํ โข ภิกฺขู เตสํ มนุสฺสานํ อุชฺฌายนฺตานํ ขิยฺยนฺตานํ วิปาเจนฺตานํ.\csromanspacer{} เย เต ภิกฺขู อปฺปิจฺฉา ฯเปฯ เต อุชฺฌายนฺติ ขิยฺยนฺติ วิปาเจนฺติ “กถํ หิ นาม ฉพฺพคฺคิยา ภิกฺขู ภิกฺขุนีหิ สทฺธิํ สํวิธาย เอกํ นาวํ อภิรุหิสฺสนฺตี”ติ ฯเปฯ.\csromanspacer{} สจฺจํ กิร ตุเมฺห ภิกฺขเว ภิกฺขุนีหิ สทฺธิํ สํวิธาย เอกํ นาวํ อภิรุหถาติ.\csromanspacer{} สจฺจํ ภควาติ.\csromanspacer{} วิครหิ พุทฺโธ ภควา ฯเปฯ.\csromanspacer{} กถํ หิ นาม ตุเมฺห โมฆปุริสา ภิกฺขุนีหิ สทฺธิํ สํวิธาย เอกํ นาวํ อภิรุหิสฺสถ.\csromanspacer{} เนตํ โมฆปุริสา อปฺปสนฺนานํ วา ปสาทาย ฯเปฯ.\csromanspacer{} เอวญฺจ ปน ภิกฺขเว อิมํ สิกฺขาปทํ อุทฺทิเสยฺยาถ\csromandash{}}

\prose{\textbf{“โย ปน ภิกฺขุ ภิกฺขุนิยา สทฺธิํ สํวิธาย เอกํ นาวํ1} \textbf{อภิรุเหยฺย อุทฺธํคามินิํ วา อโธคามินิํ วา, ปาจิตฺติยนฺ”}ติ.}

\prose{เอวญฺจิทํ ภควตา ภิกฺขูนํ สิกฺขาปทํ ปญฺญตฺตํ โหติ.}

\proseitem{187}{เตน โข ปน สมเยน สมฺพหุลา ภิกฺขู จ ภิกฺขุนิโย จ สาเกตา สาวตฺถิํ อทฺธานมคฺคปฺปฏิปนฺนา โหนฺติ.\csromanspacer{} อนฺตรามคฺเค นที ตริตพฺพา\footnote{อุตฺตริตพฺพา (สฺยา)} โหติ.\csromanspacer{} อถ โข ตา ภิกฺขุนิโย เต ภิกฺขู เอตทโวจุํ “มยมฺปิ อยฺเยหิ สทฺธิํ อุตฺตริสฺสามา”ติ.\csromanspacer{} น ภคินี กปฺปติ ภิกฺขุนิยา สทฺธิํ สํวิธาย เอกํ นาวํ อภิรุหิตุํ.\csromanspacer{} ตุเมฺห วา ปฐมํ อุตฺตรถ, มยํ วา \csromanfolio{90}อุตฺตริสฺสามาติ.\csromanspacer{} อยฺยา ภนฺเต อคฺคปุริสา, อยฺยาว ปฐมํ อุตฺตรนฺตูติ.\csromanspacer{} อถ โข ตาสํ ภิกฺขุนีนํ ปจฺฉา อุตฺตรนฺตีนํ โจรา อจฺฉินฺทิํสุ จ ทูเสสุํ จ.\csromanspacer{} อถ โข ตา ภิกฺขุนิโย สาวตฺถิํ คนฺตฺวา ภิกฺขุนีนํ เอตมตฺถํ อาโรเจสุํ.\csromanspacer{} ภิกฺขุนิโย ภิกฺขูนํ เอตมตฺถํ อาโรเจสุํ.\csromanspacer{} ภิกฺขู ภควโต เอตมตฺถํ อาโรเจสุํ.\csromanspacer{} อถ โข ภควา เอตสฺมิํ นิทาเน เอตสฺมิํ ปกรเณ ธมฺมิํ กถํ กตฺวา ภิกฺขู อามนฺเตสิ อนุชานามิ ภิกฺขเว ติริยํ ตรณาย ภิกฺขุนิยา สทฺธิํ สํวิธาย เอกํ นาวํ อภิรุหิตุํ.\csromanspacer{} เอวญฺจ ปน ภิกฺขเว อิมํ สิกฺขาปทํ อุทฺทิเสยฺยาถ\csromandash{}}

\proseitem{188}{\textbf{“โย ปน ภิกฺขุ ภิกฺขุนิยา สทฺธิํ สํวิธาย เอกํ นาวํ อภิรุเหยฺย อุทฺธํคามินิํ วา อโธคามินิํ วา อญฺญตฺร ติริยํ ตรณาย, ปาจิตฺติยนฺ”}ติ.}

\proseitem{189}{โย ปนาติ โย ยาทิโส ฯเปฯ.\csromanspacer{} \textbf{ภิกฺขู}ติ ฯเปฯ อยํ อิมสฺมิํ อตฺเถ อธิปฺเปโต ภิกฺขูติ.}

\prose{\textbf{ภิกฺขุนี} นาม อุภโตสํเฆ อุปสมฺปนฺนา.}

\prose{\textbf{สทฺธินฺ}ติ เอกโต.}

\prose{\textbf{สํวิธายา}ติ “อภิรุหาม ภคินิ อภิรุหามายฺย, อภิรุหามายฺย อภิรุหาม ภคินิ, อชฺช วา หิยฺโย วา ปุเร วา อภิรุหามา”ติ สํวิทหติ, อาปตฺติ ทุกฺกฏสฺส.}

\prose{ภิกฺขุนิยา อภิรุเฬฺห ภิกฺขุ อภิรุหติ, อาปตฺติ ปาจิตฺติยสฺส.\csromanspacer{} ภิกฺขุมฺหิ อภิรุเฬฺห ภิกฺขุนี อภิรุหติ, อาปตฺติ ปาจิตฺติยสฺส.\csromanspacer{} อุโภ วา อภิรุหนฺติ, อาปตฺติ ปาจิตฺติยสฺส.}

\prose{\textbf{อุทฺธํคามินินฺ}ติ อุชฺชวนิกาย.}

\prose{\textbf{อโธคามินินฺ}ติ โอชวนิกาย.}

\prose{\textbf{อญฺญตฺร ติริยํ ตรณายา}ติ ฐเปตฺวา ติริยํ ตรณํ.}

\prose{กุกฺกุฏสมฺปาเต คาเม คามนฺตเร คามนฺตเร อาปตฺติ ปาจิตฺติยสฺส.\csromanspacer{} อคามเก อรญฺเญ อฑฺฒโยชเน อฑฺฒโยชเน อาปตฺติ ปาจิตฺติยสฺส.}

\proseitem{190}{สํวิทหิเต สํวิทหิตสญฺญี เอกํ นาวํ อภิรุหติ อุทฺธํคามินิํ วา อโธคามินิํ วา อญฺญตฺร ติริยํ ตรณาย, อาปตฺติ \csromanfolio{91}สํวิทหิเต เวมติโก เอกํ นาวํ อภิรุหติ อุทฺธํคามินิํ วา อโธ คามินิํ วา อญฺญตฺร ติริยํ ตรณาย, อาปตฺติ ปาจิตฺติยสฺส.\csromanspacer{} สํวิทหิเต อสํวิทหิตสญฺญี เอกํ นาวํ อภิรุหติ อุทฺธํคามินิํ วา อโธคามินิํ วา อญฺญตฺร ติริยํ ตรณาย, อาปตฺติ ปาจิตฺติยสฺส.}

\prose{ภิกฺขุ สํวิทหติ ภิกฺขุนี น สํวิทหติ, อาปตฺติ ทุกฺกฏสฺส.\csromanspacer{} อสํวิทหิเต สํวิทหิตสญฺญี, อาปตฺติ ทุกฺกฏสฺส.\csromanspacer{} อสํวิทหิเต เวมติโก, อาปตฺติ ทุกฺกฏสฺส.\csromanspacer{} อสํวิทหิเต อสํวิทหิตสญฺญี, อนาปตฺติ.}

\proseitemwithcloserruled{191}{อนาปตฺติ ติริยํ ตรณาย, อสํวิทหิตฺวา อภิรุหนฺติ, ภิกฺขุนี สํวิทหติ ภิกฺขุ น สํวิทหติ, วิสงฺเกเตน อภิรุหนฺติ, อาปทาสุ อุมฺมตฺตกสฺส อาทิกมฺมิสฺสาติ.}{นาวาภิรุหนสิกฺขาปทํ นิฏฺฐิตํ อฏฺฐมํ.}
\csromanmatikaanchor{mtk.33}
\csromantocmarkref{subsection}{9. ปริปาจิตสิกฺขาปท}{mtk.33}
\bookmark[dest={mtk.33},level=2]{9. ปริปาจิตสิกฺขาปท}
\csromanheader{2}{3. \textbf{โอวาทวคฺค 9. ปริปาจิตสิกฺขาปท}}
\proseitem{192}{เตน สมเยน พุทฺโธ ภควา ราชคเห วิหรติ เวฬุวเน กลนฺทกนิวาเป.\csromanspacer{} เตน โข ปน สมเยน ถุลฺลนนฺทา ภิกฺขุนี อญฺญตรสฺส กุลสฺส กุลูปิกา โหติ นิจฺจภตฺติกา, เตน จ คหปตินา เถรา ภิกฺขู นิมนฺติตา โหนฺติ.\csromanspacer{} อถ โข ถุลฺลนนฺทา ภิกฺขุนี ปุพฺพณฺหสมยํ นิวาเสตฺวา ปตฺตจีวรมาทาย เยน ตํ กุลํ เตนุปสงฺกมิ, อุปสงฺกมิตฺวา ตํ คหปติํ เอตทโวจ “กิมิทํ คหปติ ปหูตํ ขาทนียํ โภชนียํ ปฏิยตฺตนฺ”ติ.\csromanspacer{} เถรา มยา อยฺเย นิมนฺติตาติ.\csromanspacer{} เก ปน เต คหปติ เถราติ.\csromanspacer{} อยฺโย สาริปุตฺโต อยฺโย มหาโมคฺคลฺลาโน อยฺโย มหากจฺจาโน อยฺโย มหาโกฏฺฐิโก อยฺโย มหากปฺปิโน อยฺโย มหาจุนฺโท อยฺโย อนุรุทฺโธ อยฺโย เรวโต อยฺโย อุปาลิ อยฺโย อานนฺโท อยฺโย ราหุโลติ.\csromanspacer{} กิํ ปน ตฺวํ คหปติ มหานาเค ติฏฺฐมาเน เจฏเก นิมนฺเตสีติ.}

\prose{เก ปน เต อยฺเย มหานาคาติ.\csromanspacer{} อยฺโย เทวทตฺโต อยฺโย โกกาลิโก อยฺโย กฏโมทกติสฺสโก\footnote{กฏโมรกติสฺสโก (สี) กตโมรกติสฺสโก (สฺยา)} อยฺโย ขณฺฑเทวิยา \csromanfolio{92}ปุตฺโต อยฺโย สมุทฺททตฺโตติ.\csromanspacer{} อยํ จรหิ ถุลฺลนนฺทาย ภิกฺขุนิยา อนฺตรากถา วิปฺปกตา, อถ เต เถรา ภิกฺขู ปวิสิํสุ.\csromanspacer{} สจฺจํ มหานาคา โข ตยา คหปติ นิมนฺติตาติ. “อิทาเนว โข ตฺวํ อยฺเย เจฏเก อกาสิ, อิทานิ มหานาเค”ติ ฆรโต จ นิกฺกฑฺฒิ, นิจฺจภตฺตญฺจ ปจฺฉินฺทิ.\csromanspacer{} เย เต ภิกฺขู อปฺปิจฺฉา ฯเปฯ เต อุชฺฌายนฺติ ขิยฺยนฺติ วิปาเจนฺติ “กถํ หิ นาม เทวทตฺโต ชานํ ภิกฺขุนิปริปาจิตํ ปิณฺฑปาตํ ภุญฺชิสฺสตี”ติ ฯเปฯ “สจฺจํ กิร ตฺวํ เทวทตฺต ชานํ ภิกฺขุนิปริปาจิตํ ปิณฺฑปาตํ ภุญฺชสี”ติ.\csromanspacer{} สจฺจํ ภควาติ.\csromanspacer{} วิครหิ พุทฺโธ ภควา ฯเปฯ.\csromanspacer{} กถํ หิ นาม ตฺวํ โมฆปุริส ชานํ ภิกฺขุนิปริปาจิตํ ปิณฺฑปาตํ ภุญฺชิสฺสสิ.\csromanspacer{} เนตํ โมฆปุริส อปฺปสนฺนานํ วา ปสาทาย ฯเปฯ.\csromanspacer{} เอวญฺจ ปน ภิกฺขเว อิมํ \textbf{สิกฺขาปทํ อุทฺทิเสยฺยาถ\csromandash{}}}

\prose{\textbf{“โย ปน ภิกฺขุ ชานํ ภิกฺขุนิปริปาจิตํ ปิณฺฑปาตํ ภุญฺเชยฺย, ปาจิตฺติยนฺ”}ติ.}

\prose{เอวญฺจิทํ ภควตา ภิกฺขูนํ สิกฺขาปทํ ปญฺญตฺตํ โหติ.}

\proseitem{193}{เตน โข ปน สมเยน อญฺญตโร ภิกฺขุ ราชคหา ปพฺพชิโต ญาติกุลํ อคมาสิ.\csromanspacer{} มนุสฺสา “จิรสฺสมฺปิ ภทนฺโต อาคโต”ติ สกฺกจฺจํ ภตฺตํ อกํสุ.\csromanspacer{} ตสฺส กุลสฺส กุลูปิกา ภิกฺขุนี เต มนุสฺเส เอตทโวจ “เทถยฺยสฺส อาวุโส ภตฺตนฺ”ติ.\csromanspacer{} อถ โข โส ภิกฺขุ “ภควตา ปฏิกฺขิตฺตํ ชานํ ภิกฺขุนิปริปาจิตํ ปิณฺฑปาตํ ภุญฺชิตุนฺ”ติ กุกฺกุจฺจายนฺโต น ปฏิคฺคเหสิ, นาสกฺขิ ปิณฺฑาย จริตุํ, ฉินฺนภตฺโต อโหสิ.\csromanspacer{} อถ โข โส ภิกฺขุ อารามํ คนฺตฺวา ภิกฺขูนํ เอตมตฺถํ อาโรเจสิ.\csromanspacer{} ภิกฺขู ภควโต เอตมตฺถํ อาโรเจสุํ.\csromanspacer{} อถ โข ภควา เอตสฺมิํ นิทาเน เอตสฺมิํ ปกรเณ ธมฺมิํ กถํ กตฺวา ภิกฺขู อามนฺเตสิ อนุชานามิ ภิกฺขเว ปุพฺเพ คิหิสมารมฺเภ ชานํ ภิกฺขุนิปริจิตํ ปิณฺฑปาตํ ภุญฺชิตุํ.\csromanspacer{} เอวญฺจ ปน ภิกฺขเว อิมํ สิกฺขาปทํ อุทฺทิเสยฺยาถ\csromandash{}}

\proseitem{194}{“\textbf{โย ปน ภิกฺขุ ชานํ ภิกฺขุนิปริปาจิตํ ปฺ}อิณฺฑปาตํ ภุญฺเชยฺย อญฺญตฺร ปุพฺเพ คิหิสมารมฺภา, ปาริจิตฺติยนฺ”ติ.}

\proseitem{195}{\textbf{โย ปน}ติ โย ยาทิโส ฯเปฯ.\csromanspacer{} \textbf{ภิกฺขู}ติ ฯเปฯ อยํ อิมสฺมิํ อตฺเถ อธิปฺเปโต ภิกฺขูติ.}

\prose{\csromanfolio{93}\textbf{ชานาติ} นาม สามํ วา ชานาติ, อญฺเญ วา ตสฺส อาโรเจนฺติ, สา วา อาโรเจติ.}

\prose{\textbf{ภิกฺขุนี} นาม อุภโตสํเฆ อุปสมฺปนฺนา.}

\prose{\textbf{ปริปาเจติ} นาม ปุพฺเพ อทาตุกามานํ อกตฺตุกามานํ อยฺโย ภาณโก อยฺโย พหุสฺสุโต อยฺโย สุตฺตนฺติโก อยฺโย วินยธโร อยฺโย ธมฺมกถิโก, เทถ อยฺยสฺส, กโรถ อยฺยสฺสาติ.\csromanspacer{} เอสา ปริปาเจติ นาม.}

\prose{\textbf{ปิณฺฑปาโต} นาม ปญฺจนฺนํ โภชนานํ อญฺญตรํ โภชนํ.}

\prose{\textbf{อญฺญตฺร ปุพฺเพ คิหิสมารมฺภา}ติ ฐเปตฺวา คิหิสมารมฺภํ.}

\prose{\textbf{คิหิสมารมฺโภ} นาม ญาตกา วา โหนฺติ ปวาริตา วา ปกติปฏิยตฺตํ วา.}

\prose{\textbf{อญฺญตฺร ปุพฺเพ คิหิสมารมฺภา} ภุญฺชิสฺสามีติ ปฏิคฺคณฺหาติ, อาปตฺติ ทุกฺกฏสฺส.\csromanspacer{} อชฺโฌหาเร อชฺโฌหาเร อาปตฺติ ปาจิตฺติยสฺส.}

\proseitem{196}{ปริปาจิเต ปริปาจิตสญฺญี ภุญฺชติ อญฺญตฺร ปุพฺเพ คิหิสมารมฺภา, อาปตฺติ ปาจิตฺติยสฺส.\csromanspacer{} ปริปาจิเต เวมติโก ภุญฺชติ อญฺญตฺร ปุพฺเพ คิหิสมารมฺภา, อาปตฺติ ทุกฺกฏสฺส.\csromanspacer{} ปริปาจิเต อปริปาจิตสญฺญี ภุญฺชติ อญฺญตฺร ปุพฺเพ คิหิสมารมฺภา, อนาปตฺติ.\csromanspacer{} เอกโตอุปสมฺปนฺนาย ปริปาจิตํ ภุญฺชติ อญฺญตฺร ปุพฺเพ คิหิสมารมฺภา, อาปตฺติ ทุกฺกฏสฺส.\csromanspacer{} อปริปาจิเต ปริปาจิตสญฺญี, อาปตฺติ ทุกฺกฏสฺส.\csromanspacer{} อปริปาจิเต เวมติโก, อาปตฺติ ทุกฺกฏสฺส.\csromanspacer{} อปริปาจิเต อปริปาจิตสญฺญี, อนาปตฺติ.}

\proseitemwithcloserruled{197}{อนาปตฺติ ปุพฺเพ คิหิสมารมฺเภ, สิกฺขมานา ปริปาเจติ, สามเณรี ปริปาเจติ, ปญฺจ โภชนานิ ฐเปตฺวา สพฺพตฺถ, อนาปตฺติ อุมฺมตฺตกสฺส อาทิกมฺมิกสฺสาติ.}{ปริปาจิตสิกฺขาปทํ นิฏฺฐิตํ นวมํ.}
\csromanmatikaanchor{mtk.34}
\csromantocmarkref{subsection}{10. รโหนิสชฺชสิกฺขาปท}{mtk.34}
\bookmark[dest={mtk.34},level=2]{10. รโหนิสชฺชสิกฺขาปท}
\csromanheader{2}{3. โอวาทวคฺค 10. รโหนิสชฺชสิกฺขาปท}
\proseitem{198}{เตน สมเยน พุทฺโธ ภควา สาวตฺถิยํ วิหรติ เชตวเน อนาถปิณฺฑิกสฺส อาราเม.\csromanspacer{} เตน โข ปน สมเยน อายสฺมโต \csromanfolio{94}อุทายิสฺส ปุราณทุติยิกา ภิกฺขุนีสุ ปพฺพชิตา โหติ.\csromanspacer{} สา อายสฺมโต อุทายิสฺส สนฺติเก อภิกฺขณํ อาคจฺฉติ, อายสฺมาปิ อุทายี ตสฺสา ภิกฺขุนิยา สนฺติเก อภิกฺขณํ คจฺฉติ.\csromanspacer{} เตน โข ปน สมเยน อายสฺมา อุทายี ตสฺสา ภิกฺขุนิยา สทฺธิํ เอโก เอกาย รโห นิสชฺชํ กปฺเปสิ.\csromanspacer{} เย เต ภิกฺขู อปฺปิจฺฉา ฯเปฯ เต อุชฺฌายนฺติ ขิยฺยนฺติ วิปาเจนฺติ “กถํ หิ นาม อายสฺมา อุทายี ภิกฺขุนิยา สทฺธิํ เอโก เอกาย รโห นิสชฺชํ กปฺเปสฺสตี”ติ ฯเปฯ.\csromanspacer{} สจฺจํ กิร ตฺวํ อุทายิ ภิกฺขุนิยา สทฺธิํ เอโก เอกาย รโห นิสชฺชํ กปฺเปสีติ.\csromanspacer{} สจฺจํ ภควาติ.\csromanspacer{} วิครหิ พุทฺโธ ภควา ฯเปฯ.\csromanspacer{} กถํ หิ นาม ตฺวํ โมฆปุริส ภิกฺขุนิยา สทฺธิํ เอโก เอกาย รโห นิสชฺชํ กปฺเปสฺสสิ.\csromanspacer{} เนตํ โมฆปุริส อปฺปสนฺนานํ วา ปสาทาย ฯเปฯ.\csromanspacer{} เอวญฺจ ปน ภิกฺขเว อิมํ สิกฺขาปทํ อุทฺทิเสยฺยาถ\csromandash{}}

\proseitem{199}{\textbf{“โย ปน ภิกฺขุ ภิกฺขุนิยา สทฺธิํ เอโก เอกาย รโห นิสชฺชํ กปฺเปยฺย, ปาจิตฺติยนฺ”}ติ.}

\proseitem{200}{\textbf{โย ปนา}ติ โย ยาทิโส ฯเปฯ.\csromanspacer{} \textbf{ภิกฺขู}ติ ฯเปฯ อยํ อิมสฺมิํ อตฺเถ อธิปฺเปโต ภิกฺขูติ.}

\prose{\textbf{ภิกฺขุนี} นาม อุภโตสํเฆ อุปสมฺปนฺนา.}

\prose{\textbf{สทฺธินฺ}ติ เอกโต.}

\prose{\textbf{เอโก เอกายา}ติ ภิกฺขุ เจว โหติ ภิกฺขุนี จ.}

\prose{\csromansymbolmark{*}\textbf{รโห} นาม จกฺขุสฺส รโห โสตสฺส รโห.\csromanspacer{} \textbf{จกฺขุสฺส รโห} นาม น สกฺกา โหติ อกฺขิํ วา นิขณียมาเน ภมุกํ วา อุกฺขิปียมาเน สีสํ วา อุกฺขิปียมาเน ปสฺสิตุํ.\csromanspacer{} \textbf{โสตสฺส รโห} นาม น สกฺกา โหติ ปกติกถา โสตุํ.}

\prose{\textbf{นิสชฺชํ กปฺเปยฺยา}ติ ภิกฺขุนิยา นิสินฺนาย ภิกฺขุ อุปนิสินฺโน วา โหติ อุปนิปนฺโน วา, อาปตฺติ ปาจิตฺติยสฺส.}

\prose{ภิกฺขุ นิสินฺเน ภิกฺขุนี อุปนิสินฺนา วา โหติ อุปนิปนฺนา วา, อาปตฺติ ปาจิตฺติยสฺส.\csromanspacer{} อุโภ วา นิสินฺนา โหนฺติ อุโภ วา นิปนฺนา, อาปตฺติ ปาจิตฺติยสฺส.}

\proseitem{201}{\csromanfolio{95}รโห รโหสญฺญี เอโก เอกาย นิสชฺชํ กปฺเปติ, อาปตฺติ ปาจิตฺติยสฺส.\csromanspacer{} รโห เวมติโก เอโก เอกาย นิสชฺชํ กปฺเปติ, อาปตฺติ ปาจิตฺติยสฺส.\csromanspacer{} รโห อรโหสญฺญี เอโก เอกาย นิสชฺชํ กปฺเปติ, อาปตฺติ ปาจิตฺติยสฺส.}

\prose{อรโห รโหสญฺญี, อาปตฺติ ทุกฺกฏสฺส.\csromanspacer{} อรโห เวมติโก, อาปตฺติ ทุกฺกฏสฺส.\csromanspacer{} อรโห อรโหสญฺญี, อนาปตฺติ.}

\proseitemwithcloser{202}{อนาปตฺติ โย โกจิ วิญฺญู ทุติโย โหติ, ติฏฺฐติ น นิสีทติ, อรโหเปกฺโข, อญฺญวิหิโต นิสีทติ, อุมฺมตฺตกสฺส อาทิกมฺมิกสฺสาติ.}{รโหนิสชฺชสิกฺขาปทํ นิฏฺฐิตํ ทสมํ.}
\vspace{\tipitakaparskip}
\csromancenter{โอวาทวคฺโค ตติโย.}
\csromancenter{ตสฺสุทฺทานํ}
\vspace{0.5\baselineskip}
\setlength{\csromangathapagewidth}{0pt}
\csromangathameasuregroup{อสมฺมต อตฺถงฺคตู,}{\csromangathabat{อสมฺมต อตฺถงฺคตู,}{ปสฺสยามิสทาเนน.} \\ สิพฺพติ อทฺธานํ นาวํ ภุญฺเชยฺย, เอโก เอกาย เต ทสาติ.}
\csromangathasetleft{อสมฺมต อตฺถงฺคตู,}
\csromangathagroup{\csromangathastanza{\csromangathabat{อสมฺมต อตฺถงฺคตู,}{ปสฺสยามิสทาเนน.} \\ สิพฺพติ อทฺธานํ นาวํ ภุญฺเชยฺย, เอโก เอกาย เต ทสาติ.}}

\csromanmatikaanchor{mtk.35}
\csromanmatikaanchor{mtk.36}
\csromantocmarknopage{section}{4. โภชนวคฺค}{mtk.35}
\bookmark[dest={mtk.35},level=1]{4. โภชนวคฺค}
\csromantocmarkref{subsection}{1. อาวาสถปิณฺฑสิกฺขาปท}{mtk.36}
\bookmark[dest={mtk.36},level=2]{1. อาวาสถปิณฺฑสิกฺขาปท}
\setcsromanheadh{4. โภชนวคฺค}
\csromanheader{2}{4. \textbf{โภชนวคฺค 1. อาวสถปิณฺฑสิกฺขาปท}}
\proseitem{203}{เตน สมเยน พุทฺโธ ภควา สาวตฺถิยํ วิหรติ เชตวเน อนาถปิณฺฑิกสฺส อาราเม.\csromanspacer{} เตน โข ปน สมเยน สาวตฺถิยา อวิทูเร อญฺญตรสฺส ปูคสฺส อาวสถปิณฺโฑ ปญฺญตฺโต โหติ.\csromanspacer{} ฉพฺพคฺคิยา ภิกฺขู ปุพฺพณฺหสมยํ นิวาเสตฺวา ปตฺตจีวรมาทาย สาวตฺถิํ ปิณฺฑาย ปวิสิตฺวา ปิณฺฑํ อลภมานา อาวสถํ อคมํสุ.\csromanspacer{} มนุสฺสา “จิรสฺสมฺปิ ภทนฺตา อาคตา”ติ สกฺกจฺจํ ปริวิสิํสุ.\csromanspacer{} อถ โข ฉพฺพคฺคิยา ภิกฺขู ทุติยมฺปิ ทิวสํ ฯเปฯ.\csromanspacer{} ตติยมฺปิ ทิวสํ ปุพฺพณฺหสมยํ นิวาเสตฺวา ปตฺตจีวรมาทาย สาวตฺถิํ ปิณฺฑาย ปวิสิตฺวา ปิณฺฑํ อลภมานา อาวสถํ คนฺตฺวา ภุญฺชิํสุ.\csromanspacer{} อถ โข ฉพฺพคฺคิยานํ ภิกฺขูนํ เอตทโหสิ “กิํ มยํ กริสฺสาม อารามํ คนฺตฺวา, หิยฺโยปิ อิเธว อาคนฺตพฺพํ ภวิสฺสตี”ติ ตตฺเถว อนุวสิตฺวา อนุวสิตฺวา อาวสถปิณฺฑํ ภุญฺชนฺติ.\csromanspacer{} ติตฺถิยา อปสกฺกนฺติ.\csromanspacer{} มนุสฺสา อุชฺฌายนฺติ ขิยฺยนฺติ วิปาเจนฺติ “กถํ หิ \csromanfolio{96}นาม สมณา สกฺยปุตฺติยา อนุวสิตฺวา อนุวสิตฺวา อาวสถปิณฺฑํ ภุญฺชิสฺสนฺติ.\csromanspacer{} นยิเมสญฺเญว อาวสถปิณฺโฑ ปญฺญตฺโต, สพฺเพสญฺเญว \textbf{อาวสถปิณฺโฑ ปญฺญตฺโต”ติ.}}

\prose{อสฺโสสุํ โข ภิกฺขู เตสํ มนุสฺสานํ อุชฺฌายนฺตานํ ขิยฺยนฺตานํ วิปาเจนฺตานํ.\csromanspacer{} เย เต ภิกฺขู อปฺปิจฺฉา ฯเปฯ เต อุชฺฌายนฺติ ขิยฺยนฺติ วิปาเจนฺติ “กถํ หิ นาม ฉพฺพคฺคิยา ภิกฺขู อนุวสิตฺวา อนุวสิตฺวา อาวสถปิณฺฑํ ภุญฺชิสฺสนฺตี”ติ ฯเปฯ.\csromanspacer{} สจฺจํ กิร ตุเมฺห ภิกฺขเว อนุวสิตฺวา อนุวสิตฺวา อาวสถปิณฺฑํ ภุญฺชถาติ.\csromanspacer{} สจฺจํ ภควาติ.\csromanspacer{} วิครหิ พุทฺโธ ภควา ฯเปฯ.\csromanspacer{} กถํ หิ นาม ตุเมฺห โมฆปุริสา อนุวสิตฺวา อนุวสิตฺวา อาวสถปิณฺฑํ ภุญฺชิสฺสถ.\csromanspacer{} เนตํ โมฆปุริสา อปฺปสนฺนานํ วา ปสาทาย ฯเปฯ.\csromanspacer{} เอวญฺจ ปน ภิกฺขเว อิมํ สิกฺขาปทํ อุทฺทิเสยฺยาถ\csromandash{}}

\prose{\textbf{“เอโก อาวสถปิณฺโฑ ภุญฺชิตพฺโพ, ตโต เจ อุตฺตริ ภุญฺเชยฺย, ปาจิตฺติยนฺ”}ติ.}

\prose{เอวญฺจิทํ ภควตา ภิกฺขูนํ สิกฺขาปทํ ปญฺญตฺตํ โหติ.}

\proseitem{204}{เตน โข ปน สมเยน อายสฺมา สาริปุตฺโต โกสเลสุ ชนปเท สาวตฺถิํ คจฺฉนฺโต เยน อญฺญตโร อาวสโถ เตนุปสงฺกมิ.\csromanspacer{} มนุสฺสา “จิรสฺสมฺปิ เถโร อาคโต”ติ สกฺกจฺจํ ปริวิสิํสุ.\csromanspacer{} อถ โข อายสฺมโต สาริปุตฺตสฺส ภุตฺตาวิสฺส ขโร อาพาโธ อุปฺปชฺชิ, นาสกฺขิ ตมฺหา อาวสถา ปกฺกมิตุํ.\csromanspacer{} อถ โข เต มนุสฺสา ทุติยมฺปิ ทิวสํ อายสฺมนฺตํ สาริปุตฺตํ เอตทโวจุํ “ภุญฺชถ ภนฺเต”ติ.\csromanspacer{} อถ โข อายสฺมา สาริปุตฺโต “ภควตา ปฏิกฺขิตฺตํ อนุวสิตฺวา อนุวสิตฺวา อาวสถปิณฺฑํ ภุญฺชิตุนฺ”ติ กุกฺกุจฺจายนฺโต น ปฏิคฺคเหสิ, ฉินฺนภตฺโต อโหสิ.\csromanspacer{} อถ โข อายสฺมา สาริปุตฺโต สาวตฺถิํ คนฺตฺวา ภิกฺขูนํ เอตมตฺถํ อาโรเจสิ.\csromanspacer{} ภิกฺขู ภควโต เอตมตฺถํ อาโรเจสุํ.\csromanspacer{} อถ โข ภควา เอตสฺมิํ นิทาเน เอตสฺมิํ ปกรเณ ธมฺมิํ กถํ กตฺวา ภิกฺขู อามนฺเตสิ อนุชานามิ ภิกฺขเว คิลาเนน ภิกฺขุนา อนุวสิตฺวา อนุวสิตฺวา อาวสถปิณฺฑํ ภุญฺชิตุํ.\csromanspacer{} เอวญฺจ ปน ภิกฺขเว อิมํ สิกฺขาปทํ อุทฺทิเสยฺยาถ\csromandash{}}

\proseitem{205}{\textbf{“อคิลาเนน ภิกฺขุนา เอโก อาวสถปิณฺโฑ ภุญฺชิตพฺโพ, ตโต เจ อุตฺตริ ภุญฺเชยฺย, ปาจิตฺติยนฺ”}ติ.}

\proseitem{206}{\csromanfolio{97}\textbf{อคิลาโน} นาม สกฺโกติ ตมฺหา อาวสถา ปกฺกมิตุํ.}

\prose{\textbf{คิลาโน} นาม น สกฺโกติ ตมฺหา อาวสถา ปกฺกมิตุํ.}

\prose{\textbf{อาวสถปิณฺโฑ} นาม ปญฺจนฺนํ โภชนานํ อญฺญตรํ โภชนํ สาลาย วา มณฺฑเป วา รุกฺขมูเล วา อชฺโฌกาเส วา อโนทิสฺส ยาวทตฺโถ ปญฺญตฺโต โหติ.\csromanspacer{} อคิลาเนน ภิกฺขุนา สกิํ ภุญฺชิตพฺโพ, ตโต เจ อุตฺตริ ภุญฺชิสฺสามีติ ปฏิคฺคณฺหาติ, อาปตฺติ ทุกฺกฏสฺส.\csromanspacer{} อชฺโฌหาเร อชฺโฌหาเร อาปตฺติ}

\proseitem{207}{อคิลาโน อคิลานสญฺญี ตตุตฺตริ อาวสถปิณฺฑํ ภุญฺชติ, อาปตฺติ ปาจิตฺติยสฺส.\csromanspacer{} อคิลาโน เวมติโก ตตุตฺตริ อาวสถปิณฺฑํ ภุญฺชติ, อาปตฺติ ปาจิตฺติยสฺส.\csromanspacer{} อคิลาโน คเลนสญฺญี ตตุตฺตริ อาวสถปิณฺฑํ ภุญฺชติ, อาปตฺติ}

\prose{คิลาโน อคิลานสญฺญี, อาปตฺติ ทุกฺกฏสฺส.\csromanspacer{} คิลาโน เวมติโก, อาปตฺติ ทุกฺกฏสฺส.\csromanspacer{} คิลาโน คิลานสญฺญี, อนาปตฺติ.}

\proseitemwithcloserruled{208}{อนาปตฺติ คิลานสฺส, อคิลาโน สกิํ ภุญฺชติ, คจฺฉนฺโต วา อาคจฺฉนฺโต วา ภุญฺชติ, สามิกา นิมนฺเตตฺวา โภเชนฺติ, โอทิสฺส ปญฺญตฺโต โหติ, น ยาวทตฺโถ ปญฺญตฺโต โหติ, ปญฺจ โภชนานิ ฐเปตฺวา สพฺพตฺถ, อนาปตฺติ อุมฺมตฺตกสฺส อาทิกมฺมิกสฺสาติ.}{อาวสถปิณฺฑสิกฺขาปทํ นิฏฺฐิตํ ปฐมํ.}
\csromanmatikaanchor{mtk.37}
\csromantocmarkref{subsection}{2. คณโภชนสิกฺขาปท}{mtk.37}
\bookmark[dest={mtk.37},level=2]{2. คณโภชนสิกฺขาปท}
\csromanheader{2}{4. \textbf{โภชนวคฺค 2. คณโภชนสิกฺขาปท}}
\proseitem{209}{เตน สมเยน พุทฺโธ ภควา ราชคเห วิหรติ เวฬุวเน กลนฺทกนิวาเป.\csromanspacer{} เตน โข ปน สมเยน เทวทตฺโต ปริหีนลาภสกฺกาโร สปริโส กุเลสุ วิญฺญาเปตฺวา วิญฺญาเปตฺวา ภุญฺชติ.\csromanspacer{} มนุสฺสา อุชฺฌายนฺติ ขิยฺยนฺติ วิปาเจนฺติ “กถํ หิ นาม สมณา สกฺยปุตฺติยา กุเลสุ วิญฺญาเปตฺวา วิญฺญาเปตฺวา ภุญฺชิสฺสนฺติ.\csromanspacer{} กสฺส สมฺปนฺนํ น มนาปํ, กสฺส สาทุํ น รุจฺจตี”ติ.\csromanspacer{} อสฺโสสุํ โข ภิกฺขู เตสํ มนุสฺสานํ อุชฺฌายนฺตานํ ขิยฺยนฺตานํ วิปาเจนฺตานํ.\csromanspacer{} เย เต ภิกฺขู อปฺปิจฺฉา ฯเปฯ เต \csromanfolio{98}อุชฺฌายนฺติ ขิยฺยนฺติ วิปาเจนฺติ “กถํ หิ นาม เทวทตฺโต สปริโส กุเลสุ วิญฺญาเปตฺวา วิญฺญาเปตฺวา ภุญฺชิสฺสตี”ติ ฯเปฯ.\csromanspacer{} สจฺจํ กิร ตฺวํ เทวทตฺต สปริโส กุเลสุ วิญฺญาเปตฺวา วิญฺญาเปตฺวา ภุญฺชสีติ.\csromanspacer{} สจฺจํ ภควาติ.\csromanspacer{} \textbf{วิครหิ พุทฺโธ ภควา ฯเปฯ.}\csromanspacer{}\textbf{ กถํ หิ นาม ตฺวํ โมฆปุริส} \textbf{สปริโส กุเลสุ วิญฺญาเปตฺวา วิญฺญาเปตฺวา ภุญฺชิสฺสสิ.}\csromanspacer{}\textbf{ เนตํ โมฆปุริส} \textbf{อปฺปสนฺนานํ วา ปสาทาย ฯเปฯ.}\csromanspacer{}\textbf{ เอวญฺจ ปน ภิกฺขเว อิมํ} สิกฺขาปทํ อุทฺทิเสยฺยาถ\csromandash{}}

\prose{\textbf{“คณโภชเน ปาจิตฺติยนฺ”}ติ.}

\prose{เอวญฺจิทํ ภควตา ภิกฺขูนํ สิกฺขาปทํ ปญฺญตฺตํ โหติ.}

\proseitem{210}{เตน โข ปน สมเยน มนุสฺสา คิลาเน ภิกฺขู ภตฺเตน นิมนฺเตนฺติ.\csromanspacer{} ภิกฺขู กุกฺกุจฺจายนฺตา นาธิวาเสนฺติ “ปฏิกฺขิตฺตํ ภควตา คณโภชนนฺ”ติ.\csromanspacer{} ภควโต เอตมตฺถํ อาโรเจสุํ.\csromanspacer{} อถ โข ภควา เอตสฺมิํ นิทาเน เอตสฺมิํ ปกรเณ ธมฺมิํ กถํ กตฺวา ภิกฺขู อามนฺเตสิ อนุชานามิ ภิกฺขเว คิลาเนน ภิกฺขุนา คณโภชนํ ภุญฺชิตุํ.\csromanspacer{} เอวญฺจ ปน ภิกฺขเว อิมํ สิกฺขาปทํ อุทฺทิเสยฺยาถ\csromandash{}}

\prose{“คณโภชเน อญฺญตฺร สมยา, ปาจิตฺติยํ.\csromanspacer{}\textbf{ ตตฺถายํ สมโย, คิลานสมโย, อยํ ตตฺถ สมโย”}ติ.}

\prose{เอวญฺจิทํ ภควตา ภิกฺขูนํ สิกฺขาปทํ ปญฺญตฺตํ โหติ.}

\proseitem{211}{เตน โข ปน สมเยน มนุสฺสา จีวรทานสมเย สจีวรภตฺตํ ปฏิยาเทตฺวา ภิกฺขู นิมนฺเตนฺติ “โภเชตฺวา จีวเรน อจฺฉาเทสฺสามา”ติ.\csromanspacer{} ภิกฺขู กุกฺกุจฺจายนฺตา นาธิวาเสนฺติ “ปฏิกฺขิตฺตํ ภควตา คณโภชนนฺ”ติ.\csromanspacer{} จีวรํ ปริตฺตํ อุปฺปชฺชติ.\csromanspacer{} ภควโต เอตมตฺถํ อาโรเจสุํ ฯเปฯ.\csromanspacer{} อนุชานามิ ภิกฺขเว จีวรทานสมเย คณโภชนํ ภุญฺชิตุํ.\csromanspacer{} เอวญฺจ ปน ภิกฺขเว อิมํ สิกฺขาปทํ อุทฺทิเสยฺยาถ\csromandash{}}

\prose{“คณโภชเน อญฺญตฺร สมยา ปาจิตฺติยํ.\csromanspacer{}\textbf{ ตตฺถายํ สมโย, คิลานสมโย จีวรทานสมโย, อยํ ตตฺถ สมโย”}ติ.}

\prose{เอวญฺจิทํ ภควตา ภิกฺขูนํ สิกฺขาปทํ ปญฺญตฺตํ โหติ.}

\proseitem{212}{เตน โข ปน สมเยน มนุสฺสา จิวรการเก ภิกฺขู ภตฺเตน นิมนฺเตนฺติ.\csromanspacer{} ภิกฺขู กุกฺกุจฺจายนฺตา นาธิวาเสนฺติ “ปฏิกฺขิตฺตํ ภควตา \csromanfolio{99}คณโภชนนฺ”ติ.\csromanspacer{} ภควโต เอตมตฺถํ อาโรเจสุํ ฯเปฯ.\csromanspacer{} อนุชานามิ ภิกฺขเว จีวรการสมเย คณโภชนํ ภุญฺชิตุํ.\csromanspacer{} เอวญฺจ ปน ภิกฺขเว อิมํ สิกฺขาปทํ อุทฺทิเสยฺยาถ\csromandash{}}

\prose{“คณโภชเน อญฺญตฺร สมยา ปาจิตฺติยํ.\csromanspacer{}\textbf{ ตตฺถายํ สมโย, คิลานสมโย จีวรทานสมโย จีวรการสมโย, อยํ ตตฺถ สมโย”}ติ.}

\prose{เอวญฺจิทํ ภควตา ภิกฺขูนํ สิกฺขาปทํ ปญฺญตฺตํ โหติ.}

\proseitem{213}{เตน โข ปน สมเยน ภิกฺขู มนุสฺเสหิ สทฺธิํ อทฺธานํ คจฺฉนฺติ.\csromanspacer{} อถ โข เต ภิกฺขู เต มนุสฺเส เอตทโวจุํ “มุหุตฺตํ อาวุโส อาคเมถ ปิณฺฑาย จริสฺสามา”ติ.\csromanspacer{} เต เอวมาหํสุ “อิเธว ภนฺเต ภุญฺชถา”ติ.\csromanspacer{} ภิกฺขู กุกฺกุจฺจายนฺตา น ปฏิคฺคณฺหนฺติ “ปฏิกฺขิตฺตํ ภควตา คณโภชนนฺ”ติ.\csromanspacer{} ภควโต เอตมตฺถํ อาโรเจสุํ ฯเปฯ.\csromanspacer{} อนุชานามิ ภิกฺขเว อทฺธานคมนสมเย คณโภชนํ ภุญฺชิตุํ.\csromanspacer{} เอวญฺจ ปน ภิกฺขเว อิมํ สิกฺขาปทํ อุทฺทิเสยฺยาถ\csromandash{}}

\prose{“\textbf{คณโภชเน อญฺญตฺร สมยา ปาจิตฺติยํ.}\csromanspacer{}\textbf{ ตตฺถายํ สมโย,} \textbf{คิลานสมโย จีวรทานสมยฺ}โอ จีวรการสมโย อทฺธานคมนสมโย, อยํ ตตฺถ สมโย”ติ.}

\prose{เอวญฺจิทํ ภควตา ภิกฺขูนํ สิกฺขาปทํ ปญฺญตฺตํ โหติ.}

\proseitem{214}{เตน โข ปน สมเยน ภิกฺขู มนุสฺเสหิ สทฺธิํ นาวาย คจฺฉนฺติ.\csromanspacer{} อถ โข เต ภิกฺขู เต มนุสฺเส เอตทโวจุํ “มุหุตฺตํ อาวุโส ตีรํ อุปเนถ ปิณฺฑาย จริสฺสามา”ติ.\csromanspacer{} เต เอวมาหํสุ “อิเธว ภนฺเต ภุญฺชถา”ติ.\csromanspacer{} ภิกฺขู กุกฺกุจฺจายนฺตา น ปฏิคฺคณฺหนฺติ “ปฏิกฺขิตฺตํ ภควตา คณโภชนนฺ”ติ.\csromanspacer{} ภควโต เอตมตฺถํ อาโรเจสุํ ฯเปฯ.\csromanspacer{} อนุชานามิ ภิกฺขเว นาวาภิรุหนสมเย คณโภชนํ ภุญฺชิตุํ.\csromanspacer{} เอวญฺจ ปน ภิกฺขเว อิมํ สิกฺขาปทํ อุทฺทิเสยฺยาถ\csromandash{}}

\prose{“\textbf{คณโภชเน อญฺญตฺร สมยา ปาจิตฺติยํ.}\csromanspacer{}\textbf{ ตตฺถายํ สมโย,} \textbf{คิลานสมโย จีวรทานสมยฺ}โอ จีวรการสมโย อทฺธานคมนสมโย นาวาภิรุหนสมโย, อยํ ตตฺถ สมโย”ติ.}

\prose{เอวญฺจิทํ ภควตา ภิกฺขูนํ สิกฺขาปทํ ปญฺญตฺตํ โหติ.}

\proseitem{215}{\csromanfolio{100}เตน โข ปน สมเยน ทิสาสุ วสฺสํวุฏฺฐา ภิกฺขู ราชคหํ อาคจฺฉนฺติ ภควนฺตํ ทสฺสนาย.\csromanspacer{} มนุสฺสา นานาเวรชฺชเก ภิกฺขู ปสฺสิตฺวา ภตฺเตน นิมนฺเตนฺติ.\csromanspacer{} ภิกฺขู กุกฺกุจฺจายนฺตา นาธิวาเสนฺติ “ปฏิกฺขิตฺตํ ภควตา คณโภชนนฺ”ติ.\csromanspacer{} ภควโต เอตมตฺถํ อาโรเจสุํ ฯเปฯ.\csromanspacer{} อนุชานามิ ภิกฺขเว มหาสมเย คณโภชนํ ภุญฺชิตุํ.\csromanspacer{} เอวญฺจ ปน ภิกฺขเว อิมํ สิกฺขาปทํ อุทฺทิเสยฺยาถ\csromandash{}}

\prose{“คณโภชเน อญฺญตฺร สมยา ปาจิตฺติยํ.\csromanspacer{}\textbf{ ตตฺถายํ สมโย, คิลานสมโย จีวรทานสมโย จีวรการสมโย อทฺธานคมนสมโย นาวาภิรุหนสมโย มหาสมโย, อยํ ตตฺถ สมโย”}ติ.}

\prose{เอวญฺจิทํ ภควตา ภิกฺขูนํ สิกฺขาปทํ ปญฺญตฺตํ โหติ.}

\proseitem{216}{เตน โข ปน สมเยน รญฺโญ มาคธสฺส เสนิยสฺส พิมฺพิสารสฺส ญาติสาโลหิโต อาชีวเกสุ ปพฺพชิโต โหติ.\csromanspacer{} อถ โข โส อาชีวโก เยน ราชา มาคโธ เสนิโย พิมฺพิสาโร เตนุปสงฺกมิ, อุปสงฺกมิตฺวา ราชานํ มาคธํ เสนิยํ พิมฺพิสารํ เอตทโวจ “อิจฺฉามหํ มหาราช สพฺพปาสณฺฑิกภตฺตํ กาตุนฺ”ติ.\csromanspacer{} สเจ ตฺวํ ภนฺเต พุทฺธปฺปมุขํ ภิกฺขุสํฆํ ปฐมํ โภเชยฺยาสิ, เอวํ กเรยฺยามีติ.\csromanspacer{} อถ โข โส อาชีวโก ภิกฺขูนํ สนฺติเก ทูตํ ปาเหสิ “อธิวาเสนฺตุ เม ภิกฺขู สฺวาตนาย ภตฺตนฺ”ติ.\csromanspacer{} ภิกฺขู กุกฺกุจฺจายนฺตา นาธิวาเสนฺติ “ปฏิกฺขิตฺตํ ภควตา คณโภชนนฺ”ติ.\csromanspacer{} อถ โข โส อาชีวโก เยน ภควา เตนุปสงฺกมิ, อุปสงฺกมิตฺวา ภควตา สทฺธิํ สมฺโมทิ, สมฺโมทนียํ กถํ สารณียํ วีติสาเรตฺวา เอกมนฺตํ อฏฺฐาสิ, เอกมนฺตํ ฐิโต โข โส อาชีวโก ภควนฺตํ เอตทโวจ “ภวมฺปิ โคตโม ปพฺพชิโต, อหมฺปิ ปพฺพชิโต, อรหติ ปพฺพชิโต ปพฺพชิตสฺส ปิณฺฑํ ปฏิคฺคเหตุํ, อธิวาเสตุ เม ภวํ โคตโม สฺวาตนาย ภตฺตํ สทฺธิํ ภิกฺขุสํเฆนา”ติ.\csromanspacer{} อธิวาเสสิ ภควา ตุณฺหีภาเวน.\csromanspacer{} อถ โข โส อาชีวโก ภควโต อธิวาสนํ วิทิตฺวา ปกฺกามิ.\csromanspacer{} อถ โข ภควา เอตสฺมิํ นิทาเน เอตสฺมิํ ปกรเณ ธมฺมิํ กถํ กตฺวา ภิกฺขู อามนฺเตสิ อนุชานามิ ภิกฺขเว สมณภตฺตสมเย คณโภชนํ ภุญฺชิตุํ.\csromanspacer{} เอวญฺจ ปน ภิกฺขเว อิมํ สิกฺขาปทํ อุทฺทิเสยฺยาถ\csromandash{}}

\proseitem{217}{\csromanfolio{101}“คณโภชเน อญฺญตฺร สมยา ปาจิตฺติยํ.\csromanspacer{}\textbf{ ตตฺถายํ สมโย, คิลานสมโย จีวรทานสมโย จีวรการสมโย อทฺธานคมนสมโย นาวาภิรุหนสมโย มหาสมโย สมณภตฺตสมโย, อยํ ตตฺถ สมโย”}ติ.}

\proseitem{218}{\textbf{คณโภชนํ} นาม ยตฺถ จตฺถาโร ภิกฺขู ปญฺจนฺนํ โภชนานํ อญฺญตเรน โภชเนน นิมนฺติตา ภุญฺชนฺติ, เอตํ คณโภชนานํ นาม.}

\prose{\textbf{อญฺญตฺร สมยา}ติ ฐเปตฺวา สมยํ.}

\prose{\textbf{คิลานสมโย} นาม อนฺตมโส ปาทาปิ ผลิตา\footnote{ผาลิตา (สฺยา, ก)} โหนฺติ, “คิลานสมโย”ติ ภุญฺชิตพฺพํ.}

\prose{\textbf{จีวรทานสมโย} นาม อนตฺถเต กถิเน วสฺสานสฺส ปจฺฉิโม มาโส, อตฺถเต กถิเน ปญฺจมาสา, “จีวรทานสมโย”ติ ภุญฺชิตพฺพํ.}

\prose{\textbf{จีวรการสมโย} นาม จีวเร กยิรมาเน “จีวรการสมโย”ติ ภุญฺชิตพฺพํ.}

\prose{\textbf{อทฺธานคมนสมโย} นาม อทฺธโยชนํ คจฺฉิสฺสามีติ ภุญฺชิตพฺพํ.\csromanspacer{} คจฺฉนฺเตน ภุญฺชิตพฺพํ, คเตน ภุญฺชิตพฺพํ.}

\prose{\textbf{นาวาภิรุหนสมโย} นาม นาวํ อภิรุหิสฺสามีติ ภุญฺชิตพฺพํ.\csromanspacer{} อารุเฬฺหน ภุญฺชิตพฺพํ, โอรุเฬฺหน ภุญฺชิตพฺพํ.}

\prose{\textbf{มหาสมโย} นาม ยตฺถ เทฺว ตโย ภิกฺขู ปิณฺฑาย จริตฺวา ยาเปนฺติ, จตุตฺเถ อาคเต น ยาเปนฺติ, “มหาสมโย”ติ ภุญฺชิตพฺพํ.}

\prose{\textbf{สมณภตฺตสมโย} นาม โย โกจิ ปริพฺพาชกสมาปนฺโน ภตฺตํ กโรติ, “สมณภตฺตสมโย”ติ ภุญฺชิตพฺพํ.}

\prose{อญฺญตฺร สมยา ภุญฺชิสฺสามีติ ปฏิคฺคณฺหาติ, อาปตฺติ ทุกฺกฏสฺส.\csromanspacer{} อชฺโฌหาเร อชฺโฌหาเร อาปตฺติ ปาจิตฺติยสฺส.}

\proseitem{219}{คณโภชเน คณโภชนสญฺญี อญฺญตฺร สมยา ภุญฺชติ, อาปตฺติ ปาจิตฺติยสฺส.\csromanspacer{} คณโภชเน เวมติโก อญฺญตฺร สมยา ภุญฺชติ, อาปตฺติ ปาจิตฺติยสฺส.\csromanspacer{} คณโภชเน น คณโภชนสญฺญี อญฺญตฺร สมยา ภุญฺชติ, อาปตฺติ ปาจิตฺติยสฺส.}

\prose{\csromanfolio{102}น คณโภชเน คณโภชนสญฺญี, อาปตฺติ ทุกฺกฏสฺส.\csromanspacer{} น คณโภชเน เวมติโก, อาปตฺติ ทุกฺกฏสฺส.\csromanspacer{} น คณโภชเน น คณโภชนสญฺญี, อนาปตฺติ.}

\proseitemwithcloserruled{220}{อนาปตฺติ สมเย, เทฺว ตโย เอกโต ภุญฺชนฺติ, ปิณฺฑาย จริตฺวา เอกโต สนฺนิปติตฺวา ภุญฺชนฺติ, นิจฺจภตฺตํ, สลากภตฺตํ, ปกฺขิกํ, อุโปสถิกํ, ปาฏิปทิกํ, ปญฺจ โภชนานิ ฐเปตฺวา สพฺพตฺถ, อนาปตฺติ อุมฺมตฺตกสฺส อาทิกมฺมิกสฺสาติ.}{คณโภชนสิกฺขาปทํ นิฏฺฐิตํ ทุติยํ.}
\csromanmatikaanchor{mtk.38}
\csromantocmarkref{subsection}{3. ปรมฺปรโภชนสิกฺขาปท}{mtk.38}
\bookmark[dest={mtk.38},level=2]{3. ปรมฺปรโภชนสิกฺขาปท}
\csromanheader{2}{4. โภชนวคฺค 3. ปรมฺปรโภชน\-สิกฺขาปท}
\proseitem{221}{เตน สมเยน พุทฺโธ ภควา เวสาลิยํ วิหรติ มหาวเน กูฏาคารสาลายํ.\csromanspacer{} เตน โข ปน สมเยน เวสาลิยํ ปณีตานํ ภตฺตานํ ภตฺตปฏิปาฏิ อธิฏฺฐิตา โหติ.\csromanspacer{} อถ โข อญฺญตรสฺส ทลิทฺทสฺส กมฺมการสฺส\footnote{กมฺมกรสฺส (สี)} เอตทโหสิ “น โข อิทํ โอรกํ ภวิสฺสติ ยถยิเม มนุสฺสา สกฺกจฺจํ ภตฺตํ กโรนฺติ, ยนฺนูนาหมฺปิ ภตฺตํ กเรยฺยนฺ”ติ.\csromanspacer{} อถ โข โส ทลิทฺโท กมฺมกาโร เยน กิรปติโก เตนุปสงฺกมิ, อุปสงฺกมิตฺวา ตํ กิรปติกํ เอตทโวจ “อิจฺฉามหํ อยฺยปุตฺต พุทฺธปฺปมุขสฺส ภิกฺขุสํฆสฺส ภตฺตํ กาตุํ, เทหิ เม เวตนนฺ”ติ.\csromanspacer{} โสปิ โข กิรปติโก สทฺโธ โหติ ปสนฺโน.\csromanspacer{} อถ โข โส กิรปติโก ตสฺส ทลิทฺทสฺส กมฺมการสฺส อพฺภาติเรกํ\footnote{อติเรกํ (สฺยา)} เวตนํ อทาสิ.\csromanspacer{} อถ โข โส ทลิทฺโท กมฺมกาโร เยน ภควา เตนุปสงฺกมิ, อุปสงฺกมิตฺวา ภควนฺตํ อภิวาเทตฺวา เอกมนฺตํ นิสีทิ, เอกมนฺตํ นิสินฺโน โข โส ทลิทฺโท กมฺมกาโร ภควนฺตํ เอตทโวจ “อธิวาเสตุ เม ภนฺเต ภควา สฺวาตนาย ภตฺตํ สทฺธิํ ภิกฺขุสํเฆนา”ติ. “มหา โข อาวุโส ภิกฺขุสํโฆ, ชานาหี”ติ. “โหตุ\footnote{โหตุ เม (ก)} ภนฺเต มหา ภิกฺขุสํโฆ, พหู เม พทรา ปฏิยตฺตา, พทรมิสฺเสน เปยฺยา ปริปูริสฺสนฺตี”ติ.\csromanspacer{} อธิวาเสสิ ภควา ตุณฺหีภาเวน.}

\prose{\csromanfolio{103}อถ โข โส ทลิทฺโท กมฺมกาโร ภควโต อธิวาสนํ วิทิตฺวา อุฏฺฐายาสนา ภควนฺตํ อภิวาเทตฺวา ปทกฺขิณํ กตฺวา ปกฺกามิ.\csromanspacer{} อสฺโสสุํ โข ภิกฺขู “ทลิทฺเทน กิร กมฺมกาเรน สฺวาตนาย พุทฺธปฺปมุโข ภิกฺขุสํโฆ นิมนฺติโต, พทรมิสฺเสน เปยฺยา ปริปูริสฺสนฺตี”ติ.\csromanspacer{} เต กาลสฺเสว ปิณฺฑาย จริตฺวา ภุญฺชิํสุ.\csromanspacer{} อสฺโสสุํ โข มนุสฺสา “ทลิทฺเทน กิร กมฺมกาเรน พุทฺธปฺปมุโข ภิกฺขุสํโฆ นิมนฺติโต”ติ.\csromanspacer{} เต ทลิทฺทสฺส กมฺมการสฺส ปหูตํ ขาทนียํ โภชนียํ อภิหริํสุ.\csromanspacer{} อถ โข โส ทลิทฺโท กมฺมกาโร ตสฺสา รตฺติยา อจฺจเยน ปณีตํ ขาทนียํ โภชนียํ ปฏิยาทาเปตฺวา ภควโต กาลํ อาโรจาเปสิ “กาโล ภนฺเต นิฏฺฐิตํ ภตฺตนฺ”ติ.}

\prose{อถ โข ภควา ปุพฺพณฺหสมยํ นิวาเสตฺวา ปตฺตจีวรมาทาย เยน ตสฺส ทลิทฺทสฺส กมฺมการสฺส นิเวสนํ เตนุปสงฺกมิ, อุปสงฺกมิตฺวา ปญฺญตฺเต อาสเน นิสีทิ สทฺธิํ ภิกฺขุสํเฆน.\csromanspacer{} อถ โข โส ทลิทฺโท กมฺมกาโร ภตฺตคฺเค ภิกฺขู ปริวิสติ.\csromanspacer{} ภิกฺขู เอวมาหํสุ “โถกํ อาวุโส เทหิ, โถกํ อาวุโส เทหี”ติ.\csromanspacer{} มา โข ตุเมฺห ภนฺเต อยํ ทลิทฺโท กมฺมกาโรติ โถกํ โถกํ ปฏิคฺคณฺหิตฺถ, ปหูตํ เม ขาทนียํ โภชนียํ ปฏิยตฺตํ, ปฏิคฺคณฺหถ ภนฺเต ยาวทตฺถนฺติ.\csromanspacer{} น โข มยํ อาวุโส เอตํการณา โถกํ โถกํ ปฏิคฺคณฺหาม, อปิ จ มยํ กาลสฺเสว ปิณฺฑาย จริตฺวา ภุญฺชิมฺหา, เตน มยํ โถกํ โถกํ ปฏิคฺคณฺหามาติ.}

\prose{อถ โข โส ทลิทฺโท กมฺมกาโร อุชฺฌายติ ขิยฺยติ วิปาเจติ “กถํ หิ นาม ภทนฺตา มยา นิมนฺติตา อญฺญตฺร ภุญฺชิสฺสนฺติ, น จาหํ ปฏิพโล ยาวทตฺถํ ทาตุนฺ”ติ.\csromanspacer{} อสฺโสสุํ โข ภิกฺขู ตสฺส ทลิทฺทสฺส กมฺมการสฺส อุชฺฌายนฺตสฺส ขิยฺยนฺตสฺส วิปาเจนฺตสฺส.\csromanspacer{} เย เต ภิกฺขู อปฺปิจฺฉา ฯเปฯ เต อุชฺฌายนฺติ ขิยฺยนฺติ วิปาเจนฺติ “กถํ หิ นาม ภิกฺขู อญฺญตฺร นิมนฺติตา อญฺญตฺร ภุญฺชิสฺสนฺตี”ติ ฯเปฯ.\csromanspacer{} สจฺจํ กิร ภิกฺขเว ภิกฺขู อญฺญตฺร นิมนฺติตา อญฺญตฺร ภุญฺชนฺตีติ.\csromanspacer{} สจฺจํ ภควาติ.\csromanspacer{} วิครหิ พุทฺโธ ภควา ฯเปฯ.\csromanspacer{} กถํ หิ นาม เต ภิกฺขเว โมฆปุริสา อญฺญตฺร นิมนฺติตา อญฺญตฺร ภุญฺชิสฺสนฺติ.\csromanspacer{} เนตํ ภิกฺขเว อปฺปสนฺนานํ วา ปสาทาย ฯเปฯ.\csromanspacer{} เอวญฺจ ปน ภิกฺขเว อิมํ สิกฺขาปทํ อุทฺทิเสยฺยาถ\csromandash{}}

\prose{\csromanfolio{104}\textbf{“ปรมฺปรโภชเน ปาจิตฺติยนฺ”}ติ.}

\prose{เอวญฺจิทํ ภควตา ภิกฺขูนํ สิกฺขาปทํ ปญฺญตฺตํ โหติ.}

\proseitem{222}{เตน โข ปน สมเยน อญฺญตโร ภิกฺขุ คิลาโน โหติ.\csromanspacer{} อญฺญตโร ภิกฺขุ ปิณฺฑปาตํ อาทาย เยน โส ภิกฺขุ เตนุปสงฺกมิ, อุปสงฺกมิตฺวา ตํ ภิกฺขุํ เอตทโวจ “ภุญฺชาหิ อาวุโส”ติ.\csromanspacer{} อลํ อาวุโส อตฺถิ เม ภตฺตปจฺจาสาติ.\csromanspacer{} ตสฺส ภิกฺขุโน ปิณฺฑปาโต อุสฺสูเร\footnote{อุสฺสูเรน (ก)} อาหรียิตฺถ, โส ภิกฺขุ น จิตฺตรูปํ ภุญฺชิ.\csromanspacer{} ภควโต เอตมตฺถํ อาโรเจสุํ.\csromanspacer{} อถ โข ภควา เอตสฺมิํ นิทาเน เอตสฺมิํ ปกรเณ ธมฺมิํ กถํ กตฺวา ภิกฺขู อามนฺเตสิ อนุชานามิ ภิกฺขเว คิลาเนน ภิกฺขุนา ปรมฺปรโภชนํ ภุญฺชิตุํ.\csromanspacer{} เอวญฺจ ปน ภิกฺขเว อิมํ สิกฺขาปทํ อุทฺทิเสยฺยาถ\csromandash{}}

\prose{“ปรมฺปรโภชเน อญฺญตฺร สมยา ปาจิตฺติยํ.\csromanspacer{}\textbf{ ตตฺถายํ สมโย, คิลานสมโย, อยํ ตตฺถ สมโย”}ติ.}

\prose{เอวญฺจิทํ ภควตา ภิกฺขูนํ สิกฺขาปทํ ปญฺญตฺตํ โหติ.}

\proseitem{223}{เตน โข ปน สมเยน มนุสฺสา จีวรทานสมเย สจีวรภตฺตํ ปฏิยาเทตฺวา\footnote{ปฏิยาทาเปตฺวา (อิติปิ)} ภิกฺขู นิมนฺเตนฺติ “โภเชตฺวา จีวเรน อจฺฉาเทสฺสามา”ติ.\csromanspacer{} ภิกฺขู กุกฺกุจฺจายนฺตา นาธิวาเสนฺติ “ปฏิกฺขิตฺตํ ภควตา ปรมฺปรโภชนนฺ”ติ.\csromanspacer{} จิวรํ ปริตฺตํ อุปฺปชฺชติ.\csromanspacer{} ภควโต เอตมตฺถํ อาโรเจสุํ ฯเปฯ.\csromanspacer{} อนุชานามิ ภิกฺขเว จีวรทานสมเย ปรมฺปรโภชนํ ภุญฺชิตุํ.\csromanspacer{} เอวญฺจ ปน ภิกฺขเว อิมํ สิกฺขาปทํ อุทฺทิเสยฺยาถ\csromandash{}}

\prose{“\textbf{ปรมฺปรโภชเน อญฺญตฺร สมยา ปาจิตฺติยํ.}\csromanspacer{}\textbf{ ตตฺถายํ สมโย,} คิลานสมโย จีวรทานสมโย, อยํ ตตฺถ สมโย”ติ.}

\prose{เอวญฺจิทํ ภควตา ภิกฺขูนํ สิกฺขาปทํ ปญฺญตฺตํ โหติ.}

\proseitem{224}{เตน โข ปน สมเยน มนุสฺสา จีวรการเก ภิกฺขู ภตฺเตน นิมนฺเตนฺติ.\csromanspacer{} ภิกฺขู กุกฺกุจฺจายนฺตา นาธิวาเสนฺติ “ปฏิกฺขิตฺตํ ภควตา ปรมฺปรโภชนนฺ”ติ.\csromanspacer{} ภควโต เอตมตฺถํ อาโรเจสุํ ฯเปฯ.\csromanspacer{} อนุชานามิ ภิกฺขเว จีวรการสมเย ปรมฺปรโภชนํ ภุญฺชิตุํ.\csromanspacer{} เอวญฺจ ปน ภิกฺขเว อิมํ สิกฺขาปทํ อุทฺทิเสยฺยาถ\csromandash{}}

\proseitem{225}{\csromanfolio{105}“ปรมฺปรโภชเน อญฺญตฺร สมยา ปาจิตฺติยํ.\csromanspacer{}\textbf{ ตตฺถายํ สมโย, คิลานสมโย จีวรทานสมโย จีวรการสมโย, อยํ ตตฺถ สมโย”}ติ.}

\prose{เอวญฺจิทํ ภควตา ภิกฺขูนํ สิกฺขาปทํ ปญฺญตฺตํ โหติ.}

\proseitem{226}{อถ โข ภควา ปุพฺพณฺหสมยํ นิวาเสตฺวา ปตฺตจีวรมาทาย อายสฺมตา อานนฺเทน ปจฺฉาสมเณน เยน อญฺญตรํ กุลํ เตนุปสงฺกมิ, อุปสงฺกมิตฺวา ปญฺญตฺเต อาสเน นิสีทิ.\csromanspacer{} อถ โข เต มนุสฺสา ภควโต จ อายสฺมโต จ อานนฺทสฺส โภชนํ อทํสุ.\csromanspacer{} อายสฺมา อานนฺโท กุกฺกุจฺจายนฺโต น ปฏิคฺคณฺหาติ.\csromanspacer{} คณฺหาหิ\footnote{ปติคณฺหาหิ (สี)} อานนฺทาติ.\csromanspacer{} อลํ ภควา อตฺถิ เม ภตฺตปจฺจาสาติ.\csromanspacer{} เตนหานนฺท วิกปฺเปตฺวา คณฺหาหีติ.\csromanspacer{} อถ โข ภควา เอตสฺมิํ นิทาเน เอตสฺมิํ ปกรเณ ธมฺมิํ กถํ กตฺวา ภิกฺขู อามนฺเตสิ อนุชานามิ ภิกฺขเว วิกปฺเปตฺวา\footnote{ภตฺตปจฺจาสํ วิกปฺเปตฺวา (สฺยา)} ปรมฺปรโภชนํ ภุญฺชิตุํ.\csromanspacer{} เอวญฺจ ปน ภิกฺขเว วิกปฺเปตพฺพํ “มยฺหํ ภตฺตปจฺจาสํ อิตฺถนฺนามสฺส ทมฺมี”ติ.}

\proseitem{227}{\textbf{ปรมฺปรโภชนํ} นาม ปญฺจนฺนํ โภชนานํ อญฺญตเรน โภชเนน นิมนฺติโต, ตํ ฐเปตฺวา อญฺญํ ปญฺจนฺนํ โภชนานํ อญฺญตรํ โภชนํ ภุญฺชติ, เอตํ ปรมฺปรโภชนํ นาม.}

\prose{\textbf{อญฺญตฺร สมยา}ติ ฐเปตฺวา สมยํ.}

\prose{\textbf{คิลานสมโย} นาม น สกฺโกติ เอกาสเน นิสินฺโน ยาวทตฺถํ ภุญฺชิตุํ, คิลานสมโยติ ภุญฺชิตพฺพํ.}

\prose{\textbf{จีวรทานสมโย} นาม อนตฺถเต กถิเน วสฺสานสฺส ปจฺฉิโม มาโส, อตฺถเต กถิเน ปญฺจ มาสา, จีวรทานสมโยติ ภุญฺชิตพฺพํ.}

\prose{\textbf{จีวรการสมโย} นาม จีวเร กยิรมาเน จีวรการสมโยติ ภุญฺชิตพฺพํ.}

\prose{อญฺญตฺร สมยา ภุญฺชิสฺสามีติ ปฏิคฺคณฺหาติ, อาปตฺติ ทุกฺกฏสฺส.\csromanspacer{} อชฺโฌหาเร อชฺโฌหาเร อาปตฺติ ปาจิตฺติยสฺส.}

\proseitem{228}{\csromanfolio{106}ปรมฺปรโภชเน ปรมฺปรโภชนสญฺญี อญฺญตฺร สมยา ภุญฺชติ, อาปตฺติ ปาจิตฺติยสฺส.\csromanspacer{} ปรมฺปรโภชเน เวมติโก อญฺญตฺร สมยา ภุญฺชติ, อาปตฺติ ปาจิตฺติยสฺส.\csromanspacer{} ปรมฺปรโภชเน น ปรมฺปรโภชนสญฺญี อญฺญตฺร สมยา ภุญฺชติ, อาปตฺติ ปาจิตฺติยสฺส.}

\prose{น ปรมฺปรโภชเน ปรมฺปรโภชนสญฺญี, อาปตฺติ ทุกฺกฏสฺส.\csromanspacer{} น ปรมฺปรโภชเน เวมติโก, อาปตฺติ ทุกฺกฏสฺส.\csromanspacer{} น ปรมฺปรโภชเน น ปรมฺปรโภชนสญฺญี, อนาปตฺติ.}

\proseitemwithcloserruled{229}{อนาปตฺติ สมเย, วิกปฺเปตฺวา ภุญฺชติ, เทฺว ตโย นิมนฺตเน เอกโต ภุญฺชติ, นิมนฺตนปฏิปาฏิยา ภุญฺชติ, สกเลน คาเมน นิมนฺติโต ตสฺมิํ คาเม ยตฺถ กตฺถจิ ภุญฺชติ, สกเลน ปูเคน นิมนฺติโต ตสฺมิํ ปูเค ยตฺถ กตฺถจิ ภุญฺชติ, นิมนฺติยมาโน ภิกฺขํ คเหสฺสามีติ ภณติ, นิจฺจภตฺเต สลากภตฺเต ปกฺขิเก อุโปสถิเก ปาฏิปทิเก ปญฺจ โภชนานิ ฐเปตฺวา สพฺพตฺถ, อนาปตฺติ อุมฺมตฺตกสฺส อาทิกมฺมิกสฺสาติ.}{ปรมฺปรโภชน\-สิกฺขาปทํ นิฏฺฐิตํ ตติยํ.}
\csromanmatikaanchor{mtk.39}
\csromantocmarkref{subsection}{4. กาณมาตุสิกฺขาปท}{mtk.39}
\bookmark[dest={mtk.39},level=2]{4. กาณมาตุสิกฺขาปท}
\csromanheader{2}{4. \textbf{โภชนวคฺค 4. กาณมาตุสิกฺขาปท}}
\proseitem{230}{เตน สมเยน พุทฺโธ ภควา สาวตฺถิยํ วิหรติ เชตวเน อนาถปิณฺฑิกสฺส อาราเม.\csromanspacer{} เตน โข ปน สมเยน กาณมาตา อุปาสิกา สทฺธา โหติ ปสนฺนา, กาณา คามเก อญฺญตรสฺส ปุริสสฺส ทินฺนา โหติ.\csromanspacer{} อถ โข กาณา มาตุฆรํ อคมาสิ เกนจิเทว กรณีเยน.\csromanspacer{} อถ โข กาณาย สามิโก กาณาย สนฺติเก ทูตํ ปาเหสิ “อาคจฺฉตุ กาณา, อิจฺฉามิ กาณาย อาคตนฺ”ติ.\csromanspacer{} อถ โข กาณมาตา อุปาสิกา “กิสฺมิํ วิย ริตฺตหตฺถํ คนฺตุนฺ”ติ ปูวํ\footnote{ปูปํ (ªวาทิโมคฺคลฺลาเน)} ปจิ.\csromanspacer{} ปกฺเก ปูเว อญฺญตโร ปิณฺฑจาริโก ภิกฺขุ กาณมาตาย อุปาสิกาย นิเวสนํ ปาวิสิ.\csromanspacer{} อถ โข กาณมาตา อุปาสิกา ตสฺส ภิกฺขุโน ปูวํ ทาเปสิ, โส นิกฺขมิตฺวา อญฺญสฺส อาจิกฺขิ, ตสฺสปิ ปูวํ ทาเปสิ,โส นิกฺขมิตฺวา อญฺญสฺส อาจิกฺขิ, ตสฺสปิ \csromanfolio{107}ปูวํ ทาเปสิ, ยถาปฏิยตฺตํ ปูวํ ปริกฺขยํ อคมาสิ.\csromanspacer{} ทุติยมฺปิ โข กาณาย สามิโก กาณาย สนฺติเก ทูตํ ปาเหสิ “อาคจฺฉตุ กาณา, อิจฺฉามิ กาณาย อาคตนฺ”ติ.\csromanspacer{} ทุติยมฺปิ โข กาณมาตา อุปาสิกา “กิสฺมิํ วิย ริตฺตหตฺตํ คนฺตุนฺ”ติ ปูวํ ปจิ.\csromanspacer{} ปกฺเก ปูเว อญฺญตโร ปิณฺฑจาริโก ภิกฺขุ กาณมาตาย อุปาสิกาย นิเวสนํ ปาวิสิ.\csromanspacer{} อถ โข กาณมาตา อุปาสิกา ตสฺส ภิกฺขุโน ปูวํ ทาเปสิ, โส นิกฺขมิตฺวา อญฺญสฺส อาจิกฺขิ, ตสฺสปิ ปูวํ ทาเปสิ, โส นิกฺขมิตฺวา อญฺญสฺส อาจิกฺขิ, ตสฺสปิ ปูวํ ทาเปสิ, ยถาปฏิยตฺตํ ปูวํ ปริกฺขยํ อคมาสิ.\csromanspacer{} ตติยมฺปิ โข กาณาย สามิโก กาณาย สนฺติเก ทูตํ ปาเหสิ “อาคจฺฉตุ กาณา, อิจฺฉามิ กาณาย อาคตํ, สเจ กาณา นาคจฺฉิสฺสติ, อหํ อญฺญํ ปชาปติํ อาเนสฺสามี”ติ.\csromanspacer{} ตติยมฺปิ โข กาณมาตา อุปาสิกา “กิสฺมิํ วิย ริตฺตหตฺถํ คนฺตุนฺ”ติ ปูวํ ปจิ.\csromanspacer{} ปกฺเก ปูเว อญฺญตโร ปิณฺฑจาริโก ภิกฺขุ กาณมาตาย อุปาสิกาย นิเวสนํ ปาวิสิ.\csromanspacer{} อถ โข กาณามาตา อุปาสิกา ตสฺส ภิกฺขุโน ปูวํ ทาเปสิ, โส นิกฺขมิตฺวา อญฺญสฺส อาจิกฺขิ, ตสฺสปิ ปูวํ ทาเปสิ, โส นิกฺขมิตฺวา อญฺญสฺส อาจิกฺขิ, ตสฺสปิ ปูวํ ทาเปสิ, ยถาปฏิยตฺตํ ปูวํ ปริกฺขยํ อคมาสิ.\csromanspacer{} อถ โข กาณาย สามิโก อญฺญํ ปชาปติํ อาเนสิ.}

\prose{อสฺโสสิ โข กาณา “เตน กิร ปุริเสน อญฺญา ปชาปติ อานีตา”ติ.\csromanspacer{} สา โรทนฺตี อฏฺฐาสิ.\csromanspacer{} อถ โข ภควา ปุพฺพณฺหสมยํ นิวาเสตฺวา ปตฺตจีวรมาทาย เยน กาณมาตาย อุปาสิกาย นิเวสนํ เตนุปสงฺกมิ, อุปสงฺกมิตฺวา ปญฺญตฺเต อาสเน นิสีทิ.\csromanspacer{} อถ โข กาณมาตา อุปาสิกา เยน ภควา เตนุปสงฺกมิ, อุปสงฺกมิตฺวา ภควนฺตํ อภิวาเทตฺวา เอกมนฺตํ นิสีทิ, เอกมนฺตํ นิสินฺนํ โข กาณมาตรํ อุปาสิกํ ภควา เอตทโวจ “กิสฺสายํ กาณา โรทตี”ติ.\csromanspacer{} อถ โข กาณมาตา อุปาสิกา ภควโต เอตมตฺถํ อาโรเจสิ.\csromanspacer{} อถ โข ภควา กาณมาตรํ อุปาสิกํ ธมฺมิยา กถาย สนฺทสฺเสตฺวา สมาทเปตฺวา สมุตฺเตเชตฺวา สมฺปหํเสตฺวา อุฏฺฐายาสนา ปกฺกามิ.}

\proseitem{231}{เตน โข ปน สมเยน อญฺญตโร สตฺโถ ราชคหา ปฏิยาโลกํ คนฺตุกาโม โหติ.\csromanspacer{} อญฺญตโร ปิณฺฑจาริโก ภิกฺขุ ตํ สตฺถํ ปิณฺฑาย ปาวิสิ.\csromanspacer{} อญฺญตโร อุปาสโก ตสฺส ภิกฺขุโน \csromanfolio{108}สตฺตุํ ทาเปสิ, โส นิกฺขิมิตฺวา อญฺญสฺส อาจิกฺขิ, ตสฺสปิ สตฺตุํ ทาเปสิ.\csromanspacer{} โส นิกฺขิมิตฺวา อญฺญสฺส อาจิกฺขิ, ตสฺสปิ สตฺตุํ ทาเปสิ, ยถาปฏิยตฺตํ ปาเถยฺยํ ปริกฺขยํ อคมาสิ.\csromanspacer{} อถ โข โส อุปาสโก เต มนุสฺเส เอตทโวจ “อชฺชโณฺห อยฺยา อาคเมถ, ยถาปฏิยตฺตํ ปาเถยฺยํ อยฺยานํ ทินฺนํ, ปาเถยฺยํ ปฏิยาเทสฺสามี”ติ. “นายฺโย\footnote{นายฺย (สฺยา)} สกฺกา อาคเมตุํ, ปยาโต สตฺโถ”ติ อคมํสุ.\csromanspacer{} อถ โข ตสฺส อุปาสกสฺส ปาเถยฺยํ ปฏิยาเทตฺวา \textbf{ปจฺฉา คจฺฉนฺตสฺส โจรา อจฺฉินฺทิํสุ.}\csromanspacer{}\textbf{ มนุสฺสา อุชฺฌายนฺติ ขิยฺยนฺติ} \textbf{วิปาเจนฺติ “กถํ หิ นาม สมณา สกฺยปุตฺติยา น มตฺตํ ชานิตฺวา} ปฏิคฺคเหสฺสนฺติ, อยํ อิเมสํ ทตฺวา ปจฺฉา คจฺฉนฺโต โจเรหิ อจฺฉินฺโน”ติ.\csromanspacer{} อสฺโสสุํ โข ภิกฺขู เตสํ มนุสฺสานํ อุชฺฌายนฺตานํ ขิยฺยนฺตานํ วิปาเจนฺตานํ.\csromanspacer{} อถ โข เต ภิกฺขู ภควโต เอตมตฺถํ อาโรเจสุํ.\csromanspacer{} อถ โข ภควา เอตสฺมิํ นิทาเน เอตสฺมิํ ปกรเณ ธมฺมิํ กถํ กตฺวา ภิกฺขู อามนฺเตสิ เตน หิ ภิกฺขเว ภิกฺขูนํ สิกฺขาปทํ ปญฺญเปสฺสามิ ทส อตฺถวเส ปฏิจฺจ สํฆสุฏฺฐุตาย สํฆผาสุตาย ฯเปฯ วินยานุคฺคหาย.\csromanspacer{} เอวญฺจ ปน ภิกฺขเว อิมํ สิกฺขาปทํ อุทฺทิเสยฺยาถ\csromandash{}}

\proseitem{232}{\textbf{“ภิกฺขุํ ปเนว กุลํ อุปคตํ ปูเวหิ วา มนฺเถหิ วา อภิหฏฺฐุํ ปวาเรยฺย, อากงฺขมาเนน ภิกฺขุนา ทฺวตฺติปตฺตปูรา ปฏิคฺคเหตพฺพา, ตโต เจ อุตฺตริ ปฏิคฺคเณฺหยฺย, ปาจิตฺติยํ.}\csromanspacer{} ทฺวตฺติปตฺตปูเร ปฏิคฺคเหตฺวา ตโต นีหริตฺวา ภิกฺขูหิ สทฺธิํ สํวิภชิตพฺพํ, อยํ ตตฺถ สามีจี”ติ.}

\proseitem{233}{\textbf{ภิกฺขุํ ปเนว กุล}ํ อุปคตนฺติ กุลํ นาม จตฺตาริ กุลานิ ขตฺติยกุลํ พฺราหฺมณกุลํ เวสฺสกุลํ สุทฺทกุลํ.}

\prose{\textbf{อุปคตนฺ}ติ ตตฺถ คตํ.}

\prose{\textbf{ปูวํ} นาม ยํกิญฺจิ ปเหณกตฺถาย ปฏิยตฺตํ.}

\prose{\textbf{มนฺถํ} นาม ยํกิญฺจิ ปาเถยฺยตฺถาย ปฏิยตฺตํ.}

\prose{\textbf{อภิหฏฺฐุํ ปวาเรยฺยา}ติ ยาวตกํ อิจฺฉสิ, ตาวตกํ คณฺหาหีติ.}

\prose{\textbf{อากงฺขมาเนนา}ติ อิจฺฉมาเนน.}

\prose{\textbf{ทฺวตฺติปตฺตปูรา ปฏิคฺคเหตพฺพา}ติ เทฺวตโย ปตฺตปูรา ปฏิคฺคเหตพฺพา.}

\prose{\csromanfolio{109}\textbf{ตโตเจ อุตฺตริ ปฏิคฺคเณฺหยฺยา}ติ ตตุตฺตริ ปฏิคฺคณฺหาติ, อาปตฺติ ปาจิตฺติยสฺส.}

\prose{\textbf{ทฺวตฺติปตฺตปูเร ปฏิคฺคเหตฺวา} ตโต นิกฺขมนฺเตน ภิกฺขุํ ปสฺสิตฺวา อาจิกฺขิตพฺพํ “อมุตฺร มยา ทฺวตฺติปตฺตปูรา ปฏิคฺคหิตา มา โข ตตฺถ ปฏิคฺคณฺหี”ติ.\csromanspacer{} สเจ ปสฺสิตฺวา น อาจิกฺขติ, อาปตฺติ ทุกฺกฏสฺส.\csromanspacer{} สเจ อาจิกฺขิเต ปฏิคฺคณฺหาติ, อาปตฺติ ทุกฺกฏสฺส.}

\prose{\textbf{ตโต นีหริตฺวา ภิกฺขูหิ สทฺธิํ สํวิภชิตพฺพนฺ}ติ ปฏิกฺกมนํ นีหริตฺวา สํวิภชิตพฺพํ.}

\prose{\textbf{อยํ ตตฺถ สามีจี}ติ อยํ ตตฺถ อนุธมฺมตา.}

\proseitem{234}{อติเรกทฺวตฺติปตฺตปูเร อติเรกสญฺญี ปฏิคฺคณฺหาติ, อาปตฺติ ปาจิตฺติยสฺส.\csromanspacer{} อติเรกทฺวตฺติปตฺตปูเร เวมติโก ปฏิคฺคณฺหาติ, อาปตฺติ ปาจิตฺติยสฺส.\csromanspacer{} อติเรกทฺวตฺติปตฺตปูเร อูนกสญฺญี ปฏิคฺคณฺหาติ, อาปตฺติ}

\prose{อูนกทฺวตฺติปตฺตปูเร อติเรกสญฺญี, อาปตฺติ ทุกฺกฏสฺส.\csromanspacer{} อูนกทฺวตฺติปตฺตปูเร เวมติโก, อาปตฺติ ทุกฺกฏสฺส.\csromanspacer{} อูนกทฺวตฺติปตฺตปูเร อูนกสญฺญี, อนาปตฺติ.}

\proseitemwithcloserruled{235}{อนาปตฺติ ทฺวตฺติปตฺตปูเร ปฏิคฺคณฺหาติ, อูนกทฺวตฺติปตฺตปูเร ปฏิคฺคณฺหาติ, น ปเหณกตฺถาย น ปาเถยฺยตฺถาย ปฏิยตฺตํ เทนฺติ, ปเหณกตฺถาย วา ปาเถยฺยตฺถาย วา ปฏิยตฺตเสสกํ เทนฺติ, คมเน ปฏิปฺปสฺสทฺเธ เทนฺติ, ญาตกานํ ปวาริตานํ อญฺญสฺสตฺถาย อตฺตโน ธเนน อุมฺมตฺตกสฺส อาทิกมฺมิกสฺสาติ.}{กาณมาตุสิกฺขาปทํ นิฏฺฐิตํ จตุตฺถํ.}
\csromanmatikaanchor{mtk.40}
\csromantocmarkref{subsection}{5. ปฐมปวารณาสิกฺขาปท}{mtk.40}
\bookmark[dest={mtk.40},level=2]{5. ปฐมปวารณาสิกฺขาปท}
\csromanheader{2}{4. โภชนวคฺค 5. ปฐม\-ปวารณา\-สิกฺขาปท}
\proseitem{236}{เตน สมเยน พุทฺโธ ภควา สาวตฺถิยํ วิหรติ เชตวเน อนาถปิณฺฑิกสฺส อาราเม.\csromanspacer{} เตน โข ปน สมเยน อญฺญตโร พฺราหฺมโณ ภิกฺขู นิมนฺเตตฺวา โภเชสิ.\csromanspacer{} ภิกฺขู ภุตฺตาวี ปวาริตา \csromanfolio{110}\textbf{ญาติกุลานิ คนฺตฺวา เอกจฺเจ ภุญฺชิํสุ, เอกจฺเจ ปิณฺฑปาตํ อาทาย อคมํสุ.}\csromanspacer{} \textbf{อถ โข โส พฺราหฺมโณ ปฏิวิสฺสเก}\footnote{ปฏิวิสเก (โยชนา)}\textbf{ เอตทโวจ “ภิกฺขู มยา อยฺยา} \textbf{สนฺตปฺปิตา, เอถ ตุเมฺหปิ สนฺตปฺเปสฺสามี”ติ.}\csromanspacer{}\textbf{ เต เอวมาหํสุ “กิํ ตฺวํ} อยฺโย\footnote{อยฺย (สฺยา)} อเมฺห สนฺตปฺเปสฺสสิ, เยปิ ตยา นิมนฺติตา, เตปิ อมฺหากํ ฆรานิ อาคนฺตฺวา เอกจฺเจ ภุญฺชิํสุ, เอกจฺเจ ปิณฺฑปาตํ อาทาย อคมํสู”ติ.}

\prose{อถ โข โส พฺราหฺมโณ อุชฺฌายติ ขิยฺยติ วิปาเจติ “กถํ หิ นาม ภทนฺตา อมฺหากํ ฆเร ภุญฺชิตฺวา อญฺญตฺร ภุญฺชิสฺสนฺติ, น จาหํ ปฏิพโล ยาวทตฺถํ ทาตุนฺ”ติ.\csromanspacer{} อสฺโสสุํ โข ภิกฺขู ตสฺส พฺราหฺมณสฺส อุชฺฌายนฺตสฺส ขิยฺยนฺตสฺส วิปาเจนฺตสฺส.\csromanspacer{} เย เต ภิกฺขู อปฺปิจฺฉา ฯเปฯ เต อุชฺฌายนฺติ ขิยฺยนฺติ วิปาเจนฺติ “กถํ หิ นาม ภิกฺขู ภุตฺตาวี ปวาริตา อญฺญตฺร ภุญฺชิสฺสนฺตี”ติ ฯเปฯ.\csromanspacer{} สจฺจํ กิร ภิกฺขเว ภิกฺขู ภุตฺตาวี ปวาริตา อญฺญตฺร ภุญฺชนฺตีติ.\csromanspacer{} สจฺจํ ภควาติ.\csromanspacer{} วิครหิ พุทฺโธ ภควา ฯเปฯ.\csromanspacer{} กถํ หิ นาม เต ภิกฺขเว โมฆปุริสา ภุตฺตาวี ปวาริตา อญฺญตฺร ภุญฺชิสฺสนฺติ.\csromanspacer{} เนตํ ภิกฺขเว อปฺปสนฺนานํ วา ปสาทาย ฯเปฯ.\csromanspacer{} เอวญฺจ ปน ภิกฺขเว อิมํ สิกฺขาปทํ อุทฺทิเสยฺยาถ\csromandash{}}

\prose{\textbf{“โย ปน ภิกฺขุ ภุตฺตาวี ปวาริโต ขาทนียํ วา โภชนียํ วา ขาเทยฺย วา ภุญฺเชยฺย วา, ปาจิตฺติยนฺ”}ติ.}

\prose{เอวญฺจิทํ ภควตา ภิกฺขูนํ สิกฺขาปทํ ปญฺญตฺตํ โหติ.}

\proseitem{237}{เตน โข ปน สมเยน ภิกฺขู คิลานานํ ภิกฺขูนํ ปณีเต ปิณฺฑปาเต นีหรนฺติ, คิลานา น จิตฺตรูปํ ภุญฺชนฺติ, ตานิ ภิกฺขู ฉฏฺเฏนฺติ.\csromanspacer{} อสฺโสสิ โข ภควา อุจฺจาสทฺทํ มหาสทฺทํ กาโกรวสทฺทํ, สุตฺวาน อายสฺมนฺตํ อานนฺทํ อามนฺเตสิ “กิํ นุ โข โส อานนฺท อุจฺจาสทฺโท มหาสทฺโท กาโกรวสทฺโท”ติ.\csromanspacer{} อถ โข อายสฺมา อานนฺโท ภควโต เอตมตฺถํ อาโรเจสิ.\csromanspacer{} ภุญฺเชยฺยุํ ปนานนฺท ภิกฺขู คิลานาติริตฺตนฺติ.\csromanspacer{} น ภุญฺเชยฺยุํ ภควาติ.\csromanspacer{} อถ โข ภควา เอตสฺมิํ นิทาเน เอตสฺมิํ ปกรเณ ธมฺมิํ กถํ กตฺวา ภิกฺขู อามนฺเตสิ อนุชานามิ ภิกฺขเว คิลานสฺส จ อคิลานสฺส จ อติริตฺตํ ภุญฺชิตุํ.\csromanspacer{} เอวญฺจ ปน ภิกฺขเว อติริตฺตํ กาตพฺพํ “อลเมตํ สพฺพนฺ”ติ.\csromanspacer{} เอวญฺจ ปน ภิกฺขเว อิมํ สิกฺขาปทํ อุทฺทิเสยฺยาถ\csromandash{}}

\proseitem{238}{\csromanfolio{111}\textbf{“โย ปน ภิกฺขุ ภุตฺตาวี ปวาริโต อนติริตฺตํ ขาทนียํ วา โภชนียํ วา ขาเทยฺย วา ภุญฺเชยฺย วา, ปาจิตฺติยนฺ”}ติ.}

\proseitem{239}{\textbf{โย ปนา}ติ โย ยาทิโส ฯเปฯ.\csromanspacer{} \textbf{ภิกฺขู}ติ ฯเปฯ อยํ อิมสฺมิํ อตฺเถ อธิปฺเปโต ภิกฺขูติ.}

\prose{\textbf{ภุตฺตาวี} นาม ปญฺจนฺนํ โภชนานํ อญฺญตรํ โภชนํ อนฺตมโส กุสคฺเคนปิ ภุตฺตํ โหติ.}

\prose{\textbf{ปวาริโต} นาม อสนํ ปญฺญายติ, โภชนํ ปญฺญายติ, หตฺถปาเส ฐิโต อภิหรติ, ปฏิกฺเขโป ปญฺญายติ.}

\prose{\textbf{อนติริตฺตํ} นาม อกปฺปิยกตํ โหติ, อปฺปฏิคฺคหิตกตํ โหติ, อนุจฺจาริตกตํ โหติ, อหตฺถปาเส กตํ โหติ, อภุตฺตาวินา กตํ โหติ, ภุตฺตาวินา ปวาริเตน อาสนา วุฏฺฐิเตน กตํ โหติ, “อลเมตํ สพฺพนฺ”ติ อวุตฺตํ โหติ, น คิลานาติริตฺตํ โหติ, เอตํ อนติริตฺตํ นาม.}

\prose{\textbf{อติริตฺตํ} นาม กปฺปิยกตํ โหติ, ปฏิคฺคหิตกตํ โหติ, อุจฺจาริตกตํ โหติ, หตฺถปาเส กตํ โหติ, ภุตฺตาวินา กตํ โหติ, ภุตฺตาวินา ปวาริเตน อาสนา อวุฏฺฐิเตน กตํ โหติ, “อลเมตํ สพฺพนฺ”ติ วุตฺตํ โหติ, คิลานาติริตฺตํ โหติ, เอตํ อติริตฺตํ นาม.}

\prose{\textbf{ขาทนียํ} นาม ปญฺจ โภชนานิ ยามกาลิกํ สตฺตาหกาลิกํ ยาวชีวิกํ ฐเปตฺวา อวเสสํ ขาทนียํ นาม.}

\prose{\textbf{โภชนียํ} นาม ปญฺจ โภชนานิ โอทโน กุมฺมาโส สตฺถุ มจฺโฉ มํสํ.}

\prose{“ขาทิสฺสามิ ภุญฺชิสฺสามี”ติ ปฏิคฺคณฺหาติ, อาปตฺติ ทุกฺกฏสฺส.\csromanspacer{} อชฺโฌหาเร อชฺโฌหาเร อาปตฺติ ปาจิตฺติยสฺส.}

\proseitem{240}{อนติริตฺเต อนติริตฺตสญฺญี ขาทนียํ วา โภชนียํ วา ขาทติ วา ภุญฺชติ วา, อาปตฺติ ปาจิตฺติยสฺส.\csromanspacer{} อนติริตฺเต เวมติโก ขาทนียํ วา โภชนียํ วา ขาทติ วา ภุญฺชติ วา, อาปตฺติ ปาจิตฺติยสฺส.\csromanspacer{} อนติริตฺเต อติริตฺตสญฺญี ขาทนียํ วา โภชนียํ วา ขาทติ วา ภุญฺชติ วา, อาปตฺติ}

\prose{\csromanfolio{112}ยามกาลิกํ สตฺตาหกาลิกํ ยาวชีวิกํ อาหารตฺถาย ปฏิคฺคณฺหาติ, อาปตฺติ ทุกฺกฏสฺส.\csromanspacer{} อชฺฌาหาเร อชฺฌาหาเร อาปตฺติ ทุกฺกฏสฺส.\csromanspacer{} อติริตฺเต อนติริตฺตสญฺญี, อาปตฺติ ทุกฺกฏสฺส.\csromanspacer{} อติริตฺเต เวมติโก, อาปตฺติ ทุกฺกฏสฺส.\csromanspacer{} อติริตฺเต อติริตฺตสญฺญี, อนาปตฺติ.}

\proseitemwithcloserruled{241}{อนาปตฺติ อติริตฺตํ การาเปตฺวา ภุญฺชติ, “อติริตฺตํ การาเปตฺวา ภุญฺชิสฺสามี”ติ ปฏิคฺคณฺหาติ, อญฺญสฺสตฺถาย หรนฺโต คจฺฉติ คิลานสฺส เสสกํ ภุญฺชติ, ยามกาลิกํ สตฺตาหกาลิกํ ยาวชีวิกํ สติ ปจฺจเย ปริภุญฺชติ, อุมฺมตฺตกสฺส อาทิกมฺมิกสฺสาติ.}{ปฐม\-ปวารณา\-สิกฺขาปทํ นิฏฺฐิตํ ปญฺจมํ.}
\csromanmatikaanchor{mtk.41}
\csromantocmarkref{subsection}{6. ทุติยปวารณาสิกฺขาปท}{mtk.41}
\bookmark[dest={mtk.41},level=2]{6. ทุติยปวารณาสิกฺขาปท}
\csromanheader{2}{4. โภชนวคฺค 6. ทุติย\-ปวารณา\-สิกฺขาปท}
\proseitem{242}{เตน สมเยน พุทฺโธ ภควา สาวตฺถิยํ วิหรติ เชตวเน อนาถปิณฺฑิกสฺส อาราเม.\csromanspacer{} เตน โข ปน สมเยน เทฺว ภิกฺขู โกสเลสุ ชนปเท สาวตฺถิํ อทฺธานมคฺคปฺปฏิปนฺนา โหนฺติ, เอโก ภิกฺขุ อนาจารํ อาจรติ, ทุติโย ภิกฺขุ ตํ ภิกฺขุํ เอตทโวจ “มาวุโส เอวรูปมกาสิ, เนตํ กปฺปตี”ติ, โส ตสฺมิํ อุปนนฺธิ.\csromanspacer{} อถ โข เต ภิกฺขู สาวตฺถิํ อคมํสุ.\csromanspacer{} เตน โข ปน สมเยน สาวตฺถิยํ อญฺญตรสฺส ปูคสฺส สํฆภตฺตํ โหติ, ทุติโย ภิกฺขุ ภุตฺตาวี ปวาริโต โหติ, อุปนทฺโธ\footnote{อุปนนฺโธ (ก)} ภิกฺขุ ญาติกุลํ คนฺตฺวา ปิณฺฑปาตํ อาทาย เยน โส ภิกฺขุ เตนุปสงฺกมิ, อุปสงฺกมิตฺวา ตํ ภิกฺขุํ เอตทโวจ “ภุญฺชาหิ อาวุโส”ติ.\csromanspacer{} อลํ อาวุโส ปริปุณฺโณมฺหีติ.\csromanspacer{} สุนฺทโร อาวุโส ปิณฺฑปาโต ภุญฺชาหีติ.\csromanspacer{} อถ โข โส ภิกฺขุ เตน ภิกฺขุนา นิปฺปีฬิยมาโน ตํ ปิณฺฑปาตํ ภุญฺชิ.\csromanspacer{} อุปนทฺโธ ภิกฺขุ ตํ ภิกฺขุํ เอตทโวจ “ตฺวมฺปิ\footnote{ตฺวํ หิ (สฺยา)} นาม อาวุโส มํ วตฺตพฺพํ มญฺญสิ, ยํ ตฺวํ ภุตฺตาวี ปวาริโต อนติริตฺตํ โภชนํ ภุญฺชสี”ติ.\csromanspacer{} นนุ อาวุโส อาจิกฺขิตพฺพนฺติ.\csromanspacer{} นนุ อาวุโส ปุจฺฉิตพฺพนฺติ.}

\prose{อถ โข โส ภิกฺขุ ภิกฺขูนํ เอตมตฺถํ อาโรเจสิ.\csromanspacer{} เย เต ภิกฺขู อปฺปิจฺฉา ฯเปฯ เต อุชฺฌายนฺติ ขิยฺยนฺติ วิปาเจนฺติ “กถํ หิ นาม \csromanfolio{113}ภิกฺขุ ภิกฺขุํ ภุตฺตาวิํ ปวาริตํ อนติริตฺเตน โภชเนน อภิหฏฺฐุํ ปเวเรสฺสตี”ติ ฯเปฯ.\csromanspacer{} สจฺจํ กิร ตฺวํ ภิกฺขุ ภิกฺขุํ ภุตฺตาวิํ ปวาริตํ อนติริตฺเตน โภชเนน อภิหฏฺฐุํ ปวาเรสีติ.\csromanspacer{} สจฺจํ ภควาติ.\csromanspacer{} วิครหิ พุทฺโธ ภควา ฯเปฯ.\csromanspacer{} กถํ หิ นาม ตฺวํ โมฆปุริส ภิกฺขุํ ภุตฺตาวิํ ปวาริตํ อนติริตฺเตน โภชเนน อภิหฏฺฐุํ ปวาเรสฺสสิ.\csromanspacer{} เนตํ โมฆปุริส อปฺปสนฺนานํ วา ปสาทาย ฯเปฯ.\csromanspacer{} เอวญฺจ ปน ภิกฺขเว อิมํ สิกฺขาปทํ อุทฺทิเสยฺยาถ\csromandash{}}

\proseitem{243}{\textbf{“โย ปน ภิกฺขุ ภิกฺขุํ ภุตฺตาวิํ ปวาริตํ อนติริตฺเตน ขาทนีเยน วา โภชนีเยน วา อภิหฏฺฐุํ ปวาเรยฺย ‘หนฺท ภิกฺขุ ขาท วา ภุญฺช วา’ติ ชานํ อาสาทนาเปกฺโข, ภุตฺตสฺมิํ ปาจิตฺติยนฺ”}ติ.}

\proseitem{244}{\textbf{โย ปนา}ติ โย ยาทิโส ฯเปฯ.\csromanspacer{} \textbf{ภิกฺขู}ติ ฯเปฯ อยํ อิมสฺมิํ อตฺเถ อธิปฺเปโต ภิกฺขูติ.}

\prose{\textbf{ภิกฺขุนฺ}ติ อญฺญํ ภิกฺขุํ.}

\prose{\textbf{ภุตฺตาวี} นาม ปญฺจนฺนํ โภชนานํ อญฺญตรํ โภชนํ อนฺตมโส กุสคฺเคนปิ ภุตฺตํ โหติ.}

\prose{\textbf{ปวาริโต} นาม อสนํ ปญฺญายติ, โภชนํ ปญฺญายติ, หตฺถปาเส ฐิโต อภิหรติ, ปฏิกฺเขโป ปญฺญายติ.}

\prose{\textbf{อนติริตฺตํ} นาม อกปฺปิยกตํ โหติ, อปฺปฏิคฺคหิตกตํ โหติ, อนุจฺจาริตกตํ โหติ, อหตฺถปาเส กตํ โหติ, อภุตฺตาวินา กตํ โหติ, ภุตฺตาวินา ปวาริเตน อาสนา วุฏฺฐิเตน กตํ โหติ, “อลเมตํ สพฺพนฺ”ติ อวุตฺตํ โหติ, น คิลานาติริตฺตํ โหติ, เอตํ อนติริตฺตํ นาม.}

\prose{\textbf{ขาทนียํ} นาม ปญฺจ โภชนานิ ยามกาลิกํ สตฺตาหกาลิกํ ยาวชีวิกํ ฐเปตฺวา อวเสสํ ขาทนียํ นาม.}

\prose{\textbf{โภชนิยํ} นาม ปญฺจ โภชนานิ โอทโน กุมฺมาโส สตฺตุ มจฺโฉ มํสํ.}

\prose{\textbf{อภิหฏฺฐุํ ปจาเรยฺยา}ติ ยาวตกํ อิจฺฉสิ, ตาวตกํ คณฺหาหีติ.}

\prose{\csromanfolio{114}\textbf{ชานาติ} นาม สามํ วา ชานาติ, อญฺเญ วา ตสฺส อาโรเจนฺติ, โส วา อาโรเจติ.}

\prose{\textbf{อาสาทนาเปกฺโข}ติ “อิมินา อิมํ โจเทสฺสามิ สาเรสฺสามิ, ปฏิโจเทสฺสามิ ปฏิสาเรสฺสามิ, มงฺกุํ กริสฺสามี”ติ อภิหรติ, อาปตฺติ ทุกฺกฏสฺส.\csromanspacer{} ตสฺส วจเนน “ขาทิสฺสามิ ภุญฺชิสฺสามี”ติ ปฏิคฺคณฺหาติ.\csromanspacer{} อาปตฺติ ทุกฺกฏสฺส.\csromanspacer{} อชฺโฌหาเร อชฺโฌหาเร อาปตฺติ ทุกฺกฏสฺส.\csromanspacer{} โภชนปริโยสาเน อาปตฺติ}

\proseitem{245}{ปวาริเต ปวาริตสญฺญี อนติริตฺเตน ขาทนีเยน วา โภชนีเยน วา อภิหฏฺฐุํ ปวาเรติ, อาปตฺติ ปาจิตฺติยสฺส.\csromanspacer{} ปวาริเต เวมติโก อนติริตฺเตน ขาทนีเยน วา โภชนีเยน วา อภิหฏฺฐุํ ปวาเรติ, อาปตฺติ ทุกฺกฏสฺส.\csromanspacer{} ปวาริเต อปฺปวาริตสญฺญี อนติริตฺเตน ขาทนีเยน วา โภชนีเยน วา อภิหฏฺฐุํ ปวาเรติ, อนาปตฺติ.\csromanspacer{} ยามกาลิกํ สตฺตาหกาลิกํ ยาวชีวิกํ อาหารตฺถาย อภิหรติ, อาปตฺติ ทุกฺกฏสฺส.\csromanspacer{} ตสฺส วจเนน “ขาทิสฺสามิ ภุญฺชิสฺสามี”ติ ปฏิคฺคณฺหาติ, อาปตฺติ ทุกฺกฏสฺส.\csromanspacer{} อชฺโฌหาเร อชฺโฌหาเร อาปตฺติ ทุกฺกฏสฺส.\csromanspacer{} อปฺปวาริเต ปวาริตสญฺญี, อาปตฺติ ทุกฺกฏสฺส.\csromanspacer{} อปฺปวาริเต เวมติโก, อาปตฺติ ทุกฺกฏสฺส.\csromanspacer{} อปฺปวาริเต อปฺปวาริตสญฺญี, อนาปตฺติ.}

\proseitemwithcloserruled{246}{อนาปตฺติ อติริตฺตํ การาเปตฺวา เทติ, “อติริตฺตํ การาเปตฺวา ภุญฺชาหี”ติ เทติ, อญฺญสฺสตฺถาย หรนฺโต คจฺฉาหีติ เทติ, คิลานสฺส เสสกํ เทติ, ยามกาลิกํ สตฺตาหกาลิกํ ยาวชีวิกํ สติ ปจฺจเย ปริภุญฺชาติ เทติ, อุมฺมตฺตกสฺส อาทิกมฺมิกสฺสาติ.}{ทุติย\-ปวารณา\-สิกฺขาปทํ นิฏฺฐิตํ ฉฏฺฐํ.}
\csromanmatikaanchor{mtk.42}
\csromantocmarkref{subsection}{7. วิกาลโภชนสิกฺขาปท}{mtk.42}
\bookmark[dest={mtk.42},level=2]{7. วิกาลโภชนสิกฺขาปท}
\csromanheader{2}{4. \textbf{โภชนวคฺค 7. วิกาลโภชนสิกฺขาปท}}
\proseitem{247}{เตน สมเยน พุทฺโธ ภควา ราชคเห วิหรติ เวฬุวเน กลนฺทกนิวาเป.\csromanspacer{} เตน โข ปน สมเยน ราชคเห คิรคฺคสมชฺโช โหติ.\csromanspacer{} สตฺตรสวคฺคิยา ภิกฺขู คิรคฺคสมชฺชํ ทสฺสนาย อคมํสุ.\csromanspacer{} มนุสฺสา สตฺตรสวคฺคิเย ภิกฺขู ปสฺสิตฺวา นหาเปตฺวา วิลิมฺเปตฺวา โภเชตฺวา \csromanfolio{115}ขาทนียํ อทํสุ.\csromanspacer{} สตฺตรสวคฺคิยา ภิกฺขู ขาทนียํ อาทาย อารามํ คนฺตฺวา ฉพฺพคฺคิเย ภิกฺขู เอตทโวจุํ “คณฺหาถาวุโส, ขาทนียํ ขาทถา”ติ.\csromanspacer{} กุโต ตุเมฺหหิ อาวุโส ขาทนียํ ลทฺธนฺติ.\csromanspacer{} สตฺตรสวคฺคิยา ภิกฺขู ฉพฺพคฺคิยานํ ภิกฺขูนํ เอตมตฺถํ อาโรเจสุํ.\csromanspacer{} กิํ ปน ตุเมฺห อาวุโส วิกาเล โภชนํ ภุญฺชถาติ.\csromanspacer{} เอวมาวุโสติ.\csromanspacer{} ฉพฺพคฺคิยา ภิกฺขู อุชฺฌายนฺติ ขิยฺยนฺติ วิปาเจนฺติ “กถํ หิ นาม สตฺตรสวคฺคิยา ภิกฺขู วิกาเล โภชนํ ภุญฺชิสฺสนฺตี”ติ.\csromanspacer{} อถ โข ฉพฺพคฺคิยา ภิกฺขู ภิกฺขูนํ เอตมตฺถํ อาโรเจสุํ.\csromanspacer{} เย เต ภิกฺขู อปฺปิจฺฉา ฯเปฯ เต อุชฺฌายนฺติ ขิยฺยนฺติ วิปาเจนฺติ “กถํ หิ นาม สตฺตรสวคฺคิยา ภิกฺขู วิกาเล โภชนํ ภุญฺชิสฺสนฺตี”ติ ฯเปฯ.\csromanspacer{} สจฺจํ กิร ตุเมฺห ภิกฺขเว วิกาเล โภชนํ ภุญฺชถาติ.\csromanspacer{} สจฺจํ ภควาติ.\csromanspacer{} วิครติ พุทฺโธ ภควา ฯเปฯ.\csromanspacer{} กถํ หิ นาม ตุเมฺห โมฆปุริสา วิกาเล โภชนํ ภุญฺชิสฺสถ.\csromanspacer{} เนตํ \textbf{โมฆปุริสา อปฺปสนฺนานํ วา ปสาทาย ฯเปฯ.}\csromanspacer{}\textbf{ เอวญฺจ ปน ภิกฺขเว อิมํ} \textbf{สิกฺขาปทํ อุทฺทิเสยฺยาถ\csromandash{}}}

\proseitem{248}{\textbf{“โย ปน ภิกฺขุ วิกาเล ขาทนียํ วา โภชนียํ วา ขาเทยฺย วา ภุญฺเชยฺย วา, ปาจิตฺติยนฺ”}ติ.}

\proseitem{249}{\textbf{โย ปนา}ติ โย ยาทิโส ฯเปฯ.\csromanspacer{} \textbf{ภิกฺขู}ติ ฯเปฯ อยํ อิมสฺมิํ อตฺเถ อธิปฺเปโต ภิกฺขูติ.}

\prose{\textbf{วิกาโล} นาม มชฺฌนฺหิเก วีติตตฺเต ยาว อรุณุคฺคมนา.}

\prose{\textbf{ขาทนียํ} นาม ปญฺจ โภชนานิ ยามกาลิกํ สตฺตาหกาลิกํ ยาวชีวิกํ ฐเปตฺวา อวเสสํ ขาทนียํ นาม.}

\prose{\textbf{โภชนิยํ} นาม ปญฺจ โภชนานิ โอทโน กุมฺมาโส สตฺตุ มจฺโฉ มํสํ.}

\prose{“ขาทิสฺสามิ ภุญฺชิสฺสามี”ติ ปฏิคฺคณฺหาติ, อาปตฺติ ทุกฺกฏสฺส.\csromanspacer{} อชฺโฌหาเร อชฺโฌหาเร อาปตฺติ ปาจิตฺติยสฺส.}

\proseitem{250}{วิกาเล วิกาลสญฺญี ขาทนียํ วา โภชนียํ วา ขาทติ วา ภุญฺชติ วา, อาปตฺติ ปาจิตฺติยสฺส.\csromanspacer{} วิกาเล เวมติโก ขาทนียํ วา โภชนียํ วา ขาทติ วา ภุญฺชติ วา, อาปตฺติ ปาจิตฺติยสฺส. \csromanfolio{116}วิกาเล กาลสญฺญี ขาทนียํ วา โภชนียํ วา ขาทติ วา ภุญฺชติ วา, อาปตฺติ ปาจิตฺติยสฺส.}

\prose{ยามกาลิกํ สตฺตาหกาลิกํ ยาวชีวิกํ อาหารตฺถาย ปฏิคฺคณฺหาติ, อาปตฺติ ทุกฺกฏสฺส.\csromanspacer{} อชฺโฌหาเร อชฺโฌหาเร อาปตฺติ ทุกฺกฏสฺส.\csromanspacer{} กาเล วิกาลสญฺญี, อาปตฺติ ทุกฺกฏสฺส.\csromanspacer{} กาเล เวมติโก, อาปตฺติ ทุกฺกฏสฺส.\csromanspacer{} กาเล กาลสญฺญี, อนาปตฺติ.}

\proseitemwithcloserruled{251}{อนาปตฺติ ยามกาลิกํ สตฺตาหกาลิกํ ยาวชีวิกํ สติ ปจฺจเย ปริภุญฺชติ, อุมฺมตฺตกสฺส อาทิกมฺมิกสฺสาติ.}{วิกาลโภชนสิกฺขาปทํ นิฏฺฐิตํ สตฺตมํ.}
\csromanmatikaanchor{mtk.43}
\csromantocmarkref{subsection}{8. สนฺนิธิการกสิกฺขาปท}{mtk.43}
\bookmark[dest={mtk.43},level=2]{8. สนฺนิธิการกสิกฺขาปท}
\csromanheader{2}{4. \textbf{โภชนวคฺค 8. สนฺนิธิการกสิกฺขาปท}}
\proseitem{255}{เตน สมเยน พุทฺโธ ภควา สาวตฺถิยํ วิหรติ เชตวเน อนาถปิณฺฑิกสฺส อาราเม.\csromanspacer{} เตน โข ปน สมเยน อายสฺมโต อานนฺทสฺส อุปชฺฌาโย อายสฺมา เพลฏฺฐสีโส\footnote{เพลฏฺฐิสีโส (สี), เวฬฏฺฐสีโส (สฺยา)} อรญฺเญ วิหรติ, โส ปิณฺฑาย จริตฺวา สุกฺขกุรํ อารามํ หริตฺวา สุกฺขาเปตฺวา นิกฺขิปติ, ยทา อาหาเรน อตฺโถ โหติ, ตทา อุทเกน เตเมตฺวา เตเมตฺวา ภุญฺชติ, จิเรน คามํ ปิณฺฑาย ปวิสติ.\csromanspacer{} ภิกฺขู อายสฺมนฺตํ เพลฏฺฐสีสํ เอตทโวจุํ “กิสฺส ตฺวํ อาวุโส จิเรน คามํ ปิณฺฑาย ปวิสสี”ติ.\csromanspacer{} อถ โข อายสฺมา เพลฏฺฐสีโส ภิกฺขูนํ เอตมตฺถํ อาโรเจสิ.\csromanspacer{} กิํ ปน ตฺวํ อาวุโส สนฺนิธิการกํ โภชนํ ภุญฺชสีติ.\csromanspacer{} เอวมาวุโสติ.\csromanspacer{} เย เต ภิกฺขู อปฺปิจฺฉา ฯเปฯ เต อุชฺฌายนฺติ ขิยฺยนฺติ วิปาเจนฺติ “กถํ หิ นาม อายสฺมา เพลฏฺฐสีโส สนฺนิธิการกํ โภชนํ ภุญฺชิสฺสตี”ติ ฯเปฯ.\csromanspacer{} สจฺจํ กิร ตฺวํ เพลฏฺฐสีส สนฺนิธิการกํ โภชนํ ภุญฺชสีติ.\csromanspacer{} สจฺจํ ภควาติ.\csromanspacer{} วิครหิ พุทฺโธ ภควา ฯเปฯ.\csromanspacer{} กถํ หิ นาม ตฺวํ เพลฏฺฐสีส สนฺนิธิการกํ โภชนํ ภุญฺชิสฺสสิ.\csromanspacer{} เนตํ เพลฏฺฐสีส อปฺปสนฺนานํ วา ปสาทาย ฯเปฯ.\csromanspacer{} เอวญฺจ ปน ภิกฺขเว อิมํ สิกฺขาปทํ อุทฺทิเสยฺยาถ\csromandash{}}

\proseitem{253}{\textbf{“โย ปน ภิกฺขุ สนฺนิธิการกํ ขาทนียํ วา โภชนียํ วา ขาเทยฺย วา ภุญฺเชยฺย วา, ปาจิตฺติยนฺ”}ติ.}

\proseitem{254}{\csromanfolio{117}\textbf{โย ปนา}ติ โย ยาทิโส ฯเปฯ.\csromanspacer{} \textbf{ภิกฺขู}ติ ฯเปฯ อยํ อิมสฺมิํ อตฺเถ อธิปฺเปโต ภิกฺขูติ.}

\prose{\textbf{สนฺนิธิการกํ} นาม อชฺช ปฏิคฺคหิตํ อปรชฺชุ ขาทิตํ โหติ.}

\prose{\textbf{ขาทนียํ} นาม ปญฺจ โภชนานิ ยามกาลิกํ สตฺตาหกาลิกํ ยาวชีวิกํ ฐเปตฺวา อวเสสํ ขาทนียํ นาม.}

\prose{\textbf{โภชนียํ} นาม ปญฺจ โภชนานิ โอทโน กุมฺมาโส สตฺตุ มจฺโฉ มํสํ.}

\prose{“ขาทิสฺสามิ ภุญฺชิสฺสามี”ติ ปฏิคฺคณฺหาติ, อาปตฺติ ทุกฺกฏสฺส.\csromanspacer{} อชฺโฌหาเร อชฺโฌหาเร อาปตฺติ ปาจิตฺติยสฺส.}

\proseitem{255}{สนฺนิธิการเก สนฺนิธิการกสญฺญี ขาทนียํ วา โภชนียํ วา ขาทติ วา ภุญฺชติ วา, อาปตฺติ ปาจิตฺติยสฺส.\csromanspacer{} สนฺนิธิการเก เวมติโก ขาทนียํ วา โภชนียํ วา ขาทติ วา ภุญฺชติ วา อาปตฺติ ปาจิตฺติยสฺส.\csromanspacer{} สนฺนิธิการเก อสนฺนิธิการกสญฺญี ขาทนียํ วา โภชนียํ วา ขาทติ วา ภุญฺชติ วา, อาปตฺติ ปาจิตฺติยสฺส.}

\prose{ยามกาลิกํ สตฺตาหกาลิกํ ยาวชีวิกํ อาหารตฺถาย ปฏิคฺคณฺหาติ, อาปตฺติ ทุกฺกฏสฺส.\csromanspacer{} อชฺโฌหาเร อชฺโฌหาเร อาปตฺติ ทุกฺกฏสฺส.\csromanspacer{} อสนฺนิธิการเก สนฺนิธิการกสญฺญี, อาปตฺติ ทุกฺกฏสฺส.\csromanspacer{} อสนฺนิธิการเก เวมติโก, อาปตฺติ ทุกฺกฏสฺส.\csromanspacer{} อสนฺนิธิการเก อสนฺนิธิการกสญฺญี, อนาปตฺติ.}

\proseitemwithcloser{256}{อนาปตฺติ ยาวกาลิกํ ยาวกาเล นิทหิตฺวา ภุญฺชติ, ยามกาลิกํ ยาเม นิทหิตฺวา ภุญฺชติ, สตฺตาหกาลิกํ สตฺตาหํ นิทหิตฺวา ภุญฺชติ, ยาวชีวิกํ สติ ปจฺจเย ปริภุญฺชติ, อุมฺมตฺตกสฺส อาทิกมฺมิกสฺสาติ.}{สนฺนิธิการกสิกฺขาปทํ นิฏฺฐิตํ อฏฺฐมํ.}
\csromanmatikaanchor{mtk.44}
\csromantocmarkref{subsection}{9. ปณีตโภชนสิกฺขาปท}{mtk.44}
\bookmark[dest={mtk.44},level=2]{9. ปณีตโภชนสิกฺขาปท}
\csromanheader{2}{4. \textbf{โภชนวคฺค 9. ปณีตโภชนสิกฺขาปท}}
\proseitem{257}{\csromanfolio{118}เตน สมเยน พุทฺโธ ภควา สาวตฺถิยํ วิหรติ เชตวเน อนาถปิณฺฑิกสฺส อาราเม.\csromanspacer{} เตน โข ปน สมเยน ฉพฺพคฺคิยา ภิกฺขู ปณีตโภชนานิ อตฺตโน อตฺถาย วิญฺญาเปตฺวา ภุญฺชนฺติ.\csromanspacer{} มนุสฺสา อุชฺฌายนฺติ ขิยฺยนฺติ วิปาเจนฺติ “กถํ หิ นาม สมณา สกฺยปุตฺติยา ปณีตโภชนานิ อตฺตโน อตฺถาย วิญฺญาเปตฺวา ภุญฺชิสฺสนฺติ, กสฺส สมฺปนฺนํ น มนาปํ, กสฺส สาทุํ น รุจฺจตี”ติ.\csromanspacer{} อสฺโสสุํ โข ภิกฺขู เตสํ มนุสฺสานํ อุชฺฌายนฺตานํ ขิยฺยนฺตานํ วิปาเจนฺตานํ.\csromanspacer{} เย เต ภิกฺขู อปฺปิจฺฉา ฯเปฯ เต อุชฺฌายนฺติ ขิยฺยนฺติ วิปาเจนฺติ “กถํ หิ นาม ฉพฺพคฺคิยา ภิกฺขู ปณีตโภชนานิ อตฺตโน อตฺถาย วิญฺญาเปตฺวา ภุญฺชิสฺสนฺตี”ติ ฯเปฯ.\csromanspacer{} สจฺจํ กิร ตุเมฺห ภิกฺขเว ปณีตโภชนานิ อตฺตโน อตฺถาย วิญฺญาเปตฺวา ภุญฺชถาติ.\csromanspacer{} สจฺจํ ภควาติ.\csromanspacer{} วิครหิ พุทฺโธ ภควา ฯเปฯ.\csromanspacer{} กถํ หิ นาม ตุเมฺห โมฆปุริสา ปณีตโภชนานิ อตฺตโน อตฺถาย วิญฺญาเปตฺวา ภุญฺชิสฺสถ.\csromanspacer{} เนตํ โมฆปุริสา อปฺปสนฺนานํ วา ปสาทาย ฯเปฯ.\csromanspacer{} เอวญฺจ ปน ภิกฺขเว อิมํ สิกฺขาปทํ อุทฺทิเสยฺยาถ\csromandash{}}

\prose{\textbf{“ยานิ โข ปน ตานิ ปณีตโภชนานิ, เสยฺยถิทํ, สปฺปิ นวนีตํ เตลํ มธุ ผาณิตํ มจฺโฉ มํสํ ขีรํ ทธิ.}\csromanspacer{} โย ปน ภิกฺขุ เอวรูปานิ ปณีตโภชนานิ อตฺตโน อตฺถาย วิญฺญาเปตฺวา ภุญฺเชยฺย, ปาจิตฺติยนฺ”ติ.}

\prose{เอวญฺจิทํ ภควตา ภิกฺขูนํ สิกฺขาปทํ ปญฺญตฺตํ โหติ.}

\proseitem{258}{เตน โข ปน สมเยน ภิกฺขู คิลานา โหนฺติ.\csromanspacer{} คิลานปุจฺฉกา ภิกฺขู คิลาเน ภิกฺขู เอตทโวจุํ “กจฺจาวุโส ขมนียํ, กจฺจิ ยาปนียนฺ”ติ.\csromanspacer{} ปุพฺเพ มยํ อาวุโส ปณีตโภชนานิ อตฺตโน อตฺถาย วิญฺญาเปตฺวา ภุญฺชาม, เตน โน ผาสุ โหติ, อิทานิ ปน “ภควตา ปฏิกฺขิตฺตนฺ”ติ กุกฺกุจฺจายนฺตา น วิญฺญาเปม, เตน โน น ผาสุ โหตีติ.\csromanspacer{} ภควโต เอตมตฺถํ อาโรเจสุํ.\csromanspacer{} อถ โข ภควา เอตสฺมิํ นิทาเน เอตสฺมิํ ปกรเณ ธมฺมิํ กถํ กตฺวา ภิกฺขู อามนฺเตสิ อนุชานามิ ภิกฺขเว คิลาเนน ภิกฺขุนา ปณีตโภชนานิ อตฺตโน อตฺถาย วิญฺญาเปตฺวา ภุญฺชิตุํ.\csromanspacer{} เอวญฺจ ปน ภิกฺขเว อิมํ สิกฺขาปทํ อุทฺทิเสยฺยาถ\csromandash{}}

\proseitem{259}{\csromanfolio{119}“ยานิ โข ปน ตานิ ปณีตโภชนานิ, เสยฺยถิทํ, สปฺปิ นวนีตํ เตลํ มธุ ผาณิตํ มจฺโฉ มํสํ ขีรํ ทธิ.\csromanspacer{}\textbf{ โย ปน ภิกฺขุ เอวรูปานิ ปณีตโภชนานิ อคิลาโน อตฺตโน อตฺถาย วิญฺญาเปตฺวา ภุญฺเชยฺย, ปาจิตฺติยนฺ”}ติ.}

\proseitem{260}{\textbf{ยานิ โข ปน ตานิ ป}ณีตโภชนานีติ สปฺปิ นาม โคสปฺปิ วา อชิกาสปฺปิ วา มหิํสสปฺปิ วา.\csromanspacer{} เยสํ มํสํ กปฺปติ, เตสํ สปฺปิ.}

\prose{\textbf{นวนีตํ} นาม เตสญฺเญว นวนีตํ.}

\prose{\textbf{เตลํ} นาม ติลเตลํ สาสปเตลํ มธุกเตลํ เอรณฺฑเตลํ วสาเตลํ.}

\prose{\textbf{มธุ} นาม มกฺขิกามธุ.}

\prose{\textbf{ผาณิตํ} นาม อุจฺฉุมฺหา นิพฺพตฺตํ.}

\prose{\textbf{มจฺโฉ} นาม โอทโก\footnote{อุทกจโร (สฺยา, ก)} วุจฺจติ.}

\prose{\textbf{มํสํ} นาม เยสํ มํสํ กปฺปติ, เตสํ มํสํ.}

\prose{\textbf{ขีรํ} นาม โคขีรํ วา อชิกาขีรํ วา มหิํ สขีรํ วา, เยสํ มํสํ กปฺปติ, เตสํ ขีรํ.}

\prose{\textbf{ทธิ} นาม เตสญฺเญว ทธิ.}

\prose{\textbf{โย ปนา}ติ โย ยาทิโส ฯเปฯ.\csromanspacer{} \textbf{ภิกฺขู}ติ ฯเปฯ อยํ อิมสฺมิํ อตฺเถ อธิปฺเปโต ภิกฺขูติ.}

\prose{\textbf{เอวรูปานิ ปณีตโภชนานี}ติ ตถารูปานิ ปณีตโภชนานิ.}

\prose{\textbf{อคิลาโน} นาม ยสฺส วินา ปณีตโภชนานิ ผาสุ โหติ.}

\prose{\textbf{คิลาโน} นาม ยสฺส วินา ปณีตโภชนานิ น ผาสุ โหติ.}

\prose{อคิลาโน อตฺตโน อตฺถาย วิญฺญาเปติ, ปโยเค\footnote{ปโยเค ปโยเค (ก)} ทุกฺกฏํ.\csromanspacer{} ปฏิลาเภน “ภุญฺชิสฺสามี”ติ ปฏิคฺคณฺหาติ, อาปตฺติ ทุกฺกฏสฺส.\csromanspacer{} อชฺโฌหาเร อชฺโฌหาเร อาปตฺติ}

\proseitem{261}{อคิลาโน อคิลานสญฺญี ปณีตโภชนานิ อตฺตโน อตฺถาย วิญฺญาเปตฺวา ภุญฺชติ, อาปตฺติ ปาจิตฺติยสฺส.\csromanspacer{} อคิลาโน \csromanfolio{120}เวมติโก ปณีตโภชนานิ อตฺตโน อตฺถาย วิญฺญาเปตฺวา ภุญฺชติ, อาปตฺติ ปาจิตฺติยสฺส.\csromanspacer{} อคิลาโน คิลานสญฺญี ปณีตโภชนานิ อตฺตโน อตฺถาย วิญฺญาเปตฺวา ภุญฺชติ, อาปตฺติ ปาจิตฺติยสฺส.}

\prose{คิลาโน อคิลานสญฺญี, อาปตฺติ ทุกฺกฏสฺส.\csromanspacer{} คิลาโน เวมติโก, อาปตฺติ ทุกฺกฏสฺส.\csromanspacer{} คิลาโน คิลานสญฺญี, อนาปตฺติ.}

\proseitemwithcloserruled{262}{อนาปตฺติ คิลานสฺส, คิลาโน หุตฺวา วิญฺญาเปตฺวา อคิลาโน ภุญฺชติ, คิลานสฺส เสสกํ ภุญฺชติ, ญาตกานํ ปวาริตานํ อญฺญสฺสตฺถาย อตฺตโน ธเนน อุมฺมตฺตกสฺส อาทิกมฺมิกสฺสาติ.}{ปณีตโภชนสิกฺขปทํ นิฏฺฐิตํ นวมํ.}
\csromanmatikaanchor{mtk.45}
\csromantocmarkref{subsection}{10. ทนฺตโปนสิกฺขาปท}{mtk.45}
\bookmark[dest={mtk.45},level=2]{10. ทนฺตโปนสิกฺขาปท}
\csromanheader{2}{4. \textbf{โภชนวคฺค 1}0. ทนฺตโปนสิกฺขาปท}
\proseitem{263}{เตน สมเยน พุทฺโธ ภควา เวสาลิยํ วิหรติ มหาวเน กูฏาคารสาลายํ.\csromanspacer{} เตน โข ปน สมเยน อญฺญตโร ภิกฺขุ สพฺพปํสุกูลิโก สุสาเน วิหรติ.\csromanspacer{} โส มนุสฺเสหิ ทิยฺยมานํ น อิจฺฉติ ปฏิคฺคเหตุํ, สุสาเนปิ รุกฺขมูเลปิ อุมฺมาเรปิ อยฺยโวสาฏิตกานิ สามํ คเหตฺวา ปริภุญฺชติ\footnote{ภุญฺชติ (สฺยา)}.\csromanspacer{} มนุสฺสา อุชฺฌายนฺติ ขิยฺยนฺติ วิปาเจนฺติ “กถํ หิ นาม อยํ ภิกฺขุ อมฺหากํ อยฺยโวสาฏิตกานิ สามํ คเหตฺวา ปริภุญฺชิสฺสติ, เถโรยํ ภิกฺขุ วฐโร อนุสฺสมํสํ มญฺเญ ขาทตี”ติ.\csromanspacer{} อสฺโสสุํ โข ภิกฺขู เตสํ มนุสฺสานํ อุชฺฌายนฺตานํ ขิยฺยนฺตานํ วิปาเจนฺตานํ.\csromanspacer{} เย เต ภิกฺขู อปฺปิจฺฉา ฯเปฯ เต อุชฺฌายนฺติ ขิยฺยนฺติ วิปาเจนฺติ “กถํ หิ นาม ภิกฺขุ อทินฺนํ มุขทฺวารํ อาหารํ อาหริสฺสตี”ติ ฯเปฯ.\csromanspacer{} สจฺจํ กิร ตฺวํ ภิกฺขุ อทินฺนํ มุขทฺวารํ อาหารํ อาหรสีติ.\csromanspacer{} สจฺจํ ภควาติ.\csromanspacer{} วิครหิ พุทฺโธ ภควา ฯเปฯ.\csromanspacer{} กถํ หิ นาม ตฺวํ โมฆปุริส อทินฺนํ มุขทฺวารํ อาหารํ อาหริสฺสสิ.\csromanspacer{} เนตํ โมฆปุริส อปฺปสนฺนานํ วา ปสาทาย ฯเปฯ.\csromanspacer{} เอวญฺจ ปน ภิกฺขเว อิมํ สิกฺขาปทํ อุทฺทิเสยฺยาถ\csromandash{}}

\prose{\textbf{“โย ปน ภิกฺขุ อทินฺนํ มุขทฺวารํ อาหารํ อาหเรยฺย ปาจิตฺติยนฺ”}ติ.}

\prose{เอวญฺจิทํ ภควตา ภิกฺขูนํ สิกฺขาปทํ ปญฺญตฺตํ โหติ.}

\proseitem{264}{\csromanfolio{121}เตน โข ปน สมเยน ภิกฺขู อุทกทนฺตโปเน\footnote{อุทกทนฺตโปเณ (สฺยา, ก)} กุกฺกุจฺจายนฺติ.\csromanspacer{} ภควโต เอตมตฺถํ อาโรเจสุํ ฯเปฯ.\csromanspacer{} อนุชานามิ ภิกฺขเว อุทกทนฺตโปนํ สามํ คเหตฺวา ปริภุญฺชิตุํ.\csromanspacer{} เอวญฺจ ปน ภิกฺขเว อิมํ สิกฺขาปทํ อุทฺทิเสยฺยาถ\csromandash{}}

\proseitem{265}{\textbf{“โย ปน ภิกฺขุ อทินฺนํ มุขทฺวารํ อาหารํ อาหเรยฺย อญฺญตฺร อุทกทนฺตโปนา, ปาจิตฺติยนฺ”}ติ.}

\proseitem{266}{\textbf{โย ปนา}ติ โย ยาทิโส ฯเปฯ.\csromanspacer{} \textbf{ภิกฺขู}ติ ฯเปฯ อยํ อิมสฺมิํ อตฺเถ อธิปฺเปโต ภิกฺขูติ.}

\prose{\textbf{อทินฺนํ} นาม อปฺปฏิคฺคหิตกํ วุจฺจติ.}

\prose{\textbf{ทินฺนํ} นาม กาเยน วา กายปฏิพทฺเธน วา นิสฺสคฺคิเยน วา เทนฺเต หตฺถปาเส ฐิโต กาเยน วา กายปฏิพทฺเธน วา ปฏิคฺคณฺหาติ, เอตํ ทินฺนํ นาม.}

\prose{\textbf{อาหาโร} นาม อุทกทนฺตโปนํ ฐเปตฺวา ยํกิญฺจิ อชฺโฌหรณียํ, เอโส อาหาโร นาม.}

\prose{\textbf{อญฺญตฺร อุทกทนฺตโปนา}ติ ฐเปตฺวา อุทกทนฺตโปนํ.}

\prose{“ขาทิสฺสามิ ภุญฺชิสฺสามี”ติ คณฺหาติ, อาปตฺติ ทุกฺกฏสฺส.\csromanspacer{} อชฺโฌหาเร อชฺโฌหาเร อาปตฺติ ปาจิตฺติยสฺส.}

\proseitem{267}{อปฺปฏิคฺคหิตเก อปฺปฏิคฺคหิตกสญฺญี อทินฺนํ มุขทฺวารํ อาหารํ อาหาเรติ อญฺญตฺร อุทกทนฺตโปนา, อาปตฺติ ปาจิตฺติยสฺส.\csromanspacer{} อปฺปฏิคฺคหิตเก เวมติโก อทินฺนํ มุขทฺวารํ อาหารํ อาหาเรติ อญฺญตฺร อุทกทนฺตโปนา, อาปตฺติ ปาจิตฺติยสฺส.\csromanspacer{} อปฺปฏิคฺคหิตเก ปฏิคฺคหิตกสญฺญี อทินฺนํ มุขทฺวารํ อาหารํ อาหาเรติ อญฺญตฺร อุทกทนฺตโปนา, อาปตฺติ ปาจิตฺติยสฺส.}

\prose{ปฏิคฺคหิตเก อปฺปฏิคฺคหิตกสญฺญี, อาปตฺติ ทุกฺกฏสฺส.\csromanspacer{} ปฏิคฺคหิตเก เวมติโก, อาปตฺติ ทุกฺกฏสฺส.\csromanspacer{} ปฏิคฺคหิตเก ปฏิคฺคหิตกสญฺญี, อนาปตฺติ.}

\proseitemwithcloser{268}{\csromanfolio{122}อนาปตฺติ อุทกทนฺตโปเน, จตฺตาริ มหาวิกฏานิ สติ ปจฺจเย อสติ กปฺปิยการเก สามํ คเหตฺวา ปริภุญฺชติ, อุมฺมตฺตกสฺส อาทิกมฺมิกสฺสาติ.}{ทนฺตโปนสิกฺขาปทํ นิฏฺฐิตํ ทสมํ.}
\vspace{\tipitakaparskip}
\csromancenter{โภชนวคฺโค จตุตฺโถ.}
\csromancenter{ตสฺสุทฺทานํ}
\vspace{0.5\baselineskip}
\setlength{\csromangathapagewidth}{0pt}
\csromangathameasuregroup{ปิณฺโฑ คณํ ปรํ ปูวํ,}{\csromangathabat{ปิณฺโฑ คณํ ปรํ ปูวํ,}{เทฺว จ วุตฺตา ปวารณา.} \\ วิกาเล สนฺนิธี ขีรํ, ทนฺตโปเนน เต ทสาติ.}
\csromangathasetleft{ปิณฺโฑ คณํ ปรํ ปูวํ,}
\csromangathagroup{\csromangathastanza{\csromangathabat{ปิณฺโฑ คณํ ปรํ ปูวํ,}{เทฺว จ วุตฺตา ปวารณา.} \\ วิกาเล สนฺนิธี ขีรํ, ทนฺตโปเนน เต ทสาติ.}}

\csromanmatikaanchor{mtk.46}
\csromanmatikaanchor{mtk.47}
\csromantocmarknopage{section}{5. อเจลกวคฺค}{mtk.46}
\bookmark[dest={mtk.46},level=1]{5. อเจลกวคฺค}
\csromantocmarkref{subsection}{1. อเจลกสิกฺขาปท}{mtk.47}
\bookmark[dest={mtk.47},level=2]{1. อเจลกสิกฺขาปท}
\setcsromanheadh{5. อเจลกวคฺค}
\csromanheader{2}{5. \textbf{อเจลกวคฺค 1. อเจลกสิกฺขาปท}}
\proseitem{269}{เตน สมเยน พุทฺโธ ภควา เวสาลิยํ วิหรติ มหาวเน กูฏาคารสาลายํ.\csromanspacer{} เตน โข ปน สมเยน สํฆสฺส ขาทนียํ อุสฺสนฺนํ โหติ.\csromanspacer{} อถ โข อายสฺมา อานนฺโท ภควโต เอตมตฺถํ อาโรเจสิ.\csromanspacer{} เตนหานนฺท วิฆาสาทานํ ปูวํ เทหีติ. “เอวํ ภนฺเต”ติ โข อายสฺมา อานนฺโท ภควโต ปฏิสฺสุณิตฺวา วิฆาสาเท ปฏิปาฏิยา นิสีทาเปตฺวา เอเกกํ ปูวํ เทนฺโต อญฺญตริสฺสา ปริพฺพาชิกาย เอกํ มญฺญมาโน เทฺว ปูเว อทาสิ.\csromanspacer{} สามนฺตา ปริพฺพาชิกาโย ตํ ปริพฺพาชิกํ เอตทโวจุํ “ชาโร เต เอโส สมโณ”ติ.\csromanspacer{} น เม โส สมโณ ชาโร, เอกํ มญฺญมาโน เทฺว ปูเว อทาสีติ.\csromanspacer{} ทุติยมฺปิ โข ฯเปฯ.\csromanspacer{} ตติยมฺปิ โข อายสฺมา อานนฺโท เอเกกํ ปูวํ เทนฺโต ตสฺสา เยว ปริพฺพาชิกาย เอกํ มญฺญมาโน เทฺว ปูเว อทาสิ.\csromanspacer{} สามนฺตา ปริพฺพาชิกาโย ตํ ปริพฺพาชิกํ เอตทโวจุํ “ชาโร เต เอโส สมโณ”ติ.\csromanspacer{} น เม โส สมโณ ชาโร, เอกํ มญฺญมาโน เทฺว ปูเว อทาสีติ. “ชาโร น ชาโร”ติ ภณฺฑิํสุ.\csromanspacer{} อญฺญตโรปิ อาชีวโก ปริเวสนํ อคมาสิ.\csromanspacer{} อญฺญตโร ภิกฺขุ ปหูเตน สปฺปินา โอทนํ มทฺทิตฺวา ตสฺส อาชีวกสฺส มหนฺตํ ปิณฺฑํ อทาสิ.\csromanspacer{} อถ โข โส อาชีวโก ตํ ปิณฺฑํ อาทาย อคมาสิ.\csromanspacer{} อญฺญตโร อาชีวโก ตํ อาชีวกํ เอตทโวจ “กุโต ตยา อาวุโส ปิณฺโฑ ลทฺโธ”ติ. \csromanfolio{123}ตสฺสาวุโส สมณสฺส โคตมสฺส มุณฺฑคหปติกสฺส ปริเวสนาย \textbf{ลทฺโธติ.}}

\prose{อสฺโสสุํ โข อุปาสกา เตสํ อาชีวกานํ อิมํ กถาสลฺลาปํ.\csromanspacer{} อถ โข เต อุปาสกา เยน ภควา เตนุปสงฺกมิํสุ, อุปสงฺกมิตฺวา ภควนฺตํ อภิวาเทตฺวา เอกมนฺตํ นิสีทิํสุ, เอกมนฺตํ นิสินฺนา โข เต อุปาสกา ภควนฺตํ เอตทโวจุํ “อิเม ภนฺเต ติตฺถิยา อวณฺณกามา พุทฺธสฺส อวณฺณกามา ธมฺมสฺส อวณฺณกามา สํฆสฺส, สาธุ ภนฺเต อยฺยา ติตฺถิยานํ สหตฺถา น ทเทยฺยูนฺ”ติ.\csromanspacer{} อถ โข ภควา เต อุปสเก ธมฺมิยา กถาย สนฺทสฺเสสิ สมาทเปสิ สมุตฺเตเชสิ สมฺปหํเสสิ.\csromanspacer{} อถ โข เต อุปสกา ภควตา ธมฺมิยา กถาย สนฺทสฺสิตา สมาทปิตา สมุตฺเตชิตา สมฺปหํสิตา อุฏฺฐายาสนา ภควนฺตํ อภิวาเทตฺวา ปทกฺขิณํ กตฺวา ปกฺกมิํสุ.\csromanspacer{} อถ โข ภควา เอตสฺมิํ นิทาเน เอตสฺมิํ ปกรเณ ธมฺมิํ กถํ กตฺวา ภิกฺขู อามนฺเตสิ เตน หิ ภิกฺขเว ภิกฺขูนํ สิกฺขาปทํ ปญฺญเปสฺสามิ ทส อตฺถวเส ปฏิจฺจ สํฆสุฏฺฐุตาย สํฆผาสุตาย ฯเปฯ สทฺธมฺมฏฺฐิติยา วินยานุคฺคหาย.\csromanspacer{} เอวญฺจ ปน ภิกฺขเว อิมํ สิกฺกาปทํ อุทฺทิเสยฺยาถ\csromandash{}}

\proseitem{270}{\textbf{“โย ปน ภิกฺขุ อเจลกสฺส วา ปริพฺพาชกสฺส วา ปริพฺพาชิกาย วา สหตฺถา ขาทนียํ วา โภชนียํ วา ทเทยฺย, ปาจิตฺติยนฺ”}ติ.}

\proseitem{271}{\textbf{โย ปนา}ติ โย ยาทิโส ฯเปฯ.\csromanspacer{} \textbf{ภิกฺขู}ติ ฯเปฯ อยํ อิมสฺมิํ อตฺเถ อธิปฺเปโต ภิกฺขูติ.}

\prose{\textbf{อเจลโก} นาม โย โกจิ ปริพฺพาชกสมาปนฺโน นคฺโค.}

\prose{\textbf{ปริพฺพาชโก} นาม ภิกฺขุํ จ สามเณรญฺจ ฐเปตฺวา โย โกจิ ปริพฺพชกสมาปนฺโน.}

\prose{\textbf{ปริพฺพาชิกา} นาม ภิกฺขุนิญฺจ สิกฺขมานญฺจ สามเณริญฺจ ฐเปตฺวา ยา กาจิ ปริพฺพาชิกสมาปนฺนา.}

\prose{\textbf{ขาทนียํ} นาม ปญฺจ โภชนานิ อุทกทนฺตโปนํ ฐเปตฺวา อวเสสํ ขาทนียํ นาม.}

\prose{\csromanfolio{124}\textbf{โภชนียํ} นาม ปญฺจ โภชนานิ โอทโน กุมฺมาโส สตฺตุ มจฺโฉ มํสํ.}

\prose{\textbf{ทเทยฺยา}ติ กาเยน วา กายปฏิพทฺเธน วา นิสฺสคฺคิเยน วา เทติ, อาปตฺติ}

\proseitem{272}{ติตฺถิเย ติตฺถิยสญฺญี สหตฺถา ขาทนียํ วา โภชนียํ วา เทติ, อาปตฺติ ปาจิตฺติยสฺส.\csromanspacer{} ติตฺถิเย เวมติโก สหตฺถา ขาทนียํ วา โภชนียํ วา เทติ, อาปตฺติ ปาจิตฺติยสฺส.\csromanspacer{} ติตฺถิเย อติตฺถิยสญฺญี สหตฺถา ขาทนียํ วา โภชนียํ วา เทติ, อาปตฺติ ปาจิตฺติยสฺส.}

\prose{อุทกทนฺตโปนํ เทติ, อาปตฺติ ทุกฺกฏสฺส.\csromanspacer{} อติตฺถิเย ติตฺถิยสญฺญี อาปตฺติ ทุกฺกฏสฺส.\csromanspacer{} อติตฺถิเย เวมติโก, อาปตฺติ ทุกฺกฏสฺส.\csromanspacer{} อติตฺถิเย อติตฺถิยสญฺญี อนาปตฺติ.}

\proseitemwithcloserruled{273}{อนาปตฺติ ทาเปติ น เทติ, อุปนิกฺขิปิตฺวา เทติ, พาหิราเลปํ เทติ, อุมฺมตฺตกสฺส อาทิกมฺมิกสฺสาติ.}{อเจลกสิกฺขาปทํ นิฏฺฐิตํ ปฐมํ.}
\csromanmatikaanchor{mtk.48}
\csromantocmarkref{subsection}{2. อุยฺโยชนสิกฺขาปท}{mtk.48}
\bookmark[dest={mtk.48},level=2]{2. อุยฺโยชนสิกฺขาปท}
\csromanheader{2}{5. \textbf{อเจลกวคฺค 2. อุยฺโยชนสิกฺขาปท}}
\proseitem{274}{เตน สมเยน พุทฺโธ ภควา สาวตฺถิยํ วิหรติ เชตวเน อนาถปิณฺฑิกสฺส อาราเม.\csromanspacer{} เตน โข ปน สมเยน อายสฺมา อุปนนฺโท สกฺยปุตฺโต ภาตุโน สทฺธิวิหาริกํ ภิกฺขุํ เอตทโวจ “เอหาวุโส, คามํ ปิณฺฑาย ปวิสิสฺสามา”ติ.\csromanspacer{} ตสฺส อทาเปตฺวา อุยฺโยเชสิ “คจฺฉาวุโส, น เม ตยา สทฺธิํ กถา วา นิสชฺชา วา ผาสุ โหติ, เอกกสฺส เม กถา วา นิสชฺชา วา ผาสุ โหตี”ติ.\csromanspacer{} อถ โข โส ภิกฺขุ อุปกฏฺเฐ กาเล นาสกฺขิ ปิณฺฑาย จริตุํ, ปฏิกฺกมเนปิ ภตฺตวิสคฺคํ น สมฺภาเวสิ, ฉินฺนภตฺโต อโหสิ.\csromanspacer{} อถ โข โส ภิกฺขุ อารามํ คนฺตฺวา ภิกฺขูนํ เอตมตฺถํ อาโรเจสิ.\csromanspacer{} เย เต ภิกฺขู อปฺปิจฺฉา ฯเปฯ เต อุชฺฌายนฺติ ขิยฺยนฺติ วิปาเจนฺติ “กถํ หิ นาม อายสฺมา อุปนนฺโท สกฺยปุตฺโต ภิกฺขุํ ‘เอหาวุโส, คามํ ปิณฺฑาย ปวิสิสฺสามา’ติ ตสฺส \csromanfolio{125}อทาเปตฺวา อุยฺโยเชสฺสตี”ติ ฯเปฯ.\csromanspacer{} สจฺจํ กิร ตฺวํ อุปนนฺท ภิกฺขุํ “เอหาวุโส, คามํ ปิณฺฑาย ปวิสิสฺสามา”ติ ตสฺส อทาเปตฺวา อุยฺโยเชสีติ.\csromanspacer{} สจฺจํ ภควาติ.\csromanspacer{} วิครหิ พุทฺโธ ภควา ฯเปฯ.\csromanspacer{} กถํ หิ นาม ตฺวํ โมฆปุริส ภิกฺขุํ “เอหาวุโส, คามํ ปิณฺฑาย ปวิสิสฺสามา”ติ ตสฺส อทาเปตฺวา อุยฺโยเชสฺสสิ.\csromanspacer{} เนตํ โมฆปุริส อปฺปสนฺนานํ วา ปสาทาย ฯเปฯ.\csromanspacer{} เอวญฺจ ปน ภิกฺขเว อิมํ สิกฺขาปทํ อุทฺทิเสยฺยาถ\csromandash{}}

\proseitem{275}{\textbf{“โย ปน ภิกฺขุ ภิกฺขุํ ‘เอหาวุโส, คามํ วา นิคมํ วา ปิณฺฑาย ปวิสิสฺสามา’ติ ตสฺส ทาเปตฺวา วา อทาเปตฺวา วา อุยฺโยเชยฺย ‘คจฺฉาวุโส, น เม ตยา สทฺธิํ กถา วา นิสชฺชา วา ผาสุ โหติ, เอกกสฺส เม กถา วา นิสชฺชา วา ผาสุ โหตี’ติ เอตเทว ปจฺจยํ กริตฺวา อนญฺญํ, ปาจิตฺติยนฺ”}ติ.}

\proseitem{276}{\textbf{โย ปนา}ติ โย ยาทิโส ฯเปฯ.\csromanspacer{} \textbf{ภิกฺขู}ติ ฯเปฯ อยํ อิมสฺมิํ อตฺเถ อธิปฺเปโต ภิกฺขูติ.}

\prose{\textbf{ภิกฺขุนฺ}ติ อญฺญํ ภิกฺขุํ.}

\prose{\textbf{เอหาวุโส, คามํ วา นิคมํ วา}ติ คาโมปิ นิคโมปิ นครมฺปิ คาโม เจว นิคโม จ.}

\prose{\textbf{ตสฺส ทาเปตฺวา}ติ ยาคุํ วา ภตฺตํ วา ขาทนียํ วา โภชนียํ วา ทาเปตฺวา.}

\prose{\textbf{อทาเปตฺวา}ติ น กิญฺจิ ทาเปตฺวา.}

\prose{\textbf{อุยฺโยเชยฺยา}ติ มาตุคาเมน สทฺธิํ หสิตุกาโม กีฬิตุกาโม รโห นิสีทิตุกาโม อนาจารํ อาจริตุกาโม เอวํ วเทติ “คจฺฉาวุโส, น เม ตยา สทฺธิํ กถา วา นิสชฺชา วา ผาสุ โหติ, เอกกสฺส เม กถา วา นิสชฺชา วา ผาสุ โหตี”ติ อุยฺโยเชติ, อาปตฺติ ทุกฺกฏสฺส.\csromanspacer{} ทสฺสนุปจารํ วา สวนูปจารํ วา วิชหนฺตสฺส อาปตฺติ ทุกฺกฏสฺส.\csromanspacer{} วิชหิเต อาปตฺติ}

\prose{\textbf{เอตเทว ปจฺจยํ กริตฺวา อนญฺญนฺ}ติ น อญฺโญ โกจิ ปจฺจโย โหติ อุยฺโยเชตุํ.}

\proseitem{277}{\csromanfolio{126}อุปสมฺปนฺเน อุปสมฺปนฺนสญฺญี อุยฺโยเชติ, อาปตฺติ ปาจิตฺติยสฺส.\csromanspacer{} อุปสมฺปนฺเน เวมติโก อุยฺโยเชติ, อาปตฺติ ปาจิตฺติยสฺส.\csromanspacer{} อุปสมฺปนฺเน อนุปสมฺปนฺนสญฺญี อุยฺโยเชติ, อาปตฺติ ปาจิตฺติยสฺส.}

\prose{กลิสาสนํ อาโรเปติ, อาปตฺติ ทุกฺกฏสฺส.\csromanspacer{} อนุปสมฺปนฺนํ อุยฺโยเชติ, อาปตฺติ ทุกฺกฏสฺส.\csromanspacer{} กลิสาสนํ อาโรเปติ, อาปตฺติ ทุกฺกฏสฺส.\csromanspacer{} อนุปสมฺปนฺเน อุปสมฺปนฺนสญฺญี, อาปตฺติ ทุกฺกฏสฺส.\csromanspacer{} อนุปสมฺปนฺเน เวมติโก, อาปตฺติ ทุกฺกฏสฺส.\csromanspacer{} อนุปสมฺปนฺเน อนุปสมฺปนฺนสญฺญี, อาปตฺติ ทุกฺกฏสฺส.}

\proseitemwithcloserruled{278}{อนาปตฺติ “อุโภ เอกโต น ยาเปสฺสามา”ติ อุยฺโยเชติ, “มหคฺฆํ ภณฺฑํ ปสฺสิตฺวา โลภธมฺมํ อุปฺปาเทสฺสตี”ติ อุยฺโยเชติ, “มาตุคามํ ปสฺสิตฺวา อนภิรติํ อุปฺปาเทสฺสตี”ติ อุยฺโยเชติ, “คิลานสฺส วา โอหิยฺยกสฺส วา วิหารปาลสฺส วา ยาคุํ วา ภตฺตํ วา ขาทนียํ วา โภชนียํ วา นีหรา”ติ อุยฺโยเชติ, น อนาจารํ อาจริตุกาโม, สติ กรณีเย อุยฺโยเชติ, อุมฺมตฺตกสฺส อาทิกมฺมิกสฺสาติ.}{อุยฺโยชนสิกฺขาปทํ นิฏฺฐิตํ ทุติยํ.}
\csromanmatikaanchor{mtk.49}
\csromantocmarkref{subsection}{3. สโภชนสิกฺขาปท}{mtk.49}
\bookmark[dest={mtk.49},level=2]{3. สโภชนสิกฺขาปท}
\csromanheader{2}{5. \textbf{อเจลกวคฺค 3. สโภชนสิกฺขาปท}}
\proseitem{279}{เตน สมเยน พุทฺโธ ภควา สาวตฺถิยํ วิหรติ เชตวเน อนาถปิณฺฑิกสฺส อาราเม.\csromanspacer{} เตน โข ปน สมเยน อายสฺมา อุปนนฺโท สกฺยปุตฺโต สหายกสฺส ฆรํ คนฺตฺวา ตสฺส ปชาปติยา สทฺธิํ สยนิฆเร\footnote{สยนีฆเร (สฺยา)} นิสชฺชํ กปฺเปสิ.\csromanspacer{} อถ โข โส ปุริโส เยนายสฺมา อุปนนฺโท สกฺยปุตฺโต เตนุปสงฺกมิ, อุปสงฺกมิตฺวา อายสฺมนฺตํ อุปนนฺทํ สกฺยปุตฺตํ อภิวาเทตฺวา เอกมนฺตํ นิสีทิ, เอกมนฺตํ นิสินฺโน โข โส ปุริโส ปชาปติํ เอตทโวจ “ทเทหายฺยสฺส ภิกฺขนฺ”ติ.\csromanspacer{} อถ โข สา อิตฺถี อายสฺมโต อุปนนฺทสฺส สกฺยปุตฺตสฺส ภิกฺขํ อทาสิ.\csromanspacer{} อถ โข โส ปุริโส อายสฺมนฺตํ อุปนนฺทํ สกฺยปุตฺตํ เอตทโวจ “คจฺฉถ ภนฺเต ยโต อยฺยสฺส ภิกฺขา ทินฺนา”ติ.\csromanspacer{} อถ โข สา อิตฺถี สลฺลกฺเขตฺวา ‘ปริยุฏฺฐิโต อยํ ปุริโส’ติ อยสฺมนฺตํ อุปนนฺทํ สกฺยปุตฺตํ \csromanfolio{127}เอตทโวจ “นิสีทถ ภนฺเต, มา อคมิตฺถา”ติ.\csromanspacer{} ทุติยมฺปิ โข โส ปุริโส ฯเปฯ.\csromanspacer{} ตติยมฺปิ โข โส ปุริโส อายสฺมนฺตํ อุปนนฺทํ สกฺยปุตฺตํ เอตทโวจ “คจฺฉถ ภนฺเต ยโต อยฺยสฺส ภิกฺขา ทินฺนา”ติ.\csromanspacer{} ตติยมฺปิ โข สา อิตฺถี อายสฺมนฺตํ อุปนนฺทํ สกฺยปุตฺตํ เอตทโวจ “นิสีทถ ภนฺเต, มา อคมิตฺถา”ติ.}

\prose{อถ โข โส ปุริโส นิกฺขมิตฺวา ภิกฺขู อุชฺฌาเปสิ “อยํ ภนฺเต อยฺโย อุปนนฺโท มยฺหํ ปชาปติยา สทฺธิํ สยนิฆเร นิสินฺโน.\csromanspacer{} โส มยา อุยฺโยชียมาโน น อิจฺฉติ คนฺตุํ, พหุกิจฺจา มยํ พหุกรณียา”ติ.\csromanspacer{} เย เต ภิกฺขู อปฺปิจฺฉา ฯเปฯ เต อุชฺฌายนฺติ ขิยฺยนฺติ วิปาเจนฺติ “กถํ หิ นาม อายสฺมา อุปนนฺโท สกฺยปุตฺโต สโภชเน กุเล อนุปขชฺช นิสชฺชํ สปฺเปสฺสตี”ติ ฯเปฯ.\csromanspacer{} สจฺจํ กิร ตฺวํ อุปนนฺท สโภชเน กุเล อนุปขชฺช นิสชฺชํ กปฺเปสีติ.\csromanspacer{} สจฺจํ ภควาติ.\csromanspacer{} วิครหิ พุทฺโธ ภควา ฯเปฯ.\csromanspacer{} กถํ หิ นาม ตฺวํ โมฆปุริส สโภชเน กุเล อนุปขชฺช นิสชฺชํ กปฺเปสฺสสิ.\csromanspacer{} เนตํ โมฆปุริส อปฺปสนฺนานํ วา ปสาทาย ฯเปฯ.\csromanspacer{} เอวญฺจ ปน ภิกฺขเว อิมํ สิกฺขาปทํ อุทฺทิเสยฺยาถ\csromandash{}}

\proseitem{280}{\textbf{“โย ปน ภิกฺขุ สโภชเน กุเล อนุปขชฺช นิสชฺชํ กปฺเปยฺย, ปาจิตฺติยนฺ”}ติ.}

\proseitem{281}{\textbf{โย ปนา}ติ โย ยาทิโส ฯเปฯ.\csromanspacer{} \textbf{ภิกฺขู}ติ ฯเปฯ อยํ อิมสฺมิํ อตฺเถ อธิปฺเปโต ภิกฺขูติ.}

\prose{\textbf{สโภชนํ} นาม กุลํ อิตฺถี เจว โหติ ปุริโส จ, อิตฺถี จ ปุริโส จ อุโภ อนิกฺขนฺตา โหนฺติ, อุโภ อวีตราคา.}

\prose{\textbf{อนุปขชฺชา}ติ อนุปวิสิตฺวา.}

\prose{\textbf{นิสชฺชํ กปฺเปยฺยา}ติ มหลฺลเก ฆเร ปิฏฺฐสํฆาฏสฺส\footnote{ปิฏฺฐิสงฺฆาฏสฺส (สฺยา)} หตฺถปาสํ วิชหิตฺวา นิสีทติ, อาปตฺติ ปาจิตฺติยสฺส.\csromanspacer{} ขุทฺทเก ฆเร ปิฏฺฐิวํสํ อติกฺกมิตฺวา นิสีทติ, อาปตฺติ ปาจิตฺติยสฺส.}

\proseitem{282}{สยนิฆเร สยนิฆรสญฺญี สโภชเน กุเล อนุปขชฺช นิสชฺชํ กปฺเปติ, อาปตฺติ ปาจิตฺติยสฺส.\csromanspacer{} สยนิฆเร เวมติโก สโภชเน \csromanfolio{128}\textbf{กุเล อนุปขชฺช นิสชฺชํ กปฺเปติ, อาปตฺติ ปาจิตฺติยสฺส.}\csromanspacer{}\textbf{ สยนิฆเร น} \textbf{สยนิฆรสญฺญี สโภชเน กุเล อนุปขชฺช นิสชฺชํ กปฺเปติ, อาปตฺติ}}

\prose{น สยนิฆเร สยนิฆรสญฺญี, อาปตฺติ ทุกฺกฏสฺส.\csromanspacer{} น สยนิฆเร เวมติโก, อาปตฺติ ทุกฺกฏสฺส.\csromanspacer{} น สยนิฆเร น สยนิฆรสญฺญี, อนาปตฺติ.}

\proseitemwithcloserruled{283}{อนาปตฺติ มหลฺลเก ฆเร ปิฏฺฐสํฆาฏสฺส หตฺถปาสํ อวิชหิตฺวา นิสีทติ, ขุทฺทเก ฆเร ปิฏฺฐิวํสํ อนติกฺกมิตฺวา นิสีทติ, ภิกฺขุ ทุติโย โหติ, อุโภ นิกฺขนฺตา โหนฺติ, อุโภ วีตราคา น สยนิฆเร, อุมฺมตฺตกสฺส อาทิกมฺมิกสฺสาติ.}{สโภชนสิกฺขาปทํ นิฏฺฐิตํ ตติยํ.}
\csromanmatikaanchor{mtk.50}
\csromantocmarkref{subsection}{4. รโหปฏิจฺฉนฺนสิกฺขาปท}{mtk.50}
\bookmark[dest={mtk.50},level=2]{4. รโหปฏิจฺฉนฺนสิกฺขาปท}
\csromanheader{2}{5. อเจลกวคฺค 4. รโห\-ปฏิจฺฉนฺน\-สิกฺขาปท}
\proseitem{284}{เตน สมเยน พุทฺโธ ภควา สาวตฺถิยํ วิหรติ เชตวเน อนาถปิณฺฑิกสฺส อาราเม.\csromanspacer{} เตน โข ปน สมเยน อายสฺมา อุปนนฺโท สกฺยปุตฺโต สหายกสฺส ฆรํ คนฺตฺวา ตสฺส ปชาปติยา สทฺธิํ รโห ปฏิจฺฉนฺเน อาสเน นิสชฺชํ กปฺเปสิ.\csromanspacer{} อถ โข โส ปุริโส อุชฺฌายติ ขิยฺยติ วิปาเจติ “กถํ หิ นาม อยฺโย อุปนนฺโท มยฺหํ ปชาปติยา สทฺธิํ รโห ปฏิจฺฉนฺเน อาสเน นิสชฺชํ กปฺเปสฺสตี”ติ.\csromanspacer{} อสฺโสสุํ โข ภิกฺขู ตสฺส ปุริสสฺส อุชฺฌายนฺตสฺส ขิยฺยนฺตสฺส วิปาเจนฺตสฺส.\csromanspacer{} เย เต ภิกฺขู อปฺปิจฺฉา ฯเปฯ เต อุชฺฌายนฺติ ขิยฺยนฺติ วิปาเจนฺติ “กถํ หิ นาม อายสฺมา อุปนนฺโท สกฺยปุตฺโต มาตุคาเมน สทฺธิํ รโห ปฏิจฺฉนฺเน อาสเน นิสชฺชํ กปฺเปสฺสตี”ติ ฯเปฯ.\csromanspacer{} สจฺจํ กิร ตฺวํ อุปนนฺท มาตุคาเมน สทฺธิํ รโห ปฏิจฺฉนฺเน อาสเน นิสชฺชํ กปฺเปสีติ.\csromanspacer{} สจฺจํ ภควาติ.\csromanspacer{} วิครหิ พุทฺโธ ภควา ฯเปฯ.\csromanspacer{} กถํ หิ นาม ตฺวํ โมฆปุริส มาตุคาเมน สทฺธิํ รโห ปฏิจฺฉนฺเน อาสเน นิสชฺชํ กปฺเปสฺสสิ.\csromanspacer{} เนตํ โมฆปุริส อปฺปสนฺนานํ วา ปสาทาย ฯเปฯ.\csromanspacer{} เอวญฺจ ปน ภิกฺขเว อิมํ สิกฺขาปทํ อุทฺทิเสยฺยาถ\csromandash{}}

\proseitem{285}{\csromanfolio{129}\textbf{“โย ปน ภิกฺขุ มาตุคาเมน สทฺธิํ รโห ปฏิจฺฉนฺเน อาสเน นิสชฺชํ กปฺเปยฺย, ปาจิตฺติยนฺ”}ติ.}

\proseitem{286}{\textbf{โย ปนา}ติ โย ยาทิโส ฯเปฯ.\csromanspacer{} \textbf{ภิกฺขู}ติ ฯเปฯ อยํ อิมสฺมิํ อตฺเถ อธิปฺเปโต ภิกฺขูติ.}

\prose{\textbf{มาตุคาโม} นาม มนุสฺสิตฺถี, น ยกฺขี น เปตี น ติรจฺฉานคตา, อนฺตมโส ตทหุชาตาปิ ทาริกา, ปเคว มหตฺตรี.}

\prose{\textbf{สทฺธินฺ}ติ เอกโต.}

\prose{\csromansymbolmark{*}\textbf{รโห} นาม จกฺขุสฺส รโห โสตสฺส รโห.\csromanspacer{} \textbf{จกฺขุสฺส รโห} นาม น สกฺกา โหติ อกฺขิํ วา นิขณียมาเน ภมุกํ วา อุกฺขิปียมาเน สีสํ วา อุกฺขิปียมาเน ปสฺสิตุํ.\csromanspacer{} \textbf{โสตสฺส รโห} นาม น สกฺกา โหติ ปกติกถา โสตุํ.}

\prose{\csromansymbolmark{+}\textbf{ปฏิจฺฉนฺนํ} นาม อาสนํ กุฏฺเฏน\footnote{กุฑฺเฑน (สี, สฺยา)} วา กวาเฏน วา กิลญฺเชน วา สาณิปากาเรน วา รุกฺเขน วา ถมฺเภน วา โกตฺถฬิยา\footnote{โกฏฺฐฬิยา (สี, สฺยา), โกตฺถฬิกาย (ก)} วา เยน เกนจิ ปฏิจฺฉนฺนํ โหติ.}

\prose{\textbf{นิสชฺชํ กปฺเปยฺยา}ติ มาตุคาเม นิสินฺเน ภิกฺขุ อุปนิสินฺโน วา โหติ อุปนิปนฺโน วา, อาปตฺติ ปาจิตฺติยสฺส.\csromanspacer{} ภิกฺขุ นิสินฺเน มาตุคาโม อุปนิสินฺโน วา โหติ อุปนิปนฺโน วา, อาปตฺติ ปาจิตฺติยสฺส.\csromanspacer{} อุโภ วา นิสินฺนา โหนฺติ อุโภ วา นิปนฺนา, อาปตฺติ ปาจิตฺติยสฺส.}

\proseitem{287}{มาตุคาเม มาตุคามสญฺญี รโห ปฏิจฺฉนฺเน อาสเน นิสชฺชํ กปฺเปติ, อาปตฺติ ปาจิตฺติยสฺส.\csromanspacer{} มาตุคาเม เวมติโก รโห ปฏิจฺฉนฺเน อาสเน นิสชฺชํ กปฺเปติ, อาปตฺติ ปาจิตฺติยสฺส.\csromanspacer{} มาตุคาเม อมาตุคามสญฺญี รโห ปฏิจฺฉนฺเน อาสเน นิสชฺชํ กปฺเปติ, อาปตฺติ ปาจิตฺติยสฺส.}

\prose{ยกฺขิยา วา เปติยา วา ปณฺฑเกน วา ติรจฺฉานคตาย วา มนุสฺสวิคฺคหิตฺถิยา วา สทฺธิํ รโห ปฏิจฺฉนฺเน อาสเน นิสชฺชํ กปฺเปติ, อาปตฺติ ทุกฺกฏสฺส.\csromanspacer{} อมาตุคาเม มาตุคามสญฺญี, อาปตฺติ ทุกฺกฏสฺส.\csromanspacer{} อมาตุคาเม เวมติโก, อาปฺปตฺติ ทุกฺกฏสฺส.\csromanspacer{} อมาตุคาเม อมาตุคามสญฺญี, อนาปตฺติ.}

\proseitemwithcloserruled{288}{\csromanfolio{130}อนาปตฺติ โย โกจิ วิญฺญู ปุริโส ทุติโย โหติ, ติฏฺฐติ น นิสีทติ, อรโหเปกฺโข, อญฺญวิหิโต นิสีทติ, อุมฺมตฺตกสฺส อาทิกมฺมิกสฺสาติ.}{รโห\-ปฏิจฺฉนฺน\-สิกฺขาปทํ นิฏฺฐิตํ จตุตฺถํ.}
\csromanmatikaanchor{mtk.51}
\csromantocmarkref{subsection}{5. รโหนิสชฺชสิกฺขาปท}{mtk.51}
\bookmark[dest={mtk.51},level=2]{5. รโหนิสชฺชสิกฺขาปท}
\csromanheader{2}{5. \textbf{อเจลกวคฺค 5. รโหนิสชฺชสิกฺขาปท}}
\proseitem{289}{เตน สมเยน พุทฺโธ ภควา สาวตฺถิยํ วิหรติ เชตวเน อนาถปิณฺฑิกสฺส อาราเม.\csromanspacer{} เตน โข ปนสมเยน อายสฺมา อุปนนฺโท สกฺยปุตฺโต สหายกสฺส ฆรํ คนฺตฺวา ตสฺส ปชาปติยา สทฺธิํ เอโก เอกาย รโห นิสชฺชํ กปฺเปสิ.\csromanspacer{} อถ โข โส ปุริโส อุชฺฌายติ ขิยฺยติ วิปาเจติ “กถํ หิ นาม อยฺโย อุปนนฺโท มยฺหํ ปชาปติยา สทฺธิํ เอโก เอกาย รโห นิสชฺชํ กปฺเปสฺสตี”ติ.\csromanspacer{} อสฺโสสุํ โข ภิกฺขู ตสฺส ปุริสสฺส อุชฺฌายนฺตสฺส ขิยฺยนฺตสฺส วิปาเจนฺตสฺส.\csromanspacer{} เย เต ภิกฺขู อปฺปิจฺฉา ฯเปฯ เต อุชฺฌายนฺติ ขิยฺยนฺติ วิปาเจนฺติ “กถํ หิ นาม อายสฺมา อุปนนฺโท สกฺยปุตฺโต มาตุคาเมน สทฺธิํ เอโก เอกาย รโห นิสชฺชํ กปฺเปสฺสตี”ติ ฯเปฯ.\csromanspacer{} สจฺจํ กิร ตฺวํ อุปนนฺท มาตุคาเมน สทฺธิํ เอโก เอกาย รโห นิสชฺชํ กปฺเปสีติ.\csromanspacer{} สจฺจํ ภควาติ.\csromanspacer{} วิครหิ พุทฺโธ ภควา ฯเปฯ.\csromanspacer{} กถํ หิ นาม ตฺวํ โมฆปุริส มาตุคาเมน สทฺธิํ เอโก เอกาย รโห นิสชฺชํ กปฺเปสฺสสิ.\csromanspacer{} เนตํ โมฆปุริส อปฺปสนฺนานํ วา ปสาทาย ฯเปฯ.\csromanspacer{} เอวญฺจ ปน ภิกฺขเว อิมํ สิกฺขาปทํ อุทฺทิเสยฺยาถ\csromandash{}}

\proseitem{290}{\textbf{“โย ปน ภิกฺขุ มาตุคาเมน สทฺธิํ เอโก เอกาย รโห นิสชฺชํ กปฺเปยฺย, ปาจิตฺติยนฺ”}ติ.}

\proseitem{291}{\textbf{โย ปนา}ติ โย ยาทิโส ฯเปฯ.\csromanspacer{} \textbf{ภิกฺขู}ติ ฯเปฯ อยํ อิมสฺมิํ อตฺเถ อธิปฺเปโต ภิกฺขูติ.}

\prose{\textbf{มาตุคาโม} นาม มนุสฺสิตฺถี, น ยกฺขี น เปตี น ติรจฺฉานคตา, วิญฺญู ปฏิพลา สุภาสิตทุพฺภาสิตํ ทุฏฺฐุลฺลาทุฏฺฐุลฺลํ อาชานิตุํ.}

\prose{\textbf{สทฺธินฺ}ติ เอกโต.}

\prose{\textbf{เอโก เอกายา}ติ ภิกฺขุ เจว โหติ มาตุคาโม จ.}

\prose{\csromanfolio{131}\csromansymbolfootnote{*}{94, 129 ปิฏฺเฐสุปิ; วิ 1. 287, 291 ปิฏฺเฐสุปิ จ.}\textbf{รโห} นาม จกฺขุสฺส รโห โสตสฺส รโห.\csromanspacer{} \textbf{จกฺขุสฺส รโห} นาม น สกฺกา โหติ อกฺขิํ วา นิขณียมาเน ภมุกํ วา อุกฺขิปียมาเน สีสํ วา อุกฺขิปียมาเน ปสฺสิตุํ.\csromanspacer{} \textbf{โสตสฺส รโห} นาม น สกฺกา โหติ ปกติกถา โสตุํ.}

\prose{\textbf{นิสชฺชํ กปฺเปยฺยา}ติ มาตุคาเม นิสินฺเน ภิกฺขุ อุปนิสินฺโน วา โหติ อุปนิปนฺโน วา, อาปตฺติ ปาจิตฺติยสฺส.\csromanspacer{} ภิกฺขุ นิสินฺเน มาตุคาโม อุปนิสินฺโน วา โหติ อุปนิปนฺโน วา, อาปตฺติ ปาจิตฺติยสฺส.\csromanspacer{} อุโภ วา นิสินฺนา โหนฺติ อุโภ วา นิปนฺนา, อาปตฺติ ปาจิตฺติยสฺส.}

\proseitem{292}{มาตุคาเม มาตุคามสญฺญี เอโก เอกาย รโห นิสชฺชํ กปฺเปติ, อาปตฺติ ปาจิตฺติยสฺส.\csromanspacer{} มาตุคาเม เวมติโก เอโก เอกาย รโห นิสชฺชํ กปฺเปติ, อาปตฺติ ปาจิตฺติยสฺส.\csromanspacer{} มาตุคาเม อมาตุคามสญฺญี เอโก เอกาย รโห นิสชฺชํ กปฺเปติ, อาปตฺติ ปาจิตฺติยสฺส.}

\prose{ยกฺขิยา วา เปติยา วา ปณฺฑเกน วา ติรจฺฉานคตมนุสฺสวิคฺคหิตฺถิยา วา\footnote{ติรจฺฉานคตาย วา มนุสฺสวิคฺคหิตฺถิยา วา (ก)} สทฺธิํ เอโก เอกาย รโห นิสชฺชํ กปฺเปติ, อาปตฺติ ทุกฺกฏสฺส.\csromanspacer{} อมาตุคาเม มาตุคามสญฺญี, อาปตฺติ ทุกฺกฏสฺส.\csromanspacer{} อมาตุคาเม เวมติโก, อาปตฺติ ทุกฺกฏสฺส.\csromanspacer{} อมาตุคาเม อมาตุคามสญฺญี, อนาปตฺติ.}

\proseitemwithcloserruled{293}{อนาปตฺติ โย โกจิ วิญฺญู ปุริโส ทุติโย โหติ, ติฏฺฐติ น นิสีทติ, อรโหเปกฺโข, อญฺญวิหิโต นิสีทติ, อุมฺมตฺตกสฺส อาทิกมฺมิกสฺสาติ.}{รโหนิสชฺชสิกฺขาปทํ นิฏฺฐิตํ ปญฺจมํ.}
\csromanmatikaanchor{mtk.52}
\csromantocmarkref{subsection}{6. จาริตฺตสิกฺขาปท}{mtk.52}
\bookmark[dest={mtk.52},level=2]{6. จาริตฺตสิกฺขาปท}
\csromanheader{2}{5. \textbf{อเจลกวคฺค 6. จาริตฺตสิกฺขาปท}}
\proseitem{294}{เตน สมเยน พุทฺโธ ภควา ราชคเห วิหรติ เวฬุวเน กลนฺทกนิวาเป.\csromanspacer{} เตน โข ปน สมเยน อายสฺมโต อุปนนฺทสฺส \csromanfolio{132}สกฺยปุตฺตสฺส อุปฏฺฐากกุลํ อายสฺมนฺตํ อุปนนฺทํ สกฺยปุตฺตํ ภตฺเตน นิมนฺเตสิ, อญฺเญปิ ภิกฺขู ภตฺเตน นิมนฺเตสิ.\csromanspacer{} เตน โข ปน สมเยน อายสฺมา อุปนนฺโท สกฺยปุตฺโต ปุเรภตฺตํ กุลานิ ปยิรุปาสติ.\csromanspacer{} อถ โข เต ภิกฺขู เต มนุสฺเส เอตทโวจุํ “เทถาวุโส ภตฺตนฺ”ติ.\csromanspacer{} อาคเมถ ภนฺเต ยาวายฺโย อุปนนฺโท อาคจฺฉตีติ.\csromanspacer{} ทุติยมฺปิ โข เต ภิกฺขู ฯเปฯ.\csromanspacer{} ตติยมฺปิ โข เต ภิกฺขู เต มนุสฺเส เอตทโวจุํ “เทถาวุโส ภตฺตํ ปุเร กาโล อติกฺกมตี”ติ.\csromanspacer{} ยมฺปิ มยํ ภนฺเต ภตฺตํ กริมฺหา อยฺยสฺส อุปนนฺทสฺส การณา, อาคเมถ ภนฺเต ยาวายฺโย อุปนนฺโท อาคจฺฉตีติ.}

\prose{อถ โข อายสฺมา อุปนนฺโท สกฺยปุตฺโต ปุเรภตฺตํ กุลานิ ปยิรุปาสิตฺวา ทิวา อาคจฺฉติ.\csromanspacer{} ภิกฺขู น จิตฺตรูปํ ภุญฺชิํสุ.\csromanspacer{} เย เต ภิกฺขู อปฺปิจฺฉา ฯเปฯ เต อุชฺฌายนฺติ ขิยฺยนฺติ วิปาเจนฺติ “กถํ หิ นาม อายสฺมา อุปนนฺโท สกฺยปุตฺโต นิมนฺติโต สภตฺโต สมาโน ปุเรภตฺตํ กุเลสุ จาริตฺตํ อาปชฺชิสฺสตี”ติ ฯเปฯ.\csromanspacer{} สจฺจํ กิร ตฺวํ อุปนนฺท นิมนฺติโต สภตฺโต สมาโน ปุเรภตฺตํ กุเลสุ จาริตฺตํ อาปชฺชสีติ.\csromanspacer{} สจฺจํ ภควาติ.\csromanspacer{} วิครหิ พุทฺโธ ภควา ฯเปฯ.\csromanspacer{} กถํ หิ นาม ตฺวํ โมฆปุริส นิมนฺติโต สภตฺโต สมาโน ปุเรภตฺตํ กุเลสุ จาริตฺตํ อาปชฺชิสฺสสิ.\csromanspacer{} เนตํ โมฆปุริส อปฺปสนฺนานํ วา ปสาทาย ฯเปฯ.\csromanspacer{} เอวญฺจ ปน ภิกฺขเว อิมํ สิกฺขาปทํ อุทฺทิเสยฺยาถ\csromandash{}}

\prose{\textbf{“โย ปน ภิกฺขุ นิมนฺติโต สภตฺโต สมาโน ปุเรภตฺตํ กุเลสุ จาริตฺตํ อาปชฺเชยฺย, ปาจิตฺติยนฺ”}ติ.}

\prose{เอวญฺจิทํ ภควตา ภิกฺขูนํ สิกฺขาปทํ ปญฺญตฺตํ โหติ.}

\proseitem{295}{\csromansymbolmark{*}เตน โข ปน สมเยน อายสฺมโต อุปนนฺทสฺส สกฺยปุตฺตสฺส อุปฏฺฐากกุลํ สํฆสฺสตฺถาย ขาทนียํ ปาเหสิ “อยฺยสฺส อุปนนฺทสฺส ทสฺเสตฺวา สํฆสฺส ทาตพฺพนฺ”ติ.\csromanspacer{} เตน โข ปน สมเยน อายสฺมา อุปนนฺโท สกฺยปุตฺโต คามํ ปิณฺฑาย ปวิฏฺโฐ โหติ.\csromanspacer{} อถ โข เต มนุสฺสา อารามํ คนฺตฺวา ภิกฺขู ปุจฺฉิํสุ “กหํ ภนฺเต อยฺโย อุปนนฺโท”ติ.\csromanspacer{} เอสาวุโส อายสฺมา อุปนนฺโท สกฺยปุตฺโต คามํ ปิณฺฑาย ปวิฏฺโฐติ.\csromanspacer{} อิทํ ภนฺเต ขาทนียํ อยฺยสฺส อุปนนฺทสฺส ทสฺเสตฺวา สํฆสฺส ทาตพฺพนฺติ.\csromanspacer{} ภควโต เอตมตฺถํ อาโรเจสุํ.\csromanspacer{} อถ โข ภควา เอตสฺมิํ นิทาเน \csromanfolio{133}\textbf{เอตสฺมิํ ปกรเณ ธมฺมิํ กถํ กตฺวา ภิกฺขู อามนฺเตสิ “เตน หิ} \textbf{ภิกฺขเว ปฏิคฺคเหตฺวา นิกฺขิปถ ยาว อุปนนฺโท อาคจฺฉตี”ติ.}}

\prose{อถ โข อายสฺมา อุปนนฺโท สกฺยปุตฺโต “ภควตา ปฏิกฺขิตฺตํ ปุเรภตฺตํ กุเลสุ จาริตฺตํ อาปชฺชิตุนฺ”ติ ปจฺฉาภตฺตํ กุลานิ ปยิรุปาสิตฺวา ทิวา ปฏิกฺกมิ, ขาทนียํ อุสฺสาริยิตฺถ.\csromanspacer{} เย เต ภิกฺขู อปฺปิจฺฉา, เต อุชฺฌายนฺติ ขิยฺยนฺติ วิปาเจนฺติ “กถํ หิ นาม อายสฺมา อุปนนฺโท สกฺยปุตฺโต ปจฺฉาภตฺตํ กุเลสุ จาริตฺตํ อาปชฺชิสฺสตี”ติ ฯเปฯ.\csromanspacer{} สจฺจํ กิร ตฺวํ อุปนนฺท ปจฺฉาภตฺตํ กุเลสุ จาริตฺตํ อาปชฺชสีติ.\csromanspacer{} สจฺจํ ภควาติ.\csromanspacer{} วิครหิ พุทฺโธ ภควา ฯเปฯ.\csromanspacer{} กถํ หิ นาม ตฺวํ โมฆปุริส ปจฺฉาภตฺตํ กุเลสุ จาริตฺตํ อาปชฺชิสฺสสิ.\csromanspacer{} เนตํ โมฆปุริส อปฺปสนฺนานํ วา ปสาทาย ฯเปฯ.\csromanspacer{} เอวญฺจ ปน ภิกฺขเว อิมํ สิกฺขาปทํ อุทฺทิเสยฺยาถ\csromandash{}}

\prose{\textbf{“โย ปน ภิกฺขุ นิมนฺติโต สภตฺโต สมาโน ปุเรภตฺตํ วา ปจฺฉาภตฺตํ วา กุเลสุ จาริตฺตํ อาปชฺเชยฺย, ปาจิตฺติยนฺ”}ติ.}

\prose{เอวญฺจิทํ ภควตา ภิกฺขูนํ สิกฺขาปทํ ปญฺญตฺตํ โหติ.}

\proseitem{296}{เตน โข ปน สมเยน ภิกฺขู จีวรทานสมเย กุกฺกุจฺจายนฺตา กุลานิ น ปยิรุปาสนฺติ.\csromanspacer{} จีวรํ ปริตฺตํ อุปชฺชติ.\csromanspacer{} ภควโต เอตมตฺถํ อาโรเจสุํ ฯเปฯ.\csromanspacer{} อนุชานามิ ภิกฺขเว จีวรทานสมเย กุลานิ ปยิรุปาสิตุํ.\csromanspacer{} เอวญฺจ ปน ภิกฺขเว อิมํ สิกฺขาปทํ อุทฺทิเสยฺยาถ\csromandash{}}

\prose{“\textbf{โย ปน ภิกฺขุ นิมนฺติโต สภตฺโต สมาโน ปุเรภตฺตํ วา} \textbf{ปจฺฉาภตฺตํ วา กุเลสุ จาริตฺตํ อาปชฺเชยฺ}ย อญฺญตฺร สมยา, ปาจิตฺติยํ.\csromanspacer{} ตตฺถายํ สมโย, จีวรทานสมโย, อยํ ตตฺถ สมโย”ติ.}

\prose{เอวญฺจิทํ ภควตา ภิกฺขูนํ สิกฺขาปทํ ปญฺญตฺตํ โหติ.}

\proseitem{297}{เตน โข ปน สมเยน ภิกฺขู จีวรกมฺมํ กโรนฺติ, อตฺโถ จ โหติ สูจิยาปิ สุตฺเตนปิ สตฺถเกนปิ.\csromanspacer{} ภิกฺขู กุกฺกุจฺจายนฺตา กุลานิ น ปยิรุปาสนฺติ.\csromanspacer{} ภควโต เอตมตฺถํ อาโรเจสุํ ฯเปฯ.\csromanspacer{} อนุชานามิ ภิกฺขเว จีวรการสมเย กุลานิ ปยิรุปาสิตุํ.\csromanspacer{} เอวญฺจ ปน ภิกฺขเว อิมํ สิกฺขาปทํ อุทฺทิเสยฺยาถ\csromandash{}}

\prose{“\textbf{โย ปน ภิกฺขุ นิมนฺติโต สภตฺโต สมาโน ปุเรภตฺตํ วา} \textbf{ปจฺฉาภตฺตํ วา กุเลสุ จาริตฺตํ อาปชฺเชยฺ}ย อญฺญตฺร สมยา, \csromanfolio{134}ปาจิตฺติยํ.\csromanspacer{} ตตฺถายํ สมโย, จีวรทานสมโย จีวรการสมโย, อยํ \textbf{ตตฺถ สมโย”ติ.}}

\prose{เอวญฺจิทํ ภควตา ภิกฺขูนํ สิกฺขาปทํ ปญฺญตฺตํ โหติ.}

\proseitem{298}{เตน โข ปน สมเยน ภิกฺขู คิลานา โหนฺติ, อตฺโถ จ โหติ เภสชฺเชหิ.\csromanspacer{} ภิกฺขู กุกฺกุจฺจายนฺตา กุลานิ น ปยิรุปาสนฺติ.\csromanspacer{} ภควโต เอตมตฺถํ อาโรเจสุํ ฯเปฯ.\csromanspacer{} อนุชานามิ ภิกฺขเว สนฺตํ ภิกฺขุํ อาปุจฺฉา กุลานิ ปยิรุปาสิตุํ.\csromanspacer{} เอวญฺจ ปน ภิกฺขเว อิมํ สิกฺขาปทํ อุทฺทิเสยฺยาถ\csromandash{}}

\proseitem{299}{\textbf{“โย ปน ภิกฺขุ นิมนฺติโต สภตฺโต สมาโน สนฺตํ ภิกฺขุํ อนาปุจฺฉา ปุเรภตฺตํ วา ปจฺฉาภตฺตํ วา กุเลสุ จาริตฺตํ อาปชฺเชยฺย อญฺญตฺร สมยา, ปาจิตฺติยํ.}\csromanspacer{} ตตฺถายํ สมโย, จีวรทานสมโย จีวรการสมโย, อยํ ตตฺถ สมโย”ติ.}

\proseitem{300}{\textbf{โย ปนา}ติ โย ยาทิโส ฯเปฯ.\csromanspacer{} \textbf{ภิกฺขู}ติ ฯเปฯ อยํ อิมสฺมิํ อตฺเถ อธิปฺเปโต ภิกฺขูติ.}

\prose{\textbf{นิมนฺติ}โต นาม ปญฺจนฺนํ โภชนานํ อญฺญตเรน โภชเนน นิมนฺติโต.}

\prose{\textbf{สภตฺโต} นาม เยน นิมนฺติโต, เตน สภตฺโต.}

\prose{\textbf{สนฺตํ} นาม ภิกฺขุํ สกฺกา โหติ อาปุจฺฉา ปวิสิตุํ.}

\prose{\textbf{อสนฺตํ} นาม ภิกฺขุํ น สกฺกา โหติ อาปุจฺฉา ปวิสิตุํ.}

\prose{\textbf{ปุเรภตฺตํ} นาม เยน นิมนฺติโต, ตํ อภุตฺตาวี.}

\prose{\textbf{ปจฺฉาภตฺตํ} นาม เยน นิมนฺติโต, ตํ อนฺตมโส กุสคฺเคนปิ\footnote{อรุณุคฺคมเนปิ (สี)} ภุตฺตํ โหติ.}

\prose{\textbf{กุลํ} นาม จตฺตาริ กุลานิ ขตฺติยกุลํ พฺราหฺมณกุลํ เวสฺสกุลํ สุทฺทกุลํ.}

\prose{\textbf{กุเลสุ จาริตฺตํ อาปชฺเชยฺยา}ติ อญฺญสฺส ฆรูปจารํ โอกฺกมนฺตสฺส อาปตฺติ ทุกฺกฏสฺส.\csromanspacer{} ปฐมํ ปาทํ อุมฺมารํ อติกฺกาเมติ, อาปตฺติ ทุกฺกฏสฺส.\csromanspacer{} ทุติยํ ปาทํ อติกฺกาเมติ, อาปตฺติ ปาจิตฺติยสฺส.}

\prose{\csromanfolio{135}\textbf{อญฺญตฺร สมยา}ติ ฐเปตฺวา สมยํ.}

\prose{\textbf{จีวรทานสมโย} นาม อนตฺถเต กถิเน วสฺสานสฺส ปจฺฉิโม มาโส, อตฺถเต กถิเน ปญฺจ มาสา.}

\prose{\textbf{จีวรการสมโย} นาม จีวเร กยิรมาเน.}

\proseitem{301}{นิมนฺติเต นิมนฺติตสญฺญี สนฺตํ ภิกฺขุํ อนาปุจฺฉา ปุเรภตฺตํ วา ปจฺฉาภตฺตํ วา กุเลสุ จาริตฺตํ อาปชฺชติ อญฺญตฺร สมยา, อาปตฺติ ปาจิตฺติยสฺส.\csromanspacer{} นิมนฺติเต เวมติโก สนฺตํ ภิกฺขุํ อนาปุจฺฉา ปุเรภตฺตํ วา ปจฺฉาภตฺตํ วา กุเลสุ จาริตฺตํ อาปชฺชติ อญฺญตฺร สมยา, อาปตฺติ ปาจิตฺติยสฺส.\csromanspacer{} นิมนฺติเต อนิมนฺติตสญฺญี สนฺตํ ภิกฺขุํ อนาปุจฺฉา ปุเรภตฺตํ วา ปจฺฉาภตฺตํ กุเลสุ จาริตฺตํ อาปชฺชติ อญฺญตฺร สมยา, อาปตฺติ ปาจิตฺติยสฺส.}

\prose{อนิมนฺติเต นิมนฺติตสญฺญี, อาปตฺติ ทุกฺกฏสฺส.\csromanspacer{} อนิมนฺติเต เวมติโก, อาปตฺติ ทุกฺกฏสฺส.\csromanspacer{} อนิมนฺติเต อนิมนฺติตสญฺญี, อนาปตฺติ.}

\proseitemwithcloserruled{302}{อนาปตฺติ สมเย, สนฺตํ ภิกฺขุํ อาปุจฺฉา ปวิสติ, อสนฺตํ ภิกฺขุํ อนาปุจฺฉา ปวิสติ, อญฺญสฺส ฆเรน มคฺโค โหติ, ฆรูปจาเรน มคฺโค โหติ, อนฺตรารามํ คจฺฉติ, ภิกฺขุนุปสฺสยํ คจฺฉติ, ติตฺถิยเสยฺยํ คจฺฉติ, ปฏิกฺกมนํ คจฺฉติ, ภตฺติยฆรํ คจฺฉติ, อาปทาสุ อุมฺมตฺตกสฺส อาทิกมฺมิกสฺสาติ.}{จาริตฺตสิกฺขาปทํ นิฏฺฐิตํ ฉฏฺฐํ.}
\csromanmatikaanchor{mtk.53}
\csromantocmarkref{subsection}{7. มหานามสิกฺขาปท}{mtk.53}
\bookmark[dest={mtk.53},level=2]{7. มหานามสิกฺขาปท}
\csromanheader{2}{5. \textbf{อเจลกวคฺค 7. มหานามสิกฺขาปท}}
\proseitem{303}{เตน สมเยน พุทฺโธ ภควา สกฺเกสุ วิหรติ กปิลวตฺถุสฺมิํ นิโคฺรธาราเม.\csromanspacer{} เตน โข ปน สมเยน มหานามสฺส สกฺกสฺส เภสชฺชํ อุสฺสนฺนํ โหติ.\csromanspacer{} อถ โข มหานาโม สกฺโก เยน ภควา เตนุปสงฺกมิ, อุปสงฺกมิตฺวา ภควนฺตํ อภิวาเทตฺวา เอกมนฺตํ นิสีทิ, เอกมนฺตํ นิสินฺโน โข มหานาโม สกฺโก ภควนฺตํ เอตทโวจ “อิจฺฉามหํ ภนฺเต สํฆํ จตุมาสํ\footnote{จตุมฺมาสํ (สี) จาตุมาสํ (สฺยา)} เภสชฺเชน ปวาเรตุนฺ”ติ.\csromanspacer{} สาธุ \csromanfolio{136}สาธุ มหานาม, เตน หิ ตฺวํ มหานาม สํฆํ จตุมาสํ เภสชฺเชน ปวาเรหีติ.\csromanspacer{} ภิกฺขู กุกฺกุจฺจายนฺตา นาธิวาเสนฺติ.\csromanspacer{} ภควโต เอตมตฺถํ อาโรเจสุํ ฯเปฯ.\csromanspacer{} อนุชานามิ ภิกฺขเว จตุมาสํ เภสชฺชปฺปจฺจยปวารณํ สาทิตุนฺติ.}

\proseitem{304}{เตน โข ปน สมเยน ภิกฺขู มหานามํ สกฺกํ ปริตฺตํ เภสชฺชํ วิญฺญาเปนฺติ.\csromanspacer{} ตเถว มหานามสฺส สกฺกสฺส เภสชฺชํ อุสฺสนฺนํ โหติ.\csromanspacer{} ทุติยมฺปิ โข มหานาโม สกฺโก เยน ภควา เตนุปสงฺกมิ, อุปสงฺกมิตฺวา ภควนฺตํ อภิวาเทตฺวา เอกมนฺตํ นิสีทิ, เอกมนฺตํ นิสินฺโน โข มหานาโม สกฺโก ภควนฺตํ เอตทโวจ “อิจฺฉามหํ ภนฺเต สํฆํ อปรมฺปิ จตุมาสํ เภสชฺเชน ปวาเรตุนฺ”ติ.\csromanspacer{} สาธุ สาธุ มหานาม, เตน หิ ตฺวํ มหานาม สํฆํ อปรมฺปิ จตุมาสํ เภสชฺเชน ปวาเรหีติ.\csromanspacer{} ภิกฺขู กุกฺกุจฺจายนฺตา นาธิวาเสนฺติ.\csromanspacer{} ภควโต เอตมตฺถํ อาโรเจสุํ ฯเปฯ.\csromanspacer{} อนุชานามิ ภิกฺขเว ปุน ปวารณมฺปิ สาทิตุนฺติ.}

\proseitem{305}{เตน โข ปน สมเยน ภิกฺขู มหานามํ สกฺกํ ปริตฺตํ เยว เภสชฺชํ วิญฺญาเปนฺติ.\csromanspacer{} ตเถว มหานามสฺส สกฺกสฺส เภสชฺชํ อุสฺสนฺนํ โหติ.\csromanspacer{} ตติยมฺปิ โข มหานาโม สกฺโก เยน ภควา เตนุปสงฺกมิ, อุปสงฺกมิตฺวา ภควนฺตํ อภิวาเทตฺวา เอกมนฺตํ นิสีทิ, เอกมนฺตํ นิสินฺโน โข มหานาโม สกฺโก ภควนฺตํ เอตทโวจ “อิจฺฉามหํ ภนฺเต สํฆํ ยาวชีวํ เภสชฺเชน ปวาเรตุนฺ”ติ.\csromanspacer{} สาธุ สาธุ มหานาม, เตน หิ ตฺวํ มหานาม สํฆํ ยาวชีวํ เภสชฺเชน ปวาเรหีติ.\csromanspacer{} ภิกฺขู กุกฺกุจฺจายนฺตา นาธิวาเสนฺติ.\csromanspacer{} ภควโต เอตมตฺถํ อาโรเจสุํ ฯเปฯ.\csromanspacer{} อนุชานามิ ภิกฺขเว นิจฺจปวารณมฺปิ สาทิตุนฺติ.}

\prose{เตน โข ปน สมเยน ฉพฺพคฺคิยา ภิกฺขู ทุนฺนิวตฺถา โหนฺติ ทุปฺปารุตา อนากปฺปสมฺปนฺนา.\csromanspacer{} มหานาโม สกฺโก วตฺตา โหติ “กิสฺส ตุเมฺห ภนฺเต ทุนฺนิวตฺถา ทุปฺปารุตา อนากปฺปสมฺปนฺนา, นนุ นาม ปพฺพชิเตน สุนิวตฺเถน ภวิตพฺพํ สุปารุเตน อากปฺปสมฺปนฺเนนา”ติ.\csromanspacer{} ฉพฺพคฺคิยา ภิกฺขู มหานาเม สกฺเก อุปนนฺธิํสุ.\csromanspacer{} อถ โข ฉพฺพคฺคิยานํ ภิกฺขูนํ เอตทโหสิ “เกน นุ โข มยํ อุปาเยน มหานามํ สกฺกํ มงฺกุํ กเรยฺยามา”ติ.\csromanspacer{} อถ โข ฉพฺพคฺคิยานํ ภิกฺขูนํ เอตทโหสิ “มหานาเมน โข อาวุโส สกฺเกน สํโฆ เภสชฺเชน ปวาริโต, หนฺท มยํ อาวุโส มหานามํ \csromanfolio{137}สกฺกํ สปฺปิํ วิญฺญาเปมา”ติ.\csromanspacer{} อถ โข ฉพฺพคฺคิยา ภิกฺขู เยน มหานาโม สกฺโก เตนุปสงฺกมิํสุ, อุปสงฺกมิตฺวา มหานามํ สกฺกํ เอตทโวจุํ “โทเณน อาวุโส สปฺปินา อตฺโถ”ติ.\csromanspacer{} อชฺชโณฺห ภนฺเต อาคเมถ, มนุสฺสา วชํ คตา สปฺปิํ อาหริตุํ, กาลํ อาหริสฺสถาติ\footnote{กาเล หริสฺสถาติ (สฺยา)}.}

\prose{ทุติยมฺปิ โข ฯเปฯ.\csromanspacer{} ตติยมฺปิ โข ฉพฺพคฺคิยา ภิกฺขู มหานามํ สกฺกํ เอตทโวจุํ “โทเณน อาวุโส สปฺปินา อตฺโถ”ติ.\csromanspacer{} อชฺชโณฺห ภนฺเต อาคเมถ, มนุสฺสา วชํ คตา สปฺปิํ อาหริตุํ, กาลํ อาหริสฺสถาติ1.\csromanspacer{} กิํ ปน ตยา อาวุโส อทาตุกาเมน ปวาริเตน, ยํ ตฺวํ ปวาเรตฺวา น เทสีติ.\csromanspacer{} อถ โข มหานาโม สกฺโก อุชฺฌายติ ขิยฺยติ วิปาเจติ “กถํ หิ นาม ภทนฺตา ‘อชฺชโณฺห ภนฺเต อาคเมถา’ติ วุจฺจมานา นาคเมสฺสนฺตี”ติ.\csromanspacer{} อสฺโสสุํ โข ภิกฺขู มหานามสฺส สกฺกสฺส อุชฺฌายนฺตสฺส ขิยฺยนฺตสฺส วิปาเจนฺตสฺส.\csromanspacer{} เย เต ภิกฺขู อปฺปิจฺฉา ฯเปฯ เต อุชฺฌายนฺติ ขิยฺยนฺติ วิปาเจนฺติ “กถํ หิ นาม ฉพฺพคฺคิยา ภิกฺขู มหานาเมน สกฺเกน ‘อชฺชโณฺห ภนฺเต อาคเมถา’ติ วุจฺจมานา นาคเมสฺสนฺตี”ติ ฯเปฯ.\csromanspacer{} สจฺจํ กิร ตุเมฺห ภิกฺขเว มหานาเมน สกฺเกน “อชฺชโณฺห ภนฺเต อาคเมถา”ติ วุจฺจมานา นาคเมถาติ.\csromanspacer{} สจฺจํ ภควาติ.\csromanspacer{} วิครหิ พุทฺโธ ภควา ฯเปฯ.\csromanspacer{} กถํ หิ นาม ตุเมฺห โมฆปุริสา มหานาเมน สกฺเกน “อชฺชโณฺห ภนฺเต อาคเมถา”ติ วุจฺจมานา นาคเมสฺสถ.\csromanspacer{} เนตํ โมฆปุริสา อปฺปสนฺนานํ วา ปสาทาย ฯเปฯ.\csromanspacer{} เอวญฺจ ปน ภิกฺขเว อิมํ สิกฺขาปทํ อุทฺทิเสยฺยาถ\csromandash{}}

\proseitem{306}{\textbf{“อคิลาเนน ภิกฺขุนา จตุมาสปฺปจฺจยวารณา สาทิตพฺพา อญฺญตฺร ปุนปวารณาย อญฺญตฺร นิจฺจปวารณาย, ตโต เจ อุตฺตริ สาทิเยยฺย, ปาจิตฺติยนฺ”}ติ.}

\proseitem{307}{\textbf{อคิลาเนน ภิกฺขุนา จตุ}มาสปฺปจฺจยปวารณา สาทิตพฺพาติ คิลานปฺปจฺจยปวารณา สาทิตพฺพา.}

\prose{\textbf{ปุนปวารณาปิ สาทิตพฺพา}ติ ยทา คิลาโน ภวิสฺสามิ, ตทา วิญฺญาเปสฺสามีติ.}

\prose{\textbf{นิจฺจปวารณาปิ สาทิตพฺพา}ติ ยทา คิลาโน ภวิสฺสามิ, ตทา วิญฺญาเปสฺสามีติ.}

\prose{\csromanfolio{138}\textbf{ตโต เจ อุตฺตริ สาทิเยยฺยา}ติ อตฺถิ ปวารณา เภสชฺชปริยนฺตา น รตฺติปริยนฺตา, อตฺถิ ปวารณา รตฺติปริยนฺตา น เภสชฺชปริยนฺตา, อตฺถิ ปวารณา เภสชฺชปริยนฺตา จ รตฺติปริยนฺตา จ, อตฺถิ ปวารณา เนว เภสชฺชปริยนฺตา น รตฺติปริยนฺตา.}

\prose{\textbf{เภสชฺชปริยนฺตา} นาม เภสชฺชานิ ปริคฺคหิตานิ โหนฺติ “เอตฺตเกหิ เภสชฺเชหิ ปวาเรมี”ติ.\csromanspacer{} \textbf{รตฺติปริยนฺตา} นาม รตฺติโย ปริคฺคหิตาโย โหนฺติ “เอตฺตกาสุ รตฺตีสุ ปวาเรมี”ติ.\csromanspacer{} \textbf{เภสชฺชปริยนฺตา จ รตฺติปริยนฺตา จ} นาม เภสชฺชานิ จ ปริคฺคหิตานิ โหนฺติ รตฺติโย จ ปริคฺคหิตาโย โหนฺติ “เอตฺตเกหิ เภสชฺเชหิ เอตฺตกาสุ รตฺตีสุ ปวาเรมี”ติ.\csromanspacer{} \textbf{เนว เภสชฺชปริยนฺตา น รตฺติปริยนฺตา} นาม เภสชฺชานิ จ อปริคฺคหิตานิ โหนฺติ รตฺติโย จ อปริคฺคหิตาโย โหนฺติ.}

\proseitem{308}{เภสชฺชปริยนฺเต เยหิ เภสชฺเชหิ ปวาริโต โหติ, ตานิ เภสชฺชานิ ฐเปตฺวา อญฺญานิ เภสชฺชานิ วิญฺญาเปติ, อาปตฺติ ปาจิตฺติยสฺส.\csromanspacer{} รตฺติปริยนฺเต ยาสุ รตฺตีสุ ปวาริโต โหติ, ตา รตฺติโย ฐเปตฺวา อญฺญาสุ รตฺตีสุ วิญฺญาเปติ, อาปตฺติ ปาจิตฺติยสฺส.\csromanspacer{} เภสชฺชปริยนฺเต จ รตฺติปริยนฺเต จ เยหิ เภสชฺเชหิ ปวาริโต โหติ, ตานิ เภสชฺชานิ ฐเปตฺวา ยาสุ รตฺตีสุ ปวาริโต โหติ, ตา รตฺติโย ฐเปตฺวา อญฺญานิ เภสชฺชานิ อญฺญาสุ รตฺตีสุ วิญฺญาเปติ, อาปตฺติ ปาจิตฺติยสฺส.\csromanspacer{} เนว เภสชฺชปริยนฺเต น รตฺติปริยนฺเต อนาปตฺติ.}

\proseitem{309}{น เภสชฺเชน กรณีเยน\footnote{กรณีเย (สี, สฺยา)} เภสชฺชํ วิญฺญาเปติ, อาปตฺติ ปาจิตฺติยสฺส.\csromanspacer{} อญฺเญน เภสชฺเชน กรณีเยน1 อญฺญํ เภสชฺชํ วิญฺญาเปติ, อาปตฺติ ปาจิตฺติยสฺส.\csromanspacer{} ตตุตฺตริ ตตุตฺตริสญฺญี เภสชฺชํ วิญฺญาเปติ, อาปตฺติ ปาจิตฺติยสฺส.\csromanspacer{} ตตุตฺตริ เวมติโก เภสชฺชํ วิญฺญาเปติ, อาปตฺติ ปาจิตฺติยสฺส.\csromanspacer{} ตตุตฺตริ นตตุตฺตริสญฺญี เภสชฺชํ วิญฺญาเปติ อาปตฺติ ปาจิตฺติยสฺส.}

\prose{นตตุตฺตริ ตตุตฺตริสญฺญี, อาปตฺติ ทุกฺกฏสฺส.\csromanspacer{} นตตุตฺตริ เวมติโก, อาปตฺติ ทุกฺกฏสฺส.\csromanspacer{} นตตุตฺตริ นตตุตฺตริสญฺญี, อนาปตฺติ.}

\proseitemwithcloserruled{310}{\csromanfolio{139}อนาปตฺติ เยหิ เภสชฺเชหิ ปวาริโต โหติ ตานิ เภสชฺชานิ วิญฺญาเปติ, ยาสุ รตฺตีสุ ปวาริโต โหติ ตาสุ รตฺติสุ วิญฺญาเปติ, “อิเมหิ ตยา เภสชฺเชหิ ปวาริตามฺห, อมฺหากญฺจ อิมินา จ อิมินา จ เภสชฺเชน อตฺโถ”ติ อาจิกฺขิตฺวา วิญฺญาเปติ, “ยาสุ ตยา รตฺตีสุ ปวาริตามฺห, ตาโย จ รตฺติโย วีติวตฺตา, อมฺหากญฺจ เภสชฺเชน อตฺโถ”ติ อาจิกฺขิตฺวา วิญฺญาเปติ, ญาตกานํ ปวาริตานํ อญฺญสฺสตฺถาย อตฺตโน ธเนน อุมฺมตฺตกสฺส อาทิกมฺมิกสฺสาติ.}{มหานามสิกฺขาปทํ นิฏฺฐิตํ สตฺตมํ.}
\csromanmatikaanchor{mtk.54}
\csromantocmarkref{subsection}{8. อุยฺยุตฺตเสนาสิกฺขาปท}{mtk.54}
\bookmark[dest={mtk.54},level=2]{8. อุยฺยุตฺตเสนาสิกฺขาปท}
\csromanheader{2}{5. อเจลกวคฺค 8. อุยฺยุตฺต\-เสนา\-สิกฺขาปท}
\proseitem{311}{เตน สมเยน พุทฺโธ ภควา สาวตฺถิยํ วิหรติ เชตวเน อนาถปิณฺฑิกสฺส อาราเม.\csromanspacer{} เตน โข ปน สมเยน ราชา ปเสนทิ โกสโล เสนาย อพฺภุยฺยาโต โหติ.\csromanspacer{} ฉพฺพคฺคิยา ภิกฺขู อุยฺยุตฺตํ เสนํ ทสฺสนาย อคมํสุ.\csromanspacer{} อทฺทสา โข ราชา ปเสนทิ โกสโล ฉพฺพคฺคิเย ภิกฺขู ทูรโตว อาคจฺฉนฺเต, ทิสฺวาน ปกฺโกสาเปตฺวา เอตทโวจ “กิสฺส ตุเมฺห ภนฺเต อาคตตฺถา”ติ.\csromanspacer{} มหาราชานํ มยํ ทฏฺฐุกามา\footnote{มหาราช มหาราชานํ มยํ ทฏฺฐุกามา (ก)} ติ.\csromanspacer{} กิํ ภนฺเต มํ ทิฏฺเฐน ยุทฺธาภินนฺทินํ\footnote{ยุทฺธาภินนฺทินา (ก)}, นนุ ภควา ปสฺสิตพฺโพติ.\csromanspacer{} มนุสฺสา อุชฺฌายนฺติ ขิยฺยนฺติ วิปาเจนฺติ “กถํ หิ นาม สมณา สกฺยปุตฺติยา อุยฺยุตฺตํ เสนํ ทสฺสนาย อาคจฺฉิสฺสนฺติ.\csromanspacer{} อมฺหากมฺปิ อลาภา อมฺหากมฺปิ ทุลฺลทฺธํ, เย มยํ อาชีวสฺส เหตุ ปุตฺตทารสฺส การณา เสนาย อาคจฺฉามา”ติ.\csromanspacer{} อสฺโสสุํ โข ภิกฺขู เตสํ มนุสฺสานํ อุชฺฌายนฺตานํ ขิยฺยนฺตานํ วิปาเจนฺตานํ.\csromanspacer{} เย เต ภิกฺขู อปฺปิจฺฉา ฯเปฯ เต อุชฺฌายนฺติ ขิยฺยนฺติ วิปาเจนฺติ “กถํ หิ นาม ฉพฺพคฺคิยา ภิกฺขู อุยฺยุตฺตํ เสนํ ทสฺสนาย คจฺฉิสฺสนฺตี”ติ ฯเปฯ.\csromanspacer{} สจฺจํ กิร ตุเมฺห ภิกฺขเว อุยฺยุตฺตํ เสนํ ทสฺสนาย คจฺฉถาติ.\csromanspacer{} สจฺจํ ภควาติ.\csromanspacer{} วิครหิ พุทฺโธ ภควา ฯเปฯ.\csromanspacer{} กถํ หิ นาม ตุเมฺห โมฆปุริสา อุยฺยุตฺตํ เสนํ ทสฺสนาย คจฺฉิสฺสถ.\csromanspacer{} เนตํ โมฆปุริสา อปฺปสนฺนานํ วา ปสาทาย ฯเปฯ.\csromanspacer{} เอวญฺจ ปน ภิกฺขเว อิมํ สิกฺขาปทํ อุทฺทิเสยฺยาถ\csromandash{}}

\prose{\csromanfolio{140}\textbf{“โย ปน ภิกฺขุ อุยฺยุตฺตํ เสนํ ทสฺสนาย คจฺเฉยฺย, ปาจิตฺติยนฺ”}ติ.}

\prose{เอวญฺจิทํ ภควตา ภิกฺขูนํ สิกฺขาปทํ ปญฺญตฺตํ โหติ.}

\proseitem{312}{เตน โข ปน สมเยน อญฺญตรสฺส ภิกฺขุโน มาตุโล เสนาย คิลาโน โหติ.\csromanspacer{} โส ตสฺส ภิกฺขุโน สนฺติเก ทูตํ ปาเหสิ “อหํ หิ เสนาย คิลาโน, อาคจฺฉตุ ภทนฺโต, อิจฺฉามิ ภทนฺตสฺส อาคตนฺ”ติ.\csromanspacer{} อถ โข ตสฺส ภิกฺขุโน เอตทโหสิ “ภควตา สิกฺขาปทํ ปญฺญตฺตํ ‘น อุยฺยุตฺตํ เสนํ ทสฺสนาย คนฺตพฺพนฺ’ติ, อยญฺจ เม มาตุโล เสนาย คิลาโน, กถํ นุ โข มยา ปฏิปชฺชิตพฺพนฺ”ติ.\csromanspacer{} ภควโต เอตมตฺถํ อาโรเจสิ.\csromanspacer{} อถ โข ภควา เอตสฺมิํ นิทาเน เอตสฺมิํ ปกรเณ ธมฺมิํ กถํ กตฺวา ภิกฺขู อามนฺเตสิ อนุชานามิ ภิกฺขเว ตถารูปปฺปจฺจยา เสนาย คนฺตุํ.\csromanspacer{} เอวญฺจ ปน ภิกฺขเว อิมํ สิกฺขาปทํ อุทฺทิเสยฺยาถ\csromandash{}}

\proseitem{313}{“\textbf{โย ปน ภิกฺขุ อุยฺยุตฺตํ เสนํ ทสฺสนาย คจฺฉฺ}เอยฺย อญฺญตฺร ตถารูปปฺปจฺจยา, ปาจิตฺติยนฺ”ติ.}

\proseitem{314}{\textbf{โย ปนา}ติ โย ยาทิโส ฯเปฯ.\csromanspacer{} \textbf{ภิกฺขู}ติ ฯเปฯ อยํ อิมสฺมิํ อตฺเถ อธิปฺเปโต ภิกฺขูติ.}

\prose{\textbf{อุยฺยุตฺตา} นาม เสนา คามโต นิกฺกมิตฺวา นิวิฏฺฐา วา โหติ ปยาตา วา.}

\prose{\textbf{เสนา} นาม หตฺถี อสฺสา รถา ปตฺตี.\csromanspacer{} ทฺวาทสปุริโส หตฺถี, ติปุริโส อสฺโส, จตุปุริโส รโถ, จตฺตาโร ปุริสา สรหตฺถา ปตฺติ.\csromanspacer{} ทสฺสนาย คจฺฉติ, อาปตฺติ ทุกฺกฏสฺส.\csromanspacer{} ยตฺถ ฐิโต ปสฺสติ, อาปตฺติ ปาจิตฺติยสฺส.\csromanspacer{} ทสฺสนูปจารํ วิชหิตฺวา ปุนปฺปุนํ ปสฺสติ, อาปตฺติ ปาจิตฺติยสฺส.}

\prose{\textbf{อญฺญตฺร ตถารูปปฺปจฺจยา}ติ ฐเปตฺวา ตถารูปปฺปจฺจยํ.}

\proseitem{315}{อุยฺยุตฺเต อุยฺยุตฺตสญฺญี ทสฺสนาย คจฺฉติ อญฺญตฺร ตถารูปปฺปจฺจยา, อาปตฺติ ปาจิตฺติยสฺส.\csromanspacer{} อุยฺยุตฺเต เวมติโก ทสฺสนาย คจฺฉติ อญฺญตฺร ตถารูปปฺปจฺจยา, อาปตฺติ ปาจิตฺติยสฺส.\csromanspacer{} อุยฺยุตฺเต อนุยฺยุตฺตสญฺญี ทสฺสนาย คจฺฉติ อญฺญตฺร ตถารูปปฺปจฺจยา, อาปตฺติ ปาจิตฺติยสฺส.}

\prose{\csromanfolio{141}เอกเมกํ ทสฺสนาย คจฺฉติ, อาปตฺติ ทุกฺกฏสฺส.\csromanspacer{} ยตฺถ ฐิโต ปสฺสติ, อาปตฺติ ทุกฺกฏสฺส.\csromanspacer{} ทสฺสนูปจารํ วิชหิตฺวา ปุนปฺปุนํ ปสฺสติ, อาปตฺติ ทุกฺกฏสฺส.\csromanspacer{} อนุยฺยุตฺเต อุยฺยุตฺตสญฺญี, อาปตฺติ ทุกฺกฏสฺส.\csromanspacer{} อนุยฺยุตฺเต เวมติโก, อาปตฺติ ทุกฺกฏสฺส.\csromanspacer{} อนุยฺยุตฺเต อนุยฺยุตฺตสญฺญี, อนาปตฺติ.}

\proseitemwithcloserruled{316}{อนาปตฺติ อาราเม ฐิโต ปสฺสติ, ภิกฺขุสฺส ฐิโตกาสํ วา นิสินฺโนกาสํ วา นิปนฺโนกาสํ วา อาคจฺฉติ, ปฏิปถํ คจฺฉนฺโต ปสฺสติ, ตถารูปปฺปจฺจยา อาปทาสุ อุมฺมตฺตกสฺส อาทิกมฺมิกสฺสาติ.}{อุยฺยุตฺต\-เสนา\-สิกฺขาปทํ นิฏฺฐิตํ อฏฺฐมํ.}
\csromanmatikaanchor{mtk.55}
\csromantocmarkref{subsection}{9. เสนาวาสสิกฺขาปท}{mtk.55}
\bookmark[dest={mtk.55},level=2]{9. เสนาวาสสิกฺขาปท}
\csromanheader{2}{5. \textbf{อเจลกวคฺค 9. เสนาวาสสิกฺขาปท}}
\proseitem{317}{เตน สมเยน พุทฺโธ ภควา สาวตฺถิยํ วิหรติ เชตวเน อนาถปิณฺฑิกสฺส อาราเม.\csromanspacer{} เตน โข ปน สมเยน ฉพฺพคฺคิยา ภิกฺขู สติ กรณีเย เสนํ คนฺตฺวา อติเรกติรตฺตํ เสนาย วสนฺติ.\csromanspacer{} มนุสฺสา อุชฺฌายนฺติ ขิยฺยนฺติ วิปาเจนฺติ “กถํ หิ นาม สมณา สกฺยปุตฺติยา เสนาย วสิสฺสนฺติ, อมฺหากมฺปิ อลาภา อมฺหากมฺปิ ทุลฺลทฺธํ, เย มยํ อาชีวสฺส เหตุ ปุตฺตทารสฺส การณา เสนาย ปฏิวสามา”ติ.\csromanspacer{} อสฺโสสุํ โข ภิกฺขู เตสํ มนุสฺสานํ อุชฺฌายนฺตานํ ขิยฺยนฺตานํ วิปาเจนฺตานํ.\csromanspacer{} เย เต ภิกฺขู อปฺปิจฺฉา ฯเปฯ เต อุชฺฌายนฺติ ขิยฺยนฺติ วิปาเจนฺติ “กถํ หิ นาม ฉพฺพคฺคิยา ภิกฺขู อติเรกติรตฺตํ เสนาย วสิสฺสนฺตี”ติ ฯเปฯ.\csromanspacer{} สจฺจํ กิร ตุเมฺห ภิกฺขเว อติเรกติรตฺตํ เสนาย วสถาติ.\csromanspacer{} สจฺจํ ภควาติ.\csromanspacer{} วิครหิ พุทฺโธ ภควา ฯเปฯ.\csromanspacer{} กถํ หิ นาม ตุเมฺห โมฆปุริสา อติเรกติรตฺตํ เสนาย วสิสฺสถ.\csromanspacer{} เนตํ โมฆปุริสา อปฺปสนฺนานํ วา ปสาทาย ฯเปฯ.\csromanspacer{} เอวญฺจ ปน ภิกฺขเว อิมํ สิกฺขาปทํ อุทฺทิเสยฺยาถ\csromandash{}}

\proseitem{318}{\textbf{“สิยา จ ตสฺส ภิกฺขุโน โกจิเทว ปจฺจโย เสนํ คมนาย, ทิรตฺตติรตฺตํ เตน ภิกฺขุนา เสนาย วสิตพฺพํ.}\csromanspacer{} ตโต เจ อุตฺตริ วเสยฺย, ปาจิตฺติยนฺ”ติ.}

\proseitem{319}{\textbf{สิยา จ ตสฺส ภิกฺขุโน โกจิเทว ปจฺจยฺ}โอ เสนํ คมนายาติ สิยา ปจฺจโย สิยา กรณียํ.}

\prose{\csromanfolio{142}\textbf{ทิรตฺตติรตฺตํ เตน ภิกฺขุนา เสนาย วสิตพฺพนฺ}ติ เทฺวติสฺโส รตฺติโย วสิตพฺพํ.}

\prose{\textbf{ตโต เจ อุตฺตริ วเสยฺยา}ติ จตุตฺเถ ทิวเส อตฺถงฺคเต สูริเย เสนาย วสติ, อาปตฺติ ปาจิตฺติยสฺส.}

\proseitem{320}{อติเรกติรตฺเต อติเรกสญฺญี เสนาย วสติ, อาปตฺติ ปาจิตฺติยสฺส.\csromanspacer{} อติเรกติรตฺเต เวมติโก เสนาย วสติ, อาปตฺติ ปาจิตฺติยสฺส.\csromanspacer{} อติเรกติรตฺเต อูนกสญฺญี เสนาย วสติ, อาปตฺติ ปาจิตฺติยสฺส.}

\prose{อูนกติรตฺเต อติเรกสญฺญี, อาปตฺติ ทุกฺกฏสฺส.\csromanspacer{} อูนกติรตฺเต เวมติโก, อาปตฺติ ทุกฺกฏสฺส.\csromanspacer{} อูนกติรตฺเต อูนกสญฺญี, อานาปตฺติ.}

\proseitemwithcloserruled{321}{อนาปตฺติ เทฺวติสฺโส รตฺติโย วสติ, อูนกเทฺวติสฺโส รตฺติโย วสติ, เทฺว รตฺติโย วสิตฺวา ตติยาย รตฺติยา ปุรารุณา นิกฺขมิตฺวา ปุน วสติ, คิลาโน วสติ, คิลานสฺส กรณีเยน วสติ, เสนา วา ปฏิเสนาย รุทฺธา โหติ, เกนจิ ปลิพุทฺโธ โหติ อาปทาสุ อุมฺมตฺตกสฺส อาทิกมฺมิกสฺสาติ.}{เสนาวาสสิกฺขาปทํ นิฏฺฐิตํ นวมํ.}
\csromanmatikaanchor{mtk.56}
\csromantocmarkref{subsection}{10. อุยฺโยธิกสิกฺขาปท}{mtk.56}
\bookmark[dest={mtk.56},level=2]{10. อุยฺโยธิกสิกฺขาปท}
\csromanheader{2}{5. \textbf{อเจลกวคฺค 1}0. อุยฺโยธิกสิกฺขาปท}
\proseitem{322}{เตน สมเยน พุทฺโธ ภควา สาวตฺถิยํ วิหรติ เชตวเน อนาถปิณฺฑิกสฺส อาราเม.\csromanspacer{} เตน โข ปน สมเยน ฉพฺพคฺคิยา ภิกฺขู ทิรตฺตติรตฺตํ เสนาย วสมานา อุยฺโยธิกมฺปิ พลคฺคมฺปิ เสนาพฺยูหมฺปิ อนีกทสฺสนมฺปิ คจฺฉนฺติ.\csromanspacer{} อญฺญตโรปิ ฉพฺพคฺคิโย ภิกฺขุ อุยฺโยธิกํ คนฺตฺวา กณฺเฑน ปฏิวิทฺโธ โหติ.\csromanspacer{} มนุสฺสา ตํ ภิกฺขุํ อุปฺปณฺเฑสุํ “กจฺจิ ภนฺเต สุยุทฺธํ อโหสิ, กติ เต ลกฺขานิ ลทฺธานี”ติ.\csromanspacer{} โส ภิกฺขุ เตหิ มนุสฺเสหิ อุปฺปณฺฑียมาโน มงฺกุ อโหสิ.\csromanspacer{} มนุสฺสา อุชฺฌายนฺติ ขิยฺยนฺติ วิปาเจนฺติ “กถํ หิ นาม สมณา สกฺยปุตฺติยา อุยฺโยธิกํ ทสฺสนาย อาคจฺฉิสฺสนฺติ, อมฺหากมฺปิ อลาภา อมฺหากมฺปิ ทุลฺลทฺธํ, เย มยํ อาชีวสฺส \csromanfolio{143}เหตุ ปุตฺตทารสฺส การณา อุยฺโยธิกํ อาคจฺฉามา”ติ.\csromanspacer{} อสฺโสสุํ โข ภิกฺขู เตสํ มนุสฺสานํ อุชฺฌายนฺตานํ ขิยฺยนฺตานํ วิปาเจนฺตานํ.\csromanspacer{} เย เต ภิกฺขู อปฺปิจฺฉา ฯเปฯ เต อุชฺฌายนฺติ ขิยฺยนฺติ วิปาเจนฺติ “กถํ หิ นาม ฉพฺพคฺคิยา ภิกฺขู อุยฺโยธิกํ ทสฺสนาย คจฺฉิสฺสนฺตี”ติ ฯเปฯ.\csromanspacer{} สจฺจํ กิร ตุเมฺห ภิกฺขเว อุยฺโยธิกํ ทสฺสนาย คจฺฉถาติ.\csromanspacer{} สจฺจํ ภควาติ.\csromanspacer{} วิครหิ พุทฺโธ ภควา ฯเปฯ.\csromanspacer{} กถํ หิ นาม ตุเมฺห โมฆปุริสา อุยฺโยธิกํ ทสฺสนาย คจฺฉิสฺสถ.\csromanspacer{} เนตํ โมฆปุริสา อปฺปสนฺนานํ วา ปสาทาย ฯเปฯ.\csromanspacer{} เอวญฺจ ปน ภิกฺขเว อิมํ สิกฺขาปทํ อุทฺทิเสยฺยาถ\csromandash{}}

\proseitem{323}{\textbf{“ทิรตฺตติรตฺตํ เจ ภิกฺขุ เสนาย วสมาโน อุยฺโยธิกํ วา พลคฺคํ วา เสนาพฺยูหํ วา อนีกทสฺสนํ วา คจฺเฉยฺย, ปาจิตฺติยนฺ”}ติ.}

\proseitem{324}{\textbf{ทิรตฺตติรตฺตํ เจ ภิกฺขุ เสนาย วสมฺ}อาโนติ เทฺวติสฺโส รตฺติโย วสมาโน.}

\prose{\textbf{อุยฺโยธิกํ} นาม ยตฺถ สมฺปหาโร ทิสฺสติ.}

\prose{\textbf{พลคฺคํ} นาม เอตฺตกา หตฺถี เอตฺตกา อสฺสา เอตฺตกา รถา เอตฺตกา ปตฺตี.}

\prose{\textbf{เสนาพฺยูหํ} นาม อิโต หตฺถี โหนฺตุ, อิโต อสฺสา โหนฺตุ, อิโต รถา โหนฺตุ, อิโต ปตฺติกา โหนฺตุ.}

\prose{\textbf{อนีกํ} นาม หตฺถานีกํ อสฺสานีกํ รถานีกํ ปตฺตานีกํ.\csromanspacer{} ตโย หตฺถี ปจฺฉิมํ หตฺถานีกํ, ตโย อสฺสา ปจฺฉิมํ อสฺสานีกํ, ตโย รถา ปจฺฉิมํ รถานีกํ, จตฺตาโร ปุริสา สรหตฺถา ปตฺตี ปจฺฉิมํ ปตฺตานีกํ.\csromanspacer{} ทสฺสนาย คจฺฉติ, อาปตฺติ ทุกฺกฏสฺส.\csromanspacer{} ยตฺถ ฐิโต ปสฺสติ, อาปตฺติ ปาจิตฺติยสฺส.\csromanspacer{} ทสฺสนูปจารํ วิชหิตฺวา ปุนปฺปุนํ ปสฺสติ, อาปตฺติ}

\prose{เอกเมกํ ทสฺสนาย คจฺฉติ.\csromanspacer{} อาปตฺติ ทุกฺกฏสฺส.\csromanspacer{} ยตฺถ ฐิโต ปสฺสติ, อาปตฺติ ทุกฺกฏสฺส.\csromanspacer{} ทสฺสนูปจารํ วิชหิตฺวา ปุนปฺปุนํ ปสฺสติ, อาปตฺติ ทุกฺกฏสฺส.}

\proseitemwithcloser{325}{\csromanfolio{144}อนาปตฺติ อาราเม ฐิโต ปสฺสติ, ภิกฺขุสฺส ฐิโตกาสํ วา นิสินฺโนกาสํ วา นิปนฺโนกาสํ วา อาคนฺตฺวา สมฺปหาโร ทิสฺสติ, ปฏิปถํ คจฺฉนฺโต ปสฺสติ, สติ กรณีเย คนฺตฺวา ปสฺสติ, อาปทาสุ อุมฺมตฺตกสฺส อาทิกมฺมิกสฺสาติ.}{อุยฺโยธิกสิกฺขาปทํ นิฏฺฐิตํ ทสมํ.}
\vspace{\tipitakaparskip}
\csromancenter{อเจลกวคฺโค ปญฺจโม.}
\csromancenter{ตสฺสุทฺทานํ}
\vspace{0.5\baselineskip}
\setlength{\csromangathapagewidth}{0pt}
\csromangathameasuregroup{ปูวํ กโถปนนฺทสฺส,}{\csromangathabat{ปูวํ กโถปนนฺทสฺส,}{ตยํปฏฺฐากเมว จ.} \\ มหานาโม ปเสนทิ, เสนาวิทฺโธ อิเม ทสาติ.}
\csromangathasetleft{ปูวํ กโถปนนฺทสฺส,}
\csromangathagroup{\csromangathastanza{\csromangathabat{ปูวํ กโถปนนฺทสฺส,}{ตยํปฏฺฐากเมว จ.} \\ มหานาโม ปเสนทิ, เสนาวิทฺโธ อิเม ทสาติ\footnote{อเจลกํ อุยฺโยชญฺจ, สโภชนํ ทุเว รโห. สภตฺตกญฺจ เภสชฺชํ, อุยฺยุตฺตํ เสนุยฺโยธิกํ. (สี)}.}}

\csromanmatikaanchor{mtk.57}
\csromanmatikaanchor{mtk.58}
\csromantocmarknopage{section}{6. สุราปาณวคฺค}{mtk.57}
\bookmark[dest={mtk.57},level=1]{6. สุราปาณวคฺค}
\csromantocmarkref{subsection}{1. สุราปานสิกฺขาปท}{mtk.58}
\bookmark[dest={mtk.58},level=2]{1. สุราปานสิกฺขาปท}
\setcsromanheadh{6. สุราปาณวคฺค}
\csromanheader{2}{6. \textbf{สุราปานวคฺค 1. สุราปานสิกฺขาปท}}
\proseitem{326}{เตน สมเยน พุทฺโธ ภควา เจติเยสุ จาริกํ จรมาโน เยน ภทฺทวติกา เตน ปายาสิ.\csromanspacer{} อทฺทสํสุ โข โคปาลกา ปสุปาลกา กสฺสกา ปถาวิโน ภควนฺตํ ทูรโตว อาคจฺฉนฺตํ, ทิสฺวาน ภควนฺตํ เอตทโวจุํ “มา โข ภนฺเต ภควา อมฺพติตฺถํ อคมาสิ, อมฺพติตฺเถ ภนฺเต ชฏิลสฺส อสฺสเม นาโค ปฏิวสติ อิทฺธิมา อาสิวิโส\footnote{อาสีวิโส (สี, สฺยา)} โฆรวิโส, โส ภควนฺตํ มา วิเหเฐสี”ติ.\csromanspacer{} เอวํ วุตฺเต ภควา ตุณฺหี อโหสิ.\csromanspacer{} ทุติยมฺปิ โข ฯเปฯ.\csromanspacer{} ตติยมฺปิ โข โคปาลกา ปสุปาลกา กสฺสกา ปถาวิโน ภควนฺตํ เอตทโวจุํ “มา โข ภนฺเต ภควา อมฺพติตฺถํ อคมาสิ, อมฺพติตฺเถ ภนฺเต ชฏิลสฺส อสฺสเม นาโค ปฏิวสติ อิทฺธิมา อาสิวิโส โฆรวิโส, โส ภควนฺตํ มา วิเหเฐสี”ติ.\csromanspacer{} ตติยมฺปิ โข ภควา ตุณฺหี อโหสิ.}

\prose{อถ โข ภควา อนุปุพฺเพน จาริกํ จรมาโน เยน ภทฺทวติกา ตทวสริ.\csromanspacer{} ตตฺร สุทํ ภควา ภทฺทวติกายํ วิหรติ.\csromanspacer{} อถ โข อายสฺมา สาคโต เยน อมฺพติตฺถสฺส\footnote{อมฺพติตฺถกสฺส (สี), อมฺพติตฺถํ (สฺยา)} ชฏิลสฺส อสฺสโม เตนุปสงฺกมิ \csromanfolio{145}, อุปสงฺกมิตฺวา อคฺยาคารํ ปวิสิตฺวา ติณสนฺถารกํ ปญฺญเปตฺวา นิสีทิ ปลฺลงฺกํ อาภุชิตฺวา อุชุํ กายํ ปณิธาย ปริมุขํ สติํ อุปฏฺฐเปตฺวา.\csromanspacer{} อทฺทสา โข โส นาโค อายสฺมนฺตํ สาคตํ ปวิฏฺฐํ, ทิสฺวาน ทุมฺมโน\footnote{ทุกฺขี ทุมฺมโน (สี, สฺยา)} ปธูปายิ\footnote{ปธูปาสิ (สฺยา, ก)}, อายสฺมาปิ สาคโต ปธูปายิ2.\csromanspacer{} อถ โข โส นาโค มกฺขํ อสหมาโน ปชฺชลิ, อายสฺมาปิ สาคโต เตโชธาตุํ สมาปชฺชิตฺวา ปชฺชลิ.\csromanspacer{} อถ โข อายสฺมา สาคโต ตสฺส นาคสฺส เตชสา เตชํ ปริยาทิยิตฺวา เยน ภทฺทวติกา เตนุปสงฺกมิ.\csromanspacer{} อถ โข ภควา ภทฺทวติกายํ ยถาภิรนฺตํ วิหริตฺวา เยน โกสมฺพี เตน จาริกํ ปกฺกามิ.\csromanspacer{} อสฺโสสุํ โข โกสมฺพิกา อุปาสกา “อยฺโย กิร สาคโต อมฺพติตฺถิเกน นาเคน สทฺธิํ สงฺคาเมสี”ติ.}

\prose{อถ โข ภควา อนุปุพฺเพน จาริกํ จรมาโน เยน โกสมฺพี ตทวสริ.\csromanspacer{} อถ โข โกสมฺพิกา อุปาสกา ภควโต ปจฺจุคฺคมนํ กริตฺวา เยนายสฺมา สาคโต เตนุปสงฺกมิํสุ, อุปสงฺกมิตฺวา อายสฺมนฺตํ สาคตํ อภิวาเทตฺวา เอกมนฺตํ อฏฺฐํสุ, เอกมนฺตํ ฐิตา โข โกสมฺพิกา อุปาสกา อายสฺมนฺตํ สาคตํ เอตทโวจุํ “กิํ ภนฺเต อยฺยานํ ทุลฺลภญฺจ มนาปญฺจ, กิํ ปฏิยาเทมา”ติ.\csromanspacer{} เอวํ วุตฺเต ฉพฺพคฺคิยา ภิกฺขู โกสมฺพิเก อุปาสเก เอตทโวจุํ “อตฺถาวุโส กาโปติกา นาม ปสนฺนา ภิกฺขูนํ ทุลฺลภา จ มนาปา จ, ตํ ปฏิยาเทถา”ติ.\csromanspacer{} อถ โข โกสมฺพิกา อุปาสกา ฆเร ฆเร กาโปติกํ ปสนฺนํ ปฏิยาเทตฺวา อายสฺมนฺตํ สาคตํ ปิณฺฑาย ปวิฏฺฐํ ทิสฺวาน อายสฺมนฺตํ สาคตํ เอตทโวจุํ “ปิวตุ ภนฺเต อยฺโย สาคโต กาโปติกํ ปสนฺนํ, ปิวตุ ภนฺเต อยฺโย สาคโต กาโปติกํ ปสนฺนนฺ”ติ.\csromanspacer{} อถ โข อายสฺมา สาคโต ฆเร ฆเร กาโปติกํ ปสนฺนํ ปิวิตฺวา นครมฺหา นิกฺขมนฺโต นครทฺวาเร ปริปติ.}

\prose{อถ โข ภควา สมฺพหุเลหิ ภิกฺขูหิ สทฺธิํ นครมฺหา นิกฺขมนฺโต อทฺทส อายสฺมนฺตํ สาคตํ นครทฺวาเร ปริปตนฺตํ, ทิสฺวาน ภิกฺขู อามนฺเตสิ “คณฺหถ ภิกฺขเว สาคตนฺ”ติ. “เอวํ ภนฺเต”ติ โข เต ภิกฺขู ภควโต ปฏิสฺสุณิตฺวา อายสฺมนฺตํ สาคตํ อารามํ เนตฺวา เยน ภควา เตน สีสํ กตฺวา นิปาเตสุํ.\csromanspacer{} อถ โข อายสฺมา สาคโต ปริวตฺติตฺวา เยน ภควา เตน ปาเท กริตฺวา เสยฺยํ \csromanfolio{146}\textbf{กปฺเปสิ.}\csromanspacer{}\textbf{ อถ โข ภควา ภิกฺขู อามนฺเตสิ “นนุ ภิกฺขเว ปุพฺเพ} \textbf{สาคโต ตถาคเต สคารโว อโหสิ สปฺปติสฺโส”ติ.}\csromanspacer{}\textbf{ เอวํ ภนฺเต.}\csromanspacer{}\textbf{ อปิ นุ โข} ภิกฺขเว สาคโต เอตรหิ ตถาคเต สคารโว สปฺปติสฺโสติ.\csromanspacer{} โน เหตํ ภนฺเต.\csromanspacer{} นนุ ภิกฺขเว สาคโต อมฺพติตฺถิเกน นาเคน สทฺธิํ สงฺคาเมสีติ.\csromanspacer{} เอวํ ภนฺเต.\csromanspacer{} อปิ นุ โข ภิกฺขเว สาคโต เอตรหิ ปโหติ นาเคน\footnote{qเอฑฺฑุเภนาปิ (สี, สฺยา)} สทฺธิํ สงฺคาเมตุนฺติ.\csromanspacer{} โน เหตํ ภนฺเต.\csromanspacer{} อปิ นุ โข ภิกฺขเว ตํ ปาตพฺพํ, ยํ ปิวิตฺวา วิสญฺญี อสฺสาติ.\csromanspacer{} โน เหตํ ภนฺเต.\csromanspacer{} อนนุจฺฉวิกํ ภิกฺขเว สาคตสฺส อนนุโลมิตํ อปฺปติรูปํ อสฺสามณกํ อกปฺปิยํ อกรณียํ, กถํ หิ นาม ภิกฺขเว สาคโต มชฺชํ ปิวิสฺสติ.\csromanspacer{} เนตํ ภิกฺขเว อปฺปสนฺนานํ วา ปสาทาย ฯเปฯ.\csromanspacer{} เอวญฺจ ปน ภิกฺขเว อิมํ สิกฺขาปทํ อุทฺทิเสยฺยาถ\csromandash{}}

\proseitem{327}{\textbf{“สุราเมรยปาเน ปาจิตฺติยนฺ”}ติ.}

\proseitem{328}{\textbf{สุรา} นาม ปิฏฺฐสุรา ปูวสุรา โอทนสุรา กิณฺณปกฺขิตฺตา สมฺภารสํยุตฺตา.}

\prose{\textbf{เมรโย} นาม ปุปฺผาสโว ผลาสโว มธฺวาสโว คุฬาสโว สมฺภารสํยุตฺโต.}

\prose{\textbf{ปิเวยฺยา}ติ อนฺตมโส กุสคฺเคนปิ ปิวติ, อาปตฺติ ปาจิตฺติยสฺส.}

\prose{มชฺเช มชฺชสญฺญี ปิวติ, อาปตฺติ ปาจิตฺติยสฺส.\csromanspacer{} มชฺเช เวมติโก ปิวติ, อาปตฺติ ปาจิตฺติยสฺส.\csromanspacer{} มชฺเช อมชฺชสญฺญี ปิวติ, อาปตฺติ ปาจิตฺติยสฺส.}

\prose{อมชฺเช มชฺชสญฺญี, อาปตฺติ ทุกฺกฏสฺส.\csromanspacer{} อมชฺเช เวมติโก, อาปตฺติ ทุกฺกฏสฺส.\csromanspacer{} อมชฺเช อมชฺชสญฺญี, อนาปตฺติ.}

\proseitemwithcloser{329}{อนาปตฺติ อมชฺชญฺจ โหติ มชฺชวณฺณํ มชฺชคนฺธํ มชฺชรสํ ตํ ปิวติ สูปสมฺปาเก มํสสมฺปาเก เตลสมฺปาเก อามลกผาณิเต, อมชฺชํ อริฏฺฐํ ปิวติ, อุมฺมตฺตกสฺส อาทิกมฺมิกสฺสาติ.}{สุราปานสิกฺขาปทํ นิฏฺฐิตํ ปฐมํ.}
\csromanmatikaanchor{mtk.59}
\csromantocmarkref{subsection}{2. องฺคุลิปโตทกสิกฺขาปท}{mtk.59}
\bookmark[dest={mtk.59},level=2]{2. องฺคุลิปโตทกสิกฺขาปท}
\csromanheader{2}{6. สุราปานวคฺค 2. องฺคุลิปโตทก\-สิกฺขาปท}
\proseitem{330}{\csromanfolio{147}เตน สมเยน พุทฺโธ ภควา สาวตฺถิยํ วิหรติ เชตวเน อนาถปิณฺฑิกสฺส อาราเม.\csromanspacer{} เตน โข ปน สมเยน ฉพฺพคฺคิยา ภิกฺขู สตฺตรสวคฺคิยํ ภิกฺขุํ องฺคุลิปโตทเกน หาเสสุํ.\csromanspacer{} โส ภิกฺขุ อุตฺตนฺโต อนสฺสาสโก กาลมกาสิ.\csromanspacer{} เย เต ภิกฺขู อปฺปิจฺฉา ฯเปฯ เต อุชฺฌายนฺติ ขิยฺยนฺติ วิปาเจนฺติ “กถํ หิ นาม ฉพฺพคฺคิยา ภิกฺขู ภิกฺขุํ องฺคุลิปโตทเกน หาเสสฺสนฺตี”ติ ฯเปฯ.\csromanspacer{} สจฺจํ กิร ตุเมฺห ภิกฺขเว ภิกฺขุํ องฺคุลิปโตทเกน หาเสถาติ.\csromanspacer{} สจฺจํ ภควาติ.\csromanspacer{} วิครหิ พุทฺโธ ภควา ฯเปฯ.\csromanspacer{} กถํ หิ นาม ตุเมฺห โมฆปุริสา ภิกฺขุํ องฺคุลิปโตทเกน หาเสสฺสถ.\csromanspacer{} เนตํ โมฆปุริสา อปฺปสนฺนานํ วา ปสาทาย ฯเปฯ.\csromanspacer{} เอวญฺจ ปน ภิกฺขเว อิมํ สิกฺขาปทํ อุทฺทิเสยฺยาถ\csromandash{}}

\proseitem{331}{\textbf{“องฺคุลิปโตทเก ปาจิตฺติยนฺ”}ติ.}

\proseitem{332}{\textbf{องฺคุลิปโตทโก} นาม\footnote{องฺคุลิปโตทโก นาม องฺคุลิยาปิ ตุทนฺติ (สฺยา)} อุปสมฺปนฺโน อุปสมฺปนฺนํ หสาธิปฺปาโย\footnote{หสฺสาธิปฺปาโย (สี, สฺยา)} กาเยน กายํ อามสติ, อาปตฺติ ปาจิตฺติยสฺส.\csromanspacer{} อุปสมฺปนฺเน อุปสมฺปนฺนสญฺญี องฺคุลิปโตทเกน หาเสติ, อาปตฺติ ปาจิตฺติยสฺส.\csromanspacer{} อุปสมฺปนฺเน เวมติโก องฺคุลิปโตทเกน หาเสติ, อาปตฺติ ปาจิตฺติยสฺส.\csromanspacer{} อุปสมฺปนฺเน อนุปสมฺปนฺนสญฺญี องฺคุลิปโตทเกน หาเสติ, อาปตฺติ}

\prose{กาเยน กายปฏิพทฺธํ อามสติ, อาปตฺติ ทุกฺกฏสฺส.\csromanspacer{} กายปฏิพทฺเธน กายํ อามสติ, อาปตฺติ ทุกฺกฏสฺส.\csromanspacer{} กายปฏิพทฺเธน กายปฏิพทฺธํ อามสติ, อาปตฺติ ทุกฺกฏสฺส.\csromanspacer{} นิสฺสคฺคิเยน กายํ อามสติ, อาปตฺติ ทุกฺกฏสฺส.\csromanspacer{} นิสฺสคฺคิเยน กายปฏิพทฺธํ อามสติ, อาปตฺติ ทุกฺกฏสฺส.\csromanspacer{} นิสฺสคฺคิเยน นิสฺสคฺคิยํ อามสติ, อาปตฺติ ทุกฺกฏสฺส.}

\proseitem{333}{อนุปสมฺปนฺนํ กาเยน กายํ อามสติ, อาปตฺติ ทุกฺกฏสฺส.\csromanspacer{} กาเยน กายปฏิพทฺธํ อามสติ, อาปตฺติ ทุกฺกฏสฺส.\csromanspacer{} กายปฏิพทฺเธน กายํ อามสติ, อาปตฺติ ทุกฺกฏสฺส.\csromanspacer{} กายปฏิพทฺเธน กายปฏิพทฺธํ อามสติ, อาปตฺติ ทุกฺกฏสฺส.\csromanspacer{} นิสคฺคิเยน กายํ อามสติ, อาปตฺติ \csromanfolio{148}ทุกฺกฏสฺส.\csromanspacer{} นิสฺสคฺคิเยน กายปฏิพทฺธํ อามสติ, อาปตฺติ ทุกฺกฏสฺส.\csromanspacer{} นิสฺสคฺคิเยน นิสฺสคฺคิยํ อามสติ, อาปตฺติ ทุกฺกฏสฺส.\csromanspacer{} อนุปสมฺปนฺเน อุปสมฺปนฺนสญฺญี, อาปตฺติ ทุกฺกฏสฺส.\csromanspacer{} อนุปสมฺปนฺเน เวมติโก, อาปตฺติ ทุกฺกฏสฺส.\csromanspacer{} อนุปสมฺปนฺเน อนุปสมฺปนฺนสญฺญี, อาปตฺติ ทุกฺกฏสฺส.}

\proseitemwithcloserruled{334}{อนาปตฺติ น หสาธิปฺปาโย สติ กรณีเย อามสติ, อุมฺมตฺตกสฺส อาทิกมฺมิกสฺสาติ.}{องฺคุลิปโตทก\-สิกฺขาปทํ นิฏฺฐิตํ ทุติยํ.}
\csromanmatikaanchor{mtk.60}
\csromantocmarkref{subsection}{3. หสธมฺมสิกฺขาปท}{mtk.60}
\bookmark[dest={mtk.60},level=2]{3. หสธมฺมสิกฺขาปท}
\csromanheader{2}{6. \textbf{สุราปานวคฺค 3. หสธมฺมสิกฺขาปท}}
\proseitem{335}{เตน สมเยน พุทฺโธ ภควา สาวตฺถิยํ วิหรติ เชตวเน อนาถปิณฺฑิกสฺส อาราเม.\csromanspacer{} เตน โข ปน สมเยน สตฺตรสวคฺคิยา ภิกฺขู อจิรวติยา นทิยา อุทเก กีฬนฺติ.\csromanspacer{} เตน โข ปน สมเยน ราชา ปเสนทิ โกสโล มลฺลิกาย เทวิยา สทฺธิํ อุปริปาสาทวรคโต โหติ.\csromanspacer{} อทฺทสา โข ราชา ปเสนทิ โกสโล สตฺตรสวคฺคิเย ภิกฺขู อจิรวติยา นทิยา อุทเก กีฬนฺเต, ทิสฺวาน มลฺลิกํ เทวิํ เอตทโวจ “เอเต เต มลฺลิเก อรหนฺโต อุทเก กีฬนฺตี”ติ.\csromanspacer{} นิสฺสํสยํ โข มหาราช ภควตา สิกฺขาปทํ อปญฺญตฺตํ.\csromanspacer{} เต วา ภิกฺขู อปฺปกตญฺญุโนติ.\csromanspacer{} อถ โข รญฺโญ ปเสนทิสฺส โกสลสฺส เอตทโหสิ “เกน นุ โข อหํ อุปาเยน ภควโต จ น อาโรเจยฺยํ, ภควา จ ชาเนยฺย อิเม ภิกฺขู อุทเก กีฬิตา”ติ.\csromanspacer{} อถ โข ราชา ปเสนทิ โกสโล สตฺตรสวคฺคิเย ภิกฺขู ปกฺโกสาเปตฺวา มหนฺตํ คุฬปิณฺฑํ อทาสิ “อิมํ ภนฺเต คุฬปิณฺฑํ ภควโต เทถา”ติ.\csromanspacer{} สตฺตรสวคฺคิยา ภิกฺขู ตํ คุฬปิณฺฑํ อาทาย เยน ภควา เตนุปสงฺกมิํสุ, อุปสงฺกมิตฺวา ภควนฺตํ เอตทโวจุํ “อิมํ ภนฺเต คุฬปิณฺฑํ ราชา ปเสนทิ โกสโล ภควโต เทตี”ติ.\csromanspacer{} กหํ ปน ตุเมฺห ภิกฺขเว ราชา อทฺทสาติ.\csromanspacer{} อจิรวติยา นทิยา ภควา อุทเก กีฬนฺเตติ.\csromanspacer{} วิครหิ พุทฺโธ ภควา ฯเปฯ.\csromanspacer{} กถํ หิ นาม ตุเมฺห โมฆปุริสา อุทเก กีฬิสฺสถ.\csromanspacer{} เนตํ โมฆปุริสา อปฺปสนฺนานํ วา ปสาทาย ฯเปฯ.\csromanspacer{} เอวญฺจ ปน ภิกฺขเว อิมํ สิกฺขาปทํ อุทฺทิเสยฺยาถ\csromandash{}}

\proseitem{336}{\csromanfolio{149}\textbf{“อุทเก หสธมฺเม}\footnote{หสฺสธมฺเม (สี, สฺยา)} \textbf{ปาจิตฺติยนฺ”}ติ.}

\proseitem{337}{\textbf{อุทเก หสธมฺโม} นาม อุปริโคปฺผเก อุทเก หสาธิปฺปาโย นิมุชฺชติ วา อุมฺมุชฺชติ วา ปลวติ วา, อาปตฺติ ปาจิตฺติยสฺส.}

\proseitem{338}{อุทเก หสธมฺเม หสธมฺมสญฺญี, อาปตฺติ ปาจิตฺติยสฺส.\csromanspacer{} อุทเก หสธมฺเม เวมติโก, อาปตฺติ ปาจิตฺติยสฺส.\csromanspacer{} อุทเก หสธมฺเม อหสธมฺมสญฺญี, อาปตฺติ ปาจิตฺติยสฺส.}

\prose{เหฏฺฐาโคปฺผเก อุทเก กีฬติ, อาปตฺติ ทุกฺกฏสฺส.\csromanspacer{} อุทเก นาวาย กีฬติ, อาปตฺติ ทุกฺกฏสฺส.\csromanspacer{} หตฺเถน วา ปาเทน วา กฏฺเฐน วา กฐลาย วา อุทกํ ปหรติ, อาปตฺติ ทุกฺกฏสฺส.\csromanspacer{} ภาชนคตํ อุทกํ วา กญฺชิกํ วา ขีรํ วา ตกฺกํ วา รชนํ วา ปสฺสาวํ วา จิกฺขลฺลํ วา กีฬติ, อาปตฺติ ทุกฺกฏสฺส.}

\prose{อุทเก อหสธมฺเม หสธมฺมสญฺญี, อาปตฺติ ทุกฺกฏสฺส.\csromanspacer{} อุทเก อหสธมฺเม เวมติโก, อาปตฺติ ทุกฺกฏสฺส.\csromanspacer{} อุทเก อหสธมฺเม อหสธมฺมสญฺญี, อนาปตฺติ.}

\proseitemwithcloserruled{339}{อนาปตฺติ น หสาธิปฺปาโย สติ กรณีเย อุทกํ โอตริตฺวา นิมุชฺชติ วา อุมฺมุชฺชติ วา ปลวติ วา, ปารํ คจฺฉนฺโต นิมุชฺชติ วา อุมฺมุชฺชติ วา ปลวติ วา, อาปทาสุ อุมฺมตฺตกสฺส อาทิกมฺมิกสฺสาติ.}{หสธมฺมสิกฺขาปทํ นิฏฺฐิตํ ตติยํ.}
\csromanmatikaanchor{mtk.61}
\csromantocmarkref{subsection}{4. อนาทริยสิกฺขาปท}{mtk.61}
\bookmark[dest={mtk.61},level=2]{4. อนาทริยสิกฺขาปท}
\csromanheader{2}{6. \textbf{สุราปานวคฺค 4. อนาทริยสิกฺขาปท}}
\proseitem{340}{เตน สมเยน พุทฺโธ ภควา โกสมฺพิยํ วิหรติ โฆสิตาราเม.\csromanspacer{} เตน โข ปน สมเยน อายสฺมา ฉนฺโน อนาจารํ อาจรติ.\csromanspacer{} ภิกฺขู เอวมาหํสุ “มาวุโส ฉนฺน เอวรูปํ อกาสิ, เนตํ กปฺปตี”ติ.\csromanspacer{} โส อนาทริยํ ปฏิจฺจ กโรติ เยว.\csromanspacer{} เย เต ภิกฺขู อปฺปิจฺฉา ฯเปฯ เต อุชฺฌายนฺติ ขิยฺยนฺติ วิปาเจนฺติ “กถํ หิ นาม อายสฺมา ฉนฺโน อนาทริยํ กริสฺสตี”ติ ฯเปฯ.\csromanspacer{} สจฺจํ กิร ตฺวํ ฉนฺน อนาทริยํ \csromanfolio{150}\textbf{กโรสีติ.}\csromanspacer{}\textbf{ สจฺจํ ภควาติ.}\csromanspacer{}\textbf{ วิครหิ พุทฺโธ ภควา ฯเปฯ.}\csromanspacer{}\textbf{ กถํ หิ} \textbf{นาม ตฺวํ โมฆปุริส อนาทริยํ กริสฺสสิ.}\csromanspacer{}\textbf{ เนตํ โมฆปุริส} อปฺปสนฺนานํ วา ปสาทาย ฯเปฯ.\csromanspacer{} เอวญฺจ ปน ภิกฺขเว อิมํ สิกฺขาปทํ อุทฺทิเสยฺยาถ\csromandash{}}

\proseitem{341}{\textbf{“อนาทริเย ปาจิตฺติยนฺ”}ติ.}

\proseitem{342}{\textbf{อนาทริยํ} นาม เทฺว อนาทริยานิ ปุคฺคลานาทริยญฺจ ธมฺมานาทริยญฺจ.}

\prose{\textbf{ปุคฺคลานาทริยํ} นาม อุปสมฺปนฺเนน ปญฺญตฺเตน วุจฺจมาโน “อยํ อุกฺขิตฺตโก วา วมฺภิโต วา ครหิโต วา, อิมสฺส วจนํ อกตํ ภวิสฺสตี”ติ อนาทริยํ กโรติ, อาปตฺติ ปาจิตฺติยสฺส.}

\prose{\textbf{ธมฺมานาทริยํ} นาม อุปสมฺปนฺเนน ปญฺญตฺเตน วุจฺจมาโน “กถา’ยํ นสฺเสยฺย วา วินสฺเสยฺย วา อนฺตรธาเยยฺยวา”, ตํ วา น สิกฺขิตุกาโม อนาทริยํ กโรติ, อาปตฺติ ปาจิตฺติยสฺส.}

\proseitem{343}{อุปสมฺปนฺเน อุปสมฺปนฺนสญฺญี อนาทริยํ กโรติ, อาปตฺติ ปาจิตฺติยสฺส.\csromanspacer{} อุปสมฺปนฺเน เวมติโก อนาทริยํ กโรติ, อาปตฺติ ปาจิตฺติยสฺส.\csromanspacer{} อุปสมฺปนฺเน อนุปสมฺปนฺนสญฺญี อนาทริยํ กโรติ, อาปตฺติ ปาจิตฺติยสฺส.}

\prose{อปญฺญตฺเตน วุจฺจมาโน “อิทํ น สลฺเลขาย น ธุตตฺถาย น ปาสาทิกตาย น อปจยาย น วีริยารมฺภาย สํวตฺตตี”ติ อนาทริยํ กโรติ, อาปตฺติ ทุกฺกฏสฺส.\csromanspacer{} อนุปสมฺปนฺเนน ปญฺญตฺเตน วา อปญฺญตฺเตน วา วุจฺจมาโน “อิทํ น สลฺเลขาย น ธุตตฺถาย น ปาสาทิกตาย น อปจยาย น วีริยารมฺภาย สํวตฺตตี”ติ อนาทริยํ กโรติ, อาปตฺติ ทุกฺกฏสฺส.}

\prose{อนุปสมฺปนฺเน อุปสมฺปนฺนสญฺญี, อาปตฺติ ทุกฺกฏสฺส.\csromanspacer{} อนุปสมฺปนฺเน เวมติโก, อาปตฺติ ทุกฺกฏสฺส.\csromanspacer{} อนุปสมฺปนฺเน อนุปสมฺปนฺนสญฺญี, อาปตฺติ ทุกฺกฏสฺส.}

\proseitemwithcloser{344}{อนาปตฺติ “เอวํ อมฺหากํ อาจริยานํ อุคฺคโห ปริปุจฺฉา”ติ ภณติ, อุมฺมตฺตกสฺส อาทิกมฺมิกสฺสาติ.}{อนาทริยสิกฺขาปทํ นิฏฺฐิตํ จตุตฺถํ.}
\csromanmatikaanchor{mtk.62}
\csromantocmarkref{subsection}{5. ภิํสาปนสิกฺขาปท}{mtk.62}
\bookmark[dest={mtk.62},level=2]{5. ภิํสาปนสิกฺขาปท}
\csromanheader{2}{6. \textbf{สุราปานวคฺค 5. ภิํสาปนสิกฺขาปท}}
\proseitem{345}{\csromanfolio{151}เตน สมเยน พุทฺโธ ภควา สาวตฺถิยํ วิหรติ เชตวเน อนาถปิณฺฑิกสฺส อาราเม.\csromanspacer{} เตน โข ปน สมเยน ฉพฺพคฺคิยา ภิกฺขู สตฺตรสวคฺคิเย ภิกฺขู ภิํสาเปนฺติ, เต ภิํสาปียมานา โรทนฺติ.\csromanspacer{} ภิกฺขู เอวมาหํสุ “กิสฺส ตุเมฺห อาวุโส โรทถา”ติ.\csromanspacer{} อิเม อาวุโส ฉพฺพคฺคิยา ภิกฺขู อเมฺห ภิํสาเปนฺตีติ.\csromanspacer{} เย เต ภิกฺขู อปฺปิจฺฉา ฯเปฯ เต อุชฺฌายนฺติ ขิยฺยนฺติ วิปาเจนฺติ “กถํ หิ นาม ฉพฺพคฺคิยา ภิกฺขู ภิกฺขุํ ภิํสาเปสฺสนฺตี”ติ ฯเปฯ.\csromanspacer{} สจฺจํ กิร ตุเมฺห ภิกฺขเว ภิกฺขุํ ภิํสาเปถาติ.\csromanspacer{} สจฺจํ ภควาติ.\csromanspacer{} วิครหิ พุทฺโธ ภควา ฯเปฯ.\csromanspacer{} กถํ หิ นาม ตุเมฺห โมฆปุริสา ภิกฺขุํ ภิํสาเปสฺสถ.\csromanspacer{} เนตํ โมฆปุริสา อปฺปสนฺนานํ วา ปสาทาย ฯเปฯ.\csromanspacer{} เอวญฺจ ปน ภิกฺขเว อิมํ สิกฺขาปทํ อุทฺทิเสยฺยาถ\csromandash{}}

\proseitem{346}{\textbf{“โย ปน ภิกฺขุ ภิกฺขุํ ภิํสาเปยฺย, ปาจิตฺติยนฺ”}ติ.}

\proseitem{347}{\textbf{โย ปนา}ติ โย ยาทิโส ฯเปฯ.\csromanspacer{} \textbf{ภิกฺขู}ติ ฯเปฯ อยํ อิมสฺมิํ อตฺเถ อธิปฺเปโต ภิกฺขูติ.}

\prose{\textbf{ภิกฺขุนฺ}ติ อญฺญํ ภิกฺขุํ.}

\prose{\textbf{ภิํสาเปยฺยา}ติ อุปสมฺปนฺโน อุปสมฺปนฺนํ ภิํสาเปตุกาโม รูปํ วา สทฺทํ วา คนฺธํ วา รสํ วา โผฏฺฐพฺพํ วา อุปสํหรติ, ภาเยยฺย วา โส น วา ภาเยยฺย, อาปตฺติ ปาจิตฺติยสฺส.\csromanspacer{} โจรกนฺตารํ วา วาฬกนฺตารํ วา ปิสาจกนฺตารํ วา อาจิกฺขติ, ภาเยยฺย วา โส น วา ภาเยยฺย, อาปตฺติ ปาจิตฺติยสฺส.}

\proseitem{348}{อุปสมฺปนฺเน อุปสมฺปนฺนสญฺญี ภิํสาเปติ, อาปตฺติ ปาจิตฺติยสฺส.\csromanspacer{} อุปสมฺปนฺเน เวมติโก ภิํสาเปติ, อาปตฺติ ปาจิตฺติยสฺส.\csromanspacer{} อุปสมฺปนฺเน อนุปสมฺปนฺนสญฺญี ภิํสาเปติ, อาปตฺติ ปาจิตฺติยสฺส.}

\prose{อนุปสมฺปนฺนํ ภิํสาเปตุกาโม รูปํ วา สทฺทํ วา คนฺธํ วา รสํ วา โผฏฺฐพฺพํ วา อุปสํหรติ, ภาเยยฺย วา โส น วา ภาเยยฺย, อาปตฺติ ทุกฺกฏสฺส.\csromanspacer{} โจรกนฺตารํ วา วาฬกนฺตารํ วา ปิสาจกนฺตารํ วา อาจิกฺขติ, ภาเยยฺย วา โส น วา ภาเยยฺย, อาปตฺติ \csromanfolio{152}ทุกฺกฏสฺส.\csromanspacer{} อนุปสมฺปนฺเน อุปสมฺปนฺนสญฺญี, อาปตฺติ ทุกฺกฏสฺส.\csromanspacer{} อนุปสมฺปนฺเน เวมติโก, อาปตฺติ ทุกฺกฏสฺส.\csromanspacer{} อนุปสมฺปนฺเน อนุปสมฺปนฺนสญฺญี, อาปตฺติ ทุกฺกฏสฺส.}

\proseitemwithcloserruled{349}{อนาปตฺติ น ภิํสาเปตุกาโม รูปํ วา สทฺทํ วา คนฺธํ วา รสํ วา โผฏฺฐพฺพํ วา อุปสํหรติ, โจรกนฺตารํ วา วาฬกนฺตารํ วา ปิสาจกนฺตารํ วา อาจิกฺขติ, อุมฺมตฺตกสฺส อาทิกมฺมิกสฺสาติ.}{ภิํสาปนสิกฺขาปทํ นิฏฺฐิตํ ปญฺจมํ.}
\csromanmatikaanchor{mtk.63}
\csromantocmarkref{subsection}{6. โชติกสิกฺขาปท}{mtk.63}
\bookmark[dest={mtk.63},level=2]{6. โชติกสิกฺขาปท}
\csromanheader{2}{6. \textbf{สุราปานวคฺค 6. โชติกสิกฺขาปท}}
\proseitem{350}{เตน สมเยน พุทฺโธ ภควา ภคฺเคสุ วิหรติ สุสุมารคิเร\footnote{สุํสุมารคิเร (สี, สฺยา), สํสุมารคิเร (ก)} เภสกฬาวเน มิคทาเย.\csromanspacer{} เตน โข ปน สมเยน ภิกฺขู เหมนฺติเก กาเล อญฺญตรํ มหนฺตํ สุสิรกฏฺฐํ โชติํ สมาทหิตฺวา วิสิพฺเพสุํ.\csromanspacer{} ตสฺมิํ จ สุสิเร กณฺหสปฺโป อคฺคินา สนฺตตฺโต นิกฺขมิตฺวา ภิกฺขู ปริปาเตสิ.\csromanspacer{} ภิกฺขู ตหํ ตหํ อุปธาวิํสุ.\csromanspacer{} เย เต ภิกฺขู อปฺปิจฺฉา ฯเปฯ เต อุชฺฌายนฺติ ขิยฺยนฺติ วิปาเจนฺติ “กถํ หิ นาม ภิกฺขู โชติํ สมาทหิตฺวา วิสิพฺเพสฺสนฺตี”ติ ฯเปฯ.\csromanspacer{} สจฺจํ กิร ภิกฺขเว ภิกฺขู โชติํ สมาทหิตฺวา วิสิพฺเพนฺตีติ.\csromanspacer{} สจฺจํ ภควาติ.\csromanspacer{} วิครหิ พุทฺโธ ภควา ฯเปฯ.\csromanspacer{} กถํ หิ นาม เต ภิกฺขเว โมฆปุริสา โชติํ สมาทหิตฺวา วิสิพฺเพสฺสนฺติ.\csromanspacer{} เนตํ ภิกฺขเว อปฺปสนฺนานํ วา ปสาทาย ฯเปฯ.\csromanspacer{} เอวญฺจ ปน ภิกฺขเว อิมํ สิกฺขาปทํ อุทฺทิเสยฺยาถ\csromandash{}}

\prose{\textbf{“โย ปน ภิกฺขุ วิสิพฺพนาเปกฺโข โชติํ สมาทเหยฺย วา สมาทหาเปยฺย วา, ปาจิตฺติยนฺ”}ติ.}

\prose{เอวญฺจิทํ ภควตา ภิกฺขูนํ สิกฺขาปทํ ปญฺญตฺตํ โหติ.}

\proseitem{351}{เตน โข ปน สมเยน ภิกฺขู คิลานา โหนฺติ.\csromanspacer{} คิลานปุจฺฉกา ภิกฺขู คิลาเน ภิกฺขู เอตทโวจุํ “กจฺจาวุโส ขมนียํ, กจฺจิ ยาปนียนฺ”ติ.\csromanspacer{} ปุพฺเพ มยํ อาวุโส โชติํ สมาทหิตฺวา วิสิพฺเพม, เตน โน ผาสุ โหติ, อิทานิ ปน “ภควตา ปฏิกฺขิตฺตนฺ”ติ กุกฺกุจฺจายนฺตา \csromanfolio{153}น วิสิพฺเพม, เตน โน น ผาสุ โหตีติ.\csromanspacer{} ภควโต เอตมตฺถํ อาโรเจสุํ ฯเปฯ.\csromanspacer{} อนุชานามิ ภิกฺขเว คิลาเนน ภิกฺขุนา โชติํ สมาทหิตฺวา วา สมาทหาเปตฺวา วา วิสิพฺเพตุํ.\csromanspacer{} เอวญฺจ ปน ภิกฺขเว อิมํ สิกฺขาปทํ อุทฺทิเสยฺยาถ\csromandash{}}

\prose{\textbf{“โย ปน ภิกฺขุ อคิลาโน วิสิพฺพนาเปกฺโข โชติํ สมาทเหยฺย วา สมาทหาเปยฺย วา, ปาจิตฺติยนฺ”}ติ.}

\prose{เอวญฺจิทํ ภควตา ภิกฺขูนํ สิกฺขาปทํ ปญฺญตฺตํ โหติ.}

\proseitem{352}{เตน โข ปน สมเยน ภิกฺขู ปทีเปปิ โชติเกปิ ชนฺตาฆเรปิ กุกฺกุจฺจายนฺติ.\csromanspacer{} ภควโต เอตมตฺถํ อาโรเจสุํ ฯเปฯ.\csromanspacer{} อนุชานามิ ภิกฺขเว ตถารูปปฺปจฺจยา โชติํ สมาทหิตุํ สมาทหาเปตุํ.\csromanspacer{} เอวญฺจ ปน ภิกฺขเว อิมํ สิกฺขาปทํ อุทฺทิเสยฺยาถ\csromandash{}}

\proseitem{353}{“\textbf{โย ปน ภิกฺขุ อคิลาโน วิสิพฺพนาเปกฺโข โชติํ สมาทเหยฺย วา} \textbf{สมาทหาเปยฺยฺ}อ วา อญฺญตฺร ตถารูปปฺปจฺจยา ปาจิตฺติยนฺ”ติ.}

\proseitem{354}{\textbf{โย ปนา}ติ โย ยาทิโส ฯเปฯ.\csromanspacer{} \textbf{ภิกฺขู}ติ ฯเปฯ อยํ อิมสฺมิํ อตฺเถ อธิปฺเปโต ภิกฺขูติ.}

\prose{\textbf{อคิลาโน} นาม ยสฺส วินา อคฺคินา ผาสุ โหติ.}

\prose{\textbf{คิลาโน} นาม ยสฺส วินา อคฺคินา น ผาสุ โหติ.}

\prose{\textbf{วิสิพฺพนาเปกฺโข}ติ ตปฺปิตุกาโม.\csromanspacer{} \textbf{โชติ} นาม อคฺคิ วุจฺจติ.}

\prose{\textbf{สมาทเหยฺยา}ติ สยํ สมาทหติ, อาปตฺติ ปาจิตฺติยสฺส.}

\prose{\textbf{สมาทหาเปยฺยา}ติ อญฺญํ อาณาเปติ, อาปตฺติ ปาจิตฺติยสฺส.\csromanspacer{} สกิํ อาณาตฺโต พหุกมฺปิ สมาทหติ, อาปตฺติ ปาจิตฺติยสฺส.}

\prose{\textbf{อญฺญตฺร ตถารูปปฺปจฺจยา}ติ ฐเปตฺวา ตถารูปปฺปจฺจยํ.}

\proseitem{355}{อคิลาโน อคิลานสญฺญี วิสิพฺพนาเปกฺโข โชติํ สมาทหติ วา สมาทหาเปติ วา อญฺญตฺร ตถารูปปฺปจฺจยา, อาปตฺติ ปาจิตฺติยสฺส.\csromanspacer{} อคิลาโน เวมติโก วิสิพฺพนาเปกฺโข โชติํ \csromanfolio{154}สมาทหติ วา สมาทหาเปติ วา อญฺญตฺร ตถารูปปฺปจฺจยา, อาปตฺติ ปาจิตฺติยสฺส.\csromanspacer{} อคิลาโน คิลานสญฺญี วิสิพฺพนาเปกฺโข โชติํ สมาทหติ วา สมาทหาเปติ วา อญฺญตฺร ตถารูปปฺปจฺจยา, อาปตฺติ ปาจิตฺติยสฺส.}

\prose{ปฏิลาตํ อุกฺขิปติ, อาปตฺติ ทุกฺกฏสฺส.\csromanspacer{} คิลาโน อคิลานสญฺญี, อาปตฺติ ทุกฺกฏสฺส.\csromanspacer{} คิลาโน เวมติโก, อาปตฺติ ทุกฺกฏสฺส.\csromanspacer{} คิลาโน คิลานสญฺญี, อนาปตฺติ.}

\proseitemwithcloserruled{356}{อนาปตฺติ คิลานสฺส อญฺเญน กตํ วิสิพฺเพติ, วีตจฺจิตงฺคารํ วิสิพฺเพติ, ปทีเป โชติเก ชนฺตาฆเร ตถารูปปฺปจฺจยา อาปทาสุ อุมฺมตฺตกสฺส อาทิกมฺมิกสฺสาติ.}{โชติกสิกฺขาปทํ นิฏฺฐิตํ ฉฏฺฐํ.}
\csromanmatikaanchor{mtk.64}
\csromantocmarkref{subsection}{7. นหานสิกฺขาปท}{mtk.64}
\bookmark[dest={mtk.64},level=2]{7. นหานสิกฺขาปท}
\csromanheader{2}{6. \textbf{สุราปานวคฺค 7. นหานสิกฺขาปท}}
\proseitem{357}{เตน สมเยน พุทฺโธ ภควา ราชคเห วิหรติ เวฬุวเน กลนฺทกนิวาเป.\csromanspacer{} เตน โข ปน สมเยน ภิกฺขู ตโปเท นหายนฺติ.\csromanspacer{} เตน โข ปน สมเยน\footnote{อถ โข (สี, สฺยา)} ราชา มาคโธ เสนิโย พิมฺพิสาโร “สีสํ นหายิสฺสามี”ติ ตโปทํ คนฺตฺวา “ยาวายฺยา นหายนฺตี”ติ เอกมนฺตํ ปฏิมาเนสิ.\csromanspacer{} ภิกฺขู ยาว สมนฺธการา นหายิํสุ.\csromanspacer{} อถ โข ราชา มาคโธ เสนิโย พิมฺพิสาโร วิกาเล สีสํ นหายิตฺวา นครทฺวาเร ถกิเต พหินคเร วสิตฺวา กาลสฺเสว อสมฺภินฺเนน วิเลปเนน เยน ภควา เตนุปสงฺกมิ, อุปสงฺกมิตฺวา ภควนฺตํ อภิวาเทตฺวา เอกมนฺตํ นิสีทิ.\csromanspacer{} เอกมนฺตํ นิสินฺนํ โข ราชานํ มาคธํ เสนิยํ พิมฺพิสารํ ภควา เอตทโวจ “กิสฺส ตฺวํ มหาราช กาลสฺเสว อาคโต อสมฺภินฺเนน วิเลปเนนา”ติ.\csromanspacer{} อถ โข ราชา มาคโธ เสนิโย พิมฺพิสาโร ภควโต เอตมตฺถํ อาโรเจสิ.\csromanspacer{} อถ โข ภควา ราชานํ มาคธํ เสนิยํ พิมฺพิสารํ ธมฺมิยา กถาย สนฺทสฺเสสิ สมาทเปสิ สมุตฺเตเชสิ สมฺปหํเสสิ.\csromanspacer{} อถ โข ราชา มาคโธ เสนิโย พิมฺพิสาโร ภควตา ธมฺมิยา กถาย สนฺทสฺสิโต สมาทปิโต สมุตฺเตชิโต สมฺปหํสิโต \csromanfolio{155}\textbf{อุฏฺฐายาสนา ภควนฺตํ อภิวาเทตฺวา ปทกฺขิณํ กตฺวา ปกฺกามิ.}\csromanspacer{}\textbf{ อถ} \textbf{โข ภควา เอตสฺมิํ นิทาเน เอตสฺมิํ ปกรเณ ภิกฺขุสํฆํ} \textbf{สนฺนิปาตาเปตฺวา ภิกฺขู ปฏิปุจฺฉิ “สจฺจํ กิร ภิกฺขเว ภิกฺขู} ราชานมฺปิ ปสฺสิตฺวา น มตฺตํ ชานิตฺวา นหายนฺตี”ติ.\csromanspacer{} สจฺจํ ภควาติ.\csromanspacer{} วิครหิ พุทฺโธ ภควา ฯเปฯ.\csromanspacer{} กถํ หิ นาม เต ภิกฺขเว โมฆปุริสา ราชานมฺปิ ปสฺสิตฺวา น มตฺตํ ชานิตฺวา นหายิสฺสนฺติ.\csromanspacer{} เนตํ ภิกฺขเว อปฺปสนฺนานํ วา ปสาทาย ฯเปฯ.\csromanspacer{} เอวญฺจ ปน ภิกฺขเว อิมํ สิกฺขาปทํ อุทฺทิเสยฺยาถ\csromandash{}}

\prose{\textbf{“โย ปน ภิกฺขุ โอเรนทฺธมาสํ นหาเยยฺย, ปาจิตฺติยนฺ”}ติ.}

\prose{เอวญฺจิทํ ภควตา ภิกฺขูนํ สิกฺขาปทํ ปญฺญตฺตํ โหติ.}

\proseitem{358}{เตน โข ปน สมเยน ภิกฺขู อุณฺหสมเย ปริฬาหสมเย กุกฺกุจฺจายนฺตา น นหายนฺติ.\csromanspacer{} เสทคเตน คตฺเตน สยนฺติ, จีวรมฺปิ เสนาสนมฺปิ ทุสฺสติ.\csromanspacer{} ภควโต เอตมตฺถํ อาโรเจสุํ ฯเปฯ.\csromanspacer{} อนุชานามิ ภิกฺขเว อุณฺหสมเย ปริฬาหสมเย โอเรนทฺธมาสํ นหายิตุํ.\csromanspacer{} เอวญฺจ ปน ภิกฺขเว อิมํ สิกฺขาปทํ อุทฺทิเสยฺยาถ\csromandash{}}

\prose{“\textbf{โย ปน ภิกฺขุ โอเรนทฺธมาสํ นหาเยยฺยฺ}อ อญฺญตฺร สมยา, ปาจิตฺติยํ.\csromanspacer{} ตตฺถายํ สมโย, ทิยฑฺโฒ มาโส เสโส คิมฺหานนฺติ\footnote{คิมฺหานํ (อิติปิ)} วสฺสานสฺส ปฐโม มาโส อิจฺเจเต อฑฺฒเตยฺยมาสา อุณฺหสมโย ปริฬาหสมโย, อยํ ตตฺถ สมโย”ติ.}

\prose{เอวญฺจิทํ ภควตา ภิกฺขูนํ สิกฺขาปทํ ปญฺญตฺตํ โหติ.}

\proseitem{359}{เตน โข ปน สมเยน ภิกฺขู คิลานา โหนฺติ.\csromanspacer{} คิลานปุจฺฉกา ภิกฺขู คิลาเน ภิกฺขู เอตทโวจุํ “กจฺจาวุโส ขมนียํ, กจฺจิ ยาปนียนฺ”ติ.\csromanspacer{} ปุพฺเพ มยํ อาวุโส โอเรนทฺธมาสํ นหายาม, เตน โน ผาสุ โหติ, อิทานิ ปน “ภควตา ปฏิกฺขิตฺตนฺ”ติ กุกฺกุจฺจายนฺตา น นหายาม, เตน โน น ผาสุ โหตีติ.\csromanspacer{} ภควโต เอตมตฺถํ อาโรเจสุํ ฯเปฯ.\csromanspacer{} อนุชานามิ ภิกฺขเว คิลาเนน ภิกฺขุนา โอเรนทฺธมาสํ นหายิตุํ.\csromanspacer{} เอวญฺจ ปน ภิกฺขเว อิมํ สิกฺขาปทํ อุทฺทิเสยฺยาถ\csromandash{}}

\prose{“\textbf{โย ปน ภิกฺขุ โอเรนทฺธมาสํ นหาเยยฺยฺ}อ อญฺญตฺร สมยา, ปาจิตฺติยํ.\csromanspacer{} ตตฺถายํ สมโย, ทิยฑฺโฒ มาโส เสโส คิมฺหานนฺติ1 \csromanfolio{156}\textbf{วสฺสานสฺส ปฐโม มาโส อิจฺเจเต อฑฺฒเตยฺยมาสา อุณฺหสมโย ปริฬาหสมโย คิลานสมโย, อยํ ตตฺถ สมโย”}ติ.}

\prose{เอวญฺจิทํ ภควตา ภิกฺขูนํ สิกฺขาปทํ ปญฺญตฺตํ โหติ.}

\proseitem{360}{เตน โข ปน สมเยน ภิกฺขู นวกมฺมํ กตฺวา กุกฺกุจฺจายนฺตา น นหายนฺติ.\csromanspacer{} เต เสทคเตน คตฺเตน สยนฺติ, จีวรมฺปิ เสนาสนมฺปิ ทุสฺสติ.\csromanspacer{} ภควโต เอตมตฺถํ อาโรเจสุํ ฯเปฯ.\csromanspacer{} อนุชานามิ ภิกฺขเว กมฺมสมเย โอเรนทฺธมาสํ นหายิตุํ.\csromanspacer{} เอวญฺจ ปน ภิกฺขเว อิมํ สิกฺขาปทํ อุทฺทิเสยฺยาถ\csromandash{}}

\prose{“โย ปน ภิกฺขุ โอเรนทฺธมาสํ นหาเยยฺย อญฺญตฺร สมยา, ปาจิตฺติยํ.\csromanspacer{}\textbf{ ตตฺถายํ สมโย, ทิยฑฺโฒ มาโส เสโส คิมฺหานนฺติ วสฺสานสฺส ปฐโม มาโส อิจฺเจเต อฑฺฒเตยฺยมาสา อุณฺหสมโย ปริฬาหสมโย คิลานสมโย กมฺมสมโย, อยํ ตตฺถ สมโย”}ติ.}

\prose{เอวญฺจิทํ ภควตา ภิกฺขูนํ สิกฺขาปทํ ปญฺญตฺตํ โหติ.}

\proseitem{361}{เตน โข ปน สมเยน ภิกฺขู อทฺธานํ คนฺตฺวา กุกฺกุจฺจายนฺตา น นหายนฺติ.\csromanspacer{} เต เสทคเตน คตฺเตน สยนฺติ, จีวรมฺปิ เสนาสนมฺปิ ทุสฺสติ.\csromanspacer{} ภควโต เอตมตฺถํ อาโรเจสุํ ฯเปฯ.\csromanspacer{} อนุชานามิ ภิกฺขเว อทฺธานคมนสมเย โอเรนทฺธมาสํ นหายิตุํ.\csromanspacer{} เอวญฺจ ปน ภิกฺขเว อิมํ สิกฺขาปทํ อุทฺทิเสยฺยาถ\csromandash{}}

\prose{“\textbf{โย ปน ภิกฺขุ โอเรนทฺธมาสํ นหาเยยฺย อญฺญตฺร สมยา,} \textbf{ปาจิตฺติยํ.}\csromanspacer{}\textbf{ ตตฺถายํ สมโย, ทิยฑฺโฒ มาโส เสโส คิมฺหานนฺติ วสฺสานสฺส} \textbf{ปฐโม มาโส อิจฺเจเต อฑฺฒเตยฺยมาสา อุณฺหสมโย ปริฬาหสมโย} \textbf{คิลานสมโย กมฺมสฺ}อมโย อทฺธานคมนสมโย, อยํ ตตฺถ สมโย”ติ.}

\prose{เอวญฺจิทํ ภควตา ภิกฺขูนํ สิกฺขาปทํ ปญฺญตฺตํ โหติ.}

\proseitem{362}{เตน โข ปน สมเยน สมฺพหุลา ภิกฺขู อชฺโฌกาเส จีวรกมฺมํ กโรนฺตา สรเชน วาเตน โอกิณฺณา โหนฺติ, เทโว จ โถกํ โถกํ ผุสายติ.\csromanspacer{} ภิกฺขู กุกฺกุจฺจายนฺตา น นหายนฺติ, \csromanfolio{157}กิลินฺเนน คตฺเตน สยนฺติ, จีวรมฺปิ เสนาสนมฺปิ ทุสฺสติ.\csromanspacer{} ภควโต เอตมตฺถํ อาโรเจสุํ ฯเปฯ.\csromanspacer{} อนุชานามิ ภิกฺขเว วาตวุฏฺฐิสมเย โอเรนทฺธมาสํ นหายิตุํ.\csromanspacer{} เอวญฺจ ปน ภิกฺขเว อิมํ สิกฺขาปทํ อุทฺทิเสยฺยาถ\csromandash{}}

\proseitem{363}{“โย ปน ภิกฺขุ โอเรนทฺธมาสํ นหาเยยฺย อญฺญตฺรสมยา, ปาจิตฺติยํ.\csromanspacer{}\textbf{ ตตฺถายํ สมโย, ทิยฑฺโฒ มาโส เสโส คิมฺหานนฺติ วสฺสานสฺส ปฐโม มาโส อิจฺเจเต อฑฺฒเตยฺยมาสา อุณฺหสมโย ปริฬาหสมโย คิลานสมโย กมฺมสมโย อทฺธานคมนสมโย วาตวุฏฺฐิสมโย, อยํ ตตฺถ สมโย”}ติ.}

\proseitem{364}{\textbf{โย ปนา}ติ โย ยาทิโส ฯเปฯ.\csromanspacer{} \textbf{ภิกฺขู}ติ ฯเปฯ อยํ อิมสฺมิํ อตฺเถ อธิปฺเปโต ภิกฺขูติ.}

\prose{\textbf{โอเรนทฺธมาสนฺ}ติ อูนกทฺธมาสํ.}

\prose{\textbf{นหาเยยฺยา}ติ จุณฺเณน วา มตฺติกาย วา นหายติ, ปโยเค ปโยเค ทุกฺกฏํ.\csromanspacer{} นหานปริโยสาเน อาปตฺติ ปาจิตฺติยสฺส.}

\prose{\textbf{อญฺญตฺร สมยา}ติ ฐเปตฺวา สมยํ.}

\prose{\textbf{อุณฺหสมโย} นาม ทิยฑฺโฒ มาโส เสโส คิมฺหานํ.\csromanspacer{} \textbf{ปริฬาหสมโย} นาม วสฺสานสฺส ปฐโม มาโส. “อิจฺเจเต อฑฺฒเตยฺยมาสา อุณฺหสมโย ปริฬาหสมโย”ติ นหายิตพฺพํ.}

\prose{\textbf{คิลานสมโย} นาม ยสฺส วินา นหาเนน น ผาสุ โหติ, “คิลานสมโย”ติ นหายิตพฺพํ.}

\prose{\textbf{กมฺมสมโย} นาม อนฺตมโส ปริเวณมฺปิ สมฺมฏฺฐํ โหติ, “กมฺมสมฺมโย”ติ นหายิตพฺพํ.}

\prose{\textbf{อทฺธานคมนสมโย} นาม “อทฺธโยชนํ คจฺฉิสฺสามี”ติ นหายิตพฺพํ, คจฺฉนฺเตน นหายิตพฺพํ, คเตน นหายิตพฺพํ.}

\prose{\textbf{วาตวุฏฺฐิสมโย} นาม ภิกฺขู สรเชน วาเตน โอกิณฺณา โหนฺติ, เทฺว วา ตีณิ วา อุทกผุสิตานิ กาเย ปติตานิ โหนฺติ, “วาตวุฏฺฐิสมโย”ติ นหายิตพฺพํ.}

\proseitem{365}{อูนกทฺธมาเส อูนกสญฺญี อญฺญตฺร สมยา นหายติ, อาปตฺติ ปาจิตฺติยสฺส.\csromanspacer{} อูนกทฺธมาเส เวมติโก อญฺญตฺร สมยา นหายติ, \csromanfolio{158}\textbf{อาปตฺติ ปาจิตฺติยสฺส.}\csromanspacer{}\textbf{ อูนกทฺธมาเส อติเรกสญฺญี อญฺญตฺร สมยา นหายติ,} \textbf{อาปตฺติ ปาจิตฺติยสฺส.}}

\prose{อติเรกทฺธมาเส อูนกสญฺญี, อาปตฺติ ทุกฺกฏสฺส.\csromanspacer{} อติเรกทฺธมาเส เวมติโก, อาปตฺติ ทุกฺกฏสฺส.\csromanspacer{} อติเรกทฺธมาเส อติเรกสญฺญี อนาปตฺติ.}

\proseitemwithcloserruled{366}{อนาปตฺติ สมเย, อทฺธมาสํ นหายติ, อติเรกทฺธมาสํ นหายติ, ปารํ คจฺฉนฺโต นหายติ, สพฺพปจฺจนฺติเมสุ ชนปเทสุ, อาปทาสุ อุมฺมตฺตกสฺส อาทิกมฺมิกสฺสาติ.}{นหานสิกฺขาปทํ นิฏฺฐิตํ สตฺตมํ.}
\csromanmatikaanchor{mtk.65}
\csromantocmarkref{subsection}{8. ทุพฺพณฺณกรณสิกฺขาปท}{mtk.65}
\bookmark[dest={mtk.65},level=2]{8. ทุพฺพณฺณกรณสิกฺขาปท}
\csromanheader{2}{6. \textbf{สุราปานวคฺค 8. ทุพฺพณฺณกรณสิกฺขาปท}}
\proseitem{367}{เตน สมเยน พุทฺโธ ภควา สาวตฺถิยํ วิหรติ เชตวเน อนาถปิณฺฑิกสฺส อาราเม.\csromanspacer{} เตน โข ปน สมเยน สมฺพหุลา ภิกฺขู จ ปริพฺพาชกา จ สาเกตา สาวตฺถิํ อทฺธานมคฺคปฺปฏิปนฺนา โหนฺติ.\csromanspacer{} อนฺตรามคฺเค โจรา นิกฺขมิตฺวา เต อจฺฉินฺทิํสุ.\csromanspacer{} สาวตฺถิยา ราชภฏา นิกฺขมิตฺวา เต โจเร สภณฺเฑ คเหตฺวา ภิกฺขูนํ สนฺติเก ทูตํ ปาเหสุํ “อาคจฺฉนฺตุ ภทนฺตา สกํ สกํ จีวรํ สญฺชานิตฺวา คณฺหนฺตู”ติ.\csromanspacer{} ภิกฺขู น สญฺชานนฺติ.\csromanspacer{} เต อุชฺฌายนฺติ ขิยฺยนฺติ วิปาเจนฺติ “กถํ หิ นาม ภทนฺตา อตฺตโน อตฺตโน จีวรํ น สญฺชานิสฺสนฺตี”ติ.\csromanspacer{} อสฺโสสุํ โข ภิกฺขู เตสํ มนุสฺสานํ อุชฺฌายนฺตานํ ขิยฺยนฺตานํ วิปาเจนฺตานํ.\csromanspacer{} อถ โข เต ภิกฺขู ภควโต เอตมตฺถํ อาโรเจสุํ.\csromanspacer{} อถ โข ภควา เอตสฺมิํ นิทาเน เอตสฺมิํ ปกรเณ ภิกฺขุสํฆํ สนฺนิปาตาเปตฺวา ภิกฺขูนํ ตทนุจฺฉวิกํ ตทนุโลมิกํ ธมฺมิํ กถํ กตฺวา ภิกฺขู อามนฺเตสิ เตน หิ ภิกฺขเว ภิกฺขูนํ สิกฺขาปทํ ปญฺญเปสฺสามิ ทส อตฺถวเส ปฏิจฺจ สํฆสุฏฺฐุตาย สํฆผาสุตาย ฯเปฯ สทฺธมฺมฏฺฐิติยา วินยานุคฺคหาย.\csromanspacer{} เอวญฺจ ปน ภิกฺขเว อิมํ สิกฺขาปทํ อุทฺทิเสยฺยาถ\csromandash{}}

\proseitem{368}{\textbf{“นวํ ปน ภิกฺขุนา จีวรลาเภน ติณฺณํ ทุพฺพณฺณกรณานํ อญฺญตรํ ทุพฺพณฺณกรณํ อาทาตพฺพํ นีลํ วา กทฺทมํ วา กาฬสามํ วา,} \csromanfolio{159}\textbf{อนาทา เจ ภิกฺขุ ติณฺณํ ทุพฺพณฺณกรณานํ อญฺญตรํ ทุพฺพณฺณกรณํ นวํ จีวรํ ปริภุญฺเชยฺย, ปาจิตฺติยนฺ”}ติ.}

\proseitem{369}{\textbf{นวํ} นาม อกตกปฺปํ วุจฺจติ.}

\prose{\textbf{จีวรํ} นาม ฉนฺนํ จีวรานํ อญฺญตรํ จีวรํ.}

\prose{\textbf{ติณฺณํ ทุพฺพณฺณกรณานํ อญฺญตรํ ทุพฺพณฺณกรณํ อาทาตพฺพนฺ}ติ อนฺตมโส กุสคฺเคนปิ อาทาตพฺพํ.}

\prose{\textbf{นีลํ} นาม เทฺว นีลานิ กํสนีลํ ปลาสนีลํ.}

\prose{\textbf{กทฺทโม} นาม โอทโก วุจฺจติ.}

\prose{\textbf{กาฬสามํ} นาม ยํกิญฺจิ กาฬสามกํ\footnote{กาฬกํ (สี, สฺยา)}.}

\prose{\textbf{อนาทา เจ ภิกฺขุ ติณฺณํ ทุพฺพณฺณกรณานํ อญฺญตรํ ทุพฺพณฺณกรณนฺ}ติ อนฺตมโส กุสคฺเคนปิ อนาทิยิตฺวา ติณฺณํ ทุพฺพณฺณกรณานํ อญฺญตรํ ทุพฺพณฺณกรณํ นวํ จีวรํ ปริภุญฺชติ, อาปตฺติ ปาจิตฺติยสฺส.}

\proseitem{370}{อนาทินฺเน อนาทินฺนสญฺญี ปริภุญฺชติ, อาปตฺติ ปาจิตฺติยสฺส.\csromanspacer{} อนาทินฺเน เวมติโก ปริภุญฺชติ, อาปตฺติ ปาจิตฺติยสฺส.\csromanspacer{} อนาทินฺเน อาทินฺนสญฺญี ปริภุญฺชติ, อาปตฺติ ปาจิตฺติยสฺส.}

\prose{อาทินฺเน อนาทินฺนสญฺญี, อาปตฺติ ทุกฺกฏสฺส.\csromanspacer{} อาทินฺเน เวมติโก อาปตฺติ ทุกฺกฏสฺส.\csromanspacer{} อาทินฺเน อาทินฺนสญฺญี, อนาปตฺติ.}

\proseitemwithcloserruled{371}{อนาปตฺติ อาทิยิตฺวา ปริภุญฺชติ, กปฺโป นฏฺโฐ โหติ, กปฺปกโตกาโส ชิณฺโณ โหติ, กปฺปกเตน อกปฺปกตํ สํสิพฺพิตํ โหติ, อคฺคเฬ อนุวาเต ปริภณฺเฑ อุมฺมตฺตกสฺส อาทิกมฺมิกสฺสาติ.}{ทุพฺพณฺณกรณสิกฺขาปทํ นิฏฺฐิตํ อฏฺฐมํ.}
\csromanmatikaanchor{mtk.66}
\csromantocmarkref{subsection}{9. วิกปฺปนสิกฺขาปท}{mtk.66}
\bookmark[dest={mtk.66},level=2]{9. วิกปฺปนสิกฺขาปท}
\csromanheader{2}{6. \textbf{สุราปานวคฺค 9. วิกปฺปนสิกฺขาปท}}
\proseitem{372}{เตน สมเยน พุทฺโธ ภควา สาวตฺถิยํ วิหรติ เชตวเน อนาถปิณฺฑิกสฺส อาราเม.\csromanspacer{} เตน โข ปน สมเยน อายสฺมา \csromanfolio{160}อุปนนฺโท สกฺยปุตฺโต ภาตุโน สทฺธิวิหาริกสฺส ภิกฺขุโน สามํ จีวรํ วิกปฺเปตฺวา อปฺปจฺจุทฺธารณํ\footnote{อปจฺจุทฺธารกํ (สี, สฺยา)} ปริภุญฺชติ.\csromanspacer{} อถ โข โส ภิกฺขุ ภิกฺขูนํ เอตมตฺถํ อาโรเจสิ “อยํ อาวุโส อายสฺมา อุปนนฺโท สกฺยปุตฺโต มยฺหํ จีวรํ สามํ วิกปฺเปตฺวา อปฺปจฺจุทฺธารณํ ปริภุญฺชตี”ติ.\csromanspacer{} เย เต ภิกฺขู อปฺปิจฺฉา ฯเปฯ เต อุชฺฌายนฺติ ขิยฺยนฺติ วิปาเจนฺติ “กถํ หิ นาม อายสฺมา อุปนนฺโท สกฺยปุตฺโต ภิกฺขุสฺส สามํ จีวรํ วิกปฺเปตฺวา อปฺปจฺจุทฺธารณํ ปริภุญฺชิสฺสตี”ติ ฯเปฯ.\csromanspacer{} สจฺจํ กิร ตฺวํ อุปนนฺท ภิกฺขุสฺส สามํ จีวรํ วิกปฺเปตฺวา อปฺปจฺจุทฺธารณํ ปริภุญฺชสีติ.\csromanspacer{} สจฺจํ ภควาติ.\csromanspacer{} วิครหิ พุทฺโธ ภควา ฯเปฯ.\csromanspacer{} กถํ หิ นาม ตฺวํ โมฆปุริส ภิกฺขุสฺส สามํ จีวรํ วิกปฺเปตฺวา อปฺปจฺจุทฺธารณํ ปริภุญฺชิสฺสสิ.\csromanspacer{} เนตํ โมฆปุริส อปฺปสนฺนานํ วา ปสาทาย ฯเปฯ.\csromanspacer{} เอวญฺจ ปน ภิกฺขเว อิมํ สิกฺขาปทํ อุทฺทิเสยฺยาถ\csromandash{}}

\proseitem{373}{\textbf{“โย ปน ภิกฺขุ ภิกฺขุสฺส วา ภิกฺขุนิยา วา สิกฺขมานาย วา สามเณรสฺส วา สามเณริยา วา สามํ จีวรํ วิกปฺเปตฺวา อปฺปจฺจุทฺธารณํ}1 ปริภุญฺเชยฺย, ปาจิตฺติยนฺ”ติ.}

\proseitem{374}{\textbf{โย ปนา}ติ โย ยาทิโส ฯเปฯ.\csromanspacer{} \textbf{ภิกฺขู}ติ ฯเปฯ อยํ อิมสฺมิํ อตฺเถ อธิปฺเปโต ภิกฺขูติ.}

\prose{\textbf{ภิกฺขุสฺสา}ติ อญฺญสฺส ภิกฺขุสฺส.}

\prose{\textbf{ภิกฺขุนี} นาม อุภโตสํเฆ อุปสมฺปนฺนา.}

\prose{\textbf{สิกฺขมานา} นาม เทฺว วสฺสานิ ฉสุ ธมฺเมสุ สิกฺขิตสิกฺขา.}

\prose{\textbf{สามเณโร} นาม ทสสิกฺขาปทิโก.}

\prose{\textbf{สามเณรี} นาม ทสสิกฺขาปทิกา.}

\prose{\textbf{สามนฺ}ติ สยํ วิกปฺเปตฺวา.}

\prose{\textbf{จีวรํ} นาม ฉนฺนํ จีวรานํ อญฺญตรํ จีวรํ วิกปฺปนุปคํ ปจฺฉิมํ.}

\prose{\textbf{วิกปฺปนา} นาม เทฺว วิกปฺปนา สมฺมุขาวิกปฺปนา จ ปรมฺมุขาวิกปฺปนา จ.}

\prose{\textbf{สมฺมุขาวิกปฺปนา} นาม “อิมํ จีวรํ ตุยฺหํ วิกปฺเปมิ อิตฺถนฺนามสฺส วา”ติ.}

\prose{\csromanfolio{161}\textbf{ปรมฺมุขาวิกปฺปนา} นาม “อิมํ จีวรํ วิกปฺปนตฺถาย ตุยฺหํ ทมฺมี”ติ.\csromanspacer{} เตน วตฺตพฺโพ “โก เต มิตฺโต วา สนฺทิฏฺโฐ วา”ติ.\csromanspacer{} อิตฺถนฺนาโม จ อิตฺถนฺนาโม จาติ.\csromanspacer{} เตน วตฺตพฺโพ “อหํ เตสํ ทมฺมิ, เตสํ สนฺตกํ ปริภุญฺช วา วิสฺสชฺเชหิ วา ยถาปจฺจยํ วา กโรหี”ติ.}

\prose{\textbf{อปฺปจฺจุทฺธารณํ} นาม ตสฺส วา อทินฺนํ, ตสฺส วา อวิสฺสสนฺโต ปริภุญฺชติ, อาปตฺติ ปาจิตฺติยสฺส.}

\proseitem{375}{อปฺปจฺจุทฺธารเณ อปฺปจฺจุทฺธารณสญฺญี ปริภุญฺชติ, อาปตฺติ ปาจิตฺติยสฺส.\csromanspacer{} อปฺปจฺจุทฺธารเณ เวมติโก ปริภุญฺชติ, อาปตฺติ ปาจิตฺติยสฺส.\csromanspacer{} อปฺปจฺจุทฺธารเณ ปจฺจุทฺธารณสญฺญี ปริภุญฺชติ, อาปตฺติ ปาจิตฺติยสฺส.}

\prose{อธิฏฺเฐติ วา วิสฺสชฺเชติ วา, อาปตฺติ ทุกฺกฏสฺส.\csromanspacer{} ปจฺจุทฺธารเณ อปฺปจฺจุทฺธารณสญฺญี, อาปตฺติ ทุกฺกฏสฺส.\csromanspacer{} ปจฺจุทฺธารเณ เวมติโก, อาปตฺติ ทุกฺกฏสฺส.\csromanspacer{} ปจฺจุทฺธารเณ ปจฺจุทฺธารณสญฺญี, อนาปตฺติ.}

\proseitemwithcloserruled{376}{อนาปตฺติ โส วา เทติ, ตสฺส วา วิสฺสสนฺโต ปริภุญฺชติ, อุมฺมตฺตกสฺส อาทิกมฺมิกสฺสาติ.}{วิกปฺปนสิกฺขาปทํ นิฏฺฐิตํ นวมํ.}
\csromanmatikaanchor{mtk.67}
\csromantocmarkref{subsection}{10. จีวรอปนิธานสิกฺขาปท}{mtk.67}
\bookmark[dest={mtk.67},level=2]{10. จีวรอปนิธานสิกฺขาปท}
\csromanheader{2}{6. สุราปานวคฺค 10. จีวร\-อปนิธาน\-สิกฺขาปท}
\proseitem{377}{เตน สมเยน พุทฺโธ ภควา สาวตฺถิยํ วิหรติ เชตวเน อนาถปิณฺฑิกสฺส อาราเม.\csromanspacer{} เตน โข ปน สมเยน สตฺตรสวคฺคิยา ภิกฺขู อสนฺนิหิตปริกฺขารา โหนฺติ.\csromanspacer{} ฉพฺพคฺคิยา ภิกฺขู สตฺตรสวคฺคิยานํ ภิกฺขูนํ ปตฺตมฺปิ จีวรมฺปิ อปนิเธนฺติ.\csromanspacer{} สตฺตรสวคฺคิยา ภิกฺขู ฉพฺพคฺคิเย ภิกฺขู เอตทโวจุํ “เทถาวุโส อมฺหากํ ปตฺตมฺปิ จีวรมฺปี”ติ, ฉพฺพคฺคิยา ภิกฺขู หสนฺติ, เต โรทนฺติ.\csromanspacer{} ภิกฺขู เอวมาหํสุ “กิสฺส ตุเมฺห อาวุโส โรทถา”ติ.\csromanspacer{} อิเม อาวุโส ฉพฺพคฺคิยา ภิกฺขู อมฺหากํ ปตฺตมฺปิ จีวรมฺปิ อปนิเธนฺตีติ.\csromanspacer{} เย เต ภิกฺขู อปฺปิจฺฉา ฯเปฯ เต อุชฺฌายนฺติ ขิยฺยนฺติ วิปาเจนฺติ “กถํ หิ นาม ฉพฺพคฺคิยา ภิกฺขู ภิกฺขูนํ \csromanfolio{162}ปตฺตมฺปิ จีวรมฺปิ อปนิเธสฺสนฺตี”ติ ฯเปฯ.\csromanspacer{} สจฺจํ กิร ตุเมฺห ภิกฺขเว ภิกฺขูนํ ปตฺตมฺปิ จีวรมฺปิ อปนิเธถาติ.\csromanspacer{} สจฺจํ ภควาติ.\csromanspacer{} วิครหิ พุทฺโธ ภควา ฯเปฯ.\csromanspacer{} กถํ หิ นาม ตุเมฺห โมฆปุริสา ภิกฺขูนํ ปตฺตมฺปิ จีวรมฺปิ อปนิเธสฺสถ.\csromanspacer{} เนตํ โมฆปุริสา อปฺปสนฺนานํ วา ปสาทาย ฯเปฯ.\csromanspacer{} เอวญฺจ ปน ภิกฺขเว อิมํ สิกฺขาปทํ อุทฺทิเสยฺยาถ\csromandash{}}

\proseitem{378}{\textbf{“โย ปน ภิกฺขุ ภิกฺขุสฺส ปตฺตํ วา จีวรํ วา นิสีทนํ วา สูจิฆรํ วา กายพนฺธนํ วา อปนิเธยฺย วา อปนิธาเปยฺย วา อนฺตมโส หสาเปกฺโขปิ, ปาจิตฺติยนฺ”}ติ.}

\proseitem{379}{\textbf{โย ปน}ติ โย ยาทิโส ฯเปฯ.\csromanspacer{} \textbf{ภิกฺขู}ติ ฯเปฯ อยํ อิมสฺมิํ อตฺเถ อธิปฺเปโต ภิกฺขูติ.}

\prose{\textbf{ภิกฺขุสฺสา}ติ อญฺญสฺส ภิกฺขุสฺส.}

\prose{\textbf{ปตฺโต} นาม เทฺว ปตฺตา อโยปตฺโต มตฺติกาปตฺโต.}

\prose{\textbf{จีวรํ} นาม ฉนฺนํ จีวรานํ อญฺญตรํ จีวรํ วิกปฺปนุปคํ ปจฺฉิมํ.}

\prose{\textbf{นิสีทนํ} นาม สทสํ วุจฺจติ.}

\prose{\textbf{สูจิฆรํ} นาม สสูจิกํ วา อสูจิกํ วา.}

\prose{\textbf{กายพนฺธนํ} นาม เทฺว กายพนฺธนานิ ปฏฺฏิกา สูกรนฺตกํ.}

\prose{\textbf{อปนิเธยฺย} \textbf{วา}ติ สยํ อปนิเธติ, อาปตฺติ ปาจิตฺติยสฺส.}

\prose{\textbf{อปนิธาเปยฺย} \textbf{วา}ติ อญฺญํ อาณาเปติ, อาปตฺติ ปาจิตฺติยสฺส.\csromanspacer{} สกิํ อาณตฺโต พหุกมฺปิ อปนิเธติ, อาปตฺติ ปาจิตฺติยสฺส.}

\prose{\textbf{อนฺตมโส หสาเปกฺโขปี}ติ กีฬาธิปฺปาโย.}

\proseitem{380}{อุปสมฺปนฺเน อุปสมฺปนฺนสญฺญี ปตฺตํ วา จีวรํ วา นิสีทนํ วา สูจิฆรํ วา กายพนฺธนํ วา อปนิเธติ วา อปนิธาเปติ วา อนฺตมโส หสาเปกฺโขปิ, อาปตฺติ ปาจิตฺติยสฺส.\csromanspacer{} อุปสมฺปนฺเน เวมติโก ฯเปฯ.\csromanspacer{} อุปสมฺปนฺเน อนุปสมฺปนฺนสญฺญี ปตฺตํ วา จีวรํ วา นิสีทนํ วา สูจิฆรํ วา กายพนฺธนํ วา อปนิเธติ วา อปนิธาเปติ วา อนฺตมโส หสาเปกฺโขปิ, อาปตฺติ ปาจิตฺติยสฺส.}

\prose{\csromanfolio{163}อญฺญํ ปริกฺขารํ อปนิเธติ วา อปนิธาเปติ วา อนฺตมโส หสาเปกฺโขปิ, อาปตฺติ ทุกฺกฏสฺส.\csromanspacer{} อนุปสมฺปนฺนสฺส ปตฺตํ วา จีวรํ วา อญฺญํ วา ปริกฺขารํ อปนิเธติ วา อปนิธาเปติ วา อนฺตมโส หสาเปกฺโขปิ, อาปตฺติ ทุกฺกฏสฺส.\csromanspacer{} อนุปสมฺปนฺเน อุปสมฺปนฺนสญฺญี, อาปตฺติ ทุกฺกฏสฺส.\csromanspacer{} อนุปสมฺปนฺเน เวมติโก, อาปตฺติ ทุกฺกฏสฺส.\csromanspacer{} อนุปสมฺปนฺเน อนุปสมฺปนฺนสญฺญี, อาปตฺติ ทุกฺกฏสฺส.}

\proseitemwithcloser{381}{อนาปตฺติ น หสาธิปฺปาโย, ทุนฺนิกฺขิตฺตํ ปฏิสาเมติ, “ธมฺมิํ กถํ กตฺวา ทสฺสามี”ติ ปฏิสาเมติ, อุมฺมตฺตกสฺส อาทิกมฺมิกสฺสาติ.}{จีวร\-อปนิธาน\-สิกฺขาปทํ นิฏฺฐิตํ ทสมํ.}
\vspace{\tipitakaparskip}
\csromancenter{สุราปานวคฺโค ฉฏฺโฐ.}
\csromancenter{ตสฺสุทฺทานํ}
\vspace{0.5\baselineskip}
\setlength{\csromangathapagewidth}{0pt}
\csromangathameasuregroup{สุรา องฺคุลิ หาโสจ,}{\csromangathabat{สุรา องฺคุลิ หาโสจ,}{อนาทริยญฺจ ภิํสนํ.} \\ โชตินหานทุพฺพณฺณํ, สามํ อปนิเธน จาติ.}
\csromangathasetleft{สุรา องฺคุลิ หาโสจ,}
\csromangathagroup{\csromangathastanza{\csromangathabat{สุรา องฺคุลิ หาโสจ\footnote{โตยญฺจ (อิติปิ)},}{อนาทริยญฺจ ภิํสนํ.} \\ โชตินหานทุพฺพณฺณํ, สามํ อปนิเธน จาติ.}}

\csromanmatikaanchor{mtk.68}
\csromanmatikaanchor{mtk.69}
\csromantocmarknopage{section}{7. สปฺปาณกวคฺค}{mtk.68}
\bookmark[dest={mtk.68},level=1]{7. สปฺปาณกวคฺค}
\csromantocmarkref{subsection}{1. สญฺจิจฺจสิกฺขาปท}{mtk.69}
\bookmark[dest={mtk.69},level=2]{1. สญฺจิจฺจสิกฺขาปท}
\setcsromanheadh{7. สปฺปาณกวคฺค}
\csromanheader{2}{7. \textbf{สปฺปาณกวคฺค 1. สญฺจิจฺจสิกฺขาปท}}
\proseitem{382}{เตน สมเยน พุทฺโธ ภควา สาวตฺถิยํ วิหรติ เชตวเน อนาถปิณฺฑิกสฺส อาราเม.\csromanspacer{} เตน โข ปน สมเยน อายสฺมา อุทายี อิสฺสาโส โหติ, กากา จสฺส อมนาปา โหนฺติ, โส กาเก วิชฺฌิตฺวา วิชฺฌิตฺวา สีสํ ฉินฺทิตฺวา สูเล ปฏิปาฏิยา ฐเปสิ.\csromanspacer{} ภิกฺขู เอวมาหํสุ “เกนิเม อาวุโส กากา ชีวิตา โวโรปิตา”ติ.\csromanspacer{} มยา อาวุโส, อมนาปา เม กากาติ.\csromanspacer{} เย เต ภิกฺขู อปฺปิจฺฉา ฯเปฯ เต อุชฺฌายนฺติ ขิยฺยนฺติ วิปาเจนฺติ “กถํ หิ นาม อายสฺมา อุทายี สญฺจิจฺจ ปาณํ ชีวิตา โวโรเปสฺสตี”ติ ฯเปฯ.\csromanspacer{} สจฺจํ กิร ตฺวํ อุทายิ สญฺจิจฺจ ปาณํ ชีวิตา โวโรเปสีติ.\csromanspacer{} สจฺจํ ภควาติ.\csromanspacer{} วิครหิ พุทฺโธ ภควา ฯเปฯ.\csromanspacer{} กถํ หิ นาม ตฺวํ โมฆปุริส สญฺจิจฺจ ปาณํ ชีวิตา โวโรเปสฺสสิ.\csromanspacer{} เนตํ โมฆปุริส อปฺปสนฺนานํ วา ปสาทาย ฯเปฯ.\csromanspacer{} เอวญฺจ ปน ภิกฺขเว อิมํ สิกฺขาปทํ อุทฺทิเสยฺยาถ\csromandash{}}

\proseitem{383}{\csromanfolio{164}“\textbf{โย ปน ภิกฺขุ สญฺจิจฺจ ปาณํ ชีวิตา โวโรเปยฺย, ปาจิตฺติยนฺ”ติ.}}

\proseitem{384}{\textbf{โย ปนา}ติ โย ยาทิโส ฯเปฯ.\csromanspacer{} \textbf{ภิกฺขู}ติ ฯเปฯ อยํ อิมสฺมิํ อตฺเถ อธิปฺเปโต ภิกฺขูติ.}

\prose{\textbf{สญฺจิจฺจา}ติ ชานนฺโต สญฺชานนฺโต เจจฺจ อภิวิตริตฺวา วีติกฺกโม.}

\prose{\textbf{ปาโณ} นาม ติรจฺฉานคตปาโณ วุจฺจติ.}

\prose{\textbf{ชีวิตา โวโรเปยฺยา}ติ ชีวิตินฺทฺริยํ อุปจฺฉินฺทติ อุปโรเธติ สนฺตติํ วิโกเปติ, อาปตฺติ ปาจิตฺติยสฺส.}

\proseitem{385}{ปาเณ ปาณสญฺญี ชีวิตา โวโรเปติ, อาปตฺติ ปาจิตฺติยสฺส.\csromanspacer{} ปาเณ เวมติโก ชีวิตา โวโรเปติ, อาปตฺติ ทุกฺกฏสฺส.\csromanspacer{} ปาเณ อปฺปาณสญฺญี ชีวิตา โวโรเปติ, อนาปตฺติ.\csromanspacer{} อปฺปาเณ ปาณสญฺญี อาปตฺติ ทุกฺกฏสฺส.\csromanspacer{} อปฺปาเณ เวมติโก, อาปตฺติ ทุกฺกฏสฺส.\csromanspacer{} อปฺปาเณ อปฺปาณสญฺญี, อนาปตฺติ.}

\proseitemwithcloserruled{386}{อนาปตฺติ อสญฺจิจฺจ อสฺสติยา อชานนฺตสฺส น มรณาธิปฺปายสฺส อุมฺมตฺตกสฺส อาทิกมฺมิกสฺสาติ.}{สญฺจิจฺจสิกฺขาปทํ นิฏฺฐิตํ ปฐมํ.}
\csromanmatikaanchor{mtk.70}
\csromantocmarkref{subsection}{2. สปฺปาณกสิกฺขาปท}{mtk.70}
\bookmark[dest={mtk.70},level=2]{2. สปฺปาณกสิกฺขาปท}
\csromanheader{2}{7. \textbf{สปฺปาณกวคฺค 2. สปฺปาณกสิกฺขาปท}}
\proseitem{387}{เตน สมเยน พุทฺโธ ภควา สาวตฺถิยํ วิหรติ เชตวเน อนาถปิณฺฑิกสฺส อาราเม.\csromanspacer{} เตน โข ปน สมเยน ฉพฺพคฺคิยา ภิกฺขู ชานํ สปฺปาณกํ อุทกํ ปริภุญฺชนฺติ.\csromanspacer{} เย เต ภิกฺขู อปฺปิจฺฉา ฯเปฯ เต อุชฺฌายนฺติ ขิยฺยนฺติ วิปาเจนฺติ “กถํ หิ นาม ฉพฺพคฺคิยา ภิกฺขู ชานํ สปฺปาณกํ อุทกํ ปริภุญฺชิสฺสนฺตี”ติ ฯเปฯ.\csromanspacer{} สจฺจํ กิร ตุเมฺห ภิกฺขเว ชานํ สปฺปาณกํ อุทกํ ปริภุญฺชถาติ.\csromanspacer{} สจฺจํ ภควาติ.\csromanspacer{} วิครหิ พุทฺโธ ภควา ฯเปฯ.\csromanspacer{} กถํ หิ นาม ตุเมฺห โมฆปุริสา ชานํ สปฺปาณกํ อุทกํ ปริภุญฺชิสฺสถ. \csromanfolio{165}เนตํ โมฆปุริสา อปฺปสนฺนานํ วา ปสาทาย ฯเปฯ.\csromanspacer{} เอวญฺจ ปน ภิกฺขเว อิมํ สิกฺขาปทํ อุทฺทิเสยฺยาถ\csromandash{}}

\proseitem{388}{\textbf{“โย ปน ภิกฺขุ ชานํ สปฺปาณกํ อุทกํ ปริภุญฺเชยฺย, ปาจิตฺติยนฺ”}ติ.}

\proseitem{389}{\textbf{โย ปนา}ติ โย ยาทิโส ฯเปฯ.\csromanspacer{} \textbf{ภิกฺขู}ติ ฯเปฯ อยํ อิมสฺมิํ อตฺเถ อธิปฺเปโต ภิกฺขูติ.}

\prose{\textbf{ชานาติ} นาม สามํ วา ชานาติ, อญฺเญ วา ตสฺส อาโรเจนฺติ. “สปฺปาณกนฺ”ติ ชานนฺโต “ปริโภเคน มริสฺสนฺตี”ติ ชานนฺโต ปริภุญฺชติ, อาปตฺติ ปาจิตฺติยสฺส.}

\proseitem{390}{สปฺปาณเก สปฺปาณกสญฺญี ปริภุญฺชติ, อาปตฺติ ปาจิตฺติยสฺส.\csromanspacer{} สปฺปาณเก เวมติโก ปริภุญฺชติ, อาปตฺติ ทุกฺกฏสฺส.\csromanspacer{} สปฺปาณเก อปฺปาณกสญฺญี ปริภุญฺชติ, อนาปตฺติ.\csromanspacer{} อปฺปาณเก สปฺปาณกสญฺญี, อาปตฺติ ทุกฺกฏสฺส.\csromanspacer{} อปฺปาณเก เวมติโก, อาปตฺติ ทุกฺกฏสฺส.\csromanspacer{} อปฺปาณเก อปฺปาณกสญฺญี, อนาปตฺติ.}

\proseitemwithcloserruled{391}{อนาปตฺติ “สปฺปาณกนฺ”ติ อชานนฺโต, “อปฺปาณกนฺ”ติ ชานนฺโต, “ปริโภเคน น มริสฺสนฺตี”ติ ชานนฺโต ปริภุญฺชติ, อุมฺมตฺตกสฺส อาทิกมฺมิกสฺสาติ.}{สปฺปาณกสิกฺขาปทํ นิฏฺฐิตํ ทุติยํ.}
\csromanmatikaanchor{mtk.71}
\csromantocmarkref{subsection}{3. อุกฺโกฏนสิกฺขาปท}{mtk.71}
\bookmark[dest={mtk.71},level=2]{3. อุกฺโกฏนสิกฺขาปท}
\csromanheader{2}{7. \textbf{สปฺปาณกวคฺค 3. อุกฺโกฏนสิกฺขาปท}}
\proseitem{392}{เตน สมเยน พุทฺโธ ภควา สาวตฺถิยํ วิหรติ เชตวเน อนาถปิณฺฑิกสฺส อาราเม.\csromanspacer{} เตน โข ปน สมเยน ฉพฺพคฺคิยา ภิกฺขู ชานํ ยถาธมฺมํ นิหตาธิกรณํ ปุน กมฺมาย อุกฺโกเฏนฺติ “อกตํ กมฺมํ ทุกฺกฏํ กมฺมํ ปุน กาตพฺพํ กมฺมํ อนิหตํ ทุนฺนิหตํ ปุน นิหนิตพฺพนฺ”ติ.\csromanspacer{} เย เต ภิกฺขู อปฺปิจฺฉา ฯเปฯ เต อุชฺฌายนฺติ ขิยฺยนฺติ วิปาเจนฺติ “กถํ หิ นาม ฉพฺพคฺคิยา ภิกฺขู ชานํ ยถาธมฺมํ นิหตาธิกรณํ ปุนกมฺมาย อุกฺโกเฏสฺสนฺตี \csromanfolio{166}”ติ ฯเปฯ.\csromanspacer{} สจฺจํ กิร ตุเมฺห ภิกฺขเว ชานํ ยถาธมฺมํ นิหตาธิกรณํ ปุนกมฺมาย อุกฺโกเฏถาติ.\csromanspacer{} สจฺจํ ภควาติ.\csromanspacer{} วิครหิ พุทฺโธ ภควา ฯเปฯ.\csromanspacer{} กถํ หิ นาม ตุเมฺห โมฆปุริสา ชานํ ยถาธมฺมํ นิหตาธิกรณํ ปุนกมฺมาย อุกฺโกเฏสฺสถ.\csromanspacer{} เนตํ โมฆปุริสา อปฺปสนฺนานํ วา ปสาทาย ฯเปฯ.\csromanspacer{} เอวญฺจ ปน ภิกฺขเว อิมํ สิกฺกฺกาปทํ อุทฺทิเสยฺยาถ\csromandash{}}

\proseitem{393}{\textbf{“โย ปน ภิกฺขุ ชานํ ยถาธมฺมํ นิหตาธิกรณํ ปุนกมฺมาย อุกฺโกเฏยฺย, ปาจิตฺติยนฺ”}ติ.}

\proseitem{394}{\textbf{โย ปนา}ติ โย ยาทิโส ฯเปฯ.\csromanspacer{} \textbf{ภิกฺขู}ติ ฯเปฯ อยํ อิมสฺมิํ อตฺเถ อธิปฺเปโต ภิกฺขูติ.}

\prose{\textbf{ชานาติ} นาม สามํ วา ชานาติ, อญฺเญ วา ตสฺส อาโรเจนฺติ, โส วา อาโรเจติ.}

\prose{\textbf{ยถาธมฺมํ} นาม ธมฺเมน วินเยน สตฺถุสาสเนน กตํ, เอตํ ยถาธมฺมํ นาม.}

\prose{\textbf{อธิกรณํ} นาม จตฺตาริ อธิกรณานิ วิวาทาธิกรณํ อนุวาทามิกรณํ อาปตฺตาธิกรณํ กิจฺจาธิกรณํ.}

\prose{\textbf{ปุนกมฺมาย อุกฺโกเฏยฺยา}ติ “อกตํ กมฺมํ ทุกฺกฏํ กมฺมํ ปุนกาตพฺพํ กมฺมํ อนิหตํ ทุนฺนิหตํ ปุน นิหนิตพฺพนฺ”ติ อุกฺโกเฏติ, อาปตฺติ ปาจิตฺติยสฺส.}

\proseitem{395}{ธมฺมกมฺเม ธมฺมกมฺมสญฺญี อุกฺโกเฏติ, อาปตฺติ ปาจิตฺติยสฺส.\csromanspacer{} ธมฺมกมฺเม เวมติโก อุกฺโกเฏติ, อาปตฺติ ทุกฺกฏสฺส.\csromanspacer{} ธมฺมกมฺเม อธมฺมกมฺมสญฺญี อุกฺโกเฏติ, อนาปตฺติ.\csromanspacer{} อธมฺมกมฺเม ธมฺมกมฺมสญฺญี, อาปตฺติ ทุกฺกฏสฺส.\csromanspacer{} อธมฺมกมฺเม เวมติโก, อาปตฺติ ทุกฺกฏสฺส.\csromanspacer{} อธมฺมกมฺเม อธมฺมกมฺมสญฺญี, อนาปตฺติ.}

\proseitemwithcloser{396}{อนาปตฺติ “อธมฺเมน วา วคฺเคน วา นกมฺมารหสฺส วา กมฺมํ กตนฺ”ติ ชานนฺโต อุกฺโกเฏติ, อุมฺมตฺตกสฺส อาทิกมฺมิกสฺสาติ.}{อุกฺโกฏนสิกฺขาปทํ นิฏฺฐิตํ ตติยํ.}
\csromanmatikaanchor{mtk.72}
\csromantocmarkref{subsection}{4. ทุฏฺฐุลฺลสิกฺขาปท}{mtk.72}
\bookmark[dest={mtk.72},level=2]{4. ทุฏฺฐุลฺลสิกฺขาปท}
\csromanheader{2}{7. \textbf{สปฺปาณกวคฺค 4. ทุฏฺฐุลฺลสิกฺขาปท}}
\proseitem{397}{\csromanfolio{167}เตน สมเยน พุทฺโธ ภควา สาวตฺถิยํ วิหรติ เชตวเน อนาถปิณฺฑิกสฺส อาราเม.\csromanspacer{} เตน โข ปน สมเยน อายสฺมา อุปนนฺโท สกฺยปุตฺโต สญฺเจตนิกํ สุกฺกวิสฺสฏฺฐิํ อาปตฺติํ อาปชฺชิตฺวา ภาตุโน สทฺธิวิหาริกสฺส ภิกฺขุโน อาโรเจสิ “อหํ อาวุโส สญฺเจตนิกํ สุกฺกวิสฺสฏฺฐิํ อาปตฺติํ อาปนฺโน, มา กสฺสจิ อาโรเจหี”ติ.\csromanspacer{} เตน โข ปน สมเยน อญฺญตโร ภิกฺขุ สญฺเจตนิกํ สุกฺกวิสฺสฏฺฐิํ อาปตฺติํ อาปชฺชิตฺวา สํฆํ ตสฺสา อาปตฺติยา ปริวาสํ ยาจิ, ตสฺส สํโฆ ตสฺสา อาปตฺติยา ปริวาสํ อทาสิ, โส ปริวสนฺโต ตํ ภิกฺขุํ ปสฺสิตฺวา เอตทโวจ “อหํ อาวุโส สญฺเจตนิกํ สุกฺกวิสฺสฏฺฐิํ อาปตฺติํ อาปชฺชิตฺวา สํฆํ ตสฺสา อาปตฺติยา ปริวาสํ ยาจิํ, ตสฺส เม สํโฆ ตสฺสา อาปตฺติยา ปริวาสํ อทาสิ, โสหํ ปริวสามิ, เวทิยามหํ\footnote{เวทยามหํ (สฺยา)} อาวุโส, เวทิยตีติ มํ อายสฺมา ธาเรตู”ติ.}

\prose{กิํ นุ โข อาวุโส โย อญฺโญปิ อิมํ อาปตฺติํ อาปชฺชติ, โสปิ เอวํ กโรตีติ.\csromanspacer{} เอวมาวุโสติ.\csromanspacer{} อยํ อาวุโส อายสฺมา อุปนนฺโท สกฺยปุตฺโต สญฺเจตนิกํ สุกฺกวิสฺสฏฺฐิํ อาปตฺติํ อาปชฺชิตฺวา\footnote{อาปชฺชิ (?)} โส เม อาโรเจติ “มา กสฺสจิ อาโรเจหี”ติ.\csromanspacer{} กิํ ปน ตฺวํ อาวุโส ปฏิจฺฉาเทสีติ.\csromanspacer{} เอวมาวุโสติ.\csromanspacer{} อถ โข โส ภิกฺขุ ภิกฺขูนํ เอตมตฺถํ อาโรเจสิ.\csromanspacer{} เย เต ภิกฺขู อปฺปิจฺฉา ฯเปฯ เต อุชฺฌายนฺติ ขิยฺยนฺติ วิปาเจนฺติ “กถํ หิ นาม ภิกฺขุ ภิกฺขุสฺส ชานํ ทุฏฺฐุลฺลํ อาปตฺติํ ปฏิจฺฉาเทสฺสตี”ติ ฯเปฯ.\csromanspacer{} สจฺจํ กิร ตฺวํ ภิกฺขุ ภิกฺขุสฺส ชานํ ทุฏฺฐุลฺลํ อาปตฺติํ ปฏิจฺฉาเทสีติ.\csromanspacer{} สจฺจํ ภควาติ.\csromanspacer{} วิครหิ พุทฺโธ ภควา ฯเปฯ.\csromanspacer{} กถํ หิ นาม ตฺวํ โมฆปุริส ภิกฺขุสฺส ชานํ ทุฏฺฐุลฺลํ อาปตฺติํ ปฏิจฺฉาเทสฺสสิ.\csromanspacer{} เนตํ โมฆปุริส อปฺปสนฺนานํ วา ปสาทาย ฯเปฯ.\csromanspacer{} เอวญฺจ ปน ภิกฺขเว อิมํ สิกฺขาปทํ อุทฺทิเสยฺยาถ\csromandash{}}

\proseitem{398}{“ยฺ\textbf{โอ ปน ภิกฺขุ ภิกฺขุสฺส ชานํ ทุฏฺฐุลฺลํ อาปตฺติํ ปฏิจฺฉาเทยฺย, ปาจิตฺติยนฺ”}ติ.}

\proseitem{399}{\csromanfolio{168}\textbf{โย ปนา}ติ โย ยาทิโส ฯเปฯ.\csromanspacer{} \textbf{ภิกฺขู}ติ ฯเปฯ อยํ อิมสฺมิํ อตฺเถ อธิปฺเปโต ภิกฺขูติ.}

\prose{\textbf{ภิกฺขุสฺสา}ติ อญฺญสฺส ภิกฺขุสฺส.}

\prose{\textbf{ชานาติ} นาม สามํ ชานาติ, อญฺเญ วา ตสฺส อาโรเจนฺติ, โส วา อาโรเจติ.}

\prose{\textbf{ทุฏฺฐุลฺลา} นาม อาปตฺติ จตฺตาริ จ ปาราชิกานิ เตรส จ สํฆาทิเสสา.}

\prose{\textbf{ปฏิจฺฉาเทยฺยา}ติ “อิมํ ชานิตฺวา โจเทสฺสนฺติ, สาเรสฺสนฺติ, ขุํเสสฺสนฺติ, วมฺเภสฺสนฺติ, มงฺกุํ กริสฺสนฺติ, นาโรเจสฺสามี”ติ ธุรํ นิกฺขิตฺตมตฺเต อาปตฺติ ปาจิตฺติยสฺส.}

\proseitem{400}{ทุฏฺฐุลฺลาย อาปตฺติยา ทุฏฺฐุลฺลาปตฺติสญฺญี ปฏิจฺฉาเทติ, อาปตฺติ ปาจิตฺติยสฺส.\csromanspacer{} ทุฏฺฐุลฺลาย อาปตฺติยา เวมติโก ปฏิจฺฉาเทติ, อาปตฺติ ทุกฺกฏสฺส.\csromanspacer{} ทุฏฺฐุลฺลาย อาปตฺติยา อทุฏฺฐุลฺลาปตฺติสญฺญี ปฏิจฺฉาเทติ, อาปตฺติ ทุกฺกฏสฺส.\csromanspacer{} อทุฏฺฐุลฺลํ อาปตฺติํ ปฏิจฺฉาเทติ, อาปตฺติ ทุกฺกฏสฺส.\csromanspacer{} อนุปสมฺปนฺนสฺส ทุฏฺฐุลฺลํ วา อทุฏฺฐุลฺลํ วา อชฺฌาจารํ ปฏิจฺฉาเทติ, อาปตฺติ ทุกฺกฏสฺส.\csromanspacer{} อทุฏฺฐุลฺลาย อาปตฺติยา ทุฏฺฐุลฺลาปตฺติสญฺญี, อาปตฺติ ทุกฺกฏสฺส.\csromanspacer{} อทุฏฺฐุลฺลาย อาปตฺติยา เวมติโก, อาปตฺติ ทุกฺกฏสฺส.\csromanspacer{} อทุฏฺฐุลฺลาย อาปตฺติยา อทุฏฺฐุลฺลาปตฺติสญฺญี, อาปตฺติ ทุกฺกฏสฺส.}

\proseitemwithcloser{401}{อนาปตฺติ “สํฆสฺส ภณฺฑนํ วา กลโห วา วิคฺคโห วา วิวาโท วา ภวิสฺสตี”ติ นาโรเจติ, “สํฆเภโท วา สํฆราชิ วา ภวิสฺสตี”ติ นาโรเจติ, “อยํ กกฺขโฬ ผรุโส ชีวิตนฺตรายํ วา พฺรหฺมจริยนฺตรายํ วา กริสฺสตี”ติ นาโรเจติ, อญฺเญ ปติรูเป ภิกฺขู อปสฺสนฺโต นาโรเจติ, น ฉาเทตุกาโม นาโรเจติ, “ปญฺญายิสฺสติ สเกน กมฺเมนา”ติ นาโรเจติ, อุมฺมตฺตกสฺส อาทิกมฺมิกสฺสาติ.}{ทุฏฺฐุลฺลสิกฺขาปทํ นิฏฺฐิตํ จตุตฺถํ.}
\csromanmatikaanchor{mtk.73}
\csromantocmarkref{subsection}{5. อูนวีสติวสฺสสิกฺขาปท}{mtk.73}
\bookmark[dest={mtk.73},level=2]{5. อูนวีสติวสฺสสิกฺขาปท}
\csromanheader{2}{7. สปฺปาณกวคฺค 5. อู\textbf{นวีสติวสฺสสิกฺขาปท}}
\proseitem{402}{\csromanfolio{169}เตน สมเยน พุทฺโธ ภควา ราชคเห วิหรติ เวฬุวเน กลนฺทกนิวาเป.\csromansymbolfootnote{*}{อิทํ วตฺถุ วิ 3. 108 ปิฏฺฐาทีสุปิ อาคตํ.}เตน โข ปน สมเยน ราชคเห สตฺตรสวคฺคิยา ทารกา สหายกา โหนฺติ.\csromanspacer{} อุปาลิทารโก เตสํ ปาโมกฺโข โหติ.\csromanspacer{} อถ โข อุปาลิสฺส มาตาปิตูนํ เอตทโหสิ “เกน นุ โข อุปาเยน อุปาลิ อมฺหากํ อจฺจเยน สุขญฺจ ชีเวยฺย, น จ กิลเมยฺยา”ติ.\csromanspacer{} อถ โข อุปาลิสฺส มาตาปิตูนํ เอตทโหสิ “สเจ โข อุปาลิ เลขํ สิกฺเขยฺย, เอวํ โข อุปาลิ อมฺหากํ อจฺจเยน สุขญฺจ ชีเวยฺย, น จ กิลเมยฺยา”ติ.\csromanspacer{} อถ โข อุปาลิสฺส มาตาปิตูนํ เอตทโหสิ “สเจ โข อุปาลิ เลขํ สิกฺขิสฺสติ, องฺคุลิโย ทุกฺขา ภวิสฺสนฺติ.\csromanspacer{} สเจ โข อุปาลิ คณนํ สิกฺเขยฺย, เอวํ โข อุปาลิ อมฺหากํ อจฺจเยน สุขญฺจ ชีเวยฺย, น จ กิลเมยฺยา”ติ.\csromanspacer{} อถ โข อุปาลิสฺส มาตาปิตูนํ เอตทโหสิ “สเจ โข อุปาลิ คณนํ สิกฺขิสฺสติ, อุร’สฺส ทุกฺโข ภวิสฺสติ.\csromanspacer{} สเจ โข อุปาลิ รูปํ สิกฺเขยฺย, เอวํ โข อุปาลิ อมฺหากํ อจฺจเยน สุขญฺจ ชีเวยฺย, น จ กิลเมยฺยา”ติ.\csromanspacer{} อถ โข อุปาลิสฺส มาตาปิตูนํ เอตทโหสิ “สเจ โข อุปาลิ รูปํ สิกฺขิสฺสติ, อกฺขีนิ ทุกฺขา ภวิสฺสนฺติ.\csromanspacer{} อิเม โข สมณา สกฺยปุตฺติยา สุขสีลา สุขสมาจารา สุโภชนานิ ภุญฺชิตฺวา นิวาเตสุ สยเนสุ สยนฺติ.\csromanspacer{} สเจ โข อุปาลิ สมเณสุ สกฺยปุตฺติเยสุ ปพฺพเชยฺย, เอวํ โข อุปาลิ อมฺหากํ อจฺจเยน สุขญฺจ ชีเวยฺย, น จ กิลเมยฺยา”ติ.}

\prose{อสฺโสสิ โข อุปาลิทารโก มาตาปิตูนํ อิมํ กถาสลฺลาปํ.\csromanspacer{} อถ โข อุปาลิทารโก เยน เต ทารกา เตนุปสงฺกมิ, อุปสงฺกมิตฺวา เต ทารเก เอตทโวจ “เอถ มยํ อยฺยา สมเณสุ สกฺยปุตฺติเยสุ ปพฺพชิสฺสามา”ติ.\csromanspacer{} สเจ โข ตฺวํ อยฺย ปพฺพชิสฺสสิ, เอวํ มยมฺปิ ปพฺพชิสฺสามาติ.\csromanspacer{} อถ โข เต ทารกา เอกเมกสฺส มาตาปิตโร อุปสงฺกมิตฺวา เอตทโวจุํ “อนุชานาถ มํ อคารสฺมา อนคาริยํ ปพฺพชฺชายา”ติ.\csromanspacer{} อถ โข เตสํ ทารกานํ มาตาปิตโร “สพฺเพปิเม ทารกา สมานจฺฉนฺทา กลฺยาณาธิปฺปายา”ติ อนุชานิํสุ.\csromanspacer{} เต ภิกฺขู อุปสงฺกมิตฺวา ปพฺพชฺชํ ยาจิํสุ.\csromanspacer{} เต ภิกฺขู ปพฺพาเชสุํ อุปสมฺปาเทสุํ.\csromanspacer{} เต รตฺติยา ปจฺจุสสมยํ ปจฺจุฏฺฐาย โรทนฺติ “ยาคุํ \csromanfolio{170}เทถ ภตฺตํ เทถ ขาทนียํ เทถา”ติ.\csromanspacer{} ภิกฺขู เอวมาหํสุ “อาคเมถ อาวุโส ยาว รตฺติ วิภายติ, สเจ ยาคุ ภวิสฺสติ ปิวิสฺสถ, สเจ ภตฺตํ ภวิสฺสติ ภุญฺชิสฺสถ, สเจ ขาทนียํ ภวิสฺสติ ขาทิสฺสถ, โน เจ ภวิสฺสติ ยาคุ วา ภตฺตํ วา ขาทนียํ วา ปิณฺฑาย จริตฺวา ภุญฺชิสฺสถา”ติ.\csromanspacer{} เอวํปิ โข เต ภิกฺขู ภิกฺขูหิ วุจฺจมานา โรทนฺติ เยว “ยาคุํ เทถ ภตฺตํ เทถ ขาทนียํ เทถา”ติ.\csromanspacer{} เสนาสนํ อูหทนฺติปิ อุมฺมิหนฺติปิ.}

\prose{อสฺโสสิ โข ภควา รตฺติยา ปจฺจูสสมยํ ปจฺจุฏฺฐาย ทารกสทฺทํ, สุตฺวาน อายสฺมนฺตํ อานนฺทํ อามนฺเตสิ “กิํ นุ โข โส อานนฺท ทารกสทฺโท”ติ.\csromanspacer{} อถ โข อายสฺมา อานนฺโท ภควโต เอตมตฺถํ อาโรเจสิ.\csromanspacer{} อถ โข ภควา เอตสฺมิํ นิทาเน เอตสฺมิํ ปกรเณ ภิกฺขุสํฆํ สนฺนิปาตาเปตฺวา ภิกฺขู ปฏิปุจฺฉิ “สจฺจํ กิร ภิกฺขเว ภิกฺขู ชานํ อูนวีสติวสฺสํ ปุคฺคลํ อุปสมฺปาเทนฺตี”ติ.\csromanspacer{} สจฺจํ ภควาติ.\csromanspacer{} วิครหิ พุทฺโธ ภควา ฯเปฯ.\csromanspacer{} กถํ หิ นาม เต ภิกฺขเว โมฆปุริสา ชานํ อูนวีสติวสฺสํ ปุคฺคลํ อุปสมฺปาเทสฺสนฺติ.\csromanspacer{} อูนกวีสติวสฺโส ภิกฺขเว ปุคฺคโล อกฺขโม โหติ สีตสฺส อุณฺหสฺส ชิฆจฺฉาย ปิปาสาย ฑํส\-มกส\-วาตา\-ตป\-สรีสป\-สมฺผสฺสานํ ทุรุตฺตานํ ทุราคตานํ วจนปถานํ, อุปฺปนฺนานํ สารีริกานํ เวทนานํ ทุกฺขานํ ติพฺพานํ ขรานํ กฏุกานํ อสาตานํ อมนาปานํ ปาณหรานํ อนธิวาสกชาติโก โหติ.\csromanspacer{} วีสติวสฺโสว โข ภิกฺขเว ปุคฺคโล ขโม โหติ สีตสฺส อุณฺหสฺส ฯเปฯ ปาณหรานํ อธิวาสกชาติโก โหติ.\csromanspacer{} เนตํ ภิกฺขเว อปฺปสนฺนานํ วา ปสาทาย ฯเปฯ.\csromanspacer{} เอวญฺจ ปน ภิกฺขเว อิมํ สิกฺขาปทํ อุทฺทิเสยฺยาถ\csromandash{}}

\proseitem{403}{\textbf{“โย ปน ภิกฺขุ ชานํ อูนวีสติวสฺสํ ปุคฺคลํ อุปสมฺปาเทยฺย, โส จ ปุคฺคโล อนุปสมฺปนฺโน, เต จ ภิกฺขู คารยฺหา, อิทํ ตสฺมิํ ปาจิตฺติยนฺ”}ติ.}

\proseitem{404}{\textbf{โย ปนา}ติ โย ยาทิโส ฯเปฯ.\csromanspacer{} \textbf{ภิกฺขู}ติ ฯเปฯ อยํ อิมสฺมิํ อตฺเถ อธิปฺเปโต ภิกฺขูติ.}

\prose{\textbf{ชานาติ} นาม สามํ วา ชานาติ, อญฺเญ วา ตสฺส อาโรเจนฺติ, โส วา อาโรเจติ.}

\prose{\csromanfolio{171}อู\textbf{นวีสติวสฺโส} นาม อปฺปตฺตวีสติวสฺโส.}

\prose{“อุปสมฺปาเทสฺสามี”ติ คณํ วา อาจริยํ วา ปตฺตํ วา จีวรํ วา ปริเยสติ, สีมํ วา สมฺมนฺนติ\footnote{สมฺมนติ (ก)}, อาปตฺติ ทุกฺกฏสฺส.\csromanspacer{} ญตฺติยา ทุกฺกฏํ.\csromanspacer{} ทฺวีหิ กมฺมวาจาหิ ทุกฺกฏา.\csromanspacer{} กมฺมวาจา ปริโยสาเน อุปชฺฌายสฺส อาปตฺติ ปาจิตฺติยสฺส.\csromanspacer{} คณสฺส จ อาจริยสฺส จ อาปตฺติ ทุกฺกฏสฺส.}

\proseitem{405}{อูนวีสติวสฺเส อูนวีสติวสฺสสญฺญี อุปสมฺปาเทติ, อาปตฺติ ปาจิตฺติยสฺส.\csromanspacer{} อูนวีสติวสฺเส เวมติโก อุปสมฺปาเทติ, อาปตฺติ ทุกฺกฏสฺส.\csromanspacer{} อูนวีสติวสฺเส ปริปุณฺณวีสติวสฺสสญฺญี อุปสมฺปาเทติ, อนาปตฺติ.\csromanspacer{} ปริปุณฺณวีสติวสฺเส อูนวีสติวสฺสสญฺญี.\csromanspacer{} อาปตฺติ ทุกฺกฏสฺส.\csromanspacer{} ปริปุณฺณวีสติวสฺเส เวมติโก, อาปตฺติ ทุกฺกฏสฺส.\csromanspacer{} ปริปุณฺณวีสติวสฺเส ปริปุณฺณวีสติวสฺสสญฺญี, อนาปตฺติ.}

\proseitemwithcloserruled{406}{อนาปตฺติ อูนวีสติวสฺสํ ปริปุณฺณวีสติวสฺสสญฺญี อุปสมฺปาเทติ, ปริปุณฺณวีสติวสฺสํ ปริปุณฺณวีสติวสฺสสญฺญี อุปสมฺปาเทติ, อุมฺมตฺตกสฺส อาทิกมฺมิกสฺสาติ.}{อูนวีสติวสฺสสิกฺขาปทํ นิฏฺฐิตํ ปญฺจมํ.}
\csromanmatikaanchor{mtk.74}
\csromantocmarkref{subsection}{6. เถยฺยสตฺถสิกฺขาปท}{mtk.74}
\bookmark[dest={mtk.74},level=2]{6. เถยฺยสตฺถสิกฺขาปท}
\csromanheader{2}{7. \textbf{สปฺปาณกวคฺค 6. เถยฺยสตฺถสิกฺขาปท}}
\proseitem{407}{เตน สมเยน พุทฺโธ ภควา สาวตฺถิยํ วิหรติ เชตวเน อนาถปิณฺฑิกสฺส อาราเม.\csromanspacer{} เตน โข ปน สมเยน อญฺญตโร สตฺโถ ราชคหา ปฏิยาโลกํ คนฺตุกาโม โหติ.\csromanspacer{} อญฺญตโร ภิกฺขุ เต มนุสฺเส เอตทโวจ “อหมฺปายสฺมนฺเตหิ สทฺธิํ คมิสฺสามี”ติ.\csromanspacer{} มยํ โข ภนฺเต สุงฺกํ ปริหริสฺสามาติ.\csromanspacer{} ปชานาถาวุโสติ.\csromanspacer{} อสฺโสสุํ โข กมฺมิยา\footnote{กมฺมิกา (สี, สฺยา)} “สตฺโถ กิร สุงฺกํ ปริหริสฺสตี”ติ.\csromanspacer{} เต มคฺเค ปริยุฏฺฐิํสุ.\csromanspacer{} อถ โข เต กมฺมิยา ตํ สตฺถํ คเหตฺวา อจฺฉินฺทิตฺวา ตํ ภิกฺขุํ เอตทโวจุํ “กิสฺส ตฺวํ ภนฺเต ชานํ เถยฺยสตฺเถน สทฺธิํ คจฺฉสี”ติ ปลิพุนฺเธตฺวา มุญฺจิํสุ.\csromanspacer{} อถ โข โส ภิกฺขุ สาวตฺถิํ คนฺตฺวา ภิกฺขูนํ \csromanfolio{172}เอตมตฺถํ อาโรเจสิ.\csromanspacer{} เย เต ภิกฺขู อปฺปิจฺฉา ฯเปฯ เต อุชฺฌายนฺติ ขิยฺยนฺติ วิปาเจนฺติ “กถํ หิ นาม ภิกฺขุ ชานํ เถยฺยสตฺเถน สทฺธิํ สํวิธาย เอกทฺธานมคฺคํ ปฏิปชฺชิสฺสตี”ติ ฯเปฯ.\csromanspacer{} สจฺจํ กิร ตฺวํ ภิกฺขุ ชานํ เถยฺยสตฺเถน สทฺธิํ สํวิธาย เอกทฺธานมคฺคํ ปฏิปชฺชสีติ.\csromanspacer{} สจฺจํ ภควาติ.\csromanspacer{} วิครหิ พุทฺโธ ภควา ฯเปฯ.\csromanspacer{} กถํ หิ นาม ตฺวํ โมฆปุริส ชานํ เถยฺยสตฺเถน สทฺธิํ สํวิธาย เอกทฺธานมคฺคํ ปฏิปชฺชิสฺสสิ.\csromanspacer{} เนตํ โมฆปุริส อปฺปสนฺนานํ วา ปสาทาย ฯเปฯ.\csromanspacer{} เอวญฺจ ปน ภิกฺขเว อิมํ สิกฺขาปทํ อุทฺทิเสยฺยาถ\csromandash{}}

\proseitem{408}{\textbf{“โย ปน ภิกฺขุ ชานํ เถยฺยสตฺเถน สทฺธิํ สํวิธาย เอกทฺธานมคฺคํ ปฏิปชฺเชยฺย อนฺตมโส คามนฺตรมฺปิ, ปาจิตฺติยนฺ”}ติ.}

\proseitem{409}{\textbf{โย ปนา}ติ โย ยาทิโส ฯเปฯ.\csromanspacer{} \textbf{ภิกฺขู}ติ ฯเปฯ อยํ อิมสฺมิํ อตฺเถ อธิปฺเปโต ภิกฺขูติ.}

\prose{\textbf{ชานาติ} นาม สามํ วา ชานาติ, อญฺเญ วา ตสฺส อาโรเจนฺติ, โส วา อาโรเจติ.}

\prose{\textbf{เถยฺยสตฺโถ} นาม โจรา กตกมฺมา วา โหนฺติ อกตกมฺมา วา, ราชานํ วา เถยฺยํ คจฺฉนฺติ สุงฺกํ วา ปริหรนฺติ.}

\prose{\textbf{สทฺธินฺ}ติ เอกโต.}

\prose{\textbf{สํวิธายา}ติ “คจฺฉามาวุโส คจฺฉาม ภนฺเต, คจฺฉาม ภนฺเต คจฺฉามาวุโส, อชฺช วา หิยฺโย วา ปเร วา คจฺฉามาติ” สํวิทหติ, อาปตฺติ ทุกฺกฏสฺส.}

\prose{\textbf{อนฺตมโส คามนฺตรมฺปี}ติ กุกฺกุฏสมฺปาเต คาเม คามนฺตเร คามนฺตเร อาปตฺติ ปาจิตฺติยสฺส.\csromanspacer{} อคามเก อรญฺเญ อทฺธโยชเน อทฺธโยชเน อาปตฺติ}

\proseitem{410}{เถยฺยสตฺเถ เถยฺยสตฺถสญฺญี สํวิธาย เอกทฺธานมคฺคํ ปฏิปชฺชติ อนฺตมโส คามนฺตรมฺปิ, อาปตฺติ ปาจิตฺติยสฺส.\csromanspacer{} เถยฺยสตฺเถ เวมติโก สํวิธาย เอกทฺธานมคฺคํ ปฏิปชฺชติ อนฺตมโส คามนฺตรมฺปิ, อาปตฺติ ทุกฺกฏสฺส.\csromanspacer{} เถยฺยสตฺเถ อเถยฺยสตฺถสญฺญี สํวิธาย เอกทฺธานมคฺคํ ปฏิปชฺชติ อนฺตมโส คามนฺตรมฺปิ, อนาปตฺติ.\csromanspacer{} ภิกฺขุ สํวิทหติ มนุสฺสา น สํวิทหนฺติ, \csromanfolio{173}อาปตฺติ ทุกฺกฏสฺส.\csromanspacer{} อเถยฺยสตฺเถ เถยฺยสตฺถสญฺญี, อาปตฺติ ทุกฺกฏสฺส.\csromanspacer{} อเถยฺยสตฺเถ เวมติโก, อาปตฺติ ทุกฺกฏสฺส.\csromanspacer{} อเถยฺยสตฺเถ อเถยฺยสตฺถสญฺญี, อนาปตฺติ.}

\proseitemwithcloserruled{411}{อนาปตฺติ อสํวิทหิตฺวา คจฺฉติ, มนุสฺสา สํวิทหนฺติ ภิกฺขุ น สํวิทหติ, วิสงฺเกเตน คจฺฉติ, อาปทาสุ อุมฺมตฺตกสฺส อาทิกมฺมิกสฺสาติ.}{เถยฺยสตฺถสิกฺขาปทํ นิฏฺฐิตํ ฉฏฺฐํ.}
\csromanmatikaanchor{mtk.75}
\csromantocmarkref{subsection}{7. สํวิธานสิกฺขาปท}{mtk.75}
\bookmark[dest={mtk.75},level=2]{7. สํวิธานสิกฺขาปท}
\csromanheader{2}{7. \textbf{สปฺปาณกวคฺค 7. สํวิธานสิกฺขาปท}}
\proseitem{412}{เตน สมเยน พุทฺโธ ภควา สาวตฺถิยํ วิหรติ เชตวเน อนาถปิณฺฑิกสฺส อาราเม.\csromanspacer{} เตน โข ปน สมเยน อญฺญตโร ภิกฺขุ โกสเลสุ ชนปเท สาวตฺถิํ คจฺฉนฺโต อญฺญตเรน คามทฺวาเรน อติกฺกมติ.\csromanspacer{} อญฺญตรา อิตฺถี สามิเกน สห ภณฺฑิตฺวา คามโต นิกฺขมิตฺวา ตํ ภิกฺขุํ ปสฺสิตฺวา เอตทโวจ “กหํ ภนฺเต อยฺโย คมิสฺสตี”ติ.\csromanspacer{} สาวตฺถิํ โข อหํ ภคินิ คมิสฺสามีติ.\csromanspacer{} อหํ อยฺเยน สทฺธิํ คมิสฺสามีติ.\csromanspacer{} เอยฺยาสิ ภคินีติ.\csromanspacer{} อถ โข ตสฺสา อิตฺถิยา สามิโก คามโต นิกฺขมิตฺวา มนุสฺเส ปุจฺฉิ “อปายฺโย\footnote{อปยฺยา (สี, สฺยา)} เอวรูปิํ อิตฺถิํ ปสฺเสยฺยาถา”ติ.\csromanspacer{} เอสายฺโย ปพฺพชิเตน สห คจฺฉตีติ.\csromanspacer{} อถ โข โส ปุริโส อนุพนฺธิตฺวา ตํ ภิกฺขุํ คเหตฺวา อาโกเฏตฺวา มุญฺจิ.\csromanspacer{} อถ โข โส ภิกฺขุ อญฺญตรสฺมิํ รุกฺขมูเล ปธูเปนฺโต นิสีทิ.\csromanspacer{} อถ โข สา อิตฺถี ตํ ปุริสํ เอตทโวจ “นายฺโย โส ภิกฺขุ มํ นิปฺปาเตสิ, อปิจ อหเมว เตน ภิกฺขุนา สทฺธิํ คจฺฉามิ, อการโก โส ภิกฺขุ, คจฺฉ นํ ขมาเปหี”ติ.\csromanspacer{} อถ โข โส ปุริโส ตํ ภิกฺขุํ ขมาเปสิ.\csromanspacer{} อถ โข โส ภิกฺขุ สาวตฺถิํ คนฺตฺวา ภิกฺขูนํ เอตมตฺถํ อาโรเจสิ.\csromanspacer{} เย เต ภิกฺขู อปฺปิจฺฉา ฯเปฯ เต อุชฺฌายนฺติ ขิยฺยนฺติ วิปาเจนฺติ “กถํ หิ นาม ภิกฺขุ มาตุคาเมน สทฺธิํ สํวิธาย เอกทฺธานมคฺคํ ปฏิปชฺชิสฺสตี”ติ ฯเปฯ.\csromanspacer{} สจฺจํ กิร ตฺวํ ภิกฺขุ มาตุคาเมน สทฺธิํ สํวิธาย เอกทฺธานมคฺคํ ปฏิปชฺชสีติ.\csromanspacer{} สจฺจํ ภควาติ.\csromanspacer{} วิครหิ พุทฺโธ ภควา ฯเปฯ.\csromanspacer{} กถํ หิ นาม ตฺวํ โมฆปุริส \csromanfolio{174}มาตุคาเมน สทฺธิํ สํวิธาย เอกทฺธานมคฺคํ ปฏิปชฺชิสฺสสิ.\csromanspacer{} เนตํ โมฆปุริส อปฺปสนฺนานํ วา ปสาทาย ฯเปฯ.\csromanspacer{} เอวญฺจ ปน ภิกฺขเว อิมํ สิกฺขาปทํ อุทฺทิเสยฺยาถ\csromandash{}}

\proseitem{413}{\textbf{“โย ปน ภิกฺขุ มาตุคาเมน สทฺธิํ สํวิธาย เอกทฺธานมคฺคํ ปฏิปชฺเชยฺย อนฺตมโส คามนฺตรมฺปิ ปาจิตฺติยนฺ”}ติ.}

\proseitem{414}{\textbf{โย ปนา}ติ โย ยาทิโส ฯเปฯ.\csromanspacer{} \textbf{ภิกฺขู}ติ ฯเปฯ อยํ อิมสฺมิํ อตฺเถ อธิปฺเปโต ภิกฺขูติ.}

\prose{\textbf{มาตุคาโม} นาม มนุสฺสิตฺถี, น ยกฺขี น เปตี น ติรจฺฉานคตา วิญฺญู ปฏิพลา สุภาสิตทุพฺภาสิตํ ทุฏฺฐุลฺลาทุฏฺฐุลฺลํ อาชานิตุํ.}

\prose{\textbf{สทฺธิ}นฺติ เอกโต.}

\prose{\textbf{สํวิธายา}ติ “คจฺฉาม ภคินิ คจฺฉามายฺย, คจฺฉามายฺย คจฺฉาม ภคินิ, อชฺช วา หิยฺโย วา ปเร วา คจฺฉามา”ติ สํวิทหติ, อาปตฺติ ทุกฺกฏสฺส.}

\prose{\textbf{อนฺตมโส คามนฺตรมฺปี}ติ กุกฺกุฏสมฺปาเต คาเม คามนฺตเร คามนฺตเร อาปตฺติ ปาจิตฺติยสฺส.\csromanspacer{} อคามเก อรญฺเญ อทฺธโยชเน อทฺธโยชเน อาปตฺติ}

\proseitem{415}{มาตุคาเม มาตุคามสญฺญี สํวิธาย เอกทฺธานมคฺคํ ปฏิปชฺชติ อนฺตมโส คามนฺตรมฺปิ, อาปตฺติ ปาจิตฺติยสฺส.\csromanspacer{} มาตุคาเม เวมติโก สํวิธาย เอกทฺธานมคฺคํ ปฏิปชฺชติ อนฺตมโส คามนฺตรมฺปิ, อาปตฺติ ปาจิตฺติยสฺส.\csromanspacer{} มาตุคาเม อมาตุคามสญฺญี สํวิธาย เอกทฺธานมคฺคํ ปฏิปชฺชติ อนฺตมโส คามนฺตรมฺปิ, อาปตฺติ ปาจิตฺติยสฺส.}

\prose{ภิกฺขุ สํวิทหติ มาตุคาโม น สํวิทหติ, อาปตฺติ ทุกฺกฏสฺส.\csromanspacer{} ยกฺขิยา วา เปติยา วา ปณฺฑเกน วา ติรจฺฉานคตมนุสฺสวิคฺคหิตฺถิยา วา สทฺธิํ สํวิธาย เอกทฺธานมคฺคํ ปฏิปชฺชติ อนฺตมโส คามนฺตรมฺปิ, อาปตฺติ ทุกฺกฏสฺส.\csromanspacer{} อมาตุคาเม มาตุคามสญฺญี, อาปตฺติ ทุกฺกฏสฺส.\csromanspacer{} อมาตุคาเม เวมติโก, อาปตฺติ ทุกฺกฏสฺส.\csromanspacer{} อมาตุคาเม อมาตุคามสญฺญี อนาปตฺติ.}

\proseitemwithcloserruled{416}{\csromanfolio{175}อนาปตฺติ อสํวิทหิตฺวา คจฺฉติ, มาตุคาโม สํวิทหติ ภิกฺขุ น สํวิทหติ, วิสงฺเกเตน คจฺฉติ, อาปทาสุ อุมฺมตฺตกสฺส อาทิกมฺมิกสฺสาติ.}{สํวิธานสิกฺขาปทํ นิฏฺฐิตํ สตฺตมํ.}
\csromanmatikaanchor{mtk.76}
\csromantocmarkref{subsection}{8. อริฏฺฐสิกฺขาปท}{mtk.76}
\bookmark[dest={mtk.76},level=2]{8. อริฏฺฐสิกฺขาปท}
\csromanheader{2}{7. \textbf{สปฺปาณกวคฺค 8. อริฏฺฐสิกฺขาปท}}
\proseitem{417}{\csromansymbolfootnote{*}{อิทํ วตฺถุ วิ 4. 68; ม 1. 182 ปิฏฺฐาทีสุปิ อาคตํ.}เตน สมเยน พุทฺโธ ภควา สาวตฺถิยํ วิหรติ เชตวเน อนาถปิณฺฑิกสฺส อาราเม.\csromanspacer{} เตน โข ปน สมเยน อริฏฺฐสฺส นาม ภิกฺขุโน คทฺธพาธิปุพฺพสฺส\footnote{คนฺธพาธิปุพฺพสฺส (สฺยา, ก)} เอวรูปํ ปาปกํ ทิฏฺฐิคตํ อุปฺปนฺนํ โหติ “ตถาหํ ภควตา ธมฺมํ เทสิตํ อาชานามิ, ยถา เย’เม อนฺตรายิกา ธมฺมา วุตฺตา ภควตา, เต ปฏิเสวโต นาลํ อนฺตรายายา”ติ.\csromanspacer{} อสฺโสสุํ โข สมฺพหุลา ภิกฺขู “อริฏฺฐสฺส กิร นาม ภิกฺขุโน คทฺธพาธิปุพฺพสฺส เอวรูปํ ปาปกํ ทิฏฺฐิคตํ อุปฺปนฺนํ ‘ตถาหํ ภควตา ธมฺมํ เทสิตํ อาชานามิ, ยถา เย’เม อนฺตรายิกา ธมฺมา วุตฺตา ภควตา, เต ปฏิเสวโต นาลํ อนฺตรายายา’ติ”.\csromanspacer{} อถ โข เต ภิกฺขู เยน อริฏฺโฐ ภิกฺขุ คทฺธพาธิปุพฺโพ เตนุปสงฺกมิํสุ, อุปสงฺกมิตฺวา อริฏฺฐํ ภิกฺขุํ คทฺธพาธิปุพฺพํ เอตทโวจุํ “สจฺจํ กิร เต อาวุโส อริฏฺฐ เอวรูปํ ปาปกํ ทิฏฺฐิคตํ อุปฺปนฺนํ ‘ตถาหํ ภควตา ธมฺมํ เทสิตํ อาชานามิ, ยถา เย’เม อนฺตรายิกา ธมฺมา วุตฺตา ภควตา, เต ปฏิเสวโต นาลํ อนฺตรายายา’ติ”.\csromanspacer{} เอวํพฺยาโข อหํ อาวุโส ภควตา ธมฺมํ เทสิตํ อาชานามิ, ยถา เย’เม อนฺตรายิกา ธมฺมา วุตฺตา ภควตา, เต ปฏิเสวโต นาลํ อนฺตรายายาติ.}

\prose{มา อาวุโส อริฏฺฐ เอวํ อวจ, มา ภควนฺตํ อพฺภาจิกฺขิ, น หิ สาธุ ภควโต อพฺภกฺขานํ\footnote{อพฺภาจิกฺขนํ (อิติปิ)}, น หิ ภควา เอวํ วเทยฺย, อเนกปริยาเยนา\-วุโส อริฏฺฐ อนฺตรายิกา ธมฺมา อนฺตรายิกา วุตฺตา ภควตา, อลญฺจ ปน เต ปฏิเสวโต อนฺตรายาย.\csromanspacer{} อปฺปสฺสาทา กามา วุตฺตา ภควตา พหุทุกฺขา พหุปายาสา, อาทีนโว เอตฺถ ภิยฺโย.\csromanspacer{} อฏฺฐิกงฺกลูปมา กามา วุตฺตา ภควตา พหุทุกฺขา พหุปายาสา, \csromanfolio{176}อาทีนโว เอตฺถ ภิยฺโย.\csromanspacer{} มํสเปสูปมา กามา วุตฺตา ภควตา ฯเปฯ.\csromanspacer{} ติณุกฺกูปมา กามา วุตฺตา ภควตา.\csromanspacer{} องฺคารกาสูปมา กามา วุตฺตา ภควตา.\csromanspacer{} สุปินกูปมา กามา วุตฺตา ภควตา.\csromanspacer{} ยาจิตกูปมา กามา วุตฺตา ภควตา.\csromanspacer{} รุกฺขผลูปมา กามา วุตฺตาตา.\csromanspacer{} อสิสูนูปมา กามา วุตฺตา ภควตา.\csromanspacer{} สตฺติสูลูปมา กามา วุตฺตา ภควตา.\csromanspacer{} สปฺปสิรูปมา กามา วุตฺตา ภควตา พหุทุกฺขา พหุปายาสา, อาทีนโว เอตฺถ ภิยฺโยติ.}

\prose{เอวํปิ โข อริฏฺโฐ ภิกฺขุ คทฺธพาธิปุพฺโพ เตหิ ภิกฺขูหิ วุจฺจมาโน ตเถว ตํ ปาปกํ ทิฏฺฐิคตํ ถามสา ปรามาสา อภินิวิสฺส โวหรติ “เอวํพฺยาโข อหํ อาวุโส ภควตา ธมฺมํ เทสิตํ อาชานามิ, ยถา เย’เม อนฺตรายิกา ธมฺมา วุตฺตา ภควตา, เต ปฏิเสวโต นาลํ อนฺตรายายา”ติ.\csromanspacer{} ยโต จ โข เต ภิกฺขู นาสกฺขิํสุ อริฏฺฐํ ภิกฺขุํ คทฺธพาธิปุพฺพํ เอตสฺมา ปาปกา ทิฏฺฐิคตา วิเวเจตุํ, อถ โข เต ภิกฺขู เยน ภควา เตนุปสงฺกมิํสุ, อุปสงฺกมิตฺวา ภควโต เอตมตฺถํ อาโรเจสุํ.\csromanspacer{} อถ โข ภควา เอตสฺมิํ นิทาเน เอตสฺมิํ ปกรเณ ภิกฺขุสํฆํ สนฺนิปาตาเปตฺวา อริฏฺฐํ ภิกฺขุํ คทฺธพาธิปุพฺพํ ปฏิปุจฺฉิ “สจฺจํ กิร เต อริฏฺฐ เอวรูปํ ปาปกํ ทิฏฺฐิคตํ อุปฺปนฺนํ ‘ตถาหํ ภควตา ธมฺมํ เทสิตํ อาชานามิ, ยถา เย’เม อนฺตรายิกา ธมฺมา วุตฺตา ภควตา, เต ปฏิเสวโต นาลํ อนฺตรายายา’ติ”.\csromanspacer{} เอวํพฺยาโข อหํ ภนฺเต ภควตา ธมฺมํ เทสิตํ อาชานามิ, ยถา เย’เม อนฺตรายิกา ธมฺมา วุตฺตา ภควตา, เต ปฏิเสวโต นาลํ อนฺตรายายาติ.}

\prose{กสฺส นุ โข นาม ตฺวํ โมฆปุริส มยา เอวํ ธมฺมํ เทสิตํ อาชานาสิ.\csromanspacer{} นนุ มยา โมฆปุริส อเนกปริยาเยน อนฺตรายิกา ธมฺมา อนฺตรายิกา วุตฺตา, อลญฺจ ปน เต ปฏิเสวโต อนฺตรายาย.\csromanspacer{} อปฺปสฺสาทา กามา วุตฺตา มยา พหุทุกฺขา พหุปายาสา, อาทีนโว เอตฺถ ภิยฺโย.\csromanspacer{} อฏฺฐิกงฺกลูปมา กามา วุตฺตา มยา ฯเปฯ.\csromanspacer{} มํสเปสูปมา กามา วุตฺตา มยา.\csromanspacer{} ติณุกฺกูปมา กามา วุตฺตา มยา.\csromanspacer{} องฺคารกาสูปมา กามา วุตฺตา มยา.\csromanspacer{} สุปินกูปมา กามา วุตฺตา มยา.\csromanspacer{} ยาจิตกูปมา กามา วุตฺตา มยา.\csromanspacer{} รุกฺขผลูปมา กามา วุตฺตา มยา.\csromanspacer{} อสิสูนูปมา กามา วุตฺตา มยา.\csromanspacer{} สตฺติสูลูปมา กามา วุตฺตา มยา.\csromanspacer{} สปฺปสิรูปมา \csromanfolio{177}กามา วุตฺตา มยา พหุทุกฺขา พหุปายาสา, อาทีนโว เอตฺถ ภิยฺโย.\csromanspacer{} อถ จ ปน ตฺวํ โมฆปุริส อตฺตนา ทุคฺคหิเตน อเมฺห เจว อพฺภาจิกฺขสิ, อตฺตานญฺจ ขณสิ, พหุญฺจ อปุญฺญํ ปสวสิ, ตญฺหิ เต โมฆปุริส ภวิสฺสติ ทีฆรตฺตํ อหิตาย ทุกฺขาย.\csromanspacer{} เนตํ โมฆปุริส อปฺปสนฺนานํ วา ปสาทาย ฯเปฯ.\csromanspacer{} เอวญฺจ ปน ภิกฺขเว อิมํ สิกฺขาปทํ อุทฺทิเสยฺยาถ\csromandash{}}

\proseitem{418}{“โย ปน ภิกฺขุ เอวํ วเทยฺย ‘ตถาหํ ภควตา ธมฺมํ เทสิตํ อาชานามิ, ยถา เย’เม อนฺตรายิกา ธมฺมา วุตฺตา ภควตา, เต ปฏิเสวโต นาลํ อนฺตรายายา’ติ.\csromanspacer{} โส ภิกฺขุ ภิกฺขูหิ เอวมสฺส วจนีโย ‘มายสฺมา เอวํ อวจ, มา ภควนฺตํ อพฺภาจิกฺขิ, น หิ สาธุ ภควโต อพฺภกฺขานํ, น หิ ภควา เอวํ วเทยฺย, อเนกปริยาเยนา\-วุโส อนฺตรายิกา ธมฺมา อนฺตรายิกา วุตฺตา ภควตา, อลญฺจ ปน เต ปฏิเสวโต อนฺตรายายา’ติ.\csromanspacer{} เอวญฺจ\footnote{เอวญฺจ ปน (ก)} โส ภิกฺขุ ภิกฺขูหิ วุจฺจมาโน ตเถว ปคฺคเณฺหยฺย, โส ภิกฺขุ ภิกฺขูหิ ยาวตติยํ สมนุภาสิตพฺโพ ตสฺส ปฏินิสฺสคฺคาย.\csromanspacer{} ยาวตติยญฺเจ สมนุภาสิยมาโน ตํ ปฏินิสฺสชฺเชยฺย, อิจฺเจตํ กุสลํ.\csromanspacer{} โน เจ ปฏินิสฺสชฺเชยฺย, ปาจิตฺติยนฺ”ติ.}

\proseitem{419}{\textbf{โย ปนา}ติ โย ยาทิโส ฯเปฯ.\csromanspacer{} \textbf{ภิกฺขู}ติ ฯเปฯ อยํ อิมสฺมิํ อตฺเถ อธิปฺเปโต ภิกฺขูติ.}

\prose{\textbf{เอวํ วเทยฺยา}ติ ตถาหํ ภควตา ธมฺมํ เทสิตํ อาชานามิ, ยถา เย’เม อนฺตรายิกา ธมฺมา วุตฺตา ภควตา, เต ปฏิเสวโต นาลํ อนฺตรายายาติ.}

\prose{\textbf{โส ภิกฺขู}ติ โย โส เอวํวาที ภิกฺขุ.}

\prose{ภิกฺขูหีติ อญฺเญหิ ภิกฺขูหิ.\csromanspacer{} เย ปสฺสนฺติ เย สุณนฺติ, เตหิ วตฺตพฺโพ “มายสฺมา เอวํ อวจ, มา ภควนฺตํ อพฺภาจิกฺขิ, น หิ สาธุ ภควโต อพฺภกฺขานํ, น หิ ภควา เอวํ วเทยฺย, อเนกปริยาเยนา\-วุโส อนฺตรายิกา ธมฺมา อนฺตรายิกา วุตฺตา ภควตา, อลญฺจ ปน เต ปฏิเสวโต อนฺตรายายา”ติ.\csromanspacer{} ทุติยมฺปิ วตฺตพฺโพ.\csromanspacer{} ตติยมฺปิ วตฺตพฺโพ.\csromanspacer{} สเจ ปฏินิสฺสชฺชติ, อิจฺเจตํ กุสลํ.\csromanspacer{} โน เจ ปฏินิสฺสชฺชติ, อาปตฺติ ทุกฺกฏสฺส.\csromanspacer{} สุตฺวา น วทนฺติ, อาปตฺติ ทุกฺกฏสฺส.\csromanspacer{} โส ภิกฺขุ \csromanfolio{178}สํฆมชฺฌมฺปิ อากฑฺฒิตฺวา วตฺตพฺโพ “มายสฺมา เอวํ อวจ, มา ภควนฺตํ อพฺภาจิกฺขิ, น หิ สาธุ ภควโต อพฺภกฺขานํ, น หิ ภควา เอวํ วเทยฺย, อเนกปริยาเยนา\-วุโส อนฺตรายิกา ธมฺมา อนฺตรายิกา วุตฺตา ภควตา, อลญฺจ ปน เต ปฏิเสวโต อนฺตรายายา”ติ.\csromanspacer{} ทุติยมฺปิ วตฺตพฺโพ.\csromanspacer{} ตติยมฺปิ วตฺตพฺโพ.\csromanspacer{} สเจ ปฏินิสฺสชฺชติ, อิจฺเจตํ กุสลํ.\csromanspacer{} โน เจ ปฏิสฺสชฺชติ, อาปตฺติ ทุกฺกฏสฺส.\csromanspacer{} โส ภิกฺขุ สมนุภาสิตพฺโพ.\csromanspacer{} เอวญฺจ ปน ภิกฺขเว สมนุภาสิตพฺโพ, พฺยตฺเตน ภิกฺขุนา ปฏิพเลน สํโฆ ญาเปตพฺโพ\csromandash{}}

\proseitem{420}{“สุณาตุ เม ภนฺเต สํโฆ, อิตฺถนฺนามสฺส ภิกฺขุโน เอวรูปํ ปาปกํ ทิฏฺฐิคตํ อุปฺปนฺนํ ‘ตถาหํ ภควตา ธมฺมํ เทสิตํ อาชานามิ ยถา เย’เม อนฺตรายิกา ธมฺมา วุตฺตา ภควตา, เต ปฏิเสวโต นาลํ อนฺตรายายา’ติ.\csromanspacer{} โส ตํ ทิฏฺฐิํ น ปฏินิสฺสชฺชติ, ยทิ สํฆสฺส ปตฺตกลฺลํ, สํโฆ อิตฺถนฺนามํ ภิกฺขุํ สมนุภาเสยฺย ตสฺสา ทิฏฺฐิยา ปฏินิสฺสคฺคาย, เอสา ญตฺติ.}

\prose{สุณาตุ เม ภนฺเต สํโฆ, อิตฺถนฺนามสฺส ภิกฺขุโน เอวรูปํ ปาปกํ ทิฏฺฐิคตํ อุปฺปนฺนํ ‘ตถาหํ ภควตา ธมฺมํ เทสิตํ อาชานามิ, ยถา เย’เม อนฺตรายิกา ธมฺมา วุตฺตา ภควตา, เต ปฏิเสวโต นาลํ อนฺตรายายา’ติ.\csromanspacer{} โส ตํ ทิฏฺฐิํ น ปฏินิสฺสชฺชติ, สํโฆ อิตฺถนฺนามํ ภิกฺขุํ สมนุภาสติ ตสฺสา ทิฏฺฐิยา ปฏินิสฺสคฺคาย, ยสฺสายสฺมโต ขมติ อิตฺถนฺนามสฺส ภิกฺขุโน สมนุภาสนา ตสฺสา ทิฏฺฐิยา ปฏินิสฺสคฺคาย, โส ตุณฺหสฺส, ยสฺส นกฺขมติ, โส ภาเสยฺย.}

\prose{ทุติยมฺปิ เอตมตฺถํ วทามิ ฯเปฯ.\csromanspacer{} ตติยมฺปิ เอตมตฺถํ วทามิ.\csromanspacer{} สุณาตุ เม ภนฺเต สํโฆ, อิตฺถนฺนามสฺส ภิกฺขุโน เอวรูปํ ปาปกํ ทิฏฺฐิคตํ อุปฺปนฺนํ ‘ตถาหํ ภควตา ธมฺมํ เทสิตํ อาชานามิ, ยถา เย’เม อนฺตรายิกา ธมฺมา วุตฺตา ภควตา, เต ปฏิเสวโต นาลํ อนฺตรายายา’ติ.\csromanspacer{} โส ตํ ทิฏฺฐิํ น ปฏินิสฺสชฺชติ, สํโฆ อิตฺถนฺนามํ ภิกฺขุํ สมนุภาสติ ตสฺสา ทิฏฺฐิยา ปฏินิสฺสคฺคาย, ยสฺสายสฺมโต ขมติ อิตฺถนฺนามสฺส ภิกฺขุโน สมนุภาสนา ตสฺสา ทิฏฺฐิยา ปฏินิสฺสคฺคาย, โส ตุณฺหสฺส, ยสฺส นกฺขมติ, โส ภาเสยฺย.}

\prose{\csromanfolio{179}สมนุภฏฺโฐ สํเฆน อิตฺถนฺนาโม ภิกฺขุ ตสฺสา ทิฏฺฐิยา ปฏินิสฺสคฺคาย, ขมติ สํฆสฺส, ตสฺมา ตุณฺหี, เอวเมตํ ธารยามี”ติ.}

\prose{ญตฺติยา ทุกฺกฏํ.\csromanspacer{} ทฺวีหิ กมฺมวาจาหิ ทุกฺกฏา.\csromanspacer{} กมฺมวาจาปริโยสาเน อาปตฺติ ปาจิตฺติยสฺส.}

\proseitem{421}{ธมฺมกมฺเม ธมฺมกมฺมสญฺญี น ปฏินิสฺสชฺชติ, อาปตฺติ ปาจิตฺติยสฺส.\csromanspacer{} ธมฺมกมฺเม เวมติโก น ปฏินิสฺสชฺชติ, อาปตฺติ ปาจิตฺติยสฺส.\csromanspacer{} ธมฺมกมฺเม อธมฺมกมฺมสญฺญี น ปฏินิสฺสชฺชติ, อาปตฺติ ปาจิตฺติยสฺส.}

\prose{อธมฺมกมฺเม ธมฺมกมฺมสญฺญี, อาปตฺติ ทุกฺกฏสฺส.\csromanspacer{} อธมฺมกมฺเม เวมติโก, อาปตฺติ ทุกฺกฏสฺส.\csromanspacer{} อธมฺมกมฺเม อธมฺมกมฺมสญฺญี, อาปตฺติ ทุกฺกฏสฺส.}

\proseitemwithcloserruled{422}{อนาปตฺติ อสมนุภาสนฺตสฺส ปฏินิสฺสชฺชนฺตสฺส อุมฺมตฺตกสฺสาติ.}{อริฏฺฐสิกฺขาปทํ นิฏฺฐิตํ อฏฺฐมํ.}
\csromanmatikaanchor{mtk.77}
\csromantocmarkref{subsection}{9. อุกฺขิตฺตสมฺโภคสิกฺขาปท}{mtk.77}
\bookmark[dest={mtk.77},level=2]{9. อุกฺขิตฺตสมฺโภคสิกฺขาปท}
\csromanheader{2}{7. สปฺปาณกวคฺค 9. อุกฺขิตฺต\-สมฺโภค\-สิกฺขาปท}
\proseitem{423}{เตน สมเยน พุทฺโธ ภควา สาวตฺถิยํ วิหรติ เชตวเน อนาถปิณฺฑิกสฺส อาราเม.\csromanspacer{} เตน โข ปน สมเยน ฉพฺพคฺคิยา ภิกฺขู ชานํ ตถาวาทินา อริฏฺเฐน ภิกฺขุนา\footnote{ภิกฺขุนา คทฺธพาธิปุพฺเพน (?)} อกฏานุธมฺเมน ตํ ทิฏฺฐิํ อปฺปฏินิสฺสฏฺเฐน สทฺธิํ สมฺภุญฺชนฺติปิ สํวสนฺติปิ สหาปิ เสยฺยํ กปฺเปนฺติ.\csromanspacer{} เย เต ภิกฺขู อปฺปิจฺฉา ฯเปฯ เต อุชฺฌายนฺติ ขิยฺยนฺติ วิปาเจนฺติ “กถํ หิ นาม ฉพฺพคฺคิยา ภิกฺขู ชานํ ตถาวาทินา อริฏฺเฐน ภิกฺขุนา1 อกฏานุธมฺเมน ตํ ทิฏฺฐิํ อปฺปฏินิสฺสฏฺเฐน สทฺธิํ สมฺภุญฺชิสฺสนฺติปิ สํวสิสฺสนฺติปิ สหาปิ เสยฺยํ กปฺเปสฺสนฺตี”ติ ฯเปฯ.\csromanspacer{} สจฺจํ กิร ตุเมฺห ภิกฺขเว ชานํ ตถาวาทินา อริฏฺเฐน ภิกฺขุนา1 อกฏานุธมฺเมน ตํ ทิฏฺฐิํ อปฺปฏินิสฺสฏฺเฐน สทฺธิํ สมฺภุญฺชถาปิ สํวสถาปิ สหาปิ เสยฺยํ กปฺเปถาติ.\csromanspacer{} สจฺจํ ภควาติ.\csromanspacer{} วิครหิ พุทฺโธ ภควา ฯเปฯ.\csromanspacer{} กถํ หิ นาม ตุเมฺห โมฆปุริสา ชานํ ตถาวาทินา อริฏฺเฐน ภิกฺขุนา1 อกฏานุธมฺเมน ตํ ทิฏฺฐิํ อปฺปฏินิสฺสฏฺเฐน สทฺธิํ สมฺภุญฺชิสฺสถาปิ สํวสิสฺสถาปิ สหาปิ เสยฺยํ กปฺเปสฺสถ.\csromanspacer{} เนตํ โมฆปุริสา อปฺปสนฺนานํ วา ปสาทาย ฯเปฯ.\csromanspacer{} เอวญฺจ ปน ภิกฺขเว อิมํ สิกฺขาปทํ อุทฺทิเสยฺยาถ\csromandash{}}

\proseitem{424}{\csromanfolio{180}\textbf{“โย ปน ภิกฺขุ ชานํ ตถาวาทินา ภิกฺขุนา อกฏานุธมฺเมน ตํ ทิฏฺฐิํ อปฺปฏินิสฺสฏฺเฐน สทฺธิํ สมฺภุญฺเชยฺย วา สํวเสยฺย วา สห วา เสยฺยํ กปฺเปยฺย, ปาจิตฺติยนฺ”}ติ.}

\proseitem{425}{\textbf{โย ปนา}ติ โย ยาทิโส ฯเปฯ.\csromanspacer{} \textbf{ภิกฺขู}ติ ฯเปฯ อยํ อิมสฺมิํ อตฺเถ อธิปฺเปโต ภิกฺขูติ.}

\prose{\textbf{ชานาติ} นาม สามํ วา ชานาติ, อญฺเญ วา ตสฺส อาโรเจนฺติ, โส วา อาโรเจติ.}

\prose{\textbf{ตถาวาทินา}ติ “ตถาหํ ภควตา ธมฺมํ เทสิตํ อาชานามิ, ยถา เย’เม อนฺตรายิกา ธมฺมา วุตฺตา ภควตา, เต ปฏิเสวโต นาลํ อนฺตรายายา”ติ เอวํ วาทินา.}

\prose{\textbf{อกฏานุธมฺโม} นาม อุกฺขิตฺโต อโนสาริโต.}

\prose{\textbf{ตํ ทิฏฺฐิํ อปฺปฏินิสฺสฏฺเฐน สทฺธินฺ}ติ เอตํ ทิฏฺฐิํ อปฺปฏินิสฺสฏฺเฐน สทฺธิํ.}

\prose{\textbf{สมฺภุญฺเชยฺย วา}ติ สมฺโภโคนาม เทฺว สมฺโภคา อามิสสมฺโภโค จ ธมฺมสมฺโภโค จ.\csromanspacer{} \textbf{อามิสสมฺโภโค} นาม อามิสํ เทติ วา ปฏิคฺคณฺหาติ วา, อาปตฺติ ปาจิตฺติยสฺส.\csromanspacer{} \textbf{ธมฺมสมฺโภโค} นาม อุทฺทิสติ วา อุทฺทิสาเปติ วา, ปเทน อุทฺทิสติ วา อุทฺทิสาเปติ วา, ปเท ปเท อาปตฺติ ปาจิตฺติยสฺส.\csromanspacer{} อกฺขราย อุทฺทิสติ วา อุทฺทิสาเปติ วา, อกฺขรกฺขราย อาปตฺติ ปาจิตฺติยสฺส.}

\prose{\textbf{สํวเสยฺย วา}ติ อุกฺขิตฺตเกน สทฺธิํ อุโปสถํ วา ปวารณํ วา สํฆกมฺมํ วา กโรติ, อาปตฺติ ปาจิตฺติยสฺส.}

\prose{\textbf{สห วา เสยฺยํ กปฺเปยฺยา}ติ เอกจฺฉนฺเน อุกฺขิตฺตเก นิปนฺเน ภิกฺขุ นิปชฺชติ, อาปตฺติ ปาจิตฺติยสฺส.\csromanspacer{} ภิกฺขุ นิปนฺเน อุกฺขิตฺตโก นิปชฺชติ, อาปตฺติ ปาจิตฺติยสฺส.\csromanspacer{} อุโภ วา นิปชฺชนฺติ, อาปตฺติ ปาจิตฺติยสฺส.\csromanspacer{} อุฏฺฐหิตฺวา ปุนปฺปุนํ นิปชฺชนฺติ, อาปตฺติ ปาจิตฺติยสฺส.}

\proseitem{426}{อุกฺขิตฺตเก อุกฺขิตฺตกสญฺญี สมฺภุญฺชติ วา สํวสติ วา สห วา เสยฺยํ กปฺเปติ, อาปตฺติ ปาจิตฺติยสฺส.\csromanspacer{} อุกฺขิตฺตเก เวมติโก สมฺภุญฺชติ วา สํวสติ วา สห วา เสยฺยํ กปฺเปติ, อาปตฺติ ทุกฺกฏสฺส.\csromanspacer{} อุกฺขิตฺตเก อนุกฺขิตฺตสญฺญี สมฺภุญฺชติ วา สํวสติ วา สห วา เสยฺยํ กปฺเปติ \csromanfolio{181}อนาปตฺติ.\csromanspacer{} อนุกฺขิตฺตเก อุกฺขิตฺตกสญฺญี, อาปตฺติ ทุกฺกฏสฺส.\csromanspacer{} อนุกฺขิตฺตเก เวมติโก, อาปตฺติ ทุกฺกฏสฺส.\csromanspacer{} อนุกฺขิตฺตเก อนุกฺขิตฺตกสญฺญี, อนาปตฺติ.}

\proseitemwithcloserruled{427}{อนาปตฺติ อนุกฺขิตฺโตติ ชานาติ, อุกฺขิตฺโต โอสาริโตติ ชานาติ, ตํ ทิฏฺฐิํ ปฏินิสฺสฏฺโฐติ ชานาติ, อุมฺมตฺตกสฺส อาทิกมฺมิกสฺสาติ.}{อุกฺขิตฺต\-สมฺโภค\-สิกฺขาปทํ นิฏฺฐิตํ นวมํ.}
\csromanmatikaanchor{mtk.78}
\csromantocmarkref{subsection}{10. กณฺฏกสิกฺขาปท}{mtk.78}
\bookmark[dest={mtk.78},level=2]{10. กณฺฏกสิกฺขาปท}
\csromanheader{2}{7. \textbf{สปฺปาณกวคฺค} 10. กณฺฏกสิกฺขาปท}
\proseitem{428}{เตน สมเยน พุทฺโธ ภควา สาวตฺถิยํ วิหรติ เชตวเน อนาถปิณฺฑิกสฺส อาราเม.\csromanspacer{} เตน โข ปน สมเยน กณฺฏกสฺส\footnote{กณฺฑกสฺส (สฺยา, ก)}นาม สมณุทฺเทสสฺส เอวรูปํ ปาปกํ ทิฏฺฐิคตํ อุปฺปนฺนํ โหติ “ตถาหํ ภควตา ธมฺมํ เทสิตํ อาชานามิ, ยถา เย’เม อนฺตรายิกา ธมฺมา วุตฺตา ภควตา, เต ปฏิเสวโต นาลํ อนฺตรายายา”ติ.\csromanspacer{} อสฺโสสุํ โข สมฺพหุลา ภิกฺขู “กณฺฏกสฺส นาม กิร สมณุทฺเทสสฺส เอวรูปํ ปาปกํ ทิฏฺฐิคตํ อุปฺปนฺนํ ‘ตถาหํ ภควตา ธมฺมํ เทสิตํ อาชานามิ, ยถา เยเม อนฺตรายิกา ธมฺมา วุตฺตา ภควตา, เต ปฏิเสวโต นาลํ อนฺตรายายา’ติ”.\csromanspacer{} อถ โข เต ภิกฺขู เยน กณฺฏโก สมณุทฺเทโส เตนุปสงฺกมิํสุ, อุปสงฺกมิตฺวา กณฺฏกํ สมณุทฺเทสํ เอตทโวจุํ “สจฺจํ กิร เต อาวุโส กณฺฏก เอวรูปํ ปาปกํ ทิฏฺฐิคตํ อุปฺปนฺนํ ‘ตถาหํ ภควตา ธมฺมํ เทสิตํ อาชานามิ, ยถา เย’เม อนฺตรายิกา ธมฺมา วุตฺตา ภควตา, เต ปฏิเสวโต นาลํ อนฺตรายายา’ติ”.\csromanspacer{} เอวํพฺยาโข อหํ ภนฺเต ภควตา ธมฺมํ เทสิตํ อาชานามิ, ยถา เย’เม อนฺตรายิกา ธมฺมา วุตฺตา ภควตา, เต ปฏิเสวโต นาลํ อนฺตรายายาติ.}

\prose{มา อาวุโส กณฺฏก เอวํ อวจ, มา ภควนฺตํ อพฺภาจิกฺขิ, น หิ สาธุ ภควโต อพฺภกฺขานํ, น หิ ภควา เอวํ วเทยฺย, อเนกปริยาเยนา\-วุโส กณฺฏก อนฺตรายิกา ธมฺมา อนฺตรายิกา วุตฺตา ภควตา, \csromanfolio{182}อลญฺจ ปน เต ปฏิเสวโต อนฺตรายาย.\csromanspacer{} อปฺปสฺสาทา กามา วุตฺตา ภควตา พหุทุกฺขา พหุปายาสา, อาทีนโว เอตฺถ ภิยฺโย ฯเปฯ.\csromanspacer{} เอวํปิ โข กณฺฏโก สมณุทฺเทโส เตหิ ภิกฺขูหิ วุจฺจมาโน ตเถว ตํ ปาปกํ ทิฏฺฐิคตํ ถามสา ปรามาสา อภินิวิสฺส โวหรติ “เอวํพฺยาโข อหํ ภนฺเต ภควตา ธมฺมํ เทสิตํ อาชานามิ, ยถา เย’เม อนฺตรายิกา ธมฺมา วุตฺตา ภควตา, เต ปฏิเสวโต นาลํ อนฺตรายายา”ติ.}

\prose{ยโต จ โข เต ภิกฺขู นาสกฺขิํสุ กณฺฏกํ สมณุทฺเทสํ เอตสฺมา ปาปกา ทิฏฺฐิคตา วิเวเจตุํ.\csromanspacer{} อถ โข เต ภิกฺขู เยน ภควา เตนุปสงฺกมิํสุ, อุปสงฺกมิตฺวา ภควโต เอตมตฺถํ อาโรเจสุํ.\csromanspacer{} อถ โข ภควา เอตสฺมิํ นิทาเน เอตสฺมิํ ปกรเณ ภิกฺขุสํฆํ สนฺนิปาตาเปตฺวา กณฺฏกํ สมณุทฺเทสํ ปฏิปุจฺฉิ “สจฺจํ กิร เต กณฺฏก เอวรูปํ ปาปกํ ทิฏฺฐิคตํ อุปฺปนฺนํ ‘ตถาหํ ภควตา ธมฺมํ เทสิตํ อาชานามิ, ยถา เย’เม อนฺตรายิกา ธมฺมา วุตฺตา ภควตา, เต ปฏิเสวโต นาลํ อนฺตรายายา’ติ”.\csromanspacer{} เอวํพฺยาโข อหํ ภนฺเต ภควตา ธมฺมํ เทสิตํ อาชานามิ, ยถา เย’เม อนฺตรายิกา ธมฺมา วุตฺตา ภควตา, เต ปฏิเสวโต นาลํ อนฺตรายายาติ.}

\prose{กสฺส นุ โข นาม ตฺวํ โมฆปุริส มยา เอวํ ธมฺมํ เทสิตํ อาชานาสิ.\csromanspacer{} นนุ มยา โมฆปุริส อเนกปริยาเยน อนฺตรายิกา ธมฺมา อนฺตรายิกา วุตฺตา, อลญฺจ ปน เต ปฏิเสวโต อนฺตรายาย.\csromanspacer{} อปฺปสฺสาทา กามา วุตฺตา มยา พหุทุกฺขา พหุปายาสา, อาทีนโว เอตฺถ ภิยฺโย.\csromanspacer{} อฏฺฐิกงฺกลูปมา กามา วุตฺตา มยา ฯเปฯ สปฺปสิรูปมา กามา วุตฺตา มยา พหุทุกฺขา พหุปายาสา, อาทีนโว เอตฺถ ภิยฺโย.\csromanspacer{} อถ จ ปน ตฺวํ โมฆปุริส อตฺตนา ทุคฺคหิเตน อเมฺห เจว อพฺภาจิกฺขสิ, อตฺตานญฺจ ขณสิ, พหุํ จ อปุญฺญํ ปสวสิ, ตญฺหิ เต โมฆปุริส ภวิสฺสติ ทีฆรตฺตํ อหิตาย ทุกฺขาย.\csromanspacer{} เนตํ โมฆปุริส อปฺปสนฺนานํ วา ปสาทาย ฯเปฯ ปสนฺนานญฺจ เอกจฺจานํ อญฺญถตฺตายาติ วิครหิตฺวา ฯเปฯ ธมฺมิํ กถํ กตฺวา ภิกฺขู อามนฺเตสิ.\csromanspacer{} เตน หิ ภิกฺขเว สํโฆ กณฺฏกํ สมณุทฺเทสํ นาเสตุ.\csromanspacer{} เอวญฺจ ปน ภิกฺขเว นาเสตพฺโพ “อชฺชตคฺเค เต อาวุโส กณฺฏก น เจว โส ภควา สตฺถา อปทิสิตพฺโพ, ยมฺปิ ญฺเญ สมณุทฺเทสา ลภนฺติ ภิกฺขูหิ สทฺธิํ ทิรตฺตติรตฺตํ สหเสยฺยํ, \csromanfolio{183}สาปิ เต นตฺถิ, จร ปิเร วินสฺสา”ติ.\csromanspacer{} อถ โข สํโฆ กณฺฏกํ \textbf{สมณุทฺเทสํ นาเสสิ.}}

\prose{เตน โข ปน สมเยน ฉพฺพคฺคิยา ภิกฺขู ชานํ ตถานาสิตํ กณฺฏกํ สมณุทฺเทสํ อุปลาเปนฺติปิ อุปฏฺฐาเปนฺติปิ สมฺภุญฺชนฺติปิ สหาปิ เสยฺยํ กปฺเปนฺติ.\csromanspacer{} เย เต ภิกฺขู อปฺปิจฺฉา ฯเปฯ เต อุชฺฌายนฺติ ขิยฺยนฺติ วิปาเจนฺติ “กถํ หิ นาม ฉพฺพคฺคิยา ภิกฺขู ชานํ ตถานาสิตํ กณฺฏกํ สมณุทฺเทสํ อุปลาเปสฺสนฺติปิ อุปฏฺฐาเปสฺสนฺติปิ สมฺภุญฺชิสฺสนฺติปิ สหาปิ เสยฺยํ กปฺเปสฺสนฺตี”ติ ฯเปฯ.\csromanspacer{} สจฺจํ กิร ตุเมฺห ภิกฺขเว ชานํ ตถานาสิตํ กณฺฏกํ สมณุทฺเทสํ อุปลาเปถาปิ อุปฏฺฐาเปถาปิ สมฺภุญฺชถาปิ สหาปิ เสยฺยํ กปฺเปถาติ.\csromanspacer{} สจฺจํ ภควาติ.\csromanspacer{} วิครหิ พุทฺโธ ภควา ฯเปฯ.\csromanspacer{} กถํ หิ นาม ตุเมฺห โมฆปุริสา ชานํ ตถานาสิตํ กณฺฏกํ สมณุทฺเทสํ อุปลาเปสฺสถาปิ อุปฏฺฐาเปสฺสถาปิ สมฺภุญฺชิสฺสถาปิ สหาปิ เสยฺยํ กปฺเปสฺสถ.\csromanspacer{} เนตํ โมฆปุริสา อปฺปสนฺนานํ วา ปสาทาย ฯเปฯ.\csromanspacer{} เอวญฺจ ปน ภิกฺขเว อิมํ สิกฺขาปทํ อุทฺทิเสยฺยาถ\csromandash{}}

\proseitem{429}{“สมณุทฺเทโสปิ เจ เอวํ วเทยฺย ‘ตถาหํ ภควตา ธมฺมํ เทสิตํ อาชานามิ, ยถา เย’เม อนฺตรายิกา ธมฺมา วุตฺตา ภควตา, เต ปฏิเสวโต นาลํ อนฺตรายายา’ติ.\csromanspacer{} โส สมณุทฺเทโส ภิกฺขูหิ เอวมสฺส วจนีโย ‘มาวุโส สมณุทฺเทส เอวํ อวจ, มา ภควนฺตํ อพฺภาจิกฺขิ, น หิ สาธุ ภควโต อพฺภกฺขานํ, น หิ ภควา เอวํ วเทยฺย, อเนกปริยาเยนา\-วุโส สมณุทฺเทส อนฺตรายิกา ธมฺมา อนฺตรายิกา วุตฺตา ภควตา, อลญฺจ ปน เต ปฏิเสวโต อนฺตรายายา’ติ.\csromanspacer{} เอวญฺจ\footnote{เอวญฺจ ปน (ก)} โส สมณุทฺเทโส ภิกฺขูหิ วุจฺจมาโน ตเถว ปคฺคเณฺหยฺย, โส สมณุทฺเทโส ภิกฺขูหิ เอวมสฺส วจนีโย ‘อชฺชตคฺเค เต อาวุโส สมณุทฺเทส น เจว โส ภควา สตฺถา อปทิสิตพฺโพ, ยมฺปิ จญฺเญ สมณุทฺเทสา ลภนฺติ ภิกฺขูหิ สทฺธิํ ทิรตฺตติรตฺตํ สหเสยฺยํ, สาปิ เต นตฺถิ, จร ปิเร วินสฺสา’ติ.\csromanspacer{} โย ปน ภิกฺขุ ชานํ ตถานาสิตํ สมณุทฺเทสํ อุปลาเปยฺย วา อุปฏฺฐาเปยฺย วา สมฺภุญฺเชยฺย วา สห วา เสยฺยํ กปฺเปยฺย, ปาจิตฺติยนฺ”ติ.}

\proseitem{430}{\textbf{สมณุทฺเทโส} นาม สามเณโร วุจฺจติ.}

\prose{\csromanfolio{184}\textbf{เอวํ วเทยฺยา}ติ ตถาหํ ภควตา ธมฺมํ เทสิตํ อาชานามิ, ยถา เย’เม อนฺตรายิกา ธมฺมา วุตฺตา ภควตา, เต ปฏิเสวโต นาลํ อนฺตรายายาติ.}

\prose{\textbf{โส สมณุทฺเทโส}ติ โย โส เอวํวาที สมณุทฺเทโส.}

\prose{ภิกฺขูหีติ อญฺเญหิ ภิกฺขูหิ, เย ปสฺสนฺติ เย สุณนฺติ เตหิ วตฺตพฺโพ “มา อาวุโส สมณุทฺเทส เอวํ อวจ, มา ภควนฺตํ อพฺภาจิกฺขิ, น หิ สาธุ ภควโต อพฺภกฺขานํ, น หิ ภควา เอวํ วเทยฺย, อเนกปริยาเยนา\-วุโส สมณุทฺเทส อนฺตรายิกา ธมฺมา อนฺตรายิกา วุตฺตา ภควตา, อลญฺจ ปน เต ปฏิเสวโต อนฺตรายายา”ติ.\csromanspacer{} ทุติยมฺปิ วตฺตพฺโพ.\csromanspacer{} ตติยมฺปิ วตฺตพฺโพ.\csromanspacer{} สเจ ปฏินิสฺสชฺชติ, อิจฺเจตํ กุสลํ.\csromanspacer{} โน เจ ปฏินิสฺสชฺชติ, โส สมณุทฺเทโส ภิกฺขูหิ เอวมสฺส วจนีโย “อชฺชตคฺเค เต อาวุโส สมณุทฺเทส น เจว โส ภควา สตฺถา อปทิสิตพฺโพ, ยมฺปิ จญฺเญ สมณุทฺเทสา ลภนฺติ ภิกฺขูหิ สทฺธิํ ทิรตฺตติรตฺตํ สหเสยฺยํ, สาปิ เต นตฺถิ, จร ปิเร วินสฺสา”ติ.}

\prose{\textbf{โย ปนา}ติ โย ยาทิโส ฯเปฯ.\csromanspacer{} \textbf{ภิกฺขู}ติ ฯเปฯ อยํ อิมสฺมิํ อตฺเถ อธิปฺเปโต ภิกฺขูติ.}

\prose{\textbf{ชานาติ} นาม สามํ วา ชานาติ, อญฺเญ วา ตสฺส อาโรเจนฺติ โส วา อาโรเจติ.}

\prose{\textbf{ตถานาสิตนฺ}ติ เอวํ นาสิตํ.}

\prose{\textbf{สมณุทฺเทโส} นาม สามเณโร วุจฺจติ.}

\prose{\textbf{อุปลาเปยฺย} วาติ ตสฺส ปตฺตํ \textbf{วา} จีวรํ วา อุทฺเทสํ วา ปริปุจฺฉํ วา ทสฺสามีติ อุปลาเปติ, อาปตฺติ ปาจิตฺติยสฺส.}

\prose{\textbf{อุปฏฺฐาเปยฺย วา}ติ ตสฺส จุณฺณํ วา มตฺติกํ วา ทนฺตกฏฺฐํ วา มุโขทกํ วา สาทิยติ, อาปตฺติ ปาจิตฺติยสฺส.}

\prose{\textbf{สมฺภุญฺเชยฺย วา}ติ สมฺโภโค นาม เทฺว สมฺโภคา อามิสสมฺโภโค จ ธมฺมสมฺโภโค จ.\csromanspacer{} \textbf{อามิสสมฺโภโค} นาม อามิสํ เทติ วา ปฏิคฺคณฺหาติ วา อาปตฺติ ปาจิตฺติยสฺส.\csromanspacer{} \textbf{ธมฺมสมฺโภโค} นาม อุทฺทิสติ วา อุทฺทิสาเปติ วา, ปเทน อุทฺทิสติ วา อุทฺทิสาเปติ วา, ปเท ปเท \csromanfolio{185}อาปตฺติ ปาจิตฺติยสฺส.\csromanspacer{} อกฺขราย อุทฺทิสติ วา อุทฺทิสาเปติ วา, อกฺขรกฺขราย อาปตฺติ ปาจิตฺติยสฺส.}

\prose{\textbf{สห วา เสยฺยํ กปฺเปยฺยา}ติ เอกจฺฉนฺเน นาสิตเก สมณุทฺเทเส นิปนฺเน ภิกฺขุ นิปชฺชติ, อาปตฺติ ปาจิตฺติยสฺส.\csromanspacer{} ภิกฺขุ นิปนฺเน นาสิตโก สมณุทฺเทโส นิปชฺชติ, อาปตฺติ ปาจิตฺติยสฺส.\csromanspacer{} อุโภ วา นิปชฺชนฺติ, อาปตฺติ ปาจิตฺติยสฺส.\csromanspacer{} อุฏฺฐหิตฺวา ปุนปฺปุนํ นิปชฺชนฺติ, อาปตฺติ ปาจิตฺติยสฺส.}

\proseitem{431}{นาสิตเก นาสิตกสญฺญี อุปลาเปติ วา อุปฏฺฐาเปติ วา สมฺภุญฺชติ วา สห วา เสยฺยํ กปฺเปติ, อาปตฺติ ปาจิตฺติยสฺส.\csromanspacer{} นาสิตเก เวมติโก อุปลาเปติ วา อุปฏฺฐาเปติ วา สมฺภุญฺชติ วา สห วา เสยฺยํ กปฺเปติ, อาปตฺติ ทุกฺกฏสฺส.\csromanspacer{} นาสิตเก อนาสิตกสญฺญี อุปลาเปติ วา อุปฏฺฐาเปติ วา สมฺภุญฺชติ วา สห วา เสยฺยํ กปฺเปติ, อนาปตฺติ.\csromanspacer{} อนาสิตเก นาสิตกสญฺญี, อาปตฺติ ทุกฺกฏสฺส.\csromanspacer{} อนาสิตเก เวมติโก, อาปตฺติ ทุกฺกฏสฺส.\csromanspacer{} อนาสิตเก อนาสิตกสญฺญี, อนาปตฺติ.}

\proseitemwithcloser{432}{อนาปตฺติ อนาสิตโกติ ชานาติ, ตํ ทิฏฺฐิํ ปฏินิสฺสฏฺโฐติ ชานาติ, อุมฺมตฺตกสฺส อาทิกมฺมิกสฺสาติ.}{กณฺฏกสิกฺขาปทํ นิฏฺฐิตํ ทสมํ.}
\vspace{\tipitakaparskip}
\csromancenter{สปฺปาณกวคฺโค สตฺตโม.}
\csromancenter{ตสฺสุทฺทานํ}
\vspace{0.5\baselineskip}
\setlength{\csromangathapagewidth}{0pt}
\csromangathameasuregroup{สญฺจิจฺจวธสปฺปาณํ, \\ อูนวีสติ สตฺถญฺจ,}{\csromangathabat{สญฺจิจฺจวธสปฺปาณํ,}{อุกฺโกฏํ ทุฏฺฐุลฺลฉาทนํ.} \\ \csromangathabat{อูนวีสติ สตฺถญฺจ,}{สํวิธานํ อริฏฺฐกํ.} \\ อุกฺขิตฺตํ กณฺฏกญฺเจว, ทส สิกฺขาปทา อิเมติ.}
\csromangathasetleft{สญฺจิจฺจวธสปฺปาณํ, \\ อูนวีสติ สตฺถญฺจ,}
\csromangathagroup{\csromangathastanza{\csromangathabat{สญฺจิจฺจวธสปฺปาณํ,}{อุกฺโกฏํ ทุฏฺฐุลฺลฉาทนํ.} \\ \csromangathabat{อูนวีสติ สตฺถญฺจ,}{สํวิธานํ อริฏฺฐกํ.} \\ อุกฺขิตฺตํ กณฺฏกญฺเจว, ทส สิกฺขาปทา อิเมติ.}}

\csromanmatikaanchor{mtk.79}
\csromanmatikaanchor{mtk.80}
\csromantocmarknopage{section}{8. สหธมฺมิกวคฺค}{mtk.79}
\bookmark[dest={mtk.79},level=1]{8. สหธมฺมิกวคฺค}
\csromantocmarkref{subsection}{1. สหธมฺมิกสิกฺขาปท}{mtk.80}
\bookmark[dest={mtk.80},level=2]{1. สหธมฺมิกสิกฺขาปท}
\setcsromanheadh{8. สหธมฺมิกวคฺค}
\csromanheader{2}{8. \textbf{สหธมฺมิกวคฺค 1. สหธมฺมิกสิกฺขาปท}}
\proseitem{433}{เตน สมเยน พุทฺโธ ภควา โกสมฺพิยํ วิหรติ โฆสิตาราเม.\csromanspacer{} เตน โข ปน สมเยน อายสฺมา ฉนฺโน \csromanfolio{186}อนาจารํ อาจรติ.\csromanspacer{} ภิกฺขู เอวมาหํสุ “มาวุโส ฉนฺน เอวรูปํ อกาสิ, เนตํ กปฺปตี”ติ.\csromanspacer{} โส เอวํ วเทติ “น ตาวาหํ อาวุโส เอตสฺมิํ สิกฺขาปเท สิกฺขิสฺสามิ, ยาว น อญฺญํ ภิกฺขุํ พฺยตฺตํ วินยธรํ ปริปุจฺฉามี”ติ.\csromanspacer{} เย เต ภิกฺขู อปฺปิจฺฉา ฯเปฯ เต อุชฺฌายนฺติ ขิยฺยนฺติ วิปาเจนฺติ “กถํ หิ นาม อายสฺมา ฉนฺโน ภิกฺขูหิ สหธมฺมิกํ วุจฺจมาโน เอวํ วกฺขติ น ตาวาหํ อาวุโส เอตสฺมิํ สิกฺขาปเท สิกฺขิสฺสามิ, ยาว น อญฺญํ ภิกฺขุํ พฺยตฺตํ วินยธรํ ปริปุจฺฉามี”ติ ฯเปฯ.\csromanspacer{} สจฺจํ กิร ตฺวํ ฉนฺน ภิกฺขูหิ สหธมฺมิกํ วุจฺจมาโน เอวํ วเทสิ “น ตาวาหํ อาวุโส เอตสฺมิํ สิกฺขาปเท สิกฺขิสฺสามิ, ยาว น อญฺญํ ภิกฺขุํ พฺยตฺตํ วินยธรํ ปริปุจฺฉามี”ติ.\csromanspacer{} สจฺจํ ภควาติ.\csromanspacer{} วิครหิ พุทฺโธ ภควา ฯเปฯ.\csromanspacer{} กถํ หิ นาม ตฺวํ โมฆปุริส ภิกฺขูหิ สหธมฺมิกํ วุจฺจมาโน เอวํ วกฺขสิ “น ตาวาหํ อาวุโส เอตสฺมิํ สิกฺขาปเท สิกฺขิสฺสามิ, ยาว น อญฺญํ ภิกฺขุํ พฺยตฺตํ วินยธรํ ปริปุจฺฉามี”ติ.\csromanspacer{} เนตํ โมฆปุริส อปฺปสนฺนานํ วา ปสาทาย ฯเปฯ.\csromanspacer{} เอวญฺจ ปน ภิกฺขเว อิมํ \textbf{สิกฺขาปทํ อุทฺทิเสยฺยาถ\csromandash{}}}

\proseitem{434}{\textbf{“โย ปน ภิกฺขุ ภิกฺขูหิ สหธมฺมิกํ วุจฺจมาโน เอวํ วเทยฺย ‘น ตาวาหํ อาวุโส เอตสฺมิํ สิกฺขาปเท สิกฺขิสฺสามิ, ยาว น อญฺญํ ภิกฺขุํ พฺยตฺตํ วินยธรํ ปริปุจฺฉามี’ติ, ปาจิตฺติยํ.}\csromanspacer{} สิกฺขมาเนน ภิกฺขเว ภิกฺขุนา อญฺญาตพฺพํ ปริปุจฺฉิตพฺพํ ปริปญฺหิตพฺพํ, อยํ ตตฺถ สามีจี”ติ.}

\proseitem{435}{\textbf{โย ปนา}ติ โย ยาทิโส ฯเปฯ.\csromanspacer{} \textbf{ภิกฺขู}ติ ฯเปฯ อยํ อิมสฺมิํ อตฺเถ อธิปฺเปโต ภิกฺขูติ.}

\prose{\textbf{ภิกฺขูหี}ติ อญฺเญหิ ภิกฺขูหิ.}

\prose{\textbf{สหธมฺมิกํ} นาม ยํ ภควตา ปญฺญตฺตํ สิกฺขาปทํ, เอตํ สหธมฺมิกํ นาม.\csromanspacer{} เตน วุจฺจมาโน เอวํ วเทติ\footnote{วเทยฺย (ก)} น ตาวาหํ อาวุโส เอตสฺมิํ สิกฺขาปเท สิกฺขิสฺสามิ, ยาว น อญฺญํ ภิกฺขุํ พฺยตฺตํ วินยธรํ ปริปุจฺฉามีติ\footnote{อิมานิ ปทานิ สฺยาโปตฺถเก น ทิสฺสนฺติ.} ปณฺฑิตํ พฺยตฺตํ2 เมธาวิํ พหุสฺสุตํ ธมฺมกถิกํ ปริปุจฺฉามีติ ภณติ, อาปตฺติ ปาจิตฺติยสฺส.}

\proseitem{436}{\csromanfolio{187}อุปสมฺปนฺเน อุปสมฺปนฺนสญฺญี เอวํ วเทติ, อาปตฺติ ปาจิตฺติยสฺส.\csromanspacer{} อุปสมฺปนฺเน เวมติโก เอวํ วเทติ, อาปตฺติ ปาจิตฺติยสฺส.\csromanspacer{} อุปสมฺปนฺเน อนุปสมฺปนฺนสญฺญี เอวํ วเทติ, อาปตฺติ ปาจิตฺติยสฺส.}

\prose{อปญฺญตฺเตน วุจฺจมาโน “อิทํ น สลฺเลขาย น ธุตตฺถาย น ปาสาทิกตาย น อปจยาย น วีริยารมฺภาย สํวตฺตตี”ติ เอวํ วเทติ, “น ตาวาหํ อาวุโส เอตสฺมิํ สิกฺขาปเท สิกฺขิสฺสามิ, ยาว น อญฺญํ ภิกฺขุํ พฺยตฺตํ วินยธรํ ปณฺฑิตํ เมธาวิํ พหุสฺสุตํ ธมฺมกถิกํ ปริปุจฺฉามี”ติ ภณติ, อาปตฺติ ทุกฺกฏสฺส.}

\prose{อนุปสมฺปนฺเนน ปญฺญตฺเตน วา อปญฺญตฺเตน วา วุจฺจมาโน “อิทํ น สลฺเลขาย น ธุตตฺถาย น ปาสาทิกตาย น อปจยาย น วีริยารมฺภาย สํวตฺตตี”ติ เอวํ วเทติ, “น ตาวาหํ อาวุโส เอตสฺมิํ สิกฺขาปเท สิกฺขิสฺสามิ, ยาว น อญฺญํ ภิกฺขุํ พฺยตฺตํ วินยธรํ ปณฺฑิตํ เมธาวิํ พหุสฺสุตํ ธมฺมกถิกํ ปริปุจฺฉามี”ติ ภณติ, อาปตฺติ ทุกฺกฏสฺส.\csromanspacer{} อนุปสมฺปนฺเน อุปสมฺปนฺนสญฺญี, อาปตฺติ ทุกฺกฏสฺส.\csromanspacer{} อนุปสมฺปนฺเน เวมติโก, อาปตฺติ ทุกฺกฏสฺส.\csromanspacer{} อนุปสมฺปนฺเน อนุปสมฺปนฺนสญฺญี, อาปตฺติ ทุกฺกฏสฺส.}

\prose{\textbf{สิกฺขมาเนนา}ติ สิกฺขิตุกาเมน.}

\prose{\textbf{อญฺญาตพฺพนฺ}ติ ชานิตพฺพํ.}

\prose{\textbf{ปริปุจฺฉิตพฺพนฺ}ติ อิทํ ภนฺเต กถํ อิมสฺส วา กฺวตฺโถติ.}

\prose{\textbf{ปริปญฺหิตพฺพนฺ}ติ จินฺเตตพฺพํ ตุลยิตพฺพํ.}

\prose{\textbf{อยํ ตตฺถ สามีจี}ติ อยํ ตตฺถ อนุธมฺมตา.}

\proseitemwithcloserruled{437}{อนาปตฺติ “ชานิสฺสามิ สิกฺขิสฺสามี”ติ ภณติ, อุมฺมตฺตกสฺส อาทิกมฺมิกสฺสาติ.}{สหธมฺมิกสิกฺขาปทํ นิฏฺฐิตํ ปฐมํ.}
\csromanmatikaanchor{mtk.81}
\csromantocmarkref{subsection}{2. วิเลขนสิกฺขาปท}{mtk.81}
\bookmark[dest={mtk.81},level=2]{2. วิเลขนสิกฺขาปท}
\csromanheader{2}{8. \textbf{สหธมฺมิกวคฺค 2. วิเลขนสิกฺขาปท}}
\proseitem{438}{เตน สมเยน พุทฺโธ ภควา สาวตฺถิยํ วิหรติ เชตวเน อนาถปิณฺฑิกสฺส อาราเม.\csromansymbolfootnote{*}{อิทํ วตฺถุ วิ 4. 319 ปิฏฺเฐปิ อาคตํ.}เตน โข ปน สมเยน ภควา ภิกฺขูนํ อเนกปริยาเยน วินยกถํ กเถติ, วินยสฺส วณฺณํ ภาสติ, \csromanfolio{188}วินยปริยตฺติยา วณฺณํ ภาสติ, อาทิสฺส อาทิสฺส อายสฺมโต อุปาลิสฺส วณฺณํ ภาสติ.\csromanspacer{} ภิกฺขูนํ เอตทโหสิ\footnote{ภิกฺขู ภควา (สฺยา, ก)} “ภควา โข อเนกปริยาเยน วินยกถํ กเถติ, วินยสฺส วณฺณํ ภาสติ, วินยปริยตฺติยา วณฺณํ ภาสติ, อาทิสฺส อาทิสฺส อายสฺมโต อุปาลิสฺส วณฺณํ ภาสติ, หนฺท มยํ อาวุโส อายสฺมโต อุปาลิสฺส สนฺติเก วินยํ ปริยาปุณามา”ติ, เต’ธ พหู ภิกฺขู เถรา จ นวา จ มชฺฌิมา จ อายสฺมโต อุปาลิสฺส สนฺติเก วินยํ ปริยาปุณนฺติ.}

\prose{อถ โข ฉพฺพคฺคิยานํ ภิกฺขูนํ เอตทโหสิ “เอตรหิ โข อาวุโส พหู ภิกฺขู เถรา จา นวา จ มชฺฌิมา จ อายสฺมโต อุปาลิสฺส สนฺติเก วินยํ ปริยาปุณนฺติ, สเจ อิเม วินเย ปกตญฺญุโน ภวิสฺสนฺติ, อเมฺห เยนิจฺฉกํ ยทิจฺฉกํ ยาวทิจฺฉกํ อากฑฺฒิสฺสนฺติ ปริกฑฺฒิสฺสนฺติ, หนฺท มยํ อาวุโส วินยํ วิวณฺเณมา”ติ.\csromanspacer{} อถ โข ฉพฺพคฺคิยา ภิกฺขู ภิกฺขู อุปสงฺกมิตฺวา เอวํ วทนฺติ “กิํ ปนิเมหิ ขุทฺทานุขุทฺทเกหิ สิกฺขาปเทหิ อุทฺทิฏฺเฐหิ ยาวเทว กุกฺกุจฺจาย วิเหสาย วิเลขาย สํวตฺตนฺตี”ติ.\csromanspacer{} เย เต ภิกฺขู อปฺปิจฺฉา ฯเปฯ เต อุชฺฌายนฺติ ขิยฺยนฺติ วิปาเจนฺติ “กถํ หิ นาม ฉพฺพคฺคิยา ภิกฺขู วินยํ วิวณฺเณสฺสนฺตี”ติ ฯเปฯ.\csromanspacer{} สจฺจํ กิร ตุเมฺห ภิกฺขเว วินยํ วิวณฺเณถาติ.\csromanspacer{} สจฺจํ ภควาติ.\csromanspacer{} วิครหิ พุทฺโธ ภควา ฯเปฯ.\csromanspacer{} กถํ หิ นาม ตุเมฺห โมฆปุริสา วินยํ วิวณฺเณสฺสถ.\csromanspacer{} เนตํ โมฆปุริสา อปฺปสนฺนานํ วา ปสาทาย ฯเปฯ.\csromanspacer{} เอวญฺจ ปน ภิกฺขเว อิมํ สิกฺขาปทํ อุทฺทิเสยฺยาถ\csromandash{}}

\proseitem{439}{\textbf{“โย ปน ภิกฺขุ ปาติโมกฺเข อุทฺทิสฺสมาเน เอวํ วเทยฺย ‘กิํ ปนิเมหิ ขุทฺทานุขุทฺทเกหิ สิกฺขาปเทหิ อุทฺทิฏฺเฐหิ, ยาวเทว กุกฺกุจฺจาย วิเหสาย วิเลขาย สํวตฺตนฺตี’ติ, สิกฺขาปทวิวณฺณเก ปาจิตฺติยนฺ”}ติ.}

\proseitem{440}{\textbf{โย ปนา}ติ โย ยาทิโส ฯเปฯ.\csromanspacer{} \textbf{ภิกฺขู}ติ ฯเปฯ อยํ อิมสฺมิํ อตฺเถ อธิปฺเปโต ภิกฺขูติ.}

\prose{\textbf{ปาติโมกฺเข อุทฺทิสฺสมาเน}ติ อุทฺทิสนฺเต วา อุทฺทิสาเปนฺเต วา สชฺฌายํ วา กาโรนฺเต.}

\prose{\csromanfolio{189}\textbf{เอวํ วเทยฺยา}ติ กิํ ปนิเมหิ ขุทฺทานุขุทฺทเกหิ สิกฺขาปเทหิ อุทฺทิฏฺเฐหิ, ยาวเทว กุกฺกุจฺจาย วิเหสาย วิเลขาย สํวตฺตนฺตีติ. “เย อิมํ ปริยาปุณนฺติ, เตสํ กุกฺกุจฺจํ โหติ, วิเหสา โหติ, วิเลขา โหติ.\csromanspacer{} เย อิมํ น ปริยาปุณนฺติ, เตสํ กุกฺกุจฺจํ น โหติ, วิเหสา น โหติ, วิเลขา น โหติ.\csromanspacer{} อนุทฺทิฏฺฐํ อิทํ วรํ, อนุคฺคหิตํ อิทํ วรํ, อปริยาปุฏํ อิทํ วรํ, อธาริตํ อิทํ วรํ, วินโย วา อนฺตรธายตุ, อิเม วา ภิกฺขู อปกตญฺญุโน โหนฺตู”ติ อุปสมฺปนฺนสฺส วินยํ วิวณฺเณติ, อาปตฺติ}

\proseitem{441}{อุปสมฺปนฺเน อุปสมฺปนฺนสญฺญี วินยํ วิวณฺเณติ, อาปตฺติ ปาจิตฺติยสฺส.\csromanspacer{} อุปสมฺปนฺเน เวมติโก วินยํ วิวณฺเณติ, อาปตฺติ ปาจิตฺติยสฺส.\csromanspacer{} อุปสมฺปนฺเน อนุปสมฺปนฺนสญฺญี วินยํ วิวณฺเณติ, อาปตฺติ ปาจิตฺติยสฺส.}

\prose{อญฺญํ ธมฺมํ วิวณฺเณติ, อาปตฺติ ทุกฺกฏสฺส.\csromanspacer{} อนุปสมฺปนฺนสฺส วินยํ วา อญฺญํ วา ธมฺมํ วิวณฺเณติ, อาปตฺติ ทุกฺกฏสฺส.\csromanspacer{} อนุปสมฺปนฺเน อุปสมฺปนฺนสญฺญี, อาปตฺติ ทุกฺกฏสฺส.\csromanspacer{} อนุปสมฺปนฺเน เวมติโก, อาปตฺติ ทุกฺกฏสฺส.\csromanspacer{} อนุปสมฺปนฺเน อนุปสมฺปนฺนสญฺญี, อาปตฺติ ทุกฺกฏสฺส.}

\proseitemwithcloserruled{442}{อนาปตฺติ น วิวณฺเณตุกาโม “อิงฺฆ ตฺวํ สุตฺตนฺเต วา คาถาโย วา อภิธมฺมํ วา ปริยาปุณสฺสุ, ปจฺฉา วินยํ ปริยาปุณิสฺสสี”ติ ภณติ, อุมฺมตฺตกสฺส อาทิกมฺมิกสฺสาติ.}{วิเลขนสิกฺขาปทํ นิฏฺฐิตํ ทุติยํ.}
\csromanmatikaanchor{mtk.82}
\csromantocmarkref{subsection}{3. โมหนสิกฺขาปท}{mtk.82}
\bookmark[dest={mtk.82},level=2]{3. โมหนสิกฺขาปท}
\csromanheader{2}{8. \textbf{สหธมฺมิกวคฺค 3. โมหนสิกฺขาปท}}
\proseitem{443}{เตน สมเยน พุทฺโธ ภควา สาวตฺถิยํ วิหรหิ เชตวเน อนาถปิณฺฑิกสฺส อาราเม.\csromanspacer{} เตน โข ปน สมเยน ฉพฺพคฺคิยา ภิกฺขู อนาจารํ อาจริตฺวา “อญฺญาณเกน อาปนฺนาติ ชานนฺตู”ติ ปาติโมกฺเข อุทฺทิสฺสมาเน เอวํ วทนฺติ “อิทาเนว โข มยํ ชานาม อยมฺปิ กิร ธมฺโม สุตฺตาคโต สุตฺตปริยาปนฺโน อนฺวทฺธมาสํ อุทฺเทสํ อาคจฺฉตี”ติ.\csromanspacer{} เย เต ภิกฺขู อปฺปิจฺฉา ฯเปฯ เต อุชฺฌายนฺติ ขิยฺยนฺติ วิปาเจนฺติ “กถํ หิ นาม ฉพฺพคฺคิยา ภิกฺขู ปาติโมกฺเข อุทฺทิสฺสมาเน เอวํ วกฺขนฺติ อิทาเนว โข มยํ \csromanfolio{190}ชานาม อยมฺปิ กิร ธมฺโม สุตฺตาคโต สุตฺตปริยาปนฺโน อนฺวทฺธมาสํ อุทฺเทสํ อาคจฺฉตี”ติ ฯเปฯ.\csromanspacer{} สจฺจํ กิร ตุเมฺห ภิกฺขเว ปาติโมกฺเข อุทฺทิสฺสมาเน เอวํ วเทถ “อิทาเนว โข มยํ ชานาม อยมฺปิ กิร ธมฺโม สุตฺตาคโต สุตฺตปริยาปนฺโน อนฺวทฺธมาสํ อุทฺเทสํ อาคจฺฉตี”ติ.\csromanspacer{} สจฺจํ ภควาติ.\csromanspacer{} วิครหิ พุทฺโธ ภควา ฯเปฯ กถํ หิ นาม ตุเมฺห โมฆปุริสา ปาติโมกฺเข อุทฺทิสฺสมาเน เอวํ วกฺขถ “อิทาเนว โข มยํ ชานาม อยมฺปิ กิร ธมฺโม สุตฺตาคโต สุตฺตปริยาปนฺโน อนฺวทฺธมาสํ อุทฺเทสํ อาคจฺฉตี”ติ.\csromanspacer{} เนตํ โมฆปุริสา อปฺปสนฺนานํ วา ปสาทาย ฯเปฯ.\csromanspacer{} เอวญฺจ ปน ภิกฺขเว อิมํ สิกฺขาปทํ อุทฺทิเสยฺยาถ\csromandash{}}

\proseitem{444}{\textbf{“โย ปน ภิกฺขุ อนฺวทฺธมาสํ ปาติโมกฺเข อุทฺทิสฺสมาเน เอวํ วเทยฺย ‘อิทาเนว โข อหํ ชานามิ อยมฺปิ กิร ธมฺโม สุตฺตาคโต สุตฺตปริยาปนฺโน อนฺวทฺธมาสํ อุทฺเทสํ อาคจฺฉตี’ติ.}\csromanspacer{} ตญฺเจ ภิกฺขุํ อญฺเญ ภิกฺขู ชาเนยฺยุํ นิสินฺนปุพฺพํ อิมินา ภิกฺขุนา ทฺวตฺติกฺขตฺตุํ ปาติโมกฺเข อุทฺทิสฺสมาเน, โก ปน วาโท ภิยฺโย\footnote{ภิยฺโยติ (สฺยา)}\textbf{, น จ ตสฺส ภิกฺขุโน อญฺญาณเกน มุตฺติ อตฺถิ, ยญฺจ ตตฺถ อาปตฺติํ อาปนฺโน, ตญฺจ ยถาธมฺโม กาเรตพฺโพ, อุตฺตริ จสฺส โมโห อาโรเปตพฺโพ ‘ตสฺส เต อาวุโส อลาภา, ตสฺส เต ทุลฺลทฺธํ, ยํ ตฺวํ ปาติโมกฺเข อุทฺทิสฺสมาเน น สาธุกํ อฏฺฐิํ กตฺวา}\footnote{อฏฺฐิกตฺวา (สฺยา, ก)} \textbf{มนสิ กโรสี’ติ, อิทํ ตสฺมิํ โมหนเก ปาจิตฺติยนฺ”}ติ.}

\proseitem{445}{\textbf{โย ปนา}ติ โย ยาทิโส ฯเปฯ.\csromanspacer{} \textbf{ภิกฺขู}ติ ฯเปฯ อยํ อิมสฺมิํ อตฺเถ อธิปฺเปโต ภิกฺขูติ.}

\prose{\textbf{อนฺวทฺธมาสนฺ}ติ อนุโปสถิกํ.}

\prose{\textbf{ปาติโมกฺเข อุทฺทิสฺสมาเน}ติ อุทฺทิสนฺเต.}

\prose{\textbf{เอวํ วเทยฺยา}ติ อนาจารํ อาจริตฺวา “อญฺญาณเกน อาปนฺโนติ ชานนฺตู”ติ ปาติโมกฺเข อุทฺทิสฺสมาเน เอวํ วเทติ “อิทาเนว โข อหํ ชานามิ อยมฺปิ กิร ธมฺโม สุตฺตาคโต สุตฺตปริยาปนฺโน อนฺวทฺธมาสํ อุทฺเทสํ อาคจฺฉตี”ติ, อาปตฺติ ทุกฺกฏสฺส.}

\prose{\csromanfolio{191}ตญฺเจ โมเหตุกามํ ภิกฺขุํ อญฺเญ ภิกฺขู ชาเนยฺยุํ นิสินฺนปุพฺพํ อิมินา ภิกฺขุนา ทฺวตฺติกฺขตฺตุํ ปาติโมกฺเข อุทฺทิสฺสมาเน, โก ปน วาโท ภิยฺโย, น จ ตสฺส ภิกฺขุโน อญฺญาณเกน มุตฺติ อตฺถิ, ยญฺจ ตตฺถ อาปตฺติํ อาปนฺโน, ตญฺจ ยถาธมฺโม กาเรตพฺโพ, อุตฺตริ จสฺส โมโห อาโรเปตพฺโพ.\csromanspacer{} เอวญฺจ ปน ภิกฺขเว อาโรเปตพฺโพ, พฺยตฺเตน ภิกฺขุนา ปฏิพเลน สํโฆ ญาเปตพฺโพ\csromandash{}}

\proseitem{446}{“สุณาตุ เม ภนฺเต สํโฆ, อยํ อิตฺถนฺนาโม ภิกฺขุ ปาติโมกฺเข อุทฺทิสฺสมาเน น สาธุกํ อฏฺฐิํ กตฺวา มนสิ กโรติ, ยทิ สํฆสฺส ปตฺตกลฺลํ, สํโฆ อิตฺถนฺนามสฺส ภิกฺขุโน โมหํ อาโรเปยฺย, เอสา ญตฺติ.}

\prose{สุณาตุ เม ภนฺเต สํโฆ, อยํ อิตฺถนฺนาโม ภิกฺขุ ปาติโมกฺเข อุทฺทิสฺสมาเน น สาธุกํ อฏฺฐิํ กตฺวา มนสิ กโรติ, สํโฆ อิตฺถนฺนามสฺส ภิกฺขุโน โมหํ อาโรเปติ, ยสฺสายสฺมโต ขมติ อิตฺถนฺนามสฺส ภิกฺขุโน โมหสฺส อาโรปนา, โส ตุณฺหสฺส, ยสฺส นกฺขมติ, โส ภาเสยฺย.}

\prose{อาโรปิโต สํเฆน อิตฺถนฺนามสฺส ภิกฺขุโน โมโห, ขมติ สํฆสฺส, ตสฺมา ตุณฺหี, เอวเมตํ ธารยามี”ติ.}

\prose{อนาโรปิเต โมเห โมเหติ, อาปตฺติ ทุกฺกฏสฺส.\csromanspacer{} อาโรปิเต โมเห โมเหติ, อาปตฺติ}

\proseitem{447}{ธมฺมกมฺเม ธมฺมกมฺมสญฺญี โมเหติ, อาปตฺติ ปาจิตฺติยสฺส.\csromanspacer{} ธมฺมกมฺเม เวมติโก โมเหติ, อาปตฺติ ปาจิตฺติยสฺส.\csromanspacer{} ธมฺมกมฺเม อธมฺมกมฺมสญฺญี โมเหติ, อาปตฺติ ปาจิตฺติยสฺส.}

\prose{อธมฺมกมฺเม ธมฺมกมฺมสญฺญี, อาปตฺติ ทุกฺกฏสฺส.\csromanspacer{} อธมฺมกมฺเม เวมติโก, อาปตฺติ ทุกฺกฏสฺส.\csromanspacer{} อธมฺมกมฺเม อธมฺมกมฺมสญฺญี, อาปตฺติ ทุกฺกฏสฺส.}

\proseitemwithcloser{448}{อนาปตฺติ น วิตฺถาเรน สุตํ โหติ, อูนกทฺวตฺติกฺขตฺตุํ วิตฺถาเรน สุตํ โหติ, นโมเหตุกามสฺส อุมฺมตฺตกสฺส อาทิกมฺมิกสฺสาติ.}{โมหนสิกฺขาปทํ นิฏฺฐิตํ ตติยํ.}
\csromanmatikaanchor{mtk.83}
\csromantocmarkref{subsection}{4. ปหารสิกฺขาปท}{mtk.83}
\bookmark[dest={mtk.83},level=2]{4. ปหารสิกฺขาปท}
\csromanheader{2}{8. \textbf{สหธมฺมิกวคฺค 4. ปหารสิกฺขาปท}}
\proseitem{449}{\csromanfolio{192}เตน สมเยน พุทฺโธ ภควา สาวตฺถิยํ วิหรติ เชตวเน อนาถปิณฺฑิกสฺส อาราเม.\csromanspacer{} เตน โข ปน สมเยน ฉพฺพคฺคิยา ภิกฺขู กุปิตา อนตฺตมนา สตฺตรสวคฺคิยานํ ภิกฺขูนํ ปหารํ เทนฺติ, เต โรทนฺติ.\csromanspacer{} ภิกฺขู เอวมาหํสุ “กิสฺส ตุเมฺห อาวุโส โรทถา”ติ.\csromanspacer{} อิเม อาวุโส ฉพฺพคฺคิยา ภิกฺขู กุปิตา อนตฺตมนา อมฺหากํ ปหารํ เทนฺตีติ.\csromanspacer{} เย เต ภิกฺขู อปฺปิจฺฉา ฯเปฯ เต อุชฺฌายนฺติ ขิยฺยนฺติ วิปาเจนฺติ “กถํ หิ นาม ฉพฺพคฺคิยา ภิกฺขู กุปิตา อนตฺตมนา ภิกฺขูนํ ปหารํ ทสฺสนฺตี”ติ ฯเปฯ.\csromanspacer{} สจฺจํ กิร ตุเมฺห ภิกฺขเว กุปิตา อนตฺตมนา ภิกฺขูนํ ปหารํ เทถาติ.\csromanspacer{} สจฺจํ ภควาติ.\csromanspacer{} วิครหิ พุทฺโธ ภควา ฯเปฯ.\csromanspacer{} กถํ หิ นาม ตุเมฺห โมฆปุริสา กุปิตา อนตฺตมนา ภิกฺขูนํ ปหารํ ทสฺสถ.\csromanspacer{} เนตํ โมฆปุริสา อปฺปสนฺนานํ วา ปสาทาย ฯเปฯ.\csromanspacer{} เอวญฺจ ปน ภิกฺขเว อิมํ สิกฺขาปทํ อุทฺทิเสยฺยาถ\csromandash{}}

\proseitem{450}{\textbf{“โย ปน ภิกฺขุ ภิกฺขุสฺส กุปิโต อนตฺตมโน ปหารํ ทเทยฺย, ปาจิตฺติยนฺ”}ติ.}

\proseitem{451}{\textbf{โย ปนา}ติ โย ยาทิโส ฯเปฯ.\csromanspacer{} \textbf{ภิกฺขู}ติ ฯเปฯ อยํ อิมสฺมิํ อตฺเถ อธิปฺเปโต ภิกฺขูติ.}

\prose{\textbf{ภิกฺขุสฺสา}ติ อญฺญสฺส ภิกฺขุสฺส.}

\prose{\textbf{กุปิโต อนตฺตมโน}ติ อนภิรทฺโธ อาหตจิตฺโต ขิลชาโต.}

\prose{\textbf{ปหารํ ทเทยฺยา}ติ กาเยน วา กายปฏิพทฺเธน วา นิสฺสคฺคิเยน วา อนฺตมโส อุปฺปลปตฺเตนปิ ปหารํ เทติ, อาปตฺติ ปาจิตฺติยสฺส.}

\proseitem{452}{อุปสมฺปนฺเน อุปสมฺปนฺนสญฺญี กุปิโต อนตฺตมโน ปหารํ เทติ, อาปตฺติ ปาจิตฺติยสฺส.\csromanspacer{} อุปสมฺปนฺเน เวมติโก กุปิโต อนตฺตมโน ปหารํ เทติ, อาปตฺติ ปาจิตฺติยสฺส.\csromanspacer{} อุปสมฺปนฺเน อนุปสมฺปนฺนสญฺญี กุปิโต อนตฺตมโน ปหารํ เทติ, อาปตฺติ ปาจิตฺติยสฺส.}

\prose{อนุปสมฺปนฺนสฺส กุปิโต อนตฺตมโน ปหารํ เทติ, อาปตฺติ ทุกฺกฏสฺส.\csromanspacer{} อนุปสมฺปนฺเน อุปสมฺปนฺนสญฺญี, อาปตฺติ ทุกฺกฏสฺส.\csromanspacer{} อนุปสมฺปนฺเน \csromanfolio{193}เวมติโก, อาปตฺติ ทุกฺกฏสฺส.\csromanspacer{} อนุปสมฺปนฺเน อนุปสมฺปนฺนสญฺญี, อาปตฺติ \textbf{ทุกฺกฏสฺส.}}

\proseitemwithcloserruled{453}{อนาปตฺติ เกนจิ วิเหฐียมาโน โมกฺขาธิปฺปาโย ปหารํ เทติ, อุมฺมตฺตกสฺส อาทิกมฺมิกสฺสาติ.}{ปหารสิกฺขาปทํ นิฏฺฐิตํ จตุตฺถํ.}
\csromanmatikaanchor{mtk.84}
\csromantocmarkref{subsection}{5. ตลสตฺติกสิกฺขาปท}{mtk.84}
\bookmark[dest={mtk.84},level=2]{5. ตลสตฺติกสิกฺขาปท}
\csromanheader{2}{8. \textbf{สหธมฺมิกวคฺค 5. ตลสตฺติกสิกฺขาปท}}
\proseitem{454}{เตน สมเยน พุทฺโธ ภควา สาวตฺถิยํ วิหรติ เชตวเน อนาถปิณฺฑิกสฺส อาราเม.\csromanspacer{} เตน โข ปน สมเยน ฉพฺพคฺคิยา ภิกฺขู กุปิตา อนตฺตมนา สตฺตรสวคฺคิยานํ ภิกฺขูนํ ตลสตฺติกํ อุคฺคิรนฺติ, เต ปหารสมุจฺจิตา โรทนฺติ.\csromanspacer{} ภิกฺขู เอวมาหํสุ “กิสฺส ตุเมฺห อาวุโส โรทถา”ติ.\csromanspacer{} อิเม อาวุโส ฉพฺพคฺคิยา ภิกฺขู กุปิตา อนตฺตมนา อมฺหากํ ตลสตฺติกํ อุคฺคิรนฺตีติ.\csromanspacer{} เย เต ภิกฺขู อปฺปิจฺฉา ฯเปฯ เต อุชฺฌายนฺติ ขิยฺยนฺติ วิปาเจนฺติ “กถํ หิ นาม ฉพฺพคฺคิยา ภิกฺขู กุปิตา อนตฺตมนา สตฺตรสวคฺคิยานํ ภิกฺขูนํ ตลสตฺติกํ อุคฺคิริสฺสนฺตี”ติ ฯเปฯ.\csromanspacer{} สจฺจํ กิร ตุเมฺห ภิกฺขเว กุปิตา อนตฺตมนา สตฺตรสวคฺคิยานํ ภิกฺขูนํ ตลสตฺติกํ อุคฺคิรถาติ.\csromanspacer{} สจฺจํ ภควาติ.\csromanspacer{} วิครหิ พุทฺโธ ภควา ฯเปฯ.\csromanspacer{} กถํ หิ นาม ตุเมฺห โมฆปุริสา กุปิตา อนตฺตมนา สตฺตรสวคฺคิยานํ ภิกฺขูนํ ตลสตฺติกํ อุคฺคิริสฺสถ.\csromanspacer{} เนตํ โมฆปุริสา อปฺปสนฺนานํ วา ปสาทาย ฯเปฯ.\csromanspacer{} เอวญฺจ ปน ภิกฺขเว อิมํ สิกฺขาปทํ อุทฺทิเสยฺยาถ\csromandash{}}

\proseitem{455}{\textbf{“โย ปน ภิกฺขุ ภิกฺขุสฺส กุปิโต อนตฺตมโน ตลสตฺติกํ อุคฺคิเรยฺย, ปาจิตฺติยนฺ”}ติ.}

\proseitem{456}{\textbf{โย ปนา}ติ โย ยาทิโส ฯเปฯ.\csromanspacer{} \textbf{ภิกฺขู}ติ ฯเปฯ อยํ อิมสฺมิํ อตฺเถ อธิปฺเปโต ภิกฺขูติ.}

\prose{\textbf{ภิกฺขุสฺสา}ติ อญฺญสฺส ภิกฺขุสฺส.}

\prose{\textbf{กุปิโต อนตฺตมโน}ติ อนภิรทฺโธ อาหตจิตฺโต ขิลชาโต.}

\prose{\csromanfolio{194}\textbf{ตลสตฺติกํ อุคฺคิเรยฺยา}ติ กายํ วา กายปฏิพทฺธํ วา อนฺตมโส อุปฺปลปตฺตมฺปิ อุจฺจาเรติ, อาปตฺติ ปาจิตฺติยสฺส.}

\proseitem{457}{อุปสมฺปนฺเน อุปสมฺปนฺนสญฺญี กุปิโต อนตฺตมโน ตลสตฺติกํ อุคฺคิรติ, อาปตฺติ ปาจิตฺติยสฺส.\csromanspacer{} อุปสมฺปนฺเน เวมติโก กุปิโต อนตฺตมโน ตลสตฺติกํ อุคฺคิรติ, อาปตฺติ ปาจิตฺติยสฺส.\csromanspacer{} อุปสมฺปนฺเน อนุปสมฺปนฺนสญฺญี กุปิโต อนตฺตมโน ตลสตฺติกํ อุคฺคิรติ, อาปตฺติ ปาจิตฺติยสฺส.}

\prose{อนุปสมฺปนฺนสฺส กุปิโต อนตฺตมโน ตลสตฺติกํ อุคฺคิรติ, อาปตฺติ ทุกฺกฏสฺส.\csromanspacer{} อนุปสมฺปนฺเน อุปสมฺปนฺนสญฺญี, อาปตฺติ ทุกฺกฏสฺส.\csromanspacer{} อนุปสมฺปนฺเน เวมติโก, อาปตฺติ ทุกฺกฏสฺส.\csromanspacer{} อนุปสมฺปนฺเน อนุปสมฺปนฺนสญฺญี, อาปตฺติ ทุกฺกฏสฺส.}

\proseitemwithcloserruled{458}{อนาปตฺติ เกนจิ วิเหฐียมาโน โมกฺขาธิปฺปาโย ตลสตฺติกํ อุคฺคิรติ, อุมฺมตฺตกสฺส อาทิกมฺมิกสฺสาติ.}{ตลสตฺติกสิกฺขาปทํ นิฏฺฐิตํ ปญฺจมํ.}
\csromanmatikaanchor{mtk.85}
\csromantocmarkref{subsection}{6. อมูลกสิกฺขาปท}{mtk.85}
\bookmark[dest={mtk.85},level=2]{6. อมูลกสิกฺขาปท}
\csromanheader{2}{8. \textbf{สหธมฺมิกวคฺค 6. อมูลกสิกฺขาปท}}
\proseitem{459}{เตน สมเยน พุทฺโธ ภควา สาวตฺถิยํ วิหรติ เชตวเน อนาถปิณฺฑิกสฺส อาราเม.\csromanspacer{} เตน โข ปน สมเยน ฉพฺพคฺคิยา ภิกฺขู ภิกฺขุํ อมูลเกน สํฆาทิเสเสน อนุทฺธํเสนฺติ.\csromanspacer{} เย เต ภิกฺขู อปฺปิจฺฉา ฯเปฯ เต อุชฺฌายนฺติ ขิยฺยนฺติ วิปาเจนฺติ “กถํ หิ นาม ฉพฺพคฺคิยา ภิกฺขู ภิกฺขุํ อมูลเกน สํฆาทิเสเสน อนุทฺธํเสสฺสนฺตี”ติ ฯเปฯ.\csromanspacer{} สจฺจํ กิร ตุเมฺห ภิกฺขเว ภิกฺขุํ อมูลเกน สํฆาทิเสเสน อนุทฺธํเสถาติ.\csromanspacer{} สจฺจํ ภควาติ.\csromanspacer{} วิครหิ พุทฺโธ ภควา ฯเปฯ.\csromanspacer{} กถํ หิ นาม ตุเมฺห โมฆปุริสา ภิกฺขุํ อมูลเกน สํฆาทิเสเสน อนุทฺธํเสสฺสถ.\csromanspacer{} เนตํ โมฆปุริสา อปฺปสนฺนานํ วา ปสาทาย ฯเปฯ.\csromanspacer{} เอวญฺจ ปน ภิกฺขเว อิมํ สิกฺขาปทํ อุทฺทิเสยฺยาถ\csromandash{}}

\proseitem{460}{\textbf{“โย ปน ภิกฺขุ ภิกฺขุํ อมูลเกน สํฆาทิเสเสน อนุทฺธํเสยฺย, ปาจิตฺติยนฺ”}ติ.}

\proseitem{461}{\csromanfolio{195}\textbf{โย ปนา}ติ โย ยาทิโส ฯเปฯ.\csromanspacer{} \textbf{ภิกฺขู}ติ ฯเปฯ อยํ อิมสฺมิํ อตฺเถ อธิปฺเปโต ภิกฺขูติ.}

\prose{\textbf{ภิกฺขุนฺ}ติ อญฺญํ ภิกฺขุํ.}

\prose{\textbf{อมูลกํ} นาม อทิฏฺฐํ อสฺสุตํ อปริสงฺกิตํ.}

\prose{\textbf{สํฆาทิเสเสนา}ติ เตรสนฺนํ อญฺญตเรน.}

\prose{\textbf{อนุทฺธํเสยฺยา}ติ โจเทติ วา โจทาเปติ วา, อาปตฺติ ปาจิตฺติยสฺส.}

\proseitem{462}{อุปสมฺปนฺเน อุปสมฺปนฺนสญฺญี อมูลเกน สํฆาทิเสเสน อนุทฺธํเสติ, อาปตฺติ ปาจิตฺติยสฺส.\csromanspacer{} อุปสมฺปนฺเน เวมติโก อมูลเกน สํฆาทิเสเสน อนุทฺธํเสติ, อาปตฺติ ปาจิตฺติยสฺส.\csromanspacer{} อุปสมฺปนฺเน อนุปสมฺปนฺนสญฺญี อมูลเกน สํฆาทิเสเสน อนุทฺธํเสติ, อาปตฺติ}

\prose{อาจารวิปตฺติยา วา ทิฏฺฐิวิปตฺติยา วา อนุทฺธํเสติ, อาปตฺติ ทุกฺกฏสฺส.\csromanspacer{} อนุปสมฺปนฺนํ อนุทฺธํเสติ, อาปตฺติ ทุกฺกฏสฺส.\csromanspacer{} อนุปสมฺปนฺเน อุปสมฺปนฺนสญฺญี, อาปตฺติ ทุกฺกฏสฺส.\csromanspacer{} อนุปสมฺปนฺเน เวมติโก, อาปตฺติ ทุกฺกฏสฺส.\csromanspacer{} อนุปสมฺปนฺเน อนุปสมฺปนฺนสญฺญี, อาปตฺติ ทุกฺกฏสฺส.}

\proseitemwithcloserruled{463}{อนาปตฺติ ตถาสญฺญี โจเทติ วา โจทาเปติ วา, อุมฺมตฺตกสฺส อาทิกมฺมิกสฺสาติ.}{อมูลกสิกฺขาปทํ นิฏฺฐิตํ ฉฏฺฐํ.}
\csromanmatikaanchor{mtk.86}
\csromantocmarkref{subsection}{7. สญฺจิจฺจสิกฺขาปท}{mtk.86}
\bookmark[dest={mtk.86},level=2]{7. สญฺจิจฺจสิกฺขาปท}
\csromanheader{2}{8. \textbf{สหธมฺมิกวคฺค 7. สญฺจิจฺจสิกฺขาปท}}
\proseitem{464}{เตน สมเยน พุทฺโธ ภควา สาวตฺถิยํ วิหรติ เชตวเน อนาถปิณฺฑิกสฺส อาราเม.\csromanspacer{} เตน โข ปน สมเยน ฉพฺพคฺคิยา ภิกฺขู สตฺตรสวคฺคิยานํ ภิกฺขูนํ สญฺจิจฺจ กุกฺกุจฺจํ อุปทหนฺติ “ภควตา อาวุโส สิกฺขาปทํ ปญฺญตฺตํ ‘น อูนวีสติวสฺโส ปุคฺคโล อุปสมฺปาเทตพฺโพ’ติ, ตุเมฺห จ อูนวีสติวสฺสา อุปสมฺปนฺนา, กจฺจิ โน ตุเมฺห อนุปสมฺปนฺนา”ติ.\csromanspacer{} เต โรทนฺติ.\csromanspacer{} ภิกฺขู เอวมาหํสุ “กิสฺส ตุเมฺห อาวุโส โรทถา”ติ. \csromanfolio{196}อิเม อาวุโส ฉพฺพคฺคิยา ภิกฺขู อมฺหากํ สญฺจิจฺจ กุกฺกุจฺจํ อุปทหนฺตีติ.\csromanspacer{} เย เต ภิกฺขู อปฺปิจฺฉา ฯเปฯ เต อุชฺฌายนฺติ ขิยฺยนฺติ วิปาเจนฺติ “กถํ หิ นาม ฉพฺพคฺคิยา ภิกฺขู ภิกฺขูนํ สญฺจิจฺจ กุกฺกุจฺจํ อุปทหิสฺสนฺตี”ติ ฯเปฯ.\csromanspacer{} สจฺจํ กิร ตุเมฺห ภิกฺขเว ภิกฺขูนํ สญฺจิจฺจ กุกฺกุจฺจํ อุปทหถาติ.\csromanspacer{} สจฺจํ ภควาติ.\csromanspacer{} วิครหิ พุทฺโธ ภควา ฯเปฯ.\csromanspacer{} กถํ หิ นาม ตุเมฺห โมฆปุริสา ภิกฺขูนํ สญฺจิจฺจ กุกฺกุจฺจํ อุปทหิสฺสถ.\csromanspacer{} เนตํ โมฆปุริสา อปฺปสนฺนานํ วา ปสาทาย ฯเปฯ.\csromanspacer{} เอวญฺจ ปน ภิกฺขเว อิมํ สิกฺขาปทํ อุทฺทิเสยฺยาถ\csromandash{}}

\proseitem{465}{\textbf{“โย ปน ภิกฺขุ ภิกฺขุสฺส สญฺจิจฺจ กุกฺกุจฺจํ อุปทเหยฺย}\footnote{อุปฺปาเทยฺย (อิติปิ)} \textbf{‘อิติสฺส มุหุตฺตมฺปิ อผาสุ ภวิสฺสตี’ติ เอตเทว ปจฺจยํ กริตฺวา อนญฺญํ, ปาจิตฺติยนฺ”}ติ.}

\proseitem{466}{\textbf{โย ปนา}ติ โย ยาทิโส ฯเปฯ.\csromanspacer{} \textbf{ภิกฺขู}ติ ฯเปฯ อยํ อิมสฺมิํ อตฺเถ อธิปฺเปโต ภิกฺขูติ.}

\prose{\textbf{ภิกฺขุสฺสา}ติ อญฺญสฺส ภิกฺขุสฺส.}

\prose{\textbf{สญฺจิจฺจา}ติ ชานนฺโต สญฺชานนฺโต เจจฺจ อภิวิตริตฺวา วีติกฺกโม.}

\prose{\textbf{กุกฺกุจฺจํ อุปทเหยฺยา}ติ “อูนวีสติวสฺโส มญฺเญ ตฺวํ อุปสมฺปนฺโน, วิกาเล มญฺเญ ตยา ภุตฺตํ, มชฺชํ มญฺเญ ตยา ปีตํ, มาตุคาเมน สทฺธิํ รโห มญฺเญ ตยา นิสินฺนนฺ”ติ กุกฺกุจฺจํ อุปทหติ, อาปตฺติ}

\prose{\textbf{เอตเทว ปจฺจยํ กริตฺวา อนญฺญนฺ}ติ น อญฺโญ โกจิ ปจฺจโย โหติ กุกฺกุจฺจํ อุปทหิตุํ.}

\proseitem{467}{อุปสมฺปนฺเน อุปสมฺปนฺนสญฺญี สญฺจิจฺจ กุกฺกุจฺจํ อุปทหติ, อาปตฺติ ปาจิตฺติยสฺส.\csromanspacer{} อุปสมฺปนฺเน เวมติโก สญฺจิจฺจ กุกฺกุจฺจํ อุปทหติ, อาปตฺติ ปาจิตฺติยสฺส.\csromanspacer{} อุปสมฺปนฺเน อนุปสมฺปนฺนสญฺญี สญฺจิจฺจ กุกฺกุจฺจํ อุปทหติ, อาปตฺติ ปาจิตฺติยสฺส.}

\prose{อนุปสมฺปนฺนสฺส สญฺจิจฺจ กุกฺกุจฺจํ อุปทหติ, อาปตฺติ ทุกฺกฏสฺส.\csromanspacer{} อนุปสมฺปนฺเน อุปสมฺปนฺนสญฺญี, อาปตฺติ ทุกฺกฏสฺส.\csromanspacer{} อนุปสมฺปนฺเน เวมติโก, อาปตฺติ ทุกฺกฏสฺส.\csromanspacer{} อนุปสมฺปนฺเน อนุปสมฺปนฺนสญฺญี, อาปตฺติ ทุกฺกฏสฺส.}

\proseitemwithcloserruled{468}{\csromanfolio{197}อนาปตฺติ น กุกฺกุจฺจํ อุปทหิตุกาโม “อูนวีสติวสฺโส มญฺเญ ตฺวํ อุปสมฺปนฺโน, วิกาเล มญฺเญ ตยา ภุตฺตํ, มชฺชํ มญฺเญ ตยา ปีตํ, มาตุคาเมน สทฺธิํ รโห มญฺเญ ตยา นิสินฺนํ, อิงฺฆ ชานาหิ, มา เต ปจฺฉา กุกฺกุจฺฉํ อโหสี”ติ ภณติ, อุมฺมตฺตกสฺส อาทิกมฺมิกสฺสาติ.}{สญฺจิจฺจสิกฺขาปทํ นิฏฺฐิตํ สตฺตมํ}
\csromanmatikaanchor{mtk.87}
\csromantocmarkref{subsection}{8. อุปสฺสุติสิกฺขาปท}{mtk.87}
\bookmark[dest={mtk.87},level=2]{8. อุปสฺสุติสิกฺขาปท}
\csromanheader{2}{8. \textbf{สหธมฺมิกวคฺค 8. อุปสฺสุติสิกฺขาปท}}
\proseitem{469}{เตน สมเยน พุทฺโธ ภควา สาวตฺถิยํ วิหรติ เชตวเน อนาถปิณฺฑิกสฺส อาราเม.\csromanspacer{} เตน โข ปน สมเยน ฉพฺพคฺคิยา ภิกฺขู เปสเลหิ ภิกฺขูหิ สทฺธิํ ภณฺฑนฺติ.\csromanspacer{} เปสลา ภิกฺขู เอวํ วทนฺติ “อลชฺชิโน อิเม อาวุโส ฉพฺพคฺคิยา ภิกฺขู, น สกฺกา อิเมหิ สห ภณฺฑิตุนฺ”ติ.\csromanspacer{} ฉพฺพคฺคิยา ภิกฺขู เอวํ วทนฺติ “กิสฺส ตุเมฺห อาวุโส อเมฺห อลชฺชิวาเทน ปาเปถา”ติ.\csromanspacer{} กหํ ปน ตุเมฺห อาวุโส อสฺสุตฺถาติ.\csromanspacer{} มยํ อายสฺมนฺตานํ อุปสฺสุติํ\footnote{อุปสฺสุติ (?)} ติฏฺฐมฺหาติ.\csromanspacer{} เย เต ภิกฺขู อปฺปิจฺฉา ฯเปฯ เต อุชฺฌายนฺติ ขิยฺยนฺติ วิปาเจนฺติ “กถํ หิ นาม ฉพฺพคฺคิยา ภิกฺขู ภิกฺขูนํ ภณฺฑนชาตานํ กลหชาตานํ วิวาทาปนฺนานํ อุปสฺสุติํ1 ติฏฺฐิสฺสนฺตี”ติ ฯเปฯ.\csromanspacer{} สจฺจํ กิร ตุเมฺห ภิกฺขเว ภิกฺขูนํ ภณฺฑนชาตานํ กลหชาตานํ วิวาทาปนฺนานํ อุปสฺสุติํ1 ติฏฺฐถาติ.\csromanspacer{} สจฺจํ ภควาติ.\csromanspacer{} วิครหิ พุทฺโธ ภควา ฯเปฯ.\csromanspacer{} กถํ หิ นาม ตุเมฺห โมฆปุริสา ภิกฺขูนํ ภณฺฑนชาตานํ กลหชาตานํ วิวาทาปนฺนานํ อุปสฺสุติํ1 ติฏฺฐิสฺสถ.\csromanspacer{} เนตํ โมฆปุริสา อปฺปสนฺนานํ วา ปสาทาย ฯเปฯ.\csromanspacer{} เอวญฺจ ปน ภิกฺขเว อิมํ สิกฺขาปทํ อุทฺทิเสยฺยาถ\csromandash{}}

\proseitem{470}{\textbf{“โย ปน ภิกฺขุ ภิกฺขูนํ ภณฺฑนชาตานํ กลหชาตานํ วิวาทาปนฺนานํ อุปสฺสุติํ}1 ติฏฺเฐยฺย ‘ยํ อิเม ภณิสฺสนฺติ ตํ โสสฺสามี’ติ เอตเทว ปจฺจยํ กริตฺวา อนญฺญํ, ปาจิตฺติยนฺ”ติ.}

\proseitem{471}{\textbf{โย ปนา}ติ โย ยาทิโส ฯเปฯ.\csromanspacer{} \textbf{ภิกฺขู}ติ ฯเปฯ อยํ อิมสฺมิํ อตฺเถ อธิปฺเปโต ภิกฺขูติ.}

\prose{\csromanfolio{198}\textbf{ภิกฺขูนนฺ}ติ อญฺเญสํ ภิกฺขูนํ.}

\prose{\textbf{ภณฺฑนชาตานํ กลหชาตานํ วิวาทาปนฺนานนฺ}ติ อธิกรณชาตานํ.}

\prose{\textbf{อุปสฺสุติํ ติฏฺเฐยฺยา}ติ “อิเมสํ สุตฺวา โจเทสฺสามิ สาเรสฺสามิ ปฏิโจเทสฺสามิ ปฏิสาเรสฺสามิ มงฺกู กริสฺสามี”ติ คจฺฉติ, อาปตฺติ ทุกฺกฏสฺส.\csromanspacer{} ยตฺถ ฐิโต สุณาติ, อาปตฺติ ปาจิตฺติยสฺส.\csromanspacer{} ปจฺฉโต คจฺฉนฺโต ตุริโต คจฺฉติ โสสฺสามีติ, อาปตฺติ ทุกฺกฏสฺส.\csromanspacer{} ยตฺถ ฐิโต สุณาติ, อาปตฺติ ปาจิตฺติยสฺส.\csromanspacer{} ปุรโต คจฺฉนฺโต โอหิยฺยติ โสสฺสามีติ, อาปตฺติ ทุกฺกฏสฺส.\csromanspacer{} ยตฺถ ฐิโต สุณาติ, อาปตฺติ ปาจิตฺติยสฺส.\csromanspacer{} ภิกฺขุสฺส ฐิโตกาสํ วา นิสินฺโนกาสํ วา นิปนฺโนกาสํ วา อาคนฺตฺวา มนฺเตนฺตํ อุกฺกาสิตพฺพํ, วิชานาเปตพฺพํ, โน เจ อุกฺกาเสยฺย วา วิชานาเปยฺย วา, อาปตฺติ ปาจิตฺติยสฺส.}

\prose{\textbf{เอตเทว ปจฺจยํ กริตฺวา อนญฺญนฺ}ติ น อญฺโญ โกจิ ปจฺจโย โหติ อุปสฺสุติํ ติฏฺฐิตุํ.}

\proseitem{472}{อุปสมฺปนฺเน อุปสมฺปนฺนสญฺญี อุปสฺสุติํ ติฏฺฐติ, อาปตฺติ ปาจิตฺติยสฺส.\csromanspacer{} อุปสมฺปนฺเน เวมติโก อุปสฺสุติํ ติฏฺฐติ, อาปตฺติ ปาจิตฺติยสฺส.\csromanspacer{} อุปสมฺปนฺเน อนุปสมฺปนฺนสญฺญี อุปสฺสุติํ ติฏฺฐติ, อาปตฺติ ปาจิตฺติยสฺส.}

\prose{อนุปสมฺปนฺนสฺส อุปสฺสุติํ ติฏฺฐติ, อาปตฺติ ทุกฺกฏสฺส.\csromanspacer{} อนุปสมฺปนฺเน อุปสมฺปนฺนสญฺญี, อาปตฺติ ทุกฺกฏสฺส.\csromanspacer{} อนุปสมฺปนฺเน เวมติโก, อาปตฺติ ทุกฺกฏสฺส.\csromanspacer{} อนุปสมฺปนฺเน อนุปสมฺปนฺนสญฺญี, อาปตฺติ ทุกฺกฏสฺส.}

\proseitemwithcloser{473}{อนาปตฺติ “อิเมสํ สุตฺวา โอรมิสฺสามิ วิรมิสฺสามิ วูปสมิสฺสามิ\footnote{วูปสเมสฺสามิ (สี)} อตฺตานํ ปริโมเจสฺสามี”ติ คจฺฉติ, อุมฺมตฺตกสฺส อาทิกมฺมิกสฺสาติ.}{อุปสฺสุติสิกฺขาปทํ นิฏฺฐิตํ อฏฺฐมํ.}
\csromanmatikaanchor{mtk.88}
\csromantocmarkref{subsection}{9. กมฺมปฏิพาหนสิกฺขาปท}{mtk.88}
\bookmark[dest={mtk.88},level=2]{9. กมฺมปฏิพาหนสิกฺขาปท}
\csromanheader{2}{8. สหธมฺมิกวคฺค 9. กมฺม\-ปฏิพาหน\-สิกฺขาปท}
\proseitem{474}{\csromanfolio{199}เตน สมเยน พุทฺโธ ภควา สาวตฺถิยํ วิหรติ เชตวเน อนาถปิณฺฑิกสฺส อาราเม.\csromanspacer{} เตน โข ปน สมเยน ฉพฺพคฺคิยา ภิกฺขู อนาจารํ อาจริตฺวา เอกเมกสฺส กมฺเม กยิรมาเน ปฏิกฺโกสนฺติ.\csromanspacer{} เตน โข ปน สมเยน สํโฆ สนฺนิปติโต โหติ เกนจิเทว กรณีเยน, ฉพฺพคฺคิยา ภิกฺขู จีวรกมฺมํ กโรนฺตา เอกสฺส ฉนฺทํ อทํสุ.\csromanspacer{} อถ โข สํโฆ “อยํ อาวุโส ฉพฺพคฺคิโย ภิกฺขุ เอกโก อาคโต, หนฺทสฺส มยํ กมฺมํ กโรมา”ติ ตสฺส กมฺมํ อกาสิ.\csromanspacer{} อถ โข โส ภิกฺขุ เยน ฉพฺพคฺคิยา ภิกฺขู เตนุปสงฺกมิ, ฉพฺพคฺคิยา ภิกฺขู ตํ ภิกฺขุํ เอตทโวจุํ “กิํ อาวุโส สํโฆ อกาสี”ติ.\csromanspacer{} สํโฆ เม อาวุโส กมฺมํ อกาสีติ.\csromanspacer{} น มยํ อาวุโส เอตทตฺถาย ฉนฺทํ อทมฺหา “ตุยฺหํ กมฺมํ กริสฺสตี”ติ, สเจ จ มยํ ชาเนยฺยาม “ตุยฺหํ กมฺมํ กริสฺสตี”ติ, น มยํ ฉนฺทํ ทเทยฺยามาติ.\csromanspacer{} เย เต ภิกฺขู อปฺปิจฺฉา ฯเปฯ เต อุชฺฌายนฺติ ขิยฺยนฺติ วิปาเจนฺติ “กถํ หิ นาม ฉพฺพคฺคิยา ภิกฺขู ธมฺมิกานํ กมฺมานํ ฉนฺทํ ทตฺวา ปจฺฉา ขียนธมฺมํ\footnote{ขิยฺยธมฺมํ (อิติปิ)} อาปชฺชิสฺสนฺตี”ติ ฯเปฯ.\csromanspacer{} สจฺจํ กิร ตุเมฺห ภิกฺขเว ธมฺมิกานํ กมฺมานํ ฉนฺทํ ทตฺวา ปจฺฉา ขียนธมฺมํ1 อาปชฺชถาติ.\csromanspacer{} สจฺจํ ภควาติ.\csromanspacer{} วิครหิ พุทฺโธ ภควา ฯเปฯ.\csromanspacer{} กถํ หิ นาม ตุเมฺห โมฆปุริสา ธมฺมิกานํ กมฺมานํ ฉนฺทํ ทตฺวา ปจฺฉา ขียนธมฺมํ1 อาปชฺชิสฺสถ.\csromanspacer{} เนตํ โมฆปุริสา อปฺปสนฺนานํ วา ปสาทาย ฯเปฯ.\csromanspacer{} เอวญฺจ ปน ภิกฺขเว อิมํ สิกฺขาปทํ อุทฺทิเสยฺยาถ\csromandash{}}

\proseitem{475}{\textbf{“โย ปน ภิกฺขุ ธมฺมิกานํ กมฺมานํ ฉนฺทํ ทตฺวา ปจฺฉา ขียนธมฺมํ}1 อาปชฺเชยฺย, ปาจิตฺติยนฺ”ติ.}

\proseitem{476}{\textbf{โย ปนา}ติ โย ยาทิโส ฯเปฯ.\csromanspacer{} \textbf{ภิกฺขู}ติ ฯเปฯ อยํ อิมสฺมิํ อตฺเถ อธิปฺเปโต ภิกฺขูติ.}

\prose{\textbf{ธมฺมิกํ} นาม กมฺมํ อปโลกนกมฺมํ ญตฺติ\textbf{กมฺมํ} ญตฺติทุติยกมฺมํ ญตฺติจตุตฺถกมฺมํ ธมฺเมน วินเยน สตฺถุสาสเนน กตํ, เอตํ ธมฺมิกํ นาม กมฺมํ.\csromanspacer{} ฉนฺทํ ทตฺวา ขิยฺยติ, อาปตฺติ ปาจิตฺติยสฺส.}

\proseitem{477}{\csromanfolio{200}ธมฺมกมฺเม ธมฺมกมฺมสญฺญี ฉนฺทํ ทตฺวา ขิยฺยติ, อาปตฺติ ปาจิตฺติยสฺส.\csromanspacer{} ธมฺมกมฺเม เวมติโก ฉนฺทํ ทตฺวา ขิยฺยติ, อาปตฺติ ทุกฺกฏสฺส.\csromanspacer{} ธมฺมกมฺเม อธมฺมกมฺมสญฺญี ฉนฺทํ ทตฺวา ขิยฺยติ, อนาปตฺติ.\csromanspacer{} อธมฺมกมฺเม ธมฺมกมฺมสญฺญี, อาปตฺติ ทุกฺกฏสฺส.\csromanspacer{} อธมฺมกมฺเม เวมติโก, อาปตฺติ ทุกฺกฏสฺส.\csromanspacer{} อธมฺมกมฺเม อธมฺมกมฺมสญฺญี, อนาปตฺติ.}

\proseitemwithcloserruled{478}{อนาปตฺติ “อธมฺเมน วา วคฺเคน วา น กมฺมารหสฺส วา กมฺมํ กตนฺ”ติ ชานนฺโต ขิยฺยติ, อุมฺมตฺตกสฺส อาทิกมฺมิกสฺสาติ.}{กมฺม\-ปฏิพาหน\-สิกฺขาปทํ นิฏฺฐิตํ นวมํ.}
\csromanmatikaanchor{mtk.89}
\csromantocmarkref{subsection}{10. ฉนฺทํอทตฺวาคมนสิกฺขาปท}{mtk.89}
\bookmark[dest={mtk.89},level=2]{10. ฉนฺทํอทตฺวาคมนสิกฺขาปท}
\csromanheader{2}{8. สหธมฺมิกวคฺค 10. ฉนฺทํ\-อทตฺวา\-คมน\-สิกฺขาปท}
\proseitem{479}{เตน สมเยน พุทฺโธ ภควา สาวตฺถิยํ วิหรติ เชตวเน อนาถปิณฺฑิกสฺส อาราเม.\csromanspacer{} เตน โข ปน สมเยน สํโฆ สนฺนิปติโต โหติ เกนจิเทว กรณีเยน.\csromanspacer{} ฉพฺพคฺคิยา ภิกฺขู จีวรกมฺมํ กโรนฺตา เอกสฺส ฉนฺทํ อทํสุ.\csromanspacer{} อถ โข สํโฆ “ยสฺสตฺถาย สนฺนิปติโต, ตํ กมฺมํ กริสฺสามี”ติ ญตฺติํ ฐเปสิ.\csromanspacer{} อถ โข โส ภิกฺขุ “เอวเมวิเม เอกเมกสฺส กมฺมํ กโรนฺติ, กสฺส ตุเมฺห กมฺมํ กริสฺสถา”ติ ฉนฺทํ อทตฺวา อุฏฺฐายาสนา ปกฺกามิ.\csromanspacer{} เย เต ภิกฺขู อปฺปิจฺฉา ฯเปฯ เต อุชฺฌายนฺติ ขิยฺยนฺติ วิปาเจนฺติ “กถํ หิ นาม ภิกฺขุ สํเฆ วินิจฺฉยกถาย วตฺตมานาย ฉนฺทํ อทตฺวา อุฏฺฐายาสนา ปกฺกมิสฺสตี”ติ ฯเปฯ.\csromanspacer{} สจฺจํ กิร ตฺวํ ภิกฺขุ สํเฆ วินิจฺฉยกถาย วตฺตมานาย ฉนฺทํ อทตฺวา อุฏฺฐายาสนา ปกฺกมสีติ.\csromanspacer{} สจฺจํ ภควาติ.\csromanspacer{} วิครหิ พุทฺโธ ภควา ฯเปฯ.\csromanspacer{} กถํ หิ นาม ตฺวํ โมฆปุริส สํเฆ วินิจฺฉยกถาย วตฺตมานาย ฉนฺทํ อทตฺวา อุฏฺฐายาสนา ปกฺกมิสฺสสิ.\csromanspacer{} เนตํ โมฆปุริส อปฺปสนฺนานํ วา ปสาทาย ฯเปฯ.\csromanspacer{} เอวญฺจ ปน ภิกฺขเว อิมํ สิกฺขาปทํ อุทฺทิเสยฺยาถ\csromandash{}}

\proseitem{480}{\textbf{“โย ปน ภิกฺขุ สํเฆ วินิจฺฉยกถาย วตฺตมนาย ฉนฺทํ อทตฺวา อุฏฺฐายาสนา ปกฺกเมยฺย, ปาจิตฺติยนฺ”}ติ.}

\proseitem{481}{\csromanfolio{201}\textbf{โย ปนา}ติ โย ยาทิโส ฯเปฯ.\csromanspacer{} \textbf{ภิกฺขู}ติ ฯเปฯ อยํ อิมสฺมิํ อตฺเถ อธิปฺเปโต ภิกฺขูติ.}

\prose{\textbf{สํเฆ วินิจฺฉยกถา} นาม วตฺถุ วา อาโรจิตํ โหติ อวินิจฺฉิตํ, ญตฺติ วา ฐปิตา โหติ.\csromanspacer{} กมฺมวาจา วา วิปฺปกตา โหติ.}

\prose{\textbf{ฉนฺทํ อทตฺวา อุฏฺฐายาสนา ปกฺกเมยฺยา}ติ “กถํ อิทํ กมฺมํ กุปฺปํ อสฺส วคฺคํ อสฺส น กเรยฺยา”ติ คจฺฉติ, อาปตฺติ ทุกฺกฏสฺส.\csromanspacer{} ปริสาย หตฺถปาสํ วิชหนฺตสฺส อาปตฺติ ทุกฺกฏสฺส.\csromanspacer{} วิชหิเต อาปตฺติ}

\proseitem{482}{ธมฺมกมฺเม ธมฺมกมฺมสญฺญี ฉนฺทํ อทตฺวา อุฏฺฐายาสนา ปกฺกมติ, อาปตฺติ ปาจิตฺติยสฺส.\csromanspacer{} ธมฺมกมฺเม เวมติโก ฉนฺทํ อทตฺวา อุฏฺฐายาสนา ปกฺกมติ, อาปตฺติ ทุกฺกฏสฺส.\csromanspacer{} ธมฺมกมฺเม อธมฺมกมฺมสญฺญี ฉนฺทํ อทตฺวา อุฏฺฐายาสนา ปกฺกมติ, อนาปตฺติ.\csromanspacer{} อธมฺมกมฺเม ธมฺมกมฺมสญฺญี, อาปตฺติ ทุกฺกฏสฺส.\csromanspacer{} อธมฺมกมฺเม เวมติโก, อาปตฺติ ทุกฺกฏสฺส.\csromanspacer{} อธมฺมกมฺเม อธมฺมกมฺมสญฺญี, อนาปตฺติ.}

\proseitemwithcloserruled{483}{อนาปตฺติ “สํฆสฺส ภณฺฑนํ วา กลโห วา วิคฺคโห วา วิวาโท วา ภวิสฺสตี”ติ คจฺฉติ, “สํฆเภโท วา สํฆราชิ วา ภวิสฺสตี”ติ คจฺฉติ, “อธมฺเมน วา วคฺเคน วา น กมฺมารหสฺส วา กมฺมํ กริสฺสตี”ติ คจฺฉติ, คิลาโน คจฺฉติ, คิลานสฺส กรณีเยน คจฺฉติ, อุจฺจาเรน วา ปสฺสาเวน วา ปีฬิโต คจฺฉติ, “น กมฺมํ โกเปตุกาโม ปุน ปจฺจาคมิสฺสามี”ติ คจฺฉติ, อุมฺมตฺตกสฺส อาทิกมฺมิกสฺสาติ.}{ฉนฺทํ อทตฺวา คมนสิกฺขาปทํ นิฏฺฐิตํ ทสมํ.}
\csromanmatikaanchor{mtk.90}
\csromantocmarkref{subsection}{11. ทุพฺพลสิกฺขาปท}{mtk.90}
\bookmark[dest={mtk.90},level=2]{11. ทุพฺพลสิกฺขาปท}
\csromanheader{2}{8. \textbf{สหธมฺมิกวคฺ}ค 11. ทุพฺพลสิกฺขาปท}
\proseitem{484}{เตน สมเยน พุทฺโธ ภควา ราชคเห วิหรติ เวฬุวเน กลนฺทกนิวาเป.\csromanspacer{} เตน โข ปน สมเยน อายสฺมา ทพฺโพ มลฺลปุตฺโต สํฆสฺส เสนาสนญฺจ ปญฺญเปติ, ภตฺตานิ จ อุทฺทิสติ.\csromanspacer{} โส จายสฺมา ทุพฺพลจีวโร โหติ.\csromanspacer{} เตน โข ปน สมเยน สํฆสฺส เอก จีวรํ อุปฺปนฺนํ โหติ.\csromanspacer{} อถ โข สํโฆ ตํ จีวรํ อายสฺมาโต ทพฺพสฺส มลฺลปุตฺตสฺส \csromanfolio{202}อทาสิ.\csromanspacer{} ฉพฺพคฺคิยา ภิกฺขู อุชฺฌายนฺติ ขิยฺยนฺติ วิปาเจนฺติ “ยถาสนฺถุตํ ภิกฺขู สํฆิกํ ลาภํ ปริณาเมนฺตี”ติ.\csromanspacer{} เย เต ภิกฺขู อปฺปิจฺฉา ฯเปฯ เต อุชฺฌายนฺติ ขิยฺยนฺติ วิปาเจนฺติ “กถํ หิ นาม ฉพฺพคฺคิยา ภิกฺขู สมคฺเคน สํเฆน จีวรํ ทตฺวา ปจฺฉา ขียนธมฺมํ อาปชฺชิสฺสนฺตี”ติ ฯเปฯ.\csromanspacer{} สจฺจํ กิร ตุเมฺห ภิกฺขเว สมคฺเคน สํเฆน จีวรํ ทตฺวา ปจฺฉา ขียนธมฺมํ อาปชฺชถาติ.\csromanspacer{} สจฺจํ ภควาติ.\csromanspacer{} วิครหิ พุทฺโธ ภควา ฯเปฯ.\csromanspacer{} กถํ หิ นาม ตุเมฺห โมฆปุริสา สมคฺเคน สํเฆน จีวรํ ทตฺวา ปจฺฉา ขียนธมฺมํ อาปชฺชิสฺสถ.\csromanspacer{} เนตํ โมฆปุริสา อปฺปสนฺนานํ วา ปสาทาย ฯเปฯ.\csromanspacer{} เอวญฺจ ปน ภิกฺขเว อิมํ สิกฺขาปทํ อุทฺทิเสยฺยาถ\csromandash{}}

\proseitem{485}{\textbf{“โย ปน ภิกฺขุ สมคฺเคน สํเฆน จีวรํ ทตฺวา ปจฺฉา ขียนธมฺมํ อาปชฺเชยฺย ‘ยถาสนฺถุตํ ภิกฺขู สํฆิกํ ลาภํ ปริณาเมนฺตี’ติ, ปาจิตฺติยนฺ”}ติ.}

\proseitem{486}{\textbf{โย ปนา}ติ โย ยาทิโส ฯเปฯ.\csromanspacer{} \textbf{ภิกฺขู}ติ ฯเปฯ อยํ อิมสฺมิํ อตฺเถ อธิปฺเปโต ภิกฺขูติ.}

\prose{\textbf{สมคฺโค} นาม สํโฆ สมานสํวาสโก สมานสีมายํ ฐิโต.}

\prose{\textbf{จีวรํ} นาม ฉนฺนํ จีวรานํ อญฺญตรํ จีวรํ วิกปฺปนุปคํ ปจฺฉิมํ.}

\prose{\textbf{ทตฺวา}ติ สยํ ทตฺวา.}

\prose{\textbf{ยถาสนฺถุตํ} นาม ยถามิตฺตตา ยถาสนฺทิฏฺฐตา ยถาสมฺภตฺตตา ยถาสมานุปชฺฌายกตา ยถาสมานาจริยกตา.}

\prose{\textbf{สํฆิกํ} นาม สํฆสฺส ทินฺนํ โหติ ปริจฺจตฺตํ.}

\prose{ลาโภ นาม จีวร\-ปิณฺฑปาต\-เสนาสน\-คิลานปฺปจฺจย\-เภสชฺช\-ปริกฺขารา, อนฺตมโส จุณฺณปิณฺโฑปิ ทนฺตกฏฺฐมฺปิ ทสิกสุตฺตมฺปิ.}

\prose{ปจฺฉา ขียนธมฺมํ อาปชฺเชยฺยาติ อุปสมฺปนฺนสฺส สํเฆน สมฺมตสฺส เสนาสนปญฺญาปกสฺส วา ภตฺตุทฺเทสกสฺส วา ยาคุภาชกสฺส วา ผลภาชกสฺส วา ขชฺชภาชกสฺส วา อปฺปมตฺตก\-วิสฺสชฺชกสฺส วา จีวรํ ทินฺเน ขิยฺยติ, อาปตฺติ ปาจิตฺติยสฺส.}

\proseitem{487}{ธมฺมกมฺเม ธมฺมกมฺมสญฺญี จีวรํ ทินฺเน ขิยฺยติ, อาปตฺติ ปาจิตฺติยสฺส.\csromanspacer{} ธมฺมกมฺเม เวมติโก จีวรํ ทินฺเน ขิยฺยติ, อาปตฺติ ปาจิตฺติยสฺส.\csromanspacer{} ธมฺมกมฺเม อธมฺมกมฺมสญฺญี จีวรํ ทินฺเน ขิยฺยติ, อาปตฺติ ปาจิตฺติยสฺส.}

\prose{\csromanfolio{203}อญฺญํ ปริกฺขารํ ทินฺเน ขิยฺยติ, อาปตฺติ ทุกฺกฏสฺส.\csromanspacer{} อุปสมฺปนฺนสฺส สํเฆน อสมฺมตสฺส เสนาสนปญฺญาปกสฺส วา ภตฺตุทฺเทสกสฺส วา ยาคุภาชกสฺส วา ผลภชกสฺส วา ขชฺชภาชกสฺส วา อปฺปมตฺตก\-วิสฺสชฺชกสฺส วา จีวรํ วา อญฺญํ วา ปริกฺขารํ ทินฺเน ขิยฺยติ, อาปตฺติ ทุกฺกฏสฺส.\csromanspacer{} อนุปสมฺปนฺนสฺส สํเฆน สมฺมตสฺส วา อสมฺมตสฺส วา เสนาสนปญฺญาปกสฺส วา ภตฺตุทฺเทสกสฺส วา ยาคุภาชกสฺส วา ผลภาชกสฺส วา ขชฺชภาชกสฺส วา อปฺปมตฺตก\-วิสฺสชฺชกสฺส วา จีวรํ วา อญฺญํ วา ปริกฺขารํ ทินฺเน ขิยฺยติ, อาปตฺติ ทุกฺกฏสฺส.\csromanspacer{} อธมฺมกมฺเม ธมฺมกมฺมสญฺญี, อาปตฺติ ทุกฺกฏสฺส.\csromanspacer{} อธมฺมกมฺเม เวมติโก, อาปตฺติ ทุกฺกฏสฺส, อธมฺมกมฺเม อธมฺมกมฺมสญฺญี, อนาปตฺติ.}

\proseitemwithcloserruled{488}{อนาปตฺติ ปกติยา ฉนฺทา โทสา โมหา ภยา กโรนฺตํ “กฺวตฺโถ ตสฺส ทินฺเนน ลทฺธาปิ วินิปาเตสฺสติ น สมฺมา อุปเนสฺสตี”ติ ขิยฺยติ, อุมฺมตฺตกสฺส อาทิกมฺมิกสฺสาติ.}{ทุพฺพลสิกฺขาปทํ นิฏฺฐิตํ เอกาทสมํ.}
\csromanmatikaanchor{mtk.91}
\csromantocmarkref{subsection}{12. ปริณามนสิกฺขาปท}{mtk.91}
\bookmark[dest={mtk.91},level=2]{12. ปริณามนสิกฺขาปท}
\csromanheader{2}{8. \textbf{สหธมฺมิกวคฺคฺ}อ 12. ปริณามนสิกฺขาปท}
\proseitem{489}{\csromansymbolfootnote{*}{อิทํ วตฺถุ วิ 1. 379 ปิฏฺเฐปิ อาคตํ.}เตน สมเยน พุทฺโธ ภควา สาวตฺถิยํ วิหรติ เชตวเน อนาถปิณฺฑิกสฺส อาราเม.\csromanspacer{} เตน โข ปน สมเยน สาวตฺถิยํ อญฺญตรสฺส ปูคสฺส สํฆสฺส สจีวรภตฺตํ ปฏิยตฺตํ โหติ “โภเชตฺวา จีวเรน อจฺฉาเทสฺสามา”ติ.\csromanspacer{} อถ โข ฉพฺพคฺคิยา ภิกฺขู เยน โส ปูโค เตนุปสงฺกมิํสุ, อุปสงฺกมิตฺวา ตํ ปูคํ เอตทโวจุํ “เทถาวุโส อิมานิ จีวรานิ อิเมสํ ภิกฺขูนนฺ”ติ.\csromanspacer{} น มยํ ภนฺเต ทสฺสาม, อมฺหากํ สํฆสฺส อนุวสฺสํ สจีวรภิกฺขา ปญฺญตฺตาติ.\csromanspacer{} พหู อาวุโส สํฆสฺส ทายกา, พหู สํฆสฺส ภตฺตา\footnote{ภทฺทา (ก)}, อิเม ตุเมฺห นิสฺสาย ตุเมฺห สมฺปสฺสนฺตา อิธ วิหรนฺติ, ตุเมฺห เจ อิเมสํ น ทสฺสถ, อถ โก จรหิ อิเมสํ ทสฺสติ, เทถาวุโส อิมานิ จีวรานิ อิเมสํ ภิกฺขูนนฺติ.\csromanspacer{} อถ โข โส ปูโค ฉพฺพคฺคิเยหิ ภิกฺขูหิ นิปฺปีฬิยมาโน ยถาปฏิยตฺตํ จีวรํ ฉพฺพคฺคิยานํ ภิกฺขูนํ ทตฺวา สํฆํ ภตฺเตน ปริวิสิ.\csromanspacer{} เย เต ภิกฺขู \csromanfolio{204}ชานนฺติ “สํฆสฺส สจีวรภตฺตํ ปฏิยตฺตํ, น จ ชานนฺติ ฉพฺพคฺคิยานํ ภิกฺขูนํ ทินฺนนฺ”ติ.\csromanspacer{} เต เอวมาหํสุ “โอโณเชถาวุโส สํฆสฺส จีวรนฺ”ติ.\csromanspacer{} นตฺถิ ภนฺเต ยถาปฏิยตฺตํ จีวรํ, อยฺยา ฉพฺพคฺคิยา อยฺยานํ ฉพฺพคฺคิยานํ ปริณาเมสุนฺติ.\csromanspacer{} เย เต ภิกฺขู อปฺปิจฺฉา ฯเปฯ เต อุชฺฌายนฺติ ขิยฺยนฺติ วิปาเจนฺติ “กถํ หิ นาม ฉพฺพคฺคิยา ภิกฺขู ชานํ สํฆิกํ ลาภํ ปริณตํ ปุคฺคลสฺส ปริณาเมสฺสนฺตี”ติ ฯเปฯ.\csromanspacer{} สจฺจํ กิร ตุเมฺห ภิกฺขเว ชานํ สํฆิกํ ลาภํ ปริณตํ ปุคฺคลสฺส ปริณาเมถาติ.\csromanspacer{} สจฺจํ ภควาติ.\csromanspacer{} วิครหิ พุทฺโธ ภควา ฯเปฯ.\csromanspacer{} กถํ หิ นาม ตุเมฺห โมฆปุริสา ชานํ สํฆิกํ ลาภํ ปริณตํ ปุคฺคลสฺส ปริณาเมสฺสถ.\csromanspacer{} เนตํ โมฆปุริสา อปฺปสนฺนานํ วา ปสาทาย ฯเปฯ.\csromanspacer{} เอวญฺจ ปน ภิกฺขเว อิมํ สิกฺขาปทํ อุทฺทิเสยฺยาถ\csromandash{}}

\proseitem{490}{\textbf{“โย ปน ภิกฺขุ ชานํ สํฆิกํ ลาภํ ปริณตํ ปุคฺคลสฺส ปริณาเมยฺย, ปาจิตฺติยนฺ”}ติ.}

\proseitem{491}{\textbf{โย ปนา}ติ โย ยาทิโส ฯเปฯ.\csromanspacer{} \textbf{ภิกฺขู}ติ ฯเปฯ อยํ อิมสฺมิํ อตฺเถ อธิปฺเปโต ภิกฺขูติ.}

\prose{\textbf{ชานาติ} นาม สามํ วา ชานาติ, อญฺเญ วา ตสฺส อาโรเจนฺติ, โส วา อาโรเจติ.}

\prose{\textbf{สํฆิกํ} นาม สํฆสฺส ทินฺนํ โหติ ปริจฺจตฺตํ.}

\prose{ลาโภ นาม จีวร\-ปิณฺฑปาต\-เสนาสน\-คิลานปฺปจฺจย\-เภสชฺช\-ปริกฺขารา, อนฺตมโส จุณฺณปิณฺโฑปิ ทนฺตกฏฺฐมฺปิ ทสิกสุตฺตมฺปิ.}

\prose{\textbf{ปริณตํ} นาม “ทสฺสาม กริสฺสามา”ติ วาจา ภินฺนา โหติ, ตํ ปุคฺคลสฺส ปริณาเมติ, อาปตฺติ ปาจิตฺติยสฺส.}

\proseitem{492}{ปริณเต ปริณตสญฺญี ปุคฺคลสฺส ปริณาเมติ, อาปตฺติ ปาจิตฺติยสฺส.\csromanspacer{} ปริณเต เวมติโก ปุคฺคลสฺส ปริณาเมติ, อาปตฺติ ทุกฺกฏสฺส.\csromanspacer{} ปริณเต อปริณตสญฺญี ปุคฺคลสฺส ปริณาเมติ, อนาปตฺติ.\csromanspacer{} สํฆสฺส ปริณตํ อญฺญสํฆสฺส วา เจติยสฺส วา ปริณาเมติ อาปตฺติ ทุกฺกฏสฺส.\csromanspacer{} เจติยสฺส ปริณตํ อญฺญเจติยสฺส วา สํฆสฺส วา ปุคฺคลสฺส วา ปริณาเมติ, อาปตฺติ ทุกฺกฏสฺส.\csromanspacer{} ปุคฺคลสฺส ปริณตํ อญฺญปุคฺคลสฺส วา สํฆสฺส วา เจติยสฺส วา ปริณาเมติ, อาปตฺติ \csromanfolio{205}ทุกฺกฏสฺส.\csromanspacer{} อปริณเต ปริณตสญฺญี, อาปตฺติ ทุกฺกฏสฺส.\csromanspacer{} อปริณเต เวมติโก, อาปตฺติ ทุกฺกฏสฺส.\csromanspacer{} อปริณเต อปริณตสญฺญี, อนาปตฺติ.}

\proseitemwithcloser{493}{อนาปตฺติ “กตฺถ เทมา”ติ ปุจฺฉียมาโน “ยตฺถ ตุมฺหากํ เทยฺยธมฺโม ปริโภคํ วา ลเภยฺย, ปฏิสงฺขารํ วา ลเภยฺย, จิรฏฺฐิติโก วา อสฺส, ยตฺถ วา ปน ตุมฺหากํ จิตฺตํ ปสีทติ, ตตฺถ เทถา”ติ ภณติ, อุมฺมตฺตกสฺส อาทิกมฺมิกสฺสาติ.}{ปริณามนสิกฺขาปทํ นิฏฺฐิตํ ทฺวาทสมํ.}
\vspace{\tipitakaparskip}
\csromancenter{สหธมฺมิกวคฺโค อฏฺฐโม.}
\csromancenter{ตสฺสุทฺทานํ}
\vspace{0.5\baselineskip}
\setlength{\csromangathapagewidth}{0pt}
\csromangathameasuregroup{สหธมฺม วิวณฺณญฺจ, \\ ตลสตฺติ อมูลญฺจ,}{\csromangathabat{สหธมฺม วิวณฺณญฺจ,}{โมหาปนํ ปหารกํ.} \\ \csromangathabat{ตลสตฺติ อมูลญฺจ,}{สญฺจิจฺจ จ อุปสฺสุติ.} \\ ปฏิพาหน ฉนฺทญฺจ, ทพฺพญฺจ ปริณามนนฺติ.}
\csromangathasetleft{สหธมฺม วิวณฺณญฺจ, \\ ตลสตฺติ อมูลญฺจ,}
\csromangathagroup{\csromangathastanza{\csromangathabat{สหธมฺม วิวณฺณญฺจ,}{โมหาปนํ ปหารกํ.} \\ \csromangathabat{ตลสตฺติ อมูลญฺจ,}{สญฺจิจฺจ จ อุปสฺสุติ.} \\ ปฏิพาหน ฉนฺทญฺจ, ทพฺพญฺจ ปริณามนนฺติ.}}

\csromanmatikaanchor{mtk.92}
\csromanmatikaanchor{mtk.93}
\csromantocmarknopage{section}{9. รตนวคฺค}{mtk.92}
\bookmark[dest={mtk.92},level=1]{9. รตนวคฺค}
\csromantocmarkref{subsection}{1. อนฺเตปุรสิกฺขาปท}{mtk.93}
\bookmark[dest={mtk.93},level=2]{1. อนฺเตปุรสิกฺขาปท}
\setcsromanheadh{9. รตนวคฺค}
\csromanheader{2}{9. \textbf{รตนวคฺค 1. อนฺเตปุรสิกฺขาปท}}
\proseitem{494}{เตน สมเยน พุทฺโธ ภควา สาวตฺถิยํ วิหรติ เชตวเน อนาถปิณฺฑิกสฺส อาราเม.\csromanspacer{} เตน โข ปน สมเยน ราชา ปเสนทิ โกสโล อุยฺยานปาลํ อาณาเปสิ “คจฺฉ ภเณ อุยฺยานํ โสเธหิ อุยฺยานํ คมิสฺสามา”ติ. “เอวํ เทวา”ติ โข โส อุยฺยานปาโล รญฺโญ ปเสนทิสฺส โกสลสฺส ปฏิสฺสุตฺวา อุยฺยานํ โสเธนฺโต อทฺทส ภควนฺตํ อญฺญตรสฺมิํ รุกฺขมูเล นิสินฺนํ, ทิสฺวาน เยน ราชา ปเสนทิ โกสโล เตนุปสงฺกมิ, อุปสงฺกมิตฺวา ราชานํ ปเสนทิํ โกสลํ เอตทโวจ “สุทฺธํ เทว อุยฺยานํ, อปิ จ ภควา ตตฺถ นิสินฺโน”ติ.\csromanspacer{} โหตุ ภเณ มยํ ภควนฺตํ ปยิรุปาสิสฺสามาติ.\csromanspacer{} อถ โข ราชา ปเสนทิ โกสโล อุยฺยานํ คนฺตฺวา เยน ภควา เตนุปสงฺกมิ.\csromanspacer{} เตน โข ปน สมเยน อญฺญตโร อุปาสโก ภควนฺตํ ปยิรุปาสนฺโต นิสินฺโน โหติ.\csromanspacer{} อทฺทสา โข ราชา ปเสนทิ โกสโล ตํ อุปาสกํ ภควนฺตํ ปยิรุปาสนฺตํ นิสินฺนํ, ทิสฺวาน ภีโต อฏฺฐาสิ.\csromanspacer{} อถ โข รญฺโญ ปเสนทิสฺส โกสลสฺส เอตทโหสิ “นารหตายํ ปุริโส ปาโป \csromanfolio{206}โหตุํ ยถา ภควนฺตํ ปยิรุปาสตี”ติ, เยน ภควา เตนุปสงฺกมิ, อุปสงฺกมิตฺวา ภควนฺตํ อภิวาเทตฺวา เอกมนฺตํ นิสีทิ.\csromanspacer{} อถ โข โส อุปาสโก ภควโต คารเวน ราชานํ ปเสนทิํ โกสลํ เนว อภิวาเทสิ น ปจฺจุฏฺฐาสิ.\csromanspacer{} อถ โข ราชา ปเสนทิ โกสโล อนตฺตมโน อโหสิ “กถํ หิ นามายํ ปุริโส มยิ อาคเต เนว อภิวาเทสฺสติ น ปจฺจุฏฺเฐสฺสตี”ติ.\csromanspacer{} อถ โข ภควา ราชานํ ปเสนทิํ โกสลํ อนตฺตมนํ วิทิตฺวา ราชานํ ปเสนทิํ โกสลํ เอตทโวจ “เอโส โข มหาราช อุปาสโก พหุสฺสุโต อาคตาคโม กาเมสุ วีตราโค”ติ.\csromanspacer{} อถ โข รญฺโญ ปเสนทิสฺส โกสลสฺส เอตทโหสิ “นารหตายํ อุปาสโก โอรโก โหตุํ, ภควาปิ อิมสฺส วณฺณํ ภาสตี”ติ, ตํ อุปาสกํ เอตทโวจ “วเทยฺยาสิ อุปาสก เยน อตฺโถ”ติ.\csromanspacer{} สุฏฺฐุ เทวาติ.\csromanspacer{} อถ โข ภควา ราชานํ ปเสนทิํ โกสลํ ธมฺมิยา กถายา สนฺทสฺเสสิ สมาทเปสิ สมุตฺเตเชสิ สมฺปหํเสสิ.\csromanspacer{} อถ โข ราชา ปเสนทิ โกสโล ภควตา ธมฺมิยา กถาย สนฺทสฺสิโต สมาทปิโต สมุตฺเตชิโต สมฺปหํสิโต อุฏฺฐายาสนา ภควนฺตํ อภิวาเทตฺวา ปทกฺขิณํ กตฺวา ปกฺกามิ.}

\proseitem{495}{เตน โข ปน สมเยน ราชา ปเสนทิ โกสโล อุปริปาสาทวรคโต โหติ.\csromanspacer{} อทฺทสา โข ราชา ปเสนทิ โกสโล ตํ อุปาสกํ รถิกาย\footnote{รถิยาย (อิติปิ)} ฉตฺตปาณิํ คจฺฉนฺตํ, ทิสฺวาน ปกฺโกสาเปตฺวา เอตทโวจ “ตฺวํ กิร อุปาสก พหุสฺสุโต อาคตาคโม, สาธุ อุปาสก อมฺหากํ อิตฺถาคารํ ธมฺมํ วาเจหี”ติ.\csromanspacer{} ยมหํ\footnote{ยมฺปาหํ (สี)} เทว ชานามิ อยฺยานํ วาหสา, อยฺยาว เทวสฺส อิตฺถาคารํ ธมฺมํ วาเจสฺสนฺตีติ.\csromanspacer{} อถ โข ราชา ปเสนทิ โกสโล “สจฺจํ โข อุปาสโก อาหา”ติ เยน ภควา เตนุปสงฺกมิ, อุปสงฺกมิตฺวา ภควนฺตํ อภิวาเทตฺวา เอกมนฺตํ นิสีทิ, เอกมนฺตํ นิสินฺโน โข ราชา ปเสนทิ โกสโล ภควนฺตํ เอตทโวจ “สาธุ ภนฺเต ภควา เอกํ ภิกฺขุํ อาณาเปตุ, โย อมฺหากํ อิตฺถาคารํ ธมฺมํ วาเจสฺสตี”ติ.\csromanspacer{} อถ โข ภควา ราชานํ ปเสนทิํ โกสลํ ธมฺมิยา กถาย สนฺทสฺเสสิ ฯเปฯ ปทกฺขิณํ กตฺวา ปกฺกามิ.\csromanspacer{} อถ โข ภควา อายสฺมนฺตํ อานนฺทํ อามนฺเตสิ “เตนหานนฺท รญฺโญ อิตฺถาคารํ ธมฺมํ วาเจหี”ติ.\csromanspacer{} เอวํ ภนฺเตติ โข อายสฺมา อานนฺโท \csromanfolio{207}ภควโต ปฏิสฺสุตฺวา กาเลน กาลํ ปวิสิตฺวา รญฺโญ อิตฺถาคารํ ธมฺมํ วาเจติ.\csromanspacer{} อถ โข อายสฺมา อานนฺโท ปุพฺพณฺหสมยํ นิวาเสตฺวา ปตฺตจีวรมาทาย เยน รญฺโญ ปเสนทิสฺส โกสลสฺส นิเวสนํ เตนุปสงฺกมิ.}

\proseitem{496}{เตน โข ปน สมเยน ราชา ปเสนทิ โกสโล มลฺลิกาย เทวิยา สทฺธิํ สยนคโต โหติ.\csromanspacer{} อทฺทสา โข มลฺลิกา เทวี อายสฺมนฺตํ อานนฺทํ ทูรโตว อาคจฺฉนฺตํ, ทิสฺวาน สหสา วุฏฺฐาสิ, ปีตกมฏฺฐํ ทุสฺสํ ปภสฺสิตฺถ.\csromanspacer{} อถ โข อายสฺมา อานนฺโท ตโตว ปฏินิวตฺติตฺวา อารามํ คนฺตฺวา ภิกฺขูนํ เอตมตฺถํ อาโรเจสิ.\csromanspacer{} เย เต ภิกฺขู อปฺปิจฺฉา ฯเปฯ เต อุชฺฌายนฺติ ขิยฺยนฺติ วิปาเจนฺติ “กถํ หิ นาม อายสฺมา อานนฺโท ปุพฺเพ อปฺปฏิสํวิทิโต รญฺโญ อนฺเตปุรํ ปวิสิสฺสตี”ติ ฯเปฯ.\csromanspacer{} สจฺจํ กิร ตฺวํ อานนฺท ปุพฺเพ อปฺปฏิสํวิทิโต รญฺโญ อนฺเตปุรํ ปวิสสีติ.\csromanspacer{} สจฺจํ ภควาติ.\csromanspacer{} วิครติ พุทฺโธ ภควา ฯเปฯ.\csromanspacer{} กถํ หิ นาม ตฺวํ อานนฺท ปุพฺเพ อปฺปฏิสํวิทิโต รญฺโญ อนฺเตปุรํ ปวิสิสฺสสิ.\csromanspacer{} เนตํ อานนฺท อปฺปสนฺนานํ วา ปสาทาย ฯเปฯ วิครหิตฺวา ฯเปฯ ธมฺมิํ กถํ ภิกฺขู อามนฺเตสิ\csromandash{}}

\proseitem{497}{\csromansymbolmark{*}“ทสยิเม ภิกฺขเว อาทีนวา ราชนฺเตปุรปฺปเวสเน.\csromanspacer{} กตเม ทส, อิธ ภิกฺขเว ราชา มเหสิยา สทฺธิํ นิสินฺโน โหติ, ตตฺถ ภิกฺขุ ปวิสติ, มเหสี วา ภิกฺขุํ ทิสฺวา สิตํ ปาตุกโรติ, ภิกฺขุ วา มเหสิํ ทิสฺวา สิตํ ปาตุกโรติ, ตตฺถ รญฺโญ เอวํ โหติ “อทฺธา อิเมสํ กตํ วา กริสฺสนฺติ วา”ติ, อยํ ภิกฺขเว ปฐโม อาทีนโว ราชนฺเตปุรปฺปเวสเน.}

\prose{ปุน จปรํ ภิกฺขเว ราชา พหุกิจฺโจ พหุกรณีโย อญฺญตรํ อิตฺถิํ คนฺตฺวา นสฺสรติ, สา เตน คพฺภํ คณฺหิ, ตตฺถ รญฺโญ เอวํ โหติ “น โข อิธ อญฺโญ โกจิ ปวิสติ อญฺญตฺร ปพฺพชิเตน, สิยา นุ โข ปพฺพชิตสฺส กมฺมนฺ”ติ, อยํ ภิกฺขเว ทุติโย อาทีนโว ราชนฺเตปุรปฺปเวสเน.}

\prose{ปุน จปรํ ภิกฺขเว รญฺโญ อนฺเตปุเร อญฺญตรํ รตนํ นสฺสติ, ตตฺถ รญฺโญ เอวํ โหติ “น โข อิธ อญฺโญ โกจิ ปวิสติ อญฺญตฺร ปพฺพชิเตน, สิยา นุ โข ปพฺพชิตสฺส กมฺมนฺ”ติ, อยํ ภิกฺขเว ตติโย อาทีนโว ราชนฺเตปุรปฺปเวสเน.}

\prose{\csromanfolio{208}ปุน จปรํ ภิกฺขเว รญฺโญ อนฺเตปุเร อพฺภนฺตรา คุยฺหมนฺตา พหิทฺธา สมฺเภทํ คจฺฉนฺติ, ตตฺถ รญฺโญ เอวํ โหติ “น โข อิธ อญฺโญ โกจิ ปวิสติ อญฺญตฺร ปพฺพชิเตน, สิยา นุ โข ปพฺพชิตสฺส กมฺมนฺ”ติ, อยํ ภิกฺขเว จตุตฺโถ อาทีนโว ราชนฺเตปุรปฺปเวสเน.}

\prose{ปุน จปรํ ภิกฺขเว รญฺโญ อนฺเตปุเร ปุตฺโต วา ปิตรํ ปตฺเถติ, ปิตา วา ปุตฺตํ ปตฺเถติ, เตสํ เอวํ โหติ “น โข อิธ อญฺโญ โกจิ ปวิสติ อญฺญตฺร ปพฺพชิเตน, สิยา นุ โข ปพฺพชิตสฺส กมฺมนฺ”ติ, อยํ ภิกฺขเว ปญฺจโม อาทีนโว ราชนฺเตปุรปฺปเวสเน.}

\prose{ปุน จปรํ ภิกฺขเว ราชา นีจฏฺฐานิยํ อุจฺเจ ฐาเน ฐเปติ, เยสํ ตํ อมนาปํ, เตสํ เอวํ โหติ “ราชา โข ปพฺพชิเตน สํสฏฺโฐ, สิยา นุ โข ปพฺพชิตสฺส กมฺมนฺ”ติ, อยํ ภิกฺขเว ฉฏฺโฐ อาทีนโว ราชนฺเตปุรปฺปเวสเน.}

\prose{ปุน จปรํ ภิกฺขเว ราชา อุจฺจฏฺฐานิยํ นีเจ ฐาเน ฐเปติ, เยสํ ตํ อมนาปํ, เตสํ เอวํ โหติ “ราชา โข ปพฺพชิเตน สํสฏฺโฐ, สิยา นุ โข ปพฺพชิตสฺส กมฺมนฺ”ติ, อยํ ภิกฺขเว สตฺตโม อาทีนโว ราชนฺเตปุรปฺปเวสเน.}

\prose{ปุน จปรํ ภิกฺขเว ราชา อกาเล เสนํ อุยฺโยเชติ, เยสํ ตํ อมนาปํ, เตสํ เอวํ โหติ “ราชา โข ปพฺพชิเตน สํสฏฺโฐ, สิยา นุ โข ปพฺพชิตสฺส กมฺมนฺ”ติ, อยํ ภิกฺขเว อฏฺฐโม อาทีนโว ราชนฺเต ปุรปฺปเวสเน.}

\prose{ปุน จปรํ ภิกฺขเว ราชา กาเล เสนํ อุยฺโยเชตฺวา อนฺตรามคฺคโต นิวตฺตาเปติ, เยสํ ตํ อมนาปํ, เตสํ เอวํ โหติ “ราชา โข ปพฺพชิเตน สํสฏฺโฐ, สิยา นุ โข ปพฺพชิตสฺส กมฺมนฺ”ติ, อยํ ภิกฺขเว นวโม อาทีนโว ราชนฺเตปุรปฺปเวสเน.}

\prose{ปุน จปรํ ภิกฺขเว รญฺโญ ราชนฺเตรํปุรํ หตฺถิสมฺมทฺทํ อสฺสสมฺมทฺทํ รถสมฺมทฺทํ รชนียานิ\footnote{รชฺชนียานิ (ก)} รูปสทฺทคนฺธรสโผฏฺฐพฺพานิ, ยานิ น ปพฺพชิตสฺส สารุปฺปานิ, อยํ ภิกฺขเว ทสโม อาทีนโว ราชนฺเตปุรปฺปเวสเน.\csromanspacer{} อิเม โข ภิกฺขเว ทส อาทีนวา ราชนฺเตปุรปฺปเวสเน”ติ.\csromanspacer{} อถ โข ภควา อายสฺมนฺตํ \csromanfolio{209}อานนฺทํ อเนกปริยาเยน วิครหิตฺวา ทุพฺภรตาย ฯเปฯ.\csromanspacer{} เอวญฺจ ปน ภิกฺขเว อิมํ สิกฺขาปทํ อุทฺทิเสยฺยาถ\csromandash{}}

\proseitem{498}{\textbf{“โย ปน ภิกฺขุ รญฺโญ ขตฺติยสฺส มุทฺธาวสิตฺตสฺส}\footnote{มุทฺธาภิสิตฺตสฺส (สฺยา)} \textbf{อนิกฺขนฺตราชเก อนิคฺคตรตนเก ปุพฺเพ อปฺปฏิสํวิทิโต อินฺทขีลํ อติกฺกาเมยฺย, ปาจิตฺติยนฺ”}ติ.}

\proseitem{499}{\textbf{โย ปนา}ติ โย ยาทิโส ฯเปฯ.\csromanspacer{} \textbf{ภิกฺขู}ติ ฯเปฯ อยํ อิมสฺมิํ อตฺเถ อธิปฺเปโต ภิกฺขูติ.}

\prose{\textbf{ขตฺติโย} นาม อุภโต สุชาโต โหติ มาติโต จ ปิติโต จ สํสุทฺธคหณิโก, ยาว สตฺตมา ปิตามหยุคา อกฺขิตฺโต อนุปกุฏฺโฐ ชาติวาเทน.}

\prose{\textbf{มุทฺธาวสิตฺโต} นาม ขตฺติยาภิเสเกน อภิสิตฺโต โหติ.}

\prose{\textbf{อนิกฺขนฺตราชเก}ติ ราชา สยนิฆรา อนิกฺขนฺโต โหติ.}

\prose{\textbf{อนิคฺคตรตนเก}ติ มเหสี สยนิฆรา อนิกฺขนฺตา โหติ, อุโภ วา อนิกฺขนฺตา โหนฺติ.}

\prose{\textbf{ปุพฺเพ อปฺปฏิสํวิทิโต}ติ ปุพฺเพ อนามนฺเตตฺวา.}

\prose{\textbf{อินฺทขีโล} นาม สยนิฆรสฺส อุมฺมาโร วุจฺจติ.}

\prose{สยนิฆรํ นาม ยตฺถ กตฺถจิ รญฺโญ สยนํ ปญฺญตฺตํ โหติ, อนฺตมโส สาณิปาการ\-ปริกฺขิตฺตมฺปิ.}

\prose{\textbf{อินฺทขีลํ อติกฺกาเมยฺยา}ติ ปฐมํ ปาทํ อุมฺมารํ อติกฺกาเมติ, อาปตฺติ ทุกฺกฏสฺส.\csromanspacer{} ทุติยํ ปาทํ อติกฺกาเมติ, อาปตฺติ ปาจิตฺติยสฺส.}

\proseitem{500}{อปฺปฏิสํวิทิเต อปฺปฏิสํวิทิตสญฺญี อินฺทขีลํ อติกฺกาเมติ, อาปตฺติ ปาจิตฺติยสฺส.\csromanspacer{} อปฺปฏิสํวิทิเต เวมติโก อินฺทขีลํ อติกฺกาเมติ, อาปตฺติ ปาจิตฺติยสฺส.\csromanspacer{} อปฺปฏิสํวิทิเต ปฏิสํวิทิตสญฺญี อินฺทขีลํ อติกฺกาเมติ, อาปตฺติ ปาจิตฺติยสฺส.}

\prose{ปฏิสํวิทิเต อปฺปฏิสํวิทิตสญฺญี, อาปตฺติ ทุกฺกฏสฺส.\csromanspacer{} ปฏิสํวิทิเต เวมติโก, อาปตฺติ ทุกฺกฏสฺส.\csromanspacer{} ปฏิสํวิทิเต ปฏิสํวิทิตสญฺญี, อนาปตฺติ.}

\proseitemwithcloserruled{501}{\csromanfolio{210}อนาปตฺติ ปฏิสํวิทิเต, น ขตฺติโย โหติ, น ขตฺติยาภิเสเกน อภิสิตฺโต โหติ, ราชา สยนิฆรา นิกฺขนฺโต โหติ, มเหสี สยนิฆรา นิกฺขนฺตา โหติ, อุโภ วา นิกฺขนฺตา โหนฺติ, น สยนิฆเร อุมฺมตฺตกสฺส อาทิกมฺมิกสฺสาติ.}{อนฺเตปุรสิกฺขาปทํ นิฏฺฐิตํ ปฐมํ.}
\csromanmatikaanchor{mtk.94}
\csromantocmarkref{subsection}{2. รตนสิกฺขาปท}{mtk.94}
\bookmark[dest={mtk.94},level=2]{2. รตนสิกฺขาปท}
\csromanheader{2}{9. \textbf{รตนวคฺค 2. รตนสิกฺขาปท}}
\proseitem{502}{เตน สมเยน พุทฺโธ ภควา สาวตฺถิยํ วิหรติ เชตวเน อนาถปิณฺฑิกสฺส อาราเม.\csromanspacer{} เตน โข ปน สมเยน อญฺญตโร ภิกฺขุ อจิรวติยา นทิยา นหายติ.\csromanspacer{} อญฺญตโรปิ พฺราหฺมโณ ปญฺจสตานํ ถวิกํ ถเล นิกฺขิปิตฺวา อจิรวติยา นทิยา นหายนฺโต วิสฺสริตฺวา อคมาสิ.\csromanspacer{} อถ โข โส ภิกฺขุ “ตสฺสายํ พฺราหฺมณสฺส ถวิกา มา อิธ นสฺสี”ติ อคฺคเหสิ.\csromanspacer{} อถ โข โส พฺราหฺมโณ สริตฺวา ตุริโต อาธาวิตฺวา ตํ ภิกฺขุํ เอตทโวจ “อปิ เม โภ ถวิกํ ปสฺเสยฺยาสี”ติ. “หนฺท พฺราหฺมณา”ติ อทาสิ.\csromanspacer{} อถ โข ตสฺส พฺราหฺมณสฺส เอตทโหสิ “เกน นุ โข อหํ อุปาเยน อิมสฺส ภิกฺขุโน ปุณฺณปตฺตํ น ทเทยฺยนฺ”ติ. “น เม โภ ปญฺจสตานิ, สหสฺสํ เม”ติ ปลิพุนฺเธตฺวา มุญฺจิ.\csromanspacer{} อถ โข โส ภิกฺขุ อารามํ คนฺตฺวา ภิกฺขูนํ เอตมตฺถํ อาโรเจสิ.\csromanspacer{} เย เต ภิกฺขู อปฺปิจฺฉา ฯเปฯ เต อุชฺฌายนฺติ ขิยฺยนฺติ วิปาเจนฺติ “กถํ หิ นาม ภิกฺขุ รตนํ อุคฺคเหสฺสตี”ติ ฯเปฯ.\csromanspacer{} สจฺจํ กิร ตฺวํ ภิกฺขุ รตนํ อุคฺคเหสีติ.\csromanspacer{} สจฺจํ ภควาติ.\csromanspacer{} วิครหิ พุทฺโธ ภควา ฯเปฯ.\csromanspacer{} กถํ หิ นาม ตฺวํ โมฆปุริส รตนํ อุคฺคเหสฺสสิ.\csromanspacer{} เนตํ โมฆปุริส อปฺปสนฺนานํ วา ปสาทาย ฯเปฯ.\csromanspacer{} เอวญฺจ ปน ภิกฺขเว อิมํ สิกฺขาปทํ อุทฺทิเสยฺยาถ\csromandash{}}

\prose{\textbf{“โย ปน ภิกฺขุ รตนํ วา รตนสมฺมตํ วา อุคฺคเณฺหยฺย วา อุคฺคณฺหาเปยฺย วา, ปาจิตฺติยนฺ”}ติ.}

\prose{เอวญฺจิทํ ภควตา ภิกฺขูนํ สิกฺขาปทํ ปญฺญตฺตํ โหติ.}

\proseitem{503}{เตน โข ปน สมเยน สาวตฺถิยา อุสฺสโว โหติ, มนุสฺสา อลงฺกตปฺปฏิยตฺตา อุยฺยานํ คจฺฉนฺติ, วิสาขาปิ มิคารมาตา \csromanfolio{211}อลงฺกตปฺปฏิยตฺตา “อุยฺยานํ คมิสฺสามี”ติ คามโต นิกฺขมิตฺวา “กฺยาหํ กริสฺสามิ อุยฺยานํ คนฺตฺวา ยนฺนูนาหํ ภควนฺตํ ปยิรุปาเสยฺยนฺ”ติ อาภรณํ โอมุญฺจิตฺวา อุตฺตราสงฺเคน ภณฺฑิกํ พนฺธิตฺวา ทาสิยา อทาสิ “หนฺท เช อิมํ ภณฺฑิกํ คณฺหาหี”ติ.\csromanspacer{} อถ โข วิสาขา มิคารมาตา เยน ภควา เตนุปสงฺกมิ, อุปสงฺกมิตฺวา ภควนฺตํ อภิวาเทตฺวา เอกมนฺตํ นิสีทิ, เอกมนฺตํ นิสินฺนํ โข วิสาขํ มิคารมาตรํ ภควา ธมฺมิยา กถาย สนฺทสฺเสสิ สมาทเปสิ สมุตฺเตเชสิ สมฺปหํเสสิ.\csromanspacer{} อถ โข วิสาขา มิคารมาตา ภควตา ธมฺมิยา กถาย สนฺทสฺสิตา สมาทปิตา สมุตฺเตชิตา สมฺปหํสิตา อุฏฺฐายาสนา ภควนฺตํ อภิวาเทตฺวา ปทกฺขิณํ กตฺวา ปกฺกามิ.\csromanspacer{} อถ โข สา ทาสี ตํ ภณฺฑิกํ วิสฺสริตฺวา อคมาสิ.\csromanspacer{} ภิกฺขู ปสฺสิตฺวา ภควโต เอตมตฺถํ อาโรเจสุํ. “\textbf{เตน หิ ภิกฺขเว อุคฺคเหตฺวา นิกฺขิปถา”ติ.}\csromanspacer{}\textbf{ อถ โข ภควา เอตสฺมิํ} \textbf{นิทาเน เอตสฺมิํ ปกรเณ ธมฺมิํ กถํ กตฺวา ภิกฺขู อามนฺเตสิ} \textbf{อนุชานามิ ภิกฺขเว รตนํ วา รตนสมฺมตํ วา อชฺฌาราเม อุคฺคเหตฺวา} \textbf{วา อุคฺคหาเปตฺวา วา นิกฺขิปิตุํ “ยสฺส ภวิสฺสติ, โส หริสฺสตี”ติ, เอวญฺจ} \textbf{ปน ภิกฺขเว อิมํ สิกฺขาปทํ อุทฺทิเสยฺยาถ\csromandash{}}}

\prose{\textbf{“โย ปน ภิกฺขุ รตนํ วา รตนสมฺมตํ วา อญฺญตฺร อชฺฌารามา อุคฺคเณฺหยฺย วา อุคฺคณฺหาเปยฺย วา, ปาจิตฺติยนฺ”}ติ.}

\prose{เอวญฺจิทํ ภควตา ภิกฺขูนํ สิกฺขาปทํ ปญฺญตฺตํ โหติ.}

\proseitem{504}{เตน โข ปน สมเยน กาสีสุ ชนปเท อนาถปิณฺฑิกสฺส คหปติสฺส กมฺมนฺตคาโม โหติ, เตน จ คหปตินา อนฺเตวาสี อาณตฺโต โหติ “สเจ ภทนฺตา อาคจฺฉนฺติ, ภตฺตํ กเรยฺยาสี”ติ.\csromanspacer{} เตน โข ปน สมเยน สมฺพหุลา ภิกฺขู กาสีสุ ชนปเท จาริกํ จรมานา เยน อนาถปิณฺฑิกสฺส คหปติสฺส กมฺมนฺตคาโม เตนุปสงฺกมิํสุ.\csromanspacer{} อทฺทสา โข โส ปุริโส เต ภิกฺขู ทูรโตว อาคจฺฉนฺเต, ทิสฺวาน เยน เต ภิกฺขู เตนุปสงฺกมิ, อุปสงฺกมิตฺวา เต ภิกฺขู อภิวาเทตฺวา เอตทโวจ “อธิวาเสนฺตุ ภนฺเต อยฺยา สฺวาตนาย คหปติโน ภตฺตนฺ”ติ.\csromanspacer{} อธิวาเสสุํ โข เต ภิกฺขู ตุณฺหีภาเวน.\csromanspacer{} อถ โข โส ปุริโส ตสฺสา รตฺติยา อจฺจเยน ปณีตํ ขาทนียํ โภชนียํ ปฏิยาทาเปตฺวา กาลํ อาโรจาเปตฺวา องฺคุลิมุทฺทิกํ โอมุญฺจิตฺวา เต ภิกฺขู ภตฺเตน ปริวิสิตฺวา “อยฺยา ภุญฺชิตฺวา คจฺฉนฺตุ, อหมฺปิ กมฺมนฺตํ คมิสฺสามี”ติ \csromanfolio{212}องฺคุลิมุทฺทิกํ วิสฺสริตฺวา อคมาสิ.\csromanspacer{} ภิกฺขู ปสฺสิตฺวา “สเจ มยํ คมิสฺสาม นสฺสิสฺสตายํ องฺคุลิมุทฺทิกา”ติ ตตฺเถว อจฺฉิํสุ.\csromanspacer{} อถ โข โส \textbf{ปุริโส กมฺมนฺตา อาคจฺฉนฺโต เต ภิกฺขู ปสฺสิตฺวา เอตทโวจ “กิสฺส ภนฺเต} \textbf{อยฺยา อิเธว อจฺฉนฺตี”ติ.}\csromanspacer{}\textbf{ อถ โข เต ภิกฺขู ตสฺส ปุริสสฺส เอตมตฺถํ} \textbf{อาโรเจตฺวา สาวตฺถิํ คนฺตฺวา ภิกฺขูนํ เอตมตฺถํ อาโรเจสุํ.}\csromanspacer{}\textbf{ ภิกฺขู} ภควโต เอตมตฺถํ อาโรเจสุํ.\csromanspacer{} อถ โข ภควา เอตสฺมิํ นิทาเน เอตสฺมิํ ปกรเณ ธมฺมิํ กถํ กตฺวา ภิกฺขู อามนฺเตสิ อนุชานามิ ภิกฺขเว รตนํ วา รตนสมฺมตํ วา อชฺฌาราเม วา อชฺฌาวสเถ วา อุคฺคเหตฺวา วา อุคฺคหาเปตฺวา วา นิกฺขิปิตุํ “ยสฺส ภวิสฺสติ โส หริสฺสตี”ติ.\csromanspacer{} เอวญฺจ ปน ภิกฺขเว อิมํ สิกฺขาปทํ อุทฺทิเสยฺยาถ\csromandash{}}

\proseitem{505}{“โย ปน ภิกฺขุ รตนํ วา รตนสมฺมตํ วา อญฺญตฺร อชฺฌารามา วา อชฺฌาวสถา วา อุคฺคเณฺหยฺย วา อุคฺคณฺหาเปยฺย วา, ปาจิตฺติยํ.\csromanspacer{}\textbf{ รตนํ วา ปน ภิกฺขุนา รตนสมฺมตํ วา อชฺฌาราเม วา อชฺฌาวสเถ วา อุคฺคเหตฺวา วา อุคฺคหาเปตฺวา วา นิกฺขิปิตพฺพํ ‘ยสฺส ภวิสฺสติ, โส หริสฺสตี’ติ, อยํ ตตฺถ สามีจี”}ติ.}

\proseitem{506}{\textbf{โย ปนา}ติ โย ยาทิโส ฯเปฯ.\csromanspacer{} \textbf{ภิกฺขู}ติ ฯเปฯ อยํ อิมสฺมิํ อตฺเถ อธิปฺเปโต ภิกฺขูติ.}

\prose{\textbf{รตนํ} นาม มุตฺตา มณิ เวฬุริโย สงฺโข สิลา ปวาลํ รชตํ ชาตรูปํ โลหิตงฺโก\footnote{โลหิตโก (?)} มสารคลฺลํ.}

\prose{\textbf{รตนสมฺมตํ} นาม ยํ มนุสฺสานํ อุปโภคปริโภคํ, เอตํ รตนสมฺมตํ นาม.}

\prose{\textbf{อญฺญตฺร อชฺฌารามา วา อชฺฌาวสถาวา}ติ ฐเปตฺวา อชฺฌารามํ อชฺฌาวสถํ.}

\prose{\textbf{อชฺฌาราโม} นาม ปริกฺขิตฺตสฺส อารามสฺส อนฺโต อาราโม, อปริกฺขิตฺตสฺส อุปจาโร.}

\prose{\textbf{อชฺฌาวสโถ} นาม ปริกฺขิตฺตสฺส อาวสถสฺส อนฺโต อาวสโถ, อปริกฺขิตฺตสฺส อุปจาโร.}

\prose{\csromanfolio{213}\textbf{อุคฺคเณฺหยฺยา}ติ สยํ คณฺหาติ, อาปตฺติ ปาจิตฺติยสฺส.}

\prose{\textbf{อุคฺคณฺหาเปยฺยา}ติ อญฺญํ คาหาเปติ, อาปตฺติ ปาจิตฺติยสฺส.}

\prose{\textbf{รตนํ วา ปน ภิกฺขุนา รตนสมฺมตํ วา อชฺฌาราเม วา อชฺฌาวสเถ วา อุคฺคเหตฺวา วา อุคฺคหาเปตฺวา วา นิกฺขิปิตพฺพนฺ}ติ รูเปน วา นิมิตฺเตน วา สญฺญาณํ กตฺวา นิกฺขิปิตฺวา อาจิกฺขิตพฺพํ “ยสฺส ภณฺฑํ นฏฺฐํ, โส อาคจฺฉตู”ติ.\csromanspacer{} สเจ ตตฺถ อาคจฺฉติ, โส วตฺตพฺโพ “อาวุโส กีทิสํ เต ภณฺฑนฺ”ติ.\csromanspacer{} สเจ รูเปน วา นิมิตฺเตน วา สมฺปาเทติ, ทาตพฺพํ.\csromanspacer{} โน เจ สมฺปาเทติ, “วิจินาหิ อาวุโส”ติ วตฺตพฺโพ.\csromanspacer{} ตมฺหา อาวาสา ปกฺกมนฺเตน เย ตตฺถ โหนฺติ ภิกฺขู ปติรูปา, เตสํ หตฺเถ นิกฺขิปิตฺวา ปกฺกมิตพฺพํ.\csromanspacer{} โน เจ โหนฺติ ภิกฺขู ปติรูปา, เย ตตฺถ โหนฺติ คหปติกา ปติรูปา, เตสํ หตฺเถ นิกฺขิปิตฺวา ปกฺกมิตพฺพํ.}

\prose{\textbf{อยํ ตตฺถ สามีจี}ติ อยํ ตตฺถ อนุธมฺมตา.}

\proseitemwithcloserruled{507}{อนาปตฺติ รตนํ วา รตนสมฺมตํ วา อชฺฌาราเม วา อชฺฌาวสเถ วา อุคฺคเหตฺวา วา อุคฺคหาเปตฺวา วา นิกฺขิปติ “ยสฺส ภวิสฺสติ, โส หริสฺสตี”ติ, รตนสมฺมตํ วิสฺสาสํ คณฺหาติ, ตาวกาลิกํ คณฺหาติ, ปํสุกูลสญฺญิสฺส อุมฺมตฺตกสฺส อาทิกมฺมิกสฺสาติ.}{รตนสิกฺขาปทํ นิฏฺฐิตํ ทุติยํ.}
\csromanmatikaanchor{mtk.95}
\csromantocmarkref{subsection}{3. วิกาลคามปฺปวิสนสิกฺขาปท}{mtk.95}
\bookmark[dest={mtk.95},level=2]{3. วิกาลคามปฺปวิสนสิกฺขาปท}
\csromanheader{2}{9. \textbf{รตนวคฺค 3. วิกาลคามปฺปวิสนสิกฺขาปท}}
\proseitem{508}{เตน สมเยน พุทฺโธ ภควา สาวตฺถิยํ วิหรติ เชตวเน อนาถปิณฺฑิกสฺส อาราเม.\csromanspacer{} เตน โข ปน สมเยน ฉพฺพคฺคิยา ภิกฺขู วิกาเล คามํ ปวิสิตฺวา สภายํ นิสีทิตฺวา อเนกวิหิตํ ติรจฺฉานกถํ กเถนฺติ, เสยฺยถิทํ,\csromansymbolfootnote{*}{อิมา ติรจฺฉานกถาโย ที 1. 7; ม 2. 181; สํ 3. 368; อํ 3. 358 ปิฏฺฐาทีสุปิ อาคตา.}ราชกถํ โจรกถํ มหามตฺตกถํ เสนากถํ ภยกถํ ยุทฺธกถํ อนฺนกถํ ปานกถํ วตฺถกถํ สยนกถํ มาลากถํ คนฺธกถํ ญาติกถํ ยานกถํ คามกถํ นิคมกถํ นครกถํ ชนปทกถํ อิตฺถิกถํ\footnote{อิตฺถิกถํ ปุริสกถํ (ก) มชฺฌิมปณฺณาสฏฺฐกถา 156 ปิฏฺเฐ โอโลเกตพฺพา.} สูรกถํ วิสิขากถํ กุมฺภฏฺฐานกถํ ปุพฺพเปตกถํ \csromanfolio{214}\textbf{นานตฺตกถํ โลกกฺขายิกํ สมุทฺทกฺขายิกํ อิติภวาภวกถํ อิติ} \textbf{วา.}\csromanspacer{}\textbf{ มนุสฺสา อุชฺฌายนฺติ ขิยฺยนฺติ วิปาเจนฺติ “กถํ หิ นาม สมณา} สกฺยปุตฺติยา วิกาเล คามํ ปวิสิตฺวา สภายํ นิสีทิตฺวา อเนกวิหิตํ ติรจฺฉานกถํ กเถสฺสนฺติ, เสยฺยถิทํ, ราชกถํ โจรกถํ ฯเปฯ อิติภวาภวกถํ อิติ วา เสยฺยถาปิ คิหี กามโภคิโน”ติ.}

\prose{อสฺโสสุํ โข ภิกฺขู เตสํ มนุสฺสานํ อุชฺฌายนฺตานํ ขิยฺยนฺตานํ วิปาเจนฺตานํ.\csromanspacer{} เย เต ภิกฺขู อปฺปิจฺฉา ฯเปฯ เต อุชฺฌายนฺติ ขิยฺยนฺติ วิปาเจนฺติ “กถํ หิ นาม ฉพฺพคฺคิยา ภิกฺขู วิกาเล คามํ ปวิสิตฺวา สภายํ นิสีทิตฺวา อเนกวิหิตํ ติรจฺฉานกถํ กเถสฺสนฺติ, เสยฺยถิทํ, ราชกถํ โจรกถํ ฯเปฯ อิติภวาภวกถํ อิติ วา”ติ ฯเปฯ.\csromanspacer{} สจฺจํ กิร ตุเมฺห ภิกฺขเว วิกาเล คามํ ปวิสิตฺวา สภายํ นิสีทิตฺวา อเนกวิหิตํ ติรจฺฉานกถํ กเถถ, เสยฺยถิทํ, ราชกถํ โจรกถํ ฯเปฯ อิติภวาภวกถํ อิติ วาติ.\csromanspacer{} สจฺจํ ภควาติ.\csromanspacer{} วิครหิ พุทฺโธ ภควา ฯเปฯ.\csromanspacer{} กถํ หิ นาม ตุเมฺห โมฆปุริสา วิกาเล คามํ ปวิสิตฺวา สภายํ นิสีทิตฺวา อเนกวิหิตํ ติรจฺฉานกถํ กเถสฺสถ, เสยฺยถิทํ, ราชกถํ โจรกถํ ฯเปฯ อิติภวาภวกถํ อิติ วา.\csromanspacer{} เนตํ โมฆปุริสา อปฺปสนฺนานํ วา ปสาทาย ฯเปฯ.\csromanspacer{} เอวญฺจ ปน ภิกฺขเว อิมํ สิกฺขาปทํ อุทฺทิเสยฺยาถ\csromandash{}}

\prose{\textbf{“โย ปน ภิกฺขุ วิกาเล คามํ ปวิเสยฺย, ปาจิตฺติยนฺ”}ติ.}

\prose{เอวญฺจิทํ ภควตา ภิกฺขูนํ สิกฺขาปทํ ปญฺญตฺตํ โหติ.}

\proseitem{509}{เตน โข ปน สมเยน สมฺพหุลา ภิกฺขู โกสเลสุ ชนปเท สาวตฺถิํ คจฺฉนฺตา สายํ อญฺญตรํ คามํ อุปคจฺฉิํสุ.\csromanspacer{} มนุสฺสา เต ภิกฺขู ปสฺสิตฺวา เอตทโวจุํ “ปวิสถ ภนฺเต”ติ.\csromanspacer{} อถ โข เต ภิกฺขู “ภควตา ปฏิกฺขิตฺตํ วิกาเล คามํ ปวิสิตุนฺ”ติ กุกฺกุจฺจายนฺตา น ปวิสิํสุ.\csromanspacer{} โจรา เต ภิกฺขู อจฺฉินฺทิํสุ.\csromanspacer{} อถ โข เต ภิกฺขู สาวตฺถิํ คนฺตฺวา ภิกฺขูนํ เอตมตฺถํ อาโรเจสุํ.\csromanspacer{} ภิกฺขู ภควโต เอตมตฺถํ อาโรเจสุํ.\csromanspacer{} อถ โข ภควา เอตสฺมิํ นิทาเน เอตสฺมิํ ปกรเณ ธมฺมิํ กถํ กตฺวา ภิกฺขู อามนฺเตสิ อนุชานามิ ภิกฺขเว อาปุจฺฉา วิกาเล คามํ ปวิสิตุํ.\csromanspacer{} เอวญฺจ ปน ภิกฺขเว อิมํ สิกฺขาปทํ อุทฺทิเสยฺยาถ\csromandash{}}

\prose{\csromanfolio{215}\textbf{“โย ปน ภิกฺขุ อนาปุจฺฉา วิกาเล คามํ ปวิเสยฺย, ปาจิตฺติยนฺ”}ติ.}

\prose{เอวญฺจิทํ ภควตา ภิกฺขูนํ สิกฺขาปทํ ปญฺญตฺตํ โหติ.}

\proseitem{510}{เตน โข ปน สมเยน อญฺญตโร ภิกฺขุ โกสเลสุ ชนปเท สาวตฺถิํ คจฺฉนฺโต สายํ อญฺญตรํ คามํ อุปคจฺฉิ.\csromanspacer{} มนุสฺสา ตํ ภิกฺขุํ ปสฺสิตฺวา เอตทโวจุํ “ปวิสถ ภนฺเต”ติ.\csromanspacer{} อถ โข โส ภิกฺขุ “ภควตา ปฏิกฺขิตฺตํ อนาปุจฺฉา วิกาเล คามํ ปวิสิตุนฺ”ติ กุกฺกุจฺจายนฺโต น ปาวิสิ.\csromanspacer{} โจรา ตํ ภิกฺขุํ อจฺฉินฺทิํสุ.\csromanspacer{} อถ โข โส ภิกฺขุ สาวตฺถิํ คนฺตฺวา ภิกฺขูนํ เอตมตฺถํ อาโรเจสิ.\csromanspacer{} ภิกฺขู ภควโต เอตมตฺถํ อาโรเจสุํ.\csromanspacer{} อถ โข ภควา เอตสฺมิํ นิทาเน เอตสฺมิํ ปกรเณ ธมฺมิํ กถํ กตฺวา ภิกฺขู อามนฺเตสิ อนุชานามิ ภิกฺขเว สนฺตํ ภิกฺขุํ อาปุจฺฉา วิกาเล คามํ ปวิสิตุํ.\csromanspacer{} เอวญฺจ ปน ภิกฺขเว อิมํ สิกฺขาปทํ อุทฺทิเสยฺยาถ\csromandash{}}

\prose{\textbf{“โย ปน ภิกฺขุ สนฺตํ ภิกฺขุํ อนาปุจฺฉา วิกาเล คามํ ปวิเสยฺย, ปาจิตฺติยนฺ”}ติ.}

\prose{เอวญฺจิทํ ภควตา ภิกฺขูนํ สิกฺขาปทํ ปญฺญตฺตํ โหติ.}

\proseitem{511}{เตน โข ปน สมเยน อญฺญตโร ภิกฺขุ อหินา ทฏฺโฐ โหติ.\csromanspacer{} อญฺญตโร ภิกฺขุ “อคฺคิํ อาหริสฺสามี”ติ คามํ คจฺฉติ.\csromanspacer{} อถ โข โส ภิกฺขุ “ภควตา ปฏิกฺขิตฺตํ สนฺตํ ภิกฺขุํ อนาปุจฺฉา วิกาเล คามํ ปวิสิตุนฺ”ติ กุกฺกุจฺจายนฺโต น ปาวิสิ ฯเปฯ ภควโต เอตมตฺถํ อาโรเจสุํ.\csromanspacer{} อถ โข ภควา เอตสฺมิํ นิทาเน เอตสฺมิํ ปกรเณ ธมฺมิํ กถํ กตฺวา ภิกฺขู อามนฺเตสิ อนุชานามิ ภิกฺขเว ตถารูเป อจฺจายิเก กรณีเย สนฺตํ ภิกฺขุํ อนาปุจฺฉา วิกาเล คามํ ปวิสิตุํ.\csromanspacer{} เอวญฺจ ปน ภิกฺขเว อิมํ สิกฺขาปทํ อุทฺทิเสยฺยาถ\csromandash{}}

\proseitem{512}{“\textbf{โย ปน ภิกฺขุ สนฺตํ ภิกฺขุํ อนาปุจฺฉา วิกาเล คามํ} ปวิเสยฺย อญฺญตฺร ตถารูปา อจฺจายิกา กรณียา, ปาจิตฺติยนฺ”ติ.}

\proseitem{513}{\textbf{โย ปนา}ติ โย ยาทิโส ฯเปฯ.\csromanspacer{} \textbf{ภิกฺขู}ติ ฯเปฯ อยํ อิมสฺมิํ อตฺเถ อธิปฺเปโต ภิกฺขูติ.}

\prose{\csromanfolio{216}\textbf{สนฺโต} นาม ภิกฺขุ สกฺกา โหติ อาปุจฺฉา ปวิสิตุํ.}

\prose{\textbf{อสนฺโต} นาม ภิกฺขุ น สกฺกา โหติ อาปุจฺฉา ปวิสิตุํ.}

\prose{\textbf{วิกาโล} นาม มชฺฌนฺหิเก วีติวตฺเต ยาว อรุณุคฺคมนา.}

\prose{\textbf{คามํ ปวิเสยฺยา}ติ ปริกฺขิตฺตสฺส คามสฺส ปริกฺเขปํ อติกฺกมนฺตสฺส อาปตฺติ ปาจิตฺติยสฺส.\csromanspacer{} อปริกฺขิตฺตสฺส คามสฺส อุปจารํ โอกฺกมนฺตสฺส อาปตฺติ ปาจิตฺติยสฺส.}

\prose{\textbf{อญฺญตฺร ตถารูปา อจฺจายิกา กรณียา}ติ ฐเปตฺวา ตถารูปํ อจฺจายิกํ กรณียํ.}

\proseitem{514}{วิกาเล วิกาลสญฺญี สนฺตํ ภิกฺขุํ อนาปุจฺฉา คามํ ปวิสติ อญฺญตฺร ตถารูปา อจฺจายิกา กรณียา, อาปตฺติ ปาจิตฺติยสฺส.\csromanspacer{} วิกาเล เวมติโก สนฺตํ ภิกฺขุํ อนาปุจฺฉา คามํ ปวิสติ อญฺญตฺร ตถารูปา อจฺจายิกา กรณียา, อาปตฺติ ปาจิตฺติยสฺส.\csromanspacer{} วิกาเล กาลสญฺญี สนฺตํ ภิกฺขุํ อนาปุจฺฉา คามํ ปวิสติ อญฺญตฺร ตถารูปา อจฺจายิกา กรณียา, อาปตฺติ ปาจิตฺติยสฺส.}

\prose{กาเล วิกาลสญฺญี, อาปตฺติ ทุกฺกฏสฺส.\csromanspacer{} กาเล เวมติโก, อาปตฺติ ทุกฺกฏสฺส.\csromanspacer{} กาเล กาลสญฺญี, อนาปตฺติ.}

\proseitemwithcloserruled{515}{อนาปตฺติ ตถารูเป อจฺจายิเก กรณีเย, สนฺตํ ภิกฺขุํ อาปุจฺฉา ปวิสติ, อสนฺตํ ภิกฺขุํ อนาปุจฺฉา ปวิสติ, อนฺตรารามํ คจฺฉติ, ภิกฺขุนุปสฺสยํ คจฺฉติ, ติตฺถิยเสยฺยํ คจฺฉติ, ปฏิกฺกมนํ คจฺฉติ, คาเมน มคฺโค โหติ, อาปทาสุ อุมฺมตฺตกสฺส อาทิกมฺมิกสฺสาติ.}{วิกาลคามปฺปวิสนสิกฺขาปทํ นิฏฺฐิตํ ตติยํ.}
\csromanmatikaanchor{mtk.96}
\csromantocmarkref{subsection}{4. สูจิฆรสิกฺขาปท}{mtk.96}
\bookmark[dest={mtk.96},level=2]{4. สูจิฆรสิกฺขาปท}
\csromanheader{2}{9. \textbf{รตนวคฺค 4. สูจิฆรสิกฺขาปท}}
\proseitem{516}{เตน สมเยน พุทฺโธ ภควา สกฺเกสุ วิหรติ กปิลวตฺถุสฺมิํ นิโคฺรธาราเม.\csromanspacer{} เตน โข ปน สมเยน อญฺญตเรน ทนฺตกาเรน ภิกฺขู ปวาริตา โหนฺติ “เยสํ อยฺยานํ สูจิฆเรน \csromanfolio{217}อตฺโถ, อหํ สูจิฆเรนา”ติ.\csromanspacer{} เตน โข ปน สมเยน ภิกฺขู พหู สูจิฆเร วิญฺญาเปนฺติ, เยสํ ขุทฺทกา สูจิฆรา, เต มหนฺเต สูจิฆเร วิญฺญาเปนฺติ, เยสํ มหนฺตา สูจิฆรา, เต ขุทฺทเก สูจิฆเร วิญฺญาเปนฺติ.\csromanspacer{} อถ โข โส ทนฺตกาโร ภิกฺขูนํ พหู สูจิฆเร กโรนฺโต น สกฺโกติ อญฺญํ วิกฺกายิกํ ภณฺฑํ กาตุํ, อตฺตนาปิ น ยาเปติ, ปุตฺตทาโรปิสฺส กิลมติ.\csromanspacer{} มนุสฺสา อุชฺฌายนฺติ ขิยฺยนฺติ วิปาเจนฺติ “กถํ หิ นาม สมณา สกฺยปุตฺติยา น มตฺตํ ชานิตฺวา พหู สูจิฆเร วิญฺญาเปสฺสนฺติ.\csromanspacer{} อยํ อิเมสํ พหู สูจิฆเร กโรนฺโต น สกฺโกติ อญฺญํ วิกฺกายิกํ ภณฺฑํ กาตุํ, อตฺตนาปิ น ยาเปติ, ปุตฺตทาโรปิสฺส กิลมตี”ติ.\csromanspacer{} อสฺโสสุํ โข ภิกฺขู เตสํ มนุสฺสานํ อุชฺฌายนฺตานํ ขิยฺยนฺตานํ วิปาเจนฺตานํ.\csromanspacer{} เย เต ภิกฺขู อปฺปิจฺฉา ฯเปฯ เต อุชฺฌายนฺติ ขิยฺยนฺติ วิปาเจนฺติ “กถํ หิ นาม ภิกฺขู น มตฺตํ ชานิตฺวา พหู สูจิฆเร วิญฺญาเปสฺสนฺตี”ติ ฯเปฯ.\csromanspacer{} สจฺจํ \textbf{กิร ภิกฺขเว ภิกฺขู น มตฺตํ ชานิตฺวา พหู สูจิฆเร วิญฺญาเปนฺตีติ.}\csromanspacer{} \textbf{สจฺจํ ภควาติ.}\csromanspacer{}\textbf{ วิครหิ พุทฺโธ ภควา ฯเปฯ.}\csromanspacer{}\textbf{ กถํ หิ นาม เต} ภิกฺขเว โมฆปุริสา น มตฺตํ ชานิตฺวา พหู สูจิฆเร วิญฺญาเปสฺสนฺติ.\csromanspacer{} เนตํ ภิกฺขเว อปฺปสนฺนานํ วา ปสาทาย ฯเปฯ.\csromanspacer{} เอวญฺจ ปน ภิกฺขเว อิมํ สิกฺขาปทํ อุทฺทิเสยฺยาถ\csromandash{}}

\proseitem{517}{\textbf{“โย ปน ภิกฺขุ อฏฺฐิมยํ วา ทนฺตมยํ วา วิสาณมยํ วา สูจิฆรํ การาเปยฺย, เภทนกํ ปาจิตฺติยนฺ”}ติ.}

\proseitem{518}{\textbf{โย ปนา}ติ โย ยาทิโส ฯเปฯ.\csromanspacer{} \textbf{ภิกฺขู}ติ ฯเปฯ อยํ อิมสฺมิํ อตฺเถ อธิปฺเปโต ภิกฺขูติ.}

\prose{\textbf{อฏฺฐิ} นาม ยํ กิญฺจิ อฏฺฐิ.}

\prose{\textbf{ทนฺโต} นาม หตฺถิทนฺโต วุจฺจติ.}

\prose{\textbf{วิสาณํ} นาม ยํ กิญฺจิ วิสาณํ.}

\prose{\textbf{การาเปยฺยา}ติ กโรติ วา การาเปติ วา, ปโยเค ทุกฺกฏํ.\csromanspacer{} ปฏิลาเภน ภินฺทิตฺวา ปาจิตฺติยํ เทเสตพฺพํ.}

\proseitem{519}{อตฺตนา วิปฺปกตํ อตฺตนา ปริโยสาเปติ, อาปตฺติ ปาจิตฺติยสฺส.\csromanspacer{} อตฺตนา วิปฺปกตํ ปเรหิ ปริโยสาเปติ\footnote{ปริโยสาวาเปติ (ก)}, อาปตฺติ \csromanfolio{218}ปาจิตฺติยสฺส.\csromanspacer{} ปเรหิ วิปฺปกตํ อตฺตนา ปริโยสาเปติ, อาปตฺติ ปาจิตฺติยสฺส.\csromanspacer{} ปเรหิ วิปฺปกตํ ปเรหิ ปริโยสาเปติ\footnote{ปริโยสาวาเปติ (ก)}, อาปตฺติ ปาจิตฺติยสฺส.}

\prose{อญฺญสฺสตฺถาย กโรติ วา การาเปติ วา, อาปตฺติ ทุกฺกฏสฺส.\csromanspacer{} อญฺเญน กตํ ปฏิลภิตฺวา ปริภุญฺชติ, อาปตฺติ ทุกฺกฏสฺส.}

\proseitemwithcloserruled{520}{อนาปตฺติ คณฺฐิกาย\footnote{คณฺฑิกาย (สฺยา)} อรณิเก วิเธ\footnote{วีเฐ (สี), วีเถ (สฺยา)} อญฺชนิยา อญฺชนิสลากาย วาสิชเฏ อุทกปุญฺฉนิยา อุมฺมตฺตกสฺส อาทิกมฺมิกสฺสาติ.}{สูจิฆรสิกฺขาปทํ นิฏฺฐิตํ จตุตฺถํ.}
\csromanmatikaanchor{mtk.97}
\csromantocmarkref{subsection}{5. มญฺจปีฐสิกฺขาปท}{mtk.97}
\bookmark[dest={mtk.97},level=2]{5. มญฺจปีฐสิกฺขาปท}
\csromanheader{2}{9. \textbf{รตนวคฺค 5. มญฺจปีฐสิกฺขาปท}}
\proseitem{521}{เตน สมเยน พุทฺโธ ภควา สาวตฺถิยํ วิหรติ เชตวเน อนาถปิณฺฑิกสฺส อาราเม.\csromanspacer{} เตน โข ปน สมเยน อายสฺมา อุปนนฺโท สกฺยปุตฺโต อุจฺเจ มญฺเจ สยติ.\csromanspacer{} อถ โข ภควา สมฺพหุเลหิ ภิกฺขูหิ สทฺธิํ เสนาสนจาริกํ อาหิณฺฑนฺโต เยนายสฺมโต อุปนนฺทสฺส สกฺยปุตฺตสฺส วิหาโร เตนุปสงฺกมิ.\csromanspacer{} อทฺทสา โข อายสฺมา อุปนนฺโท สกฺยปุตฺโต ภควนฺตํ ทูรโตว อาคจฺฉนฺตํ, ทิสฺวาน ภควนฺตํ เอตทโวจ “อาคจฺฉตุ เม ภนฺเต ภควา สยนํ ปสฺสตู”ติ.\csromanspacer{} อถ โข ภควา ตโตว ปฏินิวตฺติตฺวา ภิกฺขู อามนฺเตสิ “อาสยโต ภิกฺขเว โมฆปุริโส เวทิตพฺโพ”ติ.\csromanspacer{} อถ โข ภควา อายสฺมนฺตํ อุปนนฺทํ สกฺยปุตฺตํ อเนกปริยาเยน วิครหิตฺวา ทุพฺภรตาย ฯเปฯ.\csromanspacer{} เอวญฺจ ปน ภิกฺขเว อิมํ สิกฺขาปทํ อุทฺทิเสยฺยาถ\csromandash{}}

\proseitem{522}{\textbf{“นวํ ปน ภิกฺขุนา มญฺจํ วา ปีฐํ วา การยมาเนน อฏฺฐงฺคุลปาทกํ กาเรตพฺพํ สุคตงฺคุเลน อญฺญตฺร เหฏฺฐิมาย อฏนิยา, ตํ อติกฺกามยโต เฉทนกํ ปาจิตฺติยนฺ”}ติ.}

\proseitem{523}{\textbf{นวํ} นาม กรณํ อุปาทาย วุจฺจติ.}

\prose{\csromanfolio{219}\textbf{มญฺโจ} นาม จตฺตาโร มญฺจา มสารโก พุนฺทิกาพทฺโธ กุฬีรปาทโก อาหจฺจปาทโก.}

\prose{\textbf{ปีฐํ} นาม จตฺตาริ ปีฐานิ มสารกํ พุนฺทิกาพทฺธํ กุฬีรปาทกํ อาหจฺจปาทกํ.}

\prose{\textbf{การยมาเนนา}ติ กโรนฺโต วา การาเปนฺโต วา.}

\prose{\textbf{อฏฺฐงฺคุลปาทกํ กาเรตพฺพํ สุคตงฺคุเลน อญฺญตฺร เหฏฺฐิมาย อฏนิยา}ติ ฐเปตฺวา เหฏฺฐิมํ อฏนิํ, ตํ อติกฺกาเมตฺวา กโรติ วา การาเปติ วา, ปโยเค ทุกฺกฏํ, ปฏิลาเภน ฉินฺทิตฺวา ปาจิตฺติยํ เทเสตพฺพํ.}

\proseitem{524}{อตฺตนา วิปฺปกตํ อตฺตนา ปริโยสาเปติ, อาปตฺติ ปาจิตฺติยสฺส.\csromanspacer{} อตฺตนา วิปฺปกตํ ปเรหิ ปริโยสาเปติ, อาปตฺติ ปาจิตฺติยสฺส.\csromanspacer{} ปเรหิ วิปฺปกตํ อตฺตนา ปริโยสาเปติ, อาปตฺติ ปาจิตฺติยสฺส.\csromanspacer{} ปเรหิ วิปฺปกตํ ปเรหิ ปริโยสาเปติ, อาปตฺติ ปาจิตฺติยสฺส.}

\prose{อญฺญสฺสตฺถาย กโรติ วา การาเปติ วา, อาปตฺติ ทุกฺกฏสฺส.\csromanspacer{} อญฺเญน กตํ ปฏิลภิตฺวา ปริภุญฺชติ, อาปตฺติ ทุกฺกฏสฺส.}

\proseitemwithcloserruled{525}{อนาปตฺติ ปมาณิกํ กโรติ, อูนกํ กโรติ, อญฺเญน กตํ ปมาณาติกฺกนฺตํ ปฏิลภิตฺวา ฉินฺทิตฺวา ปริภุญฺชติ, อุมฺมตฺตกสฺส อาทิกมฺมิกสฺสาติ.}{มญฺจปีฐสิกฺขาปทํ นิฏฺฐิตํ ปญฺจมํ.}
\csromanmatikaanchor{mtk.98}
\csromantocmarkref{subsection}{6. ตูโลนทฺธสิกฺขาปท}{mtk.98}
\bookmark[dest={mtk.98},level=2]{6. ตูโลนทฺธสิกฺขาปท}
\csromanheader{2}{9. \textbf{รตนวคฺค 6. ตูโลนทฺธสิกฺขาปท}}
\proseitem{526}{เตน สมเยน พุทฺโธ ภควา สาวตฺถิยํ วิหรติ เชตวเน อนาถปิณฺฑิกสฺส อาราเม.\csromanspacer{} เตน โข ปน สมเยน ฉพฺพคฺคิยา ภิกฺขู มญฺจมฺปิ ปีฐมฺปิ ตูโลนทฺธํ การาเปนฺติ.\csromanspacer{} มนุสฺสา วิหารจาริกํ อาหิณฺฑนฺตา ปสฺสิตฺวา อุชฺฌายนฺติ ขิยฺยนฺติ วิปาเจนฺติ “กถํ หิ นาม สมณา สกฺยปุตฺติยา มญฺจมฺปิ ปีฐมฺปิ ตูโลนทฺธํ การาเปสฺสนฺติ เสยฺยถาปิ คิหี กามโภคิโน \csromanfolio{220}”ติ.\csromanspacer{} อสฺโสสุํ โข ภิกฺขู เตสํ มนุสฺสานํ อุชฺฌายนฺตานํ ขิยฺยนฺตานํ วิปาเจนฺตานํ.\csromanspacer{} เย เต ภิกฺขู อปฺปิจฺฉา ฯเปฯ เต อุชฺฌายนฺติ ขิยฺยนฺติ วิปาเจนฺติ “กถํ หิ นาม ฉพฺพคฺคิยา ภิกฺขู มญฺจมฺปิ ปีฐมฺปิ ตูโลนทฺธํ การาเปสฺสนฺติ”ติ ฯเปฯ.\csromanspacer{} สจฺจํ กิร ตุเมฺห ภิกฺขเว มญฺจมฺปิ ปีฐมฺปิ ตูโลนทฺธํ การาเปถาติ.\csromanspacer{} สจฺจํ ภควาติ.\csromanspacer{} วิครหิ พุทฺโธ ภควา ฯเปฯ.\csromanspacer{} กถํ หิ นาม ตุเมฺห โมฆปุริสา มญฺจมฺปิ ปีฐมฺปิ ตูโลนทฺธํ การาเปสฺสถ.\csromanspacer{} เนตํ โมฆปุริสา อปฺปสนฺนานํ วา ปสาทาย ฯเปฯ.\csromanspacer{} เอวญฺจ ปน ภิกฺขเว อิมํ สิกฺขาปทํ อุทฺทิเสยฺยาถ\csromandash{}}

\proseitem{527}{\textbf{“โย ปน ภิกฺขุ มญฺจํ วา ปีฐํ วา ตูโลนทฺธํ การาเปยฺย, อุทฺทาลนกํ ปาจิตฺติยนฺ”}ติ.}

\proseitem{528}{\textbf{โย ปนา}ติ โย ยาทิโส ฯเปฯ.\csromanspacer{} \textbf{ภิกฺขู}ติ ฯเปฯ อยํ อิมสฺมิํ อตฺเถ อธิปฺเปโต ภิกฺขูติ.}

\prose{\textbf{มญฺโจ} นาม จตฺตาโร มญฺจา มสารโก พุนฺทิกาพทฺโธ กุฬีรปาทโก อาหจฺจปาทโก.}

\prose{\textbf{ปีฐํ} นาม จตฺตาริ ปีฐานิ มสารกํ พุนฺทิกาพทฺธํ กุฬีรปาทกํ อาหจฺจปาทกํ.}

\prose{\textbf{ตูลํ} นาม ตีณิ ตูลานิ รุกฺขตูลํ ลตาตูลํ โปตกิตุลํ.}

\prose{\textbf{การาเปยฺยา}ติ กโรติ วา การาเปติ วา, ปโยเค ทุกฺกฏํ.\csromanspacer{} ปฏิลาเภน อุทฺทาเลตฺวา ปาจิตฺติยํ เทเสตพฺพํ.}

\proseitem{529}{อตฺตนา วิปฺปกตํ อตฺตนา ปริโยสาเปติ, อาปตฺติ ปาจิตฺติยสฺส.\csromanspacer{} อตฺตนา วิปฺปกตํ ปเรหิ ปริโยสาเปติ, อาปตฺติ ปาจิตฺติยสฺส.\csromanspacer{} ปเรหิ วิปฺปกตํ อตฺตนา ปริโยสาเปติ, อาปตฺติ ปาจิตฺติยสฺส.\csromanspacer{} ปเรหิ วิปฺปกตํ ปเรหิ ปริโยสาเปติ, อาปตฺติ ปาจิตฺติยสฺส.}

\prose{อญฺญสฺสตฺถาย กโรติ วา การาเปติ วา, อาปตฺติ ทุกฺกฏสฺส.\csromanspacer{} อญฺเญน กตํ ปฏิลภิตฺวา ปริภุญฺชติ, อาปตฺติ ทุกฺกฏสฺส.}

\proseitemwithcloser{530}{อนาปตฺติ อาโยเค กายพนฺธเน อํสพทฺธเก ปตฺตถวิกาย ปริสฺสาวเน พิพฺโพหนํ กโรติ, อญฺเญน กตํ ปฏิลภิตฺวา อุทฺทาเลตฺวา ปริภุญฺชติ, อุมฺมตฺตกสฺส อาทิกมฺมิกสฺสาติ.}{ตูโลนทฺธสิกฺขาปทํ นิฏฺฐิตํ ฉฏฺฐํ.}
\csromanmatikaanchor{mtk.99}
\csromantocmarkref{subsection}{7. นิสีทนสิกฺขาปท}{mtk.99}
\bookmark[dest={mtk.99},level=2]{7. นิสีทนสิกฺขาปท}
\csromanheader{2}{9. \textbf{รตนวคฺค 7. นิสีทนสิกฺขาปท}}
\proseitem{531}{\csromanfolio{221}เตน สมเยน พุทฺโธ ภควา สาวตฺถิยํ วิหรติ เชตวเน อนาถปิณฺฑิกสฺส อาราเม.\csromanspacer{} เตน โข ปน สมเยน ภควตา ภิกฺขูนํ นิสีทนํ อนุญฺญาตํ โหติ.\csromanspacer{} ฉพฺพคฺคิยา ภิกฺขู “ภควตา นิสีทนํ อนุญฺญาตนฺ”ติ อปฺปมาณิกานิ นิสีทนานิ ธาเรนฺติ, มญฺจสฺสปิ ปีฐสฺสปิ ปุรโตปิ ปจฺฉโตปิ โอลมฺเพนฺติ.\csromanspacer{} เย เต ภิกฺขู อปฺปิจฺฉา ฯเปฯ เต อุชฺฌายนฺติ ขิยฺยนฺติ วิปาเจนฺติ “กถํ หิ นาม ฉพฺพคฺคิยา ภิกฺขู อปฺปมาณิกานิ นิสีทนานิ ธาเรสฺสนฺตี”ติ ฯเปฯ.\csromanspacer{} สจฺจํ กิร ตุเมฺห ภิกฺขเว อปฺปมาณิกานิ นิสีทนานิ ธาเรถาติ.\csromanspacer{} สจฺจํ ภควาติ.\csromanspacer{} วิครหิ พุทฺโธ ภควา ฯเปฯ.\csromanspacer{} กถํ หิ นาม ตุเมฺห โมฆปุริสา อปฺปมาณิกานิ นิสีทนานิ ธาเรสฺสถ.\csromanspacer{} เนตํ โมฆปุริสา อปฺปสนฺนานํ วา ปสาทาย ฯเปฯ.\csromanspacer{} เอวญฺจ ปน ภิกฺขเว อิมํ สิกฺขาปทํ อุทฺทิเสยฺยาถ\csromandash{}}

\prose{\textbf{“นิสีทนํ ปน ภิกฺขุนา การยมาเนน ปมาณิกํ กาเรตพฺพํ, ตตฺริทํ ปมาณํ, ทีฆโส เทฺว วิทตฺถิโย สุคตวิทตฺถิยา, ติริยํ ทิยฑฺฒํ, ตํ อติกฺกามยโต เฉทนกํ ปาจิตฺติยนฺ”}ติ.}

\prose{เอวญฺจิทํ ภควตา ภิกฺขูนํ สิกฺขาปทํ ปญฺญตฺตํ โหติ.}

\proseitem{532}{เตน โข ปน สมเยน อายสฺมา อุทายี มหากาโย โหติ.\csromanspacer{} โส ภควโต ปุรโต นิสีทนํ ปญฺญเปตฺวา สมนฺตโต สมญฺฉมาโน นิสีทติ.\csromanspacer{} อถ โข ภควา อายสฺมนฺตํ อุทายิํ เอตทโวจ “กิสฺส ตฺวํ อุทายิ นิสีทนํ สมนฺตโต สมญฺฉสิ เสยฺยถาปิ ปุราณาสิโกฏฺโฐ”ติ.\csromanspacer{} ตถา หิ ปน ภนฺเต ภควตา ภิกฺขูนํ อติขุทฺทกํ นิสีทนํ อนุญฺญาตนฺติ.\csromanspacer{} อถ โข ภควา เอตสฺมิํ นิทาเน เอตสฺมิํ ปกรเณ ธมฺมิํ กถํ กตฺวา ภิกฺขู อามนฺเตสิ อนุชานามิ ภิกฺขเว นิสีทนสฺส ทสํ วิทตฺถิํ.\csromanspacer{} เอวญฺจ ปน ภิกฺขเว อิมํ สิกฺขาปทํ อุทฺทิเสยฺยาถ\csromandash{}}

\proseitem{533}{“\textbf{นิสีทนํ ปน ภิกฺขุนา การยมาเนน ปมาณิกํ กาเรตพฺพํ,} \textbf{ตตฺริทํ ปมาณํ, ทีฆโส เทฺว วิทตฺถิโย สุคตวิทตฺถิยา, ติริยํ} ทิยฑฺฒํ, ทสา วิทตฺถิ, ตํ อติกฺกามยโต เฉทนกํ ปาจิตฺติยนฺ”ติ.}

\proseitem{534}{\textbf{นิสีทนํ} นาม สทสํ วุจฺจติ.}

\prose{\csromanfolio{222}\textbf{การยมาเนนา}ติ กโรนฺโต วา การาเปนฺโต วา.\csromanspacer{} ปมาณิกํ กาเรตพฺพํ, ตตฺริทํ ปมาณํ, ทีฆโส เทฺว วิทตฺถิโย สุคตวิทตฺถิยา, ติริยํ ทิยฑฺฒํ, ทสา วิทตฺถิ, ตํ อติกฺกาเมตฺวา กโรติ วา การาเปติ วา, ปโยเค ทุกฺกฏํ.\csromanspacer{} ปฏิลาเภน ฉินฺทิตฺวา ปาจิตฺติยํ เทเสตพฺพํ.}

\proseitem{535}{อตฺตนา วิปฺปกตํ อตฺตนา ปริโยสาเปติ, อาปตฺติ ปาจิตฺติยสฺส.\csromanspacer{} อตฺตนา วิปฺปกตํ ปเรหิ ปริโยสาเปติ, อาปตฺติ ปาจิตฺติยสฺส.\csromanspacer{} ปเรหิ วิปฺปกตํ อตฺตนา ปริโยสาเปติ, อาปตฺติ ปาจิตฺติยสฺส.\csromanspacer{} ปเรหิ วิปฺปกตํ ปเรหิ ปริโยสาเปติ, อาปตฺติ ปาจิตฺติยสฺส.}

\prose{อญฺญสฺสตฺถาย กโรติ วา การาเปติ วา, อาปตฺติ ทุกฺกฏสฺส.\csromanspacer{} อญฺเญน กตํ ปฏิลภิตฺวา ปริภุญฺชติ, อาปตฺติ ทุกฺกฏสฺส.}

\proseitemwithcloserruled{536}{อนาปตฺติ ปมาณิกํ กโรติ, อูนกํ กโรติ, อญฺเญน กตํ ปมาณาติกฺกนฺตํ ปฏิลภิตฺวา ฉินฺทิตฺวา ปริภุญฺชติ, วิตานํ วา ภูมตฺถรณํ วา สาณิปาการํ วา ภิสิํ วา พิพฺโพหนํ วา กโรติ, อุมฺมตฺตกสฺส อาทิกมฺมิกสฺสาติ.}{นิสีทนสิกฺขาปทํ นิฏฺฐิตํ สตฺตมํ.}
\csromanmatikaanchor{mtk.100}
\csromantocmarkref{subsection}{8. กณฺฑุปฺปฏิจฺฉาทิสิกฺขาปท}{mtk.100}
\bookmark[dest={mtk.100},level=2]{8. กณฺฑุปฺปฏิจฺฉาทิสิกฺขาปท}
\csromanheader{2}{9. รตนวคฺค 8. กณฺฑุปฺปฏิจฺฉาทิ\-สิกฺขาปท}
\proseitem{537}{เตน สมเยน พุทฺโธ ภควา สาวตฺถิยํ วิหรติ เชตวเน อนาถปิณฺฑิกสฺส อาราเม.\csromanspacer{} เตน โข ปน สมเยน ภควตา ภิกฺขูนํ กณฺฑุปฺปฏิจฺฉาทิ อนุญฺญาตา โหติ.\csromanspacer{} ฉพฺพคฺคิยา ภิกฺขู “ภควตา กณฺฑุปฺปฏิจฺฉาทิ อนุญฺญาตา”ติ อปฺปมาณิกาโย กณฺฑุปฺปฏิจฺฉาทิโย ธาเรนฺติ, ปุรโตปิ ปจฺฉโตปิ อากฑฺฒนฺตา อาหิณฺฑนฺติ.\csromanspacer{} เย เต ภิกฺขู อปฺปิจฺฉา ฯเปฯ เต อุชฺฌายนฺติ ขิยฺยนฺติ วิปาเจนฺติ “กถํ หิ นาม ฉพฺพคฺคิยา ภิกฺขู อปฺปมาณิกาโย กณฺฑุปฺปฏิจฺฉาทิโย ธาเรสฺสนฺตี”ติ ฯเปฯ.\csromanspacer{} สจฺจํ กิร ตุเมฺห ภิกฺขเว อปฺปมาณิกาโย กณฺฑุปฺปฏิจฺฉาทิโย ธาเรถาติ.\csromanspacer{} สจฺจํ ภควาติ.\csromanspacer{} วิครหิ พุทฺโธ ภควา ฯเปฯ.\csromanspacer{} กถํ หิ นาม ตุเมฺห โมฆปุริสา อปฺปมาณิกาโย กณฺฑุปฺปฏิจฺฉาทิโย ธาเรสฺสถ.\csromanspacer{} เนตํ โมฆปุริสา \csromanfolio{223}\textbf{อปฺปสนฺนานํ วา ปสาทาย ฯเปฯ.}\csromanspacer{}\textbf{ เอวญฺจ ปน ภิกฺขเว อิมํ} สิกฺขาปทํ อุทฺทิเสยฺยาถ\csromandash{}}

\proseitem{538}{\textbf{“กณฺฑุปฺปฏิจฺฉาทิํ ปน ภิกฺขุนา การยมาเนน ปมาณิกา กาเรตพฺพา, ตตฺริทํ ปมาณํ, ทีฆโส จตสฺโส วิทตฺถิโย สุคตวิทตฺถิยา, ติริยํ เทฺว วิทตฺถิโย, ตํ อติกฺกามยโต เฉทนกํ ปาจิตฺติยนฺ”}ติ.}

\proseitem{539}{\textbf{กณฺฑุปฺปฏิจฺฉาทิ} นาม ยสฺส อโธนาภิ อุพฺภชาณุมณฺฑลํ กณฺฑุ วา ปีฬกา วา อสฺสาโว วา ถุลฺลกจฺฉุ วา อาพาโธ, ตสฺส ปฏิจฺฉาทนตฺถาย.}

\prose{\textbf{การยมาเนนา}ติ กโรนฺโต วา การาเปนฺโต วา.\csromanspacer{} ปมาณิกา กาเรตพฺพา, ตตฺริทํ ปมาณํ, ทีฆโส จตสฺโส วิทตฺถิโย สุคตวิทตฺถิยา, ติริยํ เทฺว วิทตฺถิโย, ตํ อติกฺกาเมตฺวา กโรติ วา การาเปติ วา, ปโยเค ทุกฺกฏํ.\csromanspacer{} ปฏิลาเภน ฉินฺทิตฺวา ปาจิตฺติยํ เทเสตพฺพํ.}

\proseitem{540}{อตฺตนา วิปฺปกตํ อตฺตนา ปริโยสาเปติ, อาปตฺติ ปาจิตฺติยสฺส.\csromanspacer{} อตฺตนา วิปฺปกตํ ปเรหิ ปริโยสาเปติ, อาปตฺติ ปาจิตฺติยสฺส.\csromanspacer{} ปเรหิ วิปฺปกตํ อตฺตนา ปริโยสาเปติ, อาปตฺติ ปาจิตฺติยสฺส.\csromanspacer{} ปเรหิ วิปฺปกตํ ปเรหิ ปริโยสาเปติ, อาปตฺติ ปาจิตฺติยสฺส.}

\prose{อญฺญสฺสตฺถาย กโรติ วา การาเปติ วา, อาปตฺติ ทุกฺกฏสฺส.\csromanspacer{} อญฺเญน กตํ ปฏิลภิตฺวา ปริภุญฺชติ, อาปตฺติ ทุกฺกฏสฺส.}

\proseitemwithcloser{541}{อนาปตฺติ ปมาณิกํ กโรติ, อูนกํ กโรติ, อญฺเญน กตํ ปมาณาติกฺกนฺตํ ปฏิลภิตฺวา ฉินฺทิตฺวา ปริภุญฺชติ, วิตานํ วา ภูมตฺถรณํ วา สาณิปาการํ วา ภิสิํ วา พิพฺโพหนํ วา กโรติ, อุมฺมตฺตกสฺส อาทิกมฺมิกสฺสาติ.}{กณฺฑุปฺปฏิจฺฉาทิ\-สิกฺขาปทํ นิฏฺฐิตํ อฏฺฐมํ.}
\csromanmatikaanchor{mtk.101}
\csromantocmarkref{subsection}{9. วสฺสิกสาฏิกาสิกฺขาปท}{mtk.101}
\bookmark[dest={mtk.101},level=2]{9. วสฺสิกสาฏิกาสิกฺขาปท}
\csromanheader{2}{9. รตนวคฺค 9. วสฺสิกสาฏิกา\-สิกฺขาปท}
\proseitem{542}{\csromanfolio{224}เตน สมเยน พุทฺโธ ภควา สาวตฺถิยํ วิหรติ เชตวเน อนาถปิณฺฑิกสฺส อาราเม.\csromanspacer{} เตน โข ปน สมเยน ภควตา ภิกฺขูนํ วสฺสิกสาฏิกา อนุญฺญาตา โหติ.\csromanspacer{} ฉพฺพคฺคิยา ภิกฺขู “ภควตา วสฺสิกสาฏิกา อนุญฺญาตา”ติ อปฺปมาณิกาโย วสฺสิกสาฏิกาโย ธาเรนฺติ, ปุรโตปิ ปจฺฉโตปิ อากฑฺฒนฺตา อาหิณฺฑนฺติ.\csromanspacer{} เย เต ภิกฺขู อปฺปิจฺฉา ฯเปฯ เต อุชฺฌายนฺติ ขิยฺยนฺติ วิปาเจนฺติ “กถํ หิ นาม ฉพฺพคฺคิยา ภิกฺขู อปฺปมาณิกาโย วสฺสิกสาฏิกาโย ธาเรสฺสนฺตี”ติ ฯเปฯ.\csromanspacer{} สจฺจํ กิร ตุเมฺห ภิกฺขเว อปฺปมาณิกาโย วสฺสิกสาฏิกาโย ธาเรถาติ.\csromanspacer{} สจฺจํ ภควาติ.\csromanspacer{} วิครหิ พุทฺโธ ภควา ฯเปฯ.\csromanspacer{} กถํ หิ นาม ตุเมฺห โมฆปุริสา อปฺปมาณิกาโย วสฺสิกสาฏิกาโย ธาเรสฺสถ.\csromanspacer{} เนตํ โมฆปุริสา อปฺปสนฺนานํ วา ปสาทาย ฯเปฯ.\csromanspacer{} เอวญฺจ ปน ภิกฺขเว อิมํ สิกฺขาปทํ อุทฺทิเสยฺยาถ\csromandash{}}

\proseitem{543}{\textbf{“วสฺสิกสาฏิกํ ปน ภิกฺขุนา การยมาเนน ปมาณิกา กาเรตพฺพา, ตตฺริทํ ปมาณํ, ทีฆโส ฉ วิทตฺถิโย สุคตวิทตฺถิยา, ติริยํ อฑฺฒเตยฺยา, ตํ อติกฺกามยโต เฉทนกํ ปาจิตฺติยนฺ”}ติ.}

\proseitem{544}{\textbf{วสฺสิกสาฏิกา} นาม วสฺสานสฺส จตุมาสตฺถาย.}

\prose{\textbf{การยมาเนนา}ติ กโรนฺโต วา การาเปนฺโต วา.\csromanspacer{} ปมาณิกา กาเรตพฺพา, ตตฺริทํ ปมาณํ, ทีฆโส ฉ วิทตฺถิโย สุคตวิทตฺถิยา, ติริยํ อฑฺฒเตยฺยา, ตํ อติกฺกาเมตฺวา กโรติ วา การาเปติ วา, ปโยเค ทุกฺกฏํ.\csromanspacer{} ปฏิลาเภน ฉินฺทิตฺวา ปาจิตฺติยํ เทเสตพฺพํ.}

\proseitem{545}{อตฺตนา วิปฺปกตํ อตฺตนา ปริโยสาเปติ, อาปตฺติ ปาจิตฺติยสฺส.\csromanspacer{} อตฺตนา วิปฺปกตํ ปเรหิ ปริโยสาเปติ, อาปตฺติ ปาจิตฺติยสฺส.\csromanspacer{} ปเรหิ วิปฺปกตํ อตฺตนา ปริโยสาเปติ, อาปตฺติ อาจิตฺติยสฺส.\csromanspacer{} ปเรหิ วิปฺปกตํ ปเรหิ ปริโยสาเปติ, อาปตฺติ ปาจิตฺติยสฺส.}

\prose{อญฺญสฺสตฺถาย กโรติ วา การาเปติ วา, อาปตฺติ ทุกฺกฏสฺส.\csromanspacer{} อญฺเญน กตํ ปฏิลภิตฺวา ปริภุญฺชติ, อาปตฺติ ทุกฺกฏสฺส.}

\proseitemwithcloserruled{546}{\csromanfolio{225}อนาปตฺติ ปมาณิกํ กโรติ, อูนกํ กโรติ, อญฺเญน กตํ ปมาณาติกฺกนฺตํ ปฏิลภิตฺวา ฉินฺทิตฺวา ปริภุญฺชติ, วิตานํ วา ภูมตฺถรณํ วา สาณิปาการํ วา ภิสิํ วา พิพฺโพหนํ วา กโรติ, อุมฺมตฺตกสฺส อาทิกมฺมิกสฺสาติ.}{วสฺสิกสาฏิกา\-สิกฺขาปทํ นิฏฺฐิตํ นวมํ.}
\csromanmatikaanchor{mtk.102}
\csromantocmarkref{subsection}{10. นนฺทสิกฺขาปท}{mtk.102}
\bookmark[dest={mtk.102},level=2]{10. นนฺทสิกฺขาปท}
\csromanheader{2}{9. \textbf{รตนวคฺค} 10. นนฺทสิกฺขาปท}
\proseitem{547}{เตน สมเยน พุทฺโธ ภควา สาวตฺถิยํ วิหรติ เชตวเน อนาถปิณฺฑิกสฺส อาราเม.\csromanspacer{} เตน โข ปน สมเยน อายสฺมา นนฺโท ภควโต มาตุจฺฉาปุตฺโต อภิรูโป โหติ ทสฺสนีโย ปาสาทิโก จตุรงฺคุโลมโก ภควตา\footnote{ภควโต (ก)}, โส สุคตจีวรปฺปมาณํ จีวรํ ธาเรติ.\csromanspacer{} อทฺทสํสุ โข เถรา ภิกฺขู อายสฺมนฺตํ นนฺทํ ทูรโตว อาคจฺฉนฺตํ, ทิสฺวาน “ภควา อาคจฺฉตี”ติ อาสนา วุฏฺฐหนฺติ, เต อุปคเต ชานิตฺวา\footnote{อุปคตํ สญฺชานิตฺวา (สฺยา), อุปคเต สญฺชานิตฺวา (สี)} อุชฺฌายนฺติ ขิยฺยนฺติ วิปาเจนฺติ “กถํ หิ นาม อายสฺมา นนฺโท สุคตจีวรปฺปมาณํ จีวรํ ธาเรสฺสตี”ติ ฯเปฯ.\csromanspacer{} สจฺจํ กิร ตฺวํ นนฺท สุคตจีวรปฺปมาณํ จีวรํ ธาเรสฺสตี”ติ.\csromanspacer{} สจฺจํ ภควาติ.\csromanspacer{} วิครหิ พุทฺโธ ภควา ฯเปฯ.\csromanspacer{} กถํ หิ นาม ตฺวํ นนฺท สุคตจีวรปฺปมาณํ จีวรํ ธาเรสฺสสิ.\csromanspacer{} เนตํ นนฺท อปฺปสนฺนานํ วา ปสาทาย ฯเปฯ.\csromanspacer{} เอวญฺจ ปน ภิกฺขเว อิมํ สิกฺขาปทํ อุทฺทิเสยฺยาถ\csromandash{}}

\proseitem{548}{\textbf{“โย ปน ภิกฺขุ สุคตจีวรปฺปมาณํ จีวรํ การาเปยฺย อติเรกํ วา, เฉทนกํ ปาจิตฺติยํ, ตตฺริทํ สุคตสฺส สุคตจีวรปฺปมาณํ, ทีฆโส นว วิทตฺถิโย สุคตวิทตฺถิยา, ติริยํ ฉ วิทตฺถิโย, อิทํ สุคตสฺส สุคตจีวรปฺปมาณนฺ”}ติ.}

\proseitem{549}{\textbf{โย ปนา}ติ โย ยาทิโส ฯเปฯ.\csromanspacer{} \textbf{ภิกฺขู}ติ ฯเปฯ อยํ อิมสฺมิํ อตฺเถ อธิปฺเปโต ภิกฺขูติ.}

\prose{\csromanfolio{226}\textbf{สุคตจีวรปฺปมาณํ}\footnote{สุคตจีวรํ (ก)} นาม ทีฆโส นว วิทตฺถิโย สุคตวิตฺถิยา, ติริยํ ฉ วิทตฺถิโย.}

\prose{\textbf{การาเปยฺยา}ติ กโรติ วา การาเปติ วา, ปโยเค ทุกฺกฏํ.\csromanspacer{} ปฏิลาเภน ฉินฺทิตฺวา ปาจิตฺติยํ เทเสตพฺพํ.}

\proseitem{550}{อตฺตนา วิปฺปกตํ อตฺตนา ปริโยสาเปติ, อาปตฺติ ปาจิตฺติยสฺส.\csromanspacer{} อตฺตนา วิปฺปกตํ ปเรหิ ปริโยสาเปติ, อาปตฺติ ปาจิตฺติยสฺส.\csromanspacer{} ปเรหิ วิปฺปกตํ อตฺตนา ปริโยสาเปติ, อาปตฺติ ปาจิตฺติยสฺส.\csromanspacer{} ปเรหิ วิปฺปกตํ ปเรหิ ปริโยสาเปติ, อาปตฺติ ปาจิตฺติยสฺส.}

\prose{อญฺญสฺสตฺถาย กโรติ วา การาเปติ วา, อาปตฺติ ทุกฺกฏสฺส.\csromanspacer{} อญฺเญน กตํ ปฏิลภิตฺวา ปริภุญฺชติ, อาปตฺติ ทุกฺกฏสฺส.}

\proseitemwithcloser{551}{อนาปตฺติ อูนกํ กโรติ, อญฺเญน กตํ ปฏิลภิตฺวา ฉินฺทิตฺวา ปริภุญฺชติ, วิตานํ วา ภูมตฺถรณํ วา สาณิปาการํ วา ภิสิํ วา พิพฺโพหนํ วา กโรติ, อุมฺมตฺตกสฺส อาทิกมฺมิกสฺสาติ.}{นนฺทสิกฺขาปทํ นิฏฺฐิตํ ทสมํ.}
\vspace{\tipitakaparskip}
\csromancenter{รตนวคฺโค\footnote{ราชวคฺโค (สี)} นวโม.}
\csromancenter{ตสฺสุทฺทานํ}
\vspace{0.5\baselineskip}
\setlength{\csromangathapagewidth}{0pt}
\csromangathameasuregroup{รญฺโญ จ รตนํ สนฺตํ, \\ นิสีทนญฺจ กณฺฑุญฺจ,}{\csromangathabat{รญฺโญ จ รตนํ สนฺตํ,}{สูจิ มญฺจํ จ ตูลิกํ.} \\ \csromangathabat{นิสีทนญฺจ กณฺฑุญฺจ,}{วสฺสิกา สุคเตน จาติ.}}
\csromangathasetleft{รญฺโญ จ รตนํ สนฺตํ, \\ นิสีทนญฺจ กณฺฑุญฺจ,}
\csromangathagroup{\csromangathastanza{\csromangathabat{รญฺโญ จ รตนํ สนฺตํ,}{สูจิ มญฺจํ จ ตูลิกํ.} \\ \csromangathabat{นิสีทนญฺจ กณฺฑุญฺจ,}{วสฺสิกา สุคเตน จาติ.}}}

\prosewithcloser{อุทฺทิฏฺฐา โข อายสฺมนฺโต เทฺวนวุติ ปาจิตฺติยา ธมฺมา, ตตฺถายสฺมนฺเต ปุจฺฉามิ กจฺจิตฺถ ปริสุทฺธา, ทุติยมฺปิ ปุจฺฉามิ กจฺจิตฺถ ปริสุทฺธา, ตติยมฺปิ ปุจฺฉามิ กจฺจิตฺถ ปริสุทฺธา, ปริสุทฺเธตฺถายสฺมนฺโต, ตสฺมา ตุณฺหี, เอวเมตํ ธารยามีติ.}{ขุทฺทกํ สมตฺตํ.}
\nitthitam{ปาจิตฺติยกณฺฑํ นิฏฺฐิตํ.}
\csromanmatikaanchor{mtk.103}
\csromantocmarknopage{chapter}{6. ปาฏิเทสนียกณฺฑ}{mtk.103}
\bookmark[dest={mtk.103},level=0]{6. ปาฏิเทสนียกณฺฑ}
\setcsromanheadcha{6. ปาฏิเทสนียกณฺฑ}
\setcsromanheadhclear{}
\chapterheadpageread{6. \textbf{ปาฏิเทสนียกณฺฑ}}
\csromanmatikaanchor{mtk.104}
\csromantocmarkref{section}{1. ปฐมปาฏิเทสนียสิกฺขาปท}{mtk.104}
\bookmark[dest={mtk.104},level=1]{1. ปฐมปาฏิเทสนียสิกฺขาปท}
\setcsromanheadh{1. ปฐมปาฏิเทสนียสิกฺขาปท}
\csromanheader{1}{1. ปฐม\-ปาฏิเทสนีย\-สิกฺขาปท}
\prose{\csromanfolio{227}อิเม โข ปนายสฺมนฺโต จตฺตาโร ปาฏิเทสนียา ธมฺมา อุทฺเทสํ อาคจฺฉนฺติ.}

\proseitem{552}{เตน สมเยน พุทฺโธ ภควา สาวตฺถิยํ วิหรติ เชตวเน อนาถปิณฺฑิกสฺส อาราเม.\csromanspacer{} เตน โข ปน สมเยน อญฺญตรา ภิกฺขุนี สาวตฺถิยํ ปิณฺฑาย จริตฺวา ปฏิกฺกมนกาเล อญฺญตรํ ภิกฺขุํ ปสฺสิตฺวา เอตทโวจ “หนฺทายฺย ภิกฺขํ ปฏิคฺคณฺหา”ติ. “สุฏฺฐุ ภคินี”ติ สพฺเพว อคฺคเหสิ.\csromanspacer{} สา อุปกฏฺเฐ กาเล นาสกฺขิ ปิณฺฑาย จริตุํ, ฉินฺนภตฺตา อโหสิ.\csromanspacer{} อถ โข สา ภิกฺขุนี ทุติยมฺปิ ทิวสํ ฯเปฯ.\csromanspacer{} ตติยมฺปิ ทิวสํ สาวตฺถิยํ ปิณฺฑาย จริตฺวา ปฏิกฺกมนกาเล ตํ ภิกฺขุํ ปสฺสิตฺวา เอตทโวจ “หนฺทายฺย ภิกฺขํ ปฏิคฺคณฺหา”ติ. “สุฏฺฐุ ภคินี”ติ สพฺเพว อคฺคเหสิ.\csromanspacer{} สา อุปกฏฺเฐ กาเล นาสกฺขิ ปิณฺฑาย จริตุํ, ฉินฺนภตฺตา อโหสิ.\csromanspacer{} อถ โข สา ภิกฺขุนี จตุตฺเถ ทิวเส รถิกาย ปเวเธนฺตี คจฺฉติ.\csromanspacer{} เสฏฺฐิ คหปติ รเถน ปฏิปถํ อาคจฺฉนฺโต ตํ ภิกฺขุนิํ เอตทโวจ “อเปหายฺเย”ติ.\csromanspacer{} สา โวกฺกมนฺตี ตตฺเถว ปริปติ.\csromanspacer{} เสฏฺฐิ คหปติ ตํ ภิกฺขุนิํ ขมาเปสิ “ขมาหายฺเย, มยาสิ ปาติตา”ติ.\csromanspacer{} นาหํ คหปติ ตยา ปาติตา, อปิ จ อหเมว ทุพฺพลาติ.\csromanspacer{} กิสฺส ปน ตฺวํ อยฺเย ทุพฺพลาติ.\csromanspacer{} อถ โข สา ภิกฺขุนี เสฏฺฐิสฺส คหปติสฺส เอตมตฺถํ อาโรเจสิ.\csromanspacer{} เสฏฺฐิ คหปติ ตํ ภิกฺขุนิํ ฆรํ เนตฺวา โภเชตฺวา อุชฺฌายติ ขิยฺยติ วิปาเจติ “กถํ หิ นาม ภทนฺตา ภิกฺขุนิยา หตฺถโต อามิสํ ปฏิคฺคเหสฺสนฺติ.\csromanspacer{} กิจฺฉลาโภ มาตุคาโม”ติ.}

\prose{อสฺโสสุํ โข ภิกฺขู ตสฺส เสฏฺฐิสฺส คหปติสฺส อุชฺฌายนฺตสฺส ขิยฺยนฺตสฺส วิปาเจนฺตสฺส.\csromanspacer{} เย เต ภิกฺขู อปฺปิจฺฉา ฯเปฯ เต อุชฺฌายนฺติ ขิยฺยนฺติ วิปาเจนฺติ “กถํ หิ นาม ภิกฺขุ ภิกฺขุนิยา หตฺถโต อามิสํ \csromanfolio{228}ปฏิคฺคเหสฺสตี”ติ ฯเปฯ.\csromanspacer{} สจฺจํ กิร ตฺวํ ภิกฺขุ ภิกฺขุนิยา หตฺถโต อามิสํ ปฏิคฺคเหสีติ.\csromanspacer{} สจฺจํ ภควาติ.\csromanspacer{} ญาติกา เต ภิกฺขุ อญฺญาติกาติ.\csromanspacer{} อญฺญาติกา ภควาติ.\csromanspacer{} อญฺญาตโก โมฆปุริส อญฺญาติกาย น ชานาติ ปติรูปํ วา อปฺปติรูปํ วา สนฺตํ วา อสนฺตํ วา, กถํ หิ นาม ตฺวํ โมฆปุริส อญฺญาติกาย ภิกฺขุนิยา หตฺถโต อามิสํ ปฏิคฺคเหสฺสสิ.\csromanspacer{} เนตํ โมฆปุริส อปฺปสนฺนานํ วา ปสาทาย ฯเปฯ.\csromanspacer{} เอวญฺจ ปน ภิกฺขเว อิมํ สิกฺขาปทํ อุทฺทิเสยฺยาถ\csromandash{}}

\proseitem{553}{\textbf{“โย ปน ภิกฺขุ อญฺญาติกาย ภิกฺขุนิยา อนฺตรฆรํ ปวิฏฺฐาย หตฺถโต ขาทนียํ วา โภชนียํ วา สหตฺถา ปฏิคฺคเหตฺวา ขาเทยฺย วา ภุญฺเชยฺย วา, ปฏิเทเสตพฺพํ เตน ภิกฺขุนา คารยฺหํ อาวุโส ธมฺมํ อาปชฺชิํ อสปฺปายํ ปาฏิเทสนียํ, ตํ ปฏิเทเสมี”}ติ.}

\proseitem{554}{\textbf{โย ปนา}ติ โย ยาทิโส ฯเปฯ.\csromanspacer{} \textbf{ภิกฺขู}ติ ฯเปฯ อยํ อิมสฺมิํ อตฺเถ อธิปฺเปโต ภิกฺขูติ.}

\prose{\textbf{อญฺญาติกา} นาม มาติโต วา ปิติโต วา ยาว สตฺตมา ปิตามหยุคา อสมฺพทฺธา.}

\prose{\textbf{ภิกฺขุนี} นาม อุภโตสํเฆ อุปสมฺปนฺนา.}

\prose{\textbf{อนฺตรฆรํ} นาม รถิกา พฺยูหํ สิงฺฆาฏกํ ฆรํ.}

\prose{\textbf{ขาทนียํ} นาม ปญฺจ โภชนานิ ยามกาลิกํ สตฺตาหกาลิกํ ยาวชีวิกํ ฐเปตฺวา อวเสสํ ขาทนียํ นาม.}

\prose{\textbf{โภชนียํ} นาม ปญฺจ โภชนานิ โอทโน กุมฺมาโส สตฺตุ มจฺโฉ มํสํ. “ขาทิสฺสามิ ภุญฺชิสฺสามี”ติ ปฏิคฺคณฺหาติ, อาปตฺติ ทุกฺกฏสฺส.\csromanspacer{} อชฺโฌหาเร อชฺโฌหาเร อาปตฺติ ปาฏิเทสนียสฺส.}

\proseitem{555}{อญฺญาติกาย อญฺญาติกสญฺญี อนฺตรฆรํ ปวิฏฺฐาย หตฺถโต ขาทนียํ วา โภชนียํ วา สหตฺถา ปฏิคฺคเหตฺวา ขาทติ วา ภุญฺชติ วา, อาปตฺติ ปาฏิเทสนียสฺส.\csromanspacer{} อญฺญาติกาย เวมติโก อนฺตรฆรํ ปวิฏฺฐาย หตฺถโต ขาทนียํ วา โภชนียํ สหตฺถา ปฏิคฺคเหตฺวา ขาทติ วา ภุญฺชติ วา, อาปตฺติ ปาฏิเทสนียสฺส. \csromanfolio{229}\textbf{อญฺญาติกาย ญาติกสญฺญี อนฺตฆรํ ปวิฏฺฐาย หตฺถโต ขาทนียํ วา} \textbf{โภชนียํ สหตฺถา ปฏิคเหตฺวา ขาทติ วา ภุญฺชติ วา, อาปตฺติ} \textbf{ปาฏิเทสนียสฺส.}}

\prose{ยามกาลิกํ สตฺตาหกาลิกํ ยาวชีวิกํ อาหารตฺถาย ปฏิคฺคณฺหาติ, อาปตฺติ ทุกฺกฏสฺส.\csromanspacer{} อชฺโฌหาเร อชฺโฌหาเร อาปตฺติ ทุกฺกฏสฺส.\csromanspacer{} เอกโตอุปสมฺปนฺนาย หตฺถโต ขาทนียํ วา โภชนียํ วา “ขาทิสฺสามิ ภุญฺชิสฺสามี”ติ ปฏิคฺคณฺหาติ, อาปตฺติ ทุกฺกฏสฺส.\csromanspacer{} อชฺโฌหาเร อชฺโฌหาเร อาปตฺติ ทุกฺกฏสฺส.\csromanspacer{} ญาติกาย อญฺญาติกสญฺญี, อาปตฺติ ทุกฺกฏสฺส.\csromanspacer{} ญาติกาย เวมติโก, อาปตฺติ ทุกฺกฏสฺส.\csromanspacer{} ญาติกาย ญาติกสญฺญี, อนาปตฺติ.}

\proseitemwithcloserruled{556}{อนาปตฺติ ญาติกาย, ทาเปติ น เทติ, อุปนิกฺขิปิตฺวา เทติ, อนฺตราราเม, ภิกฺขุนุปสฺสเย, ติตฺถิยเสยฺยาย, ปฏิกฺกมเน, คามโต นีหริตฺวา เทติ, ยามกาลิกํ สตฺตาหกาลิกํ ยาวชีวิกํ “สติ ปจฺจเย ปริภุญฺชา”ติ เทติ, สิกฺขมานาย สามเณริยา อุมฺมตฺตกสฺส อาทิกมฺมิกสฺสาติ.}{ปฐม\-ปาฏิเทสนีย\-สิกฺขาปทํ นิฏฺฐิตํ.}
\csromanmatikaanchor{mtk.105}
\csromantocmarkref{section}{2. ทุติยปาฏิเทสนียสิกฺขาปท}{mtk.105}
\bookmark[dest={mtk.105},level=1]{2. ทุติยปาฏิเทสนียสิกฺขาปท}
\setcsromanheadh{2. ทุติยปาฏิเทสนียสิกฺขาปท}
\csromanheader{1}{2. ทุติย\-ปาฏิเทสนีย\-สิกฺขาปท}
\proseitem{557}{เตน สมเยน พุทฺโธ ภควา ราชคเห วิหรติ เวฬุวเน กลนฺทกนิวาเป.\csromanspacer{} เตน โข ปน สมเยน ภิกฺขู กุเลสุ นิมนฺติตา ภุญฺชนฺติ, ฉพฺพคฺคิยา ภิกฺขุนิโย ฉพฺพคฺคิยานํ ภิกฺขูนํ โวสาสนฺติโย ฐิตา โหนฺติ “อิธ สูปํ เทถ อิธ โอทนํ เทถา”ติ.\csromanspacer{} ฉพฺพคฺคิยา ภิกฺขู ยาวทตฺถํ ภุญฺชนฺติ, อญฺเญ ภิกฺขู น จิตฺตรูปํ ภุญฺชนฺติ.\csromanspacer{} เย เต ภิกฺขู อปฺปิจฺฉา ฯเปฯ เต อุชฺฌายนฺติ ขิยฺยนฺติ วิปาเจนฺติ “กถํ หิ นาม ฉพฺพคฺคิยา ภิกฺขู ภิกฺขุนิโย โวสาสนฺติโย น นิวาเรสฺสนฺตี”ติ ฯเปฯ.\csromanspacer{} สจฺจํ กิร ตุเมฺห ภิกฺขเว ภิกฺขุนิโย โวสาสนฺติโย น นิวาเรถาติ.\csromanspacer{} สจฺจํ ภควาติ.\csromanspacer{} วิครหิ พุทฺโธ ภควา ฯเปฯ.\csromanspacer{} กถํ หิ นาม ตุเมฺห โมฆปุริสา ภิกฺขุนิโย โวสาสนฺติโย น นิวาเรสฺสถ.\csromanspacer{} เนตํ โมฆปุริสา อปฺปสนฺนานํ วา ปสาทาย ฯเปฯ.\csromanspacer{} เอวญฺจ ปน ภิกฺขเว อิมํ สิกฺขาปทํ อุทฺทิเสยฺยาถ\csromandash{}}

\proseitem{558}{\csromanfolio{230}\textbf{“ภิกฺขู ปเนว กุเลสุ นิมนฺติตา ภุญฺชนฺติ, ตตฺร เจ สา}\footnote{ตตฺร เจ (สฺยา)} \textbf{ภิกฺขุนี โวสาสมานรูปา ฐิตา โหติ ‘อิธ สูปํ เทถ อิธ โอทนํ เทถา’ติ, เตหิ ภิกฺขูหิ สา ภิกฺขุนี อปสาเทตพฺพา ‘อปสกฺก ตาว ภคินิ ยาว ภิกฺขู ภุญฺชนฺตี’ติ.}\csromanspacer{} เอกสฺส เจปิ\footnote{เอกสฺสปิ เจ (สี, สฺยา)} \textbf{ภิกฺขุโน น ปฏิภาเสยฺย ตํ ภิกฺขุนิํ อปสาเทตุํ ‘อปสกฺก ตาว ภคินิ ยาว ภิกฺขู ภุญฺชนฺตี’ติ, ปฏิเทเสตพฺพํ เตหิ ภิกฺขูหิ คารยฺหํ อาวุโส ธมฺมํ อาปชฺชิมฺหา อสปฺปายํ ปาฏิเทสนียํ, ตํ ปฏิเทเสมา”}ติ.}

\proseitem{559}{\textbf{ภิกฺขูปเนว กุเลสุ นิมนฺติตา ภุญฺชนฺตี}ติ กุลํ นาม จตฺตาริ กุลานิ ขตฺติยกุลํ พฺราหฺมณกุลํ เวสฺสกุลํ สุทฺทกุลํ.}

\prose{\textbf{นิมนฺติตา ภุญฺชนฺตี}ติ ปญฺจนฺนํ โภชนานํ อญฺญตเรน โภชเนน นิมนฺติตา ภุญฺชนฺติ.}

\prose{\textbf{ภิกฺขุนี} นาม อุภโตสํเฆ อุปสมฺปนฺนา.}

\prose{\textbf{โวสาสนฺตี} นาม ยถามิตฺตตา ยถาสนฺทิฏฺฐตา ยถาสมฺภตฺตตา ยถาสมานุปชฺฌายกตา ยถาสมานาจริยกตา “อิธ สูปํ เทถ อิธ โอทนํ เทถา”ติ, เอสา โวสาสนฺตี นาม.}

\prose{\textbf{เตหิ ภิกฺขูหี}ติ ภุญฺชมาเนหิ ภิกฺขูหิ.}

\prose{\textbf{สา ภิกฺขุนี}ติ ยา สา โวสาสนฺตี ภิกฺขุนี.}

\prose{เตหิ ภิกฺขุหิ สา ภิกฺขุนี อปสาเทตพฺพา “อปสกฺก ตาว ภคินิ ยาว ภิกฺขู ภุญฺชนฺตี”ติ.\csromanspacer{} เอกสฺส เจปิ2 ภิกฺขุโน อนปสาทิโต\footnote{อนปสาทิเต (สี, สฺยา)} “ขาทิสฺสามิ ภุญฺชิสฺสามี”ติ ปฏิคฺคณฺหาติ, อาปตฺติ ทุกฺกฏสฺส.\csromanspacer{} อชฺโฌหาเร อชฺโฌหาเร อาปตฺติ ปาฏิเทสนียสฺส.}

\proseitem{560}{อุปสมฺปนฺนาย อุปสมฺปนฺนสญฺญี โวสาสนฺติยา น นิวาเรติ, อาปตฺติ ปาฏิเทสนียสฺส.\csromanspacer{} อุปสมฺปนฺนาย เวมติโก โวสาสนฺติยา น นิวาเรติ, อาปตฺติ ปาฏิเทสนียสฺส.\csromanspacer{} อุปสมฺปนฺนาย อนุปสมฺปนฺนสญฺญี โวสาสนฺติยา น นิวาเรติ, อาปตฺติ ปาฏิเทสนียสฺส.}

\prose{\csromanfolio{231}เอกโตอุปสมฺปนฺนาย โวสาสนฺติยา น นิวาเรติ, อาปตฺติ ทุกฺกฏสฺส.\csromanspacer{} อนุปสมฺปนฺนาย อุปสมฺปนฺนสญฺญี, อาปตฺติ ทุกฺกฏสฺส.\csromanspacer{} อนุปสมฺปนฺนาย เวมติโก, อาปตฺติ ทุกฺกฏสฺส.\csromanspacer{} อนุปสมฺปนฺนาย อนุปสมฺปนฺนสญฺญี, อนาปตฺติ.}

\proseitemwithcloserruled{561}{อนาปตฺติ อตฺตโน ภตฺตํ ทาเปติ น เทติ, อญฺเญสํ ภตฺตํ เทติ น ทาเปติ, ยํ น ทินฺนํ ตํ ทาเปติ, ยตฺถ น ทินฺนํ ตตฺถ ทาเปติ, สพฺเพสํ สมกํ ทาเปติ, สิกฺขมานา โวสาสติ, สมเณรี โวสาสติ, ปญฺจ โภชนานิ ฐเปตฺวา สพฺพตฺถ, อนาปตฺติ อุมฺมตฺตกสฺส อาทิกมฺมิกสฺสาติ.}{ทุติย\-ปาฏิเทสนีย\-สิกฺขาปทํ นิฏฺฐิตํ.}
\csromanmatikaanchor{mtk.106}
\csromantocmarkref{section}{3. ตติยปาฏิเทสนียสิกฺขาปท}{mtk.106}
\bookmark[dest={mtk.106},level=1]{3. ตติยปาฏิเทสนียสิกฺขาปท}
\setcsromanheadh{3. ตติยปาฏิเทสนียสิกฺขาปท}
\csromanheader{1}{3. ตติยปาฏิเทสนีย\-สิกฺขาปท}
\proseitem{562}{เตน สมเยน พุทฺโธ ภควา สาวตฺถิยํ วิหรติ เชตวเน อนาถปิณฺฑิกสฺส อาราเม.\csromanspacer{} เตน โข ปน สมเยน สาวตฺถิยํ อญฺญตรํ กุลํ อุภโตปสนฺนํ โหติ.\csromanspacer{} สทฺธาย วฑฺฒติ, โภเคน หายติ, ยํ ตสฺมิํ กุเล อุปฺปชฺชติ ปุเรภตฺตํ ขาทนียํ วา โภชนียํ วา, ตํ สพฺพํ ภิกฺขูนํ วิสฺสชฺเชตฺวา อปฺเปกทา อนสิตา อจฺฉนฺติ.\csromanspacer{} มนุสฺสา อุชฺฌายนฺติ ขิยฺยนฺติ วิปาเจนฺติ “กถํ หิ นาม สมณา สกฺยปุตฺติยา น มตฺตํ ชานิตฺวา ปฏิคฺคเหสฺสนฺติ.\csromanspacer{} อิเม อิเมสํ ทตฺวา อปฺเปกทา อนสิตา อจฺฉนฺตี”ติ.\csromanspacer{} อสฺโสสุํ โข ภิกฺขู เตสํ มนุสฺสานํ อุชฺฌายนฺตานํ ขิยฺยนฺตานํ วิปาเจนฺตานํ.\csromanspacer{} อถ โข เต ภิกฺขู ภควโต เอตมตฺถํ อาโรเจสุํ.\csromanspacer{} อถ โข ภควา เอตสฺมิํ นิทาเน เอตสฺมิํ ปกรเณ ธมฺมิํ กถํ กตฺวา ภิกฺขู อามนฺเตสิ อนุชานามิ ภิกฺขเว ยํ กุลํ สทฺธาย วฑฺฒติ, โภเคน หายติ, เอวรูปสฺส กุลสฺส ญตฺติทุติเยน กมฺเมน เสกฺขสมฺมุติํ ทาตุํ.\csromanspacer{} เอวญฺจ ปน ภิกฺขเว ทาตพฺพา, พฺยตฺเตน ภิกฺขุนา ปฏิพเลน สํโฆ ญาเปตพฺโพ\csromandash{}}

\proseitem{563}{“สุณาตุ เม ภนฺเต สํโฆ, อิตฺถนฺนามํ กุลํ สทฺธาย วฑฺฒติ, โภเคน หายติ, ยทิ สํฆสฺส ปตฺตกลฺลํ, สํโฆ อิตฺถนฺนามสฺส กุลสฺส เสกฺขสมฺมุติํ ทเทยฺย, เอสา ญตฺติ.}

\prose{\csromanfolio{232}สุณาตุ เม ภนฺเต สํโฆ, อิตฺถนฺนามํ กุลํ สทฺธาย วฑฺฒติ, โภเคน หายติ, สํโฆ อิตฺถนฺนามสฺส กุลสฺส เสกฺขสมฺมุติํ เทติ, ยสฺสายสฺมโต ขมติ อิตฺถนฺนามสฺส กุลสฺส เสกฺขสมฺมุติยา ทานํ, โส ตุณฺหสฺส, ยสฺส นกฺขมติ, โส ภาเสยฺย.}

\prose{ทินฺนา สํเฆน อิตฺถนฺนามสฺส กุลสฺส เสกฺขสมฺมุติ, ขมติ สํฆสฺส, ตสฺมา ตุณฺหี, เอวเมตํ ธารยามี”ติ.}

\csromanheader{2}{เอวญฺจ ปน ภิกฺขเว อิมํ สิกฺขาปทํ อุทฺทิเสยฺยาถ\csromandash{}}
\prose{\textbf{“ยานิ โข ปน ตานิ เสกฺขสมฺมตานิ กุลานิ, โย ปน ภิกฺขุ ตถารูเปสุ เสกฺขสมฺมเตสุ กุเลสุ ขาทนียํ วา โภชนียํ วา สหตฺถา ปฏิคฺคเหตฺวา ขาเทยฺย วา ภุญฺเชยฺย วา, ปฏิเทเสตพฺพํ เตน ภิกฺขุนา คารยฺหํ อาวุโส ธมฺมํ อาปชฺชิํ อสปฺปายํ ปาฏิเทสนียํ ตํ ปฏิเทเสมี”}ติ.}

\prose{เอวญฺจิทํ ภควตา ภิกฺขูนํ สิกฺขาปทํ ปญฺญตฺตํ โหติ.}

\proseitem{546}{เตน โข ปน สมเยน สาวตฺถิยํ อุสฺสโว โหติ.\csromanspacer{} มนุสฺสา ภิกฺขู นิมนฺเตตฺวา โภเชนฺติ, ตมฺปิ โข กุลํ ภิกฺขู นิมนฺเตสิ.\csromanspacer{} ภิกฺขู กุกฺกุจฺจายนฺตา นาธิวาเสนฺติ “ปฏิกฺขิตฺตํ ภควตา เสกฺขสมฺมเตสุ กุเลสุ ขาทนียํ วา โภชนียํ วา สหตฺถา ปฏิคฺคเหตฺวา ขาทิตุํ ภุญฺชิตุนฺ”ติ.\csromanspacer{} เต อุชฺฌายนฺติ ขิยฺยนฺติ วิปาเจนฺติ “กิํ นุ โข นาม อมฺหากํ ชีวิเตน ยํ อยฺยา อมฺหากํ น ปฏิคฺคณฺหนฺตี”ติ.\csromanspacer{} อสฺโสสุํ โข ภิกฺขู เตสํ มนุสฺสานํ อุชฺฌายนฺตานํ ขิยฺยนฺตานํ วิปาเจนฺตานํ.\csromanspacer{} อถ โข เต ภิกฺขู ภควโต เอตมตฺถํ อาโรเจสุํ.\csromanspacer{} อถ โข ภควา เอตสฺมิํ นิทาเน เอตสฺมิํ ปกรเณ ธมฺมิํ กถํ กตฺวา ภิกฺขู อามนฺเตสิ อนุชานามิ ภิกฺขเว นิมนฺติเตน เสกฺขสมฺมเตสุ กุเลสุ ขาทนียํ วา โภชนียํ วา สหตฺถา ปฏิคฺคเหตฺวา ขาทิตุํ ภุญฺชิตุํ.\csromanspacer{} เอวญฺจ ปน ภิกฺขเว อิมํ สิกฺขาปทํ อุทฺทิเสยฺยาถ\csromandash{}}

\prose{“\textbf{ยานิ โข ปน ตานิ เสกฺขสมฺมตานิ กุลานิ, โย ปน ภิกฺขุ} \textbf{ตถารูเปสุ เสกฺขสมฺมตฺ}เอสุ กุเลสุ ปุพฺเพ อนิมนฺติโต ขาทนียํ วา โภชนียํ วา \textbf{สหตฺถา ปฏิคฺคเหตฺวา ขาทฺ}เอยฺย วา ภุญฺเชยฺย \csromanfolio{233}\textbf{วา, ปฏิเทเสตพฺพํ เตน ภิกฺขุนา คารยฺหํ อาวุโส ธมฺมํ อาปชฺชิํ อสปฺปายํ ปาฏิเทสนียํ, ตํ ปฏิเทเสมี”}ติ.}

\prose{เอวญฺจิทํ ภควตา ภิกฺขูนํ สิกฺขาปทํ ปญฺญตฺตํ โหติ.}

\proseitem{565}{เตน โข ปน สมเยน อญฺญตโร ภิกฺขุ ตสฺส กุลสฺส กุลูปโก โหติ.\csromanspacer{} อถ โข โส ภิกฺขุ ปุพฺพณฺหสมยํ นิวาเสตฺวา ปตฺตจีวรมาทาย เยน ตํ กุลํ เตนุปสงฺกมิ, อุปสงฺกมิตฺวา ปญฺญตฺเต อาสเน นิสีทิ.\csromanspacer{} เตน โข ปน สมเยน โส ภิกฺขุ คิลาโน โหติ.\csromanspacer{} อถ โข เต มนุสฺสา ตํ ภิกฺขุํ เอตทโวจุํ “ภุญฺชถ ภนฺเต”ติ.\csromanspacer{} อถ โข โส ภิกฺขุ “ภควตา ปฏิกฺขิตฺตํ อนิมนฺติเตน เสกฺขสมฺมเตสุ กุเลสุ ขาทนียํ วา โภชนียํ วา สหตฺถา ปฏิคฺคเหตฺวา ขาทิตุํ ภุญฺชิตุนฺ”ติ กุกฺกุจฺจายนฺโต น ปฏิคฺคเหสิ, นาสกฺขิ ปิณฺฑาย จริตุํ, ฉินฺนภตฺโต อโหสิ.\csromanspacer{} อถ โข โส ภิกฺขุ อารามํ คนฺตฺวา ภิกฺขูนํ เอตมตฺถํ อาโรเจสิ.\csromanspacer{} ภิกฺขู ภควโต เอตมตฺถํ อาโรเจสุํ.\csromanspacer{} อถ โข ภควา เอตสฺมิํ นิทาเน เอตสฺมิํ ปกรเณ ธมฺมิํ กถํ กตฺวา ภิกฺขู อามนฺเตสิ อนุชานามิ ภิกฺขเว คิลาเนน ภิกฺขุนา เสกฺขสมฺมเตสุ กุเลสุ ขาทนียํ วา โภชนียํ วา สหตฺถา ปฏิคฺคเหตฺวา ขาทิตุํ ภุญฺชิตุํ.\csromanspacer{} เอวญฺจ ปน ภิกฺขเว อิมํ สิกฺขาปทํ อุทฺทิเสยฺยาถ\csromandash{}}

\proseitem{566}{\textbf{“ยานิ โข ปน ตานิ เสกฺขสมฺมตานิ กุลานิ, โย ปน ภิกฺขุ ตถารูเปสุ เสกฺขสมฺมเตสุ กุเลสุ ปุพฺเพ อนิมนฺติโต อคิลาโน ขาทนียํ วา โภชนียํ วา สหตฺถา ปฏิคฺคเหตฺวา ขาเทยฺย วา ภุญฺเชยฺย วา, ปฏิเทเสตพฺพํ เตน ภิกฺขุนา คารยฺหํ อาวุโส ธมฺมํ อาปชฺชิํ อสปฺปายํ ปาฏิเทสนียํ, ตํ ปฏิเทเสมี”}ติ.}

\proseitem{567}{\textbf{ยานิ โข ปน ตานิ เสกฺขสมฺมตานิ กุลานฺ}อีติ เสกฺขสมฺมตํ นาม กุลํ ยํ กุลํ สทฺธาย วฑฺฒติ, โภเคน หายติ, เอวรูปสฺส กุลสฺส ญตฺติทุติเยน กมฺเมน เสกฺขสมฺมุติ ทินฺนา โหติ.}

\prose{\textbf{โย ปนา}ติ โย ยาทิโส ฯเปฯ.\csromanspacer{} \textbf{ภิกฺขู}ติ ฯเปฯ อยํ อิมสฺมิํ อตฺเถ อธิปฺปฺเปโต ภิกฺขูติ.}

\prose{\textbf{ตถารูเปสุ เสกฺขสมฺมเตสุ กุเลสู}ติ เอวรูเปสุ เสกฺขสมฺมเตสุ กุเลสุ.}

\prose{\csromanfolio{234}\textbf{อนิมนฺติโต} นาม อชฺชตนาย วา สฺวาตนาย วา อนิมนฺติโต, ฆรูปจารํ โอกฺกมนฺเต นิมนฺเตติ, เอโส อนิมนฺติโต นาม.}

\prose{\textbf{นิมนฺติโต} นาม อชฺชตนาย วา สฺวาตนาย วา นิมนฺติโต, ฆรูปจารํ อโนกฺกมนฺเต นิมนฺเตติ, เอโส นิมนฺติโต นาม.}

\prose{\textbf{อคิลาโน} นาม สกฺโกติ ปิณฺฑาย จริตุํ.}

\prose{\textbf{คิลาโน} นาม น สกฺโกติ ปิณฺฑาย จริตุํ.}

\prose{\textbf{ขาทนียํ} นาม ปญฺจ โภชนานิ ยามกาลิกํ สตฺตาหกาลิกํ ยาวชีวิกํ ฐเปตฺวา อวเสสํ ขาทนียํ นาม.}

\prose{\textbf{โภชนียํ} นาม ปญฺจ โภชนานิ โอทโน กุมฺมาโส สตฺตุ มจฺโฉ มํสํ.}

\prose{อนิมนฺติโต อคิลาโน “ขาทิสฺสามิ ภุญฺชิสฺสามี”ติ ปฏิคฺคณฺหาติ, อาปตฺติ ทุกฺกฏสฺส.\csromanspacer{} อชฺโฌหาเร อชฺโฌหาเร อาปตฺติ ปาฏิเทสนียสฺส.}

\proseitem{568}{เสกฺขสมฺมเต เสกฺขสมฺมตสญฺญี อนิมนฺติโต อคิลาโน ขาทนียํ วา โภชนียํ วา สหตฺถา ปฏิคฺคเหตฺวา ขาทติ วา ภุญฺชติ วา, อาปตฺติ ปาฏิเทสนียสฺส.\csromanspacer{} เสกฺขสมฺมเต เวมติโก ฯเปฯ.\csromanspacer{} เสกฺขสมฺมเต อเสกฺขสมฺมตสญฺญี อนิมนฺติโต อคิเลโน ขาทนียํ วา โภชนียํ วา สหตฺถา ปฏิคฺคเหตฺวา ขาทติ วา ภุญฺชติ วา, อาปตฺติ ปาฏิเทสนียสฺส.}

\prose{ยามกาลิกํ สตฺตาหกาลิกํ ยาวชีวิกํ อาหารตฺถาย ปฏิคฺคณฺหาติ, อาปตฺติ ทุกฺกฏสฺส.\csromanspacer{} อชฺโฌหาเร อชฺโฌหาเร อาปตฺติ ทุกฺกฏสฺส.\csromanspacer{} อเสกฺขสมฺมเต เสกฺขสมฺมตสญฺญี, อาปตฺติ ทุกฺกฏสฺส.\csromanspacer{} อเสกฺขสมฺมเต เวมติโก, อาปตฺติ ทุกฺกฏสฺส.\csromanspacer{} อเสกฺขสมฺมเต อเสกฺขสมฺมตสญฺญี, อนาปตฺติ.}

\proseitemwithcloser{569}{อนาปตฺติ นิมนฺติตสฺส, คิลานสฺส, นิมนฺติตสฺส วา คิลานสฺส วา เสสกํ ภุญฺชติ, อญฺเญสํ ภิกฺขา ตตฺถ ปญฺญตฺตา โหติ, ฆรโต นีหริตฺวา เทนฺติ, นิจฺจภตฺเต, สลากภตฺเต, ปกฺขิเก, อุโปสถิเก, ปาฏิปทิเก, ยามกาลิกํ สตฺตาหกาลิกํ ยาวชีวิกํ “สติ ปจฺจเย ปริภุญฺชา”ติ เทติ, อุมฺมตฺตกสฺส อาทิกมฺมิกสฺสาติ.}{ตติยปาฏิเทสนีย\-สิกฺขาปทํ นิฏฺฐิตํ.}
\csromanmatikaanchor{mtk.107}
\csromantocmarkref{section}{4. จตุตฺถปาฏิเทสนียสิกฺขาปท}{mtk.107}
\bookmark[dest={mtk.107},level=1]{4. จตุตฺถปาฏิเทสนียสิกฺขาปท}
\setcsromanheadh{4. จตุตฺถปาฏิเทสนียสิกฺขาปท}
\csromanheader{1}{4. จตุตฺถ\-ปาฏิเทสนีย\-สิกฺขาปท}
\proseitem{570}{\csromanfolio{235}เตน สมเยน พุทฺโธ ภควา สกฺเกสุ วิหรติ กปิลวตฺถุสฺมิํ นิโคฺรธาราเม.\csromanspacer{} เตน โข ปน สมเยน สากิยทาสกา อวรุทฺธา โหนฺติ.\csromanspacer{} สากิยานิโย อิจฺฉนฺติ อารญฺญเกสุ เสนาสเนสุ ภตฺตํ กาตุํ.\csromanspacer{} อสฺโสสุํ โข สากิยทาสกา “สากิยานิโย กิร อารญฺญเกสุ เสนาสเนสุ ภตฺตํ กตฺตุกามา”ติ, เต มคฺเค ปริยุฏฺฐิํสุ.\csromanspacer{} สากิยานิโย ปณีตํ ขาทนียํ โภชนียํ อาทาย อารญฺญกํ เสนาสนํ อคมํสุ.\csromanspacer{} สากิยทาสกา นิกฺขมิตฺวา สากิยานิโย อจฺฉินฺทิํสุ จ ทูเสสุํ จ.\csromanspacer{} สากิยา นิกฺขมิตฺวา เต โจเร สภณฺเฑ\footnote{สห ภณฺเฑน (ก)} คเหตฺวา อุชฺฌายนฺติ ขิยฺยนฺติ วิปาเจนฺติ “กถํ หิ นาม ภทนฺตา อาราเม โจเร ปฏิวสนฺเต นาโรเจสฺสนฺตี”ติ.\csromanspacer{} อสฺโสสุํ โข ภิกฺขู สากิยานํ อุชฺฌายนฺตานํ ขิยฺยนฺตานํ วิปาเจนฺตานํ.\csromanspacer{} อถ โข เต ภิกฺขู ภควโต เอตมตฺถํ อาโรเจสุํ.\csromanspacer{} อถ โข ภควา เอตสฺมิํ นิทาเน เอตสฺมิํ ปกรเณ ธมฺมิํ กถํ กตฺวา ภิกฺขู อามนฺเตสิ เตน หิ ภิกฺขเว ภิกฺขูนํ สิกฺขาปทํ ปญฺญเปสฺสามิ ทส อตฺถวเส ปฏิจฺจ สํฆสุฏฺฐุตาย ฯเปฯ.\csromanspacer{} เอวญฺจ ปน ภิกฺขเว อิมํ สิกฺขาปทํ อุทฺทิเสยฺยาถ\csromandash{}}

\prose{\textbf{“ยานิ โข ปน ตานิ อารญฺญกานิ เสนาสนานิ สาสงฺกสมฺมตานิ สปฺปฏิภยานิ, โย ปน ภิกฺขุ ตถารูเปสุ เสนาสเนสุ}\footnote{เสนาสเนสุ วิหรนฺโต (สี, สฺยา)} \textbf{ปุพฺเพ อปฺปฏิสํวิทิตํ ขาทนียํ วา โภชนียํ วา อชฺฌาราเม สหตฺถา ปฏิคฺคเหตฺวา ขาเทยฺย วา ภุญฺเชยฺย วา, ปฏิเทเสตพฺพํ เตน ภิกฺขุนา คารยฺหํ อาวุโส ธมฺมํ อาปชฺชิํ อสปฺปายํ ปาฏิเทสนียํ, ตํ ปฏิเทเสมี”}ติ.}

\prose{เอวญฺจิทํ ภควตา ภิกฺขูนํ สิกฺขาปทํ ปญฺญตฺตํ โหติ.}

\proseitem{571}{เตน โข ปน สมเยน อญฺญตโร ภิกฺขุ อารญฺญเกสุ เสนาสเนสุ คิลาโน โหติ, มนุสฺสา ขาทนียํ วา โภชนียํ วา อาทาย อารญฺญกํ เสนาสนํ อคมํสุ.\csromanspacer{} อถ โข เต มนุสฺสา ตํ ภิกฺขุํ เอตทโวจุํ “ภุญฺชถ ภนฺเต”ติ.\csromanspacer{} อถ โข โส ภิกฺขุ “ภควตา ปฏิกฺขิตฺตํ อารญฺญเกสุ เสนาสเนสุ ขาทนียํ วา โภชนียํ วา \csromanfolio{236}สหตฺถา ปฏิคฺคเหตฺวา ขาทิตุํ ภุญฺชิตุนฺ”ติ กุกฺกุจฺจายนฺโต น ปฏิคฺคเหสิ, นาสกฺขิ ปิณฺฑาย จริตุํ\footnote{ปวิสิตุํ (ก)}, ฉินฺนภตฺโต อโหสิ.\csromanspacer{} อถ โข โส ภิกฺขุ ภิกฺขูนํ เอตมตฺถํ อาโรเจสิ.\csromanspacer{} ภิกฺขู ภควโต เอตมตฺถํ อาโรเจสุํ.\csromanspacer{} อถ โข ภควา เอตสฺมิํ นิทาเน เอตสฺมิํ ปกรเณ ธมฺมิํ กถํ กตฺวา ภิกฺขู อามนฺเตสิ อนุชานามิ ภิกฺขเว คิลาเนน ภิกฺขุนา อารญฺญเกสุ เสนาสเนสุ ปุพฺเพ อปฺปฏิสํวิทิตํ ขาทนียํ วา โภชนียํ วา สหตฺถา ปฏิคฺคเหตฺวา ขาทิตุํ ภุญฺชิตุํ.\csromanspacer{} เอวญฺจ ปน ภิกฺขเว อิมํ สิกฺขาปทํ อุทฺทิเสยฺยาถ\csromandash{}}

\proseitem{572}{\textbf{“ยานิ โข ปน ตานิ อารญฺญกานิ เสนาสนานิ สาสงฺกสมฺมตานิ สปฺปฏิภยานิ, โย ปน ภิกฺขุ ตถารูเปสุ เสนาสเนสุ ปุพฺเพ อปฺปฏิสํวิทิตํ ขาทนียํ วา โภชนียํ วา อชฺฌาราเม สหตฺถา ปฏิคฺคเหตฺวา อคิลาโน ขาเทยฺย วา ภุญฺเชยฺย วา, ปฏิเทเสตพฺพํ เตน ภิกฺขุนา คารยฺหํ อาวุโส ธมฺมํ อาปชฺชิํ อสปฺปายํ ปาฏิเทสนียํ, ตํ ปฏิเทเสมี”}ติ.}

\proseitem{573}{\textbf{ยานิ โข ปน ตานิ อารญฺญกานิ เสนาสนานฺ}อีติ อารญฺญกํ นาม เสนาสนํ ปญฺจธนุสติกํ ปจฺฉิมํ.}

\prose{\textbf{สาสงฺกํ} นาม อาราเม อารามูปจาเร โจรานํ นิวิฏฺโฐกาโส ทิสฺสติ, ภุตฺโตกาโส ทิสฺสติ, ฐิโตกาโส ทิสฺสติ, นิสินฺโนกาโส ทิสฺสติ, นิปนฺโนกาโส ทิสฺสติ.}

\prose{\textbf{สปฺปฏิภยํ} นาม อาราเม อารามูปจาเร โจเรหิ มนุสฺสา หตา ทิสฺสนฺติ, วิลุตฺตา ทิสฺสนฺติ, อาโกฏิตา ทิสฺสนฺติ.}

\prose{\textbf{โย ปนา}ติ โย ยาทิโส ฯเปฯ.\csromanspacer{} \textbf{ภิกฺขู}ติ ฯเปฯ อยํ อิมสฺมิํ อตฺเถ อธิปฺเปโต ภิกฺขูติ.}

\prose{\textbf{ตถารูเปสุ เสนาสเนสู}ติ เอวรูเปสุ เสนาสเนสุ.}

\prose{\textbf{อปฺปฏิสํวิทิตํ} นาม ปญฺจนฺนํ ปฏิสํวิทิตํ, เอตํ อปฺปฏิสํวิทิตํ นาม.\csromanspacer{} อารามํ อารามูปจารํ ฐเปตฺวา ปฏิสํวิทิตํ, เอตํ\footnote{เอตมฺปิ (สี)} อปฺปฏิสํวิทิตํ นาม.}

\prose{\csromanfolio{237}ปฏิสํวิทิตํ นาม โย โกจิ อิตฺถี วา ปุริโส วา อารามํ อารามูปจารํ อาคนฺตฺวา อาโรเจติ “อิตฺถนฺนามสฺส ภนฺเต ขาทนียํ วา โภชนียํ วา อาหริสฺสนฺตี”ติ.\csromanspacer{} สเจ สาสงฺกํ โหติ, “สาสงฺกนฺ”ติ อาจิกฺขิตพฺพํ.\csromanspacer{} สเจ สปฺปฏิภยํ โหติ, “สปฺปฏิภยนฺ”ติ อาจิกฺขิตพฺพํ.\csromanspacer{} สเจ “โหตุ ภนฺเต อาหริยิสฺสตี”ติ ภณติ.\csromanspacer{} โจรา วตฺตพฺพา “มนุสฺสา อิธูปจรนฺติ อปสกฺกถา”ติ.\csromanspacer{} ยาคุยา ปฏิสํวิทิเต ตสฺสา ปริวาโร อาหริยฺยติ, เอตํ ปฏิสํวิทิตํ นาม.\csromanspacer{} ภตฺเตน ปฏิสํวิทิเต ตสฺส ปริวาโร อาหริยฺยติ, เอตํ ปฏิสํวิทิตํนาม.\csromanspacer{} ขาทนีเยน ปฏิสํวิทิเต ตสฺส ปริวาโร อาหริยฺยติ, เอตํ ปฏิสํวิทิตํ นาม.\csromanspacer{} กุเลน ปฏิสํวิทิเต โย ตสฺมิํ กุเล มนุสฺโส ขาทนียํ วา โภชนียํ วา อาหรติ, เอตํ ปฏิสํวิทิตํ นาม.\csromanspacer{} คาเมน ปฏิสํวิทิเต โย ตสฺมิํ คาเม มนุสฺโส ขาทนียํ วา โภชนียํ วา อาหรติ, เอตํ ปฏิสํวิทิตํ นาม.\csromanspacer{} ปูเคน ปฏิสํวิทิเต โย ตสฺมิํ ปูเค มนุสฺโส ขาทนียํ วา โภชนียํ วา อาหรติ, เอตํ ปฏิสํวิทิตํ นาม.}

\prose{\textbf{ขาทนียํ} นาม มญฺจ โภชนานิ ยามกาลิกํ สตฺตาหกาลิกํ ยาวชีวิกํ ฐเปตฺวา อวเสสํ ขาทนียํ นาม.}

\prose{\textbf{โภชนียํ} นาม ปญฺจ โภชนานิ โอทโน กุมฺมาโส สตฺตุ มจฺโฉ มํสํ.}

\prose{\textbf{อชฺฌาราโม} นาม ปริกฺขิตฺตสฺส อารามสฺส อนฺโตอาราโม.\csromanspacer{} อปริกฺขิตฺตสฺส อุปจาโร.}

\prose{\textbf{อคิลาโน} นาม สกฺโกติ ปิณฺฑาย จริตุํ 1.}

\prose{\textbf{คิลาโน} นาม น สกฺโกติ ปิณฺฑาย จริตุํ\footnote{คนฺตุํ (ก)}.}

\prose{อปฺปฏิสํวิทิตํ อคิลาโน “ขาทิสฺสามิ ภุญฺชิสฺสามี”ติ ปฏิคฺคณฺหาติ, อาปตฺติ ทุกฺกฏสฺส.\csromanspacer{} อชฺโฌหาเร อชฺโฌหาเร อาปตฺติ ปาฏิเทสนียสฺส.}

\proseitem{574}{อปฺปฏิสํวิทิเต อปฺปฏิสํวิทิตสญฺญี ขาทนียํ วา โภชนียํ อชฺฌาราเม สหตฺถา ปฏิคฺคเหตฺวา อคิลาโน ขาทติ วา ภุญฺชติ วา, อาปตฺติ ปาฏิเทสนียสฺส.\csromanspacer{} อปฺปฏิสํวิทิเต เวมติโก ขาทนียํ วา \csromanfolio{238}โภชนียํ วา อชฺฌาราเม สหตฺถา ปฏิคฺคเหตฺวา อคิลาโน ขาทติ วา ภุญฺชติ วา, อาปตฺติ ปาฏิเทสนียสฺส.\csromanspacer{} อปฺปฏิสํวิทิเต ปฏิสํวิทิตสญฺญี ขาทนียํ วา โภชนียํ วา อชฺฌาราเม สหตฺถา ปฏิคฺคเหตฺวา อคิลาโน ขาทติ วา ภุญฺชติ วา, อาปตฺติ ปาฏิเทสนียสฺส.}

\prose{ยามกาลิกํ สตฺตาหกาลิกํ ยาวชีวิกํ อาหารตฺถาย ปฏิคฺคณฺหาติ, อาปตฺติ ทุกฺกฏสฺส.\csromanspacer{} อชฺโฌหาเร อชฺโฌหาเร อาปตฺติ ทุกฺกฏสฺส.\csromanspacer{} ปฏิสํวิทิเต อปฺปฏิสํวิทิตสญฺญี, อาปตฺติ ทุกฺกฏสฺส.\csromanspacer{} ปฏิสํวิทิเต เวมติโก, อาปตฺติ ทุกฺกฏสฺส.\csromanspacer{} ปฏิสํวิทิเต ปฏิสํวิทิตสญฺญี, อนาปตฺติ.}

\proseitemwithcloser{575}{อนาปตฺติ ปฏิสํวิทิเต, คิลานสฺส, ปฏิสํวิทิเต วา คิลานสฺส วา เสสกํ ภุญฺชติ, พหาราเม ปฏิคฺคเหตฺวา อนฺโตอาราเม ภุญฺชติ, ตตฺถ ชาตกํ มูลํ วา ตจํ วา ปตฺตํ วา ปุปฺผํ วา ผลํ วา ภุญฺชติ, ยามกาลิกํ สตฺตาหกาลิกํ ยาวชีวิกํ สติ ปจฺจเย ปริภุญฺชติ, อุมฺมตฺตกสฺส อาทิกมฺมิกสฺสาติ.}{จตุตฺถ\-ปาฏิเทสนีย\-สิกฺขาปทํ นิฏฺฐิตํ.}
\prosewithcloser{อุทฺทิฏฺฐา โข อายสฺมนฺโต จตฺตาโร ปาฏิเทสนียา ธมฺมา, ตตฺถายสฺมนฺเต ปุจฺฉามิ กจฺจิตฺถ ปริสุทฺธา, ทุติยมฺปิ ปุจฺฉามิ กจฺจิตฺถ ปริสุทฺธา, ตติยมฺปิ ปุจฺฉามิ กจฺจิตฺถ ปริสุทฺธา, ปริสุทฺเธตฺถายสฺมนฺโต, ตสฺมา ตุณฺหี, เอวเมตํ ธารยามีติ.}{\textbf{ปาฏิเทสนียกณฺฑํ นิฏฺฐิตํ.}}
\csromanmatikaanchor{mtk.108}
\csromantocmarknopage{chapter}{7. เสขิยกณฺฑ}{mtk.108}
\bookmark[dest={mtk.108},level=0]{7. เสขิยกณฺฑ}
\setcsromanheadcha{7. เสขิยกณฺฑ}
\setcsromanheadhclear{}
\chapterheadpageread{7. \textbf{เสขิยกณฺฑ}}
\csromanmatikaanchor{mtk.109}
\csromantocmarkref{section}{1. ปริมณฺฑลวคฺค}{mtk.109}
\bookmark[dest={mtk.109},level=1]{1. ปริมณฺฑลวคฺค}
\setcsromanheadh{1. ปริมณฺฑลวคฺค}
\csromanheader{1}{1. \textbf{ปริมณฺฑลวคฺค}}
\prose{\csromanfolio{239}อิเม โข ปนายสฺมนฺโต เสขิยา ธมฺมา อุทฺเทสํ อาคจฺฉนฺติ.}

\proseitem{576}{เตน สมเยน พุทฺโธ ภควา สาวตฺถิยํ วิหรติ เชตวเน อนาถปิณฺฑิกสฺส อาราเม.\csromanspacer{} เตน โข ปน สมเยน ฉพฺพคฺคิยา ภิกฺขู ปุรโตปิ ปจฺฉโตปิ โอลมฺเพนฺตา นิวาเสนฺติ.\csromanspacer{} มนุสฺสา อุชฺฌายนฺติ ขิยฺยนฺติ วิปาเจนฺติ “กถํ หิ นาม สมณา สกฺยปุตฺติยา ปุรโตปิ ปจฺฉโตปิ โอลมฺเพนฺตา นิวาเสสฺสนฺติ เสยฺยถาปิ คิหี กามโภคิโน”ติ.\csromanspacer{} อสฺโสสุํ โข ภิกฺขู เตสํ มนุสฺสานํ อุชฺฌายนฺตานํ ขิยฺยนฺตานํ วิปาเจนฺตานํ.\csromanspacer{} เย เต ภิกฺขู อปฺปิจฺฉา ฯเปฯ เต อุชฺฌายนฺติ ขิยฺยนฺติ วิปาเจนฺติ “กถํ หิ นาม ฉพฺพคฺคิยา ภิกฺขู ปุรโตปิ ปจฺฉโตปิ โอลมฺเพนฺตา นิวาเสสฺสนฺตี”ติ.\csromanspacer{} อถ โข เต ภิกฺขู ฉพฺพคฺคิเย ภิกฺขู อเนกปริยาเยน วิครหิตฺวา ภควโต เอตมตฺถํ อาโรเจสุํ.\csromanspacer{} อถ โข ภควา เอตสฺมิํ นิทาเน เอตสฺมิํ ปกรเณ ภิกฺขุสํฆํ สนฺนิปาตาเปตฺวา ฉพฺพคฺคิเย ภิกฺขู ปฏิปุจฺฉิ “สจฺจํ กิร ตุเมฺห ภิกฺขเว ปุรโตปิ ปจฺฉโตปิ โอลมฺเพนฺตา นิวาเสถา”ติ.\csromanspacer{} สจฺจํ ภควาติ.\csromanspacer{} วิครหิ พุทฺโธ ภควา ฯเปฯ.\csromanspacer{} กถํ หิ นาม ตุเมฺห โมฆปุริสา ปุรโตปิ ปจฺฉโตปิ โอลมฺเพนฺตา นิวาเสสฺสถ.\csromanspacer{} เนตํ โมฆปุริสา อปฺปสนฺนานํ วา ปสาทาย ฯเปฯ.\csromanspacer{} เอวญฺจ ปน ภิกฺขเว อิมํ สิกฺขาปทํ อุทฺทิเสยฺยาถ\csromandash{}}

\prose{\textbf{“ปริมณฺฑลํ นิวาเสสฺสามีติ สิกฺขา กรณียา”}ติ.}

\prose{ปริมณฺฑลํ นิวาเสตพฺพํ นาภิมณฺฑลํ ชาณุมณฺฑลํ ปฏิจฺฉาเทนฺเตน.\csromanspacer{} โย อนาทริยํ ปฏิจฺจ ปุรโต วา ปจฺฉโต วา โอลมฺเพนฺโต นิวาเสติ, อาปตฺติ ทุกฺกฏสฺส.}

\prosewithcloser{อนาปตฺติ อสญฺจิจฺจ อสฺสติยา อชานนฺตสฺส คิลานสฺส อาปทาสุ อุมฺมตฺตกสฺส อาทิกมฺมิกสฺสาติ.}{ปฐมสิกฺขาปทํ นิฏฺฐิตํ.}
\proseitem{577}{\csromanfolio{240}เตน สมเยน พุทฺโธ ภควา สาวตฺถิยํ วิหรติ เชตวเน อนาถปิณฺฑิกสฺส อาราเม.\csromanspacer{} เตน โข ปน สมเยน ฉพฺพคฺคิยา ภิกฺขู ปุรโตปิ ปจฺฉโตปิ โอลมฺเพนฺตา ปารุปนฺติ ฯเปฯ. \csromandash{}}

\prose{\textbf{“ปริมณฺฑลํ ปารุปิสฺสามีติ สิกฺขา กรณียา”}ติ.}

\prosewithcloser{ปริมณฺฑลํ ปารุปิตพฺพํ อุโภ กณฺเณ สมํ กตฺวา.\csromanspacer{} โย อนาทริยํ ปฏิจฺจ ปุรโต วา ปจฺฉโต วา โอลมฺเพนฺโต ปารุปติ, อาปตฺติ ทุกฺกฏสฺส.\csromanspacer{} อนาปตฺติ อสญฺจิจฺจ อสฺสติยา อชานนฺตสฺส คิลานสฺส อาปทาสุ อุมฺมตฺตกสฺส อาทิกมฺมิกสฺสาติ.}{ทุติยสิกฺขาปทํ นิฏฺฐิตํ.}
\proseitem{578}{เตน สมเยน พุทฺโธ ภควา สาวตฺถิยํ วิหรติ เชตวเน อนาถปิณฺฑิกสฺส อาราเม.\csromanspacer{} เตน โข ปน สมเยน ฉพฺพคฺคิยา ภิกฺขู กายํ วิวริตฺวา อนฺตรฆเร คจฺฉนฺติ ฯเปฯ. \csromandash{}}

\prose{\textbf{“สุปฺปฏิจฺฉนฺโน อนฺตรฆเร คมิสฺสามีติ สิกฺขา กรณียา”}ติ.}

\prose{สุปฺปฏิจฺฉนฺเนน อนฺตรฆเร คนฺตพฺพํ.\csromanspacer{} โย อนาทริยํ ปฏิจฺจ กายํ วิวริตฺวา อนฺตรฆเร คจฺฉติ, อาปตฺติ ทุกฺกฏสฺส.}

\prosewithcloser{อนาปตฺติ อสญฺจิจฺจ อสฺสติยา อชานนฺตสฺส คิลานสฺส อาปทาสุ อุมฺมตฺตกสฺส อาทิกมฺมิกสฺสาติ.}{ตติยสิกฺขาปทํ นิฏฺฐิตํ.}
\proseitem{579}{เตน สมเยน พุทฺโธ ภควา สาวตฺถิยํ วิหรติ เชตวเน อนาถปิณฺฑิกสฺส อาราเม.\csromanspacer{} เตน โข ปน สมเยน ฉพฺพคฺคิยา ภิกฺขู กายํ วิวริตฺวา อนฺตรฆเร นิสีทนฺติ ฯเปฯ. \csromandash{}}

\prose{\textbf{“สุปฺปฏิจฺฉนฺโน อนฺตรฆเร นิสีทิสฺสามีติ สิกฺขา กรณียา”}ติ.}

\prose{สุปฺปฏิจฺฉนฺเนน อนฺตรฆเร นิสีทิตพฺพํ.\csromanspacer{} โย อนาทริยํ ปฏิจฺจ กายํ วิวริตฺวา อนฺตรฆเร นิสีทติ, อาปตฺติ ทุกฺกฏสฺส.}

\prosewithcloser{อนาปตฺติ อสญฺจิจฺจ อสฺสติยา อชานนฺตสฺส คิลานสฺส วาสูปคตสฺส อาปทาสุ อุมฺมตฺตกสฺส อาทิกมฺมิกสฺสาติ.}{จตุตฺถสิกฺขาปทํ นิฏฺฐิตํ.}
\proseitem{580}{\csromanfolio{241}เตน สมเยน พุทฺโธ ภควา สาวตฺถิยํ วิหรติ เชตวเน อนาถปิณฺฑิกสฺส อาราเม.\csromanspacer{} เตน โข ปน สมเยน ฉพฺพคฺคิยา ภิกฺขู หตฺถมฺปิ ปาทมฺปิ กีฬาเปนฺตา อนฺตรฆเร คจฺฉนฺติ ฯเปฯ. \csromandash{}}

\prose{\textbf{“สุสํวุโต อนฺตรฆเร คมิสฺสามีติ สิกฺขา กรณียา”}ติ.}

\prose{สุสํวุเตน อนฺตรฆเร คนฺตพฺพํ.\csromanspacer{} โย อนาทริยํ ปฏิจฺจ หตฺถํ วา ปาทํ วา กีฬาเปนฺโต อนฺตรฆเร คจฺฉติ, อาปตฺติ ทุกฺกฏสฺส.}

\prosewithcloser{อนาปตฺติ อสญฺจิจฺจ อสฺสติยา อชานนฺตสฺส คิลานสฺส อุมฺมตฺตกสฺส อาทิกมฺมิกสฺสาติ.}{ปญฺจมสิกฺขาปทํ นิฏฺฐิตํ.}
\proseitem{581}{เตน สมเยน พุทฺโธ ภควา สาวตฺถิยํ วิหรติ เชตวเน อนาถปิณฺฑิกสฺส อาราเม.\csromanspacer{} เตน โข ปน สมเยน ฉพฺพคฺคิยา ภิกฺขู หตฺถมฺปิ ปาทมฺปิ กีฬาเปนฺตา อนฺตรฆเร นิสีทนฺติ ฯเปฯ. \csromandash{}}

\prose{\textbf{“สุสํวุโต อนฺตรฆเร นิสีทิสฺสามีติ สิกฺขา กรณียา”}ติ.}

\prose{สุสํวุเตน อนฺตรฆเร นิสีทิตพฺพํ.\csromanspacer{} โย อนาทริยํ ปฏิจฺจ หตฺถํ วา ปาทํ วา กีฬาเปนฺโต อนฺตรฆเร นิสีทติ, อาปตฺติ ทุกฺกฏสฺส.}

\prosewithcloser{อนาปตฺติ อสญฺจิจฺจ อสฺสติยา อชานนฺตสฺส คิลานสฺส อุมฺมตฺตกสฺส อาทิกมฺมิกสฺสาติ.}{ฉฏฺฐสิกฺขาปทํ นิฏฺฐิตํ.}
\proseitem{582}{เตน สมเยน พุทฺโธ ภควา สาวตฺถิยํ วิหรติ เชตวเน อนาถปิณฺฑิกสฺส อาราเม.\csromanspacer{} เตน โข ปน สมเยน ฉพฺพคฺคิยา ภิกฺขู ตหํ ตหํ โอโลเกนฺตา อนฺตรฆเร คจฺฉนฺติ ฯเปฯ. \csromandash{}}

\prose{\textbf{“โอกฺขิตฺตจกฺขุ อนฺตรฆเร คมิสฺสามีติ สิกฺขา กรณียา”}ติ.}

\prose{โอกฺขิตฺตจกฺขุนา อนฺตรฆเร คนฺตพฺพํ ยุคมตฺตํ เปกฺขนฺเตน.\csromanspacer{} โย อนาทริยํ ปฏิจฺจ ตหํ ตหํ โอโลเกนฺโต อนฺตรฆเร คจฺฉติ, อาปตฺติ ทุกฺกฏสฺส.}

\prosewithcloser{อนาปตฺติ อสญฺจิจฺจ อสฺสติยา อชานนฺตสฺส คิลานสฺส อาปทาสุ อุมฺมตฺตกสฺส อาทิกมฺมิกสฺสาติ.}{สตฺตมสิกฺขาปทํ นิฏฺฐิตํ.}
\proseitem{583}{\csromanfolio{242}เตน สมเยน พุทฺโธ ภควา สาวตฺถิยํ วิหรติ เชตวเน อนาถปิณฺฑิกสฺส อาราเม.\csromanspacer{} เตน โข ปน สมเยน ฉพฺพคฺคิยา ภิกฺขู ตหํ ตหํ โอโลเกนฺตา อนฺตรฆเร นิสีทนฺติ ฯเปฯ. \csromandash{}}

\prose{\textbf{“โอกฺขิตฺตจกฺขุ อนฺตรฆเร นิสีทิสฺสามีติ สิกฺขา กรณียา”}ติ.}

\prose{โอกฺขิตฺตจกฺขุนา อนฺตรฆเร นิสีทิตพฺพํ ยุคมตฺตํ เปกฺขนฺเตน.\csromanspacer{} โย อนาทริยํ ปฏิจฺจ ตหํ ตหํ โอโลเกนฺโต อนฺตรฆเร นิสีทติ, อาปตฺติ ทุกฺกฏสฺส.}

\prosewithcloser{อนาปตฺติ อสญฺจิจฺจ อสฺสติยา อชานนฺตสฺส คิลานสฺส อาปทาสุ อุมฺมตฺตกสฺส อาทิกมฺมิกสฺสาติ.}{อฏฺฐมสิกฺขาปทํ นิฏฺฐิตํ.}
\proseitem{584}{เตน สมเยน พุทฺโธ ภควา สาวตฺถิยํ วิหรติ เชตวเน อนาถปิณฺฑิกสฺส อาราเม.\csromanspacer{} เตน โข ปน สมเยน ฉพฺพคฺคิยา ภิกฺขู อุกฺขิตฺตกาย อนฺตรฆเร คจฺฉนฺติ ฯเปฯ. \csromandash{}}

\prose{\textbf{“น อุกฺขิตฺตกาย อนฺตรฆเร คมิสฺสามีติ สิกฺขา กรณียา”}ติ.}

\prose{\textbf{น อุกฺขิตฺตกาย อนฺตรคฺ}หเร คนฺตพฺพํ.\csromanspacer{} โย อนาทริยํ ปฏิจฺจ เอกโต วา อุภโต วา อุกฺขิปิตฺวา อนฺตรฆเร คจฺฉติ, อาปตฺติ ทุกฺกฏสฺส.}

\prosewithcloser{อนาปตฺติ อสญฺจิจฺจ อสฺสติยา อชานนฺตสฺส คิลานสฺส อาปทาสุ อุมฺมตฺตกสฺส อาทิกมฺมิกสฺสาติ.}{นวมสิกฺขาปทํ นิฏฺฐิตํ.}
\proseitem{585}{เตน สมเยน พุทฺโธ ภควา สาวตฺถิยํ วิหรติ เชตวเน อนาถปิณฺฑิกสฺส อาราเม.\csromanspacer{} เตน โข ปน สมเยน ฉพฺพคฺคิยา ภิกฺขู อุกฺขิตฺตกาย อนฺตรฆเร นิสีทนฺติ ฯเปฯ. \csromandash{}}

\prose{“\textbf{น อุกฺขิตฺตกาย อนฺตรคฺ}หเร นิสีทิสฺสามีติ สิกฺขา กรณียา”ติ.}

\prose{\textbf{น อุกฺขิตฺตกาย อนฺตรคฺ}หเร นิสีทิตพฺพํ.\csromanspacer{} โย อนาทริยํ ปฏิจฺจ เอกโต วา อุภโต วา อุกฺขิปิตฺวา อนฺตรฆเร นิสีทติ, อาปตฺติ ทุกฺกฏสฺส.}

\prosewithcloser{\csromanfolio{243}อนาปตฺติ อสญฺจิจฺจ อสฺสติยา อชานนฺตสฺส คิลานสฺส วาสูปคตสฺส อาปทาสุ อุมฺมตฺตกสฺส อาทิกมฺมิกสฺสาติ.}{ทสมสิกฺขาปทํ นิฏฺฐิตํ.}
\csromanheader{2}{ปริมณฺฑลวคฺโค ปฐโม.}
\csromansectionrule
\csromanmatikaanchor{mtk.110}
\csromantocmarkref{section}{2. อุชฺชคฺฆิกวคฺค}{mtk.110}
\bookmark[dest={mtk.110},level=1]{2. อุชฺชคฺฆิกวคฺค}
\setcsromanheadh{2. อุชฺชคฺฆิกวคฺค}
\csromanheader{1}{2. \textbf{อุชฺชคฺฆิกวคฺค}}
\proseitem{586}{เตน สมเยน พุทฺโธ ภควา สาวตฺถิยํ วิหรติ เชตวเน อนาถปิณฺฑิกสฺส อาราเม.\csromanspacer{} เตน โข ปน สมเยน ฉพฺพคฺคิยา ภิกฺขู มหาหสิตํ หสนฺตา อนฺตรฆเร คจฺฉนฺติ ฯเปฯ. \csromandash{}}

\prose{\textbf{“น อุชฺชคฺฆิกาย อนฺตรฆเร คมิสฺสามีติ สิกฺขา กรณียา”}ติ.}

\prose{\textbf{น อุชฺชคฺฆิกาย อนฺตรคฺ}หเร คนฺตพฺพํ.\csromanspacer{} โย อนาทริยํ ปฏิจฺจ มหาหสิตํ หสนฺโต อนฺตรฆเร คจฺฉติ, อาปตฺติ ทุกฺกฏสฺส.}

\prosewithcloser{อนาปตฺติ อสญฺจิจฺจ อสฺสติยา อชานนฺตสฺส คิลานสฺส, หสนียสฺมิํ วตฺถุสฺมิํ มิหิตมตฺถํ กโรติ, อาปทาสุ อุมฺมตฺตกสฺส อาทิกมฺมิกสฺสาติ.}{ปฐมสิกฺขาปทํ นิฏฺฐิตํ.}
\proseitem{587}{เตน สมเยน พุทฺโธ ภควา สาวตฺถิยํ วิหรติ เชตวเน อนาถปิณฺฑิกสฺส อาราเม.\csromanspacer{} เตน โข ปน สมเยน ฉพฺพคฺคิยา ภิกฺขู มหาหสิตํ หสนฺตา อนฺตรฆเร นิสีทนฺติ ฯเปฯ. \csromandash{}}

\prose{“\textbf{น อุชฺชคฺฆิกาย อนฺตรคฺ}หเร นิสีทิสฺสามีติ สิกฺขา กรณียา”ติ.}

\prose{\textbf{น อุชฺชคฺฆิกาย อนฺตรคฺ}หเร นิสีทิตพฺพํ.\csromanspacer{} โย อนาทริยํ ปฏิจฺจ มหาหสิตํ หสนฺโต อนฺตรฆเร นิสีทติ, อาปตฺติ ทุกฺกฏสฺส.}

\prosewithcloser{อนาปตฺติ อสญฺจิจฺจ อสฺสติยา อชานนฺตสฺส คิลานสฺส, หสนียสฺมิํ วตฺถุสฺมิํ มิหิตมตฺตํ กโรติ, อาปทาสุ อุมฺมตฺตกสฺส อาทิกมฺมิกสฺสาติ.}{ทุติยสิกฺขาปทํ นิฏฺฐิตํ.}
\proseitem{588}{\csromanfolio{244}เตน สมเยน พุทฺโธ ภควา สาวตฺถิยํ วิหรติ เชตวเน อนาถปิณฺฑิกสฺส อาราเม.\csromanspacer{} เตน โข ปน สมเยน ฉพฺพคฺคิยา ภิกฺขู อุจฺจาสทฺทํ มหาสทฺทํ กโรนฺตา อนฺตรฆเร คจฺฉนฺติ ฯเปฯ. \csromandash{}}

\prose{\textbf{“อปฺปสทฺโท อนฺตรฆเร คมิสฺสามีติ สิกฺขา กรณียา”}ติ.}

\prose{อปฺปสทฺเทน อนฺตรฆเร คนฺตพฺพํ.\csromanspacer{} โย อนาทริยํ ปฏิจฺจ อุจฺจาสทฺทํ มหาสทฺทํ กโรนฺโต อนฺตรฆเร คจฺฉติ, อาปตฺติ ทุกฺกฏสฺส.}

\prosewithcloser{อนาปตฺติ อสญฺจิจฺจ อสฺสติยา อชานนฺตสฺส คิลานสฺส อาปทาสุ อุมฺมตฺตกสฺส อาทิกมฺมิกสฺสาติ.}{ตติยสิกฺขาปทํ นิฏฺฐิตํ.}
\proseitem{589}{เตน สมเยน พุทฺโธ ภควา สาวตฺถิยํ วิหรติ เชตวเน อนาถปิณฺฑิกสฺส อาราเม.\csromanspacer{} เตน โข ปน สมเยน ฉพฺพคฺคิยา ภิกฺขู อุจฺจาสทฺทํ มหาสทฺทํ กโรนฺตา อนฺตรฆเร นิสีทนฺติ ฯเปฯ. \csromandash{}}

\prose{\textbf{“อปฺปสทฺโท อนฺตรฆเร นิสีทิสฺสามีติ สิกฺขา กรณียา”}ติ.}

\prose{อปฺปสทฺเทน อนฺตรฆเร นิสีทิตพฺพํ.\csromanspacer{} โย อนาทริยํ ปฏิจฺจ อุจฺจาสทฺทํ มหาสทฺทํ กโรนฺโต อนฺตรฆเร นิสีทติ, อาปตฺติ ทุกฺกฏสฺส.}

\prosewithcloser{อนาปตฺติ อสญฺจิจฺจ ฯเปฯ อาทิกมฺมิกสฺสาติ.}{จตุตฺถสิกฺขาปทํ นิฏฺฐิตํ.}
\proseitem{590}{เตน สมเยน พุทฺโธ ภควา สาวตฺถิยํ วิหรติ เชตวเน อนาถปิณฺฑิกสฺส อาราเม.\csromanspacer{} เตน โข ปน สมเยน ฉพฺพคฺคิยา ภิกฺขู กายปฺปจาลกํ อนฺตรฆเร คจฺฉนฺติ กายํ โอลมฺเพนฺตา ฯเปฯ. \csromandash{}}

\prose{\textbf{“น กายปฺปจาลกํ อนฺตรฆเร คมิสฺสามีติ สิกฺขา กรณียา”}ติ.}

\prose{\textbf{น กายปฺปจาลกํ อนฺตร}ฆเร คนฺตพฺพํ, กายํ ปคฺคเหตฺวา คนฺตพฺพํ.\csromanspacer{} โย อนาทริยํ ปฏิจฺจ กายปฺปจาลกํ อนฺตรฆเร คจฺฉติ กายํ โอลมฺเพนฺโต, อาปตฺติ ทุกฺกฏสฺส.}

\prosewithcloser{อนาปตฺติ อสญฺจิจฺจ ฯเปฯ อาทิกมฺมิกสฺสาติ.}{ปญฺจมสิกฺขาปทํ นิฏฺฐิตํ.}
\proseitem{591}{\csromanfolio{245}เตน สมเยน พุทฺโธ ภควา สาวตฺถิยํ วิหรติ เชตวเน อนาถปิณฺฑิกสฺส อาราเม.\csromanspacer{} เตน โข ปน สมเยน ฉพฺพคฺคิยา ภิกฺขู กายปฺปจาลกํ อนฺตรฆเร นิสีทนฺติ กายํ โอลมฺเพนฺตา ฯเปฯ. \csromandash{}}

\prose{\textbf{“น กายปฺปจาลกํ อนฺตรฆเร นิสีทิสฺสามีติ สิกฺขา กรณียา”}ติ.}

\prose{\textbf{น กายปฺปจาลกํ อนฺตรคฺ}หเร นิสีทิตพฺพํ, กายํ ปคฺคเหตฺวา นิสีทิตพฺพํ.\csromanspacer{} โย อนาทริยํ ปฏิจฺจ กายปฺปจาลกํ อนฺตรฆเร นิสีทติ กายํ โอลมฺเพนฺโต, อาปตฺติ ทุกฺกฏสฺส.}

\prosewithcloser{อนาปตฺติ อสญฺจิจฺจ อสฺสติยา อชานนฺตสฺส คิลานสฺส วาสูปคตสฺส อาปทาสุ อุมฺมตฺตกสฺส อาทิกมฺมิกสฺสาติ.}{ฉฏฺฐสิกฺขาปทํ นิฏฺฐิตํ.}
\proseitem{592}{เตน สมเยน พุทฺโธ ภควา สาวตฺถิยํ วิหรติ เชตวเน อนาถปิณฺฑิกสฺส อาราเม.\csromanspacer{} เตน โข ปน สมเยน ฉพฺพคฺคิยา ภิกฺขู พาหุปฺปจาลกํ อนฺตรฆเร คจฺฉนฺติ พาหุํ โอลมฺเพนฺตา ฯเปฯ. \csromandash{}}

\prose{\textbf{“น พาหุปฺปจาลกํ อนฺตรฆเร คมิสฺสามีติ สิกฺขา กรณียา”}ติ.}

\prose{\textbf{น พาหุปฺปจาลกํ อนฺตร}ฆเร คนฺตพฺพํ, พาหุํ ปคฺคเหตฺวา คนฺตพฺพํ.\csromanspacer{} โย อนาทริยํ ปฏิจฺจ พาหุปฺปจาลกํ อนฺตรฆเร คจฺฉติ พาหุํ โอลมฺเพนฺโต, อาปตฺติ ทุกฺกฏสฺส.}

\prosewithcloser{อนาปตฺติ อสญฺจิจฺจ ฯเปฯ อาทิกมฺมิกสฺสาติ.}{สตฺตมสิกฺขาปทํ นิฏฺฐิตํ.}
\proseitem{593}{เตน สมเยน พุทฺโธ ภควา สาวตฺถิยํ วิหรติ เชตวเน อนาถปิณฺฑิกสฺส อาราเม.\csromanspacer{} เตน โข ปน สมเยน ฉพฺพคฺคิยา ภิกฺขู พาหุปฺปจาลกํ อนฺตรฆเร นิสีทนฺติ พาหุํ โอลมฺเพนฺตา ฯเปฯ. \csromandash{}}

\prose{“\textbf{น พาหุปฺปจาลกํ อนฺตร}ฆเร นิสีทิสฺสามีติ สิกฺขา กรณียา”ติ.}

\prose{\textbf{น พาหุปฺปจาลกํ อนฺตร}ฆเร นิสีทิตพฺพํ, พาหุํ ปคฺคเหตฺวา นิสีทิตพฺพํ.\csromanspacer{} โย อนาทริยํ ปฏิจฺจ พาหุปฺปจาลกํ อนฺตรฆเร นิสีทติ พาหุํ โอลมฺเพนฺโต, อาปตฺติ ทุกฺกฏสฺส.}

\prosewithcloser{\csromanfolio{246}อนาปตฺติ อสญฺจิจฺจ อสฺสติยา อชานนฺตสฺส คิลานสฺส วาสูปคตสฺส อาปทาสุ อุมฺมตฺตกสฺส อาทิกมฺมิกสฺสาติ.}{อฏฺฐมสิกฺขาปทํ นิฏฺฐิตํ.}
\proseitem{594}{เตน สมเยน พุทฺโธ ภควา สาวตฺถิยํ วิหรติ เชตวเน อนาถปิณฺฑิกสฺส อาราเม.\csromanspacer{} เตน โข ปน สมเยน ฉพฺพคฺคิยา ภิกฺขู สีสปฺปจาลกํ อนฺตรฆเร คจฺฉนฺติ สีสํ โอลมฺเพนฺตา ฯเปฯ. \csromandash{}}

\prose{\textbf{“น สีสปฺปจาลกํ อนฺตรฆเร คมิสฺสามีติ สิกฺขา กรณียา”}ติ.}

\prose{\textbf{น สีสปฺปจาลกํ อนฺตร}ฆเร คนฺตพฺพํ, สีสํ ปคฺคเหตฺวา คนฺตพฺพํ.\csromanspacer{} โย อนาทริยํ ปฏิจฺจ สีสปฺปจาลกํ อนฺตรฆเร คจฺฉติ สีสํ โอลมฺเพนฺโต, อาปตฺติ ทุกฺกฏสฺส.}

\prosewithcloser{อนาปตฺติ อสญฺจิจฺจ ฯเปฯ อาทิกมฺมิกสฺสาติ.}{นวมสิกฺขาปทํ นิฏฺฐิตํ.}
\proseitem{595}{เตน สมเยน พุทฺโธ ภควา สาวตฺถิยํ วิหรติ เชตวเน อนาถปิณฺฑิกสฺส อาราเม.\csromanspacer{} เตน โข ปน สมเยน ฉพฺพคฺคิยา ภิกฺขู สีสปฺปจาลกํ อนฺตรฆเร นิสีทนฺติ สีสํ โอลมฺเพนฺตา ฯเปฯ.}

\prose{“\textbf{น สีสปฺปจาลกํ อนฺตร}ฆเร นิสีทิสฺสามีติ สิกฺขา กรณียา”ติ.}

\prose{\textbf{น สีสปฺปจาลกํ อนฺตร}ฆเร นิสีทิตพฺพํ, สีสํ ปคฺคเหตฺวา นิสีทิตพฺพํ.\csromanspacer{} โย อนาทริยํ ปฏิจฺจ สีสปฺปจาลกํ อนฺตรฆเร นิสีทติ สีสํ โอลมฺเพนฺโต, อาปตฺติ ทุกฺกฏสฺส.}

\prosewithcloser{อนาปตฺติ อสญฺจิจฺจ อสฺสติยา อชานนฺตสฺส คิลานสฺส วาสูปคตสฺส อาปทาสุ อุมฺมตฺตกสฺส อาทิกมฺมิกสฺสาติ.}{ทสมสิกฺขาปทํ นิฏฺฐิตํ.}
\vspace{\tipitakaparskip}
\csromancenter{อุชฺชคฺฆิกวคฺโค ทุติโย.}
\csromanmatikaanchor{mtk.111}
\csromantocmarkref{section}{3. ขมฺภกตวคฺค}{mtk.111}
\bookmark[dest={mtk.111},level=1]{3. ขมฺภกตวคฺค}
\setcsromanheadh{3. ขมฺภกตวคฺค}
\csromanheader{1}{3. \textbf{ขมฺภกตวคฺค}}
\proseitem{596}{\csromanfolio{247}เตน สมเยน พุทฺโธ ภควา สาวตฺถิยํ วิหรติ เชตวเน อนาถปิณฺฑิกสฺส อาราเม.\csromanspacer{} เตน โข ปน สมเยน ฉพฺพคฺคิยา ภิกฺขู ขมฺภกตา อนฺตรฆเร คจฺฉนฺติ ฯเปฯ. \csromandash{}}

\prose{\textbf{“น ขมฺภกโต อนฺตรฆเร คมิสฺสามีติ สิกฺขา กรณียา”}ติ.}

\prose{น ขมฺภกเตน อนฺตรฆเร คนฺตพฺพํ.\csromanspacer{} โย อนาทริยํ ปฏิจฺจ เอกโต วา อุภโต วา ขมฺภํ กตฺวา อนฺตรฆเร คจฺฉติ, อาปตฺติ ทุกฺกฏสฺส.}

\prosewithcloser{อนาปตฺติ อสญฺจิจฺจ ฯเปฯ อาทิกมฺมิกสฺสาติ.}{ปฐมสิกฺขาปทํ นิฏฺฐิตํ.}
\proseitem{597}{เตน สมเยน พุทฺโธ ภควา สาวตฺถิยํ วิหรติ เชตวเน อนาถปิณฺฑิกสฺส อาราเม.\csromanspacer{} เตน โข ปน สมเยน ฉพฺพคฺคิยา ภิกฺขู ขมฺภกตา อนฺตรฆเร นิสีทนฺติ ฯเปฯ. \csromandash{}}

\prose{\textbf{“น ขมฺภกโต อนฺตรฆเร นิสีทิสฺสามีติ สิกฺขา กรณียา”}ติ.}

\prose{น ขมฺภกเตน อนฺตรฆเร นิสีทิตพฺพํ.\csromanspacer{} โย อนาทริยํ ปฏิจฺจ เอกโต วา อุภโต วา ขมฺภํ กตฺวา อนฺตรฆเร นิสีทติ, อาปตฺติ ทุกฺกฏสฺส.}

\prosewithcloser{อนาปตฺติ อสญฺจิจฺจ อสฺสติยา อชานนฺตสฺส คิลานสฺส วาสูปคตสฺส อาปทาสุ อุมฺมตฺตกสฺส อาทิกมฺมิกสฺสาติ.}{ทุติยสิกฺขาปทํ นิฏฺฐิตํ.}
\proseitem{598}{เตน สมเยน พุทฺโธ ภควา สาวตฺถิยํ วิหรติ เชตวเน อนาถปิณฺฑิกสฺส อาราเม.\csromanspacer{} เตน โข ปน สมเยน ฉพฺพคฺคิยา ภิกฺขู สสีสํ ปารุปิตฺวา อนฺตรฆเร คจฺฉนฺติ ฯเปฯ. \csromandash{}}

\prose{\textbf{“น โอคุณฺฐิโต อนฺตรฆเร คมิสฺสามีติ สิกฺขา กรณียา”}ติ.}

\prose{น โอคุณฺฐิเตน อนฺตรฆเร คนฺตพฺพํ.\csromanspacer{} โย อนาทริยํ ปฏิจฺจ สสีสํ ปารุปิตฺวา อนฺตรฆเร คจฺฉติ, อาปตฺติ ทุกฺกฏสฺส.}

\prosewithcloser{อนาปตฺติ อสญฺจิจฺจ ฯเปฯ อาทิกมฺมิกสฺสาติ.}{ตติยสิกฺขาปทํ นิฏฺฐิตํ.}
\proseitem{599}{\csromanfolio{248}เตน สมเยน พุทฺโธ ภควา สาวตฺถิยํ วิหรติ เชตวเน อนาถปิณฺฑิกสฺส อาราเม.\csromanspacer{} เตน โข ปน สมเยน ฉพฺพคฺคิยา ภิกฺขู สสีสํ ปารุปิตฺวา อนฺตรฆเร นิสีทนฺติ ฯเปฯ. \csromandash{}}

\prose{\textbf{“น โอคุณฺฐิโต อนฺตรฆเร นิสีทิสฺสามีติ สิกฺขา กรณียา”}ติ.}

\prose{น โอคุณฺฐิเตน อนฺตรฆเร นิสีทิตพฺพํ.\csromanspacer{} โย อนาทริยํ ปฏิจฺจ สสีสํ ปารุปิตฺวา อนฺตรฆเร นิสีทติ, อาปตฺติ ทุกฺกฏสฺส.}

\prosewithcloser{อนาปตฺติ อสญฺจิจฺจ อสฺสติยา อชานนฺตสฺส คิลานสฺส วาสูปคตสฺส อาปทาสุ อุมฺมตฺตกสฺส อาทิกมฺมิกสฺสาติ.}{จตุตฺถสิกฺขาปทํ นิฏฺฐิตํ.}
\proseitem{600}{เตน สมเยน พุทฺโธ ภควา สาวตฺถิยํ วิหรติ เชตวเน อนาถปิณฺฑิกสฺส อาราเม.\csromanspacer{} เตน โข ปน สมเยน ฉพฺพคฺคิยา ภิกฺขู อุกฺกุฏิกาย อนฺตรฆเร คจฺฉนฺติ ฯเปฯ. \csromandash{}}

\prose{\textbf{“น อุกฺกุฏิกาย อนฺตรฆเร คมิสฺสามีติ สิกฺขา กรณียา”}ติ.}

\prose{\textbf{น อุกฺกุฏิกาย อนฺตรฆฺ}อเร คนฺตพฺพํ.\csromanspacer{} โย อนาทริยํ ปฏิจฺจ อุกฺกุฏิกาย อนฺตรฆเร คจฺฉติ, อาปตฺติ ทุกฺกฏสฺส.}

\prosewithcloser{อนาปตฺติ อสญฺจิจฺจ ฯเปฯ อาทิกมฺมิกสฺสาติ.}{ปญฺจมสิกฺขาปทํ นิฏฺฐิตํ.}
\proseitem{601}{เตน สมเยน พุทฺโธ ภควา สาวตฺถิยํ วิหรติ เชตวเน อนาถปิณฺฑิกสฺส อาราเม.\csromanspacer{} เตน โข ปน สมเยน ฉพฺพคฺคิยา ภิกฺขู ปลฺลตฺถิกาย\footnote{ปลฺลตฺติกาย (ก)} อนฺตรฆเร นิสีทนฺติ ฯเปฯ. \csromandash{}}

\prose{\textbf{“น ปลฺลตฺถิกาย อนฺตรฆเร นิสีทิสฺสามีติ สิกฺขา กรณียา”}ติ.}

\prose{\textbf{น ปลฺลตฺถิกาย อนฺตรคฺ}หเร นิสีทิตพฺพํ.\csromanspacer{} โย อนาทริยํ ปฏิจฺจ หตฺถปลฺลตฺถิกาย วา ทุสฺสปลฺลตฺถิกาย วา อนฺตรฆเร นิสีทติ, อาปตฺติ ทุกฺกฏสฺส.}

\prosewithcloser{อนาปตฺติ อสญฺจิจฺจ อสฺสติยา อชานนฺตสฺส คิลานสฺส วาสูปคตสฺส อาปทาสุ อุมฺมตฺตกสฺส อาทิกมฺมิกสฺสาติ.}{ฉฏฺฐสิกฺขาปทํ นิฏฺฐิตํ.}
\proseitem{602}{\csromanfolio{249}เตน สมเยน พุทฺโธ ภควา สาวตฺถิยํ วิหรติ เชตวเน อนาถปิณฺฑิกสฺส อาราเม.\csromanspacer{} เตน โข ปน สมเยน ฉพฺพคฺคิยา ภิกฺขู อสกฺกจฺจํ ปิณฺฑปาตํ ปฏิคฺคณฺหนฺติ ฉฑฺเฑตุกามา วิย ฯเปฯ. \csromandash{}}

\prose{\textbf{“สกฺกจฺจํ ปิณฺฑปาตํ ปฏิคฺคเหสฺสามีติ สิกฺขา กรณียา”}ติ.}

\prose{สกฺกจฺจํ ปิณฺฑปาโต ปฏิคฺคเหตพฺโพ.\csromanspacer{} โย อนาทริยํ ปฏิจฺจ อสกฺกจฺจํ ปิณฺฑปาตํ ปฏิคฺคณฺหาติ ฉฑฺเฑตุกาโม วิย, อาปตฺติ ทุกฺกฏสฺส.}

\prosewithcloser{อนาปตฺติ อสญฺจิจฺจ ฯเปฯ อาทิกมฺมิกสฺสาติ.}{สตฺตมสิกฺขาปทํ นิฏฺฐิตํ.}
\proseitem{603}{เตน สมเยน พุทฺโธ ภควา สาวตฺถิยํ วิหรติ เชตวเน อนาถปิณฺฑิกสฺส อาราเม.\csromanspacer{} เตน โข ปน สมเยน ฉพฺพคฺคิยา ภิกฺขู ตหํ ตหํ โอโลเกนฺตา ปิณฺฑปาตํ ปฏิคฺคณฺหนฺติ, อากิรนฺเตปิ อติกฺกนฺเตปิ\footnote{อติกฺกมนฺเตปิ (สี)} น ชานนฺติ ฯเปฯ. \csromandash{}}

\prose{\textbf{“ปตฺตสญฺญี ปิณฺฑปาตํ ปฏิคฺคเหสฺสามีติ สิกฺขา กรณียา”}ติ.}

\prose{ปตฺตสญฺญินา ปิณฺฑปาโต ปฏิคฺคเหตพฺโพ.\csromanspacer{} โย อนาทริยํ ปฏิจฺจ ตหํ ตหํ โอโลเกนฺโต ปิณฺฑปาตํ ปฏิคฺคณฺหาติ, อาปตฺติ ทุกฺกฏสฺส.}

\prosewithcloser{อนาปตฺติ อสญฺจิจฺจ ฯเปฯ อาทิกมฺมิกสฺสาติ.}{อฏฺฐมสิกฺขาปทํ นิฏฺฐิตํ.}
\proseitem{604}{เตน สมเยน พุทฺโธ ภควา สาวตฺถิยํ วิหรติ เชตวเน อนาถปิณฺฑิกสฺส อาราเม.\csromanspacer{} เตน โข ปน สมเยน ฉพฺพคฺคิยา ภิกฺขู ปิณฺฑปาตํ ปฏิคฺคณฺหนฺตา สูปญฺเญว พหุํ ปฏิคฺคณฺหนฺติ ฯเปฯ. \csromandash{}}

\prose{\textbf{“สมสูปกํ ปิณฺฑปาตํ ปฏิคฺคเหสฺสามีติ สิกฺขา กรณียา”}ติ.}

\prose{\textbf{สูโป} นาม เทฺว สูปา มุคฺคสูโป มาสสูโป หตฺถหาริโย, สมสูปโก ปิณฺฑปาโต ปฏิคฺคเหตพฺโพ.\csromanspacer{} โย อนาทริยํ ปฏิจฺจ สูปญฺเญว พหุํ ปฏิคฺคณฺหาติ, อาปตฺติ ทุกฺกฏสฺส.}

\prosewithcloser{อนาปตฺติ อสญฺจิจฺจ อสฺสติยา อชานนฺตสฺส คิลานสฺส รสรเส ญาตกานํ ปวาริตานํ อญฺญสฺสตฺถาย อตฺตโน ธเนน อาปทาสุ อุมฺมตฺตกสฺส อาทิกมฺมิกสฺสาติ.}{นวมสิกฺขาปทํ นิฏฺฐิตํ.}
\proseitem{605}{\csromanfolio{250}เตน สมเยน พุทฺโธ ภควา สาวตฺถิยํ วิหรติ เชตวเน อนาถปิณฺฑิกสฺส อาราเม.\csromanspacer{} เตน โข ปน สมเยน ฉพฺพคฺคิยา ภิกฺขู ถูปีกตํ ปิณฺฑปาตํ ปฏิคฺคณฺหนฺติ ฯเปฯ. \csromandash{}}

\prose{\textbf{“สมติตฺติกํ}\footnote{สมติตฺถิกํ (ก)} \textbf{ปิณฺฑปาตํ ปฏิคฺคเหสฺสามีติ สิกฺขา กรณียา”}ติ.}

\prose{สมติตฺติโก1 ปิณฺฑปาโต ปฏิคฺคเหตพฺโพ.\csromanspacer{} โย อนาทริยํ ปฏิจฺจ ถูปีกตํ ปิณฺฑปาตํ ปฏิคฺคณฺหาติ, อาปตฺติ ทุกฺกฏสฺส.}

\prosewithcloser{อนาปตฺติ อสญฺจิจฺจ อสฺสติยา อชานนฺตสฺส อาปทาสุ อุมฺมตฺตกสฺส อาทิกมฺมิกสฺสาติ.}{ทสมสิกฺขาปทํ นิฏฺฐิตํ.}
\csromanheader{2}{ขมฺภกตวคฺโค ตติโย.}
\csromansectionrule
\csromanmatikaanchor{mtk.112}
\csromantocmarkref{section}{4. สกฺกจฺจวคฺค}{mtk.112}
\bookmark[dest={mtk.112},level=1]{4. สกฺกจฺจวคฺค}
\setcsromanheadh{4. สกฺกจฺจวคฺค}
\csromanheader{1}{4. \textbf{สกฺกจฺจวคฺค}}
\proseitem{606}{เตน สมเยน พุทฺโธ ภควา สาวตฺถิยํ วิหรติ เชตวเน อนาถปิณฺฑิกสฺส อาราเม.\csromanspacer{} เตน โข ปน สมเยน ฉพฺพคฺคิยา ภิกฺขู อสกฺกจฺจํ ปิณฺฑปาตํ ภุญฺชนฺติ อภุญฺชิตุกามา วิย ฯเปฯ. \csromandash{}}

\prose{\textbf{“สกฺกจฺจํ ปิณฺฑปาตํ ภุญฺชิสฺสามีติ สิกฺขา กรณียา”}ติ.}

\prose{สกฺกจฺจํ ปิณฺฑปาโต ภุญฺชิตพฺโพ.\csromanspacer{} โย อนาทริยํ ปฏิจฺจ อสกฺกจฺจํ ปิณฺฑปาตํ ภุญฺชติ, อาปตฺติ ทุกฺกฏสฺส.}

\prosewithcloser{อนาปตฺติ อสญฺจิจฺจ อสฺสติยา อชานนฺตสฺส คิลานสฺส อาปทาสุ อุมฺมตฺตกสฺส อาทิกมฺมิกสฺสาติ.}{ปฐมสิกฺขาปทํ นิฏฺฐิตํ.}
\proseitem{607}{เตน สมเยน พุทฺโธ ภควา สาวตฺถิยํ วิหรติ เชตวเน อนาถปิณฺฑิกสฺส อาราเม.\csromanspacer{} เตน โข ปน สมเยน ฉพฺพคฺคิยา ภิกฺขู ตหํ ตหํ โอโลเกนฺตา ปิณฺฑปาตํ ภุญฺชนฺติ, อากิรนฺเตปิ อติกฺกนฺเตปิ น ชานนฺติ ฯเปฯ. \csromandash{}}

\prose{\csromanfolio{251}\textbf{“ปตฺตสญฺญี ปิณฺฑปาตํ ภุญฺชิสฺสามีติ สิกฺขา กรณียา”}ติ.}

\prose{ปตฺตสญฺญินา ปิณฺฑปาโต ภุญฺชิตพฺโพ.\csromanspacer{} โย อนาทริยํ ปฏิจฺจ ตหํ ตหํ โอโลเกนฺโต ปิณฺฑปาตํ ภุญฺชติ, อาปตฺติ ทุกฺกฏสฺส.}

\prosewithcloser{อนาปตฺติ อสญฺจิจฺจ ฯเปฯ อาทิกมฺมิกสฺสาติ.}{ทุติยสิกฺขาปทํ นิฏฺฐิตํ.}
\proseitem{608}{เตน สมเยน พุทฺโธ ภควา สาวตฺถิยํ วิหรติ เชตวเน อนาถปิณฺฑิกสฺส อาราเม.\csromanspacer{} เตน โข ปน สมเยน ฉพฺพคฺคิยา ภิกฺขู ตหํ ตหํ โอมสิตฺวา\footnote{โอมทฺทิตฺวา (ก)} ปิณฺฑปาตํ ภุญฺชนฺติ ฯเปฯ. \csromandash{}}

\prose{\textbf{“สปทานํ ปิณฺฑปาตํ ภุญฺชิสฺสามีติ สิกฺขา กรณียา”}ติ.}

\prose{สปทานํ ปิณฺฑปาโต ภุญฺชิตพฺโพ.\csromanspacer{} โย อนาทริยํ ปฏิจฺจ ตหํ ตหํ โอมสิตฺวา1 ปิณฺฑปาตํ ภุญฺชติ, อาปตฺติ ทุกฺกฏสฺส.}

\prosewithcloser{อนาปตฺติ อสญฺจิจฺจ อสฺสติยา อชานนฺตสฺส คิลานสฺส, อญฺเญสํ เทนฺโต โอมสติ, อญฺญสฺส ภาชเน อากิรนฺโต โอมสติ, อุตฺตริภงฺเค อาปทาสุ อุมฺมตฺตกสฺส อาทิกมฺมิกสฺสาติ.}{ตติยสิกฺขาปทํ นิฏฺฐิตํ.}
\proseitem{609}{เตน สมเยน พุทฺโธ ภควา สาวตฺถิยํ วิหรติ เชตวเน อนาถปิณฺฑิกสฺส อาราเม.\csromanspacer{} เตน โข ปน สมเยน ฉพฺพคฺคิยา ภิกฺขู ปิณฺฑปาตํ ภุญฺชนฺตา สูปญฺเญว พหุํ ภุญฺชนฺติ ฯเปฯ. \csromandash{}}

\prose{\textbf{“สมสูปกํ ปิณฺฑปาตํ ภุญฺชิสฺสามีติ สิกฺขา กรณียา”}ติ.}

\prose{\textbf{สูโป} นาม เทฺว สูปา มุคฺคสูโป มาสสูโป หตฺถหาริโย.\csromanspacer{} สมสูปโก ปิณฺฑปาโต ภุญฺชิตพฺโพ.\csromanspacer{} โย อนาทริยํ ปฏิจฺจ สูปญฺเญว พหุํ ภุญฺชติ, อาปตฺติ ทุกฺกฏสฺส.}

\prosewithcloser{อนาปตฺติ อสญฺจิจฺจ อสฺสติยา อชานนฺตสฺส คิลานสฺส รสรเส ญาตกานํ ปวาริตานํ อตฺตโน ธเนน อาปทาสุ อุมฺมตฺตกสฺส อาทิกมฺมิกสฺสาติ.}{จตุตฺถสิกฺขาปทํ นิฏฺฐิตํ.}
\proseitem{610}{\csromanfolio{252}เตน สมเยน พุทฺโธ ภควา สาวตฺถิยํ วิหรติ เชตวเน อนาถปิณฺฑิกสฺส อาราเม.\csromanspacer{} เตน โข ปน สมเยน ฉพฺพคฺคิยา ภิกฺขู ถูปกโต โอมทฺทิตฺวา ปิณฺฑปาตํ ภุญฺชนฺติ ฯเปฯ. \csromandash{}}

\prose{\textbf{“น ถูปกโต โอมทฺทิตฺวา ปิณฺฑปาตํ ภุญฺชิสฺสามีติ สิกฺขา กรณียา”}ติ.}

\prose{\textbf{น ถูปกโต โอมทฺทิตฺวา ปิณฺ}ฑปาโต ภุญฺชิตพฺโพ.\csromanspacer{} โย อนาทริยํ ปฏิจฺจ ถูปกโต โอมทฺทิตฺวา ปิณฺฑปาตํ ภุญฺชติ, อาปตฺติ ทุกฺกฏสฺส.}

\prosewithcloser{อนาปตฺติ อสญฺจิจฺจ อสฺสติยา อชานนฺตสฺส คิลานสฺส, ปริตฺตเก เสเส เอกโต สํกฑฺฒิตฺวา โอมทฺทิตฺวา ภุญฺชติ, อาปทาสุ อุมฺมตฺตกสฺส อาทิกมฺมิกสฺสาติ.}{ปญฺจมสิกฺขาปทํ นิฏฺฐิตํ.}
\proseitem{611}{เตน สมเยน พุทฺโธ ภควา สาวตฺถิยํ วิหรติ เชตวเน อนาถปิณฺฑิกสฺส อาราเม.\csromanspacer{} เตน โข ปน สมเยน ฉพฺพคฺคิยา ภิกฺขู สูปมฺปิ พฺยญฺชนมฺปิ โอทเนน ปฏิจฺฉาเทนฺติ ภิยฺโยกมฺยตํ อุปาทาย.\csromandash{}}

\prose{\textbf{“น สูปํ วา พฺยญฺชนํ วา โอทเนน ปฏิจฺฉาเทสฺสามิ ภิยฺโยกมฺยตํ อุปาทายาติ สิกฺขา กรณียา”}ติ.}

\prose{\textbf{น สูปํ วา พฺยญฺชนํ วา โอทเนน} ปฏิจฺฉาเทตพฺพํ \textbf{ภิยฺโยกมฺยตํ อุปฺ}อาทาย.\csromanspacer{} โย อนาทริยํ ปฏิจฺจ สูปํ วา พฺยญฺชนํ วา โอทเนน ปฏิจฺฉาเทติ ภิยฺโยกมฺยตํ อุปาทาย, อาปตฺติ ทุกฺกฏสฺส.}

\prosewithcloser{อนาปตฺติ อสญฺจิจฺจ อสฺสติยา อชานนฺตสฺส, สามิกา ปฏิจฺฉาเทตฺวา เทนฺติ, น ภิยฺโยกมฺยตํ อุปาทาย, อาปทาสุ อุมฺมตฺตกสฺส อาทิกมฺมิกสฺสาติ.}{ฉฏฺฐสิกฺขาปทํ นิฏฺฐิตํ.}
\proseitem{612}{เตน สมเยน พุทฺโธ ภควา สาวตฺถิยํ วิหรติ เชตวเน อนาถปิณฺฑิกสฺส อาราเม.\csromanspacer{} เตน โข ปน สมเยน ฉพฺพคฺคิยา ภิกฺขู สูปมฺปิ โอทนมฺปิ อตฺตโน อตฺถาย วิญฺญาเปตฺวา ภุญฺชนฺติ.\csromanspacer{} มนุสฺสา อุชฺฌายนฺติ ขิยฺยนฺติ วิปาเจนฺติ “กถํ หิ นาม สมณา สกฺยปุตฺติยา สูปมฺปิ โอทนมฺปิ อตฺตโน อตฺถาย วิญฺญาเปตฺวา ภุญฺชิสฺสนฺติ.\csromanspacer{} กสฺส สมฺปนฺนํ น \csromanfolio{253}มนาปํ, กสฺส สาทุํ น รุจฺจตี”ติ.\csromanspacer{} อสฺโสสุํ โข ภิกฺขู เตสํ มนุสฺสานํ อุชฺฌายนฺตานํ ขิยฺยนฺตานํ วิปาเจนฺตานํ.\csromanspacer{} เย เต ภิกฺขู อปฺปิจฺฉา ฯเปฯ เต อุชฺฌายนฺติ ขิยฺยนฺติ วิปาเจนฺติ “กถํ หิ นาม \textbf{ฉพฺพคฺคิยา ภิกฺขู สูปมฺปิ โอทนมฺปิ อตฺตโน อตฺถาย วิญฺญาเปตฺวา} \textbf{ภุญฺชิสฺสนฺตี”ติ ฯเปฯ.}\csromanspacer{}\textbf{ สจฺจํ กิร ตุเมฺห ภิกฺขเว สูปมฺปิ โอทนมฺปิ อตฺตโน} อตฺถาย วิญฺญาเปตฺวา ภุญฺชถาติ.\csromanspacer{} สจฺจํ ภควาติ.\csromanspacer{} วิครหิ พุทฺโธ ภควา ฯเปฯ.\csromanspacer{} กถํ หิ นาม ตุเมฺห โมฆปุริสา สูปมฺปิ โอทนมฺปิ อตฺตโน อตฺถาย วิญฺญาเปตฺวา ภุญฺชิสฺสถ.\csromanspacer{} เนตํ โมฆปุริสา อปฺปสนฺนานํ วา ปสาทาย ฯเปฯ.\csromanspacer{} เอวญฺจ ปน ภิกฺขเว อิมํ สิกฺขาปทํ อุทฺทิเสยฺยาถ\csromandash{}}

\prose{\textbf{“น สูปํ วา โอทนํ วา อตฺตโน อตฺถาย วิญฺญาเปตฺวา ภุญฺชิสฺสามีติ สิกฺขา กรณียา”}ติ.}

\prose{เอวญฺจิทํ ภควตา ภิกฺขูนํ สิกฺขาปทํ ปญฺญตฺตํ โหติ.}

\proseitem{613}{เตน โข ปน สมเยน ภิกฺขู คิลานา โหนฺติ.\csromanspacer{} คิลานปุจฺฉกา ภิกฺขู คิลาเน ภิกฺขู เอตทโวจุํ “กจฺจาวุโส ขมนียํ, กจฺจิ ยาปนียนฺ”ติ.\csromanspacer{} ปุพฺเพ มยํ อาวุโส สูปมฺปิ โอทนมฺปิ อตฺตโน อตฺถาย วิญฺญาเปตฺวา ภุญฺชาม, เตน โน ผาสุ โหติ.\csromanspacer{} อิทานิ ปน “ภควตา ปฏิกฺขิตฺตนฺ”ติ กุกฺกุจฺจายนฺตา น วิญฺญาเปม, เตน โน น ผาสุ โหตีติ.\csromanspacer{} ภควโต เอตมตฺถํ อาโรเจสุํ ฯเปฯ.\csromanspacer{} อนุชานามิ ภิกฺขเว คิลาเนน ภิกฺขุนา สูปมฺปิ โอทนมฺปิ อตฺตโน อตฺถาย วิญฺญาเปตฺวา ภุญฺชิตุํ.\csromanspacer{} เอวญฺจ ปน ภิกฺขเว อิมํ สิกฺขาปทํ อุทฺทิเสยฺยาถ\csromandash{}}

\prose{\textbf{“น สูปํ วา โอทนํ วา อคิลาโน อตฺตโน อตฺถาย วิญฺญาเปตฺวา ภุญฺชิสฺสามีติ สิกฺขา กรณียา”}ติ.}

\prose{\textbf{น สูปํ วา โอทนํ วา} อคิลาเนน อตฺตโน อตฺถาย วิญฺญาเปตฺวา ภุญฺชิตพฺพํ.\csromanspacer{} โย อนาทริยํ ปฏิจฺจ สูปํ วา โอทนํ วา อคิลาโน อตฺตโน อตฺถาย วิญฺญาเปตฺวา ภุญฺชติ, อาปตฺติ ทุกฺกฏสฺส.}

\prosewithcloser{อนาปตฺติ อสญฺจิจฺจ อสฺสติยา อชานนฺตสฺส คิลานสฺส ญาตกานํ ปวาริตานํ อญฺญสฺสตฺถาย อตฺตโน ธเนน อาปทาสุ อุมฺมตฺตกสฺส อาทิกมฺมิกสฺสาติ.}{สตฺตมสิกฺขาปทํ นิฏฺฐิตํ.}
\proseitem{614}{\csromanfolio{254}เตน สมเยน พุทฺโธ ภควา สาวตฺถิยํ วิหรติ เชตวเน อนาถปิณฺฑิกสฺส อาราเม.\csromanspacer{} เตน โข ปน สมเยน ฉพฺพคฺคิยา ภิกฺขู อุชฺฌานสญฺญี ปเรสํ ปตฺตํ โอโลเกนฺติ ฯเปฯ.}

\prose{\textbf{“น อุชฺฌานสญฺญี ปเรสํ ปตฺตํ โอโลเกสฺสามีติ สิกฺขา กรณียา”}ติ.}

\prose{น อุชฺฌานสญฺญินา ปเรสํ ปตฺโต โอโลเกตพฺโพ.\csromanspacer{} โย อนาทริยํ ปฏิจฺจ อุชฺฌานสญฺญี ปเรสํ ปตฺตํ โอโลเกติ, อาปตฺติ ทุกฺกฏสฺส.}

\prosewithcloser{อนาปตฺติ อสญฺจิจฺจ อสฺสติยา อชานนฺตสฺส, “ทสฺสามี”ติ วา “ทาเปสฺสามี”ติ วา โอโลเกติ, น อุชฺฌานสญฺญิสฺส อาปทาสุ อุมฺมตฺตกสฺส อาทิกมฺมิกสฺสาติ.}{อฏฺฐมสิกฺขาปทํ นิฏฺฐิตํ.}
\proseitem{615}{เตน สมเยน พุทฺโธ ภควา สาวตฺถิยํ วิหรติ เชตวเน อนาถปิณฺฑิกสฺส อาราเม.\csromanspacer{} เตน โข ปน สมเยน ฉพฺพคฺคิยา ภิกฺขู มหนฺตํ กพฬํ กโรนฺติ ฯเปฯ. \csromandash{}}

\prose{\textbf{“นาติมหนฺตํ กพฬํ กริสฺสามีติ สิกฺขา กรณียา”}ติ.}

\prose{นาติมหนฺโต กพโฬ กาตพฺโพ.\csromanspacer{} โย อนาทริยํ ปฏิจฺจ มหนฺตํ กพฬํ กโรติ, อาปตฺติ ทุกฺกฏสฺส.}

\prosewithcloser{อนาปตฺติ อสญฺจิจฺจ อสฺสติยา อชานนฺตสฺส คิลานสฺส ขชฺชเก ผลาผเล อุตฺตริภงฺเค อาปทาสุ อุมฺมตฺตกสฺส อาทิกมฺมิกสฺสาติ.}{นวมสิกฺขาปทํ นิฏฺฐิตํ.}
\proseitem{616}{เตน สมเยน พุทฺโธ ภควา สาวตฺถิยํ วิหรติ เชตวเน อนาถปิณฺฑิกสฺส อาราเม.\csromanspacer{} เตน โข ปน สมเยน ฉพฺพคฺคิยา ภิกฺขู ทีฆํ อาโลปํ กโรนฺติ ฯเปฯ. \csromandash{}}

\prose{\textbf{“ปริมณฺฑลํ อาโลปํ กริสฺสามีติ สิกฺขา กรณียา”}ติ.}

\prose{ปริมณฺฑโล อาโลโป กาตพฺโพ.\csromanspacer{} โย อนาทริยํ ปฏิจฺจ ทีฆํ อาโลปํ กโรติ, อาปตฺติ ทุกฺกฏสฺส.}

\prosewithcloser{\csromanfolio{255}อนาปตฺติ อสญฺจิจฺจ อสฺสติยา อชานนฺตสฺส คิลานสฺส ขชฺชเก ผลาผเล อุตฺตริภงฺเค อาปทาสุ อุมฺมตฺตกสฺส อาทิกมฺมิกสฺสาติ.}{ทสมสิกฺขาปทํ นิฏฺฐิตํ.}
\csromanheader{2}{สกฺกจฺจวคฺโค จตุตฺโถ.}
\csromansectionrule
\csromanmatikaanchor{mtk.113}
\csromantocmarkref{section}{5. กพฬวคฺค}{mtk.113}
\bookmark[dest={mtk.113},level=1]{5. กพฬวคฺค}
\setcsromanheadh{5. กพฬวคฺค}
\csromanheader{1}{5. \textbf{กพฬวคฺค}}
\proseitem{617}{เตน สมเยน พุทฺโธ ภควา สาวตฺถิยํ วิหรติ เชตวเน อนาถปิณฺฑิกสฺส อาราเม.\csromanspacer{} เตน โข ปน สมเยน ฉพฺพคฺคิยา ภิกฺขู อนาหเฏ กพเฬ มุขทฺวารํ วิวรนฺติ ฯเปฯ. \csromandash{}}

\prose{\textbf{“น อนาหเฏ กพเฬ มุขทฺวารํ วิวริสฺสามีติ สิกฺขา กรณียา”}ติ.}

\prose{\textbf{น อนาหเฏ กพเฬ มุขทฺวฺ}อารํ วิวริตพฺพํ.\csromanspacer{} โย อนาทริยํ ปฏิจฺจ อนาหเฏ กพเฬ มุขทฺวารํ วิวรติ, อาปตฺติ ทุกฺกฏสฺส.}

\prosewithcloser{อนาปตฺติ อสญฺจิจฺจ ฯเปฯ อาทิกมฺมิกสฺสาติ.}{ปฐมสิกฺขาปทํ นิฏฺฐิตํ.}
\proseitem{618}{เตน สมเยน พุทฺโธ ภควา สาวตฺถิยํ วิหรติ เชตวเน อนาถปิณฺฑิกสฺส อาราเม.\csromanspacer{} เตน โข ปน สมเยน ฉพฺพคฺคิยา ภิกฺขู ภุญฺชมานา สพฺพํ หตฺถํ มุเข ปกฺขิปนฺติ ฯเปฯ. \csromandash{}}

\prose{\textbf{“น ภุญฺชมาโน สพฺพํ หตฺถํ มุเข ปกฺขิปิสฺสามีติ สิกฺขากรณียา”}ติ.}

\prose{น ภุญฺชมาเนน สพฺโพ หตฺโถ มุเข ปกฺขิปิตพฺโพ.\csromanspacer{} โย อนาทริยํ ปฏิจฺจ ภุญฺชมาโน สพฺพํ หตฺถํ มุเข ปกฺขิปติ, อาปตฺติ ทุกฺกฏสฺส.}

\prosewithcloser{อนาปตฺติ อสญฺจิจฺจ ฯเปฯ อาทิกมฺมิกสฺสาติ.}{ทุติยสิกฺขาปทํ นิฏฺฐิตํ.}
\proseitem{619}{เตน สมเยน พุทฺโธ ภควา สาวตฺถิยํ วิหรติ เชตวเน อนาถปิณฺฑิกสฺส อาราเม.\csromanspacer{} เตน โข ปน สมเยน ฉพฺพคฺคิยา ภิกฺขู สกพเฬน มุเขน พฺยาหรนฺติ ฯเปฯ. \csromandash{}}

\prose{\csromanfolio{256}\textbf{“น สกพเฬน มุเขน พฺยาหริสฺสามีติ สิกฺขา กรณียา”}ติ.}

\prose{\textbf{น สกพเฬน มุเขน} พฺยาหริตพฺพํ.\csromanspacer{} โย อนาทริยํ ปฏิจฺจ สกพเฬน มุเขน พฺยาหรติ, อาปตฺติ ทุกฺกฏสฺส.}

\prosewithcloser{อนาปตฺติ อสญฺจิจฺจ ฯเปฯ อาทิกมฺมิกสฺสาติ.}{ตติยสิกฺขาปทํ นิฏฺฐิตํ.}
\proseitem{620}{เตน สมเยน พุทฺโธ ภควา สาวตฺถิยํ วิหรติ เชตวเน อนาถปิณฺฑิกสฺส อาราเม.\csromanspacer{} เตน โข ปน สมเยน ฉพฺพคฺคิยา ภิกฺขู ปิณฺฑุกฺเขปกํ ภุญฺชนฺติ ฯเปฯ. \csromandash{}}

\prose{\textbf{“น ปิณฺฑุกฺเขปกํ ภุญฺชิสฺสามีติ สิกฺขา กรณียา”}ติ.}

\prose{น ปิณฺฑุกฺเขปกํ ภุญฺชิตพฺพํ.\csromanspacer{} โย อนาทริยํ ปฏิจฺจ ปิณฺฑุกฺเขปกํ ภุญฺชติ, อาปตฺติ ทุกฺกฏสฺส.}

\prosewithcloser{อนาปตฺติ อสญฺจิจฺจ อสฺสติยา อชานนฺตสฺส คิลานสฺส ขชฺชเก ผลาผเล อาปทาสุ อุมฺมตฺตกสฺส อาทิกมฺมิกสฺสาติ.}{จตุตฺถสิกฺขาปทํ นิฏฺฐิตํ.}
\proseitem{621}{เตน สมเยน พุทฺโธ ภควา สาวตฺถิยํ วิหรติ เชตวเน อนาถปิณฺฑิกสฺส อาราเม.\csromanspacer{} เตน โข ปน สมเยน ฉพฺพคฺคิยา ภิกฺขู กพฬาวจฺเฉทกํ ภุญฺชนฺติ ฯเปฯ. \csromandash{}}

\prose{\textbf{“น กพฬาวจฺเฉทกํ ภุญฺชิสฺสามีติ สิกฺขา กรณียา”}ติ.}

\prose{น กพฬาวจฺเฉทกํ ภุญฺชิตพฺพํ.\csromanspacer{} โย อนาทริยํ ปฏิจฺจ กพฬาวจฺเฉทกํ ภุญฺชติ, อาปตฺติ ทุกฺกฏสฺส.}

\prosewithcloser{อนาปตฺติ อสญฺจิจฺจ อสฺสติยา อชานนฺตสฺส คิลานสฺส ขชฺชเก ผลาผเล อุตฺตริภงฺเค อาปทาสุ อุมฺมตฺตกสฺส อาทิกมฺมิกสฺสาติ.}{ปญฺจมสิกฺขาปทํ นิฏฺฐิตํ.}
\proseitem{622}{เตน สมเยน พุทฺโธ ภควา สาวตฺถิยํ วิหรติ เชตวเน อนาถปิณฺฑิกสฺส อาราเม.\csromanspacer{} เตน โข ปน สมเยน ฉพฺพคฺคิยา ภิกฺขู อวคณฺฑการกํ ภุญฺชนฺติ ฯเปฯ. \csromandash{}}

\prose{\csromanfolio{257}\textbf{“น อวคณฺฑการกํ ภุญฺชิสฺสามีติ สิกฺขา กรณียา”}ติ.}

\prose{น อวคณฺฑการกํ ภุญฺชิตพฺพํ.\csromanspacer{} โย อนาทริยํ ปฏิจฺจ เอกโต วา อุภโต วา คณฺฑํ กตฺวา ภุญฺชติ, อาปตฺติ ทุกฺกฏสฺส.}

\prosewithcloser{อนาปตฺติ อสญฺจิจฺจ อสฺสติยา อชานนฺตสฺส คิลานสฺส ผลาผเล อาปทาสุ อุมฺมตฺตกสฺส อาทิกมฺมิกสฺสาติ.}{ฉฏฺฐสิกฺขาปทํ นิฏฺฐิตํ.}
\proseitem{623}{เตน สมเยน พุทฺโธ ภควา สาวตฺถิยํ วิหรติ เชตวเน อนาถปิณฺฑิกสฺส อาราเม.\csromanspacer{} เตน โข ปน สมเยน ฉพฺพคฺคิยา ภิกฺขู หตฺถนิทฺธุนกํ ภุญฺชนฺติ ฯเปฯ. \csromandash{}}

\prose{\textbf{“น หตฺถนิทฺธุนกํ ภุญฺชิสฺสามีติ สิกฺขา กรณียา”}ติ.}

\prose{น หตฺถนิทฺธุนกํ ภุญฺชิตพฺพํ.\csromanspacer{} โย อนาทริยํ ปฏิจฺจ หตฺถนิทฺธุนกํ ภุญฺชติ, อาปตฺติ ทุกฺกฏสฺส.}

\prosewithcloser{อนาปตฺติ อสญฺจิจฺจ อสฺสติยา อชานนฺตสฺส คิลานสฺส, กจวรํ ฉฑฺเฑนฺโต หตฺถํ นิทฺธุนาติ, อาปทาสุ อุมฺมตฺตกสฺส อาทิกมฺมิกสฺสาติ.}{สตฺตมสิกฺขาปทํ นิฏฺฐิตํ.}
\proseitem{624}{เตน สมเยน พุทฺโธ ภควา สาวตฺถิยํ วิหรติ เชตวเน อนาถปิณฺฑิกสฺส อาราเม.\csromanspacer{} เตน โข ปน สมเยน ฉพฺพคฺคิยา ภิกฺขู สิตฺถาวการกํ ภุญฺชนฺติ ฯเปฯ. \csromandash{}}

\prose{\textbf{“น สิตฺถาวการกํ ภุญฺชิสฺสามีติ สิกฺขา กรณียา”}ติ.}

\prose{น สิตฺถาวการกํ ภุญฺชิตพฺพํ.\csromanspacer{} โย อนาทริยํ ปฏิจฺจ สิตฺถาวการกํ ภุญฺชติ, อาปตฺติ ทุกฺกฏสฺส.}

\prosewithcloser{อนาปตฺติ อสญฺจิจฺจ อสฺสติยา อชานนฺตสฺส คิลานสฺส, กจวรํ ฉฑฺเฑนฺโต สิตฺถํ ฉฑฺฑยติ, อาปทาสุ อุมฺมตฺตกสฺส อาทิกมฺมิกสฺสาติ.}{อฏฺฐมสิกฺขาปทํ นิฏฺฐิตํ.}
\proseitem{625}{เตน สมเยน พุทฺโธ ภควา สาวตฺถิยํ วิหรติ เชตวเน อนาถปิณฺฑิกสฺส อาราเม.\csromanspacer{} เตน โข ปน สมเยน ฉพฺพคฺคิยา ภิกฺขู ชิวฺหานิจฺฉารกํ ภุญฺชนฺติ ฯเปฯ. \csromandash{}}

\prose{\csromanfolio{258}\textbf{“น ชิวฺหานิจฺฉารกํ ภุญฺชิสฺสามีติ สิกฺขา กรณียา”}ติ.}

\prose{น ชิวฺหานิจฺฉารกํ ภุญฺชิตพฺพํ.\csromanspacer{} โย อนาทริยํ ปฏิจฺจ ชิวฺหานิจฺฉารกํ ภุญฺชติ, อาปตฺติ ทุกฺกฏสฺส.}

\prosewithcloser{อนาปตฺติ อสญฺจิจฺจ ฯเปฯ อาทิกมฺมิกสฺสาติ.}{นวมสิกฺขาปทํ นิฏฺฐิตํ.}
\proseitem{626}{เตน สมเยน พุทฺโธ ภควา สาวตฺถิยํ วิหรติ เชตวเน อนาถปิณฺฑิกสฺส อาราเน.\csromanspacer{} เตน โข ปน สมเยน ฉพฺพคฺคิยา ภิกฺขู จปุจปุการกํ ภุญฺชนฺติ ฯเปฯ. \csromandash{}}

\prose{\textbf{“น จปุจปุการกํ ภุญฺชิสฺสามีติ สิกฺขา กรณียา”}ติ.}

\prose{น จปุจปุการกํ ภุญฺชิตพฺพํ.\csromanspacer{} โย อนาทริยํ ปฏิจฺจ จปุจปุการกํ ภุญฺชติ, อาปตฺติ ทุกฺกฏสฺส.}

\prosewithcloser{อนาปตฺติ อสญฺจิจฺจ ฯเปฯ อาทิกมฺมิกสฺสาติ.}{ทสมสิกฺขาปทํ นิฏฺฐิตํ.}
\csromanheader{2}{กพฬวคฺโค ปญฺจโม.}
\csromansectionrule
\csromanmatikaanchor{mtk.114}
\csromantocmarkref{section}{6. สุรุสุรุวคฺค}{mtk.114}
\bookmark[dest={mtk.114},level=1]{6. สุรุสุรุวคฺค}
\setcsromanheadh{6. สุรุสุรุวคฺค}
\csromanheader{1}{6. \textbf{สุรุสุรุวคฺค}}
\proseitem{627}{เตน สมเยน พุทฺโธ ภควา โกสมฺพิยํ วิหรติ โฆสิตาราเม.\csromanspacer{} เตน โข ปน สมเยน อญฺญตเรน พฺราหฺมเณน สํฆสฺส ปโยปานํ ปฏิยตฺตํ โหติ.\csromanspacer{} ภิกฺขู สุรุสุรุการกํ ขีรํ ปิวนฺติ.\csromanspacer{} อญฺญตโร นฏปุพฺพโก ภิกฺขุ เอวมาห “สพฺโพยํ มญฺเญ สํโฆ สีตีกโต”ติ.\csromanspacer{} เย เต ภิกฺขู อปฺปิจฺฉา ฯเปฯ เต อุชฺฌายนฺติ ขิยฺยนฺติ วิปาเจนฺติ “กถํ หิ นาม ภิกฺขุ สํฆํ อารพฺภ ทวํ กริสฺสตี”ติ ฯเปฯ.\csromanspacer{} สจฺจํ กิร ตฺวํ ภิกฺขุ สํฆํ อารพฺภ ทวํ อกาสีติ.\csromanspacer{} สจฺจํ ภควาติ.\csromanspacer{} วิครหิ พุทฺโธ ภควา ฯเปฯ.\csromanspacer{} กถํ หิ นาม ตฺวํ โมฆปุริส สํฆํ อารพฺภ ทวํ กริสฺสสิ.\csromanspacer{} เนตํ โมฆปุริส อปฺปสนฺนานํ วา ปสาทาย ฯเปฯ วิครหิตฺวา ฯเปฯ ธมฺมิํ กถํ กตฺวา ภิกฺขู อามนฺเตสิ “น ภิกฺขเว พุทฺธํ วา ธมฺมํ \csromanfolio{259}วา สํฆํ วา อารพฺภ ทโว กาตพฺโพ.\csromanspacer{} โย กเรยฺย, อาปตฺติ ทุกฺกฏสฺสา”ติ.\csromanspacer{} อถ โข ภควา ตํ ภิกฺขุํ อเนกปริยาเยน วิครหิตฺวา ทุพฺภรตาย ฯเปฯ.\csromanspacer{} เอวญฺจ ปน ภิกฺขเว อิมํ สิกฺขาปทํ \textbf{อุทฺทิเสยฺยาถ\csromandash{}}}

\prose{\textbf{“น สุรุสุรุการกํ ภุญฺชิสฺสามีติ สิกฺขา กรณียา”}ติ.}

\prose{น สุรุสุรุการกํ ภุญฺชิตพฺพํ.\csromanspacer{} โย อนาทริยํ ปฏิจฺจ สุรุสุรุการกํ ภุญฺชติ, อาปตฺติ ทุกฺกฏสฺส.}

\prosewithcloser{อนาปตฺติ อสญฺจิจฺจ ฯเปฯ อาทิกมฺมิกสฺสาติ.}{ปฐมสิกฺขาปทํ นิฏฺฐิตํ.}
\proseitem{628}{เตน สมเยน พุทฺโธ ภควา สาวตฺถิยํ วิหรติ เชตวเน อนาถปิณฺฑิกสฺส อาราเม.\csromanspacer{} เตน โข ปน สมเยน ฉพฺพคฺคิยา ภิกฺขู หตฺถนิลฺเลหกํ ภุญฺชนฺติ ฯเปฯ. \csromandash{}}

\prose{\textbf{“น หตฺถนิลฺเลหกํ ภุญฺชิสฺสามีติ สิกฺขา กรณียา”}ติ.}

\prose{น หตฺถนิลฺเลหกํ ภุญฺชิตพฺพํ.\csromanspacer{} โย อนาทริยํ ปฏิจฺจ หตฺถนิลฺเลหกํ ภุญฺชติ, อาปตฺติ ทุกฺกฏสฺส.}

\prosewithcloser{อนาปตฺติ อสญฺจิจฺจ ฯเปฯ อาทิกมฺมิกสฺสาติ.}{ทุติยสิกฺขาปทํ นิฏฺฐิตํ.}
\proseitem{629}{เตน สมเยน พุทฺโธ ภควา สาวตฺถิยํ วิหรติ เชตวเน อนาถปิณฺฑิกสฺส อาราเม.\csromanspacer{} เตน โข ปน สมเยน ฉพฺพคฺคิยา ภิกฺขู ปตฺตนิลฺเลหกํ ภุญฺชนฺติ ฯเปฯ. \csromandash{}}

\prose{\textbf{“น ปตฺตนิลฺเลหกํ ภุญฺชิสฺสามีติ สิกฺขา กรณียา”}ติ.}

\prose{น ปตฺตนิลฺเลหกํ ภุญฺชิตพฺพํ.\csromanspacer{} โย อนาทริยํ ปฏิจฺจ ปตฺตนิลฺเลหกํ ภุญฺชติ, อาปตฺติ ทุกฺกฏสฺส.}

\prosewithcloser{อนาปตฺติ อสญฺจิจฺจ อสฺสติยา อชานนฺตสฺส คิลานสฺส, ปริตฺตเก เสเส เอกโต สงฺกฑฺฒิตฺวา นิลฺเลหิตฺวา ภุญฺชติ, อาปทาสุ อุมฺมตฺตกสฺส อาทิกมฺมิกสฺสาติ.}{ตติยสิกฺขาปทํ นิฏฺฐิตํ.}
\proseitem{630}{\csromanfolio{260}เตน สมเยน พุทฺโธ ภควา สาวตฺถิยํ วิหรติ เชตวเน อนาถปิณฺฑิกสฺสา อาราเม.\csromanspacer{} เตน โข ปน สมเยน ฉพฺพคฺคิยา ภิกฺขู โอฏฺฐนิลฺเลหกํ ภุญฺชนฺติ ฯเปฯ. \csromandash{}}

\prose{\textbf{“น โอฏฺฐนิลฺเลหกํ ภุญฺชิสฺสามีติ สิกฺขา กรณียา”}ติ.}

\prose{น โอฏฺฐนิลฺเลหกํ ภุญฺชิตพฺพํ.\csromanspacer{} โย อนาทริยํ ปฏิจฺจ โอฏฺฐนิลฺเลหกํ ภุญฺชติ, อาปตฺติ ทุกฺกฏสฺส.}

\prosewithcloser{อนาปตฺติ อสญฺจิจฺจ ฯเปฯ อาทิกมฺมิกสฺสาติ.}{จตุตฺถสิกฺขาปทํ นิฏฺฐิตํ.}
\proseitem{613}{เตน สมเยน พุทฺโธ ภควา ภคฺเคสุ วิหรติ สุสุมารคิเร\footnote{สุํสุมารคิเร (สี, สฺยา), สํสุมารคิเร (ก)} เภสกฬาวเน มิคทาเย.\csromanspacer{} เตน โข ปน สมเยน ภิกฺขู โกกนเท\footnote{โกกนุเท (ก)} ปาสาเท สามิเสน หตฺเถน ปานียถาลกํ ปฏิคฺคณฺหนฺติ.\csromanspacer{} มนุสฺสา อุชฺฌายนฺติ ขิยฺยนฺติ วิปาเจนฺติ “กถํ หิ นาม สมณา สกฺยปุตฺติยา สามิเสน หตฺเถน ปานียถาลกํ ปฏิคฺคเหสฺสนฺติ, เสยฺยถาปิ คิหี กามโภคิโน”ติ.\csromanspacer{} อสฺโสสุํ โข ภิกฺขู เตสํ มนุสฺสานํ อุชฺฌายนฺตานํ ขิยฺยนฺตานํ วิปาเจนฺตานํ.\csromanspacer{} เย เต ภิกฺขู อปฺปิจฺฉา ฯเปฯ เต อุชฺฌายนฺติ ขิยฺยนฺติ วิปาเจนฺติ “กถํ หิ นาม ภิกฺขู สามิเสน หตฺเถน ปานียถาลกํ ปฏิคฺคเหสฺสนฺตี”ติ ฯเปฯ.\csromanspacer{} สจฺจํ กิร ภิกฺขเว ภิกฺขู สามิเสน หตฺเถน ปานียถาลกํ ปฏิคฺคณฺหนฺตีติ.\csromanspacer{} สจฺจํ ภควาติ.\csromanspacer{} วิครหิ พุทฺโธ ภควา ฯเปฯ.\csromanspacer{} กถํ หิ นาม เต ภิกฺขเว โมฆปุริสา สามิเสน หตฺเถน ปานียถาลกํ ปฏิคฺคเหสฺสนฺติ.\csromanspacer{} เนตํ ภิกฺขเว อปฺปสนฺนานํ วา ปสาทาย ฯเปฯ.\csromanspacer{} เอวญฺจ ปน ภิกฺขเว อิมํ สิกฺขาปทํ อุทฺทิเสยฺยาถ\csromandash{}}

\prose{\textbf{“น สามิเสน หตฺเถน ปานียถาลกํ ปฏิคฺคเหสฺสามีติ สิกฺขา กรณียา”}ติ.}

\prose{\textbf{น สามิเสน หตฺเถน ปา}นียถาลโก ปฏิคฺคเหตพฺโพ.\csromanspacer{} โย อนาทริยํ ปฏิจฺจ สามิเสน หตฺเถน ปานียถาลกํ ปฏิคฺคณฺหาติ, อาปตฺติ ทุกฺกฏสฺส.}

\prosewithcloser{\csromanfolio{261}อนาปตฺติ อสญฺจิจฺจ อสฺสติยา อชานนฺตสฺส คิลานสฺส, “โธวิสฺสามี”ติ วา “โธวาเปสฺสามี”ติ วา ปฏิคฺคณฺหาติ, อาปทาสุ อุมฺมตฺตกสฺส อาทิกมฺมิกสฺสาติ.}{ปญฺจมสิกฺขาปทํ นิฏฺฐิตํ.}
\proseitem{632}{เตน สมเยน พุทฺโธ ภควา ภคฺเคสุ วิหรติ สุสุมารคิเร เภสกฬาวเน มิคทาเย.\csromanspacer{} เตน โข ปน สมเยน ภิกฺขู โกกนเท ปาสาเท สสิตฺถกํ ปตฺตโธวนํ อนฺตรฆเร ฉฑฺเฑนฺติ.\csromanspacer{} มนุสฺสา อุชฺฌายนฺติ ขิยฺยนฺติ วิปาเจนฺติ “กถํ หิ นาม สมณา สกฺยปุตฺติยา สสิตฺถกํ ปตฺตโธวนํ อนฺตรฆเร ฉฑฺเฑสฺสนฺติ, เสยฺยถาปิ คิหี กามโภคิโน”ติ.\csromanspacer{} อสฺโสสุํ โข ภิกฺขู เตสํ มนุสฺสานํ อุชฺฌายนฺตานํ ขิยฺยนฺตานํ วิปาเจนฺตานํ.\csromanspacer{} เย เต ภิกฺขู อปฺปิจฺฉา ฯเปฯ เต อุชฺฌายนฺติ ขิยฺยนฺติ วิปาเจนฺติ “กถํ หิ นาม ภิกฺขู สสิตฺถกํ ปตฺตโธวนํ อนฺตรฆเร ฉฑฺเฑสฺสนฺตี”ติ ฯเปฯ.\csromanspacer{} สจฺจํ กิร ภิกฺขเว ภิกฺขู สสิตฺถกํ ปตฺตโธวนํ อนฺตรฆเร ฉฑฺเฑนฺตีติ.\csromanspacer{} สจฺจํ ภควาติ.\csromanspacer{} วิครหิ พุทฺโธ ภควา ฯเปฯ.\csromanspacer{} กถํ หิ นาม เต ภิกฺขเว โมฆปุริสา สสิตฺถกํ ปตฺตโธวนํ อนฺตรฆเร ฉฑฺเฑสฺสนฺติ.\csromanspacer{} เนตํ ภิกฺขเว อปฺปสนฺนานํ วา ปสาทาย ฯเปฯ.\csromanspacer{} เอวญฺจ ปน ภิกฺขเว อิมํ สิกฺขาปทํ อุทฺทิเสยฺยาถ\csromandash{}}

\prose{\textbf{“น สสิตฺถกํ ปตฺตโธวนํ อนฺตรฆเร ฉฑฺเฑสฺสามีติ สิกฺขา กรณียา”}ติ.}

\prose{\textbf{น สสิตฺถกํ ปตฺตโธวนํ อนฺตรฆเร จฺ}หฑฺเฑตพฺพํ.\csromanspacer{} โย อนาทริยํ ปฏิจฺจ สสิตฺถกํ ปตฺตโธวนํ อนฺตรฆเร ฉฑฺเฑติ, อาปตฺติ ทุกฺกฏสฺส.}

\prosewithcloser{อนาปตฺติ อสญฺจิจฺจ อสฺสติยา อชานนฺตสฺส คิลานสฺส, อุทฺธริตฺวา วา ภินฺทิตฺวา วา ปฏิคฺคเหตฺวา วา นีหริตฺวา วา ฉฑฺเฑติ, อาปทาสุ อุมฺมตฺตกสฺส อาทิกมฺมิกสฺสาติ.}{ฉฏฺฐสิกฺขาปทํ นิฏฺฐิตํ.}
\proseitem{633}{เตน สมเยน พุทฺโธ ภควา สาวตฺถิยํ วิหรติ เชตวเน อนาถปิณฺฑิกสฺส อาราเม.\csromanspacer{} เตน โข ปน สมเยน ฉพฺพคฺคิยา ภิกฺขู ฉตฺตปาณิสฺส ธมฺมํ เทเสนฺติ.\csromanspacer{} เย เต ภิกฺขู อปฺปิจฺฉา ฯเปฯ เต อุชฺฌายนฺติ ขิยฺยนฺติ วิปาเจนฺติ “กถํ หิ นาม ฉพฺพคฺคิยา ภิกฺขู \csromanfolio{262}ฉตฺตปาณิสฺส ธมฺมํ เทเสสฺสนฺตี”ติ ฯเปฯ.\csromanspacer{} สจฺจํ กิร ตุเมฺห ภิกฺขเว ฉตฺตปาณิสฺส ธมฺมํ เทเสถาติ.\csromanspacer{} สจฺจํ ภควาติ.\csromanspacer{} วิครหิ พุทฺโธ ภควา ฯเปฯ.\csromanspacer{} กถํ หิ นาม ตุเมฺห โมฆปุริสา ฉตฺตปาณิสฺส ธมฺมํ \textbf{เทเสสฺสถ.}\csromanspacer{}\textbf{ เนตํ โมฆปุริสา อปฺปสนฺนานํ วา ปสาทาย ฯเปฯ.}\csromanspacer{}\textbf{ เอวญฺจ} \textbf{ปน ภิกฺขเว อิมํ สิกฺขาปทํ อุทฺทิเสยฺยาถ\csromandash{}}}

\prose{\textbf{“น ฉตฺตปาณิสฺส ธมฺมํ เทเสสฺสามีติ สิกฺขา กรณียา”}ติ.}

\prose{เอวญฺจิทํ ภควตา ภิกฺขูนํ สิกฺขาปทํ ปญฺญตฺตํ โหติ.}

\proseitem{634}{เตน โข ปน สมเยน ภิกฺขู ฉตฺตปาณิสฺส คิลานสฺส ธมฺมํ เทเสตุํ กุกฺกุจฺจายนฺติ.\csromanspacer{} มนุสฺสา อุชฺฌายนฺติ ขิยฺยนฺติ วิปาเจนฺติ “กถํ หิ นาม สมณา สกฺยปุตฺติยา ฉตฺตปาณิสฺส คิลานสฺส ธมฺมํ น เทเสสฺสนฺตี”ติ.\csromanspacer{} อสฺโสสุํ โข ภิกฺขู เตสํ มนุสฺสานํ อุชฺฌายนฺตานํ ขิยฺยนฺตานํ วิปาเจนฺตานํ.\csromanspacer{} อถ โข เต ภิกฺขู ภควโต เอตมตฺถํ อาโรเจสุํ.\csromanspacer{} อถ โข ภควา เอตสฺมิํ นิทาเน เอตสฺมิํ ปกรเณ ธมฺมิํ กถํ กตฺวา ภิกฺขู อามนฺเตสิ อนุชานามิ ภิกฺขเว ฉตฺตปาณิสฺส คิลานสฺส ธมฺมํ เทเสตุํ.\csromanspacer{} เอวญฺจ ปน ภิกฺขเว อิมํ สิกฺขาปทํ อุทฺทิเสยฺยาถ\csromandash{}}

\prose{\textbf{“น ฉตฺตปาณิสฺส อคิลานสฺส ธมฺมํ เทเสสฺสามีติ สิกฺขา กรณียา”}ติ.}

\prose{\textbf{ฉตฺตํ} นาม ตีณิ ฉตฺตานิ เสตจฺฉตฺตํ กิลญฺชจฺฉตฺตํ ปณฺณจฺฉตฺตํ มณฺฑลพทฺธํ สลากพทฺธํ.}

\prose{\textbf{ธมฺโม} นาม พุทฺธภาสิโต สาวกภาสิโต อิสิภาสิโต เทวตาภาสิโต อตฺถูปสญฺหิโต ธมฺมูปสญฺหิโต.}

\prose{\textbf{เทเสยฺยา}ติ ปเทน เทเสติ, ปเท ปเท อาปตฺติ ทุกฺกฏสฺส.\csromanspacer{} อกฺขราย เทเสติ, อกฺขรกฺขราย อาปตฺติ ทุกฺกฏสฺส.\csromanspacer{} น ฉตฺตปาณิสฺส อคิลานสฺส ธมฺโม เทเสตพฺโพ.\csromanspacer{} โย อนาทริยํ ปฏิจฺจ ฉตฺตปาณิสฺส อคิลานสฺส ธมฺมํ เทเสติ, อาปตฺติ ทุกฺกฏสฺส.}

\prosewithcloser{อนาปตฺติ อสญฺจิจฺจ ฯเปฯ อาทิกมฺมิกสฺสาติ.}{สตฺตมสิกฺขาปทํ นิฏฺฐิตํ.}
\proseitem{635}{\csromanfolio{263}เตน สมเยน พุทฺโธ ภควา สาวตฺถิยํ วิหรติ เชตวเน อนาถปิณฺฑิกสฺส อาราเม.\csromanspacer{} เตน โข ปน สมเยน ฉพฺพคฺคิยา ภิกฺขู ทณฺฑปาณิสฺส ธมฺมํ เทเสนฺติ ฯเปฯ. \csromandash{}}

\prose{\textbf{“น ทณฺฑปาณิสฺส อคิลานสฺส ธมฺมํ เทเสสฺสามีติ สิกฺขา กรณียา”}ติ.}

\prose{\textbf{ทณฺโฑ} นาม มชฺฌิมสฺส ปุริสสฺส จตุหตฺโถ ทณฺโฑ.\csromanspacer{} ตโต อุกฺกฏฺโฐ อทณฺโฑ, โอมโก อทณฺโฑ.}

\prose{\textbf{น ทณฺฑปาณิสฺส อคิลานสฺส} ธมฺโม เทเสตพฺโพ.\csromanspacer{} โย อนาทริยํ ปฏิจฺจ ทณฺฑปาณิสฺส อคิลานสฺส ธมฺมํ เทเสติ, อาปตฺติ ทุกฺกฏสฺส.}

\prosewithcloser{อนาปตฺติ อสญฺจิจฺจ ฯเปฯ อาทิกมฺมิกสฺสาติ.}{อฏฺฐมสิกฺขาปทํ นิฏฺฐิตํ.}
\proseitem{636}{เตน สมเยน พุทฺโธ ภควา สาวตฺถิยํ วิหรติ เชตวเน อนาถปิณฺฑิกสฺส อาราเม.\csromanspacer{} เตน โข ปน สมเยน ฉพฺพคฺคิยา ภิกฺขู สตฺถปาณิสฺส ธมฺมํ เทเสนฺติ ฯเปฯ. \csromandash{}}

\prose{\textbf{“น สตฺถปาณิสฺส อคิลานสฺส ธมฺมํ เทเสสฺสามีติ สิกฺขา กรณียา”}ติ.}

\prose{\textbf{สตฺตํ} นาม เอกโตธารํ อุภโตธารํ ปหรณํ.}

\prose{\textbf{น สตฺถปาณิสฺส อคิลานสฺสฺ}อ ธมฺโม เทเสตพฺโพ.\csromanspacer{} โย อนาทริยํ ปฏิจฺจ สตฺถปาณิสฺส อคิลานสฺส ธมฺมํ เทเสติ, อาปตฺติ ทุกฺกฏสฺส.}

\prosewithcloser{อนาปตฺติ อสญฺจิจฺจ ฯเปฯ อาทิกมฺมิกสฺสาติ.}{นวมสิกฺขาปทํ นิฏฺฐิตํ.}
\proseitem{637}{เตน สมเยน พุทฺโธ ภควา สาวตฺถิยํ วิหรติ เชตวเน อนาถปิณฺฑิกสฺส อาราเม.\csromanspacer{} เตน โข ปน สมเยน ฉพฺพคฺคิยา ภิกฺขู อาวุธปาณิสฺส ธมฺมํ เทเสนฺติ ฯเปฯ. \csromandash{}}

\prose{\textbf{“น อาวุธปาณิสฺส อคิลานสฺส ธมฺมํ เทเสสฺสามีติ สิกฺขา กรณียา”}ติ.}

\prose{\textbf{อาวุธํ} นาม จาโป โกทณฺโฑ.}

\prose{\csromanfolio{264}น อาวุธปาณิสฺส อคิลานสฺส ธมฺโม เทเสตพฺโพ.\csromanspacer{} โย อนาทริยํ ปฏิจฺจ อาวุธปาณิสฺส อคิลานสฺส ธมฺมํ เทเสติ, อาปตฺติ ทุกฺกฏสฺส.}

\prosewithcloser{อนาปตฺติ อสญฺจิจฺจ ฯเปฯ อาทิกมฺมิกสฺสาติ.}{ทสมสิกฺขาปทํ นิฏฺฐิตํ.}
\csromanheader{2}{สุรุสุรุวคฺโค ฉฏฺโฐ.}
\csromansectionrule
\csromanmatikaanchor{mtk.115}
\csromantocmarkref{section}{7. ปาทุกวคฺค}{mtk.115}
\bookmark[dest={mtk.115},level=1]{7. ปาทุกวคฺค}
\setcsromanheadh{7. ปาทุกวคฺค}
\csromanheader{1}{7. \textbf{ปาทุกวคฺค}}
\proseitem{638}{เตน สมเยน พุทฺโธ ภควา สาวตฺถิยํ วิหรติ เชตวเน อนาถปิณฺฑิกสฺสา อาราเม.\csromanspacer{} เตน โข ปน สมเยน ฉพฺพคฺคิยา ภิกฺขู ปาทุการุฬฺหสฺส ธมฺมํ เทเสนฺติ ฯเปฯ. \csromandash{}}

\prose{\textbf{“น ปาทุการุฬฺหสฺส อคิลานสฺส ธมฺมํ เทเสสฺสามีติ สิกฺขา กรณียา”}ติ.}

\prose{\textbf{น ปาทุการุฬฺหสฺส อคิลฺ}อานสฺส ธมฺโม เทเสตพฺโพ.\csromanspacer{} โย อนาทริยํ ปฏิจฺจ อกฺกนฺตสฺส วา ปฏิมุกฺกสฺส วา โอมุกฺกสฺส วา อคิลานสฺส ธมฺมํ เทเสติ, อาปตฺติ ทุกฺกฏสฺส.}

\prosewithcloser{อนาปตฺติ อสญฺจิจฺจ ฯเปฯ อาทิกมฺมิกสฺสาติ.}{ปฐมสิกฺขาปทํ นิฏฺฐิตํ.}
\proseitem{639}{เตน สมเยน พุทฺโธ ภควา สาวตฺถิยํ วิหรติ เชตวเน อนาถปิณฺฑิกสฺส อาราเม.\csromanspacer{} เตน โข ปน สมเยน ฉพฺพคฺคิยา ภิกฺขู อุปาหนารุฬฺหสฺส ธมฺมํ เทเสนฺติ ฯเปฯ. \csromandash{}}

\prose{\textbf{“น อุปาหนารุฬฺหสฺส อคิลานสฺส ธมฺมํ เทเสสฺสามีติ สิกฺขา กรณียา”}ติ.}

\prose{\textbf{น อุปาหนารุฬฺหสฺส อคิ}ลานสฺส ธมฺโม เทเสตพฺโพ.\csromanspacer{} โย อนาทริยํ ปฏิจฺจ อกฺกนฺตสฺส วา ปฏิมุกฺกสฺส วา โอมุกฺกสฺส วา อคิลานสฺส ธมฺมํ เทเสติ, อาปตฺติ ทุกฺกฏสฺส.}

\prosewithcloser{อนาปตฺติ อสญฺจิจฺจ ฯเปฯ อาทิกมฺมิกสฺสาติ.}{ทุติยสิกฺขาปทํ นิฏฺฐิตํ.}
\proseitem{640}{\csromanfolio{265}เตน สมเยน พุทฺโธ ภควา สาวตฺถิยํ วิหรติ เชตวเน อนาถปิณฺฑิกสฺส อาราเม.\csromanspacer{} เตน โข ปน สมเยน ฉพฺพคฺคิยา ภิกฺขู ยานคตสฺส ธมฺมํ เทเสนฺติ ฯเปฯ. \csromandash{}}

\prose{\textbf{“น ยานคตสฺส อคิลานสฺส ธมฺมํ เทเสสฺสามีติ สิกฺขา กรณียา”}ติ.}

\prose{\textbf{ยานํ} นาม วยฺหํ รโถ สกฏํ สนฺทมานิกา สิวิกา ปาฏงฺกี.}

\prose{\textbf{น ยานคตสฺส อคิลานสฺส} ธมฺโม เทเสตพฺโพ.\csromanspacer{} โย อนาทริยํ ปฏิจฺจ ยานคตสฺส อคิลานสฺส ธมฺมํ เทเสติ, อาปตฺติ ทุกฺกฏสฺส.}

\prosewithcloser{อนาปตฺติ อสญฺจิจฺจ ฯเปฯ อาทิกมฺมิกสฺสาติ.}{ตติยสิกฺขาปทํ นิฏฺฐิตํ.}
\proseitem{641}{เตน สมเยน พุทฺโธ ภควา สาวตฺถิยํ วิหรติ เชตวเน อนาถปิณฺฑิกสฺส อาราเม, เตน โข ปน สมเยน ฉพฺพคฺคิยา ภิกฺขู สยนคตสฺส ธมฺมํ เทเสนฺติ ฯเปฯ. \csromandash{}}

\prose{\textbf{“น สยนคตสฺส อคิลานสฺส ธมฺมํ เทเสสฺสามีติ สิกฺขา กรณียา”}ติ.}

\prose{\textbf{น สยนคตสฺส อคิลานสฺสฺ}อ ธมฺโม เทเสตพฺโพ.\csromanspacer{} โย อนาทริยํ ปฏิจฺจ อนฺตมโส ฉมายมฺปิ นิปนฺนสฺส สยนคตสฺส อคิลานสฺส ธมฺมํ เทเสติ, อาปตฺติ ทุกฺกฏสฺส.}

\prosewithcloser{อนาปตฺติ อสญฺจิจฺจ ฯเปฯ อาทิกมฺมิกสฺสาติ.}{จตุตฺถสิกฺขาปทํ นิฏฺฐิตํ.}
\proseitem{642}{เตน สมเยน พุทฺโธ ภควา สาวตฺถิยํ วิหรติ เชตวเน อนาถปิณฺฑิกสฺสา อาราเม.\csromanspacer{} เตน โข ปน สมเยน ฉพฺพคฺคิยา ภิกฺขู ปลฺลตฺถิกาย นิสินฺนสฺส ธมฺมํ เทเสนฺติ ฯเปฯ. \csromandash{}}

\prose{\textbf{“น ปลฺลตฺถิกาย นิสินฺนสฺส อคิลานสฺส ธมฺมํ เทเสสฺสามีติ สิกฺขา กรณียา”}ติ.}

\prose{\csromanfolio{266}น ปลฺลตฺถิกาย นิสินฺนสฺส อคิลานสฺส ธมฺโม เทเสตพฺโพ.\csromanspacer{} โย อนาทริยํ ปฏิจฺจ หตฺถปลฺลตฺถิกาย วา ทุสฺสปลฺลตฺถิกาย วา นิสินฺนสฺส อคิลานสฺส ธมฺมํ เทเสติ, อาปตฺติ ทุกฺกฏสฺส.}

\prosewithcloser{อนาปตฺติ อสญฺจิจฺจ ฯเปฯ อาทิกมฺมิกสฺสาติ.}{ปญฺจมสิกฺขาปทํ นิฏฺฐิตํ.}
\proseitem{643}{เตน สมเยน พุทฺโธ ภควา สาวตฺถิยํ วิหรติ เชตวเน อนาถปิณฺฑิกสฺส อาราเม.\csromanspacer{} เตน โข ปน สมเยน ฉพฺพคฺคิยา ภิกฺขู เวฐิตสีสสฺส ธมฺมํ เทเสนฺติ ฯเปฯ. \csromandash{}}

\prose{\textbf{“น เวฐิตสีสสฺส อคิลานสฺส ธมฺมํ เทเสสฺสามีติ สิกฺขา กรณียา”}ติ.}

\prose{\textbf{เวฐิตสีโส} นาม เกสนฺตํ น ทสฺสาเปตฺวา เวฐิโต โหติ.}

\prose{\textbf{น เวฐิตสีสสฺส อคิลานสฺสฺ}อ ธมฺโม เทเสตพฺโพ.\csromanspacer{} โย อนาทริยํ ปฏิจฺจ เวฐิตสีสสฺส อคิลานสฺส ธมฺมํ เทเสติ, อาปตฺติ ทุกฺกฏสฺส.}

\prosewithcloser{อนาปตฺติ อสญฺจิจฺจ อสฺสติยา อชานนฺตสฺส คิลานสฺส, เกสนฺตํ วิวราเปตฺวา เทเสติ, อาปทาสุ อุมฺมตฺตกสฺส อาทิกมฺมิกสฺสาติ.}{ฉฏฺฐสิกฺขาปทํ นิฏฺฐิตํ.}
\proseitem{644}{เตน สมเยน พุทฺโธ ภควา สาวตฺถิยํ วิหรติ เชตวเน อนาถปิณฺฑิกสฺส อาราเม.\csromanspacer{} เตน โข ปน สมเยน ฉพฺพคฺคิยา ภิกฺขู โอคุณฺฐิตสีสสฺส ธมฺมํ เทเสนฺติ ฯเปฯ. \csromandash{}}

\prose{\textbf{“น โอคุณฺฐิตสีสสฺส อคิลานสฺส ธมฺมํ เทเสสฺสามีติ สิกฺขา กรณียา”}ติ.}

\prose{\textbf{โอคุณฺฐิตสีโส} นาม สสีสํ ปารุโต วุจฺจติ.}

\prose{\textbf{น โอคุณฺฐิตสีสสฺส อคฺ}อิลานสฺส ธมฺโม เทเสตพฺโพ.\csromanspacer{} โย อนาทริยํ ปฏิจฺจ โอคุณฺฐิตสีสสฺส อคิลานสฺส ธมฺมํ เทเสติ, อาปตฺติ ทุกฺกฏสฺส.}

\prosewithcloser{อนาปตฺติ อสญฺจิจฺจ อสฺสติยา อชานนฺตสฺส คิลานสฺส, สีสํ วิวราเปตฺวา เทเสติ, อาปทาสุ อุมฺมตฺตกสฺส อาทิกมฺมิกสฺสาติ.}{สตฺตมสิกฺขาปทํ นิฏฺฐิตํ.}
\proseitem{645}{\csromanfolio{267}เตน สมเยน พุทฺโธ ภควา สาวตฺถิยํ วิหรติ เชตวเน อนาถปิณฺฑิกสฺส อาราเม.\csromanspacer{} เตน โข ปน สมเยน ฉพฺพคฺคิยา ภิกฺขู ฉมายํ นิสีทิตฺวา อาสเน นิสินฺนสฺส ธมฺมํ เทเสนฺติ ฯเปฯ. \csromandash{}}

\prose{\textbf{“น ฉมายํ นิสีทิตฺวา อาสเน นิสินฺนสฺส อคิลานสฺส ธมฺมํ เทเสสฺสามีติ สิกฺขา กรณียา”}ติ.}

\prose{\textbf{น ฉมายํ นิสีทิตฺวา}\footnote{นิสินฺเนน (อฏฺฐกถา)}\textbf{ อาสเน นิสินฺนสฺส} อคิลานสฺส ธมฺโม เทเสตพฺโพ.\csromanspacer{} โย อนาทริยํ ปฏิจฺจ ฉมายํ นิสีทิตฺวา อาสเน นิสินฺนสฺส อคิลานสฺส ธมฺมํ เทเสติ, อาปตฺติ ทุกฺกฏสฺส.}

\prosewithcloser{อนาปตฺติ อสญฺจิจฺจ ฯเปฯ อาทิกมฺมิกสฺสาติ.}{อฏฺฐมสิกฺขาปทํ นิฏฺฐิตํ.}
\proseitem{646}{เตน สมเยน พุทฺโธ ภควา สาวตฺถิยํ วิหรติ เชตวเน อนาถปิณฺฑิกสฺส อาราเม.\csromanspacer{} เตน โข ปน สมเยน ฉพฺพคฺคิยา ภิกฺขู นีเจ อาสเน นิสีทิตฺวา อุจฺเจ อาสเน นิสินฺนสฺส ธมฺมํ เทเสนฺติ.\csromanspacer{} เย เต ภิกฺขู อปฺปิจฺฉา ฯเปฯ เต อุชฺฌายนฺติ ขิยฺยนฺติ วิปาเจนฺติ “กถํ หิ นาม ฉพฺพคฺคิยา ภิกฺขู นีเจ อาสเน นิสีทิตฺวา อุจฺเจ อาสเน นิสินฺนสฺส ธมฺมํ เทเสสฺสนฺตี”ติ ฯเปฯ.\csromanspacer{} สจฺจํ กิร ตุเมฺห ภิกฺขเว นีเจ อาสเน นิสีทิตฺวา อุจฺเจ อาสเน นิสินฺนสฺส ธมฺมํ เทเสถาติ.\csromanspacer{} สจฺจํ ภควาติ.\csromanspacer{} วิครหิ พุทฺโธ ภควา ฯเปฯ.\csromanspacer{} กถํ หิ นาม ตุเมฺห โมฆปุริสา นีเจ อาสเน นิสีทิตฺวา อุจฺเจ อาสเน นิสินฺนสฺส ธมฺมํ เทเสสฺสถ.\csromanspacer{} เนตํ โมฆปุริสา อปฺปสนฺนานํ วา ปสาทาย ฯเปฯ วิครหิตฺวา ฯเปฯ ธมฺมิํ กถํ กตฺวา ภิกฺขู อามนฺเตสิ\csromandash{}}

\proseitem{647}{“ภูตปุพฺพํ ภิกฺขเว พาราณสิยํ อญฺญตรสฺส ฉปกสฺส\footnote{ฉวกสฺส (สฺยา)} ปชาปติ คพฺภินี อโหสิ.\csromanspacer{} อถ โข ภิกฺขเว สา ฉปกี ตํ ฉปกํ เอตทโวจ ‘คพฺภินีมฺหิ อยฺยปุตฺต, อิจฺฉามิ อมฺพํ ขาทิตุนฺ’ติ.\csromanspacer{} นตฺถิ อมฺพํ\footnote{อมฺโพ (สฺยา)}, อกาโล อมฺพสฺสาติ. ‘สเจ น ลภิสฺสามิ มริสฺสามี’ติ.\csromanspacer{} เตน โข ปน สมเยน รญฺโญ อมฺโพ ธุวผโล โหติ.\csromanspacer{} อถ โข ภิกฺขเว โส ฉปโก เยน โส อมฺโพ เตนุปสงฺกมิ, อุปสงฺกมิตฺวา ตํ อมฺพํ อภิรุหิตฺวา นิลีโน อจฺฉิ.\csromanspacer{} อถ โข ภิกฺขเว ราชา ปุโรหิเตน พฺราหฺมเณน สทฺธิํ เยน โส อมฺโพ เตนุปสงฺกมิ, อุปสงฺกมิตฺวา อุจฺเจ อาสเน นิสีทิตฺวา มนฺตํ ปริยาปุณาติ.\csromanspacer{} อถ โข \csromanfolio{268}\textbf{ภิกฺขเว ตสฺส ฉปกสฺส เอตทโหสิ ‘ยาว อธมฺมิโก อยํ ราชา, ยตฺร หิ} \textbf{นาม อุจฺเจ อาสเน นิสีทิตฺวา มนฺตํ ปริยาปุณิสฺสติ.}\csromanspacer{}\textbf{ อยญฺจ พฺราหฺมโณ} \textbf{อธมฺมิโก, ยตฺร หิ นาม นีเจ อาสเน นิสีทิตฺวา อุจฺเจ อาสเน นิสินฺนสฺส} มนฺตํ วาเจสฺสติ.\csromanspacer{} อหญฺจมฺหิ อธมฺมิโก, โยหํ อิตฺถิยา การณา รญฺโญ อมฺพํ อวหรามิ.\csromanspacer{} สพฺพมิทํ จริมํ กตนฺ’ติ ตตฺเถว ปริปติ.}

\setlength{\csromangathapagewidth}{0pt}
\csromangathameasuregroup{*อุโภ อตฺถํ น ชานนฺติ, \\ โย จายํ มนฺตํ วาเจติ, \\ สาลีนํ โอทโน ภุตฺโต, \\ ตสฺมา ธมฺเม น วตฺตามิ, \\ ธิรตฺถุ ตํ ธนลาภํ, \\ ยา วุตฺติ วินิปาเตน, \\ ปริพฺพช มหาพฺรหฺเม, \\ มา ตฺวํ อธมฺโม อาจริโต,}{\csromangathabat{*อุโภ อตฺถํ น ชานนฺติ,}{อุโภ ธมฺมํ น ปสฺสเร.} \\ \csromangathabat{โย จายํ มนฺตํ วาเจติ,}{โย จาธมฺเมน’ธียติ.} \\ \csromangathabat{สาลีนํ โอทโน ภุตฺโต,}{สุจิมํสูปเสจโน.} \\ \csromangathabat{ตสฺมา ธมฺเม น วตฺตามิ,}{ธมฺโม อริเยภิ วณฺณิโต.} \\ \csromangathabat{ธิรตฺถุ ตํ ธนลาภํ,}{ยสลาภญฺจ พฺราหฺมณ.} \\ \csromangathabat{ยา วุตฺติ วินิปาเตน,}{อธมฺมจรเณน วา.} \\ \csromangathabat{ปริพฺพช มหาพฺรหฺเม,}{ปจนฺตญฺเญปิ ปาณิโน.} \\ \csromangathabat{มา ตฺวํ อธมฺโม อาจริโต,}{อสฺมา กุมฺภมิวาภิทา”ติ.}}
\csromangathasetleft{*อุโภ อตฺถํ น ชานนฺติ, \\ โย จายํ มนฺตํ วาเจติ, \\ สาลีนํ โอทโน ภุตฺโต, \\ ตสฺมา ธมฺเม น วตฺตามิ, \\ ธิรตฺถุ ตํ ธนลาภํ, \\ ยา วุตฺติ วินิปาเตน, \\ ปริพฺพช มหาพฺรหฺเม, \\ มา ตฺวํ อธมฺโม อาจริโต,}
\csromangathagroup{\csromangathastanza{\csromangathabat{\csromansymbolfootnote{*}{อิมา คาถาโย ขุ 5. 97 ปิฏฺเฐ ชาตเก อญฺญถา ทิสฺสนฺติ.}อุโภ อตฺถํ น ชานนฺติ,}{อุโภ ธมฺมํ น ปสฺสเร.} \\ \csromangathabat{โย จายํ มนฺตํ วาเจติ,}{โย จาธมฺเมน’ธียติ.}}\csromangathastanza{\csromangathabat{สาลีนํ โอทโน ภุตฺโต,}{สุจิมํสูปเสจโน.} \\ \csromangathabat{ตสฺมา ธมฺเม น วตฺตามิ,}{ธมฺโม อริเยภิ วณฺณิโต.}}\csromangathastanza{\csromangathabat{ธิรตฺถุ ตํ ธนลาภํ,}{ยสลาภญฺจ พฺราหฺมณ.} \\ \csromangathabat{ยา วุตฺติ วินิปาเตน,}{อธมฺมจรเณน วา.}}\csromangathastanza{\csromangathabat{ปริพฺพช มหาพฺรหฺเม,}{ปจนฺตญฺเญปิ ปาณิโน.} \\ \csromangathabat{มา ตฺวํ อธมฺโม อาจริโต,}{อสฺมา กุมฺภมิวาภิทา”ติ.}}}

\prose{ตทาปิ เม ภิกฺขเว อมนาปา นีเจ อาสเน นิสีทิตฺวา อุจฺเจ อาสเน นิสินฺนสฺส มนฺตํ วาเจตุํ.\csromanspacer{} กิมงฺคํ ปน เอตรหิ น อมนาปา ภวิสฺสติ นีเจ อาสเน นิสีทิตฺวา อุจฺเจ อาสเน นิสินฺนสฺส ธมฺมํ เทเสตุํ.\csromanspacer{} เนตํ ภิกฺขเว อปฺปสนฺนานํ วา ปสาทาย ฯเปฯ.\csromanspacer{} เอวญฺจ ปน ภิกฺขเว อิมํ สิกฺขาปทํ อุทฺทิเสยฺยาถ\csromandash{}}

\prose{\textbf{“น นีเจ อาสเน นิสีทิตฺวา อุจฺเจ อาสเน นิสินฺนสฺส อคิลานสฺส ธมฺมํ เทเสสฺสามีติ สิกฺขา กรณียา”}ติ.}

\prose{\textbf{น นีเจ อาสเน นิสีทิตฺวา อุจฺเจ อาสเน นิสินฺนสฺสฺ}อ อคิลานสฺส ธมฺโม เทเสตพฺโพ.\csromanspacer{} โย อนาทริยํ ปฏิจฺจ นีเจ อาสเน นิสีทิตฺวา อุจฺเจ อาสเน นิสินฺนสฺส อคิลานสฺส ธมฺมํ เทเสติ, อาปตฺติ ทุกฺกฏสฺส.}

\prosewithcloser{อนาปตฺติ อสญฺจิจฺจ ฯเปฯ อาทิกมฺมิกสฺสาติ.}{นวมสิกฺขาปทํ นิฏฺฐิตํ.}
\proseitem{648}{เตน สมเยน พุทฺโธ ภควา สาวตฺถิยํ วิหรติ เชตวเน อนาถปิณฺฑิกสฺส อาราเม.\csromanspacer{} เตน โข ปน สมเยน ฉพฺพคฺคิยา ภิกฺขู ฐิตา นิสินฺนสฺส ธมฺมํ เทเสนฺติ ฯเปฯ. \csromandash{}}

\prose{\csromanfolio{269}\textbf{“น ฐิโต นิสินฺนสฺส อคิลานสฺส ธมฺมํ เทเสสฺสามีติ สิกฺขา กรณียา”}ติ.}

\prose{น ฐิเตน นิสินฺนสฺส อคิลานสฺส ธมฺโม เทเสตพฺโพ.\csromanspacer{} โย อนาทริยํ ปฏิจฺจ ฐิโต นิสินฺนสฺส อคิลานสฺส ธมฺมํ เทเสติ, อาปตฺติ ทุกฺกฏสฺส.}

\prosewithcloser{อนาปตฺติ อสญฺจิจฺจ ฯเปฯ อาทิกมฺมิกสฺสาติ.}{ทสมสิกฺขาปทํ นิฏฺฐิตํ.}
\proseitem{649}{เตน สมเยน พุทฺโธ ภควา สาวตฺถิยํ วิหรติ เชตวเน อนาถปิณฺฑิกสฺส อาราเม.\csromanspacer{} เตน โข ปน สมเยน ฉพฺพคฺคิยา ภิกฺขู ปจฺฉโต คจฺฉนฺตา ปุรโต คจฺฉนฺตสฺส ธมฺมํ เทเสนฺติ ฯเปฯ. \csromandash{}}

\prose{\textbf{“น ปจฺฉโต คจฺฉนฺโต ปุรโต คจฺฉนฺตสฺส อคิลานสฺส ธมฺมํ เทเสสฺสามีติ สิกฺขา กรณียา”}ติ.}

\prose{น ปจฺฉโต คจฺฉนฺเตน ปุรโต คจฺฉนฺตสฺส อคิลานสฺส ธมฺโม เทเสตพฺโพ.\csromanspacer{} โย อนาทริยํ ปฏิจฺจ ปจฺฉโต คจฺฉนฺโต ปุรโต คจฺฉนฺตสฺส อคิลานสฺส ธมฺมํ เทเสติ, อาปตฺติ ทุกฺกฏสฺส.}

\prosewithcloser{อนาปตฺติ อสญฺจิจฺจ ฯเปฯ อาทิกมฺมิกสฺสาติ.}{เอกาทสมสิกฺขาปทํ นิฏฺฐิตํ.}
\proseitem{650}{เตน สมเยน พุทฺโธ ภควา สาวตฺถิยํ วิหรติ เชตวเน อนาถปิณฺฑิกสฺส อาราเม.\csromanspacer{} เตน โข ปน สมเยน ฉพฺพคฺคิยา ภิกฺขู อุปฺปเถน คจฺฉนฺตา ปเถน คจฺฉนฺตสฺส ธมฺมํ เทเสนฺติ ฯเปฯ. \csromandash{}}

\prose{\textbf{“น อุปฺปเถน คจฺฉนฺโต ปเถน คจฺฉนฺตสฺส อคิลานสฺส ธมฺมํ เทเสสฺสามีติ สิกฺขา กรณียา”}ติ.}

\prose{น อุปฺปเถน คจฺฉนฺเตน ปเถน คจฺฉนฺตสฺส อคิลานสฺส ธมฺโม เทเสตพฺโพ.\csromanspacer{} โย อนาทริยํ ปฏิจฺจ อุปฺปเถน คจฺฉนฺโต ปเถน คจฺฉนฺตสฺส อคิลานสฺส ธมฺมํ เทเสติ, อาปตฺติ ทุกฺกฏสฺส.}

\prosewithcloser{อนาปตฺติ อสญฺจิจฺจ ฯเปฯ อาทิกมฺมิกสฺสาติ.}{ทฺวาทสมสิกฺขาปทํ นิฏฺฐิตํ.}
\proseitem{651}{\csromanfolio{270}เตน สมเยน พุทฺโธ ภควา สาวตฺถิยํ วิหรติ เชตวเน อนาถปิณฺฑิกสฺส อาราเม.\csromanspacer{} เตน โข ปน สมเยน ฉพฺพคฺคิยา ภิกฺขู ฐิตา อุจฺจารมฺปิ ปสฺสาวมฺปิ กโรนฺติ ฯเปฯ. \csromandash{}}

\prose{\textbf{“น ฐิโต อคิลาโน อุจฺจารํ วา ปสฺสาวํ วา กริสฺสามีติ สิกฺขา กรณียา”}ติ.}

\prose{น ฐิเตน อคิลาเนน อุจฺจาโร วา ปสฺสาโว วา กาตพฺโพ.\csromanspacer{} โย อนาทริยํ ปฏิจฺจ ฐิโต อคิลาโน อุจฺจรํ วา ปสฺสาวํ วา กโรติ, อาปตฺติ ทุกฺกฏสฺส.}

\prosewithcloser{อนาปตฺติ อสญฺจิจฺจ ฯเปฯ อาทิกมฺมิกสฺสาติ.}{เตรสมสิกฺขาปทํ นิฏฺฐิตํ.}
\proseitem{652}{เตน สมเยน พุทฺโธ ภควา สาวตฺถิยํ วิหรติ เชตวเน อนาถปิณฺฑิกสฺส อาราเม.\csromanspacer{} เตน โข ปน สมเยน ฉพฺพคฺคิยา ภิกฺขู หริเต อุจฺจารมฺปิ ปสฺสาวมฺปิ เขฬมฺปิ กโรนฺติ ฯเปฯ. \csromandash{}}

\prose{\textbf{“น หริเต อคิลาโน อุจฺจารํ วา ปสฺสาวํ วา เขฬํ วา กริสฺสามีติ สิกฺขา กรณียา”}ติ.}

\prose{น หริเต อคิลาเนน อุจฺจาโร วา ปสฺสาโว วา เขโฬ วา กาตพฺโพ.\csromanspacer{} โย อนาทริยํ ปฏิจฺจ หริเต อคิลาโน อุจฺจารํ วา ปสฺสาวํ วา เขฬํ วา กโรติ, อาปตฺติ ทุกฺกฏสฺส.}

\prosewithcloser{อนาปตฺติ อสญฺจิจฺจ อสฺสติยา อชานนฺตสฺส คิลานสฺส, อปฺปหริเต กโต หริตํ โอตฺถรติ, อาปทาสุ อุมฺมตฺตกสฺส อาทิกมฺมิกสฺสาติ.}{จุทฺทสมสิกฺขาปทํ นิฏฺฐิตํ.}
\proseitem{653}{เตน สมเยน พุทฺโธ ภควา สาวตฺถิยํ วิหรติ เชตวเน อนาถปิณฺฑิกสฺส อาราเม.\csromanspacer{} เตน โข ปน สมเยน ฉพฺพคฺคิยา ภิกฺขู อุทเก อุจฺจารมฺปิ ปสฺสาวมฺปิ เขฬมฺปิ กโรนฺติ.\csromanspacer{} มนุสฺสา อุชฺฌายนฺติ ขิยฺยนฺติ วิปาเจนฺติ “กถํ หิ นาม สมณา สกฺยปุตฺติยา อุทเก อุจฺจารมฺปิ ปสฺสาวมฺปิ เขฬมฺปิ กริสฺสนฺติ เสยฺยถาปิ คิหี กามโภคิโน”ติ. \csromanfolio{271}อสฺโสสุํ โข ภิกฺขู เตสํ มนุสฺสานํ อุชฺฌายนฺตานํ ขิยฺยนฺตานํ วิปาเจนฺตานํ.\csromanspacer{} เย เต ภิกฺขู อปฺปิจฺฉา ฯเปฯ เต อุชฺฌายนฺติ ขิยฺยนฺติ วิปาเจนฺติ “กถํ หิ นาม ฉพฺพคฺคิยา ภิกฺขู อุทเก อุจฺจารมฺปิ ปสฺสาวมฺปิ เขฬมฺปิ กริสฺสนฺตี”ติ ฯเปฯ.\csromanspacer{} สจฺจํ กิร ตุเมฺห ภิกฺขเว อุทเก อุจฺจารมฺปิ ปสฺสาวมฺปิ เขฬมฺปิ กโรถาติ.\csromanspacer{} สจฺจํ ภควาติ.\csromanspacer{} วิครหิ พุทฺโธ ภควา ฯเปฯ.\csromanspacer{} กถํ หิ นาม ตุเมฺห โมฆปุริสา อุทเก อุจฺจารมฺปิ ปสฺสาวมฺปิ เขฬมฺปิ กริสฺสถ.\csromanspacer{} เนตํ โมฆปุริสา อปฺปสนฺนานํ วา ปสาทาย ฯเปฯ.\csromanspacer{} เอวญฺจ ปน ภิกฺขเว อิมํ สิกฺขาปทํ อุทฺทิเสยฺยาถ\csromandash{}}

\prose{\textbf{“น อุทเก อุจฺจารํ วา ปสฺสาวํ วา เขฬํ วา กริสฺสามีติ สิกฺขา กรณียา”}ติ.}

\prose{เอวญฺจิทํ ภควตา ภิกฺขูนํ สิกฺขาปทํ ปญฺญตฺตํ โหติ.}

\proseitem{654}{เตน โข ปน สมเยน คิลานา ภิกฺขู อุทเก อุจฺจารมฺปิ ปสฺสาวมฺปิ เขฬมฺปิ กาตุํ กุกฺกุจฺจายนฺติ.\csromanspacer{} ภควโต เอตมตฺถํ อาโรเจสุํ.\csromanspacer{} อถ โข ภควา เอตสฺมิํ นิทาเน เอตสฺมิํ ปกรเณ ธมฺมิํ กถํ กตฺวา ภิกฺขู อามนฺเตสิ อนุชานามิ ภิกฺขเว คิลาเนน ภิกฺขุนา อุทเก อุจฺจารมฺปิ ปสฺสาวมฺปิ เขฬมฺปิ กาตุํ.\csromanspacer{} เอวญฺจ ปน ภิกฺขเว อิมํ สิกฺขาปทํ อุทฺทิเสยฺยาถ\csromandash{}}

\prose{\textbf{“น อุทเก อคิลาโน อุจฺจารํ วา ปสฺสาวํ วา เขฬํ วา กริสฺสามีติ สิกฺขา กรณียา”}ติ.}

\prose{น อุทเก อคิลาเนน อุจฺจาโร วา ปสฺสาโว วา เขโฬ วา กาตพฺโพ.\csromanspacer{} โย อนาทริยํ ปฏิจฺจ อุทเก อคิลาโน อุจฺจารํ วา ปสฺสาวํ วา เขฬํ วา กโรติ, อาปตฺติ ทุกฺกฏสฺส.}

\prosewithcloser{อนาปตฺติ อสญฺจิจฺจ อสฺสติยา อชานนฺตสฺส คิลานสฺส, ถเล กโต อุทกํ โอตฺถรติ, อาปทาสุ อุมฺมตฺตกสฺส ขิตฺตจิตฺตสฺส เวทนาฏฺฏสฺส อาทิกมฺมิกสฺสาติ.}{ปนฺนรสมสิกฺขาปทํ นิฏฺฐิตํ.}
\vspace{\tipitakaparskip}
\csromancenter{ปาทุกวคฺโค สตฺตโม.}
\prosewithcloser{\csromanfolio{272}อุทฺทิฏฺฐา โข อายสฺมนฺโต เสขิยา ธมฺมา, ตตฺถายสฺมนฺเต ปุจฺฉามิ กจฺจิตฺถ ปริสุทฺธา, ทุติยมฺปิ ปุจฺฉามิ กจฺจิตฺถ ปริสุทฺธา, ตติยมฺปิ ปุจฺฉามิ กจฺจิตฺถ ปริสุทฺธา, ปริสุทฺเธตฺถายสฺมนฺโต, ตสฺมา ตุณฺหี, เอวเมตํ ธารยามีติ.}{เสขิยา นิฏฺฐิตา.}
\nitthitamruled{เสขิยกณฺฑํ นิฏฺฐิตํ.}
\csromanmatikaanchor{mtk.116}
\csromantocmarkref{chapter}{8. อธิกรณสมถ}{mtk.116}
\bookmark[dest={mtk.116},level=0]{8. อธิกรณสมถ}
\setcsromanheadcha{8. อธิกรณสมถ}
\setcsromanheadhclear{}
\chapterhead{8. \textbf{อธิกรณสมถ}}
\prose{อิเม โข ปนายสฺมนฺโต สตฺต อธิกรณสมถา ธมฺมา อุทฺเทสํ อาคจฺฉนฺติ.}

\proseitem{655}{อุปฺปนฺนุปฺปนฺนานํ อธิกรณานํ สมถาย วูปสมาย สมฺมุขาวินโย ทาตพฺโพ, สติวินโย ทาตพฺโพ, อมูฬฺหวินโย ทาตพฺโพ, ปฏิญฺญาย กาเรตพฺพํ, เยภุยฺยสิกา, ตสฺสปาปิยสิกา, ติณวตฺถารโกติ.}

\prosewithcloser{อุทฺทิฏฺฐา โข อายสฺมนฺโต สตฺต อธิกรณสมถา ธมฺมา, ตตฺถายสฺมนฺเต ปุจฺฉามิ กจฺจิตฺถ ปริสุทฺธา, ทุติยมฺปิ ปุจฺฉามิ กจฺจิตฺถ ปริสุทฺธา, ตติยมฺปิ ปุจฺฉามิ กจฺจิตฺถ ปริสุทฺธา, ปริสุทฺเธตฺถายสฺมนฺโต, ตสฺมา ตุณฺหี, เอวเมตํ ธารยามีติ.}{อธิกรณสมถา นิฏฺฐิตา.}
\prosewithcloser{อุทฺทิฏฺฐํ โข อายสฺมนฺโต นิทานํ, อุทฺทิฏฺฐา จตฺตาโร ปาราชิกา ธมฺมา, อุทฺทิฏฺฐา เตรส สํฆาทิเสสา ธมฺมา, อุทฺทิฏฺฐา เทฺว อนิยตา ธมฺมา, อุทฺทิฏฺฐา ติํส นิสฺสคฺคิยา ปาจิตฺติยา ธมฺมา, อุทฺทิฏฺฐา เทฺวนวุติ ปาจิตฺติยา ธมฺมา, อุทฺทิฏฺฐา จตฺตาโร ปาฏิเทสนียา ธมฺมา, อุทฺทิฏฺฐา เสขิยา ธมฺมา, อุทฺทิฏฺฐา สตฺต อธิกรณสมถา ธมฺมา, เอตฺตกํ ตสฺส ภควโต สุตฺตาคตํ สุตฺตปริยาปนฺนํ อนฺวทฺธมาสํ อุทฺเทสํ อาคจฺฉติ, ตตฺถ สพฺเพเหว สมคฺเคหิ สมฺโมทมาเนหิ อวิวทมาเนหิ สิกฺขิตพฺพนฺติ.}{\textbf{มหาวิภงฺโค นิฏฺฐิโต.}}
\csromanmatikaanchor{mtk.117}
\csromantocmarknopage{chapter}{ภิกฺขุนีวิภงฺค}{mtk.117}
\bookmark[dest={mtk.117},level=0]{ภิกฺขุนีวิภงฺค}
\setcsromanheadcha{ภิกฺขุนีวิภงฺค}
\setcsromanheadhclear{}
\chapterheadpageread{\textbf{ภิกฺขุนีวิภงฺค}}
\csromansectionrule
\namakkaram{นโม ตสฺส ภควโต อรหโต สมฺมาสมฺพุทฺธสฺส.}
\csromanmatikaanchor{mtk.118}
\csromantocmarknopage{chapter}{1. ปาราชิกกณฺฑ}{mtk.118}
\bookmark[dest={mtk.118},level=0]{1. ปาราชิกกณฺฑ}
\setcsromanheadcha{1. ปาราชิกกณฺฑ}
\setcsromanheadhclear{}
\chapterhead{1. \textbf{ปาราชิกกณฺฑ}}
\csromanmatikaanchor{mtk.119}
\csromantocmarkref{section}{1. ปฐมปาราชิกสิกฺขาปท}{mtk.119}
\bookmark[dest={mtk.119},level=1]{1. ปฐมปาราชิกสิกฺขาปท}
\setcsromanheadh{1. ปฐมปาราชิกสิกฺขาปท}
\csromanheader{1}{1. \textbf{ปฐมปาราชิก}}
\proseitem{656}{\csromanfolio{273}เตน สมเยน พุทฺโธ ภควา สาวตฺถิยํ วิหรติ เชตวเน อนาถปิณฺฑิกสฺส อาราเม.\csromanspacer{} เตน โข ปน สมเยน สาโฬฺห มิคารนตฺตา ภิกฺขุนิสํฆสฺส วิหารํ กตฺตุกาโม โหติ.\csromanspacer{} อถ โข สาโฬฺห มิคารนตฺตา ภิกฺขุนิโย อุปสงฺกมิตฺวา เอตทโวจ “อิจฺฉามหํ อยฺเย ภิกฺขุนิสํฆสฺส วิหารํ กาตุํ, เทถ เม นวกมฺมิกํ ภิกฺขุนินฺ”ติ.\csromanspacer{} เตน โข ปน สมเยน จตสฺโส ภคินิโย ภิกฺขุนีสุ ปพฺพชิตา โหนฺติ นนฺทา นนฺทวตี สุนฺทรีนนฺทา ถุลฺลนนฺทาติ.\csromanspacer{} ตาสุ สุนฺทรีนนฺทา ภิกฺขุนี ตรุณปพฺพชิตา อภิรูปา โหติ ทสฺสนียา ปาสาทิกา ปณฺฑิตา พฺยตฺตา เมธาวินี ทกฺขา อนลสา ตตฺรุปายาย วีมํสาย สมนฺนาคตา อลํ กาตุํ อลํ สํวิธาตุํ.\csromanspacer{} อถ โข ภิกฺขุนิสํโฆ สุนฺทรีนนฺทํ ภิกฺขุนิํ สมฺมนฺนิตฺวา สาฬฺหสฺส มิคารนตฺตุโน นวกมฺมิกํ อทาสิ.\csromanspacer{} เตน โข ปน สมเยน สุนฺทรีนนฺทา ภิกฺขุนี สาฬฺหสฺส มิคารนตฺตุโน นิเวสนํ อภิกฺขณํ คจฺฉติ “วาสิํ เทถ ปรสุํ\footnote{ผรสุํ (สฺยา, ก)} เทถ กุฐาริํ\footnote{กุธาริํ (ก)} เทถ กุทาลํ เทถ นิขาทนํ เทถา”ติ, สาโฬฺหปิ มิคารนตฺตา ภิกฺขุนุปสฺสยํ อภิกฺขณํ คจฺฉติ กตากตํ ชานิตุํ, เต อภิณฺหทสฺสเนน ปฏิพทฺธจิตฺตา อเหสุํ.}

\prose{อถ โข สาโฬฺห มิคารนตฺตา สุนฺทรีนนฺทํ ภิกฺขุนิํ ทูเสตุํ โอกาสํ อลภมาโน เอตเทวอตฺถาย ภิกฺขุนิสํฆสฺส ภตฺตํ อกาสิ.\csromanspacer{} อถ โข สาโฬฺห มิคารนตฺตา ภตฺตคฺเค อาสนํ ปญฺญเปนฺโต “เอตฺตกา ภิกฺขุนิโย อยฺยาย สุนฺทรีนนฺทาย วุฑฺฒตรา”ติ เอกมนฺตํ อาสนํ ปญฺญเปสิ, “เอตฺตกา นวกตรา”ติ เอกมนฺตํ อาสนํ ปญฺญเปสิ, ปฏิจฺฉนฺเน \csromanfolio{274}โอกาเส นิกูเฏ สุนฺทรีนนฺทาย ภิกฺขุนิยา อาสนํ ปญฺญเปสิ, ยถา เถรา ภิกฺขุนิโย ชาเนยฺยุํ “นวกานํ ภิกฺขุนีนํ สนฺติเก นิสินฺนา”ติ, นวกาปิ ภิกฺขุนิโย ชาเนยฺยุํ “เถรานํ ภิกฺขุนีนํ สนฺติเก นิสินฺนา”ติ.\csromanspacer{} อถ โข สาโฬฺห มิคารนตฺตา ภิกฺขุนิสํฆสฺส กาลํ อาโรจาเปสิ “กาโล อยฺเย นิฏฺฐิตํ ภตฺตนฺ”ติ.\csromanspacer{} สุนฺทรีนนฺทา ภิกฺขุนี สลฺลกฺเขตฺวา “น พหุกโต สาโฬฺห มิคารนตฺตา ภิกฺขุนิสํฆสฺส ภตฺตํ อกาสิ, มํ โส ทูเสตุกาโม, สจาหํ คมิสฺสามิ วิสฺสโร เม ภวิสฺสตี”ติ อนฺเตวาสินิํ ภิกฺขุนิํ อาณาเปสิ “คจฺฉ เม ปิณฺฑปาตํ นีหร, โย เจ มํ ปุจฺฉติ, ‘คิลานา’ติ ปฏิเวเทหี”ติ. “เอวํ อยฺเย”ติ โข สา ภิกฺขุนี สุนฺทรีนนฺทาย ภิกฺขุนิยา ปจฺจสฺโสสิ.}

\prose{เตน โข ปน สมเยน สาโฬฺห มิคารนตฺตา พหิทฺวารโกฏฺฐเก ฐิโต โหติ สุนฺทรีนนฺทํ ภิกฺขุนิํ ปฏิปุจฺฉนฺโต “กหํ อยฺเย อยฺยา สุนฺทรีนนฺทา, กหํ อยฺเย อยฺยา สุนฺทรีนนฺทา”ติ.\csromanspacer{} เอวํ วุตฺเต สุนฺทรีนนฺทาย ภิกฺขุนิยา อนฺเตวาสินี ภิกฺขุนี สาฬฺหํ มิคารนตฺตารํ เอตทโวจ “คิลานาวุโส ปิณฺฑปาตํ นีหริสฺสามี”ติ.\csromanspacer{} อถ โข สาโฬฺห มิคารนตฺตา “ยมฺปาหํ อตฺถาย\footnote{ยํปาหํ (สฺยา)} ภิกฺขุนิสํฆสฺส ภตฺตํ อกาสิํ อยฺยาย สุนฺทรีนนฺทาย การณา”ติ มนุสฺเส อาณาเปตฺวา “ภิกฺขุนิสํฆํ ภตฺเตน ปริวิสถา”ติ วตฺวา เยน ภิกฺขุนุปสฺสโย เตนุปสงฺกมิ.\csromanspacer{} เตน โข ปน สมเยน สุนฺทรีนนฺทา ภิกฺขุนี พหารามโกฏฺฐเก ฐิตา โหติ สาฬฺหํ มิคารนตฺตารํ ปติมาเนนฺตี.\csromanspacer{} อทฺทสา โข สุนฺทรีนนฺทา ภิกฺขุนี สาฬฺหํ มิคารนตฺตารํ ทูรโตว อาคจฺฉนฺตํ, ทิสฺวาน อุปสฺสยํ ปวิสิตฺวา สสีสํ ปารุปิตฺวา มญฺจเก นิปชฺชิ.\csromanspacer{} อถ โข สาโฬฺห มิคารนตฺตา เยน สุนฺทรีนนฺทา ภิกฺขุนี เตนุปสงฺกมิ, อุปสงฺกมิตฺวา สุนฺทรีนนฺทํ ภิกฺขุนิํ เอตทโวจ “กิํ เต อยฺเย อผาสุ, กิสฺส นิปนฺนาสี”ติ.\csromanspacer{} เอวเญฺหตํ อาวุโส โหติ ยา อนิจฺฉนฺตํ อิจฺฉตีติ. “กฺยาหํ ตํ อยฺเย น อิจฺฉิสฺสามิ, อปิ จาหํ โอกาสํ น ลภามิ ตํ ทูเสตุนฺ”ติ อวสฺสุโต อวสฺสุตาย สุนฺทรีนนฺทาย ภิกฺขุนิยา กายสํสคฺคํ สมาปชฺชิ.}

\prose{เตน โข ปน สมเยน อญฺญตรา ภิกฺขุนี ชราทุพฺพลา จรณคิลานา สุนฺทรีนนฺทาย ภิกฺขุนิยา อวิทูเร นิปนฺนา โหติ.\csromanspacer{} อทฺทสา โข สา \csromanfolio{275}ภิกฺขุนี สาฬฺหํ มิคารนตฺตารํ อวสฺสุตํ อวสฺสุตาย สุนฺทรีนนฺทาย ภิกฺขุนิยา กายสํสคฺคํ สมาปชฺชนฺตํ, ทิสฺวาน อุชฺฌายติ ขิยฺยติ วิปาเจติ “กถํ หิ นาม อยฺยา สุนฺทรีนนฺทา อวสฺสุตา อวสฺสุตสฺส ปุริสปุคฺคลสฺส กายสํสคฺคํ สาทิยิสฺสตี”ติ.\csromanspacer{} อถ โข สา ภิกฺขุนี \textbf{ภิกฺขุนีนํ เอตมตฺถํ อาโรเจสิ.}\csromanspacer{}\textbf{ ยา ตา ภิกฺขุนิโย อปฺปิจฺฉา สนฺตุฏฺฐา} \textbf{ลชฺชินิโย กุกฺกุจฺจิกา สิกฺขากามา, ตา อุชฺฌายนฺติ ขิยฺยนฺติ วิปาเจนฺติ “กถํ} \textbf{หิ นาม อยฺยา สุนฺทรีนนฺทา อวสฺสุตา อวสฺสุตสฺส ปุริสปุคฺคลสฺส} \textbf{กายสํสคฺคํ สาทิยิสฺสตี”ติ.}\csromanspacer{}\textbf{ อถ โข ตา ภิกฺขุนิโย ภิกฺขูนํ} \textbf{เอตมตฺถํ อาโรเจสุํ.}\csromanspacer{}\textbf{ เย เต ภิกฺขู อปฺปิจฺฉา สนฺตุฏฺฐา ลชฺชิโน กุกฺกุจฺจกา} สิกฺขากามา, เต อุชฺฌายนฺติ ขิยฺยนฺติ วิปาเจนฺติ “กถํ หิ นาม สุนฺทรีนนฺทา ภิกฺขุนี อวสฺสุตา อวสฺสุตสฺส ปุริสปุคฺคลสฺส กายสํสคฺคํ สาทิยิสฺสตี”ติ.}

\prose{อถ โข เต ภิกฺขู สุนฺทรีนนฺทํ ภิกฺขุนิํ อเนกปริยาเยน วิครหิตฺวา ภควโต เอตมตฺถํ อาโรเจสุํ.\csromanspacer{} อถ โข ภควา เอตสฺมิํ นิทาเน เอตสฺมิํ ปกรเณ ภิกฺขุสํฆํ สนฺนิปาตาเปตฺวา ภิกฺขู ปฏิปุจฺฉิ “สจฺจํ กิร ภิกฺขเว สุนฺทรีนนฺทา ภิกฺขุนี อวสฺสุตา อวสฺสุตสฺส ปุริสปุคฺคลสฺส กายสํสคฺคํ สาทิยตี”ติ\footnote{สาทิยีติ (ก)}.\csromanspacer{} สจฺจํ ภควาติ.\csromanspacer{} วิครหิ พุทฺโธ ภควา “อนนุจฺฉวิกํ ภิกฺขเว สุนฺทรีนนฺทาย ภิกฺขุนิยา อนนุโลมิกํ อปฺปติรูปํ อสฺสามณกํ อกปฺปิยํ อกรณียํ.\csromanspacer{} กถํ หิ นาม ภิกฺขเว สุนฺทรีนนฺทา ภิกฺขุนี อวสฺสุตา อวสฺสุตสฺส ปุริสปุคฺคลสฺส กายสํสคฺคํ สาทิยิสฺสติ.\csromanspacer{} เนตํ ภิกฺขเว อปฺปสนฺนานํ วา ปสาทาย ปสนฺนานํ วา ภิยฺโยภาวาย.\csromanspacer{} อถ เขฺวตํ ภิกฺขเว อปฺปสนฺนานญฺเจว อปฺปสาทาย ปสนฺนานญฺจ เอกจฺจานํ อญฺญถตฺตายา”ติ.\csromanspacer{} อถ โข ภควา สุนฺทรีนนฺทํ ภิกฺขุนิํ อเนกปริยาเยน วิครหิตฺวา ทุพฺภรตาย ทุปฺโปสตาย มหิจฺฉตาย อสนฺตุฏฺฐิตาย\footnote{อสนฺตุฏฺฐตาย (สฺยา), อสนฺตุฏฺฐิยา (ก)} สงฺคณิกาย โกสชฺชสฺส อวณฺณํ ภาสิตฺวา อเนกปริยาเยน สุภรตาย สุโปสตาย อปฺปิจฺฉสฺส สนฺตุฏฺฐสฺส\footnote{สนฺตุฏฺฐิยา (ก)} สลฺเลขสฺส ธุตสฺส ปาสาทิกสฺส อปจยสฺส วีริยารมฺภสฺส วณฺณํ ภาสิตฺวา ภิกฺขูนํ ตทนุจฺฉวิกํ ตทนุโลมิกํ ธมฺมิํ กถํ กตฺวา ภิกฺขู อามนฺเตสิ เตน หิ ภิกฺขเว ภิกฺขุนีนํ สิกฺขาปทํ ปญฺญเปสฺสามิ ทส อตฺถวเส ปฏิจฺจ สํฆสุฏฺฐุตาย สํฆผาสุตาย ทุมฺมงฺกูนํ ภิกฺขุนีนํ \csromanfolio{276}นิคฺคหาย เปสลานํ ภิกฺขุนีนํ ผาสุวิหาราย ทิฏฺฐธมฺมิกานํ อาสวานํ สํวราย สมฺปรายิกานํ อาสวานํ ปฏิฆาตาย อปฺปสนฺนานํ ปสาทาย ปสนฺนานํ ภิยฺโยภาวาย สทฺธมฺมฏฺฐิติยา วินยานุคฺคหาย.\csromanspacer{} เอวญฺจ ปน ภิกฺขเว ภิกฺขุนิโย อิมํ สิกฺขาปทํ อุทฺทิสนฺตุ\csromandash{}}

\proseitem{657}{\textbf{“ยา ปน ภิกฺขุนี อวสฺสุตา อวสฺสุตสฺส ปุริสปุคฺคลสฺส อธกฺขกํ อุพฺภชาณุมณฺฑลํ อามสนํ วา ปรามสนํ วา คหณํ วา ฉุปนํ วา ปฏิปีฬนํ วา สาทิเยยฺย, อยมฺปิ ปาราชิกา โหติ อสํวาสา อุพฺภชาณุมณฺฑลิกา”}ติ.}

\proseitem{658}{\textbf{ยา ปนา}ติ ยา ยาทิสา ยถายุตฺตา ยถาชจฺจา ยถานามา ยถาโคตฺตา ยถาสีลา ยถาวิหารินี ยถาโคจรา เถรา วา นวา วา มชฺฌิมา วา, เอสา วุจฺจติ ยา ปนาติ.}

\prose{\textbf{ภิกฺขุนี}ติ ภิกฺขิกาติ ภิกฺขุนี, ภิกฺขาจริยํ อชฺฌุปคตาติ ภิกฺขุนี ภินฺนปฏธราติ ภิกฺขุนี, สมญฺญาย ภิกฺขุนี, ปฏิญฺญาย ภิกฺขุนี, เอหิ ภิกฺขุนีติ ภิกฺขุนี, ตีหิ สรณคมเนหิ อุปสมฺปนฺนาติ ภิกฺขุนี, ภทฺรา ภิกฺขุนี, สารา ภิกฺขุนี, เสขา ภิกฺขุนี, อเสขา ภิกฺขุนี, สมคฺเคน อุภโตสํเฆน ญตฺติจตุตฺเถน กมฺเมน อกุปฺเปน ฐานารเหน อุปสมฺปนฺนาติ ภิกฺขุนี, ตตฺร ยายํ ภิกฺขุนี สมคฺเคน อุภโตสํเฆน ญตฺติจตุตฺเถน กมฺเมน อกุปฺเปน ฐานารเหน อุปสมฺปนฺนา, อยํ อิมสฺมิํ อตฺเถ อธิปฺเปตา ภิกฺขุนีติ.}

\prose{\textbf{อวสฺสุตา} นาม สารตฺตา อเปกฺขวตี ปฏิพทฺธจิตฺตา.}

\prose{\textbf{อวสฺสุโต} นาม สารตฺโต อเปกฺขวา ปฏิพทฺธจิตฺโต.}

\prose{\textbf{ปุริสปุคฺคโล} นาม มนุสฺสปุริโส น ยกฺโข น เปโต น ติรจฺฉานคโต วิญฺญู ปฏิพโล กายสํสคฺคํ สมาปชฺชิตุํ.}

\prose{\textbf{อธกฺขกนฺ}ติ เหฏฺฐกฺขกํ.}

\prose{\textbf{อุพฺภชาณุมณฺฑลนฺ}ติ อุปริชาณุมณฺฑลํ.}

\prose{\textbf{อามสนํ} นาม อามฏฺฐมตฺตํ.}

\prose{\textbf{ปรามสนํ} นาม อิโตจิโต จ สญฺโจปนํ.}

\prose{\textbf{ปคณํ} นาม คหิตมตฺตํ.}

\prose{\textbf{ฉุปนํ} นาม ผุฏฺฐมตฺตํ.}

\prose{\csromanfolio{277}\textbf{ปฏิปีฬนํ วา สาทิเยยฺยา}ติ องฺคํ คเหตฺวา นิปฺปีฬนํ สาทิยติ.}

\prose{\textbf{อยมฺปี}ติ ปุริมาโย อุปาทาย วุจฺจติ.}

\prose{\textbf{ปาราชิกา โหตี}ติ เสยฺยถาปิ นาม ปุริโส สีสจฺฉินฺโน อภพฺโพ เตน สรีรพนฺธเนน ชีวิตุํ, เอวเมว ภิกฺขุนี อวสฺสุตา อวสฺสุตสฺส ปุริสปุคฺคลสฺส อธกฺขกํ อุพฺภชาณุมณฺฑลํ อามสนํ วา ปรามสนํ วา คหณํ วา ฉุปนํ วา ปฏิปีฬนํ วา สาทิยนฺตี อสฺสมณี โหติ อสกฺยธีตา, เตน วุจฺจติ ปาราชิกา โหตีติ.}

\prose{\textbf{อสํวาสาติ สํวาโส} นาม เอกกมฺมํ เอกุทฺเทโส สมสิกฺขตา, เอโส สํวาโส นาม.\csromanspacer{} โส ตาย สทฺธิํ นตฺถิ, เตน วุจฺจติ อสํวาสาติ.}

\proseitem{659}{อุภโตอวสฺสุเต อธกฺขกํ อุพฺภชาณุมณฺฑลํ กาเยน กายํ อามสติ, อาปตฺติ ปาราชิกสฺส.\csromanspacer{} กาเยน กายปฏิพทฺธํ อามสติ, อาปตฺติ ถุลฺลจฺจยสฺส.\csromanspacer{} กายปฏิพทฺเธน กายํ อามสติ, อาปตฺติ ถุลฺลจฺจยสฺส.\csromanspacer{} กายปฏิพทฺเธน กายปฏิพทฺธํ อามสติ, อาปตฺติ ทุกฺกฏสฺส.}

\prose{นิสฺสคฺคิเยน กายํ อามสติ, อาปตฺติ ทุกฺกฏสฺส.\csromanspacer{} นิสฺสคฺคิเยน กายปฏิพทฺธํ อามสติ, อาปตฺติ ทุกฺกฏสฺส.\csromanspacer{} นิสฺสคฺคิเยน นิสฺสคฺคิยํ อามสติ, อาปตฺติ ทุกฺกฏสฺส.}

\prose{อุพฺภกฺขกํ อโธชาณุมณฺฑลํ กาเยน กายํ อามสติ, อาปตฺติ ถุลฺลจฺจยสฺส.\csromanspacer{} กาเยน กายปฏิพทฺธํ อามสติ, อาปตฺติ ทุกฺกฏสฺส.\csromanspacer{} กายปฏิพทฺเธน กายํ อามสติ, อาปตฺติ ทุกฺกฏสฺส.\csromanspacer{} กายปฏิพทฺเธน กายปฏิพทฺธํ อามสติ, อาปตฺติ ทุกฺกฏสฺส.}

\prose{นิสฺสคฺคิเยน กายํ อามสติ, อาปตฺติ ทุกฺกฏสฺส.\csromanspacer{} นิสฺสคฺคิเยน กายปฏิพทฺธํ อามสติ, อาปตฺติ ทุกฺกฏสฺส.\csromanspacer{} นิสฺสคฺคิเยน นิสฺสคฺคิยํ อามสติ, อาปตฺติ ทุกฺกฏสฺส.}

\proseitem{660}{เอกโตอวสฺสุเต อธกฺขกํ อุพฺภชาณุมณฺฑลํ กาเยน กายํ อามสติ, อาปตฺติ ถุลฺลจฺจยสฺส.\csromanspacer{} กาเยน กายปฏิพทฺธํ \csromanfolio{278}อามสติ, อาปตฺติ ทุกฺกฏสฺส.\csromanspacer{} กายปฏิพทฺเธน กายํ อามสติ, อาปตฺติ ทุกฺกฏสฺส.\csromanspacer{} กายปฏิพทฺเธน กายปฏิพทฺธํ อามสติ, อาปตฺติ ทุกฺกฏสฺส.}

\prose{นิสฺสคฺคิเยน กายํ อามสติ, อาปตฺติ ทุกฺกฏสฺส.\csromanspacer{} นิสฺสคฺคิเยน กายปฏิพทฺธํ อามสติ, อาปตฺติ ทุกฺกฏสฺส.\csromanspacer{} นิสฺสคฺคิเยน นิสฺสคฺคิยํ อามสติ, อาปตฺติ ทุกฺกฏสฺส.}

\prose{อุพฺภกฺขกํ อโธชาณุมณฺฑลํ กาเยน กายํ อามสติ, อาปตฺติ ทุกฺกฏสฺส.\csromanspacer{} กาเยน กายปฏิพทฺธํ อามสติ, อาปตฺติ ทุกฺกฏสฺส.\csromanspacer{} กายปฏิพทฺเธน กายํ อามสติ, อาปตฺติ ทุกฺกฏสฺส.\csromanspacer{} กายปฏิพทฺเธน กายปฏิพทฺธํ อามสติ, อาปตฺติ ทุกฺกฏสฺส.}

\prose{นิสฺสคฺคิเยน กายํ อามสติ, อาปตฺติ ทุกฺกฏสฺส.\csromanspacer{} นิสฺสคฺคิเยน กายปฏิพทฺธํ อามสติ, อาปตฺติ ทุกฺกฏสฺส.\csromanspacer{} นิสฺสคฺคิเยน นิสฺสคฺคิยํ อามสติ, อาปตฺติ ทุกฺกฏสฺส.}

\proseitem{661}{อุภโตอวสฺสุเต ยกฺขสฺส วา เปตสฺส วา ปณฺฑกสฺส วา ติรจฺฉานคต\-มนุสฺสวิคฺคหสฺส วา อธกฺขกํ อุพฺภชาณุมณฺฑลํ กาเยน กายํ อามสติ, อาปตฺติ ถุลฺลจฺจยสฺส.\csromanspacer{} กาเยน กายปฏิพทฺธํ อามสติ, อาปตฺติ ทุกฺกฏสฺส.\csromanspacer{} กายปฏิพทฺเธน กายํ อามสติ, อาปตฺติ ทุกฺกฏสฺส.\csromanspacer{} กายปฏิพทฺเธน กายปฏิพทฺธํ อามสติ, อาปตฺติ ทุกฺกฏสฺส.}

\prose{นิสฺสคฺคิเยน กายํ อามสติ, อาปตฺติ ทุกฺกฏสฺส.\csromanspacer{} นิสฺสคฺคิเยน กายปฏิพทฺธํ อามสติ, อาปตฺติ ทุกฺกฏสฺส.\csromanspacer{} นิสฺสคฺคิเยน นิสฺสคฺคิยํ อามสติ, อาปตฺติ ทุกฺกฏสฺส.}

\prose{อุพฺภกฺขกํ อโธชาณุมณฺฑลํ กาเยน กายํ อามสติ, อาปตฺติ ทุกฺกฏสฺส.\csromanspacer{} กาเยน กายปฏิพทฺธํ อามสติ, อาปตฺติ ทุกฺกฏสฺส.\csromanspacer{} กายปฏิพทฺเธน กายํ อามสติ, อาปตฺติ ทุกฺกฏสฺส.\csromanspacer{} กายปฏิพทฺเธน กายปฏิพทฺธํ อามสติ, อาปตฺติ ทุกฺกฏสฺส.}

\prose{นิสฺสคฺคิเยน อามสติ, อาปตฺติ ทุกฺกฏสฺส.\csromanspacer{} นิสฺสคฺคิเยน กายปฏิพทฺธํ อามสติ, อาปตฺติ ทุกฺกฏสฺส.\csromanspacer{} นิสฺสคฺคิเยน นิสฺสคฺคิยํ อามสติ, อาปตฺติ ทุกฺกฏสฺส.}

\proseitem{662}{\csromanfolio{279}เอกโตอวสฺสุเต อธกฺขกํ อุพฺภชาณุมณฺฑลํ กาเยน กายํ อามสติ, อาปตฺติ ทุกฺกฏสฺส.\csromanspacer{} กาเยน กายปฏิพทฺธํ อามสติ, อาปตฺติ ทุกฺกฏสฺส.\csromanspacer{} กายปฏิพทฺเธน กายํ อามสติ, อาปตฺติ ทุกฺกฏสฺส.\csromanspacer{} กายปฏิพทฺเธน กายปฏิพทฺธํ อามสติ, อาปตฺติ ทุกฺกฏสฺส.}

\prose{นิสฺสคฺคิเยน กายํ อามสติ, อาปตฺติ ทุกฺกฏสฺส.\csromanspacer{} นิสฺสคฺคิเยน กายปฏิพทฺธํ อามสติ, อาปตฺติ ทุกฺกฏสฺส.\csromanspacer{} นิสฺสคฺคิเยน นิสฺสคฺคิยํ อามสติ, อาปตฺติ ทุกฺกฏสฺส.}

\prose{อุพฺภกฺขกํ อโธชาณุมณฺฑลํ กาเยน กายํ อามสติ, อาปตฺติ ทุกฺกฏสฺส.\csromanspacer{} กาเยน กายปฏิพทฺธํ อามสติ, อาปตฺติ ทุกฺกฏสฺส.\csromanspacer{} กายปฏิพทฺเธน กายํ อามสติ, อาปตฺติ ทุกฺกฏสฺส.\csromanspacer{} กายปฏิพทฺเธน กายปฏิพทฺธํ อามสติ, อาปตฺติ ทุกฺกฏสฺส.}

\prose{นิสฺสคฺคิเยน กายํ อามสติ, อาปตฺติ ทุกฺกฏสฺส.\csromanspacer{} นิสฺสคฺคิเยน กายปฏิพทฺธํ อามสติ, อาปตฺติ ทุกฺกฏสฺส.\csromanspacer{} นิสฺสคฺคิเยน นิสฺสคฺคิยํ อามสติ, อาปตฺติ ทุกฺกฏสฺส.}

\proseitemwithcloserruled{663}{อนาปตฺติ อสญฺจิจฺจ อสฺสติยา อชานนฺติยา อสาทิยนฺติยา อุมฺมตฺติกาย ขิตฺตจิตฺตาย เวทนาฏฺฏาย อาทิกมฺมิกายาติ.}{ปฐมปาราชิกํ สมตฺตํ\footnote{ภิกฺขุนิวิภงฺเค สิกฺขาปทคณนา ภิกฺขูหิ อสาธารณวเสน ปกาสิตาติ เวทิตพฺพา.}.}
\csromanmatikaanchor{mtk.120}
\csromantocmarkref{section}{2. ทุติยปาราชิกสิกฺขาปท}{mtk.120}
\bookmark[dest={mtk.120},level=1]{2. ทุติยปาราชิกสิกฺขาปท}
\setcsromanheadh{2. ทุติยปาราชิกสิกฺขาปท}
\csromanheader{1}{2. \textbf{ทุติยปาราชิก}}
\proseitem{664}{เตน สมเยน พุทฺโธ ภควา สาวตฺถิยํ วิหรติ เชตวเน อนาถปิณฺฑิกสฺส อาราเม.\csromanspacer{} เตน โข ปน สมเยน สุนฺทรีนนฺทา ภิกฺขุนี สาเฬฺหน มิคารนตฺตุนา คพฺภินี โหติ.\csromanspacer{} ยาว คพฺโภ ตรุโณ อโหสิ, ตาว ฉาเทสิ, ปริปกฺเก คพฺเภ วิพฺภมิตฺวา วิชายิ.\csromanspacer{} ภิกฺขุนิโย ถุลฺลนนฺทํ ภิกฺขุนิํ เอตทโวจุํ “สุนฺทรีนนฺทา โข อยฺเย อจิรวิพฺภนฺตา วิชาตา, กจฺจิ โน สา ภิกฺขุนีเยว สมานา คพฺภินี”ติ.\csromanspacer{} เอวํ อยฺเยติ.\csromanspacer{} กิสฺส ปน ตฺวํ อยฺเย ชานํ ปาราชิกํ ธมฺมํ อชฺฌาปนฺนํ ภิกฺขุนิํ เนวตฺตนา ปฏิโจเทสิ น คณสฺส อาโรเจสีติ.\csromanspacer{} โย เอติสฺสา อวณฺโณ มเยฺหโส อวณฺโณ, ยา เอติสฺสา อกิตฺติ มเยฺหสา อกิตฺติ, โย เอติสฺสา อยโส มเยฺหโส อยโส, โย เอติสฺสา อลาโภ \csromanfolio{280}มเยฺหโส อลาโภ, กฺยาหํ อยฺเย อตฺตโน อวณฺณํ อตฺตโน อกิตฺติํ อตฺตโน อยสํ อตฺตโน อลาภํ ปเรสํ อาโรเจสฺสามีติ.\csromanspacer{} ยา ตา ภิกฺขุนิโย อปฺปิจฺฉา ฯเปฯ ตา อุชฺฌายนฺติ ขิยฺยนฺติ วิปาเจนฺติ “กถํ หิ นาม อยฺยา ถุลฺลนนฺทา ชานํ ปาราชิกํ ธมฺมํ อชฺฌาปนฺนํ ภิกฺขุนิํ เนวตฺตนา ปฏิโจเทสฺสติ, น คณสฺส อาโรเจสฺสตี”ติ.\csromanspacer{} อถ โข ตา ภิกฺขุนิโย ภิกฺขูนํ เอตมตฺถํ อาโรเจสุํ.\csromanspacer{} ภิกฺขู\footnote{เย เต ภิกฺขู ฯเปฯ (?)} ภควโต เอตมตฺถํ อาโรเจสุํ.\csromanspacer{} อถ โข ภควา เอตสฺมิํ นิทาเน เอตสฺมิํ ปกรเณ ภิกฺขุสํฆํ สนฺนิปาตาเปตฺวา ภิกฺขู ปฏิปุจฺฉิ “สจฺจํ กิร ภิกฺขเว ถุลฺลนนฺทา ภิกฺขุนี ชานํ ปาราชิกํ ธมฺมํ อชฺฌาปนฺนํ ภิกฺขุนิํ เนวตฺตนา ปฏิโจเทติ\footnote{ปฏิโจเทสิ, ... อาโรเจสีติ (ก)}, น คณสฺส อาโรเจตี”ติ2.\csromanspacer{} สจฺจํ ภควาติ.\csromanspacer{} วิครหิ พุทฺโธ ภควา ฯเปฯ.\csromanspacer{} กถํ หิ นาม ภิกฺขเว ถุลฺลนนฺทา ภิกฺขุนี ชานํ ปาราชิกํ ธมฺมํ อชฺฌาปนฺนํ ภิกฺขุนิํ เนวตฺตนา ปฏิโจเทสฺสติ, น คณสฺส อาโรเจสฺสติ.\csromanspacer{} เนตํ ภิกฺขเว อปฺปสนฺนานํ วา ปสาทาย ฯเปฯ.\csromanspacer{} เอวญฺจ ปน ภิกฺขเว ภิกฺขุนิโย อิมํ สิกฺขาปทํ อุทฺทิสนฺตุ\csromandash{}}

\proseitem{665}{\textbf{“ยา ปน ภิกฺขุนี ชานํ ปาราชิกํ ธมฺมํ อชฺฌาปนฺนํ ภิกฺขุนิํ เนวตฺตนา ปฏิโจเทยฺย น คณสฺส อาโรเจยฺย, ยทา จ สา ฐิตา วา อสฺส จุตา วา นาสิตา วา อวสฺสฏา วา, สา ปจฺฉา เอวํ วเทยฺย ‘ปุพฺเพวาหํ อยฺเย อญฺญาสิํ เอตํ ภิกฺขุนิํ เอวรูปา จ เอวรูปา จ สา ภคินีติ, โน จ โข อตฺตนา ปฏิโจเทสฺสํ, น คณสฺส อาโรเจสฺสนฺ’ติ อยมฺปิ ปาราชิกา โหติ อสํวาสา วชฺชปฺปฏิจฺฉาทิกา”}ติ.}

\proseitem{666}{\textbf{ยา ปนา}ติ ยา ยาทิสา ฯเปฯ.\csromanspacer{} \textbf{ภิกฺขุนี}ติ ฯเปฯ อยํ อิมสฺมิํ อตฺเถ อธิปฺเปตา ภิกฺขุนีติ.}

\prose{\textbf{ชานาติ} นาม สามํ วา ชานาติ, อญฺเญ วา ตสฺสา อาโรเจนฺติ, สา วา อาโรเจติ.}

\prose{\textbf{ปาราชิกํ ธมฺมํ อชฺฌาปนฺนนฺ}ติ อฏฺฐนฺนํ ปาราชิกานํ อญฺญตรํ ปาราชิกํ อชฺฌาปนฺนํ.}

\prose{\textbf{เนวตฺตนา ปฏิโจเทยฺยา}ติ น สยํ โจเทยฺย.}

\prose{\textbf{น คณสฺส อาโรเจยฺยา}ติ น อญฺญาสํ ภิกฺขุนีนํ อาโรเจยฺย.}

\prose{\csromanfolio{281}\textbf{ยทา จ สา ฐิตา วา อสฺส จุตา วา}ติ \textbf{ฐิตา} นาม สลิงฺเค ฐิตา วุจฺจติ.\csromanspacer{} \textbf{จุตา} นาม กาลงฺกตา วุจฺจติ.\csromanspacer{} \textbf{นาสิตา} นาม สยํ วา วิพฺภนฺตา โหติ อญฺเญหิ วา นาสิตา.\csromanspacer{} \textbf{อวสฺสฏา} นาม ติตฺถายตนํ สงฺกนฺตา วุจฺจติ.\csromanspacer{} สา ปจฺฉา เอวํ วเทยฺย “ปุพฺเพวาหํ อยฺเย อญฺญาสิํ เอตํ ภิกฺขุนิํ เอวรูปา จ เอวรูปา จ สา ภคินี”ติ.}

\prose{\textbf{โน จ โข อตฺตนา ปฏิโจเทสฺสนฺ}ติ สยํ วา น โจเทสฺสํ.}

\prose{\textbf{น คณสฺส อาโรเจสฺสนฺ}ติ น อญฺญาสํ ภิกฺขุนีนํ อาโรเจสฺสํ.}

\prose{\textbf{อยมฺปี}ติ ปุริมาโย อุปาทาย วุจฺจติ.}

\prose{\textbf{ปาราชิกา โหตี}ติ เสยฺยถาปิ นาม ปณฺฑุปลาโส พนฺธนา ปวุตฺโต อภพฺโพ หริตตฺถาย, เอวเมว ภิกฺขุนี ชานํ ปาราชิกํ ธมฺมํ อชฺฌาปนฺนํ ภิกฺขุนิํ เนวตฺตนา ปฏิโจเทสฺสามิ, น คณสฺส อาโรเจสฺสามีติ ธุรํ นิกฺขิตฺตมตฺเต อสฺสมณี โหติ อสกฺยธีตา, เตน วุจฺจติ ปาราชิกา โหตีติ.}

\prose{\textbf{อสํวาสา}ติ \textbf{สํวาโส} นาม เอกกมฺมํ เอกุทฺเทโส สมสิกฺขตา, เอโส สํวาโส นาม.\csromanspacer{} โส ตาย สทฺธิํ นตฺถิ, เตน วุจฺจติ อสํวาสาติ.}

\proseitemwithcloserruled{667}{อนาปตฺติ “สํฆสฺส ภณฺฑนํ วา กลโห วา วิคฺคโห วา วิวาโท วา ภวิสฺสตี”ติ นาโรเจติ, “สํฆเภโท วา สํฆราชิ วา ภวิสฺสตี”ติ นาโรเจติ, “อยํ กกฺขฬา ผรุสา ชีวิตนฺตรายํ วา พฺรหฺมจริยนฺตรายํ วา กริสฺสตี”ติ นาโรเจติ, อญฺญา ปติรูปา ภิกฺขุนิโย อปสฺสนฺตี นาโรเจติ, นจฺฉาเทตุกามา นาโรเจติ, ปญฺญายิสฺสติ สเกน กมฺเมนาติ นาโรเจติ, อุมฺมตฺติกาย ฯเปฯ อาทิกมฺมิกายาติ.}{ทุติยปาราชิกํ สมตฺตํ.}
\csromanmatikaanchor{mtk.121}
\csromantocmarkref{section}{3. ตติยปาราชิกสิกฺขาปท}{mtk.121}
\bookmark[dest={mtk.121},level=1]{3. ตติยปาราชิกสิกฺขาปท}
\setcsromanheadh{3. ตติยปาราชิกสิกฺขาปท}
\csromanheader{1}{3. \textbf{ตติยปาราชิก}}
\proseitem{668}{เตน สมเยน พุทฺโธ ภควา สาวตฺถิยํ วิหรติ เชตวเน อนาถปิณฺฑิกสฺส อาราเม.\csromanspacer{} เตน โข ปน สมเยน ถุลฺลนนฺทา ภิกฺขุนี สมคฺเคน สํเฆน อุกฺขิตฺตํ อริฏฺฐํ ภิกฺขุํ คทฺธพาธิปุพฺพํ อนุวตฺตติ. \csromanfolio{282}ยา ตา ภิกฺขุนิโย อปฺปิจฺฉา ฯเปฯ ตา อุชฺฌายนฺติ ขิยฺยนฺติ วิปาเจนฺติ “กถํ หิ นาม อยฺยา ถุลฺลนนฺทา สมคฺเคน สํเฆน อุกฺขิตฺตํ อริฏฺฐํ ภิกฺขุํ คทฺธพาธิปุพฺพํ อนุวตฺติสฺสตี”ติ ฯเปฯ.\csromanspacer{} สจฺจํ กิร ภิกฺขเว ถุลฺลนนฺทา ภิกฺขุนี สมคฺเคน สํเฆน อุกฺขิตฺตํ อริฏฺฐํ ภิกฺขุํ คทฺธพาธิปุพฺพํ อนุวตฺตตีติ.\csromanspacer{} สจฺจํ ภควาติ.\csromanspacer{} วิครหิ พุทฺโธ ภควา ฯเปฯ.\csromanspacer{} กถํ หิ นาม ภิกฺขเว ถุลฺลนนฺทา ภิกฺขุนี สมคฺเคน สํเฆน อุกฺขิตฺตํ อริฏฺฐํ ภิกฺขุํ คทฺธพาธิปุพฺพํ อนุวตฺติสฺสติ.\csromanspacer{} เนตํ ภิกฺขเว อปฺปสนฺนานํ วา ปสาทาย ฯเปฯ.\csromanspacer{} เอวญฺจ ปน ภิกฺขเว ภิกฺขุนิโย อิมํ สิกฺขาปทํ อุทฺทิสนฺตุ\csromandash{}}

\proseitem{669}{\textbf{“ยา ปน ภิกฺขุนี สมคฺเคน สํเฆน อุกฺขิตฺตํ ภิกฺขุํ ธมฺเมน วินเยน สตฺถุสาสเนน อนาทรํ อปฺปฏิการํ อกตสหายํ ตมนุวตฺเตยฺย, สา ภิกฺขุนี ภิกฺขุนีหิ เอวมสฺส วจนียา ‘เอโส โข อยฺเย ภิกฺขุ สมคฺเคน สํเฆน อุกฺขิตฺโต ธมฺเมน วินเยน สตฺถุสาสเนน อนาทโร อปฺปฏิกาโร อกตสหาโย, มายฺเย เอตํ ภิกฺขุํ อนุวตฺตี’ติ.}\csromanspacer{} เอวญฺจ สา ภิกฺขุนี ภิกฺขุนีหิ วุจฺจมานา ตเถว ปคฺคเณฺหยฺย, สา ภิกฺขุนี ภิกฺขุนีหิ ยาวตติยํ สมนุภาสิตพฺพา ตสฺส ปฏินิสฺสคฺคาย, ยาวตติยํ เจ สมนุภาสิยมานา ตํ ปฏินิสฺสชฺเชยฺย, อิจฺเจตํ กุสลํ, โน เจ ปฏินิสฺสชฺเชยฺย, อยมฺปิ ปาราชิกา โหติ อสํวาสา อุกฺขิตฺตานุวตฺติกา”ติ.}

\proseitem{670}{\textbf{ยา ปนา}ติ ยา ยาทิสา ฯเปฯ.\csromanspacer{} \textbf{ภิกฺขุนี}ติ ฯเปฯ อยํ อิมสฺมิํ อตฺเถ อธิปฺเปตา ภิกฺขุนีติ.}

\prose{\textbf{สมคฺโค} นาม \textbf{สํโฆ} สมานสํวาสโก สมานสีมายํ ฐิโต.}

\prose{\textbf{อุกฺขิตฺโต} นาม อาปตฺติยา อทสฺสเน วา อปฺปฏิกมฺเม วา อปฺปฏินิสฺสคฺเค วา\footnote{อทสฺสเนน วา อปฺปฏิกมฺเมน วา อปฺปฏินิสฺสคฺเคน วา (ก), อทสฺสเน วา อปฺปฏิกมฺเม วา ปาปิกาย ทิฏฺฐิยา อปฺปฏินิสฺสคฺเค วา (?)} อุกฺขิตฺโต.}

\prose{\textbf{ธมฺเมน วินเยนา}ติ เยน ธมฺเมน เยน วินเยน.}

\prose{\textbf{สตฺถุสาสเนนา}ติ ชินสาสเนน พุทฺธสาสเนน.}

\prose{\textbf{อนาทโร} นาม สํฆํ วา คณํ วา ปุคฺคลํ วา กมฺมํ วา นาทิยติ.}

\prose{\textbf{อปฺปฏิกาโร} นาม อุกฺขิตฺโต อโนสาริโต.}

\prose{\textbf{อกตสหาโย} นาม สมานสํวาสกา ภิกฺขู วุจฺจนฺติ สหายา.\csromanspacer{} โส เตหิ สทฺธิํ นตฺถิ, เตน วุจฺจติ อกตสหาโยติ.}

\prose{\csromanfolio{283}\textbf{ตมนุวตฺเตยฺยา}ติ ยํทิฏฺฐิโก โส โหติ ยํขนฺติโก ยํรุจิโก, สาปิ ตํทิฏฺฐิกา โหติ ตํขนฺติกา ตํรุจิกา.}

\prose{\textbf{สา ภิกฺขุนี}ติ ยา สา อุกฺขิตฺตานุวตฺติกา ภิกฺขุนี.}

\prose{\textbf{ภิกฺขุนีหี}ติ อญฺญาหิ ภิกฺขุนีหิ.\csromanspacer{} ยา ปสฺสนฺติ ยา สุณนฺติ, ตาหิ วตฺตพฺพา “เอโส โข อยฺเย ภิกฺขุ สมคฺเคน สํเฆน อุกฺขิตฺโต ธมฺเมน วินเยน สตฺถุสาสเนน อนาทโร อปฺปฏิกาโร อกตสหาโย, มายฺเย เอตํ ภิกฺขุํ อนุวตฺตี”ติ.\csromanspacer{} ทุติยมฺปิ วตฺตพฺพา.\csromanspacer{} ตติยมฺปิ วตฺตพฺพา.\csromanspacer{} สเจ ปฏินิสฺสชฺชติ, อิจฺเจตํ กุสลํ, โน เจ ปฏินิสฺสชฺชติ, อาปตฺติ ทุกฺกฏสฺส.\csromanspacer{} สุตฺวา น วทนฺติ, อาปตฺติ ทุกฺกฏสฺส.\csromanspacer{} สา ภิกฺขุนี สํฆมชฺฌมฺปิ อากฑฺฒิตฺวา วตฺตพฺพา “เอโส โข อยฺเย ภิกฺขุ สมคฺเคน สํเฆน อุกฺขิตฺโต ธมฺเมน วินเยน สตฺถุสาสเนน อนาทโร อปฺปฏิกาโร อกตสหาโย, มายฺเย เอตํ ภิกฺขุํ อนุวตฺตี”ติ.\csromanspacer{} ทุติยมฺปิ วตฺตพฺพา.\csromanspacer{} ตติยมฺปิ วตฺตพฺพา.\csromanspacer{} สเจ ปฏินิสฺสชฺชติ, อิจฺเจตํ กุสลํ, โน เจ ปฏินิสฺสชฺชติ, อาปตฺติ ทุกฺกฏสฺส.\csromanspacer{} สา ภิกฺขุนี สมนุภาสิตพฺพา.\csromanspacer{} เอวญฺจ ปน ภิกฺขเว สมนุภาสิตพฺพา, พฺยตฺตาย ภิกฺขุนิยา ปฏิพลาย สํโฆ ญาเปตพฺโพ\csromandash{}}

\proseitem{671}{“สุณาตุ เม อยฺเย สํโฆ, อยํ อิตฺถนฺนามา ภิกฺขุนี สมคฺเคน สํเฆน อุกฺขิตฺตํ ภิกฺขุํ ธมฺเมน วินเยน สตฺถุสาสเนน อนาทรํ อปฺปฏิการํ อกตสหายํ ตมนุวตฺตติ, สา ตํ วตฺถุํ น ปฏินิสฺสชฺชติ, ยทิ สํฆสฺส ปตฺตกลฺลํ, สํโฆ อิตฺถนฺนามํ ภิกฺขุนิํ สมนุภาเสยฺย ตสฺส วตฺถุสฺส ปฏินิสฺสคฺคาย, เอสา ญาตฺติ.}

\prose{สุณาตุ เม อยฺเย สํโฆ, อยํ อิตฺถนฺนามา ภิกฺขุนี สมคฺเคน สํเฆน อุกฺขิตฺตํ ภิกฺขุํ ธมฺเมน วินเยน สตฺถุสาสเนน อนาทรํ อปฺปฏิการํ อกตสหายํ ตมนุวตฺตติ, สา ตํ วตฺถุํ น ปฏินิสฺสชฺชติ, สํโฆ อิตฺถนฺนามํ ภิกฺขุนิํ สมนุภาสติ ตสฺส วตฺถุสฺส ปฏินิสฺสคฺคาย ยสฺสา อยฺยาย ขมติ อิตฺถนฺนามาย ภิกฺขุนิยา สมนุภาสนา ตสฺส วตฺถุสฺส ปฏินิสฺสคฺคาย, สา ตุณฺหสฺส, ยสฺสา นกฺขมติ, สา ภาเสยฺย.}

\prose{ทุติยมฺปิ เอตมตฺถํ วทามิ ฯเปฯ.\csromanspacer{} ตติยมฺปิ เอตมตฺถํ วทามิ ฯเปฯ.}

\prose{สมนุภฏฺฐา สํเฆน อิตฺถนฺนามา ภิกฺขุนี ตสฺส วตฺถุสฺส ปฏินิสฺสคฺคาย, ขมติ สํฆสฺส, ตสฺมา ตุณฺหี, เอวเมตํ ธารยามี”ติ.}

\prose{\csromanfolio{284}ญาตฺติยา ทุกฺกฏํ, ทฺวีหิ กมฺมวาจาติ ถุลฺลจฺจยา, กมฺมวาจา ปริโยสาเน อาปตฺติ ปาราชิกสฺส.}

\prose{\textbf{อยมฺปิ}ติ ปุริมาโย อุปาทาย วุจฺจติ.}

\prose{\textbf{ปาราชิกา โหตี}ติ เสยฺยถาปิ นาม ปุถุสิลา เทฺวธา ภินฺนา อปฺปฏิสนฺธิกา โหติ, เอวเมว ภิกฺขุนี ยาวตติยํ สมนุภาสนาย น ปฏินิสฺสชฺชนฺตี อสฺสมณี โหติ อสกฺยธีตา, เตน วุจฺจติ ปาราชิกา โหตีติ.}

\prose{\textbf{อสํวาสา}ติ \textbf{สํวาโส} นาม เอกกมฺมํ เอกุทฺเทโส สมสิกฺขตา, เอโส สํวาโส นาม.\csromanspacer{} โส ตาย สทฺธิํ นตฺถิ, เตน วุจฺจติ อสํวาสาติ.}

\proseitem{672}{ธมฺมกมฺเม ธมฺมกมฺมสญฺญา น ปฏินิสฺสชฺชติ, อาปตฺติ ปาราชิกสฺส.\csromanspacer{} ธมฺมกมฺเม เวมติกา น ปฏินิสฺสชฺชติ, อาปตฺติ ปาราชิกสฺส.\csromanspacer{} ธมฺมกมฺเม อธมฺมกมฺมสญฺญา น ปฏินิสฺสชฺชติ, อาปตฺติ ปาราชิกสฺส.}

\prose{อธมฺมกมฺเม ธมฺมกมฺมสญฺญา, อาปตฺติ ทุกฺกฏสฺส.\csromanspacer{} อธมฺมกมฺเม เวมติกา, อาปตฺติ ทุกฺกฏสฺส.\csromanspacer{} อธมฺมกมฺเม อธมฺมกมฺมสญฺญา, อาปตฺติ ทุกฺกฏสฺส.}

\proseitemwithcloserruled{673}{อนาปตฺติ อสมนุภาสนฺติยา ปฏินิสฺสชฺชนฺติยา อุมฺมตฺติกาย ฯเปฯ อาทิกมฺมิกายาติ.}{ตติยปาราชิกํ สมตฺตํ.}
\csromanmatikaanchor{mtk.122}
\csromantocmarkref{section}{4. จตุตฺถปาราชิกสิกฺขาปท}{mtk.122}
\bookmark[dest={mtk.122},level=1]{4. จตุตฺถปาราชิกสิกฺขาปท}
\setcsromanheadh{4. จตุตฺถปาราชิกสิกฺขาปท}
\csromanheader{1}{4. \textbf{จตุตฺถปาราชิก}}
\proseitem{674}{เตน สมเยน พุทฺโธ ภควา สาวตฺถิยํ วิหรติ เชตวเน อนาถปิณฺฑิกสฺส อาราเม.\csromanspacer{} เตน โข ปน สมเยน ฉพฺพคฺคิยา ภิกฺขุนิโย อวสฺสุตา อวสฺสุตสฺส ปุริสปุคฺคลสฺส หตฺถคฺคหณมฺปิ สาทิยนฺติ, สงฺฆาฏิกณฺณคฺคหณมฺปิ สาทิยนฺติ, สนฺติฏฺฐนฺติปิ, สลฺลปนฺติปิ, สงฺเกตมฺปิ คจฺฉนฺติ, ปุริสสฺสปิ อพฺภาคมนํ สาทิยนฺติ, ฉนฺนมฺปิ อนุปวิสนฺติ, กายมฺปิ ตทตฺถาย อุปสํหรนฺติ เอตสฺส อสทฺธมฺมสฺส ปฏิเสวนตฺถาย.\csromanspacer{} ยา ตา ภิกฺขุนิโย อปฺปิจฺฉา ฯเปฯ ตา อุชฺฌายนฺติ ขิยฺยนฺติ วิปาเจนฺติ “กถํ \csromanfolio{285}หิ นาม ฉพฺพคฺคิยา ภิกฺขุนิโย อวสฺสุตา อวสฺสุตสฺส ปุริสปุคฺคลสฺส หตฺถคฺคหณมฺปิ สาทิยิสฺสนฺติ, สงฺฆาฏิกณฺณคฺคหณมฺปิ สาทิยิสฺสนฺติ, สนฺติฏฺฐิสฺสนฺติปิ, สลฺลปิสฺสนฺติปิ, สงฺเกตมฺปิ คจฺฉิสฺสนฺติ, ปุริสสฺสปิ \textbf{อพฺภาคมนํ สาทิยิสฺสนฺติ, ฉนฺนมฺปิ อนุปวิสิสฺสนฺติ, กายมฺปิ ตทตฺถาย} \textbf{อุปสํหริสฺสนฺติ เอตสฺส อสทฺธมฺมสฺส ปฏิเสวนตฺถายา”ติ ฯเปฯ.}\csromanspacer{}\textbf{ สจฺจํ} กิร ภิกฺขเว ฉพฺพคฺคิยา ภิกฺขุนิโย อวสฺสุตา อวสฺสุตสฺส ปุริสปุคฺคลสฺส หตฺถคฺคหณมฺปิ สาทิยนฺติ, สงฺฆาฏิกณฺณคฺคหณมฺปิ สาทิยนฺติ, สนฺติฏฺฐนฺติปิ, สลฺลปนฺติปิ, สงฺเกตมฺปิ คจฺฉนฺติ, ปุริสสฺสปิ อพฺภาคมนํ สาทิยนฺติ, ฉนฺนมฺปิ อนุปวิสนฺติ, กายมฺปิ ตทตฺถาย อุปสํหรนฺติ เอตสฺส อสทฺธมฺมสฺส ปฏิเสวนตฺถายาติ.\csromanspacer{} สจฺจํ ภควาติ.\csromanspacer{} วิครหิ พุทฺโธ ภควา ฯเปฯ.\csromanspacer{} กถํ หิ นาม ภิกฺขเว ฉพฺพคฺคิยา ภิกฺขุนิโย อวสฺสุตา อวสฺสุตสฺส ปุริสปุคฺคลสฺส หตฺถคฺคหณมฺปิ สาทิยิสฺสนฺติ, สงฺฆาฏิกณฺณคฺคหณมฺปิ สาทิยิสฺสนฺติ, สนฺติฏฺฐิสฺสนฺติปิ, สลฺลปิสฺสนฺติปิ, สงฺเกตมฺปิ คจฺฉิสฺสนฺติ, ปุริสสฺสปิ อพฺภาคมนํ สาทิยิสฺสนฺติ, ฉนฺนมฺปิ อนุปวิสิสฺสนฺติ, กายมฺปิ ตทตฺถาย อุปสํหริสฺสนฺติ เอตสฺส อสทฺธมฺมสฺส ปฏิเสวนตฺถาย.\csromanspacer{} เนตํ ภิกฺขเว \textbf{อปฺปสนฺนานํ วา ปสาทาย ฯเปฯ.}\csromanspacer{}\textbf{ เอวญฺจ ปน ภิกฺขเว ภิกฺขุนิโย อิมํ} \textbf{สิกฺขาปทํ อุทฺทิสนฺตุ\csromandash{}}}

\proseitem{675}{\textbf{“ยา ปน ภิกฺขุนี อวสฺสุตา อวสฺสุตสฺส ปุริสปุคฺคลสฺส หตฺถคฺคหณํ วา สาทิเยยฺย สงฺฆาฏิกณฺณคฺคหณํ วา สาทิเยยฺย สนฺติฏฺเฐยฺย วา สลฺลเปยฺย วา สงฺเกตํ วา คจฺเฉยฺย ปุริสสฺส วา อพฺภาคมนํ สาทิเยยฺย ฉนฺนํ วา อนุปวิเสยฺย กายํ วา ตทตฺถาย อุปสํหเรยฺย เอตสฺส อสทฺธมฺมสฺส ปฏิเสวนตฺถาย, อยมฺปิ ปาราชิกา โหติ อสํวาสา อฏฺฐวตฺถุกา”}ติ.}

\proseitem{676}{ยา ปนาติ ยา ยาทิสา ฯเปฯ.\csromanspacer{} \textbf{ภิกฺขุนี}ติ ฯเปฯ อยํ อิมสฺมิํ อตฺเถ อธิปฺเปตา ภิกฺขุนีติ.}

\prose{\textbf{อวสฺสุตา} นาม สารตฺตา อเปกฺขวตี ปฏิพทฺธจิตฺตา.}

\prose{\textbf{อวสฺสุโต} นาม สารตฺโต อเปกฺขวา ปฏิพทฺธจิตฺโต.}

\prose{\textbf{ปุริสปุคฺคโล} นาม มนุสฺสปุริโส น ยกฺโข น เปโต น ติรจฺฉานคโต วิญฺญู ปฏิพโล กายสํสคฺคํ สมาปชฺชิตุํ.}

\prose{\csromanfolio{286}\textbf{หตฺถคฺคหณํ วา สาทิเยยฺยา}ติ หตฺโถ นาม กปฺปรํ อุปาทาย ยาว อคฺคนขา.\csromanspacer{} เอตสฺส อสทฺธมฺมสฺส ปฏิเสวนตฺถาย อุพฺภกฺขกํ อโธชาณุมณฺฑลํ คหณํ สาทิยติ, อาปตฺติ ถุลฺลจฺจยสฺส.}

\prose{\textbf{สงฺฆาฏิกณฺณคฺคหณํ วา สาทิเยยฺยา}ติ เอตสฺส อสทฺธมฺมสฺส ปฏิเสวนตฺถาย นิวตฺถํ วา ปารุตํ วา คหณํ สาทิยติ, อาปตฺติ ถุลฺลจฺจยสฺส.}

\prose{\textbf{สนฺติฏฺเฐยฺย วา}ติ เอตสฺส อสทฺธมฺมสฺส ปฏิเสวนตฺถาย ปุริสสฺส หตฺถปาเส ติฏฺฐติ, อาปตฺติ ถุลฺลจฺจยสฺส.}

\prose{\textbf{สลฺลเปยฺย วา}ติ เอตสฺส อสทฺธมฺมสฺส ปฏิเสวนตฺถาย ปุริสสฺส หตฺถปาเส ฐิตา สลฺลปติ, อาปตฺติ ถุลฺลจฺจยสฺส.}

\prose{\textbf{สงฺเกตํ วา คจฺเฉยฺยา}ติ เอตสฺส อสทฺธมฺมสฺส ปฏิเสวนตฺถาย ปุริเสน “อิตฺถนฺนามํ โอกาสํ\footnote{อิทํ ปทํ อฏฺฐกถายํ น ทิสฺสติ.} อาคจฺฉา”ติ วุตฺตา คจฺฉติ, ปเท ปเท อาปตฺติ ทุกฺกฏสฺส.\csromanspacer{} ปุริสสฺส หตฺถปาสํ โอกฺกนฺตมตฺเต อาปตฺติ ถุลฺลจฺจยสฺส.}

\prose{\textbf{ปุริสสฺส วา อพฺภาคมนํ สาทิเยยฺยา}ติ เอตสฺส อสทฺธมฺมสฺส ปฏิเสวนตฺถาย ปุริสสฺส อพฺภาคมนํ สาทิยติ, อาปตฺติ ทุกฺกฏสฺส.\csromanspacer{} หตฺถปาสํ โอกฺกนฺตมตฺเต อาปตฺติ ถุลฺลจฺจยสฺส.}

\prose{\textbf{ฉนฺนํ วา อนุปวิเสยฺยา}ติ เอตสฺส อสทฺธมฺมสฺส ปฏิเสวนตฺถาย เยน เกนจิ ปฏิจฺฉนฺนํ โอกาสํ ปวิฏฺฐมตฺเต อาปตฺติ ถุลฺลจฺจยสฺส.}

\prose{\textbf{กายํ วา ตทตฺถาย อุปสํหเรยฺยา}ติ เอตสฺส อสทฺธมฺมสฺส ปฏิเสวนตฺถาย ปุริสสฺส หตฺถปาเส ฐิตา กายํ อุปสํหรติ, อาปตฺติ ถุลฺลจฺจยสฺส.}

\prose{\textbf{อยมฺปี}ติ ปุริมาโย อุปาทาย วุจฺจติ.}

\prose{\textbf{ปาราชิกา โหตี}ติ เสยฺยถาปิ นาม ตาโล มตฺถกจฺฉินฺโน อภพฺโพ ปุน วิรุฬฺหิยา.\csromanspacer{} เอวเมว ภิกฺขุนี อฏฺฐมํ วตฺถุํ ปริปูเรนฺตี อสฺสมณี โหติ อสกฺยธีตา.\csromanspacer{} เตน วุจฺจติ ปาราชิกา โหตีติ.}

\prose{\csromanfolio{287}\textbf{อสํวาสา}ติ \textbf{สํวาโส} นาม เอกกมฺมํ เอกุทฺเทโส สมสิกฺขตา, เอโส สํวาโส นาม.\csromanspacer{} โส ตาย สทฺธิํ นตฺถิ, เตน วุจฺจติ อสํวาสาติ.}

\proseitemwithcloser{677}{อนาปตฺติ อสญฺจิจฺจ อสฺสติยา อชานนฺติยา อสาทิยนฺติยา อุมฺมตฺติกาย ขิตฺตจิตฺตาย เวทนาฏฺฏาย อาทิกมฺมิกายาติ.}{จตุตฺถปาราชิกํ สมตฺตํ.}
\prose{อุทฺทิฏฺฐา โข อยฺยาโย อฏฺฐ ปาราชิกา ธมฺมา.\csromanspacer{} เยสํ ภิกฺขุนี อญฺญตรํ วา อญฺญตรํ วา อาปชฺชิตฺวา น ลภติ ภิกฺขุนีหิ สทฺธิํ สํวาสํ ยถา ปุเร ตถา ปจฺฉา, ปาราชิกา โหติ อสํวาสา.\csromanspacer{} ตตฺถายฺยาโย ปุจฺฉามิ กจฺจิตฺถ ปริสุทฺธา, ทุติยมฺปิ ปุจฺฉามิ กจฺจิตฺถ ปริสุทฺธา, ตติยมฺปิ ปุจฺฉมิ กจฺจิตฺถ ปริสุทฺธา, ปริสุทฺเธตฺถายฺยาโย, ตสฺมา ตุณฺหี, เอวเมตํ ธารยามีติ.}

\csromanheader{2}{\textbf{ภิกฺขุนิวิภงฺเค}}
\nitthitam{\textbf{ปาราชิกกณฺฑํ นิฏฺฐิตํ.}}
\csromanmatikaanchor{mtk.123}
\csromantocmarknopage{chapter}{2. สํฆาทิเสสกณฺฑ}{mtk.123}
\bookmark[dest={mtk.123},level=0]{2. สํฆาทิเสสกณฺฑ}
\setcsromanheadcha{2. สํฆาทิเสสกณฺฑ}
\setcsromanheadhclear{}
\chapterheadpageread{2. \textbf{สํฆาทิเสสกณฺฑ}}
\csromanmatikaanchor{mtk.124}
\csromantocmarkref{section}{1. ปฐมสํฆาทิเสสสิกฺขาปท}{mtk.124}
\bookmark[dest={mtk.124},level=1]{1. ปฐมสํฆาทิเสสสิกฺขาปท}
\setcsromanheadh{1. ปฐมสํฆาทิเสสสิกฺขาปท}
\csromanheader{1}{1. \textbf{ปฐมสํฆาทิเสสสิกฺขาปท}}
\prose{\csromanfolio{288}อิเม โข ปนายฺยาโย สตฺตรส สํฆาทิเสสา ธมฺมา อุทฺเทสํ อาคจฺฉนฺติ.}

\proseitem{678}{เตน สมเยน พุทฺโธ ภควา สาวตฺถิยํ วิหรติ เชตวเน อนาถปิณฺฑิกสฺส อาราเม.\csromanspacer{} เตน โข ปน สมเยน อญฺญตโร อุปาสโก ภิกฺขุนิสํฆสฺส อุโทสิตํ\footnote{อุทฺโทสิตํ (สี, สฺยา)} ทตฺวา กาลงฺกโต โหติ, ตสฺส เทฺว ปุตฺตา โหนฺติ, เอโก อสฺสทฺโธ อปฺปสนฺโน, เอโก สทฺโธ ปสนฺโน, เต เปตฺติกํ สาปเตยฺยํ วิภชิํสุ.\csromanspacer{} อถ โข โส อสฺสทฺโธ อปฺปสนฺโน ตํ สทฺธํ ปสนฺนํ เอตทโวจ “อมฺหากํ อุโทสิโต, ตํ ภาเชมา”ติ.\csromanspacer{} เอวํ วุตฺเต โส สทฺโธ ปสนฺโน ตํ อสฺสทฺธํ อปฺปสนฺนํ เอตทโวจ “มายฺโย เอวํ อวจ, อมฺหากํ ปิตุนา ภิกฺขุนิสํฆสฺส ทินฺโน”ติ.\csromanspacer{} ทุติยมฺปิ โข โส อสฺสทฺโธ อปฺปสนฺโน ตํ สทฺธํ ปสนฺนํ เอตทโวจ “อมฺหากํ อุโทสิโต, ตํ ภาเชมา”ติ.\csromanspacer{} อถ โข โส สทฺโธ ปสนฺโน ตํ อสฺสทฺธํ อปฺปสนฺนํ เอตทโวจ “มายฺโย เอวํ อวจ, อมฺหากํ ปิตุนา ภิกฺขุนิสํฆสฺส ทินฺโน”ติ.\csromanspacer{} ตติยมฺปิ โข โส อสฺสทฺโธ อปฺปสนฺโน ตํ สทฺธํ ปสนฺนํ เอตทโวจ “อมฺหากํ อุโทสิโต, ตํ ภาเชมา”ติ.\csromanspacer{} อถ โข โส สทฺโธ ปสนฺโน “สเจ มยฺหํ ภวิสฺสติ, อหมฺปิ ภิกฺขุนิสํฆสฺส ทสฺสามี”ติ ตํ อสฺสทฺธํ อปฺปสนฺนํ เอตทโวจ “ภาเชมา”ติ.\csromanspacer{} อถ โข โส อุโทสิโต เตหิ ภาชียมาโน ตสฺส อสฺสทฺธสฺส อปฺปสนฺนสฺส ปาปุณาติ\footnote{ปาปุณิ (สฺยา)}.\csromanspacer{} อถ โข โส อสฺสทฺโธ อปฺปสนฺโน ภิกฺขุนิโย อุปสงฺกมิตฺวา เอตทโวจ “นิกฺขมถายฺเย อมฺหากํ อุโทสิโต”ติ.}

\prose{เอวํ วุตฺเต ถุลฺลนนฺทา ภิกฺขุนี ตํ ปุริสํ เอตทโวจ “มายฺโย เอวํ อวจ, ตุมฺหากํ ปิตุนา ภิกฺขุนิสํฆสฺส ทินฺโน”ติ.\csromanspacer{} ทินฺโน น ทินฺโนติ โวหาริเก มหามตฺเต ปุจฺฉิํสุ.\csromanspacer{} มหามตฺตา เอวมาหํสุ “โก อยฺเย ชานาติ ภิกฺขุนิสํฆสฺส ทินฺโน”ติ.\csromanspacer{} เอวํ วุตฺเต ถุลฺลนนฺทา ภิกฺขุนี เต มหามตฺเต เอตทโวจ “อปิ นายฺยา\footnote{อปิ นฺวยฺยา (สฺยา), อปิ นายฺโย (ก)} ตุเมฺหหิ ทิฏฺฐํ วา สุตํ วา สกฺขิํ ฐปยิตฺวา ทานํ \csromanfolio{289}ทิยฺยมานนฺ”ติ.\csromanspacer{} อถ โข เต มหามตฺตา “สจฺจํ โข อยฺยา อาหา”ติ ตํ อุโทสิตํ ภิกฺขุนิสํฆสฺส อกํสุ.\csromanspacer{} อถ โข โส ปุริโส ปราชิโต อุชฺฌายติ ขิยฺยติ วิปาเจติ “อสฺสมณิโย อิมา มุณฺฑา พนฺธกินิโย, กถํ หิ นาม อมฺหากํ อุโทสิตํ อจฺฉินฺทาเปสฺสนฺตี”ติ.\csromanspacer{} ถุลฺลนนฺทา ภิกฺขุนี มหามตฺตานํ เอตมตฺถํ อาโรเจสิ.\csromanspacer{} มหามตฺตา ตํ ปุริสํ ทณฺฑาเปสุํ.\csromanspacer{} อถ โข โส ปุริโส ทณฺฑิโต ภิกฺขุนุปสฺสยสฺส อวิทูเร อาชีวกเสยฺยํ การาเปตฺวา อาชีวเก อุยฺโยเชสิ “เอตา ภิกฺขุนิโย อจฺจาวทถา”ติ.}

\prose{ถุลฺลนนฺทา ภิกฺขุนี มหามตฺตานํ เอตมตฺถํ อาโรเจสิ.\csromanspacer{} มหามตฺตา ตํ ปุริสํ พนฺธาเปสุํ.\csromanspacer{} มนุสฺสา อุชฺฌายนฺติ ขิยฺยนฺติ วิปาเจนฺติ “ปฐมํ ภิกฺขุนิโย อุโทสิตํ อจฺฉินฺทาเปสุํ, ทุติยํ ทณฺฑาเปสุํ, ตติยํ พนฺธาเปสุํ, อิทานิ ฆาตาเปสฺสนฺตี”ติ.\csromanspacer{} อสฺโสสุํ โข ภิกฺขุนิโย เตสํ มนุสฺสานํ อุชฺฌายนฺตานํ ขิยฺยนฺตานํ วิปาเจนฺตานํ.\csromanspacer{} ยา ตา ภิกฺขุนิโย อปฺปิจฺฉา ฯเปฯ ตา อุชฺฌายนฺติ ขิยฺยนฺติ วิปาเจนฺติ “กถํ หิ นาม อยฺยา ถุลฺลนนฺทา อุสฺสยวาทิกา วิหริสฺสตี”ติ.\csromanspacer{} อถ โข ตา ภิกฺขุนิโย ภิกฺขูนํ เอตมตฺถํ อาโรเจสุํ ฯเปฯ.\csromanspacer{} สจฺจํ กิร ภิกฺขเว ถุลฺลนนฺทา ภิกฺขุนี อุสฺสยวาทิกา วิหรตีติ.\csromanspacer{} สจฺจํ ภควาติ.\csromanspacer{} วิครหิ พุทฺโธ ภควา ฯเปฯ.\csromanspacer{} กถํ หิ นาม ภิกฺขเว ถุลฺลนนฺทา ภิกฺขุนี อุสฺสยวาทิกา วิหริสฺสติ.\csromanspacer{} เนตํ ภิกฺขเว อปฺปสนฺนานํ วา ปสาทาย ฯเปฯ.\csromanspacer{} เอวญฺจ ปน ภิกฺขเว ภิกฺขุนิโย อิมํ สิกฺขาปทํ อุทฺทิสนฺตุ\csromandash{}}

\proseitem{679}{\textbf{“ยา ปน ภิกฺขุนี อุสฺสยวาทิกา วิหเรยฺย คหปตินา วา คหปติปุตฺเตน วา ทาเสน วา กมฺมกาเรน}\footnote{กมฺมกเรน (สี, สฺยา)} \textbf{วา อนฺตมโส สมณปริพฺพาชเกนาปิ, อยํ ภิกฺขุนี ปฐมาปตฺติกํ ธมฺมํ อาปนฺนา นิสฺสารณียํ สํฆาทิเสสนฺ”}ติ.}

\proseitem{680}{ยา ปนาติ ยา ยาทิสา ฯเปฯ.\csromanspacer{} \textbf{ภิกฺขุนี}ติ ฯเปฯ อยํ อิมสฺมิํ อตฺเถ อธิปฺเปตา ภิกฺขุนีติ.}

\prose{\textbf{อุสฺสยวาทิกา} นาม อฑฺฑการิกา วุจฺจติ.}

\prose{\textbf{คหปติ} นาม โย โกจิ อคารํ อชฺฌาวสติ.}

\prose{\csromanfolio{290}\textbf{คหปติปุตฺโต} นาม เย เกจิ ปุตฺตภาตโร.}

\prose{\textbf{ทาโส} นาม อนฺโตชาโต ธนกฺกีโต กรมรานีโต.}

\prose{\textbf{กมฺมกาโร} นาม ภฏโก อาหตโก.}

\prose{\textbf{สมณปริพฺพาชโก} นาม ภิกฺขุญฺจ ภิกฺขุนิญฺจ สิกฺขมานญฺจ สามเณรญฺจ สามเณริญฺจ ฐเปตฺวา โย โกจิ ปริพฺพาชกสมาปนฺโน.}

\prose{“อฑฺฑํ กริสฺสามี”ติ ทุติยํ วา ปริเยสติ คจฺฉติ วา, อาปตฺติ ทุกฺกฏสฺส.\csromanspacer{} เอกสฺส อาโรเจติ, อาปตฺติ ทุกฺกฏสฺส.\csromanspacer{} ทุติยสฺส อาโรเจติ, อาปตฺติ ถุลฺลจฺจยสฺส.\csromanspacer{} อฑฺฑปริโยสาเน อาปตฺติ สํฆาทิเสสสฺส.}

\prose{\textbf{ปฐมาปตฺติกนฺ}ติ สห วตฺถุชฺฌาจารา อาปชฺชติ อสมนุภาสนาย.}

\prose{\textbf{นิสฺสารณียนฺ}ติ สํฆมฺหา นิสฺสรียติ.}

\prose{\textbf{สํฆาทิเสสนฺ}ติ สํโฆว ตสฺสา อาปตฺติยา มานตฺตํ เทติ มูลาย ปฏิกสฺสติ อพฺเภติ, น สมฺพหุลา, น เอกา ภิกฺขุนี, เตน วุจฺจติ “สํฆาทิเสโส”ติ.\csromanspacer{} ตสฺเสว อาปตฺตินิกายสฺส นามกมฺมํ อธิวจนํ, เตนปิ วุจฺจติ “สํฆาทิเสโส”ติ.}

\proseitemwithcloserruled{681}{อนาปตฺติ มนุสฺเสหิ อากฑฺฒียมานา คจฺฉติ, อารกฺขํ ยาจติ, อโนทิสฺส อาจิกฺขติ, อุมฺมตฺติกาย ฯเปฯ อาทิกมฺมิกายาติ.}{ปฐมสํฆาทิเสสสิกฺขาปทํ นิฏฺฐิตํ.}
\csromanmatikaanchor{mtk.125}
\csromantocmarkref{section}{2. ทุติยสํฆาทิเสสสิกฺขาปท}{mtk.125}
\bookmark[dest={mtk.125},level=1]{2. ทุติยสํฆาทิเสสสิกฺขาปท}
\setcsromanheadh{2. ทุติยสํฆาทิเสสสิกฺขาปท}
\csromanheader{1}{2. \textbf{ทุติยสํฆาทิเสสสิกฺขาปท}}
\proseitem{682}{เตน สมเยน พุทฺโธ ภควา สาวตฺถิยํ วิหรติ เชตวเน อนาถปิณฺฑิกสฺส อาราเม.\csromanspacer{} เตน โข ปน สมเยน เวสาลิยํ อญฺญตรสฺส ลิจฺฉวิสฺส ปชาปติ อติจารินี โหติ.\csromanspacer{} อถ โข โส ลิจฺฉวิ ตํ อิตฺถิํ เอตทโวจ “สาธุ วิรมาหิ, อนตฺถํ โข เต กริสฺสามี”ติ.\csromanspacer{} เอวมฺปิ วุจฺจมานา นาทิยิ.\csromanspacer{} เตน โข ปน สมเยน เวสาลิยํ ลิจฺฉวิคโณ สนฺนิปติโต โหติ เกนจิเทว กรณีเยน.\csromanspacer{} อถ โข โส ลิจฺฉวิ เต ลิจฺฉวโย เอตทโวจ “เอกํ เม \csromanfolio{291}อยฺโย อิตฺถิํ อนุชานาถา”ติ.\csromanspacer{} กา นาม สาติ.\csromanspacer{} มยฺหํ ปชาปติ อติจรติ, ตํ ฆาเตสฺสามีติ.\csromanspacer{} ชานาหีติ.\csromanspacer{} อสฺโสสิ โข สา อิตฺถี “สามิโก กิร มํ ฆาเตตุกาโม”ติ วรภณฺฑํ อาทาย สาวตฺถิํ คนฺตฺวา ติตฺถิเย อุปสงฺกมิตฺวา ปพฺพชฺชํ ยาจิ.\csromanspacer{} ติตฺถิยา น อิจฺฉิํสุ ปพฺพาเชตุํ.\csromanspacer{} ภิกฺขุนิโย อุปสงฺกมิตฺวา ปพฺพชฺชํ ยาจิ.\csromanspacer{} ภิกฺขุนิโยปิ น อิจฺฉิํสุ ปพฺพาเชตุํ.\csromanspacer{} ถุลฺลนนฺทํ ภิกฺขุนิํ อุปสงฺกมิตฺวา ภณฺฑกํ ทสฺเสตฺวา ปพฺพชฺชํ ยาจิ.\csromanspacer{} ถุลฺลนนฺทา ภิกฺขุนี ภณฺฑกํ คเหตฺวา ปพฺพาเชสิ.}

\prose{อถ โข โส ลิจฺฉวิ ตํ อิตฺถิํ คเวสนฺโต สาวตฺถิํ คนฺตฺวา ภิกฺขุนีสุ ปพฺพชิตํ ทิสฺวาน เยน ราชา ปเสนทิ โกสโล เตนุปสงฺกมิ, อุปสงฺกมิตฺวา ราชานํ ปเสนทิํ โกสลํ เอตทโวจ “ปชาปติ เม เทว วรภณฺฑํ อาทาย สาวตฺถิํ อนุปฺปตฺตา, ตํ เทโว อนุชานาตู”ติ.\csromanspacer{} เตน หิ ภเณ วิจินิตฺวา อาจิกฺขาติ.\csromanspacer{} ทิฏฺฐา เทว ภิกฺขุนีสุ ปพฺพชิตาติ.\csromanspacer{} สเจ ภเณ ภิกฺขุนีสุ ปพฺพชิตา, น สา ลพฺภา กิญฺจิ กาตุํ, สฺวากฺขาโต ภควตา ธมฺโม, จรตุ พฺรหฺมจริยํ สมฺมา ทุกฺขสฺส อนฺตกิริยายาติ.\csromanspacer{} อถ โข โส ลิจฺฉวิ อุชฺฌายติ ขิยฺยติ วิปาเจติ “กถํ หิ นาม ภิกฺขุนิโย โจริํ ปพฺพาเชสฺสนฺตี”ติ.\csromanspacer{} อสฺโสสุํ โข ภิกฺขุนิโย ตสฺส ลิจฺฉวิสฺส อุชฺฌายนฺตสฺส ขิยฺยนฺตสฺส วิปาเจนฺตสฺส.\csromanspacer{} ยา ตา ภิกฺขุนิโย อปฺปิจฺฉา ฯเปฯ ตา อุชฺฌายนฺติ ขิยฺยนฺติ วิปาเจนฺติ “กถํ หิ นาม อยฺยา ถุลฺลนนฺทา โจริํ ปพฺพาเชสฺสตี”ติ.\csromanspacer{} อถ โข ตา ภิกฺขุนิโย ภิกฺขูนํ เอตมตฺถํ อาโรเจสุํ ฯเปฯ.\csromanspacer{} สจฺจํ กิร ภิกฺขเว ถุลฺลนนฺทา ภิกฺขุนี โจริํ ปพฺพาเชตีติ\footnote{ปพฺพาเชสีติ (ก)}.\csromanspacer{} สจฺจํ ภควาติ.\csromanspacer{} วิครหิ พุทฺโธ ภควา ฯเปฯ.\csromanspacer{} กถํ หิ นาม ภิกฺขเว ถุลฺลนนฺทา ภิกฺขุนี โจริํ ปพฺพาเชสฺสติ.\csromanspacer{} เนตํ ภิกฺขเว อปฺปสนฺนานํ วา ปสาทาย ฯเปฯ.\csromanspacer{} เอวญฺจ ปน ภิกฺขเว ภิกฺขุนิโย อิมํ สิกฺขาปทํ อุทฺทิสนฺตุ\csromandash{}}

\proseitem{683}{\textbf{“ยา ปน ภิกฺขุนี ชานํ โจริํ วชฺฌํ วิทิตํ อนปโลเกตฺวา ราชานํ วา สํฆํ วา คณํ วา ปูคํ วา เสณิํ วา อญฺญตฺร กปฺปา วุฏฺฐาเปยฺย, อยมฺปิ ภิกฺขุนี ปฐมาปตฺติกํ ธมฺมํ อาปนฺนา นิสฺสารณียํ สํฆาทิเสสนฺ”}ติ.}

\proseitem{684}{\textbf{ยา ปนา}ติ ยา ยาทิสา ฯเปฯ.\csromanspacer{} \textbf{ภิกฺขุนี}ติ ฯเปฯ อยํ อิมสฺมิํ อตฺเถ อธิปฺเปตา ภิกฺขุนีติ.}

\prose{\csromanfolio{292}\textbf{ชานาติ} นาม สามํ วา ชานาติ, อญฺเญ วา ตสฺสา อาโรเจนฺติ, สา วา อาโรเจติ.}

\prose{\textbf{โจรี} นาม ยา ปญฺจมาสกํ วา อติเรกปญฺจมาสกํ วา อคฺฆนกํ อทินฺนํ เถยฺยสงฺขาตํ อาทิยติ, เอสา โจรี นาม.}

\prose{\textbf{วชฺฌา} นาม ยํ กตฺวา วชฺฌปฺปตฺตา โหติ.}

\prose{\textbf{วิทิตา} นาม อญฺเญหิ มนุสฺเสหิ ญาตา โหติ “วชฺฌา เอสา”ติ.}

\prose{\textbf{อนปโลเกตฺวา}ติ อนาปุจฺฉา.}

\prose{\textbf{ราชา} นาม ยตฺถ ราชา อนุสาสติ, ราชา อปโลเกตพฺโพ.}

\prose{\textbf{สํโฆ} นาม ภิกฺขุนิสํโฆ วุจฺจติ, ภิกฺขุนิสํโฆ อปโลเกตพฺโพ.}

\prose{\textbf{คโณ} นาม ยตฺถ คโณ อนุสาสติ, คโณ อปโลเกตพฺโพ.}

\prose{\textbf{ปูโค} นาม ยตฺถ ปูโค อนุสาสติ, ปูโค อปโลเกตพฺโพ.}

\prose{\textbf{เสณิ} นาม ยตฺถ เสณิ อนุสาสติ, เสณิ อปโลเกตพฺโพ.}

\prose{\textbf{อญฺญตฺร กปฺปา}ติ ฐเปตฺวา กปฺปํ.\csromanspacer{} \textbf{กปฺปํ} นาม เทฺว กปฺปานิ ติตฺถิเยสุ วา ปพฺพชิตา โหติ อญฺญาสุ วา ภิกฺขุนีสุ ปพฺพชิตา.\csromanspacer{} อญฺญตฺร กปฺปา “วุฏฺฐาเปสฺสามี”ติ คณํ วา อาจรินิํ วา ปตฺตํ วา จีวรํ วา ปริเยสติ, สีมํ วา สมฺมนฺนติ, อาปตฺติ ทุกฺกฏสฺส.\csromanspacer{} ญตฺติยา ทุกฺกฏํ.\csromanspacer{} ทฺวีหิ กมฺมวาจาหิ ถุลฺลจฺจยา.\csromanspacer{} กมฺมวาจาปริโยสาเน อุปชฺฌายาย อาปตฺติ สํฆาทิเสสสฺส.\csromanspacer{} คณสฺส จ อาจรินิยา จ อาปตฺติ ทุกฺกฏสฺส.}

\prose{\textbf{อยมฺปี}ติ ปุริมํ อุปาทาย วุจฺจติ.}

\prose{\textbf{ปฐมาปตฺติกนฺ}ติ สห วตฺถุชฺฌาจารา อาปชฺชติ อสมนุภาสนาย.}

\prose{\textbf{นิสฺสารณียนฺ}ติ สํฆมฺหา นิสฺสารียติ.}

\prose{\textbf{สํฆาทิเสสนฺ}ติ ฯเปฯ เตนปิ วุจฺจติ สํฆาทิเสโสติ.}

\proseitem{685}{โจริยา โจริสญฺญา อญฺญตฺร กปฺปา วุฏฺฐาเปติ, อาปตฺติ สํฆาทิเสสสฺส.\csromanspacer{} โจริยา เวมติกา อญฺญตฺร กปฺปา วุฏฺฐาเปติ, \csromanfolio{293}อาปตฺติ ทุกฺกฏสฺส.\csromanspacer{} โจริยา อโจริสญฺญา อญฺญตฺร กปฺปา วุฏฺฐาเปติ, อนาปตฺติ.\csromanspacer{} อโจริยา โจริสญฺญา, อาปตฺติ ทุกฺกฏสฺส.\csromanspacer{} อโจริยา เวมติกา, อาปตฺติ ทุกฺกฏสฺส.\csromanspacer{} อโจริยา อโจริสญฺญา, อนาปตฺติ.}

\proseitemwithcloserruled{686}{อนาปตฺติ อชานนฺตี วุฏฺฐาเปติ, อปโลเกตฺวา วุฏฺฐาเปติ, กปฺปกตํ วุฏฺฐาเปติ, อุมฺมตฺติกาย อาทิกมฺมิกายาติ.}{ทุติยสํฆาทิเสสสิกฺขาปทํ นิฏฺฐิตํ.}
\csromanmatikaanchor{mtk.126}
\csromantocmarkref{section}{3. ตติยสํฆาทิเสสสิกฺขาปท}{mtk.126}
\bookmark[dest={mtk.126},level=1]{3. ตติยสํฆาทิเสสสิกฺขาปท}
\setcsromanheadh{3. ตติยสํฆาทิเสสสิกฺขาปท}
\csromanheader{1}{3. \textbf{ตติยสํฆาทิเสสสิกฺขาปท}}
\proseitem{687}{เตน สมเยน พุทฺโธ ภควา สาวตฺถิยํ วิหรติ เชตวเน อนาถปิณฺฑิกสฺส อาราเม.\csromanspacer{} เตน โข ปน สมเยน ภทฺทาย กาปิลานิยา อนฺเตวาสินี ภิกฺขุนี ภิกฺขุนีหิ สทฺธิํ ภณฺฑิตฺวา คามกํ\footnote{คามเก (สฺยา)} ญาติกุลํ อคมาสิ.\csromanspacer{} ภทฺทา กาปิลานี ตํ ภิกฺขุนิํ อปสฺสนฺตี ภิกฺขุนิโย ปุจฺฉิ “กหํ อิตฺถนฺนามา น ทิสฺสตี”ติ.\csromanspacer{} ภิกฺขุนีหิ สทฺธิํ อยฺเย ภณฺฑิตฺวา น ทิสฺสตีติ.\csromanspacer{} อมฺมา อมุกสฺมิํ คามเก เอติสฺสา ญาติกุลํ, ตตฺถ คนฺตฺวา วิจินถาติ.\csromanspacer{} ภิกฺขุนิโย ตตฺถ คนฺตฺวา ตํ ภิกฺขุนิํ ปสฺสิตฺวา เอตทโวจุํ “กิสฺส ตฺวํ อยฺเย เอกิกา อาคตา, กจฺจิสิ อปฺปธํสิตา”ติ.\csromanspacer{} อปฺปธํสิตามฺหิ อยฺเยติ.\csromanspacer{} ยา ตา ภิกฺขุนิโย อปฺปิจฺฉา ฯเปฯ ตา อุชฺฌายนฺติ ขิยฺยนฺติ วิปาเจนฺติ “กถํ หิ นาม ภิกฺขุนี เอกา คามนฺตรํ คจฺฉิสฺสตี”ติ ฯเปฯ.\csromanspacer{} สจฺจํ กิร ภิกฺขเว ภิกฺขุนี เอกา คามนฺตรํ คจฺฉตีติ.\csromanspacer{} สจฺจํ ภควาติ.\csromanspacer{} วิครหิ พุทฺโธ ภควา ฯเปฯ.\csromanspacer{} กถํ หิ นาม ภิกฺขเว ภิกฺขุนี เอกา คามนฺตรํ คจฺฉิสฺสติ.\csromanspacer{} เนตํ ภิกฺขเว อปฺปสนฺนานํ วาปสาทาย ฯเปฯ.\csromanspacer{} เอวญฺจ ปน ภิกฺขเว ภิกฺขุนิโย อิมํ สิกฺขาปทํ อุทฺทิสนฺตุ\csromandash{}}

\prose{\textbf{“ยา ปน ภิกฺขุนี เอกา คามนฺตรํ คจฺเฉยฺย, อยมฺปิ ภิกฺขุนี ปฐมาปตฺติกํ ธมฺมํ อาปนฺนา นิสฺสารณียํ สํฆาทิเสสนฺ”}ติ.}

\prose{เอวญฺจิทํ ภควตา ภิกฺขุนีนํ สิกฺขาปทํ ปญฺญตฺตํ โหติ.}

\proseitem{688}{\csromanfolio{294}เตน โข ปน สมเยน เทฺว ภิกฺขุนิโย สาเกตา สาวตฺถิํ อทฺธานมคฺคปฺปฏิปนฺนา โหนฺติ, อนฺตรามคฺเค นที ตริตพฺพา โหติ.\csromanspacer{} อถ โข ตา ภิกฺขุนิโย นาวิเก อุปสงฺกมิตฺวา เอตทโวจุํ “สาธุ โน อาวุโส ตาเรถา”ติ. “นายฺเย สกฺกา อุโภ สกิํ ตาเรตุนฺ”ติ เอโก เอกํ อุตฺตาเรสิ, อุตฺติณฺโณ อุตฺติณฺณํ ทูเสสิ, อนุตฺติณฺโณ อนุตฺติณฺณํ ทูเสสิ.\csromanspacer{} ตา ปจฺฉา สมาคนฺตฺวา ปุจฺฉิํสุ “กจฺจิสิ อยฺเย อปฺปธํสิตา”ติ.\csromanspacer{} ปธํสิตามฺหิ อยฺเย, ตฺวํ ปน อยฺเย อปฺปธํสิตาติ.\csromanspacer{} ปธํสิตามฺหิ อยฺเยติ.\csromanspacer{} อถ โข ตา ภิกฺขุนิโย สาวตฺถิํ คนฺตฺวา ภิกฺขุนีนํ เอตมตฺถํ อาโรเจสุํ.\csromanspacer{} ยา ตา ภิกฺขุนิโย อปฺปิจฺฉา ฯเปฯ ตา อุชฺฌายนฺติ ขิยฺยนฺติ วิปาเจนฺติ “กถํ หิ นาม ภิกฺขุนี เอกา นทีปารํ คจฺฉิสฺสตี”ติ.\csromanspacer{} อถ โข ตา ภิกฺขุนิโย ภิกฺขูนํ เอตมตฺถํ อาโรเจสุํ.\csromanspacer{} ภิกฺขู ภควโต เอตมตฺถํ อาโรเจสุํ ฯเปฯ.\csromanspacer{} สจฺจํ กิร ภิกฺขเว ภิกฺขุนี เอกา นทีปารํ คจฺฉตีติ.\csromanspacer{} สจฺจํ ภควาติ.\csromanspacer{} วิครหิ พุทฺโธ ภควา ฯเปฯ.\csromanspacer{} กถํ หิ นาม ภิกฺขเว ภิกฺขุนี เอกา นทีปารํ คจฺฉิสฺสติ.\csromanspacer{} เนตํ ภิกฺขเว อปฺปสนฺนานํ วา ปสาทาย ฯเปฯ.\csromanspacer{} เอวญฺจ ปน ภิกฺขเว ภิกฺขุนิโย อิมํ สิกฺขาปทํ อุทฺทิสนฺตุ\csromandash{}}

\prose{\textbf{“ยา ปน ภิกฺขุนี เอกา วา คามนฺตรํ คจฺเฉยฺย, เอกา วา นทีปารํ คจฺเฉยฺย, อยมฺปิ ภิกฺขุนี ปฐมาปตฺติกํ ธมฺมํ อาปนฺนา นิสฺสารณียํ สํฆาทิเสสนฺ”}ติ.}

\prose{เอวญฺจิทํ ภควตา ภิกฺขุนีนํ สิกฺขาปทํ ปญฺญตฺตํ โหติ.}

\proseitem{689}{เตน โข ปน สมเยน สมฺพหุลา ภิกฺขุนิโย โกสเลสุ ชนปเท สาวตฺถิํ คจฺฉนฺตา\footnote{คนฺตฺวา (ก)} สายํ อญฺญตรํ คามํ อุปคจฺฉิํสุ, ตตฺถ อญฺญตรา ภิกฺขุนี อภิรูปา โหติ ทสฺสนียา ปาสาทิกา, อญฺญตโร ปุริโส ตสฺสา ภิกฺขุนิยา สห ทสฺสเนน ปฏิพทฺธจิตฺโต โหติ.\csromanspacer{} อถ โข โส ปุริโส ตาสํ ภิกฺขุนีนํ เสยฺยํ ปญฺญเปนฺโต ตสฺสา ภิกฺขุนิยา เสยฺยํ เอกมนฺตํ ปญฺญาเปสิ.\csromanspacer{} อถ โข สา ภิกฺขุนี สลฺลกฺเขตฺวา “ปริยุฏฺฐิโต อยํ ปุริโส สเจ รตฺติํ อาคจฺฉิสฺสติ วิสฺสโร เม ภวิสฺสตี”ติ ภิกฺขุนิโย อนาปุจฺฉา อญฺญตรํ กุลํ คนฺตฺวา เสยฺยํ กปฺเปสิ.\csromanspacer{} อถ โข โส ปุริโส \csromanfolio{295}\textbf{รตฺติํ อาคนฺตฺวา ตํ ภิกฺขุนิํ คเวสนฺโต ภิกฺขุนิโย ฆฏฺเฏสิ.}\csromanspacer{} \textbf{ภิกฺขุนิโย ตํ ภิกฺขุนิํ อปสฺสนฺติโย เอวมาหํสุ “นิสฺสํสยํ โข} \textbf{สา ภิกฺขุนี ปุริเสน สทฺธิํ นิกฺขนฺตา”ติ.}}

\prose{อถ โข สา ภิกฺขุนี ตสฺสา รตฺติยา อจฺจเยน เยน ตา ภิกฺขุนิโย เตนุปสงฺกมิ.\csromanspacer{} ภิกฺขุนิโย ตํ ภิกฺขุนิํ เอตทโวจุํ “กิสฺส ตฺวํ อยฺเย ปุริเสน สทฺธิํ นิกฺขนฺตา”ติ. “นาหํ อยฺเย ปุริเสน สทฺธิํ นิกฺขนฺตา”ติ ภิกฺขุนีนํ เอตมตฺถํ อาโรเจสิ.\csromanspacer{} ยา ตา ภิกฺขุนิโย อปฺปิจฺฉา ฯเปฯ ตา อุชฺฌายนฺติ ขิยฺยนฺติ วิปาเจนฺติ “กถํ หิ นาม ภิกฺขุนี เอกา รตฺติํ วิปฺปวสิสฺสตี”ติ ฯเปฯ.\csromanspacer{} สจฺจํ กิร ภิกฺขเว ภิกฺขุนี เอกา รตฺติํ วิปฺปวสตีติ\footnote{วิปฺปวสีติ (ก)}.\csromanspacer{} สจฺจํ ภควาติ.\csromanspacer{} วิครหิ พุทฺโธ ภควา ฯเปฯ.\csromanspacer{} กถํ หิ นาม ภิกฺขเว ภิกฺขุนี เอกา รตฺติํ วิปฺปวสิสฺสติ.\csromanspacer{} เนตํ ภิกฺขเว อปฺปสนฺนานํ วา ปสาทาย ฯเปฯ.\csromanspacer{} เอวญฺจ ปน ภิกฺขเว ภิกฺขุนิโย อิมํ สิกฺขาปทํ อุทฺทิสนฺตุ\csromandash{}}

\prose{\textbf{“ยา ปน ภิกฺขุนี เอกา วา คามนฺตรํ คจฺเฉยฺย, เอกา วา นทีปารํ คจฺเฉยฺย, เอกา วา รตฺติํ วิปฺปวเสยฺย, อยมฺปิ ภิกฺขุนี ปฐมาปตฺติกํ ธมฺมํ อาปนฺนา นิสฺสารณียํ สํฆาทิเสสนฺ”}ติ.}

\prose{เอวญฺจิทํ ภควตา ภิกฺขุนีนํ สิกฺขาปทํ ปญฺญตฺตํ โหติ.}

\proseitem{690}{เตน โข ปน สมเยน สมฺพหุลา ภิกฺขุนิโย โกสเลสุ ชนปเท สาวตฺถิํ อทฺธานมคฺคปฺปฏิปนฺนา โหนฺติ, ตตฺถ อญฺญตรา ภิกฺขุนี วจฺเจน ปีฬิตา เอกิกา โอหียิตฺวา\footnote{โอหิยิตฺวา (ก)} ปจฺฉา อคมาสิ, มนุสฺสา ตํ ภิกฺขุนิํ ปสฺสิตฺวา ทูเสสุํ.\csromanspacer{} อถ โข สา ภิกฺขุนี เยน ตา ภิกฺขุนิโย เตนุปสงฺกมิ.\csromanspacer{} ภิกฺขุนิโย ตํ ภิกฺขุนิํ เอตทโวจุํ “กิสฺส ตฺวํ อยฺเย เอกิกา โอหีนา, กจฺจิสิ อปฺปธํสิตา”ติ.\csromanspacer{} ปธํสิตามฺหิ อยฺเยติ.\csromanspacer{} ยา ตา ภิกฺขุนิโย อปฺปิจฺฉา ฯเปฯ ตา อุชฺฌายนฺติ ขิยฺยนฺติ วิปาเจนฺติ “กถํ หิ นาม ภิกฺขุนี เอกา คณมฺหา โอหียิสฺสตี”ติ ฯเปฯ.\csromanspacer{} สจฺจํ กิร ภิกฺขเว ภิกฺขุนี เอกา คณมฺหา โอหียตีติ\footnote{โอหิยีติ (ก)}.\csromanspacer{} สจฺจํ ภควาติ.\csromanspacer{} วิครหิ พุทฺโธ ภควา ฯเปฯ.\csromanspacer{} กถํ หิ นาม ภิกฺขเว ภิกฺขุนี เอกา คณมฺหา โอหียิสฺสติ.\csromanspacer{} เนตํ ภิกฺขเว อปฺปสนฺนานํ วา ปสาทาย ฯเปฯ.\csromanspacer{} เอวญฺจ ปน ภิกฺขเว ภิกฺขุนิโย อิมํ สิกฺขาปทํ อุทฺทิสนฺตุ\csromandash{}}

\proseitem{691}{\csromanfolio{296}\textbf{“ยา ปน ภิกฺขุนี เอกา วา คามนฺตรํ คจฺเฉยฺย, เอกา วา นทีปารํ คจฺเฉยฺย, เอกา วา รตฺติํ วิปฺปวเสยฺย, เอกา วา คณมฺหา โอหีเยยฺย, อยมฺปิ ภิกฺขุนี ปฐมาปตฺติกํ ธมฺมํ อาปนฺนา นิสฺสารณียํ สํฆาทิเสสนฺ}”ติ.}

\proseitem{692}{\textbf{ยา ปนา}ติ ยา ยาทิสา ฯเปฯ.\csromanspacer{} \textbf{ภิกฺขุนี}ติ ฯเปฯ อยํ อิมสฺมิํ อตฺเถ อธิปฺเปตา ภิกฺขุนีติ.}

\prose{\textbf{เอกา วา คามนฺตรํ คจฺเฉยฺยา}ติ ปริกฺขิตฺตสฺส คามสฺส ปริกฺเขปํ ปฐมํ ปาทํ อติกฺกาเมนฺติยา อาปตฺติ ถุลฺลจฺจยสฺส.\csromanspacer{} ทุติยํ ปาทํ อติกฺกาเมนฺติยา อาปตฺติ สํฆาทิเสสสฺส.}

\prose{อปริกฺขิตฺตสฺส คามสฺส อุปจารํ ปฐมํ ปาทํ อติกฺกาเมนฺติยา อาปตฺติ ถุลฺลจฺจยสฺส.\csromanspacer{} ทุติยํ ปาทํ อติกฺกาเมนฺติยา อาปตฺติ สํฆาทิเสสสฺส.}

\prose{\textbf{เอกา วา นทีปารํ คจฺเฉยฺยา}ติ นที นาม ติมณฺฑลํ ปฏิจฺฉาเทตฺวา ยตฺถ กตฺถจิ อุตฺตรนฺติยา ภิกฺขุนิยา อนฺตรวาสโก เตมิยติ.\csromanspacer{} ปฐมํ ปาทํ อุตฺตรนฺติยา อาปตฺติ ถุลฺลจฺจยสฺส.\csromanspacer{} ทุติยํ ปาทํ อุตฺตรนฺติยา อาปตฺติ สํฆาทิเสสสฺส.}

\prose{\textbf{เอกา วา รตฺติํ วิปฺปวเสยฺยา}ติ สห อรุณุคฺคมนา ทุติยิกาย ภิกฺขุนิยา หตฺถปาสํ วิชหนฺติยา อาปตฺติ ถุลฺลจฺจยสฺส.\csromanspacer{} วิชหิเต อาปตฺติ สํฆาทิเสสสฺส.}

\prose{\textbf{เอกา วา คณมฺหา โอหีเยยฺยา}ติ อคามเก อรญฺเญ ทุติยิกาย ภิกฺขุนิยา ทสฺสนูปจารํ วา สวนูปจารํ วา วิชหนฺติยา อาปตฺติ ถุลฺลจฺจยสฺส.\csromanspacer{} วิชหิเต อาปตฺติ สํฆาทิเสสสฺส.}

\prose{\textbf{อยมฺปี}ติ ปุริมาโย อุปาทาย วุจฺจติ.}

\prose{\textbf{ปฐมาปตฺติกนฺ}ติ สห วตฺถุชฺฌาจารา อาปชฺชติ อสมนุภาสนาย.}

\prose{\textbf{นิสฺสารณียนฺ}ติ สํฆมฺหา นิสฺสารียติ.}

\prose{\textbf{สํฆาทิเสโส}ติ ฯเปฯ เตนปิ วุจฺจติ สํฆาทิเสโสติ.}

\proseitemwithcloser{693}{อนาปตฺติ ทุติยิกา ภิกฺขุนี ปกฺกนฺตา วา โหติ วิพฺภนฺตา วา กาลงฺกตา วา ปกฺขสงฺกนฺตา วา, อาปทาสุ อุมฺมตฺติกาย อาทิกมฺมิกายาติ.}{ตติยสํฆาทิเสสสิกฺขาปทํ นิฏฺฐิตํ.}
\csromanmatikaanchor{mtk.127}
\csromantocmarkref{section}{4. จตุตฺถสํฆาทิเสสสิกฺขาปท}{mtk.127}
\bookmark[dest={mtk.127},level=1]{4. จตุตฺถสํฆาทิเสสสิกฺขาปท}
\setcsromanheadh{4. จตุตฺถสํฆาทิเสสสิกฺขาปท}
\csromanheader{1}{4. \textbf{จตุตฺถสํฆาทิเสสสิกฺขาปท}}
\proseitem{694}{\csromanfolio{297}เตน สมเยน พุทฺโธ ภควา สาวตฺถิยํ วิหรติ เชตวเน อนาถปิณฺฑิกสฺส อาราเม.\csromanspacer{} เตน โข ปน สมเยน จณฺฑกาฬี ภิกฺขุนี ภณฺฑนการิกา โหติ กลหการิกา วิวาทการิกา ภสฺสการิกา สํเฆ อธิกรณการิกา, ถุลฺลนนฺทา ภิกฺขุนี ตสฺสา กมฺเม กรียมาเน ปฏิกฺโกสติ.\csromanspacer{} เตน โข ปน สมเยน ถุลฺลนนฺทา ภิกฺขุนี คามกํ อคมาสิ เกนจิเทว กรณีเยน.\csromanspacer{} อถ โข ภิกฺขุนิสํโฆ “ถุลฺลนนฺทา ภิกฺขุนี ปกฺกนฺตา”ติ จณฺฑกาฬิํ ภิกฺขุนิํ อาปตฺติยา อทสฺสเน\footnote{อทสฺสเนน (ก)} อุกฺขิปิ.\csromanspacer{} ถุลฺลนนฺทา ภิกฺขุนี คามเก ตํ กรณียํ ตีเรตฺวา ปุนเทว สาวตฺถิํ ปจฺจาคจฺฉิ.\csromanspacer{} จณฺฑกาฬี ภิกฺขุนี ถุลฺลนนฺทาย ภิกฺขุนิยา อาคจฺฉนฺติยา เนว อาสนํ ปญฺญเปสิ, น ปาโททกํ ปาทปีฐํ ปาทกฐลิกํ อุปนิกฺขิปิ, น ปจฺจุคฺคนฺตฺวา ปตฺตจีวรํ ปฏิคฺคเหสิ, น ปานีเยน อาปุจฺฉิ.\csromanspacer{} ถุลฺลนนฺทา ภิกฺขุนี จณฺฑกาฬิํ ภิกฺขุนิํ เอตทโวจ “กิสฺส ตฺวํ อยฺเย มยิ อาคจฺฉนฺติยา เนว อาสนํ ปญฺญเปสิ, น ปาโททกํ ปาทปีฐํ ปาทกฐลิกํ อุปนิกฺขิปิ, น ปจฺจุคฺคนฺตฺวา ปตฺตจีวรํ ปฏิคฺคเหสิ, น ปานีเยน อาปุจฺฉี”ติ.\csromanspacer{} เอวํ เหตํ อยฺเย โหติ, ยถา ตํ อนาถายาติ.\csromanspacer{} กิสฺส ปน ตฺวํ อยฺเย อนาถาติ.\csromanspacer{} อิมา มํ อยฺเย ภิกฺขุนิโย “อยํ อนาถา อปฺปญฺญาตา, นตฺถิ อิมิสฺสา กาจิ ปฏิวตฺตา”ติ อาปตฺติยา อทสฺสเน1 อุกฺขิปิํสูติ.}

\prose{ถุลฺลนนฺทา ภิกฺขุนี “พาลา เอตา อพฺยตฺตา เอตา เนว ชานนฺติ กมฺมํ วา กมฺมโทสํ วา กมฺมวิปตฺติํ วา กมฺมสมฺปตฺติํ วา, มยํ โข ชานาม กมฺมมฺปิ กมฺมโทสมฺปิ กมฺมวิปตฺติมฺปิ กมฺมสมฺปตฺติมฺปิ, มยํ โข อกตํ วา กมฺมํ กาเรยฺยาม, กตํ วา กมฺมํ โกเปยฺยามา”ติ ลหุํ ลหุํ ภิกฺขุนิสํฆํ สนฺนิปาเตตฺวา จณฺฑกาฬิํ ภิกฺขุนิํ โอสาเรสิ.\csromanspacer{} ยา ตา ภิกฺขุนิโย อปฺปิจฺฉา ฯเปฯ ตา อุชฺฌายนฺติ ขิยฺยนฺติ วิปาเจนฺติ “กถํ หิ นาม อยฺยา ถุลฺลนนฺทา สมคฺเคน สํเฆน อุกฺขิตฺตํ ภิกฺขุนิํ ธมฺเมน วินเยน สตฺถุสาสเนน อนปโลเกตฺวา การกสํฆํ อนญฺญาย คณสฺส ฉนฺทํ โอสาเรสฺสตี”ติ ฯเปฯ.\csromanspacer{} สจฺจํ กิร ภิกฺขเว ถุลฺลนนฺทา ภิกฺขุนี สมคฺเคน สํเฆน อุกฺขิตฺตํ ภิกฺขุนิํ ธมฺเมน วินเยน สตฺถุสาสเนน อนปโลเกตฺวา การกสํฆํ อนญฺญาย คณสฺส ฉนฺทํ โอสาเรตีติ\footnote{โอสาเรสีติ (ก)}. \csromanfolio{298}สจฺจํ ภควาติ.\csromanspacer{} วิครหิ พุทฺโธ ภควา ฯเปฯ.\csromanspacer{} กถํ หิ นาม ภิกฺขเว ถุลฺลนนฺทา ภิกฺขุนี สมคฺเคน สํเฆน อุกฺขิตฺตํ ภิกฺขุนิํ ธมฺเมน วินเยน สตฺถุสาสเนน อนปโลเกตฺวา การกสํฆํ อนญฺญาย คณสฺส ฉนฺทํ โอสาเรสฺสติ.\csromanspacer{} เนตํ ภิกฺขเว อปฺปสนฺนานํ วา ปสาทาย ฯเปฯ.\csromanspacer{} เอวญฺจ ปน ภิกฺขเว ภิกฺขุนิโย อิมํ สิกฺขาปทํ อุทฺทิสนฺตุ\csromandash{}}

\proseitem{695}{\textbf{“ยา ปน ภิกฺขุนี สมคฺเคน สํเฆน อุกฺขิตฺตํ ภิกฺขุนิํ ธมฺเมน วินเยน สตฺถุสาสเนน อนปโลเกตฺวา การกสํฆํ อนญฺญาย คณสฺส ฉนฺทํ โอสาเรยฺย, อยมฺปิ ภิกฺขุนี ปฐมาปตฺติกํ ธมฺมํ อาปนฺนา นิสฺสารณียํ สํฆาทิเสสนฺ”}ติ.}

\proseitem{696}{\textbf{ยา ปนา}ติ ยาทิสา ฯเปฯ.\csromanspacer{} \textbf{ภิกฺขุนี}ติ ฯเปฯ อยํ อิมสฺมิํ อตฺเถ อธิปฺเปตา ภิกฺขุนีติ.}

\prose{\textbf{สมคฺโค} นาม \textbf{สํโฆ} สมานสํวาสโก สมานสีมายํ ฐิโต.}

\prose{\textbf{อุกฺขิตฺตา} นาม อาปตฺติยา อทสฺสเน วา อปฺปฏิกมฺเม วา อปฺปฏินิสฺสคฺเค วา\footnote{อทสฺสเนน วา อปฺปฏิกมฺเมน วา อปฺปฏินิสฺสคฺเคน วา (ก), อทสฺสเน วา อปฺปฏิกมฺเม วา ปาปิกาย ทิฏฺฐิยา อปฺปฏินิสฺสคฺเค วา (?)} อุกฺขิตฺตา.}

\prose{\textbf{ธมฺเมน วินเยนา}ติ เยน ธมฺเมน เยน วินเยน.}

\prose{\textbf{สตฺถุสาสเนนา}ติ ชินสาสเนน พุทฺธสาสเนน.}

\prose{\textbf{อนปโลเกตฺวา การกสํฆนฺ}ติ กมฺมการกสํฆํ อนาปุจฺฉา.}

\prose{\textbf{อนญฺญาย คณสฺส ฉนฺทนฺ}ติ คณสฺส ฉนฺทํ อชานิตฺวา.}

\prose{“โอสาเรสฺสามี”ติ คณํ วา ปริเยสติ, สีมํ วา สมฺมนฺนติ, อาปตฺติ ทุกฺกฏสฺส.\csromanspacer{} ญตฺติยา ทุกฺกฏํ.\csromanspacer{} ทฺวีหิ กมฺมวาจาหิ ถุลฺลจฺจยา.\csromanspacer{} กมฺมวาจาปริโยสาเน อาปตฺติ สํฆาทิเสสสฺส.}

\prose{\textbf{อยมฺปี}ติ ปุริมาโย อุปาทาย วุจฺจติ.}

\prose{\textbf{ปฐมาปตฺติกนฺ}ติ สห วตฺถุชฺฌาจารา อาปชฺชติ อสมนุภาสนาย.}

\prose{\textbf{นิสฺสารณียนฺ}ติ สํฆมฺหา นิสฺสารียติ.}

\prose{\textbf{สํฆาทิเสโส}ติ ฯเปฯ เตนปิ วุจฺจติ สํฆาทิเสโสติ.}

\proseitem{697}{\csromanfolio{299}ธมฺมกมฺเม ธมฺมกมฺมสญฺญา โอสาเรติ, อาปตฺติ สํฆาทิเสสสฺส.\csromanspacer{} ธมฺมกมฺเม เวมติกา โอสาเรติ, อาปตฺติ สํฆาทิเสสสฺส.\csromanspacer{} ธมฺมกมฺเม อธมฺมกมฺมสญฺญา โอสาเรติ, อาปตฺติ สํฆาทิเสสสฺส.}

\prose{อธมฺมกมฺเม ธมฺมกมฺมสญฺญา, อาปตฺติ ทุกฺกฏสฺส.\csromanspacer{} อธมฺมกมฺเม เวมติกา, อาปตฺติ ทุกฺกฏสฺส.\csromanspacer{} อธมฺมกมฺเม อธมฺมกมฺมสญฺญา, อาปตฺติ ทุกฺกฏสฺส.}

\proseitemwithcloserruled{698}{อนาปตฺติ กมฺมการกสํฆํ อปโลเกตฺวา โอสาเรติ, คณสฺส ฉนฺทํ ชานิตฺวา โอสาเรติ, วตฺเต วตฺตนฺติํ โอสาเรติ, อสนฺเต กมฺมการกสํเฆ โอสาเรติ, อุมฺมตฺติกาย อาทิกมฺมิกายาติ.}{จตุตฺถสํฆาทิเสสสิกฺขาปทํ นิฏฺฐิตํ.}
\csromanmatikaanchor{mtk.128}
\csromantocmarkref{section}{5. ปญฺจมสํฆาทิเสสสิกฺขาปท}{mtk.128}
\bookmark[dest={mtk.128},level=1]{5. ปญฺจมสํฆาทิเสสสิกฺขาปท}
\setcsromanheadh{5. ปญฺจมสํฆาทิเสสสิกฺขาปท}
\csromanheader{1}{5. \textbf{ปญฺจมสํฆาทิเสสสิกฺขาปท}}
\proseitem{699}{เตน สมเยน พุทฺโธ ภควา สาวตฺถิยํ วิหรติ เชตวเน อนาถปิณฺฑิกสฺส อาราเม.\csromanspacer{} เตน โข ปน สมเยน สุนฺทรีนนฺทา ภิกฺขุนี อภิรูปา โหติ ทสฺสนียา ปาสาทิกา.\csromanspacer{} มนุสฺสา ภตฺตคฺเค สุนฺทรีนนฺทํ ภิกฺขุนิํ ปสฺสิตฺวา อวสฺสุตา อวสฺสุตาย สุนฺทรีนนฺทาย ภิกฺขุนิยา อคฺคมคฺคานิ โภชนานิ เทนฺติ.\csromanspacer{} สุนฺทรีนนฺทา ภิกฺขุนี ยาวทตฺถํ ภุญฺชติ, อญฺญา ภิกฺขุนิโย น จิตฺตรูปํ ลภนฺติ.\csromanspacer{} ยา ตา ภิกฺขุนิโย อปฺปิจฺฉา ฯเปฯ ตา อุชฺฌายนฺติ ขิยฺยนฺติ วิปาเจนฺติ “กถํ หิ นาม อยฺยา สุนฺทรีนนฺทา อวสฺสุตา อวสฺสุตสฺส ปุริสปุคฺคลสฺส หตฺถโต ขาทนียํ วา โภชนียํ วา สหตฺถา ปฏิคฺคเหตฺวา ขาทิสฺสติ ภุญฺชิสฺสตี”ติ ฯเปฯ.\csromanspacer{} สจฺจํ กิร ภิกฺขเว สุนฺทรีนนฺทา ภิกฺขุนี อวสฺสุตา อวสฺสุตสฺส ปุริสปุคฺคลสฺส หตฺถโต ขาทนียํ วา โภชนียํ วา สหตฺถา ปฏิคฺคเหตฺวา ขาทติ ภุญฺชตีติ.\csromanspacer{} สจฺจํ ภควาติ.\csromanspacer{} วิครหิ พุทฺโธ ภควา ฯเปฯ.\csromanspacer{} กถํ หิ นาม ภิกฺขเว สุนฺทรีนนฺทา ภิกฺขุนี อวสฺสุตา อวสฺสุตสฺส ปุริสปุคฺคลสฺส หตฺถโต ขาทนียํ วา โภชนียํ วา สหตฺถา ปฏิคฺคเหตฺวา ขาทิสฺสติ ภุญฺชิสฺสติ.\csromanspacer{} เนตํ ภิกฺขเว อปฺปสนฺนานํ วา ปสาทาย ฯเปฯ.\csromanspacer{} เอวญฺจ ปน ภิกฺขเว ภิกฺขุนิโย อิมํ สิกฺขาปทํ อุทฺทิสนฺตุ\csromandash{}}

\proseitem{700}{\csromanfolio{300}\textbf{“ยา ปน ภิกฺขุนี อวสฺสุตา อวสฺสุตสฺส ปุริสปุคฺคลสฺส หตฺถโต ขาทนียํ วา โภชนียํ วา สหตฺถา ปฏิคฺคเหตฺวา ขาเทยฺย วา ภุญฺเชยฺย วา, อยมฺปิ ภิกฺขุนี ปฐมาปตฺติกํ ธมฺมํ อาปนฺนา นิสฺสารณียํ สํฆาทิเสสนฺ”}ติ.}

\proseitem{701}{\textbf{ยา ปนา}ติ ยา ยาทิสา ฯเปฯ.\csromanspacer{} \textbf{ภิกฺขุนี}ติ ฯเปฯ อยํ อิมสฺมิํ อตฺเถ อธิปฺเปตา ภิกฺขุนีติ.}

\prose{\textbf{อวสฺสุตา} นาม สารตฺตา อเปกฺขวตี ปฏิพทฺธจิตฺตา.}

\prose{\textbf{อวสฺสุโต} นาม สารตฺโต อเปกฺขวา ปฏิพทฺธจิตฺโต.}

\prose{\textbf{ปุริสปุคฺคโล} นาม มนุสฺสปุริโส น ยกฺโข น เปโต น ติรจฺฉานคโต วิญฺญู ปฏิพโล สารชฺชิตุํ.}

\prose{\textbf{ขาทนียํ} นาม ปญฺจ โภชนานิ อุทกทนฺตโปนํ ฐเปตฺวา อวเสสํ ขาทนียํ นาม.}

\prose{\textbf{โภชนียํ} นาม ปญฺจ โภชนานิ โอทโน กุมฺมาโส สตฺตุ มจฺโฉ มํสํ.}

\prose{“ขาทิสฺสามิ ภุญฺชิสฺสามี”ติ ปฏิคฺคณฺหาติ, อาปตฺติ ถุลฺลจฺจยสฺส.\csromanspacer{} อชฺโฌหาเร อชฺโฌหาเร อาปตฺติ สํฆาทิเสสสฺส.}

\prose{\textbf{อยมฺปี}ติ ปุริมาโย อุปาทาย วุจฺจติ.}

\prose{\textbf{ปฐมาปตฺติกนฺ}ติ สห วตฺถุชฺฌาจารา อาปชฺชติ อสมนุภาสนาย.}

\prose{\textbf{นิสฺสารณียนฺ}ติ สํฆมฺหา นิสฺสารียติ.}

\prose{\textbf{สํฆาทิเสโส}ติ ฯเปฯ เตนปิ วุจฺจติ สํฆาทิเสโสติ.}

\prose{อุทกทนฺตโปนํ ปฏิคฺคณฺหาติ, อาปตฺติ ทุกฺกฏสฺส.\csromanspacer{} เอกโตอวสฺสุเต “ขาทิสฺสามิ ภุญฺชิสฺสามี”ติ ปฏิคฺคณฺหาติ, อาปตฺติ ทุกฺกฏสฺส.}

\proseitem{702}{อชฺโฌหาเร อชฺโฌหาเร อาปตฺติ ถุลฺลจฺจยสฺส.\csromanspacer{} อุทกทนฺตโปนํ ปฏิคฺคณฺหาติ, อาปตฺติ ทุกฺกฏสฺส.\csromanspacer{} อุภโตอวสฺสุเต ยกฺขสฺส วา เปตสฺส วา ปณฺฑกสฺส วา ติรจฺฉานคต\-มนุสฺสวิคฺคหสฺส วา หตฺถโต “ขาทิสฺสามิ ภุญฺชิสฺสามี”ติ ปฏิคฺคณฺหาติ, อาปตฺติ ทุกฺกฏสฺส.\csromanspacer{} อชฺโฌหาเร \csromanfolio{301}อชฺโฌหาเร อาปตฺติ ถุลฺลจฺจยสฺส.\csromanspacer{} อุทกทนฺตโปนํ ปฏิคฺคณฺหาติ, อาปตฺติ ทุกฺกฏสฺส.\csromanspacer{} เอกโต อวสฺสุเต “ขาทิสฺสามิ ภุญฺชิสฺสามี”ติ ปฏิคฺคณฺหาติ, อาปตฺติ ทุกฺกฏสฺส.\csromanspacer{} อชฺโฌหาเร อชฺโฌหาเร อาปตฺติ ทุกฺกฏสฺส.\csromanspacer{} อุทกทนฺตโปนํ ปฏิคฺคณฺหาติ, อาปตฺติ ทุกฺกฏสฺส.}

\proseitemwithcloserruled{703}{อนาปตฺติ อุภโตอนวสฺสุตา โหนฺติ, “อนวสฺสุโต”ติ ชานนฺตี ปฏิคฺคณฺหาติ, อุมฺมตฺติกาย อาทิกมฺมิกายาติ.}{ปญฺจมสํฆาทิเสสสิกฺขาปทํ นิฏฺฐิตํ.}
\csromanmatikaanchor{mtk.129}
\csromantocmarkref{section}{6. ฉฏฺฐสํฆาทิเสสสิกฺขาปท}{mtk.129}
\bookmark[dest={mtk.129},level=1]{6. ฉฏฺฐสํฆาทิเสสสิกฺขาปท}
\setcsromanheadh{6. ฉฏฺฐสํฆาทิเสสสิกฺขาปท}
\csromanheader{1}{6. \textbf{ฉฏฺฐสํฆาทิเสสสิกฺขาปท}}
\proseitem{704}{เตน สมเยน พุทฺโธ ภควา สาวตฺถิยํ วิหรติ เชตวเน อนาถปิณฺฑิกสฺส อาราเม.\csromanspacer{} เตน โข ปน สมเยน สุนฺทรีนนฺทา ภิกฺขุนี อภิรูปา โหติ ทสฺสนียา ปาสาทิกา, มนุสฺสา ภตฺตคฺเค สุนฺทรีนนฺทํ ภิกฺขุนิํ ปสฺสิตฺวา อวสฺสุตา สุนฺทรีนนฺทาย ภิกฺขุนิยา อคฺคมคฺคานิ โภชนานิ เทนฺติ, สุนฺทรีนนฺทา ภิกฺขุนี กุกฺกุจฺจายนฺตี น ปฏิคฺคณฺหาติ.\csromanspacer{} อนนฺตริกา ภิกฺขุนี สุนฺทรีนนฺทํ ภิกฺขุนิํ เอตทโวจ “กิสฺส ตฺวํ อยฺเย น ปฏิคฺคณฺหาสี”ติ.\csromanspacer{} อวสฺสุตา อยฺเยติ.\csromanspacer{} ตฺวํ ปน อยฺเย อวสฺสุตาติ.\csromanspacer{} นาหํ อยฺเย อวสฺสุตาติ.\csromanspacer{} กิํ เต อยฺเย เอโส ปุริสปุคฺคโล กริสฺสติ อวสฺสุโต วา อนวสฺสุโต วา, ยโต ตฺวํ อนวสฺสุตา, อิงฺฆ อยฺเย ยํ เต เอโส ปุริสปุคฺคโล เทติ ขาทนียํ วา โภชนียํ วา, ตํ ตฺวํ สหตฺถา ปฏิคฺคเหตฺวา ขาท วา ภุญฺช วาติ.}

\prose{ยา ตา ภิกฺขุนิโย อปฺปิจฺฉา ฯเปฯ ตา อุชฺฌายนฺติ ขิยฺยนฺติ วิปาเจนฺติ “กถํ หิ นาม ภิกฺขุนี เอวํ วกฺขติ ‘กิํ เต อยฺเย เอโส ปุริสปุคฺคโล กริสฺสติ อวสฺสุโต วา อนวสฺสุโต วา, ยโต ตฺวํ อนวสฺสุตา, อิงฺฆ อยฺเย ยํ เต เอโส ปุริสปุคฺคโล เทติ ขาทนียํ วา โภชนียํ วา, ตํ ตฺวํ สหตฺถา ปฏิคฺคเหตฺวา ขาท วา ภุญฺช วา’ติ” ฯเปฯ.\csromanspacer{} สจฺจํ กิร ภิกฺขเว ภิกฺขุนี เอวํ วเทติ “กิํ เต อยฺเย เอโส ปุริสปุคฺคโล กริสฺสติ อวสฺสุโต วา อนวสฺสุโต วา, ยโต ตฺวํ อนวสฺสุตา, อิงฺฆ อยฺเย ยํ เต เอโส ปุริสปุคฺคโล เทติ ขาทนียํ วา โภชนียํ วา, ตํ ตฺวํ สหตฺถา ปฏิคฺคเหตฺวา ขาท วา ภุญฺช วา”ติ.\csromanspacer{} สจฺจํ ภควาติ. \csromanfolio{302}วิครหิ พุทฺโธ ภควา ฯเปฯ.\csromanspacer{} กถํ หิ นาม ภิกฺขเว ภิกฺขุนี เอวํ วกฺขติ “กิํ เต อยฺเย เอโส ปุริสปุคฺคโล กริสฺสติ อวสฺสุโต วา อนวสฺสุโต วา, ยโต ตฺวํ อนวสฺสุตา, อิงฺฆ อยฺเย ยํ เต เอโส ปุริสปุคฺคโล เทติ ขาทนียํ วา โภชนียํ วา, ตํ ตฺวํ สหตฺถา ปฏิคฺคเหตฺวา ขาท วา ภุญฺช วา”ติ.\csromanspacer{} เนตํ ภิกฺขเว อปฺปสนฺนานํ วา ปสาทาย ฯเปฯ.\csromanspacer{} เอวญฺจ ปน ภิกฺขเว ภิกฺขุนิโย อิมํ สิกฺขาปทํ อุทฺทิสนฺตุ\csromandash{}}

\proseitem{705}{\textbf{“ยา ปน ภิกฺขุนี เอวํ วเทยฺย ‘กิํ เต อยฺเย เอโส ปุริสปุคฺคโล กริสฺสติ อวสฺสุโต วา อนวสฺสุโต วา, ยโต ตฺวํ อนวสฺสุตา, อิงฺฆ อยฺเย ยํ เต เอโส ปุริสปุคฺคโล เทติ ขาทนียํ วา โภชนียํ วา, ตํ ตฺวํ สหตฺถา ปฏิคฺคเหตฺวา ขาท วา ภุญฺช วา’ติ, อยมฺปิ ภิกฺขุนี ปฐมาปตฺติกํ ธมฺมํ อาปนฺนา นิสฺสารณียํ สํฆาทิเสสนฺ”}ติ.}

\proseitem{706}{\textbf{ยา ปนา}ติ ยา ยาทิสา ฯเปฯ.\csromanspacer{} \textbf{ภิกฺขุนี}ติ ฯเปฯ อยํ อิมสฺมิํ อตฺเถ อธิปฺเปตา ภิกฺขุนีติ.}

\prose{\textbf{เอวํ วเทยฺยา}ติ “กิํ เต อยฺเย เอโส ปุริสปุคฺคโล กริสฺสติ อวสฺสุโต วา อนวสฺสุโต วา, ยโต ตฺวํ อนวสฺสุตา, อิงฺฆ อยฺเย ยํ เต เอโส ปุริสปุคฺคโล เทติ ขาทนียํ วา โภชนียํ วา, ตํ ตฺวํ สหตฺถา ปฏิคฺคเหตฺวา ขาท วา ภุญฺช วา”ติ อุยฺโยเชติ อาปตฺติ ทุกฺกฏสฺส.\csromanspacer{} ตสฺสา วจเนน “ขาทิสฺสามิ ภุญฺชิสฺสามี”ติ ปฏิคฺคณฺหาติ, อาปตฺติ ทุกฺกฏสฺส.\csromanspacer{} อชฺโฌหาเร อชฺโฌหาเร อาปตฺติ ถุลฺลจฺจยสฺส.\csromanspacer{} โภชนปริโยสาเน อาปตฺติ สํฆาทิเสสสฺส.}

\prose{\textbf{อยมฺปี}ติ ปุริมาโย อุปาทาย วุจฺจติ.}

\prose{\textbf{ปฐมาปตฺติกนฺ}ติ สห วตฺถุชฺฌาจารา อาปชฺชติ อสมนุภาสนาย.}

\prose{\textbf{นิสฺสารณียนฺ}ติ สํฆมฺหา นิสฺสารียติ.}

\prose{\textbf{สํฆาทิเสโส}ติ ฯเปฯ เตนปิ วุจฺจติ สํฆาทิเสโสติ.}

\prose{อุทกทนฺตโปนํ “ปฏิคฺคณฺหา”ติ อุยฺโยเชติ, อาปตฺติ ทุกฺกฏสฺส.\csromanspacer{} ตสฺสา วจเนน “ขาทิสฺสามิ ภุญฺชิสฺสามี”ติ ปฏิคฺคณฺหาติ, อาปตฺติ ทุกฺกฏสฺส.}

\proseitem{707}{\csromanfolio{303}เอกโต อวสฺสุเต ยกฺขสฺส วา เปตสฺส วา ปณฺฑกสฺส วา ติรจฺฉานคต\-มนุสฺสวิคฺคหสฺส วา หตฺถโต ขาทนียํ วา โภชนียํ วา “ขาท วา ภุญฺช วา”ติ อุยฺโยเชติ, อาปตฺติ ทุกฺกฏสฺส.\csromanspacer{} ตสฺสา วจเนน “ขาทิสฺสามิ ภุญฺชิสฺสามี”ติ ปฏิคฺคณฺหาติ, อาปตฺติ ทุกฺกฏสฺส.\csromanspacer{} อชฺโฌหาเร อชฺโฌหาเร อาปตฺติ ทุกฺกฏสฺส.\csromanspacer{} โภชนปริโยสาเน อาปตฺติ ถุลฺลจฺจยสฺส.\csromanspacer{} อุทกทนฺตโปนํ “ปฏิคฺคณฺหา”ติ อุยฺโยเชติ, อาปตฺติ ทุกฺกฏสฺส.\csromanspacer{} ตสฺสา วจเนน “ขาทิสฺสามิ ภุญฺชิสฺสามี”ติ ปฏิคฺคณฺหาติ, อาปตฺติ ทุกฺกฏสฺส.}

\proseitemwithcloserruled{708}{อนาปตฺติ “อนวสฺสุโต”ติ ชานนฺตี อุยฺโยเชติ, “กุปิตา น ปฏิคฺคณฺหาตี”ติ อุยฺโยเชติ, “กุลานุทฺทยตาย น ปฏิคฺคณฺหาตี”ติ อุยฺโยเชติ, อุมฺมตฺติกาย อาทิกมฺมิกายาติ.}{ฉฏฺฐสํฆาทิเสสสิกฺขาปทํ นิฏฺฐิตํ.}
\csromanmatikaanchor{mtk.130}
\csromantocmarkref{section}{7. สตฺตมสํฆาทิเสสสิกฺขาปท}{mtk.130}
\bookmark[dest={mtk.130},level=1]{7. สตฺตมสํฆาทิเสสสิกฺขาปท}
\setcsromanheadh{7. สตฺตมสํฆาทิเสสสิกฺขาปท}
\csromanheader{1}{7. \textbf{สตฺตมสํฆาทิเสสสิกฺขาปท}}
\proseitem{709}{เตน สมเยน พุทฺโธ ภควา สาวตฺถิยํ วิหรติ เชตวเน อนาถปิณฺฑิกสฺส อาราเม.\csromanspacer{} เตน โข ปน สมเยน จณฺฑกาฬี ภิกฺขุนี ภิกฺขุนีหิ สทฺธิํ ภณฺฑิตฺวา กุปิตา อนตฺตมนา เอวํ วเทติ “พุทฺธํ ปจฺจาจิกฺขามิ, ธมฺมํ ปจฺจาจิกฺขามิ, สํฆํ ปจฺจาจิกฺขามิ, สิกฺขํ ปจฺจาจิกฺขามิ, กินฺนุมาว สมณิโย ยา สมณิโย สกฺยธีตโร, สนฺตญฺญาปิ สมณิโย ลชฺชินิโย กุกฺกุจฺจิกา สิกฺขากามา, ตาสาหํ สนฺติเก พฺรหฺมจริยํ จริสฺสามี”ติ.\csromanspacer{} ยา ตา ภิกฺขุนิโย อปฺปิจฺฉา ฯเปฯ ตา อุชฺฌายนฺติ ขิยฺยนฺติ วิปาเจนฺติ “กถํ หิ นาม อยฺยา จณฺฑกาฬี\footnote{จณฺฑกาฬี ภิกฺขุนี (ก)} กุปิตา อนตฺตมนา เอวํ วกฺขติ พุทฺธํ ปจฺจาจิกฺขามิ ฯเปฯ สิกฺขํ ปจฺจาจิกฺขามิ, กินฺนุมาว สมณิโย ยา สมณิโย สกฺยธีตโร, สนฺตญฺญาปิ สมณิโย ลชฺชินิโย กุกฺกุจฺจิกา สิกฺขากามา, ตสาหํ สนฺติเก พฺรหฺมจริยํ จริสฺสามี”ติ ฯเปฯ.\csromanspacer{} สจฺจํ กิร ภิกฺขเว จณฺฑกาฬี ภิกฺขุนี กุปิตา อนตฺตมนา เอวํ วเทติ “พุทฺธํ ปจฺจาจิกฺขามิ ฯเปฯ สิกฺขํ ปจฺจาจิกฺขามิ, กินฺนุมาว สมณิโย ยา สมณิโย สกฺยธีตโร, สนฺตญฺญาปิ สมณีโย ลชฺชินิโย \csromanfolio{304}กุกฺกุจฺจิกา สิกฺขากามา, ตาสาหํ สนฺติเก พฺรหฺมจริยํ จริสฺสามี”ติ.\csromanspacer{} สจฺจํ ภควาติ.\csromanspacer{} วิครหิ พุทฺโธ ภควา ฯเปฯ.\csromanspacer{} กถํ หิ นาม ภิกฺขเว จณฺฑกาฬี ภิกฺขุนี กุปิตา อนตฺตมนา เอวํ วกฺขติ “พุทฺธํ ปจฺจาจิกฺขามิ ฯเปฯ สิกฺขํ ปจฺจาจิกฺขา- มิ, กินฺนุมาว สมณิโย ยา สมณิโย สกฺยธีตโร, สนฺตญฺญาปิ สมณิโย ลชฺชินิโย กุ- กฺกุจฺจิกา สิกฺขากามา, ตาสาหํ สนฺติเก พฺรหฺมจริยํ จริสฺสามี”ติ.\csromanspacer{} เนตํ ภิกฺขเว อปฺปสนฺนานํ วา ปสาทาย ฯเปฯ.\csromanspacer{} เอวญฺจ ปน ภิกฺขเว ภิกฺขุนิโย อิมํ สิกฺขาปทํ อุทฺทิสนฺตุ\csromandash{}}

\proseitem{710}{“ยา ปน ภิกฺขุนี กุปิตา อนตฺตมนา เอวํ วเทยฺย ‘พุทฺธํ ปจฺจาจิกฺขามิ, ธมฺมํ ปจฺจาจิกฺขามิ, สํฆํ ปจฺจาจิกฺขามิ, สิกฺขํ ปจฺจาจิกฺขามิ, กินฺนุมาว สมณิโย ยา สมณิโย สกฺยธีตโร, สนฺตญฺญาปิ สมณิโย ลชฺชินิโย กุกฺกุจฺจิกา สิกฺขากามา, ตาสาหํ สนฺติเก พฺรหฺมจริยํ จริสฺสามี’ติ.\csromanspacer{}\textbf{ สา ภิกฺขุนี ภิกฺขุนีหิ เอวมสฺส วจนียา ‘มายฺเย กุปิตา อนตฺตมนา เอวํ อวจ พุทฺธํ ปจฺจาจิกฺขามิ, ธมฺมํ ปจฺจาจิกฺขามิ, สํฆํ ปจฺจาจิกฺขามิ, สิกฺขํ ปจฺจาจิกฺขามิ, กินฺนุมาว สมณิโย ยา สมณิโย สกฺยธีตโร, สนฺตญฺญาปิ สมณิโย ลชฺชินิโย กุกฺกุจฺจิกา สิกฺขากามา, ตาสาหํ สนฺติเก พฺรหฺมจริยํ จริสฺสามีติ, อภิรมายฺเย สฺวากฺขาโต ธมฺโม, จร พฺรหฺมจริยํ สมฺมา ทุกฺขสฺส อนฺตกิริยายา’ติ, เอวญฺจ สา ภิกฺขุนี ภิกฺขุนีหิ วุจฺจมานา ตเถว ปคฺคเณฺหยฺย, สา ภิกฺขุนี ภิกฺขุนีหิ ยาวตติยํ สมนุภาสิตพฺพา ตสฺส ปฏินิสฺสคฺคาย, ยาวตติยญฺเจ สมนุภาสียมานา ตํ ปฏินิสฺสชฺเชยฺย, อิจฺเจตํ กุสลํ, โน เจ ปฏินิสฺสชฺเชยฺย, อยมฺปิ ภิกฺขุนี ยาวตติยกํ ธมฺมํ อาปนฺนา นิสฺสารณียํ สํฆาทิเสสนฺ”}ติ.}

\proseitem{711}{\textbf{ยา ปนา}ติ ยา ยาทิสา ฯเปฯ.\csromanspacer{} \textbf{ภิกฺขุนี}ติ ฯเปฯ อยํ อิมสฺมิํ อตฺเถ อธิปฺเปตา ภิกฺขุนีติ.}

\prose{\textbf{กุปิตา อนตฺตมนา}ติ อนภิรทฺธา อาหตจิตฺตา ขิลชาตา.}

\prose{\textbf{เอวํ วเทยฺยา}ติ “พุทฺธํ ปจฺจาจิกฺขามิ ฯเปฯ สิกฺขํ ปจฺจาจิกฺขามิ, กินฺนุมาว สมณิโย ยา สมณิโย สกฺยธีตโร, สนฺตญฺญาปิ สมณิโย ลชฺชินิโย กุกฺกุจฺจิกา สิกฺขากามา, ตาสาหํ สนฺติเก พฺรหฺมจริยํ จริสฺสามี”ติ.}

\prose{\csromanfolio{305}\textbf{สา ภิกฺขุนี}ติ ยา สา เอวํวาทินี ภิกฺขุนี.\csromanspacer{} \textbf{ภิกฺขุนีหี}ติ อญฺญาหิ ภิกฺขุนีหิ.}

\prose{ยา ปสฺสนฺติ, ยา สุณนฺติ, ตาหิ วตฺตพฺพา “มายฺเย กุปิตา อนตฺตมนา เอวํ อวจ ‘พุทฺธํ ปจฺจาจิกฺขามิ ฯเปฯ สิกฺขํ ปจฺจาจิกฺขามิ, กินฺนุมาว สมณิโย ยา สมาณิโย สกฺยธีตโร, สนฺตญฺญาปิ สมณิโย ลชฺชินิโย กุกฺกุจฺจิกา สิกฺขากามา, ตาสาหํ สนฺติเก พฺรหฺมจริยํ จริสฺสามี’ติ, อภิรมายฺเย, สฺวากฺขาโต ธมฺโม, จร พฺรหฺมจริยํ สมฺมา ทุกฺขสฺส อนฺตกิริยายา”ติ.\csromanspacer{} ทุติยมฺปิ วตฺตพฺพา.\csromanspacer{} ตติยมฺปิ วตฺตพฺพา.\csromanspacer{} สเจ ปฏินิสฺสชฺชติ, อิจฺเจตํ กุสลํ, โน เจ ปฏินิสฺสชฺชติ, อาปตฺติ ทุกฺกฏสฺส.\csromanspacer{} สุตฺวา น วทนฺติ, อาปตฺติ ทุกฺกฏสฺส.\csromanspacer{} สา ภิกฺขุนี สํฆมชฺฌมฺปิ อากฑฺฒิตฺวา วตฺตพฺพา “มายฺเย กุปิตา อนตฺตมนา เอวํ อวจ ‘พุทฺธํ ปจฺจาจิกฺขามิ ฯเปฯ สิกฺขํ ปจฺจาจิกฺขามิ, กินฺนุมาว สมณิโย ยา สมณิโย สกฺยธีตโร, สนฺตญฺญาปิ สมณิโย ลชฺชินิโย กุกฺกุจฺจิกา สิกฺขากามา, ตาสาหํ สนฺติเก พฺรหฺมจริยํ จริสฺสามี’ติ.\csromanspacer{} อภิรมายฺเย, สฺวากฺขาโต ธมฺโม, จร พฺรหฺมจริยํ สมฺมา ทุกฺขสฺส อนฺตกิริยายา”ติ.\csromanspacer{} ทุติยมฺปิ วตฺตพฺพา.\csromanspacer{} ตติยมฺปิ วตฺตพฺพา.\csromanspacer{} สเจ ปฏินิสฺสชฺชติ, อิจฺเจตํ กุสลํ, โน เจ ปฏินิสฺสชฺชติ, อาปตฺติ ทุกฺกฏสฺส.\csromanspacer{} สา ภิกฺขุนี สมนุภาสิตพฺพา, เอวญฺจ ปน ภิกฺขเว สมนุภาสิตพฺพา, พฺยตฺตาย ภิกฺขุนิยา ปฏิพลาย สํโฆ ญาเปตพฺโพ\csromandash{}}

\proseitem{712}{“สุณาตุ เม อยฺเย สํโฆ, อยํ อิตฺถนฺนามา ภิกฺขุนี กุปิตา อนตฺตมนา เอวํ วเทติ ‘พุทฺธํ ปจฺจาจิกฺขามิ, ธมฺมํ ปจฺจาจิกฺขามิ, สํฆํ ปจฺจาจิกฺขามิ, สิกฺขํ ปจฺจาจิกฺขามิ, กินฺนุมาว สมณิโย ยา สมณิโย สกฺยธีตโร, สนฺตญฺญาปิ สมณิโย ลชฺชินิโย กุกฺกุจฺจิกา สิกฺขากามา, ตาสาหํ สนฺติเก พฺรหฺมจริยํ จริสฺสามี’ติ, สา ตํ วตฺถุํ น ปฏินิสฺสชฺชติ, ยทิ สํฆสฺส ปตฺตกลฺลํ, สํโฆ อิตฺถนฺนามํ ภิกฺขุนิํ สมนุภาเสยฺย ตสฺส วตฺถุสฺส ปฏินิสฺสคฺคาย, เอสา ญตฺติ.}

\prose{สุณาตุ เม อยฺเย สํโฆ, อยํ อิตฺถนฺนามา ภิกฺขุนี กุปิตา อนตฺตมนา เอวํ วเทติ ‘พุทฺธํ ปจฺจาจิกฺขามิ ฯเปฯ สิกฺขํ ปจฺจาจิกฺขามิ, กินฺนุมาว สมณิโย ยา สมณิโย สกฺยธีตโร, สนฺตญฺญาปิ สมณิโย ลชฺชินิโย กุกฺกุจฺจิกา สิกฺขากามา, ตาสาหํ สนฺติเก พฺรหฺมจริยํ \csromanfolio{306}\textbf{จริสฺสามี’ติ, สา ตํ วตฺถุํ น ปฏินิสฺสชฺชติ, สํโฆ อิตฺถนฺนามํ} ภิกฺขุนิํ สมนุภาสติ ตสฺส วตฺถุสฺส ปฏินิสฺสคฺคาย, ยสฺสา อยฺยาย ขมติ อิตฺถนฺนามาย ภิกฺขุนิยา สมนุภาสนา ตสฺส วตฺถุสฺส ปฏินิสฺสคฺคาย, สา ตุณฺหสฺส, ยสฺสา นกฺขมติ, สา ภาเสยฺย.}

\prose{ทุติยมฺปิ เอตมตฺถํ วทามิ ฯเปฯ.\csromanspacer{} ตติยมฺปิ เอตมตฺถํ วทามิ ฯเปฯ.}

\prose{สมนุภฏฺฐา สํเฆน อิตฺถนฺนามา ภิกฺขุนี ตสฺส วตฺถุสฺส ปฏินิสฺสคฺคาย, ขมติ สํฆสฺส, ตสฺมา ตุณฺหี, เอวเมตํ ธารยามี”ติ.}

\prose{ญตฺติยา ทุกฺกฏํ, ทฺวีหิ กมฺมวาจาหิ ถุลฺลจฺจยา, กมฺมวาจาปริโยสาเน อาปตฺติ สํฆาทิเสสสฺส.\csromanspacer{} สํฆาทิเสสํ อชฺฌาปชฺชนฺติยา ญตฺติยา ทุกฺกฏํ ทฺวีหิ กมฺมวาจาหิ ถุลฺลจฺจยา ปฏิปฺปสฺสมฺภนฺติ.}

\prose{\textbf{อยมฺปี}ติ ปุริมาโย อุปาทาย วุจฺจติ.}

\prose{\textbf{ยาวตติยกนฺ}ติ ยาวตติยํ สมนุภาสนาย อาปชฺชติ, น สห วตฺถุชฺฌาจารา.}

\prose{\textbf{นิสฺสรณียนฺ}ติ สํฆมฺหา นิสฺสารียติ.}

\prose{\textbf{สํฆาทิเสโส}ติ ฯเปฯ เตนปิ วุจฺจติ สํฆาทิเสโสติ.}

\proseitem{713}{ธมฺมกมฺเม ธมฺมกมฺมสญฺญา น ปฏินิสฺสชฺชติ, อาปตฺติ สํฆาทิเสสสฺส.\csromanspacer{} ธมฺมกมฺเม เวมติกา น ปฏินิสฺสชฺชติ, อาปตฺติ สํฆาทิเสสสฺส.\csromanspacer{} ธมฺมกมฺเม อธมฺมกมฺมสญฺญา น ปฏินิสฺสชฺชติ, อาปตฺติ สํฆาทิเสสสฺส.}

\prose{อธมฺมกมฺเม ธมฺมกมฺมสญฺญา, อาปตฺติ ทุกฺกฏสฺส.\csromanspacer{} อธมฺมกมฺเม เวมติกา, อาปตฺติ ทุกฺกฏสฺส.\csromanspacer{} อธมฺมกมฺเม อธมฺมกมฺมสญฺญา, อาปตฺติ ทุกฺกฏสฺส.}

\proseitemwithcloser{714}{อนาปตฺติ อสมนุภาสนฺติยา ปฏินิสฺสชฺชนฺติยา อุมฺมตฺติกาย อาทิกมฺมิกายาติ.}{สตฺตมสํฆาทิเสสสิกฺขาปทํ นิฏฺฐิตํ.}
\csromanmatikaanchor{mtk.131}
\csromantocmarkref{section}{8. อฏฺฐมสํฆาทิเสสสิกฺขาปท}{mtk.131}
\bookmark[dest={mtk.131},level=1]{8. อฏฺฐมสํฆาทิเสสสิกฺขาปท}
\setcsromanheadh{8. อฏฺฐมสํฆาทิเสสสิกฺขาปท}
\csromanheader{1}{8. \textbf{อฏฺฐมสํฆาทิเสสสิกฺขาปท}}
\proseitem{715}{\csromanfolio{307}เตน สมเยน พุทฺโธ ภควา สาวตฺถิยํ วิหรติ เชตวเน อนาถปิณฺฑิกสฺส อาราเม.\csromanspacer{} เตน โข ปน สมเยน จณฺฑกาฬี ภิกฺขุนี กิสฺมิญฺจิเทว อธิกรเณ ปจฺจากตา กุปิตา อนตฺตมนา เอวํ วเทติ “ฉนฺทคามินิโย จ ภิกฺขุนิโย โทสคามินิโย จ ภิกฺขุนิโย โมหคามินิโย จ ภิกฺขุนิโย ภยคามินิโย จ ภิกฺขุนิโย”ติ.\csromanspacer{} ยา ตา ภิกฺขุนิโย อปฺปิจฺฉา ฯเปฯ ตา อุชฺฌายนฺติ ขิยฺยนฺติ วิปาเจนฺติ “กถํ หิ นาม อยฺยา จณฺฑกาฬี กิสฺมิญฺจิเทว อธิกรเณ ปจฺจากตา กุปิตา อนตฺตมนา เอวํ วกฺขติ ‘ฉนฺทคามินิโย จ ภิกฺขุนิโย ฯเปฯ ภยคามินิโย จ ภิกฺขุนิโย’ติ ฯเปฯ.\csromanspacer{} สจฺจํ กิร ภิกฺขเว จณฺฑกาฬี ภิกฺขุนี กิสฺมิญฺจิเทว อธิกรเณ ปจฺจากตา กุปิตา อนตฺตมนา เอวํ วเทติ “ฉนฺทคามินิโย จ ภิกฺขุนิโย ฯเปฯ ภยคามินิโย จ ภิกฺขุนิโย”ติ.\csromanspacer{} สจฺจํ ภควาติ.\csromanspacer{} วิครหิ พุทฺโธ ภควา ฯเปฯ.\csromanspacer{} กถํ หิ นาม ภิกฺขเว จณฺฑกาฬี ภิกฺขุนี กิสฺมิญฺจิเทว อธิกรเณ ปจฺจากตา กุปิตา อนตฺตมนา เอวํ วกฺขติ “ฉนฺทคามินิโย จ ภิกฺขุนิโย ฯเปฯ ภยคามินิโย จ ภิกฺขุนิโย”ติ.\csromanspacer{} เนตํ ภิกฺขเว อปฺปสนฺนานํ วา ปสาทาย ฯเปฯ.\csromanspacer{} เอวญฺจ ปน ภิกฺขเว ภิกฺขุนิโย อิมํ สิกฺขาปทํ อุทฺทิสนฺตุ\csromandash{}}

\proseitem{716}{\textbf{“ยา ปน ภิกฺขุนี กิสฺมิญฺจิเทว อธิกรเณ ปจฺจากตา กุปิตา อนตฺตมนา เอวํ วเทยฺย ‘ฉนฺทคามินิโย จ ภิกฺขุนิโย โทสคามินิโย จ ภิกฺขุนิโย โมหคามินิโย จ ภิกฺขุนิโย ภยคามินิโย จ ภิกฺขุนิโย’ติ, สา ภิกฺขุนี ภิกฺขุนีหิ เอวมสฺส วจนียา ‘มายฺเย กิสฺมิญฺจิเทว อธิกรเณ ปจฺจากตา กุปิตา อนตฺตมนา เอวํ อวจ ฉนฺทคามินิโย จ ภิกฺขุนิโย โทสคามินิโย จ ภิกฺขุนิโย โมหคามินิโย จ ภิกฺขุนิโย ภยคามินิโย จ ภิกฺขุนิโยติ, อยฺยา โข ฉนฺทาปิ คจฺเฉยฺย, โทสาปิ คจฺเฉยฺย, โมหาปิ คจฺเฉยฺย, ภยาปิ คจฺเฉยฺยา’ติ.}\csromanspacer{} เอวญฺจ สา ภิกฺขุนี ภิกฺขุนีหิ วุจฺจมานา ตเถว ปคฺคเณฺหยฺย, สา ภิกฺขุนี ภิกฺขุนีหิ ยาวตติยํ สมนุภาสิตพฺพา ตสฺส ปฏินิสฺสคฺคาย, ยาวตติยญฺเจ สมนุภาสียมานา ตํ ปฏินิสฺสชฺเชยฺย, อิจฺเจตํ กุสลํ, โน เจ ปฏินิสฺสชฺเชยฺย, อยมฺปิ ภิกฺขุนี ยาวตติยกํ ธมฺมํ อาปนฺนา นิสฺสารณียํ สํฆาทิเสสนฺ”ติ.}

\proseitem{717}{\csromanfolio{308}\textbf{ยา ปน}ติ ยา ยาทิสา ฯเปฯ.\csromanspacer{} \textbf{ภิกฺขุนี}ติ ฯเปฯ อยํ อิมสฺมิํ อตฺเถ อธิปฺเปตา ภิกฺขุนีติ.}

\prose{\textbf{กิสฺมิญฺจิเทว อธิกรเณ}ติ อธิกรณํ นาม จตฺตาริ อธิกรณานิ วิวาทาธิกรณํ อนุวาทาธิกรณํ อาปตฺตาธิกรณํ กิจฺจาธิกรณํ.}

\prose{\textbf{ปจฺจากตา} นาม ปราชิตา วุจฺจติ.}

\prose{\textbf{กุปิตา อนตฺตมนา}ติ อนภิรทฺธา อาหตจิตฺตา ขิลชาตา.}

\prose{\textbf{เอวํ วเทยฺยา}ติ “ฉนฺทคามินิโย จ ภิกฺขุนิโย ฯเปฯ ภยคามินิโย จ ภิกฺขุนิโย”ติ.}

\prose{\textbf{สา ภิกฺขุนี}ติ ยา สา เอวํวาทินี ภิกฺขุนี.}

\prose{\textbf{ภิกฺขุนีหี}ติ อญฺญาหิ ภิกฺขุนีหิ.}

\prose{ยา ปสฺสนฺติ, ยา สุณนฺติ, ตาหิ วตฺตพฺพา “มายฺเย กิสฺมิญฺจิเทว อธิกรเณ ปจฺจากตา กุปิตา อนตฺตมนา เอวํ อวจ ‘ฉนฺทคามินิโย จ ภิกฺขุนิโย ฯเปฯ ภยคามินิโย จ ภิกฺขุนิโย’ติ, อยฺยา โข ฉนฺทาปิ คจฺเฉยฺย ฯเปฯ ภยาปิ คจฺเฉยฺยา”ติ.\csromanspacer{} ทุติยมฺปิ วตฺตพฺพา.\csromanspacer{} ตติยมฺปิ วตฺตพฺพา.\csromanspacer{} สเจ ปฏินิสฺสชฺชติ, อิจฺเจตํ กุสลํ, โน เจ ปฏินิสฺสชฺชติ, อาปตฺติ ทุกฺกฏสฺส.\csromanspacer{} สุตฺวา น วทนฺติ, อาปตฺติ ทุกฺกฏสฺส.\csromanspacer{} สา ภิกฺขุนี สํฆมชฺฌมฺปิ อากฑฺฒิตฺวา วตฺตพฺพา “มายฺเย กิสฺมิญฺจิเทว อธิกรเณ ปจฺจากตา กุปิตา อนตฺตมนา เอวํ อวจ ‘ฉนฺทคามินิโย จ ภิกฺขุนิโย ฯเปฯ ภยคามินิโย จ ภิกฺขุนิโย’ติ, อยฺยา โข ฉนฺทาปิ คจฺเฉยฺย ฯเปฯ ภยาปิ คจฺเฉยฺยา”ติ.\csromanspacer{} ทุติยมฺปิ วตฺตพฺพา.\csromanspacer{} ตติยมฺปิ วตฺตพฺพา.\csromanspacer{} สเจ ปฏินิสฺสชฺชติ, อิจฺเจตํ กุสลํ, โน เจ ปฏินิสฺสชฺชติ, อาปตฺติ ทุกฺกฏสฺส.\csromanspacer{} สา ภิกฺขุนี สมนุภาสิตพฺพา, เอวญฺจ ปน ภิกฺขเว สมนุภาสิตพฺพา, พฺยตฺตาย ภิกฺขุนิยา ปฏิพลาย สํโฆ ญาเปตพฺโพ\csromandash{}}

\proseitem{718}{“สุณาตุ เม อยฺเย สํโฆ, อยํ อิตฺถนฺนามา ภิกฺขุนี กิสฺมิญฺจิเทว อธิกรเณ ปจฺจากตา กุปิตา อนตฺตมนา เอวํ วเทติ ‘ฉนฺทคามินิโย จ ภิกฺขุนิโย ฯเปฯ ภยคามินิโย จ ภิกฺขุนิโย’ติ, สา ตํ วตฺถุํ น ปฏินิสฺสชฺชติ, ยทิ สํฆสฺส ปตฺตกลฺลํ, สํโฆ อิตฺถนฺนามํ ภิกฺขุนิํ สมนุภาเสยฺย ตสฺส วตฺถุสฺส ปฏินิสฺสคฺคาย, เอสา ญตฺติ.}

\prose{\csromanfolio{309}สุณาตุ เม อยฺเย สํโฆ, อยํ อิตฺถนฺนามา ภิกฺขุนี กิสฺมิญฺจิเทว อธิกรเณ ปจฺจากตา กุปิตา อนตฺตมนา เอวํ วเทติ ‘ฉนฺทคามินิโย จ ภิกฺขุนิโย ฯเปฯ ภยคามินิโย จ ภิกฺขุนิโย’ติ, สา ตํ วตฺถุํ น ปฏินิสฺสชฺชติ, สํโฆ อิตฺถนฺนามํ ภิกฺขุนิํ สมนุภาสติ ตสฺส วตฺถุสฺส ปฏินิสฺสคฺคาย, ยสฺสา อยฺยาย ขมติ อิตฺถนฺนามาย ภิกฺขุนิยา สมนุภาสนา ตสฺส วตฺถุสฺส ปฏินิสฺสคฺคาย, สา ตุณฺหสฺส, ยสฺสา นกฺขมติ, สา ภาเสยฺย.}

\prose{ทุติยมฺปิ เอตมตฺถํ วทามิ ฯเปฯ.\csromanspacer{} ตติยมฺปิ เอตมตฺถํ วทามิ ฯเปฯ.}

\prose{สมนุภฏฺฐา สํเฆน อิตฺถนฺนามา ภิกฺขุนี ตสฺส วตฺถุสฺส ปฏินิสฺสคฺคาย, ขมติ สํฆสฺส, ตสฺมา ตุณฺหี, เอวเมตํ ธารยามี”ติ.}

\prose{ญตฺติยา ทุกฺกฏํ, ทฺวีหิ กมฺมวาจาหิ ถุลฺลจฺจยา, กมฺมวาจาปริโยสาเน อาปตฺติ สํฆาทิเสสสฺส.\csromanspacer{} สํฆาทิเสสํ อชฺฌาปชฺชนฺติยา ญตฺติยา ทุกฺกฏํ ทฺวีหิ กมฺมวาจาหิ ถุลฺลจฺจยา ปฏิปฺปสฺสมฺภนฺติ.}

\prose{\textbf{อยมฺปี}ติ ปุริมาโย อุปาทาย วุจฺจติ.}

\prose{\textbf{ยาวตติยกนฺ}ติ ยาวตติยํ สมนุภาสนาย อาปชฺชติ, น สห วตฺถุชฺฌาจารา.}

\prose{\textbf{นิสฺสรณียนฺ}ติ สํฆมฺหา นิสฺสารียติ.}

\prose{\textbf{สํฆาทิเสโส}ติ ฯเปฯ เตนปิ วุจฺจติ สํฆาทิเสโสติ.}

\proseitem{719}{ธมฺมกมฺเม ธมฺมกมฺมสญฺญา น ปฏินิสฺสชฺชติ, อาปตฺติ สํฆาทิเสสสฺส.\csromanspacer{} ธมฺมกมฺเม เวมติกา น ปฏินิสฺสชฺชติ, อาปตฺติ สํฆาทิเสสสฺส.\csromanspacer{} ธมฺมกมฺเม อธมฺมกมฺมสญฺญา น ปฏินิสฺสชฺชติ, อาปตฺติ สํฆาทิเสสสฺส.}

\prose{อธมฺมกมฺเม ธมฺมกมฺมสญฺญา, อาปตฺติ ทุกฺกฏสฺส.\csromanspacer{} อธมฺมกมฺเม เวมติกา, อาปตฺติ ทุกฺกฏสฺส.\csromanspacer{} อธมฺมกมฺเม อธมฺมกมฺมสญฺญา, อาปตฺติ ทุกฺกฏสฺส.}

\proseitemwithcloser{720}{อนาปตฺติ อสมนุภาสนฺติยา ปฏินิสฺสชฺชนฺติยา อุมฺมตฺติกาย อาทิกมฺมิกายาติ.}{อฏฺฐมสํฆาทิเสสสิกฺขาปทํ นิฏฺฐิตํ.}
\csromanmatikaanchor{mtk.132}
\csromantocmarkref{section}{9. นวมสํฆาทิเสสสิกฺขาปท}{mtk.132}
\bookmark[dest={mtk.132},level=1]{9. นวมสํฆาทิเสสสิกฺขาปท}
\setcsromanheadh{9. นวมสํฆาทิเสสสิกฺขาปท}
\csromanheader{1}{9. \textbf{นวมสํฆาทิเสสสิกฺขาปท}}
\proseitem{721}{\csromanfolio{310}เตน สมเยน พุทฺโธ ภควา สาวตฺถิยํ วิหรติ เชตวเน อนาถปิณฺฑิกสฺส อาราเม.\csromanspacer{} เตน โข ปน สมเยน ถุลฺลนนฺทาย ภิกฺขุนิยา อนฺเตวาสิกา ภิกฺขุนิโย สํสฏฺฐา วิหรนฺติ ปาปาจารา ปาปสทฺทา ปาปสิโลกา ภิกฺขุนิสํฆสฺส วิเหสิกา อญฺญมญฺญิสฺสา วชฺชปฺปฏิจฺฉาทิกา.\csromanspacer{} ยา ตา ภิกฺขุนิโย อปฺปิจฺฉา ฯเปฯ ตา อุชฺฌายนฺติ ขิยฺยนฺติ วิปาเจนฺติ “กถํ หิ นาม ภิกฺขุนิโย สํสฏฺฐา วิหริสฺสนฺติ ปาปาจารา ปาปสทฺทา ปาปสิโลกา ภิกฺขุนิสํฆสฺส วิเหสิกา อญฺญมญฺญิสฺสา วชฺชปฺปฏิจฺฉาทิกา”ติ ฯเปฯ.\csromanspacer{} สจฺจํ กิร ภิกฺขเว ภิกฺขุนิโย สํสฏฺฐา วิหรนฺติ ปาปาจารา ปาปสทฺทา ปาปสิโลกา ภิกฺขุนิสํฆสฺส วิเหสิกา อญฺญมญฺญิสฺสา วชฺชปฺปฏิจฺฉาทิกาติ.\csromanspacer{} สจฺจํ ภควาติ.\csromanspacer{} วิครหิ พุทฺโธ ภควา ฯเปฯ.\csromanspacer{} กถํ หิ นาม ภิกฺขเว ภิกฺขุนิโย สํสฏฺฐา วิหริสฺสนฺติ ปาปาจารา ปาปสทฺทา ปาปสิโลกา ภิกฺขุนิสํฆสฺส วิเหสิกา อญฺญมญฺญิสฺสา วชฺชปฺปฏิจฺฉาทิกา.\csromanspacer{} เนตํ ภิกฺขเว อปฺปสนฺนานํ วา ปสาทาย ฯเปฯ.\csromanspacer{} เอวญฺจ ปน ภิกฺขเว ภิกฺขุนิโย อิมํ สิกฺขาปทํ อุทฺทิสนฺตุ\csromandash{}}

\proseitem{722}{\textbf{“ภิกฺขุนิโย ปเนว สํสฏฺฐา วิหรนฺติ ปาปาจารา ปาปสทฺทา ปาปสิโลกา ภิกฺขุนิสํฆสฺส วิเหสิกา อญฺญมญฺญิสฺสา วชฺชปฺปฏิจฺฉาทิกา, ตา ภิกฺขุนิโย ภิกฺขุนีหิ เอวมสฺสุ วจนียา ‘ภคินิโย โข สํสฏฺฐา วิหรนฺติ ปาปาจารา ปาปสทฺทา ปาปสิโลกา ภิกฺขุนิสํฆสฺส วิเหสิกา อญฺญมญฺญิสฺสา วชฺชปฺปฏิจฺฉาทิกา, วิวิจฺจถายฺเย, วิเวกญฺเญว ภคินีนํ สํโฆ วณฺเณตี’ติ.}\csromanspacer{} เอวญฺจ ตา ภิกฺขุนิโย ภิกฺขุนีหิ วุจฺจมานา ตเถว ปคฺคเณฺหยฺยุํ, ตา ภิกฺขุนิโย ภิกฺขุนีหิ ยาวตติยํ สมนุภาสิตพฺพา ตสฺส ปฏินิสฺสคฺคาย, ยาวตติยญฺเจ สมนุภาสียมานา ตํ ปฏินิสฺสชฺเชยฺยุํ, อิจฺเจตํ กุสลํ, โน เจ ปฏินิสฺสชฺเชยฺยุํ, อิมาปิ ภิกฺขุนิโย ยาวตติยกํ ธมฺมํ อาปนฺนา นิสฺสารณียํ สํฆาทิเสสนฺ”ติ.}

\proseitem{723}{\textbf{ภิกฺขุนิโย ปเนวา}ติ อุปสมฺปนฺนาโย วุจฺจนฺติ.}

\prose{\textbf{สํสฏฺฐา วิหรนฺตี}ติ สํสฏฺฐา นาม อนนุโลมิเกน กายิกวาจสิเกน สํสฏฺฐา วิหรนฺติ.}

\prose{\csromanfolio{311}\textbf{ปาปาจารา}ติ ปาปเกน อาจาเรน สมนฺนาคตา.}

\prose{\textbf{ปาปสทฺทา}ติ ปาปเกน กิตฺติสทฺเทน อพฺภุคฺคตา.}

\prose{\textbf{ปาปสิโลกา}ติ ปาปเกน มิจฺฉาชีเวน ชีวิตํ กปฺเปนฺติ.}

\prose{\textbf{ภิกฺขุนิสํฆสฺส วิเหสิกา}ติ อญฺญมญฺญิสฺสา กมฺเม กรียมาเน ปฏิกฺโกสนฺติ.}

\prose{\textbf{อญฺญมญฺญิสฺสา วชฺชปฺปฏิจฺฉาทิกา}ติ อญฺญมญฺญํ วชฺชํ ปฏิจฺฉาเทนฺติ.}

\prose{\textbf{ตา ภิกฺขุนิโย}ติ ยา ตา สํสฏฺฐา ภิกฺขุนิโย.}

\prose{\textbf{ภิกฺขุนีหี}ติ อญฺญาหิ ภิกฺขุนีหิ.}

\prose{ยา ปสฺสนฺติ, ยา สุณนฺติ, ตาหิ วตฺตพฺพา “ภคินิโย โข สํสฏฺฐา วิหรนฺติ ปาปาจารา ปาปสทฺทา ปาปสิโลกา ภิกฺขุนิสํฆสฺส วิเหสิกา อญฺญมญฺญิสฺสา วชฺชปฺปฏิจฺฉาทิกา, วิวิจฺจถายฺเย, วิเวกญฺเญว ภคินีนํ สํโฆ วณฺเณตี”ติ.\csromanspacer{} ทุติยมฺปิ วตฺตพฺพา.\csromanspacer{} ตติยมฺปิ วตฺตพฺพา.\csromanspacer{} สเจ ปฏินิสฺสชฺชนฺติ, อิจฺเจตํ กุสลํ, โน เจ ปฏินิสฺสชฺชนฺติ, อาปตฺติ ทุกฺกฏสฺส.\csromanspacer{} สุตฺวา น วทนฺติ, อาปตฺติ ทุกฺกฏสฺส.\csromanspacer{} ตา ภิกฺขุนิโย สํฆมชฺฌมฺปิ อากฑฺฒิตฺวา วตฺตพฺพา “ภคินิโย โข สํสฏฺฐา วิหรนฺติ ปาปาจารา ปาปสทฺทา ปาปสิโลกา ภิกฺขุนิสํฆสฺส วิเหสิกา อญฺญมญฺญิสฺสา วชฺชปฺปฏิจฺฉาทิกา, วิวิจฺจถายฺเย, วิเวกญฺเญว ภคินีนํ สํโฆ วณฺเณตี”ติ.\csromanspacer{} ทุติยมฺปิ วตฺตพฺพา.\csromanspacer{} ตติยมฺปิ วตฺตพฺพา.\csromanspacer{} สเจ ปฏินิสฺสชฺชนฺติ, อิจฺเจตํ กุสลํ, โน เจ ปฏินิสฺสชฺชนฺติ, อาปตฺติ ทุกฺกฏสฺส.\csromanspacer{} ตา ภิกฺขุนิโย สมนุภาสิตพฺพา, เอวญฺจ ปน ภิกฺขเว สมนุภาสิตพฺพา, พฺยตฺตาย ภิกฺขุนิยา ปฏิพลาย สํโฆ ญาเปตพฺโพ\csromandash{}}

\proseitem{724}{“สุณาตุ เม อยฺเย สํโฆ, อิตฺถนฺนามา จ อิตฺถนฺนามา จ ภิกฺขุนิโย สํสฏฺฐา วิหรนฺติ ปาปาจารา ปาปสทฺทา ปาปสิโลกา ภิกฺขุนิสํฆสฺส วิเหสิกา อญฺญมญฺญิสฺสา วชฺชปฺปฏิจฺฉาทิกา, ตา ตํ วตฺถุํ น ปฏินิสฺสชฺชนฺติ, ยทิ สํฆสฺส ปตฺตกลฺลํ, สํโฆ อิตฺถนฺนามญฺจ อิตฺถนฺนามญฺจ ภิกฺขุนิโย สมนุภาเสยฺย ตสฺส วตฺถุสฺส ปฏินิสฺสคฺคาย, เอสา ญตฺติ.}

\prose{สุณาตุ เม อยฺเย สํโฆ, อิตฺถนฺนามา จ อิตฺถนฺนามา จ ภิกฺขุนิโย สํสฏฺฐา วิหรนฺติ ปาปาจารา ปาปสทฺทา ปาปสิโลกา ภิกฺขุนิสํฆสฺส วิเหสิกา อญฺญมญฺญิสฺสา วชฺชปฺปฏิจฺฉาทิกา, ตา ตํ วตฺถุํ น ปฏินิสฺสชฺชนฺติ, \csromanfolio{312}\textbf{สํโฆ อิตฺถนฺนามญฺจ อิตฺถนฺนามญฺจ ภิกฺขุนิโย สมนุภาสติ ตสฺส} \textbf{วตฺถุสฺส ปฏินิสฺสคฺคาย, ยสฺสา อยฺยาย ขมติ อิตฺถนฺนามาย จ} อิตฺถนฺนามาย จ ภิกฺขุนีนํ สมนุภาสนา ตสฺส วตฺถุสฺส ปฏินิสฺสคฺคาย, สา ตุณฺหสฺส, ยสฺสา นกฺขมติ, สา ภาเสยฺย.}

\prose{ทุติยมฺปิ เอตมตฺถํ วทามิ ฯเปฯ.\csromanspacer{} ตติยมฺปิ เอตมตฺถํ วทามิ ฯเปฯ.}

\prose{สมนุภฏฺฐา สํเฆน อิตฺถนฺนามา จ อิตฺถนฺนามา จ ภิกฺขุนิโย ตสฺส วตฺถุสฺส ปฏินิสฺสคฺคาย, ขมติ สํฆสฺส, ตสฺมา ตุณฺหี, เอวเมตํ ธารยามี”ติ.}

\prose{ญตฺติยา ทุกฺกฏํ, ทฺวีหิ กมฺมวาจาหิ ถุลฺลจฺจยา, กมฺมวาจาปริโยสาเน อาปตฺติ สํฆาทิเสสสฺส.\csromanspacer{} สํฆาทิเสสํ อชฺฌาปชฺชนฺตีนํ ญตฺติยา ทุกฺกฏํ ทฺวีหิ กมฺมวาจาติ ถุลฺลจฺจยา ปฏิปฺปสฺสมฺภนฺติ, เทฺว ติสฺโส เอกโต สมนุภาสิตพฺพา, ตตุตฺตริ น สมนุภาสิตพฺพา.}

\prose{\textbf{อิมาปิ ภิกฺขุนิโย}ติ ปุริมาโย อุปาทาย วุจฺจนฺติ.}

\prose{\textbf{ยาวตติยกนฺ}ติ ยาวรติยํ สมนุภาสนาย อาปชฺชนฺติ, น สห วตฺถุชฺฌาจารา.}

\prose{\textbf{นิสฺสารณียนฺ}ติ สํฆมฺหา นิสฺสารียติ.}

\prose{\textbf{สํฆาทิเสโส}ติ ฯเปฯ เตนปิ วุจฺจติ สํฆาทิเสโสติ.}

\proseitem{725}{ธมฺมกมฺเม ธมฺมกมฺมสญฺญา น ปฏินิสฺสชฺชนฺติ, อาปตฺติ สํฆาทิเสสสฺส.\csromanspacer{} ธมฺมกมฺเม เวมติกา น ปฏินิสฺสชฺชนฺติ, อาปตฺติ สํฆาทิเสสสฺส.\csromanspacer{} ธมฺมกมฺเม อธมฺมกมฺมสญฺญา น ปฏินิสฺสชฺชนฺติ, อาปตฺติ สํฆาทิเสสสฺส.}

\prose{อธมฺมกมฺเม ธมฺมกมฺมสญฺญา, อาปตฺติ ทุกฺกฏสฺส.\csromanspacer{} อธมฺมกมฺเม เวมติกา, อาปตฺติ ทุกฺกฏสฺส.\csromanspacer{} อธมฺมกมฺเม อธมฺมกมฺมสญฺญา, อาปตฺติ ทุกฺกฏสฺส.}

\proseitemwithcloser{726}{อนาปตฺติ อสมนุภาสนฺตีนํ ปฏินิสฺสชฺชนฺตีนํ อุมฺมตฺติกานํ อาทิกมฺมิกานนฺติ.}{นวมสํฆาทิเสสสิกฺขาปทํ นิฏฺฐิตํ.}
\csromanmatikaanchor{mtk.133}
\csromantocmarkref{section}{10. ทสมสํฆาทิเสสสิกฺขาปท}{mtk.133}
\bookmark[dest={mtk.133},level=1]{10. ทสมสํฆาทิเสสสิกฺขาปท}
\setcsromanheadh{10. ทสมสํฆาทิเสสสิกฺขาปท}
\csromanheader{1}{10. \textbf{ทสมสํฆาทิเสสสิกฺขาปท}}
\proseitem{727}{\csromanfolio{313}เตน สมเยน พุทฺโธ ภควา สาวตฺถิยํ วิหรติ เชตวเน อนาถปิณฺฑิกสฺส อาราเม.\csromanspacer{} เตน โข ปน สมเยน ถุลฺลนนฺทา ภิกฺขุนี สํเฆน สมนุภฏฺฐา ภิกฺขุนิโย เอวํ วเทติ “สํสฏฺฐาว อยฺเย ตุเมฺห วิหรถ, มา ตุเมฺห นานา วิหริตฺถ, สนฺติ สํเฆ อญฺญาปิ ภิกฺขุนิโย เอวาจารา เอวํสทฺทา เอวํสิโลกา ภิกฺขุนิสํฆสฺส วิเหสิกา อญฺญมญฺญิสฺสา วชฺชปฺปฏิจฺฉาทิกา, ตา สํโฆ น กิญฺจิ อาห, ตุมฺหญฺเญว\footnote{ตุเมฺหเยว (สฺยา)} สํโฆ อุญฺญาย ปริภเวน อกฺขนฺติยา เวภสฺสิยา\footnote{เวภสฺสา (สี, สฺยา)} ทุพฺพลฺยา เอวมาห ‘ภคินิโย โข สํสฏฺฐา วิหรนฺติ ปาปาจารา ปาปสทฺทา ปาปสิโลกา ภิกฺขุนิสํฆสฺส วิเหสิกา อญฺญมญฺญิสฺสา วชฺชปฺปฏิจฺฉาทิกา, วิวิจฺจถายฺเย, วิเวกญฺเญว ภคินีนํ สํโฆ วณฺเณตี’ติ”.\csromanspacer{} ยา ตา ภิกฺขุนิโย อปฺปิจฺฉา ฯเปฯ ตา อุชฺฌายนฺติ ขิยฺยนฺติ วิปาเจนฺติ “กถํ หิ นาม อยฺยา ถุลฺลนนฺทา สํเฆน สมนุภฏฺฐา ภิกฺขุนิโย เอวํ วกฺขติ ‘สํสฏฺฐาว อยฺเย ตุเมฺห วิหรถ, มา ตุเมฺห นานา วิหริตฺถ, สนฺติ สํเฆ อญฺญาปิ ภิกฺขุนิโย ฯเปฯ วิวิจฺจถายฺเย, วิเวกญฺเญว ภคินีนํ สํโฆ วณฺเณตี’ติ” ฯเปฯ.\csromanspacer{} สจฺจํ กิร ภิกฺขเว ถุลฺลนนฺทา ภิกฺขุนี สํเฆน สมนุภฏฺฐา ภิกฺขุนิโย เอวํ วเทติ “สํสฏฺฐาว อยฺเย ตุเมฺห วิหรถ, มา ตุเมฺห นานา วิหรตฺถ, สนฺติ สํเฆ อญฺญาปิ ภิกฺขุนิโย เอวาจารา เอวํสทฺทา เอวํสิโลกา ภิกฺขุนิสํฆสฺส วิเหสิกา อญฺญมญฺญิสฺสา วชฺชปฺปฏิจฺฉาทิกา, ตา สํโฆ น กิญฺจิ อาห, ตุมฺหญฺเญว สํโฆ อุญฺญาย ปริภเวน อกฺขนฺติยา เวภสฺสิยา ทุพฺพลฺยา เอวมาห ‘ภคินิโย โข สํสฏฺฐา วิหรนฺติ ปาปาจารา ปาปสทฺทา ปาปสิโลกา ภิกฺขุนิสํฆสฺส วิเหสิกา อญฺญมญฺญิสฺสา วชฺชปฺปฏิจฺฉาทิกา, วิวิจฺจถายฺเย, วิเวกญฺเญว ภคินีนํ สํโฆ วณฺเณตี’ติ”.\csromanspacer{} สจฺจํ ภควาติ.\csromanspacer{} วิครหิ พุทฺโธ ภควา ฯเปฯ.\csromanspacer{} กถํ หิ นาม ภิกฺขเว ถุลฺลนนฺทา ภิกฺขุนี สํเฆน สมนุภฏฺฐา ภิกฺขุนิโย เอวํ วกฺขติ “สํสฏฺฐาว อยฺเย ตุเมฺห วิหรถ, มา ตุเมฺห นานา วิหริตฺถ, สนฺติ สํเฆ อญฺญาปิ ภิกฺขุนิโย ฯเปฯ วิวิจฺจถายฺเย, วิเวกญฺเญว ภคินีนํ สํโฆ วณฺเณตี”ติ.\csromanspacer{} เนตํ ภิกฺขเว อปฺปสนฺนานํ วา ปสาทาย ฯเปฯ.\csromanspacer{} เอวญฺจ ปน ภิกฺขเว ภิกฺขุนิโย อิมํ สิกฺขาปทํ อุทฺทิสนฺตุ\csromandash{}}

\proseitem{728}{\csromanfolio{314}“ยา ปน ภิกฺขุนี เอวํ วเทยฺย ‘สํสฏฺฐาว อยฺเย ตุเมฺห วิหรถ, มา ตุเมฺห นานา วิหริตฺถ, สนฺติ สํเฆ อญฺญาปิ ภิกฺขุนิโย เอวาจารา เอวํสทฺทา เอวํสิโลกา ภิกฺขุนิสํฆสฺส วิเหสิกา อญฺญมญฺญิสฺสา วชฺชปฺปฏิจฺฉาทิกา, ตา สํโฆ น กิญฺจิ อาห, ตุมฺหญฺเญว สํโฆ อุญฺญาย ปริภเวน อกฺขนฺติยา เวภสฺสิยา ทุพฺพลฺยา เอวมาห ภคินิโย โข สํสฏฺฐา วิหรนฺติ ปาปาจารา ปาปสทฺทา ปาปสิโลกา ภิกฺขุนิสํฆสฺส วิเหสิกา อญฺญมญฺญิสฺสา วชฺชปฺปฏิจฺฉาทิกา, วิวิจฺจถายฺเย, วิเวกญฺเญว ภคินีนํ สํโฆ วณฺเณตี’ติ.\csromanspacer{}\textbf{ สา ภิกฺขุนี ภิกฺขุนีหิ เอวมสฺส วจนียา ‘มา อยฺเย เอวํ อวจ สํสฏฺฐาว อยฺเย ตุเมฺห วิหรถ, มา ตุเมฺห นานา วิหริตฺถ, สนฺติ สํเฆ อญฺญาปิ ภิกฺขุนิโย เอวาจารา เอวํสทฺทา เอวํสิโลกา ภิกฺขุนิสํฆสฺส วิเหสิกา, อญฺญมญฺญิสฺสา วชฺชปฺปฏิจฺฉาทิกา, ตา สํโฆ น กิญฺจิ อาห, ตุมฺหญฺเญว สํโฆ อุญฺญาย ปริภเวน อกฺขนฺติยา เวภสฺสิยา ทุพฺพลฺยา เอวมาห ภคินิโย โข สํสฏฺฐา วิหรนฺติ ปาปาจารา ปาปสทฺทา ปาปสิโลกา ภิกฺขุนิสํฆสฺส วิเหสิกา อญฺญมญฺญิสฺสา วชฺชปฺปฏิจฺฉาทิกา, วิวิจฺจถายฺเย, วิเวกญฺเญว ภคินีนํ สํโฆ วณฺเณตี’ติ.}\csromanspacer{} เอวญฺจ สา ภิกฺขุนี ภิกฺขุนีหิ วุจฺจมานา ตเถว ปคฺคเณฺหยฺย, สา ภิกฺขุนี ภิกฺขุนีหิ ยาวตติยํ สมนุภาสิตพฺพา ตสฺส ปฏินิสฺสคฺคาย, ยาวตติยญฺเจ สมนุภาสียมานา ตํ ปฏินิสฺสชฺเชยฺย, อิจฺเจตํ กุสลํ, โน เจ ปฏินิสฺสชฺเชยฺย, อยมฺปิ ภิกฺขุนี ยาวตติยกํ ธมฺมํ อาปนฺนา นิสฺสารณียํ สํฆาทิเสสนฺ”ติ.}

\proseitem{729}{\textbf{ยา ปนา}ติ ยา ยาทิสา ฯเปฯ.\csromanspacer{} \textbf{ภิกฺขุนี}ติ ฯเปฯ อยํ อิมสฺมิํ อตฺเถ อธิปฺเปตา ภิกฺขุนีติ.}

\prose{\textbf{เอวํ วเทยฺยา}ติ “สํสฏฺฐาว อยฺเย ตุเมฺห วิหรถ, มา ตุเมฺห นานา วิหริตฺถ, สนฺติ สํเฆ อญฺญาปิ ภิกฺขุนิโย เอวาจารา เอวํสทฺทา เอวํสิโลกา ภิกฺขุนิสํฆสฺส วิเหสิกา อญฺญมญฺญิสฺสา วชฺชปฺปฏิจฺฉาทิกา, ตา สํโฆ น กิญฺจิ อาห”.}

\prose{\textbf{ตุมฺหญฺเญว สํโฆ อุญฺญายา}ติ อวญฺญาย.}

\prose{\textbf{ปริภเวนา}ติ ปาริภพฺยตา.}

\prose{\textbf{อกฺขนฺติยา}ติ โกเปน.}

\prose{\csromanfolio{315}\textbf{เวภสฺสิยา}ติ วิภสฺสีกตา\footnote{วิภสฺสิกตาย (สี)}.}

\prose{\textbf{ทุพฺพลฺยา}ติ อปกฺขตา.}

\prose{เอวมาห “ภคินิโย โข สํสฏฺฐา วิหรนฺติ ปาปาจารา ปาปสทฺทา ปาปสิโลกา ภิกฺขุนิสํฆสฺส วิเหสิกา อญฺญมญฺญิสฺสา วชฺชปฺปฏิจฺฉาทิกา, วิวิจฺจถายฺเย, วิเวกญฺเญว ภคินีนํ สํโฆ วณฺเณตี”ติ.}

\prose{\textbf{สา ภิกฺขุนี}ติ ยา สา เอวํวาทินี ภิกฺขุนี.}

\prose{\textbf{ภิกฺขุนีหี}ติ อญฺญาหิ ภิกฺขุนีหิ.}

\prose{ยา ปสฺสนฺติ, ยา สุณนฺติ, ตาหิ วตฺตพฺพา “มายฺเย เอวํ อวจ ‘สํสฏฺฐาว อยฺเย ตุเมฺห วิหรถ, มา ตุเมฺห นานา วิหริตฺถ, สนฺติ สํเฆ อญฺญาปิ ภิกฺขุนิโย ฯเปฯ วิวิจฺจถายฺเย, วิเวกญฺเญว ภคินีนํ สํโฆ วณฺเณตี’ติ”.\csromanspacer{} ทุติยมฺปิ วตฺตพฺพา.\csromanspacer{} ตติยมฺปิ วตฺตพฺพา.\csromanspacer{} สเจ ปฏินิสฺสชฺชติ, อิจฺเจตํ กุสลํ, โน เจ ปฏินิสฺสชฺชติ, อาปตฺติ ทุกฺกฏสฺส.\csromanspacer{} สุตฺวา น วทนฺติ, อาปตฺติ ทุกฺกฏสฺส.\csromanspacer{} สา ภิกฺขุนี สํฆมชฺฌมฺปิ อากฑฺฒิตฺวา วตฺตพฺพา “มายฺเย เอวํ อวจ ‘สํสฏฺฐาว อยฺเย ตุเมฺห วิหรถ, มา ตุเมฺห นานา วิหริตฺถ, สนฺติ สํเฆ อญฺญาปิ ภิกฺขุนิโย ฯเปฯ วิวิจฺจถายฺเย, วิเวกญฺเญว ภคินีนํ สํโฆ วณฺเณตี’ติ”.\csromanspacer{} ทุติยมฺปิ วตฺตพฺพา.\csromanspacer{} ตติยมฺปิ วตฺตพฺพา.\csromanspacer{} สเจ ปฏินิสฺสชฺชติ, อิจฺเจตํ กุสลํ, โน เจ ปฏินิสฺสชฺชติ, อาปตฺติ ทุกฺกฏสฺส.\csromanspacer{} สา ภิกฺขุนี สมนุภาสิตพฺพา, เอวญฺจ ปน ภิกฺขเว สมนุภาสิตพฺพา, พฺยตฺตาย ภิกฺขุนิยา ปฏิพลาย สํโฆ ญาเปตพฺโพ\csromandash{}}

\proseitem{730}{สุณาตุ เม อยฺเย สํโฆ, อยํ อิตฺถนฺนามา ภิกฺขุนี สํเฆน สมนุภฏฺฐา ภิกฺขุนิโย เอวํ วเทติ ‘สํสฏฺฐาว อยฺเย ตุเมฺห วิหรถ, มา ตุเมฺห นานา วิหริตฺถ, สนฺติ สํเฆ อญฺญาปิ ภิกฺขุนิโย เอวาจารา เอวํสทฺทา เอวํสิโลกา ภิกฺขุนิสํฆสฺส วิเหสิกา อญฺญฺมญฺญิสฺสา วชฺชปฺปฏิจฺฉาทิกา, ตา สํโฆ น กิญฺจิ อาห, ตุมฺหญฺเญว สํโฆ อุญฺญาย ปริภเวน อกฺขนฺติยา เวภสฺสิยา ทุพฺพลฺยา เอวมาห ‘ภคินิโย โข สํสฏฺฐา วิหรนฺติ ปาปาจารา ปาปสทฺทา ปาปสิโลกา ภิกฺขุนิสํฆสฺส วิเหสิกา อญฺญมญฺญิสฺสา วชฺชปฺปฏิจฺฉาทิกา, วิวิจฺจถายฺเย, วิเวกญฺเญว ภคินีนํ สํโฆ วณฺเณตี’ติ’.\csromanspacer{} สา ตํ วตฺถุํ น ปฏินิสฺสชฺชติ, ยทิ สํฆสฺส ปตฺตกลฺลํ, สํโฆ อิตฺถนฺนามํ ภิกฺขุนิํ สมนุภาเสยฺย ตสฺส วตฺถุสฺส ปฏินิสฺสคฺคาย, เอสา ญตฺติ.}

\prose{\csromanfolio{316}สุณาตุ เม อยฺเย สํโฆ, อยํ อิตฺถนฺนามา ภิกฺขุนี สํเฆน สมนุภฏฺฐา ภิกฺขุนิโย เอวํ วเทติ ‘สํสฏฺฐาว อยฺเย ตุเมฺห วิหรถ, มา ตุเมฺห นานา วิหริตฺถ, สนฺติ สํเฆ อญฺญาปิ ภิกฺขุนิโย เอวาจารา เอวํสทฺทา เอวํสิโลกา ภิกฺขุนิสํฆสฺส วิเหสิกา อญฺญมญฺญิสฺสา วชฺชปฺปฏิจฺฉาทิกา, ตา สํโฆ น กิญฺจิ อาห, ตุมฺหญฺเญว สํโฆ อุญฺญาย ปริภเวน อกฺขนฺติยา เวภสฺสิยา ทุพฺพลฺยา เอวมาห ‘ภคินิโย โข สํสฏฺฐา วิหรนฺติ ปาปาจารา ปาปสทฺทา ปาปสิโลกา ภิกฺขุนิสํฆสฺส วิเหสิกา อญฺญมญฺญิสฺสา วชฺชปฺปฏิจฺฉาทิกา, วิวิจฺจถายฺเย, วิเวกญฺเญว ภคินีนํ สํโฆ วณฺเณตี’ติ’.\csromanspacer{} สา ตํ วตฺถุํ น ปฏินิสฺสชฺชติ, สํโฆ อิตฺถนฺนามํ ภิกฺขุนิํ สมนุภาสติ ตสฺส วตฺถุสฺส ปฏินิสฺสคฺคาย, ยสฺสา อยฺยาย ขมติ อิตฺถนฺนามาย ภิกฺขุนิยา สมนุภาสนา ตสฺส วตฺถุสฺส ปฏินิสฺสคฺคาย, สา ตุณฺหสฺส, ยสฺสา นกฺขมติ, สา ภาเสยฺย.}

\prose{ทุติยมฺปิ เอตมตฺถํ วทามิ ฯเปฯ.\csromanspacer{} ตติยมฺปิ เอตมตฺถํ วทามิ ฯเปฯ.}

\prose{สมนุภฏฺฐา สํเฆน อิตฺถนฺนามา ภิกฺขุนี ตสฺส วตฺถุสฺส ปฏินิสฺสคฺคาย ขมติ สํฆสฺส, ตสฺมา ตุณฺหี, เอวเมตํ ธารยามี”ติ.}

\prose{ญตฺติยา ทุกฺกฏํ, ทฺวีหิ กมฺมวาจาหิ ถุลฺลจฺจยา, กมฺมวาจาปริโยสาเน อาปตฺติ สํฆาทิเสสสฺส.\csromanspacer{} สํฆาทิเสสํ อชฺฌาปชฺชนฺติยา ญตฺติยา ทุกฺกฏํ ทฺวีหิ กมฺมวาจาหิ ถุลฺลจฺจยา ปฏิปฺปสฺสมฺภนฺติ.}

\prose{\textbf{อยมฺปี}ติ ปุริมาโย อุปาทาย วุจฺจติ.}

\prose{\textbf{ยาวตติยกนฺ}ติ ยาวตติยํ สมนุภาสนาย อาปชฺชติ, น สห วตฺถุชฺฌาจารา.}

\prose{\textbf{นิสฺสารณียนฺ}ติ สํฆมฺหา นิสฺสารียติ.}

\prose{\textbf{สํฆาทิเสโส}ติ สํโฆว ตสฺสา อาปตฺติยา มานตฺตํ เทติ มูลาย ปฏิกสฺสติ อพฺเภติ, น สมฺพหุลา, น เอกา ภิกฺขุนี, เตน วุจฺจติ “สํฆาทิเสโส”ติ.\csromanspacer{} ตสฺเสว อาปตฺตินิกายสฺส นามกมฺมํ อธิวจนํ, เตนปิ วุจฺจติ “สํฆาทิเสโส”ติ.}

\proseitem{731}{ธมฺมกมฺเม ธมฺมกมฺมสญฺญา น ปฏินิสฺสชฺชติ, อาปตฺติ สํฆาทิเสสสฺส.\csromanspacer{} ธมฺมกมฺเม เวมติกา น ปฏินิสฺสชฺชติ, อาปตฺติ สํฆาทิเสสสฺส.\csromanspacer{} ธมฺมกมฺเม อธมฺมกมฺมสญฺญา น ปฏินิสฺสชฺชติ, อาปตฺติ สํฆาทิเสสสฺส.}

\prose{\csromanfolio{317}อธมฺมกมฺเม ธมฺมกมฺมสญฺญา, อาปตฺติ ทุกฺกฏสฺส.\csromanspacer{} อธมฺมกมฺเม เวมติกา, อาปตฺติ ทุกฺกฏสฺส.\csromanspacer{} อธมฺมกมฺเม อธมฺมกมฺมสญฺญา, อาปตฺติ ทุกฺกฏสฺส.}

\proseitemwithcloser{732}{อนาปตฺติ อสมนุภาสนฺติยา ปฏินิสฺสชฺชนฺติยา อุมฺมตฺติกาย อาทิกมฺมิกายาติ.}{ทสมสํฆาทิเสสสิกฺขาปทํ นิฏฺฐิตํ.}
\prose{อุทฺทิฏฺฐา โข อยฺยาโย สตฺตรส สํฆาทิเสสา ธมฺมา นว ปฐมาปตฺติกา อฏฺฐ ยาวตติยกา, เยสํ ภิกฺขุนี อญฺญตรํ วา อญฺญตรํ วา อาปชฺชติ, ตาย ภิกฺขุนิยา อุภโตสํเฆ ปกฺขมานตฺตํ จริตพฺพํ.\csromanspacer{} จิณฺณมานตฺตา ภิกฺขุนี ยตฺถ สิยา วีสติคโณ ภิกฺขุนิสํโฆ, ตตฺถ (\csromansymbolfootnote{( )}{(กตฺถจิ นตฺถิ)}สา ภิกฺขุนี) อพฺเภตพฺพา.\csromanspacer{} เอกายปิ เจ อูโน วีสติคโณ ภิกฺขุนิสํโฆ ตํ ภิกฺขุนิํ อพฺเภยฺย, สา จ ภิกฺขุนี อนพฺภิตา, ตา จ ภิกฺขุนิโย คารยฺหา, อยํ ตตฺถ สามีจิ.}

\prosewithcloser{ตตฺถายฺยาโย ปุจฺฉามิ กจฺจิตฺถ ปริสุทฺธา, ทุติยมฺปิ ปุจฺฉามิ กจฺจิตฺถ ปริสุทฺธา, ตติยมฺปิ ปุจฺฉามิ กจฺจิตฺถ ปริสุทฺธา, ปริสุทฺเธตฺถายฺยาโย, ตสฺมา ตุณฺหี, เอวเมตํ ธารยามีติ.}{สตฺตรสกํ นิฏฺฐิตํ.}
\csromanheader{2}{\textbf{ภิกฺขุนิวิภงฺเค}}
\nitthitam{\textbf{สํฆาทิเสสกณฺฑํ นิฏฺฐิตํ.}}
\csromanmatikaanchor{mtk.134}
\csromantocmarknopage{chapter}{3. นิสฺสคฺคิยกณฺฑ}{mtk.134}
\bookmark[dest={mtk.134},level=0]{3. นิสฺสคฺคิยกณฺฑ}
\setcsromanheadcha{3. นิสฺสคฺคิยกณฺฑ}
\setcsromanheadhclear{}
\chapterheadpageread{3. \textbf{นิสฺสคฺคิยกณฺฑ}}
\csromanmatikaanchor{mtk.135}
\csromanmatikaanchor{mtk.136}
\csromantocmarknopage{section}{1. ปตฺตวคฺค}{mtk.135}
\bookmark[dest={mtk.135},level=1]{1. ปตฺตวคฺค}
\csromantocmarkref{subsection}{1. ปฐมสิกฺขาปท}{mtk.136}
\bookmark[dest={mtk.136},level=2]{1. ปฐมสิกฺขาปท}
\setcsromanheadh{1. ปตฺตวคฺค}
\csromanheader{2}{1. \textbf{ปตฺตวคฺค 1. ปฐมสิกฺขาปท}}
\prose{\csromanfolio{318}อิเม โข ปนายฺยาโย ติํส นิสฺสคฺคิยา ปาจิตฺติยา ธมฺมา อุทฺเทสํ อาคจฺฉนฺติ.}

\proseitem{733}{เตน สมเยน พุทฺโธ ภควา สาวตฺถิยํ วิหรติ เชตวเน อนาถปิณฺฑิกสฺส อาราเม.\csromanspacer{} เตน โข ปน สมเยน ฉพฺพคฺคิยา ภิกฺขุนิโย พหู ปตฺเต สนฺนิจยํ กโรนฺติ.\csromanspacer{} มนุสฺสา วิหารจาริกํ อาหิณฺฑนฺตา ปสฺสิตฺวา อุชฺฌายนฺติ ขิยฺยนฺติ วิปาเจนฺติ “กถํ หิ นาม ภิกฺขุนิโย พหู ปตฺเต สนฺนิจยํ กริสฺสนฺติ, ปตฺตวาณิชฺชํ วา ภิกฺขุนิโย กริสฺสนฺติ, อามตฺติกาปณํ วา ปสาเรสฺสนฺตี”ติ.\csromanspacer{} อสฺโสสุํ โข ภิกฺขุนิโย เตสํ มนุสฺสานํ อุชฺฌายนฺตานํ ขิยฺยนฺตานํ วิปาเจนฺตานํ.\csromanspacer{} ยา ตา ภิกฺขุนิโย อปฺปิจฺฉา ฯเปฯ ตา อุชฺฌายนฺติ ขิยฺยนฺติ วิปาเจนฺติ “กถํ หิ นาม ฉพฺพคฺคิยา ภิกฺขุนิโย ปตฺตสนฺนิจยํ กริสฺสนฺตี”ติ ฯเปฯ.\csromanspacer{} สจฺจํ กิร ภิกฺขเว ฉพฺพคฺคิยา ภิกฺขุนิโย ปตฺตสนฺนิจยํ กโรนฺตีติ.\csromanspacer{} สจฺจํ ภควาติ.\csromanspacer{} วิครหิ พุทฺโธ ภควา ฯเปฯ.\csromanspacer{} กถํ หิ นาม ภิกฺขเว ฉพฺพคฺคิยา ภิกฺขุนิโย ปตฺตสนฺนิจยํ กริสฺสนฺติ.\csromanspacer{} เนตํ ภิกฺขเว อปฺปสนฺนานํ วา ปสาทาย ฯเปฯ.\csromanspacer{} เอวญฺจ ปน ภิกฺขเว ภิกฺขุนิโย อิมํ สิกฺขาปทํ อุทฺทิสนฺตุ\csromandash{}}

\proseitem{734}{\textbf{“ยา ปน ภิกฺขุนี ปตฺตสนฺนิจยํ กเรยฺย, นิสฺสคฺคิยํ ปาจิตฺติยนฺ”}ติ.}

\proseitem{735}{\textbf{ยา ปนา}ติ ยา ยาทิสา ฯเปฯ.\csromanspacer{} \textbf{ภิกฺขุนี}ติ ฯเปฯ อยํ อิมสฺมิํ อตฺเถ อธิปฺเปตา ภิกฺขุนีติ.}

\prose{\textbf{ปตฺโต} นาม เทฺว ปตฺตา อโยปตฺโต มตฺติกาปตฺโต.\csromanspacer{} ตโย ปตฺตสฺส วณฺณา อุกฺกฏฺโฐ ปตฺโต มชฺฌิโม ปตฺโต โอมโก ปตฺโต.\csromanspacer{} \textbf{อุกฺกฏฺโฐ} นาม \textbf{ปตฺโต} อฑฺฒาฬฺหโกทนํ คณฺหาติ จตุภาคํ ขาทนํ ตทุปิยํ พฺยญฺชนํ.\csromanspacer{} \textbf{มชฺฌิโม} นาม \textbf{ปตฺโต} นาฬิโกทนํ คณฺหาติ จตุภาคํ ขาทนํ ตทุปิยํ พฺยญฺชนํ.\csromanspacer{} \textbf{โอมโก} นาม ปตฺโต ปตฺโถทนํ คณฺหาติ จตุภาคํ ขาทนํ ตทุปิยํ พฺยญฺชนํ.\csromanspacer{} ตโต อุกฺกฏฺโฐ อ\textbf{ปตฺโต} โอมโก อปตฺโต.}

\prose{\textbf{สนฺนิจยํ กเรยฺยา}ติ อนธิฏฺฐิโต อวิกปฺปิโต.}

\prose{\csromanfolio{319}\textbf{นิสฺสคฺคิโย โหตี}ติ สห อรุณุคฺคมนา นิสฺสคฺคิโย โหติ.\csromanspacer{} นิสฺสชฺชิตพฺโพ สํฆสฺส วา คณสฺส วา เอกภิกฺขุนิยา วา.\csromanspacer{} เอวญฺจ ปน ภิกฺขเว นิสฺสชฺชิตพฺโพ, ตาย ภิกฺขุนิยา สํฆํ อุปสงฺกมิตฺวา เอกํสํ อุตฺตราสงฺคํ กริตฺวา วุฑฺฒานํ ภิกฺขุนีนํ ปาเท วนฺทิตฺวา อุกฺกุฏิกํ นิสีทิตฺวา อญฺชลิํ ปคฺคเหตฺวา เอวมสฺส วจนีโย “อยํ เม อยฺเย ปตฺโต รตฺตาติกฺกนฺโต นิสฺสคฺคิโย, อิมาหํ สํฆสฺส นิสฺสชฺชามี”ติ, นิสฺสชฺชิตฺวา อาปตฺติ เทเสตพฺพา, พฺยตฺตาย ภิกฺขุนิยา ปฏิพลาย อาปตฺติ ปฏิคฺคเหตพฺพา, นิสฺสฏฺฐปตฺโต ทาตพฺโพ\csromandash{}}

\prose{“สุณาตุ เม อยฺเย สํโฆ, อยํ ปตฺโต อิตฺถนฺนามาย ภิกฺขุนิยา นิสฺสคฺคิโย สํฆสฺส นิสฺสฏฺโฐ, ยทิ สํฆสฺส ปตฺตกลฺลํ, สํโฆ อิมํ ปตฺตํ อิตฺถนฺนามาย ภิกฺขุนิยา ทเทยฺยา”ติ.}

\prose{ตาย ภิกฺขุนิยา สมฺพหุลา ภิกฺขุนิโย อุปสงฺกมิตฺวา เอกํสํ อุตฺตราสงฺคํ กริตฺวา วุฑฺฒานํ ภิกฺขุนีนํ ปาเท วนฺทิตฺวา อุกฺกุฏิกํ นิสีทิตฺวา อญฺชลิํ ปคฺคเหตฺวา เอวมสฺสุ วจนียา “อยํ เม อยฺยาโย ปตฺโต รตฺตาติกฺกนฺโต นิสฺสคฺคิโย, อิมาหํ อยฺยานํ นิสฺสชฺชามี”ติ, นิสฺสชฺชิตฺวา อาปตฺติ เทเสตพฺพา, พฺยตฺตาย ภิกฺขุนิยา ปฏิพลาย อาปตฺติ ปฏิคฺคเหตพฺพา, นิสฺสฏฺฐปตฺโต ทาตพฺโพ\csromandash{}}

\prose{“สุณนฺตุ เม อยฺยาโย, อยํ ปตฺโต อิตฺถนฺนามาย ภิกฺขุนิยา นิสฺสคฺคิโย อยฺยานํ นิสฺสฏฺโฐ, ยทิ อยฺยานํ ปตฺตกลฺลํ, อยฺยาโย อิมํ ปตฺตํ อิตฺถนฺนามาย ภิกฺขุนิยา ทเทยฺยุนฺ”ติ.}

\prose{ตาย ภิกฺขุนิยา เอกํ ภิกฺขุนิํ อุปสงฺกมิตฺวา เอกํสํ อุตฺตราสงฺคํ กริตฺวา อุกฺกุฏิกํ นิสีทิตฺวา อญฺชลิํ ปคเหตฺวา เอวมสฺส วจนียา “อยํ เม อยฺเย ปตฺโต รตฺตาติกฺกนฺโต นิสฺสคฺคิโย, อิมาหํ อยฺยาย นิสฺสชฺชามี”ติ, นิสฺสชฺชิตฺวา อาปตฺติ เทเสตพฺพา, ตาย ภิกฺขุนิยา อาปตฺติ ปฏิคฺคเหตพฺพา, นิสฺสฏฺฐปตฺโต ทาตพฺโพ “อิมํ ปตฺตํ อยฺยาย ทมฺมี”ติ.}

\proseitem{736}{รตฺตาติกฺกนฺเต อติกฺกนฺตสญฺญา, นิสฺสคฺคิยํ ปาจิตฺติยํ.\csromanspacer{} รตฺตาติกฺกนฺเต เวมติกา, นิสฺสคฺคิยํ ปาจิตฺติยํ.\csromanspacer{} รตฺตาติกฺกนฺเต อนติกฺกนฺตสญฺญา, นิสฺสคฺคิยํ ปาจิตฺติยํ.\csromanspacer{} อนธิฏฺฐิเต อธิฏฺฐิตสญฺญา, นิสฺสคฺคิยํ ปาจิตฺติยํ.\csromanspacer{} อวิกปฺปิเต วิกปฺปิตสญฺญา, นิสฺสคฺคิยํ ปาจิตฺติยํ. \csromanfolio{320}อวิสฺสชฺชิเต วิสฺสชฺชิตสญฺญา, นิสฺสคฺคิยํ ปาจฺจิตฺติยํ.\csromanspacer{} อนฏฺเฐ นฏฺฐสญฺญา.\csromanspacer{} อวินฏฺเฐ วินฏฺฐสญฺญา.\csromanspacer{} อภินฺเน ภินฺนสญฺญา.\csromanspacer{} อวิลุตฺเต วิลุตฺตสญฺญา, นิสฺสคฺคิยํ ปาจิตฺติยํ.}

\prose{นิสฺสคฺคิยํ ปตฺตํ อนิสฺสชฺชิตฺวา ปริภุญฺชติ, อาปตฺติ ทุกฺกฏสฺส.\csromanspacer{} รตฺตานติกฺกนฺเต อติกฺกนฺตสญฺญา, อาปตฺติ ทุกฺกฏสฺส.\csromanspacer{} รตฺตานติกฺกนฺเต เวมติกา, อาปตฺติ ทุกฺกฏสฺส.\csromanspacer{} รตฺตานติกฺกนฺเต อนติกฺกนฺตสญฺญา, อนาปตฺติ.}

\proseitem{737}{อนาปตฺติ อนฺโต อรุเณ อธิฏฺเฐติ วิกปฺเปติ วิสฺสชฺเชติ นสฺสติ วินสฺสติ ภิชฺชติ อจฺฉินฺทิตฺวา คณฺหนฺติ วิสฺสาสํ คณฺหนฺติ, อุมฺมตฺติกาย อาทิกมฺมิกายาติ.}

\prosewithcloserruled{เตน โข ปน สมเยน ฉพฺพคฺคิยา ภิกฺขุนิโย นิสฺสฏฺฐปตฺตํ น เทนฺติ.\csromanspacer{} ภควโต เอตมตฺถํ อาโรเจสุํ.\csromanspacer{} น ภิกฺขเว นิสฺสฏฺฐปตฺโต น ทาตพฺโพ, ยา น ทเทยฺย, อาปตฺติ ทุกฺกฏสฺสาติ.}{ปฐมสิกฺขาปทํ นิฏฺฐิตํ.}
\csromanmatikaanchor{mtk.137}
\csromantocmarkref{subsection}{2. ทุติยสิกฺขาปท}{mtk.137}
\bookmark[dest={mtk.137},level=2]{2. ทุติยสิกฺขาปท}
\csromanheader{2}{1. \textbf{ปตฺตวคฺค 2. ทุติยสิกฺขาปท}}
\proseitem{738}{เตน สมเยน พุทฺโธ ภควา สาวตฺถิยํ วิหรติ เชตวเน อนาถปิณฺฑิกสฺส อาราเม.\csromanspacer{} เตน โข ปน สมเยน สมฺพหุลา ภิกฺขุนิโย คามกาวาเส วสฺสํวุฏฺฐา\footnote{วสฺสํวุตฺถา (สี, สฺยา)} สาวตฺถิํ อคมํสุ วตฺตสมฺปนฺนา อิริยาปถสมฺปนฺนา ทุจฺโจฬา ลูขจีวรา.\csromanspacer{} อุปาสกา ตา ภิกฺขุนิโย ปสฺสิตฺวา “อิมา ภิกฺขุนิโย วตฺตสมฺปนฺนา อิริยาปถสมฺปนฺนา ทุจฺโจฬา ลูขจีวรา, อิมา ภิกฺขุนิโย อจฺฉินฺนา ภวิสฺสนฺตี”ติ ภิกฺขุนิสํฆสฺส อกาลจีวรํ อทํสุ.\csromanspacer{} ถุลฺลนนฺทา ภิกฺขุนี “อมฺหากํ กถินํ อตฺถตํ กาลจีวรนฺ”ติ อธิฏฺฐหิตฺวา ภาชาเปสิ.\csromanspacer{} อุปาสกา ตา ภิกฺขุนิโย ปสฺสิตฺวา เอตทโวจุํ “อปยฺยาหิ จีวรํ ลทฺธนฺ”ติ.\csromanspacer{} น มยํ อาวุโส จีวรํ ลภาม, อยฺยา ถุลฺลนนฺทา “อมฺหากํ กถินํ อตฺถตํ กาลจีวรนฺ”ติ อธิฏฺฐหิตฺวา ภาชาเปสีติ.\csromanspacer{} อุปาสกา อุชฺฌายนฺติ ขิยฺยนฺติ วิปาเจนฺติ “กถํ หิ นาม อยฺยา ถุลฺลนนฺทา อกาลจีวรํ ‘กาลจีวรนฺ’ติ อธิฏฺฐหิตฺวา ภาชาเปสฺสตี”ติ. \csromanfolio{321}อสฺโสสุํ โข ภิกฺขุนิโย เตสํ อุปาสกานํ อุชฺฌายนฺตานํ ขิยฺยนฺตานํ วิปาเจนฺตานํ.\csromanspacer{} ยา ตา ภิกฺขุนิโย อปฺปิจฺฉา ฯเปฯ ตา อุชฺฌายนฺติ ขิยฺยนฺติ วิปาเจนฺติ “กถํ หิ นาม อยฺยา ถุลฺลนนฺทา อกาลจีวรํ ‘กาลจีวรนฺ’ติ อธิฏฺฐหิตฺวา ภาชาเปสฺสตี”ติ.\csromanspacer{} อถ โข ตา ภิกฺขุนิโย ภิกฺขูนํ เอตมตฺถํ อาโรเจสุํ.\csromanspacer{} ภิกฺขู ภควโต เอตมตฺถํ อาโรเจสุํ ฯเปฯ.\csromanspacer{} สจฺจํ กิร ภิกฺขเว ถุลฺลนนฺทา ภิกฺขุนี อกาลจีวรํ “กาลจีวรนฺ”ติ อธิฏฺฐหิตฺวา ภาชาเปตีติ\footnote{ภาชาเปสีติ (ก)}.\csromanspacer{} สจฺจํ ภควาติ.\csromanspacer{} วิครหิ “พุทฺโธ ภควา ฯเปฯ.\csromanspacer{} กถํ หิ นาม ภิกฺขเว ถุลฺลนนฺทา ภิกฺขุนี อกาลจีวรํ “กาลจีวรนฺ”ติ อธิฏฺฐหิตฺวา ภาชาเปสฺสติ.\csromanspacer{} เนตํ ภิกฺขเว อปฺปสนฺนานํ วา ปสาทาย ฯเปฯ.\csromanspacer{} เอวญฺจ ปน ภิกฺขเว ภิกฺขุนิโย อิมํ สิกฺขาปทํ อุทฺทิสนฺตุ\csromandash{}}

\proseitem{739}{\textbf{“ยา ปน ภิกฺขุนี อกาลจีวรํ ‘กาลจีวรนฺ’ติ อธิฏฺฐหิตฺวา ภาชาเปยฺย, นิสฺสคฺคิยํ ปาจิตฺติยนฺ}”ติ.}

\proseitem{740}{\textbf{ยา ปนา}ติ ยา ยาทิสา ฯเปฯ.\csromanspacer{} \textbf{ภิกฺขุนี}ติ ฯเปฯ อยํ อิมสฺมิํ อตฺเถ อธิปฺเปตา ภิกฺขุนีติ.}

\prose{\textbf{อกาลจีวรํ} นาม อนตฺถเต กถิเน เอกาทสมาเส อุปฺปนฺนํ, อตฺถเต กถิเน สตฺตมาเส อุปฺปนฺนํ, กาเลปิ อาทิสฺส ทินฺนํ, เอตํ อกาลจีวรํ นาม.}

\prose{อกาลจีวรํ “กาลจีวรนฺ”ติ อธิฏฺฐหิตฺวา ภาชาเปติ, ปโยเค ทุกฺกฏํ.\csromanspacer{} ปฏิลาเภน นิสฺสคฺคิยํ โหติ.\csromanspacer{} นิสฺสชฺชิตพฺพํ สํฆสฺส วา คณสฺส วา เอกภิกฺขุนิยา วา.\csromanspacer{} เอวญฺจ ปน ภิกฺขเว นิสฺสชฺชิตพฺพํ ฯเปฯ “อิทํ เม อยฺเย อกาลจีวรํ ‘กาลจีวรนฺ’ติ อธิฏฺฐหิตฺวา ภาชาปิตํ นิสฺสคฺคิยํ, อิมาหํ สํฆสฺส นิสฺสชฺชามี”ติ ฯเปฯ ทเทยฺยาติ ฯเปฯ ทเทยฺยุนฺติ ฯเปฯ อยฺยาย ทมฺมีติ.}

\proseitem{741}{อกาลจีวเร อกาลจีวรสญฺญา “กาลจีวรนฺ”ติ อธิฏฺฐหิตฺวา ภาชาเปติ, นิสฺสคฺคิยํ ปาจิตฺติยํ.\csromanspacer{} อกาลจีวเร เวมติกา “กาลจีวรนฺ”ติ อธิฏฺฐหิตฺวา ภาชาเปติ, อาปตฺติ ทุกฺกฏสฺส.\csromanspacer{} อกาลจีวเร กาลจีวรสญฺญา “กาลจีวรนฺ”ติ อธิฏฺฐหิตฺวา \csromanfolio{322}\textbf{ภาชาเปติ, อนาปตฺติ.}\csromanspacer{}\textbf{ กาลจีวเร อกาลจีวรสญฺญา, อาปตฺติ ทุกฺกฏสฺส.}\csromanspacer{} \textbf{กาลจีวเร เวมติกา, อาปตฺติ ทุกฺกฏสฺส.}\csromanspacer{}\textbf{ กาลจีวเร กาลจีวรสญฺญา,} \textbf{อนาปตฺติ.}}

\proseitemwithcloserruled{742}{อนาปตฺติ อกาลจีวรํ กาลจีวรสญฺญา ภาชาเปติ, กาลจีวรํ กาลจีวรสญฺญา ภาชาเปติ, อุมฺมตฺติกาย อาทิกมฺมิกายาติ.}{ทุติยสิกฺขาปทํ นิฏฺฐิตํ.}
\csromanmatikaanchor{mtk.138}
\csromantocmarkref{subsection}{3. ตติยสิกฺขาปท}{mtk.138}
\bookmark[dest={mtk.138},level=2]{3. ตติยสิกฺขาปท}
\csromanheader{2}{1. \textbf{ปตฺตวคฺค 3. ตติยสิกฺขาปท}}
\proseitem{743}{เตน สมเยน พุทฺโธ ภควา สาวตฺถิยํ วิหรติ เชตวเน อนาถปิณฺฑิกสฺส อาราเม.\csromanspacer{} เตน โข ปน สมเยน ถุลฺลนนฺทา ภิกฺขุนี อญฺญตราย ภิกฺขุนิยา สทฺธิํ จีวรํ ปริวตฺเตตฺวา ปริภุญฺชิ.\csromanspacer{} อถ โข สา ภิกฺขุนี ตํ จีวรํ สงฺฆริตฺวา\footnote{สํหริตฺวา (ก)} นิกฺขิปิ.\csromanspacer{} ถุลฺลนนฺทา ภิกฺขุนี ตํ ภิกฺขุนิํ เอตทโวจ “ยํ เต อยฺเย มยา สทฺธิํ จีวรํ ปริวตฺติตํ, กหํ ตํ จีวรนฺ”ติ.\csromanspacer{} อถ โข สา ภิกฺขุนี ตํ จีวรํ นีหริตฺวา ถุลฺลนนฺทาย ภิกฺขุนิยา ทสฺเสสิ.\csromanspacer{} ถุลฺลนนฺทา ภิกฺขุนี ตํ ภิกฺขุนิํ เอตทโวจ “หนฺทายฺเย ตุยฺหํ จีวรํ, อาหร เม’ตํ จีวรํ, ยํ ตุยฺหํ ตุยฺหเมเวตํ, ยํ มยฺหํ มยฺหเมเวตํ, อาหร เม’ตํ, สกํ ปจฺจาหรา”ติ อจฺฉินฺทิ.\csromanspacer{} อถ โข สา ภิกฺขุนี ภิกฺขุนีนํ เอตมตฺถํ อาโรเจสิ.\csromanspacer{} ยา ตา ภิกฺขุนิโย อปฺปิจฺฉา ฯเปฯ ตา อุชฺฌายนฺติ ขิยฺยนฺติ วิปาเจนฺติ “กถํ หิ นาม อยฺยา ถุลฺลนนฺทา ภิกฺขุนิยา สทฺธิํ จีวรํ ปริวตฺเตตฺวา อจฺฉินฺทิสฺสตี”ติ.\csromanspacer{} อถ โข ตา ภิกฺขุนิโย ภิกฺขูนํ เอตมตฺถํ อาโรเจสุํ.\csromanspacer{} ภิกฺขู ภควโต เอตมตฺถํ อาโรเจสุํ ฯเปฯ.\csromanspacer{} สจฺจํ กิร ภิกฺขเว ถุลฺลนนฺทา ภิกฺขุนี ภิกฺขุนิยา สทฺธิํ จีวรํ ปริวตฺเตตฺวา อจฺฉินฺทตีติ\footnote{อจฺฉินฺทีติ (ก)}.\csromanspacer{} สจฺจํ ภควาติ.\csromanspacer{} วิครหิ พุทฺโธ ภควา ฯเปฯ.\csromanspacer{} กถํ หิ นาม ภิกฺขเว ถุลฺลนนฺทา ภิกฺขุนี ภิกฺขุนิยา สทฺธิํ จีวรํ ปริวตฺเตตฺวา อจฺฉินฺทิสฺสติ.\csromanspacer{} เนตํ ภิกฺขเว อปฺปสนฺนานํ วา ปสาทาย ฯเปฯ.\csromanspacer{} เอวญฺจ ปน ภิกฺขเว ภิกฺขุนิโย อิมํ สิกฺขาปทํ อุทฺทิสนฺตุ\csromandash{}}

\proseitem{744}{\csromanfolio{323}\textbf{“ยา ปน ภิกฺขุนี ภิกฺขุนิยา สทฺธิํ จีวรํ ปริวตฺเตตฺวา สา ปจฺฉา เอวํ วเทยฺย ‘หนฺทายฺเย ตุยฺหํ จีวรํ, อาหร เม’ตํ จีวรํ, ยํ ตุยฺหํ ตุยฺหเมเวตํ, ยํ มยหํ มยฺหเมเวตํ, อาหร เม’ตํ, สกํ ปจฺจาหรา”ติ อจฺฉินฺเทยฺย วา อจฺฉินฺทาเปยฺย วา, นิสฺสคฺคิยํ ปาจิตฺติยนฺ”}ติ.}

\proseitem{745}{\textbf{ยา ปนา}ติ ยา ยาทิสา ฯเปฯ.\csromanspacer{} \textbf{ภิกฺขุนี}ติ ฯเปฯ อยํ อิมสฺมิํ อตฺเถ อธิปฺเปตา ภิกฺขุนีติ.}

\prose{\textbf{ภิกฺขุนิยา สทฺธินฺ}ติ อญฺญาย ภิกฺขุนิยา สทฺธิํ.}

\prose{\textbf{จีวรํ} นาม ฉนฺนํ จีวรานํ อญฺญตรํ จีวรํ วิกปฺปนุปคํ ปจฺฉิมํ.}

\prose{\textbf{ปริวตฺเตตฺวา}ติ ปริตฺเตน วา วิปุลํ, วิปุเลน วา ปริตฺตํ.}

\prose{\textbf{อจฺฉินฺเทยฺยา}ติ สยํ อจฺฉินฺทติ, นิสฺสคฺคิยํ ปาจิตฺติยํ.}

\prose{\textbf{อจฺฉินฺทาเปยฺยา}ติ อญฺญํ อาณาเปติ, อาปตฺติ ทุกฺกฏสฺส.\csromanspacer{} สกิํ อาณตฺตา พหุกมฺปิ อจฺฉินฺทติ, นิสฺสคฺคิยํ โหติ.\csromanspacer{} นิสฺสชฺชิตพฺพํ สํฆสฺส วา คณสฺส วา เอกภิกฺขุนิยา วา.\csromanspacer{} เอวญฺจา ปน ภิกฺขเว นิสฺสชฺชิตพฺพํ ฯเปฯ “อิทํ เม อยฺเย จีวรํ ภิกฺขุนิยา สทฺธิํ ปริวตฺเตตฺวา อจฺฉินฺนํ นิสฺสคฺคิยํ, อิมาหํ สํฆสฺส นิสฺสชฺชามี”ติ ฯเปฯ ทเทยฺยาติ ฯเปฯ ทเทยฺยุนฺติ ฯเปฯ อยฺยาย ทมฺมีติ.}

\proseitem{746}{อุปสมฺปนฺนาย อุปสมฺปนฺนสญฺญา จีวรํ ปริวตฺเตตฺวา อจฺฉินฺทติ วา อจฺฉินฺทาเปติ วา, นิสฺสคฺคิยํ ปาจิตฺติยํ.\csromanspacer{} อุปสมฺปนฺนาย เวมติกา จีวรํ ปริวตฺเตตฺวา อจฺฉินฺทติ วา อจฺฉินฺทาเปติ วา, นิสฺสคฺคิยํ ปาจิตฺติยํ.\csromanspacer{} อุปสมฺปนฺนาย อนุปสมฺปนฺนสญฺญา จีวรํ ปริวตฺเตตฺวา อจฺฉินฺทติ วา อจฺฉินฺทาเปติ วา, นิสฺสคฺคิยํ ปาจิตฺติยํ.}

\prose{อญฺญํ ปริกฺขารํ ปริวตฺเตตฺวา อจฺฉินฺทติ วา อจฺฉินฺทาเปติ วา, อาปตฺติ ทุกฺกฏสฺส.\csromanspacer{} อนุปสมฺปนฺนาย สทฺธิํ จีวรํ วา อญฺญํ วา ปริกฺขารํ ปริวตฺเตตฺวา อจฺฉินฺทติ วา อจฺฉินฺทาเปติ วา, อาปตฺติ ทุกฺกฏสฺส.\csromanspacer{} อนุปสมฺปนฺนาย อุปสมฺปนฺนสญฺญา, อาปตฺติ ทุกฺกฏสฺส.\csromanspacer{} อนุปสมฺปนฺนาย เวมติกา, อาปตฺติ ทุกฺกฏสฺส.\csromanspacer{} อนุปสมฺปนฺนาย อนุปสมฺปนฺนสญฺญา, อาปตฺติ ทุกฺกฏสฺส.}

\proseitemwithcloser{747}{อนาปตฺติ สา วา เทติ, ตสฺสา วา วิสฺสสนฺตี คณฺหาติ, อุมฺมตฺติกาย อาทิกมฺมิกายาติ.}{ตติยสิกฺขาปทํ นิฏฺฐิตํ.}
\csromanmatikaanchor{mtk.139}
\csromantocmarkref{subsection}{4. จตุตฺถสิกฺขาปท}{mtk.139}
\bookmark[dest={mtk.139},level=2]{4. จตุตฺถสิกฺขาปท}
\csromanheader{2}{1. \textbf{ปตฺตวคฺค 4. จตุตฺถสิกฺขาปท}}
\proseitem{748}{\csromanfolio{324}เตน สมเยน พุทฺโธ ภควา สาวตฺถิยํ วิหรติ เชตวเน อนาถปิณฺฑิกสฺส อาราเม.\csromanspacer{} เตน โข ปน สมเยน ถุลฺลนนฺทา ภิกฺขุนี คิลานา โหติ.\csromanspacer{} อถ โข อญฺญตโร อุปาสโก เยน ถุลฺลนนฺทา ภิกฺขุนี เตนุปสงฺกมิ, อุปสงฺกมิตฺวา ถุลฺลนนฺทํ ภิกฺขุนิํ เอตทโวจ “กิํ เต อยฺเย อผาสุ, กิํ อาหรียตู”ติ.\csromanspacer{} สปฺปินา เม อาวุโส อตฺโถติ.\csromanspacer{} อถ โข โส อุปาสโก อญฺญตรสฺส อาปณิกสฺส ฆรา กหาปณสฺส สปฺปิํ อาหริตฺวา ถุลฺลนนฺทาย ภิกฺขุนิยา อทาสิ.\csromanspacer{} ถุลฺลนนฺทา ภิกฺขุนี เอวมาห “น เม อาวุโส สปฺปินา อตฺโถ, เตเลน เม อตฺโถ”ติ.\csromanspacer{} อถ โข โส อุปาสโก เยน โส อาปณิโก เตนุปสงฺกมิ, อุปสงฺกมิตฺวา ตํ อาปณิกํ เอตทโวจ “น กิรายฺโย อยฺยาย สปฺปินา อตฺโถ, เตเลน อตฺโถ, หนฺท เต สปฺปิํ, เตลํ เม เทหี”ติ.\csromanspacer{} สเจ มยํ อยฺโย วิกฺกีตํ ภณฺฑํ ปุน อาทิยิสฺสาม\footnote{อาหริสฺสาม (ก)}, กทา อมฺหากํ ภณฺฑํ วิกฺกายิสฺสติ\footnote{วิกฺกียิสฺสติ (?)}, สปฺปิสฺส กเยน สปฺปิ หฏํ, เตลสฺส กยํ อาหร, เตลํ หริสฺสสีติ.\csromanspacer{} อถ โข โส อุปาสโก อุชฺฌายติ ขิยฺยติ วิปาเจติ “กถํ หิ นาม อยฺยา ถุลฺลนนฺทา อญฺญํ วิญฺญาเปตฺวา อญฺญํ วิญฺญาเปสฺสตี”ติ.\csromanspacer{} อสฺโสสุํ โข ภิกฺขุนิโย ตสฺส อุปาสกสฺส อุชฺฌายนฺตสฺส ขิยฺยนฺตสฺส วิปาเจนฺตสฺส.\csromanspacer{} ยา ตา ภิกฺขุนิโย อปฺปิจฺฉา ฯเปฯ ตา อุชฺฌายนฺติ ขิยฺยนฺติ วิปาเจนฺติ ฯเปฯ.\csromanspacer{} อถ โข ตา ภิกฺขุนิโย ภิกฺขูนํ เอตมตฺถํ อาโรเจสุํ.\csromanspacer{} ภิกฺขู ภควโต เอตมตฺถํ อาโรเจสุํ ฯเปฯ.\csromanspacer{} สจฺจํ กิร ภิกฺขเว ถุลฺลนนฺทา ภิกฺขุนี อญฺญํ วิญฺญาเปตฺวา อญฺญํ วิญฺญาเปตีติ\footnote{วิญฺญาเปสีติ (ก)}.\csromanspacer{} สจฺจํ ภควาติ.\csromanspacer{} วิครหิ พุทฺโธ ภควา ฯเปฯ.\csromanspacer{} กถํ หิ นาม ภิกฺขเว ถุลฺลนนฺทา ภิกฺขุนี อญฺญํ วิญฺญาเปตฺวา อญฺญํ วิญฺญาเปสฺสติ.\csromanspacer{} เนตํ ภิกฺขเว อปฺปสนฺนานํ วา ปสาทาย ฯเปฯ.\csromanspacer{} เอวญฺจ ปน ภิกฺขเว ภิกฺขุนิโย อิมํ สิกฺขาปทํ อุทฺทิสนฺตุ\csromandash{}}

\proseitem{749}{\textbf{“ยา ปน ภิกฺขุนี อญฺญํ วิญฺญาเปตฺวา อญฺญํ วิญฺญาเปยฺย, นิสฺสคฺคิยํ ปาจิตฺติยนฺ}”ติ.}

\proseitem{750}{\textbf{ยา ปนา}ติ ยา ยาทิสา ฯเปฯ.\csromanspacer{} \textbf{ภิกฺขุนี}ติ ฯเปฯ อยํ อิมสฺมิํ อตฺเถ อธิปฺเปตา ภิกฺขุนีติ.}

\prose{\csromanfolio{325}\textbf{อญฺญํ วิญฺญาเปตฺวา}ติ ยํ กิญฺจิ วิญฺญาเปตฺวา.}

\prose{\textbf{อญฺญํ วิญฺญาเปยฺยา}ติ ตํ ฐเปตฺวา อญฺญํ วิญฺญาเปติ, ปโยเค ทุกฺกฏํ.\csromanspacer{} ปฏิลาเภน นิสฺสคฺคิยํ โหติ.\csromanspacer{} นิสฺสชฺชิตพฺพํ สํฆสฺส วา คณสฺส วา เอกภิกฺขุนิยา วา.\csromanspacer{} เอวญฺจ ปน ภิกฺขเว นิสฺสชฺชิตพฺพํ ฯเปฯ “อิทํ เม อยฺเย อญฺญํ วิญฺญาเปตฺวา อญฺญํ วิญฺญาปิตํ นิสฺสคฺคิยํ, อิมาหํ สํฆสฺส นิสฺสชฺชามีติ ฯเปฯ ทเทยฺยาติ ฯเปฯ ทเทยฺยุนฺติ ฯเปฯ อยฺยาย ทมฺมี”ติ.}

\proseitem{751}{อญฺเญ อญฺญสญฺญา อญฺญํ วิญฺญาเปติ นิสฺสคฺคิยํ ปาจิตฺติยํ.\csromanspacer{} อญฺเญ เวมติกา อญฺญํ วิญฺญาเปติ, นิสฺสคฺคิยํ ปาจิตฺติยํ.\csromanspacer{} อญฺเญ อนญฺญสญฺญา อญฺญํ วิญฺญาเปติ, นิสฺสคฺคิยํ ปาจิตฺติยํ.}

\prose{อนญฺเญ อญฺญสญฺญา อนญฺญํ วิญฺญาเปติ, อาปตฺติ ทุกฺกฏสฺส.\csromanspacer{} อนญฺเญ เวมติกา อนญฺญํ วิญฺญาเปติ, อาปตฺติ ทุกฺกฏสฺส.\csromanspacer{} อนญฺเญ อนญฺญสญฺญา, อนาปตฺติ.}

\proseitem{752}{อนาปตฺติ ตญฺเญว\footnote{ตญฺเจว (สฺยา)} วิญฺญาเปติ, อญฺญญฺจ วิญฺญาเปติ, อานิสํสํ ทสฺเสตฺวา วิญฺญาเปติ, อุมฺมตฺติกาย อาทิกมฺมิกายาติ.\csromanspacer{} จตุตฺถสิกฺขาปทํ นิฏฺฐิตํ.}
\csromansectionrule

\csromanmatikaanchor{mtk.140}
\csromantocmarkref{subsection}{5. ปญฺจมสิกฺขาปท}{mtk.140}
\bookmark[dest={mtk.140},level=2]{5. ปญฺจมสิกฺขาปท}
\csromanheader{2}{1. \textbf{ปตฺตวคฺค 5. ปญฺจมสิกฺขาปท}}
\proseitem{753}{เตน สมเยน พุทฺโธ ภควา สาวตฺถิยํ วิหรติ เชตวเน อนาถปิณฺฑิกสฺส อาราเม.\csromanspacer{} เตน โข ปน สมเยน ถุลฺลนนฺทา ภิกฺขุนี คิลานา โหติ.\csromanspacer{} อถ โข อญฺญตโร อุปาสโก เยน ถุลฺลนนฺทา ภิกฺขุนี เตนุปสงฺกมิ, อุปสงฺกมิตฺวา ถุลฺลนนฺทํ ภิกฺขุนิํ เอตทโวจ “กจฺจิ อยฺเย ขมนียํ, กจฺจิ ยาปนียนฺ”ติ.\csromanspacer{} น เม อาวุโส ขมนียํ, น ยาปนียนฺติ.\csromanspacer{} อมุกสฺส อยฺเย อาปณิกสฺส ฆเร กหาปณํ นิกฺขิปิสฺสามิ, ตโต ยํ อิจฺเฉยฺยาสิ, ตํ อาหราเปยฺยาสีติ.\csromanspacer{} ถุลฺลนนฺทา ภิกฺขุนี อญฺญตรํ สิกฺขมานํ อาณาเปสิ “คจฺฉ สิกฺขมาเน อมุกสฺส อาปณิกสฺส ฆรา กหาปณสฺส เตลํ อาหรา”ติ.\csromanspacer{} อถ โข สา สิกฺขมานา ตสฺส อาปณิกสฺส ฆรา กหาปณสฺส \csromanfolio{326}เตลํ อาหริตฺวา ถุลฺลนนฺทาย ภิกฺขุนิยา อทาสิ.\csromanspacer{} ถุลฺลนนฺทา ภิกฺขุนี เอวมาห “น เม สิกฺขมาเน เตเลน อตฺโถ, สปฺปินา เม อตฺโถ”ติ.\csromanspacer{} อถ โข สา สิกฺขมานา เยน โส อาปณิโก เตนุปสงฺกมิ, อุปสงฺกมิตฺวา ตํ อาปณิกํ เอตทโวจ “น กิร อาวุโส อยฺยาย เตเลน อตฺโถ, สปฺปินา อตฺโถ, หนฺท เต เตลํ, สปฺปิํ เม เทหี”ติ.\csromanspacer{} สเจ มยํ อยฺเย วิกฺกีตํ ภณฺฑํ ปุน อาทิยิสฺสาม, กทา อมฺหากํ ภณฺฑํ วิกฺกายิสฺสติ, เตลสฺส กเยน เตลํ หฏํ, สปฺปิสฺส กยํ อาหร, สปฺปิํ หริสฺสสีติ.\csromanspacer{} อถ โข สา สิกฺขมานา โรทนฺตี อฏฺฐาสิ.\csromanspacer{} ภิกฺขุนิโย ตํ สิกฺขมานํ เอตทโวจุํ “กิสฺส ตฺวํ สิกฺขมาเน โรทสี”ติ.\csromanspacer{} อถ โข สา สิกฺขมานา ภิกฺขุนีนํ เอตมตฺถํ อาโรเจสิ.\csromanspacer{} ยา ตา ภิกฺขุนิโย อปฺปิจฺฉา ฯเปฯ ตา อุชฺฌายนฺติ ขิยฺยนฺติ วิปาเจนฺติ “กถํ หิ นาม อยฺยา ถุลฺลนนฺทา อญฺญํ เจตาเปตฺวา อญฺญํ เจตาเปสฺสตี”ติ ฯเปฯ.\csromanspacer{} สจฺจํ กิร ภิกฺขเว ถุลฺลนนฺทา ภิกฺขุนี อญฺญํ เจตาเปตฺวา อญฺญํ เจตาเปตี’ติ\footnote{เจตาเปสีติ (ก)}.\csromanspacer{} \textbf{สจฺจํ ภควาติ.}\csromanspacer{}\textbf{ วิครหิ พุทฺโธ ภควา ฯเปฯ.}\csromanspacer{}\textbf{ กถํ หิ นาม} \textbf{ภิกฺขเว ถุลฺลนนฺทา ภิกฺขุนี อญฺญํ เจตาเปตฺวา อญฺญํ เจตาเปสฺสติ.}\csromanspacer{} เนตํ ภิกฺขเว อปฺปสนฺนานํ วา ปสาทาย ฯเปฯ.\csromanspacer{} เอวญฺจ ปน ภิกฺขเว ภิกฺขุนิโย อิมํ สิกฺขาปทํ อุทฺทิสนฺตุ\csromandash{}}

\proseitem{754}{\textbf{“ยา ปน ภิกฺขุนี อญฺญํ เจตาเปตฺวา อญฺญํ เจตาเปยฺย, นิสฺสคฺคิยํ ปาจิตฺติยนฺ}”ติ.}

\proseitem{755}{\textbf{ยา ปนา}ติ ยา ยาทิสา ฯเปฯ.\csromanspacer{} \textbf{ภิกฺขุนี}ติ ฯเปฯ อยํ อิมสฺมิํ อตฺเถ อธิปฺเปตา ภิกฺขุนีติ.}

\prose{\textbf{อญฺญํ เจตาเปตฺวา}ติ ยํ กิญฺจิ เจตาเปตฺวา.}

\prose{\textbf{อญฺญํ เจตาเปยฺยา}ติ ตํ ฐเปตฺวา อญฺญํ เจตาเปติ, ปโยเค ทุกฺกฏํ.\csromanspacer{} ปฏิลาเภน นิสฺสคฺคิยํ โหติ.\csromanspacer{} นิสฺสชฺชิตพฺพํ สํฆสฺส วา คณสฺส วา เอกภิกฺขุนิยา วา.\csromanspacer{} เอวญฺจ ปน ภิกฺขเว นิสฺสชฺชิตพฺพํ. “อิทํ เม อยฺเย อญฺญํ เจตาเปตฺวา อญฺญํ เจตาปิตํ นิสฺสคฺคิยํ, อิมาหํ สํฆสฺส นิสฺสชฺชามี”ติ ฯเปฯ ทเทยฺยาติ ฯเปฯ ทเทยฺยุนฺติ ฯเปฯ อยฺยาย ทมฺมีติ.}

\proseitem{756}{\csromanfolio{327}อญฺเญ อญฺญสญฺญา อญฺญํ เจตาเปติ, นิสฺสคฺคิยํ ปาจิตฺติยํ.\csromanspacer{} อญฺเญ เวมติกา อญฺญํ เจตาเปติ, นิสฺสคฺคิยํ ปาจิตฺติยํ.\csromanspacer{} อญฺเญ อนญฺญสญฺญา อญฺญํ เจตาเปติ, นิสฺสคฺคิยํ ปาจิตฺติยํ.}

\prose{อนญฺเญ อญฺญสญฺญา อนญฺญํ เจตาเปติ, อาปตฺติ ทุกฺกฏสฺส.\csromanspacer{} อนญฺเญ เวมติกา อนญฺญํ เจตาเปติ, อาปตฺติ ทุกฺกฏสฺส.\csromanspacer{} อนญฺเญ อนญฺญสญฺญา อนาปตฺติ.}

\proseitem{757}{อนาปตฺติ ตญฺเญว\footnote{ตญฺเจว (สฺยา)} เจตาเปติ, อญฺญญฺจ เจตาเปติ, อานิสํสํ ทสฺเสตฺวา เจตาเปติ, อุมฺมตฺติกาย อาทิกมฺมิกายาติ.\csromanspacer{} ปญฺจมสิกฺขาปทํ นิฏฺฐิตํ.}
\csromansectionrule

\csromanmatikaanchor{mtk.141}
\csromantocmarkref{subsection}{6. ฉฏฺฐสิกฺขาปท}{mtk.141}
\bookmark[dest={mtk.141},level=2]{6. ฉฏฺฐสิกฺขาปท}
\csromanheader{2}{1. \textbf{ปตฺตวคฺค 6. ฉฏฺฐสิกฺขาปท}}
\proseitem{758}{เตน สมเยน พุทฺโธ ภควา สาวตฺถิยํ วิหรติ เชตวเน อนาถปิณฺฑิกสฺส อาราเม.\csromanspacer{} เตน โข ปน สมเยน อุปาสกา ภิกฺขุนิสํฆสฺส จีวรตฺถาย ฉนฺทกํ สงฺฆริตฺวา อญฺญตรสฺส ปาวาริกสฺส ฆเร ปริกฺขารํ นิกฺขิปิตฺวา ภิกฺขุนิโย อุปสงฺกมิตฺวา เอตทโวจุํ “อมุกสฺส อยฺเย ปาวาริกสฺส ฆเร จีวรตฺถาย ปริกฺขาโร นิกฺขิตฺโต, ตโต จีวรํ อาหราเปตฺวา ภาเชถา”ติ.\csromanspacer{} ภิกฺขุนิโย เตน ปริกฺขาเรน เภสชฺชํ เจตาเปตฺวา ปริภุญฺชิํสุ.\csromanspacer{} อุปาสกา ชานิตฺวา อุชฺฌายนฺติ ขิยฺยนฺติ วิปาเจนฺติ “กถํ หิ นาม ภิกฺขุนิโย อญฺญทตฺถิเกน ปริกฺขาเรน อญฺญุทฺทิสิเกน สํฆิเกน อญฺญํ เจตาเปสฺสนฺตี”ติ.\csromanspacer{} อสฺโสสุํ โข ภิกฺขุนิโย เตสํ อุปาสกานํ อุชฺฌายนฺตานํ ขิยฺยนฺตานํ วิปาเจนฺตานํ.\csromanspacer{} ยา ตา ภิกฺขุนิโย อปฺปิจฺฉา ฯเปฯ ตา อุชฺฌายนฺติ ขิยฺยนฺติ วิปาเจนฺติ “กถํ หิ นาม ภิกฺขุนิโย อญฺญทตฺถิเกน ปริกฺขาเรน อญฺญุทฺทิสิเกน สํฆิเกน อญฺญํ เจตาเปสฺสนฺตี”ติ ฯเปฯ สจฺจํ กิร ภิกฺขเว ภิกฺขุนิโย อญฺญทตฺถิเกน ปริกฺขาเรน อญฺญุทฺทิสิเกน สํฆิเกน อญฺญํ เจตาเปนฺตีติ.\csromanspacer{} สจฺจํ ภควาติ.\csromanspacer{} วิครหิ พุทฺโธ ภควา ฯเปฯ.\csromanspacer{} กถํ หิ นาม ภิกฺขเว ภิกฺขุนิโย \csromanfolio{328}อญฺญทตฺถิเกน ปริกฺขาเรน อญฺญุทฺทิสิเกน สํฆิเกน อญฺญํ เจตาเปสฺสนฺติ.\csromanspacer{} เนตํ ภิกฺขเว อปฺปสนฺนานํ วา ปสาทาย ฯเปฯ.\csromanspacer{} เอวญฺจ ปน ภิกฺขเว ภิกฺขุนิโย อิมํ สิกฺขาปทํ อุทฺทิสนฺตุ\csromandash{}}

\proseitem{759}{\textbf{“ยา ปน ภิกฺขุนี อญฺญทตฺถิเกน ปริกฺขาเรน อญฺญุทฺทิสิเกน สํฆิเกน อญฺญํ เจตาเปยฺย, นิสฺสคฺคิยํ ปาจิตฺติยนฺ}”ติ.}

\proseitem{760}{\textbf{ยา ปนา}ติ ยา ยาทิสา ฯเปฯ.\csromanspacer{} \textbf{ภิกฺขุนี}ติ ฯเปฯ อยํ อิมสฺมิํ อตฺเถ อธิปฺเปตา ภิกฺขุนีติ.}

\prose{\textbf{อญฺญทตฺถิเกน ปริกฺขาเรน อญฺญุทฺทิสิเกนา}ติ อญฺญสฺสตฺถาย ทินฺเนน.}

\prose{\textbf{สํฆิเกนา}ติ สํฆสฺส, น คณสฺส, น เอกภิกฺขุนิยา.}

\prose{\textbf{อญฺญํ เจตาเปยฺยา}ติ ยํอตฺถาย ทินฺนํ, ตํ ฐเปตฺวา อญฺญํ เจตาเปติ, ปโยเค ทุกฺกฏํ.\csromanspacer{} ปฏิลาเภน นิสฺสคฺคิยํ โหติ.\csromanspacer{} นิสฺสชฺชิตพฺพํ สํฆสฺส วา คณสฺส วา เอกภิกฺขุนิยา วา.\csromanspacer{} เอวญฺจ ปน ภิกฺขเว นิสฺสชฺชิตพฺพํ ฯเปฯ “อิทํ เม อยฺเย อญฺญทตฺถิเกน ปริกฺขาเรน อญฺญุทฺทิสิเกน สํฆิเกน อญฺญํ เจตาปิตํ นิสฺสคฺคิยํ, อิมาหํ สํฆสฺส นิสฺสชฺชามี”ติ ฯเปฯ ทเทยฺยาติ ฯเปฯ ทเทยฺยุนฺติ ฯเปฯ อยฺยาย ทมฺมีติ.}

\proseitem{761}{อญฺญทตฺถิเก อญฺญทตฺถิกสญฺญา อญฺญํ เจตาเปติ, นิสฺสคฺคิยํ ปาจิตฺติยํ.\csromanspacer{} อญฺญทตฺถิเก เวมติกา อญฺญํ เจตาเปติ, นิสฺสคฺคิยํ ปาจิตฺติยํ.\csromanspacer{} อญฺญทตฺถิเก อนญฺญทตฺถิกสญฺญา อญฺญํ เจตาเปติ, นิสฺสคฺคิยํ ปาจิตฺติยํ.\csromanspacer{} นิสฺสฏฺฐํ ปฏิลภิตฺวา ยถาทาเน อุปเนตพฺพํ.}

\prose{อนญฺญทตฺถิเก อญฺญทตฺถิกสญฺญา, อาปตฺติ ทุกฺกฏสฺส.\csromanspacer{} อนญฺญทตฺติเก เวมติกา, อาปตฺติ ทุกฺกฏสฺส.\csromanspacer{} อนญฺญทตฺถิเก อนญฺญทตฺถิกสญฺญา, อนาปตฺติ.}

\proseitem{762}{อนาปตฺติ เสสกํ อุปเนติ, สามิเก อปโลเกตฺวา อุปเนติ, อาปทาสุ อุมฺมตฺติกาย อาทิกมฺมิกายาติ.\csromanspacer{} ฉฏฺฐสิกฺขาปทํ นิฏฺฐิตํ.}

\csromanmatikaanchor{mtk.142}
\csromantocmarkref{subsection}{7. สตฺตมสิกฺขาปท}{mtk.142}
\bookmark[dest={mtk.142},level=2]{7. สตฺตมสิกฺขาปท}
\csromanheader{2}{1. \textbf{ปตฺตวคฺค 7. สตฺตมสิกฺขาปท}}
\proseitem{763}{\csromanfolio{329}เตน สมเยน พุทฺโธ ภควา สาวตฺถิยํ วิหรติ เชตวเน อนาถปิณฺฑิกสฺส อาราเม.\csromanspacer{} เตน โข ปน สมเยน อุปาสกา ภิกฺขุนิสํฆสฺส จีวรตฺถาย ฉนฺทกํ สงฺฆริตฺวา อญฺญตรสฺส ปาวาริกสฺส ฆเร ปริกฺขารํ นิกฺขิปิตฺวา ภิกฺขุนิโย อุปสงฺกมิตฺวา เอตทโวจุํ “อมุกสฺส อยฺเย ปาวาริกสฺส ฆเร จีวรตฺถาย ปริกฺขาโร นิกฺขิตฺโต, ตโต จีวรํ อาหราเปตฺวา ภาเชถา”ติ.\csromanspacer{} ภิกฺขุนิโย เตน จ ปริกฺขาเรน สยมฺปิ ยาจิตฺวา เภสชฺชํ เจตาเปตฺวา ปริภุญฺชิํสุ.\csromanspacer{} อุปาสกา ชานิตฺวา อุชฺฌายนฺติ ขิยฺยนฺติ วิปาเจนฺติ “กถํ หิ นาม ภิกฺขุนิโย อญฺญทตฺถิเกน ปริกฺขาเรน อญฺญุทฺทิสิเกน สํฆิเกน สญฺญาจิเกน อญฺญํ เจตาเปสฺสนฺตี”ติ ฯเปฯ.\csromanspacer{} สจฺจํ กิร ภิกฺขเว ภิกฺขุนิโย อญฺญทตฺถิเกน ปริกฺขาเรน อญฺญุทฺทิสิเกน สํฆิเกน สญฺญาจิเกน อญฺญํ เจตาเปนฺตีติ.\csromanspacer{} สจฺจํ ภควาติ.\csromanspacer{} วิครหิ พุทฺโธ ภควา ฯเปฯ.\csromanspacer{} กถํ หิ นาม ภิกฺขเว ภิกฺขุนิโย อญฺญทตฺถิเกน ปริกฺขาเรน อญฺญุทฺทิสิเกน สํฆิเกน สญฺญาจิเกน อญฺญํ เจตาเปสฺสนฺติ.\csromanspacer{} เนตํ ภิกฺขเว อปฺปสนฺนานํ วา ปสาทาย ฯเปฯ.\csromanspacer{} เอวญฺจ ปน ภิกฺขเว ภิกฺขุนิโย อิมํ สิกฺขาปทํ อุทฺทิสนฺตุ\csromandash{}}

\proseitem{764}{\textbf{“ยา ปน ภิกฺขุนี อญฺญทตฺถิเกน ปริกฺขาเรน อญฺญุทฺทิสิเกน สํฆิเกน สญฺญาจิเกน อญฺญํ เจตาเปยฺย, นิสฺสคฺคิยํ ปาจิตฺติยนฺ}”ติ.}

\proseitem{765}{\textbf{ยา ปนา}ติ ยา ยาทิสา ฯเปฯ.\csromanspacer{} \textbf{ภิกฺขุนี}ติ ฯเปฯ อยํ อิมสฺมิํ อตฺเถ อธิปฺเปตา ภิกฺขุนีติ.}

\prose{\textbf{อญฺญทตฺถิเกน ปริกฺขาเรน อญฺญุทฺทิสิเกนา}ติ อญฺญสฺสตฺถาย ทินฺเนน.}

\prose{\textbf{สํฆิเกนา}ติ สํฆสฺส, น คณสฺส, น เอกภิกฺขุนิยา.}

\prose{\textbf{สญฺญาจิเกนา}ติ สยํ ยาจิตฺวา.}

\prose{\textbf{อญฺญํ เจตาเปยฺยา}ติ ยํอตฺถาย ทินฺนํ, ตํ ฐเปตฺวา อญฺญํ เจตาเปติ, ปโยเค ทุกฺกฏํ.\csromanspacer{} ปฏิลาเภน นิสฺสคฺคิยํ โหติ.\csromanspacer{} นิสฺสชฺชิตพฺพํ สํฆสฺส วา คณสฺส วา เอกภิกฺขุนิยา วา.\csromanspacer{} เอวญฺจ ปน ภิกฺขเว นิสฺสชฺชิตพฺพํ ฯเปฯ “อิทํ เม อยฺเย อญฺญทตฺถิเกน ปริกฺขาเรน \csromanfolio{330}อญฺญุทฺทิสิเกน สํฆิเกน สญฺญาจิเกน อญฺญํ เจตาปิตํ นิสฺสคฺคิยํ, อิมาหํ สํฆสฺส นิสฺสชฺชามี”ติ ฯเปฯ ทเทยฺยาติ ฯเปฯ ทเทยฺยุนฺติ ฯเปฯ อยฺยาย ทมฺมีติ.}

\proseitem{766}{อญฺญทตฺถิเก อญฺญทตฺถิกสญฺญา อญฺญํ เจตาเปติ, นิสฺสคฺคิยํ ปาจิตฺติยํ.\csromanspacer{} อญฺญทตฺถิเก เวมติกา อญฺญํ เจตาเปติ, นิสฺสคฺคิยํ ปาจิตฺติยํ.\csromanspacer{} อญฺญทตฺถิเก อนญฺญทตฺถิกสญฺญา อญฺญํ เจตาเปติ, นิสฺสคฺคิยํ ปาจิตฺติยํ.\csromanspacer{} นิสฺสฏฺฐํ ปฏิลภิตฺวา ยถาทาเน อุปเนตพฺพํ.}

\prose{อนญฺญทตฺถิเก อญฺญทตฺถิกสญฺญา, อาปตฺติ ทุกฺกฏสฺส.\csromanspacer{} อนญฺญทตฺถิเก เวมติกา, อาปตฺติ ทุกฺกฏสฺส.\csromanspacer{} อนญฺญทตฺถิเก อนญฺญทตฺถิกสญฺญา, อนาปตฺติ.}

\proseitem{767}{อนาปตฺติ เสสกํ อุปเนติ, สามิเก อปโลเกตฺวา อุปเนติ, อาปทาสุ อุมฺมตฺติกาย อาทิกมฺมิกายาติ.\csromanspacer{} สตฺตมสิกฺขาปทํ นิฏฺฐิตํ.}
\csromansectionrule

\csromanmatikaanchor{mtk.143}
\csromantocmarkref{subsection}{8. อฏฺฐมสิกฺขาปท}{mtk.143}
\bookmark[dest={mtk.143},level=2]{8. อฏฺฐมสิกฺขาปท}
\csromanheader{2}{1. \textbf{ปตฺตวคฺค 8. อฏฺฐมสิกฺขาปท}}
\proseitem{768}{เตน สมเยน พุทฺโธ ภควา สาวตฺถิยํ วิหรติ เชตวเน อนาถปิณฺฑิกสฺส อาราเม.\csromanspacer{} เตน โข ปน สมเยน อญฺญตรสฺส ปูคสฺส ปริเวณวาสิกา ภิกฺขุนิโย ยาคุยา กิลมนฺติ.\csromanspacer{} อถ โข โส ปูโค ภิกฺขุนีนํ ยาคุอตฺถาย ฉนฺทกํ สงฺฆริตฺวา อญฺญตรสฺส อาปณิกสฺส ฆเร ปริกฺขารํ นิกฺขิปิตฺวา ภิกฺขุนิโย อุปสงฺกมิตฺวา เอตทโวจ “อมุกสฺส อยฺเย อาปณิกสฺส ฆเร ยาคุอตฺถาย ปริกฺขาโร นิกฺขิตฺโต, ตโต ตณฺฑุลํ อาหราเปตฺวา ยาคุํ ปจาเปตฺวา ปริภุญฺชถา”ติ.\csromanspacer{} ภิกฺขุนิโย เตน ปริกฺขาเรน เภสชฺชํ เจตาเปตฺวา ปริภุญฺชิํสุ.\csromanspacer{} โส ปูโค ชานิตฺวา อุชฺฌายติ ขิยฺยติ วิปาเจติ “กถํ หิ นาม ภิกฺขุนิโย อญฺญทตฺถิเกน ปริกฺขาเรน อญฺญุทฺทิสิเกน มหาชนิเกน อญฺญํ เจตาเปสฺสนฺตี”ติ ฯเปฯ.\csromanspacer{} สจฺจํ กิร ภิกฺขเว ภิกฺขุนิโย อญฺญทตฺถิเกน ปริกฺขาเรน อญฺญุทฺทิสิเกน \csromanfolio{331}\textbf{มหาชนิเกน อญฺญํ เจตาเปนฺตีติ.}\csromanspacer{}\textbf{ สจฺจํ ภควาติ.}\csromanspacer{}\textbf{ วิครหิ พุทฺโธ} \textbf{ภควา ฯเปฯ.}\csromanspacer{}\textbf{ กถํ หิ นาม ภิกฺขเว ภิกฺขุนิโย อญฺญทตฺถิเกน} ปริกฺขาเรน อญฺญุทฺทิสิเกน มหาชนิเกน อญฺญํ เจตาเปสฺสนฺติ.\csromanspacer{} เนตํ ภิกฺขเว อปฺปสนฺนานํ วา ปสาทาย ฯเปฯ.\csromanspacer{} เอวญฺจ ปน ภิกฺขเว ภิกฺขุนิโย อิมํ สิกฺขาปทํ อุทฺทิสนฺตุ\csromandash{}}

\proseitem{769}{\textbf{“ยา ปน ภิกฺขุนี อญฺญทตฺถิเกน ปริกฺขาเรน อญฺญุทฺทิสิเกน มหาชนิเกน อญฺญํ เจตาเปยฺย, นิสฺสคฺคิยํ ปาจิตฺติยนฺ}”ติ.}

\proseitem{770}{\textbf{ยา ปนา}ติ ยา ยาทิสา ฯเปฯ.\csromanspacer{} \textbf{ภิกฺขุนี}ติ ฯเปฯ อยํ อิมสฺมิํ อตฺเถ อธิปฺเปตา ภิกฺขุนีติ.}

\prose{\textbf{อญฺญทตฺถิเกน ปริกฺขาเรน อญฺญุทฺทิสิเกนา}ติ อญฺญสฺสตฺถาย ทินฺเนน.}

\prose{\textbf{มหาชนิเกนา}ติ คณสฺส, น สํฆสฺส, น เอกภิกฺขุนิยา.}

\prose{\textbf{อญฺญํ เจตาเปยฺยา}ติ ยํอตฺถาย ทินฺนํ, ตํ ฐเปตฺวา อญฺญํ เจตาเปติ, ปโยเค ทุกฺกฏํ.\csromanspacer{} ปฏิลาเภน นิสฺสคฺคิยํ โหติ.\csromanspacer{} นิสฺสชฺชิตพฺพํ สํฆสฺส วา คณสฺส วา เอกภิกฺขุนิยา วา.\csromanspacer{} เอวญฺจ ปน ภิกฺขเว นิสฺสชฺชิตพฺพํ ฯเปฯ “อิทํ เม อยฺเย อญฺญทตฺถิเกน ปริกฺขาเรน อญฺญุทฺทิสิเกน มหาชนิเกน อญฺญํ เจตาปิตํ นิสฺสคฺคิยํ, อิมาหํ สํฆสฺส นิสฺสชฺชามี”ติ ฯเปฯ ทเทยฺยาติ ฯเปฯ ทเทยฺยุนฺติ ฯเปฯ อยฺยาย ทมฺมีติ.}

\proseitem{771}{อญฺญทตฺถิเก อญฺญทตฺถิกสญฺญา อญฺญํ เจตาเปติ, นิสฺสคฺคิยํ ปาจิตฺติยํ.\csromanspacer{} อญฺญทตฺถิเก เวมติกา อญฺญํ เจตาเปติ, นิสฺสคฺคิยํ ปาจิตฺติยํ.\csromanspacer{} อญฺญทตฺถิเก อนญฺญทตฺถิกสญฺญา อญฺญํ เจตาเปติ, นิสฺสคฺคิยํ ปาจิตฺติยํ.\csromanspacer{} นิสฺสฏฺฐํ ปฏิลภิตฺวา ยถาทาเน อุปเนตพฺพํ.}

\prose{อนญฺญทตฺถิเก อญฺญทตฺถิกสญฺญา, อาปตฺติ ทุกฺกฏสฺส.\csromanspacer{} อนญฺญทตฺถิเก เวมติกา, อาปตฺติ ทุกฺกฏสฺส.\csromanspacer{} อนญฺญทตฺถิเก อนญฺญทตฺถิกสญฺญา, อนาปตฺติ.}

\proseitem{772}{อนาปตฺติ เสสกํ อุปเนติ, สามิเก อปโลเกตฺวา อุปเนติ, อาปทาสุ อุมฺมตฺติกาย อาทิกมฺมิกายาติ.\csromanspacer{} อฏฺฐมสิกฺขาปทํ นิฏฺฐิตํ.}

\csromanmatikaanchor{mtk.144}
\csromantocmarkref{subsection}{9. นวมสิกฺขาปท}{mtk.144}
\bookmark[dest={mtk.144},level=2]{9. นวมสิกฺขาปท}
\csromanheader{2}{1. \textbf{ปตฺตวคฺค 9. นวมสิกฺขาปท}}
\proseitem{773}{\csromanfolio{332}เตน สมเยน พุทฺโธ ภควา สาวตฺถิยํ วิหรติ เชตวเน อนาถปิณฺฑิกสฺส อาราเม.\csromanspacer{} เตน โข ปน สมเยน อญฺญตรสฺส ปูคสฺส ปริเวณวาสิกา ภิกฺขุนิโย ยาคุยา กิลมนฺติ.\csromanspacer{} อถ โข โส ปูโค ภิกฺขุนีนํ ยาคุอตฺถาย ฉนฺทกํ สงฺฆริตฺวา อญฺญตรสฺส อาปณิกสฺส ฆเร ปริกฺขารํ นิกฺขิปิตฺวา ภิกฺขุนิโย อุปสงฺกมิตฺวา เอตทโวจ “อมุกสฺส อยฺเย อาปณิกสฺส ฆเร ยาคุอตฺถาย ปริกฺขาโร นิกฺขิตฺโต, ตโต ตณฺฑุเล อาหราเปตฺวา ยาคุํ ปจาเปตฺวา ปริภุญฺชถา”ติ ภิกฺขุนิโย เตน จ ปริกฺขาเรน สยมฺปิ ยาจิตฺวา เภสชฺชํ เจตาเปตฺวา ปริภุญฺชิํสุ.\csromanspacer{} โส ปูโค ชานิตฺวา อุชฺฌายติ ขิยฺยติ วิปาเจติ “กถํ หิ นาม ภิกฺขุนิโย อญฺญทตฺถิเกน ปริกฺขาเรน อญฺญุทฺทิสิเกน มหาชนิเกน สญฺญาจิเกน อญฺญํ เจตาเปสฺสนฺตี”ติ ฯเปฯ.\csromanspacer{} สจฺจํ กิร ภิกฺขเว ภิกฺขุนิโย อญฺญทตฺถิเกน ปริกฺขาเรน อญฺญุทฺทิสิเกน มหาชนิเกน สญฺญาจิเกน อญฺญํ เจตาเปนฺตีติ.\csromanspacer{} สจฺจํ ภควาติ.\csromanspacer{} วิครหิ พุทฺโธ ภควา ฯเปฯ.\csromanspacer{} กถํ หิ นาม ภิกฺขเว ภิกฺขุนิโย อญฺญทตฺถิเกน ปริกฺขาเรน อญฺญุทฺทิสิเกน มหาชนิเกน สญฺญาจิเกน อญฺญํ เจตาเปสฺสนฺติ.\csromanspacer{} เนตํ ภิกฺขเว อปฺปสนฺนานํ วา ปสาทาย ฯเปฯ.\csromanspacer{} เอวญฺจ ปน ภิกฺขเว ภิกฺขุนิโย อิมํ สิกฺขาปทํ อุทฺทิสนฺตุ\csromandash{}}

\proseitem{774}{\textbf{“ยา ปน ภิกฺขุนี อญฺญทตฺถิเกน ปริกฺขาเรน อญฺญุทฺทิสิเกน มหาชนิเกน สญฺญาจิเกน อญฺญํ เจตาเปยฺย, นิสฺสคฺคิยํ ปาจิตฺติยนฺ}”ติ.}

\proseitem{775}{\textbf{ยา ปนา}ติ ยา ยาทิสา ฯเปฯ.\csromanspacer{} \textbf{ภิกฺขุนี}ติ ฯเปฯ อยํ อิมสฺมิํ อตฺเถ อธิปฺเปตา ภิกฺขุนีติ.}

\prose{\textbf{อญฺญทตฺถิเกน ปริกฺขาเรน อญฺญุทฺทิสิเกนา}ติ อญฺญสฺสตฺถาย ทินฺเนน.}

\prose{\textbf{มหาชนิเกนา}ติ คณสฺส, น สํฆสฺส, น เอกภิกฺขุนิยา.}

\prose{\textbf{สญฺญาจิเกนา}ติ สยํ ยาจิตฺวา.}

\prose{\textbf{อญฺญํ เจตาเปยฺยา}ติ ยํอตฺถาย ทินฺนํ, ตํ ฐเปตฺวา อญฺญํ เจตาเปติ, ปโยเค ทุกฺกฏํ.\csromanspacer{} ปฏิลาเภน นิสฺสคฺคิยํ โหติ.\csromanspacer{} นิสฺสชฺชิตพฺพํ สํฆสฺส \csromanfolio{333}วา คณสฺส วา เอกภิกฺขุนิยา วา.\csromanspacer{} เอวญฺจ ปน ภิกฺขเว นิสฺสชฺชิตพฺพํ ฯเปฯ “อิทํ เม อยฺเย อญฺญทตฺถิเกน ปริกฺขาเรน อญฺญุทฺทิสิเกน มหาชนิเกน สญฺญาจิเกน อญฺญํ เจตาปิตํ นิสฺสคฺคิยํ, อิมาหํ สํฆสฺส นิสฺสชฺชามี”ติ ฯเปฯ ทเทยฺยาติ ฯเปฯ ทเทยฺยุนฺติ ฯเปฯ อยฺยาย ทมฺมีติ.}

\proseitem{776}{อญฺญทตฺถิเก อญฺญทตฺถิกสญฺญา อญฺญํ เจตาเปติ, นิสฺสคฺคิยํ ปาจิตฺติยํ.\csromanspacer{} อญฺญทตฺถิเก เวมติกา อญฺญํ เจตาเปติ, นิสฺสคฺคิยํ ปาจิตฺติยํ.\csromanspacer{} อญฺญทตฺถิเก อนญฺญทตฺถิกสญฺญา อญฺญํ เจตาเปติ, นิสฺสคฺคิยํ ปาจิตฺติยํ.\csromanspacer{} นิสฺสฏฺฐํ ปฏิลภิตฺวา ยถาทาเน อุปเนตพฺพํ.}

\prose{อนญฺญทตฺถิเก อญฺญทตฺถิกสญฺญา, อาปตฺติ ทุกฺกฏสฺส.\csromanspacer{} อนญฺญทตฺถิเก เวมติกา, อาปตฺติ ทุกฺกฏสฺส.\csromanspacer{} อนญฺญทตฺถิเก อนญฺญทตฺถิกสญฺญา, อนาปตฺติ.}

\proseitem{777}{อนาปตฺติ เสสกํ อุปเนติ, สามิเก อปโลเกตฺวา อุปเนติ, อาปทาสุ อุมฺมตฺติกาย อาทิกมฺมิกายาติ.\csromanspacer{} นวมสิกฺขาปทํ นิฏฺฐิตํ.}
\csromansectionrule

\csromanmatikaanchor{mtk.145}
\csromantocmarkref{subsection}{10. ทสมสิกฺขาปท}{mtk.145}
\bookmark[dest={mtk.145},level=2]{10. ทสมสิกฺขาปท}
\csromanheader{2}{1. \textbf{ปตฺตวคฺค 1}0. ทสมสิกฺขาปท}
\proseitem{778}{เตน สมเยน พุทฺโธ ภควา สาวตฺถิยํ วิหรติ เชตวเน อนาถปิณฺฑิกสฺส อาราเม.\csromanspacer{} เตน โข ปน สมเยน ถุลฺลนนฺทา ภิกฺขุนี พหุสฺสุตา โหติ ภาณิกา วิสารทา ปฏฺฏา ธมฺมิํ กถํ กาตุํ.\csromanspacer{} พหู มนุสฺสา ถุลฺลนนฺทํ ภิกฺขุนิํ ปยิรุปาสนฺติ.\csromanspacer{} เตน โข ปน สมเยน ถุลฺลนนฺทาย ภิกฺขุนิยา ปริเวณํ อุนฺทฺริยติ\footnote{อุทฺรียติ (สฺยา)}.\csromanspacer{} มนุสฺสา ถุลฺลนนฺทํ ภิกฺขุนิํ เอตทโวจุํ “กิสฺสิทํ เต อยฺเย ปริเวณํ อุนฺทฺริยตี”ติ.\csromanspacer{} นตฺถาวุโส ทายกา, นตฺถิ การกาติ.\csromanspacer{} อถ โข เต มนุสฺสา ถุลฺลนนฺทาย ภิกฺขุนิยา ปริเวณตฺถาย ฉนฺทกํ สงฺฆริตฺวา ถุลฺลนนฺทาย ภิกฺขุนิยา ปริกฺขารํ อทํสุ.\csromanspacer{} ถุลฺลนนฺทา ภิกฺขุนี เตน จ ปริกฺขาเรน สยมฺปิ ยาจิตฺวา เภสชฺชํ เจตาเปตฺวา ปริภุญฺชิ.\csromanspacer{} มนุสฺสา ชานิตฺวา อุชฺฌายนฺติ ขิยฺยนฺติ วิปาเจนฺติ “กถํ หิ นาม อยฺยา ถุลฺลนนฺทา อญฺญทตฺถิเกน ปริกฺขาเรน อญฺญุทฺทิสิเกน ปุคฺคลิเกน สญฺญาจิเกน อญฺญํ เจตาเปสฺสตี”ติ ฯเปฯ \csromanfolio{334}\textbf{สจฺจํ กิร ภิกฺขเว ถุลฺลนนฺทา ภิกฺขุนี อญฺญทตฺถิเกน ปริกฺขาเรน} \textbf{อญฺญุทฺทิสิเกน ปุคฺคลิเกน สญฺญาจิเกน อญฺญํ เจตาเปตีติ.}\csromanspacer{}\textbf{ สจฺจํ} ภควาติ.\csromanspacer{} วิครหิ พุทฺโธ ภควา ฯเปฯ.\csromanspacer{} กถํ หิ นาม ภิกฺขเว ถุลฺลนนฺทา ภิกฺขุนี อญฺญทตฺถิเกน ปริกฺขาเรน อญฺญุทฺทิสิเกน ปุคฺคลิเกน สญฺญาจิเกน อญฺญํ เจตาเปสฺสติ.\csromanspacer{} เนตํ ภิกฺขเว อปฺปสนฺนานํ วา ปสาทาย ฯเปฯ.\csromanspacer{} เอวญฺจ ปน ภิกฺขเว ภิกฺขุนิโย อิมํ สิกฺขาปทํ อุทฺทิสนฺตุ\csromandash{}}

\proseitem{779}{\textbf{“ยา ปน ภิกฺขุนี อญฺญทตฺถิเกน ปริกฺขาเรน อญฺญุทฺทิสิเกน ปุคฺคลิเกน สญฺญาจิเกน อญฺญํ เจตาเปยฺย, นิสฺสคฺคิยํ ปาจิตฺติยนฺ”}ติ.}

\proseitem{780}{\textbf{ยา ปนา}ติ ยา ยาทิสา ฯเปฯ.\csromanspacer{} \textbf{ภิกฺขุนี}ติ ฯเปฯ อยํ อิมสฺมิํ อตฺเถ อธิปฺเปตา ภิกฺขุนีติ.}

\prose{\textbf{อญฺญทตฺถิเกน ปริกฺขาเรน อญฺญุทฺทิสิเกนา}ติ อญฺญสฺสตฺถาย ทินฺเนน.}

\prose{\textbf{ปุคฺคลิเกนา}ติ เอกาย ภิกฺขุนิยา, น สํฆสฺส, น คณสฺส.}

\prose{\textbf{สญฺญาจิเกนา}ติ สยํ ยาจิตฺวา.}

\prose{\textbf{อญฺญํ เจตาเปยฺยา}ติ ยํอตฺถาย ทินฺนํ, ตํ ฐเปตฺวา อญฺญํ เจตาเปติ, ปโยเค ทุกฺกฏํ.\csromanspacer{} ปฏิลาเภน นิสฺสคฺคิยํ โหติ.\csromanspacer{} นิสฺสชฺชิตพฺพํ สํฆสฺส วา คณสฺส วา เอกภิกฺขุนิยา วา.\csromanspacer{} เอวญฺจ ปน ภิกฺขเว นิสฺสชฺชิตพฺพํ ฯเปฯ “อิทํ เม อยฺเย อญฺญทตฺถิเกน ปริกฺขาเรน อญฺญุทฺทิสิเกน ปุคฺคลิเกน สญฺญาจิเกน อญฺญํ เจตาปิตํ นิสฺสคฺคิยํ, อิมาหํ สํฆสฺส นิสฺสชฺชามี”ติ ฯเปฯ ทเทยฺยาติ ฯเปฯ ทเทยฺยุนฺติ ฯเปฯ อยฺยาย ทมฺมีติ.}

\proseitem{781}{อญฺญทตฺถิเก อญฺญทตฺถิกสญฺญา อญฺญํ เจตาเปติ, นิสฺสคฺคิยํ ปาจิตฺติยํ.\csromanspacer{} อญฺญทตฺถิเก เวมติกา อญฺญํ เจตาเปติ, นิสฺสคฺคิยํ ปาจิตฺติยํ.\csromanspacer{} อญฺญทตฺถิเก อนญฺญทตฺถิกสญฺญา อญฺญํ เจตาเปติ, นิสฺสคฺคิยํ ปาจิตฺติยํ.\csromanspacer{} นิสฺสฏฺฐํ ปฏิลภิตฺวา ยถาทาเน อุปเนตพฺพํ.}

\prose{อนญฺญทตฺถิเก อญฺญทตฺถิกสญฺญา, อาปตฺติ ทุกฺกฏสฺส.\csromanspacer{} อนญฺญทตฺถิเก เวมติกา, อาปตฺติ ทุกฺกฏสฺส.\csromanspacer{} อนญฺญทตฺถิเก อนญฺญทตฺถิกสญฺญา, อนาปตฺติ.}

\proseitem{782}{อนาปตฺติ เสสกํ อุปเนติ, สามิเก อปโลเกตฺวา อุปเนติ, อาปทาสุ อุมฺมตฺติกาย อาทิกมฺมิกายาติ.\csromanspacer{} ทสมสิกฺขาปทํ นิฏฺฐิตํ.}

\csromanmatikaanchor{mtk.146}
\csromantocmarkref{subsection}{11. เอกาทสมสิกฺขาปท}{mtk.146}
\bookmark[dest={mtk.146},level=2]{11. เอกาทสมสิกฺขาปท}
\csromanheader{2}{1. ปตฺตวคฺค 11. เอกาทสมสิกฺขาปท}
\proseitem{783}{\csromanfolio{335}เตน สมเยน พุทฺโธ ภควา สาวตฺถิยํ วิหรติ เชตวเน อนาถปิณฺฑิกสฺส อาราเม.\csromanspacer{} เตน โข ปน สมเยน ถุลฺลนนฺทา ภิกฺขุนี พหุสฺสุตา โหติ ภาณิกา วิสารทา ปฏฺฏา ธมฺมิํ กถํ กาตุํ.\csromanspacer{} อถ โข ราชา ปเสนทิ โกสโล สีตกาเล มหคฺฆํ กมฺพลํ ปารุปิตฺวา เยน ถุลฺลนนฺทา ภิกฺขุนี เตนุปสงฺกมิ, อุปสงฺกมิตฺวา ถุลฺลนนฺทํ ภิกฺขุนิํ อภิวาเทตฺวา เอกมนฺตํ นิสีทิ, เอกมนฺตํ นิสินฺนํ โข ราชานํ ปเสนทิํ โกสลํ ถุลฺลนนฺทา ภิกฺขุนี ธมฺมิยา กถาย สนฺทสฺเสสิ สมาทเปสิ สมุตฺเตเชสิ สมฺปหํเสสิ.\csromanspacer{} อถ โข ราชา ปเสนทิ โกสโล ถุลฺลนฺทาย ภิกฺขุนิยา ธมฺมิยา กถาย สนฺทสฺสิโต สมาทปิโต สมุตฺเตชิโต สมฺปหํสิโต ถุลฺลนนฺทํ ภิกฺขุนิํ เอตทโวจ ‘วเทยฺยาสิ อยฺเย เยน อตฺโถ”ติ.\csromanspacer{} สเจ เม ตฺวํ มหาราช ทาตุกาโมสิ, อิมํ กมฺพลํ เทหีติ.\csromanspacer{} อถ โข ราชา ปเสนทิ โกสโล ถุลฺลนนฺทาย ภิกฺขุนิยา กมฺพลํ ทตฺวา อุฏฺฐายาสนา ถุลฺลนนฺทํ ภิกฺขุนิํ อภิวาเทตฺวา ปทกฺขิณํ กตฺวา ปกฺกามิ.\csromanspacer{} มนุสฺสา อุชฺฌายนฺติ ขิยฺยนฺติ วิปาเจนฺติ “มหิจฺฉา อิมา ภิกฺขุนิโย อสนฺตุฏฺฐา, กถํ หิ นาม ราชานํ กมฺพลํ วิญฺญาเปสฺสนฺตี”ติ.\csromanspacer{} อสฺโสสุํ โข ภิกฺขุนิโย เตสํ มนุสฺสานํ อุชฺฌายนฺตานํ ขิยฺยนฺตานํ วิปาเจนฺตานํ.\csromanspacer{} ยา ตา ภิกฺขุนิโย อปฺปิจฺฉา ฯเปฯ ตา อุชฺฌายนฺติ ขิยฺยนฺติ วิปาเจนฺติ “กถํ หิ นาม อยฺยา ถุลฺลนนฺทา ราชานํ กมฺพลํ วิญฺญาเปสฺสตี”ติ ฯเปฯ.\csromanspacer{} สจฺจํ กิร ภิกฺขเว ถุลฺลนนฺทา ภิกฺขุนี ราชานํ กมฺพลํ วิญฺญาเปตีติ\footnote{วิญฺญาเปสีติ (ก)}.\csromanspacer{} สจฺจํ ภควาติ.\csromanspacer{} วิครหิ พุทฺโธ ภควา ฯเปฯ.\csromanspacer{} กถํ หิ นาม ภิกฺขเว ถุลฺลนนฺทา ภิกฺขุนี ราชานํ กมฺพลํ วิญฺญาเปสฺสติ.\csromanspacer{} เนตํ ภิกฺขเว อปฺปสนฺนานํ วา ปสาทาย ฯเปฯ.\csromanspacer{} เอวญฺจ ปน ภิกฺขเว ภิกฺขุนิโย อิมํ สิกฺขาปทํ อุทฺทิสนฺตุ\csromandash{}}

\proseitem{784}{\textbf{“ครุปาวุรณํ}\footnote{ปาปุรณํ (สี, สฺยา)} \textbf{ปน ภิกฺขุนิยา เจตาเปนฺติยา จตุกฺกํสปรมํ เจตาเปตพฺพํ, ตโต เจ อุตฺตริ เจตาเปยฺย, นิสฺสคฺคิยํ ปาจิตฺติยนฺ”}ติ.}

\proseitem{785}{\csromanfolio{336}\textbf{ครุปาวุรณํ} นาม ยํ กิญฺจิ สีตกาเล ปาวุรณํ.}

\prose{\textbf{เจตาเปนฺติยา}ติ วิญฺญาเปนฺติยา.}

\prose{\textbf{จตุกฺกํสปรมํ เจตาเปตพฺพนฺ}ติ โสฬสกหาปณคฺฆนกํ เจตาเปตพฺพํ.}

\prose{\textbf{ตโต เจ อุตฺตริ เจตาเปยฺยา}ติ ตตุตฺตริ วิญฺญาเปติ ปโยเค ทุกฺกฏํ.\csromanspacer{} ปฏิลาเภน นิสฺสคฺคิยํ โหติ.\csromanspacer{} นิสฺสชฺชิตพฺพํ สํฆสฺส วา คณสฺส วา เอกภิกฺขุนิยา วา.\csromanspacer{} เอวญฺจ ปน ภิกฺขเว นิสฺสชฺชิตพฺพํ ฯเปฯ “อิทํ เม อยฺเย ครุปาวุรณํ อติเรกจตุกฺกํสปรมํ เจตาปิตํ นิสฺสคฺคิยํ, อิมาหํ สํฆสฺส นิสฺสชฺชามี”ติ ฯเปฯ ทเทยฺยาติ ฯเปฯ ทเทยฺยุนฺติ ฯเปฯ อยฺยาย ทมฺมีติ.}

\proseitem{786}{อติเรกจตุกฺกํเส อติเรกสญฺญา เจตาเปติ, นิสฺสคฺคิยํ ปาจิตฺติยํ.\csromanspacer{} อติเรกจตุกฺกํเส เวมติกา เจตาเปติ, นิสฺสคฺคิยํ ปาจิตฺติยํ.\csromanspacer{} อติเรกจตุกฺกํเส อูนกสญฺญา เจตาเปติ, นิสฺสคฺคิยํ ปาจิตฺติยํ.}

\prose{อูนกจตุกฺกํเส อติเรกสฺสญฺญา, อาปตฺติ ทุกฺกฏสฺส.\csromanspacer{} อูนกจตุกฺกํเส เวมติกา, อาปตฺติ ทุกฺกฏสฺส.\csromanspacer{} อูนกจตุกฺกํเส อูนกสญฺญา, อนาปตฺติ.}

\proseitem{787}{อนาปตฺติ จตุกฺกํสปรมํ เจตาเปติ, อูนกจตุกฺกํสปรมํ เจตาเปติ, ญาตกานํ, ปวาริตานํ, อญฺญสฺสตฺถาย, อตฺตโน ธเนน, มหคฺฆํ เจตาเปตุกามสฺส อปฺปคฺฆํ เจตาเปติ, อุมฺมตฺติกาย อาทิกมฺมิกายาติ.\csromanspacer{} เอกาทสมสิกฺขาปทํ นิฏฺฐิตํ.}
\csromansectionrule

\csromanmatikaanchor{mtk.147}
\csromantocmarkref{subsection}{12. ทฺวาทสมสิกฺขาปท}{mtk.147}
\bookmark[dest={mtk.147},level=2]{12. ทฺวาทสมสิกฺขาปท}
\csromanheader{2}{1. ปตฺตวคฺค 12. ทฺวาทสมสิกฺขาปท}
\proseitem{788}{เตน สมเยน พุทฺโธ ภควา สาวตฺถิยํ วิหรติ เชตวเน อนาถปิณฺฑิกสฺส อาราเม.\csromanspacer{} เตน โข ปน สมเยน ถุลฺลนนฺทา ภิกฺขุนี พหุสฺสุตา โหติ ภาณิกา วิสารทา ปฏฺฏา ธมฺมิํ กถํ กาตุํ.\csromanspacer{} อถ โข ราชา ปเสนทิ โกสโล อุณฺหกาเล มหคฺฆํ โขมํ ปารุปิตฺวา เยน ถุลฺลนนฺทา ภิกฺขุนี เตนุปสงฺกมิ, อุปสงฺกมิตฺวา \csromanfolio{337}ถุลฺลนนฺทํ ภิกฺขุนิํ อภิวาเทตฺวา เอกมนฺตํ นิสีทิ, เอกมนฺตํ นิสินฺนํ โข ราชานํ ปเสนทิํ โกสลํ ถุลฺลนนฺทา ภิกฺขุนี ธมฺมิยา กถาย สนฺทสฺเสสิ สมาทเปสิ สมุตฺเตเชสิ สมฺปหํเสสิ.\csromanspacer{} อถ โข ราชา ปเสนทิ โกสโล ถุลฺลนนฺทาย ภิกฺขุนิยา ธมฺมิยา กถาย สนฺทสฺสิโต สมาทปิโต สมุตฺเตชิโต สมฺปหํสิโต ถุลฺลนนฺทํ ภิกฺขุนิํ เอตทโวจ “วเทยฺยาสิ อยฺเย เยน อตฺโถ”ติ.\csromanspacer{} สเจ เม ตฺวํ มหาราช ทาตุกาโมสิ, อิมํ โขมํ เทหีติ.\csromanspacer{} อถ โข ราชา ปเสนทิ โกสโล ถุลฺลนนฺทาย ภิกฺขุนิยา โขมํ ทตฺวา อุฏฺฐายาสนา ถุลฺลนนฺทํ ภิกฺขุนิํ อภิวาเทตฺวา ปทกฺขิณํ กตฺวา ปกฺกามิ.\csromanspacer{} มนุสฺสา อุชฺฌายนฺติ ขิยฺยนฺติ วิปาเจนฺติ “มหิจฺฉา อิมา ภิกฺขุนิโย อสนฺตุฏฺฐา, กถํ หิ นาม ราชานํ โขมํ วิญฺญาเปสฺสนฺตี”ติ.\csromanspacer{} อสฺโสสุํ โข ภิกฺขุนิโย เตสํ มนุสฺสานํ อุชฺฌายนฺตานํ ขิยฺยนฺตานํ วิปาเจนฺตานํ.\csromanspacer{} ยา ตา ภิกฺขุนิโย อปฺปิจฺฉา ฯเปฯ ตา อุชฺฌายนฺติ ขิยฺยนฺติ วิปาเจนฺติ “กถํ หิ นาม อยฺยา ถุลฺลนนฺทา ราชานํ โขมํ วิญฺญาเปสฺสตี”ติ ฯเปฯ.\csromanspacer{} สจฺจํ กิร ภิกฺขเว ถุลฺลนนฺทา ภิกฺขุนี ราชานํ โขมํ วิญฺญาเปตีติ\footnote{วิญฺญาเปสีติ (ก)}.\csromanspacer{} สจฺจํ ภควาติ.\csromanspacer{} วิครหิ พุทฺโธ ภควา ฯเปฯ.\csromanspacer{} กถํ หิ นาม ภิกฺขเว ถุลฺลนนฺทา ภิกฺขุนี ราชานํ โขมํ วิญฺญาเปสฺสติ.\csromanspacer{} เนตํ ภิกฺขเว อปฺปสนฺนานํ วา ปสาทาย ฯเปฯ.\csromanspacer{} เอวญฺจ ปน ภิกฺขเว ภิกฺขุนิโย อิมํ สิกฺขาปทํ อุทฺทิสนฺตุ\csromandash{}}

\proseitem{789}{\textbf{“ลหุปาวุรณํ ปน ภิกฺขุนิยา เจตาเปนฺติยา อฑฺฒเตยฺยกสปรมํ เจตาเปตพฺพํ, ตโต เจ อุตฺตริ เจตาเปยฺย, นิสฺสคฺคิยํ ปาจิตฺติยนฺ”}ติ.}

\proseitem{790}{\textbf{ลหุปาวุรณํ} นาม ยํ กิญฺจิ อุณฺหกาเล ปาวุรณํ.}

\prose{\textbf{เจตาเปนฺติยา}ติ วิญฺญาเปนฺติยา.}

\prose{\textbf{อฑฺฒเตยฺยกํสปรมํ เจตาเปตพฺพนฺ}ติ ทสกหาปณคฺฆนกํ เจตาเปตพฺพํ.}

\prose{\textbf{ตโต เจ อุตฺตริ เจตาเปยฺยา}ติ ตตุตฺตริ วิญฺญาเปติ, ปโยเค ทุกฺกฏํ.\csromanspacer{} ปฏิลาเภน นิสฺสคฺคิยํ โหติ.\csromanspacer{} นิสฺสชฺชิตพฺพํ สํฆสฺส วา คณสฺส วา เอกภิกฺขุนิยา วา.\csromanspacer{} เอวญฺจ ปน ภิกฺขเว นิสฺสชฺชิตพฺพํ ฯเปฯ “อิทํ เม อยฺเย ลหุปาวุรณํ อติเรก อฑฺฒเตยฺยกํสปรมํ เจตาปิตํ นิสฺสคฺคิยํ, \csromanfolio{338}อิมาหํ สํฆสฺส นิสฺสชฺชามี”ติ ฯเปฯ ทเทยฺยาติ ฯเปฯ ทเทยฺยุนฺติ ฯเปฯ อยฺยาย ทมฺมีติ.}

\proseitem{791}{อติเรกอฑฺฒเตยฺยกํเส อติเรกสญฺญา เจตาเปติ, นิสฺสคฺคิยํ ปาจิตฺติยํ.\csromanspacer{} อติเรก อฑฺฒเตยฺยกํเส เวมติกา เจตาเปติ, นิสฺสคฺคิยํ ปาจิตฺติยํ.\csromanspacer{} อติเรก อฑฺฒเตยฺยกํเส อูนกสญฺญา เจตาเปติ, นิสฺสคฺคิยํ ปาจิตฺติยํ.}

\prose{อูนกอฑฺฒเตยฺยกํเส อติเรกสญฺญา, อาปตฺติ ทุกฺกฏสฺส.\csromanspacer{} อูนก อฑฺฒเตยฺยกํเส เวมติกา, อาปตฺติ ทุกฺกฏสฺส.\csromanspacer{} อูนกอฑฺฒเตยฺยกํเส อูนกสญฺญา, อนาปตฺติ.}

\proseitemwithcloser{792}{อนาปตฺติ อฑฺฒเตยฺยกํสปรมํ เจตาเปติ, อูนก อฑฺฒเตยฺยกํสปรมํ เจตาเปติ, ญาตกานํ, ปวาริตานํ, อญฺญสฺสตฺถาย, อตฺตโน ธเนน, มหคฺฆํ เจตาเปตุกามสฺส อปฺปคฺฆํ เจตาเปติ, อุมฺมตฺติกาย อาทิกมฺมิกายาติ.}{ทฺวาทสมสิกฺขาปทํ นิฏฺฐิตํ.}
\prose{อุทฺทิฏฺฐา โข อยฺยาโย ติํส นิสฺสคฺคิยา ปาจิตฺติยา ธมฺมา.\csromanspacer{} ตตฺถายฺยาโย ปุจฺฉามิ กจฺจิตฺถ ปริสุทฺธา, ทุติยมฺปิ ปุจฺฉามิ กจฺจิตฺถ ปริสุทฺธา ตติยมฺปิ ปุจฺฉามิ กจฺจิตฺถ ปริสุทฺธา, ปริสุทฺเธตฺถายฺยาโย, ตสฺมา ตุณฺหี, เอวเมตํ ธารยามีติ.}

\csromanheader{2}{\textbf{ภิกฺขุนิวิภงฺเค}}
\nitthitam{\textbf{นิสฺสคฺคิยกณฺฑํ นิฏฺฐิตํ.}}
\chapterheadpageread{4. \textbf{ปาจิตฺติยกณฺฑ}}
\csromanmatikaanchor{mtk.148}
\csromanmatikaanchor{mtk.149}
\csromanmatikaanchor{mtk.150}
\csromantocmarknopage{chapter}{4. ปาจิตฺติยกณฺฑ}{mtk.148}
\bookmark[dest={mtk.148},level=0]{4. ปาจิตฺติยกณฺฑ}
\csromantocmarknopage{section}{1. ลสุณวคฺค}{mtk.149}
\bookmark[dest={mtk.149},level=1]{1. ลสุณวคฺค}
\csromantocmarkref{subsection}{1. ปฐมสิกฺขาปท}{mtk.150}
\bookmark[dest={mtk.150},level=2]{1. ปฐมสิกฺขาปท}
\setcsromanheadcha{4. ปาจิตฺติยกณฺฑ}
\setcsromanheadhclear{}
\setcsromanheadh{1. ลสุณวคฺค}
\csromanheader{2}{1. \textbf{ลสุณวคฺค 1. ปฐมสิกฺขาปท}}
\prose{\csromanfolio{339}อิเม โข ปนายฺยาโย ฉสฏฺฐิสตา ปาจิตฺติยา ธมฺมา อุทฺเทสํ อาคจฺฉนฺติ.}

\proseitem{793}{เตน สมเยน พุทฺโธ ภควา สาวตฺถิยํ วิหรติ เชตวเน อนาถปิณฺฑิกสฺส อาราเม.\csromanspacer{} เตน โข ปน สมเยน อญฺญตเรน อุปาสเกน ภิกฺขุนิสํโฆ ลสุเณน ปวาริโต โหติ “ยาสํ อยฺยานํ ลสุเณน อตฺโถ, อหํ ลสุเณนา”ติ.\csromanspacer{} เขตฺตปาโล จ อาณตฺโต โหติ “สเจ ภิกฺขุนิโย อาคจฺฉนฺติ, เอกเมกาย ภิกฺขุนิยา เทฺวตโย ภณฺฑิเก เทหี”ติ.\csromanspacer{} เตน โข ปน สมเยน สาวตฺถิยํ อุสฺสโว โหติ, ยถาภตํ ลสุณํ ปริกฺขยํ อคมาสิ.\csromanspacer{} ภิกฺขุนิโย ตํ อุปาสกํ อุปสงฺกมิตฺวา เอตทโวจุํ “ลสุเณน อาวุโส อตฺโถ”ติ.\csromanspacer{} นตฺถายฺเย ยถาภตํ ลสุณํ ปริกฺขีณํ, เขตฺตํ คจฺฉถาติ.\csromanspacer{} ถุลฺลนนฺทา ภิกฺขุนี เขตฺตํ คนฺตฺวา น มตฺตํ ชานิตฺวา พหุํ ลสุณํ หราเปสิ.\csromanspacer{} เขตฺตปาโล อุชฺฌายติ ขิยฺยติ วิปาเจติ “กถํ หิ นาม ภิกฺขุนิโย น มตฺตํ ชานิตฺวา พหุํ ลสุณํ หราเปสฺสนฺตี”ติ.\csromanspacer{} อสฺโสสุํ โข ภิกฺขุนิโย ตสฺส เขตฺตปาลสฺส อุชฺฌายนฺตสฺส ขิยฺยนฺตสฺส วิปาเจนฺตสฺส.\csromanspacer{} ยา ตา ภิกฺขุนิโย อปฺปิจฺฉา ฯเปฯ ตา อุชฺฌายนฺติ ขิยฺยนฺติ วิปาเจนฺติ “กถํ หิ นาม อยฺยา ถุลฺลนนฺทา น มตฺตํ ชานิตฺวา พหุํ ลสุณํ หราเปสฺสตี”ติ ฯเปฯ.\csromanspacer{} สจฺจํ กิร ภิกฺขเว ถุลฺลนนฺทา ภิกฺขุนี น มตฺตํ ชานิตฺวา พหุํ ลสุณํ หราเปตีติ\footnote{หราเปสีติ (ก)}.\csromanspacer{} สจฺจํ ภควาติ.\csromanspacer{} วิครหิ พุทฺโธ ภควา ฯเปฯ.\csromanspacer{} กถํ หิ นาม ภิกฺขเว ถุลฺลนนฺทา ภิกฺขุนี น มตฺตํ ชานิตฺวา พหุํ ลสุณํ หราเปสฺสติ.\csromanspacer{} เนตํ ภิกฺขเว อปฺปสนฺนานํ วา ปสาทาย ฯเปฯ ธมฺมิํ กถํ กตฺวา ภิกฺขู อามนฺเตสิ\csromandash{}}

\prose{ภูตปุพฺพํ ภิกฺขเว ถุลฺลนนฺทา ภิกฺขุนี อญฺญตรสฺส พฺราหฺมณสฺส ปชาปติ อโหสิ, ติสฺโส จ ธีตโร นนฺทา นนฺทวตี สุนฺทรีนนฺทา.\csromanspacer{} อถ โข ภิกฺขเว โส พฺราหฺมโณ กาลงฺกตฺวา อญฺญตรํ หํสโยนิํ อุปปชฺชิ, ตสฺส สพฺพโสวณฺณมยา ปตฺตา อเหสุํ, โส ตาสํ เอเกกํ ปตฺตํ เทติ.\csromanspacer{} อถ โข ภิกฺขเว ถุลฺลนนฺทา ภิกฺขุนี \csromanfolio{340}“\textbf{อยํ หํโส อมฺหากํ เอเกกํ ปตฺตํ เทตี”ติ ตํ หํสราชํ คเหตฺวา} \textbf{นิปฺปตฺตํ อกาสิ, ตสฺส ปุน ชายมานา ปตฺตา เสตา สมฺปชฺชิํสุ.}\csromanspacer{}\textbf{ ตทาปิ} ภิกฺขเว ถุลฺลนนฺทา ภิกฺขุนี อติโลเภน สุวณฺณา ปริหีนา, อิทานิ ลสุณา ปริหายิสฺสตีติ.}

\setlength{\csromangathapagewidth}{0pt}
\csromangathameasuregroup{*ยํ ลทฺธํ เตน ตุฏฺฐพฺพํ, \\ หํสราชํ คเหตฺวาน,}{\csromangathabat{*ยํ ลทฺธํ เตน ตุฏฺฐพฺพํ,}{อติโลโภ หิ ปาปโก.} \\ \csromangathabat{หํสราชํ คเหตฺวาน,}{สุวณฺณา ปริหายถาติ.}}
\csromangathasetleft{*ยํ ลทฺธํ เตน ตุฏฺฐพฺพํ, \\ หํสราชํ คเหตฺวาน,}
\csromangathagroup{\csromangathastanza{\csromangathabat{\csromansymbolfootnote{*}{ขุ 5. 32 ปิฏฺเฐ สุวณฺณหํสชาตเกปิ.}ยํ ลทฺธํ เตน ตุฏฺฐพฺพํ,}{อติโลโภ หิ ปาปโก.} \\ \csromangathabat{หํสราชํ คเหตฺวาน,}{สุวณฺณา ปริหายถาติ.}}}

\prose{อถ โข ภควา ถุลฺลนนฺทํ อเนกปริยาเยน วิครหิตฺวา ทุพฺภรตาย ฯเปฯ.\csromanspacer{} เอวญฺจ ปน ภิกฺขเว ภิกฺขุนิโย อิมํ สิกฺขาปทํ อุทฺทิสนฺตุ\csromandash{}}

\proseitem{794}{\textbf{“ยา ปน ภิกฺขุนี ลสุณํ ขาเทยฺย, ปาจิตฺติยนฺ”}ติ.}

\proseitem{795}{\textbf{ยา ปนา}ติ ยา ยาทิสา ฯเปฯ.\csromanspacer{} \textbf{ภิกฺขุนี}ติ ฯเปฯ อยํ อิมสฺมิํ อตฺเถ อธิปฺเปตา ภิกฺขุนีติ.}

\prose{\textbf{ลสุณํ} นาม มาคธกํ วุจฺจติ.}

\prose{“ขาทิสฺสามี”ติ ปฏิคฺคณฺหาติ, อาปตฺติ ทุกฺกฏสฺส.\csromanspacer{} อชฺโฌหาเร อชฺโฌหาเร อาปตฺติ ปาจิตฺติยสฺส.}

\proseitem{796}{ลสุเณ ลสุณสญฺญา ขาทติ, อาปตฺติ ปาจิตฺติยสฺส.\csromanspacer{} ลสุเณ เวมติกา ขาทติ, อาปตฺติ ปาจิตฺติยสฺส.\csromanspacer{} ลสุเณ อลสุณสญฺญา ขาทติ, อาปตฺติ ปาจิตฺติยสฺส.}

\prose{อลสุเณ ลสุณสญฺญา ขาทติ, อาปตฺติ ทุกฺกฏสฺส.\csromanspacer{} อลสุเณ เวมติกา ขาทติ, อาปตฺติ ทุกฺกฏสฺส.\csromanspacer{} อลสุเณ อลสุณสญฺญา ขาทติ, อนาปตฺติ.}

\proseitem{797}{อนาปตฺติ ปลณฺฑุเก ภญฺชนเก หรีตเก จาปลสุเณ สูปสมฺปาเก มํสสมฺปาเก เตลสมฺปาเก สาฬเว อุตฺตริภงฺเค อุมฺมตฺติกาย อาทิกมฺมิกายาติ.\csromanspacer{} ลสุณสิกฺขาปทํ ปฐมํ นิฏฺฐิตํ.}

\csromanmatikaanchor{mtk.151}
\csromantocmarkref{subsection}{2. ทุติยสิกฺขาปท}{mtk.151}
\bookmark[dest={mtk.151},level=2]{2. ทุติยสิกฺขาปท}
\csromanheader{2}{1. \textbf{ลสุณวคฺค 2. ทุติยสิกฺขาปท}}
\proseitem{798}{\csromanfolio{341}เตน สมเยน พุทฺโธ ภควา สาวตฺถิยํ วิหรติ เชตวเน อนาถปิณฺฑิกสฺส อาราเม.\csromanspacer{} เตน โข ปน สมเยน ฉพฺพคฺคิยา ภิกฺขุนิโย สมฺพาเธ โลมํ สํหราเปตฺวา อจิรวติยา นทิยา เวสิยาหิ สทฺธิํ นคฺคา เอกติตฺเถ นหายนฺติ.\csromanspacer{} เวสิยา อุชฺฌายนฺติ ขิยฺยนฺติ วิปาเจนฺติ “กถํ หิ นาม ภิกฺขุนิโย สมฺพาเธ โลมํ สํหราเปสฺสนฺติ เสยฺยถาปิ คิหินิโย กามโภคินิโย”ติ.\csromanspacer{} อสฺโสสุํ โข ภิกฺขุนิโย ตาสํ เวสิยานํ อุชฺฌายนฺตีนํ ขิยฺยนฺตีนํ วิปาเจนฺตีนํ.\csromanspacer{} ยา ตา ภิกฺขุนิโย อปฺปิจฺฉา ฯเปฯ ตา อุชฺฌายนฺติ ขิยฺยนฺติ วิปาเจนฺติ “กถํ หิ นาม ฉพฺพคฺคิยา ภิกฺขุนิโย สมฺพาเธ โลมํ สํหราเปสฺสนฺตี”ติ ฯเปฯ.\csromanspacer{} สจฺจํ กิร ภิกฺขเว ฉพฺพคฺคิยา ภิกฺขุนิโย สมฺพาเธ โลมํ สํหราเปนฺตีติ.\csromanspacer{} สจฺจํ ภควาติ.\csromanspacer{} วิครหิ พุทฺโธ ภควา ฯเปฯ.\csromanspacer{} กถํ หิ นาม ภิกฺขเว ฉพฺพคฺคิยา ภิกฺขุนิโย สมฺพาเธ โลมํ สํหราเปสฺสนฺติ.\csromanspacer{} เนตํ ภิกฺขเว อปฺปสนฺนานํ วา ปสาทาย ฯเปฯ.\csromanspacer{} เอวญฺจ ปน ภิกฺขเว ภิกฺขุนิโย อิมํ สิกฺขาปทํ อุทฺทิสนฺตุ\csromandash{}}

\proseitem{799}{\textbf{“ยา ปน ภิกฺขุนี สมฺพาเธ โลมํ สํหราเปยฺย, ปาจิตฺติยนฺ”}ติ.}

\proseitem{800}{\textbf{ยา ปนา}ติ ยา ยาทิสา ฯเปฯ.\csromanspacer{} \textbf{ภิกฺขุนี}ติ ฯเปฯ อยํ อิมสฺมิํ อตฺเถ อธิปฺเปตา ภิกฺขุนีติ.}

\prose{\textbf{สมฺพาโธ} นาม อุโภ อุปกจฺฉกา มุตฺตกรณํ.}

\prose{\textbf{สํหราเปยฺยา}ติ เอกมฺปิ โลมํ สํหราเปติ, อาปตฺติ ปาจิตฺติยสฺส.\csromanspacer{} พหุเกปิ โลเม สํหราเปติ, อาปตฺติ ปาจิตฺติยสฺส.}

\proseitemwithcloserruled{801}{อนาปตฺติ อาพาธปจฺจยา อุมฺมตฺติกาย อาทิกมฺมิกายาติ.}{ทุติยสิกฺขาปทํ นิฏฺฐิตํ.}
\csromanmatikaanchor{mtk.152}
\csromantocmarkref{subsection}{3. ตติยสิกฺขาปท}{mtk.152}
\bookmark[dest={mtk.152},level=2]{3. ตติยสิกฺขาปท}
\csromanheader{2}{1. \textbf{ลสุณวคฺค 3. ตติยสิกฺขาปท}}
\proseitem{802}{เตน สมเยน พุทฺโธ ภควา สาวตฺถิยํ วิหรติ เชตวเน อนาถปิณฺฑิกสฺส อาราเม.\csromanspacer{} เตน โข ปน สมเยน เทฺว ภิกฺขุนิโย \csromanfolio{342}อนภิรติยา ปีฬิตา โอวรกํ ปวิสิตฺวา ตลฆาตกํ กโรนฺติ.\csromanspacer{} ภิกฺขุนิโย เตน สทฺเทน อุปธาวิตฺวา ตา ภิกฺขุนิโย เอตทโวจุํ “กิสฺส ตุเมฺห อยฺเย ปุริเสน สทฺธิํ สมฺปทุสฺสถา”ติ. “น มยํ อยฺเย ปุริเสน สทฺธิํ สมฺปทุสฺสามา”ติ ภิกฺขุนีนํ เอตมตฺถํ อาโรเจสุํ.\csromanspacer{} ยา ตา ภิกฺขุนิโย อปฺปิจฺฉา ฯเปฯ ตา อุชฺฌายนฺติ ขิยฺยนฺติ วิปาเจนฺติ “กถํ หิ นาม ภิกฺขุนิโย ตลฆาตกํ กริสฺสนฺตี”ติ ฯเปฯ.\csromanspacer{} สจฺจํ กิร ภิกฺขเว ภิกฺขุนิโย ตลฆาตกํ กโรนฺตีติ.\csromanspacer{} สจฺจํ ภควาติ.\csromanspacer{} วิครหิ พุทฺโธ ภควา ฯเปฯ.\csromanspacer{} กถํ หิ นาม ภิกฺขเว ภิกฺขุนิโย ตลฆาตกํ กริสฺสนฺติ.\csromanspacer{} เนตํ ภิกฺขเว อปฺปสนฺนานํ วา ปสาทาย ฯเปฯ.\csromanspacer{} เอวญฺจ ปน ภิกฺขเว ภิกฺขุนิโย อิมํ สิกฺขาปทํ อุทฺทิสนฺตุ\csromandash{}}

\proseitem{803}{\textbf{“ตลฆาตเก ปาจิตฺติยนฺ”}ติ.}

\proseitem{804}{\textbf{ตลฆาตกํ} นาม สมฺผสฺสํ สาทิยนฺตี อนฺตมโส อุปฺปลปตฺเตนปิ มุตฺตกรเณ ปหารํ เทติ, อาปตฺติ ปาจิตฺติยสฺส.}

\proseitemwithcloserruled{805}{อนาปตฺติ อาพาธปจฺจยา อุมฺมตฺติกาย อาทิกมฺมิกายาติ.}{ตติยสิกฺขาปทํ นิฏฺฐิตํ.}
\csromanmatikaanchor{mtk.153}
\csromantocmarkref{subsection}{4. จตุตฺถสิกฺขาปท}{mtk.153}
\bookmark[dest={mtk.153},level=2]{4. จตุตฺถสิกฺขาปท}
\csromanheader{2}{1. \textbf{ลสุณวคฺค 4. จตุตฺถสิกฺขาปท}}
\proseitem{806}{เตน สมเยน พุทฺโธ ภควา สาวตฺถิยํ วิหรติ เชตวเน อนาถปิณฺฑิกสฺส อาราเม.\csromanspacer{} เตน โข ปน สมเยน อญฺญตรา ปุราณราโชโรธา ภิกฺขุนีสุ ปพฺพชิตา\footnote{อญฺญตโร ปุราณราโชโรโธ ภิกฺขุนีสุ ปพฺพชิโต (สฺยา)} โหติ.\csromanspacer{} อญฺญตรา ภิกฺขุนี อนภิรติยา ปีฬิตา เยน สา ภิกฺขุนี เตนุปสงฺกมิ, อุปสงฺกมิตฺวา ตํ ภิกฺขุนิํ เอตทโวจ “ราชา โข อยฺเย ตุเมฺห จิราจิรํ คจฺฉติ, กถํ ตุเมฺห ธาเรถา”ติ.\csromanspacer{} ชตุมฏฺฐเกน อยฺเยติ.\csromanspacer{} กิํ เอตํ อยฺเย ชตุมฏฺฐกนฺติ.\csromanspacer{} อถ โข สา ภิกฺขุนี ตสฺสา ภิกฺขุนิยา ชตุมฏฺฐกํ อาจิกฺขิ.\csromanspacer{} อถ โข สา ภิกฺขุนี ชตุมฏฺฐกํ อาทิยิตฺวา โธวิตุํ วิสฺสริตฺวา เอกมนฺตํ ฉฑฺเฑสิ.\csromanspacer{} ภิกฺขุนิโย มกฺขิกาหิ สมฺปริกิณฺณํ \csromanfolio{343}ปสฺสิตฺวา เอวมาหํสุ “กสฺสิทํ กมฺมนฺ”ติ.\csromanspacer{} สา เอวมาห “มยฺหิทํ กมฺมนฺ”ติ.\csromanspacer{} ยา ตา ภิกฺขุนิโย อปฺปิจฺฉา ฯเปฯ ตา อุชฺฌายนฺติ ขิยฺยนฺติ วิปาเจนฺติ “กถํ หิ นาม ภิกฺขุนี ชตุมฏฺฐกํ อาทิยิสฺสตี”ติ ฯเปฯ.\csromanspacer{} สจฺจํ กิร ภิกฺขเว ภิกฺขุนี ชตุมฏฺฐกํ อาทิยตีติ\footnote{อาทิยีติ (ก)}.\csromanspacer{} สจฺจํ ภควาติ.\csromanspacer{} วิครหิ พุทฺโธ ภควา ฯเปฯ.\csromanspacer{} กถํ หิ นาม ภิกฺขเว ภิกฺขุนี ชตุมฏฺฐกํ อาทิยิสฺสติ.\csromanspacer{} เนตํ ภิกฺขเว อปฺปสนฺนานํ วา ปสาทาย ฯเปฯ.\csromanspacer{} เอวญฺจ ปน ภิกฺขเว ภิกฺขุนิโย อิมํ สิกฺขาปทํ อุทฺทิสนฺตุ\csromandash{}}

\proseitem{807}{\textbf{“ชตุมฏฺฐเก}\footnote{ชตุมฏฺฏเก (สี)} \textbf{ปาจิตฺติยนฺ”}ติ.}

\proseitem{808}{ชตุมฏฺฐกํ นาม ชตุมยํ กฏฺฐมยํ ปิฏฺฐมยํ มตฺติกามยํ.}

\prose{\textbf{อาทิเยยฺยา}ติ สมฺผสฺสํ สาทิยนฺตี อนฺตมโส อุปฺปลปตฺตมฺปิ มุตฺตกรณํ ปเวเสติ, อาปตฺติ ปาจิตฺติยสฺส.}

\proseitemwithcloserruled{809}{อนาปตฺติ อาพาธปจฺจยา อุมฺมตฺติกาย อาทิกมฺมิกายาติ.}{จตุตฺถสิกฺขาปทํ นิฏฺฐิตํ.}
\csromanmatikaanchor{mtk.154}
\csromantocmarkref{subsection}{5. ปญฺจมสิกฺขาปท}{mtk.154}
\bookmark[dest={mtk.154},level=2]{5. ปญฺจมสิกฺขาปท}
\csromanheader{2}{1. \textbf{ลสุณวคฺค 5. ปญฺจมสิกฺขาปท}}
\proseitem{810}{เตน สมเยน พุทฺโธ ภควา สกฺเกสุ วิหรติ กปิลวตฺถุสฺมิํ นิโคฺรธาราเม.\csromanspacer{} อถ โข มหาปชาปติ โคตมี เยน ภควา เตนุปสงฺกมิ, อุปสงฺกมิตฺวา ภควนฺตํ อภิวาเทตฺวา อโธวาเต อฏฺฐาสิ “ทุคฺคนฺโธ ภควา มาตุคาโม”ติ.\csromanspacer{} อถ โข ภควา “อาทิยนฺตุ โข ภิกฺขุนิโย อุทกสุทฺธิกนฺ”ติ มหาปชาปติํ โคตมิํ ธมฺมิยา กถาย สนฺทสฺเสสิ สมาทเปสิ สมุตฺเตเชสิ สมฺปหํเสสิ.\csromanspacer{} อถ โข มหาปชาปติ โคตมี ภควตา ธมฺมิยา กถาย สนฺทสฺสิตา สมาทปิตา สมุตฺเตชิตา สมฺปหํสิตา ภควนฺตํ อภิวาเทตฺวา ปทกฺขิณํ กตฺวา ปกฺกามิ.\csromanspacer{} อถ โข ภควา เอตสฺมิํ นิทาเน เอตสฺมิํ ปกรเณ ธมฺมิํ กถํ กตฺวา ภิกฺขู อามนฺเตสิ “อนุชานามิ ภิกฺขเว ภิกฺขุนีนํ อุทกสุทฺธิกนฺ”ติ.\csromanspacer{} เตน โข ปน สมเยน อญฺญตรา ภิกฺขุนี \csromanfolio{344}“\textbf{ภควตา อุทกสุทฺธิกา อนุญฺญาตา”ติ อติคมฺภีรํ อุทกสุทฺธิกํ} \textbf{อาทิยนฺตี มุตฺตกรเณ วณํ อกาสิ.}\csromanspacer{}\textbf{ อถ โข สา ภิกฺขุนี ภิกฺขุนีนํ} เอตมตฺถํ อาโรเจติ.\csromanspacer{} ยา ตา ภิกฺขุนิโย อปฺปิจฺฉา ฯเปฯ ตา อุชฺฌายนฺติ ขิยฺยนฺติ วิปาเจนฺติ “กถํ หิ นาม ภิกฺขุนี อติคมฺภีรํ อุทกสุทฺธิกํ อาทิยิสฺสตี”ติ ฯเปฯ.\csromanspacer{} สจฺจํ กิร ภิกฺขเว ภิกฺขุนี อติคมฺภีรํ อุทกสุทฺธิกํ อาทิยตีติ\footnote{อาทิยีติ (ก)}.\csromanspacer{} สจฺจํ ภควาติ.\csromanspacer{} วิครหิ พุทฺโธ ภควา ฯเปฯ.\csromanspacer{} กถํ หิ นาม ภิกฺขเว ภิกฺขุนี อติคมฺภีรํ อุทกสุทฺธิกํ อาทิยิสฺสติ.\csromanspacer{} เนตํ ภิกฺขเว อปฺปสนฺนานํ วา ปสาทาย ฯเปฯ.\csromanspacer{} เอวญฺจ ปน ภิกฺขเว ภิกฺขุนิโย อิมํ สิกฺขาปทํ อุทฺทิสนฺตุ\csromandash{}}

\proseitem{811}{\textbf{“อุทกสุทฺธิกํ ปน ภิกฺขุนิยา อาทิยมานาย ทฺวงฺคุลปพฺพปรมํ อาทาตพฺพํ, ตํ อติกฺกาเมนฺติยา ปาจิตฺติยนฺ”}ติ.}

\proseitem{812}{\textbf{อุทกสุทฺธิกํ} นาม มุตฺตกรณสฺส โธวนา วุจฺจติ.}

\prose{\textbf{อาทิยมานายา}ติ โธวนฺติยา.}

\prose{\textbf{ทฺวงฺคุลปพฺพปรมํ อาทาตพฺพนฺ}ติ ทฺวีสุ องฺคุเลสุ เทฺว ปพฺพปรมา อาทาตพฺพา.}

\prose{\textbf{ตํ อติกฺกาเมนฺติยา}ติ สมฺผสฺสํ สาทิยนฺตี อนฺตมโส เกสคฺคมตฺตมฺปิ อติกฺกาเมติ, อาปตฺติ ปาจิตฺติยสฺส.}

\proseitem{813}{อติเรกทฺวงฺคุลปพฺเพ อติเรกสญฺญา อาทิยติ, อาปตฺติ ปาจิตฺติยสฺส.\csromanspacer{} อติเรกทฺวงฺคุลปพฺเพ เวมติกา อาทิยติ, อาปตฺติ ปาจิตฺติยสฺส.\csromanspacer{} อติเรกทฺวงฺคุลปพฺเพ อูนกสญฺญา อาทิยติ, อาปตฺติ ปาจิตฺติยสฺส.}

\prose{อูนกทฺวงฺคุลปพฺเพ อติเรกสญฺญา, อาปตฺติ ทุกฺกฏสฺส.\csromanspacer{} อูนกทฺวงฺคุลปพฺเพ เวมติกา, อาปตฺติ ทุกฺกฏสฺส, อูนกทฺวงฺคุลปพฺเพ อูนกสญฺญา, อนาปตฺติ.}

\proseitemwithcloser{814}{อนาปตฺติ ทฺวงฺคุลปพฺพปรมํ อาทิยติ, อูนกทฺวงฺคุลปพฺพปรมํ อาทิยติ, อาพาธปจฺจยา อุมฺมตฺติกาย อาทิกมฺมิกายาติ.}{ปญฺจมสิกฺขาปทํ นิฏฺฐิตํ.}
\csromanmatikaanchor{mtk.155}
\csromantocmarkref{subsection}{6. ฉฏฺฐสิกฺขาปท}{mtk.155}
\bookmark[dest={mtk.155},level=2]{6. ฉฏฺฐสิกฺขาปท}
\csromanheader{2}{1. \textbf{ลสุณวคฺค 6. ฉฏฺฐสิกฺขาปท}}
\proseitem{815}{\csromanfolio{345}เตน สมเยน พุทฺโธ ภควา สาวตฺถิยํ วิหรติ เชตวเน อนาถปิณฺฑิกสฺส อาราเม.\csromanspacer{} เตน โข ปน สมเยน อาโรหนฺโต นาม มหามตฺโต ภิกฺขูสุ ปพฺพชิโต โหติ, ตสฺส ปุราณทุติยิกา ภิกฺขุนีสุ ปพฺพชิตา โหติ.\csromanspacer{} เตน โข ปน สมเยน โส ภิกฺขุ ตสฺสา ภิกฺขุนิยา สนฺติเก ภตฺตวิสฺสคฺคํ กโรติ.\csromanspacer{} อถ โข สา ภิกฺขุนี ตสฺส ภิกฺขุโน ภุญฺชนฺตสฺส ปานีเยน จ วิธูปเนน จ อุปติฏฺฐิตฺวา อจฺจาวทติ.\csromanspacer{} อถ โข โส ภิกฺขุ ตํ ภิกฺขุนิํ อปสาเทติ “มา ภคินิ เอวรูปํ อกาสิ, เนตํ กปฺปตี”ติ. “ปุพฺเพ มํ ตฺวํ เอวญฺจ เอวญฺจ กโรสิ, อิทานิ เอตฺตกํ น สหสี”ติ ปานียถาลกํ มตฺถเก อาสุมฺภิตฺวา วิธูปเนน ปหารํ อทาสิ.\csromanspacer{} ยา ตา ภิกฺขุนิโย อปฺปิจฺฉา ฯเปฯ ตา อุชฺฌายนฺติ ขิยฺยนฺติ วิปาเจนฺติ “กถํ หิ นาม ภิกฺขุนี ภิกฺขุสฺส ปหารํ ทสฺสตี”ติ ฯเปฯ.\csromanspacer{} สจฺจํ กิร ภิกฺขเว ภิกฺขุนี ภิกฺขุสฺส ปหารํ อทาสีติ.\csromanspacer{} สจฺจํ ภควาติ.\csromanspacer{} วิครหิ พุทฺโธ ภควา ฯเปฯ.\csromanspacer{} กถํ หิ นาม ภิกฺขเว ภิกฺขุนี ภิกฺขุสฺส ปหารํ ทสฺสติ.\csromanspacer{} เนตํ ภิกฺขเว อปฺปสนฺนานํ วา ปสาทาย ฯเปฯ.\csromanspacer{} เอวญฺจ ปน ภิกฺขเว ภิกฺขุนิโย อิมํ สิกฺขาปทํ อุทฺทิสนฺตุ\csromandash{}}

\proseitem{816}{\textbf{“ยา ปน ภิกฺขุนี ภิกฺขุสฺส ภุญฺชนฺตสฺส ปานีเยน วา วิธูปเนน วา อุปติฏฺเฐยฺย, ปาจิตฺติยนฺ”}ติ.}

\proseitem{817}{\textbf{ยา ปนา}ติ ยา ยาทิสา ฯเปฯ.\csromanspacer{} \textbf{ภิกฺขุนี}ติ ฯเปฯ อยํ อิมสฺมิํ อตฺเถ อธิปฺเปตา ภิกฺขุนีติ.}

\prose{\textbf{ภิกฺขุสฺสา}ติ อุปสมฺปนฺนสฺส.}

\prose{\textbf{ภุญฺชนฺตสฺสา}ติ ปญฺจนฺนํ โภชนานํ อญฺญตรํ โภชนํ ภุญฺชนฺตสฺส.}

\prose{\textbf{ปานียํ} นาม ยํ กิญฺจิ ปานียํ.}

\prose{\textbf{วิธูปนํ} นาม ยา กาจิ พีชนี.}

\prose{\textbf{อุปติฏฺเฐยฺยา}ติ หตฺถปาเส ติฏฺฐติ, อาปตฺติ ปาจิตฺติยสฺส.}

\proseitem{818}{อุปสมฺปนฺเน อุปสมฺปนฺนสญฺญา ปานีเยน วา วิธูปเนน วา อุปติฏฺฐติ, อาปตฺติ ปาจิตฺติยสฺส.\csromanspacer{} อุปสมฺปนฺเน เวมติกา ปานีเยน วา วิธูปเนน วา อุปติฏฺฐติ, อาปตฺติ ปาจิตฺติยสฺส.\csromanspacer{} อุปสมฺปนฺเน อนุปสมฺปนฺนสญฺญา ปานีเยน วา วิธูปเนน วา อุปติฏฺฐติ, อาปตฺติ}

\prose{\csromanfolio{346}หตฺถปาสํ วิชหิตฺวา อุปติฏฺฐติ, อาปตฺติ ทุกฺกฏสฺส.\csromanspacer{} ขาทนียํ ขาทนฺตสฺส อุปติฏฺฐติ อาปตฺติ ทุกฺกฏสฺส.\csromanspacer{} อนุปสมฺปนฺนสฺส อุปติฏฺฐติ, อาปตฺติ ทุกฺกฏสฺส.\csromanspacer{} อนุปสมฺปนฺเน อุปสมฺปนฺนสญฺญา, อาปตฺติ ทุกฺกฏสฺส.\csromanspacer{} อนุปสมฺปนฺเน เวมติกา, อาปตฺติ ทุกฺกฏสฺส.\csromanspacer{} อนุปสมฺปนฺเน อนุปสมฺปนฺนสญฺญา, อาปตฺติ ทุกฺกฏสฺส.}

\proseitemwithcloserruled{819}{อนาปตฺติ เทติ, ทาเปติ, อนุปสมฺปนฺนํ อาณาเปติ, อุมฺมตฺติกาย อาทิกมฺมิกายาติ.}{ฉฏฺฐสิกฺขาปทํ นิฏฺฐิตํ.}
\csromanmatikaanchor{mtk.156}
\csromantocmarkref{subsection}{7. สตฺตมสิกฺขาปท}{mtk.156}
\bookmark[dest={mtk.156},level=2]{7. สตฺตมสิกฺขาปท}
\csromanheader{2}{1. \textbf{ลสุณวคฺค 7. สตฺตมสิกฺขาปท}}
\proseitem{820}{เตน สมเยน พุทฺโธ ภควา สาวตฺถิยํ วิหรติ เชตวเน อนาถปิณฺฑิกสฺส อาราเม.\csromanspacer{} เตน โข ปน สมเยน ภิกฺขุนิโย สสฺสกาเล อามกธญฺญํ วิญฺญาเปตฺวา นครํ อติหรนฺติ, ทฺวารฏฺฐาเน “เทถายฺเย ภาคนฺ”ติ ปลิพุนฺเธตฺวา มุญฺจิํสุ.\csromanspacer{} อถ โข ตา ภิกฺขุนิโย อุปสฺสยํ คนฺตฺวา ภิกฺขุนีนํ เอตมตฺถํ อาโรเจสุํ.\csromanspacer{} ยา ตา ภิกฺขุนิโย อปฺปิจฺฉา ฯเปฯ ตา อุชฺฌายนฺติ ขิยฺยนฺติ วิปาเจนฺติ “กถํ หิ นาม ภิกฺขุนิโย อามกธญฺญํ วิญฺญาเปสฺสนฺตี”ติ ฯเปฯ.\csromanspacer{} สจฺจํ กิร ภิกฺขเว ภิกฺขุนิโย อามกธญฺญํ วิญฺญาเปนฺตีติ.\csromanspacer{} สจฺจํ ภควาติ.\csromanspacer{} วิครหิ พุทฺโธ ภควา ฯเปฯ.\csromanspacer{} กถํ หิ นาม ภิกฺขเว ภิกฺขุนิโย อามกธญฺญํ วิญฺญาเปสฺสนฺติ.\csromanspacer{} เนตํ ภิกฺขเว อปฺปสนฺนานํ วา ปสาทาย ฯเปฯ.\csromanspacer{} เอวญฺจ ปน ภิกฺขเว ภิกฺขุนิโย อิมํ สิกฺขาปทํ อุทฺทิสนฺตุ\csromandash{}}

\proseitem{821}{\textbf{“ยา ปน ภิกฺขุนี อามกธญฺญํ วิญฺญตฺวา วา วิญฺญาเปตฺวา วา ภชฺชิตฺวา วา ภชฺชาเปตฺวา วา โกฏฺเฏตฺวา วา โกฏฺฏาเปตฺวา วา ปจิตฺวา วา ปจาเปตฺวา วา ภุญฺเชยฺย, ปาจิตฺติยนฺ”}ติ.}

\proseitem{822}{\textbf{ยา ปนา}ติ ยา ยาทิสา ฯเปฯ.\csromanspacer{} \textbf{ภิกฺขุนี}ติ ฯเปฯ อยํ อิมสฺมิํ อตฺเถ อธิปฺเปตา ภิกฺขุนีติ.}

\prose{\textbf{อามกธญฺญํ} นาม สาลิ วีหิ ยโว โคธุโม กงฺคุ วรโก กุทฺรุสโก.}

\prose{\csromanfolio{347}\textbf{วิญฺญตฺวา}ติ สยํ วิญฺญตฺวา.\csromanspacer{} \textbf{วิญฺญาเปตฺวา}ติ อญฺญํ วิญฺญาเปตฺวา.}

\prose{\textbf{ภชฺชิตฺวา}ติ สยํ ภชฺชิตฺวา.\csromanspacer{} \textbf{ภชฺชาเปตฺวา}ติ อญฺญํ ภชฺชาเปตฺวา.}

\prose{\textbf{โกฏฺเฏตฺวา}ติ สยํ โกฏฺเฏตฺวา.\csromanspacer{} \textbf{โกฏฺฏาเปตฺวา}ติ อญฺญํ โกฏฺฏาเปตฺวา.}

\prose{\textbf{ปจิตฺวา}ติ สยํ ปจิตฺวา.\csromanspacer{} \textbf{ปจาเปตฺวา}ติ อญฺญํ ปจาเปตฺวา.}

\prose{“ภุญฺชิสฺสามี”ติ ปฏิคฺคณฺหาติ, อาปตฺติ ทุกฺกฏสฺส.\csromanspacer{} อชฺโฌหาเร อชฺโฌหาเร อาปตฺติ ปาจิตฺติยสฺส.}

\proseitemwithcloserruled{823}{อนาปตฺติ อาพาธปจฺจยา, อปรณฺณํ วิญฺญาเปติ, อุมฺมตฺติกาย อาทิกมฺมิกายาติ.}{สตฺตมสิกฺขาปทํ นิฏฺฐิตํ.}
\csromanmatikaanchor{mtk.157}
\csromantocmarkref{subsection}{8. อฏฺฐมสิกฺขาปท}{mtk.157}
\bookmark[dest={mtk.157},level=2]{8. อฏฺฐมสิกฺขาปท}
\csromanheader{2}{1. \textbf{ลสุณวคฺค 8. อฏฺฐมสิกฺขาปท}}
\proseitem{824}{เตน สมเยน พุทฺโธ ภควา สาวตฺถิยํ วิหรติ เชตวเน อนาถปิณฺฑิกสฺส อาราเม.\csromanspacer{} เตน โข ปน สมเยน อญฺญตโร พฺราหฺมโณ นิพฺพิฏฺฐราชภโฏ “ตญฺเญว ภฏปถํ ยาจิสฺสามี”ติ สีสํ นหายิตฺวา ภิกฺขุนุปสฺสยํ นิสฺสาย ราชกุลํ คจฺฉติ.\csromanspacer{} อญฺญตรา ภิกฺขุนี กฏาเห วจฺจํ กตฺวา ติโรกุฏฺเฏ ฉฑฺเฑนฺตี ตสฺส พฺราหฺมณสฺส มตฺถเก อาสุมฺภิ.\csromanspacer{} อถ โข โส พฺราหฺมโณ อุชฺฌายติ ขิยฺยนฺติ วิปาเจติ “อสฺสมณิโย อิมา มุณฺฑา พนฺธกินิโย, กถํ หิ นาม คูถกฏาหํ มตฺถเก อาสุมฺภิสฺสนฺติ, อิมาสํ อุปสฺสยํ ฌาเปสฺสามี”ติ อุมฺมุกํ คเหตฺวา อุปสฺสยํ ปวิสติ.\csromanspacer{} อญฺญตโร อุปาสโก อุปสฺสยา นิกฺขมนฺโต อทฺทส ตํ พฺราหฺมณํ อุมฺมุกํ คเหตฺวา อุปสฺสยํ ปวิสนฺตํ, ทิสฺวาน ตํ พฺราหฺมณํ เอตทโวจ “กิสฺส ตฺวํ โภ อุมฺมุกํ คเหตฺวา อุปสฺสยํ ปวิสสี”ติ.\csromanspacer{} อิมา มํ โภ มุณฺฑา พนฺธกินิโย คูถกฏาหํ มตฺถเก อาสุมฺภิํสุ, อิมาสํ อุปสฺสยํ ฌาเปสฺสามีติ.\csromanspacer{} คจฺฉ โภ พฺราหฺมณ, มงฺคลํ เอตํ, สหสฺสํ ลจฺฉสิ ตญฺจ ภฏปถนฺติ.\csromanspacer{} อถ โข โส พฺราหฺมโณ สีสํ นหายิตฺวา ราชกุลํ คนฺตฺวา สหสฺสํ อลตฺถ ตญฺจ ภฏปถํ.\csromanspacer{} อถ โข โส อุปาสโก อุปสฺสยํ ปวิสิตฺวา ภิกฺขุนีนํ \csromanfolio{348}\textbf{เอตมตฺถํ อาโรเจตฺวา ปริภาสิ.}\csromanspacer{}\textbf{ ยา ตา ภิกฺขุนิโย อปฺปิจฺฉา ฯเปฯ ตา} อุชฺฌายนฺติ ขิยฺยนฺติ วิปาเจนฺติ “กถํ หิ นาม ภิกฺขุนิโย อุจฺจารํ ติโรกุฏฺเฏ ฉฑฺเฑสฺสนฺตี”ติ ฯเปฯ.\csromanspacer{} สจฺจํ กิร ภิกฺขเว ภิกฺขุนิโย อุจฺจารํ ติโรกุฏฺเฏ ฉฑฺเฑนฺตีติ.\csromanspacer{} สจฺจํ ภควาติ.\csromanspacer{} วิครหิ พุทฺโธ ภควาติ ฯเปฯ.\csromanspacer{} กถํ หิ นาม ภิกฺขเว ภิกฺขุนิโย อุจฺจารํ ติโรกุฏฺเฏ ฉฑฺเฑสฺสนฺติ.\csromanspacer{} เนตํ ภิกฺขเว อปฺปสนฺนานํ วา ปสาทาย ฯเปฯ.\csromanspacer{} เอวญฺจ ปน ภิกฺขเว ภิกฺขุนิโย อิมํ สิกฺขาปทํ อุทฺทิสนฺตุ\csromandash{}}

\proseitem{825}{\textbf{“ยา ปน ภิกฺขุนี อุจฺจารํ วา ปสฺสาวํ วา สงฺการํ วา วิฆาสํ วา ติโรกุฏฺเฏ วา ติโรปากาเร วา ฉฑฺเฑยฺย วา ฉฑฺฑาเปยฺย วา, ปาจิตฺติยนฺ”}ติ.}

\proseitem{826}{\textbf{ยา ปนา}ติ ยา ยาทิสา ฯเปฯ.\csromanspacer{} \textbf{ภิกฺขุนี}ติ ฯเปฯ อยํ อิมสฺมิํ อตฺเถ อธิปฺเปตา ภิกฺขุนีติ.}

\prose{\textbf{อุจฺจาโร} นาม คูโถ วุจฺจติ.\csromanspacer{} \textbf{ปสฺสาโว} นาม มุตฺตํ วุจฺจติ.}

\prose{\textbf{สงฺการํ} นาม กจวรํ วุจฺจติ.}

\prose{\textbf{วิฆาสํ} นาม จลกานิ วา อฏฺฐิกานิ วา อุจฺฉิฏฺโฐทกํ วา.}

\prose{\textbf{กุฏฺโฏ} นาม ตโย กุฏฺฏา อิฏฺฐกากุฏฺโฏ สิลากุฏฺโฏ ทารุกุฏฺโฏ.}

\prose{\textbf{ปากาโร} นาม ตโย ปาการา อิฏฺฐกาปากาโร สิลาปากาโร ทารุปากาโร.}

\prose{\textbf{ติโรกุฏฺเฏ}ติ กุฏฺฏสฺส ปรโต.\csromanspacer{} \textbf{ติโรปากาเร}ติ ปาการสฺส ปรโต.}

\prose{\textbf{ฉฑฺเฑยฺยา}ติ สยํ ฉฑฺเฑติ, อาปตฺติ ปาจิตฺติยสฺส.\csromanspacer{} \textbf{ฉฑฺฑาเปยฺยา}ติ อญฺญํ อาณาเปติ, อาปตฺติ ทุกฺกฏสฺส.\csromanspacer{} สกิํ อาณตฺตา พหุกมฺปิ ฉฑฺเฑติ, อาปตฺติ}

\proseitemwithcloser{827}{อนาปตฺติ โอโลเกตฺวา ฉฑฺเฑติ, อวฬญฺเช ฉฑฺเฑติ, อุมฺมตฺติกาย อาทิกมฺมิกายาติ.}{อฏฺฐมสิกฺขาปทํ นิฏฺฐิตํ.}
\csromanmatikaanchor{mtk.158}
\csromantocmarkref{subsection}{9. นวมสิกฺขาปท}{mtk.158}
\bookmark[dest={mtk.158},level=2]{9. นวมสิกฺขาปท}
\csromanheader{2}{1. \textbf{ลสุณวคฺค 9. นวมสิกฺขาปท}}
\proseitem{828}{\csromanfolio{349}เตน สมเยน พุทฺโธ ภควา สาวตฺถิยํ วิหรติ เชตวเน อนาถปิณฺฑิกสฺส อาราเม.\csromanspacer{} เตน โข ปน สมเยน อญฺญตรสฺส พฺราหฺมณสฺส ภิกฺขุนุปสฺสยํ นิสฺสาย ยวเขตฺตํ โหติ.\csromanspacer{} ภิกฺขุนิโย อุจฺจารมฺปิ ปสฺสาวมฺปิ สงฺการมฺปิ วิฆาสมฺปิ เขตฺเต ฉฑฺเฑนฺติ.\csromanspacer{} อถ โข โส พฺราหฺมโณ อุชฺฌายติ ขิยฺยนฺติ วิปาเจติ “กถํ หิ นาม ภิกฺขุนิโย อมฺหากํ ยวเขตฺตํ ทูเสสฺสนฺตี”ติ.\csromanspacer{} อสฺโสสุํ โข ภิกฺขุนิโย ตสฺส พฺรหฺมณสฺส อุชฺฌายนฺตสฺส ขิยฺยนฺตสฺส วิปาเจนฺตสฺส.\csromanspacer{} ยา ตา ภิกฺขุนิโย อปฺปิจฺฉา ฯเปฯ ตา อุชฺฌายนฺติ ขิยฺยนฺติ วิปาเจนฺติ “กถํ หิ นาม ภิกฺขุนิโย อุจฺจารมฺปิ ปสฺสาวมฺปิ สงฺการมฺปิ วิฆาสมฺปิ หริเต ฉฑฺเฑสฺสนฺตี”ติ ฯเปฯ.\csromanspacer{} สจฺจํ กิร ภิกฺขเว ภิกฺขุนิโย อุจฺจารมฺปิ ปสฺสาวมฺปิ สงฺการมฺปิ วิฆาสมฺปิ หริเต ฉฑฺเฑนฺตีติ.\csromanspacer{} สจฺจํ ภควาติ.\csromanspacer{} วิครหิ พุทฺโธ ภควา ฯเปฯ.\csromanspacer{} กถํ หิ นาม ภิกฺขเว ภิกฺขุนิโย อุจฺจารมฺปิ ปสฺสาวมฺปิ สงฺการมฺปิ วิฆาสมฺปิ หริเต ฉฑฺเฑสฺสนฺติ.\csromanspacer{} เนตํ ภิกฺขเว อปฺปสนฺนานํ วา ปสาทาย ฯเปฯ.\csromanspacer{} เอวญฺจ ปน ภิกฺขเว ภิกฺขุนิโย อิมํ สิกฺขาปทํ อุทฺทิสนฺตุ\csromandash{}}

\proseitem{829}{\textbf{“ยา ปน ภิกฺขุนี อุจฺจารํ วา ปสฺสาวํ วา สงฺการํ วา วิฆาสํ วา หริเต ฉฑฺเฑยฺย วา ฉฑฺฑาเปยฺย วา, ปาจิตฺติยนฺ”}ติ.}

\proseitem{830}{\textbf{ยา ปนา}ติ ยา ยาทิสา ฯเปฯ.\csromanspacer{} \textbf{ภิกฺขุนี}ติ ฯเปฯ อยํ อิมสฺมิํ อตฺเถ อธิปฺเปตา ภิกฺขุนีติ.}

\prose{\textbf{อุจฺจาโร} นาม คูโถ วุจฺจติ.\csromanspacer{} \textbf{ปสฺสาโว} นาม มุตฺตํ วุจฺจติ.}

\prose{\textbf{สงฺการํ} นาม กจวรํ วุจฺจติ.}

\prose{\textbf{วิฆาสํ} นาม จลกานิ วา อฏฺฐิกานิ วา อุจฺฉิฏฺโฐทกํ วา.}

\prose{\textbf{หริตํ} นาม ปุพฺพณฺณํ อปรณฺณํ, ยํ มนุสฺสานํ อุปโภคปริโภคํ โรปิมํ.}

\prose{\textbf{ฉฑฺเฑยฺยา}ติ สยํ ฉฑฺเฑติ, อาปตฺติ ปาจิตฺติยสฺส.}

\prose{\textbf{ฉฑฺฑาเปยฺยา}ติ อญฺญํ อาณาเปติ, อาปตฺติ ทุกฺกฏสฺส.\csromanspacer{} สกิํ อาณตฺตา พหุกมฺปิ ฉฑฺเฑติ, อาปตฺติ ปาจิตฺติยสฺส.}

\proseitem{831}{หริเต หริตสญฺญา ฉฑฺเฑติ วา ฉฑฺฑาเปติ วา อาปตฺติ ปาจิตฺติยสฺส.\csromanspacer{} หริเต เวมติกา ฉฑฺเฑติ จา ฉฑฺฑาเปติ วา, อาปตฺติ \csromanfolio{350}ปาจิตฺติยสฺส.\csromanspacer{} หริเต อหริตสญฺญา ฉฑฺเฑติ วา ฉฑฺฑาเปติ วา อาปตฺติ}

\prose{อหริเต หริตสญฺญา, อาปตฺติ ทุกฺกฏสฺส.\csromanspacer{} อหริเต เวมติกา, อาปตฺติ ทุกฺกฏสฺส.\csromanspacer{} อหริเต อหริตสญฺญา, อนาปตฺติ.}

\proseitemwithcloserruled{832}{อนาปตฺติ โอโลเกตฺวา ฉฑฺเฑติ, เขตฺตมริยาเท\footnote{ฉฑฺฑิตเขตฺเต (?) อฏฺฐกถา โอโลเกตพฺพา.} ฉฑฺเฑติ, สามิเก อาปุจฺฉิตฺวา อปโลเกตฺวา ฉฑฺเฑติ, อุมฺมตฺติกาย อาทิกมฺมิกายาติ.}{นวมสิกฺขาปทํ นิฏฺฐิตํ.}
\csromanmatikaanchor{mtk.159}
\csromantocmarkref{subsection}{10. ทสมสิกฺขาปท}{mtk.159}
\bookmark[dest={mtk.159},level=2]{10. ทสมสิกฺขาปท}
\csromanheader{2}{1. \textbf{ลสุณวคฺค} 10. ทสมสิกฺขาปท}
\proseitem{833}{เตน สมเยน พุทฺโธ ภควา ราชคเห วิหรติ เวฬุวเน กลนฺทกนิวาเป.\csromanspacer{} เตน โข ปน สมเยน ราชคเห คิรคฺคสมชฺโช โหติ.\csromanspacer{} ฉพฺพคฺคิยา ภิกฺขุนิโย คิรคฺคสมชฺชํ ทสฺสนาย อคมํสุ.\csromanspacer{} มนุสฺสา อุชฺฌายนฺติ ขิยฺยนฺติ วิปาเจนฺติ “กถํ หิ นาม ภิกฺขุนิโย นจฺจมฺปิ คีตมฺปิ วาทิตมฺปิ ทสฺสนาย คจฺฉิสฺสนฺติ เสยฺยถาปิ คิหินิโย กามโภคินิโย”ติ.\csromanspacer{} อสฺโสสุํ โข ภิกฺขุนิโย เตสํ มนุสฺสานํ อุชฺฌายนฺตานํ ขิยฺยนฺตานํ วิปาเจนฺตานํ.\csromanspacer{} ยา ตา ภิกฺขุนิโย อปฺปิจฺฉา ฯเปฯ ตา อุชฺฌายนฺติ ขิยฺยนฺติ วิปาเจนฺติ “กถํ หิ นาม ฉพฺพคฺคิยา ภิกฺขุนิโย นจฺจมฺปิ คีตมฺปิ วาทิตมฺปิ ทสฺสนาย คจฺฉิสฺสนฺตี”ติ ฯเปฯ.\csromanspacer{} สจฺจํ กิร ภิกฺขเว ฉพฺพคฺคิยา ภิกฺขุนิโย นจฺจมฺปิ คีตมฺปิ วาทิตมฺปิ ทสฺสนาย คจฺฉนฺตีติ.\csromanspacer{} สจฺจํ ภควาติ.\csromanspacer{} วิครหิ พุทฺโธ ภควา ฯเปฯ.\csromanspacer{} กถํ หิ นาม ภิกฺขเว ฉพฺพคฺคิยา ภิกฺขุนิโย นจฺจมฺปิ คีตมฺปิ วาทิตมฺปิ ทสฺสนาย คจฺฉิสฺสนฺติ.\csromanspacer{} เนตํ ภิกฺขเว อปฺปสนฺนานํ วา ปสาทาย ฯเปฯ.\csromanspacer{} เอวญฺจ ปน ภิกฺขเว ภิกฺขุนิโย อิมํ สิกฺขาปทํ อุทฺทิสนฺตุ\csromandash{}}

\proseitem{834}{\textbf{“ยา ปน ภิกฺขุนี นจฺจํ วา คีตํ วา วาทิตํ วา ทสฺสนาย คจฺเฉยฺย, ปาจิตฺติยนฺ”}ติ.}

\proseitem{835}{\textbf{ยา ปนา}ติ ยา ยาทิสา ฯเปฯ.\csromanspacer{} \textbf{ภิกฺขุนี}ติ ฯเปฯ อยํ อิมสฺมิํ อตฺเถ อธิปฺเปตา ภิกฺขุนีติ.}

\prose{\csromanfolio{351}\textbf{นจฺจํ} นาม ยํ กิญฺจิ นจฺจํ.\csromanspacer{} \textbf{คีตํ} นาม ยํ กิญฺจิ คีตํ.\csromanspacer{} \textbf{วาทิตํ} นาม ยํ กิญฺจิ วาทิตํ.}

\proseitem{836}{ทสฺสนาย คจฺฉติ, อาปตฺติ ทุกฺกฏสฺส.\csromanspacer{} ยตฺถ ฐิตา ปสฺสติ วา สุณาติ วา, อาปตฺติ ปาจิตฺติยสฺส.\csromanspacer{} ทสฺสนูปจารํ วิชหิตฺวา ปุนปฺปุนํ ปสฺสติ วา สุณาติ วา, อาปตฺติ ปาจิตฺติยสฺส.\csromanspacer{} เอกเมกํ ทสฺสนาย คจฺฉติ, อาปตฺติ ทุกฺกฏสฺส.\csromanspacer{} ยตฺถ ฐิตา ปสฺสติ วา สุณาติ วา, อาปตฺติ ปาจิตฺติยสฺส.\csromanspacer{} ทสฺสนูปจารํ วิชหิตฺวา ปุนปฺปุนํ ปสฺสติ วา สุณาติ วา, อาปตฺติ}

\proseitemwithcloser{837}{อนาปตฺติ อาราเม ฐิตา ปสฺสติ วา สุณาติ วา, ภิกฺขุนิยา ฐิโตกาสํ วา นิสินฺโนกาสํ วา นิปนฺโนกาสํ วา อาคนฺตฺวา นจฺจนฺติ วา คายนฺติ วา วาเทนฺติ วา, ปฏิปถํ คจฺฉนฺตี ปสฺสติ วา สุณาติ วา, สติ กรณีเย คนฺตฺวา ปสฺสติ วา สุณาติ วา, อาปทาสุ อุมฺมตฺติกาย อาทิกมฺมิกายาติ.}{ทสมสิกฺขาปทํ นิฏฺฐิตํ.}
\csromanheader{2}{ลสุณวคฺโค ปฐโม.}
\csromansectionrule
\csromanmatikaanchor{mtk.160}
\csromanmatikaanchor{mtk.161}
\csromantocmarknopage{section}{2. อนฺธการวคฺค}{mtk.160}
\bookmark[dest={mtk.160},level=1]{2. อนฺธการวคฺค}
\csromantocmarkref{subsection}{1. ปฐมสิกฺขาปท}{mtk.161}
\bookmark[dest={mtk.161},level=2]{1. ปฐมสิกฺขาปท}
\setcsromanheadh{2. อนฺธการวคฺค}
\csromanheader{2}{2. \textbf{อนฺธการวคฺค 1. ปฐมสิกฺขาปท}}
\proseitem{838}{เตน สมเยน พุทฺโธ ภควา สาวตฺถิยํ วิหรติ เชตวเน อนาถปิณฺฑิกสฺส อาราเม.\csromanspacer{} เตน โข ปน สมเยน ภทฺทาย กาปิลานิยา อนฺเตวาสินิยา ภิกฺขุนิยา ญาตโก ปุริโส คามกา สาวตฺถิํ อคมาสิ เกนจิเทว กรณีเยน.\csromanspacer{} อถ โข สา ภิกฺขุนี เตน ปุริเสน สทฺธิํ รตฺตนฺธกาเร อปฺปทีเป เอเกเนกา สนฺติฏฺฐติปิ สลฺลปติปิ.\csromanspacer{} ยา ตา ภิกฺขุนิโย อปฺปิจฺฉา ฯเปฯ ตา อุชฺฌายนฺติ ขิยฺยนฺติ วิปาเจนฺติ “กถํ หิ นาม ภิกฺขุนี รตฺตนฺธกาเร อปฺปทีเป ปุริเสน สทฺธิํ เอเกเนกา สนฺติฏฺฐิสฺสติปิ สลฺลปิสฺสติปี”ติ ฯเปฯ.\csromanspacer{} สจฺจํ กิร ภิกฺขเว ภิกฺขุนี รตฺตนฺธกาเร อปฺปทีเป ปุริเสน สทฺธิํ เอเกเนกา สนฺติฏฺฐติปิ สลฺลปติปีติ.\csromanspacer{} สจฺจํ ภควาติ.\csromanspacer{} วิครหิ พุทฺโธ ภควา ฯเปฯ.\csromanspacer{} กถํ หิ นาม \csromanfolio{352}ภิกฺขเว ภิกฺขุนี รตฺตนฺธกาเร อปฺปทีเป ปุริเสน สทฺธิํ เอเกเนกา สนฺติฏฺฐิสฺสติปิ สลฺลปิสฺสติปิ.\csromanspacer{} เนตํ ภิกฺขเว อปฺปสนฺนานํ วา ปสาทาย - ป-.\csromanspacer{} เอวญฺจ ปน ภิกฺขเว ภิกฺขุนิโย อิมํ สิกฺขาปทํ อุทฺทิสนฺตุ\csromandash{}}

\proseitem{839}{\textbf{“ยา ปน ภิกฺขุนี รตฺตนฺธกาเร อปฺปทีเป ปุริเสน สทฺธิํ เอเกเนกา สนฺติฏฺเฐยฺย วา สลฺลเปยฺย วา, ปาจิตฺติยนฺ”}ติ.}

\proseitem{840}{\textbf{ยา ปนา}ติ ยา ยาทิสา ฯเปฯ.\csromanspacer{} \textbf{ภิกฺขุนี}ติ ฯเปฯ อยํ อิมสฺมิํ อตฺเถ อธิปฺเปตา ภิกฺขุนีติ.}

\prose{\textbf{รตฺตนฺธกาเร}ติ โอคฺคเต สูริเย.\csromanspacer{} \textbf{อปฺปทีเป}ติ อนาโลเก.}

\prose{\textbf{ปุริโส} นาม มนุสฺสปุริโส น ยกฺโข น เปโต น ติรจฺฉานคโต วิญฺญู ปฏิพโล สนฺติฏฺฐิตุํ สลฺลปิตุํ.}

\prose{\textbf{สทฺธินฺ}ติ เอกโต\textbf{.}\csromanspacer{} เอเกเนกาติ ปุริโส เจว โหติ ภิกฺขุนี จ.}

\prose{\textbf{สนฺติฏฺเฐยฺย วา}ติ ปุริสสฺส หตฺถปาเส ติฏฺฐติ, อาปตฺติ ปาจิตฺติยสฺส.}

\prose{\textbf{สลฺลเปยฺย วา}ติ ปุริสสฺส หตฺถปาเส ฐิตา สลฺลปติ, อาปตฺติ ปาจิตฺติยสฺส.}

\prose{หตฺถปาสํ วิชหิตฺวา สนฺติฏฺฐติ วา สลฺลปติ วา, อาปตฺติ ทุกฺกฏสฺส.\csromanspacer{} ยกฺเขน วา เปเตน วา ปณฺฑเกน วา ติรจฺฉานคตมนุสฺสวิคฺคเหน วา สทฺธิํ สนฺติฏฺฐติ วา สลฺลปติ วา, อาปตฺติ ทุกฺกฏสฺส.}

\proseitemwithcloserruled{841}{อนาปตฺติ โย โกจิ วิญฺญู ทุติโย\footnote{ยา กาจิ วิญฺญู ทุติยา (สฺยา)} โหติ, อรโหเปกฺขา, อญฺญวิหิตา สนฺติฏฺฐติ วา สลฺลปติ วา, อุมฺมตฺติกาย อาทิกมฺมิกายาติ.}{ปฐมสิกฺขาปทํ นิฏฺฐิตํ.}
\csromanmatikaanchor{mtk.162}
\csromantocmarkref{subsection}{2. ทุติยสิกฺขาปท}{mtk.162}
\bookmark[dest={mtk.162},level=2]{2. ทุติยสิกฺขาปท}
\csromanheader{2}{2. \textbf{อนฺธการวคฺค 2. ทุติยสิกฺขาปท}}
\proseitem{842}{เตน สมเยน พุทฺโธ ภควา สาวตฺถิยํ วิหรติ เชตวเน อนาถปิณฺฑิกสฺส อาราเม.\csromanspacer{} เตน โข ปน สมเยน ภทฺทาย กาปิลานิยา อนฺเตวาสินิยา ภิกฺขุนิยา ญาตโก ปุริโส คามกา \csromanfolio{353}สาวตฺถิํ อคมาสิ เกนจิเทว กรณีเยน.\csromanspacer{} อถ โข สา ภิกฺขุนี “ภควตา ปฏิกฺขิตฺตํ รตฺตนฺธกาเร อปฺปทีเป ปุริเสน สทฺธิํ เอเกเนกา สนฺติฏฺฐิตุํ สลฺลปิตุนฺ”ติ เตเนว ปุริเสน สทฺธิํ ปฏิจฺฉนฺเน โอกาเส เอเกเนกา สนฺติฏฺฐติปิ สลฺลปติปิ.\csromanspacer{} ยา ตา ภิกฺขุนิโย อปฺปิจฺฉา ฯเปฯ ตา อุชฺฌายนฺติ ขิยฺยนฺติ วิปาเจนฺติ “กถํ หิ นาม ภิกฺขุนี ปฏิจฺฉนฺเน โอกาเส ปุริเสน สทฺธิํ เอเกเนกา สนฺติฏฺฐิสฺสติปิ สลฺลปิสฺสติปี”ติ ฯเปฯ.\csromanspacer{} \textbf{สจฺจํ กิร ภิกฺขเว ภิกฺขุนี ปฏิจฺฉนฺเน โอกาเส ปุริเสน สทฺธิํ} \textbf{เอเกเนกา สนฺติฏฺฐติปิ สลฺลปติปีติ.}\csromanspacer{}\textbf{ สจฺจํ ภควาติ.}\csromanspacer{}\textbf{ วิครหิ พุทฺโธ} ภควา ฯเปฯ.\csromanspacer{} กถํ หิ นาม ภิกฺขเว ภิกฺขุนี ปฏิจฺฉนฺเน โอกาเส ปุริเสน สทฺธิํ เอเกเนกา สนฺติฏฺฐิสฺสติปิ สลฺลปิสฺสติปิ.\csromanspacer{} เนตํ ภิกฺขเว อปฺปสนฺนานํ วา ปสาทาย ฯเปฯ.\csromanspacer{} เอวญฺจ ปน ภิกฺขเว ภิกฺขุนิโย อิมํ สิกฺขาปทํ อุทฺทิสนฺตุ\csromandash{}}

\proseitem{843}{\textbf{“ยา ปน ภิกฺขุนี ปฏิจฺฉนฺเน โอกาเส ปุริเสน สทฺธิํ เอเกเนกา สนฺติฏฺเฐยฺย วา สลฺลเปยฺย วา, ปาจิตฺติยนฺ”}ติ.}

\proseitem{844}{\textbf{ยา ปนา}ติ ยา ยาทิสา ฯเปฯ.\csromanspacer{} \textbf{ภิกฺขุนี}ติ ฯเปฯ อยํ อิมสฺมิํ อตฺเถ อธิปฺเปตา ภิกฺขุนีติ.}

\prose{\textbf{ปฏิจฺฉนฺโน} นาม โอกาโส กุฏฺเฏน วา กวาเฏน วา กิลญฺเชน วา สาณิปากาเรน วา รุกฺเขน วา ถมฺเภน วา โกตฺถฬิยา วา เยน เกนจิ ปฏิจฺฉนฺโน โหติ.}

\prose{\textbf{ปุริโส} นาม มนุสฺสปุริโส น ยกฺโข น เปโต น ติรจฺฉานคโต วิญฺญู ปฏิพโล สนฺติฏฺฐิตุํ สลฺลปิตุํ.}

\prose{\textbf{สทฺธินฺ}ติ เอกโต.\csromanspacer{} \textbf{เอเกเนกา}ติ ปุริโส เจว โหติ ภิกฺขุนี จ.}

\prose{\textbf{สนฺติฏฺเฐยฺย วา}ติ ปุริสสฺส หตฺถปาเส ติฏฺฐติ, อาปตฺติ ปาจิตฺติยสฺส.}

\prose{\textbf{สลฺลเปยฺย วา}ติ ปุริสสฺส หตฺถปาเส ฐิตา สลฺลปติ, อาปตฺติ ปาจิตฺติยสฺส.}

\prose{หตฺถปาสํ วิชหิตฺวา สนฺติฏฺฐติ วา สลฺลปติ วา, อาปตฺติ ทุกฺกฏสฺส.\csromanspacer{} ยกฺเขน วา เปเตน วา ปณฺฑเกน วา ติรจฺฉานคตมนุสฺสวิคฺคเหน วา สทฺธิํ สนฺติฏฺฐติ วา สลฺลปติ วา, อาปตฺติ ทุกฺกฏสฺส.}

\proseitemwithcloserruled{845}{\csromanfolio{354}อนาปตฺติ โย โกจิ วิญฺญู ทุติโย โหติ, อรโหเปกฺขา, อญฺญวิหิตา สนฺติฏฺฐติ วา สลฺลปติ วา, อุมฺมตฺติกาย อาทิกมฺมิกายาติ.}{ทุติยสิกฺขาปทํ นิฏฺฐิตํ.}
\csromanmatikaanchor{mtk.163}
\csromantocmarkref{subsection}{3. ตติยสิกฺขาปท}{mtk.163}
\bookmark[dest={mtk.163},level=2]{3. ตติยสิกฺขาปท}
\csromanheader{2}{2. \textbf{อนฺธการวคฺค 3. ตติยสิกฺขาปท}}
\proseitem{846}{เตน สมเยน พุทฺโธ ภควา สาวตฺถิยํ วิหรติ เชตวเน อนาถปิณฺฑิกสฺส อาราเม.\csromanspacer{} เตน โข ปน สมเยน ภทฺทาย กาปิลานิยา อนฺเตวาสินิยา ภิกฺขุนิยา ญาตโก ปุริโส คามกา สาวตฺถิํ อคมาสิ เกนจิเทว กรณีเยน.\csromanspacer{} อถ โข สา ภิกฺขุนี “ภควตา ปฏิกฺขิตฺตํ ปฏิจฺฉนฺเน โอกาเส ปุริเสน สทฺธิํ เอเกเนกา สนฺติฏฺฐิตุํ สลฺลปิตุนฺ”ติ เตเนว ปุริเสน สทฺธิํ อชฺโฌกาเส เอเกเนกา สนฺติฏฺฐติปิ สลฺลปติปิ.\csromanspacer{} ยา ตา ภิกฺขุนิโย อปฺปิจฺจา ฯเปฯ ตา อุชฺฌายนฺติ ขิยฺยนฺติ วิปาเจนฺติ “กถํ หิ นาม ภิกฺขุนี อชฺโฌกาเส ปุริเสน สทฺธิํ เอเกเนกา สนฺติฏฺฐิสฺสติปิ สลฺลปิสฺสติปี”ติ ฯเปฯ.\csromanspacer{} สจฺจํ กิร ภิกฺขเว ภิกฺขุนี อชฺโฌกาเส ปุริเสน สทฺธิํ เอเกเนกา สนฺติฏฺฐติปิ สลฺลปติปีติ.\csromanspacer{} สจฺจํ ภควาติ.\csromanspacer{} วิครหิ พุทฺโธ ภควา ฯเปฯ.\csromanspacer{} กถํ หิ นาม ภิกฺขเว ภิกฺขุนี อชฺโฌกาเส ปุริเสน สทฺธิํ เอเกเนกา สนฺติฏฺฐิสฺสติปิ สลฺลปิสฺสติปิ.\csromanspacer{} เนตํ ภิกฺขเว อปฺปสนฺนานํ วา ปสาทาย ฯเปฯ.\csromanspacer{} เอวญฺจ ปน ภิกฺขเว ภิกฺขุนิโย อิมํ สิกฺขาปทํ อุทฺทิสนฺตุ\csromandash{}}

\proseitem{847}{\textbf{“ยา ปน ภิกฺขุนี อชฺโฌกาเส ปุริเสน สทฺธิํ เอเกเนกา สนฺติฏฺเฐยฺย วา สลฺลเปยฺย วา, ปาจิตฺติยนฺ”}ติ.}

\proseitem{848}{\textbf{ยา ปนา}ติ ยา ยาทิสา ฯเปฯ.\csromanspacer{} \textbf{ภิกฺขุนี}ติ ฯเปฯ อยํ อิมสฺมิํ อตฺเถ อธิปฺเปตา ภิกฺขุนีติ.}

\prose{\textbf{อชฺโฌกาโส} นาม อปฺปฏิจฺฉนฺโน โหติ กุฏฺเฏน วา กวาเฏน วา กิลญฺเชน วา สาณิปากาเรน วา รุกฺเขน วา ถมฺเภน วา โกตฺถฬิยา วา, เยน เกนจิ อปฺปฏิจฺฉนฺโน โหติ.}

\prose{\csromanfolio{355}\textbf{ปุริโส} นาม มนุสฺสปุริโส น ยกฺโข น เปโต น ติรจฺฉานคโต วิญฺญู ปฏิพโล สนฺติฏฺฐิตุํ สลฺลปิตุํ.}

\prose{\textbf{สทฺธินฺ}ติ เอกโต.\csromanspacer{} \textbf{เอเกเนกา}ติ ปุริโส เจว โหติ ภิกฺขุนี จ.}

\prose{\textbf{สนฺติฏฺเฐยฺย วา}ติ ปุริสสฺส หตฺถปาเส ติฏฺฐติ, อาปตฺติ ปาจิตฺติยสฺส.}

\prose{\textbf{สลฺลเปยฺย วา}ติ ปุริสสฺส หตฺถปาเส ฐิตา สลฺลปติ, อาปตฺติ ปาจิตฺติยสฺส.}

\prose{หตฺถปาสํ วิชหิตฺวา สนฺติฏฺฐติ วา สลฺลปติ วา, อาปตฺติ ทุกฺกฏสฺส.\csromanspacer{} ยกฺเขน วา เปเตน วา ปณฺฑเกน วา ติรจฺฉานคตมนุสฺสวิคฺคเหน วา สทฺธิํ สนฺติฏฺฐติ วา สลฺลปติ วา, อาปตฺติ ทุกฺกฏสฺส.}

\proseitemwithcloserruled{849}{อนาปตฺติ โย โกจิ วิญฺญู ทุติโย โหติ, อรโหเปกฺขา, อญฺญวิหิตา สนฺติฏฺฐติ วา สลฺลปติ วา, อุมฺมตฺติกาย อาทิกมฺมิกายาติ.}{ตติยสิกฺขาปทํ นิฏฺฐิตํ.}
\csromanmatikaanchor{mtk.164}
\csromantocmarkref{subsection}{4. จตุตฺถสิกฺขาปท}{mtk.164}
\bookmark[dest={mtk.164},level=2]{4. จตุตฺถสิกฺขาปท}
\csromanheader{2}{2. \textbf{อนฺธการวคฺค 4. จตุตฺถสิกฺขาปท}}
\proseitem{850}{เตน สมเยน พุทฺโธ ภควา สาวตฺถิยํ วิหรติ เชตวเน อนาถปิณฺฑิกสฺส อาราเม.\csromanspacer{} เตน โข ปน สมเยน ถุลฺลนนฺทา ภิกฺขุนี รถิกายปิ พฺยูเหปิ สิงฺฆาฏเกปิ ปุริเสน สทฺธิํ เอเกเนกา สนฺติฏฺฐติปิ สลฺลปติปิ นิกณฺณิกมฺปิ ชปฺเปติ ทุติยิกมฺปิ ภิกฺขุนิํ อุยฺโยเชติ.\csromanspacer{} ยา ตา ภิกฺขุนิโย อปฺปิจฺฉา ฯเปฯ ตา อุชฺฌายนฺติ ขิยฺยนฺติ วิปาเจนฺติ “กถํ หิ นาม อยฺยา ถุลฺลนนฺทา รถิกายปิ พฺยูเหปิ สิงฺฆาฏเกปิ ปุริเสน สทฺธิํ เอเกเนกา สนฺติฏฺฐิสฺสติปิ สลฺลปิสฺสติปิ นิกณฺณิกมฺปิ ชปฺปิสฺสติ ทุติยิกมฺปิ ภิกฺขุนิํ อุยฺโยเชสฺสตี”ติ ฯเปฯ.\csromanspacer{} สจฺจํ กิร ภิกฺขเว ถุลฺลนนฺทา ภิกฺขุนี รถิกายปิ พฺยูเหปิ สิงฺฆาฏเกปิ ปุริเสน สทฺธิํ เอเกเนกา สนฺติฏฺฐติปิ สลฺลปติปิ นิกณฺณิกมฺปิ ชปฺเปติ ทุติยิกมฺปิ ภิกฺขุนิํ อุยฺโยเชตีติ.\csromanspacer{} สจฺจํ ภควาติ.\csromanspacer{} วิครหิ พุทฺโธ ภควา ฯเปฯ.\csromanspacer{} กถํ หิ นาม ภิกฺขเว ถุลฺลนนฺทา ภิกฺขุนี รถิกายปิ พฺยูเหปิ สิงฺฆาฏเกปิ ปุริเสน สทฺธิํ เอเกเนกา สนฺติฏฺฐิสฺสติปิ สลฺลปิสฺสติปิ นิกณฺณิกมฺปิ ชปฺปิสฺสติ \csromanfolio{356}\textbf{ทุติยิกมฺปิ ภิกฺขุนิํ อุยฺโยเชสฺสติ.}\csromanspacer{}\textbf{ เนตํ ภิกฺขเว อปฺปสนฺนานํ วา} \textbf{ปสาทาย ฯเปฯ.}\csromanspacer{}\textbf{ เอวญฺจ ปน ภิกฺขเว ภิกฺขุนิโย อิมํ สิกฺขาปทํ} อุทฺทิสนฺตุ\csromandash{}}

\proseitem{851}{\textbf{“ยา ปน ภิกฺขุนี รถิกาย วา พฺยูเห วา สิงฺฆาฏเก วา ปุริเสน สทฺธิํ เอเกเนกา สนฺติฏฺเฐยฺย วา สลฺลเปยฺย วา นิกณฺณิกํ วา ชปฺเปยฺย ทุติยิกํ วา ภิกฺขุนิํ อุยฺโยเชยฺย, ปาจิตฺติยนฺ”}ติ.}

\proseitem{852}{\textbf{ยา ปนา}ติ ยา ยาทิสา ฯเปฯ.\csromanspacer{} \textbf{ภิกฺขุนี}ติ ฯเปฯ อยํ อิมสฺมิํ อตฺเถ อธิปฺเปตา ภิกฺขุนีติ.}

\prose{\textbf{รถิกา} นาม รจฺฉา วุจฺจติ.\csromanspacer{} \textbf{พฺยูหํ} นาม เยเนว ปวิสนฺติ, เตเนว นิกฺขมนฺติ.\csromanspacer{} \textbf{สิงฺฆาฏโก} นาม จจฺจรํ วุจฺจติ.}

\prose{\textbf{ปุริโส} นาม มนุสฺสปุริโส น ยกฺโข น เปโต น ติรจฺฉานคโต วิญฺญู ปฏิพโล สนฺติฏฺฐิตุํ สลฺลปิตุํ.}

\prose{\textbf{สทฺธินฺ}ติ เอกโต.\csromanspacer{} \textbf{เอเกเนกา}ติ ปุริโส เจว โหติ ภิกฺขุนี จ.}

\prose{\textbf{สนฺติฏฺเฐยฺย วา}ติ ปุริสสฺส หตฺถปาเส ติฏฺฐติ, อาปตฺติ ปาจิตฺติยสฺส.}

\prose{\textbf{สลฺลเปยฺย วา}ติ ปุริสสฺส หตฺถปาเส ฐิตา สลฺลปติ, อาปตฺติ ปาจิตฺติยสฺส.}

\prose{\textbf{นิกณฺณิกํ วา ชปฺเปยฺยา}ติ ปุริสสฺส อุปกณฺณเก อาโรเจติ, อาปตฺติ}

\prose{\textbf{ทุติยิกํ วา ภิกฺขุนิํ อุยฺโยเชยฺยา}ติ อนาจารํ อาจริตุกามา ทุติยิกมฺปิ ภิกฺขุนิํ อุยฺโยเชติ, อาปตฺติ ทุกฺกฏสฺส.\csromanspacer{} ทสฺสนูปจารํ วา สวนูปจารํ วา วิชหนฺติยา อาปตฺติ ทุกฺกฏสฺส.\csromanspacer{} วิชหิเต อาปตฺติ ปาจิตฺติยสฺส.}

\prose{หตฺถปาสํ วิชหิตฺวา สนฺติฏฺฐติ วา สลฺลปติ วา, อาปตฺติ ทุกฺกฏสฺส.\csromanspacer{} ยกฺเขน วา เปเตน วา ปณฺฑเกน วา ติรจฺฉานคตมนุสฺสวิคฺคเหน วา สทฺธิํ สนฺติฏฺฐติ วา สลฺลปติ วา อาปตฺติ ทุกฺกฏสฺส.}

\proseitemwithcloser{853}{อนาปตฺติ โย โกจิ วิญฺญู ทุติโย โหติ, อรโหเปกฺขา, อญฺญวิหิตา สนฺติฏฺฐติ วา สลฺลปติ วา, น อนาจารํ อาจริตุกามา, สติ กรณีเย ทุติยิกํ ภิกฺขุนิํ อุยฺโยเชติ, อุมฺมตฺติกาย อาทิกมฺมิกายาติ.}{จตุตฺถสิกฺขาปทํ นิฏฺฐิตํ.}
\csromanmatikaanchor{mtk.165}
\csromantocmarkref{subsection}{5. ปญฺจมสิกฺขาปท}{mtk.165}
\bookmark[dest={mtk.165},level=2]{5. ปญฺจมสิกฺขาปท}
\csromanheader{2}{2. \textbf{อนฺธการวคฺค 5. ปญฺจมสิกฺขาปท}}
\proseitem{854}{\csromanfolio{357}เตน สมเยน พุทฺโธ ภควา สาวตฺถิยํ วิหรติ เชตวเน อนาถปิณฺฑิกสฺส อาราเม.\csromanspacer{} เตน โข ปน สมเยน อญฺญตรา ภิกฺขุนี อญฺญตรสฺส กุลสฺส กุลูปิกา โหติ นิจฺจภตฺติกา.\csromanspacer{} อถ โข สา ภิกฺขุนี ปุพฺพณฺหสมยํ นิวาเสตฺวา ปตฺตจีวรมาทาย เยน ตํ กุลํ เตนุปสงฺกมิ, อุปสงฺกมิตฺวา อาสเน นิสีทิตฺวา สามิเก อนาปุจฺฉา ปกฺกามิ.\csromanspacer{} ตสฺส กุลสฺส ทาสี ฆรํ สมฺมชฺชนฺตี ตํ อาสนํ ภาชนนฺตริกาย ปกฺขิปิ.\csromanspacer{} มนุสฺสา ตํ อาสนํ อปสฺสนฺตา ตํ ภิกฺขุนิํ เอตทโวจุํ “กหํ ตํ อยฺเย อาสนนฺ”ติ.\csromanspacer{} นาหํ ตํ อาวุโส อาสนํ ปสฺสามีติ. “เทถายฺเย ตํ อาสนนฺ”ติ ปริภาสิตฺวา นิจฺจภตฺตํ ปจฺฉินฺทิํสุ.\csromanspacer{} อถ โข เต มนุสฺสา ฆรํ โสเธนฺตา ตํ อาสนํ ภาชนนฺตริกาย ปสฺสิตฺวา ตํ ภิกฺขุนิํ ขมาเปตฺวา นิจฺจภตฺตํ ปฏฺฐเปสุํ.\csromanspacer{} อถ โข สา ภิกฺขุนี ภิกฺขุนีนํ เอตมตฺถํ อาโรเจสิ.\csromanspacer{} ยา ตา ภิกฺขุนิโย อปฺปิจฺฉา ฯเปฯ ตา อุชฺฌายนฺติ ขิยฺยนฺติ วิปาเจนฺติ “กถํ หิ นาม ภิกฺขุนี ปุเรภตฺตํ กุลานิ อุปสงฺกมิตฺวา อาสเน นิสีทิตฺวา สามิเก อนาปุจฺฉา ปกฺกมิสฺสตี”ติ ฯเปฯ .\csromanspacer{} สจฺจํ กิร ภิกฺขเว ภิกฺขุนี ปุเรภตฺตํ กุลานิ อุปสงฺกมิตฺวา อาสเน นิสีทิตฺวา สามิเก อนาปุจฺฉา ปกฺกมตีติ\footnote{ปกฺกมีติ (ก)}.\csromanspacer{} สจฺจํ ภควาติ.\csromanspacer{} วิครหิ พุทฺโธ ภควา ฯเปฯ.\csromanspacer{} กถํ หิ นาม ภิกฺขเว ภิกฺขุนี ปุเรภตฺตํ กุลานิ อุปสงฺกมิตฺวา อาสเน นิสีทิตฺวา สามิเก อนาปุจฺฉา ปกฺกมิสฺสติ.\csromanspacer{} เนตํ ภิกฺขเว อปฺปสนฺนานํ วา ปสาทาย ฯเปฯ.\csromanspacer{} เอวญฺจ ปน ภิกฺขเว ภิกฺขุนิโย อิมํ สิกฺขาปทํ อุทฺทิสนฺตุ\csromandash{}}

\proseitem{855}{\textbf{“ยา ปน ภิกฺขุนี ปุเรภตฺตํ กุลานิ อุปสงฺกมิตฺวา อาสเน นิสีทิตฺวา สามิเก อนาปุจฺฉา ปกฺกเมยฺย, ปาจิตฺติยนฺ”}ติ.}

\proseitem{856}{\textbf{ยา ปนา}ติ ยา ยาทิสา ฯเปฯ.\csromanspacer{} \textbf{ภิกฺขุนี}ติ ฯเปฯ อยํ อิมสฺมิํ อตฺเถ อธิปฺเปตา ภิกฺขุนีติ.}

\prose{\textbf{ปุเรภตฺตํ} นาม อรุณุคฺคมนํ อุปาทาย ยาว มชฺฌนฺหิกา.}

\prose{\textbf{กุลํ} นาม จตฺตาริ กุลานิ ขตฺติยกุลํ พฺราหฺมณกุลํ เวสฺสกุลํ สุทฺทกุลํ.\csromanspacer{} \textbf{อุปสงฺกมิตฺวา}ติ ตตฺถ คนฺตฺวา.}

\prose{\csromanfolio{358}\textbf{อาสนํ} นาม ปลฺลงฺกสฺส โอกาโส วุจฺจติ.\csromanspacer{} \textbf{นิสีทิตฺวา}ติ ตสฺมิํ นิสีทิตฺวา.}

\prose{\textbf{สามิเก อนาปุจฺฉา ปกฺกเมยฺยา}ติ โย ตสฺมิํ กุเล มนุสฺโส วิญฺญู, ตํ อนาปุจฺฉา อโนวสฺสกํ อติกฺกาเมนฺติยา อาปตฺติ ปาจิตฺติยสฺส.\csromanspacer{} อชฺโฌกาเส อุปจารํ อติกฺกาเมนฺติยา อาปตฺติ ปาจิตฺติยสฺส.}

\proseitem{857}{อนาปุจฺฉิเต อนาปุจฺฉิตสญฺญา ปกฺกมติ, อาปตฺติ ปาจิตฺติยสฺส.\csromanspacer{} อนาปุจฺฉิเต เวมติกา ปกฺกมติ, อาปตฺติ ปาจิตฺติยสฺส.\csromanspacer{} อนาปุจฺฉิเต อาปุจฺฉิตสญฺญา ปกฺกมติ, อาปตฺติ ปาจิตฺติยสฺส.}

\prose{ปลฺลงฺกสฺส อโนกาเส อาปตฺติ ทุกฺกฏสฺส.\csromanspacer{} อาปุจฺฉิเต อนาปุจฺฉิตสญฺญา, อาปตฺติ ทุกฺกฏสฺส.\csromanspacer{} อาปุจฺฉิเต เวมติกา, อาปตฺติ ทุกฺกฏสฺส.\csromanspacer{} อาปุจฺฉิเต อาปุจฺฉิตสญฺญา, อนาปตฺติ.}

\proseitemwithcloserruled{858}{อนาปตฺติ อาปุจฺฉา คจฺฉติ, อสํหาริเม คิลานาย อาปทาสุ อุมฺมตฺติกาย อาทิกมฺมิกายาติ.}{ปญฺจมสิกฺขาปทํ นิฏฺฐิตํ.}
\csromanmatikaanchor{mtk.166}
\csromantocmarkref{subsection}{6. ฉฏฺฐสิกฺขาปท}{mtk.166}
\bookmark[dest={mtk.166},level=2]{6. ฉฏฺฐสิกฺขาปท}
\csromanheader{2}{2. \textbf{อนฺธการวคฺค 6. ฉฏฺฐสิกฺขาปท}}
\proseitem{859}{เตน สมเยน พุทฺโธ ภควา สาวตฺถิยํ วิหรติ เชตวเน อนาถปิณฺฑิกสฺส อาราเม.\csromanspacer{} เตน โข ปน สมเยน ถุลฺลนนฺทา ภิกฺขุนี ปจฺฉาภตฺตํ กุลานิ อุปสงฺกมิตฺวา สามิเก อนาปุจฺฉา อาสเน อภินิสีทติปิ อภินิปชฺชติปิ.\csromanspacer{} มนุสฺสา ถุลฺลนนฺทํ ภิกฺขุนิํ หิรียมานา อาสเน เนว อภินิสีทนฺติ น อภินิปชฺชนฺติ.\csromanspacer{} มนุสฺสา อุชฺฌายนฺติ ขิยฺยนฺติ วิปาเจนฺติ “กถํ หิ นาม อยฺยา ถุลฺลนนฺทา ปจฺฉาภตฺตํ กุลานิ อุปสงฺกมิตฺวา สามิเก อนาปุจฺฉา อาสเน อภินิสีทิสฺสติปิ อภินิปชฺชิสฺสติปี”ติ.\csromanspacer{} อสฺโสสุํ โข ภิกฺขุนิโย เตสํ มนุสฺสานํ อุชฺฌายนฺตานํ ขิยฺยนฺตานํ วิปาเจนฺตานํ.\csromanspacer{} ยา ตา ภิกฺขุนิโย อปฺปิจฺฉา ฯเปฯ ตา อุชฺฌายนฺติ ขิยฺยนฺติ วิปาเจนฺติ “กถํ หิ นาม อยฺยา ถุลฺลนนฺทา ปจฺฉาภตฺตํ \csromanfolio{359}กุลานิ อุปสงฺกมิตฺวา สามิเก อนาปุจฺฉา อาสเน อภินิสีทิสฺสติปิ อภินิปชฺชิสฺสติปี”ติ ฯเปฯ.\csromanspacer{} สจฺจํ กิร ภิกฺขเว ถุลฺลนนฺทา ภิกฺขุนี ปจฺฉาภตฺตํ กุลานิ อุปสงฺกมิตฺวา สามิเก อนาปุจฺฉา อาสเน อภินิสีทติปิ อภินิปชฺชติปีติ.\csromanspacer{} สจฺจํ ภควาติ.\csromanspacer{} วิครหิ พุทฺโธ ภควา ฯเปฯ.\csromanspacer{} กถํ หิ นาม ภิกฺขเว ถุลฺลนนฺทา ภิกฺขุนี ปจฺฉาภตฺตํ กุลานิ อุปสงฺกมิตฺวา สามิเก อนาปุจฺฉา อาสเน อภินิสีทิสฺสติปิ อภินิปชฺชิสฺสติปิ.\csromanspacer{} เนตํ ภิกฺขเว อปฺปสนฺนานํ วา ปสาทาย ฯเปฯ.\csromanspacer{} เอวญฺจ ปน ภิกฺขเว ภิกฺขุนิโย อิมํ สิกฺขาปทํ อุทฺทิสนฺตุ\csromandash{}}

\proseitem{860}{\textbf{“ยา ปน ภิกฺขุนี ปจฺฉาภตฺตํ กุลานิ อุปสงฺกมิตฺวา สามิเก อนาปุจฺฉา อาสเน อภินิสีเทยฺย วา อภินิปชฺเชยฺย วา, ปาจิตฺติยนฺ”}ติ.}

\proseitem{861}{\textbf{ยา ปนา}ติ ยา ยาทิสา ฯเปฯ.\csromanspacer{} \textbf{ภิกฺขุนี}ติ ฯเปฯ อยํ อิมสฺมิํ อตฺเถ อธิปฺเปตา ภิกฺขุนีติ.}

\prose{\textbf{ปจฺฉาภตฺตํ} นาม มชฺฌนฺหิเก วีติวตฺเต ยาว อตฺถงฺคเต สูริเย.}

\prose{\textbf{กุลํ} นาม จตฺตาริ กุลานิ ขตฺติยกุลํ พฺราหฺมณกุลํ เวสฺสกุลํ สุทฺทกุลํ.\csromanspacer{} \textbf{อุปสงฺกมิตฺวา}ติ ตตฺถ คนฺตฺวา.}

\prose{\textbf{สามิเก อนาปุจฺฉา}ติ โย ตสฺมิํ กุเล มนุสฺโส สามิโก ทาตุํ ตํ อนาปุจฺฉา.}

\prose{\textbf{อาสนํ} นาม ปลฺลงฺกสฺส โอกาโส วุจฺจติ.}

\prose{\textbf{อภินิสีเทยฺยา}ติ ตสฺมิํ อภินิสีทติ, อาปตฺติ ปาจิตฺติยสฺส.}

\prose{\textbf{อภินิปชฺเชยฺยา}ติ ตสฺมิํ อภินิปชฺชติ, อาปตฺติ ปาจิตฺติยสฺส.}

\proseitem{862}{อนาปุจฺฉิเต อนาปุจฺฉิตสญฺญา อาสเน อภินิสีทติ วา อภินิปชฺชติ วา, อาปตฺติ ปาจิตฺติยสฺส.\csromanspacer{} อนาปุจฺฉิเต เวมติกา อาสเน อภินิสีทติ วา อภินิปชฺชติ วา, อาปตฺติ ปาจิตฺติยสฺส.\csromanspacer{} อนาปุจฺฉิเต อาปุจฺฉิตสญฺญา อาสเน อภินิสีทติ วา อภินิปชฺชติ วา, อาปตฺติ ปาจิตฺติยสฺส.}

\prose{\csromanfolio{360}ปลฺลงฺกสฺส อโนกาเส อาปตฺติ ทุกฺกฏสฺส.\csromanspacer{} อาปุจฺฉิเต อนาปุจฺฉิตสญฺญา, อาปตฺติ ทุกฺกฏสฺส.\csromanspacer{} อาปุจฺฉิเต เวมติกา, อาปตฺติ ทุกฺกฏสฺส.\csromanspacer{} อาปุจฺฉิเต อาปุจฺฉิตสญฺญา, อนาปตฺติ.}

\proseitemwithcloserruled{863}{อนาปตฺติ อาปุจฺฉา อาสเน อภินิสีทติ วา อภินิปชฺชติ วา, ธุวปญฺญตฺเต คิลานาย อาปทาสุ อุมฺมตฺติกาย อาทิกมฺมิกายาติ.}{ฉฏฺฐสิกฺขาปทํ นิฏฺฐิตํ.}
\csromanmatikaanchor{mtk.167}
\csromantocmarkref{subsection}{7. สตฺตมสิกฺขาปท}{mtk.167}
\bookmark[dest={mtk.167},level=2]{7. สตฺตมสิกฺขาปท}
\csromanheader{2}{2. \textbf{อนฺธการวคฺค 7. สตฺตมสิกฺขาปท}}
\proseitem{864}{เตน สมเยน พุทฺโธ ภควา สาวตฺถิยํ วิหรติ เชตวเน อนาถปิณฺฑิกสฺส อาราเม.\csromanspacer{} เตน โข ปน สมเยน สมฺพหุลา ภิกฺขุนิโย โกสเลสุ ชนปเท สาวตฺถิํ คจฺฉนฺติโย สายํ อญฺญตรํ คามํ อุปคนฺตฺวา อญฺญตรํ พฺราหฺมณกุลํ อุปสงฺกมิตฺวา โอกาสํ ยาจิํสุ.\csromanspacer{} อถ โข สา พฺราหฺมณี ตา ภิกฺขุนิโย เอตทโวจ “อาคเมถ อยฺเย ยาว พฺราหฺมโณ อาคจฺฉตี”ติ.\csromanspacer{} ภิกฺขุนิโย “ยาว พฺราหฺมโณ อาคจฺฉตี”ติ เสยฺยํ สนฺถริตฺวา เอกจฺจา นิสีทิํสุ เอกจฺจา นิปชฺชิํสุ.\csromanspacer{} อถ โข โส พฺราหฺมโณ รตฺติํ อาคนฺตฺวา ตํ พฺราหฺมณิํ เอตทโวจ “กา อิมา”ติ.\csromanspacer{} ภิกฺขุนิโย อยฺยาติ. “นิกฺกฑฺฒถ อิมา มุณฺฑา พนฺธกินิโย”ติ ฆรโต นิกฺกฑฺฒาเปสิ.\csromanspacer{} อถ โข ตา ภิกฺขุนิโย สาวตฺถิํ คนฺตฺวา ภิกฺขุนีนํ เอตมตฺถํ อาโรเจสุํ.\csromanspacer{} ยา ตา ภิกฺขุนิโย อปฺปิจฺฉา ฯเปฯ ตา อุชฺฌายนฺติ ขิยฺยนฺติ วิปาเจนฺติ “กถํ หิ นาม ภิกฺขุนิโย วิกาเล กุลานิ อุปสงฺกมิตฺวา สามิเก อนาปุจฺฉา เสยฺยํ สนฺถริตฺวา อภินิสีทิสฺสนฺติปิ อภินิปชฺชิสฺสนฺติปี”ติ ฯเปฯ.\csromanspacer{} สจฺจํ กิร ภิกฺขเว ภิกฺขุนิโย วิกาเล กุลานิ อุปสงฺกมิตฺวา สามิเก อนาปุจฺฉา เสยฺยํ สนฺถริตฺวา อภินิสีทนฺติปิ อภินิปชฺชนฺติปีติ.\csromanspacer{} สจฺจํ ภควาติ.\csromanspacer{} วิครหิ พุทฺโธ ภควา ฯเปฯ.\csromanspacer{} กถํ หิ นาม ภิกฺขเว ภิกฺขุนิโย วิกาเล กุลานิ อุปสงฺกมิตฺวา สามิเก อนาปุจฺฉา เสยฺยํ สนฺถริตฺวา อภินิสีทิสฺสนฺติปิ อภินิปชฺชิสฺสนฺติปิ.\csromanspacer{} เนตํ ภิกฺขเว อปฺปสนฺนานํ วา ปสาทาย ฯเปฯ.\csromanspacer{} เอวญฺจ ปน ภิกฺขเว ภิกฺขุนิโย อิมํ สิกฺขาปทํ อุทฺทิสนฺตุ\csromandash{}}

\proseitem{865}{\csromanfolio{361}\textbf{“ยา ปน ภิกฺขุนี วิกาเล กุลานิ อุปสงฺกมิตฺวา สามิเก อนาปุจฺฉา เสยฺยํ สนฺถริตฺวา วา สนฺถราเปตฺวา วา อภินิสีเทยฺย วา อภินิปชฺเชยฺย วา, ปาจิตฺติยนฺ”}ติ.}

\proseitem{866}{\textbf{ยา ปนา}ติ ยา ยาทิสา ฯเปฯ.\csromanspacer{} \textbf{ภิกฺขุนี}ติ ฯเปฯ อยํ อิมสฺมิํ อตฺเถ อธิปฺเปตา ภิกฺขุนีติ.}

\prose{\textbf{วิกาโล} นาม อตฺถงฺคเต สูริเย ยาว อรุณุคฺคมนา.}

\prose{\textbf{กุลํ} นาม จตฺตาริ กุลานิ ขตฺติยกุลํ พฺราหฺมณกุลํ เวสฺสกุลํ สุทฺทกุลํ.\csromanspacer{} \textbf{อุปสงฺกมิตฺวา}ติ ตตฺถ คนฺตฺวา.}

\prose{\textbf{สามิเก อนาปุจฺฉา}ติ โย ตสฺมิํ กุเล มนุสฺโส สามิโก ทาตุํ, ตํ อนาปุจฺฉา.}

\prose{\textbf{เสยฺยํ} นาม อนฺตมโส ปณฺณสนฺถาโรปิ.\csromanspacer{} \textbf{สนฺถริตฺวา}ติ สยํ สนฺถริตฺวา.\csromanspacer{} \textbf{สนฺถราเปตฺวา}ติ อญฺญํ สนฺถราเปตฺวา.\csromanspacer{} \textbf{อภินิสีเทยฺยา}ติ ตสฺมิํ อภินิสีทติ, อาปตฺติ ปาจิตฺติยสฺส.\csromanspacer{} \textbf{อภินิปชฺเชยฺยา}ติ ตสฺมิํ อภินิปชฺชติ, อาปตฺติ}

\proseitem{867}{อนาปุจฺฉิเต อนาปุจฺฉิตสญฺญา เสยฺยํ สนฺถริตฺวา วา สนฺถาราเปตฺวา วา อภินิสีทติ วา อภินิปชฺชติ วา, อาปตฺติ ปาจิตฺติยสฺส.\csromanspacer{} อนาปุจฺฉิเต เวมติกา เสยฺยํ สนฺถริตฺวา วา สนฺถราเปตฺวา วา อภินิสีทติ วา อภินิปชฺชติ วา, อาปตฺติ ปาจิตฺติยสฺส.\csromanspacer{} อนาปุจฺฉิเต อาปุจฺฉิตสญฺญา เสยฺยํ สนฺถริตฺวา วา สนฺถราเปตฺวา วา อภินิสีทติ วา อภินิปชฺชติ วา, อาปตฺติ}

\prose{อาปุจฺฉิเต อนาปุจฺฉิตสญฺญา, อาปตฺติ ทุกฺกฏสฺส.\csromanspacer{} อาปุจฺฉิเต เวมติกา, อาปตฺติ ทุกฺกฏสฺส.\csromanspacer{} อาปุจฺฉิเต อาปุจฺฉิตสญฺญา, อนาปตฺติ.}

\proseitemwithcloser{868}{อนาปตฺติ อาปุจฺฉา เสยฺยํ สนฺถริตฺวา วา สนฺถราเปตฺวา วา อภินิสีทติ วา อภินิปชฺชติ วา, คิลานาย อาปทาสุ อุมฺมตฺติกาย อาทิกมฺมิกายาติ.}{สตฺตมสิกฺขาปทํ นิฏฺฐิตํ.}
\csromanmatikaanchor{mtk.168}
\csromantocmarkref{subsection}{8. อฏฺฐมสิกฺขาปท}{mtk.168}
\bookmark[dest={mtk.168},level=2]{8. อฏฺฐมสิกฺขาปท}
\csromanheader{2}{2. \textbf{อนฺธการวคฺค 8. อฏฺฐมสิกฺขาปท}}
\proseitem{869}{\csromanfolio{362}เตน สมเยน พุทฺโธ ภควา สาวตฺถิยํ วิหรติ เชตวเน อนาถปิณฺฑิกสฺส อาราเม.\csromanspacer{} เตน โข ปน สมเยน ภทฺทาย กาปิลานิยา อนฺเตวาสินี ภิกฺขุนี ภทฺทํ กาปิลานิํ สกฺกจฺจํ อุปฏฺเฐติ.\csromanspacer{} ภทฺทา กาปิลานี ภิกฺขุนิโย เอตทโวจ “อยํ มํ อยฺเย ภิกฺขุนี สกฺกจฺจํ อุปฏฺเฐติ, อิมิสฺสาหํ จีวรํ ทสฺสามี”ติ.\csromanspacer{} อถ โข สา ภิกฺขุนี ทุคฺคหิเตน ทูปธาริเตน ปรํ อุชฺฌาเปสิ “อหํ กิรายฺเย อยฺยํ น สกฺกจฺจํ อุปฏฺเฐมิ, น กิร เม อยฺยา จีวรํ ทสฺสตี”ติ.\csromanspacer{} ยา ตา ภิกฺขุนิโย อปฺปิจฺฉา ฯเปฯ ตา อุชฺฌายนฺติ ขิยฺยนฺติ วิปาเจนฺติ “กถํ หิ นาม ภิกฺขุนี ทุคฺคหิเตน ทูปธาริเตน ปรํ อุชฺฌาเปสฺสตี”ติ ฯเปฯ.\csromanspacer{} สจฺจํ กิร ภิกฺขเว ภิกฺขุนี ทุคฺคหิเตน ทูปธาริเตน ปรํ อุชฺฌาเปตีติ\footnote{อุชฺฌาเปสีติ (ก)}.\csromanspacer{} สจฺจํ ภควาติ.\csromanspacer{} วิครหิ พุทฺโธ ภควา ฯเปฯ.\csromanspacer{} กถํ หิ นาม ภิกฺขเว ภิกฺขุนี ทุคฺคหิเตน ทูปธาริเตน ปรํ อุชฺฌาเปสฺสติ.\csromanspacer{} เนตํ ภิกฺขเว อปฺปสนฺนานํ วา ปสาทาย ฯเปฯ.\csromanspacer{} เอวญฺจ ปน ภิกฺขเว ภิกฺขุนิโย อิมํ สิกฺขาปทํ อุทฺทิสนฺตุ\csromandash{}}

\proseitem{870}{\textbf{“ยา ปน ภิกฺขุนี ทุคฺคหิเตน ทูปธาริเตน ปรํ อุชฺฌาเปยฺย, ปาจิตฺติยนฺ”}ติ.}

\proseitem{871}{\textbf{ยา ปนา}ติ ยา ยาทิสา ฯเปฯ.\csromanspacer{} \textbf{ภิกฺขุนี}ติ ฯเปฯ อยํ อิมสฺมิํ อตฺเถ อธิปฺเปตา ภิกฺขุนีติ.}

\prose{\textbf{ทุคฺคหิเตนา}ติ อญฺญถา คหิเตน\footnote{อุคฺคหิเตน (ก)}.}

\prose{\textbf{ทูปธาริเตนา}ติ อญฺญถา อุปธาริเตน.}

\prose{\textbf{ปรนฺ}ติ อุปสมฺปนฺนํ อุชฺฌาเปติ, อาปตฺติ ปาจิตฺติยสฺส.}

\proseitem{872}{อุปสมฺปนฺนาย อุปสมฺปนฺนสญฺญา อุชฺฌาเปติ, อาปตฺติ ปาจิตฺติยสฺส.\csromanspacer{} อุปสมฺปนฺนาย เวมติกา อุชฺฌาเปติ, อาปตฺติ ปาจิตฺติยสฺส.\csromanspacer{} อุปสมฺปนฺนาย อนุปสมฺปนฺนสญฺญา อุชฺฌาเปติ, อาปตฺติ ปาจิตฺติยสฺส.}

\prose{\csromanfolio{363}อนุปสมฺปนฺนํ อุชฺฌาเปติ, อาปตฺติ ทุกฺกฏสฺส.\csromanspacer{} อนุปสมฺปนฺนาย อุปสมฺปนฺนสญฺญา, อาปตฺติ ทุกฺกฏสฺส.\csromanspacer{} อนุปสมฺปนฺนาย เวมติกา, อาปตฺติ ทุกฺกฏสฺส.\csromanspacer{} อนุปสมฺปนฺนาย อนุปสมฺปนฺนสญฺญา, อาปตฺติ ทุกฺกฏสฺส.}

\proseitemwithcloserruled{873}{อนาปตฺติ อุมฺมตฺติกาย อาทิกมฺมิกายาติ.}{อฏฺฐมสิกฺขาปทํ นิฏฺฐิตํ.}
\csromanmatikaanchor{mtk.169}
\csromantocmarkref{subsection}{9. นวมสิกฺขาปท}{mtk.169}
\bookmark[dest={mtk.169},level=2]{9. นวมสิกฺขาปท}
\csromanheader{2}{2. \textbf{อนฺธการวคฺค 9. นวมสิกฺขาปท}}
\proseitem{874}{เตน สมเยน พุทฺโธ ภควา สาวตฺถิยํ วิหรติ เชตวเน อนาถปิณฺฑิกสฺส อาราเม.\csromanspacer{} เตน โข ปน สมเยน ภิกฺขุนิโย อตฺตโน ภณฺฑกํ อปสฺสนฺติโย จณฺฑกาฬิํ ภิกฺขุนิํ เอตทโวจุํ “อปายฺเย อมฺหากํ ภณฺฑกํ ปสฺเสยฺยาสี”ติ.\csromanspacer{} จณฺฑกาฬี ภิกฺขุนี อุชฺฌายติ ขิยฺยติ วิปาเจติ “อหเมว นูน โจรี, อหเมว นูน อลชฺชินี, ยา อยฺยาโย อตฺตโน ภณฺฑกํ อปสฺสนฺติโย, ตา มํ เอวมาหํสุ ‘อปายฺเย อมฺหากํ ภณฺฑกํ ปสฺเสยฺยาสี’ติ.\csromanspacer{} สจาหํ อยฺเย ตุมฺหากํ ภณฺฑกํ คณฺหามิ, อสฺสมณี โหมิ, พฺรหฺมจริยา จวามิ, นิรยํ อุปปชฺชามิ, ยา ปน มํ อภูเตน เอวมาห, สาปิ อสฺสมณี โหตุ, พฺราหฺมจริยา จวตุ, นิรยํ อุปปชฺชตู”ติ.\csromanspacer{} ยา ตา ภิกฺขุนิโย อปฺปิจฺฉา ฯเปฯ ตา อุชฺฌายนฺติ ขิยฺยนฺติ วิปาเจนฺติ “กถํ หิ นาม อยฺยา จณฺฑกาฬี อตฺตานมฺปิ ปรมฺปิ นิรเยนปิ พฺรหฺมจริเยนปิ อภิสปิสฺสตี”ติ ฯเปฯ.\csromanspacer{} สจฺจํ กิร ภิกฺขเว จณฺฑกาฬี ภิกฺขุนี อตฺตานมฺปิ ปรมฺปิ นิรเยนปิ พฺรหฺมจริเยนปิ อภิสปตีติ.\csromanspacer{} สจฺจํ ภควาติ.\csromanspacer{} วิครหิ พุทฺโธ ภควา ฯเปฯ.\csromanspacer{} กถํ หิ นาม ภิกฺขเว จณฺฑกาฬี ภิกฺขุนี อตฺตานมฺปิ ปรมฺปิ นิรเยนปิ พฺรหฺมจริเยนปิ อภิสปิสฺสติ.\csromanspacer{} เนตํ ภิกฺขเว อปฺปสนฺนานํ วา ปสาทาย ฯเปฯ.\csromanspacer{} เอวญฺจ ปน ภิกฺขเว ภิกฺขุนิโย อิมํ สิกฺขาปทํ อุทฺทิสนฺตุ\csromandash{}}

\proseitem{875}{\textbf{“ยา ปน ภิกฺขุนี อตฺตานํ วา ปรํ วา นิรเยน วา พฺรหฺมจริเยน วา อภิสเปยฺย, ปาจิตฺติยนฺ”}ติ.}

\proseitem{876}{\textbf{ยา ปนา}ติ ยา ยาทิสา ฯเปฯ.\csromanspacer{} \textbf{ภิกฺขุนี}ติ ฯเปฯ อยํ อิมสฺมิํ อตฺเถ อธิปฺเปตา ภิกฺขุนีติ.}

\prose{\csromanfolio{364}\textbf{อตฺตานนฺ}ติ ปจฺจตฺตํ.\csromanspacer{} \textbf{ปรนฺ}ติ อุปสมฺปนฺนํ.\csromanspacer{} นิรเยน วา พฺราหฺมจริเยน วา อภิสปติ, อาปตฺติ ปาจิตฺติยสฺส.}

\proseitem{877}{อุปสมฺปนฺนาย อุปสมฺปนฺนสญฺญา นิรเยน วา พฺรหฺมจริเยน วา อภิสปติ, อาปตฺติ ปาจิตฺติยสฺส.\csromanspacer{} อุปสมฺปนฺนาย เวมติกา นิรเยน วา พฺรหฺมจริเยน วา อภิสปติ, อาปตฺติ ปาจิตฺติยสฺส.\csromanspacer{} อุปสมฺปนฺนาย อนุปสมฺปนฺนสญฺญา นิรเยน วา พฺรหฺมจริเยน วา อภิสปติ, อาปตฺติ}

\prose{ติรจฺฉานโยนิยา วา เปตฺติวิสเยน วา มนุสฺสโทภคฺเคน วา อภิสปติ, อาปตฺติ ทุกฺกฏสฺส.\csromanspacer{} อนุปสมฺปนฺนํ อภิสปติ, อาปตฺติ ทุกฺกฏสฺส.\csromanspacer{} อนุปสมฺปนฺนาย อุปสมฺปนฺนสญฺญา, อาปตฺติ ทุกฺกฏสฺส.\csromanspacer{} อนุปสมฺปนฺนาย เวมติกา, อาปตฺติ ทุกฺกฏสฺส.\csromanspacer{} อนุปสมฺปนฺนาย อนุปสมฺปนฺนสญฺญา, อาปตฺติ ทุกฺกฏสฺส.}

\proseitemwithcloserruled{878}{อนาปตฺติ อตฺถปุเรกฺขาราย ธมฺมปุเรกฺขาราย อนุสาสนิ\-ปุเรกฺขาราย อุมฺมตฺติกาย อาทิกมฺมิกายาติ.}{นวมสิกฺขาปทํ นิฏฺฐิตํ.}
\csromanmatikaanchor{mtk.170}
\csromantocmarkref{subsection}{10. ทสมสิกฺขาปท}{mtk.170}
\bookmark[dest={mtk.170},level=2]{10. ทสมสิกฺขาปท}
\csromanheader{2}{2. \textbf{อนฺธการวคฺคฺ}อ 10. ทสมสิกฺขาปท}
\proseitem{879}{เตน สมเยน พุทฺโธ ภควา สาวตฺถิยํ วิหรติ เชตวเน อนาถปิณฺฑิกสฺส อาราเม.\csromanspacer{} เตน โข ปน สมเยน จณฺฑกาฬี ภิกฺขุนี ภิกฺขุนีหิ สทฺธิํ ภณฺฑิตฺวา อตฺตานํ วธิตฺวา วธิตฺวา โรทติ.\csromanspacer{} ยา ตา ภิกฺขุนิโย อปฺปิจฺฉา ฯเปฯ ตา อุชฺฌายนฺติ ขิยฺยนฺติ วิปาเจนฺติ “กถํ หิ นาม อยฺยา จณฺฑกาฬี อตฺตานํ วธิตฺวา วธิตฺวา โรทิสฺสตี”ติ ฯเปฯ.\csromanspacer{} สจฺจํ กิร ภิกฺขเว จณฺฑกาฬี ภิกฺขุนี อตฺตานํ วธิตฺวา วธิตฺวา โรทตีติ.\csromanspacer{} สจฺจํ ภควาติ.\csromanspacer{} วิครหิ พุทฺโธ ภควา ฯเปฯ.\csromanspacer{} กถํ หิ นาม ภิกฺขเว จณฺฑกาฬี ภิกฺขุนี อตฺตานํ วธิตฺวา วธิตฺวา โรทิสฺสติ.\csromanspacer{} เนตํ ภิกฺขเว อปฺปสนฺนานํ วา ปสาทาย ฯเปฯ.\csromanspacer{} เอวญฺจ ปน ภิกฺขเว ภิกฺขุนิโย อิมํ สิกฺขาปทํ อุทฺทิสนฺตุ\csromandash{}}

\proseitem{880}{\csromanfolio{365}\textbf{“ยา ปน ภิกฺขุนี อตฺตานํ วธิตฺวา วธิตฺวา โรเทยฺย, ปาจิตฺติยนฺ”}ติ.}

\proseitem{881}{\textbf{ยา ปนา}ติ ยา ยาทิสา ฯเปฯ.\csromanspacer{} \textbf{ภิกฺขุนี}ติ ฯเปฯ อยํ อิมสฺมิํ อตฺเถ อธิปฺเปตา ภิกฺขุนีติ.}

\prose{\textbf{อตฺตานนฺ}ติ ปจฺจตฺตํ.\csromanspacer{} วธิตฺวา วธิตฺวา โรทติ, อาปตฺติ ปาจิตฺติยสฺส.\csromanspacer{} วธติ น โรทติ, อาปตฺติ ทุกฺกฏสฺส.\csromanspacer{} โรทติ น วธติ, อาปตฺติ ทุกฺกฏสฺส.}

\proseitemwithcloser{882}{อนาปตฺติ ญาติพฺยสเนน วา โภคพฺยสเนน วา โรคพฺยสเนน วา ผุฏฺฐา โรทติ น วธติ, อุมฺมตฺติกาย อาทิกมฺมิกายาติ.}{ทสมสิกฺขาปทํ นิฏฺฐิตํ.}
\csromanheader{2}{อนฺธการวคฺโค ทุติโย.}
\csromansectionrule
\csromanmatikaanchor{mtk.171}
\csromanmatikaanchor{mtk.172}
\csromantocmarknopage{section}{3. นคฺควคฺค}{mtk.171}
\bookmark[dest={mtk.171},level=1]{3. นคฺควคฺค}
\csromantocmarkref{subsection}{1. ปฐมสิกฺขาปท}{mtk.172}
\bookmark[dest={mtk.172},level=2]{1. ปฐมสิกฺขาปท}
\setcsromanheadh{3. นคฺควคฺค}
\csromanheader{2}{3. \textbf{นคฺควคฺค 1. ปฐมสิกฺขาปท}}
\proseitem{883}{เตน สมเยน พุทฺโธ ภควา สาวตฺถิยํ วิหรติ เชตวเน อนาถปิณฺฑิกสฺส อาราเม.\csromanspacer{} เตน โข ปน สมเยน สมฺพหุลา\csromansymbolfootnote{*}{วิ 3. 407 ปิฏฺเฐปิ.}ภิกฺขุนิโย อจิรวติยา นทิยา เวสิยาหิ สทฺธิํ นคฺคา เอกติตฺเถ นหายนฺติ.\csromanspacer{} เวสิยา ตา ภิกฺขุนิโย อุปฺปณฺเฑสุํ “กิํ นุ โข นาม ตุมฺหากํ อยฺเย ทหรานํ\footnote{ทหรานํ ทหารานํ (สี)} พฺรหฺมจริยํ จิณฺเณน, นนุ นาม กามา ปริภุญฺชิตพฺพา, ยทา ชิณฺณา ภวิสฺสถ, ตทา พฺรหฺมจริยํ จริสฺสถ, เอวํ ตุมฺหากํ อุโภ อตฺถา ปริคฺคหิตา ภวิสฺสนฺตี”ติ.\csromanspacer{} ภิกฺขุนิโย เวสิยาหิ อุปฺปณฺฑิยมานา มงฺกู อเหสุํ.\csromanspacer{} อถ โข ตา ภิกฺขุนิโย อุปสฺสยํ คนฺตฺวา ภิกฺขุนีนํ เอตมตฺถํ อาโรเจสุํ.\csromanspacer{} ภิกฺขุนิโย ภิกฺขูนํ เอตมตฺถํ อาโรเจสุํ.\csromanspacer{} ภิกฺขู ภควโต เอตมตฺถํ อาโรเจสุํ.\csromanspacer{} อถ โข ภควา เอตสฺมิํ นิทาเน เอตสฺมิํ ปกรเณ ธมฺมิํ กถํ กตฺวา ภิกฺขู อามนฺเตสิ เตน หิ ภิกฺขเว ภิกฺขุนีนํ สิกฺขาปทํ ปญฺญเปสฺสามิ ทส อตฺถวเส ปฏิจฺจ สํฆสุฏฺฐุตาย ฯเปฯ วินยานุคฺคหาย.\csromanspacer{} เอวญฺจ ปน ภิกฺขเว ภิกฺขุนิโย อิมํ สิกฺขาปทํ อุทฺทิสนฺตุ\csromandash{}}

\proseitem{884}{\csromanfolio{366}\textbf{“ยา ปน ภิกฺขุนี นคฺคา นหาเยยฺย, ปาจิตฺติยนฺ”}ติ.}

\proseitem{885}{\textbf{ยา ปนา}ติ ยา ยาทิสา ฯเปฯ.\csromanspacer{} \textbf{ภิกฺขุนี}ติ ฯเปฯ อยํ อิมสฺมิํ อตฺเถ อธิปฺเปตา ภิกฺขุนีติ.}

\prose{\textbf{นคฺคา นหาเยยฺยา}ติ อนิวตฺถา วา อปารุตา วา นหายติ, ปโยเค ทุกฺกฏํ.\csromanspacer{} นหานปริโยสาเน อาปตฺติ ปาจิตฺติยสฺส.}

\proseitemwithcloserruled{886}{อนาปตฺติ อจฺฉินฺนจีวริกาย วา, นฏฺฐจีวริกาย วา, อาปทาสุ อุมฺมตฺติกาย อาทิกมฺมิกายาติ.}{ปฐมสิกฺขาปทํ นิฏฺฐิตํ.}
\csromanmatikaanchor{mtk.173}
\csromantocmarkref{subsection}{2. ทุติยสิกฺขาปท}{mtk.173}
\bookmark[dest={mtk.173},level=2]{2. ทุติยสิกฺขาปท}
\csromanheader{2}{3. \textbf{นคฺควคฺค 2. ทุติยสิกฺขาปท}}
\proseitem{887}{เตน สมเยน พุทฺโธ ภควา สาวตฺถิยํ วิหรติ เชตวเน อนาถปิณฺฑิกสฺส อาราเม.\csromanspacer{} เตน โข ปน สมเยน ภควตา ภิกฺขุนีนํ\csromansymbolfootnote{*}{วิ 3. 408 ปิฏฺเฐ.}อุทกสาฏิกา อนุญฺญาตา โหติ.\csromanspacer{} ฉพฺพคฺคิยา ภิกฺขุนิโย “ภควตา อุทกสาฏิกา อนุญฺญาตา”ติ อปฺปมาณิกาโย อุทกสาฏิกาโย ธาเรสุํ\footnote{ธาเรนฺติ (221, 222, 224 ปิฏฺเฐสุ)}, ปุรโตปิ ปจฺฉโตปิ อากฑฺฒนฺตา อาหิณฺฑนฺติ.\csromanspacer{} ยา ตา ภิกฺขุนิโย อปฺปิจฺฉา ฯเปฯ ตา อุชฺฌายนฺติ ขิยฺยนฺติ วิปาเจนฺติ “กถํ หิ นาม ฉพฺพคฺคิยา ภิกฺขุนิโย อปฺปมาณิกาโย อุทกสาฏิกาโย ธาเรสฺสนฺตี”ติ ฯเปฯ.\csromanspacer{} สจฺจํ กิร ภิกฺขเว ฉพฺพคฺคิยา ภิกฺขุนิโย อปฺปมาณิกาโย อุทกสาฏิกาโย ธาเรนฺตีติ.\csromanspacer{} สจฺจํ ภควาติ.\csromanspacer{} วิครหิ พุทฺโธ ภควา ฯเปฯ.\csromanspacer{} กถํ หิ นาม ภิกฺขเว ฉพฺพคฺคิยา ภิกฺขุนิโย อปฺปมาณิกาโย อุทกสาฏิกาโย ธาเรสฺสนฺติ.\csromanspacer{} เนตํ ภิกฺขเว อปฺปสนฺนานํ วา ปสาทาย ฯเปฯ.\csromanspacer{} เอวญฺจ ปน ภิกฺขเว ภิกฺขุนิโย อิมํ สิกฺขาปทํ อุทฺทิสนฺตุ\csromandash{}}

\proseitem{888}{\textbf{“อุทกสาฏิกํ ปน ภิกฺขุนิยา การยมานาย ปมาณิกา กาเรตพฺพา, ตตฺริทํ ปมาณํ, ทีฆโส จตสฺโส วิทตฺถิโย สุคตวิทตฺถิยา, ติริยํ เทฺว วิทตฺถิโย, ตํ อติกฺกาเมนฺติยา เฉทนกํ ปาจิตฺติยนฺ”}ติ.}

\proseitem{889}{\csromanfolio{367}\textbf{อุทกสาฏิกา} นาม ยาย นิวตฺถา\footnote{นิวตฺถาย (สฺยา)} นหายติ.}

\prose{\textbf{การยมานายา}ติ กโรนฺติยา วา การาเปนฺติยา วา.}

\prose{ปมาณิกา กาเรตพฺพา, ตตฺริทํ ปมาณํ, ทีฆโส จตสฺโส วิทตฺถิโย สุคตวิทตฺถิยา, ติริยํ เทฺว วิทตฺถิโย.\csromanspacer{} ตํ อติกฺกาเมตฺวา กโรติ วา การาเปติ วา, ปโยเค ทุกฺกฏํ.\csromanspacer{} ปฏิลาเภน ฉินฺทิตฺวา ปาจิตฺติยํ เทเสตพฺพํ.}

\proseitem{890}{อตฺตนา วิปฺปกตํ อตฺตนา ปริโยสาเปติ, อาปตฺติ ปาจิตฺติยสฺส.\csromanspacer{} อตฺตนา วิปฺปกตํ ปเรหิ ปริโยสาเปติ, อาปตฺติ ปาจิตฺติยสฺส.\csromanspacer{} ปเรหิ วิปฺปกตํ อตฺตนา ปริโยสาเปติ, อาปตฺติ ปาจิตฺติยสฺส.\csromanspacer{} ปเรหิ วิปฺปกตํ ปเรหิ ปริโยสาเปติ, อาปตฺติ ปาจิตฺติยสฺส.}

\prose{อญฺญสฺสตฺถาย กโรติ วา การาเปติ วา, อาปตฺติ ทุกฺกฏสฺส.\csromanspacer{} อญฺเญน กตํ ปฏิลภิตฺวา ปริภุญฺชติ, อาปตฺติ ทุกฺกฏสฺส.}

\proseitemwithcloserruled{891}{อนาปตฺติ ปมาณิกํ กโรติ, อูนกํ กโรติ, อญฺเญน กตํ ปมาณาติกฺกนฺตํ ปฏิลภิตฺวา ฉินฺทิตฺวา ปริภุญฺชติ, วิตานํ วา ภูมตฺถรณํ วา สาณิปาการํ วา ภิสิํ วา พิพฺโพหนํ วา กโรติ, อุมฺมตฺติกาย อาทิกมฺมิกายาติ.}{ทุติยสิกฺขาปทํ นิฏฺฐิตํ.}
\csromanmatikaanchor{mtk.174}
\csromantocmarkref{subsection}{3. ตติยสิกฺขาปท}{mtk.174}
\bookmark[dest={mtk.174},level=2]{3. ตติยสิกฺขาปท}
\csromanheader{2}{3. \textbf{นคฺควคฺค 3. ตติยสิกฺขาปท}}
\proseitem{892}{เตน สมเยน พุทฺโธ ภควา สาวตฺถิยํ วิหรติ เชตวเน อนาถปิณฺฑิกสฺส อาราเม.\csromanspacer{} เตน โข ปน สมเยน อญฺญตริสฺสา ภิกฺขุนิยา มหคฺเฆ จีวรทุสฺเส จีวรํ ทุกฺกฏํ โหติ ทุสฺสิพฺพิตํ.\csromanspacer{} ถุลฺลนนฺทา ภิกฺขุนี ตํ ภิกฺขุนิํ เอตทโวจ “สุนฺทรํ โข อิทํ เต อยฺเย จีวรทุสฺสํ, จีวรญฺจ โข ทุกฺกฏํ ทุสฺสิพฺพิตนฺ”ติ.\csromanspacer{} วิสิพฺเพมิ อยฺเย, สิพฺพิสฺสสีติ.\csromanspacer{} อามายฺเย สิพฺพิสฺสามีติ.\csromanspacer{} อถ โข สา ภิกฺขุนี ตํ จีวรํ วิสิพฺเพตฺวา ถุลฺลนนฺทาย \csromanfolio{368}ภิกฺขุนิยา อทาสิ.\csromanspacer{} ถุลฺลนนฺทา ภิกฺขุนี “สิพฺพิสฺสามิ สิพฺพิสฺสามี”ติ เนว สิพฺพติ, น สิพฺพาปนาย อุสฺสุกฺกํ กโรติ.\csromanspacer{} อถ โข สา ภิกฺขุนี ภิกฺขุนีนํ เอตมตฺถํ อาโรเจสิ.\csromanspacer{} ยา ตา ภิกฺขุนิโย อปฺปิจฺฉา ฯเปฯ ตา อุชฺฌายนฺติ ขิยฺยนฺติ วิปาเจนฺติ “กถํ หิ นาม อยฺยา ถุลฺลนนฺทา ภิกฺขุนิยา จีวรํ วิสิพฺพาเปตฺวา เนว สิพฺพิสฺสติ น สิพฺพาปนาย อุสฺสุกํ กริสฺสตี”ติ ฯเปฯ.\csromanspacer{} สจฺจํ กิร ภิกฺขเว ถุลฺลนนฺทา ภิกฺขุนี ภิกฺขุนิยา จีวรํ วิสิพฺพาเปตฺวา เนว สิพฺพติ, น สิพฺพาปนาย อุสฺสุกฺกํ กโรตีติ.\csromanspacer{} สจฺจํ ภควาติ.\csromanspacer{} วิครหิ พุทฺโธ ภควา ฯเปฯ.\csromanspacer{} กถํ หิ นาม ภิกฺขเว ถุลฺลนนฺทา ภิกฺขุนี ภิกฺขุนิยา จีวรํ วิสิพฺพาเปตฺวา เนว สิพฺพิสฺสติ, น สิพฺพาปนาย อุสฺสุกํ กริสฺสติ.\csromanspacer{} เนตํ ภิกฺขเว อปฺปสนฺนานํ วา ปสาทาย ฯเปฯ.\csromanspacer{} เอวญฺจ ปน ภิกฺขเว ภิกฺขุนิโย อิมํ สิกฺขาปทํ อุทฺทิสนฺตุ\csromandash{}}

\proseitem{893}{\textbf{“ยา ปน ภิกฺขุนี ภิกฺขุนิยา จีวรํ วิสิพฺเพตฺวา วา วิสิพฺพาเปตฺวา วา สา ปจฺฉา อนนฺตรายิกินี เนว สิพฺเพยฺย, น สิพฺพาปนาย อุสฺสุกฺกํ กเรยฺย อญฺญตฺร จตูหปญฺจาหา, ปาจิตฺติยนฺ”}ติ.}

\proseitem{894}{\textbf{ยา ปนา}ติ ยา ยาทิสา ฯเปฯ.\csromanspacer{} \textbf{ภิกฺขุนี}ติ ฯเปฯ อยํ อิมสฺมิํ อตฺเถ อธิปฺเปตา ภิกฺขุนีติ.}

\prose{\textbf{ภิกฺขุนิยา}ติ อญฺญาย ภิกฺขุนิยา.}

\prose{\textbf{จีวรํ} นาม ฉนฺนํ จีวรานํ อญฺญตรํ จีวรํ.}

\prose{\textbf{วิสิพฺเพตฺวา}ติ สยํ วิสิพฺเพตฺวา.\csromanspacer{} \textbf{วิสิพฺพาเปตฺวา}ติ อญฺญํ วิสิพฺพาเปตฺวา.}

\prose{\textbf{สา ปจฺฉา อนนฺตรายิกินี}ติ อสติ อนฺตราเย.}

\prose{\textbf{เนว สิพฺเพยฺยา}ติ น สยํ สิพฺเพยฺย.\csromanspacer{} \textbf{น สิพฺพาปนาย อุสฺสุกฺกํ กเรยฺยา}ติ น อญฺญํ อาณาเปยฺย.}

\prose{\textbf{อญฺญตฺร จตูหปญฺจาหา}ติ ฐเปตฺวา จตูหปญฺจาหํ. “เนว สิพฺพิสฺสามิ น สิพฺพาปนาย อุสฺสุกฺกํ กริสฺสามี”ติ ธุรํ นิกฺขิตฺตมตฺเต อาปตฺติ}

\proseitem{895}{อุปสมฺปนฺนาย อุปสมฺปนฺนสญฺญา จีวรํ วิสิพฺเพตฺวา วา วิสิพฺพาเปตฺวา วา สา ปจฺฉา อนนฺตรายิกินี เนว สิพฺพติ, น สิพฺพาปนาย อุสฺสุกฺกํ กโรติ อญฺญตฺร จตูหปญฺจาหา, อาปตฺติ ปาจิตฺติยสฺส.\csromanspacer{} อุปสมฺปนฺนาย \csromanfolio{369}เวมติกา จีวรํ วิสิพฺเพตฺวา วา วิสิพฺพาเปตฺวา วา สา ปจฺฉา อนนฺตรายิกินี เนว สิพฺพติ, น สิพฺพาปนาย อุสฺสุกฺกํ กโรติ อญฺญตฺร จตูหปญฺจาหา, อาปฺปติ ปาจิตฺติยสฺส.\csromanspacer{} อุปสมฺปนฺนาย อนุปสมฺปนฺนสญฺญา จีวรํ วิสิพฺเพตฺวา วา วิสิพฺพาเปตฺวา วา สา ปจฺฉา อนนฺตรายิกินี เนว สิพฺพติ, น สิพฺพาปนาย อุสฺสุกฺกํ กโรติ อญฺญตฺร จตูหปญฺจาหา, อาปตฺติ ปาจิตฺติยสฺส.}

\prose{อญฺญํ ปริกฺขารํ วิสิพฺเพตฺวา วา วิสิพฺพาเปตฺวา วา สา ปจฺฉา อนนฺตรายิกินี เนว สิพฺพติ, น สิพฺพาปนาย อุสฺสุกฺกํ กโรติ อญฺญตฺร จตูหปญฺจาหา, อาปตฺติ ทุกฺกฏสฺส.\csromanspacer{} อนุปสมฺปนฺนาย จีวรํ วา อญฺญํ วา ปริกฺขารํ วิสิพฺเพตฺวา วา วิสิพฺพาเปตฺวา วา สา ปจฺฉา อนนฺตรายิกินี เนว สิพฺพติ, น สิพฺพาปนาย อุสฺสุกฺกํ กโรติ อญฺญตฺร จตูหปญฺจาหา, อาปตฺติ ทุกฺกฏสฺส.\csromanspacer{} อนุปสมฺปนฺนาย อุปสมฺปนฺนสญฺญา, อาปตฺติ ทุกฺกฏสฺส.\csromanspacer{} อนุปสมฺปนฺนาย เวมติกา, อาปตฺติ ทุกฺกฏสฺส.\csromanspacer{} อนุปสมฺปนฺนาย อนุปสมฺปนฺนสญฺญา, อาปตฺติ ทุกฺกฏสฺส.}

\proseitemwithcloserruled{896}{อนาปตฺติ สติ อนฺตราเย, ปริเยสิตฺวา น ลภติ, กโรนฺตี จตูหปญฺจาหํ อติกฺกาเมติ, คิลานาย อาปทาสุ อุมฺมตฺติกาย อาทิกมฺมิกายาติ.}{ตติยสิกฺขาปทํ นิฏฺฐิตํ.}
\csromanmatikaanchor{mtk.175}
\csromantocmarkref{subsection}{4. จตุตฺถสิกฺขาปท}{mtk.175}
\bookmark[dest={mtk.175},level=2]{4. จตุตฺถสิกฺขาปท}
\csromanheader{2}{3. \textbf{นคฺควคฺค 4. จตุตฺถสิกฺขาปท}}
\proseitem{897}{เตน สมเยน พุทฺโธ ภควา สาวตฺถิยํ วิหรติ เชตวเน อนาถปิณฺฑิกสฺส อาราเม.\csromanspacer{} เตน โข ปน สมเยน ภิกฺขุนิโย ภิกฺขุนีนํ หตฺเถ จีวรํ นิกฺขิปิตฺวา สนฺตรุตฺตเรน ชนปทจาริกํ ปกฺกมนฺติ.\csromanspacer{} ตานิ จีวรานิ จิรํ นิกฺขิตฺตานิ กณฺณกิตานิ โหนฺติ.\csromanspacer{} ตานิ ภิกฺขุนิโย โอตาเปนฺติ.\csromanspacer{} ภิกฺขุนิโย ตา ภิกฺขุนิโย เอตทโวจุํ “กสฺสิมานิ อยฺเย จีวรานิ กณฺณกิตานี”ติ.\csromanspacer{} อถ โข ตา ภิกฺขุนิโย ภิกฺขุนีนํ เอตมตฺถํ อาโรเจสุํ.\csromanspacer{} ยา ตา ภิกฺขุนิโย อปฺปิจฺฉา ฯเปฯ ตา อุชฺฌายนฺติ ขิยฺยนฺติ วิปาเจนฺติ “กถํ หิ นาม ภิกฺขุนิโย ภิกฺขุนีนํ หตฺเถ จีวรํ นิกฺขิปิตฺวา สนฺตรุตฺตเรน ชนปทจาริกํ ปกฺกมิสฺสนฺตี”ติ ฯเปฯ.\csromanspacer{} สจฺจํ กิร ภิกฺขเว ภิกฺขุนิโย ภิกฺขุนีนํ หตฺเถ จีวรํ นิกฺขิปิตฺวา สนฺตรุตฺตเรน \csromanfolio{370}ชนปทจาริกํ ปกฺกมนฺตีติ.\csromanspacer{} สจฺจํ ภควาติ.\csromanspacer{} วิครหิ พุทฺโธ ภควา ฯเปฯ.\csromanspacer{} กถํ หิ นาม ภิกฺขเว ภิกฺขุนิโย ภิกฺขุนีนํ หตฺเถ จีวรํ นิกฺขิปิตฺวา สนฺตรุตฺตเรน ชนปทจาริกํ ปกฺกมิสฺสนฺติ.\csromanspacer{} เนตํ ภิกฺขเว อปฺปสนฺนานํ วา ปสาทาย ฯเปฯ.\csromanspacer{} เอวญฺจ ปน ภิกฺขเว ภิกฺขุนิโย อิมํ สิกฺขาปทํ อุทฺทิสนฺตุ\csromandash{}}

\proseitem{898}{\textbf{“ยา ปน ภิกฺขุนี ปญฺจาหิกํ สงฺฆาฏิจารํ}\footnote{สงฺฆาฏิวารํ (สฺยา)} \textbf{อติกฺกาเมยฺย, ปาจิตฺติยนฺ”}ติ.}

\proseitem{899}{\textbf{ยา ปนา}ติ ยา ยาทิสา ฯเปฯ.\csromanspacer{} \textbf{ภิกฺขุนี}ติ ฯเปฯ อยํ อิมสฺมิํ อตฺเถ อธิปฺเปตา ภิกฺขุนีติ.}

\prose{\textbf{ปญฺจาหิกํ สงฺฆาฏิจารํ อติกฺกาเมยฺยา}ติ ปญฺจมํ ทิวสํ ปญฺจ จีวรานิ เนว นิวาเสติ น ปารุปติ น โอตาเปติ ปญฺจมํ ทิวสํ อติกฺกาเมติ, อาปตฺติ ปาจิตฺติยสฺส.}

\proseitem{900}{ปญฺจาหาติกฺกนฺเต อติกฺกนฺตสญฺญา, อาปตฺติ ปาจิตฺติยสฺส.\csromanspacer{} ปญฺจาหาติกฺกนฺเต เวมติกา, อาปตฺติ ปาจิตฺติยสฺส.\csromanspacer{} ปญฺจาหาติกฺกนฺเต อนติกฺกนฺตสญฺญา, อาปตฺติ ปาจิตฺติยสฺส.}

\prose{ปญฺจาหานติกฺกนฺเต อติกฺกนฺตสญฺญา, อาปตฺติ ทุกฺกฏสฺส.\csromanspacer{} ปญฺจาหานติกฺกนฺเต เวมติกา, อาปตฺติ ทุกฺกฏสฺส.\csromanspacer{} ปญฺจาหานติกฺกนฺเต อนติกฺกนฺตสญฺญา, อนาปตฺติ.}

\proseitemwithcloserruled{901}{อนาปตฺติ ปญฺจมํ ทิวสํ ปญฺจ จีวรานิ นิวาเสติ วา ปารุปติ วา โอตาเปติ วา, คิลานาย อาปทาสุ อุมฺมตฺติกาย อาทิกมฺมิกายาติ.}{จตุตฺถสิกฺขาปทํ นิฏฺฐิตํ.}
\csromanmatikaanchor{mtk.176}
\csromantocmarkref{subsection}{5. ปญฺจมสิกฺขาปท}{mtk.176}
\bookmark[dest={mtk.176},level=2]{5. ปญฺจมสิกฺขาปท}
\csromanheader{2}{3. \textbf{นคฺควคฺค 5. ปญฺจมสิกฺขาปท}}
\proseitem{902}{เตน สมเยน พุทฺโธ ภควา สาวตฺถิยํ วิหรติ เชตวเน อนาถปิณฺฑิกสฺส อาราเม.\csromanspacer{} เตน โข ปน สมเยน อญฺญตรา ภิกฺขุนี ปิณฺฑาย จริตฺวา อลฺลจีวรํ ปตฺถริตฺวา วิหารํ ปาวิสิ.\csromanspacer{} อญฺญตรา ภิกฺขุนี \csromanfolio{371}\textbf{ตํ จีวรํ ปารุปิตฺวา คามํ ปิณฺฑาย ปาวิสิ.}\csromanspacer{}\textbf{ สา นิกฺขมิตฺวา ภิกฺขุนิโย} \textbf{ปุจฺฉิ “อปายฺเย มยฺหํ จีวรํ ปสฺเสยฺยาถา”ติ.}\csromanspacer{}\textbf{ ภิกฺขุนิโย ตสฺสา} ภิกฺขุนิยา เอตมตฺถํ อาโรเจสุํ.\csromanspacer{} อถ โข สา ภิกฺขุนี อุชฺฌายติ ขิยฺยนฺติ วิปาเจติ “กถํ หิ นาม ภิกฺขุนี มยฺหํ จีวรํ อนาปุจฺฉา ปารุปิสฺสตี”ติ.\csromanspacer{} อถ โข สา ภิกฺขุนี ภิกฺขุนีนํ เอตมตฺถํ อาโรเจสิ.\csromanspacer{} ยา ตา ภิกฺขุนิโย อปฺปิจฺฉา ฯเปฯ ตา อุชฺฌายนฺติ ขิยฺยนฺติ วิปาเจนฺติ “กถํ หิ นาม ภิกฺขุนี ภิกฺขุนิยา จีวรํ อนาปุจฺฉา ปารุปิสฺสตี”ติ ฯเปฯ.\csromanspacer{} สจฺจํ กิร ภิกฺขเว ภิกฺขุนี ภิกฺขุนิยา จีวรํ อนาปุจฺฉา ปารุปตีติ\footnote{ปารุปีติ (ก)}.\csromanspacer{} สจฺจํ ภควาติ.\csromanspacer{} วิครหิ พุทฺโธ ภควา ฯเปฯ.\csromanspacer{} กถํ หิ นาม ภิกฺขเว ภิกฺขุนี ภิกฺขุนิยา จีวรํ อนาปุจฺฉา ปารุปิสฺสติ.\csromanspacer{} เนตํ ภิกฺขเว อปฺปสนฺนานํ วา ปสาทาย ฯเปฯ.\csromanspacer{} เอวญฺจ ปน ภิกฺขเว ภิกฺขุนิโย อิมํ สิกฺขาปทํ อุทฺทิสนฺตุ\csromandash{}}

\proseitem{903}{\textbf{“ยา ปน ภิกฺขุนี จีวรสงฺกมนียํ ธาเรยฺย, ปาจิตฺติยนฺ”}ติ.}

\proseitem{904}{\textbf{ยา ปนา}ติ ยา ยาทิสา ฯเปฯ.\csromanspacer{} \textbf{ภิกฺขุนี}ติ ฯเปฯ อยํ อิมสฺมิํ อตฺเถ อธิปฺเปตา ภิกฺขุนีติ.}

\prose{\textbf{จีวรสงฺกมนียํ} นาม อุปสมฺปนฺนาย ปญฺจนฺนํ จีวรานํ อญฺญตรํ จีวรํ ตสฺสา วา อทินฺนํ ตํ วา อนาปุจฺฉา นิวาเสติ วา ปารุปติ วา, อาปตฺติ ปาจิตฺติยสฺส.}

\proseitem{905}{อุปสมฺปนฺนาย อุปสมฺปนฺนสญฺญา จีวรสงฺกมนียํ ธาเรติ, อาปตฺติ ปาจิตฺติยสฺส.\csromanspacer{} อุปสมฺปนฺนาย เวมติกา จีวรสงฺกมนียํ ธาเรติ, อาปตฺติ ปาจิตฺติยสฺส.\csromanspacer{} อุปสมฺปนฺนาย อนุปสมฺปนฺนสญฺญา จีวรสงฺกมนียํ ธาเรติ, อาปตฺติ ปาจิตฺติยสฺส.}

\prose{อนุปสมฺปนฺนาย จีวรสงฺกมนียํ ธาเรติ, อาปตฺติ ทุกฺกฏสฺส.\csromanspacer{} อนุปสมฺปนฺนาย อุปสมฺปนฺนสญฺญา, อาปตฺติ ทุกฺกฏสฺส.\csromanspacer{} อนุปสมฺปนฺนาย เวมติกา, อาปตฺติ ทุกฺกฏสฺส.\csromanspacer{} อนุปสมฺปนฺนาย อนุปสมฺปนฺนสญฺญา, อาปตฺติ ทุกฺกฏสฺส.}

\proseitemwithcloser{906}{อนาปตฺติ สา วา เทติ, ตํ วา อาปุจฺฉา นิวาเสติ วา ปารุปติ วา, อจฺฉินฺนจีวริกาย นฏฺฐจีวริกาย อาปทาสุ อุมฺมตฺติกาย อาทิกมฺมิกายาติ.}{ปญฺจมสิกฺขาปทํ นิฏฺฐิตํ.}
\csromanmatikaanchor{mtk.177}
\csromantocmarkref{subsection}{6. ฉฏฺฐสิกฺขาปท}{mtk.177}
\bookmark[dest={mtk.177},level=2]{6. ฉฏฺฐสิกฺขาปท}
\csromanheader{2}{3. \textbf{นคฺควคฺค 6. ฉฏฺฐสิกฺขาปท}}
\proseitem{907}{\csromanfolio{372}เตน สมเยน พุทฺโธ ภควา สาวตฺถิยํ วิหรติ เชตวเน อนาถปิณฺฑิกสฺส อาราเม.\csromanspacer{} เตน โข ปน สมเยน ถุลฺลนนฺทาย ภิกฺขุนิยา อุปฏฺฐากกุลํ ถุลฺลนนฺทํ ภิกฺขุนิํ เอตทโวจ “ภิกฺขุนิสํฆสฺส อยฺเย จีวรํ ทสฺสามา”ติ.\csromanspacer{} ถุลฺลนนฺทา ภิกฺขุนี “ตุเมฺห พหุกิจฺจา พหุกรณียา”ติ อนฺตรายํ อกาสิ.\csromanspacer{} เตน โข ปน สมเยน ตสฺส กุลสฺส ฆรํ ฑยฺหติ.\csromanspacer{} เต อุชฺฌายนฺติ ขิยฺยนฺติ วิปาเจนฺติ “กถํ หิ นาม อยฺยา ถุลฺลนนฺทา อมฺหากํ เทยฺยธมฺมํ อนฺตรายํ กริสฺสติ.\csromanspacer{} อุภเยนามฺห ปริพาหิรา\footnote{ปริหีนา (สี)} โภเคหิ จ ปุญฺเญน จา”ติ.\csromanspacer{} อสฺโสสุํ โข ภิกฺขุนิโย เตสํ มนุสฺสานํ อุชฺฌายนฺตานํ ขิยฺยนฺตานํ วิปาเจนฺตานํ.\csromanspacer{} ยา ตา ภิกฺขุนิโย อปฺปิจฺฉา ฯเปฯ ตา อุชฺฌายนฺติ ขิยฺยนฺติ วิปาเจนฺติ “กถํ หิ นาม อยฺยา ถุลฺลนนฺทา คณสฺส จีวรลาภํ อนฺตรายํ กริสฺสตี”ติ ฯเปฯ.\csromanspacer{} สจฺจํ กิร ภิกฺขเว ถุลฺลนนฺทา ภิกฺขุนี คณสฺส จีวรลาภํ อนฺตรายํ อกาสีติ.\csromanspacer{} สจฺจํ ภควาติ.\csromanspacer{} วิครหิ พุทฺโธ ภควา ฯเปฯ.\csromanspacer{} กถํ หิ นาม ภิกฺขเว ถุลฺลนนฺทา ภิกฺขุนี คณสฺส จีวรลาภํ อนฺตรายํ กริสฺสติ.\csromanspacer{} เนตํ ภิกฺขเว อปฺปสนฺนานํ วา ปสาทาย ฯเปฯ.\csromanspacer{} เอวญฺจ ปน ภิกฺขเว ภิกฺขุนิโย อิมํ สิกฺขาปทํ อุทฺทิสนฺตุ\csromandash{}}

\proseitem{908}{\textbf{“ยา ปน ภิกฺขุนี คณสฺส จีวรลาภํ อนฺตรายํ กเรยฺย, ปาจิตฺติยนฺ”}ติ.}

\proseitem{909}{\textbf{ยา ปนา}ติ ยา ยาทิสา ฯเปฯ.\csromanspacer{} \textbf{ภิกฺขุนี}ติ ฯเปฯ อยํ อิมสฺมิํ อตฺเถ อธิปฺเปตา ภิกฺขุนีติ.}

\prose{\textbf{คโณ} นาม ภิกฺขุนิสํโฆ วุจฺจติ.}

\prose{\textbf{จีวรํ} นาม ฉนฺนํ จีวรานํ อญฺญตรํ จีวรํ วิกปฺปนุปคํ ปจฺฉิมํ.}

\prose{\textbf{อนฺตรายํ กเรยฺยา}ติ “กถํ อิเม จีวรํ น\footnote{อิมํ จีวรํ (ก)} ทเทยฺยุนฺ”ติ อนฺตรายํ กโรติ, อาปตฺติ ปาจิตฺติยสฺส.\csromanspacer{} อญฺญํ ปริกฺขารํ อนฺตรายํ กโรติ, อาปตฺติ ทุกฺกฏสฺส.\csromanspacer{} สมฺพหุลานํ ภิกฺขุนีนํ วา เอกภิกฺขุนิยา วา อนุปสมฺปนฺนาย วา จีวรํ วา อญฺญํ วา ปริกฺขารํ อนฺตรายํ กโรติ, อาปตฺติ ทุกฺกฏสฺส.}

\proseitemwithcloserruled{910}{\csromanfolio{373}อนาปตฺติ อานิสํสํ ทสฺเสตฺวา นิวาเรติ, อุมฺมตฺติกาย อาทิกมฺมิกายาติ.}{ฉฏฺฐสิกฺขาปทํ นิฏฺฐิตํ.}
\csromanmatikaanchor{mtk.178}
\csromantocmarkref{subsection}{7. สตฺตมสิกฺขาปท}{mtk.178}
\bookmark[dest={mtk.178},level=2]{7. สตฺตมสิกฺขาปท}
\csromanheader{2}{3. \textbf{นคฺควคฺค 7. สตฺตมสิกฺขาปท}}
\proseitem{911}{เตน สมเยน พุทฺโธ ภควา สาวตฺถิยํ วิหรติ เชตวเน อนาถปิณฺฑิกสฺส อาราเม.\csromanspacer{} เตน โข ปน สมเยน ภิกฺขุนิสํฆสฺส อกาลจีวรํ อุปฺปนฺนํ โหติ.\csromanspacer{} อถ โข ภิกฺขุนิสํโฆ ตํ จีวรํ ภาเชตุกาโม สนฺนิปติ.\csromanspacer{} เตน โข ปน สมเยน ถุลฺลนนฺทาย ภิกฺขุนิยา อนฺเตวาสินิโย ภิกฺขุนิโย ปกฺกนฺตา โหนฺติ.\csromanspacer{} ถุลฺลนนฺทา ภิกฺขุนี ตา ภิกฺขุนิโย เอตทโวจ “อยฺเย ภิกฺขุนิโย ปกฺกนฺตา, น ตาว จีวรํ ภาชียิสฺสตี”ติ จีวรวิภงฺคํ ปฏิพาหิ.\csromanspacer{} ภิกฺขุนิโย น ตาว จีวรํ ภาชียิสฺสตีติ ปกฺกมิํสุ\footnote{วิปกฺกมิํสุ (สฺยา)}.\csromanspacer{} ถุลฺลนนฺทา ภิกฺขุนี อนฺเตวาสินีสุ ภิกฺขุนีสุ อาคตาสุ ตํ จีวรํ ภาชาเปสิ.\csromanspacer{} ยา ตา ภิกฺขุนิโย อปฺปิจฺฉา - ป- ตา อุชฺฌายนฺติ ขิยฺยนฺติ วิปาเจนฺติ “กถํ หิ นาม อยฺยา ถุลฺลนนฺทา ธมฺมิกํ จีวรวิภงฺคํ ปฏิพาหิสฺสตี”ติ ฯเปฯ.\csromanspacer{} สจฺจํ กิร ภิกฺขเว ถุลฺลนนฺทา ภิกฺขุนี ธมฺมิกํ จีวรวิภงฺคํ ปฏิพาหตีติ\footnote{ปฏิพาหีติ (ก)}.\csromanspacer{} สจฺจํ ภควาติ.\csromanspacer{} วิครหิ พุทฺโธ ภควา ฯเปฯ.\csromanspacer{} กถํ หิ นาม ภิกฺขเว ถุลฺลนนฺทา ภิกฺขุนี ธมฺมิกํ จีวรวิภงฺคํ ปฏิพาหิสฺสติ.\csromanspacer{} เนตํ ภิกฺขเว อปฺปสนฺนานํ วา ปสาทาย ฯเปฯ.\csromanspacer{} เอวญฺจ ปน ภิกฺขเว ภิกฺขุนิโย อิมํ สิกฺขาปทํ อุทฺทิสนฺตุ\csromandash{}}

\proseitem{912}{\textbf{“ยา ปน ภิกฺขุนี ธมฺมิกํ จีวรวิภงฺคํ ปฏิพาเหยฺย, ปาจิตฺติยนฺ”}ติ.}

\proseitem{913}{\textbf{ยา ปนา}ติ ยา ยาทิสา ฯเปฯ.\csromanspacer{} \textbf{ภิกฺขุนี}ติ ฯเปฯ อยํ อิมสฺมิํ อตฺเถ อธิปฺเปตา ภิกฺขุนีติ.}

\prose{\textbf{ธมฺมิโก} นาม จีวรวิภงฺโค สมคฺโค ภิกฺขุนิสํโฆ สนฺนิปติตฺวา ภาเชติ.}

\prose{\textbf{ปฏิพาเหยฺยา}ติ “กถํ อิมํ จีวรํ น ภาเชยฺยา”ติ\footnote{จีวรํ ภาเชยฺยาติ (ก)} ปฏิพาหติ, อาปตฺติ}

\proseitem{914}{\csromanfolio{374}ธมฺมิเก ธมฺมิกสญฺญา ปฏิพาหติ, อาปตฺติ ปาจิตฺติยสฺส.\csromanspacer{} ธมฺมิเก เวมติกา ปฏิพาหติ, อาปตฺติ ทุกฺกฏสฺส.\csromanspacer{} ธมฺมิเก อธมฺมิกสญฺญา ปฏิพาหติ, อนาปตฺติ.\csromanspacer{} อธมฺมิเก ธมฺมิกสญฺญา, อาปตฺติ ทุกฺกฏสฺส.\csromanspacer{} อธมฺมิเก เวมติกา, อาปตฺติ ทุกฺกฏสฺส.\csromanspacer{} อธมฺมิเก อธมฺมิกสญฺญา, อนาปตฺติ.}

\proseitemwithcloserruled{915}{อนาปตฺติ อานิสํสํ ทสฺเสตฺวา ปฏิพาหติ, อุมฺมตฺติกาย อาทิกมฺมิกายาติ.}{สตฺตมสิกฺขาปทํ นิฏฺฐิตํ.}
\csromanmatikaanchor{mtk.179}
\csromantocmarkref{subsection}{8. อฏฺฐมสิกฺขาปท}{mtk.179}
\bookmark[dest={mtk.179},level=2]{8. อฏฺฐมสิกฺขาปท}
\csromanheader{2}{3. \textbf{นคฺควคฺค 8. อฏฺฐมสิกฺขาปท}}
\proseitem{916}{เตน สมเยน พุทฺโธ ภควา สาวตฺถิยํ วิหรติ เชตวเน อนาถปิณฺฑิกสฺส อาราเม.\csromanspacer{} เตน โข ปน สมเยน ถุลฺลนนฺทา ภิกฺขุนี นฏานมฺปิ นฏกานมฺปิ ลงฺฆกานมฺปิ โสกชฺฌายิกานมฺปิ\footnote{โสกสายิกานมฺปิ (สี)} กุมฺภถูณิกานมฺปิ\footnote{กุมฺภถุณิกานมฺปิ (ก)} สมณจีวรํ เทติ “มยฺหํ ปริสติ\footnote{ปริสติํ (สี)} วณฺณํ ภาสถา”ติ.\csromanspacer{} นฏาปิ นฏกาปิ ลงฺฆกาปิ โสกชฺฌายิกาปิ กุมฺภถูณิกาปิ ถุลฺลนนฺทาย ภิกฺขุนิยา ปริสติ วณฺณํ ภาสนฺติ “อยฺยา ถุลฺลนนฺทา พหุสฺสุตา ภาณิกา วิสารทา ปฏฺฏา ธมฺมิํ กถํ กาตุํ, เทถ อยฺยาย, กโรถ อยฺยายา”ติ.\csromanspacer{} ยา ตา ภิกฺขุนิโย อปฺปิจฺฉา, ตา อุชฺฌายนฺติ ขิยฺยนฺติ วิปาเจนฺติ “กถํ หิ นาม อยฺยา ถุลฺลนนฺทา อคาริกสฺส สมณจีวรํ ทสฺสตี”ติ ฯเปฯ.\csromanspacer{} สจฺจํ กิร ภิกฺขเว ถุลฺลนนฺทา ภิกฺขุนี อคาริกสฺส สมณจีวรํ เทตีติ.\csromanspacer{} สจฺจํ ภควาติ.\csromanspacer{} วิครหิ พุทฺโธ ภควา ฯเปฯ.\csromanspacer{} กถํ หิ นาม ภิกฺขเว ถุลฺลนนฺทา ภิกฺขุนี อคาริกสฺส สมณจีวรํ ทสฺสติ.\csromanspacer{} เนตํ ภิกฺขเว อปฺปสนฺนานํ วา ปสาทาย ฯเปฯ.\csromanspacer{} เอวญฺจ ปน ภิกฺขเว ภิกฺขุนิโย อิมํ สิกฺขาปทํ อุทฺทิสนฺตุ\csromandash{}}

\proseitem{917}{\textbf{“ยา ปน ภิกฺขุนี อคาริกสฺส วา ปริพฺพาชกสฺส วา ปริพฺพาชิกาย วา สมณจีวรํ ทเทยฺย, ปาจิตฺติยนฺ”}ติ.}

\proseitem{918}{\textbf{ยา ปนา}ติ ยา ยาทิสา ฯเปฯ.\csromanspacer{} \textbf{ภิกฺขุนี}ติ ฯเปฯ อยํ อิมสฺมิํ อตฺเถ อธิปฺเปตา ภิกฺขุนีติ.}

\prose{\csromanfolio{375}\textbf{อคาริโก} นาม โย โกจิ อคารํ อชฺฌาวสติ.}

\prose{\textbf{ปริพฺพาชโก} นาม ภิกฺขุญฺจ สามเณรญฺจ ฐเปตฺวา โย โกจิ ปริพฺพาชกสมาปนฺโน.}

\prose{\textbf{ปริพฺพาชิกา} นาม ภิกฺขุนิญฺจ สิกฺขมานญฺจ สามเณริญฺจ ฐเปตฺวา ยา กาจิ ปริพฺพาชิกสมาปนฺนา.}

\prose{\textbf{สมณจีวรํ} นาม กปฺปกตํ วุจฺจติ.\csromanspacer{} เทติ, อาปตฺติ ปาจิตฺติยสฺส.}

\proseitemwithcloserruled{919}{อนาปตฺติ มาตาปิตูนํ เทติ, ตาวกาลิกํ เทติ อุมฺมตฺติกาย อาทิกมฺมิกายาติ.}{อฏฺฐมสิกฺขาปทํ นิฏฺฐิตํ.}
\csromanmatikaanchor{mtk.180}
\csromantocmarkref{subsection}{9. นวมสิกฺขาปท}{mtk.180}
\bookmark[dest={mtk.180},level=2]{9. นวมสิกฺขาปท}
\csromanheader{2}{3. \textbf{นคฺควคฺค 9. นวมสิกฺขาปท}}
\proseitem{920}{เตน สมเยน พุทฺโธ ภควา สาวตฺถิยํ วิหรติ เชตวเน อนาถปิณฺฑิกสฺส อาราเม.\csromanspacer{} เตน โข ปน สมเยน ถุลฺลนนฺทาย ภิกฺขุนิยา อุปฏฺฐากกุลํ ถุลฺลนนฺทํ ภิกฺขุนิํ เอตทโวจ “สเจ มยํ อยฺเย สกฺโกม, ภิกฺขุนิสํฆสฺส จีวรํ ทสฺสามา”ติ.\csromanspacer{} เตน โข ปน สมเยน วสฺสํวุฏฺฐา ภิกฺขุนิโย จีวรํ ภาเชตุกามา สนฺนิปติํสุ.\csromanspacer{} ถุลฺลนนฺทา ภิกฺขุนี ตา ภิกฺขุนิโย เอตทโวจ “อาคเมถ อยฺเย, อตฺถิ ภิกฺขุนิสํฆสฺส จีวรปจฺจาสา”ติ.\csromanspacer{} ภิกฺขุนิโย ถุลฺลนนฺทํ ภิกฺขุนิํ เอตทโวจุํ “คจฺฉายฺเย, ตํ จีวรํ ชานาหี”ติ.\csromanspacer{} ถุลฺลนนฺทา ภิกฺขุนี เยน ตํ กุลํ เตนุปสงฺกมิ, อุปสงฺกมิตฺวา เต มนุสฺเส เอตทโวจ “เทถาวุโส ภิกฺขุนิสํฆสฺส จีวรนฺ”ติ.\csromanspacer{} น มยํ อยฺเย สกฺโกม ภิกฺขุนิสํฆสฺส จีวรํ ทาตุนฺติ.\csromanspacer{} ถุลฺลนนฺทา ภิกฺขุนี ภิกฺขุนีนํ เอตมตฺถํ อาโรเจสิ.\csromanspacer{} ยา ตา ภิกฺขุนิโย อปฺปิจฺฉา ฯเปฯ ตา อุชฺฌายนฺติ ขิยฺยนฺติ วิปาเจนฺติ “กถํ หิ นาม อยฺยา ถุลฺลนนฺทา ทุพฺพลจีวรปจฺจาสาย จีวรกาลสมยํ อติกฺกาเมสฺสตี”ติ ฯเปฯ.\csromanspacer{} สจฺจํ กิร ภิกฺขเว ถุลฺลนนฺทา ภิกฺขุนี ทุพฺพลจีวรปจฺจาสาย จีวรกาลสมยํ อติกฺกาเมตีติ\footnote{อติกฺกาเมสีติ (ก)}.\csromanspacer{} สจฺจํ ภควาติ.\csromanspacer{} วิครหิ พุทฺโธ ภควา ฯเปฯ.\csromanspacer{} กถํ หิ นาม ภิกฺขเว ถุลฺลนนฺทา \csromanfolio{376}\textbf{ภิกฺขุนี ทุพฺพลจีวรปจฺจาสาย จีวรกาลสมยํ อติกฺกาเมสฺสติ.}\csromanspacer{} เนตํ ภิกฺขเว อปฺปสนฺนานํ วา ปสาทาย ฯเปฯ.\csromanspacer{} เอวญฺจ ปน ภิกฺขเว ภิกฺขุนิโย อิมํ สิกฺขาปทํ อุทฺทิสนฺตุ\csromandash{}}

\proseitem{921}{\textbf{“ยา ปน ภิกฺขุนี ทุพฺพลจีวรปจฺจาสาย จีวรกาลสมยํ อติกฺกาเมยฺย, ปาจิตฺติยนฺ”}ติ.}

\proseitem{922}{\textbf{ยา ปนา}ติ ยา ยาทิสา ฯเปฯ.\csromanspacer{} \textbf{ภิกฺขุนี}ติ ฯเปฯ อยํ อิมสฺมิํ อตฺเถ อธิปฺเปตา ภิกฺขุนีติ.}

\prose{\textbf{ทุพฺพลจีวรปจฺจาสา} นาม “สเจ มยํ สกฺโกม, ทสฺสาม กริสฺสามา”ติ วาจา ภินฺนา โหติ.}

\prose{\textbf{จีวรกาลสมโย} นาม อนตฺถเต กถิเน วสฺสานสฺส ปจฺฉิโม มาโส, อตฺถเต กถิเน ปญฺจมาสา.}

\prose{\textbf{จิวรกาลสมยํ อติกฺกาเมยฺยา}ติ อนตฺถเต กถิเน วสฺสานสฺส ปจฺฉิมํ ทิวสํ อติกฺกาเมติ, อาปตฺติ ปาจิตฺติยสฺส.\csromanspacer{} อตฺถเต กถิเน กถินุทฺธารทิวสํ อติกฺกาเมติ, อาปตฺติ ปาจิตฺติยสฺส.}

\proseitem{923}{ทุพฺพลจีวเร ทุพฺพลจีวรสญฺญา จีวรกาลสมยํ อติกฺกาเมติ, อาปตฺติ ปาจิตฺติยสฺส.\csromanspacer{} ทุพฺพลจีวเร เวมติกา จีวรกาลสมยํ อติกฺกาเมติ, อาปตฺติ ทุกฺกฏสฺส.\csromanspacer{} ทุพฺพลจีวเร อทุพฺพลจีวรสญฺญา จีวรกาลสมยํ อติกฺกาเมติ, อนาปตฺติ.}

\prose{อทุพฺพลจีวเร ทุพฺพลจีวรสญฺญา, อาปตฺติ ทุกฺกฏสฺส.\csromanspacer{} อทุพฺพลจีวเร เวมติกา, อาปตฺติ ทุกฺกฏสฺส.\csromanspacer{} อทุพฺพลจีวเร อทุพฺพลจีวรสญฺญา, อนาปตฺติ.}

\proseitemwithcloser{924}{อนาปตฺติ อานิสํสํ ทสฺเสตฺวา นิวาเรติ, อุมฺมตฺติกาย อาทิกมฺมิกายาติ.}{นวมสิกฺขาปทํ นิฏฺฐิตํ.}
\csromanmatikaanchor{mtk.181}
\csromantocmarkref{subsection}{10. ทสมสิกฺขาปท}{mtk.181}
\bookmark[dest={mtk.181},level=2]{10. ทสมสิกฺขาปท}
\csromanheader{2}{3. \textbf{นคฺควคฺค 1}0. ทสมสิกฺขาปท}
\proseitem{925}{\csromanfolio{377}เตน สมเยน พุทฺโธ ภควา สาวตฺถิยํ วิหรติ เชตวเน อนาถปิณฺฑิกสฺส อาราเม.\csromanspacer{} เตน โข ปน สมเยน อญฺญตเรน อุปาสเกน สํฆํ อุทฺทิสฺส วิหาโร การาปิโต โหติ.\csromanspacer{} โส ตสฺส วิหารสฺส มเห อุภโตสํฆสฺส อกาลจีวรํ ทาตุกาโม โหติ.\csromanspacer{} เตน โข ปน สมเยน อุภโตสํฆสฺส กถินํ อตฺถตํ โหติ.\csromanspacer{} อถ โข โส อุปาสโก สํฆํ อุปสงฺกมิตฺวา กถินุทฺธารํ ยาจิ.\csromanspacer{} ภควโต เอตมตฺถํ อาโรเจสุํ.\csromanspacer{} อถ โข ภควา เอตสฺมิํ นิทาเน เอตสฺมิํ ปกรเณ ธมฺมิํ กถํ กตฺวา ภิกฺขู อามนฺเตสิ อนุชานามิ ภิกฺขเว กถินํ อุทฺธริตุํ, เอวญฺจ ปน ภิกฺขเว กถินํ อุทฺธริตพฺพํ.\csromanspacer{} พฺยตฺเตน ภิกฺขุนา ปฏิพเลน สํโฆ ญาเปตพฺโพ\csromandash{}}

\proseitem{926}{“สุณาตุ เม ภนฺเต สํโฆ, ยทิ สํฆสฺส ปตฺตกลฺลํ, สํโฆ กถินํ อุทฺธเรยฺย, เอสา ญตฺติ.}

\prose{สุณาตุ เม ภนฺเต สํโฆ, สํโฆ กถินํ อุทฺธรติ, ยสฺสายสฺมโต ขมติ กถินสฺส อุทฺธาโร, โส ตุณฺหสฺส, ยสฺส นกฺขมติ, โส ภาเสยฺย.}

\prose{อุพฺภตํ สํเฆน กถินํ, ขมติ สํฆสฺส, ตสฺมา ตุณฺหี, เอวเมตํ ธารยามี”ติ.}

\proseitem{927}{อถ โข โส อุปาสโก ภิกฺขุนิสํฆํ อุปสงฺกมิตฺวา กถินุทฺธารํ ยาจิ.\csromanspacer{} ถุลฺลนนฺทา ภิกฺขุนี “จีวรํ อมฺหากํ ภวิสฺสตี”ติ กถินุทฺธารํ ปฏิพาหิ.\csromanspacer{} อถ โข โส อุปาสโก อุชฺฌายติ ขิยฺยติ วิปาเจติ “กถํ หิ นาม ภิกฺขุนิโย อมฺหากํ กถินุทฺธารํ น ทสฺสนฺตี”ติ.\csromanspacer{} อสฺโสสุํ โข ภิกฺขุนิโย ตสฺส อุปาสกสฺส อุชฺฌายนฺตสฺส ขิยฺยนฺตสฺส วิปาเจนฺตสฺส.\csromanspacer{} ยา ตา ภิกฺขุนิโย อปฺปิจฺฉา ฯเปฯ ตา อุชฺฌายนฺติ ขิยฺยนฺติ วิปาเจนฺติ “กถํ หิ นาม อยฺยา ถุลฺลนนฺทา ธมฺมิกํ กถินุทฺธารํ ปฏิพาหิสฺสตี”ติ ฯเปฯ.\csromanspacer{} สจฺจํ กิร ภิกฺขเว ถุลฺลนนฺทา ภิกฺขุนี ธมฺมิกํ กถินุทฺธารํ ปฏิพาหตีติ\footnote{ปฏิพาหีติ (ก)}.\csromanspacer{} สจฺจํ ภควาติ.\csromanspacer{} วิครหิ พุทฺโธ ภควา ฯเปฯ.\csromanspacer{} กถํ หิ นาม ภิกฺขเว ถุลฺลนนฺทา ภิกฺขุนี ธมฺมิกํ กถินุทฺธารํ \csromanfolio{378}ปฏิพาหิสฺสติ.\csromanspacer{} เนตํ ภิกฺขเว อปฺปสนฺนานํ วา ปสาทาย ฯเปฯ.\csromanspacer{} เอวญฺจ ปน ภิกฺขเว ภิกฺขุนิโย อิมํ สิกฺขาปทํ อุทฺทิสนฺตุ\csromandash{}}

\proseitem{928}{\textbf{“ยา ปน ภิกฺขุนี ธมฺมิกํ กถินุทฺธารํ ปฏิพาเหยฺย, ปาจิตฺติยนฺ”}ติ.}

\proseitem{929}{\textbf{ยา ปนา}ติ ยา ยาทิสา ฯเปฯ.\csromanspacer{} \textbf{ภิกฺขุนี}ติ ฯเปฯ อยํ อิมสฺมิํ อตฺเถ อธิปฺเปตา ภิกฺขุนีติ.}

\prose{\textbf{ธมฺมิโก} นาม กถินุทฺธาโร สมคฺโค ภิกฺขุนิสํโฆ สนฺนิปติตฺวา อุทฺธรติ.}

\prose{\textbf{ปฏิพาเหยฺยา}ติ “กถํ อิทํ กถินํ น อุทฺธเรยฺยา”ติ\footnote{กถินํ อุทฺธเรยฺยาติ (ก)} ปฏิพาหติ, อาปตฺติ ปาจิตฺติยสฺส.}

\proseitem{930}{ธมฺมิเก ธมฺมิกสญฺญา ปฏิพาหติ, อาปตฺติ ปาจิตฺติยสฺส.\csromanspacer{} ธมฺมิเก เวมติกา ปฏิพาหติ, อาปตฺติ ทุกฺกฏสฺส.\csromanspacer{} ธมฺมิเก อธมฺมิกสญฺญา ปฏิพาหติ, อนาปตฺติ.\csromanspacer{} อธมฺมิเก ธมฺมิกสญฺญา, อาปตฺติ ทุกฺกฏสฺส.\csromanspacer{} อธมฺมิเก เวมติกา, อาปตฺติ ทุกฺกฏสฺส.\csromanspacer{} อธมฺมิเก อธมฺมิกสญฺญา, อนาปตฺติ.}

\proseitemwithcloser{931}{อนาปตฺติ อานิสํสํ ทสฺเสตฺวา ปฏิพาหติ, อุมฺมตฺติกาย อาทิกมฺมิกายาติ.}{ทสมสิกฺขาปทํ นิฏฺฐิตํ.}
\csromanheader{2}{นคฺควคฺโค ตติโย.}
\csromansectionrule
\csromanmatikaanchor{mtk.182}
\csromanmatikaanchor{mtk.183}
\csromantocmarknopage{section}{4. ตุวฏฺฏวคฺค}{mtk.182}
\bookmark[dest={mtk.182},level=1]{4. ตุวฏฺฏวคฺค}
\csromantocmarkref{subsection}{1. ปฐมสิกฺขาปท}{mtk.183}
\bookmark[dest={mtk.183},level=2]{1. ปฐมสิกฺขาปท}
\setcsromanheadh{4. ตุวฏฺฏวคฺค}
\csromanheader{2}{4. \textbf{ตุวฏฺฏวคฺค 1. ปฐมสิกฺขาปท}}
\proseitem{932}{เตน สมเยน พุทฺโธ ภควา สาวตฺถิยํ วิหรติ เชตวเน อนาถปิณฺฑิกสฺส อาราเม.\csromanspacer{} เตน โข ปน สมเยน ภิกฺขุนิโย เทฺว เอกมญฺเจ ตุวฏฺเฏนฺติ.\csromanspacer{} มนุสฺสา วิหารจาริกํ อาหิณฺฑนฺตา ปสฺสิตฺวา \csromanfolio{379}อุชฺฌายนฺติ ขิยฺยนฺติ วิปาเจนฺติ “กถํ หิ นาม ภิกฺขุนิโย เทฺว เอกมญฺเจ ตุวฏฺเฏสฺสนฺติ เสยฺยถาปิ คิหินิโย กามโภคินิโย”ติ.\csromanspacer{} อสฺโสสุํ โข ภิกฺขุนิโย เตสํ มนุสฺสานํ อุชฺฌายนฺตานํ ขิยฺยนฺตานํ วิปาเจนฺตานํ.\csromanspacer{} ยา ตา ภิกฺขุนิโย อปฺปิจฺฉา ฯเปฯ ตา อุชฺฌายนฺติ ขิยฺยนฺติ วิปาเจนฺติ “กถํ หิ นาม ภิกฺขุนิโย เทฺว เอกมญฺเจ ตุวฏฺเฏสฺสนฺตี”ติ ฯเปฯ.\csromanspacer{} สจฺจํ กิร ภิกฺขเว ภิกฺขุนิโย เทฺว เอกมญฺเจ ตุวฏฺเฏนฺตีติ.\csromanspacer{} สจฺจํ ภควาติ.\csromanspacer{} วิครหิ พุทฺโธ ภควา ฯเปฯ.\csromanspacer{} กถํ หิ นาม ภิกฺขเว ภิกฺขุนิโย เทฺว เอกมญฺเจ ตุวฏฺเฏสฺสนฺติ.\csromanspacer{} เนตํ ภิกฺขเว อปฺปสนฺนานํ \textbf{วา ปสาทาย ฯเปฯ.}\csromanspacer{}\textbf{ เอวญฺจ ปน ภิกฺขเว ภิกฺขุนิโย อิมํ สิกฺขาปทํ} \textbf{อุทฺทิสนฺตุ\csromandash{}}}

\proseitem{933}{\textbf{“ยา ปน ภิกฺขุนิโย เทฺว เอกมญฺเจ ตุวฏฺเฏยฺยุํ, ปาจิตฺติยนฺ”}ติ.}

\proseitem{934}{\textbf{ยา ปนา}ติ ยา ยาทิสา ฯเปฯ.\csromanspacer{} \textbf{ภิกฺขุนิโย}ติ อุปสมฺปนฺนาโย วุจฺจนฺติ.}

\prose{\textbf{เทฺว เอกมญฺเจ ตุวฏฺเฏยฺยุนฺ}ติ เอกาย นิปนฺนาย อปรา นิปชฺชติ, อาปตฺติ ปาจิตฺติยสฺส.\csromanspacer{} อุโภ วา นิปชฺชนฺติ, อาปตฺติ ปาจิตฺติยสฺส.\csromanspacer{} อุฏฺฐหิตฺวา ปุนปฺปุนํ นิปชฺชนฺติ, อาปตฺติ ปาจิตฺติยสฺส.}

\proseitemwithcloserruled{935}{อนาปตฺติ เอกาย นิปนฺนาย อปรา นิสีทติ, อุโภ วา นิสีทนฺติ, อุมฺมตฺติกานํ อาทิกมฺมิกานนฺติ.}{ปฐมสิกฺขาปทํ นิฏฺฐิตํ.}
\csromanmatikaanchor{mtk.184}
\csromantocmarkref{subsection}{2. ทุติยสิกฺขาปท}{mtk.184}
\bookmark[dest={mtk.184},level=2]{2. ทุติยสิกฺขาปท}
\csromanheader{2}{4. \textbf{ตุวฏฺฏวคฺค 2. ทุติยสิกฺขาปท}}
\proseitem{936}{เตน สมเยน พุทฺโธ ภควา สาวตฺถิยํ วิหรติ เชตวเน อนาถปิณฺฑิกสฺส อาราเม.\csromanspacer{} เตน โข ปน สมเยน ภิกฺขุนิโย เทฺว เอกตฺถรณปาวุรณา ตุวฏฺเฏนฺติ.\csromanspacer{} มนุสฺสา วิหารจาริกํ อาหิณฺฑนฺตา ปสฺสิตฺวา อุชฺฌายนฺติ ขิยฺยนฺติ วิปาเจนฺติ “กถํ หิ นาม ภิกฺขุนิโย เทฺว เอกตฺถรณปาวุรณา ตุวฏฺเฏสฺสนฺติ เสยฺยถาปิ คิหินิโย กามโภคินิโย”ติ.\csromanspacer{} อสฺโสสุํ โข ภิกฺขุนิโย เตสํ มนุสฺสานํ \csromanfolio{380}\textbf{อุชฺฌายนฺตานํ ขิยฺยนฺตานํ วิปาเจนฺตานํ.}\csromanspacer{}\textbf{ ยา ตา ภิกฺขุนิโย อปฺปิจฺฉา} \textbf{ฯเปฯ ตา อุชฺฌายนฺติ ขิยฺยนฺติ วิปาเจนฺติ “กถํ หิ นาม ภิกฺขุนิโย เทฺว} เอกตฺถรณปาวุรณา ตุวฏฺเฏสฺสนฺตี”ติ ฯเปฯ.\csromanspacer{} สจฺจํ กิร ภิกฺขเว ภิกฺขุนิโย เทฺว เอกตฺถรณปาวุรณา ตุวฏฺเฏนฺตีติ.\csromanspacer{} สจฺจํ ภควาติ.\csromanspacer{} วิครหิ พุทฺโธ ภควา ฯเปฯ.\csromanspacer{} กถํ หิ นาม ภิกฺขเว ภิกฺขุนิโย เทฺว เอกตฺถรณปาวุรณา ตุวฏฺเฏสฺสนฺติ.\csromanspacer{} เนตํ ภิกฺขเว อปฺปสนฺนานํ วา ปสาทาย ฯเปฯ.\csromanspacer{} เอวญฺจ ปน ภิกฺขเว ภิกฺขุนิโย อิมํ สิกฺขาปทํ อุทฺทิสนฺตุ\csromandash{}}

\proseitem{937}{\textbf{“ยา ปน ภิกฺขุนิโย เทฺว เอกตฺถรณปาวุรณา ตุวฏฺเฏยฺยุํ, ปาจิตฺติยนฺ”}ติ.}

\proseitem{938}{\textbf{ยา ปนา}ติ ยา ยาทิสา ฯเปฯ.\csromanspacer{} \textbf{ภิกฺขุนิโย}ติ อุปสมฺปนฺนาโย วุจฺจนฺติ.}

\prose{\textbf{เทฺว เอกตฺถรณปาวุรณา ตุวฏฺเฏยฺยุนฺ}ติ ตญฺเญว อตฺถริตฺวา ตญฺเญว ปารุปนฺติ, อาปตฺติ ปาจิตฺติยสฺส.}

\proseitem{939}{เอกตฺถรณปาวุรเณ เอกตฺถรณปาวุรณ\-สญฺญา ตุวฏฺเฏนฺติ, อาปตฺติ ปาจิตฺติยสฺส.\csromanspacer{} เอกตฺถรณปาวุรเณ เวมติกา ตุวฏฺเฏนฺติ, อาปตฺติ ปาจิตฺติยสฺส.\csromanspacer{} เอกตฺถรณปาวุรเณ นานตฺถรณปาวุรณ\-สญฺญา ตุวฏฺเฏนฺติ, อาปตฺติ}

\prose{เอกตฺถรเณ นานาปาวุรเณ\footnote{นานาปาวุรณสญฺญา (ก)}, อาปตฺติ ทุกฺกฏสฺส.\csromanspacer{} นานตฺถรเณ เอกปาวุรเณ\footnote{เอกปาวุรณสญฺญา (ก)}, อาปตฺติ ทุกฺกฏสฺส.\csromanspacer{} นานตฺถรณปาวุรเณ เอกตฺถรณปาวุรณ\-สญฺญา, อาปตฺติ ทุกฺกฏสฺส.\csromanspacer{} นานตฺถรณปาวุรเณ เวมติกา, อาปตฺติ ทุกฺกฏสฺส.\csromanspacer{} นานตฺถรณปาวุรเณ นานตฺถรณปาวุรณ\-สญฺญา, อนาปตฺติ.}

\proseitemwithcloser{940}{อนาปตฺติ ววตฺถานํ ทสฺเสตฺวา นิปชฺชนฺติ, อุมฺมตฺติกานํ อาทิกมฺมิกานนฺติ.}{ทุติยสิกฺขาปทํ นิฏฺฐิตํ.}
\csromanmatikaanchor{mtk.185}
\csromantocmarkref{subsection}{3. ตติยสิกฺขาปท}{mtk.185}
\bookmark[dest={mtk.185},level=2]{3. ตติยสิกฺขาปท}
\csromanheader{2}{4. \textbf{ตุวฏฺฏวคฺค 3. ตติยสิกฺขาปท}}
\proseitem{941}{\csromanfolio{381}เตน สมเยน พุทฺโธ ภควา สาวตฺถิยํ วิหรติ เชตวเน อนาถปิณฺฑิกสฺส อาราเม.\csromanspacer{} เตน โข ปน สมเยน ถุลฺลนนฺทา ภิกฺขุนี พหุสฺสุตา โหติ ภาณิกา วิสารทา ปฏฺฏา ธมฺมิํ กถํ กาตุํ.\csromanspacer{} ภทฺทาปิ กาปิลานี พหุสฺสุตา โหติ ภาณิกา วิสารทา ปฏฺฏา ธมฺมิํ กถํ กาตุํ อุฬารสมฺภาวิตา.\csromanspacer{} มนุสฺสา “อยฺยา ภทฺทา กาปิลานี พหุสฺสุตา ภาณิกา วิสารทา ปฏฺฏา ธมฺมิํ กถํ กาตุํ อุฬารสมฺภาวิตา”ติ ภทฺทํ กาปิลานิํ ปฐมํ ปยิรุปาสิตฺวา ปจฺฉา ถุลฺลนนฺทํ ภิกฺขุนิํ ปยิรุปาสนฺติ.\csromanspacer{} ถุลฺลนนฺทา ภิกฺขุนี อิสฺสาปกตา “อิมา กิร อปฺปิจฺฉา สนฺตุฏฺฐา ปวิวิตฺตา อสํสฏฺฐา, ยา อิมา สญฺญตฺติพหุลา วิญฺญตฺติพหุลา วิหรนฺตี”ติ ภทฺทาย กาปิลานิยา ปุรโต จงฺกมติปิ ติฏฺฐติปิ นิสีทติปิ เสยฺยมฺปิ กปฺเปติ อุทฺทิสติปิ อุทฺทิสาเปติปิ สชฺฌายมฺปิ กโรติ.\csromanspacer{} ยา ตา ภิกฺขุนิโย อปฺปิจฺฉา ฯเปฯ ตา อุชฺฌายนฺติ ขิยฺยนฺติ วิปาเจนฺติ “กถํ หิ นาม อยฺยา ถุลฺลนนฺทา อยฺยาย ภทฺทาย กาปิลานิยา สญฺจิจฺจ อผาสุํ กริสฺสตี”ติ ฯเปฯ.\csromanspacer{} สจฺจํ กิร ภิกฺขเว ถุลฺลนนฺทา ภิกฺขุนี ภทฺทาย กาปิลานิยา สญฺจิจฺจ อผาสุํ กโรตีติ.\csromanspacer{} สจฺจํ ภควาติ.\csromanspacer{} วิครหิ พุทฺโธ ภควา ฯเปฯ.\csromanspacer{} กถํ หิ นาม ภิกฺขเว ถุลฺลนนฺทา ภิกฺขุนี ภทฺทาย กาปิลานิยา สญฺจิจฺจ อผาสุํ กริสฺสติ.\csromanspacer{} เนตํ ภิกฺขเว อปฺปสนฺนานํ วา ปสาทาย ฯเปฯ.\csromanspacer{} เอวญฺจ ปน ภิกฺขเว ภิกฺขุนิโย อิมํ สิกฺขาปทํ อุทฺทิสนฺตุ\csromandash{}}

\proseitem{942}{\textbf{“ยา ปน ภิกฺขุนี ภิกฺขุนิยา สญฺจิจฺจ อผาสุํ กเรยฺย, ปาจิตฺติยนฺ”}ติ.}

\proseitem{943}{\textbf{ยา ปนา}ติ ยา ยาทิสา ฯเปฯ.\csromanspacer{} \textbf{ภิกฺขุนี}ติ ฯเปฯ อยํ อิมสฺมิํ อตฺเถ อธิปฺเปตา ภิกฺขุนีติ.}

\prose{\textbf{ภิกฺขุนิยา}ติ อญฺญาย ภิกฺขุนิยา.}

\prose{\textbf{สญฺจิจฺจา}ติ ชานนฺตี สญฺชานนฺตี เจจฺจ อภิวิตริตฺวา วีติกฺกโม.}

\prose{\textbf{อผาสุํ กเรยฺยา}ติ “อิมินา อิมิสฺสา อผาสุ ภวิสฺสตี”ติ อนาปุจฺฉา ปุรโต จงฺกมติ วา ติฏฺฐติ วา นิสีทติ วา เสยฺยํ วา กปฺเปติ อุทฺทิสติ วา อุทฺทิสาเปติ วา สชฺฌายํ วา กโรติ, อาปตฺติ ปาจิตฺติยสฺส.}

\proseitem{944}{\csromanfolio{382}อุปสมฺปนฺนาย อุปสมฺปนฺนสญฺญา สญฺจิจฺจ อผาสุํ กโรติ, อาปตฺติ ปาจิตฺติยสฺส.\csromanspacer{} อุปสมฺปนฺนาย เวมติกา สญฺจิจฺจ อผาสุํ กโรติ, อาปตฺติ ปาจิตฺติยสฺส.\csromanspacer{} อุปสมฺปนฺนาย อนุปสมฺปนฺนสญฺญา สญฺจิจฺจ อผาสุํ กโรติ, อาปตฺติ ปาจิตฺติยสฺส.}

\prose{อนุปสมฺปนฺนาย สญฺจิจฺจ อผาสุํ กโรติ, อาปตฺติ ทุกฺกฏสฺส.\csromanspacer{} อนุปสมฺปนฺนาย อุปสมฺปนฺนสญฺญา, อาปตฺติ ทุกฺกฏสฺส.\csromanspacer{} อนุปสมฺปนฺนาย เวมติกา, อาปตฺติ ทุกฺกฏสฺส.\csromanspacer{} อนุปสมฺปนฺนาย อนุปสมฺปนฺนสญฺญา, อาปตฺติ ทุกฺกฏสฺส.}

\proseitemwithcloserruled{945}{อนาปตฺติ น อผาสุํ กตฺตุกามา อาปุจฺฉา ปุรโต จงฺกมติ วา ติฏฺฐติ วา นิสีทติ วา เสยฺยํ วา กปฺเปติ อุทฺทิสติ วา อุทฺทิสาเปติ วา สชฺฌายํ วา กโรติ, อุมฺมตฺติกาย อาทิกมฺมิกายาติ.}{ตติยสิกฺขาปทํ นิฏฺฐิตํ.}
\csromanmatikaanchor{mtk.186}
\csromantocmarkref{subsection}{4. จตุตฺถสิกฺขาปท}{mtk.186}
\bookmark[dest={mtk.186},level=2]{4. จตุตฺถสิกฺขาปท}
\csromanheader{2}{4. \textbf{ตุวฏฺฏวคฺค 4. จตุตฺถสิกฺขาปท}}
\proseitem{946}{เตน สมเยน พุทฺโธ ภควา สาวตฺถิยํ วิหรติ เชตวเน อนาถปิณฺฑิกสฺส อาราเม.\csromanspacer{} เตน โข ปน สมเยน ถุลฺลนนฺทา ภิกฺขุนี ทุกฺขิตํ สหชีวินิํ เนว อุปฏฺเฐติ, น อุปฏฺฐาปนาย อุสฺสุกฺกํ กโรติ.\csromanspacer{} ยา ตา ภิกฺขุนิโย อปฺปิจฺฉา ฯเปฯ ตา อุชฺฌายนฺติ ขิยฺยนฺติ วิปาเจนฺติ “กถํ หิ นาม อยฺยา ถุลฺลนนฺทา ทุกฺขิตํ สหชีวินิํ เนว อุปฏฺเฐสฺสติ, น อุปฏฺฐาปนาย อุสฺสุกฺกํ กริสฺสตี”ติ ฯเปฯ.\csromanspacer{} สจฺจํ กิร ภิกฺขเว ถุลฺลนนฺทา ภิกฺขุนี ทุกฺขิตํ สหชีวินิํ เนว อุปฏฺเฐติ น อุปฏฺฐาปนาย อุสฺสุกฺกํ กโรตีติ.\csromanspacer{} สจฺจํ ภควาติ.\csromanspacer{} วิครหิ พุทฺโธ ภควา ฯเปฯ.\csromanspacer{} กถํ หิ นาม ภิกฺขเว ถุลฺลนนฺทา ภิกฺขุนี ทุกฺขิตํ สหชีวินิํ เนว อุปฏฺเฐสฺสติ, น อุปฏฺฐาปนาย อุสฺสุกฺกํ กริสฺสติ.\csromanspacer{} เนตํ ภิกฺขเว อปฺปสนฺนานํ วา ปสาทาย ฯเปฯ.\csromanspacer{} เอวญฺจ ปน ภิกฺขเว ภิกฺขุนิโย อิมํ สิกฺขาปทํ อุทฺทิสนฺตุ\csromandash{}}

\proseitem{947}{\textbf{“ยา ปน ภิกฺขุนี ทุกฺขิตํ สหชีวินิํ เนว อุปฏฺเฐยฺย น อุปฏฺฐาปนาย อุสฺสุกฺกํ กเรยฺย, ปาจิตฺติยนฺ”}ติ.}

\proseitem{948}{\csromanfolio{383}\textbf{ยา ปนา}ติ ยา ยาทิสา ฯเปฯ.\csromanspacer{} \textbf{ภิกฺขุนี}ติ ฯเปฯ อยํ อิมสฺมิํ อตฺเถ อธิปฺเปตา ภิกฺขุนีติ.}

\prose{\textbf{ทุกฺขิตา} นาม คิลานา วุจฺจติ.\csromanspacer{} \textbf{สหชีวินี} นาม สทฺธิวิหารินี วุจฺจติ.}

\prose{\textbf{เนว อุปฏฺเฐยฺยา}ติ น สยํ อุปฏฺเฐยฺย.}

\prose{\textbf{น อุปฏฺฐาปนาย อุสฺสุกฺกํ กเรยฺยา}ติ น อญฺญํ อาณาเปยฺย.}

\prose{“เนว อุปฏฺเฐสฺสามิ, น อุปฏฺฐาปนาย อุสฺสุกฺกํ กริสฺสามี”ติ ธุรํ นิกฺขิตฺตมตฺเต อาปตฺติ ปาจิตฺติยสฺส.\csromanspacer{} อนฺเตวาสินิํ วา อนุปสมฺปนฺนํ วา เนว อุปฏฺเฐติ, น อุปฏฺฐาปนาย อุสฺสุกฺกํ กโรติ, อาปตฺติ ทุกฺกฏสฺส.}

\proseitemwithcloserruled{949}{อนาปตฺติ สติ อนฺตราเย, ปริเยสิตฺวา น ลภติ, คิลานาย อาปทาสุ อุมฺมตฺติกาย อาทิกมฺมิกายาติ.}{จตุตฺถสิกฺขาปทํ นิฏฺฐิตํ.}
\csromanmatikaanchor{mtk.187}
\csromantocmarkref{subsection}{5. ปญฺจมสิกฺขาปท}{mtk.187}
\bookmark[dest={mtk.187},level=2]{5. ปญฺจมสิกฺขาปท}
\csromanheader{2}{4. \textbf{ตุวฏฺฏวคฺค 5. ปญฺจมสิกฺขาปท}}
\proseitem{950}{เตน สมเยน พุทฺโธ สาวตฺถิยํ วิหรติ เชตวเน อนาถปิณฺฑิกสฺส อาราเม.\csromanspacer{} เตน โข ปน สมเยน ภทฺทา กาปิลานี สาเกเต วสฺสํ อุปคตา โหติ.\csromanspacer{} สา เกนจิเทว กรณีเยน อุพฺพาฬฺหา ถุลฺลนนฺทาย ภิกฺขุนิยา สนฺติเก ทูตํ ปาเหสิ “สเจ เม อยฺยา ถุลฺลนนฺทา อุปสฺสยํ ทเทยฺย, อาคจฺเฉยฺยามหํ สาวตฺถินฺ”ติ.\csromanspacer{} ถุลฺลนนฺทา ภิกฺขุนี เอวมาห “อาคจฺฉตุ ทสฺสามี”ติ.\csromanspacer{} อถ โข ภทฺทา กาปิลานี สาเกตา สาวตฺถิํ อคมาสิ.\csromanspacer{} ถุลฺลนนฺทา ภิกฺขุนี ภทฺทาย กาปิลานิยา อุปสฺสยํ อทาสิ.\csromanspacer{} เตน โข ปน สมเยน ถุลฺลนนฺทา ภิกฺขุนี พหุสฺสุตา โหติ ภาณิกา วิสารทา ปฏฺฏา ธมฺมิํ กถํ กาตุํ.\csromanspacer{} ภทฺทาปิ กาปิลานี พหุสฺสุตา\footnote{พหุสฺสุตาเยว (ก)} โหติ ภาณิกา วิสารทา ปฏฺฏฺตา ธมฺมิํ กถํ กาตุํ อุฬารสมฺภาวิตา.\csromanspacer{} มนุสฺสา “อยฺยา ภทฺทา กาปิลานี พหุสฺสุตา ภาณิกา วิสารทา ปฏฺฏา ธมฺมิํ กถํ กาตุํ อุฬารสมฺภาวิตา”ติ ภทฺทํ กาปิลานิํ ปฐมํ ปยิรุปาสิตฺวา ปจฺฉา ถุลฺลนนฺทํ ภิกฺขุนิํ ปยิรุปาสนฺติ.\csromanspacer{} ถุลฺลนนฺทา ภิกฺขุนี อิสฺสาปกตา “อิมา \csromanfolio{384}กิร อปฺปิจฺฉา สนฺตุฏฺฐา ปวิวิตฺตา อสํสฏฺฐา, ยา อิมา สญฺญตฺติพหุลา วิญฺญตฺติพหุลา วิหรนฺตี”ติ กุปิตา อนตฺตมนา ภทฺทํ กาปิลานิํ อุปสฺสยา นิกฺกฑฺฒิ.\csromanspacer{} ยา ตา ภิกฺขุนิโย อปฺปิจฺฉา ฯเปฯ ตา อุชฺฌายนฺติ ขิยฺยนฺติ วิปาเจนฺติ “กถํ หิ นาม อยฺยา ถุลฺลนนฺทา อยฺยาย ภทฺทาย กาปิลานิยา อุปสฺสยํ ทตฺวา กุปิตา อนตฺตมนา นิกฺกฑฺฒิสฺสตี”ติ ฯเปฯ.\csromanspacer{} สจฺจํ กิร ภิกฺขเว ถุลฺลนนฺทา ภิกฺขุนี ภทฺทาย กาปิลานิยา อุปสฺสยํ ทตฺวา กุปิตา อนตฺตมนา นิกฺกฑฺฒตีติ\footnote{นิกฺกฑฺฒีติ (ก)}.\csromanspacer{} สจฺจํ ภควาติ.\csromanspacer{} วิครหิ พุทฺโธ ภควา ฯเปฯ.\csromanspacer{} กถํ หิ นาม ภิกฺขเว ถุลฺลนนฺทา ภิกฺขุนี ภทฺทาย กาปิลานิยา อุปสฺสยํ ทตฺวา กุปิตา อนตฺตมนา นิกฺกฑฺฒิสฺสติ.\csromanspacer{} เนตํ ภิกฺขเว อปฺปสนฺนานํ วา ปสาทาย ฯเปฯ.\csromanspacer{} เอวญฺจ ปน ภิกฺขเว ภิกฺขุนิโย อิมํ สิกฺขาปทํ อุทฺทิสนฺตุ\csromandash{}}

\proseitem{951}{\textbf{“ยา ปน ภิกฺขุนี ภิกฺขุนิยา อุปสฺสยํ ทตฺวา กุปิตา อนตฺตมนา นิกฺกฑฺเฒยฺย วา นิกฺกฑฺฒาเปยฺย วา, ปาจิตฺติยนฺ”}ติ.}

\proseitem{952}{\textbf{ยา ปนา}ติ ยา ยาทิสา ฯเปฯ.\csromanspacer{} \textbf{ภิกฺขุนี}ติ ฯเปฯ อยํ อิมสฺมิํ อตฺเถ อธิปฺเปตา ภิกฺขุนีติ.}

\prose{\textbf{ภิกฺขุนิยา}ติ อญฺญาย ภิกฺขุนิยา.\csromanspacer{} \textbf{อุปสฺสโย} นาม กวาฏพทฺโธ วุจฺจติ.\csromanspacer{} \textbf{ทตฺวา}ติ สยํ ทตฺวา.\csromanspacer{} \textbf{กุปิตา อนตฺตมนา}ติ อนภิรทฺธา อาหตจิตฺตา ขิลชาตา.}

\prose{\textbf{นิกฺกฑฺเฒยฺยา}ติ คพฺเภ คเหตฺวา ปมุขํ นิกฺกฑฺฒติ, อาปตฺติ ปาจิตฺติยสฺส.\csromanspacer{} ปมุเข คเหตฺวา พหิ นิกฺกฑฺฒติ, อาปตฺติ ปาจิตฺติยสฺส.\csromanspacer{} เอเกน ปโยเคน พหุเกปิ ทฺวาเร อติกฺกาเมติ, อาปตฺติ ปาจิตฺติยสฺส.}

\prose{\textbf{นิกฺกฑฺฒาเปยฺยา}ติ อญฺญํ อาณาเปติ, อาปตฺติ ปาจิตฺติยสฺส\footnote{เอตฺถ “อาปตฺติ ทุกฺกฏสฺสา”ติ สพฺพตฺถ ทิสฺสติ, ภิกฺขุวิภงฺเค ปน อีทิเสสุ อาณตฺติสหคตสิกฺขาปเทสุ “อญฺญํ อาณาเปติ อาปตฺติ ปาจิตฺติยสฺสา”ติ อาคตตฺตา ตมนุวตฺติตฺวา วิโสธิตมิทํ.}.\csromanspacer{} สกิํ อาณตฺตา พหุเกปิ ทฺวาเร อติกฺกาเมติ, อาปตฺติ ปาจิตฺติยสฺส.}

\proseitem{935}{อุปสมฺปนฺนาย อุปสมฺปนฺนสญฺญา อุปสฺสยํ ทตฺวา กุปิตา อนตฺตมนา นิกฺกฑฺฒติ วา นิกฺกฑฺฒาเปติ วา, อาปตฺติ ปาจิตฺติยสฺส.\csromanspacer{} อุปสมฺปนฺนาย เวมติกา อุปสฺสยํ ทตฺวา กุปิตา อนตฺตมนา นิกฺกฑฺฒติ \csromanfolio{385}วา นิกฺกฑฺฒาเปติ วา, อาปตฺติ ปาจิตฺติยสฺส.\csromanspacer{} อุปสมฺปนฺนาย อนุปสมฺปนฺนสญฺญา อุปสฺสยํ ทตฺวา กุปิตา อนตฺตมนา นิกฺกฑฺฒติ วา นิกฺกฑฺฒาเปติ วา, อาปตฺติ ปาจิตฺติยสฺส.}

\prose{ตสฺสา ปริกฺขารํ นิกฺกฑฺฒติ วา นิกฺกฑฺฒาเปติ วา, อาปตฺติ ทุกฺกฏสฺส.\csromanspacer{} อกวาฏพทฺธา นิกฺกฑฺฒติ วา นิกฺกฑฺฒาเปติ วา, อาปตฺติ ทุกฺกฏสฺส.\csromanspacer{} ตสฺสา ปริกฺขารํ นิกฺกฑฺฒติ วา นิกฺกฑฺฒาเปติ วา, อาปตฺติ ทุกฺกฏสฺส.\csromanspacer{} อนุปสมฺปนฺนํ กวาฏพทฺธา วา อกวาฏพทฺธา วา นิกฺกฑฺฒติ วา นิกฺกฑฺฒาเปติ วา, อาปตฺติ ทุกฺกฏสฺส.\csromanspacer{} ตสฺสา ปริกฺขารํ นิกฺกฑฺฒติ วา นิกฺกฑฺฒาเปติ วา, อาปตฺติ ทุกฺกฏสฺส.\csromanspacer{} อนุปสมฺปนฺนาย อุปสมฺปนฺนสญฺญา, อาปตฺติ ทุกฺกฏสฺส.\csromanspacer{} อนุปสมฺปนฺนาย เวมติกา, อาปตฺติ ทุกฺกฏสฺส.\csromanspacer{} อนุปสมฺปนฺนาย อนุปสมฺปนฺนสญฺญา, อาปตฺติ ทุกฺกฏสฺส.}

\proseitemwithcloserruled{954}{อนาปตฺติ อลชฺชินิํ นิกฺกฑฺฒติ วา นิกฺกฑฺฒาเปติ วา, ตสฺสา ปริกฺขารํ นิกฺกฑฺฒติ วา นิกฺกฑฺฒาเปติ วา, อุมฺมตฺติกํ นิกฺกฑฺฒติ วา นิกฺกฑฺฒาเปติ วา, ตสฺสา ปริกฺขารํ นิกฺกฑฺฒติ วา นิกฺกฑฺฒาเปติ วา, ภณฺฑนการิกํ กลหการิกํ วิวาทการิกํ ภสฺสการิกํ สํเฆ อธิกรณการิกํ นิกฺกฑฺฒติ วา นิกฺกฑฺฒาเปติ วา, ตสฺสา ปริกฺขารํ นิกฺกฑฺฒติ วา นิกฺกฑฺฒาเปติ วา, อนฺเตวาสินิํ วา สทฺธิวิหารินิํ วา น สมฺมา วตฺตนฺติํ นิกฺกฑฺฒติ วา นิกฺกฑฺฒาเปติ วา, ตสฺสา ปริกฺขารํ นิกฺกฑฺฒติ วา นิกฺกฑฺฒาเปติ วา, อุมฺมตฺติกาย อาทิกมฺมิกายาติ.}{ปญฺจมสิกฺขาปทํ นิฏฺฐิตํ.}
\csromanmatikaanchor{mtk.188}
\csromantocmarkref{subsection}{6. ฉฏฺฐสิกฺขาปท}{mtk.188}
\bookmark[dest={mtk.188},level=2]{6. ฉฏฺฐสิกฺขาปท}
\csromanheader{2}{4. \textbf{ตุวฏฺฏวคฺค 6. ฉฏฺฐสิกฺขาปท}}
\proseitem{955}{เตน สมเยน พุทฺโธ ภควา สาวตฺถิยํ วิหรติ เชตวเน อนาถปิณฺฑิกสฺส อาราเม.\csromanspacer{} เตน โข ปน สมเยน จณฺฑกาฬี ภิกฺขุนี สํสฏฺฐา วิหรติ คหปตินาปิ คหปติปุตฺเตนปิ.\csromanspacer{} ยา ตา ภิกฺขุนิโย อปฺปิจฺฉา ฯเปฯ ตา อุชฺฌายนฺติ ขิยฺยนฺติ วิปาเจนฺติ “กถํ หิ นาม อยฺยา จณฺฑกาฬี สํสฏฺฐา วิหริสฺสติ คหปตินาปิ คหปติปุตฺเตนปี”ติ ฯเปฯ.\csromanspacer{} สจฺจํ กิร ภิกฺขเว จณฺฑกาฬี ภิกฺขุนี สํสฏฺฐา วิหรติ คหปตินาปิ \csromanfolio{386}คหปติปุตฺเตนปีติ.\csromanspacer{} สจฺจํ ภควาติ.\csromanspacer{} วิครหิ พุทฺโธ ภควา ฯเปฯ.\csromanspacer{} กถํ หิ นาม ภิกฺขเว จณฺฑกาฬี ภิกฺขุนี สํสฏฺฐา วิหริสฺสติ คหปตินาปิ คหปติปุตฺเตนปิ.\csromanspacer{} เนตํ ภิกฺขเว อปฺปสนฺนานํ วา ปสาทาย ฯเปฯ.\csromanspacer{} เอวญฺจ ปน ภิกฺขเว ภิกฺขุนิโย อิมํ สิกฺขาปทํ อุทฺทิสนฺตุ\csromandash{}}

\proseitem{956}{\textbf{“ยา ปน ภิกฺขุนี สํสฏฺฐา วิหเรยฺย คหปตินา วา คหปติปุตฺเตน วา, สา ภิกฺขุนี ภิกฺขุนีหิ เอวมสฺส วจนียา ‘มายฺเย สํสฏฺฐา วิหริ คหปตินาปิ คหปติปุตฺเตนปิ, วิวิจฺจายฺเย วิเวกญฺเญว ภคินิยา สํโฆ วณฺเณตี’ติ.}\csromanspacer{} เอวญฺจ\footnote{เอวญฺจ ปน (ก)} \textbf{สา ภิกฺขุนี ภิกฺขุนีหิ วุจฺจมานา ตเถว ปคฺคเณฺหยฺย, สา ภิกฺขุนี ภิกฺขุนีหิ ยาวตติยํ สมนุภาสิตพฺพา ตสฺส ปฏินิสฺสคฺคาย, ยาวตติยญฺเจ สมนุภาสียมานา ตํ ปฏินิสฺสชฺเชยฺย, อิจฺเจตํ กุสลํ, โน เจ ปฏินิสฺสชฺเชยฺย, ปาจิตฺติยนฺ”}ติ.}

\proseitem{957}{\textbf{ยา ปนา}ติ ยา ยาทิสา ฯเปฯ.\csromanspacer{} \textbf{ภิกฺขุนี}ติ ฯเปฯ อยํ อิมสฺมิํ อตฺเถ อธิปฺเปตา ภิกฺขุนีติ.}

\prose{\textbf{สํสฏฺฐา} นาม อนนุโลมิเกน กายิกวาจสิเกน สํสฏฺฐา.}

\prose{\textbf{คหปติ} นาม โย โกจิ อคารํ อชฺฌาวสติ.}

\prose{\textbf{คหปติปุตฺโต} นาม เย เกจิ\footnote{โย โกจิ (ก)} ปุตฺตภาตโร.}

\prose{\textbf{สา ภิกฺขุนี}ติ ยา สา สํสฏฺฐา ภิกฺขุนี.}

\prose{\textbf{ภิกฺขุนีหี}ติ อญฺญาหิ ภิกฺขุนีหิ.\csromanspacer{} ยา ปสฺสนฺติ ยา สุณนฺติ, ตาหิ วตฺตพฺพา “มายฺเย สํสฏฺฐา วิหริ คหปตินาปิ คหปติปุตฺเตนปิ,วิวิจฺจายฺเย วิเวกญฺเญว ภคินิยา สํโฆ วณฺเณตีติ”.\csromanspacer{} ทุติยมฺปิ วตฺตพฺพา.\csromanspacer{} ตติยมฺปิ วตฺตพฺพา.\csromanspacer{} สเจ ปฏินิสฺสชฺชติ, อิจฺเจตํ กุสลํ, โน เจ ปฏินิสฺสชฺชติ, อาปตฺติ ทุกฺกฏสฺส.\csromanspacer{} สุตฺวา น วทนฺติ, อาปตฺติ ทุกฺกฏสฺส.\csromanspacer{} สา ภิกฺขุนี สํฆมชฺฌมฺปิ อากฑฺฒิตฺวา วตฺตพฺพา “มายฺเย สํสฏฺฐา วิหริ คหปตินาปิ คหปติปุตฺเตนปิ, วิวิจฺจายฺเย วิเวกญฺเญว ภคินิยา สํโฆ วณฺเณตี”ติ.\csromanspacer{} ทุติยมฺปิ วตฺตพฺพา.\csromanspacer{} ตติยมฺปิ วตฺตพฺพา.\csromanspacer{} สเจ ปฏินิสฺสชฺชติ, อิจฺเจตํ กุสลํ, โน เจ ปฏินิสฺสชฺชติ, อาปตฺติ ทุกฺกฏสฺส.\csromanspacer{} สา ภิกฺขุนี สมนุภาสิตพฺพา, เอวญฺจ ปน ภิกฺขเว สมนุภาสิตพฺพา.\csromanspacer{} พฺยตฺตาย ภิกฺขุนิยา ปฏิพลาย สํโฆ ญาเปตพฺโพ\csromandash{}}

\proseitem{958}{\csromanfolio{387}“สุณาตุ เม อยฺเย สํโฆ, อยํ อิตฺถนฺนามา ภิกฺขุนี สํสฏฺฐา วิหรติ คหปตินาปิ คหปติปุตฺเตนปิ, สา ตํ วตฺถุํ น ปฏินิสฺสชฺชติ, ยทิ สํฆสฺส ปตฺตกลฺลํ, สํโฆ อิตฺถนฺนามํ ภิกฺขุนิํ สมนุภาเสยฺย ตสฺส วตฺถุสฺส ปฏินิสฺสคฺคาย, เอสา ญตฺติ.}

\prose{สุณาตุ เม อยฺเย สํโฆ, อยํ อิตฺถนฺนามา ภิกฺขุนี สํสฏฺฐา วิหรติ คหปตินาปิ คหปติปุตฺเตนปิ, สา ตํ วตฺถุํ น ปฏินิสฺสชฺชติ, สํโฆ อิตฺถนฺนามํ ภิกฺขุนิํ สมนุภาสติ ตสฺส วตฺถุสฺส ปฏินิสฺสคฺคาย, ยสฺสา อยฺยาย ขมติ อิตฺถนฺนามาย ภิกฺขุนิยา สมนุภาสนา ตสฺส วตฺถุสฺส ปฏินิสฺสคฺคาย, สา ตุณฺหสฺส, ยสฺสา นกฺขมติ, สา ภาเสยฺย.}

\prose{ทุติยมฺปิ เอตมตฺถํ วทามิ ฯเปฯ.\csromanspacer{} ตติยมฺปิ เอตมตฺถํ วทามิ ฯเปฯ.}

\prose{สมนุภฏฺฐา สํเฆน อิตฺถนฺนามา ภิกฺขุนี ตสฺส วตฺถุสฺส ปฏินิสฺสคฺคาย, ขมติ สํฆสฺส, ตสฺมา ตุณฺหี, เอวเมตํ ธารยามี”ติ.}

\prose{ญตฺติยา ทุกฺกฏํ.\csromanspacer{} ทฺวีหิ กมฺมวาจาหิ ทุกฺกฏา.\csromanspacer{} กมฺมวาจาปริโยสาเน อาปตฺติ ปาจิตฺติยสฺส.}

\proseitem{959}{ธมฺมกมฺเม ธมฺมกมฺมสญฺญา น ปฏินิสฺสชฺชติ, อาปตฺติ ปาจิตฺติยสฺส.\csromanspacer{} ธมฺมกมฺเม เวมติกา น ปฏินิสฺสชฺชติ, อาปตฺติ ปาจิตฺติยสฺส.\csromanspacer{} ธมฺมกมฺเม อธมฺมกมฺมสญฺญา น ปฏินิสฺสชฺชติ, อาปตฺติ ปาจิตฺติยสฺส.}

\prose{อธมฺมกมฺเม ธมฺมกมฺมสญฺญา, อาปตฺติ ทุกฺกฏสฺส.\csromanspacer{} อธมฺมกมฺเม เวมติกา, อาปตฺติ ทุกฺกฏสฺส.\csromanspacer{} อธมฺมกมฺเม อธมฺมกมฺมสญฺญา, อาปตฺติ ทุกฺกฏสฺส.}

\proseitemwithcloserruled{960}{อนาปตฺติ อสมนุภาสนฺติยา ปฏินิสฺสชฺชนฺติยา อุมฺมตฺติกาย อาทิกมฺมิกายาติ.}{ฉฏฺฐสิกฺขาปทํ นิฏฺฐิตํ.}
\csromanmatikaanchor{mtk.189}
\csromantocmarkref{subsection}{7. สตฺตมสิกฺขาปท}{mtk.189}
\bookmark[dest={mtk.189},level=2]{7. สตฺตมสิกฺขาปท}
\csromanheader{2}{4. \textbf{ตุวฏฺฏวคฺค 7. สตฺตมสิกฺขาปท}}
\proseitem{961}{เตน สมเยน พุทฺโธ ภควา สาวตฺถิยํ วิหรติ เชตวเน อนาถปิณฺฑิกสฺส อาราเม.\csromanspacer{} เตน โข ปน สมเยน ภิกฺขุนิโย อนฺโตรฏฺเฐ สาสงฺกสมฺมเต สปฺปฏิภเย อสตฺถิกา จาริกํ จรนฺติ, \csromanfolio{388}\textbf{ธุตฺตา ทูเสนฺติ.}\csromanspacer{}\textbf{ ยา ตา ภิกฺขุนิโย อปฺปิจฺฉา ฯเปฯ ตา อุชฺฌายนฺติ ขิยฺยนฺติ} วิปาเจนฺติ “กถํ หิ นาม ภิกฺขุนิโย อนฺโตรฏฺเฐ สาสงฺกสมฺมเต สปฺปฏิภเย อสตฺถิกา จาริกํ จริสฺสนฺตี”ติ ฯเปฯ.\csromanspacer{} สจฺจํ กิร ภิกฺขเว ภิกฺขุนิโย อนฺโตรฏฺเฐ สาสงฺกสมฺมเต สปฺปฏิภเย อสตฺถิกา จาริกํ จรนฺตีติ.\csromanspacer{} สจฺจํ ภควาติ.\csromanspacer{} วิครหิ พุทฺโธ ภควา ฯเปฯ.\csromanspacer{} กถํ หิ นาม ภิกฺขเว ภิกฺขุนิโย อนฺโตรฏฺเฐ สาสงฺกสมฺมเต สปฺปฏิภเย อสตฺถิกา จาริกํ จริสฺสนฺติ.\csromanspacer{} เนตํ ภิกฺขเว อปฺปสนฺนานํ วา ปสาทาย ฯเปฯ.\csromanspacer{} เอวญฺจ ปน ภิกฺขเว ภิกฺขุนิโย อิมํ สิกฺขาปทํ อุทฺทิสนฺตุ\csromandash{}}

\proseitem{962}{\textbf{“ยา ปน ภิกฺขุนี อนฺโตรฏฺเฐ สาสงฺกสมฺมเต สปฺปฏิภเย อสตฺถิกา จาริกํ จเรยฺย, ปาจิตฺติยนฺ”}ติ.}

\proseitem{963}{\textbf{ยา ปนา}ติ ยา ยาทิสา ฯเปฯ.\csromanspacer{} \textbf{ภิกฺขุนี}ติ ฯเปฯ อยํ อิมสฺมิํ อตฺเถ อธิปฺเปตา ภิกฺขุนีติ.}

\prose{\textbf{อนฺโตรฏฺเฐ}ติ ยสฺส วิชิเต วิหรติ, ตสฺส รฏฺเฐ.}

\prose{\textbf{สาสงฺกํ} นาม ตสฺมิํ มคฺเค โจรานํ นิวิฏฺโฐกาโส ทิสฺสติ, ภุตฺโตกาโส ทิสฺสติ, ฐิโตกาโส ทิสฺสติ, นิสินฺโนกาโส ทิสฺสติ, นิปนฺโนกาโส ทิสฺสติ.}

\prose{\textbf{สปฺปฏิภยํ} นาม ตสฺมิํ มคฺเค โจเรหิ มนุสฺสา หตา ทิสฺสนฺติ, วิลุตฺตา ทิสฺสนฺติ, อาโกฏิตา ทิสฺสนฺติ.}

\prose{\textbf{อสตฺถิกา} นาม วินา สตฺเถน.}

\prose{\textbf{จาริกํ จเรยฺยา}ติ กุกฺกุฏสมฺปาเต คาเม คามนฺตเร คามนฺตเร อาปตฺติ ปาจิตฺติยสฺส.\csromanspacer{} อคามเก อรญฺเญ อทฺธโยชเน อทฺธโยชเน อาปตฺติ ปาจิตฺติยสฺส.}

\proseitemwithcloser{964}{อนาปตฺติ สตฺเถน สห คจฺฉติ, เขเม อปฺปฏิภเย คจฺฉติ, อาปทาสุ อุมฺมตฺติกาย อาทิกมฺมิกายติ.}{สตฺตมสิกฺขาปทํ นิฏฺฐิตํ.}
\csromanmatikaanchor{mtk.190}
\csromantocmarkref{subsection}{8. อฏฺฐมสิกฺขาปท}{mtk.190}
\bookmark[dest={mtk.190},level=2]{8. อฏฺฐมสิกฺขาปท}
\csromanheader{2}{4. \textbf{ตุวฏฺฏวคฺค 8. อฏฺฐมสิกฺขาปท}}
\proseitem{965}{\csromanfolio{389}เตน สมเยน พุทฺโธ ภควา สาวตฺถิยํ วิหรติ เชตวเน อนาถปิณฺฑิกสฺส อาราเม.\csromanspacer{} เตน โข ปน สมเยน ภิกฺขุนิโย ติโรรฏฺเฐ สาสงฺกสมฺมเต สปฺปฏิภเย อสตฺถิกา จาริกํ จรนฺติ, ธุตฺตา ทูเสนฺติ.\csromanspacer{} ยา ตา ภิกฺขุนิโย อปฺปิจฺฉา ฯเปฯ ตา อุชฺฌายนฺติ ขิยฺยนฺติ วิปาเจนฺติ “กถํ หิ นาม ภิกฺขุนิโย ติโรรฏฺเฐ สาสงฺกสมฺมเต สปฺปฏิภเย อสตฺถิกา จาริกํ จริสฺสนฺตี”ติ ฯเปฯ.\csromanspacer{} สจฺจํ กิร ภิกฺขเว ภิกฺขุนิโย ติโรรฏฺเฐ สาสงฺกสมฺมเต สปฺปฏิภเย อสตฺถิกา จาริกํ จรนฺตีติ.\csromanspacer{} สจฺจํ ภควาติ.\csromanspacer{} วิครหิ พุทฺโธ ภควา ฯเปฯ.\csromanspacer{} กถํ หิ นาม ภิกฺขเว ภิกฺขุนิโย ติโรรฏฺเฐ สาสงฺกสมฺมเต สปฺปฏิภเย อสตฺถิกา จาริกํ จริสฺสนฺติ.\csromanspacer{} เนตํ ภิกฺขเว อปฺปสนฺนานํ วา ปสาทาย ฯเปฯ.\csromanspacer{} เอวญฺจ ปน ภิกฺขเว ภิกฺขุนิโย อิมํ สิกฺขาปทํ อุทฺทิสนฺตุ\csromandash{}}

\proseitem{966}{\textbf{“ยา ปน ภิกฺขุนี ติโรรฏฺเฐ สาสงฺกสมฺมเต สปฺปฏิภเย อสตฺถิกา จาริกํ จเรยฺย, ปาจิตฺติยนฺ”}ติ.}

\proseitem{967}{\textbf{ยา ปนา}ติ ยา ยาทิสา ฯเปฯ.\csromanspacer{} \textbf{ภิกฺขุนี}ติ ฯเปฯ อยํ อิมสฺมิํ อตฺเถ อธิปฺเปตา ภิกฺขุนีติ.}

\prose{\textbf{ติโรรฏฺเฐ}ติ ยสฺส วิชิเต วิหรติ, ตํ ฐเปตฺวา อญฺญสฺส รฏฺเฐ.}

\prose{\textbf{สาสงฺกํ} นาม ตสฺมิํ มคฺเค โจรานํ นิวิฏฺโฐกาโส ทิสฺสติ, ภุตฺโตกาโส ทิสฺสติ, ฐิโตกาโส ทิสฺสติ, นิสินฺโนกาโส ทิสฺสติ, นิปนฺโนกาโส ทิสฺสติ.}

\prose{\textbf{สปฺปฏิภยํ} นาม ตสฺมิํ มคฺเค โจเรหิ มนุสฺสา หตา ทิสฺสนฺติ, วิลุตฺตา ทิสฺสนฺติ, อาโกฏิตา ทิสฺสนฺติ.}

\prose{\textbf{อสตฺถิกา} นาม วินา สตฺเถน.}

\prose{\textbf{จาริกํ จเรยฺยา}ติ กุกฺกุฏสมฺปาเต คาเม คามนฺตเร คามนฺตเร อาปตฺติ ปาจิตฺติยสฺส.\csromanspacer{} อคามเก อรญฺเญ อทฺธโยชเน อทฺธโยชเน อาปตฺติ ปาจิตฺติยสฺส.}

\proseitemwithcloser{968}{อนาปตฺติ สตฺเถน สห คจฺฉติ, เขเม อปฺปฏิภเย คจฺฉติ, อาปทาสุ อุมฺมตฺติกาย อาทิกมฺมิกายาติ.}{อฏฺฐมสิกฺขาปทํ นิฏฺฐิตํ.}
\csromanmatikaanchor{mtk.191}
\csromantocmarkref{subsection}{9. นวมสิกฺขาปท}{mtk.191}
\bookmark[dest={mtk.191},level=2]{9. นวมสิกฺขาปท}
\csromanheader{2}{4. \textbf{ตุวฏฺฏวคฺค 9. นวมสิกฺขาปท}}
\proseitem{969}{\csromanfolio{390}เตน สมเยน พุทฺโธ ภควา ราชคเห วิหรติ เวฬุวเน กลนฺทกนิวาเป.\csromanspacer{} เตน โข ปน สมเยน ภิกฺขุนิโย อนฺโตวสฺสํ จาริกํ จรนฺติ.\csromanspacer{} มนุสฺสา อุชฺฌายนฺติ ขิยฺยนฺติ วิปาเจนฺติ “กถํ หิ นาม ภิกฺขุนิโย อนฺโตวสฺสํ จาริกํ จริสฺสนฺติ หริตานิ ติณานิ จ สมฺมทฺทนฺตา เอกินฺทฺริยํ ชีวํ วิเหเฐนฺตา พหู ขุทฺทเก ปาเณ สงฺฆาตํ อาปาเทนฺตา”ติ.\csromanspacer{} อสฺโสสุํ โข ภิกฺขุนิโย เตสํ มนุสฺสานํ อุชฺฌายนฺตานํ ขิยฺยนฺตานํ วิปาเจนฺตานํ.\csromanspacer{} ยา ตา ภิกฺขุนิโย อปฺปิจฺฉา ฯเปฯ ตา อุชฺฌายนฺติ ขิยฺยนฺติ วิปาเจนฺติ “กถํ หิ นาม ภิกฺขุนิโย อนฺโตวสฺสํ จาริกํ จริสฺสนฺตี”ติ - ป-.\csromanspacer{} สจฺจํ กิร ภิกฺขเว ภิกฺขุนิโย อนฺโตวสฺสํ จาริกํ จรนฺตีติ.\csromanspacer{} สจฺจํ ภควาติ.\csromanspacer{} วิครหิ พุทฺโธ ภควา ฯเปฯ.\csromanspacer{} กถํ หิ นาม ภิกฺขเว ภิกฺขุนิโย อนฺโตวสฺสํ จาริกํ จริสฺสนฺติ.\csromanspacer{} เนตํ ภิกฺขเว อปฺปสนฺนานํ วา ปสาทาย ฯเปฯ.\csromanspacer{} เอวญฺจ ปน ภิกฺขเว ภิกฺขุนิโย อิมํ สิกฺขาปทํ อุทฺทิสนฺตุ\csromandash{}}

\proseitem{970}{\textbf{“ยา ปน ภิกฺขุนี อนฺโตวสฺสํ จาริกํ จเรยฺย, ปาจิตฺติยนฺ”}ติ.}

\proseitem{971}{\textbf{ยา ปนา}ติ ยา ยาทิสา ฯเปฯ.\csromanspacer{} \textbf{ภิกฺขุนี}ติ ฯเปฯ อยํ อิมสฺมิํ อตฺเถ อธิปฺเปตา ภิกฺขุนีติ.}

\prose{\textbf{อนฺโตวสฺสนฺ}ติ ปุริมํ วา เตมาสํ ปจฺฉิมํ วา เตมาสํ อวสิตฺวา.}

\prose{\textbf{จาริกํ} \textbf{จเรยฺยา}ติ กุกฺกุฏสมฺปาเต คาเม คามนฺตเร คามนฺตเร อาปตฺติ ปาจิตฺติยสฺส.\csromanspacer{} อคามเก อรญฺเญ อทฺธโยชเน อทฺธโยชเน อาปตฺติ ปาจิตฺติยสฺส.}

\proseitemwithcloserruled{972}{อนาปตฺติ สตฺตาหกรณีเยน คจฺฉติ, เกนจิ อุพฺพาฬฺหา คจฺฉติ, อาปทาสุ อุมฺมตฺติกาย อาทิกมฺมิกายาติ.}{นวมสิกฺขาปทํ นิฏฺฐิตํ.}
\csromanmatikaanchor{mtk.192}
\csromantocmarkref{subsection}{10. ทสมสิกฺขาปท}{mtk.192}
\bookmark[dest={mtk.192},level=2]{10. ทสมสิกฺขาปท}
\csromanheader{2}{4. \textbf{ตุวฏฺฏวคฺค} 10. ทสมสิกฺขาปท}
\proseitem{973}{เตน สมเยน พุทฺโธ ภควา ราชคเห วิหรติ เวฬุวเน กลนฺทกนิวาเป.\csromanspacer{} เตน โข ปน สมเยน ภิกฺขุนิโย ตตฺเถว \csromanfolio{391}ราชคเห วสฺสํ วสนฺติ ตตฺถ เหมนฺตํ ตตฺถ คิมฺหํ.\csromanspacer{} มนุสฺสา อุชฺฌายนฺติ ขิยฺยนฺติ วิปาเจนฺติ “อาหุนฺทริกา ภิกฺขุนีนํ ทิสา อนฺธการา, น อิมาสํ ทิสา ปกฺขายนฺตี”ติ.\csromanspacer{} อสฺโสสุํ โข ภิกฺขุนิโย เตสํ มนุสฺสานํ อุชฺฌายนฺตานํ ขิยฺยนฺตานํ วิปาเจนฺตานํ.\csromanspacer{} อถ โข ตา ภิกฺขุนิโย ภิกฺขูนํ เอตมตฺถํ อาโรเจสุํ.\csromanspacer{} ภิกฺขู ภควโต เอตมตฺถํ อาโรเจสุํ.\csromanspacer{} อถ โข ภควา เอตสฺมิํ นิทาเน เอตสฺมิํ ปกรเณ ธมฺมิํ กถํ กตฺวา ภิกฺขู อามนฺเตสิ เตน หิ ภิกฺขเว ภิกฺขุนีนํ สิกฺขาปทํ ปญฺญเปสฺสามิ ทส อตฺถวเส ปฏิจฺจ สํฆสุฏฺฐุตาย ฯเปฯ สทฺธมฺมฏฺฐิติยา วินยานุคฺคหาย.\csromanspacer{} เอวญฺจ ปน ภิกฺขเว ภิกฺขุนิโย อิมํ สิกฺขาปทํ อุทฺทิสนฺตุ\csromandash{}}

\proseitem{974}{\textbf{“ยา ปน ภิกฺขุนี วสฺสํวุฏฺฐา จาริกํ น ปกฺกเมยฺย อนฺตมโส ฉปฺปญฺจโยชนานิปิ, ปาจิตฺติยนฺ”}ติ.}

\proseitem{975}{\textbf{ยา ปนา}ติ ยา ยาทิสา ฯเปฯ.\csromanspacer{} \textbf{ภิกฺขุนี}ติ ฯเปฯ อยํ อิมสฺมิํ อตฺเถ อธิปฺเปตา ภิกฺขุนีติ.}

\prose{\textbf{วสฺสํวุฏฺฐา} นาม ปุริมํ วา เตมาสํ ปจฺฉิมํ วา เตมาสํ วุฏฺฐา.}

\prose{“จาริกํ น ปกฺกมิสฺสามิ อนฺตมโส ฉปฺปญฺจโยชนานิปี”ติ ธุรํ นิกฺขิตฺตมตฺเต อาปตฺติ ปาจิตฺติยสฺส.}

\proseitemwithcloser{976}{อนาปตฺติ สติ อนฺตราเย, ปริเยสิตฺวา ทุติยิกํ ภิกฺขุนิํ น ลภติ, คิลานาย อาปทาสุ อุมฺมตฺติกาย อาทิกมฺมิกายาติ.}{ทสมสิกฺขาปทํ นิฏฺฐิตํ.}
\csromanheader{2}{ตุวฏฺฏวคฺโค จตุตฺโถ.}
\csromansectionrule
\csromanmatikaanchor{mtk.193}
\csromanmatikaanchor{mtk.194}
\csromantocmarknopage{section}{5. จิตฺตาคารวคฺค}{mtk.193}
\bookmark[dest={mtk.193},level=1]{5. จิตฺตาคารวคฺค}
\csromantocmarkref{subsection}{1. ปฐมสิกฺขาปท}{mtk.194}
\bookmark[dest={mtk.194},level=2]{1. ปฐมสิกฺขาปท}
\setcsromanheadh{5. จิตฺตาคารวคฺค}
\csromanheader{2}{5. \textbf{จิตฺตาคารวคฺค 1. ปฐมสิกฺขาปท}}
\proseitem{977}{เตน สมเยน พุทฺโธ ภควา สาวตฺถิยํ วิหรติ เชตวเน อนาถปิณฺฑิกสฺส อาราเม.\csromanspacer{} เตน โข ปน สมเยน รญฺโญ ปเสนทิสฺส โกสลสฺส อุยฺยาเน จิตฺตาคาเร ปฏิภานจิตฺตํ \csromanfolio{392}กตํ โหติ.\csromanspacer{} พหู มนุสฺสา จิตฺตาคารํ ทสฺสนาย คจฺฉนฺติ.\csromanspacer{} ฉพฺพคฺคิยาปิ ภิกฺขุนิโย จิตฺตาคารํ ทสฺสนาย อคมํสุ.\csromanspacer{} มนุสฺสา อุชฺฌายนฺติ ขิยฺยนฺติ วิปาเจนฺติ “กถํ หิ นาม ภิกฺขุนิโย จิตฺตาคารํ ทสฺสนาย คจฺฉิสฺสนฺติ เสยฺยถาปิ คิหินิโย กามโภคินิโย”ติ.\csromanspacer{} อสฺโสสุํ โข ภิกฺขุนิโย เตสํ มนุสฺสานํ อุชฺฌายนฺตานํ ขิยฺยนฺตานํ วิปาเจนฺตานํ.\csromanspacer{} ยา ตา ภิกฺขุนิโย อปฺปิจฺฉา ฯเปฯ ตา อุชฺฌายนฺติ ขิยฺยนฺติ วิปาเจนฺติ “กถํ หิ \textbf{นาม ฉพฺพคฺคิยา ภิกฺขุนิโย จิตฺตาคารํ ทสฺสนาย คจฺฉิสฺสนฺตี”ติ ฯเปฯ.}\csromanspacer{} \textbf{สจฺจํ กิร ภิกฺขเว ฉพฺพคฺคิยา ภิกฺขุนิโย จิตฺตาคารํ ทสฺสนาย} คจฺฉนฺตีติ.\csromanspacer{} สจฺจํ ภควาติ.\csromanspacer{} วิครหิ พุทฺโธ ภควา ฯเปฯ.\csromanspacer{} กถํ หิ นาม ภิกฺขเว ฉพฺพคฺคิยา ภิกฺขุนิโย จิตฺตาคารํ ทสฺสนาย คจฺฉิสฺสนฺติ.\csromanspacer{} เนตํ ภิกฺขเว อปฺปสนฺนานํ วา ปสาทาย ฯเปฯ.\csromanspacer{} เอวญฺจ ปน ภิกฺขเว ภิกฺขุนิโย อิมํ สิกฺขาปทํ อุทฺทิสนฺตุ\csromandash{}}

\proseitem{978}{\textbf{“ยา ปน ภิกฺขุนี ราชาคารํ วา จิตฺตาคารํ วา อารามํ วา อุยฺยานํ วา โปกฺขรณิํ วา ทสฺสนาย คจฺเฉยฺย, ปาจิตฺติยนฺ”}ติ.}

\proseitem{979}{\textbf{ยา ปนา}ติ ยา ยาทิสา ฯเปฯ.\csromanspacer{} \textbf{ภิกฺขุนี}ติ ฯเปฯ อยํ อิมสฺมิํ อตฺเถ อธิปฺเปตา ภิกฺขุนีติ.}

\prose{\textbf{ราชาคารํ} นาม ยตฺถ กตฺถจิ รญฺโญ กีฬิตุํ รมิตุํ กตํ โหติ.}

\prose{\textbf{จิตฺตาคารํ} นาม ยตฺถ กตฺถจิ มนุสฺสานํ กีฬิตุํ รมิตุํ กตํ โหติ.}

\prose{\textbf{อาราโม} นาม ยตฺถ กตฺถจิ มนุสฺสานํ กีฬิตุํ รมิตุํ กโต โหติ.}

\prose{\textbf{อุยฺยานํ} นาม ยตฺถ กตฺถจิ มนุสฺสานํ กีฬิตุํ รมิตุํ กตํ โหติ.}

\prose{\textbf{โปกฺขรณี} นาม ยตฺถ กตฺถจิ มนุสฺสานํ กีฬิตุํ รมิตุํ กตา โหติ.}

\proseitem{980}{ทสฺสนาย คจฺฉติ, อาปตฺติ ทุกฺกฏสฺส.\csromanspacer{} ยตฺถ ฐิตา ปสฺสติ, อาปตฺติ ปาจิตฺติยสฺส.\csromanspacer{} ทสฺสนูปจารํ วิชหิตฺวา ปุนปฺปุนํ ปสฺสติ, อาปตฺติ ปาจิตฺติยสฺส.\csromanspacer{} เอกเมกํ ทสฺสนาย คจฺฉติ, อาปตฺติ ทุกฺกฏสฺส.\csromanspacer{} ยตฺถ ฐิตา ปสฺสติ, อาปตฺติ ปาจิตฺติยสฺส.\csromanspacer{} ทสฺสนูปจารํ วิชหิตฺวา ปุนปฺปุนํ ปสฺสติ, อาปตฺติ ปาจิตฺติยสฺส.}

\proseitemwithcloserruled{981}{\csromanfolio{393}อนาปตฺติ อาราเม ฐิตา ปสฺสติ, คจฺฉนฺตี วา อาคจฺฉนฺตี วา ปสฺสติ, สติ กรณีเย คนฺตฺวา ปสฺสติ, อาปทาสุ อุมฺมตฺติกาย อาทิกมฺมิกายาติ.}{ปฐมสิกฺขาปทํ นิฏฺฐิตํ.}
\csromanmatikaanchor{mtk.195}
\csromantocmarkref{subsection}{2. ทุติยสิกฺขาปท}{mtk.195}
\bookmark[dest={mtk.195},level=2]{2. ทุติยสิกฺขาปท}
\csromanheader{2}{5. \textbf{จิตฺตาคารวคฺค 2. ทุติยสิกฺขาปท}}
\proseitem{982}{เตน สมเยน พุทฺโธ ภควา สาวตฺถิยํ วิหรติ เชตวเน อนาถปิณฺฑิกสฺส อาราเม.\csromanspacer{} เตน โข ปน สมเยน ภิกฺขุนิโย อาสนฺทิมฺปิ ปลฺลงฺกมฺปิ ปริภุญฺชนฺติ.\csromanspacer{} มนุสฺสา วิหารจาริกํ อาหิณฺฑนฺตา ปสฺสิตฺวา อุชฺฌายนฺติ ขิยฺยนฺติ วิปาเจนฺติ “กถํ หิ นาม ภิกฺขุนิโย อาสนฺทิมฺปิ ปลฺลงฺกมฺปิ ปริภุญฺชิสฺสนฺติ เสยฺยถาปิ คิหินิโย กามโภคินิโย”ติ.\csromanspacer{} อสฺโสสุํ โข ภิกฺขุนิโย เตสํ มนุสฺสานํ อุชฺฌายนฺตานํ ขิยฺยนฺตานํ วิปาเจนฺตานํ.\csromanspacer{} ยา ตา ภิกฺขุนิโย อปฺปิจฺฉา ฯเปฯ ตา อุชฺฌายนฺติ ขิยฺยนฺติ วิปาเจนฺติ “กถํ หิ นาม ภิกฺขุนิโย อาสนฺทิมฺปิ ปลฺลงฺกมฺปิ ปริภุญฺชิสฺสนฺตี”ติ ฯเปฯ.\csromanspacer{} สจฺจํ กิร ภิกฺขเว ภิกฺขุนิโย อาสนฺทิมฺปิ ปลฺลงฺกมฺปิ ปริภุญฺชนฺตีติ.\csromanspacer{} สจฺจํ ภควาติ.\csromanspacer{} วิครหิ พุทฺโธ ภควา - ป-.\csromanspacer{} กถํ หิ นาม ภิกฺขเว ภิกฺขุนิโย อาสนฺทิมฺปิ ปลฺลงฺกมฺปิ ปริภุญฺชิสฺสนฺติ.\csromanspacer{} เนตํ ภิกฺขเว อปฺปสนฺนานํ วา ปสาทาย ฯเปฯ.\csromanspacer{} เอวญฺจ ปน ภิกฺขเว ภิกฺขุนิโย อิมํ สิกฺขาปทํ อุทฺทิสนฺตุ\csromandash{}}

\proseitem{983}{\textbf{“ยา ปน ภิกฺขุนี อาสนฺทิํ วา ปลฺลงฺกํ วา ปริภุญฺเชยฺย, ปาจิตฺติยนฺ”}ติ.}

\proseitem{984}{\textbf{ยา ปนา}ติ ยา ยาทิสา ฯเปฯ.\csromanspacer{} \textbf{ภิกฺขุนี}ติ ฯเปฯ อยํ อิมสฺมิํ อตฺเถ อธิปฺเปตา ภิกฺขุนีติ.}

\prose{\textbf{อาสนฺที} นาม อติกฺกนฺตปฺปมาณา วุจฺจติ.}

\prose{\textbf{ปลฺลงฺโก} นาม อาหริเมหิ\footnote{“อสํหาริเมหิ” อิติ กงฺขาวิตรณิยา สเมติ.} วาเฬหิ กโต โหติ.}

\prose{\textbf{ปริภุญฺเชยฺยา}ติ ตสฺมิํ อภินิสีทติ วา อภินิปชฺชติ วา, อาปตฺติ}

\proseitemwithcloserruled{985}{\csromanfolio{394}อนาปตฺติ อาสนฺทิยา ปาเท ฉินฺทิตฺวา ปริภุญฺชติ, ปลฺลงฺกสฺส วาเฬ ภินฺทิตฺวา ปริภุญฺชติ, อุมฺมตฺติกาย อาทิกมฺมิกายาติ.}{ทุติยสิกฺขาปทํ นิฏฺฐิตํ.}
\csromanmatikaanchor{mtk.196}
\csromantocmarkref{subsection}{3. ตติยสิกฺขาปท}{mtk.196}
\bookmark[dest={mtk.196},level=2]{3. ตติยสิกฺขาปท}
\csromanheader{2}{5. \textbf{จิตฺตาคารวคฺค 3. ตติยสิกฺขาปท}}
\proseitem{986}{เตน สมเยน พุทฺโธ ภควา สาวตฺถิยํ วิหรติ เชตวเน อนาถปิณฺฑิกสฺส อาราเม.\csromanspacer{} เตน โข ปน สมเยน ฉพฺพคฺคิยา ภิกฺขุนิโย สุตฺตํ กนฺตนฺติ.\csromanspacer{} มนุสฺสา วิหารจาริกํ อาหิณฺฑนฺตา ปสฺสิตฺวา อุชฺฌายนฺติ ขิยฺยนฺติ วิปาเจนฺติ “กถํ หิ นาม ภิกฺขุนิโย สุตฺตํ กนฺติสฺสนฺติ เสยฺยถาปิ คิหินิโย กามโภคินิโย”ติ.\csromanspacer{} อสฺโสสุํ โข ภิกฺขุนิโย เตสํ มนุสฺสานํ อุชฺฌายนฺตานํ ขิยฺยนฺตานํ วิปาเจนฺตานํ.\csromanspacer{} ยา ตา ภิกฺขุนิโย อปฺปิจฺฉา ฯเปฯ ตา อุชฺฌายนฺติ ขิยฺยนฺติ วิปาเจนฺติ “กถํ หิ นาม ฉพฺพคฺคิยา ภิกฺขุนิโย สุตฺตํ กนฺติสฺสนฺตี”ติ - ป-.\csromanspacer{} สจฺจํ กิร ภิกฺขเว ฉพฺพคฺคิยา ภิกฺขุนิโย สุตฺตํ กนฺตนฺตีติ.\csromanspacer{} สจฺจํ ภควาติ.\csromanspacer{} วิครหิ พุทฺโธ ภควา ฯเปฯ.\csromanspacer{} กถํ หิ นาม ภิกฺขเว ฉพฺพคฺคิยา ภิกฺขุนิโย สุตฺตํ กนฺติสฺสนฺติ.\csromanspacer{} เนตํ ภิกฺขเว อปฺปสนฺนานํ วา ปสาทาย ฯเปฯ.\csromanspacer{} เอวญฺจ ปน ภิกฺขเว ภิกฺขุนิโย อิมํ สิกฺขาปทํ อุทฺทิสนฺตุ\csromandash{}}

\proseitem{987}{\textbf{“ยา ปน ภิกฺขุนี สุตฺตํ กนฺเตยฺย, ปาจิตฺติยนฺ”}ติ.}

\proseitem{988}{\textbf{ยา ปนา}ติ ยา ยาทิสา ฯเปฯ.\csromanspacer{} \textbf{ภิกฺขุนี}ติ ฯเปฯ อยํ อิมสฺมิํ อตฺเถ อธิปฺเปตา ภิกฺขุนีติ.}

\prose{\textbf{สุตฺตํ} นาม ฉ สุตฺตานิ โขมํ กปฺปาสิกํ โกเสยฺยํ กมฺพลํ สาณํ ภงฺคํ.}

\prose{\textbf{กนฺเตยฺยา}ติ สยํ กนฺตติ, ปโยเค ทุกฺกฏํ.\csromanspacer{} อุชฺชวุชฺชเว อาปตฺติ}

\proseitemwithcloser{989}{อนาปตฺติ กนฺติตสุตฺตํ กนฺตติ, อุมฺมตฺติกาย อาทิกมฺมิกายาติ.}{ตติยสิกฺขาปทํ นิฏฺฐิตํ.}
\csromanmatikaanchor{mtk.197}
\csromantocmarkref{subsection}{4. จตุตฺถสิกฺขาปท}{mtk.197}
\bookmark[dest={mtk.197},level=2]{4. จตุตฺถสิกฺขาปท}
\csromanheader{2}{5. \textbf{จิตฺตาคารวคฺค 4. จตุตฺถสิกฺขาปท}}
\proseitem{990}{\csromanfolio{395}เตน สมเยน พุทฺโธ ภควา สาวตฺถิยํ วิหรติ เชตวเน อนาถปิณฺฑิกสฺส อาราเม.\csromanspacer{} เตน โข ปน สมเยน ภิกฺขุนิโย คิหิเวยฺยาวจฺจํ กโรนฺติ.\csromanspacer{} ยา ตา ภิกฺขุนิโย อปฺปิจฺฉา ฯเปฯ ตา อุชฺฌายนฺติ ขิยฺยนฺติ วิปาเจนฺติ “กถํ หิ นาม ภิกฺขุนิโย คิหิเวยฺยาวจฺจํ กริสฺสนฺตี”ติ ฯเปฯ.\csromanspacer{} สจฺจํ กิร ภิกฺขเว ภิกฺขุนิโย คิหิเวยฺยาวจฺจํ กโรนฺตีติ.\csromanspacer{} สจฺจํ ภควาติ.\csromanspacer{} วิครหิ พุทฺโธ ภควา ฯเปฯ.\csromanspacer{} กถํ หิ นาม ภิกฺขเว ภิกฺขุนิโย คิหิเวยฺยาวจฺจํ กริสฺสนฺติ.\csromanspacer{} เนตํ ภิกฺขเว อปฺปสนฺนานํ วา ปสาทาย ฯเปฯ.\csromanspacer{} เอวญฺจ ปน ภิกฺขเว ภิกฺขุนิโย อิมํ สิกฺขาปทํ อุทฺทิสนฺตุ\csromandash{}}

\proseitem{991}{\textbf{“ยา ปน ภิกฺขุนี คิหิเวยฺยาวจฺจํ กเรยฺย, ปาจิตฺติยนฺ”}ติ.}

\proseitem{992}{\textbf{ยา ปนา}ติ ยา ยาทิสา ฯเปฯ.\csromanspacer{} \textbf{ภิกฺขุนี}ติ ฯเปฯ อยํ อิมสฺมิํ อตฺเถ อธิปฺเปตา ภิกฺขุนีติ.}

\prose{\textbf{คิหิเวยฺยาวจฺจํ} นาม อคาริกสฺส ยาคุํ วา ภตฺตํ วา ขาทนียํ วา ปจติ, สาฏกํ วา เวฐนํ วา โธวติ, อาปตฺติ ปาจิตฺติยสฺส.}

\proseitemwithcloserruled{993}{อนาปตฺติ ยาคุปาเน, สํฆภตฺเต, เจติยปูชาย, อตฺตโน เวยฺยาวจฺจกรสฺส ยาคุํ วา ภตฺตํ วา ขาทนียํ วา ปจติ สาฏกํ วา เวฐนํ วา โธวติ, อุมฺมตฺติกาย อาทิกมฺมิกายาติ.}{จตุตฺถสิกฺขาปทํ นิฏฺฐิตํ.}
\csromanmatikaanchor{mtk.198}
\csromantocmarkref{subsection}{5. ปญฺจมสิกฺขาปท}{mtk.198}
\bookmark[dest={mtk.198},level=2]{5. ปญฺจมสิกฺขาปท}
\csromanheader{2}{5. \textbf{จิตฺตาคารวคฺค 5. ปญฺจมสิกฺขาปท}}
\proseitem{994}{เตน สมเยน พุทฺโธ ภควา สาวตฺถิยํ วิหรติ เชตวเน อนาถปิณฺฑิกสฺส อาราเม.\csromanspacer{} เตน โข ปน สมเยน อญฺญตรา ภิกฺขุนี ถุลฺลนนฺทํ ภิกฺขุนิํ อุปสงฺกมิตฺวา เอตทโวจ “เอหายฺเย, อิมํ อธิกรณํ วูปสเมหี”ติ.\csromanspacer{} ถุลฺลนนฺทา ภิกฺขุนี “สาธู”ติ ปฏิสฺสุณิตฺวา เนว วูปสเมติ, น วูปสมาย อุสฺสุกฺกํ กโรติ.\csromanspacer{} อถ โข สา \csromanfolio{396}ภิกฺขุนี ภิกฺขุนีนํ เอตมตฺถํ อาโรเจสิ.\csromanspacer{} ยา ตา ภิกฺขุนิโย อปฺปิจฺฉา - ป- ตา อุชฺฌายนฺติ ขิยฺยนฺติ วิปาเจนฺติ “กถํ หิ นาม อยฺยา ถุลฺลนนฺทา ภิกฺขุนิยา ‘เอหายฺเย อิมํ อธิกรณํ วูปสเมหี’ติ วุจฺจมานา ‘สาธู’ติ ปฏิสฺสุณิตฺวา เนว วูปสเมสฺสติ, น วูปสมาย อุสฺสุกฺกํ กริสฺสตี”ติ ฯเปฯ.\csromanspacer{} สจฺจํ กิร ภิกฺขเว ถุลฺลนนฺทา ภิกฺขุนี ภิกฺขุนิยา “เอหายฺเย, อิมํ อธิกรณํ วูปสเมหี”ติ วุจฺจมานา “สาธู”ติ ปฏิสฺสุณิตฺวา เนว วูปสเมติ, น วูปสมาย อุสฺสุกฺกํ กโรตีติ.\csromanspacer{} สจฺจํ ภควาติ.\csromanspacer{} วิครหิ พุทฺโธ ภควา ฯเปฯ.\csromanspacer{} กถํ หิ นาม ภิกฺขเว ถุลฺลนนฺทา ภิกฺขุนี ภิกฺขุนิยา “เอหายฺเย, อิมํ อธิกรณํ วูปสเมหี”ติ วุจฺจมานา “สาธู”ติ ปฏิสฺสุณิตฺวา เนว วูปสเมสฺสติ, น วูปสมาย อุสฺสุกฺกํ กริสฺสติ.\csromanspacer{} เนตํ ภิกฺขเว อปฺปสนฺนานํ วา ปสาทาย ฯเปฯ.\csromanspacer{} เอวญฺจ ปน ภิกฺขเว ภิกฺขุนิโย อิมํ สิกฺขาปทํ อุทฺทิสนฺตุ\csromandash{}}

\proseitem{995}{\textbf{“ยา ปน ภิกฺขุนี ภิกฺขุนิยา ‘เอหายฺเย อิมํ อธิกรณํ วูปสเมหี’ติ วุจฺจมานา ‘สาธู’ติ ปฏิสฺสุณิตฺวา สา ปจฺฉา อนนฺตรายิกินี เนว วูปสเมยฺย น วูปสมาย อุสฺสุกฺกํ กเรยฺย, ปาจิตฺติยนฺ”}ติ.}

\proseitem{996}{\textbf{ยา ปนา}ติ ยา ยาทิสา ฯเปฯ.\csromanspacer{} \textbf{ภิกฺขุนี}ติ ฯเปฯ อยํ อิมสฺมิํ อตฺเถ อธิปฺเปตา ภิกฺขุนีติ.}

\prose{\textbf{ภิกฺขุนิยา}ติ อญฺญาย ภิกฺขุนิยา.}

\prose{\textbf{อธิกรณ} นาม จตฺตาริ อธิกรณานิ วิวาทาธิกรณํ อนุวาทาธิกรณํ อาปตฺตาธิกรณํ กิจฺจาธิกรณํ.}

\prose{\textbf{เอหายฺเย อิมํ อธิกรณํ วูปสเมหี}ติ เอหายฺเย อิมํ อธิกรณํ วินิจฺเฉหิ.}

\prose{\textbf{สา ปจฺฉา อนนฺตรายิกินี}ติ อสติ อนฺตราเย.}

\prose{\textbf{เนว วูปสเมยฺยา}ติ น สยํ วูปสเมยฺย.}

\prose{\textbf{น วูปสมาย อุสฺสุกฺกํ กเรยฺยา}ติ น อญฺญํ อาณาเปยฺย, “เนว วูปสเมสฺสามิ, น วูปสมาย อุสฺสุกฺกํ กริสฺสามี”ติ ธุรํ นิกฺขิตฺตมตฺเต อาปตฺติ ปาจิตฺติยสฺส.}

\proseitem{997}{\csromanfolio{397}อุปสมฺปนฺนาย อุปสมฺปนฺนสญฺญา อธิกรณํ เนว วูปสเมติ, น วูปสมาย อุสฺสุกฺกํ กโรติ, อาปตฺติ ปาจิตฺติยสฺส.\csromanspacer{} อุปสมฺปนฺนาย เวมติกา อธิกรณํ เนว วูปสเมติ, น วูปสมาย อุสฺสุกฺกํ กโรติ, อาปตฺติ ปาจิตฺติยสฺส.\csromanspacer{} อุปสมฺปนฺนาย อนุปสมฺปนฺนสญฺญา อธิกรณํ เนว วูปสเมติ, น วูปสมาย อุสฺสุกฺกํ กโรติ, อาปตฺติ}

\prose{อนุปสมฺปนฺนาย อธิกรณํ เนว วูปสเมติ, น วูปสมาย อุสฺสุกฺกํ กโรติ, อาปตฺติ ทุกฺกฏสฺส.\csromanspacer{} อนุปสมฺปนฺนาย อุปสมฺปนฺนสญฺญา, อาปตฺติ ทุกฺกฏสฺส.\csromanspacer{} อนุปสมฺปนฺนาย เวมติกา, อาปตฺติ ทุกฺกฏสฺส.\csromanspacer{} อนุปสมฺปนฺนาย อนุปสมฺปนฺนสญฺญา, อาปตฺติ ทุกฺกฏสฺส.}

\proseitemwithcloserruled{998}{อนาปตฺติ สติ อนฺตราเย, ปริเยสิตฺวา น ลภติ, คิลานาย อาปทาสุ อุมฺมตฺติกาย อาทิกมฺมิกายาติ.}{ปญฺจมสิกฺขาปทํ นิฏฺฐิตํ.}
\csromanmatikaanchor{mtk.199}
\csromantocmarkref{subsection}{6. ฉฏฺฐสิกฺขาปท}{mtk.199}
\bookmark[dest={mtk.199},level=2]{6. ฉฏฺฐสิกฺขาปท}
\csromanheader{2}{5. \textbf{จิตฺตาคารวคฺค 6. ฉฏฺฐสิกฺขาปท}}
\proseitem{999}{เตน สมเยน พุทฺโธ ภควา สาวตฺถิยํ วิหรติ เชตวเน อนาถปิณฺฑิกสฺส อาราเม.\csromanspacer{} เตน โข ปน สมเยน ถุลฺลนนฺทา ภิกฺขุนี นฏานมฺปิ นฏกานมฺปิ ลงฺฆกานมฺปิ โสกชฺฌายิกานมฺปิ กุมฺภถูณิกานมฺปิ สหตฺถา ขาทนียํ โภชนียํ เทติ “มยฺหํ ปริสติ วณฺณํ ภาสถา”ติ.\csromanspacer{} นฏาปิ นฏกาปิ ลงฺฆกาปิ โสกชฺฌายิกาปิ กุมฺภถูณิกาปิ ถุลฺลนนฺทาย ภิกฺขุนิยา ปริสติ วณฺณํ ภาสนฺติ “อยฺยา ถุลฺลนนฺทา พหุสฺสุตา ภาณิกา วิสารทา ปฏฺฏา ธมฺมิํ กถํ กาตุํ, เทถ อยฺยาย, กโรถ อยฺยายา”ติ.\csromanspacer{} ยา ตา ภิกฺขุนิโย อปฺปิจฺฉา ฯเปฯ ตา อุชฺฌายนฺติ ขิยฺยนฺติ วิปาเจนฺติ “กถํ หิ นาม อยฺยา ถุลฺลนนฺทา อคาริกสฺส สหตฺถา ขาทนียํ โภชนียํ ทสฺสตี”ติ ฯเปฯ.\csromanspacer{} สจฺจํ กิร ภิกฺขเว ถุลฺลนนฺทา ภิกฺขุนี อคาริกสฺส สหตฺถา ขาทนียํ โภชนียํ เทตีติ.\csromanspacer{} สจฺจํ ภควาติ.\csromanspacer{} วิครหิ พุทฺโธ ภควา ฯเปฯ.\csromanspacer{} กถํ หิ นาม ภิกฺขเว ถุลฺลนนฺทา ภิกฺขุนี อคาริกสฺส สหตฺถา ขาทนียํ โภชนียํ ทสฺสติ.\csromanspacer{} เนตํ ภิกฺขเว อปฺปสฺสนฺนานํ วา ปสาทาย ฯเปฯ.\csromanspacer{} เอวญฺจ ปน ภิกฺขเว ภิกฺขุนิโย อิมํ สิกฺขาปทํ อุทฺทิสนฺตุ\csromandash{}}

\proseitem{1000}{\csromanfolio{398}\textbf{“ยา ปน ภิกฺขุนี อคาริกสฺส วา ปริพฺพาชกสฺส วา ปริพฺพาชิกายวา สหตฺถา ขาทนียํ วา โภชนียํ วา ทเทยฺย, ปาจิตฺติยนฺ”}ติ.}

\proseitem{1001}{\textbf{ยา ปนา}ติ ยา ยาทิสา ฯเปฯ.\csromanspacer{} \textbf{ภิกฺขุนี}ติ ฯเปฯ อยํ อิมสฺมิํ อตฺเถ อธิปฺเปตา ภิกฺขุนีติ.}

\prose{\textbf{อคาริโก} นาม โย โกจิ อคารํ อชฺฌาวสติ.}

\prose{\textbf{ปริพฺพาชโก} นาม ภิกฺขุญฺจ สามเณรญฺจ ฐเปตฺวา โย โกจิ ปริพฺพาชกสมาปนฺโน.}

\prose{\textbf{ปริพฺพาชิกา} นาม ภิกฺขุนิญฺจ สิกฺขมานญฺจ สามเณริญฺจ ฐเปตฺวา ยา กาจิ ปริพฺพาชิกสมาปนฺนา.}

\prose{\textbf{ขาทนียํ} นาม ปญฺจ โภชนานิ อุทกทนฺตโปนํ ฐเปตฺวา อวเสสํ ขาทนียํ นาม.}

\prose{\textbf{โภชนียํ} นาม ปญฺจ โภชนานิ โอทโน กุมฺมาโส สตฺตุ มจฺโฉ มํสํ.}

\prose{\textbf{ทเทยฺยา}ติ กาเยน วา กายปฏิพทฺเธน วา นิสฺสคฺคิเยน วา เทติ, อาปตฺติ ปาจิตฺติยสฺส.\csromanspacer{} อุทกทนฺตโปนํ เทติ, อาปตฺติ ทุกฺกฏสฺส.}

\proseitemwithcloserruled{1002}{อนาปตฺติ ทาเปติ น เทติ, อุปนิกฺขิปิตฺวา เทติ, พาหิราเลปํ เทติ, อุมฺมตฺติกาย อาทิกมฺมิกายาติ.}{ฉฏฺฐสิกฺขาปทํ นิฏฺฐิตํ.}
\csromanmatikaanchor{mtk.200}
\csromantocmarkref{subsection}{7. สตฺตมสิกฺขาปท}{mtk.200}
\bookmark[dest={mtk.200},level=2]{7. สตฺตมสิกฺขาปท}
\csromanheader{2}{5. \textbf{จิตฺตาคารวคฺค 7. สตฺตมสิกฺขาปท}}
\proseitem{1003}{เตน สมเยน พุทฺโธ ภควา สาวตฺถิยํ วิหรติ เชตวเน อนาถปิณฺฑิกสฺส อาราเม.\csromanspacer{} เตน โข ปน สมเยน ถุลฺลนนฺทา ภิกฺขุนี อาวสถจีวรํ อนิสฺสชฺชิตฺวา ปริภุญฺชติ.\csromanspacer{} อญฺญา อุตุนิโย ภิกฺขุนิโย น ลภนฺติ.\csromanspacer{} ยา ตา ภิกฺขุนิโย อปฺปิจฺฉา ฯเปฯ ตา อุชฺฌายนฺติ ขิยฺยนฺติ วิปาเจนฺติ “กถํ หิ นาม อยฺยา ถุลฺลนนฺทา อาวสถจีวรํ \csromanfolio{399}\textbf{อนิสฺสชฺชิตฺวา ปริภุญฺชิสฺสตี”ติ ฯเปฯ.}\csromanspacer{}\textbf{ สจฺจํ กิร ภิกฺขเว ถุลฺลนนฺทา} ภิกฺขุนี อาวสถจีวรํ อนิสฺสชฺชิตฺวา ปริภุญฺชตีติ.\csromanspacer{} สจฺจํ ภควาติ.\csromanspacer{} วิครหิ พุทฺโธ ภควา ฯเปฯ.\csromanspacer{} กถํ หิ นาม ภิกฺขเว ถุลฺลนนฺทา ภิกฺขุนี อาวสถจีวรํ อนิสฺสชฺชิตฺวา ปริภุญฺชิสฺสติ.\csromanspacer{} เนตํ ภิกฺขเว อปฺปสนฺนานํ วา ปสาทาย ฯเปฯ.\csromanspacer{} เอวญฺจ ปน ภิกฺขเว ภิกฺขุนิโย อิมํ สิกฺขาปทํ อุทฺทิสนฺตุ\csromandash{}}

\proseitem{1004}{\textbf{“ยา ปน ภิกฺขุนี อาวสถจีวรํ อนิสฺสชฺชิตฺวา ปริภุญฺเชยฺย, ปาจิตฺติยนฺ”}ติ.}

\proseitem{1005}{\textbf{ยา ปนา}ติ ยา ยาทิสา ฯเปฯ.\csromanspacer{} \textbf{ภิกฺขุนี}ติ ฯเปฯ อยํ อิมสฺมิํ อตฺเถ อธิปฺเปตา ภิกฺขุนีติ.}

\prose{\textbf{อาวสถจีวรํ} นาม “อุตุนิโย ภิกฺขุนิโย ปริภุญฺชนฺตู”ติ ทินฺนํ โหติ.}

\prose{\textbf{อนิสฺสชฺชิตฺวา ปริภุญฺเชยฺยา}ติ เทฺวติสฺโส รตฺติโย ปริภุญฺชิตฺวา จตุตฺถทิวเส โธวิตฺวา ภิกฺขุนิยา วา สิกฺขมานาย วา สามเณริยา วา อนิสฺสชฺชิตฺวา ปริภุญฺชติ, อาปตฺติ ปาจิตฺติยสฺส.}

\proseitem{1006}{อนิสฺสชฺชิเต อนิสฺสชฺชิตสญฺญา ปริภุญฺชติ, อาปตฺติ ปาจิตฺติยสฺส.\csromanspacer{} อนิสฺสชฺชิเต เวมติกา ปริภุญฺชติ, อาปตฺติ ปาจิตฺติยสฺส.\csromanspacer{} อนิสฺสชฺชิเต นิสฺสชฺชิตสญฺญา ปริภุญฺชติ, อาปตฺติ ปาจิตฺติยสฺส.}

\prose{นิสฺสชฺชิเต อนิสฺสชฺชิตสญฺญา, อาปตฺติ ทุกฺกฏสฺส.\csromanspacer{} นิสฺสชฺชิเต เวมติกา, อาปตฺติ ทุกฺกฏสฺส.\csromanspacer{} นิสฺสชฺชิเต นิสฺสชฺชิตสญฺญา, อนาปตฺติ.}

\proseitemwithcloser{1007}{อนาปตฺติ นิสฺสชฺชิตฺวา ปริภุญฺชติ, ปุน ปริยาเยน ปริภุญฺชติ, อญฺญา อุตุนิโย ภิกฺขุนิโย น โหนฺติ, อจฺฉินฺนจีวริกาย นฏฺฐจีวริกาย อาปทาสุ อุมฺมตฺติกาย อาทิกมฺมิกายาติ.}{สตฺตมสิกฺขาปทํ นิฏฺฐิตํ.}
\csromanmatikaanchor{mtk.201}
\csromantocmarkref{subsection}{8. อฏฺฐมสิกฺขาปท}{mtk.201}
\bookmark[dest={mtk.201},level=2]{8. อฏฺฐมสิกฺขาปท}
\csromanheader{2}{5. \textbf{จิตฺตาคารวคฺค 8. อฏฺฐมสิกฺขาปท}}
\proseitem{1008}{\csromanfolio{400}เตน สมเยน พุทฺโธ ภควา สาวตฺถิยํ วิหรติ เชตวเน อนาถปิณฺฑิกสฺส อาราเม.\csromanspacer{} เตน โข ปน สมเยน ถุลฺลนนฺทา ภิกฺขุนี อาวสถํ อนิสฺสชฺชิตฺวา จาริกํ ปกฺกามิ.\csromanspacer{} เตน โข ปน สมเยน ถุลฺลนนฺทาย ภิกฺขุนิยา อาวสโถ ฑยฺหติ.\csromanspacer{} ภิกฺขุนิโย เอวมาหํสุ “หนฺทายฺเย ภณฺฑกํ นีหรามา”ติ.\csromanspacer{} เอกจฺจา เอวมาหํสุ “น มยํ อยฺเย นีหริสฺสาม, ยํ กิญฺจิ นฏฺฐํ สพฺพํ อเมฺห อภิยุญฺชิสฺสตี”ติ.\csromanspacer{} ถุลฺลนนฺทา ภิกฺขุนี ปุนเทว ตํ อาวสถํ ปจฺจาคนฺตฺวา ภิกฺขุนิโย ปุจฺฉิ “อปายฺเย ภณฺฑกํ นีหริตฺถา”ติ.\csromanspacer{} น มยํ อยฺเย นีหริมฺหาติ.\csromanspacer{} ถุลฺลนนฺทา ภิกฺขุนี อุชฺฌายติ ขิยฺยติ วิปาเจติ “กถํ หิ นาม ภิกฺขุนิโย อาวสเถ ฑยฺหมาเน ภณฺฑกํ น นีหริสฺสนฺตี”ติ.\csromanspacer{} ยา ตา ภิกฺขุนิโย อปฺปิจฺฉา ฯเปฯ ตา อุชฺฌายนฺติ ขิยฺยนฺติ วิปาเจนฺติ “กถํ หิ นาม อยฺยา ถุลฺลนนฺทา อาวสถํ อนิสฺสชฺชิตฺวา จาริกํ ปกฺกมิสฺสตี”ติ ฯเปฯ.\csromanspacer{} สจฺจํ กิร ภิกฺขเว ถุลฺลนนฺทา ภิกฺขุนี อาวสถํ อนิสฺสชฺชิตฺวา จาริกํ ปกฺกมตีติ\footnote{ปกฺกมีติ (ก)}.\csromanspacer{} สจฺจํ ภควาติ.\csromanspacer{} วิครหิ พุทฺโธ ภควา ฯเปฯ.\csromanspacer{} กถํ หิ นาม ภิกฺขเว ถุลฺลนนฺทา ภิกฺขุนี อาวสถํ อนิสฺสชฺชิตฺวา จาริกํ ปกฺกมิสฺสติ.\csromanspacer{} เนตํ ภิกฺขเว อปฺปสนฺนานํ วา ปสาทาย ฯเปฯ.\csromanspacer{} เอวญฺจ ปน ภิกฺขเว ภิกฺขุนิโย อิมํ สิกฺขาปทํ อุทฺทิสนฺตุ\csromandash{}}

\proseitem{1009}{\textbf{“ยา ปน ภิกฺขุนี อาวสถํ อนิสฺสชฺชิตฺวา จาริกํ ปกฺกเมยฺย, ปาจิตฺติยนฺ”}ติ.}

\proseitem{1010}{\textbf{ยา ปนา}ติ ยา ยาทิสา ฯเปฯ.\csromanspacer{} \textbf{ภิกฺขุนี}ติ ฯเปฯ อยํ อิมสฺมิํ อตฺเถ อธิปฺเปตา ภิกฺขุนีติ.}

\prose{\textbf{อาวสโถ} นาม กวาฏพทฺโธ วุจฺจติ.}

\prose{\textbf{อนิสฺสชฺชิตฺวา จาริกํ ปกฺกเมยฺยา}ติ ภิกฺขุนิยา วา สิกฺขมานาย วา สามเณริยา วา อนิสฺสชฺชิตฺวา ปริกฺขิตฺตสฺส อาวสถสฺส ปริกฺเขปํ อติกฺกาเมนฺติยา อาปตฺติ ปาจิตฺติยสฺส.\csromanspacer{} อปริกฺขิตฺตสฺส อาวสถสฺส อุปจารํ อติกฺกาเมนฺติยา อาปตฺติ ปาจิตฺติยสฺส.}

\proseitem{1011}{อนิสฺสชฺชิเต อนิสฺสชฺชิตสญฺญา ปกฺกมติ, อาปตฺติ ปาจิตฺติยสฺส.\csromanspacer{} อนิสฺสชฺชิเต เวมติกา ปกฺกมติ, อาปตฺติ ปาจิตฺติยสฺส.\csromanspacer{} อนิสฺสชฺชิเต นิสฺสชฺชิตสญฺญา ปกฺกมติ, อาปตฺติ ปาจิตฺติยสฺส.}

\prose{\csromanfolio{401}อกวาฏพทฺธํ อนิสฺสชฺชิตฺวา ปกฺกมติ, อาปตฺติ ทุกฺกฏสฺส.\csromanspacer{} นิสฺสชฺชิเต อนิสฺสชฺชิตสญฺญา, อาปตฺติ ทุกฺกฏสฺส.\csromanspacer{} นิสฺสชฺชิเต เวมติกา, อาปตฺติ ทุกฺกฏสฺส.\csromanspacer{} นิสฺสชฺชิเต นิสฺสชฺชิตสญฺญา, อนาปตฺติ.}

\proseitemwithcloserruled{1012}{อนาปตฺติ นิสฺสชฺชิตฺวา ปกฺกมติ, สติ อนฺตราเย, ปริเยสิตฺวา น ลภติ, คิลานาย อาปทาสุ อุมฺมตฺติกาย อาทิกมฺมิกายาติ.}{อฏฺฐมสิกฺขาปทํ นิฏฺฐิตํ.}
\csromanmatikaanchor{mtk.202}
\csromantocmarkref{subsection}{9. นวมสิกฺขาปท}{mtk.202}
\bookmark[dest={mtk.202},level=2]{9. นวมสิกฺขาปท}
\csromanheader{2}{5. \textbf{จิตฺตาคารวคฺค 9. นวมสิกฺขาปท}}
\proseitem{1013}{เตน สมเยน พุทฺโธ ภควา สาวตฺถิยํ วิหรติ เชตวเน อนาถปิณฺฑิกสฺส อาราเม.\csromanspacer{} เตน โข ปน สมเยน ฉพฺพคฺคิยา ภิกฺขุนิโย ติรจฺฉานวิชฺชํ ปริยาปุณนฺติ.\csromanspacer{} มนุสฺสา อุชฺฌายนฺติ ขิยฺยนฺติ วิปาเจนฺติ “กถํ หิ นาม ภิกฺขุนิโย ติรจฺฉานวิชฺชํ ปริยาปุณิสฺสนฺติ เสยฺยถาปิ คิหินิโย กามโภคินิโย”ติ.\csromanspacer{} อสฺโสสุํ โข ภิกฺขุนิโย เตสํ มนุสฺสานํ อุชฺฌายนฺตานํ วิปาเจนฺตานํ.\csromanspacer{} ยา ตา ภิกฺขุนิโย อปฺปิจฺฉา ฯเปฯ ตา อุชฺฌายนฺติ ขิยฺยนฺติ วิปาเจนฺติ “กถํ หิ นาม ฉพฺพคฺคิยา ภิกฺขุนิโย ติรจฺฉานวิชฺชํ ปริยาปุณิสฺสนฺตี”ติ ฯเปฯ.\csromanspacer{} สจฺจํ กิร ภิกฺขเว ฉพฺพคฺคิยา ภิกฺขุนิโย ติรจฺฉานวิชฺชํ ปริยาปุณนฺตีติ.\csromanspacer{} สจฺจํ ภควาติ.\csromanspacer{} วิครหิ พุทฺโธ ภควา ฯเปฯ.\csromanspacer{} กถํ หิ นาม ภิกฺขเว ฉพฺพคฺคิยา ภิกฺขุนิโย ติรจฺฉานวิชฺชํ ปริยาปุณิสฺสนฺติ.\csromanspacer{} เนตํ ภิกฺขเว อปฺปสนฺนานํ วา ปสาทาย ฯเปฯ.\csromanspacer{} เอวญฺจ ปน ภิกฺขเว ภิกฺขุนิโย อิมํ สิกฺขาปทํ อุทฺทิสนฺตุ\csromandash{}}

\proseitem{1014}{\textbf{“ยา ปน ภิกฺขุนี ติรจฺฉานวิชฺชํ ปริยาปุเณยฺย, ปาจิตฺติยนฺ”}ติ.}

\proseitem{1015}{\textbf{ยา ปนา}ติ ยา ยาทิสา ฯเปฯ.\csromanspacer{} \textbf{ภิกฺขุนี}ติ ฯเปฯ อยํ อิมสฺมิํ อตฺเถ อธิปฺเปตา ภิกฺขุนีติ.}

\prose{\textbf{ติรจฺฉานวิชฺชา}\footnote{ติรจฺฉานวิชฺชํ (ก)} นาม ยํ กิญฺจิ พาหิรกํ อนตฺถสํหิตํ.}

\prose{\csromanfolio{402}\textbf{ปริยาปุเณยฺยา}ติ ปเทน ปริยาปุณาติ, ปเท ปเท อาปตฺติ ปาจิตฺติยสฺส.\csromanspacer{} อกฺขราย ปริยาปุณาติ, อกฺขรกฺขราย อาปตฺติ ปาจิตฺติยสฺส.}

\proseitemwithcloserruled{1016}{อนาปตฺติ เลขํ ปริยาปุณาติ, ธารณํ ปริยาปุณาติ คุตฺตตฺถาย ปริตฺตํ ปริยาปุณาติ, อุมฺมตฺติกาย อาทิกมฺมิกายาติ.}{นวมสิกฺขาปทํ นิฏฺฐิตํ.}
\csromanmatikaanchor{mtk.203}
\csromantocmarkref{subsection}{10. ทสมสิกฺขาปท}{mtk.203}
\bookmark[dest={mtk.203},level=2]{10. ทสมสิกฺขาปท}
\csromanheader{2}{5. \textbf{จิตฺตาคารวคฺคฺ}อ 10. ทสมสิกฺขาปท}
\proseitem{1017}{เตน สมเยน พุทฺโธ ภควา สาวตฺถิยํ วิหรติ เชตวเน อนาถปิณฺฑิกสฺส อาราเม.\csromanspacer{} เตน โข ปน สมเยน ฉพฺพคฺคิยา ภิกฺขุนิโย ติรจฺฉานวิชฺชํ วาเจนฺติ.\csromanspacer{} มนุสฺสา อุชฺฌายนฺติ ขิยฺยนฺติ วิปาเจนฺติ “กถํ หิ นาม ภิกฺขุนิโย ติรจฺฉานวิชฺชํ วาเจสฺสนฺติ เสยฺยถาปิ คิหินิโย กามโภคินิโย”ติ.\csromanspacer{} อสฺโสสุํ โข ภิกฺขุนิโย เตสํ มนุสฺสานํ อุชฺฌายนฺตานํ ขิยฺยนฺตานํ วิปาเจนฺตานํ.\csromanspacer{} ยา ตา ภิกฺขุนิโย อปฺปิจฺฉา ฯเปฯ ตา อุชฺฌายนฺติ ขิยฺยนฺติ วิปาเจนฺติ “กถํ หิ นาม ฉพฺพคฺคิยา ภิกฺขุนิโย ติรจฺฉานวิชฺชํ วาเจสฺสนฺตี”ติ ฯเปฯ.\csromanspacer{} สจฺจํ กิร ภิกฺขเว ฉพฺพคฺคิยา ภิกฺขุนิโย ติรจฺฉานวิชฺชํ วาเจนฺตีติ.\csromanspacer{} สจฺจํ ภควาติ.\csromanspacer{} วิครหิ พุทฺโธ ภควา ฯเปฯ.\csromanspacer{} กถํ หิ นาม ภิกฺขเว ฉพฺพคฺคิยา ภิกฺขุนิโย ติรจฺฉานวิชฺชํ วาเจสฺสนฺติ.\csromanspacer{} เนตํ ภิกฺขเว อปฺปสนฺนานํ วา ปสาทาย ฯเปฯ.\csromanspacer{} เอวญฺจ ปน ภิกฺขเว ภิกฺขุนิโย อิมํ สิกฺขาปทํ อุทฺทิสนฺตุ\csromandash{}}

\proseitem{1018}{\textbf{“ยา ปน ภิกฺขุนี ติรจฺฉานวิชฺชํ วาเจยฺย, ปาจิตฺติยนฺ”}ติ.}

\proseitem{1019}{\textbf{ยา ปนา}ติ ยา ยาทิสา ฯเปฯ.\csromanspacer{} \textbf{ภิกฺขุนี}ติ ฯเปฯ อยํ อิมสฺมิํ อตฺเถ อธิปฺเปตา ภิกฺขุนีติ.}

\prose{\textbf{ติรจฺฉานวิชฺชา}\footnote{ติรจฺฉานวิชฺชํ (ก)} นาม ยํ กิญฺจิ พาหิรกํ อนตฺถสํหิตํ.}

\prose{\textbf{วาเจยฺยา}ติ ปเทน วาเจติ, ปเท ปเท อาปตฺติ ปาจิตฺติยสฺส.\csromanspacer{} อกฺขราย วาเจติ, อกฺขรกฺขราย อาปตฺติ ปาจิตฺติยสฺส.}

\proseitemwithcloser{1020}{\csromanfolio{403}อนาปตฺติ เลขํ วาเจติ, ธารณํ วาเจติ, คุตฺตตฺถาย ปริตฺตํ วาเจติ, อุมฺมตฺติกาย อาทิกมฺมิกายาติ.}{ทสมสิกฺขาปทํ นิฏฺฐิตํ.}
\csromanheader{2}{จิตฺตาคารวคฺโค ปญฺจโม.}
\csromansectionrule
\csromanmatikaanchor{mtk.204}
\csromanmatikaanchor{mtk.205}
\csromantocmarknopage{section}{6. อารามวคฺค}{mtk.204}
\bookmark[dest={mtk.204},level=1]{6. อารามวคฺค}
\csromantocmarkref{subsection}{1. ปฐมสิกฺขาปท}{mtk.205}
\bookmark[dest={mtk.205},level=2]{1. ปฐมสิกฺขาปท}
\setcsromanheadh{6. อารามวคฺค}
\csromanheader{2}{6. \textbf{อารามวคฺค 1. ปฐมสิกฺขาปท}}
\proseitem{1021}{เตน สมเยน พุทฺโธ ภควา สาวตฺถิยํ วิหรติ เชตวเน อนาถปิณฺฑิกสฺส อาราเม.\csromanspacer{} เตน โข ปน สมเยน สมฺพหุลา ภิกฺขู คามกาวาเส เอกจีวรา จีวรกมฺมํ กโรนฺติ.\csromanspacer{} ภิกฺขุนิโย อนาปุจฺฉา อารามํ ปวิสิตฺวา เยน เต ภิกฺขู เตนุปสงฺกมิํสุ.\csromanspacer{} ภิกฺขู อุชฺฌายนฺติ ขิยฺยนฺติ วิปาเจนฺติ “กถํ หิ นาม ภิกฺขุนิโย อนาปุจฺฉา อารามํ ปวิสิสฺสนฺตี”ติ ฯเปฯ.\csromanspacer{} สจฺจํ กิร ภิกฺขเว ภิกฺขุนิโย อนาปุจฺฉา อารามํ ปวิสนฺตีติ.\csromanspacer{} สจฺจํ ภควาติ.\csromanspacer{} วิครหิ พุทฺโธ ภควา ฯเปฯ.\csromanspacer{} กถํ หิ นาม ภิกฺขเว ภิกฺขุนิโย อนาปุจฺฉา อารามํ ปวิสิสฺสนฺติ.\csromanspacer{} เนตํ ภิกฺขเว อปฺปสนฺนานํ วา ปสาทาย ฯเปฯ.\csromanspacer{} เอวญฺจ ปน ภิกฺขเว ภิกฺขุนิโย อิมํ สิกฺขาปทํ อุทฺทิสนฺตุ\csromandash{}}

\prose{\textbf{“ยา ปน ภิกฺขุนี อนาปุจฺฉา อารามํ ปวิเสยฺย, ปาจิตฺติยนฺ”}ติ.}

\prose{เอวญฺจิทํ ภควตา ภิกฺขุนีนํ สิกฺขาปทํ ปญฺญตฺตํ โหติ.}

\proseitem{1022}{เตน โข ปน สมเยน เต ภิกฺขู ตมฺหา อาวาสา ปกฺกมิํสุ.\csromanspacer{} ภิกฺขุนิโย “อยฺยา ปกฺกนฺตา”ติ อารามํ นาคมํสุ.\csromanspacer{} อถ โข เต ภิกฺขู ปุนเทว ตํ อาวาสํ ปจฺจาคจฺฉิํสุ.\csromanspacer{} ภิกฺขุนิโย “อยฺยา อาคตา”ติ อาปุจฺฉา อารามํ ปวิสิตฺวา เยน เต ภิกฺขู เตนุปสงฺกมิํสุ, อุปสงฺกมิตฺวา เต ภิกฺขู อภิวาเทตฺวา เอกมนฺตํ อฏฺฐํสุ, เอกมนฺตํ ฐิตา โข ตา ภิกฺขุนิโย เต ภิกฺขู เอตทโวจุํ “กิสฺส ตุเมฺห ภคินิโย อารามํ เนว สมฺมชฺชิตฺถ, น ปานียํ ปริโภชนียํ อุปฏฺฐาปิตฺถา”ติ.\csromanspacer{} ภควตา อยฺยา สิกฺขาปทํ ปญฺญตฺตํ\footnote{ปญฺญตฺตํ โหติ (ก)} “น อนาปุจฺฉา อาราโม ปวิสิตพฺโพ”ติ.\csromanspacer{} เตน มยํ น อาคมิมฺหาติ.\csromanspacer{} ภควโต เอตมตฺถํ อาโรเจสุํ ฯเปฯ.\csromanspacer{} อนุชานามิ ภิกฺขเว สนฺตํ ภิกฺขุํ อาปุจฺฉา \csromanfolio{404}อารามํ ปวิสิตุํ.\csromanspacer{} เอวญฺจ ปน ภิกฺขเว ภิกฺขุนิโย อิมํ สิกฺขาปทํ อุทฺทิสนฺตุ\csromandash{}}

\prose{\textbf{“ยา ปน ภิกฺขุนี สนฺตํ ภิกฺขุํ อนาปุจฺฉา อารามํ ปวิเสยฺย ปาจิตฺติยนฺ”}ติ.}

\prose{เอวญฺจิทํ ภควตา ภิกฺขุนีนํ สิกฺขาปทํ ปญฺญตฺตํ โหติ.}

\proseitem{1023}{เตน โข ปน สมเยน เต ภิกฺขู ตมฺหา อาวาสา ปกฺกมิตฺวา ปุนเทว ตํ อาวาสํ ปจฺจาคจฺฉิํสุ.\csromanspacer{} ภิกฺขุนิโย “อยฺยา ปกฺกนฺตา”ติ อนาปุจฺฉา อารามํ ปวิสิํสุ.\csromanspacer{} ตาสํ กุกฺกุจฺจํ อโหสิ “ภควตา สิกฺขาปทํ ปญฺญตฺตํ ‘น สนฺตํ ภิกฺขุํ อนาปุจฺฉา อาราโม ปวิสิตพฺโพ’ติ, มยญฺจมฺหา สนฺตํ ภิกฺขุํ อนาปุจฺฉา อารามํ ปวิสิมฺหา, กจฺจิ นุ โข มยํ ปาจิตฺติยํ อาปตฺติํ อาปนฺนา”ติ ฯเปฯ.\csromanspacer{} ภควโต เอตมตฺถํ อาโรเจสุํ.\csromanspacer{} อถ โข ภควา เอตสฺมิํ นิทาเน เอตสฺมิํ ปกรเณ ธมฺมิํ กถํ กตฺวา ภิกฺขู อามนฺเตสิ ฯเปฯ.\csromanspacer{} เอวญฺจ ปน ภิกฺขเว ภิกฺขุนิโย อิมํ สิกฺขาปทํ อุทฺทิสนฺตุ\csromandash{}}

\proseitem{1024}{\textbf{“ยา ปน ภิกฺขุนี ชานํ สภิกฺขุกํ อารามํ อนาปุจฺฉา ปวิเสยฺย, ปาจิตฺติยนฺ”}ติ.}

\proseitem{1025}{\textbf{ยา ปนา}ติ ยา ยาทิสา ฯเปฯ.\csromanspacer{} \textbf{ภิกฺขุนี}ติ ฯเปฯ อยํ อิมสฺมิํ อตฺเถ อธิปฺเปตา ภิกฺขุนีติ.}

\prose{\textbf{ชานาติ} นาม สามํ วา ชานาติ, อญฺเญ วา ตสฺสา อาโรเจนฺติ, เต วา อาโรเจนฺติ.}

\prose{\textbf{สภิกฺขุโก} นาม อาราโม ยตฺถ ภิกฺขู รุกฺขมูเลปิ วสนฺติ.}

\prose{\textbf{อนาปุจฺฉา อารามํ ปวิเสยฺยา}ติ ภิกฺขุํ วา สามเณรํ วา อารามิกํ วา อนาปุจฺฉา ปริกฺขิตฺตสฺส อารามสฺส ปริกฺเขปํ อติกฺกาเมนฺติยา อาปตฺติ ปาจิตฺติยสฺส.\csromanspacer{} อปริกฺขิตฺตสฺส อารามสฺส อุปจารํ โอกฺกมนฺติยา อาปตฺติ}

\proseitem{1026}{สภิกฺขุเก สภิกฺขุกสญฺญา สนฺตํ ภิกฺขุํ อนาปุจฺฉา อารามํ ปวิสติ, อาปตฺติ ปาจิตฺติยสฺส.\csromanspacer{} สภิกฺขุเก เวมติกา สนฺตํ ภิกฺขุํ \csromanfolio{405}อนาปุจฺฉา อารามํ ปวิสติ, อาปตฺติ ทุกฺกฏสฺส.\csromanspacer{} สภิกฺขุเก อภิกฺขุกสญฺญา สนฺตํ ภิกฺขุํ อนาปุจฺฉา อารามํ ปวิสติ, อนาปตฺติ.}

\prose{อภิกฺขุเก สภิกฺขุกสญฺญา, อาปตฺติ ทุกฺกฏสฺส.\csromanspacer{} อภิกฺขุเก เวมติกา, อาปตฺติ ทุกฺกฏสฺส.\csromanspacer{} อภิกฺขุเก อภิกฺขุกสญฺญา, อนาปตฺติ.}

\proseitemwithcloserruled{1027}{อนาปตฺติ สนฺตํ ภิกฺขุํ อาปุจฺฉา ปวิสติ, อสนฺตํ ภิกฺขุํ อนาปุจฺฉา ปวิสติ, สีสานุโลกิกา คจฺฉติ, ยตฺถ ภิกฺขุนิโย สนฺนิปติตา โหนฺติ ตตฺถ คจฺฉติ, อาราเมน มคฺโค โหติ, คิลานาย อาปทาสุ อุมฺมตฺติกาย อาทิกมฺมิกายาติ.}{ปฐมสิกฺขาปทํ นิฏฺฐิตํ.}
\csromanmatikaanchor{mtk.206}
\csromantocmarkref{subsection}{2. ทุติยสิกฺขาปท}{mtk.206}
\bookmark[dest={mtk.206},level=2]{2. ทุติยสิกฺขาปท}
\csromanheader{2}{6. \textbf{อารามวคฺค 2. ทุติยสิกฺขาปท}}
\proseitem{1028}{เตน สมเยน พุทฺโธ ภควา เวสาลิยํ วิหรติ มหาวเน กูฏาคารสาลายํ.\csromanspacer{} เตน โข ปน สมเยน อายสฺมโต อุปาลิสฺส อุปชฺฌาโย อายสฺมา กปฺปิตโก สุสาเน วิหรติ.\csromanspacer{} เตน โข ปน สมเยน ฉพฺพคฺคิยานํ ภิกฺขุนีนํ มหตฺตรา\footnote{มหนฺตตรา (สี)} ภิกฺขุนี กาลงฺกตา โหติ.\csromanspacer{} ฉพฺพคฺคิยา ภิกฺขุนิโย ตํ ภิกฺขุนิํ นีหริตฺวา อายสฺมโต กปฺปิตกสฺส วิหารสฺส อวิทูเร ฌาเปตฺวา ถูปํ กตฺวา คนฺตฺวา ตสฺมิํ ถูเป โรทนฺติ.\csromanspacer{} อถ โข อายสฺมา กปฺปิตโก เตน สทฺเทน อุพฺพาโฬฺห ตํ ถูปํ ภินฺทิตฺวา ปกิเรสิ.\csromanspacer{} ฉพฺพคฺคิยา ภิกฺขุนิโย “อิมินา กปฺปิตเกน อมฺหากํ อยฺยาย ถูโป ภินฺโน, หนฺท นํ ฆาเตมา”ติ มนฺเตสุํ.\csromanspacer{} อญฺญตรา ภิกฺขุนี อายสฺมโต อุปาลิสฺส เอตมตฺถํ อาโรเจสิ.\csromanspacer{} อายสฺมา อุปาลิ อายสฺมโต กปฺปิตกสฺส เอตมตฺถํ อาโรเจสิ.\csromanspacer{} อถ โข อายสฺมา กปฺปิตโก วิหารา นิกฺขมิตฺวา นิลีโน อจฺฉิ.\csromanspacer{} อถ โข ฉพฺพคฺคิยา ภิกฺขุนิโย เยนายสฺมโต กปฺปิตกสฺส วิหาโร เตนุปสงฺกมิํสุ, อุปสงฺกมิตฺวา อายสฺมโต กปฺปิตกสฺส วิหารํ ปาสาเณหิ จ เลฑฺฑูหิ จ โอตฺถราเปตฺวา “มโต กปฺปิตโก”ติ ปกฺกมิํสุ.}

\prose{\csromanfolio{406}อถ โข อายสฺมา กปฺปิตโก ตสฺสา รตฺถิยา อจฺจเยน ปุพฺพณฺหสมยํ นิวาเสตฺวา ปตฺตจีวรมาทาย เวสาลิํ ปิณฺฑาย ปาวิสิ.\csromanspacer{} อทฺทสํสุ โข ฉพฺพคฺคิยา ภิกฺขุนิโย อายสฺมนฺตํ กปฺปิตกํ ปิณฺฑาย จรนฺตํ, ทิสฺวาน เอวมาหํสุ “อยํ กปฺปิตโก ชีวติ, โก นุ โข อมฺหากํ มนฺตํ สํหรี”ติ.\csromanspacer{} อสฺโสสุํ โข ฉพฺพคฺคิยา ภิกฺขุนิโย “อยฺเยน กิร อุปาลินา อมฺหากํ มนฺโต สํหโฏ”ติ.\csromanspacer{} ตา อายสฺมนฺตํ อุปาลิํ อกฺโกสิํสุ “กถํ หิ นาม อยํ กาสาวโฏ มลมชฺชโน นิหีนชจฺโจ อมฺหากํ มนฺตํ สํหริสฺสตี”ติ.\csromanspacer{} ยา ตา ภิกฺขุนิโย อปฺปิจฺฉา ฯเปฯ ตา อุชฺฌายนฺติ ขิยฺยนฺติ วิปาเจนฺติ “กถํ หิ นาม ฉพฺพคฺคิยา ภิกฺขุนิโย อยฺยํ อุปาลิํ อกฺโกสิสฺสนฺตี”ติ ฯเปฯ.\csromanspacer{} สจฺจํ กิร ภิกฺขเว ฉพฺพคฺคิยา ภิกฺขุนิโย อุปาลิํ อกฺโกสนฺตีติ.\csromanspacer{} สจฺจํ ภควาติ.\csromanspacer{} วิครหิ พุทฺโธ ภควา ฯเปฯ.\csromanspacer{} กถํ หิ นาม ภิกฺขเว ฉพฺพคฺคิยา ภิกฺขุนิโย อุปาลิํ อกฺโกสิสฺสนฺติ.\csromanspacer{} เนตํ ภิกฺขเว อปฺปสนฺนานํ วา ปสาทาย ฯเปฯ.\csromanspacer{} เอวญฺจ ปน ภิกฺขเว ภิกฺขุนิโย อิมํ สิกฺขาปทํ อุทฺทิสนฺตุ\csromandash{}}

\proseitem{1029}{\textbf{“ยา ปน ภิกฺขุนี ภิกฺขุํ อกฺโกเสยฺย วา ปริภาเสยฺย วา, ปาจิตฺติยนฺ”}ติ.}

\proseitem{1030}{\textbf{ยา ปนา}ติ ยา ยาทิสา ฯเปฯ.\csromanspacer{} \textbf{ภิกฺขุนี}ติ ฯเปฯ อยํ อิมสฺมิํ อตฺเถ อธิปฺเปตา ภิกฺขุนีติ.}

\prose{\textbf{ภิกฺขุนฺ}ติ อุปสมฺปนฺนํ.\csromanspacer{} \textbf{อกฺโกเสยฺย วา}ติ ทสหิ วา อกฺโกสวตฺถูหิ อกฺโกสติ เอเตสํ วา อญฺญตเรน, อาปตฺติ ปาจิตฺติยสฺส.}

\prose{\textbf{ปริภาเสยฺย} \textbf{วา}ติ ภยํ อุปทํเสติ, อาปตฺติ ปาจิตฺติยสฺส.}

\proseitem{1031}{อุปสมฺปนฺเน อุปสมฺปนฺนสญฺญา อกฺโกสติ วา ปริภาสติ วา, อาปตฺติ ปาจิตฺติยสฺส.\csromanspacer{} อุปสมฺปนฺเน เวมติกา อกฺโกสติ วา ปริภาสติ วา, อาปตฺติ ปาจิตฺติยสฺส.\csromanspacer{} อุปสมฺปนฺเน อนุปสมฺปนฺนสญฺญา อกฺโกสติ วา ปริภาสติ วา, อาปตฺติ ปาจิตฺติยสฺส.}

\prose{อนุปสมฺปนฺนํ อกฺโกสติ วา ปริภาสติ วา, อาปตฺติ ทุกฺกฏสฺส.\csromanspacer{} อนุปสมฺปนฺเน อุปสมฺปนฺนสญฺญา, อาปตฺติ ทุกฺกฏสฺส.\csromanspacer{} อนุปสมฺปนฺเน เวมติกา, อาปตฺติ ทุกฺกฏสฺส.\csromanspacer{} อนุปสมฺปนฺเน อนุปสมฺปนฺนสญฺญา, อาปตฺติ ทุกฺกฏสฺส.}

\proseitemwithcloserruled{1032}{\csromanfolio{407}อนาปตฺติ อตฺถปุเรกฺขาราย ธมฺมปุเรกฺขาราย อนุสาสนิ\-ปุเรกฺขาราย อุมฺมตฺติกาย อาทิกมฺมิกายาติ.}{ทุติยสิกฺขาปทํ นิฏฺฐิตํ.}
\csromanmatikaanchor{mtk.207}
\csromantocmarkref{subsection}{3. ตติยสิกฺขาปท}{mtk.207}
\bookmark[dest={mtk.207},level=2]{3. ตติยสิกฺขาปท}
\csromanheader{2}{6. \textbf{อารามวคฺค 3. ตติยสิกฺขาปท}}
\proseitem{1033}{เตน สมเยน พุทฺโธ ภควา สาวตฺถิยํ วิหรติ เชตวเน อนาถปิณฺฑิกสฺส อาราเม.\csromanspacer{} เตน โข ปน สมเยน จณฺฑกาฬี ภิกฺขุนี ภณฺฑนการิกา โหติ กลหการิกา วิวาทการิกา ภสฺสการิกา สํเฆ อธิกรณการิกา.\csromanspacer{} ถุลฺลนนฺทา ภิกฺขุนี ตสฺสา กมฺเม กรียมาเน ปฏิกฺโกสติ.\csromanspacer{} เตน โข ปน สมเยน ถุลฺลนนฺทา ภิกฺขุนี คามกํ อคมาสิ เกนจิเทว กรณีเยน.\csromanspacer{} อถ โข ภิกฺขุนิสํโฆ “ถุลฺลนนฺทา ภิกฺขุนี ปกฺกนฺตา”ติ จณฺฑกาฬิํ ภิกฺขุนิํ อาปตฺติยา อทสฺสเน อุกฺขิปิ.\csromanspacer{} ถุลฺลนนฺทา ภิกฺขุนี คามเก ตํ กรณียํ ตีเรตฺวา ปุนเทว สาวตฺถิํ ปจฺจาคจฺฉิ.\csromanspacer{} จณฺฑกาฬี ภิกฺขุนี ถุลฺลนนฺทาย ภิกฺขุนิยา อาคจฺฉนฺติยา เนว อาสนํ ปญฺญเปสิ, น ปาโททกํ ปาทปีฐํ ปาทกฐลิกํ อุปนิกฺขิปิ, น ปจฺจุคฺคนฺตฺวา ปตฺตจีวรํ ปฏิคฺคเหสิ, น ปานีเยน อาปุจฺฉิ.\csromanspacer{} ถุลฺลนนฺทา ภิกฺขุนี จณฺฑกาฬิํ ภิกฺขุนิํ เอตทโวจ “กิสฺส ตฺวํ อยฺเย มยิ อาคจฺฉนฺติยา เนว อาสนํ ปญฺญเปสิ, น ปาโททกํ ปาทปีฐํ ปาทกฐลิกํ อุปนิกฺขิปิ, น ปจฺจุคฺคนฺตฺวา ปตฺตจีวรํ ปฏิคฺคเหสิ, น ปานีเยน อาปุจฺฉี”ติ.\csromanspacer{} เอวํ เหตํ อยฺเย โหติ, ยถา ตํ อนาถายาติ.\csromanspacer{} กิสฺส ปน ตฺวํ อยฺเย อนาถาติ.\csromanspacer{} อิมา มํ อยฺเย ภิกฺขุนิโย “อยํ อนาถา อปฺปญฺญาตา, นตฺถิ อิมิสฺสา กาจิ ปติวตฺตา”ติ อาปตฺติยา อทสฺสเน อุกฺขิปิํสูติ.\csromanspacer{} ถุลฺลนนฺทา ภิกฺขุนี “พาลา เอตา, อพฺยตฺตา เอตา, เนตา ชานนฺติ กมฺมํ วา กมฺมโทสํ วา กมฺมวิปตฺติํ วา กมฺมสมฺปตฺติํ วา”ติ จณฺฑีกตา คณํ ปริภาสิ.\csromanspacer{} ยา ตา ภิกฺขุนิโย อปฺปิจฺฉา ฯเปฯ ตา อุชฺชฺชายนฺติ ขิยฺยนฺติ วิปาเจนฺติ “กถํ หิ นาม อยฺยา ถุลฺลนนฺทา จณฺฑีกตา คณํ ปริภาสิสฺสตี”ติ ฯเปฯ.\csromanspacer{} สจฺจํ กิร ภิกฺขเว ถุลฺลนนฺทา ภิกฺขุนี จณฺฑีกตา คณํ ปริภาสตีติ\footnote{ปริภาสีติ (ก)}.\csromanspacer{} สจฺจํ ภควาติ.\csromanspacer{} วิครหิ พุทฺโธ ภควา ฯเปฯ.\csromanspacer{} กถํ หิ นาม ภิกฺขเว ถุลฺลนนฺทา ภิกฺขุนี จณฺฑีกตา \csromanfolio{408}คณํ ปริภาสิสฺสติ.\csromanspacer{} เนตํ ภิกฺขเว อปฺปสนฺนานํ วา ปสาทาย ฯเปฯ.\csromanspacer{} เอวญฺจ ปน ภิกฺขเว ภิกฺขุนิโย อิมํ สิกฺขาปทํ อุทฺทิสนฺตุ\csromandash{}}

\proseitem{1034}{\textbf{“ยา ปน ภิกฺขุนี จณฺฑีกตา คณํ ปริภาเสยฺย, ปาจิตฺติยนฺ”}ติ.}

\proseitem{1035}{\textbf{ยา ปนา}ติ ยา ยาทิสา ฯเปฯ.\csromanspacer{} \textbf{ภิกฺขุนี}ติ ฯเปฯ อยํ อิมสฺมิํ อตฺเถ อธิปฺเปตา ภิกฺขุนีติ.}

\prose{\textbf{จณฺฑีกตา} นาม โกธนา วุจฺจติ.}

\prose{\textbf{คโณ} นาม ภิกฺขุนิสํโฆ วุจฺจติ.}

\prose{\textbf{ปริภาเสยฺยา}ติ “พาลา เอตา, อพฺยตฺตา เอตา, เนตา ชานนฺติ กมฺมํ วา กมฺมโทสํ วา กมฺมวิปตฺติํ วา กมฺมสมฺปตฺติํ วา”ติ ปริภาสติ, อาปตฺติ ปาจิตฺติยสฺส.\csromanspacer{} สมฺพหุลา ภิกฺขุนิโย วา เอกํ ภิกฺขุนิํ วา อนุปสมฺปนฺนํ วา ปริภาสติ, อาปตฺติ ทุกฺกฏสฺส.}

\proseitemwithcloserruled{1036}{อนาปตฺติ อตฺถปุเรกฺขาราย ธมฺมปุเรกฺขาราย อนุสาสนิ\-ปุเรกฺขาราย อุมฺมตฺติกาย อาทิกมฺมิกายาติ.}{ตติยสิกฺขาปทํ นิฏฺฐิตํ.}
\csromanmatikaanchor{mtk.208}
\csromantocmarkref{subsection}{4. จตุตฺถสิกฺขาปท}{mtk.208}
\bookmark[dest={mtk.208},level=2]{4. จตุตฺถสิกฺขาปท}
\csromanheader{2}{6. \textbf{อารามวคฺค 4. จตุตฺถสิกฺขาปท}}
\proseitem{1037}{เตน สมเยน พุทฺโธ ภควา สาวตฺถิยํ วิหรติ เชตวเน อนาถปิณฺฑิกสฺส อาราเม.\csromanspacer{} เตน โข ปน สมเยน อญฺญตโร พฺราหฺมโณ ภิกฺขุนิโย นิมนฺเตตฺวา โภเชสิ.\csromanspacer{} ภิกฺขุนิโย ภุตฺตาวี\footnote{ภุตฺตาวินี (ก)} ปวาริตา ญาติกุลานิ คนฺตฺวา เอกจฺจา ภุญฺชิํสุ, เอกจฺจา ปิณฺฑปาตํ อาทาย อคมํสุ.\csromanspacer{} อถ โข โส พฺราหฺมโณ ปฏิวิสฺสเก เอตทโวจ “ภิกฺขุนิโย มยา อยฺยา สนฺตปฺปิตา, เอถ ตุเมฺหปิ สนฺตปฺเปสฺสามี”ติ.\csromanspacer{} เต เอวมาหํสุ “กิํ ตฺวํ อยฺโย อเมฺห สนฺตปฺเปสฺสสิ, ยาปิ ตยา นิมนฺติตา, ตาปิ อมฺหากํ ฆรานิ อาคนฺตฺวา เอกจฺจา ภุญฺชิํสุ, เอกจฺจา \csromanfolio{409}ปิณฺฑปาตํ อาทาย อคมํสู”ติ.\csromanspacer{} อถ โข โส พฺราหฺมโณ อุชฺฌายติ ขิยฺยติ วิปาเจติ “กถํ หิ นาม ภิกฺขุนิโย อมฺหากํ ฆเร ภุญฺชิตฺวา อญฺญตฺร ภุญฺชิสฺสนฺติ, น จาหํ ปฏิพโล ยาวทตฺถํ ทาตุนฺ”ติ.\csromanspacer{} อสฺโสสุํ โข ภิกฺขุนิโย ตสฺส พฺราหฺมณสฺส อุชฺฌายนฺตสฺส ขิยฺยนฺตสฺส วิปาเจนฺตสฺส.\csromanspacer{} ยา ตา ภิกฺขุนิโย อปฺปิจฺฉา ฯเปฯ ตา อุชฺฌายนฺติ ขิยฺยนฺติ วิปาเจนฺติ “กถํ หิ นาม ภิกฺขุนิโย ภุตฺตาวี\footnote{ภุตฺตาวินี (ก)} ปวาริตา อญฺญตฺร ภุญฺชิสฺสนฺตี”ติ ฯเปฯ.\csromanspacer{} สจฺจํ กิร ภิกฺขเว ภิกฺขุนิโย ภุตฺตาวี ปวาริตา อญฺญตฺร ภุญฺชนฺตีติ.\csromanspacer{} สจฺจํ ภควาติ.\csromanspacer{} วิครหิ พุทฺโธ ภควา ฯเปฯ.\csromanspacer{} กถํ หิ นาม ภิกฺขเว ภิกฺขุนิโย ภุตฺตาวี ปวาริตา อญฺญตฺร ภุญฺชิสฺสนฺติ.\csromanspacer{} เนตํ ภิกฺขเว อปฺปสนฺนานํ วา ปสาทาย ฯเปฯ.\csromanspacer{} เอวญฺจ ปน ภิกฺขเว ภิกฺขุนิโย อิมํ สิกฺขาปทํ อุทฺทิสนฺตุ\csromandash{}}

\proseitem{1038}{\textbf{“ยา ปน ภิกฺขุนี นิมนฺติตา วา ปวาริตา วา ขาทนียํ วา โภชนียํ วา ขาเทยฺย วา ภุญฺเชยฺย วา, ปาจิตฺติยนฺ”}ติ.}

\proseitem{1039}{\textbf{ยา ปนา}ติ ยา ยาทิสา ฯเปฯ.\csromanspacer{} \textbf{ภิกฺขุนี}ติ ฯเปฯ อยํ อิมสฺมิํ อตฺเถ อธิปฺเปตา ภิกฺขุนีติ.}

\prose{\textbf{นิมนฺติตา} นาม ปญฺจนฺนํ โภชนานํ อญฺญตเรน โภชเนน นิมนฺติตา.}

\prose{\textbf{ปวาริตา} นาม อสนํ ปญฺญายติ, โภชนํ ปญฺญายติ, หตฺถปาเส ฐิตา อภิหรติ, ปฏิกฺเขโป ปญฺญายติ.}

\prose{\textbf{ขาทนียํ} นาม ปญฺจ โภชนานิ ยาคุํ ยามกาลิกํ สตฺตาหกาลิกํ ยาวชีวิกํ ฐเปตฺวา อวเสสํ ขาทนียํ นาม.}

\prose{\textbf{โภชนียํ} นาม ปญฺจ โภชนานิ โอทโน กุมฺมาโส สตฺตุ มจฺโฉ มํสํ.}

\prose{“ขาทิสฺสามิ ภุญฺชิสฺสามี”ติ ปฏิคฺคณฺหาติ, อาปตฺติ ทุกฺกฏสฺส.\csromanspacer{} อชฺโฌหาเร อชฺโฌหาเร อาปตฺติ ปาจิตฺติยสฺส.}

\proseitem{1040}{นิมนฺติเต นิมนฺติตสญฺญา ขาทนียํ วา โภชนียํ วา ขาทติ วา ภุญฺชติ วา, อาปตฺติ ปาจิตฺติยสฺส.\csromanspacer{} นิมนฺติเต เวมติกา ขาทนียํ วา \csromanfolio{410}โภชนียํ วา ขาทติ วา ภุญฺชติ วา, อาปตฺติ ปาจิตฺติยสฺส.\csromanspacer{} นิมนฺติเต อนิมนฺติตสญฺญา ขาทนียํ วา โภชนียํ วา ขาทติ วา ภุญฺชติ วา, อาปตฺติ ปาจิตฺติยสฺส.}

\prose{ยามกาลิกํ สตฺตาหกาลิกํ ยาวชีวิกํ อาหารตฺถาย ปฏิคฺคณฺหาติ, อาปตฺติ ทุกฺกฏสฺส.\csromanspacer{} อชฺโฌหาเร อชฺโฌหาเร อาปตฺติ ทุกฺกฏสฺส ฯเปฯ.}

\proseitemwithcloserruled{1041}{อนาปตฺติ นิมนฺติตา อปฺปวาริตา ยาคุํ ปิวติ สามิเก อปโลเกตฺวา ภุญฺชติ, ยามกาลิกํ สตฺตาหกาลิกํ ยาวชีวิกํ สติ ปจฺจเย ปริภุญฺชติ, อุมฺมตฺติกาย อาทิกมฺมิกายาติ.}{จตุตฺถสิกฺขาปทํ นิฏฺฐิตํ.}
\csromanmatikaanchor{mtk.209}
\csromantocmarkref{subsection}{5. ปญฺจมสิกฺขาปท}{mtk.209}
\bookmark[dest={mtk.209},level=2]{5. ปญฺจมสิกฺขาปท}
\csromanheader{2}{6. \textbf{อารามวคฺค 5. ปญฺจมสิกฺขาปท}}
\proseitem{1042}{เตน สมเยน พุทฺโธ ภควา สาวตฺถิยํ วิหรติ เชตวเน อนาถปิณฺฑิกสฺส อาราเม.\csromanspacer{} เตน โข ปน สมเยน อญฺญตรา ภิกฺขุนี สาวตฺถิยํ อญฺญตริสฺสา วิสิขาย ปิณฺฑาย จรมานา เยน อญฺญตรํ กุลํ เตนุปสงฺกมิ, อุปสงฺกมิตฺวา ปญฺญตฺเต อาสเน นิสีทิ.\csromanspacer{} อถ โข เต มนุสฺสา ตํ ภิกฺขุนิํ โภเชตฺวา เอตทโวจุํ “อญฺญาปิ อยฺเย ภิกฺขุนิโย อาคจฺฉนฺตู”ติ.\csromanspacer{} อถ โข สา ภิกฺขุนี “กถํ หิ นาม\footnote{กถํ อญฺญา (สฺยา)} ภิกฺขุนิโย นาคจฺเฉยฺยุนฺ”ติ ภิกฺขุนิโย อุปสงฺกมิตฺวา เอตทโวจ “อมุกสฺมิํ อยฺเย โอกาเส วาฬา สุนขา จณฺโฑ พลีพทฺโท จิกฺขลฺโล โอกาโส, มา โข ตตฺถ อคมิตฺถา”ติ.\csromanspacer{} อญฺญตราปิ ภิกฺขุนี ตสฺสา วิสิขาย ปิณฺฑาย จรมานา เยน ตํ กุลํ เตนุปสงฺกมิ, อุปสงฺกมิตฺวา ปญฺญตฺเต อาสเน นิสีทิ.\csromanspacer{} อถ โข เต มนุสฺสา ตํ ภิกฺขุนิํ โภเชตฺวา เอตทโวจุํ “กิสฺส อยฺเย ภิกฺขุนิโย น อาคจฺฉนฺตี”ติ.\csromanspacer{} อถ โข สา ภิกฺขุนี เตสํ มนุสฺสานํ เอตมตฺถํ อาโรเจสิ.\csromanspacer{} มนุสฺสา อุชฺฌายนฺติ ขิยฺยนฺติ วิปาเจนฺติ “กถํ หิ นาม ภิกฺขุนี กุลํ มจฺฉรายิสฺสตี”ติ ฯเปฯ.\csromanspacer{} สฺสจฺจํ กิร ภิกฺขเว ภิกฺขุนี \csromanfolio{411}กุลํ มจฺฉรายตีติ.\csromanspacer{} สจฺจํ ภควาติ.\csromanspacer{} วิครหิ พุทฺโธ ภควา ฯเปฯ.\csromanspacer{} กถํ หิ นาม ภิกฺขเว ภิกฺขุนี กุลํ มจฺฉรายิสฺสติ.\csromanspacer{} เนตํ ภิกฺขเว อปฺปสนฺนานํ วา ปสาทาย ฯเปฯ.\csromanspacer{} เอวญฺจ ปน ภิกฺขเว \textbf{ภิกฺขุนิโย อิมํ สิกฺขาปทํ อุทฺทิสนฺตุ\csromandash{}}}

\proseitem{1043}{\textbf{“ยา ปน ภิกฺขุนี กุลมจฺฉรินี อสฺส, ปาจิตฺติยนฺ”}ติ.}

\proseitem{1044}{\textbf{ยา ปนา}ติ ยา ยาทิสา ฯเปฯ.\csromanspacer{} \textbf{ภิกฺขุนี}ติ ฯเปฯ อยํ อิมสฺมิํ อตฺเถ อธิปฺเปตา ภิกฺขุนีติ.}

\prose{\textbf{กุลํ} นาม จตฺตาริ กุลานิ ขตฺติยกุลํ พฺราหฺมณกุลํ เวสฺสกุลํ สุทฺทกุลํ.}

\prose{\textbf{มจฺฉรินี อสฺสา}ติ “กถํ ภิกฺขุนิโย นาคจฺเฉยฺยุนฺ”ติ ภิกฺขุนีนํ สนฺติเก กุลสฺส อวณฺณํ ภาสติ, อาปตฺติ ปาจิตฺติยสฺส.\csromanspacer{} กุลสฺส วา สนฺติเก ภิกฺขุนีนํ อวณฺณํ ภาสติ, อาปตฺติ ปาจิตฺติยสฺส.}

\proseitemwithcloserruled{1045}{อนาปตฺติ กุลํ น มจฺฉรายนฺตี สนฺตํ เยว อาทีนวํ อาจิกฺขติ, อุมฺมตฺติกาย อาทิกมฺมิกายาติ.}{ปญฺจมสิกฺขาปทํ นิฏฺฐิตํ.}
\csromanmatikaanchor{mtk.210}
\csromantocmarkref{subsection}{6. ฉฏฺฐสิกฺขาปท}{mtk.210}
\bookmark[dest={mtk.210},level=2]{6. ฉฏฺฐสิกฺขาปท}
\csromanheader{2}{6. \textbf{อารามวคฺค 6. ฉฏฺฐสิกฺขาปท}}
\proseitem{1046}{เตน สมเยน พุทฺโธ ภควา สาวตฺถิยํ วิหรติ เชตวเน อนาถปิณฺฑิกสฺส อาราเม.\csromanspacer{} เตน โข ปน สมเยน สมฺพหุลา ภิกฺขุนิโย คามกาวาเส วสฺสํวุฏฺฐา สาวตฺถิํ อคมํสุ.\csromanspacer{} ภิกฺขุนิโย ตา ภิกฺขุนิโย เอตทโวจุํ “กตฺถายฺยาโย วสฺสํวุฏฺฐา, กจฺจิ โอวาโท อิทฺโธ อโหสี”ติ.\csromanspacer{} นตฺถายฺเย ตตฺถ ภิกฺขู, กุโต โอวาโท อิทฺโธ ภวิสฺสตีติ.\csromanspacer{} ยา ตา ภิกฺขุนิโย อปฺปิจฺฉา ฯเปฯ ตา อุชฺฌายนฺติ ขิยฺยนฺติ วิปาเจนฺติ “กถํ หิ นาม ภิกฺขุนิโย อภิกฺขุเก อาวาเส วสฺสํ วสิสฺสนฺตี”ติ ฯเปฯ.\csromanspacer{} สจฺจํ กิร ภิกฺขเว ภิกฺขุนิโย อภิกฺขุเก อาวาเส วสฺสํ วสนฺตีติ.\csromanspacer{} สจฺจํ ภควาติ.\csromanspacer{} วิครหิ พุทฺโธ ภควาติ ฯเปฯ.\csromanspacer{} กถํ หิ นาม \csromanfolio{412}ภิกฺขเว ภิกฺขุนิโย อภิกฺขุเก อาวาเส วสฺสํ วสิสฺสนฺติ.\csromanspacer{} เนตํ ภิกฺขเว อปฺปสนฺนานํ วา ปสาทาย ฯเปฯ.\csromanspacer{} เอวญฺจ ปน ภิกฺขเว ภิกฺขุนิโย อิมํ สิกฺขาปทํ อุทฺทิสนฺตุ\csromandash{}}

\proseitem{1047}{\textbf{“ยา ปน ภิกฺขุนี อภิกฺขุเก อาวาเส วสฺสํ วเสยฺย, ปาจิตฺติยนฺ”}ติ.}

\proseitem{1048}{\textbf{ยา ปนา}ติ ยา ยาทิสา ฯเปฯ.\csromanspacer{} \textbf{ภิกฺขุนี}ติ ฯเปฯ อยํ อิมสฺมิํ อตฺเถ อธิปฺเปตา ภิกฺขุนีติ.}

\prose{\textbf{อภิกฺขุโก} นาม อาวาโส น สกฺกา โหติ โอวาทาย วา สํวาสาย วา คนฺตุํ. “วสฺสํ วสิสฺสามี”ติ เสนาสนํ ปญฺญเปติ, ปานียํ ปริโภชนียํ อุปฏฺฐเปติ, ปริเวณํ สมฺมชฺชติ, อาปตฺติ ทุกฺกฏสฺส.\csromanspacer{} สห อรุณุคฺคมนา อาปตฺติ ปาจิตฺติยสฺส.}

\proseitemwithcloserruled{1049}{อนาปตฺติ วสฺสุปคตา ภิกฺขู ปกฺกนฺตา วา โหนฺติ, วิพฺภนฺตา วา, กาลงฺกตา วา, ปกฺขสงฺกนฺตา วา, อาปทาสุ อุมฺมตฺติกาย อาทิกมฺมิกายาติ.}{ฉฏฺฐสิกฺขาปทํ นิฏฺฐิตํ.}
\csromanmatikaanchor{mtk.211}
\csromantocmarkref{subsection}{7. สตฺตมสิกฺขาปท}{mtk.211}
\bookmark[dest={mtk.211},level=2]{7. สตฺตมสิกฺขาปท}
\csromanheader{2}{6. \textbf{อารามวคฺค 7. สตฺตมสิกฺขาปท}}
\proseitem{1050}{เตน สมเยน พุทฺโธ ภควา สาวตฺถิยํ วิหรติ เชตวเน อนาถปิณฺฑิกสฺส อาราเม.\csromanspacer{} เตน โข ปน สมเยน สมฺพหุลา ภิกฺขุนิโย คามกาวาเส วสฺสํวุฏฺฐา สาวตฺถิํ อคมํสุ.\csromanspacer{} ภิกฺขุนิโย ตา ภิกฺขุนิโย เอตทโวจุํ “กตฺถายฺยาโย วสฺสํวุฏฺฐา, กตฺถ\footnote{กจฺจิ (สฺยา)} ภิกฺขุสํโฆ ปวาริโต”ติ.\csromanspacer{} น มยํ อยฺเย ภิกฺขุสํฆํ ปวาเรมาติ.\csromanspacer{} ยา ตา ภิกฺขุนิโย อปฺปิจฺฉา ฯเปฯ ตา อุชฺฌายนฺติ ขิยฺยนฺติ วิปาเจนฺติ “กถํ หิ นาม ภิกฺขุนิโย วสฺสํวุฏฺฐา ภิกฺขุสํฆํ น ปวาเรสฺสนฺตี”ติ ฯเปฯ.\csromanspacer{} สจฺจํ กิร ภิกฺขเว ภิกฺขุนิโย วสฺสํวุฏฺฐา ภิกฺขุสํฆํ น ปวาเรนฺตีติ.\csromanspacer{} สจฺจํ ภควาติ.\csromanspacer{} วิครหิ พุทฺโธ ภควา ฯเปฯ.\csromanspacer{} กถํ หิ นาม ภิกฺขเว ภิกฺขุนิโย \csromanfolio{413}\textbf{วสฺสํวุฏฺฐา ภิกฺขุสํฆํ น ปวาเรสฺสนฺติ.}\csromanspacer{}\textbf{ เนตํ ภิกฺขเว} \textbf{อปฺปสนฺนานํ วา ปสาทาย ฯเปฯ.}\csromanspacer{}\textbf{ เอวญฺจ ปน ภิกฺขเว ภิกฺขุนิโย อิมํ} \textbf{สิกฺขาปทํ อุทฺทิสนฺตุ\csromandash{}}}

\proseitem{1051}{\textbf{“ยา ปน ภิกฺขุนี วสฺสํวุฏฺฐา อุภโตสํเฆ ตีหิ ฐาเนหิ น ปวาเรยฺย ทิฏฺเฐน วา สุเตน วา ปริสงฺกาย วา, ปาจิตฺติยนฺ”}ติ.}

\proseitem{1052}{\textbf{ยา ปนา}ติ ยา ยาทิสา ฯเปฯ.\csromanspacer{} \textbf{ภิกฺขุนี}ติ ฯเปฯ อยํ อิมสฺมิํ อตฺเถ อธิปฺเปตา ภิกฺขุนีติ.}

\prose{\textbf{วสฺสํวุฏฺฐา} นาม ปุริมํ วา เตมาสํ ปจฺฉิมํ วา เตมาสํ วุฏฺฐา. “อุภโตสํเฆ ตีหิ ฐาเนหิ น ปวาเรสฺสามิ ทิฏฺเฐน วา สุเตน วา ปริสงฺกาย วา”ติ ธุรํ นิกฺขิตฺตมตฺเต อาปตฺติ ปาจิตฺติยสฺส.}

\proseitemwithcloserruled{1053}{อนาปตฺติ สติ อนฺตราเย, ปริเยสิตฺวา น ลภติ, คิลานาย อาปทาสุ อุมฺมตฺติกาย อาทิกมฺมิกายาติ.}{สตฺตมสิกฺขาปทํ นิฏฺฐิตํ.}
\csromanmatikaanchor{mtk.212}
\csromantocmarkref{subsection}{8. อฏฺฐมสิกฺขาปท}{mtk.212}
\bookmark[dest={mtk.212},level=2]{8. อฏฺฐมสิกฺขาปท}
\csromanheader{2}{6. \textbf{อารามวคฺค 8. อฏฺฐมสิกฺขาปท}}
\proseitem{1054}{เตน สมเยน พุทฺโธ ภควา สกฺเกสุ วิหรติ กปิลวตฺถุสฺมิํ นิโคฺรธาราเม.\csromanspacer{} เตน โข ปน สมเยน ฉพฺพคฺคิยา ภิกฺขู ภิกฺขุนุปสฺสยํ อุปสงฺกมิตฺวา ฉพฺพคฺคิยา ภิกฺขุนิโย โอวทนฺติ.\csromanspacer{} ภิกฺขุนิโย ฉพฺพคฺคิยา ภิกฺขุนิโย เอตทโวจุํ “เอถายฺเย โอวาทํ คมิสฺสามา”ติ.\csromanspacer{} ยมฺปิ มยํ อยฺเย คจฺเฉยฺยาม โอวาทสฺส การณา, อยฺยา ฉพฺพคฺคิยา อิเธว อาคนฺตฺวา อเมฺห โอวทนฺตีติ.\csromanspacer{} ยา ตา ภิกฺขุนิโย อปฺปิจฺฉา ฯเปฯ ตา อุชฺฌายนฺติ ขิยฺยนฺติ วิปาเจนฺติ “กถํ หิ นาม ฉพฺพคฺคิยา ภิกฺขุนิโย โอวาทํ น คจฺฉิสฺสนฺตี”ติ ฯเปฯ.\csromanspacer{} สจฺจํ กิร ภิกฺขเว ฉพฺพคฺคิยา ภิกฺขุนิยา โอวาทํ น คจฺฉนฺตีติ.\csromanspacer{} สจฺจํ ภควาติ.\csromanspacer{} วิครหิ พุทฺโธ ภควา ฯเปฯ.\csromanspacer{} กถํ หิ นาม ภิกฺขเว ฉพฺพคฺคิยา ภิกฺขุนิยา โอวาทํ น คจฺฉิสฺสนฺติ.\csromanspacer{} เนตํ ภิกฺขเว อปฺปสนฺนานํ วา ปสาทาย ฯเปฯ.\csromanspacer{} เอวญฺจ ปน ภิกฺขเว ภิกฺขุนิยา อิมํ สิกฺขาปทํ อุทฺทิสนฺตุ\csromandash{}}

\proseitem{1055}{\csromanfolio{414}\textbf{“ยา ปน ภิกฺขุนี โอวาทาย วา สํวาสาย วา น คจฺเฉยฺย, ปาจิตฺติยนฺ”}ติ.}

\proseitem{1056}{\textbf{ยา ปนา}ติ ยา ยาทิสา ฯเปฯ.\csromanspacer{} \textbf{ภิกฺขุนี}ติ ฯเปฯ อยํ อิมสฺมิํ อตฺเถ อธิปฺเปตา ภิกฺขุนีติ.}

\prose{\textbf{โอวาโท} นาม อฏฺฐ ครุธมฺมา.}

\prose{\textbf{สํวาโส} นาม เอกกมฺมํ เอกุทฺเทโส สมสิกฺขตา. “โอวาทาย วา สํวาสาย วา น คจฺฉิสฺสามี”ติ ธุรํ นิกฺขิตฺตมตฺเต อาปตฺติ ปาจิตฺติยสฺส.}

\proseitemwithcloserruled{1057}{อนาปตฺติ สติ อนฺตราเย, ปริเยสิตฺวา ทุติยิกํ ภิกฺขุนิํ น ลภติ, คิลานาย อาปทาสุ อุมฺมตฺติกาย อาทิกมฺมิกายาติ.}{อฏฺฐมสิกฺขาปทํ นิฏฺฐิตํ.}
\csromanmatikaanchor{mtk.213}
\csromantocmarkref{subsection}{9. นวมสิกฺขาปท}{mtk.213}
\bookmark[dest={mtk.213},level=2]{9. นวมสิกฺขาปท}
\csromanheader{2}{6. \textbf{อารามวคฺค 9. นวมสิกฺขาปท}}
\proseitem{1058}{เตน สมเยน พุทฺโธ ภควา สาวตฺถิยํ วิหรติ เชตวเน อนาถปิณฺฑิกสฺส อาราเม.\csromanspacer{} เตน โข ปน สมเยน ภิกฺขุนิโย อุโปสถมฺปิ น ปุจฺฉนฺติ, โอวาทมฺปิ น ยาจนฺติ.\csromanspacer{} ภิกฺขู อุชฺฌายนฺติ ขิยฺยนฺติ วิปาเจนฺติ “กถํ หิ นาม ภิกฺขุนิโย อุโปสถมฺปิ น ปุจฺฉิสฺสนฺติ, โอวาทมฺปิ น ยาจิสฺสนฺตี”ติ ฯเปฯ.\csromanspacer{} สจฺจํ กิร ภิกฺขเว ภิกฺขุนิโย “อุโปสถมฺปิ น ปุจฺฉนฺติ, โอวาทมฺปิ น ยาจนฺตี”ติ.\csromanspacer{} สจฺจํ ภควาติ.\csromanspacer{} วิครหิ พุทฺโธ ภควา ฯเปฯ.\csromanspacer{} กถํ หิ นาม ภิกฺขเว ภิกฺขุนิโย อุโปสถมฺปิ น ปุจฺฉิสฺสนฺติ, โอวาทมฺปิ น ยาจิสฺสนฺติ.\csromanspacer{} เนตํ ภิกฺขเว อปฺปสนฺนานํ วา ปสาทาย ฯเปฯ.\csromanspacer{} เอวญฺจ ปน ภิกฺขเว ภิกฺขุนิโย อิมํ สิกฺขาปทํ อุทฺทิสนฺตุ\csromandash{}}

\proseitem{1059}{\textbf{“อนฺวทฺธมาสํ ภิกฺขุนิยา ภิกฺขุสํฆโต เทฺว ธมฺมา ปจฺจาสีสิตพฺพา อุโปสถปุจฺฉกญฺจ โอวาทูปสงฺกมนญฺจ, ตํ อติกฺกาเมนฺติยา ปาจิตฺติยนฺ”}ติ.}

\proseitem{1060}{\csromanfolio{415}\textbf{อนฺวทฺธมาสนฺ}ติ อนุโปสถิกํ.\csromanspacer{}\textbf{ อุโปสโถ} นาม เทฺว อุโปสถา จาตุทฺทสิโก จ ปนฺนรสิโก จ.}

\prose{\textbf{โอวาโท} นาม อฏฺฐ ครุธมฺมา. “อุโปสถมฺปิ น ปุจฺฉิสฺสามิ โอวาทมฺปิ น ยาจิสฺสามี”ติ ธุรํ นิกฺขิตฺตมตฺเต อาปตฺติ ปาจิตฺติยสฺส.}

\proseitemwithcloserruled{1061}{อนาปตฺติ สติ อนฺตราเย, ปริเยสิตฺวา ทุติยิกํ ภิกฺขุนิํ น ลภติ, คิลานาย อาปทาสุ อุมฺมตฺติกาย อาทิกมฺมิกายาติ.}{นวมสิกฺขาปทํ นิฏฺฐิตํ.}
\csromanmatikaanchor{mtk.214}
\csromantocmarkref{subsection}{10. ทสมสิกฺขาปท}{mtk.214}
\bookmark[dest={mtk.214},level=2]{10. ทสมสิกฺขาปท}
\csromanheader{2}{6. \textbf{อารามวคฺค 1}0. ทสมสิกฺขาปท}
\proseitem{1062}{เตน สมเยน พุทฺโธ ภควา สาวตฺถิยํ วิหรติ เชตวเน อนาถปิณฺฑิกสฺส อาราเม.\csromanspacer{} เตน โข ปน สมเยน อญฺญตรา ภิกฺขุนี ปสาเข ชาตํ คณฺฑํ ปุริเสน สทฺธิํ เอเกเนกา เภทาเปสิ.\csromanspacer{} อถ โข โส ปุริโส ตํ ภิกฺขุนิํ ทูเสตุํ อุปกฺกมิ, สา วิสฺสรมกาสิ.\csromanspacer{} ภิกฺขุนิโย อุปธาวิตฺวา ตํ ภิกฺขุนิํ เอตทโวจุํ “กิสฺส ตฺวํ อยฺเย วิสฺสรมกาสี”ติ.\csromanspacer{} อถ โข สา ภิกฺขุนี ภิกฺขุนีนํ เอตมตฺถํ อาโรเจสิ.\csromanspacer{} ยา ตา ภิกฺขุนิโย อปฺปิจฺฉา ฯเปฯ ตา อุชฺชฺยายนฺติ ขิยฺยนฺติ วิปาเจนฺติ “กถํ หิ นาม ภิกฺขุนี ปสาเข ชาตํ คณฺฑํ ปุริเสน สทฺธิํ เอเกเนกา เภทาเปสฺสตี”ติ ฯเปฯ.\csromanspacer{} สจฺจํ กิร ภิกฺขเว ภิกฺขุนี ปสาเข ชาตํ คณฺฑํ ปุริเสน สทฺธิํ เอเกเนกา เภทาเปตีติ\footnote{เภทาเปสีติ (ก)}.\csromanspacer{} สจฺจํ ภควาติ.\csromanspacer{} วิครหิ พุทฺโธ ภควา ฯเปฯ.\csromanspacer{} กถํ หิ นาม ภิกฺขเว ภิกฺขุนี ปสาเข ชาตํ คณฺฑํ ปุริเสน สทฺธิํ เอเกเนกา เภทาเปสฺสติ.\csromanspacer{} เนตํ ภิกฺขเว อปฺปสนฺนานํ วา ปสาทาย ฯเปฯ.\csromanspacer{} เอวญฺจ ปน ภิกฺขเว ภิกฺขุนิโย อิมํ สิกฺขาปทํ อุทฺทิสนฺตุ\csromandash{}}

\proseitem{1063}{\textbf{“ยา ปน ภิกฺขุนี ปสาเข ชาตํ คณฺฑํ วา รุธิตํ วา อนปโลเกตฺวา สํฆํ วา คณํ วา ปุริเสน สทฺธิํ เอเกเนกา เภทาเปยฺย วา ผาลาเปยฺย วา โธวาเปยฺย วา อาลิมฺปาเปยฺย วา พนฺธาเปยฺย วา โมจาเปยฺย วา, ปาจิตฺติยนฺ”}ติ.}

\proseitem{1064}{\csromanfolio{416}\textbf{ยา ปนา}ติ ยา ยาทิสา ฯเปฯ.\csromanspacer{} \textbf{ภิกฺขุนี}ติ ฯเปฯ อยํ อิมสฺมิํ อตฺเถ อธิปฺเปตา ภิกฺขุนีติ.}

\prose{\textbf{ปสาขํ} นาม อโธนาภิ อุพฺภชาณุมณฺฑลํ.\csromanspacer{}\textbf{ ชาต}นฺติ ตตฺถ ชาตํ.\csromanspacer{} \textbf{คณฺโฑ} นาม โย โกจิ คณฺโฑ.\csromanspacer{} \textbf{รุธิตํ} นาม ยํ กิญฺจิ วณํ.\csromanspacer{} \textbf{อนปโลเกตฺวา}ติ อนาปุจฺฉา.\csromanspacer{} \textbf{สํโฆ} นาม ภิกฺขุนิสํโฆ วุจฺจติ.\csromanspacer{} \textbf{คโณ} นาม สมฺพหุลา ภิกฺขุนิโย วุจฺจนฺติ.\csromanspacer{} \textbf{ปุริโส} นาม มนุสฺสปุริโส น ยกฺโข น เปโต น ติรจฺฉานคโต วิญฺญู ปฏิพโล ทูเสตุํ.\csromanspacer{} \textbf{สทฺธิ}นฺติ เอกโต.\csromanspacer{} \textbf{เอเกเนกา}ติ ปุริโส เจว โหติ ภิกฺขุนี จ.}

\proseitem{1065}{“ภินฺทา”ติ อาณาเปติ, อาปตฺติ ทุกฺกฏสฺส.\csromanspacer{} ภินฺเน อาปตฺติ ปาจิตฺติยสฺส. “ผาเลหี”ติ อาณาเปติ, อาปตฺติ ทุกฺกฏสฺส.\csromanspacer{} ผาลิเต อาปตฺติ ปาจิตฺติยสฺส. “โธวา”ติ อาณาเปติ, อาปตฺติ ทุกฺกฏสฺส.\csromanspacer{} โธวิเต\footnote{โธเต (สฺยา)} อาปตฺติ ปาจิตฺติยสฺส. “อาลิมฺปา”ติ อาณาเปติ, อาปตฺติ ทุกฺกฏสฺส.\csromanspacer{} ลิตฺเต\footnote{อาลิตฺเต (สฺยา)} อาปตฺติ ปาจิตฺติยสฺส. “พนฺธาหี”ติ อาณาเปติ, อาปตฺติ ทุกฺกฏสฺส.\csromanspacer{} พทฺเธ อาปตฺติ ปาจิตฺติยสฺส. “โมเจหี”ติ อาณาเปติ, อาปตฺติ ทุกฺกฏสฺส.\csromanspacer{} มุตฺเต อาปตฺติ ปาจิตฺติยสฺส.}

\proseitemwithcloser{1066}{อนาปตฺติ อปโลเกตฺวา เภทาเปติ วา ผาลาเปติ วา โธวาเปติ วา อาลิมฺปาเปติ วา พนฺธาเปติ วา โมจาเปติ วา, ยา กาจิ วิญฺญู ทุติยิกา\footnote{ยา กาจิ วิญฺญู ทุติยา (สฺยา), โย โกจิ วิญฺญู ทุติโย (ก)} โหติ, อุมฺมตฺติกาย อาทิกมฺมิกายาติ.}{ทสมสิกฺขาปทํ นิฏฺฐิตํ.}
\csromanheader{2}{อารามวคฺโค ฉฏฺโฐ.}
\csromansectionrule
\csromanmatikaanchor{mtk.215}
\csromanmatikaanchor{mtk.216}
\csromantocmarknopage{section}{7. คพฺภินีวคฺค}{mtk.215}
\bookmark[dest={mtk.215},level=1]{7. คพฺภินีวคฺค}
\csromantocmarkref{subsection}{1. ปฐมสิกฺขาปท}{mtk.216}
\bookmark[dest={mtk.216},level=2]{1. ปฐมสิกฺขาปท}
\setcsromanheadh{7. คพฺภินีวคฺค}
\csromanheader{2}{7. \textbf{คพฺภินีวคฺค 1. ปฐมสิกฺขาปท}}
\proseitem{1067}{เตน สมเยน พุทฺโธ ภควา สาวตฺถิยํ วิหรติ เชตวเน อนาถปิณฺฑิกสฺส อาราเม.\csromanspacer{} เตน โข ปน สมเยน ภิกฺขุนิโย คพฺภินิํ วุฏฺฐาเปนฺติ.\csromanspacer{} สา ปิณฺฑาย จรติ.\csromanspacer{} มนุสฺสา เอวมาหํสุ “เทถายฺยาย \csromanfolio{417}\textbf{ภิกฺขํ, ครุภารา}\footnote{ครุคพฺภา (สฺยา)}\textbf{ อยฺยา”ติ.}\csromanspacer{}\textbf{ มนุสฺสา อุชฺฌายนฺติ ขิยฺยนฺติ วิปาเจนฺติ} “\textbf{กถํ หิ นาม ภิกฺขุนิโย คพฺภินิํ วุฏฺฐาเปสฺสนฺตี”ติ.}\csromanspacer{}\textbf{ อสฺโสสุํ โข} ภิกฺขุนิโย เตสํ มนุสฺสานํ อุชฺฌายนฺตานํ ขิยฺยนฺตานํ วิปาเจนฺตานํ.\csromanspacer{} ยา ตา ภิกฺขุนิโย อปฺปิจฺฉา ฯเปฯ ตา อุชฺฌายนฺติ ขิยฺยนฺติ วิปาเจนฺติ “กถํ หิ นาม ภิกฺขุนิโย คพฺภินิํ วุฏฺฐาเปสฺสนฺตี”ติ ฯเปฯ.\csromanspacer{} สจฺจํ กิร ภิกฺขเว ภิกฺขุนิโย คพฺภินิํ วุฏฺฐาเปนฺตีติ.\csromanspacer{} สจฺจํ ภควาติ.\csromanspacer{} วิครหิ พุทฺโธ ภควา ฯเปฯ.\csromanspacer{} กถํ หิ นาม ภิกฺขเว ภิกฺขุนิโย คพฺภินิํ วุฏฺฐาเปสฺสนฺติ.\csromanspacer{} เนตํ ภิกฺขเว อปสนฺนานํ วา ปสาทาย ฯเปฯ.\csromanspacer{} เอวญฺจ ปน ภิกฺขเว ภิกฺขุนิโย อิมํ สิกฺขาปทํ อุทฺทิสนฺตุ\csromandash{}}

\proseitem{1068}{\textbf{“ยา ปน ภิกฺขุนี คพฺภินิํ วุฏฺฐาเปยฺย, ปาจิตฺติยนฺ”}ติ.}

\proseitem{1069}{\textbf{ยา ปนา}ติ ยา ยาทิสา ฯเปฯ.\csromanspacer{} \textbf{ภิกฺขุนี}ติ ฯเปฯ อยํ อิมสฺมิํ อตฺเถ อธิปฺเปตา ภิกฺขุนีติ.}

\prose{\textbf{คพฺภินี} นาม อาปนฺนสตฺตา วุจฺจติ.\csromanspacer{} \textbf{วุฏฺฐาเปยฺยา}ติ อุปสมฺปาเทยฺย.}

\prose{“วุฏฺฐาเปสฺสามี”ติ คณํ วา อาจรินิํ วา ปตฺตํ วา จีวรํ วา ปริเยสติ, สีมํ วา สมฺมนฺนติ, อาปตฺติ ทุกฺกฏสฺส.\csromanspacer{} ญตฺติยา ทุกฺกฏํ.\csromanspacer{} ทฺวีหิ กมฺมวาจาหิ ทุกฺกฏา.\csromanspacer{} กมฺมวาจาปริโยสาเน อุปชฺฌายาย อาปตฺติ ปาจิตฺติยสฺส.\csromanspacer{} คณสฺส จ อาจรินิยา จ อาปตฺติ ทุกฺกฏสฺส.}

\proseitem{1070}{คพฺภินิยา คพฺภินิสญฺญา วุฏฺฐาเปติ, อาปตฺติ ปาจิตฺติยสฺส.\csromanspacer{} คพฺภินิยา เวมติกา วุฏฺฐาเปติ, อาปตฺติ ทุกฺกฏสฺส.\csromanspacer{} คพฺภินิยา อคพฺภินิสญฺญา วุฏฺฐาเปติ, อนาปตฺติ.\csromanspacer{} อคพฺภินิยา คพฺภินิสญฺญา, อาปตฺติ ทุกฺกฏสฺส.\csromanspacer{} อคพฺภินิยา เวมติกา, อาปตฺติ ทุกฺกฏสฺส.\csromanspacer{} อคพฺภินิยา อคพฺภินิสญฺญา, อนาปตฺติ.}

\proseitemwithcloser{1071}{อนาปตฺติ คพฺภินิํ อคพฺภินิสญฺญา วุฏฺฐาเปติ, อคพฺภินิํ อคพฺภินิสญฺญา วุฏฺฐาเปติ, อุมฺมตฺติกาย อาทิกมฺมิกายาติ.}{ปฐมสิกฺขาปทํ นิฏฺฐิตํ.}
\csromanmatikaanchor{mtk.217}
\csromantocmarkref{subsection}{2. ทุติยสิกฺขาปท}{mtk.217}
\bookmark[dest={mtk.217},level=2]{2. ทุติยสิกฺขาปท}
\csromanheader{2}{7. \textbf{คพฺภินีวคฺค 2. ทุติยสิกฺขาปท}}
\proseitem{1072}{\csromanfolio{418}เตน สมเยน พุทฺโธ ภควา สาวตฺถิยํ วิหรติ เชตวเน อนาถปิณฺฑิกสฺส อาราเม.\csromanspacer{} เตน โข ปน สมเยน ภิกฺขุนิโย ปายนฺติํ วุฏฺฐาเปนฺติ.\csromanspacer{} สา ปิณฺฑาย จรติ.\csromanspacer{} มนุสฺสา เอวมาหํสุ “เทถายฺยาย ภิกฺขํ, สทุติยิกา อยฺยา”ติ.\csromanspacer{} มนุสฺสา อุชฺฌายนฺติ ขิยฺยนฺติ วิปาเจนฺติ “กถํ หิ นาม ภิกฺขุนิโย ปายนฺติํ วุฏฺฐาเปสฺสนฺตี”ติ.\csromanspacer{} อสฺโสสุํ โข ภิกฺขุนิโย เตสํ มนุสฺสานํ อุชฺฌายนฺตานํ ขิยฺยนฺตานํ วิปาเจนฺตานํ.\csromanspacer{} ยา ตา ภิกฺขุนิโย อปฺปิจฺฉา ฯเปฯ ตา อุชฺฌายนฺติ ขิยฺยนฺติ วิปาเจนฺติ “กถํ หิ นาม ภิกฺขุนิโย ปายนฺติํ วุฏฺฐาเปสฺสนฺตี”ติ ฯเปฯ.\csromanspacer{} สจฺจํ กิร ภิกฺขเว ภิกฺขุนิโย ปายนฺติํ วุฏฺฐาเปนฺตีติ.\csromanspacer{} สจฺจํ ภควาติ.\csromanspacer{} วิครหิ พุทฺโธ ภควา ฯเปฯ.\csromanspacer{} กถํ หิ นาม ภิกฺขเว ภิกฺขุนิโย ปายนฺติํ วุฏฺฐาเปสฺสนฺติ.\csromanspacer{} เนตํ ภิกฺขเว อปฺปสนฺนานํ วา ปสาทาย ฯเปฯ.\csromanspacer{} เอวญฺจ ปน ภิกฺขเว ภิกฺขุนิโย อิมํ สิกฺขาปทํ อุทฺทิสนฺตุ\csromandash{}}

\proseitem{1073}{\textbf{“ยา ปน ภิกฺขุนี ปายนฺติํ วุฏฺฐาเปยฺย, ปาจิตฺติยนฺ”}ติ.}

\proseitem{1074}{\textbf{ยา ปนา}ติ ยา ยาทิสา ฯเปฯ.\csromanspacer{} \textbf{ภิกฺขุนี}ติ ฯเปฯ อยํ อิมสฺมิํ อตฺเถ อธิปฺเปตา ภิกฺขุนีติ.}

\prose{\textbf{ปายนฺตี} นาม มาตา วา โหติ\footnote{โหตุ (ก)} ธาติ\footnote{ธาตี (สี, สฺยา)} วา.}

\prose{\textbf{วุฏฺฐาเปยฺยา}ติ อุปสมฺปาเทยฺย.}

\prose{“วุฏฺฐาเปสฺสามี”ติ คณํ วา อาจรินิํ วา ปตฺตํ วา จีวรํ วา ปริเยสติ, สีมํ วา สมฺมนฺนติ, อาปตฺติ ทุกฺกฏสฺส.\csromanspacer{} ญตฺติยา ทุกฺกฏํ.\csromanspacer{} ทฺวีหิ กมฺมวาจาหิ ทุกฺกฏา.\csromanspacer{} กมฺมวาจาปริโยสาเน อุปชฺฌายาย อาปตฺติ ปาจิตฺติยสฺส.\csromanspacer{} คณสฺส จ อาจรินิยา จ อาปตฺติ ทุกฺกฏสฺส.}

\proseitem{1075}{ปายนฺติยา ปายนฺติสญฺญา วุฏฺฐาเปติ, อาปตฺติ ปาจิตฺติยสฺส.\csromanspacer{} ปายนฺติยา เวมติกา วุฏฺฐาเปติ, อาปตฺติ ทุกฺกฏสฺส.\csromanspacer{} ปายนฺติยา อปายนฺติสญฺญา วุฏฺฐาเปติ, อนาปตฺติ.\csromanspacer{} อปายนฺติยา ปายนฺติสญฺญา, อาปตฺติ ทุกฺกฏสฺส.\csromanspacer{} อปายนฺติยา เวมติกา, อาปตฺติ ทุกฺกฏสฺส.\csromanspacer{} อปายนฺติยา อปายนฺติสญฺญา, อนาปตฺติ.}

\proseitemwithcloserruled{1076}{\csromanfolio{419}อนาปตฺติ ปายนฺติํ อปายนฺติสญฺญา วุฏฺฐาเปติ, อปายนฺติํ อปายนฺติสญฺญา วุฏฺฐาเปติ, อุมฺมตฺติกาย อาทิกมฺมิกายาติ.}{ทุติยสิกฺขาปทํ นิฏฺฐิตํ.}
\csromanmatikaanchor{mtk.218}
\csromantocmarkref{subsection}{3. ตติยสิกฺขาปท}{mtk.218}
\bookmark[dest={mtk.218},level=2]{3. ตติยสิกฺขาปท}
\csromanheader{2}{7. \textbf{คพฺภินีวคฺค 3. ตติยสิกฺขาปท}}
\proseitem{1077}{เตน สมเยน พุทฺโธ ภควา สาวตฺถิยํ วิหรติ เชตวเน อนาถปิณฺฑิกสฺส อาราเม.\csromanspacer{} เตน โข ปน สมเยน ภิกฺขุนิโย เทฺว วสฺสานิ ฉสุ ธมฺเมสุ อสิกฺขิตสิกฺขํ สิกฺขมานํ วุฏฺฐาเปนฺติ, ตา พาลา โหนฺติ อพฺยตฺตา, น ชานนฺติ กปฺปิยํ วา อกปฺปิยํ วา.\csromanspacer{} ยา ตา ภิกฺขุนิโย อปฺปิจฺฉา ฯเปฯ ตา อุชฺฌายนฺติ ขิยฺยนฺติ วิปาเจนฺติ “กถํ หิ นาม ภิกฺขุนิโย เทฺว วสฺสานิ ฉสุ ธมฺเมสุ อสิกฺขิตสิกฺขํ สิกฺขมานํ วุฏฺฐาเปสฺสนฺตี”ติ ฯเปฯ.\csromanspacer{} สจฺจํ กิร ภิกฺขเว ภิกฺขุนิโย เทฺว วสฺสานิ ฉสุ ธมฺเมสุ อสิกฺขิตสิกฺขํ สิกฺขมานํ วุฏฺฐาเปนฺตี”ติ.\csromanspacer{} สจฺจํ ภควาติ.\csromanspacer{} วิครหิ พุทฺโธ ภควา ฯเปฯ.\csromanspacer{} กถํ หิ นาม ภิกฺขเว ภิกฺขุนิโย เทฺว วสฺสานิ ฉสุ ธมฺเมสุ อสิกฺขิตสิกฺขํ สิกฺขมานํ วุฏฺฐาเปสฺสนฺตี.\csromanspacer{} เนตํ ภิกฺขเว อปฺปสนฺนานํ วา ปสาทาย ฯเปฯ วิครหิตฺวา ฯเปฯ ธมฺมิํ กถํ กตฺวา ภิกฺขู อามนฺเตสิ อนุชานามิ ภิกฺขเว สิกฺขมานาย เทฺว วสฺสานิ ฉสุ ธมฺเมสุ สิกฺขาสมฺมุติํ ทาตุํ.\csromanspacer{} เอวญฺจ ปน ภิกฺขเว ทาตพฺพา, ตาย สิกฺขมานาย สํฆํ อุปสงฺกมิตฺวา เอกํสํ อุตฺตราสงฺคํ กริตฺวา ภิกฺขุนีนํ ปาเท วนฺทิตฺวา อุกฺกุฏิกํ นิสีทิตฺวา อญฺชลิํ ปคฺคเหตฺวา เอวมสฺส วจนีโย “อหํ อยฺเย อิตฺถนฺนามา อิตฺถนฺนามาย อยฺยาย สิกฺขมานา สํฆํ เทฺว วสฺสานิ ฉสุ ธมฺเมสุ สิกฺขาสมฺมุติํ ยาจามี”ติ.\csromanspacer{} ทุติยมฺปิ ยาจิตพฺพา.\csromanspacer{} ตติยมฺปิ ยาจิตพฺพา.\csromanspacer{} พฺยตฺตาย ภิกฺขุนิยา ปฏิพลาย สํโฆ ญาเปตพฺโพ\csromandash{}}

\proseitem{1078}{“สุณาตุ เม อยฺเย สํโฆ, อยํ อิตฺถนฺนามา อิตฺถนฺนามาย อยฺยาย สิกฺขมานา สํฆํ เทฺว วสฺสานิ ฉสุ ธมฺเมสุ สิกฺขาสมฺมุติํ ยาจติ, ยทิ สํฆสฺส ปตฺตกลฺลํ, สํโฆ อิตฺถนฺนามาย สิกฺขมานาย เทฺว วสฺสานิ ฉสุ ธมฺเมสุ สิกฺขาสมฺมุติํ ทเทยฺย, เอสา ญตฺติ.}

\prose{\csromanfolio{420}สุณาตุ เม อยฺเย สํโฆ, อยํ อิตฺถนฺนามา อิตฺถนฺนามาย อยฺยาย สิกฺขมานา สํฆํ เทฺว วสฺสานิ ฉสุ ธมฺเมสุ สิกฺขาสมฺมุติํ ยาจติ, สํโฆ อิตฺถนฺนามาย สิกฺขมานาย เทฺว วสฺสานิ ฉสุ ธมฺเมสุ สิกฺขาสมฺมุติํ เทติ, ยสฺสา อยฺยาย ขมติ อิตฺถนฺนามาย สิกฺขมานาย เทฺว วสฺสานิ ฉสุ ธมฺเมสุ สิกฺขาสมฺมุติยา ทานํ, สา ตุณฺหสฺส, ยสฺสา นกฺขมติ, สา ภาเสยฺย.}

\prose{ทินฺนา สํเฆน อิตฺถนฺนามาย สิกฺขมานาย เทฺว วสฺสานิ ฉสุ ธมฺเมสุ สิกฺขาสมฺมุติ, ขมติ สํฆสฺส, ตสฺมา ตุณฺหี, เอวเมตํ ธารยามี”ติ.}

\proseitem{1079}{สา สิกฺขมานา “เอวํ วเทหี”ติ วตฺตพฺพา “ปาณาติปาตา เวรมณิํ เทฺว วสฺสานิ อวีติกฺกมฺม สมาทานํ สมาทิยามิ.\csromanspacer{} อทินฺนทานา เวรมณิํ เทฺว วสฺสานิ อวีติกฺกมฺม สมาทานํ สมาทิยามิ.\csromanspacer{} อพฺรหฺมจริยา เวรมณิํ เทฺว วสฺสานิ อวีติกฺกมฺม สมาทานํ สมาทิยามิ.\csromanspacer{} มุสาวาทา เวรมณิํ เทฺว วสฺสานิ อวีติกฺกมฺม สมาทานํ สมาทิยามิ.\csromanspacer{} สุราเมรยมชฺชปฺปมาทฏฺฐานา เวรมณิํ เทฺว วสฺสานิ อวีติกฺกมฺม สมาทานํ สมาทิยามิ.\csromanspacer{} วิกาลโภชนา เวรมณิํ เทฺว วสฺสานิ อวีติกฺกมฺม สมาทานํ สมาทิยามี”ติ.}

\prose{อถ โข ภควา ตา ภิกฺขุนิโย อเนกปริยาเยน วิครหิตฺวา ทุพฺภรตาย ฯเปฯ.\csromanspacer{} เอวญฺจ ปน ภิกฺขเว ภิกฺขุนิโย อิมํ สิกฺขาปทํ อุทฺทิสนฺตุ\csromandash{}}

\proseitem{1080}{\textbf{“ยา ปน ภิกฺขุนี เทฺว วสฺสานิ ฉสุ ธมฺเมสุ อสิกฺขิตสิกฺขํ สิกฺขมานํ วุฏฺฐาเปยฺย, ปาจิตฺติยนฺ”}ติ.}

\proseitem{1081}{\textbf{ยา ปนา}ติ ยา ยาทิยา ฯเปฯ.\csromanspacer{} \textbf{ภิกฺขุนีตฺ}อิ ฯเปฯ อยํ อิมสฺมิํ อตฺเถ อธิปฺเปตา ภิกฺขุนีติ.}

\prose{\textbf{เทฺว วสฺสานี}ติ เทฺว สํวจฺฉรานิ.\csromanspacer{} \textbf{อสิกฺขิตสิกฺขา} นาม สิกฺขา วา น ทินฺนา โหติ, ทินฺนา วา สิกฺขา กุปิตา.\csromanspacer{} \textbf{วุฏฺฐาเปยฺยา}ติ อุปสมฺปาเทยฺย.}

\prose{“วุฏฺฐาเปสฺสามี”ติ คณํ วา อาจรินิํ วา ปตฺตํ วา จีวรํ วา ปริเยสติ, สีมํ วา สมฺมนฺนติ, อาปตฺติ ทุกฺกฏสฺส.\csromanspacer{} ญตฺติยา ทุกฺกฏํ.\csromanspacer{} ทฺวีหิ กมฺมวาจาหิ ทุกฺกฏา.\csromanspacer{} กมฺมวาจาปริโยสาเน อุปชฺฌายาย อาปตฺติ ปาจิตฺติยสฺส.\csromanspacer{} คณสฺส จ อาจรินิยา จ อาปตฺติ ทุกฺกฏสฺส.}

\proseitem{1082}{\csromanfolio{421}ธมฺมกมฺเม ธมฺมกมฺมสญฺญา วุฏฺฐาเปติ, อาปตฺติ ปาจิตฺติยสฺส.\csromanspacer{} ธมฺมกมฺเม เวมติกา วุฏฺฐาเปติ, อาปตฺติ ปาจิตฺติยสฺส.\csromanspacer{} ธมฺมกมฺเม อธมฺมกมฺมสญฺญา วุฏฺฐาเปติ, อาปตฺติ ปาจิตฺติยสฺส.}

\prose{อธมฺมกมฺเม ธมฺมกมฺมสญฺญา, อาปตฺติ ทุกฺกฏสฺส.\csromanspacer{} อธมฺมกมฺเม เวมติกา, อาปตฺติ ทุกฺกฏสฺส.\csromanspacer{} อธมฺมกมฺเม อธมฺมกมฺมสญฺญา, อาปตฺติ ทุกฺกฏสฺส.}

\proseitemwithcloserruled{1083}{อนาปตฺติ เทฺว วสฺสานิ ฉสุ ธมฺเมสุ สิกฺขิตสิกฺขํ สิกฺขมานํ วุฏฺฐาเปติ, อุมฺมตฺติกาย อาทิกมฺมิกายาติ.}{ตติยสิกฺขาปทํ นิฏฺฐิตํ.}
\csromanmatikaanchor{mtk.219}
\csromantocmarkref{subsection}{4. จตุตฺถสิกฺขาปท}{mtk.219}
\bookmark[dest={mtk.219},level=2]{4. จตุตฺถสิกฺขาปท}
\csromanheader{2}{7. \textbf{คพฺภินีวคฺค 4. จตุตฺถสิกฺขาปท}}
\proseitem{1084}{เตน สมเยน พุทฺโธ ภควา สาวตฺถิยํ วิหรติ เชตวเน อนาถปิณฺฑิกสฺส อาราเม.\csromanspacer{} เตน โข ปน สมเยน ภิกฺขุนิโย เทฺว วสฺสานิ ฉสุ ธมฺเมสุ สิกฺขิตสิกฺขํ สิกฺขมานํ สํเฆน อสมฺมตํ วุฏฺฐาเปนฺติ.\csromanspacer{} ภิกฺขุนิโย เอวมาหํสุ “เอถ สิกฺขมานา อิมํ ชานาถ อิมํ เทถ อิมํ อาหรถ อิมินา อตฺโถ อิมํ กปฺปิยํ กโรถา”ติ.\csromanspacer{} ตา เอวมาหํสุ “น มยํ อยฺเย สิกฺขมานา, ภิกฺขุนิโย มยนฺ”ติ.\csromanspacer{} ยา ตา ภิกฺขุนิโย อปฺปิจฺฉา ฯเปฯ ตา อุชฺฌายนฺติ ขิยฺยนฺติ วิปาเจนฺติ “กถํ หิ นาม ภิกฺขุนิโย เทฺว วสฺสานิ ฉสุ ธมฺเมสุ สิกฺขิตสิกฺขํ สิกฺขมานํ สํเฆน อสมฺมตํ วุฏฺฐาเปสฺสนฺตี”ติ ฯเปฯ.\csromanspacer{} สจฺจํ กิร ภิกฺขเว ภิกฺขุนิโย เทฺว วสฺสานิ ฉสุ ธมฺเมสุ สิกฺขิตสิกฺขํ สิกฺขมานํ สํเฆน อสมฺมตํ วุฏฺฐาเปนฺตีติ.\csromanspacer{} สจฺจํ ภควาติ.\csromanspacer{} วิครหิ พุทฺโธ ภควา ฯเปฯ.\csromanspacer{} กถํ หิ นาม ภิกฺขเว ภิกฺขุนิโย เทฺว วสฺสานิ ฉสุ ธมฺเมสุ สิกฺขิตสิกฺขํ สิกฺขมานํ สํเฆน อสมฺมตํ วุฏฺฐาเปสฺสนฺติ.\csromanspacer{} เนตํ ภิกฺขเว อปฺปสนฺนานํ วา ปสาทาย ฯเปฯ วิครหิตฺวา ฯเปฯ ธมฺมิํ กถํ กตฺวา ภิกฺขู อามนฺเตสิ อนุชานามิ ภิกฺขเว เทฺว วสฺสานิ ฉสุ ธมฺเมสุ สิกฺขิตสิกฺขาย สิกฺขมานาย วุฏฺฐานสมฺมุติํ ทาตุํ.\csromanspacer{} เอวญฺจ ปน ภิกฺขเว ทาตพฺพา, ตาย เทฺว วสฺสานิ ฉสุ ธมฺเมสุ สิกฺขิตสิกฺขาย สิกฺขมานาย สํฆํ อุปสงฺกมิตฺวา เอกํสํ อุตฺตราสงฺคํ กริตฺวา ภิกฺขุนีนํ ปาเท วนฺทิตฺวา \csromanfolio{422}อุกฺกุฏิกํ นิสีทิตฺวา อญฺชลิํ ปคฺคเหตฺวา เอวมสฺส วจนีโย “อหํ อยฺเย อิตฺถนฺนามา อิตฺถนฺนามาย อยฺยาย เทฺว วสฺสานิ ฉสุ ธมฺเมสุ สิกฺขิตสิกฺขา สิกฺขมานา สํฆํ วุฏฺฐานสมฺมุติํ ยาจามี”ติ.\csromanspacer{} ทุติยมฺปิ ยาจิตพฺพา.\csromanspacer{} ตติยมฺปิ ยาจิตพฺพา.\csromanspacer{} พฺยตฺตาย ภิกฺขุนิยา ปฏิพลาย สํโฆ ญาเปตพฺโพ\csromandash{}}

\proseitem{1085}{“สุณาตุ เม อยฺเย สํโฆ, อยํ อิตฺถนฺนามา อิตฺถนฺนามาย อยฺยาย เทฺว วสฺสานิ ฉสุ ธมฺเมสุ สิกฺขิตสิกฺขา สิกฺขมานา สํฆํ วุฏฺฐานสมฺมุติํ ยาจติ, ยทิ สํฆสฺส ปตฺตกลฺลํ, สํโฆ อิตฺถนฺนามาย เทฺว วสฺสานิ ฉสุ ธมฺเมสุ สิกฺขิตสิกฺขาย สิกฺขมานาย วุฏฺฐานสมฺมุติํ ทเทยฺย, เอสา ญตฺติ.}

\prose{สุณาตุ เม อยฺเย สํโฆ, อยํ อิตฺถนฺนามา อิตฺถนฺนามาย อยฺยาย เทฺว วสฺสานิ ฉสุ ธมฺเมสุ สิกฺขิตสิกฺขา สิกฺขมานา สํฆํ วุฏฺฐานสมฺมุติํ ยาจติ, สํโฆ อิตฺถนฺนามาย เทฺว วสฺสานิ ฉสุ ธมฺเมสุ สิกฺขิตสิกฺขาย สิกฺขมานาย วุฏฺฐานสมฺมุติํ เทติ, ยสฺสา อยฺยาย ขมติ อิตฺถนามาย เทฺว วสฺสานิ ฉสุ ธมฺเมสุ สิกฺขิตสิกฺขาย สิกฺขมานาย วุฏฺฐานสมฺมุติยา ทานํ, สา ตุณฺหสฺส, ยสฺสา นกฺขมติ, สา ภาเสยฺย.}

\prose{ทินฺนา สํเฆน อิตฺถนฺนามาย เทฺว วสฺสานิ ฉสุ ธมฺเมสุ สิกฺขิตสิกฺขาย สิกฺขมานาย วุฏฺฐานสมฺมุติ, ขมติ สํฆสฺส, ตสฺมา ตุณฺหี, เอวเมตํ ธารยามี”ติ.}

\prose{อถ โข ภควา ตา ภิกฺขุนิโย อเนกปริยาเยน วิครหิตฺวา ทุพฺภรตาย ฯเปฯ.\csromanspacer{} เอวญฺจ ปน ภิกฺขเว ภิกฺขุนิโย อิมํ สิกฺขาปทํ อุทฺทิสนฺตุ\csromandash{}}

\proseitem{1086}{\textbf{“ยา ปน ภิกฺขุนี เทฺว วสฺสานิ ฉสุ ธมฺเมสุ สิกฺขิตสิกฺขํ สิกฺขมานํ สํเฆน อสมฺมตํ วุฏฺฐาเปยฺย, ปาจิตฺติยนฺ”}ติ.}

\proseitem{1087}{\textbf{ยา ปนา}ติ ยา ยาทิสา ฯเปฯ.\csromanspacer{} \textbf{ภิกฺขุนี}ติ ฯเปฯ อยํ อิมสฺมิํ อตฺเถ อธิปฺเปตา ภิกฺขุนีติ.}

\prose{\csromanfolio{423}\textbf{เทฺว วสฺสานี}ติ เทฺว สํวจฺฉรานิ.\csromanspacer{} \textbf{สิกฺขิตสิกฺขา} นาม ฉสุ ธมฺเมสุ สิกฺขิตสิกฺขา.\csromanspacer{} \textbf{อสมฺมตา} นาม ญตฺติทุติเยน กมฺเมน วุฏฺฐานสมฺมุติ น ทินฺนา โหติ.\csromanspacer{} \textbf{วุฏฺฐาเปยฺยา}ติ อุปสมฺปาเทยฺย.}

\prose{“วุฏฺฐาเปสฺสามี”ติ คณํ วา อาจรินิํ วา ปตฺตํ วา จีวรํ วา ปริเยสติ, สีมํ วา สมฺมนฺนติ, อาปตฺติ ทุกฺกฏสฺส.\csromanspacer{} ญตฺติยา ทุกฺกฏํ.\csromanspacer{} ทฺวีหิ กมฺมวาจาหิ ทุกฺกฏา.\csromanspacer{} กมฺมวาจาปริโยสาเน อุปชฺฌายาย อาปตฺติ ปาจิตฺติยสฺส.\csromanspacer{} คณสฺส จ อาจรินิยา จ อาปตฺติ ทุกฺกฏสฺส.}

\proseitem{1088}{ธมฺมกมฺเม ธมฺมกมฺมสญฺญา วุฏฺฐาเปติ, อาปตฺติ ปาจิตฺติยสฺส.\csromanspacer{} ธมฺมกมฺเม เวมติกา วุฏฺฐาเปติ, อาปตฺติ ปาจิตฺติยสฺส.\csromanspacer{} ธมฺมกมฺเม อธมฺมกมฺมสญฺญา วุฏฺฐาเปติ, อาปตฺติ ปาจิตฺติยสฺส.}

\prose{อธมฺมกมฺเม ธมฺมกมฺมสญฺญา, อาปตฺติ ทุกฺกฏสฺส.\csromanspacer{} อธมฺมกมฺเม เวมติกา, อาปตฺติ ทุกฺกฏสฺส.\csromanspacer{} อธมฺมกมฺเม อธมฺมกมฺมสญฺญา, อาปตฺติ ทุกฺกฏสฺส.}

\proseitemwithcloserruled{1089}{อนาปตฺติ เทฺว วสฺสานิ ฉสุ ธมฺเมสุ สิกฺขิตสิกฺขํ สิกฺขมานํ สํเฆน สมฺมตํ วุฏฺฐาเปติ, อุมฺมตฺติกาย อาทิกมฺมิกายาติ.}{จตุตฺถสิกฺขาปทํ นิฏฺฐิตํ.}
\csromanmatikaanchor{mtk.220}
\csromantocmarkref{subsection}{5. ปญฺจมสิกฺขาปท}{mtk.220}
\bookmark[dest={mtk.220},level=2]{5. ปญฺจมสิกฺขาปท}
\csromanheader{2}{7. \textbf{คพฺภินีวคฺค 5. ปญฺจมสิกฺขาปท}}
\proseitem{1090}{เตน สมเยน พุทฺโธ ภควา สาวตฺถิยํ วิหรติ เชตวเน อนาถปิณฺฑิกสฺส อาราเม.\csromanspacer{} เตน โข ปน สมเยน ภิกฺขุนิโย อูนทฺวาทสวสฺสํ คิหิคตํ วุฏฺฐาเปนฺติ, ตา อกฺขมา โหนฺติ สีตสฺส อุณฺหสฺส ชิฆจฺฉาย ปิปาสาย ฑํส\-มกส\-วาตา\-ตป\-สรีสป\-สมฺผสฺสานํ, ทุรุตฺตานํ ทุราคตานํ วจนปถานํ, อุปฺปนฺนานํ สารีริกานํ เวทนานํ ทุกฺขานํ ติพฺพานํ ขรานํ กฏุกานํ อสาตานํ อมนาปานํ ปาณหรานํ อนธิวาสกชาติกา โหนฺติ.\csromanspacer{} ยา ตา ภิกฺขุนิโย อปฺปิจฺฉา ฯเปฯ ตา อุชฺฌายนฺติ ขิยฺยนฺติ วิปาเจนฺติ “กถํ หิ นาม ภิกฺขุนิโย อูนทฺวาทสวสฺสํ คิหิคตํ วุฏฺฐาเปสฺสนฺตี”ติ ฯเปฯ.\csromanspacer{} สจฺจํ กิร ภิกฺขเว ภิกฺขุนิโย อูนทฺวาทสวสฺสํ คิหิคตํ วุฏฺฐาเปนฺตีติ.\csromanspacer{} สจฺจํ ภควาติ. \csromanfolio{424}วิครหิ พุทฺโธ ภควา ฯเปฯ.\csromanspacer{} กถํ หิ นาม ภิกฺขเว ภิกฺขุนิโย อูนทฺวาทสวสฺสํ คิหิคตํ วุฏฺฐาเปสฺสนฺติ, อูนทฺวาทสวสฺสา\footnote{อูนทฺวาทสวสฺสา หิ (ก)} ภิกฺขเว คิหิคตา อกฺขมา โหติ สีตสฺส อุณฺหสฺส ชิฆจฺฉาย ปิปาสาย ฑํส\-มกส\-วาตา\-ตป\-สรีสป\-สมฺผสฺสานํ ทุรุตฺตานํ ทุราคตานํ วจนปถานํ, อุปฺปนฺนานํ สารีริกานํ เวทนานํ ทุกฺขานํ ติพฺพานํ ขรานํ กฏุกานํ อสาตานํ อมนาปานํ ปาณหรานํ อนธิวาสกชาติกา โหติ, ทฺวาทสวสฺสาว\footnote{ทฺวาทสวสฺสา จ (สฺยา, ก)} โข ภิกฺขเว คิหิคตา ขมา โหติ สีตสฺส อุณฺหสฺส ชิฆจฺฉาย ปิปาสาย ฑํส\-มกส\-วาตา\-ตป\-สรีสป\-สมฺผสฺสานํ ทุรุตฺตานํ ทุราคตานํ วจนปถานํ อุปฺปนฺนานํ สารีริกานํ เวทนานํ ทุกฺขานํ ติพฺพานํ ขรานํ กฏุกานํ อสาตานํ อมนาปานํ ปาณหรานํ อธิวาสกชาติกา โหติ.\csromanspacer{} เนตํ ภิกฺขเว อปฺปสนฺนานํ วา ปสาทาย ฯเปฯ.\csromanspacer{} เอวญฺจ ปน ภิกฺขเว ภิกฺขุนิโย อิมํ สิกฺขาปทํ อุทฺทิสนฺตุ\csromandash{}}

\proseitem{1091}{\textbf{“ยา ปน ภิกฺขุนี อูนทฺวาทสวสฺสํ คิหิคตํ วุฏฺฐาเปยฺย, ปาจิตฺติยนฺ”}ติ.}

\proseitem{1092}{\textbf{ยา ปนา}ติ ยา ยาทิสา ฯเปฯ.\csromanspacer{} \textbf{ภิกฺขุนี}ติ ฯเปฯ อยํ อิมสฺมิํ อตฺเถ อธิปฺเปตา ภิกฺขุนีติ.}

\prose{อู\textbf{นทฺวาทสวสฺสา} นาม อปฺปตฺตทฺวาทสวสฺสา.\csromanspacer{} \textbf{คิหิคตา} นาม ปุริสนฺตรคตา วุจฺจติ.\csromanspacer{} \textbf{วุฏฺฐาเปยฺยา}ติ อุปสมฺปาเทยฺย.}

\prose{“วุฏฺฐาเปสฺสามี”ติ คณํ วา อาจรินิํ วา ปตฺตํ วา จีวรํ วา ปริเยสติ, สีมํ วา สมฺมนฺนติ, อาปตฺติ ทุกฺกฏสฺส.\csromanspacer{} ญตฺติยา ทุกฺกฏํ.\csromanspacer{} ทฺวีหิ กมฺมวาจาหิ ทุกฺกฏา.\csromanspacer{} กมฺมวาจาปริโยสาเน อุปชฺฌายาย อาปตฺติ ปาจิตฺติยสฺส.\csromanspacer{} คณสฺส จ อาจรินิยา จ อาปตฺติ ทุกฺกฏสฺส.}

\proseitem{1093}{อูนทฺวาทสวสฺสาย อูนทฺวาทสวสฺสสญฺญา วุฏฺฐาเปติ, อาปตฺติ ปาจิตฺติยสฺส.\csromanspacer{} อูนทฺวาทสวสฺสาย เวมติกา วุฏฺฐาเปติ, อาปตฺติ ทุกฺกฏสฺส.\csromanspacer{} อูนทฺวาทสวสฺสาย ปริปุณฺณสญฺญา วุฏฺฐาเปติ, อนาปตฺติ.}

\prose{\csromanfolio{425}ปริปุณฺณ\-ทฺวาทสวสฺสาย อูนทฺวาทสวสฺสสญฺญา, อาปตฺติ ทุกฺกฏสฺส.\csromanspacer{} ปริปุณฺณ\-ทฺวาทสวสฺสาย เวมติกา, อาปตฺติ ทุกฺกฏสฺส.\csromanspacer{} ปริปุณฺณ\-ทฺวาทสวสฺสาย ปริปุณฺณสญฺญา, อนาปตฺติ.}

\proseitemwithcloserruled{1094}{อนาปตฺติ อูนทฺวาทสวสฺสํ ปริปุณฺณสญฺญา วุฏฺฐาเปติ, ปริปุณฺณทฺวาทสวสฺสํ ปริปุณฺณสญฺญา วุฏฺฐาเปติ, อุมฺมตฺติกาย อาทิกมฺมิกายาติ.}{ปญฺจมสิกฺขาปทํ นิฏฺฐิตํ.}
\csromanmatikaanchor{mtk.221}
\csromantocmarkref{subsection}{6. ฉฏฺฐสิกฺขาปท}{mtk.221}
\bookmark[dest={mtk.221},level=2]{6. ฉฏฺฐสิกฺขาปท}
\csromanheader{2}{7. \textbf{คพฺภินีวคฺค 6. ฉฏฺฐสิกฺขาปท}}
\proseitem{1095}{เตน สมเยน พุทฺโธ ภควา สาวตฺถิยํ วิหรติ เชตวเน อนาถปิณฺฑิกสฺส อาราเม.\csromanspacer{} เตน โข ปน สมเยน ภิกฺขุนิโย ปริปุณฺณทฺวาทสวสฺสํ คิหิคตํ เทฺว วสฺสานิ ฉสุ ธมฺเมสุ อสิกฺขิตสิกฺขํ วุฏฺฐาเปนฺติ, ตา พาลา โหนฺติ อพฺยตฺตา, น ชานนฺติ กปฺปิยํ วา อกปฺปิยํ วา.\csromanspacer{} ยา ตา ภิกฺขุนิโย อปฺปิจฺฉา ฯเปฯ ตา อุชฺฌายนฺติ ขิยฺยนฺติ วิปาเจนฺติ “กถํ หิ นาม ภิกฺขุนิโย ปริปุณฺณทฺวาทสวสฺสํ คิหิคตํ เทฺว วสฺสานิ ฉสุ ธมฺเมสุ อสิกฺขิตสิกฺขํ วุฏฺฐาเปสฺสนฺตี”ติ ฯเปฯ.\csromanspacer{} สจฺจํ กิร ภิกฺขเว ภิกฺขุนิโย ปริปุณฺณทฺวาทสวสฺสํ คิหิคตํ เทฺว วสฺสานิ ฉสุ ธมฺเมสุ อสิกฺขิตสิกฺขํ วุฏฺฐาเปนฺตีติ.\csromanspacer{} สจฺจํ ภควาติ.\csromanspacer{} วิครหิ พุทฺโธ ภควา ฯเปฯ.\csromanspacer{} กถํ หิ นาม ภิกฺขเว ภิกฺขุนิโย ปริปุณฺณทฺวาทสวสฺสํ คิหิคตํ เทฺว วสฺสานิ ฉสุ ธมฺเมสุ อสิกฺขิตสิกฺขํ วุฏฺฐาเปสฺสนฺติ.\csromanspacer{} เนตํ ภิกฺขเว อปฺปสนฺนานํ วา ปสาทาย ฯเปฯ วิครหิตฺวา ฯเปฯ ธมฺมิํ กถํ กตฺวา ภิกฺขู อามนฺเตสิ อนุชานามิ ภิกฺขเว ปริปุณฺณ\-ทฺวาทสวสฺสาย คิหิคตาย เทฺว วสฺสานิ ฉสุ ธมฺเมสุ สิกฺขาสมฺมุติํ ทาตุํ.\csromanspacer{} เอวญฺจ ปน ภิกฺขเว ทาตพฺพา, ตาย ปริปุณฺณ\-ทฺวาทสวสฺสาย คิหิคตาย สํฆํ อุปสงฺกมิตฺวา เอกํสํ อุตฺตราสงฺคํ กริตฺวา ภิกฺขุนีนํ ปาเท วนฺทิตฺวา อุกฺกุฏิกํ นิสีทิตฺวา อญฺชลิํ ปคฺคเหตฺวา เอวมสฺส วจนีโย “อหํ อยฺเย อิตฺถนฺนามา อิตฺถนฺนามาย อยฺยาย ปริปุณฺณทฺวาทสวสฺสา คิหิคตา สํฆํ เทฺว วสฺสานิ ฉสุ ธมฺเมสุ สิกฺขาสมฺมุติํ ยาจามี”ติ.\csromanspacer{} ทุติยมฺปิ ยาจิตพฺพา.\csromanspacer{} ตติยมฺปิ ยาจิตพฺพา.\csromanspacer{} พฺยตฺตาย ภิกฺขุนิยา ปฏิพลาย สํโฆ ญาเปตพฺโพ\csromandash{}}

\proseitem{1096}{\csromanfolio{426}“สุณาตุ เม อยฺเย สํโฆ, อยํ อิตฺถนฺนามา อิตฺถนฺนามาย อยฺยาย ปริปุณฺณทฺวาทสวสฺสา คิหิคตา สํฆํ เทฺว วสฺสานิ ฉสุ ธมฺเมสุ สิกฺขาสมฺมุติํ ยาจติ, ยทิ สํฆสฺส ปตฺตกลฺลํ, สํโฆ อิตฺถนฺนามาย ปริปุณฺณ\-ทฺวาทสวสฺสาย คิหิคตาย เทฺว วสฺสานิ ฉสุ ธมฺเมสุ สิกฺขาสมฺมุติํ ทเทยฺย, เอสา ญตฺติ.}

\prose{สุณาตุ เม อยฺเย สํโฆ, อยํ อิตฺถนฺนามา อิตฺถนฺนามาย อยฺยาย ปริปุณฺณทฺวาทสวสฺสา คิหิคตา สํฆํ เทฺว วสฺสานิ ฉสุ ธมฺเมสุ สิกฺขาสมฺมุติํ ยาจติ, สํโฆ อิตฺถนฺนามาย ปริปุณฺณ\-ทฺวาทสวสฺสาย คิหิคตาย เทฺว วสฺสานิ ฉสุ ธมฺเมสุ สิกฺขาสมฺมุติํ เทติ, ยสฺสา อยฺยาย ขมติ อิตฺถนฺนามาย ปริปุณฺณ\-ทฺวาทสวสฺสาย คิหิคตาย เทฺว วสฺสานิ ฉสุ ธมฺเมสุ สิกฺขาสมฺมุติยา ทานํ, สา ตุณฺหสฺส, ยสฺสา นกฺขมติ, สา ภาเสยฺย.}

\prose{ทินฺนา สํเฆน อิตฺถนฺนามาย ปริปุณฺณ\-ทฺวาทสวสฺสาย คิหิคตาย เทฺว วสฺสานิ ฉสุ ธมฺเมสุ สิกฺขาสมฺมุติ, ขมติ สํฆสฺส, ตสฺมา ตุณฺหี, เอวเมตํ ธารยามี”ติ.}

\prose{สา ปริปุณฺณทฺวาทสวสฺสา คิหิคตา เอวํ วเทหีติ วตฺตพฺพา “ปาณาติปาตา เวรมณิํ เทฺว วสฺสานิ อวีติกฺกมฺม สมาทานํ สมาทิยามิ ฯเปฯ วิกาลโภชนา เวรมณิํ เทฺว วสฺสานิ อวีติกฺกมฺม สมาทานํ สมาทิยามี”ติ.}

\prose{อถ โข ภควา ตา ภิกฺขุนิโย อเนกปริยาเยน วิครหิตฺวา ทุพฺภรตาย ฯเปฯ.\csromanspacer{} เอวญฺจ ปน ภิกฺขเว ภิกฺขุนิโย อิมํ สิกฺขาปทํ อุทฺทิสนฺตุ\csromandash{}}

\proseitem{1097}{\textbf{“ยา ปน ภิกฺขุนี ปริปุณฺณทฺวาทสวสฺสํ คิหิคตํ เทฺว วสฺสานิ ฉสุ ธมฺเมสุ อสิกฺขิตสิกฺขํ วุฏฺฐาเปยฺย, ปาจิตฺติยนฺ”}ติ.}

\proseitem{1098}{\textbf{ยา ปนา}ติ ยา ยาทิสา ฯเปฯ.\csromanspacer{} \textbf{ภิกฺขุนี}ติ ฯเปฯ อยํ อิมสฺมิํ อตฺเถ อธิปฺเปตา ภิกฺขุนีติ.}

\prose{\textbf{ปริปุณฺณทฺวาทสวสฺสา} นาม ปตฺตทฺวาทสวสฺสา.\csromanspacer{} \textbf{คิหิคตา} นาม ปุริสนฺตรคตา วุจฺจติ.\csromanspacer{} \textbf{เทฺว วสฺสานี}ติ เทฺว สํวจฺฉรานิ.\csromanspacer{} \textbf{อสิกฺขิตสิกฺขา} นาม สิกฺขา วา น ทินฺนา โหติ, ทินฺนา วา สิกฺขา กุปิตา.\csromanspacer{} \textbf{วุฏฺฐาเปยฺยา}ติ อุปสมฺปาเทยฺย.}

\prose{\csromanfolio{427}“วุฏฺฐาเปสฺสามี”ติ คณํ วา อาจรินิํ วา ปตฺตํ วา จีวรํ วา ปริเยสติ, สีมํ วา สมฺมนฺนติ, อาปตฺติ ทุกฺกฏสฺส.\csromanspacer{} ญตฺติยา ทุกฺกฏํ.\csromanspacer{} ทฺวีหิ กมฺมวาจาหิ ทุกฺกฏา.\csromanspacer{} กมฺมวาจาปริโยสาเน อุปชฺฌายาย อาปตฺติ ปาจิตฺติยสฺส.\csromanspacer{} คณสฺส จ อาจรินิยา จ อาปตฺติ ทุกฺกฏสฺส.}

\proseitem{1099}{ธมฺมกมฺเม ธมฺมกมฺมสญฺญา วุฏฺฐาเปติ, อาปตฺติ ปาจิตฺติยสฺส.\csromanspacer{} ธมฺมกมฺเม เวมติกา วุฏฺฐาเปติ, อาปตฺติ ปาจิตฺติยสฺส.\csromanspacer{} ธมฺมกมฺเม อธมฺมกมฺมสญฺญา วุฏฺฐาเปติ, อาปตฺติ ปาจิตฺติยสฺส.}

\prose{อธมฺมกมฺเม ธมฺมกมฺมสญฺญา วุฏฺฐาเปติ, อาปตฺติ ทุกฺกฏสฺส.\csromanspacer{} อธมฺมกมฺเม เวมติกา วุฏฺฐาเปติ, อาปตฺติ ทุกฺกฏสฺส.\csromanspacer{} อธมฺมกมฺเม อธมฺมกมฺมสญฺญา วุฏฺฐาเปติ, อาปตฺติ ทุกฺกฏสฺส.}

\proseitemwithcloserruled{1100}{อนาปตฺติ ปริปุณฺณทฺวาทสวสฺสํ คิหิคตํ เทฺว วสฺสานิ ฉสุ ธมฺเมสุ สิกฺขิตสิกฺขํ วุฏฺฐาเปติ, อุมฺมตฺติกาย อาทิกมฺมิกายาติ.}{ฉฏฺฐสิกฺขาปทํ นิฏฺฐิตํ.}
\csromanmatikaanchor{mtk.222}
\csromantocmarkref{subsection}{7. สตฺตมสิกฺขาปท}{mtk.222}
\bookmark[dest={mtk.222},level=2]{7. สตฺตมสิกฺขาปท}
\csromanheader{2}{7. \textbf{คพฺภินีวคฺค 7. สตฺตมสิกฺขาปท}}
\proseitem{1101}{เตน สมเยน พุทฺโธ ภควา สาวตฺถิยํ วิหรติ เชตวเน อนาถปิณฺฑิกสฺส อาราเม.\csromanspacer{} เตน โข ปน สมเยน ภิกฺขุนิโย ปริปุณฺณทฺวาทสวสฺสํ คิหิคตํ เทฺว วสฺสานิ ฉสุ ธมฺเมสุ สิกฺขิตสิกฺขํ สํเฆน อสมฺมตํ วุฏฺฐาเปนฺติ.\csromanspacer{} ภิกฺขุนิโย เอวมาหํสุ “เอถ สิกฺขมานา อิมํ ชานาถ อิมํ เทถ อิมํ อาหรถ อิมินา อตฺโถ อิมํ กปฺปิยํ กโรถา”ติ.\csromanspacer{} ตา เอวมาหํสุ “น มยํ อยฺเย สิกฺขมานา, ภิกฺขุนิโย มยนฺ”ติ.\csromanspacer{} ยา ตา ภิกฺขุนิโย อปฺปิจฺฉา ฯเปฯ ตา อุชฺฌายนฺติ ขิยฺยนฺติ วิปาเจนฺติ “กถํ หิ นาม ภิกฺขุนิโย ปริปุณฺณทฺวาทสวสฺสํ คิหิคตํ เทฺว วสฺสานิ ฉสุ ธมฺเมสุ สิกฺขิตสิกฺขํ สํเฆน อสมฺมตํ วุฏฺฐาเปสฺสนฺตี”ติ ฯเปฯ.\csromanspacer{} สจฺจํ กิร ภิกฺขเว ภิกฺขุนิโย ปริปุณฺณทฺวาทสวสฺสํ คิหิคตํ เทฺว วสฺสานิ ฉสุ ธมฺเมสุ สิกฺขิตสิกฺขํ สํเฆน อสมฺมตํ วุฏฺฐาเปนฺตีติ.\csromanspacer{} สจฺจํ ภควาติ.\csromanspacer{} วิครหิ พุทฺโธ ภควา ฯเปฯ.\csromanspacer{} กถํ หิ นาม ภิกฺขเว ภิกฺขุนิโย ปริปุณฺณทฺวาทสวสฺสํ คิหิคตํ เทฺว วสฺสานิ ฉสุ ธมฺเมสุ สิกฺขิตสิกฺขํ \csromanfolio{428}สํเฆน อสมฺมตํ วุฏฺฐาเปสฺสนฺติ.\csromanspacer{} เนตํ ภิกฺขเว อปฺปสนฺนานํ วา ปสาทาย ฯเปฯ วิครหิตฺวา ฯเปฯ ธมฺมิํ กถํ กตฺวา ภิกฺขู อามนฺเตสิ อนุชานามิ ภิกฺขเว ปริปุณฺณ\-ทฺวาทสวสฺสาย คิหิคตาย เทฺว วสฺสานิ ฉสุ ธมฺเมสุ สิกฺขิตสิกฺขาย วุฏฺฐานสมฺมุติํ ทาตุํ.\csromanspacer{} เอวญฺจ ปน ภิกฺขเว ทาตพฺพา, ตาย ปริปุณฺณ\-ทฺวาทสวสฺสาย คิหิคตาย เทฺว วสฺสฺสานิ ฉสุ ธมฺเมสุ สิกฺขิตสิกฺขาย สํฆํ อุปสงฺกมิตฺวา เอกํสํ อุตฺตราสงฺคํ กริตฺวา ภิกฺขุนีนํ ปาเท วนฺทิตฺวา อุกฺกุฏิกํ นิสีทิตฺวา อญฺชลิํ ปคฺคเหตฺวา เอวมสฺส วจนีโย “อหํ อยฺเย อิตฺถนฺนามา อิตฺถนฺนามาย อยฺยาย ปริปุณฺณทฺวาทสวสฺสา คิหิคตา เทฺว วสฺสานิ ฉสุ ธมฺเมสุ สิกฺขิตสิกฺขา สํฆํ วุฏฺฐานสมฺมุติํ ยาจามีติ.\csromanspacer{} ทุติยมฺปิ ยาจิตพฺพา.\csromanspacer{} ตติยมฺปิ ยาจิตพฺพา.\csromanspacer{} พฺยตฺตาย ภิกฺขุนิยา ปฏิพลาย สํโฆ ญาเปตพฺโพ\csromandash{}}

\proseitem{1102}{“สุณาตุ เม อยฺเย สํโฆ, อยํ อิตฺถนฺนามา อิตฺถนฺนามาย อยฺยาย ปริปุณฺณทฺวาทสวสฺสา คิหิคตา เทฺว วสฺสานิ ฉสุ ธมฺเมสุ สิกฺขิตสิกฺขา สํฆํ วุฏฺฐานสมฺมุติํ ยาจติ, ยทิ สํฆสฺส ปตฺตกลฺลํ, สํโฆ อิตฺถนฺนามาย ปริปุณฺณ\-ทฺวาทสวสฺสาย คิหิคตาย เทฺว วสฺสานิ ฉสุ ธมฺเมสุ สิกฺขิตสิกฺขาย วฺòฏฺถานสมฺมุติํ ทเทยฺย, เอสา ญตฺติ.}

\prose{สุณาตุ เม อยฺเย สํโฆ, อยํ อิตฺถนฺนามา อิตฺถนฺนามาย อยฺยาย ปริปุณฺณทฺวาทสวสฺสา คิหิคตา เทฺว วสฺสานิ ฉสุ ธมฺเมสุ สิกฺขิตสิกฺขา สํฆํ วุฏฺฐานสมฺมุติํ ยาจติ, สํโฆ อิตฺถนฺนามาย ปริปุณฺณ\-ทฺวาทสวสฺสาย คิหิคตาย เทฺว วสฺสานิ ฉสุ ธมฺเมสุ สิกฺขิตสิกฺขาย วุฏฺฐานสมฺมุติํ เทติ, ยสฺสา อยฺยาย ขมติ อิตฺถนฺนามาย ปริปุณฺณ\-ทฺวาทสวสฺสาย คิหิคตาย เทฺว วสฺสานิ ฉสุ ธมฺเมสุ สิกฺขิตสิกฺขาย วุฏฺฐานสมฺมุติยา ทานํ, สา ตุณฺหสฺส, ยสฺสา นกฺขมติ, สา ภาเสยฺย.}

\prose{ทินฺนา สํเฆน อิตฺถนฺนามาย ปริปุณฺณ\-ทฺวาทสวสฺสาย คิหิคตาย เทฺว วสฺสานิ ฉสุ ธมฺเมสุ สิกฺขิตสิกฺขาย วุฏฺฐานสมฺมุติ, ขมติ สํฆสฺส, ตสฺมา ตุณฺหี, เอวเมตํ ธารยามี”ติ.}

\prose{อถ โข ภควา ตา ภิกฺขุนิโย อเนกปริยาเยน วิครหิตฺวา ทุพฺภรตาย ฯเปฯ.\csromanspacer{} เอวญฺจ ปน ภิกฺขเว ภิกฺขุนิโย อิมํ สิกฺขาปทํ อุทฺทิสนฺตุ\csromandash{}}

\proseitem{1103}{\csromanfolio{429}\textbf{“ยา ปน ภิกฺขุนี ปริปุณฺณทฺวาทสวสฺสํ คิหิคตํ เทฺว วสฺสานิ ฉสุ ธมฺเมสุ สิกฺขิตสิกฺขํ สํเฆน อสมฺมตํ วุฏฺฐาเปยฺย, ปาจิตฺติยนฺ”}ติ.}

\proseitem{1104}{\textbf{ยา ปนา}ติ ยา ยาทิสา ฯเปฯ.\csromanspacer{} \textbf{ภิกฺขุนี}ติ ฯเปฯ อยํ อิมสฺมิํ อตฺเถ อธิปฺเปตา ภิกฺขุนีติ.}

\prose{\textbf{ปริปุณฺณทฺวาทสวสฺสา} นาม ปตฺตทฺวาทสวสฺสา.\csromanspacer{} \textbf{คิหิคตา} นาม ปุริสนฺตรคตา วุจฺจติ.\csromanspacer{} \textbf{เทฺว วสฺสานี}ติ เทฺว สํวจฺฉรานิ.\csromanspacer{} \textbf{สิกฺขิตสิกฺขา} นาม ฉสุ ธมฺเมสุ สิกฺขิตสิกฺขา.\csromanspacer{} \textbf{อสมฺมตา} นาม ญตฺติทุติเยน กมฺเมน วุฏฺฐานสมฺมุติ น ทินฺนา โหติ.\csromanspacer{} \textbf{วุฏฺฐาเปยฺยา}ติ อุปสมฺปาเทยฺย.}

\prose{“วุฏฺฐาเปสฺสามี”ติ คณํ วา อาจรินิํ วา ปตฺตํ วา จีวรํ วา ปริเยสติ, สีมํ วา สมฺปนฺนติ, อาปตฺติ ทุกฺกฏสฺส.\csromanspacer{} ญตฺติยา ทุกฺกฏํ.\csromanspacer{} ทฺวีหิ กมฺมวาจาหิ ทุกฺกฏา.\csromanspacer{} กมฺมวาจาปริโยสาเน อุปชฺฌายาย อาปตฺติ ปาจิตฺติยสฺส.\csromanspacer{} คณสฺส จ อาจรินิยา จ อาปตฺติ ทุกฺกฏสฺส.}

\proseitem{1105}{ธมฺมกมฺเม ธมฺมกมฺมสญฺญา วุฏฺฐาเปติ, อาปตฺติ ปาจิตฺติยสฺส.\csromanspacer{} ธมฺมกมฺเม เวมติกา วุฏฺฐาเปติ, อาปตฺติ ปาจิตฺติยสฺส.\csromanspacer{} ธมฺมกมฺเม อธมฺมกมฺมสญฺญา วุฏฺฐาเปติ, อาปตฺติ ปาจิตฺติยสฺส.}

\prose{อธมฺมกมฺเม ธมฺมกมฺมสญฺญา, อาปตฺติ ทุกฺกฏสฺส.\csromanspacer{} อธมฺมกมฺเม เวมติกา, อาปตฺติ ทุกฺกฏสฺส.\csromanspacer{} อธมฺมกมฺเม อธมฺมกมฺมสญฺญา, อาปตฺติ ทุกฺกฏสฺส.}

\proseitemwithcloserruled{1106}{อนาปตฺติ ปริปุณฺณทฺวาทสวสฺสํ คิหิคตํ เทฺว วสฺสานิ ฉสุ ธมฺเมสุ สิกฺขิตสิกฺขํ สํเฆน สมฺมตํ วุฏฺฐาเปติ, อุมฺมตฺติกาย อาทิกมฺมิกายาติ.}{สตฺตมสิกฺขาปทํ นิฏฺฐิตํ.}
\csromanmatikaanchor{mtk.223}
\csromantocmarkref{subsection}{8. อฏฺฐมสิกฺขาปท}{mtk.223}
\bookmark[dest={mtk.223},level=2]{8. อฏฺฐมสิกฺขาปท}
\csromanheader{2}{7. \textbf{คพฺภินีวคฺค 8. อฏฺฐมสิกฺขาปท}}
\proseitem{1107}{เตน สมเยน พุทฺโธ ภควา สาวตฺถิยํ วิหรติ เชตวเน อนาถปิณฺฑิกสฺส อาราเม.\csromanspacer{} เตน โข ปน สมเยน ถุลฺลนนฺทา ภิกฺขุนี สหชีวินิํ วุฏฺฐาเปตฺวา เทฺว วสฺสานิเนว เนว อนุคฺคณฺหาติ น อนุคฺคณฺหาเปติ, \csromanfolio{430}ตา พาลา โหนฺติ อพฺยตฺตา, น ชานนฺติ กปฺปิยํ วา อกปฺปิยํ วา.\csromanspacer{} ยา ตา ภิกฺขุนิโย อปฺปิจฺฉา ฯเปฯ ตา อุชฺฌายนฺติ ขิยฺยนฺติ วิปาเจนฺติ “กถํ หิ นาม อยฺยา ถุลฺลนนฺทา สหชีวินิํ วุฏฺฐาเปตฺวา เทฺว วสฺสานิ เนว อนุคฺคณฺหิสฺสติ น อนุคฺคณฺหาเปสฺสตี”ติ ฯเปฯ.\csromanspacer{} สจฺจํ กิร ภิกฺขเว ถุลฺลนนฺทา ภิกฺขุนี สหชีวินิํ วุฏฺฐาเปตฺวา เทฺว วสฺสานิ เนว อนุคฺคณฺหาติ น อนุคฺคณฺหาเปตีติ.\csromanspacer{} สจฺจํ ภควาติ.\csromanspacer{} วิครหิ พุทฺโธ ภควา ฯเปฯ.\csromanspacer{} กถํ หิ นาม ภิกฺขเว ถุลฺลนนฺทา ภิกฺขุนี สหชีวินิํ วุฏฺฐาเปตฺวา เทฺว วสฺสานิ เนว อนุคฺคณฺหิสฺสติ น อนุคฺคณฺหาเปสฺสติ.\csromanspacer{} เนตํ ภิกฺขเว อปฺปสนฺนานํ วา ปสาทาย ฯเปฯ.\csromanspacer{} เอวญฺจ ปน ภิกฺขเว ภิกฺขุนิโย อิมํ สิกฺขาปทํ อุทฺทิสนฺตุ\csromandash{}}

\proseitem{1108}{\textbf{“ยา ปน ภิกฺขุนี สหชีวินิํ วุฏฺฐาเปตฺวา เทฺว วสฺสานิ เนว อนุคฺคเณฺหยฺย น อนุคฺคณฺหาเปยฺย, ปาจิตฺติยนฺ”}ติ.}

\proseitem{1109}{\textbf{ยา ปนา}ติ ยา ยาทิสา ฯเปฯ.\csromanspacer{} \textbf{ภิกฺขุนี}ติ ฯเปฯ อยํ อิมสฺมิํ อตฺเถ อธิปฺเปตา ภิกฺขุนีติ.}

\prose{\textbf{สหชีวินี} นาม สทฺธิวิหารินี วุจฺจติ.\csromanspacer{} \textbf{วุฏฺฐาเปตฺวา}ติ อุปสมฺปาเทตฺวา.\csromanspacer{} \textbf{เทฺว วสฺสานี}ติ เทฺว สํวจฺฉรานิ.}

\prose{\textbf{เนว อนุคฺคเณฺหยฺยา}ติ น สยํ อนุคฺคเณฺหยฺย อุทฺเทเสน ปริปุจฺฉาย โอวาเทน อนุสาสนิยา.}

\prose{\textbf{น อนุคฺคณฺหาเปยฺยา}ติ น อญฺญํ อาณาเปยฺย. “เทฺว วสฺสานิ เนว อนุคฺคณฺหิสฺสามิ น อนุคฺคณฺหาเปสฺสามี”ติ ธุรํ นิกฺขิตฺตมตฺเต อาปตฺติ}

\proseitemwithcloserruled{1110}{อนาปตฺติ สติ อนฺตราเย, ปริเยสิตฺวา น ลภติ, คิลานาย อาปทาสุ อุมฺมตฺติกาย อาทิกมฺมิกายาติ.}{อฏฺฐมสิกฺขาปทํ นิฏฺฐิตํ.}
\csromanmatikaanchor{mtk.224}
\csromantocmarkref{subsection}{9. นวมสิกฺขาปท}{mtk.224}
\bookmark[dest={mtk.224},level=2]{9. นวมสิกฺขาปท}
\csromanheader{2}{7. \textbf{คพฺภินีวคฺค 9. นวมสิกฺขาปท}}
\proseitem{1111}{เตน สมเยน พุทฺโธ ภควา สาวตฺถิยํ วิหรติ เชตวเน อนาถปิณฺฑิกสฺส อาราเม.\csromanspacer{} เตน โข ปน สมเยน ภิกฺขุนิโย \csromanfolio{431}วุฏฺฐาปิตํ ปวตฺตินิํ เทฺว วสฺสานิ นานุพทฺธนฺติ, ตา พาลา โหนฺติ อพฺยตฺตา, น ชานนฺติ กปฺปิยํ วา อกปฺปิยํ วา.\csromanspacer{} ยา ตา ภิกฺขุนิโย อปฺปิจฺฉา ฯเปฯ ตา อุชฺฌายนฺติ ขิยฺยนฺติ วิปาเจนฺติ “กถํ หิ นาม ภิกฺขุนิโย วุฏฺฐาปิตํ ปวตฺตินิํ เทฺว วสฺสานิ นานุพนฺธิสฺสนฺตี”ติ ฯเปฯ.\csromanspacer{} สจฺจํ กิร ภิกฺขเว ภิกฺขุนิโย วุฏฺฐาปิตํ ปวตฺตินิํ เทฺว วสฺสานิ นานุพนฺธนฺตีติ.\csromanspacer{} สจฺจํ ภควาติ.\csromanspacer{} วิครหิ พุทฺโธ ภควา ฯเปฯ.\csromanspacer{} กถํ หิ นาม ภิกฺขเว ภิกฺขุนิโย วุฏฺฐาปิตํ ปวตฺตินิํ เทฺว วสฺสานิ นานุพนฺธิสฺสนฺติ.\csromanspacer{} เนตํ ภิกฺขเว อปฺปสนฺนานํ วา ปสาทาย ฯเปฯ.\csromanspacer{} เอวญฺจ ปน ภิกฺขเว \textbf{ภิกฺขุนิโย อิมํ สิกฺขาปทํ อุทฺทิสนฺตุ\csromandash{}}}

\proseitem{1112}{\textbf{“ยา ปน ภิกฺขุนี วุฏฺฐาปิตํ ปวตฺตินิํ เทฺว วสฺสานิ นานุพนฺเธยฺย, ปาจิตฺติยนฺ”}ติ.}

\proseitem{1113}{\textbf{ยา ปนา}ติ ยา ยาทิสา ฯเปฯ.\csromanspacer{} \textbf{ภิกฺขุนี}ติ ฯเปฯ อยํ อิมสฺมิํ อตฺเถ อธิปฺเปตา ภิกฺขุนีติ.}

\prose{\textbf{วุฏฺฐาปิตนฺ}ติ อุปสมฺปาทิตํ.\csromanspacer{} \textbf{ปวตฺตินี} นาม อุปชฺฌายา วุจฺจติ.\csromanspacer{} \textbf{เทฺว วสฺสานี}ติ เทฺว สํวจฺฉรานิ.\csromanspacer{} \textbf{นานุพนฺเธยฺยา}ติ น สยํ อุปฏฺฐเหยฺย. “เทฺว วสฺสานิ นานุพนฺธิสฺสามี”ติ ธุรํ นิกฺขิตฺตมตฺเต อาปตฺติ ปาจิตฺติยสฺส.}

\proseitemwithcloserruled{1114}{อนาปตฺติ อุปชฺฌายา พาลา วา โหติ อลชฺชินี วา, คิลานาย อาปทาสุ อุมฺมตฺติกาย อาทิกมฺมิกายาติ.}{นวมสิกฺขาปทํ นิฏฺฐิตํ.}
\csromanmatikaanchor{mtk.225}
\csromantocmarkref{subsection}{10. ทสมสิกฺขาปท}{mtk.225}
\bookmark[dest={mtk.225},level=2]{10. ทสมสิกฺขาปท}
\csromanheader{2}{7. \textbf{คพฺภินีวคฺค} 10. ทสมสิกฺขาปท}
\proseitem{1115}{เตน สมเยน พุทฺโธ ภควา สาวตฺถิยํ วิหรติ เชตวเน อนาถปิณฺฑิกสฺส อาราเม.\csromanspacer{} เตน โข ปน สมเยน ถุลฺลนนฺทา ภิกฺขุนี สหชีวินิํ วุฏฺฐาเปตฺวา เนว วูปกาเสติ น วูปกาสาเปติ, สามิโก อคฺคเหสิ.\csromanspacer{} ยา ตา ภิกฺขุนิโย อปฺปิจฺฉา ฯเปฯ ตา อุชฺฌายนฺติ ขิยฺยนฺติ วิปาเจนฺติ “กถํ หิ นาม อยฺยา ถุลฺลนนฺทา สหชีวินิํ วุฏฺฐาเปตฺวา เนว วูปกาเสสฺสติ น วูปกาสาเปสฺสติ, สามิโก \csromanfolio{432}อคฺคเหสิ, สจายํ ภิกฺขุนี ปกฺกนฺตา อสฺส น จ สามิโก คเณฺหยฺยา”ติ ฯเปฯ.\csromanspacer{} สจฺจํ กิร ภิกฺขเว ถุลฺลนนฺทา ภิกฺขุนี สหชีวินิํ วุฏฺฐาเปตฺวา เนว วูปกาเสติ\footnote{...สิ (ก)} น วูปกาสาเปติ1, สามิโก อคฺคเหสีติ.\csromanspacer{} สจฺจํ ภควาติ.\csromanspacer{} วิครหิ พุทฺโธ ภควา ฯเปฯ.\csromanspacer{} กถํ หิ นาม ภิกฺขเว ถุลฺลนนฺทา ภิกฺขุนี สหชีวินิํ วุฏฺฐาเปตฺวา เนว วูปกาเสสฺสติ น วูปกาสาเปสฺสติ, สามิโก อคฺคเหสิ.\csromanspacer{} เนตํ ภิกฺขเว อปฺปสนฺนานํ วา ปสาทาย ฯเปฯ.\csromanspacer{} เอวญฺจ ปน ภิกฺขเว ภิกฺขุนิโย อิมํ สิกฺขาปทํ อุทฺทิสนฺตุ\csromandash{}}

\proseitem{1116}{\textbf{“ยา ปน ภิกฺขุนี สหชีวินิํ วุฏฺฐาเปตฺวา เนว วูปกาเสยฺย นวูปกาสาเปยฺย อนฺตมโส ฉปฺปญฺจโยชนานิปิ, ปาจิตฺติยนฺ”}ติ.}

\proseitem{1117}{\textbf{ยา ปนา}ติ ยา ยาทิสา ฯเปฯ.\csromanspacer{} \textbf{ภิกฺขุนี}ติ ฯเปฯ อยํ อิมสฺมิํ อตฺเถ อธิปฺเปตา ภิกฺขุนีติ.}

\prose{\textbf{สหชีวินี} นาม สทฺธิวิหารินี วุจฺจติ.\csromanspacer{} \textbf{วุฏฺฐาเปตฺวา}ติ อุปสมฺปาเทตฺวา.\csromanspacer{} \textbf{เนว วูปกาเสยฺยา}ติ น สยํ วูปกาเสยฺย.\csromanspacer{} \textbf{น วูปกาสาเปยฺยา}ติ น อญฺญํ อาณาเปยฺย. “เนว วูปกาเสสฺสามิ น วูปกาสาเปสฺสามิ อนฺตมโส ฉปฺปญฺจโยชนานิปี”ติ ธุรํ นิกฺขิตฺตมตฺเต อาปตฺติ ปาจิตฺติยสฺส.}

\proseitemwithcloser{1118}{อนาปตฺติ สติ อนฺตราเย, ปริเยสิตฺวา ทุติยิกํ ภิกฺขุนิํ น ลภติ, คิลานาย อาปทาสุ อุมฺมตฺติกาย อาทิกมฺมิกายาติ.}{ทสมสิกฺขาปทํ นิฏฺฐิตํ.}
\csromanheader{2}{คพฺภินิวคฺโค สตฺตโม.}
\csromansectionrule
\csromanmatikaanchor{mtk.226}
\csromanmatikaanchor{mtk.227}
\csromantocmarknopage{section}{8. กุมารีภูตวคฺค}{mtk.226}
\bookmark[dest={mtk.226},level=1]{8. กุมารีภูตวคฺค}
\csromantocmarkref{subsection}{1. ปฐมสิกฺขาปท}{mtk.227}
\bookmark[dest={mtk.227},level=2]{1. ปฐมสิกฺขาปท}
\setcsromanheadh{8. กุมารีภูตวคฺค}
\csromanheader{2}{8. \textbf{กุมารีภูตวคฺค 1. ปฐมสิกฺขาปท}}
\proseitem{1119}{เตน สมเยน พุทฺโธ ภควา สาวตฺถิยํ วิหรติ เชตวเน อนาถปิณฺฑิกสฺส อาราเม.\csromanspacer{} เตน โข ปน สมเยน ภิกฺขุนิโย อูนวีสติวสฺสํ กุมาริภูตํ วุฏฺฐาเปนฺติ, ตา อกฺขมา โหนฺติ สีตสฺส \csromanfolio{433}อุณฺหสฺส ชิฆจฺฉาย ปิปาสาย ฑํส\-มกส\-วาตา\-ตป\-สรีสป\-สมฺผสฺสานํ ทุรุตฺตานํ ทุราคตานํ วจนปถานํ, อุปฺปนฺนานํ สารีริกานํ เวทนานํ ทุกฺขานํ ติพฺพานํ ขรานํ กฏุกานํ อสาตานํ อมนาปานํ ปาณหรานํ อนธิวาสกชาติกา โหนฺติ.\csromanspacer{} ยา ตา ภิกฺขุนิโย อปฺปิจฺฉา ฯเปฯ ตา อุชฺฌายนฺติ ขิยฺยนฺติ วิปาเจนฺติ “กถํ หิ นาม ภิกฺขุนิโย อูนวีสติวสฺสํ กุมาริภูตํ วุฏฺฐาเปสฺสนฺตี”ติ ฯเปฯ.\csromanspacer{} สจฺจํ กิร ภิกฺขเว ภิกฺขุนิโย อูนวีสติวสฺสํ กุมาริภูตํ วุฏฺฐาเปนฺตีติ.\csromanspacer{} สจฺจํ ภควาติ.\csromanspacer{} วิครหิ พุทฺโธ ภควา ฯเปฯ.\csromanspacer{} กถํ หิ นาม ภิกฺขเว ภิกฺขุนิโย อูนวีสติวสฺสํ กุมาริภูตํ วุฏฺฐาเปสฺสนฺติ, อูนวีสติวสฺสา ภิกฺขเว กุมาริภูตา อกฺขมา โหติ สีตสฺส อุณฺหสฺส ฯเปฯ ปาณหรานํ อนธิวาสกชาติกา โหติ, วีสติวสฺสาว โข ภิกฺขเว กุมาริภูตา ขมา โหติ สีตสฺส อุณฺหสฺส ฯเปฯ ปาณหรานํ อธิวาสกชาติกา โหติ.\csromanspacer{} เนตํ ภิกฺขเว อปฺปสนฺนานํ วา ปสาทาย ฯเปฯ.\csromanspacer{} เอวญฺจ ปน ภิกฺขเว ภิกฺขุนิโย อิมํ สิกฺขาปทํ อุทฺทิสนฺตุ\csromandash{}}

\proseitem{1120}{\textbf{“ยา ปน ภิกฺขุนี อูนวีสติวสฺสํ กุมาริภูตํ วุฏฺฐาเปยฺย, ปาจิตฺติยนฺ”}ติ.}

\proseitem{1121}{\textbf{ยา ปนา}ติ ยา ยาทิสา ฯเปฯ.\csromanspacer{} \textbf{ภิกฺขุนี}ติ ฯเปฯ อยํ อิมสฺมิํ อตฺเถ อธิปฺเปตา ภิกฺขุนีติ.}

\prose{อู\textbf{นวีสติวสฺสา} นาม อปฺปตฺตวีสติวสฺสา.\csromanspacer{} \textbf{กุมาริภูตา} นาม สามเณรี วุจฺจติ.\csromanspacer{} \textbf{วุฏฺฏฺถาเปยฺยา}ติ อุปสมฺปาเทยฺย.}

\prose{“วุฏฺฐาเปสฺสามี”ติ คณํ วา อาจรินิํ วา ปตฺตํ วา จีวรํ วา ปริเยสติ, สีมํ วา สมฺมนฺนติ, อาปตฺติ ทุกฺกฏสฺส.\csromanspacer{} ญตฺติยา ทุกฺกฏํ.\csromanspacer{} ทฺวีหิ กมฺมวาจาหิ ทุกฺกฏา.\csromanspacer{} กมฺมวาจาปริโยสาเน อุปชฺฌายาย อาปตฺติ ปาจิตฺติยสฺส.\csromanspacer{} คณสฺส จ อาจรินิยา จ อาปตฺติ ทุกฺกฏสฺส.}

\proseitem{1122}{อูนวีสติวสฺสาย อูนวีสติวสฺสสญฺญา วุฏฺฐาเปติ, อาปตฺติ ปาจิตฺติยสฺส.\csromanspacer{} อูนวีสติวสฺสาย เวมติกา วุฏฺฐาเปติ, อาปตฺติ ทุกฺกฏสฺส.\csromanspacer{} อูนวีสติวสฺสาย ปริปุณฺณสญฺญา วุฏฺฐาเปติ, อนาปตฺติ.\csromanspacer{} ปริปุณฺณวีสติวสฺสาย อูนวีสติวสฺสสญฺญา, อาปตฺติ ทุกฺกฏสฺส.\csromanspacer{} ปริปุณฺณวีสติวสฺสาย \csromanfolio{434}เวมติกา, อาปตฺติ ทุกฺกฏสฺส.\csromanspacer{} \textbf{ปริปุณฺณวีสติวสฺสา}ย ปริปุณฺณสญฺญา, อนาปตฺติ.}

\proseitemwithcloserruled{1123}{อนาปตฺติ อูนวีสติวสฺสํ ปริปุณฺณสญฺญา วุฏฺฐาเปติ, ปริปุณฺณวีสติวสฺสํ ปริปุณฺณสญฺญา วุฏฺฐาเปติ, อุมฺมตฺติกาย อาทิกมฺมิกายาติ.}{ปฐมสิกฺขาปทํ นิฏฺฐิตํ.}
\csromanmatikaanchor{mtk.228}
\csromantocmarkref{subsection}{2. ทุติยสิกฺขาปท}{mtk.228}
\bookmark[dest={mtk.228},level=2]{2. ทุติยสิกฺขาปท}
\csromanheader{2}{8. \textbf{กุมารีภูตวคฺค 2. ทุติยสิกฺขาปท}}
\proseitem{1124}{เตน สมเยน พุทฺโธ ภควา สาวตฺถิยํ วิหรติ เชตวเน อนาถปิณฺฑิกสฺส อาราเม.\csromanspacer{} เตน โข ปน สมเยน ภิกฺขุนิโย ปริปุณฺณวีสติวสฺสํ กุมาริภูตํ เทฺว วสฺสานิ ฉสุ ธมฺเมสุ อสิกฺขิตสิกฺขํ วุฏฺฐาเปนฺติ.\csromanspacer{} ตา พาลา โหนฺติ อพฺยตฺตา, น ชานนฺติ กปฺปิยํ วา อกปฺปิยํ วา.\csromanspacer{} ยา ตา ภิกฺขุนิโย อปฺปิจฺฉา ฯเปฯ ตา อุชฺฌายนฺติ ขิยฺยนฺติ วิปาเจนฺติ “กถํ หิ นาม ภิกฺขุนิโย ปริปุณฺณวีสติวสฺสํ กุมาริภูตํ เทฺว วสฺสานิ ฉสุ ธมฺเมสุ อสิกฺขิตสิกฺขํ วุฏฺฐาเปสฺสนฺตี”ติ ฯเปฯ.\csromanspacer{} สจฺจํ กิร ภิกฺขเว ภิกฺขุนิโย ปริปุณฺณวีสติวสฺสํ กุมาริภูตํ เทฺว วสฺสานิ ฉสุ ธมฺเมสุ อสิกฺขิตสิกฺขํ วุฏฺฐาเปนฺตีติ.\csromanspacer{} สจฺจํ ภควาติ.\csromanspacer{} วิครหิ พุทฺโธ ภควา ฯเปฯ.\csromanspacer{} กถํ หิ นาม ภิกฺขเว ภิกฺขุนิโย ปริปุณฺณวีสติวสฺสํ กุมาริภูตํ เทฺว วสฺสานิ ฉสุ ธมฺเมสุ อสิกฺขิตสิกฺขํ วุฏฺฐาเปสฺสนฺติ.\csromanspacer{} เนตํ ภิกฺขเว อปฺปสนฺนานํ วา ปสาทาย ฯเปฯ วิครหิตฺวา ฯเปฯ ธมฺมิํ กถํ กตฺวา ภิกฺขู อามนฺเตสิ อนุชานามิ ภิกฺขเว อฏฺฐารสวสฺสาย กุมาริภูตาย เทฺว วสฺสานิ ฉสุ ธมฺเมสุ สิกฺขาสมฺมุติํ ทาตุํ.\csromanspacer{} เอวญฺจ ปน ภิกฺขเว ทาตพฺพา, ตาย อฏฺฐารสวสฺสาย กุมาริภูตาย สํฆํ อุปสงฺกมิตฺวา เอกํสํ อุตฺตราสงฺคํ กริตฺวา ภิกฺขุนีนํ ปาเท วนฺทิตฺวา อุกฺกุฏิกํ นิสีทิตฺวา อญฺชลิํ ปคฺคเหตฺวา เอวมสฺส วจนีโย “อหํ อยฺเย อิตฺถนฺนามา อิตฺถนฺนามาย อยฺยาย อฏฺฐารสวสฺสา กุมาริภูตา สํฆํ เทฺว วสฺสานิ ฉสุ ธมฺเมสุ สิกฺขาสมฺมุติํ ยาจามี”ติ.\csromanspacer{} ทุติยมฺปิ ยาจิตพฺพา.\csromanspacer{} ตติยมฺปิ ยาจิตพฺพา.\csromanspacer{} พฺยตฺตาย ภิกฺขุนิยา ปฏิพลาย สํโฆ ญาเปตพฺโพ\csromandash{}}

\proseitem{1125}{\csromanfolio{435}“สุณาตุ เม อยฺเย สํโฆ, อยํ อิตฺถนฺนามา อิตฺถนฺนามาย อยฺยาย อฏฺฐารสวสฺสา กุมาริภูตา สํฆํ เทฺว วสฺสานิ ฉสุ ธมฺเมสุ สิกฺขาสมฺมุติํ ยาจติ, ยทิ สํฆสฺส ปตฺตกลฺลํ, สํโฆ อิตฺถนฺนามาย อฏฺฐารสวสฺสาย กุมาริภูตาย เทฺว วสฺสานิ ฉสุ ธมฺเมสุ สิกฺขาสมฺมุติํ ทเทยฺย, เอสา ญตฺติ.}

\prose{สุณาตุ เม อยฺเย สํโฆ, อยํ อิตฺถนฺนามา อิตฺถนฺนามาย อยฺยาย อฏฺฐารสวสฺสา กุมาริภูตา สํฆํ เทฺว วสฺสานิ ฉสุ ธมฺเมสุ สิกฺขาสมฺมุติํ ยาจติ, สํโฆ อิตฺถนฺนามาย อฏฺฐารสวสฺสาย กุมาริภูตาย เทฺว วสฺสานิ ฉสุ ธมฺเมสุ สิกฺขาสมฺมุติํ เทติ, ยสฺสา อยฺยาย ขมติ อิตฺถนฺนามาย อฏฺฐารสวสฺสาย กุมาริภูตาย เทฺว วสฺสานิ ฉสุ ธมฺเมสุ สิกฺขาสมฺมุติยา ทานํ, สา ตุณฺหสฺส, ยสฺสา นกฺขมติ, สา ภาเสยฺย.}

\prose{ทินฺนา สํเฆน อิตฺถนฺนามาย อฏฺฐารสวสฺสาย กุมาริภูตาย เทฺว วสฺสานิ ฉสุ ธมฺเมสุ สิกฺขาสมฺมุติ, ขมติ สํฆสฺส, ตสฺมา ตุณฺหี, เอวเมตํ ธารยามี”ติ.}

\prose{สา อฏฺฐารสวสฺสา กุมาริภูตา เอวํ วเทหีติ วตฺตพฺพา “ปาณาติปาตา เวรมณิํ เทฺว วสฺสานิ อวีติกฺกมฺม สมาทานํ สมาทิยามิ ฯเปฯ วิกาลโภชนา เวรมณิํ เทฺว วสฺสานิ อวีติกฺกมฺม สมาทานํ สมาทิยามี”ติ.}

\prose{อถ โข ภควา ตา ภิกฺขุนิโย อเนกปริยาเยน วิครหิตฺวา ทุพฺภรตาย ฯเปฯ.\csromanspacer{} เอวญฺจ ปน ภิกฺขเว ภิกฺขุนิโย อิมํ สิกฺขาปทํ อุทฺทิสนฺตุ\csromandash{}}

\proseitem{1126}{\textbf{“ยา ปน ภิกฺขุนี ปริปุณฺณวีสติวสฺสํ กุมาริภูตํ เทฺว วสฺสานิ ฉสุ ธมฺเมสุ อสิกฺขิตสิกฺขํ วุฏฺฐาเปยฺย, ปาจิตฺติยนฺ”}ติ.}

\proseitem{1127}{\textbf{ยา ปนา}ติ ยา ยาทิสา ฯเปฯ.\csromanspacer{} \textbf{ภิกฺขุนี}ติ ฯเปฯ อยํ อิมสฺมิํ อตฺเถ อธิปฺเปตา ภิกฺขุนีติ.}

\prose{\textbf{ปริปุณฺณวีสติวสฺสา} นาม ปตฺตวีสติวสฺสา.\csromanspacer{} \textbf{กุมาริภูตา} นาม สามเณรี วุจฺจติ.\csromanspacer{} \textbf{เทฺว วสฺสานี}ติ เทฺว สํวจฺฉรานิ.\csromanspacer{} \textbf{อสิกฺขิตสิกฺขา} นาม สิกฺขา วา น ทินฺนา โหติ, ทินฺนา วา สิกฺขา กุปิตา.\csromanspacer{} \textbf{วุฏฺฐาเปยฺยา}ติ อุปสมฺปาเทยฺย.}

\prose{\csromanfolio{436}“วุฏฺฐาเปสฺสามี”ติ คณํ วา อาจรินิํ วา ปตฺตํ วา จีวรํ วา ปริเยสติ, สีมํ วา สมฺมนฺนติ, อาปตฺติ ทุกฺกฏสฺส.\csromanspacer{} ญตฺติยา ทุกฺกฏํ.\csromanspacer{} ทฺวีหิ กมฺมวาจาหิ ทุกฺกฏา.\csromanspacer{} กมฺมวาจาปริโยสาเน อุปชฺฌายาย อาปตฺติ ปาจิตฺติยสฺส.\csromanspacer{} คณสฺส จ อาจรินิยา จ อาปตฺติ ทุกฺกฏสฺส.}

\proseitem{1128}{ธมฺมกมฺเม ธมฺมกมฺมสญฺญา วุฏฺฐาเปติ, อาปตฺติ ปาจิตฺติยสฺส.\csromanspacer{} ธมฺมกมฺเม เวมติกา วุฏฺฐาเปติ, อาปตฺติ ปาจิตฺติยสฺส.\csromanspacer{} ธมฺมกมฺเม อธมฺมกมฺมสญฺญา วุฏฺฐาเปติ, อาปตฺติ ปาจิตฺติยสฺส.}

\prose{อธมฺมกมฺเม ธมฺมกมฺมสญฺญา, อาปตฺติ ทุกฺกฏสฺส.\csromanspacer{} อธมฺมกมฺเม เวมติกา, อาปตฺติ ทุกฺกฏสฺส.\csromanspacer{} อธมฺมกมฺเม อธมฺมกมฺมสญฺญา, อาปตฺติ ทุกฺกฏสฺส.}

\proseitemwithcloserruled{1129}{อนาปตฺติ ปริปุณฺณวีสติวสฺสํ กุมาริภูตํ เทฺว วสฺสานิ ฉสุ ธมฺเมสุ สิกฺขิตสิกฺขํ วุฏฺฐาเปติ, อุมฺมตฺติกาย อาทิกมฺมิกายาติ.}{ทุติยสิกฺขาปทํ นิฏฺฐิตํ.}
\csromanmatikaanchor{mtk.229}
\csromantocmarkref{subsection}{3. ตติยสิกฺขาปท}{mtk.229}
\bookmark[dest={mtk.229},level=2]{3. ตติยสิกฺขาปท}
\csromanheader{2}{8. \textbf{กุมารีภูตวคฺค 3. ตติยสิกฺขาปท}}
\proseitem{1130}{เตน สมเยน พุทฺโธ ภควา สาวตฺถิยํ วิหรติ เชตวเน อนาถปิณฺฑิกสฺส อาราเม.\csromanspacer{} เตน โข ปน สมเยน ภิกฺขุนิโย ปริปุณฺณวีสติวสฺสํ กุมาริภูตํ เทฺว วสฺสานิ ฉสุ ธมฺเมสุ สิกฺขิตสิกฺขํ สํเฆน อสมฺมตํ วุฏฺฐาเปนฺติ.\csromanspacer{} ภิกฺขุนิโย เอวมาหํสุ “เอถ สิกฺขมานา อิมํ ชานาถ อิมํ เทถ อิมํ อาหรถ อิมินา อตฺโถ อิมํ กปฺปิยํ กโรถา”ติ.\csromanspacer{} ตา เอวมาหํสุ “น มยํ อยฺเย สิกฺขมานา, ภิกฺขุนิโย มยนฺ”ติ.\csromanspacer{} ยา ตา ภิกฺขุนิโย อปฺปิจฺฉา ฯเปฯ ตา อุชฺฌายนฺติ ขิยฺยนฺติ วิปาเจนฺติ “กถํ หิ นาม ภิกฺขุนิโย ปริปุณฺณวีสติวสฺสํ กุมาริภูตํ เทฺว วสฺสานิ ฉสุ ธมฺเมสุ สิกฺขิตสิกฺขํ สํเฆน อสมฺมตํ วุฏฺฐาเปสฺสนฺตี”ติ ฯเปฯ.\csromanspacer{} สจฺจํ กิร ภิกฺขเว ภิกฺขุนิโย ปริปุณฺณวีสติวสฺสํ กุมาริภูตํ เทฺว วสฺสานิ ฉสุ ธมฺเมสุ สิกฺขิตสิกฺขํ สํเฆน อสมฺมตํ วุฏฺฐาเปนฺตีติ.\csromanspacer{} สจฺจํ ภควาติ.\csromanspacer{} วิครหิ พุทฺโธ ภควา ฯเปฯ.\csromanspacer{} กถํ หิ นาม ภิกฺขเว ภิกฺขุนิโย ปริปุณฺณวีสติวสฺสํ กุมาริภูตํ เทฺว วสฺสานิ ฉสุ ธมฺเมสุ สิกฺขิตสิกฺขํ สํเฆน อสมฺมตํ วุฏฺฐาเปสฺสนฺติ.\csromanspacer{} เนตํ \csromanfolio{437}\textbf{ภิกฺขเว อปฺปสนฺนานํ วา ปสาทาย ฯเปฯ วิครหิตฺวา ฯเปฯ ธมฺมิํ} \textbf{กถํ กตฺวา ภิกฺขู อามนฺเตสิ อนุชานามิ ภิกฺขเว} \textbf{ปริปุณฺณวีสติวสฺสาย กุมาริภูตาย เทฺว วสฺสานิ ฉสุ ธมฺเมสุ} \textbf{สิกฺขิตสิกฺขาย วุฏฺฐานสมฺมุติํ ทาตุํ.}\csromanspacer{}\textbf{ เอวญฺจ ปน ภิกฺขเว} ทาตพฺพา, ตาย ปริปุณฺณวีสติวสฺสาย กุมาริภูตาย เทฺว วสฺสานิ ฉสุ ธมฺเมสุ สิกฺขิตสิกฺขาย สํฆํ อุปสงฺกมิตฺวา เอกํสํ อุตฺตราสงฺคํ กริตฺวา ภิกฺขุนีนํ ปาเท วนฺทิตฺวา อุกฺกุฏิกํ นิสีทิตฺวา อญฺชลิํ ปคฺคเหตฺวา เอวมสฺส วจนีโย “อหํ อยฺเย อิตฺถนฺนามา อิตฺถนฺนามาย อยฺยาย ปริปุณฺณวีสติวสฺสา กุมาริภูตา เทฺว วสฺสานิ ฉสุ ธมฺเมสุ สิกฺขิตสิกฺขา สํฆํ วุฏฺฐานสมฺมุติํ ยาจามี”ติ.\csromanspacer{} ทุติยมฺปิ ยาจิตพฺพา.\csromanspacer{} ตติยมฺปิ ยาจิตพฺพา.\csromanspacer{} พฺยตฺตาย ภิกฺขุนิยา ปฏิพลาย สํโฆ ญาเปตพฺโพ\csromandash{}}

\proseitem{1131}{“สุณาตุ เม อยฺเย สํโฆ, อยํ อิตฺถนฺนามา อิตฺถนฺนามาย อยฺยาย ปริปุณฺณวีสติวสฺสา กุมาริภูตา เทฺว วสฺสานิ ฉสุ ธมฺเมสุ สิกฺขิตสิกฺขา สํฆํ วุฏฺฐานสมฺมุติํ ยาจติ, ยทิ สํฆสฺส ปตฺตกลฺลํ, สํโฆ อิตฺถนฺนามาย ปริปุณฺณวีสติวสฺสาย กุมาริภูตาย เทฺว วสฺสานิ ฉสุ ธมฺเมสุ สิกฺขิตสิกฺขาย วุฏฺฐานสมฺมุติํ ทเทยฺย, เอสา ญตฺติ.}

\prose{สุณาตุ เม อยฺเย สํโฆ, อยํ อิตฺถนฺนามา อิตฺถนฺนามาย อยฺยาย ปริปุณฺณวีสติวสฺสา กุมาริภูตา เทฺว วสฺสานิ ฉสุ ธมฺเมสุ สิกฺขิตสิกฺขา สํฆํ วุฏฺฐานสมฺมุติํ ยาจติ, สํโฆ อิตฺถนฺนามาย ปริปุณฺณวีสติวสฺสาย กุมาริภูตาย เทฺว วสฺสานิ ฉสุ ธมฺเมสุ สิกฺขิตสิกฺขาย วุฏฺฐานสมฺมุติํ เทติ, ยสฺสา อยฺยาย ขมติ อิตฺถนฺนามาย ปริปุณฺณวีสติวสฺสาย กุมาริภูตาย เทฺว วสฺสานิ ฉสุ ธมฺเมสุ สิกฺขิตสิกฺขาย วุฏฺฐานสมฺมุติยา ทานํ, สา ตุณฺหสฺส, ยสฺสา นกฺขมติ, สา ภาเสยฺย.}

\prose{ทินฺนา สํเฆน อิตฺถนฺนามาย ปริปุณฺณวีสติวสฺสาย กุมาริภูตาย เทฺว วสฺสานิ ฉสุ ธมฺเมสุ สิกฺขิตสิกฺขาย วุฏฺฐานสมฺมุติ, ขมติ สํฆสฺส, ตสฺมา ตุณฺหี, เอวเมตํ ธารยามี”ติ.}

\prose{อถ โข ภควา ตา ภิกฺขุนิโย อเนกปริยาเยน วิครหิตฺวา ทุพฺภรตาย ฯเปฯ.\csromanspacer{} เอวญฺจ ปน ภิกฺขเว ภิกฺขุนิโย อิมํ สิกฺขาปทํ อุทฺทิสนฺตุ\csromandash{}}

\proseitem{1132}{\csromanfolio{438}\textbf{“ยา ปน ภิกฺขุนี ปริปุณฺณวีสติวสฺสํ กุมาริภูตํ เทฺว วสฺสานิ ฉสุ ธมฺเมสุ สิกฺขิตสิกฺขํ สํเฆน อสมฺมตํ วุฏฺฐาเปยฺย, ปาจิตฺติยนฺ”}ติ.}

\proseitem{1133}{\textbf{ยา ปนา}ติ ยา ยาทิสา ฯเปฯ.\csromanspacer{} \textbf{ภิกฺขุนี}ติ ฯเปฯ อยํ อิมสฺมิํ อตฺเถ อธิปฺเปตา ภิกฺขุนีติ.}

\prose{\textbf{ปริปุณฺณวีสติวสฺสา} นาม ปตฺตวีสติวสฺสา.\csromanspacer{} \textbf{กุมาริภูตา} นาม สามเณรี วุจฺจติ.\csromanspacer{} \textbf{เทฺว วสฺสานี}ติ เทฺว สํวจฺฉรานิ.\csromanspacer{} \textbf{สิกฺขิตสิกฺขา} นาม ฉสุ ธมฺเมสุ สิกฺขิตสิกฺขา.\csromanspacer{} \textbf{อสมฺมตา} นาม ญตฺติทุติเยน กมฺเมน วุฏฺฐานสมฺมุติ น ทินฺนา โหติ.\csromanspacer{} \textbf{วุฏฺฐาเปยฺยา}ติ อุปสมฺปาเทยฺย.}

\prose{“วุฏฺฐาเปสฺสามี”ติ คณํ วา อาจรินิํ วา ปตฺตํ วา จีวรํ วา ปริเยสติ, สีมํ วา สมฺมนฺนติ, อาปตฺติ ทุกฺกฏสฺส.\csromanspacer{} ญตฺติยา ทุกฺกฏํ.\csromanspacer{} ทฺวีหิ กมฺมวาจาหิ ทุกฺกฏา.\csromanspacer{} กมฺมวาจาปริโยสาเน อุปชฺฌายาย อาปตฺติ ปาจิตฺติยสฺส.\csromanspacer{} คณสฺส จ อาจรินิยา จ อาปตฺติ ทุกฺกฏสฺส.}

\proseitem{1134}{ธมฺมกมฺเม ธมฺมกมฺมสญฺญา วุฏฺฐาเปติ, อาปตฺติ ปาจิตฺติยสฺส.\csromanspacer{} ธมฺมกมฺเม เวมติกา วุฏฺฐาเปติ, อาปตฺติ ปาจิตฺติยสฺส.\csromanspacer{} ธมฺมกมฺเม อธมฺมกมฺมสญฺญา วุฏฺฐาเปติ, อาปตฺติ ปาจิตฺติยสฺส.}

\prose{อธมฺมกมฺเม ธมฺมกมฺมสญฺญา, อาปตฺติ ทุกฺกฏสฺส.\csromanspacer{} อธมฺมกมฺเม เวมติกา, อาปตฺติ ทุกฺกฏสฺส.\csromanspacer{} อธมฺมกมฺเม อธมฺมกมฺมสญฺญา, อาปตฺติ ทุกฺกฏสฺส.}

\proseitemwithcloserruled{1135}{อนาปตฺติ ปริปุณฺณวีสติวสฺสํ กุมาริภูตํ เทฺว วสฺสานิ ฉสุ ธมฺเมสุ สิกฺขิตสิกฺขํ สํเฆน สมฺมตํ วุฏฺฐาเปติ, อุมฺมตฺติกาย อาทิกมฺมิกายาติ.}{ตติยสิกฺขาปทํ นิฏฺฐิตํ.}
\csromanmatikaanchor{mtk.230}
\csromantocmarkref{subsection}{4. จตุตฺถสิกฺขาปท}{mtk.230}
\bookmark[dest={mtk.230},level=2]{4. จตุตฺถสิกฺขาปท}
\csromanheader{2}{8. \textbf{กุมารีภูตวคฺค 4. จตุตฺถสิกฺขาปท}}
\proseitem{1136}{เตน สมเยน พุทฺโธ ภควา สาวตฺถิยํ วิหรติ เชตวเน อนาถปิณฺฑิกสฺส อาราเม.\csromanspacer{} เตน โข ปน สมเยน ภิกฺขุนิโย อูนทฺวาทสวสฺสา วุฏฺฐาเปนฺติ, ตา พาลา โหนฺติ อพฺยตฺตา, น ชานนฺติ \csromanfolio{439}กปฺปิยํ วา อกปฺปิยํ วา.\csromanspacer{} สทฺธิวิหารินิโยปิ พาลา โหนฺติ อพฺยตฺตา, น ชานนฺติ กปฺปิยํ วา อกปฺปิยํ วา.\csromanspacer{} ยา ตา ภิกฺขุนิโย อปฺปิจฺฉา ฯเปฯ ตา อุชฺฌายนฺติ ขิยฺยนฺติ วิปาเจนฺติ “กถํ หิ นาม ภิกฺขุนิโย อูนทฺวาทสวสฺสา วุฏฺฐาเปสฺสนฺตี”ติ ฯเปฯ.\csromanspacer{} สจฺจํ กิร ภิกฺขเว ภิกฺขุนิโย อูนทฺวาทสวสฺสา วุฏฺฐาเปนฺตีติ.\csromanspacer{} สจฺจํ ภควาติ.\csromanspacer{} วิครหิ พุทฺโธ ภควา ฯเปฯ.\csromanspacer{} กถํ หิ นาม ภิกฺขเว ภิกฺขุนิโย อูนทฺวาทสวสฺสา วุฏฺฐาเปสฺสนฺติ.\csromanspacer{} เนตํ ภิกฺขเว อปฺปสนฺนานํ วา ปสาทาย ฯเปฯ.\csromanspacer{} เอวญฺจ ปน ภิกฺขเว ภิกฺขุนิโย อิมํ สิกฺขาปทํ อุทฺทิสนฺตุ\csromandash{}}

\proseitem{1137}{\textbf{“ยา ปน ภิกฺขุนี อูนทฺวาทสวสฺสา วุฏฺฐาเปยฺย, ปาจิตฺติยนฺ”}ติ.}

\proseitem{1138}{\textbf{ยา ปนา}ติ ยา ยาทิสา ฯเปฯ.\csromanspacer{} \textbf{ภิกฺขุนี}ติ ฯเปฯ อยํ อิมสฺมิํ อตฺเถ อธิปฺเปตา ภิกฺขุนีติ.}

\prose{อู\textbf{นทฺวาทสวสฺสา} นาม อปฺปตฺตทฺวาทสวสฺสา.}

\prose{\textbf{วุฏฺฐาเปยฺยา}ติ อุปสมฺปาเทยฺย.}

\prose{“วุฏฺฐาเปสฺสามี”ติ คณํ อาจรินิํ วา ปตฺตํ วา จีวรํ วา ปริเยสติ, สีมํ วา สมฺมนฺนติ, อาปตฺติ ทุกฺกฏสฺส.\csromanspacer{} ญตฺติยา ทุกฺกฏํ.\csromanspacer{} ทฺวีหิ กมฺมวาจาหิ ทุกฺกฏา.\csromanspacer{} กมฺมวาจาปริโยสาเน อุปชฺฌายาย อาปตฺติ ปาจิตฺติยสฺส.\csromanspacer{} คณสฺส จ อาจรินิยา จ อาปตฺติ ทุกฺกฏสฺส.}

\proseitemwithcloserruled{1139}{อนาปตฺติ ปริปุณฺณทฺวาทสวสฺสา วุฏฺฐาเปติ, อุมฺมตฺติกาย อาทิกมฺมิกายาติ.}{จตุตฺถสิกฺขาปทํ นิฏฺฐิตํ.}
\csromanmatikaanchor{mtk.231}
\csromantocmarkref{subsection}{5. ปญฺจมสิกฺขาปท}{mtk.231}
\bookmark[dest={mtk.231},level=2]{5. ปญฺจมสิกฺขาปท}
\csromanheader{2}{8. \textbf{กุมารีภูตวคฺค 5. ปญฺจมสิกฺขาปท}}
\proseitem{1140}{เตน สมเยน พุทฺโธ ภควา สาวตฺถิยํ วิหรติ เชตวเน อนาถปิณฺฑิกสฺส อาราเม.\csromanspacer{} เตน โข ปน สมเยน ภิกฺขุนิโย ปริปุณฺณทฺวาทสวสฺสา สํเฆน อสมฺมตา วุฏฺฐาเปนฺติ, ตา พาลา โหนฺติ อพฺยตฺตา, น ชานนฺติ กปฺปิยํ วา อกปฺปิยํ วา.\csromanspacer{} สทฺธิวิหารินิโยปิ พาลา \csromanfolio{440}โหนฺติ อพฺยตฺตา, น ชานนฺติ กปฺปิยํ วา อกปฺปิยํ วา.\csromanspacer{} ยา ตา ภิกฺขุนิโย อปฺปิจฺฉา ฯเปฯ ตา อุชฺฌายนฺติ ขิยฺยนฺติ วิปาเจนฺติ “กถํ หิ นาม ภิกฺขุนิโย ปริปุณฺณทฺวาทสวสฺสา สํเฆน อสมฺมตา วุฏฺฐาเปสฺสนฺตี”ติ ฯเปฯ.\csromanspacer{} สจฺจํ กิร ภิกฺขเว ภิกฺขุนิโย ปริปุณฺณทฺวาทสวสฺสา สํเฆน อสมฺมตา วุฏฺฐาเปนฺตีติ.\csromanspacer{} สจฺจํ ภควาติ.\csromanspacer{} วิครหิ พุทฺโธ ภควา ฯเปฯ.\csromanspacer{} กถํ หิ นาม ภิกฺขเว ภิกฺขุนิโย ปริปุณฺณทฺวาทสวสฺสา สํเฆน อสมฺมตา วุฏฺฐาเปสฺสนฺติ.\csromanspacer{} เนตํ ภิกฺขเว อปฺปสนฺนานํ วา ปสาทาย ฯเปฯ วิครหิตฺวา ฯเปฯ ธมฺมิํ กถํ กตฺวา ภิกฺขู อามนฺเตสิ อนุชานามิ ภิกฺขเว ปริปุณฺณ\-ทฺวาทสวสฺสาย ภิกฺขุนิยา วุฏฺฐาปนสมฺมุติํ ทาตุํ.\csromanspacer{} เอวญฺจ ปน ภิกฺขเว ทาตพฺพา, ตาย ปริปุณฺณ\-ทฺวาทสวสฺสาย ภิกฺขุนิยา สํฆํ อุปสงฺกมิตฺวา เอกํสํ อุตฺตราสงฺคํ กริตฺวา วุฑฺฒานํ ภิกฺขุนีนํ ปาเท วนฺทิตฺวา อุกฺกุฏิกํ นิสีทิตฺวา อญฺชลิํ ปคฺคเหตฺวา เอวมสฺส วจนีโย “อหํ อยฺเย อิตฺถนฺนามา ปริปุณฺณทฺวาทสวสฺสา ภิกฺขุนี สํฆํ วุฏฺฐาปนสมฺมุติํ ยาจามี”ติ.\csromanspacer{} ทุติยมฺปิ ยาจิตพฺพา.\csromanspacer{} ตติยมฺปิ ยาจิตพฺพา.\csromanspacer{} สา ภิกฺขุนี สํเฆน ปริจฺฉินฺทิตพฺพา “พฺยตฺตายํ ภิกฺขุนี ลชฺชินี”ติ.\csromanspacer{} สเจ พาลา จ โหติ อลชฺชินี จ, น ทาตพฺพา.\csromanspacer{} สเจ พาลา โหติ ลชฺชินี, น ทาตพฺพา.\csromanspacer{} สเจ พฺยตฺตา โหติ อลชฺชินี, น ทาตพฺพา.\csromanspacer{} สเจ พฺยตฺตา จ โหติ ลชฺชินี จ, ทาตพฺพา.\csromanspacer{} เอวญฺจ ปน ภิกฺขเว ทาตพฺพา, พฺยตฺตาย ภิกฺขุนิยา ปฏิพลาย สํโฆ ญาเปตพฺโพ\csromandash{}}

\proseitem{1141}{“สุณาตุ เม อยฺเย สํโฆ, อยํ อิตฺถนฺนามา ปริปุณฺณทฺวาทสวสฺสา ภิกฺขุนี สํฆํ วุฏฺฐาปนสมฺมุติํ ยาจติ, ยทิ สํฆสฺส ปตฺตกลฺลํ, สํโฆ อิตฺถนฺนามาย ปริปุณฺณ\-ทฺวาทสวสฺสาย ภิกฺขุนิยา วุฏฺฐาปนสมฺมุติํ ทเทยฺย, เอสา ญตฺติ.}

\prose{สุณาตุ เม อยฺเย สํโฆ, อยํ อิตฺถนฺนามา ปริปุณฺณทฺวาทสวสฺสา ภิกฺขุนี สํฆํ วุฏฺฐาปนสมฺมุติํ ยาจติ, สํโฆ อิตฺถนฺนามาย ปริปุณฺณ\-ทฺวาทสวสฺสาย ภิกฺขุนิยา วุฏฺฐาปนสมฺมุติํ เทติ, ยสฺสา อยฺยาย ขมติ อิตฺถนฺนามาย ปริปุณฺณ\-ทฺวาทสวสฺสาย ภิกฺขุนิยา วุฏฺฐาปนสมฺมุติยา ทานํ, สา ตุณฺหสฺส, ยสฺสา นกฺขมติ, สา ภาเสยฺย.}

\prose{ทินฺนา สํเฆน อิตฺถนฺนามาย ปริปุณฺณ\-ทฺวาทสวสฺสาย ภิกฺขุนิยา วุฏฺฐาปนสมฺมุติ, ขมติ สํฆสฺส, ตสฺมา ตุณฺหี, เอวเมตํ ธารยามี”ติ.}

\prose{\csromanfolio{441}อถ โข ภควา ตา ภิกฺขุนิโย อเนกปริยาเยน วิครหิตฺวา ทุพฺภรตาย ฯเปฯ.\csromanspacer{} เอวญฺจ ปน ภิกฺขเว ภิกฺขุนิโย อิมํ สิกฺขาปทํ อุทฺทิสนฺตุ\csromandash{}}

\proseitem{1142}{\textbf{“ยา ปน ภิกฺขุนี ปริปุณฺณทฺวาทสวสฺสา สํเฆน อสมฺมตา วุฏฺฐาเปยฺย, ปาจิตฺติยนฺ”}ติ.}

\proseitem{1134}{\textbf{ยา ปนา}ติ ยา ยาทิสา ฯเปฯ.\csromanspacer{} \textbf{ภิกฺขุนี}ติ ฯเปฯ อยํ อิมสฺมิํ อตฺเถ อธิปฺเปตา ภิกฺขุนีติ.}

\prose{\textbf{ปริปุณฺณทฺวาทสวสฺสา} นาม ปตฺตทฺวาทสวสฺสา.}

\prose{\textbf{อสมฺมตา} นาม ญตฺติทุติเยน กมฺเมน วุฏฺฐาปนสมฺมุติ น ทินฺนา โหติ.}

\prose{\textbf{วุฏฺฐาเปยฺยา}ติ อุปสมฺปาเทยฺย.}

\prose{“วุฏฺฐาเปสฺสามี”ติ คณํ วา อาจรินิํ วา ปตฺตํ วา จีวรํ วา ปริเยสติ, สีมํ วา สมฺมนฺนติ, อาปตฺติ ทุกฺกฏสฺส.\csromanspacer{} ญตฺติยา ทุกฺกฏํ.\csromanspacer{} ทฺวีหิ กมฺมวาจาหิ ทุกฺกฏา.\csromanspacer{} กมฺมวาจาปริโยสาเน อุปชฺฌายาย อาปตฺติ ปาจิตฺติยสฺส.\csromanspacer{} คณสฺส จ อาจรินิยา จ อาปตฺติ ทุกฺกฏสฺส.}

\proseitem{1144}{ธมฺมกมฺเม ธมฺมกมฺมสญฺญา วุฏฺฐาเปติ, อาปตฺติ ปาจิตฺติยสฺส.\csromanspacer{} ธมฺมกมฺเม เวมติกา วุฏฺฐาเปติ, อาปตฺติ ปาจิตฺติยสฺส.\csromanspacer{} ธมฺมกมฺเม อธมฺมกมฺมสญฺญา วุฏฺฐาเปติ, อาปตฺติ ปาจิตฺติยสฺส.}

\prose{อธมฺมกมฺเม ธมฺมกมฺมสญฺญา, อาปตฺติ ทุกฺกฏสฺส.\csromanspacer{} อธมฺมกมฺเม เวมติกา อาปตฺติ ทุกฺกฏสฺส.\csromanspacer{} อธมฺมกมฺเม อธมฺมกมฺมสญฺญา, อาปตฺติ ทุกฺกฏสฺส.}

\proseitemwithcloserruled{1145}{อนาปตฺติ ปริปุณฺณทฺวาทสวสฺสา สํเฆน สมฺมตา วุฏฺฐาเปติ, อุมฺมตฺติกาย อาทิกมฺมิกายาติ.}{ปญฺจมสิกฺขาปทํ นิฏฺฐิตํ.}
\csromanmatikaanchor{mtk.232}
\csromantocmarkref{subsection}{6. ฉฏฺฐสิกฺขาปท}{mtk.232}
\bookmark[dest={mtk.232},level=2]{6. ฉฏฺฐสิกฺขาปท}
\csromanheader{2}{8. \textbf{กุมารีภูตวคฺค 6. ฉฏฺฐสิกฺขาปท}}
\proseitem{1146}{เตน สมเยน พุทฺโธ ภควา สาวตฺถิยํ วิหรติ เชตวเน อนาถปิณฺฑิกสฺส อาราเม.\csromanspacer{} เตน โข ปน สมเยน จณฺฑกาฬี \csromanfolio{442}\textbf{ภิกฺขุนี ภิกฺขุนิสํฆํ อุปสงฺกมิตฺวา วุฏฺฐาปนสมฺมุติํ ยาจติ.}\csromanspacer{} \textbf{อถ โข ภิกฺขุนิสํโฆ จณฺฑกาฬิํ ภิกฺขุนิํ ปริจฺฉินฺทิตฺวา “อลํ} ตาว เต อยฺเย วุฏฺฐาปิเตนา”ติ วุฏฺฐาปนสมฺมุติํ น อทาสิ.\csromanspacer{} จณฺฑกาฬี ภิกฺขุนี สาธูติ ปฏิสฺสุณิ.\csromanspacer{} เตน โข ปน สมเยน ภิกฺขุนิสํโฆ อญฺญาสํ ภิกฺขุนีนํ วุฏฺฐาปนสมฺมุติํ เทติ.\csromanspacer{} จณฺฑกาฬี ภิกฺขุนี อุชฺฌายติ ขิยฺยติ วิปาเจติ “อหเมว นูน พาลา, อหเมว นูน อลชฺชินี, ยํ สํโฆ อญฺญาสํ ภิกฺขุนีนํ วุฏฺฐาปนสมฺมุติํ เทติ, มยฺหเมว น เทตี”ติ.\csromanspacer{} ยา ตา ภิกฺขุนิโย อปฺปิจฺฉา ฯเปฯ ตา อุชฺฌายนฺติ ขิยฺยนฺติ วิปาเจนฺติ “กถํ หิ นาม อยฺยา จณฺฑกาฬี ‘อลํ ตาว เต อยฺเย วุฏฺฐาปิเตนา’ติ วุจฺจมานา สาธูติ ปฏิสฺสุณิตฺวา ปจฺฉา ขียนธมฺมํ อาปชฺชิสฺสตี”ติ ฯเปฯ.\csromanspacer{} สจฺจํ กิร ภิกฺขเว จณฺฑกาฬี ภิกฺขุนี “อลํ ตาว เต อยฺเย วุฏฺฐาปิเตนา”ติ วุจฺจมานา สาธูติ ปฏิสฺสุณิตฺวา ปจฺฉา ขียนธมฺมํ อาปชฺชตีติ\footnote{อาปชฺชีติ (ก)}.\csromanspacer{} สจฺจํ ภควาติ.\csromanspacer{} วิครหิ พุทฺโธ ภควา ฯเปฯ.\csromanspacer{} กถํ หิ นาม ภิกฺขเว จณฺฑกาฬี ภิกฺขุนี “อลํ ตาว เต อยฺเย วุฏฺฐาปิเตนา”ติ วุจฺจมานา สาธูติ ปฏิสฺสุณิตฺวา ปจฺฉา ขียนธมฺมํ อาปชฺชิสฺสติ.\csromanspacer{} เนตํ ภิกฺขเว อปฺปสนฺนานํ วา ปสาทาย ฯเปฯ.\csromanspacer{} เอวญฺจ ปน ภิกฺขเว ภิกฺขุนิโย อิมํ สิกฺขาปทํ อุทฺทิสนฺตุ\csromandash{}}

\proseitem{1147}{\textbf{“ยา ปน ภิกฺขุนี ‘อลํ ตาว เต อยฺเย วุฏฺฐาปิเตนา’ติ วุจฺจมานา สาธูติ ปฏิสฺสุณิตฺวา ปจฺฉา ขียนธมฺมํ อาปชฺเชยฺย, ปาจิตฺติยนฺ”}ติ.}

\proseitem{1148}{\textbf{ยา ปนา}ติ ยา ยาทิสา ฯเปฯ.\csromanspacer{} \textbf{ภิกฺขุนี}ติ ฯเปฯ อยํ อิมสฺมิํ อตฺเถ อธิปฺเปตา ภิกฺขุนีติ.}

\prose{\textbf{อลํ ตาว เต อยฺเย วุฏฺฐาปิเตนา}ติ อลํ ตาว เต อยฺเย อุปสมฺปาทิเตน.\csromanspacer{} สาธูติ ปฏิสฺสุณิตฺวา ปจฺฉา ขียนธมฺมํ อาปชฺชติ, อาปตฺติ ปาจิตฺติยสฺส.}

\proseitemwithcloser{1149}{อนาปตฺติ ปกติยา ฉนฺทา โทสา โมหา ภยา กโรนฺตํ ขิยฺยติ, อุมฺมตฺติกาย อาทิกมฺมิกายาติ.}{ฉฏฺฐสิกฺขาปทํ นิฏฺฐิตํ.}
\csromanmatikaanchor{mtk.233}
\csromantocmarkref{subsection}{7. สตฺตมสิกฺขาปท}{mtk.233}
\bookmark[dest={mtk.233},level=2]{7. สตฺตมสิกฺขาปท}
\csromanheader{2}{8. \textbf{กุมารีภูตวคฺค 7. สตฺตมสิกฺขาปท}}
\proseitem{1150}{\csromanfolio{443}เตน สมเยน พุทฺโธ ภควา สาวตฺถิยํ วิหรติ เชตวเน อนาถปิณฺฑิกสฺส อาราเม.\csromanspacer{} เตน โข ปน สมเยน อญฺญตรา สิกฺขมานา ถุลฺลนนฺทํ ภิกฺขุนิํ อุปสงฺกมิตฺวา อุปสมฺปทํ ยาจิ.\csromanspacer{} ถุลฺลนนฺทา ภิกฺขุนี ตํ สิกฺขมานํ “สเจ เม ตฺวํ อยฺเย จีวรํ ทสฺสสิ, เอวาหํ ตํ วุฏฺฐาเปสฺสามี”ติ วตฺวา เนว วุฏฺฐาเปติ, น วุฏฺฐาปนาย อุสฺสุกฺกํ กโรติ.\csromanspacer{} อถ โข สา สิกฺขมานา ภิกฺขุนีนํ เอตมตฺถํ อาโรเจสิ.\csromanspacer{} ยา ตา ภิกฺขุนิโย อปฺปิจฺฉา ฯเปฯ ตา อุชฺฌายนฺติ ขิยฺยนฺติ วิปาเจนฺติ “กถํ หิ นาม อยฺยา ถุลฺลนนฺทา สิกฺขมานํ ‘สเจ เม ตฺวํ อยฺเย จีวรํ ทสฺสสิ, เอวาหํ ตํ วุฏฺฐาเปสฺสามี’ติ วตฺวา เนว วุฏฺฐาเปสฺสติ, น วุฏฺฐาปนาย อุสฺสุกฺกํ กริสฺสตี”ติ ฯเปฯ.\csromanspacer{} สจฺจํ กิร ภิกฺขเว ถุลฺลนนฺทา ภิกฺขุนี สิกฺขมานํ “สเจ เม ตฺวํ อยฺเย จีวรํ ทสฺสสิ, เอวาหํ ตํ วุฏฺฐาเปสฺสามี”ติ วตฺวา เนว วุฏฺฐาเปติ, น วุฏฺฐาปนาย อุสฺสุกฺกํ กโรตีติ.\csromanspacer{} สจฺจํ ภควาติ.\csromanspacer{} วิครหิ พุทฺโธ ภควา ฯเปฯ.\csromanspacer{} กถํ หิ นาม ภิกฺขเว ถุลฺลนนฺทา ภิกฺขุนี สิกฺขมานํ “สเจ เม ตฺวํ อยฺเย จีวรํ ทสฺสสิ, เอวาหํ ตํ วุฏฺฐาเปสฺสามี”ติ วตฺวา เนว วุฏฺฐาเปสฺสติ, น วุฏฺฐาปนาย อุสฺสุกฺกํ กริสฺสติ.\csromanspacer{} เนตํ ภิกฺขเว อปฺปสนฺนานํ วา ปสาทาย ฯเปฯ.\csromanspacer{} เอวญฺจ ปน ภิกฺขเว ภิกฺขุนิโย อิมํ สิกฺขาปทํ อุทฺทิสนฺตุ\csromandash{}}

\proseitem{1151}{\textbf{“ยา ปน ภิกฺขุนี สิกฺขมานํ ‘สเจ เม ตฺวํ อยฺเย จีวรํ ทสฺสสิ, เอวาหํ ตํ วุฏฺฐาเปสฺสามี’ติ วตฺวา สา ปจฺฉา อนนฺตรายิกินี เนว วุฏฺฐาเปยฺย น วุฏฺฐาปนาย อุสฺสุกฺกํ กเรยฺย, ปาจิตฺติยนฺ”ติ}.}

\proseitem{1152}{\textbf{ยา ปนา}ติ ยา ยาทิสา ฯเปฯ.\csromanspacer{} \textbf{ภิกฺขุนี}ติ ฯเปฯ อยํ อิมสฺมิํ อตฺเถ อธิปฺเปตา ภิกฺขุนีติ.}

\prose{\textbf{สิกฺขมานา} นาม เทฺว วสฺสานิ ฉสุ ธมฺเมสุ สิกฺขิตสิกฺขา.}

\prose{\textbf{สเจ เม ตฺวํ อยฺเย จีวรํ ทสฺสสิ, เอวาหํ ตํ วุฏฺฐาเปสฺสามี}ติ เอวาหํ ตํ อุปสมฺปาเทสฺสามิ.}

\prose{\textbf{สา ปจฺฉา อนนฺตรายิกินี}ติ อสติ อนฺตราเย.}

\prose{\textbf{เนว วุฏฺฐาเปยฺยา}ติ น สยํ วุฏฺฐาเปยฺย.}

\prose{\csromanfolio{444}\textbf{น วุฏฺฐาปนาย อุสฺสุกฺกํ กเรยฺยา}ติ น อญฺญํ อาณาเปยฺย. “เนว วุฏฺฐาเปสฺสามิ น วุฏฺฐาปนาย อุสฺสุกฺกํ กริสฺสามี”ติ ธุรํ นิกฺขิตฺตมตฺเต อาปตฺติ ปาจิตฺติยสฺส.}

\proseitemwithcloserruled{1153}{อนาปตฺติ สติ อนฺตราเย, ปริเยสิตฺวา น ลภติ, คิลานาย อาปทาสุ อุมฺมตฺติกาย อาทิกมฺมิกายาติ.}{สตฺตมสิกฺขาปทํ นิฏฺฐิตํ.}
\csromanmatikaanchor{mtk.234}
\csromantocmarkref{subsection}{8. อฏฺฐมสิกฺขาปท}{mtk.234}
\bookmark[dest={mtk.234},level=2]{8. อฏฺฐมสิกฺขาปท}
\csromanheader{2}{8. \textbf{กุมารีภูตวคฺค 8. อฏฺฐมสิกฺขาปท}}
\proseitem{1154}{เตน สมเยน พุทฺโธ ภควา สาวตฺถิยํ วิหรติ เชตวเน อนาถปิณฺฑิกสฺส อาราเม.\csromanspacer{} เตน โข ปน สมเยน อญฺญตรา สิกฺขมานา ถุลฺลนนฺทํ ภิกฺขุนิํ อุปสงฺกมิตฺวา อุปสมฺปทํ ยาจิ.\csromanspacer{} ถุลฺลนนฺทา ภิกฺขุนี ตํ สิกฺขมานํ “สเจ มํ ตฺวํ อยฺเย เทฺว วสฺสานิ อนุพนฺธิสฺสสิ, เอวาหํ ตํ วุฏฺฐาเปสฺสามี”ติ วตฺวา เนว วุฏฺฐาเปติ, น วุฏฺฐาปนาย อุสฺสุกฺกํ กโรติ.\csromanspacer{} อถ โข สา สิกฺขมานา ภิกฺขุนีนํ เอตมตฺถํ อาโรเจสิ.\csromanspacer{} ยา ตา ภิกฺขุนิโย อปฺปิจฺฉา ฯเปฯ ตา อุชฺฌายนฺติ ขิยฺยนฺติ วิปาเจนฺติ “กถํ หิ นาม อยฺยา ถุลฺลนนฺทา สิกฺขมานํ ‘สเจ มํ ตฺวํ อยฺเย เทฺว วสฺสานิ อนุพนฺธิสฺสสิ, เอวาหํ ตํ วุฏฺฐาเปสฺสามี’ติ วตฺวา เนว วุฏฺฐาเปสฺสติ, น วุฏฺฐาปนาย อุสฺสุกฺกํ กริสฺสตี”ติ ฯเปฯ.\csromanspacer{} สจฺจํ กิร ภิกฺขเว ถุลฺลนนฺทา ภิกฺขุนี สิกฺขมานํ “สเจ มํ ตฺวํ อยฺเย เทฺว วสฺสานิ อนุพนฺธิสฺสสิ, เอวาหํ ตํ วุฏฺฐาเปสฺสามี”ติ วตฺวา เนว วุฏฺฐาเปติ, น วุฏฺฐาปนาย อุสฺสุกฺกํ กโรตีติ.\csromanspacer{} สจฺจํ ภควาติ.\csromanspacer{} วิครหิ พุทฺโธ ภควา ฯเปฯ.\csromanspacer{} กถํ หิ นาม ภิกฺขเว ถุลฺลนนฺทา ภิกฺขุนี สิกฺขมานํ “สเจ มํ ตฺวํ อยฺเย เทฺว วสฺสานิ อนุพนฺธิสฺสสิ, เอวาหํ ตํ วุฏฺฐาเปสฺสามี”ติ วตฺวา เนว วุฏฺฐาเปสฺสติ, น วุฏฺฐาปนาย อุสฺสุกฺกํ กริสฺสติ.\csromanspacer{} เนตํ ภิกฺขเว อปฺปสนฺนานํ วา ปสาทาย ฯเปฯ.\csromanspacer{} เอวญฺจ ปน ภิกฺขเว ภิกฺขุนิโย อิมํ สิกฺขาปทํ อุทฺทิสนฺตุ\csromandash{}}

\proseitem{1155}{\textbf{“ยา ปน ภิกฺขุนี สิกฺขมานํ ‘สเจ มํ ตฺวํ อยฺเย เทฺว วสฺสานิ อนุพนฺธิสฺสสิ, เอวาหํ ตํ วุฏฺฐาเปสฺสามี’ติ วตฺวา สา ปจฺฉา} \csromanfolio{445}\textbf{อนนฺตรายิกินี เนว วุฏฺฐาเปยฺย, น วุฏฺฐาปนาย อุสฺสุกฺกํ กเรยฺย, ปาจิตฺติยนฺ”}ติ.}

\proseitem{1156}{\textbf{ยา ปนา}ติ ยา ยาทิสา ฯเปฯ.\csromanspacer{} \textbf{ภิกฺขุนี}ติ ฯเปฯ อยํ อิมสฺมิํ อตฺเถ อธิปฺเปตา ภิกฺขุนีติ.}

\prose{\textbf{สิกฺขมานา} นาม เทฺว วสฺสานิ ฉสุ ธมฺเมสุ สิกฺขิตสิกฺขา.}

\prose{\textbf{สเจ มํ ตฺวํ อยฺเย เทฺว วสฺสานิ อนุพนฺธิสฺสสี}ติ เทฺว สํวจฺฉรานิ อุปฏฺฐหิสฺสสิ.}

\prose{\textbf{เอวาหํ ตํ วุฏฺฐาเปสฺสามี}ติ เอวาหํ ตํ อุปสมฺปาเทสฺสามิ.}

\prose{\textbf{สา ปจฺฉา อนนฺตรายิกินี}ติ อสติ อนฺตราเย.}

\prose{\textbf{เนว วุฏฺฐาเปยฺยา}ติ น สยํ วุฏฺฐาเปยฺย.}

\prose{\textbf{น วุฏฺฐาปนาย อุสฺสุกฺกํ กเรยฺยา}ติ น อญฺญํ อาณาเปยฺย. “เนว วุฏฺฐาเปสฺสามิ น วุฏฺฐาปนาย อุสฺสุกฺกํ กริสฺสามี”ติ ธุรํ นิกฺขิตฺตมตฺเต อาปตฺติ ปาจิตฺติยสฺส.}

\proseitemwithcloserruled{1157}{อนาปตฺติ สติ อนฺตราเย, ปริเยสิตฺวา น ลภติ, คิลานาย อาปทาสุ อุมฺมตฺติกาย อาทิกมฺมิกายาติ.}{อฏฺฐมสิกฺขาปทํ นิฏฺฐิตํ.}
\csromanmatikaanchor{mtk.235}
\csromantocmarkref{subsection}{9. นวมสิกฺขาปท}{mtk.235}
\bookmark[dest={mtk.235},level=2]{9. นวมสิกฺขาปท}
\csromanheader{2}{8. \textbf{กุมารีภูตวคฺค 9. นวมสิกฺขาปท}}
\proseitem{1158}{เตน สมเยน พุทฺโธ ภควา สาวตฺถิยํ วิหรติ เชตวเน อนาถปิณฺฑิกสฺส อาราเม.\csromanspacer{} เตน โข ปน สมเยน ถุลฺลนนฺทา ภิกฺขุนี ปุริสสํสฏฺฐํ กุมารกสํสฏฺฐํ จณฺฑิํ โสกาวาสํ\footnote{โสกวสฺสํ (สี) โสกาวสฺสํ (สฺยา)} จณฺฑกาฬิํ สิกฺขมานํ วุฏฺฐาเปติ.\csromanspacer{} ยา ตา ภิกฺขุนิโย อปฺปิจฺฉา ฯเปฯ ตา อุชฺฌายนฺติ ขิยฺยนฺติ วิปาเจนฺติ “กถํ หิ นาม อยฺยา ถุลฺลนนฺทา ปุริสสํสฏฺฐํ กุมารกสํสฏฺฐํ จณฺฑิํ โสกาวาสํ1 จณฺฑกาฬิํ สิกฺขมานํ วุฏฺฐาเปสฺสตี”ติ ฯเปฯ.\csromanspacer{} สจฺจํ กิร ภิกฺขเว ถุลฺลนนฺทา ภิกฺขุนี ปุริสสํสฏฺฐํ กุมารกสํสฏฺฐํ จณฺฑิํ โสกาวาสํ1 จณฺฑกาฬิํ สิกฺขมานํ วุฏฺฐาเปตีติ\footnote{วุฏฺฐาเปสีติ (ก)}.\csromanspacer{} สจฺจํ ภควาติ.\csromanspacer{} วิครหิ \csromanfolio{446}พุทฺโธ ภควา ฯเปฯ.\csromanspacer{} กถํ หิ นาม ภิกฺขเว ถุลฺลนนฺทา ภิกฺขุนี ปุริสสํสฏฺฐํ กุมารกสํสฏฺฐํ จณฺฑิํ โสกาวาสํ จณฺฑกาฬิํ สิกฺขมานํ วุฏฺฐาเปสฺสติ.\csromanspacer{} เนตํ ภิกฺขเว อปฺปสนฺนานํ วา ปสาทาย ฯเปฯ.\csromanspacer{} เอวญฺจ ปน ภิกฺขเว ภิกฺขุนิโย อิมํ สิกฺขาปทํ อุทฺทิสนฺตุ\csromandash{}}

\proseitem{1159}{\textbf{“ยา ปน ภิกฺขุนี ปุริสสํสฏฺฐํ กุมารกสํสฏฺฐํ จณฺฑิํ โสกาวาสํ สิกฺขมานํ วุฏฺฐาเปยฺย, ปาจิตฺติยนฺ”}ติ.}

\proseitem{1160}{\textbf{ยา ปนา}ติ ยา ยาทิสา ฯเปฯ.\csromanspacer{} \textbf{ภิกฺขุนี}ติ ฯเปฯ อยํ อิมสฺมิํ อตฺเถ อธิปฺเปตา ภิกฺขุนีติ.}

\prose{\textbf{ปุริโส} นาม ปตฺตวีสติวสฺโส.\csromanspacer{} \textbf{กุมารโก} นาม อปฺปตฺตวีสติวสฺโส.\csromanspacer{} \textbf{สํสฏฺฐา} นาม อนนุโลมิเกน กายิกวาจสิเกน สํสฏฺฐา.\csromanspacer{} \textbf{จณฺฑี} นาม โกธนา วุจฺจติ.}

\prose{\textbf{โสกาวาสา} นาม ปเรสํ ทุกฺขํ อุปฺปาเทติ, โสกํ อาวิสติ.\csromanspacer{} \textbf{สิกฺขมานา} นาม เทฺว วสฺสานิ ฉสุ ธมฺเมสุ สิกฺขิตสิกฺขา.\csromanspacer{} \textbf{วุฏฺฐาเปยฺยา}ติ อุปสมฺปาเทยฺย.}

\prose{“วุฏฺฐาเปสฺสามี”ติ คณํ วา อาจรินิํ วา ปตฺตํ วา จีวรํ วา ปริเยสติ, สีมํ วา สมฺมนฺนติ, อาปตฺติ ทุกฺกฏสฺส.\csromanspacer{} ญตฺติยา ทุกฺกฏํ.\csromanspacer{} ทฺวีหิ กมฺมวาจาหิ ทุกฺกฏา.\csromanspacer{} กมฺมวาจาปริโยสาเน อุปชฺฌายาย อาปตฺติ ปาจิตฺติยสฺส.\csromanspacer{} คณสฺส จ อาจรินิยา จ อาปตฺติ ทุกฺกฏสฺส.}

\proseitemwithcloserruled{1161}{อนาปตฺติ อชานนฺตี วุฏฺฐาเปติ, อุมฺมตฺติกาย อาทิกมฺมิกายาติ.}{นวมสิกฺขาปทํ นิฏฺฐิตํ.}
\csromanmatikaanchor{mtk.236}
\csromantocmarkref{subsection}{10. ทสมสิกฺขาปท}{mtk.236}
\bookmark[dest={mtk.236},level=2]{10. ทสมสิกฺขาปท}
\csromanheader{2}{8. \textbf{กุมารีภูตวคฺ}ค 10. ทสมสิกฺขาปท}
\proseitem{1162}{เตน สมเยน พุทฺโธ ภควา สาวตฺถิยํ วิหรติ เชตวเน อนาถปิณฺฑิกสฺส อาราเม.\csromanspacer{} เตน โข ปน สมเยน ถุลฺลนนฺทา ภิกฺขุนี มาตาปิตูหิปิ สามิเกนปิ อนนุญฺญาตํ สิกฺขมานํ วุฏฺฐาเปติ.\csromanspacer{} มาตาปิตโรปิ สามิโกปิ อุชฺฌายนฺติ ขิยฺยนฺติ วิปาเจนฺติ “กถํ หิ \csromanfolio{447}นาม อยฺยา ถุลฺลนนฺทา อเมฺหหิ อนนุญฺญาตํ สิกฺขมานํ วุฏฺฐาเปสฺสตี”ติ.\csromanspacer{} อสฺโสสุํ โข ภิกฺขุนิโย มาตาปิตูนํปิ สามิกสฺสปิ อุชฺฌายนฺตานํ ขิยฺยนฺตานํ วิปาเจนฺตานํ.\csromanspacer{} ยา ตา ภิกฺขุนิโย อปฺปิจฺฉา ฯเปฯ ตา อุชฺฌายนฺติ ขิยฺยนฺติ วิปาเจนฺติ “กถํ หิ นาม อยฺยา ถุลฺลนนฺทา มาตาปิตูหิปิ สามิเกนปิ อนนุญฺญาตํ สิกฺขมานํ วุฏฺฐาเปสฺสตี”ติ ฯเปฯ.\csromanspacer{} สจฺจํ กิร ภิกฺขเว ถุลฺลนนฺทา ภิกฺขุนี มาตาปิตูหิปิ สามิเกนปิ อนนุญฺญาตํ สิกฺขมานํ วุฏฺฐาเปตีติ.\csromanspacer{} สจฺจํ ภควาติ.\csromanspacer{} วิครหิ พุทฺโธ ภควา ฯเปฯ.\csromanspacer{} กถํ หิ นาม ภิกฺขเว ถุลฺลนนฺทา ภิกฺขุนี มาตาปิตูหิปิ สามิเกนปิ อนนุญฺญาตํ สิกฺขมานํ วุฏฺฐาเปสฺสติ.\csromanspacer{} เนตํ ภิกฺขเว อปฺปสนฺนานํ วา ปสาทาย ฯเปฯ.\csromanspacer{} เอวญฺจ ปน ภิกฺขเว ภิกฺขุนิโย อิมํ สิกฺขาปทํ อุทฺทิสนฺตุ\csromandash{}}

\proseitem{1163}{\textbf{“ยา ปน ภิกฺขุนี มาตาปิตูหิ วา สามิเกน วา อนนุญฺญาตํ สิกฺขมานํ วุฏฺฐาเปยฺย, ปาจิตฺติยนฺ”}ติ.}

\proseitem{1164}{\textbf{ยา ปนา}ติ ยา ยาทิสา ฯเปฯ.\csromanspacer{} \textbf{ภิกฺขุนี}ติ ฯเปฯ อยํ อิมสฺมิํ อตฺเถ อธิปฺเปตา ภิกฺขุนีติ.}

\prose{\textbf{มาตาปิตโร} นาม ชนกา วุจฺจนฺติ.\csromanspacer{} \textbf{สามิโก} นาม เยน ปริคฺคหิตา โหติ.\csromanspacer{} \textbf{อนนุญฺญาตา}ติ อนาปุจฺฉา.\csromanspacer{} \textbf{สิกฺขมานา} นาม เทฺว วสฺสานิ ฉสุ ธมฺเมสุ สิกฺขิตสิกฺขา.\csromanspacer{} \textbf{วุฏฺฐาเปยฺยา}ติ อุปสมฺปาเทยฺย.}

\prose{“วุฏฺฐาเปสฺสามี”ติ คณํ วา อาจรินิํ วา ปตฺตํ วา จีวรํ วา ปริเยสติ, สีมํ วา สมฺมนฺนติ, อาปตฺติ ทุกฺกฏสฺส.\csromanspacer{} ญตฺติยา ทุกฺกฏํ.\csromanspacer{} ทฺวีหิ กมฺมวาจาหิ ทุกฺกฏา.\csromanspacer{} กมฺมวาจาปริโยสาเน อุปชฺฌายาย อาปตฺติ ปาจิตฺติยสฺส.\csromanspacer{} คณสฺส จ อาจรินิยา จ อาปตฺติ ทุกฺกฏสฺส.}

\proseitemwithcloserruled{1165}{อนาปตฺติ อชานนฺตี วุฏฺฐาเปติ, อปโลเกตฺวา วุฏฺฐาเปติ, อุมฺมตฺติกาย อาทิกมฺมิกายาติ.}{ทสมสิกฺขาปทํ นิฏฺฐิตํ.}
\csromanmatikaanchor{mtk.237}
\csromantocmarkref{subsection}{11. เอกาทสมสิกฺขาปท}{mtk.237}
\bookmark[dest={mtk.237},level=2]{11. เอกาทสมสิกฺขาปท}
\csromanheader{2}{8. \textbf{กุมารีภูตวคฺคฺ}อ 11. เอกาทสมสิกฺขาปท}
\proseitem{1166}{เตน สมเยน พุทฺโธ ภควา ราชคเห วิหรติ เวฬุวเน กลนฺทกนิวาเป.\csromanspacer{} เตน โข ปน สมเยน ถุลฺลนนฺทา ภิกฺขุนี “สิกฺขมานํ \csromanfolio{448}วุฏฺฐาเปสฺสามี”ติ เถเร ภิกฺขู สนฺนิปาเตตฺวา ปหูตํ ขาทนียํ โภชนียํ ปสฺสิตฺวา “น ตาวาหํ อยฺยา สิกฺขมานํ วุฏฺฐาเปสฺสามี”ติ เถเร ภิกฺขู อุยฺโยเชตฺวา เทวทตฺตํ โกกาลิกํ กฏโมทกติสฺสกํ ขณฺฑเทวิยา ปุตฺตํ สมุทฺททตฺตํ สนฺนิปาเตตฺวา สิกฺขมานํ วุฏฺฐาเปสิ.\csromanspacer{} ยา ตา ภิกฺขุนิโย อปฺปิจฺฉา ฯเปฯ ตา อุชฺฌายนฺติ ขิยฺยนฺติ วิปาเจนฺติ “กถํ หิ นาม อยฺยา ถุลฺลนนฺทา ปาริวาสิกฉนฺททาเนน สิกฺขมานํ วุฏฺฐาเปสฺสตี”ติ ฯเปฯ.\csromanspacer{} สจฺจํ กิร ภิกฺขเว ถุลฺลนนฺทา ภิกฺขุนี ปาริวาสิกฉนฺททาเนน สิกฺขมานํ วุฏฺฐาเปตีติ\footnote{วุฏฺฐาเปสีติ (ก)}.\csromanspacer{} สจฺจํ ภควาติ.\csromanspacer{} วิครหิ พุทฺโธ ภควา ฯเปฯ.\csromanspacer{} กถํ หิ นาม ภิกฺขเว ถุลฺลนนฺทา ภิกฺขุนี ปาริวาสิกฉนฺททาเนน สิกฺขมานํ วุฏฺฐาเปสฺสติ.\csromanspacer{} เนตํ ภิกฺขเว อปฺปสนฺนานํ วา ปสาทาย ฯเปฯ.\csromanspacer{} เอวญฺจ ปน ภิกฺขเว ภิกฺขุนิโย อิมํ สิกฺขาปทํ อุทฺทิสนฺตุ.}

\proseitem{1167}{\textbf{“ยา ปน ภิกฺขุนี ปาริวาสิกฉนฺททาเนน สิกฺขมานํ วุฏฺฐาเปยฺย, ปาจิตฺติยนฺ”}ติ.}

\proseitem{1168}{\textbf{ยา ปนา}ติ ยา ยาทิสา ฯเปฯ.\csromanspacer{} พฺ\textbf{หิกฺขุนี}ติ ฯเปฯ อยํ อิมสฺมิํ อตฺเถ อธิปฺเปตา ภิกฺขุนีติ.}

\prose{\textbf{ปาริวาสิกฉนฺททาเนนา}ติ วุฏฺฐิตาย ปริสาย.\csromanspacer{} \textbf{สิกฺขมานา} นาม เทฺว วสฺสานิ ฉสุ ธมฺเมสุ สิกฺขิตสิกฺขา.\csromanspacer{} \textbf{วุฏฺฐาเปยฺยา}ติ อุปสมฺปาเทยฺย.}

\prose{“วุฏฺฐาเปสฺสามี”ติ คณํ วา อาจรินิํ วา ปตฺตํ วา จีวรํ วา ปริเยสติ, สีมํ วา สมฺมนฺนติ, อาปตฺติ ทุกฺกฏสฺส.\csromanspacer{} ญตฺติยา ทุกฺกฏํ.\csromanspacer{} ทฺวีหิ กมฺมวาจาหิ ทุกฺกฏา.\csromanspacer{} กมฺมวาจาปริโยสาเน อุปชฺฌายาย อาปตฺติ ปาจิตฺติยสฺส.\csromanspacer{} คณสฺส จ อาจรินิยา จ อาปตฺติ ทุกฺกฏสฺส.}

\proseitemwithcloserruled{1169}{อนาปตฺติ อวุฏฺฐิตาย ปริสาย วุฏฺฐาเปติ, อุมฺมตฺติกาย อาทิกมฺมิกายาติ.}{เอกาทสมสิกฺกาปทํ นิฏฺฐิตํ.}
\csromanmatikaanchor{mtk.238}
\csromantocmarkref{subsection}{12. ทฺวาทสมสิกฺขาปท}{mtk.238}
\bookmark[dest={mtk.238},level=2]{12. ทฺวาทสมสิกฺขาปท}
\csromanheader{2}{8. \textbf{กุมารีภูตวคฺคฺ}อ 12. ทฺวาทสมสิกฺขาปท}
\proseitem{1170}{เตน สมเยน พุทฺโธ ภควา สาวตฺถิยํ วิหรติ เชตวเน อนาถปิณฺฑิกสฺส อาราเม.\csromanspacer{} เตน โข ปน สมเยน ภิกฺขุนิโย \csromanfolio{449}อนุวสฺสํ วุฏฺฐาเปนฺติ, อุปสฺสโย น สมฺมติ.\csromanspacer{} มนุสฺสา\footnote{มนุสฺสา วิหารจาริกํ อาหิณฺฑนฺตา ปสฺสิตฺวา (สี)} อุชฺฌายนฺติ ขิยฺยนฺติ วิปาเจนฺติ “กถํ หิ นาม ภิกฺขุนิโย อนุวสฺสํ วุฏฺฐาเปสฺสนฺติ อุปสฺสโย น สมฺมตี”ติ.\csromanspacer{} อสฺโสสุํ โข ภิกฺขุนิโย เตสํ มนุสฺสานํ อุชฺฌายนฺตานํ ขิยฺยนฺตานํ วิปาเจนฺตานํ, ยา ตา ภิกฺขุนิโย อปฺปิจฺฉา ฯเปฯ ตา อุชฺฌายนฺติ ขิยฺยนฺติ วิปาเจนฺติ “กถํ หิ นาม ภิกฺขุนิโย อนุวสฺสํ วุฏฺฐาเปสฺสนฺตี”ติ ฯเปฯ.\csromanspacer{} สจฺจํ กิร ภิกฺขเว ภิกฺขุนิโย อนุวสฺสํ วุฏฺฐาเปนฺตีติ.\csromanspacer{} สจฺจํ ภควาติ.\csromanspacer{} วิครหิ พุทฺโธ ภควา ฯเปฯ.\csromanspacer{} กถํ หิ นาม ภิกฺขเว ภิกฺขุนิโย อนุวสฺสํ วุฏฺฐาเปสฺสนฺติ.\csromanspacer{} เนตํ ภิกฺขเว อปฺปสนฺนานํ วา ปสาทาย ฯเปฯ.\csromanspacer{} เอวญฺจ ปน ภิกฺขเว \textbf{ภิกฺขุนิโย อิมํ สิกฺขาปทํ อุทฺทิสนฺตุ\csromandash{}}}

\proseitem{1171}{\textbf{“ยา ปน ภิกฺขุนี อนุวสฺสํ วุฏฺฐาเปยฺย, ปาจิตฺติยนฺ”}ติ.}

\proseitem{1172}{\textbf{ยา ปนา}ติ ยา ยาทิสา ฯเปฯ.\csromanspacer{} \textbf{ภิกฺขุนี}ติ ฯเปฯ อยํ อิมสฺมิํ อตฺเถ อธิปฺเปตา ภิกฺขุนีติ.}

\prose{\textbf{อนุวสฺส}นฺติ อนุสํวจฺฉรํ.\csromanspacer{} \textbf{วุฏฺฐาเปยฺยา}ติ อุปสมฺปาเทยฺย.}

\prose{“วุฏฺฐาเปสฺสามี”ติ คณํ วา อาจรินิํ วา ปตฺตํ วา จีวรํ วา ปริเยสติ, สีมํ วา สมฺมนฺนติ, อาปตฺติ ทุกฺกฏสฺส.\csromanspacer{} ญตฺติยา ทุกฺกฏํ.\csromanspacer{} ทฺวีหิ กมฺมวาจาหิ ทุกฺกฏา.\csromanspacer{} กมฺมวาจาปริโยสาเน อุปชฺฌายาย อาปตฺติ ปาจิตฺติยสฺส.\csromanspacer{} คณสฺส จ อาจรินิยา จ อาปตฺติ ทุกฺกฏสฺส.}

\proseitemwithcloserruled{1173}{อนาปตฺติ เอกนฺตริกํ วุฏฺฐาเปติ, อุมฺมตฺติกาย อาทิกมฺมิกายาติ.}{ทฺวาทสมสิกฺขาปทํ นิฏฺฐิตํ.}
\csromanmatikaanchor{mtk.239}
\csromantocmarkref{subsection}{13. เตรสมสิกฺขาปท}{mtk.239}
\bookmark[dest={mtk.239},level=2]{13. เตรสมสิกฺขาปท}
\csromanheader{2}{8. \textbf{กุมารีภูตวคฺคฺ}อ 13. เตรสมสิกฺขาปท}
\proseitem{1174}{เตน สมเยน พุทฺโธ ภควา สาวตฺถิยํ วิหรติ เชตวเน อนาถปิณฺฑิกสฺส อาราเม.\csromanspacer{} เตน โข ปน สมเยน ภิกฺขุนิโย เอกํ วสฺสํ เทฺว วุฏฺฐาเปนฺติ, อุปสฺสโย ตเถว น สมฺมติ.\csromanspacer{} มนุสฺสา1 ตเถว \csromanfolio{450}อุชฺฌายนฺติ ขิยฺยนฺติ วิปาเจนฺติ “กถํ หิ นาม ภิกฺขุนิโย เอกํ วสฺสํ เทฺว วุฏฺฐาเปสฺสนฺติ, อุปสฺสโย ตเถว น สมฺมตี”ติ.\csromanspacer{} อสฺโสสุํ โข ภิกฺขุนิโย เตสํ มนุสฺสานํ อุชฺฌายนฺตานํ ขิยฺยนฺตานํ วิปาเจนฺตานํ.\csromanspacer{} ยา ตา ภิกฺขุนิโย อปฺปิจฺฉา ฯเปฯ ตา อุชฺฌายนฺติ ขิยฺยนฺติ วิปาเจนฺติ “กถํ หิ นาม ภิกฺขุนิโย เอกํ วสฺสํ เทฺว วุฏฺฐาเปสฺสนฺตี”ติ ฯเปฯ.\csromanspacer{} สจฺจํ กิร ภิกฺขเว ภิกฺขุนิโย เอกํ วสฺสํ เทฺว วุฏฺฐาเปนฺตีติ.\csromanspacer{} สจฺจํ ภควาติ.\csromanspacer{} วิครหิ พุทฺโธ ภควา ฯเปฯ.\csromanspacer{} กถํ หิ นาม ภิกฺขเว ภิกฺขุนิโย เอกํ วสฺสํ เทฺว วุฏฺฐาเปสฺสนฺติ.\csromanspacer{} เนตํ ภิกฺขเว อปฺปสนฺนานํ วา ปสาทาย ฯเปฯ.\csromanspacer{} เอวญฺจ ปน ภิกฺขเว ภิกฺขุนิโย อิมํ สิกฺขาปทํ อุทฺทิสนฺตุ\csromandash{}}

\proseitem{1175}{\textbf{“ยา ปน ภิกฺขุนี เอกํ วสฺสํ เทฺว วุฏฺฐาเปยฺย, ปาจิตฺติยนฺ}”ติ.}

\proseitem{1176}{\textbf{ยา ปนา}ติ ยา ยาทิสา ฯเปฯ.\csromanspacer{} \textbf{ภิกฺขุนี}ติ ฯเปฯ อยํ อิมสฺมิํ อตฺเถ อธิปฺเปตา ภิกฺขุนีติ.}

\prose{\textbf{เอกํ วสฺสนฺ}ติ เอกํ สํวจฺฉรํ.\csromanspacer{} \textbf{เทฺว วุฏฺฐาเปยฺยา}ติ เทฺว อุปสมฺปาเทยฺย.}

\prose{“เทฺว วุฏฺฐาเปสฺสามี”ติ คณํ วา อาจรินิํ วา ปตฺตํ วา จีวรํ วา ปริเยสติ, สีมํ วา สมฺมนฺนติ, อาปตฺติ ทุกฺกฏสฺส.\csromanspacer{} ญตฺติยา ทุกฺกฏํ.\csromanspacer{} ทฺวีหิ กมฺมวาจาหิ ทุกฺกฏา.\csromanspacer{} กมฺมวาจาปริโยสาเน อุปชฺฌายาย อาปตฺติ ปาจิตฺติยสฺส.\csromanspacer{} คณสฺส จ อาจรินิยา จ อาปตฺติ ทุกฺกฏสฺส.}

\proseitemwithcloser{1177}{อนาปตฺติ เอกนฺตริกํ เอกํ วุฏฺฐาเปติ, อุมฺมตฺติกาย อาทิกมฺมิกายาติ.}{เตรสมสิกฺขาปทํ นิฏฺฐิตํ.}
\csromanheader{2}{กุมาริภูตวคฺโค อฏฺฐโม.}
\csromansectionrule
\csromanmatikaanchor{mtk.240}
\csromanmatikaanchor{mtk.241}
\csromantocmarknopage{section}{9. ฉตฺตุปาหนวคฺค}{mtk.240}
\bookmark[dest={mtk.240},level=1]{9. ฉตฺตุปาหนวคฺค}
\csromantocmarkref{subsection}{1. ปฐมสิกฺขาปท}{mtk.241}
\bookmark[dest={mtk.241},level=2]{1. ปฐมสิกฺขาปท}
\setcsromanheadh{9. ฉตฺตุปาหนวคฺค}
\csromanheader{2}{9. \textbf{ฉตฺตุปาหนวคฺค 1. ปฐมสิกฺขาปท}}
\proseitem{1178}{เตน สมเยน พุทฺโธ ภควา สาวตฺถิยํ วิหรติ เชตวเน อนาถปิณฺฑิกสฺส อาราเม.\csromanspacer{} เตน โข ปน สมเยน ฉพฺพคฺคิยา ภิกฺขุนิโย ฉตฺตุปาหนํ ธาเรนฺติ.\csromanspacer{} มนุสฺสา อุชฺฌายนฺติ ขิยฺยนฺติ วิปาเจนฺติ}

\prose{\csromanfolio{451}“กถํ หิ นาม ภิกฺขุนิโย ฉตฺตุปาหนํ ธาเรสฺสนฺติ, เสยฺยถาปิ คิหินิโย กามโภคินิโย”ติ.\csromanspacer{} อสฺโสสุํ โข ภิกฺขุนิโย เตสํ มนุสฺสานํ อุชฺฌายนฺตานํ ขิยฺยนฺตานํ วิปาเจนฺตานํ.\csromanspacer{} ยา ตา ภิกฺขุนิโย อปฺปิจฺฉา ฯเปฯ ตา อุชฺฌายนฺติ ขิยฺยนฺติ วิปาเจนฺติ “กถํ หิ นาม ฉพฺพคฺคิยา ภิกฺขุนิโย ฉตฺตุปาหนํ ธาเรสฺสนฺตี”ติ ฯเปฯ.\csromanspacer{} สจฺจํ กิร ภิกฺขเว ฉพฺพคฺคิยา ภิกฺขุนิโย ฉตฺตุปาหนํ ธาเรนฺตีติ.\csromanspacer{} สจฺจํ ภควาติ.\csromanspacer{} วิครหิ พุทฺโธ ภควา ฯเปฯ.\csromanspacer{} กถํ หิ นาม ภิกฺขเว ฉพฺพคฺคิยา ภิกฺขุนิโย ฉตฺตุปาหนํ ธาเรสฺสนฺติ.\csromanspacer{} เนตํ ภิกฺขเว อปฺปสนฺนานํ วา ปสาทาย ฯเปฯ.\csromanspacer{} เอวญฺจ ปน ภิกฺขเว ภิกฺขุนิโย อิมํ สิกฺขาปทํ อุทฺทิสนฺตุ\csromandash{}}

\prose{\textbf{“ยา ปน ภิกฺขุนี ฉตฺตุปาหนํ ธาเรยฺย, ปาจิตฺติยนฺ”}ติ.}

\prose{เอวญฺจิทํ ภควตา ภิกฺขุนีนํ สิกฺขาปทํ ปญฺญตฺตํ โหติ.}

\proseitem{1179}{เตน โข ปน สมเยน อญฺญตรา ภิกฺขุนี คิลานา โหติ, ตสฺสา วินา ฉตฺตุปาหนํ น ผาสุ โหติ ฯเปฯ.\csromanspacer{} ภควโต เอตมตฺถํ อาโรเจสุํ ฯเปฯ.\csromanspacer{} อนุชานามิ ภิกฺขเว คิลานาย ภิกฺขุนิยา ฉตฺตุปาหนํ.\csromanspacer{} เอวญฺจ ปน ภิกฺขเว ภิกฺขุนิโย อิมํ สิกฺขาปทํ อุทฺทิสนฺตุ\csromandash{}}

\proseitem{1180}{\textbf{“ยา ปน ภิกฺขุนี อคิลานา ฉตฺตุปาหนํ ธาเรยฺย, ปาจิตฺติยนฺ”}ติ.}

\proseitem{1181}{\textbf{ยา ปนา}ติ ยา ยาทิสา ฯเปฯ.\csromanspacer{} \textbf{ภิกฺขุนี}ติ ฯเปฯ อยํ อิมสฺมิํ อตฺเถ อธิปฺเปตา ภิกฺขุนีติ.}

\prose{\textbf{อคิลานา} นาม ยสฺสา วินา ฉตฺตุปาหนํ ผาสุ โหติ.}

\prose{\textbf{คิลานา} นาม ยสฺสา วินา ฉตฺตุปาหนํ น ผาสุ โหติ.}

\prose{\textbf{ฉตฺตํ} นาม ตีณิ ฉตฺตานิ เสตจฺฉตฺตํ กิลญฺชจฺฉตฺตํ ปณฺณจฺฉตฺตํ มณฺฑลพทฺธํ สลากพทฺธํ.\csromanspacer{} \textbf{ธาเรยฺยา}ติ สกิํปิ ธาเรติ, อาปตฺติ}

\proseitem{1182}{อคิลานา อคิลานสญฺญา ฉตฺตุปาหนํ ธาเรติ, อาปตฺติ ปาจิตฺติยสฺส.\csromanspacer{} อคิลานา เวมติกา ฉตฺตุปาหนํ ธาเรติ, อาปตฺติ ปาจิตฺติยสฺส.\csromanspacer{} อคิลานา คิลานสญฺญา ฉตฺตุปาหนํ ธาเรติ, อาปตฺติ ปาจิตฺติยสฺส. \csromanfolio{452}\textbf{ฉตฺตํ ธาเรติ น อุปาหนํ, อาปตฺติ ทุกฺกฏสฺส.}\csromanspacer{}\textbf{ อุปาหนํ ธาเรติ น} \textbf{ฉตฺตํ, อาปตฺติ ทุกฺกฏสฺส.}\csromanspacer{}\textbf{ คิลานา อคิลานสญฺญา, อาปตฺติ ทุกฺกฏสฺส.}\csromanspacer{}\textbf{ คิลานา} เวมติกา, อาปตฺติ ทุกฺกฏสฺส.\csromanspacer{} คิลานา คิลานสญฺญา, อนาปตฺติ.}

\proseitemwithcloserruled{1183}{อนาปตฺติ คิลานาย, อาราเม อารามูปจาเร ธาเรติ, อาปทาสุ อุมฺมตฺติกาย อาทิกมฺมิกายาติ.}{ปฐมสิกฺขาปทํ นิฏฺฐิตํ.}
\csromanmatikaanchor{mtk.242}
\csromantocmarkref{subsection}{2. ทุติยสิกฺขาปท}{mtk.242}
\bookmark[dest={mtk.242},level=2]{2. ทุติยสิกฺขาปท}
\csromanheader{2}{9. \textbf{ฉตฺตุปาหนวคฺค 2. ทุติยสิกฺขาปท}}
\proseitem{1184}{เตน สมเยน พุทฺโธ ภควา สาวตฺถิยํ วิหรติ เชตวเน อนาถปิณฺฑิกสฺส อาราเม.\csromanspacer{} เตน โข ปน สมเยน ฉพฺพคฺคิยา ภิกฺขุนิโย ยาเนน ยายนฺติ.\csromanspacer{} มนุสฺสา อุชฺฌายนฺติ ขิยฺยนฺติ วิปาเจนฺติ “กถํ หิ นาม ภิกฺขุนิโย ยาเนน ยายิสฺสนฺติ เสยฺยาถาปิ คิหินิโย กามโภคินิโย”ติ.\csromanspacer{} อสฺโสสุํ โข ภิกฺขุนิโย เตสํ มนุสฺสานํ อุชฺฌายนฺตานํ ขิยฺยนฺตานํ วิปาเจนฺตานํ.\csromanspacer{} ยา ตา ภิกฺขุนิโย อปฺปิจฺฉา ฯเปฯ ตา อุชฺฌายนฺติ ขิยฺยนฺติ วิปาเจนฺติ “กถํ หิ นาม ฉพฺพคฺคิยา ภิกฺขุนิโย ยาเนน ยายิสฺสนฺตี”ติ ฯเปฯ.\csromanspacer{} สจฺจํ กิร ภิกฺขเว ฉพฺพคฺคิยา ภิกฺขุนิโย ยาเนน ยายนฺตีติ.\csromanspacer{} สจฺจํ ภควาติ.\csromanspacer{} วิครหิ พุทฺโธ ภควา ฯเปฯ.\csromanspacer{} กถํ หิ นาม ภิกฺขเว ฉพฺพคฺคิยา ภิกฺขุนิโย ยาเนน ยายิสฺสนฺติ.\csromanspacer{} เนตํ ภิกฺขเว อปฺปสนฺนานํ วา ปสาทาย ฯเปฯ.\csromanspacer{} เอวญฺจ ปน ภิกฺขเว ภิกฺขุนิโย อิมํ สิกฺขาปทํ อุทฺทิสนฺตุ\csromandash{}}

\prose{\textbf{“ยา ปน ภิกฺขุนี ยาเนน ยาเยยฺย, ปาจิตฺติยนฺ”}ติ.}

\prose{เอวญฺจิทํ ภควตา ภิกฺขุนีนํ สิกฺขาปทํ ปญฺญตฺตํ โหติ.}

\proseitem{1185}{เตน โข ปน สมเยน อญฺญตรา ภิกฺขุนี คิลานา โหติ, น สกฺโกติ ปทสา คนฺตุํ ฯเปฯ.\csromanspacer{} ภควโต เอตมตฺถํ อาโรเจสุํ ฯเปฯ.\csromanspacer{} อนุชานามิ ภิกฺขเว คิลานาย ภิกฺขุนิยา ยานํ.\csromanspacer{} เอวญฺจ ปน ภิกฺขเว ภิกฺขุนิโย อิมํ สิกฺขาปทํ อุทฺทิสนฺตุ\csromandash{}}

\proseitem{1186}{\csromanfolio{453}\textbf{“ยา ปน ภิกฺขุนี อคิลานา ยาเนน ยาเยยฺย, ปาจิตฺติยนฺ”}ติ.}

\proseitem{1187}{\textbf{ยา ปนา}ติ ยา ยาทิสา ฯเปฯ.\csromanspacer{} \textbf{ภิกฺขุนี}ติ ฯเปฯ อยํ อิมสฺมิํ อตฺเถ อธิปฺเปตา ภิกฺขุนีติ.}

\prose{\textbf{อคิลานา} นาม สกฺโกติ ปทสา คนฺตุํ.}

\prose{\textbf{คิลานา} นาม น สกฺโกติ ปทสา คนฺตุํ.}

\prose{\textbf{ยานํ} นาม วยฺหํ รโถ สกฏํ สนฺทมานิกา สิวิกา ปาฏงฺกี.}

\prose{\textbf{ยาเยยฺยา}ติ สกิํปิ ยาเนน ยายติ, อาปตฺติ ปาจิตฺติยสฺส.}

\proseitem{1188}{อคิลานา อคิลานสญฺญา ยาเนน ยายติ, อาปตฺติ ปาจิตฺติยสฺส.\csromanspacer{} อคิลานา เวมติกา ยาเนน ยายติ, อาปตฺติ ปาจิตฺติยสฺส.\csromanspacer{} อคิลานา คิลานสญฺญา ยาเนน ยายติ, อาปตฺติ ปาจิตฺติยสฺส.}

\prose{คิลานา อคิลานสญฺญา, อาปตฺติ ทุกฺกฏสฺส.\csromanspacer{} คิลานา เวมติกา, อาปตฺติ ทุกฺกฏสฺส.\csromanspacer{} คิลานา คิลานสญฺญา, อนาปตฺติ.}

\proseitemwithcloserruled{1189}{อนาปตฺติ คิลานาย อาปทาสุ อุมฺมตฺติกาย อาทิกมฺมิกายาติ.}{ทุติยสิกฺขาปทํ นิฏฺฐิตํ.}
\csromanmatikaanchor{mtk.243}
\csromantocmarkref{subsection}{3. ตติยสิกฺขาปท}{mtk.243}
\bookmark[dest={mtk.243},level=2]{3. ตติยสิกฺขาปท}
\csromanheader{2}{9. \textbf{ฉตฺตุปาหนวคฺค 3. ตติยสิกฺขาปท}}
\proseitem{1190}{เตน สมเยน พุทฺโธ ภควา สาวตฺถิยํ วิหรติ เชตวเน อนาถปิณฺฑิกสฺส อาราเม.\csromanspacer{} เตน โข ปน สมเยน อญฺญตรา ภิกฺขุนี อญฺญตริสฺสา อิตฺถิยา กุลูปิกา โหติ.\csromanspacer{} อถ โข สา อิตฺถี ตํ ภิกฺขุนิํ เอตทโวจ “หนฺทายฺเย อิมํ สงฺฆาณิํ อมุกาย นาม อิตฺถิยา เทหี”ติ.\csromanspacer{} อถ โข สา ภิกฺขุนี “สจาหํ ปตฺเตน อาทาย คจฺฉามิ วิสฺสโร เม ภวิสฺสตี”ติ ปฏิมุญฺจิตฺวา อคมาสิ.\csromanspacer{} ตสฺสา รถิกาย สุตฺตเก ฉินฺเน วิปฺปกิริยิํสุ.\csromanspacer{} มนุสฺสา อุชฺฌายนฺติ ขิยฺยนฺติ วิปาเจนฺติ “กถํ \csromanfolio{454}หิ นาม ภิกฺขุนิโย สงฺฆาณิํ ธาเรสฺสนฺติ เสยฺยถาปิ คิหินิโย กามโภคินิโย”ติ.\csromanspacer{} อสฺโสสุํ โข ภิกฺขุนิโย เตสํ มนุสฺสานํ อุชฺฌายนฺตานํ ขิยฺยนฺตานํ วิปาเจนฺตานํ.\csromanspacer{} ยา ตา ภิกฺขุนิโย อปฺปิจฺฉา \textbf{ฯเปฯ ตา อุชฺฌายนฺติ ขิยฺยนฺติ วิปาเจนฺติ “กถํ หิ นาม ภิกฺขุนี สงฺฆาณิํ} \textbf{ธาเรสฺสตี”ติ ฯเปฯ.}\csromanspacer{}\textbf{ สจฺจํ กิร ภิกฺขเว ภิกฺขุนี สงฺฆาณิํ ธาเรตีติ}\footnote{ธาเรสีติ (ก)}\textbf{.}\csromanspacer{} สจฺจํ ภควาติ.\csromanspacer{} วิครหิ พุทฺโธ ภควา ฯเปฯ.\csromanspacer{} กถํ หิ นาม ภิกฺขเว ภิกฺขุนี สงฺฆาณิํ ธาเรสฺสติ.\csromanspacer{} เนตํ ภิกฺขเว อปฺปสนฺนานํ วา ปสาทาย ฯเปฯ.\csromanspacer{} เอวญฺจ ปน ภิกฺขเว ภิกฺขุนิโย อิมํ สิกฺขาปทํ อุทฺทิสนฺตุ\csromandash{}}

\proseitem{1191}{\textbf{“ยา ปน ภิกฺขุนี สงฺฆาณิํ ธาเรยฺย, ปาจิตฺติยนฺ”}ติ.}

\proseitem{1192}{\textbf{ยา ปนา}ติ ยา ยาทิสา ฯเปฯ.\csromanspacer{} \textbf{ภิกฺขุนี}ติ ฯเปฯ อยํ อิมสฺมิํ อตฺเถ อธิปฺเปตา ภิกฺขุนีติ.}

\prose{\textbf{สงฺฆาณิ} นาม ยา กาจิ กฏูปคา.}

\prose{\textbf{ธาเรยฺยา}ติ สกิํปิ ธาเรติ, อาปตฺติ ปาจิตฺติยสฺส.}

\proseitemwithcloserruled{1193}{อนาปตฺติ อาพาธปฺปจฺจยา, กฏิสุตฺตกํ ธาเรติ, อุมฺมตฺติกาย อาทิกมฺมิกายาติ.}{ตติยสิกฺขาปทํ นิฏฺฐิตํ.}
\csromanmatikaanchor{mtk.244}
\csromantocmarkref{subsection}{4. จตุตฺถสิกฺขาปท}{mtk.244}
\bookmark[dest={mtk.244},level=2]{4. จตุตฺถสิกฺขาปท}
\csromanheader{2}{9. \textbf{ฉตฺตุปาหนวคฺค 4. จตุตฺถสิกฺขาปท}}
\proseitem{1194}{เตน สมเยน พุทฺโธ ภควา สาวตฺถิยํ วิหรติ เชตวเน อนาถปิณฺฑิกสฺส อาราเม.\csromanspacer{} เตน โข ปน สมเยน ฉพฺพคฺคิยา ภิกฺขุนิโย อิตฺถาลงฺการํ ธาเรนฺติ.\csromanspacer{} มนุสฺสา อุชฺฌายนฺติ ขิยฺยนฺติ วิปาเจนฺติ “กถํ หิ นาม ภิกฺขุนิโย อิตฺถาลงฺการํ ธาเรสฺสนฺติ เสยฺยถาปิ คิหินิโย กามโภคินิโย”ติ.\csromanspacer{} อสฺโสสุํ โข ภิกฺขุนิโย เตสํ มนุสฺสานํ อุชฺฌายนฺตานํ ขิยฺยนฺตานํ วิปาเจนฺตานํ.\csromanspacer{} ยา ตา ภิกฺขุนิโย อปฺปิจฺฉา ฯเปฯ ตา อุชฺฌายนฺติ ขิยฺยนฺติ วิปาเจนฺติ “กถํ หิ นาม ฉพฺพคฺคิยา ภิกฺขุนิโย อิตฺถาลงฺการํ ธาเรสฺสนฺตี”ติ ฯเปฯ.\csromanspacer{} สจฺจํ กิร ภิกฺขเว ฉพฺพคฺคิยา ภิกฺขุนิโย อิตฺถาลงฺการํ ธาเรนฺตีติ.\csromanspacer{} สจฺจํ ภควาติ.\csromanspacer{} วิครหิ พุทฺโธ ภควา ฯเปฯ.\csromanspacer{} กถํ หิ นาม ภิกฺขเว ฉพฺพคฺคิยา ภิกฺขุนิโย \csromanfolio{455}อิตฺถาลงฺการํ ธาเรสฺสนฺติ.\csromanspacer{} เนตํ ภิกฺขเว อปฺปสนฺนานํ วา ปสาทาย ฯเปฯ.\csromanspacer{} เอวญฺจ ปน ภิกฺขเว ภิกฺขุนิโย อิมํ สิกฺขาปทํ อุทฺทิสนฺตุ\csromandash{}}

\proseitem{1195}{\textbf{“ยา ปน ภิกฺขุนี อิตฺถาลงฺการํ ธาเรยฺย, ปาจิตฺติยนฺ}”ติ.}

\proseitem{1196}{\textbf{ยา ปนา}ติ ยา ยาทิสา ฯเปฯ.\csromanspacer{} \textbf{ภิกฺขุนี}ติ ฯเปฯ อยํ อิมสฺมิํ อตฺเถ อธิปฺเปตา ภิกฺขุนีติ.}

\prose{\textbf{อิตฺถาลงฺกาโร} นาม สีสูปโค คีวูปโค หตฺถูปโค ปาทูปโค กฏูปโค.\csromanspacer{} \textbf{ธาเรยฺยา}ติ สกิํปิ ธาเรติ, อาปตฺติ ปาจิตฺติยสฺส.}

\proseitemwithcloserruled{1197}{อนาปตฺติ อาพาธปฺปจฺจยา, อุมฺมตฺติกาย อาทิกมฺมิกายาติ.}{จตุตฺถสิกฺขาปทํ นิฏฺฐิตํ.}
\csromanmatikaanchor{mtk.245}
\csromantocmarkref{subsection}{5. ปญฺจมสิกฺขาปท}{mtk.245}
\bookmark[dest={mtk.245},level=2]{5. ปญฺจมสิกฺขาปท}
\csromanheader{2}{9. \textbf{ฉตฺตุปาหนวคฺค 5. ปญฺจมสิกฺขาปท}}
\proseitem{1198}{เตน สมเยน พุทฺโธ ภควา สาวตฺถิยํ วิหรติ เชตวเน อนาถปิณฺฑิกสฺส อาราเม.\csromanspacer{} เตน โข ปน สมเยน ฉพฺพคฺคิยา ภิกฺขุนิโย คนฺธวณฺณเกน นหายนฺติ.\csromanspacer{} มนุสฺสา อุชฺฌายนฺติ ขิยฺยนฺติ วิปาเจนฺติ “กถํ หิ นาม ภิกฺขุนิโย คนฺธวณฺณเกน นหายิสฺสนฺติ เสยฺยถาปิ คิหินิโย กามโภคินิโย”ติ.\csromanspacer{} อสฺโสสุํ โข ภิกฺขุนิโย เตสํ มนุสฺสานํ อุชฺฌายนฺตานํ ขิยฺยนฺตานํ วิปาเจนฺตานํ.\csromanspacer{} ยา ตา ภิกฺขุนิโย อปฺปิจฺฉา ฯเปฯ ตา อุชฺฌายนฺติ ขิยฺยนฺติ วิปาเจนฺติ “กถํ หิ นาม ฉพฺพคฺคิยา ภิกฺขุนิโย คนฺธวณฺณเกน นหายิสฺสนฺตี”ติ ฯเปฯ.\csromanspacer{} สจฺจํ กิร ภิกฺขเว ฉพฺพคฺคิยา ภิกฺขุนิโย คนฺธวณฺณเกน นหายนฺตีติ.\csromanspacer{} สจฺจํ ภควาติ.\csromanspacer{} วิครหิ พุทฺโธ ภควา ฯเปฯ.\csromanspacer{} กถํ หิ นาม ภิกฺขเว ฉพฺพคฺคิยา ภิกฺขุนิโย คนฺธวณฺณเกน นหายิสฺสนฺติ.\csromanspacer{} เนตํ ภิกฺขเว อปฺปสนฺนานํ วา ปสาทาย ฯเปฯ.\csromanspacer{} เอวญฺจ ปน ภิกฺขเว ภิกฺขุนิโย อิมํ สิกฺขาปทํ อุทฺทิสนฺตุ\csromandash{}}

\proseitem{1199}{\textbf{“ยา ปน ภิกฺขุนี คนฺธวณฺณเกน นหาเยยฺย, ปาจิตฺติยนฺ”}ติ.}

\proseitem{1200}{\textbf{ยา ปนา}ติ ยา ยาทิสา ฯเปฯ.\csromanspacer{} \textbf{ภิกฺขุนี}ติ ฯเปฯ อยํ อิมสิํ อตฺเถ อธิปฺเปตา ภิกฺขุนีติ.}

\prose{\csromanfolio{456}\textbf{คนฺโธ} นาม โย โกจิ คนฺโธ.\csromanspacer{} \textbf{วณฺณกํ} นาม ยํ กิญฺจิ วณฺณกํ.\csromanspacer{} \textbf{นหาเยยฺยา}ติ นหายติ, ปโยเค ทุกฺกฏํ.\csromanspacer{} นหานปริโยสาเน อาปตฺติ}

\proseitemwithcloserruled{1201}{อนาปตฺติ อาพาธปฺปจฺจยา, อุมฺมตฺติกาย อาทิกมฺมิกายาติ.}{ปญฺจมสิกฺขาปทํ นิฏฺฐิตํ.}
\csromanmatikaanchor{mtk.246}
\csromantocmarkref{subsection}{6. ฉฏฺฐสิกฺขาปท}{mtk.246}
\bookmark[dest={mtk.246},level=2]{6. ฉฏฺฐสิกฺขาปท}
\csromanheader{2}{9. \textbf{ฉตฺตุปาหนวคฺค 6. ฉฏฺฐสิกฺขาปท}}
\proseitem{1202}{เตน สมเยน พุทฺโธ ภควา สาวตฺถิยํ วิหรติ เชตวเน อนาถปิณฺฑิกสฺส อาราเม.\csromanspacer{} เตน โข ปน สมเยน ฉพฺพคฺคิยา ภิกฺขุนิโย วาสิตเกน ปิญฺญาเกน นหายนฺติ.\csromanspacer{} มนุสฺสา อุชฺฌายนฺติ ขิยฺยนฺติ วิปาเจนฺติ “กถํ หิ นาม ภิกฺขุนิโย วาสิตเกน ปิญฺญาเกน นหายิสฺสนฺติ เสยฺยถาปิ คิหินิโย กามโภคินิโย”ติ.\csromanspacer{} อสฺโสสุํ โข ภิกฺขุนิโย เตสํ มนุสฺสานํ อุชฺฌายนฺตานํ ขิยฺยนฺตานํ วิปาเจนฺตานํ.\csromanspacer{} ยา ตา ภิกฺขุนิโย อปฺปิจฺฉา ฯเปฯ ตา อุชฺฌายนฺติ ขิยฺยนฺติ วิปาเจนฺติ “กถํ หิ นาม ฉพฺพคฺคิยา ภิกฺขุนิโย วาสิตเกน ปิญฺญาเกน นหายิสฺสนฺตี”ติ ฯเปฯ.\csromanspacer{} สจฺจํ กิร ภิกฺขเว ฉพฺพคฺคิยา ภิกฺขุนิโย วาสิตเกน ปิญฺญาเกน นหายนฺตีติ.\csromanspacer{} สจฺจํ ภควาติ.\csromanspacer{} วิครหิ พุทฺโธ ภควา ฯเปฯ.\csromanspacer{} กถํ หิ นาม ภิกฺขเว ฉพฺพคฺคิยา ภิกฺขุนิโย วาสิตเกน ปิญฺญาเกน นหายิสฺสนฺติ.\csromanspacer{} เนตํ ภิกฺขเว อปฺปสนฺนานํ วา ปสาทาย ฯเปฯ.\csromanspacer{} เอวญฺจ ปน ภิกฺขเว ภิกฺขุนิโย อิมํ สิกฺขาปทํ อุทฺทิสนฺตุ\csromandash{}}

\proseitem{1203}{\textbf{“ยา ปน ภิกฺขุนี วา สิตเกน ปิญฺญาเกน นหาเยยฺย, ปาจิตฺติยนฺ”}ติ.}

\proseitem{1204}{\textbf{ยา ปนา}ติ ยา ยาทิสา ฯเปฯ.\csromanspacer{} \textbf{ภิกฺขุนี}ติ ฯเปฯ อยํ อิมสฺมิํ อตฺเถ อธิปฺเปตา ภิกฺขุนีติ.}

\prose{\textbf{วาสิตกํ} นาม ยํ กิญฺจิ คนฺธวาสิตกํ.\csromanspacer{} \textbf{ปิญฺญากํ} นาม ติลปิฏฺฐํ วุจฺจติ.\csromanspacer{} \textbf{นหาเยยฺยา}ติ นหายติ, ปโยเค ทุกฺกฏํ.\csromanspacer{} นหานปริโยสาเน อาปตฺติ ปาจิตฺติยสฺส.}

\proseitemwithcloserruled{1205}{\csromanfolio{457}อนาปตฺติ อาพาธปฺปจฺจยา, ปกติปิญฺญาเกน นหายติ, อุมฺมตฺติกาย อาทิกมฺมิกายาติ.}{ฉฏฺฐสิกฺขาปทํ นิฏฺฐิตํ.}
\csromanmatikaanchor{mtk.247}
\csromantocmarkref{subsection}{7. สตฺตมสิกฺขาปท}{mtk.247}
\bookmark[dest={mtk.247},level=2]{7. สตฺตมสิกฺขาปท}
\csromanheader{2}{9. \textbf{ฉตฺตุปาหนวคฺค 7. สตฺตมสิกฺขาปท}}
\proseitem{1206}{เตน สมเยน พุทฺโธ ภควา สาวตฺถิยํ วิหรติ เชตวเน อนาถปิณฺฑิกสฺส อาราเม.\csromanspacer{} เตน โข ปน สมเยน ภิกฺขุนิโย ภิกฺขุนิยา อุมฺมทฺทาเปนฺติปิ ปริมทฺทาเปนฺติปิ.\csromanspacer{} มนุสฺสา วิหารจาริกํ อาหิณฺฑนฺตา ปสฺสิตฺวา อุชฺฌายนฺติ ขิยฺยนฺติ วิปาเจนฺติ “กถํ หิ นาม ภิกฺขุนิโย ภิกฺขุนิยา อุมฺมทฺทาเปสฺสนฺติปิ ปริมทฺทาเปสฺสนฺติปิ เสยฺยถาปิ คิหินิโย กามโภคินิโย”ติ.\csromanspacer{} อสฺโสสุํ โข ภิกฺขุนิโย เตสํ มนุสฺสานํ อุชฺฌายนฺตานํ ขิยฺยนฺตานํ วิปาเจนฺตานํ.\csromanspacer{} ยา ตา ภิกฺขุนิโย อปฺปิจฺฉา ฯเปฯ ตา อุชฺฌายนฺติ ขิยฺยนฺติ วิปาเจนฺติ “กถํ หิ นาม ภิกฺขุนิโย ภิกฺขุนิยา อุมฺมทฺทาเปสฺสนฺติปิ ปริมทฺทาเปสฺสนฺติปี”ติ ฯเปฯ.\csromanspacer{} สจฺจํ กิร ภิกฺขเว ภิกฺขุนิโย ภิกฺขุนิยา อุมฺมทฺทาเปนฺติปิ ปริมทฺทาเปนฺติปีติ.\csromanspacer{} สจฺจํ ภควาติ.\csromanspacer{} วิครหิ พุทฺโธ ภควา ฯเปฯ.\csromanspacer{} กถํ หิ นาม ภิกฺขเว ภิกฺขุนิโย ภิกฺขุนิยา อุมฺมทฺทาเปสฺสนฺติปิ ปริมทฺทาเปสฺสนฺติปิ.\csromanspacer{} เนตํ ภิกฺขเว อปฺปสนฺนานํ วา ปสาทาย ฯเปฯ.\csromanspacer{} เอวญฺจ ปน ภิกฺขเว ภิกฺขุนิโย อิมํ สิกฺขาปทํ อุทฺทิสนฺตุ\csromandash{}}

\proseitem{1207}{\textbf{“ยา ปน ภิกฺขุนี ภิกฺขุนิยา อุมฺมทฺทาเปยฺย วา ปริมทฺทาเปยฺย วา, ปาจิตฺติยนฺ”}ติ.}

\proseitem{1208}{\textbf{ยา ปนา}ติ ยา ยาทิสา ฯเปฯ.\csromanspacer{} \textbf{ภิกฺขุนี}ติ ฯเปฯ อยํ อิมสฺมิํ อตฺเถ อธิปฺเปตา ภิกฺขุนีติ.}

\prose{\textbf{ภิกฺขุนิยา}ติ อญฺญาย ภิกฺขุนิยา.\csromanspacer{} \textbf{อุมฺมทฺทาเปยฺย} วาติ อุมฺมทฺทาเปติ\footnote{อุพฺพฏฺฏาเปติ (สี)}, อาปตฺติ ปาจิตฺติยสฺส.\csromanspacer{} ปริมทฺทาเปยฺย \textbf{วา}ติ สมฺพาหาเปติ, อาปตฺติ}

\proseitemwithcloserruled{1209}{\csromanfolio{458}อนาปตฺติ คิลานาย อาปทาสุ อุมฺมตฺติกาย อาทิกมฺมิกายาติ.}{สตฺตมสิกฺขาปทํ นิฏฺฐิตํ.}
\csromanmatikaanchor{mtk.248}
\csromantocmarkref{subsubsection}{8-9-10. อฏฺฐมนวมทสมสิกฺขาปท}{mtk.248}
\bookmark[dest={mtk.248},level=3]{8-9-10. อฏฺฐมนวมทสมสิกฺขาปท}
\csromanheader{3}{9. \textbf{ฉตฺตุปาหนวคฺค 8-9-}10. อฏฺฐม นวม ทสมสิกฺขาปท}
\proseitem{1210}{เตน สมเยน พุทฺโธ ภควา สาวตฺถิยํ วิหรติ เชตวเน อนาถปิณฺฑิกสฺส อาราเม.\csromanspacer{} เตน โข ปน สมเยน ภิกฺขุนิโย สิกฺขมานาย อุมฺมทฺทาเปนฺติปิ ปริมทฺทาเปนฺติปิ ฯเปฯ สามเณริยา อุมฺมทฺทาเปนฺติปิ ปริมทฺทาเปนฺติปิ ฯเปฯ คิหินิยา อุมฺมทฺทาเปนฺติปิ ปริมทฺทาเปนฺติปิ.\csromanspacer{} มนุสฺสา วิหารจาริกํ อาหิณฺฑนฺตา ปสฺสิตฺวา อุชฺฌายนฺติ ขิยฺยนฺติ วิปาเจนฺติ “กถํ หิ นาม ภิกฺขุนิโย คิหินิยา อุมฺมทฺทาเปสฺสนฺติปิ ปริมทฺทาเปสฺสนฺติปิ เสยฺยถาปิ คิหินิโย กามโภคินิโย”ติ.\csromanspacer{} อสฺโสสุํ โข ภิกฺขุนิโย เตสํ มนุสฺสานํ อุชฺฌายนฺตานํ ขิยฺยนฺตานํ วิปาเจนฺตานํ.\csromanspacer{} ยา ตา ภิกฺขุนิโย อปฺปิจฺฉา ฯเปฯ ตา อุชฺฌายนฺติ ขิยฺยนฺติ วิปาเจนฺติ “กถํ หิ นาม ภิกฺขุนิโย คิหินิยา อุมฺมทฺทาเปสฺสนฺติปิ ปริมทฺทาเปสฺสนฺติปี”ติ ฯเปฯ.\csromanspacer{} สจฺจํ กิร ภิกฺขเว ภิกฺขุนิโย คิหินิยา อุมฺมทฺทาเปนฺติปิ ปริมทฺทาเปนฺติปีติ.\csromanspacer{} สจฺจํ ภควาติ.\csromanspacer{} วิครหิ พุทฺโธ ภควา ฯเปฯ.\csromanspacer{} กถํ หิ นาม ภิกฺขเว ภิกฺขุนิโย คิหินิยา อุมฺมทฺทาเปสฺสนฺติปิ ปริมทฺทาเปสฺสนฺติปิ.\csromanspacer{} เนตํ ภิกฺขเว อปฺปสนฺนานํ วา ปสาทาย ฯเปฯ.\csromanspacer{} เอวญฺจ ปน ภิกฺขเว ภิกฺขุนิโย อิมํ สิกฺขาปทํ อุทฺทิสนฺตุ\csromandash{}}

\proseitem{1211}{\textbf{“ยา ปน ภิกฺขุนี (สิกฺขมานาย ฯเปฯ สามเณริยา ฯเปฯ) คิหินิยา อุมฺมทฺทาเปยฺย วา ปริมทฺทาเปยฺย วา, ปาจิตฺติยนฺ”}ติ.}

\proseitem{1212}{\textbf{ยา ปนา}ติ ยา ยาทิสา ฯเปฯ.\csromanspacer{} \textbf{ภิกฺขุนี}ติ ฯเปฯ อยํ อิมสฺมิํ อตฺเถ อธิปฺเปตา ภิกฺขุนีติ.}

\prose{\textbf{สิกฺขมานา} นาม เทฺว วสฺสานิ ฉสุ ธมฺเมสุ สิกฺขิตสิกฺขา.}

\prose{\textbf{สามเณรี} นาม ทสสิกฺขาปทิกา.}

\prose{\textbf{คิหินี} นาม อคารินี วุจฺจติ.}

\prose{\csromanfolio{459}\textbf{อุมฺมทฺทาเปยฺย วา}ติ อุมฺมทฺทาเปติ, อาปตฺติ ปาจิตฺติยสฺส.}

\prose{\textbf{ปริมทฺทาเปยฺย} \textbf{วา}ติ สมฺพาหาเปติ, อาปตฺติ ปาจิตฺติยสฺส.}

\proseitemwithcloserruled{1213}{อนาปตฺติ คิลานาย\footnote{อาพาธปฺปจฺจยา (ก)} อาปทาสุ อุมฺมตฺติกาย อาทิกมฺมิกายาติ.}{อฏฺฐม นวม ทสมสิกฺขาปทานิ นิฏฺฐิตานิ.}
\csromanmatikaanchor{mtk.249}
\csromantocmarkref{subsection}{11. เอกาทสมสิกฺขาปท}{mtk.249}
\bookmark[dest={mtk.249},level=2]{11. เอกาทสมสิกฺขาปท}
\csromanheader{2}{9. \textbf{ฉตฺตุปาหนวคฺคฺ}อ 11. เอกาทสมสิกฺขาปท}
\proseitem{1214}{เตน สมเยน พุทฺโธ ภควา สาวตฺถิยํ วิหรติ เชตวเน อนาถปิณฺฑิกสฺส อาราเม.\csromanspacer{} เตน โข ปน สมเยน ภิกฺขุนิโย ภิกฺขุสฺส ปุรโต อนาปุจฺฉา อาสเน นิสีทนฺติ.\csromanspacer{} ภิกฺขู อุชฺฌายนฺติ ขิยฺยนฺติ วิปาเจนฺติ “กถํ หิ นาม ภิกฺขุนิโย ภิกฺขุสฺส ปุรโต อนาปุจฺฉา อาสเน นิสีทิสฺสนฺตี”ติ ฯเปฯ.\csromanspacer{} สจฺจํ กิร ภิกฺขเว ภิกฺขุนิโย ภิกฺขุสฺส ปุรโต อนาปุจฺฉา อาสเน นิสีทนฺตีติ.\csromanspacer{} สจฺจํ ภควาติ.\csromanspacer{} วิครหิ พุทฺโธ ภควา ฯเปฯ.\csromanspacer{} กถํ หิ นาม ภิกฺขเว ภิกฺขุนิโย ภิกฺขุสฺส ปุรโต อนาปุจฺฉา อาสเน นิสีทิสฺสนฺติ.\csromanspacer{} เนตํ ภิกฺขเว อปฺปสนฺนานํ วา ปสาทาย ฯเปฯ.\csromanspacer{} เอวญฺจ ปน ภิกฺขเว ภิกฺขุนิโย อิมํ สิกฺขาปทํ อุทฺทิสนฺตุ\csromandash{}}

\proseitem{1215}{\textbf{“ยา ปน ภิกฺขุนี ภิกฺขุสฺส ปุรโต อนาปุจฺฉา อาสเน นิสีเทยฺย, ปาจิตฺติยนฺ”}ติ.}

\proseitem{1216}{\textbf{ยา ปนา}ติ ยา ยาทิสา ฯเปฯ.\csromanspacer{} \textbf{ภิกฺขุนีตฺ}อิ ฯเปฯ อยํ อิมสฺมิํ อตฺเถ อธิปฺเปตา ภิกฺขุนีติ.}

\prose{\textbf{ภิกฺขุสฺส ปุรโต}ติ อุปสมฺปนฺนสฺส ปุรโต.\csromanspacer{} \textbf{อนาปุจฺฉา}ติ อนปโลเกตฺวา.\csromanspacer{} \textbf{อาสเน นิสีเทยฺยา}ติ อนฺตมโส ฉมายปิ นิสีทติ, อาปตฺติ ปาจิตฺติยสฺส.}

\proseitem{1217}{อนาปุจฺฉิเต อนาปุจฺฉิตสญฺญา อาสเน นิสีทติ, อาปตฺติ ปาจิตฺติยสฺส.\csromanspacer{} อนาปุจฺฉิเต เวมติกา อาสเน นิสีทติ, อาปตฺติ \csromanfolio{460}ปาจิตฺติยสฺส.\csromanspacer{} อนาปุจฺฉิเต อาปุจฺฉิตสญฺญา อาสเน นิสีทติ, อาปตฺติ ปาจิตฺติยสฺส.}

\prose{อาปุจฺฉิเต อนาปุจฺฉิตสญฺญา, อาปตฺติ ทุกฺกฏสฺส.\csromanspacer{} อาปุจฺฉิเต เวมติกา, อาปตฺติ ทุกฺกฏสฺส.\csromanspacer{} อาปุจฺฉิเต อาปุจฺฉิตสญฺญา, อนาปตฺติ.}

\proseitemwithcloserruled{1218}{อนาปตฺติ อาปุจฺฉา อาสเน นิสีทติ, คิลานาย อาปทาสุ อุมฺมตฺติกาย อาทิกมฺมิกายาติ.}{เอกาทสมสิกฺขาปทํ นิฏฺฐิตํ.}
\csromanmatikaanchor{mtk.250}
\csromantocmarkref{subsection}{12. ทฺวาทสมสิกฺขาปท}{mtk.250}
\bookmark[dest={mtk.250},level=2]{12. ทฺวาทสมสิกฺขาปท}
\csromanheader{2}{9. \textbf{ฉตฺตุปาหนวคฺคฺ}อ 12. ทฺวาทสมสิกฺขาปท}
\proseitem{1219}{เตน สมเยน พุทฺโธ ภควา สาวตฺถิยํ วิหรติ เชตวเน อนาถปิณฺฑิกสฺส อาราเม.\csromanspacer{} เตน โข ปน สมเยน ภิกฺขุนิโย อโนกาสกตํ ภิกฺขุํ ปญฺหํ ปุจฺฉนฺติ.\csromanspacer{} ภิกฺขู อุชฺฌายนฺติ ขิยฺยนฺติ วิปาเจนฺติ “กถํ หิ นาม ภิกฺขุนิโย อโนกาสกตํ ภิกฺขุํ ปญฺหํ ปุจฺฉิสฺสนฺตี”ติ ฯเปฯ.\csromanspacer{} สจฺจํ กิร ภิกฺขเว ภิกฺขุนิโย อโนกาสกตํ ภิกฺขุํ ปญฺหํ ปุจฺฉนฺตีติ.\csromanspacer{} สจฺจํ ภควาติ.\csromanspacer{} วิครหิ พุทฺโธ ภควา ฯเปฯ.\csromanspacer{} กถํ หิ นาม ภิกฺขเว ภิกฺขุนิโย อโนกาสกตํ ภิกฺขุํ ปญฺหํ ปุจฺฉิสฺสนฺติ.\csromanspacer{} เนตํ ภิกฺขเว อปฺปสนฺนานํ วา ปสาทาย ฯเปฯ.\csromanspacer{} เอวญฺจ ปน ภิกฺขเว ภิกฺขุนิโย อิมํ สิกฺขาปทํ อุทฺทิสนฺตุ\csromandash{}}

\proseitem{1220}{\textbf{“โย ปน ภิกฺขุนี อโนกาสกตํ ภิกฺขุํ ปญฺหํ ปุจฺเฉยฺย, ปาจิตฺติยนฺ”}ติ.}

\proseitem{1221}{\textbf{ยา ปนา}ติ ยา ยาทิสา ฯเปฯ.\csromanspacer{} \textbf{ภิกฺขุนี}ติ ฯเปฯ อยํ อิมสฺมิํ อตฺเถ อธิปฺเปตา ภิกฺขุนีติ.}

\prose{\textbf{อโนกาสกตนฺ}ติ อนาปุจฺฉา.\csromanspacer{} \textbf{ภิกฺขุ}นฺติ อุปสมฺปนฺนํ.\csromanspacer{} \textbf{ปญฺหํ ปุจฺเฉยฺยา}ติ สุตฺตนฺเต โอกาสํ การาเปตฺวา วินยํ วา อภิธมฺมํ วา ปุจฺฉติ, อาปตฺติ ปาจิตฺติยสฺส.\csromanspacer{} วินเย โอกาสํ การาเปตฺวา สุตฺตนฺตํ วา อภิธมฺมํ วา ปุจฺฉติ, อาปตฺติ ปาจิตฺติยสฺส.\csromanspacer{} อภิธมฺเม โอกาสํ การาเปตฺวา สุตฺตนฺตํ วา วินยํ วา ปุจฺฉติ, อาปตฺติ ปาจิตฺติยสฺส.}

\proseitem{1222}{\csromanfolio{461}อนาปุจฺฉิเต อนาปุจฺฉิตสญฺญา ปญฺหํ ปุจฺฉติ, อาปตฺติ ปาจิตฺติยสฺส.\csromanspacer{} อนาปุจฺฉิเต เวมติกา ปญฺหํ ปุจฺฉติ, อาปตฺติ ปาจิตฺติยสฺส.\csromanspacer{} อนาปุจฺฉิเต อาปุจฺฉิตสญฺญา ปญฺหํ ปุจฺฉติ, อาปตฺติ ปาจิตฺติยสฺส.}

\prose{อาปุจฺฉิเต อนาปุจฺฉิตสญฺญา, อาปตฺติ ทุกฺกฏสฺส.\csromanspacer{} อาปุจฺฉิเต เวมติกา, อาปตฺติ ทุกฺกฏสฺส.\csromanspacer{} อาปุจฺฉิเต อาปุจฺฉิตสญฺญา, อนาปตฺติ.}

\proseitemwithcloserruled{1223}{อนาปตฺติ โอกาสํ การาเปตฺวา ปุจฺฉติ, อโนทิสฺส โอกาสํ การาเปตฺวา ยตฺถ กตฺถจิ ปุจฺฉติ, อุมฺมตฺติกาย อาทิกมฺมิกายาติ.}{ทฺวาทสมสิกฺขาปทํ นิฏฺฐิตํ.}
\csromanmatikaanchor{mtk.251}
\csromantocmarkref{subsection}{13. เตรสมสิกฺขาปท}{mtk.251}
\bookmark[dest={mtk.251},level=2]{13. เตรสมสิกฺขาปท}
\csromanheader{2}{9. \textbf{ฉตฺตุปาหนวคฺ}ค 13. เตรสมสิกฺขาปท}
\proseitem{1224}{เตน สมเยน พุทฺโธ ภควา สาวตฺถิยํ วิหรติ เชตวเน อนาถปิณฺฑิกสฺส อาราเม.\csromanspacer{} เตน โข ปน สมเยน อญฺญตรา ภิกฺขุนี อสํกจฺจิกา\footnote{อสงฺกจฺฉิกา (สฺยา)} คามํ ปิณฺฑาย ปาวิสิ, ตสฺสา รถิกาย วาตมณฺฑลิกา สงฺฆาฏิโย อุกฺขิปิํสุ.\csromanspacer{} มนุสฺสา อุกฺกุฏฺฐิํ อกํสุ “สุนฺทรา อยฺยาย ถนุทรา”ติ.\csromanspacer{} สา ภิกฺขุนี เตหิ มนุสฺเสหิ อุปฺปณฺฑิยมานา มงฺกุ อโหสิ.\csromanspacer{} อถ โข สา ภิกฺขุนี อุปสฺสยํ คนฺตฺวา ภิกฺขุนีนํ เอตมตฺถํ อาโรเจสิ.\csromanspacer{} ยา ตา ภิกฺขุนิโย อปฺปิจฺฉา ฯเปฯ ตา อุชฺฌายนฺติ ขิยฺยนฺติ วิปาเจนฺติ “กถํ หิ นาม ภิกฺขุนี อสํกจฺจิกา คามํ ปวิสิสฺสตี”ติ ฯเปฯ.\csromanspacer{} สจฺจํ กิร ภิกฺขเว ภิกฺขุนี อสํกจฺจิกา คามํ ปาวิสีติ.\csromanspacer{} สจฺจํ ภควาติ.\csromanspacer{} วิครหิ พุทฺโธ ภควา ฯเปฯ.\csromanspacer{} กถํ หิ นาม ภิกฺขเว ภิกฺขุนี อสํกจฺจิกา คามํ ปวิสิสฺสติ.\csromanspacer{} เนตํ ภิกฺขเว อปฺปสนฺนานํ วา ปสาทาย ฯเปฯ.\csromanspacer{} เอวญฺจ ปน ภิกฺขเว ภิกฺขุนิโย อิมํ สิกฺขาปทํ อุทฺทิสนฺตุ\csromandash{}}

\proseitem{1225}{\textbf{“ยา ปน ภิกฺขุนี อสํกจฺจิกา}1 \textbf{คามํ ปวิเสยฺย, ปาจิตฺติยนฺ”}ติ.}

\proseitem{1226}{\csromanfolio{462}\textbf{ยา ปนา}ติ ยา ยาทิสา ฯเปฯ.\csromanspacer{} \textbf{ภิกฺขุนี}ติ ฯเปฯ อยํ อิมสฺมิํ อตฺเถ อธิปฺเปตา ภิกฺขุนีติ.}

\prose{\textbf{อสํกจฺจิกา}ติ วินา สํกจฺจิกํ.}

\prose{\textbf{สํกจฺจิกํ} นาม อธกฺขกํ อุพฺภนาภิ, ตสฺส ปฏิจฺฉาทนตฺถาย.}

\prose{\textbf{คามํ ปวิเสยฺยา}ติ ปริกฺขิตฺตสฺส คามสฺส ปริกฺเขปํ อติกฺกาเมนฺติยา อาปตฺติ ปาจิตฺติยสฺส.\csromanspacer{} อปริกฺขิตฺตสฺส คามสฺส อุปจารํ โอกฺกมนฺติยา อาปตฺติ}

\proseitemwithcloser{1227}{อนาปตฺติ อจฺฉินฺนจีวริกาย นฏฺฐจีวริกาย คิลานาย อสฺสติยา อชานนฺติยา อาปทาสุ อุมฺมตฺติกาย อาทิกมฺมิกายาติ.}{เตรสมสิกฺขาปทํ นิฏฺฐิตํ.}
\vspace{\tipitakaparskip}
\csromancenter{ฉตฺตุปาหนวคฺโค นวโม.}
\prosewithcloser{อุทฺทิฏฺฐา โข อยฺยาโย ฉสฏฺฐิสตา ปาจิตฺติยา ธมฺมา, ตตฺถายฺยาโย ปุจฺฉามิ กจฺจิตฺถ ปริสุทฺธา, ทุติยมฺปิ ปุจฺฉามิ กจฺจิตฺถ ปริสุทฺธา, ตติยมฺปิ ปุจฺฉามิ กจฺจิตฺถ ปริสุทฺธา, ปริสุทฺเธตฺถายฺยาโย, ตสฺมา ตุณฺหี, เอวเมตํ ธารยามีติ.}{ขุทฺทกํ สมตฺตํ.}
\csromanheader{2}{\textbf{ภิกฺขุนิวิภงฺเค}}
\nitthitam{\textbf{ปาจิตฺติยกณฺฑํ นิฏฺฐิตํ.}}
\chapterheadpageread{5. \textbf{ปาฏิเทสนียกณฺฑ}}
\csromanmatikaanchor{mtk.252}
\csromanmatikaanchor{mtk.253}
\csromantocmarknopage{chapter}{5. ปาฏิเทสนียกณฺฑ}{mtk.252}
\bookmark[dest={mtk.252},level=0]{5. ปาฏิเทสนียกณฺฑ}
\csromantocmarkref{section}{1. ปฐมปาฏิเทสนียสิกฺขาปท}{mtk.253}
\bookmark[dest={mtk.253},level=1]{1. ปฐมปาฏิเทสนียสิกฺขาปท}
\setcsromanheadcha{5. ปาฏิเทสนียกณฺฑ}
\setcsromanheadhclear{}
\setcsromanheadh{1. ปฐมปาฏิเทสนียสิกฺขาปท}
\csromanheader{1}{1. ปฐม\-ปาฏิเทสนีย\-สิกฺขาปท}
\prose{\csromanfolio{463}อิเม โข ปนายฺยาโย อฏฺฐ ปาฏิเทสนียา ธมฺมา อุทฺเทสํ อาคจฺฉนฺติ.}

\proseitem{1228}{เตน สมเยน พุทฺโธ ภควา สาวตฺถิยํ วิหรติ เชตวเน อนาถปิณฺฑิกสฺส อาราเม.\csromanspacer{} เตน โข ปน สมเยน ฉพฺพคฺคิยา ภิกฺขุนิโย สปฺปิํ วิญฺญาเปตฺวา ภุญฺชนฺติ.\csromanspacer{} มนุสฺสา อุชฺฌายนฺติ ขิยฺยนฺติ วิปาเจนฺติ “กถํ หิ นาม ภิกฺขุนิโย สปฺปิํ วิญฺญาเปตฺวา ภุญฺชิสฺสนฺติ, กสฺส สมฺปนฺนํ น มนาปํ, กสฺส สาทุํ น รุจฺจตี”ติ.\csromanspacer{} อสฺโสสุํ โข ภิกฺขุนิโย เตสํ มนุสฺสานํ อุชฺฌายนฺตานํ ขิยฺยนฺตานํ วิปาเจนฺตานํ.\csromanspacer{} ยา ตา ภิกฺขุนิโย อปฺปิจฺฉา ฯเปฯ ตา อุชฺฌายนฺติ ขิยฺยนฺติ วิปาเจนฺติ “กถํ หิ นาม ฉพฺพคฺคิยา ภิกฺขุนิโย สปฺปิํ วิญฺญาเปตฺวา ภุญฺชิสฺสนฺตี”ติ ฯเปฯ.\csromanspacer{} สจฺจํ กิร ภิกฺขเว ฉพฺพคฺคิยา ภิกฺขุนิโย สปฺปิํ วิญฺญาเปตฺวา ภุญฺชนฺตีติ.\csromanspacer{} สจฺจํ ภควาติ.\csromanspacer{} วิครหิ พุทฺโธ ภควา ฯเปฯ.\csromanspacer{} กถํ หิ นาม ภิกฺขเว ฉพฺพคฺคิยา ภิกฺขุนิโย สปฺปิํ วิญฺญาเปตฺวา ภุญฺชิสฺสนฺติ.\csromanspacer{} เนตํ ภิกฺขเว อปฺปสนฺนานํ วา ปสาทาย ฯเปฯ.\csromanspacer{} เอวญฺจ ปน ภิกฺขเว ภิกฺขุนิโย อิมํ สิกฺขาปทํ อุทฺทิสนฺตุ\csromandash{}}

\prose{\textbf{“ยา ปน ภิกฺขุนี สปฺปิํ วิญฺญาเปตฺวา ภุญฺเชยฺย, ปฏิเทเสตพฺพํ ตาย ภิกฺขุนิยา คารยฺหํ อยฺเย ธมฺมํ อาปชฺชิํ อสปฺปายํ ปาฏิเทสนียํ ตํ ปฏิเทเสมี”}ติ.}

\prose{เอวญฺจิทํ ภควตา ภิกฺขุนีนํ สิกฺขาปทํ ปญฺญตฺตํ โหติ.}

\proseitem{1229}{เตน โข ปน สมเยน ภิกฺขุนิโย คิลานา โหนฺติ.\csromanspacer{} คิลานปุจฺฉิกา ภิกฺขุนิโย คิลานา ภิกฺขุนิโย เอตทโวจุํ “กจฺจิ อยฺเย ขมนียํ, กจฺจิ ยาปนียนฺ”ติ.\csromanspacer{} ปุพฺเพ มยํ อยฺเย สปฺปิํ วิญฺญาเปตฺวา \csromanfolio{464}ภุญฺชาม, เตน โน ผาสุ โหติ, อิทานิ ปน “ภควตา ปฏิกฺขิตฺตนฺ”ติ กุกฺกุจฺจายนฺตา น วิญฺญาเปม, เตน โน น ผาสุ โหตีติ ฯเปฯ.\csromanspacer{} ภควโต เอตมตฺถํ อาโรเจสุํ ฯเปฯ.\csromanspacer{} อนุชานามิ ภิกฺขเว คิลานาย ภิกฺขุนิยา สปฺปิํ วิญฺญาเปตฺวา ภุญฺชิตุํ.\csromanspacer{} เอวญฺจ ปน ภิกฺขเว ภิกฺขุนิโย อิมํ สิกฺขาปทํ อุทฺทิสนฺตุ\csromandash{}}

\proseitem{1230}{\textbf{“ยา ปน ภิกฺขุนี อคิลานา สปฺปิํ วิญฺญาเปตฺวา ภุญฺเชยฺย, ปฏิเทเสตพฺพํ ตาย ภิกฺขุนิยา คารยฺหํ อยฺเย ธมฺมํ อาปชฺชิํ อสปฺปายํ ปาฏิเทสนียํ ตํ ปฏิเทเสมี”}ติ.}

\proseitem{1231}{\textbf{ยา ปนา}ติ ยา ยาทิสา ฯเปฯ.\csromanspacer{} \textbf{ภิกฺขุนี}ติ ฯเปฯ อยํ อิมสฺมิํ อตฺเถ อธิปฺเปตา ภิกฺขุนีติ.}

\prose{\textbf{อคิลานา} นาม ยสฺสา วินา สปฺปินา ผาสุ โหติ.}

\prose{\textbf{คิลานา} นาม ยสฺสา วินา สปฺปินา น ผาสุ โหติ.}

\prose{\textbf{สปฺปิ} นาม โคสปฺปิ วา อชิกาสปฺปิ วา มหิํสสปฺปิ วา, เยสํ มํสํ กปฺปติ เตสํ สปฺปิ.}

\prose{อคิลานา อตฺตโน อตฺถาย วิญฺญาเปติ, ปโยเค ทุกฺกฏํ.\csromanspacer{} ปฏิลาเภน “ภุญฺชิสฺสามี”ติ, ปฏิคฺคณฺหาติ, อาปตฺติ ทุกฺกฏสฺส.\csromanspacer{} อชฺโฌหาเร อชฺโฌหาเร อาปตฺติ ปาฏิเทสนียสฺส.}

\proseitem{1232}{อคิลานา อคิลานสญฺญา สปฺปิํ วิญฺญาเปตฺวา ภุญฺชติ, อาปตฺติ ปาฏิเทสนียสฺส.\csromanspacer{} อคิลานา เวมติกา สปฺปิํ วิญฺญาเปตฺวา ภุญฺชติ, อาปตฺติ ปาฏิเทสนียสฺส.\csromanspacer{} อคิลานา คเลนสญฺญา สปฺปิํ วิญฺญาเปตฺวา ภุญฺชติ, อาปตฺติ ปาฏิเทสนียสฺส.}

\prose{คิลานา อคิลานสญฺญา, อาปตฺติ ทุกฺกฏสฺส.\csromanspacer{} คิลานา เวมติกา, อาปตฺติ ทุกฺกฏสฺส.\csromanspacer{} คิลานา คิลานสญฺญา, อนาปตฺติ.}

\proseitemwithcloserruled{1233}{\csromanfolio{465}อนาปตฺติ คิลานาย, คิลานา หุตฺวา วิญฺญาเปตฺวา อคิลานา ภุญฺชติ, คิลานาย เสสกํ ภุญฺชติ, ญาตกานํ ปวาริตานํ อญฺญสฺสตฺถาย อตฺตโน ธเนน อุมฺมตฺติกาย อาทิกมฺมิกายาติ.}{ปฐม\-ปาฏิเทสนีย\-สิกฺขาปทํ นิฏฺฐิตํ.}
\csromanmatikaanchor{mtk.254}
\csromantocmarkref{section}{2. ทุติยาทิปาฏิเทสนียสิกฺขาปท}{mtk.254}
\bookmark[dest={mtk.254},level=1]{2. ทุติยาทิปาฏิเทสนียสิกฺขาปท}
\setcsromanheadh{2. ทุติยาทิปาฏิเทสนียสิกฺขาปท}
\csromanheader{1}{2. \textbf{ทุติยาทิปาฏิเทสนียสิกฺขาปท}}
\proseitem{1234}{เตน สมเยน พุทฺโธ ภควา สาวตฺถิยํ วิหรติ เชตวเน อนาถปิณฺฑิกสฺส อาราเม.\csromanspacer{} เตน โข ปน สมเยน ฉพฺพคฺคิยา ภิกฺขุนิโย เตลํ วิญฺญาเปตฺวา ภุญฺชนฺติ ฯเปฯ.\csromanspacer{} มธุํ วิญฺญาเปตฺวา ภุญฺชนฺติ ฯเปฯ.\csromanspacer{} ผาณิตํ วิญฺญาเปตฺวา ภุญฺชนฺติ ฯเปฯ.\csromanspacer{} มจฺฉํ วิญฺญาเปตฺวา ภุญฺชนฺติ - ป-.\csromanspacer{} มํสํ วิญฺญาเปตฺวา ภุญฺชนฺติ ฯเปฯ.\csromanspacer{} ขีรํ วิญฺญาเปตฺวา ภุญฺชนฺติ ฯเปฯ.\csromanspacer{} ทธิํ วิญฺญาเปตฺวา ภุญฺชนฺติ.\csromanspacer{} มนุสฺสา อุชฺฌายนฺติ ขิยฺยนฺติ วิปาเจนฺติ “กถํ หิ นาม ภิกฺขุนิโย ทธิํ วิญฺญาเปตฺวา ภุญฺชิสฺสนฺติ, กสฺส สมฺปนฺนํ น มนาปํ, กสฺส สาทุํ น รุจฺจตี”ติ.\csromanspacer{} อสฺโสสุํ โข ภิกฺขุนิโย เตสํ มนุสฺสานํ อุชฺฌายนฺตานํ ขิยฺยนฺตานํ วิปาเจนฺตานํ.\csromanspacer{} ยา ตา ภิกฺขุนิโย อปฺปิจฺฉา ฯเปฯ ตา อุชฺฌายนฺติ ขิยฺยนฺติ วิปาเจนฺติ “กถํ หิ นาม ฉพฺพคฺคิยา ภิกฺขุนิโย ทธิํ วิญฺญาเปตฺวา ภุญฺชิสฺสนฺตี”ติ ฯเปฯ.\csromanspacer{} สจฺจํ กิร ภิกฺขเว ฉพฺพคฺคิยา ภิกฺขุนิโย ทธิํ วิญฺญาเปตฺวา ภุญฺชนฺตีติ.\csromanspacer{} สจฺจํ ภควาติ.\csromanspacer{} วิครหิ พุทฺโธ ภควา ฯเปฯ.\csromanspacer{} กถํ หิ นาม ภิกฺขเว ฉพฺพคฺคิยา ภิกฺขุนิโย ทธิํ วิญฺญาเปตฺวา ภุญฺชิสฺสนฺติ.\csromanspacer{} เนตํ ภิกฺขเว อปฺปสนฺนานํ วา ปสาทาย ฯเปฯ.\csromanspacer{} เอวญฺจ ปน ภิกฺขเว ภิกฺขุนิโย อิมํ สิกฺขาปทํ อุทฺทิสนฺตุ\csromandash{}}

\prose{\textbf{“ยา ปน ภิกฺขุนี ทธิํ วิญฺญาเปตฺวา ภุญฺเชยฺย, ปฏิเทเสตพฺพํ ตาย ภิกฺขุนิยา คารยฺหํ อยฺเย ธมฺมํ อาปชฺชิํ อสปฺปายํ ปาฏิเทสนียํ ตํ ปฏิเทเสมี”}ติ.}

\prose{เอวญฺจิทํ ภควตา ภิกฺขุนีนํ สิกฺขาปทํ ปญฺญตฺตํ โหติ.}

\proseitem{1235}{เตน โข ปน สมเยน ภิกฺขุนิโย คิลานา โหนฺติ.\csromanspacer{} คิลานปุจฺฉิกา ภิกฺขุนิโย คิลานา ภิกฺขุนิโย เอตทโวจุํ “กจฺจิ อยฺเย ขมนียํ, กจฺจิ ยาปนียนฺ”ติ.\csromanspacer{} ปุพฺเพ มยํ อยฺเย ทธิํ วิญฺญาเปตฺวา \csromanfolio{466}ภุญฺชิมฺหา, เตน โน ผาสุ โหติ, อิทานิ ปน “ภควตา ปฏิกฺขิตฺตนฺ”ติ กุกฺกุจฺจายนฺตา น วิญฺญาเปม, เตน โน น ผาสุ โหตีติ ฯเปฯ.\csromanspacer{} ภควโต เอตมตฺถํ อาโรเจสุํ ฯเปฯ.\csromanspacer{} อนุชานามิ ภิกฺขเว คิลานาย ภิกฺขุนิยา ทธิํ วิญฺญาเปตฺวา ภุญฺชิตุํ.\csromanspacer{} เอวญฺจ ปน ภิกฺขเว ภิกฺขุนิโย อิมํ สิกฺขาปทํ อุทฺทิสนฺตุ\csromandash{}}

\proseitem{1236}{\textbf{“ยา ปน ภิกฺขุนี อคิลานา (เตลํ ฯเปฯ มธุํ ฯเปฯ ผาณิตํ ฯเปฯ มจฺฉํ ฯเปฯ มํสํ ฯเปฯ ขีรํ ฯเปฯ) ทธิํ วิญฺญาเปตฺวา ภุญฺเชยฺย, ปฏิเทเสตพฺพํ ตาย ภิกฺขุนิยา คารยฺหํ อยฺเย ธมฺมํ อาปชฺชิํ อสปฺปายํ ปาฏิเทสนียํ ตํ ปฏิเทเสมี”}ติ.}

\proseitem{1237}{\textbf{ยา ปนา}ติ ยา ยาทิสา ฯเปฯ.\csromanspacer{} \textbf{ภิกฺขุนี}ติ ฯเปฯ อยํ อิมสฺมิํ อตฺเถ อธิปฺเปตา ภิกฺขุนีติ.}

\prose{\textbf{อคิลานา} นาม ยสฺสา วินา ทธินา ผาสุ โหติ.}

\prose{\textbf{คิลานา} นาม ยสฺสา วินา ทธินา น ผาสุ โหติ.}

\prose{\textbf{เตลํ} นาม ติลเตลํ สาสปเตลํ มธุกเตลํ เอรณฺฑเตลํ วสาเตลํ.}

\prose{\textbf{มธุ} นาม มกฺขิกามธุ.\csromanspacer{} \textbf{ผาณิตํ} นาม อุจฺฉุมฺหา นิพฺพตฺตํ.\csromanspacer{} \textbf{มจฺโฉ} นาม โอทโก วุจฺจติ.\csromanspacer{} \textbf{มํสํ} นาม เยสํ มํสํ กปฺปติ เตสํ มํสํ.\csromanspacer{} \textbf{ขีรํ} นาม โคขีรํ วา อชิกาขีรํ วา มหิํสขีรํ วา, เยสํ มํสํ กปฺปติ เตสํ ขีรํ.\csromanspacer{} \textbf{ทธิ} นาม เตสญฺเญว ทธิ.}

\prose{อคิลานา อตฺตโน อตฺถาย วิญฺญาเปติ, ปโยเค ทุกฺกฏํ.\csromanspacer{} ปฏิลาเภน “ภุญฺชิสฺสามี”ติ ปฏิคฺคณฺหาติ, อาปตฺติ ทุกฺกฏสฺส.\csromanspacer{} อชฺโฌหาเร อชฺโฌหาเร อาปตฺติ ปาฏิเทสนียสฺส.}

\proseitem{1238}{อคิลานา อคิลานสญฺญา ทธิํ วิญฺญาเปตฺวา ภุญฺชติ, อาปตฺติ ปาฏิเทสนียสฺส.\csromanspacer{} อคิลานา เวมติกา ทธิํ วิญฺญาเปตฺวา ภุญฺชติ, อาปตฺติ ปาฏิเทสนียสฺส.\csromanspacer{} อคิลานา คเลนสญฺญา ทธิํ วิญฺญาเปตฺวา ภุญฺชติ, อาปตฺติ ปาฏิเทสนียสฺส.}

\prose{\csromanfolio{467}คิลานา อคิลานสญฺญา, อาปตฺติ ทุกฺกฏสฺส.\csromanspacer{} คิลานา เวมติกา อาปตฺติ ทุกฺกฏสฺส.\csromanspacer{} คิลานา คิลานสญฺญา, อนาปตฺติ.}

\proseitemwithcloser{1239}{อนาปตฺติ คิลานาย, คิลานา หุตฺวา วิญฺญาเปตฺวา อคิลานา ภุญฺชติ, คิลานาย เสสกํ ภุญฺชติ, ญาตกานํ ปวาริตานํ อญฺญสฺสตฺถาย อตฺตโน ธเนน อุมฺมตฺติกาย อาทิกมฺมิกายาติ.}{อฏฺฐม\-ปาฏิเทสนีย\-สิกฺขาปทํ นิฏฺฐิตํ.}
\prose{อุทฺทิฏฺฐา โข อยฺยาโย อฏฺฐ ปาฏิเทสนียา ธมฺมา, ตตฺถายฺยาโย ปุจฺฉามิ กจฺจิตฺถ ปริสุทฺธา, ทุติยมฺปิ ปุจฺฉามิ กจฺจิตฺถ ปริสุทฺธา, ตติยมฺปิ ปุจฺฉามิ กจฺจิตฺถ ปริสุทฺธา, ปริสุทฺเธตฺถายฺยาโย, ตสฺมา ตุณฺหี, เอวเมตํ ธารยามีติ.}

\csromanheader{2}{\textbf{ภิกฺขุนิวิภงฺเค}}
\nitthitam{\textbf{ปาฏิเทสนียกณฺฑํ นิฏฺฐิตํ.}}
\chapterheadpageread{6. \textbf{เสขิยกณฺฑ}}
\csromanmatikaanchor{mtk.255}
\csromanmatikaanchor{mtk.256}
\csromantocmarknopage{chapter}{6. เสขิยกณฺฑ}{mtk.255}
\bookmark[dest={mtk.255},level=0]{6. เสขิยกณฺฑ}
\csromantocmarkref{section}{1. ปริมณฺฑลวคฺค}{mtk.256}
\bookmark[dest={mtk.256},level=1]{1. ปริมณฺฑลวคฺค}
\setcsromanheadcha{6. เสขิยกณฺฑ}
\setcsromanheadhclear{}
\setcsromanheadh{1. ปริมณฺฑลวคฺค}
\csromanheader{1}{1. \textbf{ปริมณฺฑลวคฺค}}
\prose{\csromanfolio{468}อิเม โข ปนายฺยาโย เสขิยา ธมฺมา อุทฺเทสํ อาคจฺฉนฺติ.}

\proseitem{1240}{เตน สมเยน พุทฺโธ ภควา สาวตฺถิยํ วิหรติ เชตวเน อนาถปิณฺฑิกสฺส อาราเม.\csromanspacer{} เตน โข ปน สมเยน ฉพฺพคฺคิยา ภิกฺขุนิโย ปุรโตปิ ปจฺฉโตปิ โอลมฺเพนฺตี นิวาเสนฺติ.\csromanspacer{} มนุสฺสา อุชฺฌายนฺติ ขิยฺยนฺติ วิปาเจนฺติ “กถํ หิ นาม ภิกฺขุนิโย ปุรโตปิ ปจฺฉโตปิ โอลมฺเพนฺตี นิวาเสสฺสนฺติ เสยฺยถาปิ คิหินิโย กามโภคินิโย”ติ.\csromanspacer{} อสฺโสสุํ โข ภิกฺขุนิโย เตสํ มนุสฺสานํ อุชฺฌายนฺตานํ ขิยฺยนฺตานํ วิปาเจนฺตานํ.\csromanspacer{} ยา ตา ภิกฺขุนิโย อปฺปิจฺฉา ฯเปฯ ตา อุชฺฌายนฺติ ขิยฺยนฺติ วิปาเจนฺติ “กถํ หิ นาม ฉพฺพคฺคิยา ภิกฺขุนิโย ปุรโตปิ ปจฺฉโตปิ โอลมฺเพนฺตี นิวาเสสฺสนฺตี”ติ ฯเปฯ.\csromanspacer{} สจฺจํ กิร ภิกฺขเว ฉพฺพคฺคิยา ภิกฺขุนิโย ปุรโตปิ ปจฺฉโตปิ โอลมฺเพนฺตี นิวาเสนฺตีติ.\csromanspacer{} สจฺจํ ภควาติ.\csromanspacer{} วิครหิ พุทฺโธ ภควา ฯเปฯ.\csromanspacer{} กถํ หิ นาม ภิกฺขเว ฉพฺพคฺคิยา ภิกฺขุนิโย ปุรโตปิ ปจฺฉโตปิ โอลมฺเพนฺตี นิวาเสสฺสนฺติ.\csromanspacer{} เนตํ ภิกฺขเว อปฺปสนฺนานํ วา ปสาทาย ฯเปฯ.\csromanspacer{} เอวญฺจ ปน ภิกฺขเว ภิกฺขุนิโย อิมํ สิกฺขาปทํ อุทฺทิสนฺตุ\csromandash{}}

\prose{\textbf{“ปริมณฺฑลํ นิวาเสสฺสามีติ สิกฺขา กรณียา”}ติ.}

\prose{ปริมณฺฑลํ นิวาเสตพฺพํ นาภิมณฺฑลํ ชาณุมณฺฑลํ ปฏิจฺฉาเทนฺติยา.\csromanspacer{} ยา อนาทริยํ ปฏิจฺจ ปุรโต วา ปจฺฉโต วา โอลมฺเพนฺตี นิวาเสติ, อาปตฺติ ทุกฺกฏสฺส.}

\prose{อนาปตฺติ อสญฺจิจฺจ อสฺสติยา อชานนฺติยา คิลานาย อาปทาสุ อุมฺมตฺติกาย อาทิกมฺมิกายาติ ฯเปฯ. (สํขิตฺตํ)}
\csromansectionrule

\csromanmatikaanchor{mtk.257}
\csromantocmarkref{subsection}{7. ปาทุกวคฺค}{mtk.257}
\bookmark[dest={mtk.257},level=2]{7. ปาทุกวคฺค}
\csromanheader{2}{7. \textbf{ปาทุกวคฺค}}
\proseitem{1241}{เตน สมเยน พุทฺโธ ภควา สาวตฺถิยํ วิหรติ เชตวเน อนาถปิณฺฑิกสฺส อาราเม.\csromanspacer{} เตน โข ปน สมเยน ฉพฺพคฺคิยา ภิกฺขุนิโย อุทเก อุจฺจารมฺปิ ปสฺสาวมฺปิ เขฬมฺปิ กโรนฺติ.\csromanspacer{} มนุสฺสา อุชฺฌายนฺติ ขิยฺยนฺติ วิปาเจนฺติ “กถํ หิ นาม ภิกฺขุนิโย อุทเก \csromanfolio{469}อุจฺจารมฺปิ ปสฺสาวมฺปิ เขฬมฺปิ กริสฺสนฺติ เสยฺยถาปิ คิหินิโย กามโภคินิโย”ติ.\csromanspacer{} อสฺโสสุํ โข ภิกฺขุนิโย เตสํ มนุสฺสานํ อุชฺฌายนฺตานํ ขิยฺยนฺตานํ วิปาเจนฺตานํ.\csromanspacer{} ยา ตา ภิกฺขุนิโย อปฺปิจฺฉา, ตา อุชฺฌายนฺติ ขิยฺยนฺติ วิปาเจนฺติ “กถํ หิ นาม ฉพฺพคฺคิยา ภิกฺขุนิโย อุทเก อุจฺจารมฺปิ ปสฺสาวมฺปิ เขฬมฺปิ กริสฺสนฺตี”ติ.\csromanspacer{} อถ โข ภิกฺขุนิโย ภิกฺขูนํ เอตมตฺถํ อาโรเจสุํ.\csromanspacer{} ภิกฺขู\footnote{เย เต ภิกฺขู ฯเปฯ (?)} ภควโต เอตมตฺถํ อาโรเจสุํ.\csromanspacer{} อถ โข ภควา ฯเปฯ ภิกฺขู ปฏิปุจฺฉิ “สจฺจํ กิร ภิกฺขเว ฉพฺพคฺคิยา ภิกฺขุนิโย อุทเก อุจฺจารมฺปิ ปสฺสาวมฺปิ เขฬมฺปิ กโรนฺตี”ติ.\csromanspacer{} สจฺจํ ภควาติ.\csromanspacer{} วิครหิ พุทฺโธ ภควา ฯเปฯ.\csromanspacer{} กถํ หิ นาม ภิกฺขเว ฉพฺพคฺคิยา ภิกฺขุนิโย อุทเก อุจฺจารมฺปิ ปสฺสาวมฺปิ เขฬมฺปิ กริสฺสนฺติ.\csromanspacer{} เนตํ ภิกฺขเว อปฺปสนฺนานํ วา ปสาทาย ฯเปฯ.\csromanspacer{} เอวญฺจ ปน ภิกฺขเว ภิกฺขุนิโย อิมํ สิกฺขาปทํ อุทฺทิสนฺตุ\csromandash{}}

\prose{\textbf{“น อุทเก อุจฺจารํ วา ปสฺสาวํ วา เขฬํ วา กริสฺสามีติ สิกฺขา กรณียา”}ติ.}

\prose{เอวญฺจิทํ ภควตา ภิกฺขุนีนํ สิกฺขาปทํ ปญฺญตฺตํ โหติ.}

\prose{เตน โข ปน สมเยน คิลานา ภิกฺขุนิโย อุทเก อุจฺจารมฺปิ ปสฺสาวมฺปิ เขฬมฺปิ กาตุํ กุกฺกุจฺจายนฺติ.\csromanspacer{} ภควโต เอตมตฺถํ อาโรเจสุํ ฯเปฯ.\csromanspacer{} อนุชานามิ ภิกฺขเว คิลานาย ภิกฺขุนิยา อุทเก อุจฺจารมฺปิ ปสฺสาวมฺปิ เขฬมฺปิ กาตุํ.\csromanspacer{} เอวญฺจ ปน ภิกฺขเว ภิกฺขุนิโย อิมํ สิกฺขาปทํ อุทฺทิสนฺตุ\csromandash{}}

\prose{\textbf{“น อุทเก อคิลานา อุจฺจารํ วา ปสฺสาวํ วา เขฬํ วา กริสฺสามีติ สิกฺขา กรณียา”}ติ.}

\prose{น อุทเก อคิลานาย อุจฺจาโร วา ปสฺสาโว วา เขโฬ วา กาตพฺโพ.\csromanspacer{} ยา อนาทริยํ ปฏิจฺจ อุทเก อคิลานา อุจฺจารํ วา ปสฺสาวํ วา เขฬํ วา กโรติ, อาปตฺติ ทุกฺกฏสฺส.}

\prosewithcloser{อนาปตฺติ อสญฺจิจฺจ อสฺสติยา อชานนฺติยา คิลานาย ถเล กโต อุทกํ โอตฺถรติ, อาปทาสุ อุมฺมตฺติกาย ขิตฺตจิตฺตาย เวทนาฏฺฏาย อาทิกมฺมิกายาติ.}{ปนฺนรสมสิกฺขาปทํ นิฏฺฐิตํ.}
\vspace{\tipitakaparskip}
\csromancenter{ปาทุกวคฺโค สตฺตโม.}
\prosewithcloserruled{\csromanfolio{470}อุทฺทิฏฺฐา โข อยฺยาโย เสขิยา ธมฺมา, ตตฺถายฺยาโย ปุจฺฉามิ กจฺจิตฺถ ปริสุทฺธา, ทุติยมฺปิ ปุจฺฉามิ กจฺจิตฺถ ปริสุทฺธา, ตติยมฺปิ ปุจฺฉามิ กจฺจิตฺถ ปริสุทฺธา, ปริสุทฺเธตฺถายฺยาโย, ตสฺมา ตุณฺหี, เอวเมตํ ธารยามีติ.}{เสขิยกณฺฑํ นิฏฺฐิตํ.}
\csromanmatikaanchor{mtk.258}
\csromantocmarkref{chapter}{7. อธิกรณสมถ}{mtk.258}
\bookmark[dest={mtk.258},level=0]{7. อธิกรณสมถ}
\setcsromanheadcha{7. อธิกรณสมถ}
\setcsromanheadhclear{}
\chapterhead{7. \textbf{อธิกรณสมถ}}
\prose{อิเม โข ปนายฺยาโย สตฺต อธิกรณสมถา ธมฺมา อุทฺเทสํ อาคจฺฉนฺติ.}

\proseitem{1242}{อุปฺปนฺนุปฺปนฺนานํ อธิกรณานํ สมถาย วูปสมาย สมฺมุขาวินโย ทาตพฺโพ, สติวินโย ทาตพฺโพ, อมูฬฺหวินโย ทาตพฺโพ, ปฏิญฺญาย กาเรตพฺพํ, เยภุยฺยสิกา, ตสฺสปาปิยสิกา, ติณวตฺถารโกติ.}

\prosewithcloser{อุทฺทิฏฺฐา โข อยฺยาโย สตฺต อธิกรณสมถา ธมฺมา, ตตฺถายฺยาโย ปุจฺฉามิ กจฺจิตฺถ ปริสุทฺธา, ทุติยมฺปิ ปุจฺฉามิ กจฺจิตฺถ ปริสุทฺธา, ตติยมฺปิ ปุจฺฉามิ กจฺจิตฺถ ปริสุทฺธา, ปริสุทฺเธตฺถายฺยาโย, ตสฺมา ตุณฺหี, เอวเมตํ ธารยามีติ.}{อธิกรณสมถา นิฏฺฐิตา.}
\prosewithcloser{อุทฺทิฏฺฐํ โข อยฺยาโย นิทานํ, อุทฺทิฏฺฐา อฏฺฐ ปาราชิกา ธมฺมา, อุทฺทิฏฺฐา สตฺตรส สํฆาทิเสสา ธมฺมา, อุทฺทิฏฺฐา ติํส นิสฺสคฺคิยา ปาจิตฺติยา ธมฺมา, อุทฺทิฏฺฐา ฉสฏฺฐิสตา ปาจิตฺติยา ธมฺมา, อุทฺทิฏฺฐา อฏฺฐ ปาฏิเทสนียา ธมฺมา, อุทฺทิฏฺฐา เสขิยา ธมฺมา, อุทฺทิฏฺฐา สตฺต อธิกรณสมถา ธมฺมา, เอตฺตกํ ตสฺส ภควโต สุตฺตาคตํ สุตฺตปริยาปนฺนํ อนฺวทฺธมาสํ อุทฺเทสํ อาคจฺฉติ, ตตฺถ สพฺพาเหว สมคฺคาหิ สมฺโมทมานาหิ อวิวทมานาหิ สิกฺขิตพฺพนฺติ.}{\textbf{ภิกฺขุนิวิภงฺโค นิฏฺฐิโต.}}
\nitthitam{\textbf{ปาจิตฺติยปาฬิ นิฏฺฐิตา.}}
